\documentclass[10pt, letterpaper]{article}

% Inhaltsverzeichnis für Pakettypen (nur für Übersicht im Header, wird nicht im Dokument angezeigt)
% 1. Seitenlayout und Ränder
% 2. Sprache und Zeichensatz
% 3. Mathematik und Theorem-Umgebungen
% 4. Eigene Makros
% 5. Diagramme und Grafiken
% 6. Tabellen und Aufzählungen
% 7. Inhaltsverzeichnis
% 8. Abschnittsüberschriften
% 9. Abstrakt-Umgebung
% 10. Todos/Notizen
% 11. Rahmen/Box-Umgebungen
% 12. Python-Integration
% 13. Literaturverwaltung
% 14. Hyperlinks
% 15. Absatzeinstellungen
% 16. Umgebungen
% 17  Graphik
% 18  Extra
% 00. Titel und Autor

\usepackage{slashed}

% --- 1. Seitenlayout und Ränder ---
\usepackage[margin=3cm]{geometry}

% --- 2. Sprache und Zeichensatz ---
\usepackage[english]{babel}
\usepackage[T1]{fontenc}
\usepackage[utf8]{inputenc}

% --- 3. Mathematik und Theorem-Umgebungen ---
\usepackage{amsmath, amssymb, amsthm}
\usepackage{mathrsfs}
\DeclareMathOperator{\WF}{WF}

% --- 4. Eigene Makros ---
\usepackage{xcolor}
\newcommand{\SKP}{\langle\cdot,\cdot\rangle}
\newcommand{\R}{\mathbb{R}}
\newcommand{\N}{\mathbb{N}}
\newcommand{\Q}{\mathbb{Q}}
\newcommand{\Z}{\mathbb{Z}}
\newcommand{\C}{\mathbb{C}}
\newcommand{\entwurf}[1]{\textcolor{red}{#1}}
\newcommand{\GL}{\mathrm{GL}}
\newcommand{\SL}{\mathrm{SL}}
\newcommand{\SO}{\mathrm{SO}}
\newcommand{\SU}{\mathrm{SU}}
\newcommand{\Sp}{\mathrm{Sp}}
\newcommand{\Ogroup}{\mathrm{O}}
\newcommand{\U}{\mathrm{U}}

% --- 5. Diagramme und Grafiken ---
\usepackage{graphicx}
\graphicspath{{../../Library/Images/}}
\usepackage[export]{adjustbox}
\usepackage{tikz}
\usetikzlibrary{decorations.pathreplacing, arrows.meta, positioning}
\usepackage{tikz-cd}

% --- 6. Tabellen und Aufzählungen ---
\usepackage{enumitem}
\setlist[itemize]{left=0.5cm}

\newenvironment{romanenum}[1][]
  {%
    \ifx&#1&
    \else
      \textbf{#1}\quad
    \fi
    \begin{enumerate}[label=\roman*)]
  }
  {%
    \end{enumerate}%
  }

% --- 7. Inhaltsverzeichnis ---
\usepackage{tocloft}
\renewcommand{\cftsecfont}{\footnotesize}
\renewcommand{\cftsubsecfont}{\footnotesize}
\renewcommand{\cftsubsubsecfont}{\footnotesize}
\renewcommand{\cftsecpagefont}{\footnotesize}
\renewcommand{\cftsubsecpagefont}{\footnotesize}
\renewcommand{\cftsubsubsecpagefont}{\footnotesize}
\usepackage{etoc}

% --- 8. Abschnittsüberschriften ---
\usepackage{titlesec}
\titleformat{\section}{\normalfont\large\bfseries}{\thesection}{1em}{}
\titleformat{\subsection}{\normalfont\normalsize\bfseries}{\thesubsection}{0.5em}{}
\titleformat{\subsubsection}{\normalfont\normalsize\bfseries}{\thesubsubsection}{0.5em}{}
\setcounter{secnumdepth}{4}

% --- 9. Abstrakt-Umgebung ---
\usepackage{changepage}
\renewenvironment{abstract}
  {
    \begin{adjustwidth}{1.5cm}{1.5cm}
    \small
    \textsc{Abstract. –}%
  }
  {
    \end{adjustwidth}
  }

% --- 10. Todos/Notizen ---
\usepackage{todonotes}
% Hellblau definieren
\definecolor{hellblau}{rgb}{0.3,0.6,1}
\newenvironment{note}{\par\begingroup\color{hellblau}\itshape}{\par\endgroup}

% --- 11. Rahmen/Box-Umgebungen ---
\usepackage{mdframed}
\usepackage{tcolorbox}
\colorlet{shadecolor}{gray!25}

\newenvironment{customTheorem}
  {\vspace{10pt}%
   \begin{mdframed}[
     backgroundcolor=gray!20,
     linewidth=0pt,
     innertopmargin=10pt,
     innerbottommargin=10pt,
     skipabove=\dimexpr\topsep+\ht\strutbox\relax,
     skipbelow=\topsep,
   ]}
  {\end{mdframed}
   \vspace{10pt}%
  }

% --- 12. Python-Integration ---
% (Deaktiviert in dieser Version, aktiviere bei Bedarf)
% \usepackage{pythontex}
% \usepackage[makestderr]{pythontex}

% --- 13. Literaturverwaltung ---
\usepackage{csquotes}
%\usepackage[backend=biber, style=alphabetic, citestyle=alphabetic]{biblatex}
%\addbibresource{bibliography.bib}

% --- 14. Hyperlinks ---
\usepackage{hyperref}
\hypersetup{
  colorlinks   = true,
  urlcolor     = blue,
  linkcolor    = blue,
  citecolor    = blue,
  frenchlinks  = true
}

% --- 15. Absatzeinstellungen ---
\usepackage[parfill]{parskip}
\sloppy

% --- 16. Umgebungen ---
\usepackage{thmtools}

\newcommand{\CustomHeading}[3]{%
  \par\medskip\noindent%
  \textbf{#1 #2} \textnormal{(#3)}.\enskip%
}

\newenvironment{DEF}[2]{\begin{unitbox}\CustomHeading{Definition}{#1}{#2}}{\end{unitbox}}
\newenvironment{PROP}[2]{\begin{unitbox}\CustomHeading{Proposition}{#1}{#2}}{\end{unitbox}}
\newenvironment{THEO}[2]{\begin{unitbox}\CustomHeading{Theorem}{#1}{#2}}{\end{unitbox}}
\newenvironment{LEM}[2]{\begin{unitbox}\CustomHeading{Lemma}{#1}{#2}}{\end{unitbox}}
\newenvironment{KORO}[2]{\begin{unitbox}\CustomHeading{Corollar}{#1}{#2}}{\end{unitbox}}
\newenvironment{REM}[2]{\begin{unitbox}\CustomHeading{Remark}{#1}{#2}}{\end{unitbox}}
\newenvironment{EXA}[2]{\begin{unitbox}\CustomHeading{Example}{#1}{#2}}{\end{unitbox}}
\newenvironment{STUD}[2]{\begin{unitbox}\CustomHeading{Study}{#1}{#2}}{\end{unitbox}}
\newenvironment{CONC}[2]{\begin{unitbox}\CustomHeading{Concept}{#1}{#2}}{\end{unitbox}}
\newenvironment{OTH}[2]{\begin{unitbox}\CustomHeading{Other}{#1}{#2}}{\end{unitbox}}
\newenvironment{EXE}[2]{\begin{unitbox}\CustomHeading{Exercise}{#1}{#2}}{\end{unitbox}}
\newenvironment{MOT}[2]{\begin{unitbox}\CustomHeading{Motivation}{#1}{#2}}{\end{unitbox}}
\newenvironment{PROOF}[2]{\begin{unitbox}\CustomHeading{Proof}{#1}{#2}}{\end{unitbox}}

% --- Unit Umgebung für Source-Inhalte ---
\usepackage{mdframed}
\newmdenv[
  linewidth=1pt,
  topline=false,
  bottomline=false,
  rightline=false,
  leftmargin=0cm,
  rightmargin=0cm,
  skipabove=10pt,
  skipbelow=10pt,
  innertopmargin=0.5\baselineskip,
  innerbottommargin=0.5\baselineskip,
  backgroundcolor=gray!10,
  linecolor=gray
]{unitbox}

\newenvironment{unit}[1]
  {\begin{unitbox}\textbf{Unit #1}\par\smallskip}
  {\end{unitbox}}

% --- 17. ... ---

% --- 18. Extras ---
\usepackage{stmaryrd}
\usepackage{bbold}  % falls du \mathbb{1} nutzen willst

% --- 00. Titel und Autor ---
\title{Mein Titel}
\author{Tim Jaschik}
\date{\today}

\begin{document}

\maketitle
\rule{\textwidth}{0.5pt}
\begin{abstract}
Kurze Beschreibung …
\end{abstract}
\rule{\textwidth}{0.5pt}
\vspace{0.5cm}

\tableofcontents

\pagebreak


\section{Funktorieller Zugang zur Lie-Theorie}

Wir schreiben:
\begin{itemize}
  \item \textbf{Kategorie der Mannigfaltigkeiten:} 
  Die Kategorie der glatten, endlichdimensionalen Mannigfaltigkeiten (ohne Rand) mit glatten Abbildungen.
  \item \textbf{Kategorie der Lie-Gruppen:} 
  Objekte sind Lie-Gruppen, Morphismen sind glatte Gruppenhomomorphismen.
  \item \textbf{Kategorie der Lie-Algebren:}
  Objekte sind endlichdimensionale reelle Lie-Algebren, Morphismen sind Lie-Algebrenhomomorphismen.
  \item \textbf{Kategorie der Vektorräume:}
  Objekte sind endlichdimensionale reelle Vektorräume, Morphismen sind lineare Abbildungen.
\end{itemize}



Sei \( \mathcal{C} \) eine Kategorie mit endlichen Produkten.  
Ein \emph{Gruppenobjekt} in \( \mathcal{C} \) besteht aus:
\[
(G,\, \mu,\, \iota,\, e),
\]
wobei
\[
\mu : G \times G \to G, \quad
\iota : G \to G, \quad
e : * \to G
\]
Morphismen in \( \mathcal{C} \) sind, die die Diagrammform der Gruppenaxiome erfüllen:

\begin{enumerate}
  \item \textbf{Assoziativität:}
  \[
  \begin{tikzcd}
    (G \times G) \times G \arrow[r,"{\mu \times \mathrm{id}_G}"] \arrow[d,"{\cong}"'] 
      & G \times G \arrow[d,"\mu"] \\
    G \times (G \times G) \arrow[r,"{\mathrm{id}_G \times \mu}"'] & G
  \end{tikzcd}
  \]

  \item \textbf{Neutrales Element:}
  \[
  \begin{tikzcd}
    G \arrow[r,"{\langle \mathrm{id}_G,\, e \circ ! \rangle}"] \arrow[dr,"\mathrm{id}_G"'] 
      & G \times G \arrow[d,"\mu"] \\
      & G
  \end{tikzcd}
  \qquad
  \begin{tikzcd}
    G \arrow[r,"{\langle e \circ !,\, \mathrm{id}_G \rangle}"] \arrow[dr,"\mathrm{id}_G"'] 
      & G \times G \arrow[d,"\mu"] \\
      & G
  \end{tikzcd}
  \]
  wobei \( !: G \to * \) die eindeutige Abbildung zum terminalen Objekt ist.

  \item \textbf{Inversion:}
  \[
  \begin{tikzcd}
    G \arrow[r,"{\langle \mathrm{id}_G,\, \iota \rangle}"] \arrow[drr,"e \circ !"'] 
      & G \times G \arrow[r,"\mu"] & G \\
      && G \arrow[u,equals]
  \end{tikzcd}
  \]
\end{enumerate}



Eine \emph{Lie-Gruppe} ist ein Gruppenobjekt
\[
(G,\, \mu,\, \iota,\, e)
\]
in der Kategorie der glatten Mannigfaltigkeiten.  
Das bedeutet:
\begin{itemize}
  \item \( G \) ist eine glatte Mannigfaltigkeit.
  \item Die Abbildungen
  \[
  \mu : G \times G \to G, \qquad
  \iota : G \to G, \qquad
  e : * \to G
  \]
  sind glatt.
  \item Sie erfüllen die Assoziativitäts-, Neutralitäts- und Inversionsdiagramme aus der Definition oben.
\end{itemize}

Morphismen zwischen Lie-Gruppen sind glatte Gruppenhomomorphismen \(\phi : G \to H\), welche die Strukturabbildungen respektieren, d.\,h.
\[
\phi \circ \mu_G = \mu_H \circ (\phi \times \phi), \quad
\phi \circ \iota_G = \iota_H \circ \phi, \quad
\phi(e_G) = e_H.
\]



Der \emph{Tangentialfunktor}
\[
T : \text{Kategorie der Mannigfaltigkeiten} \longrightarrow 
    \text{Kategorie der Mannigfaltigkeiten}
\]
ordnet zu:
\[
M \longmapsto TM, 
\qquad 
f : M \to N \longmapsto Tf : TM \to TN,
\]
wobei \( Tf \) das Differential von \( f \) ist.  
Er erfüllt:
\[
T(\mathrm{id}_M) = \mathrm{id}_{TM}, 
\qquad 
T(g \circ f) = Tg \circ Tf,
\]
und ist \textbf{produktverträglich}:
\[
T(M \times N) \cong TM \times TN.
\]



Der \emph{Tangentialfunktor}
\[
T : \text{Kategorie der Mannigfaltigkeiten} \longrightarrow 
    \text{Kategorie der Mannigfaltigkeiten}
\]
ordnet zu:
\[
M \longmapsto TM, 
\qquad 
f : M \to N \longmapsto Tf : TM \to TN,
\]
wobei \( Tf \) das Differential von \( f \) ist.  
Er erfüllt:
\[
T(\mathrm{id}_M) = \mathrm{id}_{TM}, 
\qquad 
T(g \circ f) = Tg \circ Tf,
\]
und ist \textbf{produktverträglich}:
\[
T(M \times N) \cong TM \times TN.
\]


Sei \((G, \mu, \iota, e)\) ein Gruppenobjekt in der Kategorie der Mannigfaltigkeiten.  
Dann ist
\[
(TG,\, T\mu,\, T\iota,\, Te)
\]
wieder ein Gruppenobjekt in dieser Kategorie.  
Insbesondere ist \( TG \) kanonisch eine Lie-Gruppe.


Wendet man den Tangentialfunktor \(T\) auf alle Strukturabbildungen und Diagramme des Gruppenobjekts an, so bleiben die Diagramme aufgrund der Funktorialität und Produktverträglichkeit von \(T\) kommutativ. Damit erfüllen die abgeleiteten Abbildungen die Gruppenaxiome in der Kategorie der Mannigfaltigkeiten.



Unter der Identifikation \( T(G \times G) \cong TG \times TG \) ist die Multiplikation auf \( TG \) gegeben durch
\[
(v_g) \cdot (w_h)
\;:=\;
T_{(g,h)}\mu\, (v_g, w_h)
\;\in\;
T_{gh}G.
\]
Die Inversion lautet \( (v_g)^{-1} := T_g\iota(v_g) \),
und das neutrale Element ist der Nullvektor \( 0_e \in T_e G \).



Für $g \in G$ sei die linke Translation
\[
L_g : G \longrightarrow G,\qquad L_g(h) := gh.
\]
Ihr Differential in $h$ wird mit $(L_g)_{*}\big|_h : T_h G \to T_{gh} G$ bezeichnet.



Linke Parallelisierung der Tangentialräume

Für jedes $g \in G$ ist $(L_g)_{*}\big|_e : T_e G \to T_g G$ ein linearer Isomorphismus. Die Abbildung
\[
\lambda : TG \longrightarrow G \times T_e G,\qquad
\lambda(v_g) := \bigl(g,\ (L_{g^{-1}})_{*}\big|_g\, v_g \bigr),
\]
ist glatt und bijektiv. Ihre inverse Abbildung lautet
\[
\lambda^{-1}(g,X) := (L_g)_{*}\big|_e\, X \;\in\; T_g G.
\]



Die Abbildung $(L_g)_{*}\big|_e$ ist ein Isomorphismus, da $L_g$ ein Diffeomorphismus ist. Für die Definition von $\lambda$ gilt:
\[
(L_{g^{-1}})_{*}\big|_g : T_g G \longrightarrow T_e G
\]
ist ebenfalls ein linearer Isomorphismus. Offensichtlich ist $\lambda$ glatt (sie besteht aus glatten Bündelmorphismen) und bijektiv mit inverser Abbildung wie angegeben. Die Umkehrung prüft man direkt:
\[
\lambda\big( (L_g)_{*}\big|_e X \big) = \bigl(g,\ (L_{g^{-1}})_{*}\big|_g (L_g)_{*}\big|_e X \bigr) = (g,X),
\]
sowie umgekehrt $\lambda^{-1}\bigl(g,(L_{g^{-1}})_{*} v_g\bigr)=(L_g)_{*}(L_{g^{-1}})_{*}v_g=v_g$.



Kanonische Trivialisierung

Die Abbildung $\lambda$ identifiziert $TG$ als \emph{triviales Vektorbündel} über $G$:
\[
TG \cong G \times T_e G,
\]
wobei die Faser über $g$ durch $(L_g)_{*}\big|_e : T_e G \to T_g G$ kanonisch identifiziert wird. Diese Trivialisierung ist \emph{linksinvariant} und wird häufig als \emph{Maurer–Cartan-Trivialisierung} des Tangentialbündels bezeichnet.





Konjugation und adjungierte Wirkung

Für $g \in G$ sei die Konjugationsabbildung
\[
c_g : G \longrightarrow G,\qquad c_g(h) := g h g^{-1}.
\]
Die \emph{adjungierte Wirkung} ist definiert durch das Differential in der Identität:
\[
\mathrm{Ad}_g := \bigl(T_e c_g\bigr) : T_e G \longrightarrow T_e G.
\]




Darstellungseigenschaften der adjungierten Wirkung

Die Abbildung $g \mapsto \mathrm{Ad}_g$ ist ein glatter Gruppenhomomorphismus
\[
\mathrm{Ad} : G \longrightarrow \mathrm{GL}(T_e G),\qquad
\mathrm{Ad}_{gh} = \mathrm{Ad}_g \circ \mathrm{Ad}_h,\quad \mathrm{Ad}_e = \mathrm{id}_{T_e G}.
\]


Für $g,h \in G$ gilt $c_{gh} = c_g \circ c_h$. Differenziation in $e$ liefert
\[
T_e c_{gh} = T_e(c_g \circ c_h) = (T_e c_g) \circ (T_e c_h),
\]
also $\mathrm{Ad}_{gh} = \mathrm{Ad}_g \circ \mathrm{Ad}_h$. Weiter ist $c_e = \mathrm{id}_G$, also $\mathrm{Ad}_e=\mathrm{id}$.


Adjungierte Lie-Algebra-Wirkung
Die Ableitung der adjungierten Wirkung in der Identität liefert
\[
\mathrm{ad} : T_e G \longrightarrow \mathfrak{gl}(T_e G),\qquad
\mathrm{ad}_X := \left.\frac{d}{dt}\right|_{t=0} \mathrm{Ad}_{\exp(tX)}.
\]
Die Lie-Klammer auf $T_e G$ ist gegeben durch
\[
[X,Y] := \mathrm{ad}_X(Y).
\]



Gruppenoperation auf $TG$ in der Trivialisierung

Unter der Identifikation $\lambda : TG \xrightarrow{\cong} G \times T_e G$ ist die (in Teil~I via $T\mu$ definierte) Multiplikation auf $TG$ durch
\[
(g,X)\cdot(h,Y) \;=\; \bigl(gh,\ X + \mathrm{Ad}_g(Y)\bigr)
\]
gegeben. Das neutrale Element ist $(e,0)$ und die Inversion lautet
\[
(g,X)^{-1} \;=\; \bigl(g^{-1},\, -\,\mathrm{Ad}_{g^{-1}}(X)\bigr).
\]



Seien $v_g \in T_g G$, $w_h \in T_h G$. Per Definition der Gruppenstruktur auf $TG$ gilt
\[
v_g \cdot w_h \;=\; T_{(g,h)}\mu \,(v_g,w_h) \;\in\; T_{gh}G,
\]
mit $\mu(g,h)=gh$. Wenden wir auf beiden Seiten $\lambda$ an und benutzen die Definition $\lambda(v_g)=(g,(L_{g^{-1}})_{*}v_g)$, so genügt es, die zweite Komponente zu berechnen:
\[
(L_{(gh)^{-1}})_{*}\, \bigl( T_{(g,h)}\mu \,(v_g,w_h) \bigr)
\;\in\; T_e G.
\]
Die Kettenregel und die Identität $\mu = L_g \circ R_h$ an der Stelle $(g,h)$ (besser: die bekannte Bilinearitätszerlegung des Differentials der Multiplikation) liefern
\[
T_{(g,h)}\mu \,(v_g,w_h) \;=\; T_h(L_g)\,w_h \;+\; T_g(R_h)\,v_g.
\]
Damit
\[
(L_{(gh)^{-1}})_{*}\, T_{(g,h)}\mu \,(v_g,w_h)
= (L_{(gh)^{-1}})_{*}\, T_h(L_g)\,w_h \;+\; (L_{(gh)^{-1}})_{*}\, T_g(R_h)\,v_g.
\]
Für den ersten Term:
\[
(L_{(gh)^{-1}})_{*}\, T_h(L_g) = (L_{h^{-1}})_{*}\,(L_{g^{-1}})_{*}\, T_h(L_g) = (L_{h^{-1}})_{*},
\]
da $(L_{g^{-1}})_{*}\, T_h(L_g)=\mathrm{id}$ nach Funktorialität. Somit ist der erste Beitrag $(L_{h^{-1}})_{*}\, w_h$.

Für den zweiten Term:
\[
(L_{(gh)^{-1}})_{*}\, T_g(R_h)\,v_g
= \bigl( (L_{h^{-1}})_{*} \circ (L_{g^{-1}})_{*} \bigr)\, T_g(R_h)\,v_g
= (L_{h^{-1}})_{*}\, \bigl( (L_{g^{-1}})_{*} \circ T_g(R_h) \bigr)\, v_g.
\]
Die Kettenregel in Gestalt der Konjugation $c_g = L_g \circ R_{g^{-1}}$ impliziert
\[
(L_{g^{-1}})_{*}\big|_g \circ T_g(R_h) = T_e(c_{g}) \circ (L_{g^{-1}})_{*}\big|_g = \mathrm{Ad}_g \circ (L_{g^{-1}})_{*}\big|_g.
\]
Setzt man dies ein, ergibt sich
\[
(L_{(gh)^{-1}})_{*}\, T_g(R_h)\,v_g
= (L_{h^{-1}})_{*}\, \bigl( \mathrm{Ad}_g \big( (L_{g^{-1}})_{*} v_g \big) \bigr).
\]
In Summe:
\[
(L_{(gh)^{-1}})_{*}\, \bigl( T_{(g,h)}\mu \,(v_g,w_h) \bigr)
= (L_{h^{-1}})_{*}\, w_h \;+\; (L_{h^{-1}})_{*}\, \bigl( \mathrm{Ad}_g ( (L_{g^{-1}})_{*} v_g ) \bigr).
\]
Wenden wir nun $\lambda$ auf $v_g$ und $w_h$ an, also $X=(L_{g^{-1}})_{*} v_g$, $Y=(L_{h^{-1}})_{*} w_h$, so erhalten wir die zweite Komponente
\[
Y \;+\; \mathrm{Ad}_g(X).
\]
Beachtet man die Reihenfolge, in der die Faktoren auftreten, erhält man unter $\lambda$ schließlich
\[
\lambda(v_g \cdot w_h) \;=\; \bigl( gh,\ X + \mathrm{Ad}_g(Y) \bigr),
\]
wobei $\lambda(v_g)=(g,X)$ und $\lambda(w_h)=(h,Y)$ ist. Dies zeigt die angegebene Multiplikationsformel. Neutrales Element und Inversion folgen durch direkte Rechnung aus $T\iota$ und $Te$ bzw.\ durch Lösen der Gleichung $(g,X)\cdot(g,X)^{-1}=(e,0)$.


Rechts-trivialisierte Form

Verwendet man statt linker die rechte Trivialisierung $(R_g)_{*}$, erhält man analog die Formel
\[
(g,X)\cdot(h,Y) = \bigl(gh,\ \mathrm{Ad}_{h^{-1}}(X) + Y \bigr).
\]
Beide Beschreibungen sind äquivalent und durch eine Faser-Automorphie ineinander überführbar.




Semidirektes Produkt

Sei $G$ eine Lie-Gruppe und $\rho : G \to \mathrm{GL}(V)$ eine glatte Darstellung auf einem endlichdimensionalen reellen Vektorraum $V$. Das \emph{semidirekte Produkt} $G \ltimes_{\rho} V$ ist die Mannigfaltigkeit $G \times V$ mit Multiplikation
\[
(g,X)\cdot(h,Y) := \bigl(gh,\ X + \rho(g)Y \bigr),
\]
Inversion $(g,X)^{-1} := \bigl(g^{-1}, -\rho(g^{-1})X\bigr)$ und neutralem Element $(e,0)$.


$TG$ als semidirektes Produkt

Unter der kanonischen Trivialisierung $\lambda : TG \to G \times T_e G$ ist $TG$ isomorph zur semidirekten Produktgruppe
\[
G \ltimes_{\mathrm{Ad}} T_e G,
\]
wobei die Darstellung $\rho$ die adjungierte Wirkung $\mathrm{Ad}$ ist. Der Isomorphismus ist ein Lie-Gruppen-Isomorphismus, natürlich in $G$.



Dies ist unmittelbare Umschreibung der Multiplikationsformel der vorigen Proposition mit $\rho=\mathrm{Ad}$.


Kompatibilität mit Morphismen

Sei $\phi : G \to H$ ein glatter Gruppenhomomorphismus zwischen Lie-Gruppen. Dann ist
\[
T\phi : TG \longrightarrow TH
\]
ein Lie-Gruppenhomomorphismus. Unter den Trivialisierungen $\lambda_G$ und $\lambda_H$ korrespondiert $T\phi$ zur Abbildung
\[
\Phi : G \times T_e G \longrightarrow H \times T_e H,\qquad
\Phi(g,X) := \bigl(\phi(g),\ d\phi\big|_e (X)\bigr),
\]
und diese ist ein Homomorphismus von semidirekten Produkten, da
\[
d\phi\big|_e \circ \mathrm{Ad}_g = \mathrm{Ad}_{\phi(g)} \circ d\phi\big|_e.
\]



Die Funktorialität des Tangentialfunktors liefert die Homomorphieeigenschaft von $T\phi$; die angegebene Formel erhält man durch Einsetzen der Trivialisierungen und die Tatsache, dass $\phi$ die Konjugation respektiert, also $c_{\phi(g)} \circ \phi = \phi \circ c_g$, woraus nach Differentiation in $e$ die Ad-Kompatibilität folgt.




Lie-Algebra von $TG$

Die Lie-Algebra der Lie-Gruppe $TG \cong G \ltimes_{\mathrm{Ad}} T_e G$ ist kanonisch isomorph zu
\[
T_e G \;\oplus\; T_e G,
\]
mit Klammer
\[
[(X_1,X_2),(Y_1,Y_2)] \;=\;
\bigl(\ [X_1,Y_1]\ ,\ [X_1,Y_2] + [X_2,Y_1]\ \bigr),
\]
wobei auf der rechten Seite die Lie-Klammer von $G$ in $T_e G$ verwendet wird.



Allgemein gilt für ein semidirektes Produkt $G \ltimes_{\rho} V$ mit Darstellung $\rho$ und abelscher Lie-Algebra $V$:
\[
\mathfrak{g}\ltimes_{\mathrm{d}\rho} V,\qquad
[(X,u),(Y,v)] = \bigl([X,Y],\, (\mathrm{d}\rho)_X v - (\mathrm{d}\rho)_Y u \bigr).
\]
Hier ist $V = T_e G$ und $\rho=\mathrm{Ad}$, also $(\mathrm{d}\rho)_X = \mathrm{ad}_X$, und $V$ trägt die triviale abelsche Klammer. Daher
\[
[(X_1,X_2),(Y_1,Y_2)] = \bigl([X_1,Y_1],\, \mathrm{ad}_{X_1} Y_2 - \mathrm{ad}_{Y_1} X_2 \bigr)
= \bigl([X_1,Y_1],\, [X_1,Y_2] + [X_2,Y_1] \bigr).
\]




In diesem Abschnitt analysieren wir die Gruppenmultiplikation einer Lie-Gruppe $G$ in einer Umgebung der Einheit mittels ihrer Jet-Entwicklung. Dadurch wird klar, dass die Lie-Klammer auf $\mathfrak g = T_e G$ genau dem antisymmetrischen Anteil des zweiten Jets der Multiplikation entspricht.



$k$-Jet

Seien $M,N$ glatte Mannigfaltigkeiten und $p\in M$, $q=f(p)\in N$. 
Zwei glatte Abbildungen $f,g:M\to N$ heißen \emph{$k$-äquivalent in $p$}, falls ihre partiellen Ableitungen bis Ordnung $k$ in einer Karte um $p$ übereinstimmen. 

Die Äquivalenzklasse von $f$ in $p$ wird als
\[
j^k_p f \in J^k(M,N)_{(p,q)}
\]
bezeichnet und heißt \emph{$k$-Jet} von $f$ in $p$. 



Für $k=1$ ist $j^1_p f$ vollständig durch das Differential $T_p f$ bestimmt. 
Für $k=2$ treten zusätzlich die zweiten partiellen Ableitungen auf, die die erste Nichtlinearität von $f$ erfassen.




Jets der Gruppenstruktur

Für eine Lie-Gruppe $(G,\mu,\iota,e)$ bezeichnen wir
\[
j^k_{(e,e)}\mu \in J^k(G\times G,G)_{((e,e),e)},\qquad
j^k_e\iota \in J^k(G,G)_{(e,e)}
\]
als die $k$-Jet-Daten der Gruppenstruktur in der Einheit. 
Diese Jets kodieren, wie die Multiplikation in einer Umgebung der Einheit aussieht.



\begin{itemize}
  \item Der $0$-Jet enthält nur die Information $\mu(e,e)=e$.
  \item Der $1$-Jet ist das Differential $T_{(e,e)}\mu:\mathfrak g\oplus\mathfrak g\to\mathfrak g$.
  \item Der $2$-Jet enthält die ersten Nichtlinearitäten, insbesondere die Nichtkommutativität.
\end{itemize}




Lokale Darstellung der Multiplikation

In einer Exponentialumgebung von $e$ gibt es eine glatte Abbildung 
\[
Z : \mathfrak g \times \mathfrak g \to \mathfrak g
\]
so dass
\[
\exp(X)\exp(Y)=\exp(Z(X,Y)).
\]
Diese Abbildung ist eindeutig in einer Umgebung von $(0,0)$ und erfüllt $Z(0,0)=0$.


Beweisskizze

Die Exponentialabbildung ist ein lokaler Diffeomorphismus $ \exp : \mathfrak g \supset U \to V \subset G$ mit $U,V$ Umgebungen von $0,e$. Damit können wir $\mu$ auf $U\times U$ transportieren:
\[
Z := \log \circ \mu \circ (\exp\times\exp).
\]
Die Glattheit folgt aus der Glattheit von $\mu$ und $\exp$. 



Baker--Campbell--Hausdorff-Reihe

Die formale Taylorreihe von $Z(X,Y)$ um $(0,0)$ wird als \emph{Baker--Campbell--Hausdorff (BCH)-Reihe} bezeichnet:
\[
Z(X,Y) = X + Y + \tfrac{1}{2}[X,Y]
+ \tfrac{1}{12}[X,[X,Y]] - \tfrac{1}{12}[Y,[X,Y]] + \cdots.
\]




Erster Jet

Der erste Jet $j^1_{(e,e)}\mu$ entspricht der linearen Abbildung
\[
T_{(e,e)}\mu(X,Y) = X + Y,
\]
d.\,h.\ die linearisierte Multiplikation ist kommutativ und identifiziert sich mit der Vektorraumaddition auf $\mathfrak g$.



Setze $g=\exp(tX)$, $h=\exp(sY)$ und benutze $Z(X,Y)=X+Y+O(2)$. Dann gilt
\[
\mu(\exp(tX),\exp(sY)) = \exp(tX)\exp(sY)
 = \exp(tX + sY + O(t^2,s^2,ts)).
\]
Differentiation nach $t,s$ in $0$ liefert $T_{(e,e)}\mu(X,Y)=X+Y$.



Zweiter Jet und antisymmetrischer Anteil

Der antisymmetrische Teil der zweiten Ableitung von $Z$ in $(0,0)$ ist
\[
\operatorname{Alt}\,D^2 Z(0,0)(X,Y)
= \tfrac{1}{2}[X,Y].
\]



Die BCH-Reihe zeigt, dass der Term $\tfrac12[X,Y]$ der erste nichtlineare (quadratische) Beitrag von $Z$ ist. 
Die Symmetrieeigenschaften der partiellen Ableitungen implizieren, dass der symmetrische Teil der zweiten Ableitung verschwindet, sodass die antisymmetrische Komponente $\tfrac12[X,Y]$ ergibt.



Lie-Klammer aus dem zweiten Jet

Die Lie-Klammer auf $\mathfrak g$ ist die eindeutig bestimmte antisymmetrische Bilinearform
\[
[X,Y] := 2\,\operatorname{Alt}\,D^2 Z(0,0)(X,Y),
\]
die die Jacobi-Identität erfüllt. 
Damit ist die Lie-Algebra $(\mathfrak g,[\cdot,\cdot])$ vollständig durch die $2$-Jet-Daten der Gruppenmultiplikation bestimmt.



Gruppenkommutator

Für $g,h\in G$ sei
\[
\mathcal{C}(g,h):= ghg^{-1}h^{-1}.
\]


Differential des Kommutators

In Exponentialkoordinaten gilt
\[
\mathcal{C}(\exp(tX),\exp(sY))
= \exp\!\big(ts[X,Y] + O(t^2,s^2,ts^2,t^2s)\big).
\]
Insbesondere folgt
\[
[X,Y]
= \left.\frac{\partial^2}{\partial t\,\partial s}\right|_{(0,0)}
\log\!\mathcal{C}(\exp(tX),\exp(sY)).
\]



Einsetzen der BCH-Reihe in $\exp(tX)\exp(sY)\exp(-tX)\exp(-sY)$ und Entwickeln nach $t,s$ zeigt, dass der niedrigste nichttriviale Term $ts[X,Y]$ ist. 


Interpretation

Der Kommutator misst die Nichtkommutativität der Gruppe, 
während die Lie-Klammer die \emph{infinitesimale Nichtkommutativität} darstellt, d.\,h.\ die linearisierte Abweichung von der Kommutativität in zweiter Ordnung.





Infinitesimale Struktur

Sei $(G,\mu,\iota,e)$ eine Lie-Gruppe. 
Die \emph{infinitesimale Struktur} von $G$ ist das Paar
\[
(T_e G,\ [\cdot,\cdot]),
\]
wobei
\begin{itemize}
  \item $T_e G$ die linearisierte (erste Jet-)Struktur der Mannigfaltigkeit in der Einheit beschreibt,
  \item die Lie-Klammer $[\cdot,\cdot]$ der antisymmetrische Anteil des zweiten Jets der Multiplikation ist.
\end{itemize}
Diese Struktur hängt nur von den $2$-Jet-Daten der Multiplikation in der Einheit ab und bestimmt $G$ bis zu lokaler Isomorphie in der Umgebung von $e$.


Der Lie-Funktor

Die Zuweisung
\[
\mathrm{Lie}: 
\text{Kategorie der Lie-Gruppen} 
\longrightarrow 
\text{Kategorie der Lie-Algebren},\qquad
G \mapsto (T_e G,[\cdot,\cdot]),
\]
\[
\phi:G\to H \longmapsto d\phi_e : T_e G \to T_e H
\]
ist ein kovarianter Funktor.




Für ein Homomorphismus $\phi$ gilt $\phi\circ\mu_G = \mu_H\circ(\phi\times\phi)$ und $\phi\circ c_g = c_{\phi(g)}\circ\phi$. Differentiation in der Einheit zeigt
\[
d\phi_e\,[X,Y] = [\,d\phi_e X,\, d\phi_e Y\,],
\]
also ist $d\phi_e$ ein Lie-Algebrenhomomorphismus. Funktorialität folgt aus der Kettenregel.



\[
\begin{tikzcd}[column sep=large, row sep=large]
  \text{Lie-Gruppen} \arrow[r,"{\text{Differentiation in }e}"] \arrow[d,hook] 
    & \text{Lie-Algebren} \arrow[d,hook] \\
  \text{glatte Mannigfaltigkeiten} \arrow[r,"{T_e}"] & \text{Vektorräume}
\end{tikzcd}
\]


Die Infinitesimalisierung ist damit nichts anderes als die Anwendung des Tangentialfunktors auf das Gruppenobjekt und die Auswertung in der Einheit. Der erste Jet liefert den linearen Anteil, der zweite Jet den antisymmetrischen Anteil, welcher die Lie-Klammer erzeugt.



Die Lie-Algebra einer Lie-Gruppe ist somit das einfachste nichttriviale „lineare Modell“ der Gruppenstruktur: 
Sie kodiert alle Informationen bis zur ersten Abweichung von Linearität und Kommutativität. 
Höhere Jets enthalten zusätzliche formale Information (z.\,B.\ Krümmungstermen), aber keine weitere lineare Struktur. 

\[
\boxed{
\text{Infinitesimale Struktur von } G 
= (T_e G,\ [\cdot,\cdot]) 
= \big(j^1_{e}(G),\ \operatorname{Alt}\,j^2_{(e,e)}\mu\big).
}
\]


Aus Teil III wissen wir: Die $1$-Jet-Daten der Mannigfaltigkeit $G$ in $e$ liefern die Vektorraumstruktur auf $T_eG$, und der antisymmetrische Anteil des $2$-Jets der Multiplikation $\mu$ in $(e,e)$ liefert die Lie-Klammer $[\cdot,\cdot]$ auf $T_eG$. In diesem Teil zeigen wir, warum diese \emph{lokalen} Daten am Einselement ausreichen, um die \emph{infinitesimale Struktur auf allen Tangentialräumen} $T_gG$ zu bestimmen, und präzisieren den Mechanismus der Übertragung.



Linke Translation und kanonische Isomorphismen der Fasern

Für jedes $g\in G$ ist die linke Translation $L_g:G\to G$, $L_g(h)=gh$, ein Diffeomorphismus. Ihr Differential in $e$,
\[
(L_g)_{*}\big|_e \;:\; T_eG \longrightarrow T_gG,
\]
ist ein kanonischer linearer Isomorphismus. Die Familie $\{(L_g)_{*}\big|_e\}_{g\in G}$ trivialisiert das Tangentialbündel zu
\[
TG \;\cong\; G \times T_eG,\qquad
v_g \longmapsto \bigl(g,\ (L_{g^{-1}})_{*}\big|_g v_g\bigr).
\]


Begriffliche Bedeutung

Diese Trivialisierung ist \emph{linksinvariant} und kanonisch: Es wird \emph{keine} Zusatzwahl (wie eine Verbindung) getroffen. Sie kodiert, dass die Tangentialräume überall als Kopien von $T_eG$ identifiziert werden, über die Gruppengeometrie selbst (nicht bloß über eine Chart).




Transportierte Faserklammer

Sei $(T_eG,[\cdot,\cdot])$ die in Teil~III konstruierte Lie-Algebra. Für jedes $g\in G$ definieren wir auf $T_gG$ eine bilineare Operation durch
\[
[u_g,v_g]_{(g)} \;:=\; (L_g)_{*}\big|_e \,\Big[\,(L_{g^{-1}})_{*}\big|_g u_g\;,\; (L_{g^{-1}})_{*}\big|_g v_g\,\Big] \;\in\; T_gG.
\]


Kohärenz der übertragenen Klammern

Für jedes $g\in G$ ist $[\cdot,\cdot]_{(g)}$ eine Lie-Klammer (bilinear, antisymmetrisch, Jacobi-Identität). Zudem ist sie \emph{kompatibel} mit Translationen:
\[
(L_h)_{*}\big|_g \,\big[u_g,v_g\big]_{(g)} 
\;=\; 
\Big[(L_h)_{*}\big|_g u_g,\ (L_h)_{*}\big|_g v_g\Big]_{(hg)}.
\]



Die Eigenschaften folgen unmittelbar aus der Tatsache, dass $(L_g)_{*}\big|_e$ ein linearer Isomorphismus ist und Lie-Klammer-Eigenschaften unter Isomorphismen erhalten bleiben. Die Kompatibilität ist Natur der Pushforwards: $(L_h)_{*}\circ(L_g)_{*}=(L_{hg})_{*}$.



Damit ist präzisiert, was es heißt, dass die \emph{infinitesimale Information in $e$} auf \emph{alle anderen Tangentialräume} übertragen werden kann: Die linke Translation liefert kanonische Isomorphismen der Fasern und transportiert die Lie-Klammer funktoriell auf jede Faser.



Linksinvariante Vektorfelder

Zu $X\in T_eG$ gehört das linksinvariante Vektorfeld $\widetilde{X}$ durch
\[
\widetilde{X}(g) \;:=\; (L_g)_{*}\big|_e X \;\in\; T_gG.
\]
Die Abbildung $X\mapsto \widetilde{X}$ ist ein linearer Isomorphismus zwischen $T_eG$ und dem Raum der linksinvarianten Vektorfelder auf $G$.



Klammererhalt der Linksinvarianz

Für $X,Y\in T_eG$ gilt
\[
[\widetilde{X},\widetilde{Y}](g) \;=\; (L_g)_{*}\big|_e \,[X,Y].
\]
Insbesondere ist $[\widetilde{X},\widetilde{Y}]=\widetilde{[X,Y]}$.



Dies ist die klassische Naturverhaltensformel der Lie-Klammer unter Pushforwards und die Definition der linksinvarianten Felder; vgl.\ $ (L_g)_{*}[\cdot,\cdot] = [ (L_g)_{*}(\cdot),(L_g)_{*}(\cdot) ]$.



Rekonstruktion aus $T_eG$

Die Daten $(T_eG,[\cdot,\cdot])$ bestimmen \emph{alle} linksinvarianten Vektorfelder und deren Klammer global. Umgekehrt bestimmt die Klammer der linksinvarianten Vektorfelder die Lie-Klammer in $T_eG$ durch Auswertung in $e$.



Linke Maurer–Cartan-Form

Die \emph{linke Maurer–Cartan-Form} ist die $T_eG$-wertige $1$-Form
\[
\omega^L \;:\; TG \longrightarrow T_eG,\qquad
\omega^L_g \;:=\; (L_{g^{-1}})_{*}\big|_g \;:\; T_gG \to T_eG.
\]
Sie ist linksinvariant: $(L_h)^{*}\omega^L = \omega^L$.




Maurer–Cartan-Gleichung

Die Form $\omega^L$ erfüllt die Strukturformel
\[
d\omega^L \;+\; \tfrac{1}{2}\,[\omega^L,\omega^L] \;=\; 0,
\]
wobei $[\omega^L,\omega^L]$ der durch die Lie-Klammer auf $T_eG$ induzierte bilineare Ausdruck ist.



$\omega^L$ realisiert die oben beschriebene Übertragung: Für $u_g\in T_gG$ ist $\omega^L_g(u_g)$ genau der „zurücktransportierte“ Vektor in $T_eG$. Die Maurer–Cartan-Gleichung kodiert die Konsistenz dieser Parallelisierung mit der Lie-Klammer (sie ist die Differentialform-Version der Jacobi-Identität).




Genügsamkeit des Jets in der Einheit


Die $1$- und $2$-Jet-Daten der Multiplikation in $(e,e)$ (also die Vektorraumstruktur auf $T_eG$ und die Lie-Klammer $[\cdot,\cdot]$) bestimmen:
\begin{enumerate}
  \item die Klammer aller linksinvarianten Vektorfelder,
  \item die faserweisen Klammern $[\cdot,\cdot]_{(g)}$ auf allen $T_gG$ mittels linkem Transport,
  \item die adjungierte Wirkung $g\mapsto \mathrm{Ad}_g$ (als linearisierte Konjugation),
  \item und (lokal um $e$) die Gruppenmultiplikation bis zu höheren Ordnungen (Baker–Campbell–Hausdorff).
\end{enumerate}



(1) folgt aus $[\widetilde{X},\widetilde{Y}]=\widetilde{[X,Y]}$.  
(2) folgt aus Definition der transportierten Klammer via $(L_g)_{*}$.  
(3) folgt aus $\mathrm{Ad}_g = T_e( h\mapsto ghg^{-1} )$ und der durch die Lie-Algebra bestimmten Klammer; die Beziehung $\mathrm{ad}_X=\left.\frac{d}{dt}\right|_{0}\mathrm{Ad}_{\exp(tX)}$ bindet beides.  
(4) folgt aus der BCH-Entwicklung $Z(X,Y)=X+Y+\frac12[X,Y]+\cdots$.



Begriff „Übertragung auf alle Tangentialräume“

Formal bedeutet dies: Es existieren \emph{kanonische Isomorphismen der Fasern} $(L_g)_{*}\big|_e:T_eG\to T_gG$, mit denen jede in $T_eG$ definierte algebraische Struktur (insbesondere $[\cdot,\cdot]$) \emph{funktoriell} in \emph{jedem} Punkt $g$ als Struktur auf $T_gG$ reproduziert wird. Diese Kohärenz ist genau die linke Invarianz.




Einparameteruntergruppen und Exponentialabbildung

Für jedes $X\in T_eG$ existiert eine eindeutig bestimmte Einparameteruntergruppe $\gamma_X:\mathbb{R}\to G$ mit $\gamma_X(0)=e$ und $\dot{\gamma}_X(0)=X$. Die Exponentialabbildung ist $\exp(X):=\gamma_X(1)$. Die Relationen
\[
\exp(tX)\exp(sY) = \exp\bigl(Z(tX,sY)\bigr),
\qquad
Z = \text{BCH}(X,Y),
\]
zeigen, wie die (lokal) globale Multiplikation aus $(T_eG,[\cdot,\cdot])$ rekonstruiert wird.


Konzeptuelle Quintessenz

Die Lie-Algebra $(T_eG,[\cdot,\cdot])$ ist die \emph{komprimierte} infinitesimale Information; linke Translation und die Maurer–Cartan-Form sind die \emph{Transportmechanismen}, die diese Information \emph{ohne Zusatzwahl} auf alle anderen Punkte und Tangentialräume ausbreiten. 


\begin{itemize}
  \item Die Jet-Charakterisierung zeigt, dass die \emph{niedrigste nichttriviale} infinitesimale Information der Gruppenmultiplikation im Punkt $(e,e)$ exakt die Lie-Klammer auf $T_eG$ ist.
  \item Die linke Translation liefert kanonische Isomorphismen $T_eG\cong T_gG$ und transportiert die Klammer auf jede Faser, kompatibel zur Geometrie.
  \item Linksinvariante Vektorfelder sind die globalisierte Form der Daten in $T_eG$; ihre Klammer ist die globalisierte Lie-Klammer.
  \item Die Maurer–Cartan-Form realisiert diesen Transport differenzialgeometrisch und erfüllt die Maurer–Cartan-Gleichung als Integrabilitätsbedingung.
\end{itemize}






\pagebreak


\section{Funktoren der globalen Analysis}


\begin{DEF}{GA-FUN-01}{Tangentialfunktor}
Sei $\mathbf{Man}$ die Kategorie glatter Mannigfaltigkeiten mit glatten Abbildungen.  
Der \emph{Tangentialfunktor}
\[
T:\mathbf{Man}\longrightarrow \mathbf{VBun}
\]
weist jedem Objekt $M$ das Tangentialbündel $TM\xrightarrow{\;\pi\;}M$ zu und jedem Morphismus
$f:M\to N$ den Bündelmorphismus
\[
Tf:TM\to TN,\qquad v_p\mapsto T_p f\cdot v_p,
\]
wobei $T_pf:T_pM\to T_{f(p)}N$ die totale Ableitung ist.
\medskip

\textbf{Eigenschaften.}
\begin{enumerate}
\item $T(\mathrm{id}_M)=\mathrm{id}_{TM}$.
\item $T(g\circ f)=Tg\circ Tf$.
\item $T$ ist kovariant und mit Faserprojektionen verträglich: $\pi_N\circ Tf=f\circ \pi_M$.
\end{enumerate}
\end{DEF}

\begin{DEF}{GA-FUN-02}{Kotangentialfunktor}
Der \emph{Kotangentialfunktor}
\[
T^{*}:\mathbf{Man}\longrightarrow \mathbf{VBun}
\]
ordnet $M$ das Kotangentialbündel $T^{*}M\to M$ zu und $f:M\to N$ den Bündelmorphismus
\[
T^{*}f:T^{*}N\to T^{*}M,\qquad \alpha_{f(p)}\mapsto (T_pf)^{*}\alpha_{f(p)},
\]
also die duale Abbildung $(T_pf)^{*}:T_{f(p)}^{*}N\to T_p^{*}M$ punktweise.
\medskip

\textbf{Eigenschaften.} $T^{*}$ ist kontravariant auf der Ebene der Fasern, aber als Bündelmorphismus \emph{kovariant} über $f$:
$\pi_M\circ T^{*}f = f^{-1}\circ \pi_N$ (als Abbildung der Gesamt\-räume).
Funktorial: $T^{*}(\mathrm{id})=\mathrm{id}$, $T^{*}(g\circ f)=T^{*}f\circ T^{*}g$.
\end{DEF}


\begin{DEF}{GA-FUN-03}{Tensorbündel-Funktoren}
Für $r,s\in\mathbb{N}$ definiert der \emph{Tensorfunktor}
\[
T^{r}_{\;s}:\mathbf{Man}\longrightarrow \mathbf{VBun},\qquad 
M\longmapsto T^{r}_{\;s}(TM):=\underbrace{TM\otimes\cdots\otimes TM}_{r}\otimes
\underbrace{T^{*}M\otimes\cdots\otimes T^{*}M}_{s},
\]
und für $f:M\to N$
\[
T^{r}_{\;s}(f):T^{r}_{\;s}(TM)\longrightarrow T^{r}_{\;s}(TN)
\]
als die faserweise funktorielle Kombination von $Tf$ und $T^{*}f$:
kontravariant in den $s$ kotangentialen und kovariant in den $r$ tangentialen Argumenten.
\medskip

\textbf{Eigenschaften.} $T^{r}_{\;s}$ ist ein kovarianter Funktor, verträglich mit Tensorprodukt, Dualbildung und Pullback/Pushforward auf der Faser.
\end{DEF}


\begin{DEF}{GA-FUN-04}{Jet-Funktor}
Für $k\in\mathbb{N}$ ordnet der \emph{$k$-Jet-Funktor}
\[
J^{k}:\mathbf{Man}\longrightarrow \mathbf{FibBun}
\]
der Mannigfaltigkeit $M$ das Bündel $J^{k}(M,N)$ der $k$-Jets glatter Abbildungen nach $N$ (bei festem Ziel $N$) zu; für $f:M\to M'$ gibt es den induzierten Morphismus
\[
J^{k}f: J^{k}(M,N)\to J^{k}(M',N),\qquad j^{k}_{p}\phi\mapsto j^{k}_{f(p)}(\phi\circ f^{-1})
\]
(geeignet präzisiert über geeignete Kategorien/Unterkategorien, etwa lokale Diffeomorphismen).
\medskip

\textbf{Eigenschaften.} $J^{k}$ bewahrt Komposition und Identitäten; er kodiert Ableitungen bis zur Ordnung $k$ und ist grundlegend in der Geometrie partieller Differentialgleichungen.
\end{DEF}

\begin{DEF}{GA-FUN-05}{Differentialformen-Funktor}
Für jedes $k\ge 0$ definiert
\[
\Omega^{k}:\mathbf{Man}\longrightarrow \mathbf{CAlg}^{\mathrm{op}}
\]
den kontravarianten Funktor, der $M\mapsto \Omega^{k}(M)$ und für $f:M\to N$ den Pullback
\[
f^{*}:\Omega^{k}(N)\to \Omega^{k}(M)
\]
zuordnet, gegeben durch $(f^{*}\eta)_p(v_{1},\dots,v_{k})=\eta_{f(p)}(T_pf\,v_{1},\dots,T_pf\,v_{k})$.
Für $\Omega=\bigoplus_{k}\Omega^{k}$ ist $f^{*}$ ein Ringhomomorphismus bzgl.\ $\wedge$.
\end{DEF}


\begin{DEF}{GA-FUN-06}{Algebrafunktor glatter Funktionen}
Der kontravarante \emph{Funktional-Funktor}
\[
C^{\infty}:\mathbf{Man}\longrightarrow \mathbf{CAlg}^{\mathrm{op}},
\qquad M\mapsto C^{\infty}(M),
\]
ordnet $f:M\to N$ den Pullback $f^{*}:C^{\infty}(N)\to C^{\infty}(M)$, $f^{*}(h)=h\circ f$, zu.
Er ist treu und spiegelt (viel von) der glatten Struktur.
\end{DEF}


\begin{DEF}{GA-FUN-07}{Lie-Algebra-Funktor}
Sei $\mathbf{LieGrp}$ die Kategorie der Lie-Gruppen und glatten Gruppenhomomorphismen.  
Der \emph{Lie-Algebra-Funktor}
\[
\mathfrak{L}:\mathbf{LieGrp}\longrightarrow \mathbf{LieAlg},\qquad 
G\mapsto \mathfrak{g}:=T_{e}G,
\]
ordnet einem Homomorphismus $\Phi:G\to H$ die abgeleitete Algebraabbildung
\[
d\Phi:\mathfrak{g}\to \mathfrak{h},\qquad v\mapsto T_{e}\Phi\cdot v,
\]
zu.  
\textbf{Eigenschaften.} $\mathfrak{L}$ ist ein (starker) Funktor: $d(\Psi\circ\Phi)=d\Psi\circ d\Phi$, $d(\mathrm{id})=\mathrm{id}$; der Lie-Klammer verträglich über linksinvariante Felder.
\end{DEF}


\begin{DEF}{GA-FUN-08}{Rahmenbündel-Funktor}
Der \emph{Frame-Funktor}
\[
\mathrm{Fr}:\mathbf{Man}_{n}\longrightarrow \mathbf{Prin}_{GL(n)}
\]
ordnet einer $n$-dimensionalen Mannigfaltigkeit $M$ ihr \emph{Rahmenbündel} $\mathrm{Fr}(M)\to M$ zu, dessen Faser am Punkt $p$ die Menge linearer Isomorphismen $u:\mathbb{R}^{n}\to T_{p}M$ ist. Eine glatte Abbildung $f:M\to N$ induziert
\[
\mathrm{Fr}(f):\mathrm{Fr}(M)\to \mathrm{Fr}(N),\qquad u\mapsto T_pf\circ u,
\]
wobei $GL(n)$ rechts regulär wirkt. $\mathrm{Fr}$ ist kovariant in $\mathbf{Prin}_{GL(n)}$.
\end{DEF}


\begin{DEF}{GA-FUN-09}{Assoziierter-Bündel-Funktor}
Sei $\mathbf{Prin}_{G}$ die Kategorie glatter $G$-Hauptbündel und $G\to GL(V)$ eine Darstellung $\rho$.  
Der Funktor
\[
\mathrm{Assoc}_{\rho}:\mathbf{Prin}_{G}\longrightarrow \mathbf{VBun},\qquad 
P\mapsto P\times_{\rho} V
\]
ordnet einem Hauptbündel $P\to M$ das assoziierte Vektorbündel zu (Orbitraum des $G$-Rechts\-wirkens $(p,v)\cdot g=(pg,\rho(g^{-1})v)$).  
Ein Morphismus $F:P\to P'$ über $f:M\to M'$ induziert
\[
\mathrm{Assoc}_{\rho}(F):P\times_{\rho}V\to P'\times_{\rho}V,\quad [p,v]\mapsto [F(p),v].
\]
\end{DEF}


\begin{DEF}{GA-FUN-10}{Pullback-Funktor für Vektorbündel}
Für $f:M\to N$ definiert der \emph{Bündel-Pullback-Funktor}
\[
f^{*}:\mathbf{VBun}(N)\longrightarrow \mathbf{VBun}(M),\qquad 
E\xrightarrow{\;\pi\;}N\ \mapsto\ f^{*}E:=\{(p,e)\in M\times E:\ f(p)=\pi(e)\}
\]
mit Projektion $(p,e)\mapsto p$.  
Ein Bündelmorphismus $\Phi:E\to E'$ über $\mathrm{id}_{N}$ induziert $f^{*}\Phi:f^{*}E\to f^{*}E'$ durch $(p,e)\mapsto (p,\Phi(e))$.  
\textbf{Eigenschaften.} Funktorialität in $E$ und $f$; Faserweise ist $(f^{*}E)_p\simeq E_{f(p)}$.
\end{DEF}

\begin{DEF}{GA-FUN-11}{(Faser-)Pushforward/Integration entlang der Faser}
Sei $\pi:E\to B$ eine glatte surjektive Submersion mit kompakten Fasern, $E,B$ orientiert.  
Dann definiert \emph{Integration entlang der Faser} einen Funktor
\[
\pi_{*}:\Omega^{k+\dim\mathrm{Faser}}(E)\longrightarrow \Omega^{k}(B),
\]
eindeutig charakterisiert durch
\[
\int_{B}\beta\wedge \pi_{*}\alpha=\int_{E}\pi^{*}\beta\wedge \alpha\qquad(\beta\in\Omega^{\mathrm{top}-k}(B)).
\]
Für einen glatten, properen Morphismus zwischen geeigneten Faserbündeln ist $\pi_{*}$ kovariant und verträglich mit $d$ (bis auf Vorzeichen, Stokes).
\end{DEF}


\begin{DEF}{GA-FUN-12}{de Rham-Komplex- und Kohomologiefunktor}
Der \emph{de Rham-Funktor}
\[
\mathcal{R}:\mathbf{Man}^{\mathrm{op}}\longrightarrow \mathbf{DGA},\qquad 
M\mapsto (\Omega^{\bullet}(M),d,\wedge)
\]
ordnet Pullbacks $f^{*}$ zu, die mit $d$ kommutieren: $d\circ f^{*}=f^{*}\circ d$.  
Durch Quotientieren nach exakten Formen entsteht der \emph{de Rham-Kohomologiefunktor}
\[
H^{\bullet}_{\mathrm{dR}}:\mathbf{Man}^{\mathrm{op}}\longrightarrow \mathbf{GrVect},\qquad 
M\mapsto H^{\bullet}_{\mathrm{dR}}(M),
\]
mit $f^{*}:H^{\bullet}_{\mathrm{dR}}(N)\to H^{\bullet}_{\mathrm{dR}}(M)$.
\end{DEF}


\begin{DEF}{GA-FUN-13}{Hodge-\texorpdfstring{$\star$}{*}-Funktor (metrisch orientiert)}
Für eine orientierte Riemannsche Mannigfaltigkeit $(M,g)$ definiert der \emph{Hodge-Stern-Funktor}
\[
\star_{g}:\Omega^{k}(M)\longrightarrow \Omega^{n-k}(M)
\]
punktweise als den eindeutigen Isomorphismus mit
\[
\alpha\wedge \star_{g}\beta = \langle \alpha,\beta\rangle_{g}\,\mathrm{vol}_{g}.
\]
Für einen isometrischen Diffeomorphismus $f:(M,g)\to (N,h)$ gilt $f^{*}\circ \star_{h}=\star_{g}\circ f^{*}$; damit ist $\star$ natürlich in der Kategorie orientierter Riemannscher Mannigfaltigkeiten und Isometrien.
\end{DEF}


\begin{DEF}{GA-FUN-14}{Laplace–Beltrami-Zuordnung}
Der \emph{Laplace–Beltrami-Funktor} ordnet einer orientierten Riemannschen Mannigfaltigkeit $(M,g)$ den (formell) selbstadjungierten elliptischen Operator
\[
\Delta_{g}:=\mathrm{d}\,\delta+\delta\,\mathrm{d}:\Omega^{k}(M)\to\Omega^{k}(M),
\qquad \delta:=(-1)^{nk+n+1}\star\,\mathrm{d}\,\star,
\]
zu.  
Ein isometrischer Diffeomorphismus $f:(M,g)\to(N,h)$ erfüllt $f^{*}\circ \Delta_{h}=\Delta_{g}\circ f^{*}$; damit entsteht ein Funktor in die Kategorie linearer Operatoren modulo unitärer Äquivalenz.
\end{DEF}

\begin{DEF}{GA-FUN-15}{Index-Funktor (Atiyah–Singer)}
Sei $\mathbf{Ell}$ die Kategorie elliptischer, Fredholm-Operatoren zwischen Vektorbündelsektionen über kompakten Mannigfaltigkeiten (Morphismen: Bündel\-isomorphismen/Konkordanz).  
Der \emph{Index-Funktor}
\[
\mathrm{ind}:\mathbf{Ell}\longrightarrow \mathbf{Ab},
\qquad D\mapsto \dim\ker D-\dim\mathrm{coker}\,D\in\mathbb{Z},
\]
ist invariant unter homotopen Deformationen von $D$ und faktorisert durch die topologische $K$-Theorie des Baseraums.
\end{DEF}


\begin{DEF}{GA-FUN-16}{Integrationsfunktor (Volumenformen)}
Für eine orientierte, kompakte $n$-Mannigfaltigkeit $M$ definiert
\[
\int_{M}:\Omega^{n}(M)\longrightarrow \mathbb{R},\qquad \omega\mapsto \text{das Integral von }\omega,
\]
einen natürlichen, kovarianten Funktor auf der Unterkategorie orientierter Diffeomorphismen, d.\,h.\ $\int_{M} f^{*}\eta=\int_{N}\eta$ für orientierungsbewahrende Diffeomorphismen $f:M\to N$.
\end{DEF}


\section{Lie-Theorie und mathematische Physik}




\begin{DEF}{LIE-L09-08-01}{Forminvarianz von Operatorgleichungen unter einer Lie-Gruppenwirkung}
Sei $G$ eine Lie-Gruppe mit Darstellung 
\[
U: G \to \mathrm{Aut}(V)
\]
auf einem topologischen Vektorraum $V$ (typischerweise $V = C^\infty(M)$, $L^2(M)$ oder $\mathcal{S}(M)$, 
je nach Regularität der Felder oder Zustände).  
Sei ferner 
\[
E: \mathrm{Dom}(E) \subseteq V \to V
\]
ein (linearer) Operator mit dichtem Definitionsbereich $\mathrm{Dom}(E)$ in $V$ 
(z.\,B. ein Differentialoperator, ein Pseudodifferentialoperator, 
oder der Generator einer zeitlichen Entwicklung).

Dann heißt die Gleichung
\[
E[\psi] = 0
\]
\emph{forminvariant unter $G$}, wenn für alle $g \in G$ gilt:
\[
E[U(g)\psi] = U(g)\,E[\psi]
\quad \text{für alle } \psi \in \mathrm{Dom}(E).
\]
Äquivalent:
\[
E\,U(g) = U(g)\,E \quad \text{auf } \mathrm{Dom}(E).
\]

Ist $E$ dicht definiert und die Darstellung $U$ stark stetig, 
so ist dies äquivalent zur infinitesimalen Invarianzbedingung auf der Lie-Algebra-Ebene:
\[
[dU(X), E] = 0, \qquad \forall X \in \mathfrak{g},
\]
wobei $dU(X)$ der infinitesimale Generator der Darstellung $U$ ist:
\[
dU(X) := \left.\frac{d}{dt}\right|_{t=0} U(\exp(tX)).
\]
\end{DEF}








\begin{DEF}{LIE-L09-08-02}{Lorentz-Invarianz von Differentialgleichungen}
Sei $E$ ein linearer Differentialoperator auf $\mathcal{S}(\mathbb{R}^{1,3}, \mathbb{C}^N)$ 
oder $C^\infty(\mathbb{R}^{1,3}, \mathbb{C}^N)$, 
und sei $G = SO(1,3)$ die Lorentz-Gruppe 
mit Darstellung
\[
U(\Lambda)\psi(x) := D(\Lambda)\psi(\Lambda^{-1}x),
\]
wobei $D: SO(1,3) \to GL(\mathbb{C}^N)$ eine (möglicherweise nichttriviale) endliche oder unendliche Darstellung ist
(z.\,B. Vektordarstellung, Spinordarstellung, Tensor-Darstellung).

Dann heißt die Gleichung
\[
E[\psi] = 0
\]
\emph{Lorentz-invariant}, wenn für alle $\Lambda \in SO(1,3)$ gilt:
\[
E[U(\Lambda)\psi] = U(\Lambda)\,E[\psi],
\]
d.\,h. $E$ kommutiert mit der Lorentz-Wirkung. 
Äquivalent gilt auf der infinitesimalen Ebene:
\[
[dU(X),E] = 0 \quad \forall X \in \mathfrak{so}(1,3).
\]
\end{DEF}





\begin{EXA}{LIE-L09-08-03}{Forminvarianz unter Translationen}
Sei $M=\mathbb{R}^{1,3}$ und $G=(\mathbb{R}^{1,3},+)$ die Translationsgruppe mit Darstellung
\[
(U(a)\psi)(x) = \psi(x-a).
\]
Für den Operator
\[
E = \Box + m^2, \qquad \Box = \eta^{\mu\nu}\partial_\mu\partial_\nu,
\]
gilt:
\[
E[U(a)\psi](x) = (\Box + m^2)\psi(x-a) = (U(a)E\psi)(x).
\]
Also ist $E[\psi]=0$ forminvariant unter Raum-Zeit-Translationen.
\end{EXA}



\begin{EXA}{LIE-L09-08-04}{Forminvarianz unter Rotationen}
Sei $M=\mathbb{R}^3$ und $G=SO(3)$ mit Darstellung
\[
(U(R)\psi)(\mathbf{x}) = \psi(R^{-1}\mathbf{x}).
\]
Für den Operator
\[
E = -\Delta + V(|\mathbf{x}|),
\]
gilt:
\[
E[U(R)\psi](\mathbf{x}) = -\Delta_\mathbf{x}[\psi(R^{-1}\mathbf{x})] + V(|\mathbf{x}|)\psi(R^{-1}\mathbf{x})
= (U(R)E\psi)(\mathbf{x}),
\]
da der Laplace-Operator $\Delta$ und das Potential $V(|\mathbf{x}|)$ rotationsinvariant sind.
Somit ist die Schrödingergleichung mit zentralem Potential forminvariant unter $SO(3)$.
\end{EXA}



\begin{EXA}{LIE-L09-08-05}{Forminvarianz unter Zeittranslationen}
Sei $G=(\mathbb{R},+)$ mit Darstellung
\[
(U(t)\psi)(x) = e^{-iHt/\hbar}\psi(x).
\]
Dann gilt für den Operator
\[
E = i\hbar \frac{\partial}{\partial t} - H,
\]
dass $E[U(t)\psi]=U(t)E[\psi]$ gilt, da $U(t)$ die Zeitentwicklung selbst beschreibt.
\end{EXA}





\begin{EXA}{LIE-L09-08-06}{Klein--Gordon-Gleichung}
Die Klein--Gordon-Gleichung
\[
E[\phi] := (\Box + m^2)\phi = 0, \qquad 
\Box = \eta^{\mu\nu}\partial_\mu\partial_\nu,
\]
ist Lorentz-invariant.
Für die Darstellung
\[
(U(\Lambda)\phi)(x) = \phi(\Lambda^{-1}x)
\]
gilt
\[
E[U(\Lambda)\phi](x)
= \eta^{\mu\nu}\partial_\mu\partial_\nu \phi(\Lambda^{-1}x) + m^2\phi(\Lambda^{-1}x)
= (U(\Lambda)E\phi)(x),
\]
da $\eta^{\mu\nu}\Lambda_\mu^{\ \rho}\Lambda_\nu^{\ \sigma} = \eta^{\rho\sigma}$.
\end{EXA}



\begin{EXA}{LIE-L09-08-07}{Dirac-Gleichung}
Die Dirac-Gleichung
\[
E[\psi] := (i\gamma^\mu \partial_\mu - m)\psi = 0
\]
ist Lorentz-invariant, wenn $\psi$ sich nach der Spinor-Darstellung transformiert:
\[
(U(\Lambda)\psi)(x) = S(\Lambda)\psi(\Lambda^{-1}x),
\]
wobei $S(\Lambda)$ die Spinor-Darstellung von $SO(1,3)$ ist, die
\[
S^{-1}(\Lambda)\gamma^\mu S(\Lambda) = \Lambda^\mu_{\ \nu}\gamma^\nu
\]
erfüllt.
Dann gilt:
\[
E[U(\Lambda)\psi](x)
= i\gamma^\mu \partial_\mu \big(S(\Lambda)\psi(\Lambda^{-1}x)\big) - mS(\Lambda)\psi(\Lambda^{-1}x)
= S(\Lambda)(i\gamma^\mu \partial_\mu - m)\psi(\Lambda^{-1}x)
= (U(\Lambda)E\psi)(x).
\]
\end{EXA}



\begin{EXA}{LIE-L09-08-08}{Maxwell-Gleichungen}
Die homogenen und inhomogenen Maxwell-Gleichungen
\[
\partial_{[\lambda}F_{\mu\nu]} = 0, 
\qquad 
\partial_\mu F^{\mu\nu} = j^\nu,
\]
sind Lorentz-invariant, wenn das elektromagnetische Feld $F_{\mu\nu}$ als antisymmetrischer Tensor transformiert:
\[
F'_{\mu\nu}(x') = \Lambda_\mu^{\ \rho}\Lambda_\nu^{\ \sigma}F_{\rho\sigma}(x),
\quad x' = \Lambda x.
\]
Da $\partial_\mu \mapsto \Lambda_\mu^{\ \rho}\partial_\rho$ und die Metrik invariant bleibt,
erhalten die Maxwell-Gleichungen dieselbe Form in allen Lorentz-Systemen.
\end{EXA}




\begin{DEF}{LIE-L09-08-09}{Invarianz und Darstellungsstruktur des Lösungsraums}
Sei $G$ eine Lie-Gruppe mit Darstellung
\[
U: G \to \mathrm{Aut}(V),
\]
wobei $V$ ein topologischer Vektorraum ist (z.\,B.\ $C^\infty(M)$, $\mathcal{S}(M)$ oder $L^2(M)$).  
Sei ferner
\[
E : \mathrm{Dom}(E) \subseteq V \to V
\]
ein linearer Operator mit dichtem Definitionsbereich $\mathrm{Dom}(E)$.

Angenommen, $E$ ist \emph{forminvariant} unter der Wirkung von $G$, d.\,h.
\[
E\,U(g) = U(g)\,E
\qquad \forall g \in G.
\]

Dann gilt:

\begin{enumerate}
\item Der Lösungsraum
\[
\ker(E) := \{ \psi \in \mathrm{Dom}(E) \mid E[\psi] = 0 \}
\]
ist \emph{invariant unter der Wirkung von $G$}, d.\,h.
\[
U(g)\psi \in \ker(E)
\qquad \forall g \in G, \ \forall \psi \in \ker(E).
\]
\item Die Einschränkung der Darstellung
\[
U|_{\ker(E)} : G \to \mathrm{Aut}(\ker(E)), \qquad
U|_{\ker(E)}(g) := U(g)|_{\ker(E)}
\]
definiert eine Darstellung von $G$ auf dem Lösungsraum.
\item Damit wirkt die Symmetriegruppe $G$ auf dem Lösungsraum $\ker(E)$, und
\[
\ker(E)
\]
ist ein $G$-invarianter Unterraum von $V$, der sich (falls $V$ unitär ist) häufig als direkte Summe irreduzibler Darstellungen zerlegen lässt:
\[
\ker(E) \cong \bigoplus_{\alpha \in \mathcal{I}} V_\alpha,
\]
wobei jedes $V_\alpha$ eine irreduzible Darstellung von $G$ ist.
\end{enumerate}

Anschaulich bedeutet dies:
Ist $\psi$ eine Lösung der Gleichung $E[\psi]=0$, so ist für jedes Gruppenelement $g\in G$ auch $U(g)\psi$ eine Lösung.  
Damit bildet die Symmetriegruppe $G$ auf dem Raum der Lösungen eine Darstellung, d.\,h. eine strukturtreue lineare Abbildung ihrer Symmetrieoperationen.
\end{DEF}





\begin{DEF}{LIE-L09-08-10}{Orbits und Orbitraum im Lösungsraum einer forminvarianten Gleichung}
Sei $G$ eine Lie-Gruppe mit Darstellung 
\[
U: G \to \mathrm{Aut}(V),
\]
und sei $E : \mathrm{Dom}(E) \subseteq V \to V$ ein linearer Operator, 
der forminvariant unter $G$ ist, also $EU(g) = U(g)E$ für alle $g \in G$.

Dann wirkt $G$ auf dem Lösungsraum
\[
\ker(E) = \{ \psi \in \mathrm{Dom}(E) \mid E[\psi] = 0 \}
\]
durch
\[
G \times \ker(E) \to \ker(E), \qquad (g,\psi) \mapsto U(g)\psi.
\]

Für eine Lösung $\psi \in \ker(E)$ heißt
\[
\mathcal{O}_\psi := \{ U(g)\psi \mid g \in G \}
\]
der \emph{Orbit von $\psi$} unter der Gruppenwirkung $U$.

Der \emph{Stabilisator} (Isotropiegruppe) von $\psi$ ist
\[
G_\psi := \{ g \in G \mid U(g)\psi = \psi \}.
\]

Der \emph{Orbitraum} der Wirkung von $G$ auf dem Lösungsraum ist der Quotientenraum
\[
\ker(E)/G := \{ \mathcal{O}_\psi \mid \psi \in \ker(E) \},
\]
der die Menge der durch $G$-Äquivalenz bestimmten physikalisch unterscheidbaren Lösungen beschreibt.
\end{DEF}



\begin{DEF}{LIE-L09-08-11}{Eigenwertproblem und Symmetrieinvarianz}
Sei $G$ eine Lie-Gruppe mit Darstellung 
\[
U: G \to \mathrm{Aut}(V),
\]
und sei $E : \mathrm{Dom}(E) \subseteq V \to V$ ein linearer Operator, der \emph{forminvariant} unter $G$ ist:
\[
E\,U(g) = U(g)\,E \quad \forall g \in G.
\]

Dann gilt:
\begin{enumerate}
\item Für jedes Eigenpaar $(\lambda, \psi)$, d.\,h.\ $E[\psi] = \lambda \psi$, ist $U(g)\psi$ ebenfalls Eigenvektor zum selben Eigenwert:
\[
E[U(g)\psi] = \lambda\,U(g)\psi.
\]
\item Damit ist der Eigenraum
\[
V_\lambda := \{ \psi \in \mathrm{Dom}(E) \mid E[\psi] = \lambda \psi \}
\]
invariant unter $G$:
\[
U(g)V_\lambda \subseteq V_\lambda.
\]
\item Die Einschränkung $U|_{V_\lambda}: G \to \mathrm{Aut}(V_\lambda)$ definiert eine Darstellung von $G$ auf dem Eigenraum $V_\lambda$.
\end{enumerate}
\end{DEF}



\begin{DEF}{LIE-L09-08-12}{Orbits und physikalische Äquivalenz}
Sei $G$ eine Lie-Gruppe, die auf dem Lösungsraum
\[
\ker(E)=\{\psi\mid E[\psi]=0\}
\]
durch $U(g)\psi$ wirkt. Für $\psi\in\ker(E)$ sei
\[
\mathcal{O}_\psi=\{U(g)\psi\mid g\in G\}
\]
der Orbit.

\begin{enumerate}
\item \textbf{Physikalische Symmetrien:}  
Wenn $G$ eine externe Symmetriegruppe des Gesetzes ist (z.\,B.\ Raumzeit-, Rotations- oder Lorentzgruppe),
dann beschreiben alle $\psi'\in\mathcal{O}_\psi$ \emph{verschiedene Repräsentationen desselben physikalischen Prozesses}.
Der Orbit strukturiert den Lösungsraum in Bahnen von Zuständen mit denselben invarianten Größen
(Energie, Impuls, Spin, Ladung, …).

\item \textbf{Eichsymmetrien:}  
Wenn $G$ eine Eichgruppe ist (d.\,h.\ die Transformationen wirken nur auf die Repräsentation, nicht auf beobachtbare Größen),
dann beschreiben alle $\psi'\in\mathcal{O}_\psi$ \emph{denselben physikalischen Zustand}.
Die Menge der physikalisch unterscheidbaren Lösungen ist dann der Quotient
\[
\mathcal{M} = \ker(E)/G,
\]
der \emph{Modulraum der Lösungen modulo Gauge}.
\end{enumerate}
\end{DEF}



\begin{DEF}{LIE-L09-08-13}{Physikalische Symmetrien vs. Eichsymmetrien}
Sei $\mathcal{C}$ der Konfigurationsraum eines physikalischen Systems
(z.\,B.\ Felder, Wellenfunktionen, Metriken),
und $G$ eine Lie-Gruppe, die auf $\mathcal{C}$ wirkt:
\[
\Phi: G\times \mathcal{C}\to \mathcal{C},\qquad (g,\psi)\mapsto g\cdot\psi.
\]
Sei ferner $\mathcal{E}[\psi]=0$ das dynamische Gesetz
(z.\,B.\ Feld- oder Wellengleichung) auf $\mathcal{C}$.

\begin{enumerate}
\item \textbf{Physikalische Symmetrie (externe Symmetrie).}\\[2mm]
$G$ heißt \emph{physikalische Symmetriegruppe} des Gesetzes, wenn
\[
\mathcal{E}[g\cdot\psi]=0\ \Longleftrightarrow\ \mathcal{E}[\psi]=0,
\]
aber die Wirkungen $g\cdot\psi$ und $\psi$ beschreiben \emph{physikalisch verschiedene Konfigurationen},
die jedoch demselben Gesetz unterliegen.
Der Lösungsraum $\ker(\mathcal{E})$ wird dann in $G$-Orbits zerlegt,
deren Punkte $\psi$ und $g\cdot\psi$ mathematisch verschieden,
aber physikalisch korrespondierend sind.
\[
G \curvearrowright \ker(\mathcal{E})\quad\text{mit}\quad \mathcal{O}_\psi=\{g\cdot\psi\mid g\in G\}.
\]
Hier bleibt der physikalische Zustand identisch nur bis auf Wahl von Koordinaten
oder Bezugssystem; es wird \emph{nicht} nach $G$ quotientiert.

\item \textbf{Eichsymmetrie (interne Redundanz).}\\[2mm]
$G$ heißt \emph{Eichgruppe} des Systems, wenn
\[
\mathcal{E}[g\cdot\psi]=0\ \Longleftrightarrow\ \mathcal{E}[\psi]=0
\quad\text{und}\quad g\cdot\psi\ \text{repräsentiert dieselbe physikalische Konfiguration}.
\]
Dann ist die Wirkung von $G$ \emph{redundant}:
verschiedene Punkte $\psi,\psi'\in\mathcal{C}$, die in demselben Orbit liegen,
beschreiben denselben physikalischen Zustand.
Der physikalische Konfigurationsraum ist der Quotient
\[
\mathcal{C}_\text{phys} = \mathcal{C}/G, \qquad 
\mathcal{M}_\text{phys} = \ker(\mathcal{E})/G.
\]
\end{enumerate}
\end{DEF}





\begin{DEF}{LIE-L09-08-14}{Strukturelle Lösungen von Operatorgleichungen}
\begin{enumerate}
\item \textbf{Intuitive Ebene: Lösung als Funktion auf Raum–Zeit}

Naive Lösung einer Operatorgleichung: Sei $M\simeq\mathbb{R}^{1,3}$ eine Koordinatenraumzeit mit Karte $x=(x^\mu)$
und $V\simeq\mathbb{C}^N$ ein fester Darstellungsraum.
Eine \emph{naive Lösung} einer linearen Operatorgleichung
\[
E[\psi]=0
\]
ist eine Abbildung $\psi\colon M\to V$, die $E[\psi](x)=0$ für alle $x\in M$ erfüllt, wobei
$E$ ein (lokaler) Differentialoperator ist, typischerweise $E=\sum_{|\alpha|\le m}A_\alpha(x)\partial^\alpha$ mit $A_\alpha(x)\in\mathrm{End}(V)$.

Grenze des naiven Bildes: Die Darstellung einer Lösung als Funktion $\psi(x)$ hängt von
\emph{(i)} der Wahl von Koordinaten auf $M$ und
\emph{(ii)} der Wahl einer Basis von $V$ ab.
Ein Wechsel des Beobachters (Koordinaten) oder des internen Bezugs (Basis in $V$)
ändert die \emph{Darstellung} derselben physikalischen Lösung.

\item \textbf{Struktureller Übergang: Bündel, Gruppenwirkung, Darstellung}

Prinzipal- und assoziiertes Bündel: Sei $M$ eine glatte Raumzeitmannigfaltigkeit, $G$ eine Lie-Gruppe und
$P\overset{\pi}{\to} M$ ein $G$-Prinzipalbündel. Sei
$D\colon G\to GL(W)$ eine (endliche) Darstellung auf einem komplexen Vektorraum $W$.
Das \emph{assoziierte Vektorbündel} ist
\[
E := P\times_D W \;\;\xrightarrow{\;\;\pi\;\;}\; M,
\]
mit Faser $E_x\simeq W$ für $x\in M$.

Schnitt als „koordinatenfreie“ Lösung: Ein \emph{(glatter) Schnitt} des Bündels $E$ ist eine glatte Abbildung
\[
\psi \in \Gamma(E) := \{\, \psi\colon M\to E \mid \pi\circ\psi=\mathrm{id}_M \,\}.
\]
Eine \emph{strukturelle Lösung} einer Gleichung ist ein Schnitt $\psi\in\Gamma(E)$, der eine
bündelwertige Gleichung $E[\psi]=0$ erfüllt (s.\,u. zur Kovarianz von $E$).

Symmetriegruppe und ihre Wirkung: Sei $G$ eine Lie-Gruppe, die
\emph{(i)} auf $M$ durch Diffeomorphismen $a\colon M\to M$ (externe Symmetrie) und
\emph{(ii)} auf den Fasern durch $D\colon G\to GL(W)$ (interne Struktur)
wirkt.
Die induzierte Wirkung auf Schnitten lautet
\[
U(g)\psi \;\;:\;\; x \longmapsto D(g)\,\psi\bigl(g^{-1}\!\cdot x\bigr)\;\in E_x.
\]

Koordinaten- und Darstellungswechsel: Die Information, \emph{wie} eine Lösung unter Beobachter- oder Basiswechsel transformiert,
ist \textbf{nicht} in $\psi(x)$ allein kodiert, sondern in der Paarung
\[
\bigl(\;D\colon G\to GL(W),\;\; U(g)\psi(x)=D(g)\psi(g^{-1}\!x)\;\bigr).
\]
Dies macht die „Abhängigkeit vom Koordinatensystem“ als \emph{strukturierte Transformationsregel} explizit.

\item \textbf{Kovarianter Operator und Forminvarianz}

Kovariante Ableitung und kovarianter Operator: Sei $E=P\times_D W$ ein assoziiertes Bündel und $\nabla$ eine zu $D$ kompatible kovariante Ableitung auf $E$ (d.\,h.\ eine Wahl von Zusammenhang, die die $G$-Wirkung respektiert).
Ein linearer Differentialoperator $E$ heißt \emph{kovariant}, wenn er aus $\nabla$ und $G$-invarianten
Tensoren (z.\,B.\ Metrik $\eta$, Gamma-Matrizen $\gamma^\mu$ in geeigneter Darstellung) gebaut ist.


Forminvarianz / Kovarianz: Die Gleichung $E[\psi]=0$ heißt \emph{forminvariant} (kovariant) unter $G$, wenn
\[
E[U(g)\psi] \;=\; U(g)\,E[\psi]\qquad\forall\, g\in G,\;\psi\in\Gamma(E).
\]
Äquivalent gilt infinitesimal für $X\in\mathfrak{g}$:
\[
[dU(X),E]=0.
\]


Skalar- und Spinorfelder:
\begin{enumerate}
\item \textbf{Skalarfeld:} $W=\mathbb{C}$, $D(g)=1$. Mit $\square=\eta^{\mu\nu}\nabla_\mu\nabla_\nu$ ist
$E=\square+m^2$ Lorentz-kovariant: $E[U(\Lambda)\phi]=U(\Lambda)E[\phi]$.
\item \textbf{Dirac-Spinor:} $W\simeq\mathbb{C}^4$, $D(\Lambda)=S(\Lambda)$ (Spinor-Darstellung). Mit
$E=i\gamma^\mu\nabla_\mu - m$ und $S^{-1}(\Lambda)\gamma^\mu S(\Lambda)=\Lambda^\mu_{\;\nu}\gamma^\nu$
ist $E$ Lorentz-kovariant.
\end{enumerate}


\item \textbf{Lösungsraum, Orbits, Beobachter- und Eichaspekt}

Lösungsraum und Gruppenwirkung: Der \emph{Lösungsraum} ist
\[
\mathcal{S} := \ker(E) \;=\; \{\, \psi\in\Gamma(E)\mid E[\psi]=0 \,\}.
\]
Die Darstellung $U\colon G\to\mathrm{Aut}(\Gamma(E))$ schränkt sich zu einer Wirkung
\[
U\colon G\times \mathcal{S}\to \mathcal{S},\qquad (g,\psi)\mapsto U(g)\psi
\]
ein (wegen Forminvarianz).

Orbits und physikalische Interpretation: Für $\psi\in\mathcal{S}$ ist der \emph{Orbit}
\[
\mathcal{O}_\psi \;:=\;\{\, U(g)\psi\mid g\in G \,\}.
\]
\begin{itemize}
\item \textbf{Externe (physikalische) Symmetrien:} $\psi$ und $U(g)\psi$ sind
verschiedene \emph{Darstellungen} desselben physikalischen Gesetzes (z.\,B.\ anderes Inertialsystem).
Der Quotient $\mathcal{S}/G$ wird \emph{nicht} als physikalischer Zustandsraum genommen.
\item \textbf{Eichsymmetrien (Redundanz):} $\psi$ und $U(g)\psi$ sind \emph{physikalisch identisch}.
Der \emph{physikalische Lösungsraum} ist der \emph{Modulraum}
\[
\mathcal{M}_\mathrm{phys}\;:=\;\mathcal{S}/G.
\]
\end{itemize}

\item \textbf{Wo die „Beobachter-/Darstellungsinformation“ formal liegt}

Träger der Transformationsinformation: Die Information, \emph{wie} sich eine Lösung unter Koordinaten-/Beobachterwechseln und internen Basiswechseln verhält, liegt in den Daten
\[
\bigl(M,\;P\to M,\;D\colon G\to GL(W),\;\nabla,\;U(g)\psi(x)=D(g)\psi(g^{-1}\!x)\bigr).
\]
Insbesondere:
\begin{enumerate}
\item $D$ kodiert die \emph{interne} Transformationsregel (Spin, Isospin, Farbe).
\item Die Wirkung von $G$ auf $M$ kodiert den \emph{externen} Koordinatenwechsel (z.\,B.\ Lorentz).
\item $\nabla$ (Zusammenhang) koppelt beide Ebenen zu einem \emph{kovarianten} Operator $E$.
\end{enumerate}
Begründung (skizziert): Für $U(g)\psi(x)=D(g)\psi(g^{-1}x)$ folgt aus der Kovarianz $E[U(g)\psi]=U(g)E[\psi]$,
dass $E$ aus $\nabla$ und $G$-invarianten Tensoren zusammengesetzt sein muss. Damit ist die
Transformationsregel von $\psi$ unter Wechsel der Karte (Wirkung auf $M$) und der internen Basis (Wirkung durch $D$)
struktureller Bestandteil der Definition von $(E,\psi)$, nicht ein äußerliches Anhängsel.


\item \textbf{Spektral- und Darstellungsstruktur}: Sei $A$ ein (im Wesentlichen) selbstadjungierter Observable-Operator auf $\Gamma(E)$ mit $[A,U(g)]=0$.
Für einen Eigenwert $\lambda$ ist der Eigenraum
\[
V_\lambda:=\{\psi\mid A\psi=\lambda\psi\}
\]
invariant unter $U$, d.\,h.\ $U(g)V_\lambda\subseteq V_\lambda$.
Somit trägt jeder $V_\lambda$ eine Darstellung von $G$; eine (maximale) Familie kommutierender Observablen (CSCO)
zerlegt $V_\lambda$ weiter in gemeinsame Eigenräume (Gewichte, Projektionen etc.).


\item \textbf{Zusammenfassung als „koordinatenfreies“ Lösungsobjekt}: Eine \emph{strukturelle Lösung} einer kovarianten Gleichung ist die Datenmenge
\[
\Bigl(M,\;P\to M,\;D,\;E,\;\nabla,\;\psi\in\Gamma(P\times_D W)\Bigr),
\]
so dass $E[\psi]=0$ und $E$ unter der durch $D$ und die $G$-Wirkung auf $M$ induzierten Darstellung $U$ forminvariant ist.
\emph{Physikalische Information} besteht aus
\begin{enumerate}
\item der \textbf{konkreten Konfiguration} $\psi$ (als Schnitt),
\item ihren \textbf{Transformationsgesetzen} $(U,D)$ (Beobachter-/Darstellungsebene),
\item der \textbf{kovarianten Dynamik} $E$ (Gesetzesstruktur).
\end{enumerate}
Bei Eichgruppen ist nur die Klasse $[\psi]\in\Gamma(E)/G$ physikalisch.
\end{enumerate}
\end{DEF}


\begin{DEF}{LIE-L09-08-15}{Struktureller Zugang Klein--Gordon Gleichung}
\begin{enumerate}

\item \textbf{Intuitive Ebene: Lösung als Funktion auf Raum–Zeit}

Naive Lösung einer Operatorgleichung:  
Sei $M\simeq\mathbb{R}^{1,3}$ mit globaler Karte $x=(x^\mu)$ und Minkowski-Metrik $\eta=\mathrm{diag}(-,+,+,+)$.  
Setze $V\simeq\mathbb{C}$ (Skalarwert).  
Die \emph{Klein–Gordon-Gleichung} lautet
\[
E[\phi]\;=\;(\Box + m^2)\,\phi \;=\;0, 
\qquad \Box := \eta^{\mu\nu}\partial_\mu\partial_\nu = -\partial_t^2+\Delta.
\]
Eine \emph{naive Lösung} ist eine Funktion $\phi:M\to\mathbb{C}$ mit $(\Box+m^2)\phi=0$.

Grenze des naiven Bildes:  
Die Darstellung $\phi(x)$ hängt von  
\emph{(i)} der Wahl von Koordinaten $x^\mu$ (Bezugssystem)  
und \emph{(ii)} einer (hier trivialen) Wahl des Werteraums ab.  
Lorentz-Boosts oder -Drehungen ändern die \emph{Darstellung} derselben physikalischen Lösung.


\item \textbf{Struktureller Übergang: Bündel, Gruppenwirkung, Darstellung}

Bündel und Darstellung:  
Für Skalarfelder ist das assoziierte Bündel trivial:
\[
E = M\times\mathbb{C} \;\xrightarrow{\;\pi\;}\; M.
\]
Die relevante Symmetriegruppe ist die \emph{Poincaré-Gruppe}
\[
\mathcal{P} = \mathbb{R}^{1,3}\rtimes SO^+(1,3),
\]
mit Elementen $(a,\Lambda)$ (Translation $a$, Lorentz-Transformation $\Lambda$).  
Die (externe) Wirkung auf $M$ ist $x\mapsto \Lambda x + a$.  
Für ein \emph{Skalarfeld} ist die Faserwirkung trivial: $D(\Lambda)\equiv 1$.

Wirkung auf Schnitten (Pullback-Darstellung):  
Für $\phi\in\Gamma(E)\simeq C^\infty(M,\mathbb{C})$ definiert
\[
\bigl(U(a,\Lambda)\phi\bigr)(x) = \phi\bigl(\Lambda^{-1}(x-a)\bigr).
\]
Dies ist die natürliche Pullback-Wirkung der Poincaré-Gruppe auf Skalarfeldern.

Koordinaten- und Darstellungswechsel:  
Die Information, \emph{wie} sich $\phi$ beim Beobachterwechsel transformiert, liegt
\emph{nicht} in $\phi(x)$, sondern in der Paarung $(U,\,D\equiv 1)$:
\[
U(a,\Lambda)\phi(x)=\phi\bigl(\Lambda^{-1}(x-a)\bigr).
\]


\item \textbf{Kovarianter Operator und Forminvarianz}

Kovarianter Operator:  
Im flachen Minkowski-Raum ist die Levi-Civita-Ableitung $\nabla_\mu=\partial_\mu$.  
Der Klein–Gordon-Operator ist
\[
E = \Box + m^2 = \eta^{\mu\nu}\nabla_\mu\nabla_\nu + m^2,
\]
zusammengesetzt aus $G$-invarianten Daten (Metrik $\eta$, trivialer Zusammenhang).

Forminvarianz / Kovarianz:  
Für alle $(a,\Lambda)\in\mathcal{P}$ und alle $\phi$ gilt
\[
E\bigl[U(a,\Lambda)\phi\bigr] = U(a,\Lambda)\,E[\phi].
\]
Beweis (Rechnung):  
Mit $(U(a,\Lambda)\phi)(x)=\phi(\Lambda^{-1}(x-a))$ ist
\[
\partial_\mu (U\phi)(x) 
= (\Lambda^{-1})_\mu{}^{\nu}\, (\partial_\nu\phi)\bigl(\Lambda^{-1}(x-a)\bigr),
\]
\[
\partial_\mu\partial_\rho (U\phi)(x) 
= (\Lambda^{-1})_\mu{}^{\nu}(\Lambda^{-1})_\rho{}^{\sigma}\,
(\partial_\nu\partial_\sigma\phi)\bigl(\Lambda^{-1}(x-a)\bigr).
\]
Daher
\[
\Box(U\phi)(x)
= \eta^{\mu\rho}(\Lambda^{-1})_\mu{}^{\nu}(\Lambda^{-1})_\rho{}^{\sigma}\,
(\partial_\nu\partial_\sigma\phi)\bigl(\Lambda^{-1}(x-a)\bigr)
= \eta^{\nu\sigma}(\partial_\nu\partial_\sigma\phi)\bigl(\Lambda^{-1}(x-a)\bigr),
\]
weil $\eta^{\mu\rho}(\Lambda^{-1})_\mu{}^{\nu}(\Lambda^{-1})_\rho{}^{\sigma}
=\eta^{\nu\sigma}$ für $\Lambda\in SO^+(1,3)$.  
Damit
\[
E[U\phi](x)=\bigl(\Box+m^2\bigr)\phi\bigl(\Lambda^{-1}(x-a)\bigr) = \bigl(U E[\phi]\bigr)(x).
\]

Infinitesimale Generatoren:  
Mit
\[
P_\mu := i\,\partial_\mu, 
\qquad 
M_{\mu\nu} := i\,(x_\mu\partial_\nu - x_\nu\partial_\mu)
\]
gilt
\[
[P_\mu,\,E]=0, \qquad [M_{\mu\nu},\,E]=0.
\]


\item \textbf{Lösungsraum, Orbits, Beobachter- und Eichaspekt}

Lösungsraum und Gruppenwirkung:  
\[
\mathcal{S} := \ker(E) = \{\, \phi\in C^\infty(M)\mid (\Box+m^2)\phi=0 \,\}.
\]
Die Darstellung $U:\mathcal{P}\to\mathrm{Aut}(C^\infty(M))$ schränkt sich zu
\[
U:\mathcal{P}\times\mathcal{S}\to \mathcal{S},\qquad (g,\phi)\mapsto U(g)\phi
\]
ein (wegen Forminvarianz).

Orbits:  
Für $\phi\in\mathcal{S}$ ist der Orbit
\[
\mathcal{O}_\phi := \{\, U(g)\phi \mid g\in\mathcal{P} \,\}.
\]
\emph{Interpretation:} Punkte eines Orbits sind \emph{verschiedene Darstellungen}
derselben Physik in unterschiedlichen Inertialsystemen (externe Symmetrie).  
Für die Klein–Gordon-Theorie gibt es keine Eichredundanz des Feldes selbst,
d.h. man quotientiert hier nicht nach $\mathcal{P}$.


\item \textbf{Spektral- und Darstellungsstruktur}

Gleichzeitige Diagonalisierung der Translationen:  
Die Generatoren $P_\mu$ kommutieren und können simultan (verallgemeinert) diagonalisiert werden.  
Die ebenen Wellen
\[
\phi_p(x) = e^{-i\,p\cdot x}, \qquad p\cdot x := \eta_{\mu\nu}p^\mu x^\nu,
\]
erfüllen
\[
P_\mu \phi_p = p_\mu\,\phi_p, \qquad
E[\phi_p]=(-p^2 + m^2)\,\phi_p = 0,
\]
also $p^2=\eta_{\mu\nu}p^\mu p^\nu = m^2$ (Massenschale).

Massenschale und Lorentz-Orbits:  
\[
\mathcal{H}_m = \{\, p\in\mathbb{R}^{1,3} \mid p^2=m^2,\ p^0>0 \,\}.
\]
Die Lorentzgruppe wirkt transitiv auf $\mathcal{H}_m$, 
der Stabilisator ist $SO(3)$ („little group“).

CSCO und Casimir:  
Der Poincaré-Casimir $P^2=\eta^{\mu\nu}P_\mu P_\nu$ wirkt auf $\ker(E)$ als Multiplikation mit $m^2$.  
Ein CSCO ist z.B.\ $\{P^2,\mathbf{P}\}$; auf der Massenschale liefert $\mathbf{P}$ die kontinuierlichen Label, 
$P^2$ fixiert die physikalische Teilchenart (Masse).  
Für ein \emph{Skalarfeld} wirkt die little group trivial (keine Spin-Entartung).

Fourierdarstellung von Lösungen:  
Jede (geeignete) Lösung besitzt die Darstellung
\[
\phi(x) = \int_{\mathcal{H}_m} \widetilde{\phi}(p)\,e^{-ip\cdot x}\,d\mu_m(p), 
\qquad
d\mu_m(p) = \frac{d^3\mathbf{p}}{(2\pi)^3\,2\omega_{\mathbf{p}}},\ 
\omega_{\mathbf{p}}=\sqrt{m^2+\|\mathbf{p}\|^2}.
\]
Unter $(a,\Lambda)\in\mathcal{P}$ transformiert $\widetilde{\phi}$ nach
\[
\bigl(\widetilde{U(a,\Lambda)\phi}\bigr)(p)
= e^{-ip\cdot a}\,\widetilde{\phi}(\Lambda^{-1}p),
\]
was die Einheitlichkeit der Massenschale (Lorentz-Orbits) und die Translationsphasen (Impuls-Eigenwerte) zeigt.


\item \textbf{Wo die „Beobachter-/Darstellungsinformation“ formal liegt}

Träger der Transformationsinformation:  
Die Information, \emph{wie} sich eine Lösung unter Koordinaten-/Beobachterwechseln verhält, liegt in den Daten
\[
\bigl(M,\;E=M\times\mathbb{C},\;D\equiv 1,\;\eta,\;U(a,\Lambda)\phi(x)=\phi(\Lambda^{-1}(x-a))\bigr).
\]
Insbesondere:
\begin{enumerate}
\item $D$ kodiert die (hier triviale) \emph{interne} Transformationsregel.
\item Die Wirkung von $\mathcal{P}$ auf $M$ kodiert den \emph{externen} Koordinatenwechsel (Lorentz + Translation).
\item Die Metrik $\eta$ koppelt beide Ebenen zu einem \emph{kovarianten} Operator $E=\Box+m^2$.
\end{enumerate}
Damit ist die Transformationsregel von $\phi$ unter Wechsel der Karte (Wirkung auf $M$)
struktureller Bestandteil der Definition von $(E,\phi)$, nicht ein äußerliches Anhängsel.


\item \textbf{Zusammenfassung als „koordinatenfreies“ Lösungsobjekt}

Eine \emph{strukturelle Lösung} der Klein–Gordon-Gleichung ist die Datenmenge
\[
\Bigl(M=\mathbb{R}^{1,3},\;\eta,\;E=M\times\mathbb{C},\;D\equiv 1,\;U:\mathcal{P}\curvearrowright\Gamma(E),\;
E=\Box+m^2,\;\phi\in\Gamma(E)\Bigr),
\]
so dass $E[\phi]=0$ und $E$ unter der durch $U$ induzierten Darstellung forminvariant ist:
\[
E[U(g)\phi]=U(g)E[\phi].
\]
\emph{Physikalische Information} besteht aus
\begin{enumerate}
\item der \textbf{konkreten Konfiguration} $\phi$ (Schnitt eines trivialen Bündels),
\item ihren \textbf{Transformationsgesetzen} $U(a,\Lambda)$ (Beobachter-/Koordinatenwechsel),
\item der \textbf{kovarianten Dynamik} $E=\Box+m^2$ (Poincaré-Kovarianz).
\end{enumerate}
Orbits in $\ker(E)$ sind Darstellungen derselben Physik in verschiedenen Inertialsystemen.
Eine Eichreduktion tritt für das Skalarfeld nicht auf.
\end{enumerate}
\end{DEF}





\begin{DEF}{LIE-L09-08-16}{Struktureller Zugang Dirac Gleichung}
\begin{enumerate}

\item \textbf{Intuitive Ebene: Lösung als Funktion auf Raum–Zeit}

Naive Lösung einer Operatorgleichung:  
Sei $M\simeq\mathbb{R}^{1,3}$ mit globaler Karte $x=(x^\mu)$ und Minkowski-Metrik $\eta=\mathrm{diag}(-,+,+,+)$.  
Setze $V\simeq\mathbb{C}^4$ (Spinorraum).  
Die \emph{Dirac-Gleichung} lautet
\[
E[\psi] = (i\gamma^\mu\partial_\mu - m)\psi = 0,
\]
wobei die $\gamma^\mu$ die Clifford-Relationen erfüllen:
\[
\{\gamma^\mu,\gamma^\nu\} = \gamma^\mu\gamma^\nu+\gamma^\nu\gamma^\mu = 2\eta^{\mu\nu}\mathbf{1}_4.
\]
Eine \emph{naive Lösung} ist eine Abbildung $\psi:M\to\mathbb{C}^4$ mit $(i\gamma^\mu\partial_\mu - m)\psi=0$.

Grenze des naiven Bildes:  
Die Darstellung $\psi(x)$ hängt von  
\emph{(i)} der Wahl von Koordinaten $x^\mu$ (Bezugssystem),  
\emph{(ii)} der Wahl einer Basis des Spinorraums $\mathbb{C}^4$,  
und \emph{(iii)} der Wahl eines konkreten $\gamma$-Matrix-Systems ab.  
Ein Wechsel der Lorentz-Basis (oder der Spinbasis) ändert die \emph{Darstellung} derselben physikalischen Lösung.


\item \textbf{Struktureller Übergang: Bündel, Gruppenwirkung, Darstellung}

Spinstruktur und assoziiertes Bündel:  
Die Dirac-Spinoren leben als Schnitte eines \emph{Spinorbündels}
\[
S = P_{\mathrm{Spin}}\times_D \mathbb{C}^4 \;\xrightarrow{\;\pi\;}\; M,
\]
das aus einem \emph{Spin-Prinzipalbündel} $P_{\mathrm{Spin}}\to M$ mit Strukturgruppe
\[
\mathrm{Spin}(1,3)\simeq SL(2,\mathbb{C})
\]
und der fundamentalen Darstellung $D:\mathrm{Spin}(1,3)\to GL(\mathbb{C}^4)$ assoziiert ist.

Schnitt als koordinatenfreie Lösung:  
Ein \emph{Dirac-Spinor} ist ein glatter Schnitt $\psi\in\Gamma(S)$.
Die Dirac-Gleichung ist dann eine Bündelgleichung $E[\psi]=0$ mit
\[
E = i\gamma^\mu\nabla_\mu - m,
\]
wobei $\nabla_\mu$ die zum Spinbündel kompatible \emph{kovariante Ableitung} ist.

Symmetriegruppe und Wirkung:  
Die Poincaré-Gruppe $\mathcal{P}=\mathbb{R}^{1,3}\rtimes SO^+(1,3)$ wirkt auf $M$ über $x\mapsto \Lambda x + a$.
Ihre zweifache Überlagerung $\widetilde{\mathcal{P}}=\mathbb{R}^{1,3}\rtimes \mathrm{Spin}(1,3)$ wirkt auf dem Spinorbündel durch
\[
(U(a,\Lambda)\psi)(x) = S(\Lambda)\,\psi(\Lambda^{-1}(x-a)),
\]
wobei $S(\Lambda)\in \mathrm{Spin}(1,3)$ die $4\times4$-Matrix der Spinor-Darstellung ist.

Koordinaten- und Darstellungswechsel:  
Die Information, \emph{wie} sich $\psi$ unter Lorentz-Transformationen verändert, liegt
nicht in $\psi(x)$ selbst, sondern in der Paarung $(U,S)$:
\[
U(a,\Lambda)\psi(x) = S(\Lambda)\psi(\Lambda^{-1}(x-a)).
\]


\item \textbf{Kovarianter Operator und Forminvarianz}

Kovarianter Operator:  
Im Minkowski-Raum ist $\nabla_\mu=\partial_\mu$ und
\[
E = i\gamma^\mu\partial_\mu - m.
\]
Die $\gamma^\mu$-Matrizen transformieren gemäß
\[
S^{-1}(\Lambda)\gamma^\mu S(\Lambda) = \Lambda^\mu_{\;\nu}\gamma^\nu.
\]

Forminvarianz / Kovarianz:  
Für alle $(a,\Lambda)\in\widetilde{\mathcal{P}}$ und alle $\psi$ gilt
\[
E[U(a,\Lambda)\psi] = U(a,\Lambda)E[\psi].
\]
Beweis (skizziert):  
Aus der Transformationsregel
\[
(U(a,\Lambda)\psi)(x)=S(\Lambda)\psi(\Lambda^{-1}(x-a))
\]
folgt
\[
\partial_\mu(U\psi)(x)
= S(\Lambda)\,(\Lambda^{-1})_\mu{}^{\nu}\,(\partial_\nu\psi)(\Lambda^{-1}(x-a)),
\]
und damit
\[
E[U\psi](x)
= S(\Lambda)\,\bigl(i\gamma^\mu(\Lambda^{-1})_\mu{}^{\nu}\partial_\nu - m\bigr)\psi(\Lambda^{-1}(x-a))
= S(\Lambda)\,(i\gamma^\nu\partial_\nu - m)\psi(\Lambda^{-1}(x-a))
= (UE\psi)(x),
\]
weil $S^{-1}(\Lambda)\gamma^\mu S(\Lambda) = \Lambda^\mu_{\;\nu}\gamma^\nu$.

Infinitesimale Generatoren:  
\[
P_\mu = i\partial_\mu, 
\qquad 
M_{\mu\nu} = i(x_\mu\partial_\nu - x_\nu\partial_\mu) + \frac{i}{4}[\gamma_\mu,\gamma_\nu],
\]
wobei der zweite Term die \emph{Spinor-Darstellung} der Lorentz-Algebra ist.
Diese Generatoren erfüllen
\[
[P_\mu,E]=0,\qquad [M_{\mu\nu},E]=0.
\]


\item \textbf{Lösungsraum, Orbits, Beobachter- und Eichaspekt}

Lösungsraum und Gruppenwirkung:  
\[
\mathcal{S} := \ker(E) = \{\, \psi\in C^\infty(M,\mathbb{C}^4) \mid (i\gamma^\mu\partial_\mu - m)\psi=0 \,\}.
\]
Die Darstellung $U:\widetilde{\mathcal{P}}\to\mathrm{Aut}(C^\infty(M,\mathbb{C}^4))$ schränkt sich zu
\[
U:\widetilde{\mathcal{P}}\times\mathcal{S}\to\mathcal{S},\qquad (g,\psi)\mapsto U(g)\psi
\]
ein (wegen Forminvarianz).

Orbits:  
Für $\psi\in\mathcal{S}$ ist der Orbit
\[
\mathcal{O}_\psi = \{\, U(g)\psi \mid g\in\widetilde{\mathcal{P}} \,\}.
\]
\emph{Interpretation:} Punkte eines Orbits sind verschiedene \emph{Darstellungen}
derselben Physik in unterschiedlichen Inertialsystemen.  
Da die Dirac-Gleichung keine Eichredundanz besitzt, 
bleibt der gesamte Lösungsraum $\ker(E)$ physikalisch relevant.


\item \textbf{Spektral- und Darstellungsstruktur}

Ebenen Wellen und Spektralbedingungen:  
Setze
\[
\psi_p(x) = u(p)\,e^{-ip\cdot x}, \qquad p\cdot x = \eta_{\mu\nu}p^\mu x^\nu.
\]
Einsetzen in die Dirac-Gleichung liefert
\[
(i\gamma^\mu p_\mu - m)u(p)=0.
\]
Nichttriviale Lösungen $u(p)$ existieren genau für $p^2=m^2$.  
Die zugehörigen $u(p)$ bilden zweidimensionale Spinor-Unterräume („Spin up/down“).

Massenschale und Spinstruktur:  
Die zulässigen Vierimpulse liegen auf der Massenschale
\[
\mathcal{H}_m = \{p\in\mathbb{R}^{1,3} \mid p^2=m^2,\ p^0>0\},
\]
die ein Lorentz-Orbit ist.  
Die little group des Vierimpulses $p=(m,0,0,0)$ ist $SO(3)$,  
auf deren zweifacher Überlagerung $\mathrm{Spin}(3)\simeq SU(2)$
die Spinor-Darstellung $D$ als irreduzible $2$-dimensionale Darstellung wirkt.

Casimir-Operatoren:  
Die beiden Poincaré-Casimirs
\[
P^2 = \eta^{\mu\nu}P_\mu P_\nu,\qquad 
W^\mu W_\mu = -m^2 s(s+1),
\]
mit dem Pauli–Lubanski-Vektor
\[
W^\mu = \frac{1}{2}\epsilon^{\mu\nu\rho\sigma}P_\nu M_{\rho\sigma},
\]
klassifizieren die Darstellungen.  
Für die Dirac-Gleichung gilt $P^2=m^2$ und $s=\frac{1}{2}$.

Fourier-Darstellung von Lösungen:  
Jede (geeignete) Lösung kann geschrieben werden als
\[
\psi(x) = \int_{\mathcal{H}_m} 
\sum_{s=\pm\frac{1}{2}} \bigl(
a_s(p)\,u_s(p)\,e^{-ip\cdot x} + b_s^\ast(p)\,v_s(p)\,e^{ip\cdot x}
\bigr)\, d\mu_m(p),
\]
mit Spinoren $u_s(p)$, $v_s(p)$, die
\[
(\gamma^\mu p_\mu - m)u_s(p)=0, \qquad (\gamma^\mu p_\mu + m)v_s(p)=0
\]
erfüllen.  
Unter $(a,\Lambda)\in\widetilde{\mathcal{P}}$ transformieren die Fourierkoeffizienten nach
\[
a_s(p) \mapsto a_{s'}(\Lambda^{-1}p)\,D_{s's}(\Lambda),
\]
mit $D(\Lambda)\in SU(2)$ als Spin-$\frac{1}{2}$-Darstellung der little group.


\item \textbf{Wo die „Beobachter-/Darstellungsinformation“ formal liegt}

Träger der Transformationsinformation:  
Die Information, \emph{wie} sich eine Lösung unter Koordinaten- und Spinbasiswechseln verhält, liegt in den Daten
\[
\bigl(M,\;S=P_{\mathrm{Spin}}\times_D\mathbb{C}^4,\;
D:\mathrm{Spin}(1,3)\to GL(\mathbb{C}^4),\;
\gamma^\mu,\;
U(a,\Lambda)\psi(x)=S(\Lambda)\psi(\Lambda^{-1}(x-a))\bigr).
\]
Insbesondere:
\begin{enumerate}
\item $D$ kodiert die \emph{interne} Spinorstruktur (Spin-$\frac{1}{2}$).
\item Die Wirkung von $\widetilde{\mathcal{P}}$ auf $M$ kodiert den \emph{externen} Koordinatenwechsel.
\item Die $\gamma^\mu$ koppeln Raumzeitindizes und Spinorindizes und sichern Kovarianz.
\item Der Operator $E=i\gamma^\mu\partial_\mu - m$ ist \emph{forminvariant} unter dieser kombinierten Wirkung.
\end{enumerate}
Damit ist die Transformationsregel von $\psi$ unter Wechsel des Beobachters
struktureller Bestandteil der Definition von $(E,\psi)$.


\item \textbf{Zusammenfassung als „koordinatenfreies“ Lösungsobjekt}

Eine \emph{strukturelle Lösung} der Dirac-Gleichung ist die Datenmenge
\[
\Bigl(
M=\mathbb{R}^{1,3},\;
\eta,\;
S=P_{\mathrm{Spin}}\times_D\mathbb{C}^4,\;
D:\mathrm{Spin}(1,3)\to GL(\mathbb{C}^4),\;
U:\widetilde{\mathcal{P}}\curvearrowright\Gamma(S),\;
E=i\gamma^\mu\partial_\mu - m,\;
\psi\in\Gamma(S)
\Bigr),
\]
so dass $E[\psi]=0$ und $E$ unter der durch $(U,D)$ induzierten Darstellung forminvariant ist:
\[
E[U(g)\psi]=U(g)E[\psi].
\]
\emph{Physikalische Information} besteht aus
\begin{enumerate}
\item der \textbf{Konfiguration} $\psi$ (Schnitt des Spinorbündels),
\item ihren \textbf{Transformationsgesetzen} $S(\Lambda)$ (Spinor- und Beobachterwechsel),
\item der \textbf{kovarianten Dynamik} $E=i\gamma^\mu\partial_\mu - m$ (Poincaré-Kovarianz).
\end{enumerate}
Orbits in $\ker(E)$ sind Darstellungen derselben Physik in verschiedenen Inertialsystemen,
gegliedert nach Masse $m$ und Spin $s=\tfrac{1}{2}$.
\end{enumerate}
\end{DEF}





% =========================================
% Skript: Lie-Gruppen, Flüsse und Zustandsbeschreibungen
% =========================================

\subsection{Lie-Gruppen, Flüsse und die Entwicklung von Zustandsbeschreibungen}

\subsection*{Grundobjekte}




\subsection*{Flüsse in der Lie-Gruppe und induzierte Flüsse}

Linksinvariantes Vektorfeld und Gruppenfluss

Zu $X\in\mathfrak g$ gehört das linksinvariante Vektorfeld $\widetilde X$ auf $G$ mit
$\widetilde X(g)=(L_g)_*X$. Sein Fluss ist die Einparameteruntergruppe
\[
\gamma_X(t)\;=\;\exp(tX),\qquad \dot\gamma_X(t)=(L_{\gamma_X(t)})_*X,\ \gamma_X(0)=e.
\]



Induzierter Fluss auf $M$

Der \emph{induzierte Fluss} auf $M$ durch $X$ ist
\[
\Phi_X(t,p)\;:=\;\exp(tX)\cdot p,
\]
und löst die ODE
\[
\frac{d}{dt}\Phi_X(t,p)\;=\;X_M\big(\Phi_X(t,p)\big),\qquad \Phi_X(0,p)=p.
\]


Gegenüberstellung: lokal vs. global

$X\in\mathfrak g$ kodiert die \emph{infinitesimale} Änderungsrichtung der Beschreibung,
der Fluss $t\mapsto\exp(tX)$ die \emph{integrierte, globale} Transformation in $G$;
$\Phi_X$ ist die entsprechende integrierte Bewegung der \emph{Beschreibungen} auf $M$.


\subsection*{Zeitunabhängige vs. zeitabhängige Generatoren}

Zeitunabhängiger Generator

Ist $X\in\mathfrak g$ konstant, so ist die Zeitentwicklung
\[
g(t)=\exp(tX),\qquad p(t)=g(t)\cdot p_0=\exp(tX)\cdot p_0
\]
eine Einparameteruntergruppe in $G$ bzw. deren induzierter Fluss auf $M$.



Zeitabhängiger Generator und geordnete Exponentialabbildung


Für eine glatte Kurve $X(t)\in\mathfrak g$ ist $g(t)\in G$ die Lösung der
Lie-Gruppen-DGL
\[
\dot g(t)=(L_{g(t)})_*X(t),\qquad g(0)=e,
\]
formal geschrieben als
\[
g(t)\;=\;\mathcal{P}\exp\!\Big(\int_0^t X(s)\,ds\Big),
\]
wobei $\mathcal P$ die Zeitordnung bezeichnet. Der induzierte Fluss ist $p(t)=g(t)\cdot p_0$.




Stationäre (zeitunabhängige) vs. zeitabhängige Zustände

\begin{itemize}
\item \emph{Zeitunabhängige Zustände:} $p(t)=\exp(tX)\cdot p_0$ mit konstantem $X$;
in der Quantenmechanik: $\psi(t)=e^{-\,\frac{i}{\hbar}t\hat H}\psi(0)$, 
\emph{stationäre} Zustände sind Eigenzustände von $\hat H$ (nur Phasenentwicklung).
\item \emph{Zeitabhängige Zustände/Generatoren:} $X(t)$ nicht konstant; Entwicklung via $\mathcal P\exp$,
z.\,B. zeitabhängige Hamiltonoperatoren, zeitabhängige Eichpotenziale.
\end{itemize}


\subsection*{Invarianz, Orbits und Erhaltung}


Invariante Größen

Eine Funktion $f:M\to\mathbb{R}$ (Observable) heiße $G$-invariant, wenn
$f(g\cdot p)=f(p)$ für alle $g\in G$. Äquivalent: $X_M f=0$ für alle $X\in\mathfrak g$.




Orbit und physikalischer Inhalt


Der $G$-Orbit eines Zustands $p\in M$ ist $\mathcal{O}_p:=\{g\cdot p\mid g\in G\}$.
Die physikalisch äquivalente Beschreibungsklasse eines Zustands ist seine Orbitklasse $[p]=\mathcal{O}_p$.




Noether-Idee, Casimirs

Invarianz $\Rightarrow$ Erhaltung: Aus $X_M f=0$ folgen Konstanten der Bewegung entlang $\Phi_X$.
Bei Darstellungen quantisiert: Casimir-Operatoren liefern Quantenzahlen (z.\,B. $J^2$ bei Rotationen).


\subsection*{Beispiele}

\paragraph{(1) Klassische Mechanik als Fluss in $\mathrm{Diff}(M)$.}
$M$ Phasenraum, $\omega$ symplektisch, $H:M\to\mathbb{R}$.
Hamiltonsches Vektorfeld $X_H$ mit $\iota_{X_H}\omega=dH$.
Fluss $\Phi_t=\exp(tX_H)$, Zustand $z(t)=\Phi_t(z_0)$, $\dot z=X_H(z)$.
Dies realisiert die Zeitentwicklung als Kurve in der Lie-Gruppe $\mathrm{Diff}(M)$.

\paragraph{(2) Quantenmechanik als Kurve in $U(\mathcal H)$.}
$G=U(\mathcal H)$, $\mathfrak g=\mathfrak u(\mathcal H)$.
Zeitunabhängig: $U(t)=e^{-\,\frac{i}{\hbar}t\hat H}$, $\psi(t)=U(t)\psi(0)$, $i\hbar\dot\psi=\hat H\psi$.
Zeitabhängig: $U(t)=\mathcal P\exp\!\big(-\frac{i}{\hbar}\int_0^t \hat H(s)\,ds\big)$.

\paragraph{(3) Rotationen und Drehimpuls/Spin.}
$G=SO(3)$ (bzw. $SU(2)$), $\mathfrak g=\mathfrak{so}(3)$ mit $[J_i,J_j]=\varepsilon_{ijk}J_k$.
Fluss: $R_{\hat n}(\theta)=\exp(\theta\,\hat n\cdot J)$, induzierte Wirkung auf $M$ (Raum oder Spinraum).
Invariante: $J^2$; Orbitklassen sind Orientationen desselben physikalischen Zustands.

\paragraph{(4) $U(1)$-Eichung (Ladung).}
$G=U(1)$, $\mathfrak g\simeq i\mathbb R$. Fluss: $e^{i\theta Q}$,
Wirkung auf Wellenfunktionen als globale Phase, in Feldtheorien lokal via Verbindung $A$:
Paralleltransport $P\exp\int A$, Krümmung $F=dA+\tfrac12[A,A]$.


\subsection*{Bemerkung zur Interpretation}
Die Lie-Gruppe $G$ ist der Raum der zulässigen \emph{Beschreibungstransformationen}.
Flüsse in $G$ kodieren die \emph{Prozessstruktur} solcher Transformationen.
Über die Wirkung $a:G\times M\to M$ werden daraus Flüsse der \emph{Zustandsbeschreibung}.
Physikalischer Inhalt liegt in Invarianzen (Orbitklassen, Casimirs) oder in gezieltem Symmetriebruch
(zeit-/raumabhängige Generatoren, anisotrope Hamiltonians).









[Lie-Gruppenwirkungen, Flüsse und die Entwicklung von Zustandsbeschreibungen]

Sei $G$ eine Lie-Gruppe mit Lie-Algebra $\mathfrak g = T_eG$ und $\mathcal{H}$ ein komplexer Hilbertraum.

\textbf{(1) Zustände und Beschreibungen.}
Ein \emph{(reiner) Zustand} eines physikalischen Systems wird mathematisch durch einen normierten Vektor
$\psi \in \mathcal H$, $\|\psi\|=1$, \emph{nicht eindeutig}, sondern nur bis auf eine komplexe Phase bestimmt.
Daher ist $\psi$ keine physikalische Entität, sondern eine \emph{Beschreibung} eines physikalischen Zustands.
Die Menge aller Beschreibungen ist die Einheitsphäre
\[
S(\mathcal H)=\{\psi\in\mathcal H\mid \|\psi\|=1\}.
\]
Der physikalische Zustand selbst entspricht der Äquivalenzklasse (dem \emph{Strahl})
\[
[\psi]=\{\,e^{i\theta}\psi \mid \theta\in\mathbb R\,\}\subset S(\mathcal H),
\]
also einem Punkt im projektiven Hilbertraum $\mathbb P(\mathcal H)$.

\textbf{(2) Lie-Gruppenwirkung auf dem Hilbertraum.}
Eine \emph{Darstellung} der Lie-Gruppe $G$ auf $\mathcal H$ ist ein glatter Gruppenhomomorphismus
\[
U:G\longrightarrow \mathcal U(\mathcal H),\qquad g\mapsto U(g),
\]
mit $U(g_1g_2)=U(g_1)U(g_2)$.
$G$ beschreibt den Raum aller \emph{zulässigen Beschreibungstransformationen}.
Die Wirkung auf Zustandsbeschreibungen ist
\[
a:G\times S(\mathcal H)\to S(\mathcal H),\qquad (g,\psi)\mapsto U(g)\psi.
\]


\emph{(a) Grundidee.}
Eine \emph{Darstellung} einer Lie-Gruppe $G$ auf einem komplexen Hilbertraum $\mathcal H$
ist eine Vorschrift, die jedem Gruppenelement $g\in G$ eine unitäre (oder projektiv unitäre) Abbildung
$U(g):\mathcal H\to\mathcal H$ zuordnet, sodass die Gruppenstruktur erhalten bleibt:
\[
U(g_1g_2)=U(g_1)U(g_2),\qquad U(e)=\mathbb 1.
\]
Die Abbildung
\[
U:G\longrightarrow \mathcal U(\mathcal H),\qquad g\mapsto U(g)
\]
heißt \emph{unitäre Darstellung} von $G$ auf $\mathcal H$.

Die physikalische Bedeutung:  
$G$ beschreibt die Symmetriegruppe der Theorie (z.\,B.\ Translationen, Rotationen, Lorentz-, Gauge- oder Zeitentwicklung),
und $U(g)$ beschreibt, \emph{wie sich Zustandsbeschreibungen im Hilbertraum unter dieser Symmetrie verändern}.
Symmetrien in der Physik werden somit als unitäre (oder projektiv unitäre) Transformationen auf Zuständen realisiert.

\emph{(b) Starke Stetigkeit und infinitesimale Generatoren.}
Eine Darstellung $U$ heißt \emph{stetig} (bzw.\ stark stetig), wenn für alle $\psi\in\mathcal H$
die Abbildung $g\mapsto U(g)\psi$ stetig ist.
Nach dem \emph{Theorem von Stone} (bzw.\ von Stone–von Neumann)
besitzt jede stark stetige, unitäre Darstellung eine zugehörige Lie-Algebra-Darstellung:
\[
dU:\mathfrak g\longrightarrow \mathfrak u(\mathcal H),\qquad
X\mapsto dU(X),
\]
wobei $dU(X)$ ein (im Wesentlichen selbstadjungierter) Operator ist, der durch
\[
dU(X)\psi := \left.\frac{d}{dt}\right|_{t=0} U(\exp tX)\psi
\]
definiert ist.  
Die zugehörigen \emph{Generatoren}
\[
\hat X := i\hbar\,dU(X)
\]
sind die physikalischen Observablen, die den infinitesimalen Symmetrierichtungen entsprechen.
Damit gilt
\[
U(\exp tX)=e^{-\,\frac{i}{\hbar}t\hat X}.
\]

\emph{(c) Zusammenhang mit physikalischen Größen.}
Für eine Lie-Gruppe $G$ von Symmetrien sind die Generatoren $\hat X_i$
die fundamentalen Observablen des Systems:
\begin{itemize}
\item $G=\mathbb R^3$ (Translationen) $\Rightarrow$ $\hat X_i = \hat p_i$ (Impulsoperatoren).
\item $G=SO(3)$ (Rotationen) $\Rightarrow$ $\hat X_i = \hat J_i$ (Drehimpulsoperatoren).
\item $G=U(1)$ (Phasenrotationen) $\Rightarrow$ $\hat X = \hat Q$ (Ladungsoperator).
\item $G=\mathbb R$ (Zeitentwicklung) $\Rightarrow$ $\hat X = \hat H$ (Hamiltonoperator).
\end{itemize}
Alle physikalisch relevanten Größen erscheinen also als Generatoren unitärer Darstellungen
von kontinuierlichen Symmetriegruppen.




\emph{Setting.}
Sei $G$ eine Lie-Gruppe mit Lie-Algebra $\mathfrak g=T_eG$. Sei $(\mathcal H,\langle\cdot,\cdot\rangle)$ ein komplexer Hilbertraum.
Eine (stark stetige) unitäre Darstellung ist ein Homomorphismus
\[
U: G \longrightarrow \mathcal U(\mathcal H),\qquad U(g_1g_2)=U(g_1)U(g_2),\quad U(e)=\mathbb 1.
\]
Das \emph{Differential} (Lie-Algebra-Darstellung) ist
\[
dU:\mathfrak g \longrightarrow \mathfrak u(\mathcal H),\qquad
dU(X)\psi=\left.\frac{d}{dt}\right|_{0}U(\exp tX)\psi,\quad \psi\in\mathcal D,
\]
wobei $\mathcal D$ eine dichte, $dU(\mathfrak g)$-invariante Domäne ist.
Die (physikalischen) Generatoren sind $\hat X := i\hbar\,dU(X)$ (im Wesentlichen selbstadjungig).

\emph{Die drei natürlichen Kommutativitäts-Beziehungen.}
\begin{enumerate}
\item[\textbf{(G)}] \textbf{„Exponential-Brücke“ (Gruppen–Operatoren).} Für alle $X\in\mathfrak g$ und $t\in\mathbb R$ gilt
\[
\boxed{ \quad U(\exp tX) \;=\; \exp\!\big( t\,dU(X)\big) \quad }
\]
(Stone–Theorem / Hille–Yosida für stark stetige unitäre Gruppen).\\
\emph{Diagramm (Exponential-Naturalität):}
\[
\begin{tikzcd}[column sep=large,row sep=large]
\mathfrak g \arrow{r}{\exp} \arrow{d}{dU}
& G \arrow{d}{U}
\\
\mathfrak u(\mathcal H) \arrow{r}{\exp}
& \mathcal U(\mathcal H)
\end{tikzcd}
\qquad\text{(kommutativ)}
\]
Die Kommutativität bedeutet: „\emph{erst exponenzieren in $G$, dann darstellen}“ = „\emph{erst differenzieren, dann Operator exponenzieren}“.

\item[\textbf{(A)}] \textbf{„Klammer-Brücke“ (Algebra–Operatoren).} Für alle $X,Y\in\mathfrak g$ gilt
\[
\boxed{\quad dU([X,Y]) \;=\; [\,dU(X),\,dU(Y)\,] \quad}
\]
(Algebrahomomorphismus).\\
\emph{Diagramm (BCH/Bracket-Kompatibilität):}
\[
\begin{tikzcd}[column sep=large,row sep=large]
\mathfrak g\times \mathfrak g \arrow{r}{[\ ,\ ]} \arrow{d}{dU\times dU}
& \mathfrak g \arrow{d}{dU}
\\
\mathfrak u(\mathcal H)\times \mathfrak u(\mathcal H) \arrow{r}{[\ ,\ ]}
& \mathfrak u(\mathcal H)
\end{tikzcd}
\qquad\text{(kommutativ)}
\]
Dies ist die infinitesimale Form der BCH-Kompatibilität: die erste nichttriviale Abweichung 2.~Ordnung (Kommutator) wird respektiert.

\item[\textbf{(Ad)}] \textbf{„Konjugations-Brücke“ (Adjungierte Wirkung).} Für alle $g\in G$, $X\in\mathfrak g$ gilt
\[
\boxed{\quad dU(\mathrm{Ad}_g X) \;=\; \mathrm{Ad}_{U(g)}\big(dU(X)\big)
\;=\; U(g)\,dU(X)\,U(g)^{-1} \quad}
\]
\emph{Diagramm (Ad-Naturalität):}
\[
\begin{tikzcd}[column sep=large,row sep=large]
\mathfrak g \arrow{r}{\mathrm{Ad}_g} \arrow{d}{dU}
& \mathfrak g \arrow{d}{dU}
\\
\mathfrak u(\mathcal H) \arrow{r}{\mathrm{Ad}_{U(g)}}
& \mathfrak u(\mathcal H)
\end{tikzcd}
\qquad\text{(kommutativ)}
\]
Dies ist die formale Aussage der \emph{Kovarianz}: „Konjugation in $G$“ entspricht „Heisenberg-Konjunktion“ auf Operatorseite.
\end{enumerate}

\emph{Ein zusammenfassendes „Hauptdiagramm“.}
\[
\begin{tikzcd}[column sep=huge,row sep=huge]
\mathfrak g \arrow{r}{\exp} \arrow{d}{dU}
  \arrow[bend right=20]{dd}[swap]{\mathrm{Ad}_g}
& G \arrow{d}{U}
  \arrow[bend left=20]{dd}{\mathrm{Ad}_g}
\\
\mathfrak u(\mathcal H) \arrow{r}{\exp}
  \arrow[bend right=20]{dd}[swap]{\mathrm{Ad}_{U(g)}}
& \mathcal U(\mathcal H)
  \arrow[bend left=20]{dd}{\mathrm{Ad}_{U(g)}}
\\[1ex]
\mathfrak g \arrow{r}{\exp} \arrow{d}{dU}
& G \arrow{d}{U}
\\
\mathfrak u(\mathcal H) \arrow{r}{\exp}
& \mathcal U(\mathcal H)
\end{tikzcd}
\]
Alle Vierecke (Exponential- und Ad-Squares) sowie die Bracket-Kompatibilität kommutieren (letztere als separates Quadrat auf der Algebren-Ebene).

\emph{Konsequenzen (physikalische Lesart).}
\begin{itemize}
\item \emph{Generatoren $\leftrightarrow$ Observablen:} $X\in\mathfrak g \mapsto \hat X=i\hbar\,dU(X)$ ist die „lineare Kodierung“ der Symmetrie als messbare Größe.
\item \emph{Flüsse:} $g(t)=\exp(tX)$ auf Gruppenseite $\leftrightarrow$ $U(g(t))=\exp\big(t\,dU(X)\big)$ auf Operatorseite. 
      \,„Lokale Richtungen“ (in $\mathfrak g$) integrieren zu „globalen Prozessen“ (in $G$ \& $\mathcal U(\mathcal H)$).
\item \emph{Kovarianz/Heisenberg-Bild:} $U(g)\,\hat O\,U(g)^{-1}$ ist die Ad-Wirkung; sie entspricht $\mathrm{Ad}_g$ auf $G$.
      Invarianz $[\hat O,dU(\mathfrak g)]=0$ heißt: $\hat O$ ist entlang aller Symmetrieflüsse konstant (Noether).
\item \emph{BCH und Kommutator:} Die 2.-Ordnung-Information der Gruppenmultiplikation (BCH-Term $\tfrac12[X,Y]$) wird auf Operatorseite exakt durch 
      $[dU(X),dU(Y)]$ reproduziert. Das erklärt die Rolle des Kommutators als \emph{erste nichttriviale} Korrektur.
\end{itemize}

\emph{Zeitabhängige Verallgemeinerung (Pfadordnung).}
Für eine glatte Kurve $X(t)\in\mathfrak g$ gilt
\[
g(t)=\mathcal P\exp\!\Big(\int_0^t X(s)\,ds\Big),\qquad
U(g(t))=\mathcal P\exp\!\Big(\int_0^t dU\big(X(s)\big)\,ds\Big),
\]
sodass die „Exponential-Brücke“ \emph{pfadgeordnet} weiter gilt. Nichtkommutativität von $X(s)$ erzeugt die bekannte Zeitordnung.

\emph{Projektive Darstellungen (Wigner) und zentrale Erweiterungen.}
In Quantenmechanik erscheinen oft \emph{projektiv unitäre} Darstellungen $\overline U:G\to \mathcal U(\mathcal H)/U(1)$.
Dann hebt sich $\overline U$ zu einer echten unitären Darstellung $U:\widetilde G\to\mathcal U(\mathcal H)$ der \emph{zentralen Erweiterung} $\widetilde G$
(anomale Phasen $\leadsto$ zentrale Ladung). Alle obigen Quadrate kommutieren dann auf $\widetilde G$.

\emph{Domänen- und Selbstadjungigkeitsbemerkung.}
Die Operatoren $dU(X)$ sind im Allgemeinen unbeschränkt. 
Wählt man eine gemeinsame dichte, invariante Domäne $\mathcal D$ (z.\,B.\ Schwarzraum),
so gelten die Kommutator- und Ad-Relationen streng auf $\mathcal D$; die selbstadjungigen Abschlüsse liefern die physikalischen Observablen
$\hat X=i\hbar\,dU(X)$ mit wohldefinierter Exponential-Dynamik.

\emph{Mini-Beispiel (Rotation).}
$G=SO(3)$, $\mathfrak g=\mathfrak{so}(3)$, $U:SO(3)\to\mathcal U(\mathcal H)$,
$dU(J_i)=\frac{1}{i\hbar}\hat J_i$, $U(\exp\theta\,\hat n\cdot J)=\exp\!\big(-\frac{i}{\hbar}\theta\,\hat n\cdot\hat{\mathbf J}\big)$.
Für $g\in SO(3)$: $dU(\mathrm{Ad}_g J_i)=U(g)\,dU(J_i)\,U(g)^{-1}$ (Achsenrotation der Generatoren),
und $[dU(J_i),dU(J_j)]=dU([J_i,J_j])=\frac{1}{(i\hbar)^2}[\hat J_i,\hat J_j]$.

\emph{Zusammenfassung in einer Box.}
\[
\boxed{
\begin{aligned}
&\text{(Exponential)}\quad U\!\circ\!\exp \;=\; \exp\!\circ dU,\\
&\text{(Klammer)}\quad dU([X,Y]) \;=\; [\,dU(X),dU(Y)\,],\\
&\text{(Ad-Kovarianz)}\quad dU(\mathrm{Ad}_g X) \;=\; \mathrm{Ad}_{U(g)}(dU(X)).
\end{aligned}
}
\]
Diese drei Relationen sind der \emph{kommunikative Nerv} zwischen Geometrie (Lie-Gruppe/-Algebra) und Operatoren (Darstellung): 
Sie garantieren, dass jede geometrische Aussage über $G,\mathfrak g$ unmittelbar eine operatorielle Aussage über $U(G),dU(\mathfrak g)$ besitzt — und umgekehrt.




\textbf{(2.0) Die Lie-Algebra der unitären Gruppe $\mathcal U(\mathcal H)$}

\emph{(1) Die unitäre Gruppe.}
Sei $(\mathcal H,\langle\cdot,\cdot\rangle)$ ein komplexer Hilbertraum (endlich- oder unendlichdimensional).
Die \emph{unitäre Gruppe} ist
\[
\mathcal U(\mathcal H)
:=\{\,U\in\mathfrak B(\mathcal H)\mid U^\dagger U=\mathbb 1\,\}.
\]
Sie ist eine Banach-Lie-Gruppe (für die starke Operator-Topologie sogar Fréchet-Lie-Gruppe)
mit Gruppenoperation $(U_1,U_2)\mapsto U_1U_2$ und Identität $\mathbb 1$.

\emph{(2) Tangentialraum an der Identität.}
Die \emph{Lie-Algebra} $\mathfrak u(\mathcal H)$ ist der Tangentialraum $T_{\mathbb 1}\mathcal U(\mathcal H)$:
\[
\mathfrak u(\mathcal H)
=\Bigl\{A=\dot U(0)\,\Big|\,
U:\mathbb R\to\mathcal U(\mathcal H)\text{ glatt, }U(0)=\mathbb 1\Bigr\}.
\]
Für eine glatte Kurve $t\mapsto U(t)\in\mathcal U(\mathcal H)$ mit $U(0)=\mathbb 1$ gilt
\[
U(t)^\dagger U(t)=\mathbb 1.
\]
Differenzieren bei $t=0$ liefert
\[
(\dot U(0))^\dagger+\dot U(0)=0,
\]
also
\[
\boxed{\;\mathfrak u(\mathcal H)=\{A\in\mathfrak B(\mathcal H)\mid A^\dagger=-A\}\;}
\]
— die Menge aller \emph{skew-adjungierten (anti-hermiteschen)} Operatoren.

\emph{(3) Lie-Klammer.}
Die Lie-Klammer auf $\mathfrak u(\mathcal H)$ ist der Kommutator
\[
[A,B]=AB-BA.
\]
Man überprüft:
\[
A^\dagger=-A,\;B^\dagger=-B \;\Rightarrow\; [A,B]^\dagger=-[A,B],
\]
also ist $\mathfrak u(\mathcal H)$ unter dieser Klammer abgeschlossen und bildet eine reelle Lie-Algebra.

\emph{(4) Zusammenhang mit selbstadjungierten Observablen.}
Physikalisch arbeitet man lieber mit \emph{selbstadjungierten} Operatoren $\hat A=i\hbar\,A$.
Durch Multiplikation mit $i\hbar$ erhält man eine isomorphe reelle Lie-Algebra
\[
\boxed{
\mathfrak{herm}(\mathcal H)
=\{\,\hat A\in\mathfrak B(\mathcal H)\mid \hat A^\dagger=\hat A\,\},
\qquad
[A,B]_{\text{QM}}=\frac{1}{i\hbar}[\,\hat A,\hat B\,].
}
\]
Damit gilt
\[
\mathfrak u(\mathcal H)
\simeq
\frac{1}{i\hbar}\,\mathfrak{herm}(\mathcal H),
\qquad
U(\exp tA)=\exp(tA)=\exp\!\left(-\frac{i}{\hbar}t\,\hat A\right).
\]

\emph{(5) Zusammenhang mit Darstellungen.}
Für eine unitäre Darstellung $U:G\to\mathcal U(\mathcal H)$ einer Lie-Gruppe $G$ gilt:
\[
dU:\mathfrak g\longrightarrow \mathfrak u(\mathcal H),
\qquad
X\mapsto dU(X),
\]
und äquivalent
\[
\hat X=i\hbar\,dU(X)\in\mathfrak{herm}(\mathcal H).
\]
Dann
\[
U(\exp tX)=\exp\big(t\,dU(X)\big)=\exp\!\left(-\frac{i}{\hbar}t\,\hat X\right),
\]
wobei $\hat X$ die physikalische Observable (Generator) zur Symmetrie $X$ ist.

\emph{(6) Strukturdiagramm.}
\[
\begin{tikzcd}[column sep=huge,row sep=large]
\mathfrak g \arrow{r}{dU} \arrow{d}{X\mapsto i\hbar X}
& \mathfrak u(\mathcal H) \arrow{d}{A\mapsto i\hbar A}
\\
i\hbar\,\mathfrak g \arrow{r}{X\mapsto \hat X=i\hbar\,dU(X)}
& \mathfrak{herm}(\mathcal H)
\end{tikzcd}
\]
Mit Klammern:
\[
[dU(X),dU(Y)]=dU([X,Y])
\quad\Leftrightarrow\quad
[\hat X,\hat Y]=i\hbar\,\widehat{[X,Y]}.
\]
Dies ist die \emph{quantenmechanische Lie-Klammerrelation}:
Die Kommutatoren der Observablen spiegeln die Lie-Algebra-Struktur der zugrunde liegenden Symmetriegruppe wider.

\emph{(7) Endlichdimensionale Spezialfälle.}
Für $\dim\mathcal H=n<\infty$:
\[
\mathcal U(\mathcal H)=U(n)
=\{U\in GL(n,\mathbb C)\mid U^\dagger U=\mathbb 1\},\qquad
\mathfrak u(\mathcal H)=\mathfrak u(n)
=\{A\in M_n(\mathbb C)\mid A^\dagger=-A\}.
\]
Schreibt man $A=iH$ mit $H=H^\dagger$, so ist
\[
[H_1,H_2]_{\text{QM}}=\frac{1}{i\hbar}[H_1,H_2].
\]

\emph{(8) Unendlichdimensionale Bemerkungen.}
Im unendlichdimensionalen Fall ($\mathcal H=L^2(\mathbb R^3)$):
\begin{itemize}
\item $\mathcal U(\mathcal H)$ ist eine Banach-Lie-Gruppe in der starken Operator-Topologie.
\item $\mathfrak u(\mathcal H)$ enthält alle \emph{beschränkten} skew-adjungierten Operatoren.
\item Physikalisch relevante Generatoren $\hat A$ sind meist unbeschränkt;
  sie erzeugen stark stetige, aber nicht normstetige Gruppen:
  \[
  U(t)=e^{-\,\frac{i}{\hbar}t\hat A}.
  \]
  Nach dem Satz von Stone gilt: zu jedem selbstadjungierten Operator $\hat A$ existiert eine eindeutige
  stark stetige unitäre Einparametergruppe dieser Form.
\end{itemize}

\emph{(9) Zusammenfassung.}
\[
\boxed{
\begin{aligned}
&\mathcal U(\mathcal H)=\{U\mid U^\dagger U=\mathbb 1\},\\[2pt]
&\mathfrak u(\mathcal H)=T_{\mathbb 1}\mathcal U(\mathcal H)
 =\{A\mid A^\dagger=-A\},\\[2pt]
&\text{Klammer: } [A,B]=AB-BA,\\[2pt]
&\text{Physikalische Generatoren: } \hat A=i\hbar\,A\text{ (selbstadjungiert)},\\[2pt]
&U(\exp tX)=\exp(t\,dU(X))
 =\exp\!\left(-\,\frac{i}{\hbar}t\,\hat X\right),\\[2pt]
&[\hat X,\hat Y]=i\hbar\,\widehat{[X,Y]}.
\end{aligned}
}
\]

\emph{(10) Interpretation.}
\begin{itemize}
\item $\mathcal U(\mathcal H)$: Raum aller \emph{reversiblen Beschreibungstransformationen}.
\item $\mathfrak u(\mathcal H)$: \emph{Infinitesimale Änderungen} dieser Transformationen.
\item $i\hbar\,\mathfrak u(\mathcal H)$: Raum der \emph{Observablen}, d.\,h.\ der Richtungen,
  entlang derer sich Beschreibungen entwickeln (Energie, Impuls, Drehimpuls, \dots).
\end{itemize}




\textbf{(2.1) Zustände, Beschreibungen, Observablen und Messungen im Darstellungsrahmen.}

\emph{(a) Beschreibungsebenen.}

In einer unitären Darstellung $(\mathcal H,U)$ einer Lie-Gruppe $G$ gilt:

\begin{itemize}
\item[\textbf{(1)}] \emph{Beschreibungen von Zuständen:}
Jeder Vektor $\psi\in\mathcal H$ beschreibt eine mögliche \emph{Zustandsbeschreibung} des Systems.
Die Wahl von $\psi$ ist jedoch nicht eindeutig: Skalierung mit einer Phase $e^{i\theta}$ liefert dieselbe physikalische Situation.
Daher sind physikalisch relevante Beschreibungen Strahlen
\[
[\psi]=\{\,e^{i\theta}\psi \mid \theta\in\mathbb R\,\}
\]
im projektiven Hilbertraum $\mathbb P(\mathcal H)$.

\item[\textbf{(2)}] \emph{Physikalische Zustände:}
Physikalische Zustände sind Äquivalenzklassen von Beschreibungen unter denjenigen Transformationen,
die keine physikalische Veränderung darstellen:
\[
[\psi]_\mathrm{phys}=\{\,U(g)\psi \mid g\in G_\mathrm{gauge}\,\},
\]
wobei $G_\mathrm{gauge}\subseteq G$ die interne (redundante) Symmetriegruppe ist.
Damit ist der physikalische Zustandsraum ein Quotient
\[
\mathcal M_\mathrm{phys}=\mathbb P(\mathcal H)/G_\mathrm{gauge}.
\]
\end{itemize}

\emph{(b) Observablen und Generatoren.}

Eine \emph{Observable} ist ein (im Wesentlichen) selbstadjungierter Operator $\hat A$ auf $\mathcal H$.  
Die Observablen bilden eine *nichtkommutative Algebra* $\mathcal A$:
\[
\mathcal A:=\{\hat A\in\mathfrak B(\mathcal H)\mid \hat A=\hat A^\dagger\}.
\]

\begin{itemize}
\item Jeder Lie-Algebra-Generator $\hat X=i\hbar\,dU(X)$ ist eine Observable,
  deren Exponential $\exp(-\frac{i}{\hbar}t\hat X)=U(\exp tX)$
  eine kontinuierliche Symmetrietransformation erzeugt.
\item Kommutierende Observablen entsprechen gleichzeitig bestimmbaren Größen:
  $[\hat A,\hat B]=0$  $\Rightarrow$ gemeinsame Eigenbasen.
\item Invariante Observablen unter $U(G)$ sind physikalisch \emph{erhaltende Größen}:
  $U(g)\hat A\,U(g)^{-1}=\hat A$ oder infinitesimal $[dU(X),\hat A]=0$.
\end{itemize}

\emph{(c) Messung.}

\begin{itemize}
\item Eine Messung der Observable $\hat A$ an einer Zustandsbeschreibung $\psi$
liefert ein reelles Ergebnis $\lambda$ aus dem Spektrum $\mathrm{spec}(\hat A)$.
\item Die \emph{Wahrscheinlichkeit}, bei Zustand $\psi$ den Wert $\lambda$ zu messen, ist
\[
P_\psi(\lambda)=\langle\psi,E_\lambda\psi\rangle,
\]
wobei $\{E_\lambda\}$ die Spektralzerlegung von $\hat A$ ist
($\hat A=\int\lambda\,dE_\lambda$).
\item Der Erwartungswert der Observable in $\psi$ ist
\[
\langle\hat A\rangle_\psi = \langle\psi,\hat A\psi\rangle.
\]
\item Nach der Messung kollabiert die Zustandsbeschreibung (idealisiert)
auf einen Eigenvektor $\psi_\lambda$ mit $\hat A\psi_\lambda=\lambda\psi_\lambda$,
also auf eine \emph{stationäre Richtung} des Flusses, der durch $\hat A$ erzeugt wird.
\end{itemize}

\emph{(d) Orbits, Invarianz und physikalische Identität.}

\begin{itemize}
\item Der Orbit eines Zustands unter der Symmetriegruppe
\[
\mathcal O_\psi=\{U(g)\psi\mid g\in G\}
\]
beschreibt alle \emph{Darstellungen desselben physikalischen Sachverhalts} unter verschiedenen Beobachter- oder Koordinatenwahlen.
\item Observablen, die auf einem Orbit denselben Erwartungswert liefern,
\[
\langle U(g)\psi,\hat A\,U(g)\psi\rangle=\langle\psi,\hat A\psi\rangle\ \forall g,
\]
heißen \emph{invariant}; sie messen genuin physikalische Größen.
\item Gauge-Orbits (Untergruppe $G_\mathrm{gauge}$) sind echte Redundanzen:
Zustände auf demselben Gauge-Orbit repräsentieren dasselbe physikalische System.
\end{itemize}

\emph{(e) Verknüpfung mit dem Diagramm (Lie $\leftrightarrow$ Operator).}

Das zentrale Diagramm
\[
\begin{tikzcd}[column sep=huge,row sep=large]
\mathfrak g \arrow{r}{\exp} \arrow{d}{dU}
& G \arrow{d}{U}
\\
\mathfrak u(\mathcal H) \arrow{r}{\exp}
& \mathcal U(\mathcal H)
\end{tikzcd}
\]
verbindet geometrische und physikalische Ebenen:

\begin{itemize}
\item Auf Gruppenebene: $g=\exp(tX)$ beschreibt eine Transformation der \emph{Beschreibung}.
\item Auf Operatorseite: $U(g)=\exp(-\frac{i}{\hbar}t\hat X)$ ist der zugehörige physikalische Prozess (z.\,B.\ Zeitentwicklung).
\item Der Operator $\hat X$ ist damit die messbare Größe, deren „Fluss“ die Transformation erzeugt.
\end{itemize}

\emph{(f) Zusammenspiel von Zuständen, Observablen und Symmetrie.}

\begin{itemize}
\item \emph{Zustandsbeschreibung:} $\psi\in\mathcal H$ (abhängig von Darstellung/Basis).
\item \emph{Physikalischer Zustand:} Äquivalenzklasse $[\psi]\in\mathbb P(\mathcal H)/G_\mathrm{gauge}$.
\item \emph{Observable:} selbstadjungierter Operator $\hat A$, der eine Richtung in der Lie-Algebra der Messoperationen festlegt.
\item \emph{Messung:} Spektralzerlegung $\hat A=\int\lambda\,dE_\lambda$;
  Wahrscheinlichkeit durch $|\langle\psi_\lambda,\psi\rangle|^2$.
\item \emph{Invarianz:} $[\hat A, dU(X)]=0$ $\Rightarrow$ Erhaltungsgröße unter Fluss $\Phi_X(t,\psi)$.
\item \emph{Orbit:} Menge der Beschreibungen eines physikalisch identischen Zustands.
\item \emph{Physikalische Entwicklung:} Pfad $t\mapsto U(\exp tX)\psi$ in $\mathcal H$;
  Projektion $[\psi(t)]$ auf $\mathbb P(\mathcal H)$ ist die Entwicklung des physikalischen Zustands.
\end{itemize}

\emph{(g) Zusammenfassung in einer formalen Übersicht.}

\[
\boxed{
\begin{array}{lll}
\textbf{Mathematisches Objekt} & \textbf{Symbol / Raum} & \textbf{Physikalische Bedeutung}\\[4pt]
\hline
\text{Lie-Gruppe} & G & \text{Raum der Beschreibungstransformationen}\\[3pt]
\text{Lie-Algebra} & \mathfrak g & \text{Infinitesimale Richtungen der Veränderung}\\[3pt]
\text{Darstellung} & U:G\to\mathcal U(\mathcal H) & \text{Regel, wie Beschreibungen sich transformieren}\\[3pt]
\text{Generatoren} & \hat X=i\hbar\,dU(X) & \text{Observablen / messbare Größen}\\[3pt]
\text{Hilbertraum} & \mathcal H & \text{Raum der Zustandsbeschreibungen}\\[3pt]
\text{Projektiver Raum} & \mathbb P(\mathcal H) & \text{Raum der physikalischen Zustände}\\[3pt]
\text{Orbit} & \mathcal O_\psi=U(G)\psi & \text{Darstellungen desselben physikalischen Zustands}\\[3pt]
\text{Observable} & \hat A=\hat A^\dagger & \text{Messbare Größe (Selbstadjungiertheit)}\\[3pt]
\text{Messung} & \langle\psi,\hat A\psi\rangle & \text{Erwartungswert / Ergebnisstatistik}\\[3pt]
\text{Invarianz} & [\hat A,dU(\mathfrak g)]=0 & \text{Erhaltungsgröße (Symmetrie)}\\[3pt]
\end{array}
}
\]

\emph{(h) Interpretationshinweis.}
Die mathematische Struktur trennt sauber:
\begin{itemize}
\item[$\circ$] \emph{Geometrie:} $G$, $\mathfrak g$ — Raum der Transformationen.
\item[$\circ$] \emph{Algebra:} $\mathcal U(\mathcal H)$, $\mathfrak u(\mathcal H)$ — Operatorstruktur der Beschreibung.
\item[$\circ$] \emph{Physik:} Invariante Relationen in $\mathbb P(\mathcal H)$ (Orbitstruktur, Spektren, Erhaltungsgrößen).
\end{itemize}
Damit ist der physikalische Zustand eine \emph{Geometrieklasse eines Vektors},
und eine Messung ist eine \emph{Operation im Darstellungsraum}, deren Ergebnisstruktur
(Spektrum, Wahrscheinlichkeiten) durch die Lie-Algebra- und Spektralstruktur der Observablen bestimmt ist.




\emph{(d) Beispiel 1: Translationen in der Quantenmechanik}

\textbf{Daten und Darstellung.}
Sei $G=\mathbb{R}^3$ die (additive) Translationsgruppe, $\mathfrak g=\mathbb{R}^3$.
Auf $\mathcal H=L^2(\mathbb R^3)$ definiert
\[
\bigl(U(\mathbf a)\psi\bigr)(\mathbf r)\;=\;\psi(\mathbf r-\mathbf a),\qquad
\mathbf a,\mathbf r\in\mathbb R^3,
\]
eine unitäre Darstellung.

\smallskip
\textbf{Repräsentationseigenschaft und Unitarität.}
Für $\mathbf a,\mathbf b\in\mathbb R^3$ gilt
\[
U(\mathbf a)U(\mathbf b)\psi(\mathbf r)
=\psi\bigl(\mathbf r-(\mathbf a+\mathbf b)\bigr)
=U(\mathbf a+\mathbf b)\psi(\mathbf r),\qquad U(\mathbf 0)=\mathrm{id}.
\]
Weiter
\[
\langle U(\mathbf a)\psi,U(\mathbf a)\phi\rangle
=\int_{\mathbb R^3}\overline{\psi(\mathbf r-\mathbf a)}\,\phi(\mathbf r-\mathbf a)\,d^3\mathbf r
=\langle \psi,\phi\rangle,
\]
also $U(\mathbf a)$ unitär und $U(\mathbf a)^{-1}=U(-\mathbf a)=U(\mathbf a)^\dagger$.
Die Abbildung $\mathbf a\mapsto U(\mathbf a)$ ist stark stetig.

\smallskip
\textbf{(2) Infinitesimale Generatoren (Stone).}

Einparametergruppen und starke Stetigkeit. Für $i\in\{1,2,3\}$ sei $U_i(t):=U(t\hat e_i)$ die Einparametergruppe der
Translationen längs der $i$-ten Koordinate auf $\mathcal H=L^2(\mathbb R^3)$,
\[
\bigl(U_i(t)\psi\bigr)(\mathbf r)=\psi(\mathbf r-t\hat e_i).
\]
Dann ist $U_i:\mathbb R\to\mathcal U(\mathcal H)$ (stark) stetig, denn
$\|U_i(t)\psi-\psi\|_2\to0$ folgt aus der Stetigkeit von Verschiebungen in $L^2$.

\smallskip
Stone-Theorem (skizziert). Für jede stark stetige unitäre Einparametergruppe $(U_i(t))_{t\in\mathbb R}$
existiert ein eindeutig bestimmter (maximal) schiefadjungierter Generator $A_i$ mit
\[
U_i(t)=e^{tA_i},\qquad A_i^\dagger=-A_i,
\]
wobei $A_i$ dicht definiert ist und der \emph{starke Ableitungsgrenzwert}
\[
A_i\psi=\lim_{t\to0}\frac{U_i(t)\psi-\psi}{t}
\]
auf einer dichten Domäne $\mathcal D(A_i)$ existiert.

\smallskip
Ableitung auf einem dichten Kern. Wähle als Kern $\mathcal S(\mathbb R^3)$ (Schwarzraum). Für $\psi\in\mathcal S(\mathbb R^3)$ gilt
\[
\frac{d}{dt}\Big|_{t=0}\bigl(U_i(t)\psi\bigr)(\mathbf r)
=\frac{d}{dt}\Big|_{t=0}\psi\bigl(\mathbf r-t\hat e_i\bigr)
=-\,\partial_{r_i}\psi(\mathbf r),
\]
und die Ableitung darf (dominiert) unter das Integral gezogen werden, also
\[
A_i\psi=-\,\partial_{r_i}\psi\qquad\text{auf }\mathcal S(\mathbb R^3).
\]
Damit ist $A_i$ auf $\mathcal S$ als Differentialoperator identifiziert.

\smallskip
Symmetrie und essenzielle Selbstadjungiertheit der Impulsoperatoren. Definiere die (formalen) Impulsoperatoren
\[
\hat p_i:=-\,i\hbar\,A_i=-\,i\hbar\,\partial_{r_i}\qquad\text{auf }\mathcal S(\mathbb R^3).
\]
Für $\phi,\psi\in\mathcal S(\mathbb R^3)$ liefert partielle Integration
\[
\langle \phi,\hat p_i\psi\rangle
=\int \overline{\phi}(-i\hbar\,\partial_{r_i}\psi)\,d^3\mathbf r
=\int (i\hbar\,\partial_{r_i}\overline{\phi})\,\psi\,d^3\mathbf r
=\langle \hat p_i\phi,\psi\rangle,
\]
also ist $\hat p_i$ \emph{symmetrisch} auf $\mathcal S$.
Mittels Fourier-Transformation $\mathcal F:L^2(\mathbb R^3)\to L^2(\mathbb R^3)$,
\[
\tilde\psi(\mathbf p)=\frac{1}{(2\pi\hbar)^{3/2}}\int e^{-\,\frac{i}{\hbar}\mathbf p\cdot\mathbf r}\psi(\mathbf r)\,d^3\mathbf r,
\]
gilt
\[
\bigl(\mathcal F\,\hat p_i\,\mathcal F^{-1}\tilde\psi\bigr)(\mathbf p)=p_i\,\tilde\psi(\mathbf p),
\]
also ist der Abschluss von $\hat p_i$ unitär äquivalent zur Multiplikation mit $p_i$.
Damit ist $\hat p_i$ \emph{essentiell selbstadjungiert} auf $\mathcal S$ und sein selbstadjungiger Abschluss hat Domäne
\[
\mathcal D(\hat p_i)=\Bigl\{\psi\in L^2(\mathbb R^3)\;\Big|\;p_i\,\tilde\psi\in L^2(\mathbb R^3)\Bigr\}.
\]

\smallskip
Weyl-Form und Exponentialdarstellung. Aus Stone folgt $U_i(t)=e^{tA_i}=e^{-\,\frac{i}{\hbar}t\,\hat p_i}$.
Für $\mathbf a=(a_1,a_2,a_3)$ und kommutierende Komponenten $[\hat p_i,\hat p_j]=0$ erhält man
\[
U(\mathbf a)=\prod_{i=1}^3 e^{-\,\frac{i}{\hbar}a_i\hat p_i}
=\exp\!\Bigl(-\,\frac{i}{\hbar}\sum_{i=1}^3 a_i\hat p_i\Bigr)
=\exp\!\Bigl(-\,\frac{i}{\hbar}\,\mathbf a\cdot\hat{\mathbf p}\Bigr).
\]
Direkt in Ortsdarstellung:
\[
\bigl(U(\mathbf a)\psi\bigr)(\mathbf r)
=\psi(\mathbf r-\mathbf a),
\]
was mit der Exponentialdarstellung konsistent ist, da die $\hat p_i$ vertauschen.

\smallskip 
Kan. Vertauschungsrelationen und Kovarianz. Die Ortsoperatoren $(\hat r_i\psi)(\mathbf r)=r_i\psi(\mathbf r)$ erfüllen
\[
[\hat r_i,\hat p_j]=i\hbar\,\delta_{ij}\,\mathbb 1,\qquad [\hat p_i,\hat p_j]=0.
\]
Für die adjoint-Wirkung gilt
\[
U(\mathbf a)\,\hat{\mathbf r}\,U(\mathbf a)^\dagger=\hat{\mathbf r}+\mathbf a,\qquad
U(\mathbf a)\,\hat{\mathbf p}\,U(\mathbf a)^\dagger=\hat{\mathbf p},
\]
also ist $\hat{\mathbf p}$ der Generator translationskovarianter Verschiebungen, während $\hat{\mathbf r}$ entsprechend verschoben wird.

\smallskip Domänenhinweis und Varianten. Auf kompakten Räumen (z.\,B.\ $3$-Torus) werden die zugehörigen Generatoren durch Spektralbedingungen (diskretes Spektrum) bestimmt; auf Intervallen $\subset\mathbb R$ spielen Randbedingungen (selbstadjungige Erweiterungen) eine Rolle.
Auf $\mathbb R^3$ liefert $\mathcal S(\mathbb R^3)$ einen gemeinsamen dichten Kern für alle $\hat p_i$.

\smallskip
Fazit. Die Einparametergruppen $U_i(t)$ sind stark stetig und unitär; ihre Stone-Generatoren sind $A_i=-\partial_{r_i}$ (auf $\mathcal S$). Die physikalisch relevanten selbstadjungigen Operatoren sind
\[
\hat p_i=-\,i\hbar\,\partial_{r_i},\qquad
U(\mathbf a)=\exp\!\Bigl(-\,\frac{i}{\hbar}\,\mathbf a\cdot\hat{\mathbf p}\Bigr),
\]
womit Translationen \emph{exakt} als Exponential der Impulsoperatoren realisiert sind.

\smallskip
\textbf{Selbstadjungiertheit und Vertauschungsrelationen.}
Auf dem dichten Kern $\mathcal S(\mathbb R^3)$ sind die $\hat p_i$ symmetrisch, essenziell selbstadjungiert
(ihre Abschlüsse sind selbstadjungiert).
Außerdem
\[
[\hat p_i,\hat p_j]=0,\qquad
[\hat r_i,\hat p_j]=i\hbar\,\delta_{ij}\,\mathbb 1,
\]
wobei $(\hat r_i\psi)(\mathbf r)=r_i\psi(\mathbf r)$ der Ortsoperator ist.
Die zweite Relation folgt direkt aus der obigen Ableitung.

\smallskip
\textbf{Heisenberg-Kovarianz (Translationskovarianz).}
Für jede Observablen $ \hat O$ gilt im Heisenberg-Bild
\[
\hat O\;\mapsto\;\hat O(\mathbf a):=U(\mathbf a)\,\hat O\,U(\mathbf a)^\dagger.
\]
Insbesondere
\[
U(\mathbf a)\,\hat{\mathbf r}\,U(\mathbf a)^\dagger
= \hat{\mathbf r}+\mathbf a,\qquad
U(\mathbf a)\,\hat{\mathbf p}\,U(\mathbf a)^\dagger=\hat{\mathbf p}.
\]
Also „verschiebt“ $U(\mathbf a)$ den Ortsoperator und lässt den Impuls invariant — genau die physikalische Bedeutung einer Translation.

\smallskip
\textbf{Spektralbild (Impulsraum) und verallgemeinerte Eigenfunktionen.}
Mit der Fourier-Transformation $\mathcal F:L^2(\mathbb R^3)\to L^2(\mathbb R^3)$
\[
\tilde\psi(\mathbf p)=\frac{1}{(2\pi\hbar)^{3/2}}\int e^{-\,\frac{i}{\hbar}\mathbf p\cdot \mathbf r}\,\psi(\mathbf r)\,d^3\mathbf r
\]
wird $\hat{\mathbf p}$ zur Multiplikation:
\[
(\mathcal F\,\hat{\mathbf p}\,\mathcal F^{-1}\tilde\psi)(\mathbf p)=\mathbf p\,\tilde\psi(\mathbf p).
\]
Die verallgemeinerten Eigenfunktionen von $\hat{\mathbf p}$ sind eben die ebenen Wellen
$ \varphi_{\mathbf p}(\mathbf r)=(2\pi\hbar)^{-3/2} e^{\frac{i}{\hbar}\mathbf p\cdot \mathbf r}$ mit
$\hat{\mathbf p}\,\varphi_{\mathbf p}=\mathbf p\,\varphi_{\mathbf p}$ (im Distributionensinn).
Im Impulsraum wirkt $U(\mathbf a)$ als Phasenfaktor:
\[
(\widetilde{U(\mathbf a)\psi})(\mathbf p)=e^{-\,\frac{i}{\hbar}\mathbf a\cdot\mathbf p}\,\tilde\psi(\mathbf p).
\]

\smallskip
\textbf{Zusammenhang mit Dynamik (Translationsinvarianz).}
Ist ein Hamiltonoperator $\hat H$ translationsinvariant, d.\,h.\ $[\,\hat H,\,\hat{\mathbf p}\,]=0$,
so kommutiert $U(\mathbf a)$ mit der Zeitentwicklung $e^{-\,\frac{i}{\hbar}t\hat H}$:
\[
U(\mathbf a)\,e^{-\,\frac{i}{\hbar}t\hat H}=e^{-\,\frac{i}{\hbar}t\hat H}\,U(\mathbf a).
\]
Dann ist der Impuls erhalten (Noether) und der Lösungsraum zerfällt in Impuls-Eigenräume.
Für das freie Teilchen $\hat H=\frac{\hat{\mathbf p}^2}{2m}$ gilt
\(
e^{-\,\frac{i}{\hbar}t\hat H}\varphi_{\mathbf p}=e^{-\,\frac{i}{\hbar}\frac{\|\mathbf p\|^2}{2m}t}\,\varphi_{\mathbf p}.
\)

\smallskip
\textbf{Flussbild (Geometrie).}
Die Kurve $\mathbf a(t)=t\,\mathbf a_0$ in $G=\mathbb R^3$ erzeugt den Fluss
\[
\psi(t)=U(\mathbf a(t))\psi_0
=\exp\!\Big(-\,\frac{i}{\hbar}t\,\mathbf a_0\cdot\hat{\mathbf p}\Big)\psi_0,
\]
der in Ortsdarstellung die Zustandsbeschreibung starr verschiebt:
$\psi(t,\mathbf r)=\psi_0(\mathbf r-t\mathbf a_0)$.
Infinitesimal ist die „Geschwindigkeit“ des Flusses gerade $-\,\mathbf a_0\cdot\nabla$,
also der Generator $\mathbf a_0\cdot\hat{\mathbf p}/(i\hbar)$.

\smallskip
\textbf{Zusammenfassung.}
\[
\boxed{
\begin{aligned}
&U(\mathbf a)\psi(\mathbf r)=\psi(\mathbf r-\mathbf a),\quad
U(\mathbf a)=\exp\!\Big(-\frac{i}{\hbar}\mathbf a\cdot\hat{\mathbf p}\Big),\\
&\hat p_i=-\,i\hbar\,\partial_{r_i},\quad [\hat r_i,\hat p_j]=i\hbar\,\delta_{ij},\quad
U(\mathbf a)\,\hat{\mathbf r}\,U(\mathbf a)^\dagger=\hat{\mathbf r}+\mathbf a,\\
&\text{Impulsraum: } (\widetilde{U(\mathbf a)\psi})(\mathbf p)=e^{-\,\frac{i}{\hbar}\mathbf a\cdot\mathbf p}\tilde\psi(\mathbf p),\quad
[\,\hat H,\,\hat{\mathbf p}\,]=0\ \Rightarrow\ \text{Impulserhaltung.}
\end{aligned}
}
\]

\emph{(e) Beispiel 2: Rotationen und Spin.}
\[
G=SO(3)\text{ oder }SU(2),\qquad \mathfrak g=\mathfrak{so}(3)\simeq\mathbb R^3.
\]
Darstellung auf $\mathcal H=L^2(S^2)$ oder $\mathbb C^{2s+1}$:
\[
U(R)\psi(\mathbf r)=\psi(R^{-1}\mathbf r),
\]
bzw. in der Spin-$s$-Darstellung
\[
U(R)=\exp\!\left(-\frac{i}{\hbar}\theta\,\hat{\mathbf n}\cdot\hat{\mathbf J}\right),
\qquad
[\hat J_i,\hat J_j]=i\hbar\,\varepsilon_{ijk}\hat J_k.
\]
Die Eigenzustände $\hat J^2|j,m\rangle=\hbar^2j(j+1)|j,m\rangle$ und $\hat J_z|j,m\rangle=\hbar m|j,m\rangle$
bilden eine Basis der irreduziblen Darstellung von Dimension $2j+1$.
Dies sind die fundamentalen Bausteine der Darstellungstheorie kompakter Gruppen:
alle Darstellungen von $SU(2)$ zerfallen in direkte Summen solcher irreduzibler Spin-$j$-Darstellungen.

\emph{(f) Beispiel 3: Lorentz-Gruppe und Spinor-Darstellung.}
\[
G=SL(2,\mathbb C)\simeq\text{Lorentz}^\uparrow_+,\qquad \mathfrak g=\mathfrak{sl}(2,\mathbb C).
\]
Die Dirac-Darstellung $D:G\to GL(\mathbb C^4)$ erfüllt
\[
S(\Lambda)\gamma^\mu S^{-1}(\Lambda)=\Lambda^\mu_{\ \nu}\gamma^\nu,
\]
und erzeugt die Lorentz-kovariante Transformation der Dirac-Gleichung
\[
(i\gamma^\mu\partial_\mu - m)\psi=0,\qquad \psi'(x)=S(\Lambda)\psi(\Lambda^{-1}x).
\]
Die Spinor-Darstellung ist projektiv unitär und zeigt, wie interne und externe (Raumzeit-)Transformationen
in einer einheitlichen Lie-Gruppenstruktur wirken.

\emph{(g) Beispiel 4: Phasen- und Eichgruppen.}
Für $G=U(1)$ ist
\[
U(\theta)\psi = e^{i q\theta/\hbar}\psi,
\]
wobei $q$ die Ladung ist.
In einer Eichfeldtheorie mit Verbindung $A_\mu(x)$ wirkt lokal
\[
U(\theta(x))\psi(x)=e^{i q\theta(x)/\hbar}\psi(x),
\]
und die kovariante Ableitung $D_\mu=\partial_\mu - i q A_\mu/\hbar$
stellt sicher, dass die Gleichung $E[\psi]=0$ kovariant bleibt.
Damit ist die lokale Darstellung der Eichgruppe der mathematische Ausdruck
für die Redundanz (Gauge-Freiheit) der Zustandsbeschreibung.


\emph{(h) Beispiel 5: Zeitentwicklung eines beschränkten Wellenpakets}

\emph{(1) Ausgangsrahmen.}
Wir betrachten ein Teilchen im $\mathbb R^3$, dessen Hilbertraum
\[
\mathcal H=L^2(\mathbb R^3),\qquad
\langle\psi,\phi\rangle=\int_{\mathbb R^3}\overline{\psi(\mathbf r)}\phi(\mathbf r)\,d^3\mathbf r
\]
ist.
Die Lie-Gruppe der Raum-Zeit-Translationen ist
\[
G=\mathbb R_t\times\mathbb R^3_x,
\]
mit Generatoren
\[
\hat H\quad\text{(Zeittranslation)}\qquad
\hat{\mathbf P}=(\hat P_x,\hat P_y,\hat P_z)\quad\text{(Raumtranslationen)}.
\]
Diese Operatoren erfüllen
\[
[\hat P_i,\hat P_j]=0,\qquad [\hat H,\hat P_i]=0
\]
(freies Teilchen), und wirken durch
\[
\hat P_i=-\,i\hbar\,\frac{\partial}{\partial x_i},\qquad
\hat H=\frac{\hat{\mathbf P}^2}{2m}=-\frac{\hbar^2}{2m}\nabla^2.
\]
Die Zeitentwicklung ist die Einparametergruppe
\[
U(t)=e^{-\,\frac{i}{\hbar}t\hat H}.
\]

\emph{(2) Anfangszustand (beschränktes Wellenpaket).}
Wähle zur Zeit $t=0$ eine normierte Superposition ebener Wellen:
\[
\psi(\mathbf r,0)
=\frac{1}{(2\pi\hbar)^{3/2}}\int_{\mathbb R^3}\!\!a(\mathbf p)\,e^{\frac{i}{\hbar}\mathbf p\cdot\mathbf r}\,d^3\mathbf p,
\]
wobei $a(\mathbf p)$ ein stark lokalisiertes Amplitudenprofil im Impulsraum ist,
typischerweise ein Gaußpaket um $\mathbf p_0$:
\[
a(\mathbf p)=A_0\,\exp\!\left(-\frac{|\mathbf p-\mathbf p_0|^2}{2\sigma_p^2}\right),
\]
mit $\sigma_p\ll|\mathbf p_0|$.

\emph{(3) Zeitentwicklung im Impulsraum.}
Die Schrödinger-Gleichung
\[
i\hbar\,\partial_t\psi(\mathbf r,t)=\hat H\psi(\mathbf r,t)
\]
entspricht im Impulsraum der einfachen Phasenentwicklung
\[
\tilde\psi(\mathbf p,t)
=e^{-\,\frac{i}{\hbar}\frac{p^2}{2m}t}\,\tilde\psi(\mathbf p,0)
=e^{-\,\frac{i}{\hbar}\frac{p^2}{2m}t}\,a(\mathbf p).
\]
Dies ist die Wirkung des Flusses
\[
U(t)=\exp\!\Big(-\frac{i}{\hbar}t\frac{\hat{\mathbf P}^2}{2m}\Big)
\]
in der Darstellung $dU(\hat H)=\frac{\hat{\mathbf P}^2}{2m}$.

\emph{(4) Rücktransformation in Ortsdarstellung.}
\[
\psi(\mathbf r,t)
=\frac{1}{(2\pi\hbar)^{3/2}}\int_{\mathbb R^3}
a(\mathbf p)\,
\exp\!\Big[\frac{i}{\hbar}\big(\mathbf p\cdot\mathbf r-\frac{p^2}{2m}t\big)\Big]
\,d^3\mathbf p.
\]
Für ein enges Gaußpaket ($\sigma_p\ll|\mathbf p_0|$) kann man die Phase um $\mathbf p_0$ entwickeln:
\[
E(\mathbf p)=\frac{p^2}{2m}\approx E_0+\mathbf v_g\cdot(\mathbf p-\mathbf p_0),
\quad
\mathbf v_g=\nabla_\mathbf p E|_{\mathbf p_0}=\frac{\mathbf p_0}{m}.
\]
Dann ergibt sich
\[
\psi(\mathbf r,t)
\approx e^{\frac{i}{\hbar}(\mathbf p_0\cdot\mathbf r - E_0 t)}\;
\phi(\mathbf r-\mathbf v_g t),
\]
mit
\[
\phi(\mathbf r)
=\frac{1}{(2\pi\hbar)^{3/2}}\int
a(\mathbf p_0+\mathbf q)\,e^{\frac{i}{\hbar}\mathbf q\cdot\mathbf r}\,d^3\mathbf q.
\]

\emph{(5) Interpretation.}
\begin{itemize}
\item $\phi(\mathbf r)$ beschreibt die räumliche Hüllkurve („Form“ des Pakets),
  die sich mit konstanter Geschwindigkeit $\mathbf v_g$ verschiebt:
  \[
  |\psi(\mathbf r,t)|^2=|\phi(\mathbf r-\mathbf v_g t)|^2.
  \]
\item Der Faktor $e^{\frac{i}{\hbar}(\mathbf p_0\cdot\mathbf r - E_0 t)}$
  ist eine schnelle interne Phase, assoziiert mit der „Trägerwelle“ der Energie $E_0$.
\item $\mathbf v_g=\nabla_\mathbf p E(\mathbf p_0)$ ist die \emph{Gruppengeschwindigkeit}:
  die Geschwindigkeit, mit der die Energie und das Maximum der Wahrscheinlichkeit sich fortpflanzen.
\item Der Fluss in $\mathcal U(\mathcal H)$,
  \[
  U(t)=\exp\!\Big(-\frac{i}{\hbar}t\,\frac{\hat{\mathbf P}^2}{2m}\Big),
  \]
  ist die exponentielle Darstellung der Lie-Algebra der Zeittranslationen.
  Seine Wirkung auf den Zustand ist der „Transport“ des Wellenpakets entlang der Raumzeit-Flüsse.
\end{itemize}

\emph{(6) Geometrisch-Lie-theoretische Deutung.}
\begin{itemize}
\item Der Raum der Beschreibungen $\psi(\mathbf r,t)$ ist $\mathcal H$;
  die Zeitentwicklung $U(t)$ ist eine Einparameter-Untergruppe von $\mathcal U(\mathcal H)$.
\item Der Generator $dU(\partial_t)= -\,\frac{i}{\hbar}\hat H$ liegt in $\mathfrak u(\mathcal H)$.
\item Der Fluss $\psi(t)=U(t)\psi(0)$ ist die Kurve in $\mathcal H$, die der Integralfluss des Vektorfeldes
  $X_\psi=-\,\frac{i}{\hbar}\hat H\psi$ ist.
\item In der projektiven Raumdarstellung $\mathbb P(\mathcal H)$ beschreibt der Fluss den physikalischen Zustand (Phase eliminiert).
\end{itemize}

\emph{(7) Zusammenfassung.}
\[
\boxed{
\begin{aligned}
&\text{Ein beschränktes Wellenpaket } \psi(\mathbf r,0)
=\int a(\mathbf p)\,e^{\frac{i}{\hbar}\mathbf p\cdot\mathbf r}\,d^3\mathbf p,\\[2pt]
&\text{entwickelt sich gemäß } U(t)=e^{-\,\frac{i}{\hbar}t\hat H},\quad \hat H=\tfrac{\hat{\mathbf P}^2}{2m},\\[2pt]
&\Rightarrow\ \psi(\mathbf r,t)\approx e^{\frac{i}{\hbar}(\mathbf p_0\cdot\mathbf r-E_0t)}\phi(\mathbf r-\mathbf v_g t),\quad
\mathbf v_g=\frac{\mathbf p_0}{m}.
\end{aligned}
}
\]
Das Wellenpaket bleibt formstabil (bis auf Dispersionskorrekturen) und bewegt sich mit Gruppengeschwindigkeit $\mathbf v_g$
in Richtung des mittleren Impulses $\mathbf p_0$.


\emph{(1) Grunddiagramm.}

\[
\begin{tikzcd}[column sep=huge,row sep=large]
\mathfrak g \arrow{r}{\exp} \arrow{d}{dU}
& G \arrow{d}{U}
\\
\mathfrak u(\mathcal H) \arrow{r}{\exp}
& \mathcal U(\mathcal H)
\end{tikzcd}
\]
\[
\Downarrow
\]
\[
\text{„Geometrische Zeittranslation“} \;\longmapsto\;
\text{„Unitäre Zeitentwicklung“}.
\]

\emph{(2) Geometrische Seite.}

\begin{itemize}
\item Gruppe der Zeittranslationen:
  \[
  G=\mathbb R_t,\qquad \text{Multiplikation: } (t_1,t_2)\mapsto t_1+t_2.
  \]
\item Lie-Algebra:
  \[
  \mathfrak g=\mathbb R,\qquad
  X=\frac{\partial}{\partial t},\qquad [X,X]=0.
  \]
\item Exponentialabbildung:
  \[
  \exp(tX)=t.
  \]
\end{itemize}

\emph{(3) Darstellungsseite.}

\begin{itemize}
\item Darstellung:
  \[
  U:G\to\mathcal U(\mathcal H),\qquad U(t)=e^{-\,\frac{i}{\hbar}t\hat H}.
  \]
\item Abgeleitete Darstellung:
  \[
  dU:\mathfrak g\to\mathfrak u(\mathcal H),\qquad dU(X)=-\,\frac{i}{\hbar}\hat H.
  \]
\item Generator:
  \[
  \hat H\in\mathfrak{herm}(\mathcal H)
  \quad\text{(selbstadjungierter Operator, Energie).}
  \]
\item Infinitesimale Lie-Algebra-Wirkung:
  \[
  \frac{d}{dt}U(t)\psi=-\,\frac{i}{\hbar}\hat H\,U(t)\psi.
  \]
\end{itemize}

Damit ist die Schrödinger-Gleichung exakt die Darstellung der Lie-Algebra:
\[
i\hbar\,\partial_t\psi=\hat H\psi
\quad\Longleftrightarrow\quad
\partial_t\psi=dU(X)\psi.
\]

\emph{(4) Der Fluss auf dem Hilbertraum.}

Das Wellenpaket $\psi(t)=U(t)\psi_0$ ist eine Kurve im Hilbertraum $\mathcal H$:
\[
t\mapsto \psi(t)=e^{-\,\frac{i}{\hbar}t\hat H}\psi_0.
\]
Diese Kurve ist der Integralfluss des (linksinvarianten) Vektorfeldes
\[
V_\psi=-\,\frac{i}{\hbar}\hat H\psi
\]
auf $\mathcal H$, also der Lie-Gruppenfluss induziert durch $dU(X)$.

Im projektiven Raum $\mathbb P(\mathcal H)$ (physikalische Zustände) ist der Fluss dieselbe Kurve modulo globaler Phase:
\[
t\mapsto [\psi(t)].
\]

\emph{(5) Der Orbit in $\mathcal H$.}

\[
\mathcal O_{\psi_0}
=\{U(t)\psi_0\mid t\in\mathbb R\}
\]
ist die Bahn der Zeitentwicklung — eine \emph{Einparameter-Untergruppe}-Orbit des Zustands $\psi_0$ unter $U(G)$.
In einem unitaren Bild ist dieser Orbit eine geschlossene Kurve, falls das Spektrum von $\hat H$ diskret ist (periodische Entwicklung),
oder eine nichtgeschlossene Helix bei kontinuierlichem Spektrum.

\emph{(6) Beispiel: Freies Wellenpaket.}

Für $\mathcal H=L^2(\mathbb R^3)$ und $\hat H=\tfrac{\hat{\mathbf P}^2}{2m}$ gilt:
\[
U(t)=\exp\!\Big(-\frac{i}{\hbar}t\,\frac{\hat{\mathbf P}^2}{2m}\Big),\qquad
dU(X)=-\frac{i}{\hbar}\,\frac{\hat{\mathbf P}^2}{2m}.
\]
Das Wellenpaket ist
\[
\psi(t)=U(t)\psi_0,\quad
\psi(\mathbf r,t)
=\frac{1}{(2\pi\hbar)^{3/2}}\int
a(\mathbf p)e^{\frac{i}{\hbar}(\mathbf p\cdot\mathbf r-\frac{p^2}{2m}t)}\,d^3\mathbf p.
\]
→ Das ist die konkrete Realisierung der Exponentialabbildung im Diagramm.

\emph{(7) Gruppengeschwindigkeit als differentielle Struktur.}

Die Gruppengeschwindigkeit ist die Ableitung der Energie nach dem Impuls,
also der Differentialoperator der Lie-Algebra-Wirkung:
\[
\mathbf v_g=\nabla_\mathbf p E(\mathbf p_0)
=\frac{d}{d\mathbf p}\Big|_{\mathbf p_0}\!\bigl(\tfrac{p^2}{2m}\bigr)
=\frac{\mathbf p_0}{m}.
\]
In der Diagrammsprache: 
\[
E:\mathfrak g\to\mathbb R,\quad \mathbf p\mapsto\langle\psi_0,\hat H\psi_0\rangle,
\]
und $\mathbf v_g$ ist der Pushforward der Energie unter der Darstellungsdifferentialabbildung $dU$.

\emph{(8) Diagramm für das Wellenpaket.}

\[
\begin{tikzcd}[column sep=huge,row sep=large]
\mathfrak g=\mathbb R\langle X\rangle \arrow{r}{\exp} \arrow{d}{dU}
& G=\mathbb R_t \arrow{d}{U(t)=e^{-\,\frac{i}{\hbar}t\hat H}} \arrow{r}{\text{Wirkung auf }M=\mathbb R^3}
& \mathbf r\mapsto \mathbf r+\mathbf v_g t
\\
\mathfrak u(\mathcal H)=\mathbb R\langle -\,\tfrac{i}{\hbar}\hat H\rangle
\arrow{r}{\exp}
& \mathcal U(\mathcal H)
\arrow{r}{\text{Wirkung auf }\mathcal H}
& \psi(\mathbf r,0)\mapsto\psi(\mathbf r,t)
\end{tikzcd}
\]
Dieses Diagramm fasst alle Ebenen zusammen:
\begin{enumerate}
\item Oben: Geometrische Zeittranslation auf der Raumzeit.
\item Mitte: Lie-Algebra–Darstellungsabbildung $dU$ als Übersetzung von Geometrie in Operatorik.
\item Unten: Wirkung auf Zustandsbeschreibungen (Hilbertraum).
\end{enumerate}

\emph{(9) Zusammenfassung.}
\[
\boxed{
\begin{aligned}
&\text{Geometrische Zeittranslation } \exp(tX)
\;\xrightarrow{U}\;
\text{Unitäre Zeitentwicklung } e^{-\,\frac{i}{\hbar}t\hat H}\\[2pt]
&\text{Generator } X\in\mathfrak g
\;\xrightarrow{dU}\;
dU(X)=-\,\frac{i}{\hbar}\hat H\in\mathfrak u(\mathcal H)\\[2pt]
&\text{Orbit } \psi(t)=U(t)\psi_0
\text{ beschreibt die Bewegung des Wellenpakets.}\\[2pt]
&\text{Gruppengeschwindigkeit } \mathbf v_g
=\nabla_\mathbf p E(\mathbf p_0)
\text{ ist die Differentialstruktur des Flusses in Impulsraum.}
\end{aligned}
}
\]


\emph{Zeitentwicklung als induzierter Fluss von Zustandsbeschreibungen}

\textbf{(1) Strukturelle Daten.}
Sei $G=\mathbb R_t$ die Lie-Gruppe der Zeittranslationen mit Lie-Algebra
\[
\mathfrak g=\mathbb R\langle X\rangle,\qquad [X,X]=0,\qquad \exp(tX)=t.
\]
Sei $(\mathcal H,\langle\cdot,\cdot\rangle)$ ein komplexer Hilbertraum
und $U:G\to\mathcal U(\mathcal H)$ eine stark stetige unitäre Darstellung,
deren abgeleitete Darstellung
\[
dU:\mathfrak g\longrightarrow\mathfrak u(\mathcal H),
\qquad dU(X)=-\,\frac{i}{\hbar}\hat H
\]
den (im Wesentlichen) selbstadjungierten Hamiltonoperator $\hat H$ definiert.

\textbf{(2) Aktiver Fluss im Hilbertraum.}
Die zugehörige Einparametergruppe
\[
U(t)=\exp(t\,dU(X))
=\exp\!\Big(-\frac{i}{\hbar}t\,\hat H\Big)
\]
erzeugt auf $\mathcal H$ einen \emph{unitären Fluss}
\[
\Phi_t:\mathcal H\to\mathcal H,\qquad
\Phi_t(\psi_0)=U(t)\psi_0.
\]
Die Kurve $\psi(t):=\Phi_t(\psi_0)$ erfüllt die \emph{Schrödinger-Gleichung}
\[
\dot\psi(t)=dU(X)\psi(t)
=-\,\frac{i}{\hbar}\hat H\psi(t).
\]

\textbf{(3) Projektiver (physikalischer) Fluss.}
Da globale Phasen physikalisch irrelevant sind,
projiziert $\Phi_t$ zu einem Fluss auf dem projektiven Raum
\[
[\Phi_t]:\mathbb P(\mathcal H)\longrightarrow \mathbb P(\mathcal H),
\qquad
[\psi]\longmapsto [U(t)\psi].
\]
Die Bahn
\[
\mathcal O_{[\psi_0]}
=\{[U(t)\psi_0]\mid t\in\mathbb R\}
\subset\mathbb P(\mathcal H)
\]
ist die \emph{physikalische Zeitentwicklung} des Zustands.

\textbf{(4) Darstellung in einer Ortsbasis.}
Wählt man eine feste Orthonormalbasis $|\mathbf r\rangle\in\mathcal H$
(\emph{Ortsdarstellung}), so sind die Koordinaten der Kurve $\psi(t)$ gegeben durch
\[
\psi(t,\mathbf r)
:=\langle \mathbf r|\psi(t)\rangle
=\langle \mathbf r|U(t)|\psi_0\rangle.
\]
Dies ist die \emph{Darstellungsfunktion} der Zustandsbeschreibung
zur Zeit $t$ im gewählten Koordinatensystem.

\textbf{(5) Passive Interpretation.}
Für $U(t)=e^{-\,\frac{i}{\hbar}t\hat H}$ gilt
\[
\psi(t,\mathbf r)
=\psi_0(\mathbf r-\mathbf v_g t)
\qquad(\text{freies Teilchen}),
\]
so dass die Zeitentwicklung im Ortsraum als Translation der Koordinatenfunktion erscheint.
Formal gilt
\[
\psi(t,\mathbf r)
=(U(t)\psi_0)(\mathbf r)
=\psi_0\bigl(\Phi_{-t}(\mathbf r)\bigr),
\]
wobei $\Phi_{t}:\mathbf r\mapsto \mathbf r+\mathbf v_g t$
der geometrische Fluss der Raumtranslation ist.
Dies ist eine \emph{passive Beschreibung} derselben Transformation:
\[
\text{Aktiv: } \psi(t)=U(t)\psi_0,
\qquad
\text{Passiv: } \psi_t(\mathbf r)=\psi_0(\Phi_{-t}(\mathbf r)).
\]
Beide sind Darstellungen desselben Flusses auf unterschiedlichen Ebenen.

\textbf{(6) Zusammenfassung.}
\[
\boxed{
\begin{aligned}
&\text{Aktiv:}\quad \psi(t)=U(t)\psi_0
\text{ beschreibt den Orbit im Hilbertraum,}\\[2pt]
&\text{Passiv:}\quad \psi_t(\mathbf r)=\psi_0(\mathbf r-\mathbf v_g t)
\text{ beschreibt die gleiche Entwicklung in fixen Koordinaten.}\\[4pt]
&\text{Physikalisch:}\quad [\psi(t)]\in\mathbb P(\mathcal H)
\text{ ist der eigentliche Zustand; die „Bewegung“}\\
&\text{des Wellenpakets ist die Veränderung der Darstellung des Zustands unter dem Fluss }U(t).
\end{aligned}
}
\]




\emph{(h) Irreduzible Darstellungen und physikalische Elementarobjekte.}
Eine Darstellung $(\mathcal H,U)$ heißt \emph{irreduzibel}, wenn kein nichttrivialer abgeschlossener Unterraum
$V\subset \mathcal H$ invariant unter allen $U(g)$ ist:
\[
U(g)V\subseteq V\ \forall g\in G \ \Rightarrow\ V=\{0\}\text{ oder }V=\mathcal H.
\]
Physikalisch entsprechen irreduzible Darstellungen den \emph{Elementarteilchenzuständen}:
Jede irreduzible Darstellung der Poincaré-Gruppe
wird durch Masse $m$ und Spin $s$ charakterisiert (nach Wigner 1939).
Diese $(m,s)$-Labels sind die fundamentalen invarianten Parameter (Casimirs)
der Darstellung.

\emph{(i) Zusammenfassung.}
\[
\boxed{
\begin{aligned}
&\text{Darstellung }U:G\to\mathcal U(\mathcal H)\text{ verbindet Symmetrie und Zustandsbeschreibung.}\\
&\text{Generatoren } \hat X=i\hbar\,dU(X)\text{ sind Observablen (Impuls, Energie, Drehimpuls,...).}\\
&\text{Einparameteruntergruppen } e^{tX}\text{ erzeugen Flüsse (Zeit, Translation, Rotation).}\\
&\text{Irreduzible Darstellungen } \Rightarrow \text{elementare physikalische Systeme.}\\
&\text{Kovarianz } [dU(X),E]=0 \text{ sichert Forminvarianz physikalischer Gleichungen.}
\end{aligned}
}
\]

\textbf{(3) Physikalische Zustände und Orbits.}
Zwei Beschreibungen $\psi,\psi'\in S(\mathcal H)$ heißen \emph{physikalisch äquivalent}, 
wenn es ein $g\in G$ mit $\psi'=U(g)\psi$ gibt (und ggf. eine globale Phase $e^{i\theta}$).
Der \emph{Orbit}
\[
\mathcal O_\psi := \{\, e^{i\theta}U(g)\psi \mid g\in G,\ \theta\in\mathbb R \,\}
\]
repräsentiert dann den \emph{physikalischen Zustand}.
Jeder Orbit fasst alle Beschreibungen zusammen, die dasselbe physikalische System repräsentieren.
Die physikalisch relevanten Objekte sind somit Punkte in der Orbitenmannigfaltigkeit
$S(\mathcal H)/(\mathbb R\times G)$ bzw. im projektiven Hilbertraum $\mathbb P(\mathcal H)$.

\textbf{(4) Infinitesimale Struktur und Generatoren.}
Das Differential von $U$ an der Identität
\[
dU:\mathfrak g\longrightarrow \mathfrak u(\mathcal H)
\]
ordnet jedem $X\in\mathfrak g$ einen (selbstadjungierten) Operator 
$\hat X=i\hbar\,dU(X)$ zu, der den \emph{Generator der infinitesimalen Transformation} in Richtung $X$ bildet:
\[
U(\exp tX)=e^{-\,\frac{i}{\hbar}t\hat X}.
\]
$\hat X$ beschreibt, wie sich eine Zustandsbeschreibung unter einer infinitesimalen Änderung
der Darstellungsvorschrift transformiert.

\textbf{(5) Globale Flüsse in der Lie-Gruppe.}
Zu $X\in\mathfrak g$ gehört das linksinvariante Vektorfeld $\widetilde X(g)=(L_g)_*X$
mit Fluss
\[
\gamma_X:\mathbb R\to G,\qquad \gamma_X(0)=e,\quad \dot\gamma_X(t)=(L_{\gamma_X(t)})_*X.
\]
Die Lösung $\gamma_X(t)=\exp(tX)$ ist eine \emph{Einparameter-Untergruppe} von $G$
und beschreibt die \emph{Prozessstruktur einer kontinuierlichen Transformation}
im Raum der Beschreibungstransformationen.

\textbf{(6) Induzierte Flüsse auf dem Zustandsraum.}
Über die Darstellung $U$ induziert $\gamma_X$ den \emph{unitären Fluss auf $S(\mathcal H)$}
\[
\Phi_X(t,\psi):=U(\exp tX)\psi,\qquad \psi\in S(\mathcal H),
\]
der die Differentialgleichung
\[
\frac{d}{dt}\Phi_X(t,\psi)=-\frac{i}{\hbar}\,\hat X\,\Phi_X(t,\psi),
\qquad \Phi_X(0,\psi)=\psi
\]
erfüllt.
Dieser Fluss beschreibt die \emph{zeitliche oder parametrische Entwicklung der Beschreibung}
bei kontinuierlicher Anwendung der durch $X$ erzeugten Transformation.

\textbf{(7) Zeitabhängige Generatoren.}
Für zeitabhängige $X(t)\in\mathfrak g$ gilt
\[
\dot g(t)=(L_{g(t)})_*X(t),\quad g(0)=e,
\]
mit Lösung
\[
U(t)=\mathcal P\exp\!\left(-\frac{i}{\hbar}\int_0^t \hat X(s)\,ds\right),\qquad
\psi(t)=U(t)\psi(0),
\]
was die allgemeine Form einer \emph{zeitabhängigen Zustandsbeschreibung} ist.

\textbf{(8) Invarianz, Dynamik und physikalischer Gehalt.}
\begin{itemize}
\item Eine Observable $\hat O$ ist $G$-invariant, wenn $U(g)^\dagger \hat O\,U(g)=\hat O$
oder äquivalent $[\hat O,\hat X]=0$ für alle $X\in\mathfrak g$.
Erwartungswerte solcher Observablen bleiben entlang jedes Flusses $\Phi_X(t,\psi)$ konstant:
\[
\frac{d}{dt}\langle \Phi_X(t,\psi),\hat O\,\Phi_X(t,\psi)\rangle=0.
\]
\item Die Punkte desselben Orbits $\mathcal O_\psi$ beschreiben \emph{dasselbe physikalische System}
unter verschiedenen Darstellungen.
\item Wird die Symmetrie gebrochen (etwa durch ein äußeres Feld, sodass $[\hat H,\hat X]\neq0$),
dann ist der Fluss $\Phi_X(t,\psi)$ keine reine Beschreibungstransformation mehr,
sondern eine \emph{physikalische Entwicklung} im Sinne dynamischer Veränderung.
\end{itemize}

\textbf{(9) Zusammenfassung.}
\[
\boxed{
\begin{aligned}
&\text{Lie-Gruppe } G: \text{Raum der zulässigen Beschreibungstransformationen,}\\
&\text{Lie-Algebra } \mathfrak g: \text{infinitesimale Richtungen dieser Änderungen,}\\
&\text{Fluss } \gamma_X(t)=\exp(tX): \text{kontinuierliche Transformationsgeschichte,}\\
&\text{Darstellung } U: G\to\mathcal U(\mathcal H): \text{Übertragung auf Beschreibungen,}\\
&\text{Orbit } \mathcal O_\psi: \text{physikalischer Zustand (Äquivalenzklasse von Beschreibungen).}\\
&\text{Genuin physikalisch: Invarianzstrukturen, Relationen zwischen Orbits, Symmetriebrüche.}
\end{aligned}
}
\]

\textbf{(10) Kommutierende Operatoren und die Beschreibung physikalischer Systeme.}

Ein physikalisches System wird in der Regel durch eine Menge von Observablen beschrieben,
die durch selbstadjungierte Operatoren $\hat A_i$ auf $\mathcal H$ repräsentiert werden.
Die algebraische Struktur dieser Operatoren bestimmt, welche physikalischen Größen
\emph{gleichzeitig bestimmbar} und \emph{invariant} sind.

\emph{(a) Kommutierende Operatoren.}
Zwei Observablen $\hat A,\hat B$ heißen \emph{kompatibel} oder \emph{kommutierend}, wenn
\[
[\hat A,\hat B]=0.
\]
Dann besitzen sie ein gemeinsames System von Eigenzuständen, das heißt:
es existiert eine Basis $\{\psi_\alpha\}\subset\mathcal H$ mit
\[
\hat A\psi_\alpha = a_\alpha\psi_\alpha,\qquad \hat B\psi_\alpha=b_\alpha\psi_\alpha.
\]
Diese Situation definiert eine \emph{kommutative Familie} von Operatoren,
deren Spektren die gleichzeitig messbaren Werte der entsprechenden Observablen darstellen.

Kommutierende Operatoren sind die quantenmechanische Formulierung
von \emph{klassisch verträglichen Größen}, die in denselben Koordinaten simultan definiert sind.
Ihre gemeinsame Spektralzerlegung strukturiert den Hilbertraum in
Unterräume, die jeweils eine \emph{vollständige Beschreibung eines möglichen Zustands} liefern.

\emph{(b) Rolle in der Gruppensymmetrie.}
Ist eine Observable $\hat A$ invariant unter der Wirkung von $G$
($U(g)^\dagger\hat A\,U(g)=\hat A$ oder $[\hat A,\hat X]=0$ für alle $X\in\mathfrak g$),
so kommutiert sie mit allen Generatoren der Symmetrie.
In diesem Fall beschreibt $\hat A$ eine physikalische Größe,
deren Messwert entlang aller Flüsse $\Phi_X(t,\psi)$ konstant bleibt:
\[
\frac{d}{dt}\langle \Phi_X(t,\psi),\hat A\,\Phi_X(t,\psi)\rangle = 0.
\]
Somit sind kommutierende Operatoren die \emph{quantitativen Repräsentanten von Invarianzstrukturen}
unter der Lie-Gruppenwirkung.

\emph{(c) Maximale kommutierende Familien.}
Eine Menge $\{\hat A_1,\ldots,\hat A_n\}$ selbstadjungierter Operatoren heißt \emph{maximal kommutierend},
wenn kein weiterer unabhängiger Operator $\hat B$ existiert, der mit allen $\hat A_i$ kommutiert.
Eine solche Familie definiert ein gemeinsames Eigenbasissystem $\{\psi_\lambda\}$, das den Hilbertraum vollständig
in simultane Eigenräume zerlegt. Physikalisch entsprechen diese Eigenzustände den \emph{stationären oder stabilen Beschreibungen}
des Systems, deren Entwicklung durch die entsprechenden Flüsse trivial (Phasenentwicklung) ist.

\emph{(d) Beispiel (Hamilton-Formalismus).}
Der Hamiltonoperator $\hat H$ ist Generator der Zeitentwicklung.
Ein Zustand $\psi$ erfüllt das Eigenwertproblem
\[
\hat H\psi = E\psi
\]
genau dann, wenn seine zeitliche Entwicklung stationär ist:
\[
\psi(t) = e^{-\,\frac{i}{\hbar}Et}\psi(0)=e^{-\,\frac{i}{\hbar}Et}\psi.
\]
Alle Erwartungswerte $\langle\psi(t),\hat O\psi(t)\rangle$
bleiben konstant für Observablen $\hat O$, die mit $\hat H$ kommutieren.
Daher sind Eigenzustände von $\hat H$ die \emph{zeitinvarianten Beschreibungen}
unter der durch $\hat H$ erzeugten Einparametergruppe.
Die Eigenwerte $E$ sind die \emph{charakteristischen Invarianten}
dieser Gruppe, formal die Spektralparameter des Flusses $\Phi_H(t,\psi)$.

\emph{(e) Eigenwertprobleme in der Symmetrieanalyse.}
Allgemein liefert das Eigenwertproblem
\[
\hat X\,\psi = \lambda\psi,\qquad X\in\mathfrak g,
\]
die stationären Beschreibungen bezüglich der durch $X$ erzeugten Transformation.
Physikalisch bedeutet das:
- $\psi$ ist \emph{invariant bis auf eine Phasenänderung} unter der durch $X$ erzeugten Symmetrie,
- $\lambda$ ist der \emph{Messwert} der zugehörigen physikalischen Größe (z.\,B.\ Impuls, Drehimpuls, Energie),
- das Spektrum $\{\lambda\}$ beschreibt die möglichen invarianten Quantitäten.

Damit sind Eigenwertprobleme die mathematische Form der Aussage,
dass ein Zustand entlang eines gegebenen Gruppenflusses
nur eine triviale (phasenartige) Veränderung erfährt — also „stationär“ bleibt.

\emph{(f) Verknüpfung mit Orbits.}
Zwei Eigenzustände mit demselben Satz von Eigenwerten kommutierender Observablen
liegen auf demselben Orbit der Symmetriegruppe und repräsentieren denselben physikalischen Zustand.
Unterschiedliche Eigenwertkombinationen definieren verschiedene Orbitklassen.
Das Spektrum der kommutierenden Operatoren bestimmt somit die topologische und algebraische
Struktur der Orbitenmannigfaltigkeit, d.\,h.\ die Raumstruktur der physikalischen Zustände selbst.

\textbf{(11) Zusammenfassung.}
\[
\boxed{
\begin{aligned}
&\text{Kommutierende Operatoren} \;\Rightarrow\; \text{gleichzeitig messbare Größen, Invarianzen;}\\
&\text{Eigenwertprobleme} \;\Rightarrow\; \text{stationäre Beschreibungen (invariante Richtungen der Flüsse);}\\
&\text{Eigenwerte} \;\Rightarrow\; \text{charakteristische Parameter physikalischer Zustände (Quantenzahlen);}\\
&\text{Orbitstruktur} \;\Rightarrow\; \text{physikalische Zustände als Äquivalenzklassen gemeinsamer Eigenräume.}
\end{aligned}
}
\]





\textbf{(12) Lie-Gruppen, Generatoren und Invarianz physikalischer Gleichungen.}

\emph{(a) Dynamische Gleichungen als Flüsse auf dem Zustandsraum.}
Sei $G$ eine Lie-Gruppe, die auf dem (komplexen) Zustandsraum $\mathcal H$ mit Darstellung $U:G\to\mathcal U(\mathcal H)$ wirkt.
Ein Operator $E:\mathcal D(E)\subset\mathcal H\to\mathcal H$ (Differential- oder Integraloperator)
definiert eine physikalische Gleichung
\[
E[\psi]=0,
\]
deren Lösungsmenge $\mathcal S=\ker(E)$ den Raum der zulässigen Beschreibungen bildet.
Für jeden Generator $\hat X=i\hbar\,dU(X)$, $X\in\mathfrak g$, erzeugt der Fluss
\[
\Phi_X(t,\psi)=U(\exp tX)\psi
\]
eine Einparametergruppe von Transformationen auf $\mathcal S$.
Die infinitesimale Invarianzbedingung lautet
\[
[dU(X),E]=0 \qquad\Longleftrightarrow\qquad E[U(\exp tX)\psi]=U(\exp tX)E[\psi].
\]
Dann sind alle Punkte des Flusses $\Phi_X(t,\psi)$ wieder Lösungen derselben Gleichung.

\emph{(b) Bedeutung der Invarianz.}
Die Bedingung $[dU(X),E]=0$ besagt:
\begin{itemize}
\item Der Operator $E$ ist \emph{kovariant} unter der durch $X$ erzeugten Transformation — 
  die zugrunde liegende physikalische Gleichung ist forminvariant.
\item Der Fluss $\Phi_X(t)$ erzeugt keine neue Dynamik, sondern verändert nur die \emph{Darstellung} der Lösung.
\item Wird die Invarianz verletzt ($[dU(X),E]\neq0$), wird der Fluss dynamisch:
  die Transformation ändert den physikalischen Zustand selbst (Symmetriebruch, äußeres Feld, Wechselwirkung).
\end{itemize}

\emph{(c) Kommutierende Generatoren und Integrabilität.}
Besitzen mehrere unabhängige Generatoren $X_i\in\mathfrak g$ kommutierende Repräsentanten $[\hat X_i,\hat X_j]=0$,
so sind die zugehörigen Flüsse integrierbar:
\[
U(\exp t_i X_i)\,U(\exp t_j X_j)=U(\exp t_j X_j)\,U(\exp t_i X_i),
\]
und der gemeinsame Fluss ist wohldefiniert:
\[
\Phi_{\mathbf t}(\psi)=e^{-\,\frac{i}{\hbar}\sum_i t_i\hat X_i}\psi.
\]
Physikalisch bedeutet das:
\begin{itemize}
\item Alle Größen $\hat X_i$ sind \emph{kompatibel} und gleichzeitig messbar.
\item Der Raum $\mathcal S$ lässt sich in gemeinsame Eigenräume zerlegen.
\item Diese Eigenräume sind invariant unter der Wirkung der Symmetriegruppe.
\end{itemize}

\emph{(d) Eigenwertprobleme und stationäre Lösungen.}
Ein Eigenwertproblem
\[
E[\psi]=\lambda\psi
\]
oder im Spezialfall $\hat H\psi=E\psi$
definiert Zustände, die unter der zugehörigen Einparametergruppe stationär bleiben:
\[
U(t)\psi=e^{-\,\frac{i}{\hbar}\lambda t}\psi.
\]
Damit sind Eigenzustände die \emph{stationären Richtungen} der Flüsse,
und die Eigenwerte $\lambda$ sind die invarianten Kennzahlen (Spektralparameter) dieser Bewegung.
In einem kommutativen Algebrazusammenhang $\{\hat X_i\}$ liefern gemeinsame Eigenwerte
$\lambda_i$ ein vollständiges Koordinatensystem auf der Raumstruktur der Orbitklassen physikalischer Zustände.

\emph{(e) Invarianz der Feldgleichungen.}
Für Differentialoperatoren $E[\psi]=0$ (etwa Klein–Gordon, Dirac, Maxwell, Yang–Mills)
bedeutet die Kovarianz
\[
E[U(g)\psi]=U(g)E[\psi],\qquad g\in G,
\]
dass die Gleichung die Struktur der zugrunde liegenden Raumzeit- oder Eichsymmetrie bewahrt.
Diese Bedingung bestimmt den Aufbau des Operators $E$:
\begin{itemize}
\item In der Relativistik erzwingt sie Lorentz-Kovarianz 
  ($E$ aus $\eta_{\mu\nu},\,\gamma^\mu$).
\item In der Eichtheorie erzwingt sie die Verwendung kovarianter Ableitungen $\nabla_\mu$
  mit Verbindung $A_\mu$.
\item In der Quantenmechanik garantiert sie die Erhaltung der Wahrscheinlichkeit
  (Unitarität des Flusses).
\end{itemize}

\emph{(f) Rolle physikalischer Gleichungen im Symmetriegefüge.}
Eine physikalische Gleichung identifiziert die \emph{zulässigen Beschreibungen} innerhalb des gesamten Hilbertraums:
\[
E[\psi]=0 \;\;\Longrightarrow\;\; \psi\in \mathcal S\subset\mathcal H.
\]
Die Lie-Gruppenstruktur legt fest, wie diese Beschreibungen unter Transformationen verknüpft werden;
Invarianz garantiert, dass die Menge der Lösungen $\mathcal S$ durch die Gruppenwirkung stabil bleibt:
\[
U(g)\mathcal S=\mathcal S.
\]
Damit verbindet die Invarianzstruktur das \emph{Gesetz} (Operatorgleichung) mit der \emph{Symmetrie} (Gruppenwirkung)
und legt fest, welche Gleichungen überhaupt physikalisch konsistent sind.

\textbf{(13) Zusammenfassung.}
\[
\boxed{
\begin{aligned}
&\text{Operatorgleichungen } E[\psi]=0 \text{ definieren dynamische oder stationäre Gesetze.}\\
&\text{Invarianz } [dU(X),E]=0 \text{ garantiert, dass die Gleichung Symmetrien von }G\text{ respektiert.}\\
&\text{Kommutierende Generatoren } \Rightarrow \text{ integrierbare und gleichzeitig beobachtbare Größen.}\\
&\text{Eigenwertprobleme } \Rightarrow \text{ stationäre Richtungen der Flüsse, charakteristische Spektren.}\\
&\text{Lösungsräume } \mathcal S\text{ sind durch Gruppenwirkung stabil, Orbits bilden physikalische Zustände.}\\
&\text{Gesetz, Symmetrie und Invarianz bilden die dreifache Struktur:}\\
&\qquad\text{(i) dynamischer Operator }E, \quad (ii)\text{ Symmetriegruppe }G, \quad (iii)\text{ Invarianzrelation }[E,dU(\mathfrak g)]=0.
\end{aligned}
}
\]



\textbf{(14) Allgemeine Kurven in Lie-Gruppen und ihre induzierten Flüsse.}

\emph{(a) Motivation.}
Einparameter-Untergruppen $\exp(tX)$ beschreiben homogene, zeitunabhängige Prozesse,
bei denen die Richtung der Transformation in der Lie-Algebra $\mathfrak g$ konstant bleibt.
In vielen physikalischen Situationen — etwa bei zeitabhängigen Hamiltonoperatoren,
nicht-stationären Feldern oder zusammengesetzten Transformationen — variiert die Richtung der Veränderung mit der Zeit.
Dazu benötigt man allgemeine glatte Kurven
\[
g:\mathbb R \to G,\qquad g(0)=e,\qquad g(t)\in G,
\]
die \emph{keine Untergruppen} sind, sondern von einer zeitabhängigen Lie-Algebra-Kurve $X(t)\in\mathfrak g$ erzeugt werden.

\emph{(b) Differentialgleichung in der Lie-Gruppe.}
Für eine glatte Kurve $g(t)\in G$ existiert (lokal eindeutig) eine Lie-Algebra-Kurve $X(t)\in\mathfrak g$, so dass
\[
\dot g(t) = (L_{g(t)})_*\,X(t),\qquad g(0)=e.
\]
Dies definiert den \emph{linksinvarianten Geschwindigkeitsvektor}
\[
X(t) := (L_{g(t)^{-1}})_*\dot g(t)\in\mathfrak g.
\]
Umgekehrt wird $g(t)$ durch $X(t)$ eindeutig bestimmt:
\[
g(t)=\mathcal{P}\exp\!\left(\int_0^t X(s)\,ds\right),
\]
wobei $\mathcal P$ die Zeitordnung bezeichnet.
Für konstantes $X$ reduziert sich dies zu $g(t)=\exp(tX)$.

\emph{(c) Induzierte Flüsse auf dem Zustandsraum.}
Über eine Darstellung $U:G\to\mathcal U(\mathcal H)$ induziert $g(t)$ einen Fluss von Zuständen
\[
\Phi(t,\psi) = U(g(t))\psi.
\]
Differenziation ergibt die \emph{zeitabhängige Evolutionsgleichung}
\[
\frac{d}{dt}\Phi(t,\psi) = -\,\frac{i}{\hbar}\,\hat X(t)\,\Phi(t,\psi),\qquad 
\hat X(t)=i\hbar\,dU(X(t)).
\]
Damit ist jedes zeitabhängige physikalische Gesetz (z.\,B.\ $i\hbar\dot\psi=\hat H(t)\psi$)
die Darstellung eines Flusses, der durch eine allgemeine Kurve in $G$ erzeugt wird.

\emph{(d) Physikalische Bedeutung.}
\begin{itemize}
\item Allgemeine Kurven $g(t)$ beschreiben \emph{nicht-stationäre oder zusammengesetzte Entwicklungen}:
      zeitabhängige Hamiltonoperatoren, beschleunigte Beobachter, zeitvariable Eichtransformationen usw.
\item Das zugehörige $X(t)$ beschreibt zu jedem Zeitpunkt die \emph{momentane Richtung der Veränderung} im Raum der Beschreibungstransformationen.
\item Der Fluss $\Phi(t,\psi)$ ist die \emph{integrierte Entwicklung der Zustandsbeschreibung} entlang dieser sukzessiven infinitesimalen Richtungen.
\end{itemize}

\emph{(e) Beispiel (zeitabhängige Dynamik).}
In der Quantenmechanik:
\[
i\hbar\frac{d}{dt}\psi(t)=\hat H(t)\psi(t)
\quad\Longleftrightarrow\quad
g(t)=\mathcal P\exp\!\left(-\frac{i}{\hbar}\int_0^t \hat H(s)\,ds\right),
\qquad \psi(t)=g(t)\psi(0).
\]
Hier ist $G=U(\mathcal H)$, $\mathfrak g=\mathfrak u(\mathcal H)$, und $X(t)=-\frac{i}{\hbar}\hat H(t)$.
Ist $\hat H$ zeitunabhängig, so reduziert sich $g(t)=e^{-\,\frac{i}{\hbar}t\hat H}$ auf eine Einparameter-Untergruppe.

\emph{(f) Nichtkommutative Zusammensetzung und geometrische Struktur.}
Für zwei zeitabhängige Kurven $X_1(t),X_2(t)$ mit nichtkommutierenden Werten $[X_1(t),X_2(s)]\neq0$
ist die zusammengesetzte Bewegung nicht einfach additiv:
\[
\mathcal P\exp\!\left(\int_0^t (X_1+X_2)\,ds\right)
\;\neq\;
\mathcal P\exp\!\left(\int_0^t X_1\,ds\right)
\mathcal P\exp\!\left(\int_0^t X_2\,ds\right).
\]
Die Nichtkommutativität erzeugt eine \emph{Krümmungsstruktur} im Raum der Beschreibungstransformationen;
sie ist formal das, was in der Eichgeometrie als Feldstärke $F=dA+\tfrac12[A,A]$ erscheint.
Damit beschreiben allgemeine Kurven $g(t)$ \emph{gekrümmte Pfade im Raum der Symmetrien}, deren lokale Nichtkommutativität
physikalisch Wechselwirkung oder Kopplung ausdrückt.

\emph{(g) Zusammenfassende Interpretation.}
\begin{itemize}
\item Einparameter-Untergruppen ($X=$ konstant): stationäre, homogene Prozesse (freie Dynamik, reine Symmetrieflüsse).
\item Allgemeine Kurven $g(t)$: nichtstationäre, kombinierte oder feldabhängige Prozesse (zeitabhängige Dynamik, Wechselwirkungen, Transport).
\item Die Lie-Algebra-Kurve $X(t)$ liefert die lokale Richtung der Beschreibung,
  der Pfad $g(t)$ ihre integrierte Geschichte,
  und die Darstellung $U(g(t))$ die konkrete Entwicklung der Zustandsbeschreibung.
\item Physikalisch liegt hier der Übergang von \emph{lokaler Struktur} (infinitesimale Generatoren) 
  zu \emph{globaler Dynamik} (integrierte Prozesse).
\end{itemize}

\textbf{(15) Zusammenfassung.}
\[
\boxed{
\begin{aligned}
&\text{Allgemeine Kurven } g(t)\in G \text{ repräsentieren sukzessive Transformationen;}\\
&X(t)=(L_{g(t)^{-1}})_*\dot g(t) \text{ ist die momentane Richtung in } \mathfrak g;\\
&\Phi(t,\psi)=U(g(t))\psi \text{ beschreibt die fortlaufende Entwicklung der Beschreibung; }\\
&\text{Nichtkommutativität der } X(t) \text{ erzeugt geometrische Krümmung (Wechselwirkungen).}\\
&\text{Stationäre Prozesse: } X=\text{konstant},\quad\text{allgemeine Dynamik: } X=X(t). 
\end{aligned}
}
\]


[Aufspaltung entarteter Eigenräume durch kommutierende Operatoren]

\textbf{(1) Ausgangssituation.}
Sei $G$ eine Lie-Gruppe mit Darstellung
\[
U:G\longrightarrow \mathcal U(\mathcal H),
\]
und $\mathfrak g$ ihre Lie-Algebra mit abgeleiteter Darstellung
\[
dU:\mathfrak g\longrightarrow \mathfrak u(\mathcal H),
\qquad
X\mapsto dU(X).
\]
Sei $\hat H=i\hbar\,dU(X_H)$ der selbstadjungierte Operator,
der die Zeittranslation (oder allgemein eine beobachtbare Symmetrie) generiert.

Nehmen wir an, $\hat H$ habe entartete Eigenräume:
\[
\hat H\psi_{E,a}=E\,\psi_{E,a},\qquad a=1,\ldots, n_E.
\]
Dann ist jeder Energieeigenraum
\[
V_E:=\operatorname{span}\{\psi_{E,a}\}_{a=1}^{n_E}
\]
eine \emph{invariante Unterdarstellung} von $U$ — ein Unterraum, auf dem $G$ (bzw.\ $U(G)$) noch nicht vollständig reduziert ist.

\textbf{(2) Lie-theoretische Sicht.}
Die Entartung bedeutet: die Einschränkung von $U$ auf den Einparameter-Fluss $\exp(tX_H)$ ist nicht irreduzibel;
es existieren nichttriviale Invarianten in $V_E$.

Die Aufspaltung erfolgt, indem man zusätzliche Operatoren $\{\hat A_i\}$ betrachtet, die mit $\hat H$ kommutieren:
\[
[\hat H,\hat A_i]=0,\qquad [\hat A_i,\hat A_j]=0.
\]
Diese Operatoren bilden eine abelsche Unteralgebra
\[
\mathfrak a:=\operatorname{span}\{dU(X_i)\}\subset \mathfrak u(\mathcal H),
\]
die mit $dU(X_H)$ kommutiert:
\[
[dU(X_H),\mathfrak a]=0.
\]
Damit wird die Darstellung auf $\mathcal H$ durch eine größere Untergruppe $A:=\exp(\mathfrak a)$ gemeinsam diagonalisierbar.

\textbf{(3) Diagrammatische Interpretation.}
\[
\begin{tikzcd}[column sep=huge,row sep=large]
\mathfrak g \arrow{r}{dU} \arrow[dashed,swap]{dr}{\text{Reduktion auf }\mathfrak a}
& \mathfrak u(\mathcal H)
  \arrow{r}{[\ ,\ ]=0 \text{ auf }\mathfrak a}
  & \text{Gemeinsame Spektralunterräume}
\\
\mathfrak a\subset\mathfrak g \arrow{ur}{dU|_{\mathfrak a}} \arrow{r}{\exp}
& A=\exp(\mathfrak a)
  \arrow{ur}{U|_A}
\end{tikzcd}
\]

- Der Schritt $\mathfrak g\to\mathfrak a$ bedeutet:  
  Auswahl einer \emph{maximal abelschen Unteralgebra} (MAU) innerhalb der Symmetrie.  
- Deren Exponential $A=\exp(\mathfrak a)$ ist eine abelsche Untergruppe von $G$ (meist ein Torus).
- Die Darstellung $U|_A$ ist vollständig diagonalisierbar; ihre Eigenwerte sind die \emph{quantenzahlen}.

\textbf{(4) Physikalische Bedeutung.}
\begin{itemize}
\item Entartung eines Spektrums $\hat H\psi=E\psi$ heißt:
  es existiert eine Rest-Symmetriegruppe $G_E\subset G$, deren Wirkung den Eigenraum $V_E$ invariant lässt.
\item Kommutierende Operatoren $\{\hat A_i\}$ sind die Generatoren dieser Restgruppe $A\subset G_E$.
\item Ihre gleichzeitige Diagonalisierung liefert eine feinere Zerlegung
  \[
  V_E = \bigoplus_{\alpha} V_{E,\alpha},
  \]
  wobei $\alpha$ ein \emph{Gewichtsvektor} (Quantenzahl) der gemeinsamen abelschen Unteralgebra ist.
\item Diese Zerlegung entspricht der Aufspaltung des Orbits von $U(G)$ in kleinere Orbits von $U(A)$ innerhalb des Eigenraums.
\end{itemize}

\textbf{(5) Top-down Erklärung.}
Von oben betrachtet (Diagramm-Ebene):
\[
\begin{tikzcd}[column sep=huge,row sep=large]
\text{Lie-Algebra } \mathfrak g \arrow[dashed]{r}{\text{MAU } \mathfrak a} \arrow{d}{dU}
& \text{Abelsche Unteralgebra } \mathfrak a \arrow{d}{dU|_{\mathfrak a}}
\\
\text{Operator-Algebra } \mathfrak u(\mathcal H) \arrow[dashed]{r}{\text{Kommutative Unteralgebra}}
& \text{Selbstadjungierte Kommutanten } \{\hat H,\hat A_i\}
\end{tikzcd}
\]

Die Reduktion
\[
\mathfrak g \longrightarrow \mathfrak a
\]
ist die Lie-theoretische Entsprechung der
„Suche nach einem vollständigen Satz kommutierender Observablen“ (CSCO):
eine maximal abelsche Unteralgebra der Operatorenalgebra bestimmt eine simultane Spektralzerlegung.

\textbf{(6) Formale Zusammenfassung.}
\[
\boxed{
\begin{aligned}
&\text{Entartung:} &[\hat H,\hat A]=0\ \forall A\in\mathfrak g_E
 &\Rightarrow\text{reduzierte Symmetriegruppe }G_E,\\[2pt]
&\text{Aufspaltung:} &\mathfrak a\subset\mathfrak g_E\text{ maximal abelsch}
 &\Rightarrow \text{gemeinsame Eigenräume }V_{E,\alpha},\\[2pt]
&\text{Physikalisch:}&
(E,\alpha)\text{ bilden vollständige Satz von Quantenzahlen.}
\end{aligned}
}
\]

\textbf{(7) Beispiel.}
Für ein Zentralpotential:
\[
\hat H=\frac{\hat{\mathbf p}^2}{2m}+V(r),
\]
gilt $[\hat H,\hat L^2]=[\hat H,\hat L_z]=0$.
Dann ist
\[
\mathfrak a = \operatorname{span}\{dU(X_{L^2}),dU(X_{L_z})\}
\]
abelsch, und die gemeinsame Spektralzerlegung liefert
\[
|E,\ell,m\rangle,
\]
die Aufspaltung der entarteten Energieniveaus $E$ nach Drehimpuls-Quantenzahlen $(\ell,m)$.

\textbf{(8) Interpretative Formulierung.}
Im Diagramm-Setting bedeutet das:
\[
\text{Ein entarteter Eigenraum ist eine nichttriviale Bahn eines Restflusses in } \mathcal U(\mathcal H),
\]
und das Einführen kommutierender Operatoren
entspricht dem \emph{Zerlegen dieses Orbits in Unterorbits} der abelschen Untergruppe $A$.



\pagebreak



[Formale Struktur der dynamischen Seite der Quantenmechanik]

\textbf{(1) Motivation und Rolle.}
Die \emph{dynamische Seite} der Quantenmechanik modelliert die \emph{Gesetze der Veränderung}
von Zuständen und Observablen, unabhängig von der Wahl der Beschreibung oder Symmetrie.
Während die \emph{Symmetrie-Seite} beschreibt, \emph{welche Transformationen} eine Theorie
invariant lassen, beschreibt die \emph{Dynamik-Seite}, \emph{wie Zustände sich entwickeln}.

Diese Dynamik entsteht aus einer \emph{Lie-Struktur}, die bereits klassisch vorhanden ist,
und wird durch die Quantisierung in die Operatorwelt übertragen.

\medskip
\textbf{(2) Klassischer struktureller Rahmen.}

\begin{enumerate}[label=\textbf{(C\arabic*)}]
\item \textbf{Phasenraum:}
      Sei $(\mathcal P,\omega)$ eine symplektische Mannigfaltigkeit,
      typischerweise $\mathcal P=T^*Q$ für einen Konfigurationsraum $Q$.
      \[
      \omega = \sum_{i=1}^n dq^i \wedge dp_i.
      \]
      $(\mathcal P,\omega)$ ist der Raum der \emph{klassischen Zustände}
      (vollständige Information über Ort und Impuls).
\item \textbf{Observablen:}
      Eine Observable ist eine glatte Funktion $f\in C^\infty(\mathcal P,\mathbb R)$.
      Sie ordnet jedem Zustand $\xi\in\mathcal P$ einen Messwert $f(\xi)$ zu.
\item \textbf{Poisson-Klammer:}
      Auf $C^\infty(\mathcal P)$ ist durch die symplektische Form
      die \emph{Poisson-Klammer}
      \[
      \{f,g\}
      =\omega(X_f,X_g)
      =\sum_{i=1}^n
      \left(
        \frac{\partial f}{\partial q^i}\frac{\partial g}{\partial p_i}
        -\frac{\partial f}{\partial p_i}\frac{\partial g}{\partial q^i}
      \right)
      \]
      definiert.
      $(C^\infty(\mathcal P),\{\ ,\ \})$ ist eine Lie-Algebra.
\item \textbf{Hamiltonscher Vektorfeldfluss:}
      Jede Observable $f$ erzeugt über $\omega$ ein Vektorfeld $X_f$:
      \[
      \iota_{X_f}\omega = df.
      \]
      Dieses $X_f$ beschreibt die infinitesimale Bewegung auf $\mathcal P$, erzeugt
      durch $f$, mit Fluss
      \[
      \Phi^f_t = \exp(tX_f):\mathcal P\to\mathcal P.
      \]
\item \textbf{Zeitentwicklung:}
      Für den Hamiltonian $H\in C^\infty(\mathcal P)$ gilt:
      \[
      \dot f = \{H,f\},\qquad \dot\xi = X_H(\xi),
      \]
      d.h.\ der Fluss $\Phi^H_t$ bestimmt die Dynamik des Systems auf $\mathcal P$.
\item \textbf{Lie-Struktur der Dynamik:}
      Der Zusammenhang
      \[
      [X_f,X_g]=X_{\{f,g\}}
      \]
      zeigt, dass die Zuordnung
      $f\mapsto X_f$ ein Lie-Algebrahomomorphismus ist:
      \[
      (C^\infty(\mathcal P),\{\ ,\ \})
      \longrightarrow
      (\mathfrak X_{\mathrm{Ham}}(\mathcal P),[\ ,\ ]).
      \]
\end{enumerate}

\medskip
\textbf{Interpretation:}
\begin{itemize}
\item Elemente von $C^\infty(\mathcal P)$ sind \emph{Messgrößen / Observablen}.
\item Die Poisson-Klammer ist die infinitesimale Beschreibung der \emph{Nichtkommutativität der Flüsse}.
\item Die Flüsse $\Phi^f_t$ sind \emph{kanonische Transformationen} (symplektische Diffeomorphismen),
      die das physikalische System in sich selbst bewegen.
\item Die Dynamik eines Systems ist also nicht einfach „Zeitentwicklung einer Zahl“,
      sondern \emph{Fluss einer Funktion auf einem symplektischen Raum}.
\end{itemize}

\medskip
\textbf{(3) Physikalische Bedeutung.}
\begin{itemize}
\item Jede Observable $f$ definiert eine „Richtung der Veränderung“ im Raum der Zustände — 
      ihr Hamiltonsches Vektorfeld.
\item Die Poisson-Klammer $\{f,g\}$ misst, wie stark die durch $f$ und $g$ erzeugten Flüsse \emph{nicht kommutieren};
      sie ist also die „Infinitesimalform der Wechselwirkung“.
\item Der Hamiltonian $H$ erzeugt den Fluss, der die \emph{Zeitentwicklung der Zustände} beschreibt:
      \[
      \dot\xi = X_H(\xi),\quad \xi(t)=\Phi^H_t(\xi_0).
      \]
\end{itemize}

\medskip
\textbf{(4) Quantisierung als strukturelle Übersetzung.}

Die Quantisierung ist ein funktorieller Übergang zwischen der
\emph{klassischen Lie-Algebra der Observablen}
und der \emph{Operator-Lie-Algebra der Quantenmechanik}:

\[
\begin{tikzcd}[column sep=huge,row sep=large]
(C^\infty(\mathcal P),\{\ ,\ \})
\arrow{r}{f\mapsto X_f}
\arrow{d}{\mathcal Q:f\mapsto\hat f}
& \mathfrak X_{\mathrm{Ham}}(\mathcal P)
\arrow{r}{\exp}
& \mathrm{Ham}(\mathcal P)
\\
(\mathfrak u(\mathcal H),\tfrac{1}{i\hbar}[\ ,\ ])
\arrow{r}{A\mapsto \mathrm{ad}_A}
& \mathrm{Der}(\mathfrak u(\mathcal H))
\arrow{r}{\exp}
& \mathcal U(\mathcal H)
\end{tikzcd}
\]

Dabei gilt:
\[
\boxed{
\mathcal Q:\{f,g\}\;\mapsto\;\frac{1}{i\hbar}[\hat f,\hat g].
}
\]

\begin{enumerate}[label=\textbf{(Q\arabic*)}]
\item \textbf{Zuordnung:}
      $f\in C^\infty(\mathcal P)\mapsto\hat f$ (selbstadjungierter Operator) auf $\mathcal H$.
\item \textbf{Flüsse:}
      $\Phi_t^f=\exp(tX_f)$ wird zu einem unitären Fluss
      $U_f(t)=\exp(-\,\frac{i}{\hbar}t\hat f)$ auf $\mathcal H$.
\item \textbf{Zeitentwicklung:}
      $H\in C^\infty(\mathcal P)\mapsto \hat H$ mit Schrödinger-Gleichung
      $i\hbar\,\partial_t\psi=\hat H\psi$.
\item \textbf{Klammern:}
      $\{f,g\}\leftrightarrow\tfrac{1}{i\hbar}[\hat f,\hat g]$ stellt die Lie-Struktur sicher.
\end{enumerate}

\medskip
\textbf{(5) Physikalische Interpretation der Quantisierung.}

\begin{itemize}
\item Sie ersetzt den \emph{Raum der Zustände} $(\mathcal P,\omega)$ durch einen \emph{Hilbertraum} $(\mathcal H,\langle\cdot,\cdot\rangle)$.
\item Sie ersetzt Observablen $f$ (als reelle Funktionen) durch Operatoren $\hat f$,
      die nun \emph{nicht-kommutativ} sind.
\item Sie ersetzt klassische Flüsse $\Phi^f_t$ durch unitäre Operatorflüsse $U_f(t)$.
\item Die Poisson-Klammer — das Maß der Nichtkommutativität von Flüssen — wird
      durch den Kommutator ersetzt, der die Nichtkommutativität von Operatoren misst.
\item Damit werden die \emph{dynamischen Gesetze} selbst in dieselbe Lie-Struktur
      eingebettet, die auch die \emph{Symmetrien} tragen:
      \[
      \{\,\cdot\,,\,\cdot\,\} \longleftrightarrow \tfrac{1}{i\hbar}[\,\cdot\,,\,\cdot\,].
      \]
\end{itemize}

\medskip
\textbf{(6) Gemeinsamer struktureller Raum.}

Beide Seiten – klassisch und quantisiert – leben in einem \emph{Lie-Algebraischen Raum der Flüsse}:

\[
\begin{tikzcd}[column sep=huge,row sep=large]
\text{Klassisch: } (C^\infty(\mathcal P),\{\ ,\ \})
\arrow{r}{f\mapsto X_f}
& (\mathfrak X_{\mathrm{Ham}}(\mathcal P),[\ ,\ ]) \\
\text{Quantisiert: } (\mathfrak u(\mathcal H),\tfrac{1}{i\hbar}[\ ,\ ])
\arrow{r}{A\mapsto \mathrm{ad}_A}
& (\mathrm{Der}(\mathfrak u(\mathcal H)),[\ ,\ ])
\end{tikzcd}
\]
%
In beiden Fällen sind die Objekte:
\begin{itemize}
\item \emph{Generatoren} (Observablen oder Operatoren),
\item \emph{Flüsse} (Hamiltonsche bzw.\ unitäre Flüsse),
\item \emph{Lie-Klammer}, die die infinitesimale Kompositionsstruktur kodiert.
\end{itemize}

\medskip
\textbf{(7) Physikalischer Gehalt der dynamischen Struktur.}

\begin{itemize}
\item Die Lie-Struktur der Dynamik ist die mathematische Kodierung der
      \emph{Gesetzmäßigkeit der Veränderung von Zuständen}.
\item Ihre Quantisierung macht diese Gesetzmäßigkeit kompatibel mit dem probabilistischen
      (Hilbertraum-)Rahmen.
\item Die klassische Dynamik lebt auf einem kommutativen Algebra von Observablen,
      die Quantenmechanik auf einer nichtkommutativen;
      der Übergang beschreibt also den \emph{Verlust der simultanen Messbarkeit}.
\item Die Nichtkommutativität (Kommutator/Poisson) ist nicht technisches Artefakt,
      sondern das \emph{Infinitesimalgesetz} der Dynamik selbst.
\end{itemize}

\medskip
\textbf{(8) Fazit.}
Die \emph{dynamische Seite} der Quantenmechanik besteht aus:
\[
\boxed{
\begin{aligned}
&\text{Objekte: } f\in C^\infty(\mathcal P) \leftrightarrow \hat f\in\mathfrak{herm}(\mathcal H),\\[3pt]
&\text{Operation: } \{f,g\}\leftrightarrow\tfrac{1}{i\hbar}[\hat f,\hat g],\\[3pt]
&\text{Flüsse: } \Phi^f_t=\exp(tX_f)\leftrightarrow U_f(t)=e^{-\,\frac{i}{\hbar}t\hat f},\\[3pt]
&\text{Gesetz: } \dot f = \{H,f\}\leftrightarrow i\hbar\,\dot{\hat f}=[\hat f,\hat H].
\end{aligned}
}
\]
Damit wird Dynamik als \emph{Lie-Flussgesetz} realisiert,
dessen strukturelle Gestalt unter der Quantisierung erhalten bleibt.



\pagebreak

[Top–Down: Strukturelle Ziele der Quantenmechanik und die zwei Lie-Algebren]

\textbf{Zielbild.}
Eine Quantenmechanik (QM) einer gegebenen physikalischen Theorie soll vier \emph{strukturelle} Anforderungen erfüllen:
\begin{enumerate}[label=\textbf{Z\arabic*}.]
\item \textbf{Zustandsraum und Wahrscheinlichkeit} (Kinematik):
      Ein komplexer Hilbertraum $\mathcal H$ (bzw.\ projektiver Raum $\mathbb P(\mathcal H)$) mit Born-Regel.
\item \textbf{Dynamik} (Gesetz der Zeitentwicklung):
      Eine reversible Zeitentwicklung als \emph{unitärer Fluss} $U(t)=e^{-\,\frac{i}{\hbar}t\hat H}$ (Stone),
      erzeugt durch einen selbstadjungierten Hamiltonoperator $\hat H$.
\item \textbf{Symmetrieprinzip} (Kovarianz/Erhaltungsgrößen):
      Eine (projektiv) unitäre Darstellung $U:G\to\mathcal U(\mathcal H)$ der externen Symmetriegruppe $G$
      (Raum-/Zeit-/interne Symmetrien) mit $U(\exp X)=e^{\,dU(X)}$.
\item \textbf{Observablen und Messung}:
      Observablen als selbstadjungierte Operatoren $\hat A$ (Spektraltheorie),
      kompatibel mit (Z1)–(Z3).
\end{enumerate}

\medskip
\textbf{Zwei Lie–Algebren, zwei Rollen.}
Diese vier Ziele trennen sich \emph{notwendig} in zwei Lie–algebraische Ebenen, die zusammenwirken:

\begin{center}
\begin{tabular}{lll}

\textbf{Ebene} & \textbf{Lie–Algebra} & \textbf{Rolle}\\

\emph{(I) Symmetrie/Kinematik} & $(\mathfrak g,[\ ,\ ])$ & externe Symmetrien der Theorie (Rotation, Translation, Lorentz, Gauge)\\
\emph{(II) Dynamik/Observablen} & $(\mathfrak g_{\mathrm{dyn}},\{\ ,\ \})$ & interne Flussstruktur (Poisson/Heisenberg), d.\,h.\ Bewegungsgesetz\\

\end{tabular}
\end{center}

\textbf{(I) Symmetrie–Lie–Algebra.}
Die Daten $(G,\mathfrak g)$ fixieren \emph{wie Beschreibungen transformieren}:
\[
\begin{tikzcd}[column sep=large,row sep=large]
\mathfrak g \arrow{r}{\exp} \arrow{d}{dU}
& G \arrow{d}{U}
\\
\mathfrak u(\mathcal H) \arrow{r}{\exp}
& \mathcal U(\mathcal H)
\end{tikzcd}
\quad\text{(kommutativ)}.
\]
Hier liefern $dU(X)$ die \emph{kinematischen Generatoren} (Impuls, Drehimpuls, Ladung).
Sie kodieren \emph{Forminvarianz} (Kovarianz) und – über Noether – Erhaltungsgrößen, aber \emph{nicht} die konkrete Dynamik.

\textbf{(II) Dynamische Lie–Algebra.}
Die \emph{Bewegung} eines Systems ist klassisch ein \emph{Hamiltonscher Fluss} auf dem Phasenraum $(\mathcal P,\omega)$:
\[
(C^\infty(\mathcal P),\{\ ,\ \})\quad\text{mit}\quad \dot f=\{H,f\},\ X_H=\{H,\cdot\},\ \Phi_t=\exp(tX_H).
\]
Für endlich viele Freiheitsgrade spannt die \emph{Heisenberg–Lie–Algebra}
\[
\mathfrak h_n=\mathrm{span}\{q^i,p_i,\mathbf 1\},\qquad \{q^i,p_j\}=\delta^i_j
\]
die \emph{minimale} dynamische Struktur auf (kanonische Paare).
In QFT tritt die \emph{kontinuierliche} Version (CCR/CAR) an die Stelle.

\medskip
\textbf{Warum beide Ebenen \emph{separat} notwendig sind.}
\begin{itemize}
\item \emph{Symmetrie} beantwortet: Welche Transformationen verändern nur die \emph{Beschreibung} (Beobachterwechsel, interne Phasen), ohne Physik zu ändern?  
      — Sie liefert $U(G)$, Casimirs, Auswahlregeln, aber keine Bewegungsgleichung.
\item \emph{Dynamik} beantwortet: Wie entwickelt sich ein Zustand \emph{in der Zeit}?  
      — Sie verlangt ein \emph{Gesetz des Flusses} (Hamilton/Heisenberg), also eine zweite Lie–Klammer (Poisson/Kommutator).
\end{itemize}

\medskip
\textbf{Natürlicher Zusammenhang: Quantisierung als Lie–Funktor.}
Die Brücke zwischen (II) und der Operatorseite ist die \emph{Quantisierung}
\[
\mathcal Q:\ (C^\infty(\mathcal P),\{\ ,\ \})\ \longrightarrow\ \big(\mathfrak u(\mathcal H),\tfrac{1}{i\hbar}[\ ,\ ]\big),
\qquad
\frac{1}{i\hbar}[\hat f,\hat g]\ \approx\ \widehat{\{f,g\}}.
\]
Damit wird \emph{jeder klassische Fluss} zu einem \emph{unitären Fluss}:
\[
\Phi^f_t=\exp(tX_f)\quad\leadsto\quad U_f(t)=\exp\!\Big(-\frac{i}{\hbar}t\,\hat f\Big).
\]
Besonders: $\hat H$ erzeugt die Zeitentwicklung $U(t)=e^{-\,\frac{i}{\hbar}t\hat H}$ (Stone).

\medskip
\textbf{Kompatibilität mit der Symmetrie–Ebene.}
Die Dynamik muss \emph{kovariant} sein:
\[
U(g)\,e^{-\,\frac{i}{\hbar}t\hat H}\,U(g)^{-1}
\;=\;
e^{-\,\frac{i}{\hbar}t\,\hat H_g},
\qquad
\hat H_g:=U(g)\hat H U(g)^{-1}.
\]
Bei echter Symmetrie gilt $\hat H_g=\hat H$ und $[dU(X),\hat H]=0$ für $X\in\mathfrak g$ (Erhaltungssatz).
Somit verknüpft die Darstellung $U$ die \emph{äußeren} Generatoren $dU(\mathfrak g)$ mit der \emph{inneren} Flussstruktur $\hat H$.

\medskip
\textbf{Heisenberg–Gruppe als gemeinsamer Träger der Dynamik.}
Die exponentielle Weyl–Form der CCR erfordert die \emph{zentrale Erweiterung} $\exp(\mathfrak h_n)=H_n$.
Die Schrödinger–Darstellung $U:H_n\to\mathcal U(\mathcal H)$ realisiert
\[
dU(q^i)=\frac{1}{i\hbar}\hat q^i,\qquad dU(p_i)=\frac{1}{i\hbar}\hat p_i,\qquad
[\hat q^i,\hat p_j]=i\hbar\,\delta^i_j\,\mathbf 1,
\]
und ist (Stone–von Neumann) bis auf Unitäräquivalenz eindeutig: die kanonische Quantisierung ist dadurch \emph{strukturell fixiert}.

\medskip
\textbf{Gesamtdiagramm (Synthese).}
\[
\begin{tikzcd}[column sep=huge,row sep=large]
\text{Klassik/Dynamik: } C^\infty(\mathcal P) \arrow{r}{f\mapsto X_f} \arrow{d}{\mathcal Q:f\mapsto\hat f}
& \mathfrak X_{\mathrm{Ham}}(\mathcal P) \arrow{r}{\exp}
& \mathrm{Ham}(\mathcal P)
\\[2pt]
\text{Quant/Observablen: } \mathfrak u(\mathcal H) \arrow{r}{[\ ,\ ]/(i\hbar)} \arrow{d}{\exp}
& \mathfrak u(\mathcal H) \arrow{r}{\exp}
& \mathcal U(\mathcal H)
\\[2pt]
\text{Symmetrie/Kinematik: } \mathfrak g \arrow{r}{\exp} \arrow{u}{dU}
& G \arrow{u}{U} \arrow{r}{\text{Wirkung}}
& \mathcal H
\end{tikzcd}
\]
Oben: \emph{dynamische} Flüsse (Poisson) $\to$ Unten Mitte: \emph{unitäre} Flüsse (Kommutator);  
unten: \emph{symmetrische} Darstellungen (Kovarianz).  
Die Quantisierung $\mathcal Q$ ist die vertikale Brücke, $dU,U$ koppeln Symmetrie und Dynamik.

\medskip
\textbf{Schlussfolgerung.}
\begin{enumerate}[label=\textbf{S\arabic*}.]
\item Die \emph{Symmetrie–Lie–Algebra} $(\mathfrak g,[\ ,\ ])$ ist für \emph{Kinematik/Kovarianz} notwendig (Transformationsgesetz, Erhaltung).
\item Die \emph{dynamische Lie–Algebra} $(\mathfrak g_{\mathrm{dyn}},\{\ ,\ \})$ ist für \emph{Bewegungsgesetze} notwendig (Flüsse).
\item Die \emph{Quantisierung} ist der \emph{natürliche} Lie–funktorielle Übergang, der (II) in Operatorflüsse der Ebene (I) überträgt:
      \[
      \{\,\cdot\,,\,\cdot\,\}\;\longmapsto\;\tfrac{1}{i\hbar}[\,\cdot\,,\,\cdot\,],\qquad
      \Phi_t\;\longmapsto\;e^{-\,\frac{i}{\hbar}t(\cdot)}.
      \]
\item Die \emph{Notwendigkeit der Trennung} (Kinematik vs.\ Dynamik) folgt aus den \emph{unterschiedlichen Trägern}
      (Beschreibungsänderungen vs.\ Bewegungsflüsse);
      der \emph{natürliche Zusammenhang} folgt aus der gemeinsamen Lie–Struktur und ihrer funktoriellen Übersetzung.
\end{enumerate}



\pagebreak

[Dynamische Struktur am Beispiel des harmonischen Oszillators]

\textbf{(1) Klassische Struktur.}

Sei $Q=\mathbb R$, Phasenraum $\mathcal P=T^*Q\simeq\mathbb R^2$ mit Koordinaten $(q,p)$ und
symplektischer Form $\omega=dq\wedge dp$.
Die Poisson-Klammer ist
\[
\{f,g\}=\frac{\partial f}{\partial q}\frac{\partial g}{\partial p}
        -\frac{\partial f}{\partial p}\frac{\partial g}{\partial q}.
\]
%
Hamiltonfunktion:
\[
H(q,p)=\frac{p^2}{2m}+\frac{1}{2}m\omega_0^2 q^2.
\]

\textbf{(2) Hamiltonscher Fluss.}

Hamiltonsches Vektorfeld:
\[
X_H=\{H,\cdot\}
= \frac{\partial H}{\partial p}\,\frac{\partial}{\partial q}
  -\frac{\partial H}{\partial q}\,\frac{\partial}{\partial p}
= \frac{p}{m}\frac{\partial}{\partial q}
  -m\omega_0^2 q\,\frac{\partial}{\partial p}.
\]
%
Bewegungsgleichungen:
\[
\dot q=\frac{p}{m},\qquad
\dot p=-m\omega_0^2 q.
\]
%
Lösung (Hamilton-Fluss):
\[
\Phi_t(q,p)
=\Bigl(q\cos(\omega_0 t)+\frac{p}{m\omega_0}\sin(\omega_0 t),\;
       p\cos(\omega_0 t)-m\omega_0 q\sin(\omega_0 t)\Bigr).
\]
%
Das ist eine \emph{Rotation} im $(q,p)$-Phasenraum.
Die Trajektorien sind Ellipsen, Energie $H(q,p)=\text{const}$.

\textbf{(3) Lie-Struktur der Observablen.}

Die elementaren Observablen $q,p$ erfüllen:
\[
\{q,p\}=1,\quad
\{H,q\}=-\frac{p}{m},\quad
\{H,p\}=m\omega_0^2 q.
\]
%
Die Unteralgebra
\[
\mathfrak h_1=\mathrm{span}\{q,p,\mathbf 1\},\qquad
\{q,p\}=1,
\]
ist die \emph{Heisenberg-Lie-Algebra} (klassisch, abelsche Erweiterung).
Ihr Exponentialbild ist die \emph{Heisenberg-Gruppe} $H_1$,
deren Elemente $(q,p,s)$ wirken auf $\mathcal P$ als kanonische Transformationen.

\textbf{(4) Quantisierung.}

Quantisierung $\mathcal Q: (C^\infty(\mathcal P),\{\ ,\ \})\to(\mathfrak u(\mathcal H),\tfrac{1}{i\hbar}[\ ,\ ])$:
\[
\mathcal Q(q)=\hat q,\quad
\mathcal Q(p)=\hat p,\quad
\frac{1}{i\hbar}[\hat q,\hat p]=1.
\]
%
Im Ortsraum $\mathcal H=L^2(\mathbb R,dq)$:
\[
(\hat q\psi)(q)=q\psi(q),\qquad
(\hat p\psi)(q)=-i\hbar\frac{d}{dq}\psi(q).
\]
%
Hamiltonoperator:
\[
\hat H
=\frac{\hat p^2}{2m}+\frac{1}{2}m\omega_0^2\hat q^2.
\]
%
Damit:
\[
i\hbar\dot\psi(t)=\hat H\psi(t),\qquad
\psi(t)=e^{-\,\frac{i}{\hbar}t\hat H}\psi(0).
\]

\textbf{(5) Operator-Lie-Struktur und Fluss.}

\begin{itemize}
\item \emph{Kommutatoren:}
      \[
      [\hat q,\hat p]=i\hbar,\qquad
      [\hat H,\hat q]=-\frac{i\hbar}{m}\hat p,\qquad
      [\hat H,\hat p]=i\hbar m\omega_0^2\hat q.
      \]
      Dies entspricht exakt der klassischen Struktur mit
      $\{\ ,\ \}\leftrightarrow[\ ,\ ]/(i\hbar)$.
\item \emph{Heisenberg-Bild:}
      Für Operatoren $\hat q(t),\hat p(t)$ gilt
      \[
      i\hbar\dot{\hat q}=[\hat q,\hat H]=\frac{\hbar}{m}\hat p,\qquad
      i\hbar\dot{\hat p}=[\hat p,\hat H]=-i\hbar m\omega_0^2\hat q,
      \]
      d.h.
      \[
      \ddot{\hat q}+\omega_0^2\hat q=0,
      \]
      also derselbe Fluss wie klassisch.
\end{itemize}

\textbf{(6) Lie-algebraische und Gruppenstruktur.}

Die Operatoren
\[
\hat a=\sqrt{\frac{m\omega_0}{2\hbar}}\Bigl(\hat q+\frac{i}{m\omega_0}\hat p\Bigr),
\quad
\hat a^\dagger=\sqrt{\frac{m\omega_0}{2\hbar}}\Bigl(\hat q-\frac{i}{m\omega_0}\hat p\Bigr)
\]
erfüllen
\[
[\hat a,\hat a^\dagger]=1.
\]
%
In diesen Variablen
\[
\hat H=\hbar\omega_0\Bigl(\hat a^\dagger\hat a+\tfrac{1}{2}\Bigr).
\]
%
Die $\hat a,\hat a^\dagger$ spannen die \emph{Heisenberg-Lie-Algebra}
\[
\mathfrak h_1=\mathrm{span}\{\hat a,\hat a^\dagger,\mathbf 1\},
\qquad
[\hat a,\hat a^\dagger]=\mathbf 1,
\]
deren Exponentialgruppe $H_1$ durch die Weyl-Operatoren
\[
W(\alpha)=e^{\alpha \hat a^\dagger-\alpha^*\hat a}
\]
realisiert wird.
$H_1$ wirkt unitar auf $\mathcal H$,
und der Hamiltonfluss ist eine Einparameter-Untergruppe:
\[
U_H(t)=e^{-\,\frac{i}{\hbar}t\hat H}=e^{-\,i\omega_0 t(\hat a^\dagger\hat a+\tfrac{1}{2})}.
\]

\textbf{(7) Zusammenfassung im Diagramm.}

\[
\begin{tikzcd}[column sep=huge,row sep=large]
\text{Klassisch: } & (q,p)\in T^*\mathbb R
\arrow{r}{\Phi_t}
& (q_t,p_t)
\\
\text{Quantisiert: } & (\hat q,\hat p)\in\mathfrak u(\mathcal H)
\arrow{r}{U_H(t)=e^{-\,\frac{i}{\hbar}t\hat H}}
& (\hat q(t),\hat p(t))=U_H(t)(\hat q,\hat p)U_H(t)^{-1}
\end{tikzcd}
\]
%
Beide Flüsse sind durch dieselbe Lie-Struktur bestimmt:
\[
\{H,\cdot\}\quad\longleftrightarrow\quad\frac{1}{i\hbar}[\hat H,\cdot].
\]
%
Die Quantisierung überführt also die \emph{klassische Rotationsdynamik im Phasenraum}
in eine \emph{unitäre Rotation im Hilbertraum}, 
erzeugt durch denselben Lie-Generator.

\textbf{(8) Physikalische Interpretation.}
\begin{itemize}
\item Klassisch: Bewegung einer Masse $m$ im Potential $V=\tfrac{1}{2}m\omega_0^2 q^2$; 
      Trajektorien sind Kreise in $(q,p)$.
\item Quantisiert: Zustände $\psi_n(q)$ bilden eine Orthonormalbasis von Energieeigenzuständen:
      \[
      \hat H\psi_n=\hbar\omega_0(n+\tfrac{1}{2})\psi_n.
      \]
      Die Zeitentwicklung $U_H(t)$ ist die unitäre Darstellung
      der Einparameter-Untergruppe der „Phasenrotationen“ im Fockraum.
\item Strukturell: dieselbe Lie-Struktur (Heisenberg-Algebra) trägt sowohl
      die klassische Dynamik (Poisson-Flüsse) als auch die quantisierte (Kommutator-Flüsse).
\end{itemize}

\textbf{(9) Fazit.}
\[
\boxed{
\begin{aligned}
&\text{Klassisch: } (C^\infty(T^*\mathbb R),\{\ ,\ \}),\quad
H=\tfrac{p^2}{2m}+\tfrac{1}{2}m\omega_0^2q^2,\\
&\text{Quantisiert: } (\mathfrak u(L^2(\mathbb R)),\tfrac{1}{i\hbar}[\ ,\ ]),\quad
\hat H=\tfrac{\hat p^2}{2m}+\tfrac{1}{2}m\omega_0^2\hat q^2,\\
&\text{Lie-Brücke: } \{H,\cdot\}\leftrightarrow\tfrac{1}{i\hbar}[\hat H,\cdot],\\
&\text{Fluss: } \Phi_t=\exp(tX_H)\leftrightarrow U_H(t)=e^{-\,\frac{i}{\hbar}t\hat H}.
\end{aligned}
}
\]

\pagebreak



[Dynamische Struktur eines reellen Skalarfeldes als verallgemeinerte Lie-Struktur]

\textbf{(1) Klassischer Ausgangspunkt.}

Ein reelles Skalarfeld $\phi:\mathbb R^{1,3}\to\mathbb R$ ist ein Feld auf der Raumzeit
$M=\mathbb R^{1,3}$ mit Lagrangedichte
\[
\mathscr L
=\tfrac{1}{2}\,(\partial_\mu\phi)(\partial^\mu\phi)
-\tfrac{1}{2}\,m^2\phi^2.
\]
Die Euler–Lagrange–Gleichung ist
\[
(\Box+m^2)\phi(x)=0,
\]
also die Klein–Gordon–Gleichung.

\textbf{(2) Kanonische Variablen und Hamiltonstruktur.}

Wähle einen festen Zeitschnitt $t=\mathrm{const}$.
Dann sind
\[
\phi(\mathbf x):=\phi(t,\mathbf x),\qquad
\Pi(\mathbf x):=\frac{\partial\mathscr L}{\partial(\partial_0\phi)}=\dot\phi(t,\mathbf x)
\]
\emph{kanonisch konjugierte Felder.}
%
Die Hamiltondichte lautet
\[
\mathscr H(\mathbf x)
=\tfrac{1}{2}\Big(\Pi^2+(\nabla\phi)^2+m^2\phi^2\Big),
\]
und die Hamiltonfunktion (Gesamtenergie)
\[
H[\phi,\Pi]=\int_{\mathbb R^3}\mathrm d^3x\;\mathscr H(\mathbf x).
\]

\textbf{(3) Unendlichdimensionale Poisson-Struktur.}

Die \emph{kanonische Poisson-Klammer} für Funktionale $F[\phi,\Pi]$,
$G[\phi,\Pi]$ auf dem Phasenraum
\[
\mathcal P=\{(\phi,\Pi)\mid \phi,\Pi:\mathbb R^3\to\mathbb R\}
\]
ist definiert durch
\[
\{F,G\}
=\int_{\mathbb R^3}\mathrm d^3x\,
\left(
\frac{\delta F}{\delta\phi(\mathbf x)}\frac{\delta G}{\delta\Pi(\mathbf x)}
-\frac{\delta F}{\delta\Pi(\mathbf x)}\frac{\delta G}{\delta\phi(\mathbf x)}
\right).
\]
%
Daraus folgen die fundamentalen Relationen:
\[
\{\phi(\mathbf x),\phi(\mathbf y)\}=0,\quad
\{\Pi(\mathbf x),\Pi(\mathbf y)\}=0,\quad
\{\phi(\mathbf x),\Pi(\mathbf y)\}=\delta^3(\mathbf x-\mathbf y).
\]
%
Damit ist
\[
(C^\infty_{\mathrm{loc}}(\mathcal P),\{\ ,\ \})
\]
eine \emph{(unendlichdimensionale) Poisson-Lie-Algebra.}

\textbf{(4) Dynamik als Hamiltonscher Fluss.}

Die Zeitentwicklung wird durch den Hamiltonian $H$ erzeugt:
\[
\dot\phi(\mathbf x,t)=\{\phi(\mathbf x),H\}=\Pi(\mathbf x,t),\qquad
\dot\Pi(\mathbf x,t)=\{\Pi(\mathbf x),H\}=\nabla^2\phi(\mathbf x,t)-m^2\phi(\mathbf x,t),
\]
also die Klein–Gordon–Gleichung
\[
(\partial_t^2-\nabla^2+m^2)\phi=0.
\]

\textbf{(5) Quantisierung: Übergang zur Operator-Lie-Struktur.}

Die Quantisierung
\[
\mathcal Q:
(C^\infty(\mathcal P),\{\ ,\ \})
\;\longrightarrow\;
(\mathfrak u(\mathcal H),\tfrac{1}{i\hbar}[\ ,\ ])
\]
besteht formal aus:

\begin{enumerate}[label=\textbf{(Q\arabic*)}]
\item \emph{Zuordnung:}
      \[
      \phi(\mathbf x)\mapsto\hat\phi(\mathbf x),\qquad
      \Pi(\mathbf x)\mapsto\hat\Pi(\mathbf x),
      \]
      d.h.\ Felder werden zu Operatorwerten.
\item \emph{Kommutatorregeln (kanonische Vertauschungsrelationen):}
      \[
      [\hat\phi(\mathbf x,t),\hat\phi(\mathbf y,t)]=0,\quad
      [\hat\Pi(\mathbf x,t),\hat\Pi(\mathbf y,t)]=0,\quad
      [\hat\phi(\mathbf x,t),\hat\Pi(\mathbf y,t)]=i\hbar\,\delta^3(\mathbf x-\mathbf y).
      \]
      Dies ist die kontinuierliche Version der Heisenberg-Algebra.
\item \emph{Hamiltonoperator:}
      \[
      \hat H(t)
      =\int\mathrm d^3x\;\hat{\mathscr H}(\mathbf x)
      =\tfrac{1}{2}\int\mathrm d^3x\,
      \Big(\hat\Pi^2+(\nabla\hat\phi)^2+m^2\hat\phi^2\Big).
      \]
\item \emph{Zeitentwicklung:}
      \[
      i\hbar\,\frac{\partial\hat\phi}{\partial t}
      =[\,\hat\phi,\hat H\,],\qquad
      i\hbar\,\frac{\partial\hat\Pi}{\partial t}
      =[\,\hat\Pi,\hat H\,],
      \]
      was wieder die Operatorform der Klein–Gordon–Gleichung ergibt.
\end{enumerate}

\textbf{(6) Strukturelle Interpretation.}

\begin{itemize}
\item Der klassische Phasenraum $\mathcal P$ ist unendlichdimensional;
      seine Punkte sind Feldkonfigurationen $(\phi,\Pi)$.
\item Die Poisson-Klammer auf $\mathcal P$ liefert eine unendlichdimensionale
      \emph{Lie-Algebra der Observablen} (lokale Funktionale).
\item Die Quantisierung ersetzt diese durch eine
      \emph{kontinuierliche Heisenberg-Algebra}
      mit $\delta^3(\mathbf x-\mathbf y)$ als strukturellem Kern:
      \[
      [\hat q_x,\hat p_y]=i\hbar\,\delta^3(x-y).
      \]
\item Damit wird jeder Punkt $\mathbf x$ formal zu einem
      unabhängigen Oszillator mit
      \[
      H=\int\mathrm d^3x\,H_{\mathrm{osc}}(\mathbf x).
      \]
      Die Feldtheorie ist also die unendliche Kopplung von harmonischen Oszillatoren.
\item Der unitäre Fluss
      \[
      U_H(t)=\exp\!\Big(-\frac{i}{\hbar}t\hat H\Big)
      \]
      ist die direkte Verallgemeinerung des Flusses $U(t)=e^{-\,\frac{i}{\hbar}tH}$
      des eindimensionalen Oszillators.
\end{itemize}

\textbf{(7) Strukturelle Einordnung in die Gesamtarchitektur.}

\[
\begin{tikzcd}[column sep=huge,row sep=large]
\text{Klassische Felddynamik: } (C^\infty(\mathcal P),\{\ ,\ \})
\arrow{r}{f\mapsto X_f}
\arrow{d}{\mathcal Q}
& \mathfrak X_{\mathrm{Ham}}(\mathcal P)\arrow{r}{\exp}
& \mathrm{Ham}(\mathcal P)
\\
\text{Quantisierte Felder: } (\mathfrak u(\mathcal H),\tfrac{1}{i\hbar}[\ ,\ ])
\arrow{r}{A\mapsto \mathrm{ad}_A}
& \mathrm{Der}(\mathfrak u(\mathcal H))\arrow{r}{\exp}
& \mathcal U(\mathcal H)
\end{tikzcd}
\]
%
Die vertikale Abbildung $\mathcal Q$ ist formal identisch mit derjenigen des
harmonischen Oszillators; sie wirkt jetzt nur auf eine unendliche Menge von
kanonischen Freiheitsgraden, also ein kontinuierliches Tensorprodukt von
Heisenberg-Algebren.

\textbf{(8) Physikalische Bedeutung.}

\begin{itemize}
\item Die Felder $\hat\phi(\mathbf x)$ und $\hat\Pi(\mathbf x)$ sind Operatoren
      im Fockraum, die jedem Raumzeitpunkt eine Observablenstruktur zuordnen.
\item Die Gleichung
      \[
      [\,\hat\phi(\mathbf x),\hat\Pi(\mathbf y)\,]=i\hbar\delta^3(\mathbf x-\mathbf y)
      \]
      bedeutet, dass die Theorie lokal dieselbe infinitesimale
      Lie-Struktur besitzt wie die Quantenmechanik, aber über einen kontinuierlichen Indexraum.
\item Die Quantisierung ist somit keine neue Regel,
      sondern die \emph{Erweiterung der Poisson--Kommutator-Brücke} auf unendlich viele Freiheitsgrade.
\end{itemize}

\textbf{(9) Fazit.}
\[
\boxed{
\begin{aligned}
&\text{Klassische Seite: } 
(C^\infty(\mathcal P),\{\ ,\ \}),\;
\mathcal P=\{(\phi,\Pi)\},\;
\{\,\phi(\mathbf x),\Pi(\mathbf y)\,\}=\delta^3(\mathbf x-\mathbf y);\\[3pt]
&\text{Quantisierte Seite: } 
(\mathfrak u(\mathcal H),\tfrac{1}{i\hbar}[\ ,\ ]),\;
[\,\hat\phi(\mathbf x),\hat\Pi(\mathbf y)\,]=i\hbar\delta^3(\mathbf x-\mathbf y);\\[3pt]
&\text{Flussgesetz: } 
\dot F=\{H,F\}
\;\longleftrightarrow\;
i\hbar\dot{\hat F}=[\hat F,\hat H];\\[3pt]
&\text{Dynamische Struktur: }
\text{lokale Heisenberg-Algebra,}
\quad
\text{Hamiltonfluss }U_H(t)=e^{-\,\frac{i}{\hbar}t\hat H}.
\end{aligned}
}
\]


[Formale Konstruktion der Quantisierung eines reellen Skalarfeldes]

\textbf{(1) Klassische Struktur.}

\begin{enumerate}[label=\textbf{(C\arabic*)}]
\item \textbf{Feldraum.}
      Das reelle Skalarfeld ist eine glatte Abbildung
      \[
      \phi: M=\mathbb R^{1,3}\longrightarrow\mathbb R,
      \qquad
      x=(t,\mathbf x).
      \]
      Der Raum aller glatten Felder wird bezeichnet mit
      \[
      \mathcal F := C^\infty(M,\mathbb R).
      \]

\item \textbf{Lagrangedichte und Lagrange-Funktional.}
      \[
      \mathscr L(\phi,\partial_\mu\phi)
      =\tfrac{1}{2}(\partial_\mu\phi)(\partial^\mu\phi)
       -\tfrac{1}{2}m^2\phi^2,
      \qquad
      L[\phi]=\int_{\mathbb R^3}\mathrm d^3x\,\mathscr L.
      \]

\item \textbf{Phasenraum.}
      Wähle einen Zeitschnitt $t=\mathrm{const}$.
      Dann definiere den unendlichdimensionalen Phasenraum
      \[
      \mathcal P:=T^*\mathcal C,\qquad
      \mathcal C:=C^\infty(\mathbb R^3,\mathbb R),
      \]
      also den Kotangentialbündelraum über den Raum der Feldkonfigurationen.
      Ein Punkt in $\mathcal P$ ist ein Paar
      \[
      (\phi,\Pi),\quad
      \phi:\mathbb R^3\to\mathbb R,\quad
      \Pi:\mathbb R^3\to\mathbb R.
      \]

\item \textbf{Symplektische Form und Poisson-Klammer.}
      Auf $\mathcal P$ ist die (formale) symplektische Form
      \[
      \Omega((\delta\phi_1,\delta\Pi_1),(\delta\phi_2,\delta\Pi_2))
      =\int_{\mathbb R^3}\mathrm d^3x\,
       (\delta\phi_1\,\delta\Pi_2-\delta\phi_2\,\delta\Pi_1),
      \]
      und daraus folgt die Poisson-Klammer für Funktionale $F,G:\mathcal P\to\mathbb R$:
      \[
      \{F,G\}
      =\int_{\mathbb R^3}\mathrm d^3x\,
       \left(
         \frac{\delta F}{\delta\phi(\mathbf x)}\frac{\delta G}{\delta\Pi(\mathbf x)}
        -\frac{\delta F}{\delta\Pi(\mathbf x)}\frac{\delta G}{\delta\phi(\mathbf x)}
       \right).
      \]
      Damit ist $(C^\infty_{\mathrm{loc}}(\mathcal P),\{\ ,\ \})$
      eine (unendlichdimensionale) Lie-Algebra.

\item \textbf{Hamiltonfunktion.}
      \[
      H[\phi,\Pi]
      =\int_{\mathbb R^3}\mathrm d^3x\,\tfrac{1}{2}
        \left(\Pi^2+(\nabla\phi)^2+m^2\phi^2\right).
      \]
      Das erzeugt den Hamiltonschen Fluss:
      \[
      \dot\phi(\mathbf x,t)=\{\phi(\mathbf x),H\}=\Pi(\mathbf x,t),\qquad
      \dot\Pi(\mathbf x,t)=\{\Pi(\mathbf x),H\}=\nabla^2\phi(\mathbf x,t)-m^2\phi(\mathbf x,t).
      \]
      Dies entspricht $(\Box+m^2)\phi=0$.
\end{enumerate}

\medskip
\textbf{(2) Abstrakte Definition der Quantisierung.}

Die \emph{Quantisierung} ist eine Abbildung
\[
\mathcal Q:\ (C^\infty_{\mathrm{loc}}(\mathcal P),\{\ ,\ \})
\longrightarrow
(\mathfrak{Op},\tfrac{1}{i\hbar}[\ ,\ ]),
\]
wobei:
\begin{itemize}
\item $\mathfrak{Op}$ eine *-Algebra (meist eine C\(^*\)-Algebra oder von-Neumann-Algebra) von Operatoren ist,
      die auf einem Hilbertraum $\mathcal H$ wirkt;
\item $\mathcal Q$ die fundamentalen Zuordnungen
      \[
      \phi(\mathbf x)\longmapsto\hat\phi(\mathbf x),\qquad
      \Pi(\mathbf x)\longmapsto\hat\Pi(\mathbf x)
      \]
      erfüllt, sowie die Bedingung
      \[
      \frac{1}{i\hbar}[\hat F,\hat G]\approx\widehat{\{F,G\}}
      \]
      auf einer geeigneten dichten Domäne in $\mathcal H$.
\end{itemize}

\textbf{(3) Fundierende Relationen: kanonische Vertauschungsregeln (CCR).}

Aus der Poisson-Struktur folgt, dass die Operatoren $\hat\phi,\hat\Pi$
die CCR (Canonical Commutation Relations) erfüllen:
\[
[\hat\phi(\mathbf x,t),\hat\phi(\mathbf y,t)]=0,\quad
[\hat\Pi(\mathbf x,t),\hat\Pi(\mathbf y,t)]=0,\quad
[\hat\phi(\mathbf x,t),\hat\Pi(\mathbf y,t)]=i\hbar\,\delta^3(\mathbf x-\mathbf y).
\]
Diese bilden eine \emph{kontinuierliche Heisenberg-Algebra}:
\[
\mathfrak h_\infty
=\mathrm{span}\{\hat\phi(f),\hat\Pi(g),\mathbf 1\mid f,g\in C_0^\infty(\mathbb R^3)\},
\quad
[\hat\phi(f),\hat\Pi(g)]=i\hbar\int f(\mathbf x)g(\mathbf x)\,\mathrm d^3x.
\]

\textbf{(4) Darstellung auf einem Hilbertraum (Fockraum).}

Sei $\mathcal S(\mathbb R^3)$ der Schwartz-Raum.
Die Operatoren werden in Fouriermoden zerlegt:
\[
\hat\phi(\mathbf x)
=\int\frac{\mathrm d^3k}{(2\pi)^3\sqrt{2\omega_{\mathbf k}}}
 \left(\hat a_{\mathbf k}e^{i\mathbf k\cdot\mathbf x}
      +\hat a_{\mathbf k}^\dagger e^{-i\mathbf k\cdot\mathbf x}\right),
\qquad
\omega_{\mathbf k}=\sqrt{\mathbf k^2+m^2}.
\]
%
Damit gilt
\[
[\hat a_{\mathbf k},\hat a_{\mathbf k'}^\dagger]=(2\pi)^3\delta^3(\mathbf k-\mathbf k').
\]
%
Der zugehörige Hilbertraum ist der symmetrische Fockraum
\[
\mathcal H=\bigoplus_{n=0}^\infty L^2_{\mathrm{sym}}((\mathbb R^3)^n).
\]
Die Erzeuger– und Vernichteroperatoren $(\hat a_{\mathbf k},\hat a_{\mathbf k}^\dagger)$
wirken auf $\mathcal H$ gemäß
\[
\hat a_{\mathbf k}|n_1,\dots,n_{\mathbf k},\dots\rangle
=\sqrt{n_{\mathbf k}}\,|n_{\mathbf k}-1\rangle,\qquad
\hat a_{\mathbf k}^\dagger|n_{\mathbf k}\rangle
=\sqrt{n_{\mathbf k}+1}\,|n_{\mathbf k}+1\rangle.
\]

\textbf{(5) Hamiltonoperator und Fluss.}

Der quantisierte Hamiltonoperator lautet
\[
\hat H
=\int\frac{\mathrm d^3k}{(2\pi)^3}\,\hbar\omega_{\mathbf k}\,
 \hat a_{\mathbf k}^\dagger\hat a_{\mathbf k} + E_0,
\]
mit Nullpunktenergie $E_0=\frac{1}{2}\int\frac{\mathrm d^3k}{(2\pi)^3}\hbar\omega_{\mathbf k}$.
%
Der zugehörige unitäre Fluss auf $\mathcal H$ ist
\[
U_H(t)=\exp\!\Big(-\frac{i}{\hbar}t\hat H\Big),
\]
und die Operatoren $\hat\phi(\mathbf x,t)$ entwickeln sich nach
\[
\hat\phi(\mathbf x,t)=U_H(t)\hat\phi(\mathbf x,0)U_H(t)^{-1}.
\]

\textbf{(6) Formale Abbildung zwischen den Strukturräumen.}

\[
\begin{tikzcd}[column sep=huge,row sep=large]
\text{Klassische Poisson-Algebra} &
(C^\infty_{\mathrm{loc}}(\mathcal P),\{\ ,\ \})
\arrow{r}{f\mapsto X_f} \arrow{d}{\mathcal Q}
& \mathfrak X_{\mathrm{Ham}}(\mathcal P)
\\
\text{Quantisierte Operator-Algebra} &
(\mathfrak h_\infty,\tfrac{1}{i\hbar}[\ ,\ ])
\arrow{r}{A\mapsto\mathrm{ad}_A}
& \mathrm{Der}(\mathfrak h_\infty)
\end{tikzcd}
\]

Vertikal: Quantisierungsfunktor $\mathcal Q$  
— er ersetzt die kommutative Algebra glatter Funktionale durch die nichtkommutative *-Algebra der Operatoren.  
Horizontal: Lie-Klammern definieren infinitesimale Flüsse auf beiden Ebenen.

\textbf{(7) Physikalische Bedeutung.}

\begin{itemize}
\item \emph{Klassisch:} Der Punkt $(\phi,\Pi)\in\mathcal P$ beschreibt eine konkrete Feldkonfiguration (Amplitude + „Geschwindigkeit“).  
\item \emph{Quantisiert:} Zustände sind Wellenfunktionen im Fockraum, Operatoren $\hat\phi,\hat\Pi$ sind nichtkommutative Observablen.  
\item Die unendlichdimensionale Poisson-Struktur wird zur kontinuierlichen Heisenberg-Algebra.  
\item Die Quantisierung ist also formal dieselbe Lie-Funktor-Abbildung wie in der Quantenmechanik,  
      nur dass der Indexraum von endlichdimensional (\(i=1,\dots,n\)) zu kontinuierlich (\(\mathbf x\in\mathbb R^3\)) erweitert ist.
\end{itemize}

\textbf{(8) Schlussformel (Strukturzusammenfassung).}

\[
\boxed{
\begin{aligned}
&\text{Klassische Seite: } (C^\infty_{\mathrm{loc}}(\mathcal P),\{\ ,\ \}),\quad
\{\,\phi(\mathbf x),\Pi(\mathbf y)\,\}=\delta^3(\mathbf x-\mathbf y),\\[3pt]
&\text{Quantisierte Seite: } (\mathfrak h_\infty,\tfrac{1}{i\hbar}[\ ,\ ]),\quad
[\,\hat\phi(\mathbf x),\hat\Pi(\mathbf y)\,]=i\hbar\delta^3(\mathbf x-\mathbf y),\\[3pt]
&\text{Funktorielle Brücke: } 
\mathcal Q:\{f,g\}\longmapsto\tfrac{1}{i\hbar}[\hat f,\hat g],\\[3pt]
&\text{Trägerraum der Dynamik: } \mathcal H=\text{symmetrischer Fockraum über }L^2(\mathbb R^3),\\[3pt]
&\text{Fluss: } U_H(t)=e^{-\,\frac{i}{\hbar}t\hat H},\qquad
i\hbar\dot{\hat F}=[\hat F,\hat H].
\end{aligned}
}
\]


\pagebreak

[CCR-\(C^*\)-Algebra und Fock-Darstellung für das reelle Skalarfeld]

\textbf{(1) Symplektischer Präraum der klassischen Daten.}
Sei \(S\) ein reeller lokalkonvexer Vektorraum (z.\,B.\ glatte, kompakt getragene Cauchy-Daten
oder glatte Testfunktionen modulo Lösungen) und
\[
\sigma:S\times S\to\mathbb R
\]
eine stetige, alternierende, nicht-entartete Bilinearform (symplektische Form).
Im Minkowski-Fall kann man wählen:
\begin{itemize}
\item \(S=\mathcal D(\mathbb R^3,\mathbb R)\oplus\mathcal D(\mathbb R^3,\mathbb R)\) (Cauchy-Daten \((\phi,\Pi)\) am festen \(t\)),
\item
\(
\sigma\big((\phi,\Pi),(\phi',\Pi')\big)=\int_{\mathbb R^3}\!\big(\phi\,\Pi'-\phi'\,\Pi\big)\,d^3x.
\)
\end{itemize}

\textbf{(2) Weyl-\(*\)-Algebra (algebraische CCR).}
Die \emph{Weyl-\(*\)-Algebra} \(\mathcal W(S,\sigma)\) ist die universelle \(*\)-Algebra,
erzeugt von formalen Einheiten \(W(f)\) (\(f\in S\)) mit den \emph{Weyl-Relationen}
\[
\boxed{\
W(f)W(g)=e^{-\,\tfrac{i}{2}\sigma(f,g)}\,W(f+g),\qquad
W(f)^*=W(-f),\qquad W(0)=\mathbf 1.\
}
\]
Die universelle \(C^*\)-Algebra der CCR ist die Vollendung der \(*\)-Algebra bzgl.\ der
größten \(C^*\)-Halbnorm, für die alle \emph{regularen} \(*\)-Darstellungen stetig sind.
Wir schreiben diese Vollendung wieder als \(\mathfrak{A}_{\mathrm{CCR}}(S,\sigma)\).

\textbf{(3) Zustände und GNS-Darstellung.}
Ein \emph{Zustand} ist ein positives, normiertes Funktional \(\omega:\mathfrak{A}_{\mathrm{CCR}}\to\mathbb C\).
Die \emph{GNS-Konstruktion} liefert aus \((\mathfrak{A}_{\mathrm{CCR}},\omega)\) eine
Hilbertraumdarstellung \(\pi_\omega:\mathfrak{A}_{\mathrm{CCR}}\to\mathcal B(\mathcal H_\omega)\)
mit zyklischem Vektor \(\Omega_\omega\) (dem „Vakuum“), so dass
\(\omega(A)=\langle\Omega_\omega,\pi_\omega(A)\Omega_\omega\rangle\).

\textbf{(4) Quasi-freie (gaußsche) Zustände.}
Ein Zustand \(\omega\) heißt \emph{quasi-frei}, wenn er durch seine Zweipunktfunktion
\(\Lambda:S\times S\to\mathbb C\) bestimmt ist
und alle \(n\)-Punkt-Funktionen Wick-artig (Isserlis) aus \(\Lambda\) zusammengesetzt werden.
Konsistenzbedingungen:
\[
\mathrm{Im}\,\Lambda(f,g)=\tfrac{1}{2}\,\sigma(f,g),
\qquad
\Lambda(f,f)\ge 0,
\]
sowie geeignete Stetigkeit. Dann ist
\(\omega(W(f))=\exp\!\big(-\tfrac{1}{2}\Lambda(f,f)\big)\).

\textbf{(5) Einteilchenstruktur und komplexe Struktur.}
Eine (reelle) \emph{Einteilchenstruktur} ist ein komplexer Hilbertraum \(\mathfrak h\) mit Skalarprodukt
\(\langle\cdot,\cdot\rangle_{\mathfrak h}\) und eine reell-lineare, dicht-bildende Abbildung
\(K:S\to\mathfrak h\) derart, dass
\[
\boxed{\
\sigma(f,g)=2\,\mathrm{Im}\,\langle Kf,Kg\rangle_{\mathfrak h},\qquad
\mu(f,g):=2\,\mathrm{Re}\,\langle Kf,Kg\rangle_{\mathfrak h}\ \text{positiv definit}.\
}
\]
Äquivalent: Wahl einer \emph{komplexen Struktur} \(J:S\to S\) mit
\(J^2=-\mathrm{id}\), \(\sigma(J\cdot,J\cdot)=\sigma(\cdot,\cdot)\) und
\(\mu(\cdot,\cdot):=\sigma(\cdot,J\cdot)\) positiv, und setze dann
\(\mathfrak h\) als Komplettierung von \(S\) bzgl.\ \(\langle f,g\rangle_{\mathfrak h}=\tfrac{1}{2}\big(\mu(f,g)+i\sigma(f,g)\big)\)
und \(K=\mathrm{inkl}\).

\textbf{(6) Symmetrischer Fockraum und Feldoperatoren.}
Der \emph{Fockraum} über \(\mathfrak h\) ist
\[
\mathcal F_+(\mathfrak h)=\bigoplus_{n=0}^\infty \mathfrak h^{\otimes_s n}.
\]
Auf \(\mathcal F_+(\mathfrak h)\) existieren Erzeuger/Vernichter \(\,a^\dagger(\xi),a(\xi)\) für \(\xi\in\mathfrak h\)
mit \([a(\xi),a^\dagger(\eta)]=\langle\xi,\eta\rangle_{\mathfrak h}\).
Die (geschmierte) \emph{Feldoperatorfamilie} ist
\[
\boxed{\
\phi(f):=\frac{1}{\sqrt{2}}\big(a^\dagger(Kf)+a(Kf)\big),\qquad f\in S,\
}
\]
definiert auf dem endlichen Teilchenzahl-Domänenkern, essenziell selbstadjungiert für reelles \(f\).
Dann ist die \emph{Weyl-Darstellung}
\[
\pi(W(f)):=e^{\,i\phi(f)}\in\mathcal U(\mathcal F_+(\mathfrak h)),
\]
und die \emph{CCR} halten in geschmierter Form:
\[
[\phi(f),\phi(g)]=i\,\sigma(f,g)\,\mathbf 1,\qquad
\pi(W(f))\,\pi(W(g))=e^{-\,\tfrac{i}{2}\sigma(f,g)}\,\pi(W(f+g)).
\]
Der Vektor \(\Omega\) (Vakuum, \(a(\xi)\Omega=0\)) definiert einen reinen quasi-freien Zustand
\(\omega(A)=\langle\Omega,\pi(A)\Omega\rangle\).

\textbf{(7) Klein–Gordon-Fall (Minkowski): kanonische Wahl von \(\mathfrak h\).}
Für \((\Box+m^2)\phi=0\) mit positivem Energiekegel wähle
\[
\mathfrak h=L^2\big(\mathbb R^3,\tfrac{d^3k}{(2\pi)^3\,2\omega_{\mathbf k}}\big),
\qquad \omega_{\mathbf k}=\sqrt{|\mathbf k|^2+m^2},
\]
und definiere \(K\) als Projektion der Cauchy-Daten auf die \emph{positive Frequenz}:
\[
K:\ (\phi,\Pi)\mapsto \widetilde{K(\phi,\Pi)}(\mathbf k)
=\frac{1}{\sqrt{2}}\Big(\sqrt{\omega_{\mathbf k}}\,\tilde\phi(\mathbf k)
+\frac{i}{\sqrt{\omega_{\mathbf k}}}\,\tilde\Pi(\mathbf k)\Big).
\]
Dann ist \(\sigma(f,g)=2\,\mathrm{Im}\,\langle Kf,Kg\rangle\) erfüllt und die zugehörige Fock-Darstellung
ist \emph{Poincaré-invariant} und \emph{energiepositiv}.

\textbf{(8) Zeitentwicklung und Implementierung.}
Sei \(T_t:S\to S\) die klassisch-symplektische Evolution (hier: Klein–Gordon-Fluss auf Cauchy-Daten).
\(\{T_t\}_{t\in\mathbb R}\) ist eine symplektische Einparametergruppe:
\(\sigma(T_tf,T_tg)=\sigma(f,g)\).
\begin{itemize}
\item Ist \(T_t\) \emph{komplex linear} bzgl.\ der gewählten Einteilchenstruktur (\(KT_t=e^{-i h t}K\) für einen
selbstadjungierten \(h\) auf \(\mathfrak h\)), dann implementiert
\[
U(t)=\Gamma\big(e^{-iht}\big)=e^{-\,i\,d\Gamma(h)\,t}
\]
die Zeitentwicklung auf \(\mathcal F_+(\mathfrak h)\) und
\(\phi(T_tf)=U(t)\phi(f)U(t)^{-1}\).
\item Allgemein (zeitabhängige Hintergründe) gilt das \emph{Shale–Stinespring-Kriterium} für Implementierbarkeit:
\(T\) ist implementierbar in der Fock-Darstellung genau dann, wenn der „Bogoliubov-\(\beta\)-Teil“
bzgl.\ \(\mathfrak h\) Hilbert–Schmidt ist.
\end{itemize}

\textbf{(9) Eindeutigkeit und Nicht-Eindeutigkeit.}
\begin{itemize}
\item In \emph{endlicher} Dimension impliziert Stone–von Neumann die Einzigkeit (bis Unitäräquivalenz) der irreduziblen CCR-Darstellung.
\item In \emph{unendlicher} Dimension existieren i.\,A.\ \emph{viele unitar nichtäquivalente} Darstellungen.
   Physikalische Zusatzkriterien (Poincaré-Invarianz, Positivität der Energie, Hadamard-Bedingung, KMS/thermisch)
   selektieren bevorzugte quasi-freie Zustände und damit bevorzugte Fock-Darstellungen.
\end{itemize}

\textbf{(10) Fazit (wohl\-definierte Prozedur).}
Die Quantisierung des reellen Skalarfeldes ist wohldefiniert als:
\[
\boxed{\
(S,\sigma)\ \longmapsto\ \mathfrak{A}_{\mathrm{CCR}}(S,\sigma)\ \xrightarrow{\ \text{quasi-freier Zustand }\omega\ }\ 
\big(\pi_\omega,\mathcal H_\omega,\Omega_\omega\big)\
}
\]
und — bei Wahl einer Einteilchenstruktur \((\mathfrak h,K)\) — explizit als Fockdarstellung
\[
\pi(W(f))=e^{\,i\phi(f)},\qquad
\phi(f)=\tfrac{1}{\sqrt{2}}\big(a^\dagger(Kf)+a(Kf)\big),
\]
mit Zeitentwicklung \(U(t)=\Gamma(e^{-iht})\) (falls implementierbar).
Damit wird die algebraische CCR-Struktur \emph{unabhängig von Operator-Domänenproblemen} definiert
und erst anschließend durch GNS/Fock konkret realisiert.


\pagebreak

[Positionierung der Quantisierungsdiagramme im klassischen Lie-Gruppen-Setting]

\textbf{(1) Zwei komplementäre Rollen von Lie-Gruppen.}

In der mathematischen Struktur physikalischer Theorien treten Lie-Gruppen in zwei
fundamental verschiedenen, aber strukturell verknüpften Rollen auf:

\begin{center}
\begin{tabular}{lll}

\textbf{Ebene} & \textbf{Rolle der Gruppe} & \textbf{Beispielhafte Gruppen} \\

\emph{Symmetrie-Ebene (I)} & externe Transformationsgruppen (Raum, Zeit, interne) & 
$SO(3)$, $\mathbb R_t$, $SU(2)$, Poincaré, Galilei \\
\emph{Dynamische Ebene (II)} & interne Bewegungsgruppen (Flüsse im Phasenraum) &
Hamiltonsche Diffeomorphismen, Heisenberg-Gruppe $H_n$ \\

\end{tabular}
\end{center}

\begin{itemize}
\item Auf Ebene (I) beschreibt $G$ Symmetrien, die auf Zustandsräumen oder Feldräumen wirken.
\item Auf Ebene (II) beschreibt eine Poisson- oder Heisenberg-Struktur die interne
  Dynamik des Systems, also Flüsse auf dem Phasenraum.
\end{itemize}

Die Quantisierung fungiert als Brücke zwischen diesen Ebenen.

%------------------------------------------------------------
\textbf{(2) Symmetrie-Ebene (I): Darstellungen von Lie-Gruppen.}

Eine glatte Darstellung
\[
U:G\longrightarrow \mathcal U(\mathcal H),
\qquad
dU:\mathfrak g\longrightarrow \mathfrak u(\mathcal H),
\]
erfüllt das \emph{Darstellungsdiagramm}
\[
\begin{tikzcd}[column sep=large,row sep=large]
\mathfrak g \arrow{r}{\exp} \arrow{d}{dU}
& G \arrow{d}{U}
\\
\mathfrak u(\mathcal H) \arrow{r}{\exp}
& \mathcal U(\mathcal H)
\end{tikzcd}
\qquad
\text{(kommutativ)}.
\]
%
Hier:
\begin{itemize}
\item $G$ – Lie-Gruppe der physikalischen Symmetrien
  (Rotation, Translation, Zeitentwicklung, Eichgruppe),
\item $\mathfrak g=T_eG$ – Lie-Algebra der Generatoren
  (z.\,B. $H$, $P_i$, $J_i$),
\item $\mathcal U(\mathcal H)$ – unitäre Gruppe auf dem Hilbertraum der Zustände,
\item $\mathfrak u(\mathcal H)$ – Lie-Algebra der antihermiteschen Generatoren
  ($i\hbar\,\mathfrak u(\mathcal H)$ entspricht den selbstadjungierten Observablen).
\end{itemize}
%
Das Diagramm beschreibt die Exponential-Korrespondenz
zwischen infinitesimalen Generatoren und endlichen Transformationen
sowie deren Umsetzung als unitäre Operationen auf Zuständen:
\[
U(\exp(tX))=\exp(t\,dU(X)),\qquad X\in\mathfrak g.
\]

%------------------------------------------------------------
\textbf{(3) Dynamische Ebene (II): Poisson- und Flussstruktur.}

Der klassische Phasenraum $(\mathcal P,\omega)$ (typisch $\mathcal P=T^*Q$)
trägt die Poisson-Struktur
\[
\{f,g\}
=\omega(X_f,X_g)
=\sum_i\left(
\frac{\partial f}{\partial q^i}\frac{\partial g}{\partial p_i}
-
\frac{\partial f}{\partial p_i}\frac{\partial g}{\partial q^i}
\right).
\]
Damit wird $(C^\infty(\mathcal P),\{\ ,\ \})$ zur Lie-Algebra der Observablen.
Jede Observable $f$ erzeugt ein Hamiltonsches Vektorfeld
$X_f=\{f,\cdot\}$ und einen Fluss $\Phi^f_t=\exp(tX_f)$:
\[
[X_f,X_g]=X_{\{f,g\}}.
\]
%
Analog zum obigen Diagramm gilt formal
\[
\begin{tikzcd}[column sep=large,row sep=large]
C^\infty(\mathcal P) \arrow{r}{f\mapsto X_f} \arrow{d}{f\mapsto\hat f}
& \mathfrak X_\text{Ham}(\mathcal P) \arrow{r}{\exp}
& \mathrm{Ham}(\mathcal P)\\
\mathfrak u(\mathcal H) \arrow{r}{[\ ,\ ]/(i\hbar)}
& \mathfrak u(\mathcal H) \arrow{r}{\exp}
& \mathcal U(\mathcal H)
\end{tikzcd}
\]
wobei die vertikale Abbildung die Quantisierung ist.

%------------------------------------------------------------
\textbf{(4) Brücke zwischen beiden Ebenen: die Quantisierung.}

Quantisierung ist ein (formal) Lie-funktorieller Übergang:
\[
\mathcal Q:(C^\infty(\mathcal P),\{\ ,\ \})
\longrightarrow
(\mathfrak u(\mathcal H),\tfrac{1}{i\hbar}[\ ,\ ]).
\]
Die Forderung
\[
\frac{1}{i\hbar}[\hat f,\hat g]\approx\widehat{\{f,g\}}
\]
stellt die strukturelle Kompatibilität sicher.

Damit entspricht:
\begin{center}
\begin{tabular}{lll}

\textbf{Klassische Seite} & \textbf{Quantisierte Seite} & \textbf{Interpretation}\\

Observable $f$ & Operator $\hat f$ & Repräsentant einer Größe\\
Hamiltonscher Fluss $\Phi^f_t$ & Unitärer Fluss $e^{-\,\frac{i}{\hbar}t\hat f}$ & Zeitentwicklung / Symmetriefluss\\
Poisson-Klammer $\{f,g\}$ & Kommutator $[\hat f,\hat g]/(i\hbar)$ & Lie-Struktur-Erhalt\\

\end{tabular}
\end{center}

Die Heisenberg-Gruppe $H_n$ (bzw. ihre kontinuierliche Version in QFT)
ist die \emph{zentrale Erweiterung} der abelschen Gruppe $\mathbb R^{2n}$,
die sicherstellt, dass die kanonischen Relationen auch für endliche Transformationen
(Weyl-Form) exponentiell konsistent bleiben.

%------------------------------------------------------------
\textbf{(5) Zusammenführung der Diagramme.}

Die Symmetrie-Ebene und die Dynamik-Ebene verschmelzen in der Quantentheorie:
\[
\begin{tikzcd}[column sep=huge,row sep=large]
\text{(klassische Flüsse)} &
C^\infty(\mathcal P)\arrow{r}{\{\ ,\ \}} \arrow{d}{\text{Quantisierung}}
& \mathfrak X_\text{Ham}(\mathcal P)\arrow{r}{\exp}\arrow{d}{dU}
& \mathrm{Ham}(\mathcal P)\arrow{d}{U}
\\
\text{(unitäre Flüsse)} &
\mathfrak u(\mathcal H)\arrow{r}{\frac{1}{i\hbar}[\ ,\ ]}\arrow{d}{\exp}
& \mathfrak u(\mathcal H)\arrow{r}{\exp}\arrow{d}{\exp}
& \mathcal U(\mathcal H)\arrow[equal]{d}
\\
&
\mathcal U(\mathcal H)\arrow[equal]{r}
& \mathcal U(\mathcal H)
& \mathcal U(\mathcal H)
\end{tikzcd}
\]
Das oberste Diagramm beschreibt die klassische Dynamik,
das untere die quantisierte Dynamik.
Vertikal verbinden sie die Quantisierung und die Darstellung $dU,U$.


%------------------------------------------------------------
\textbf{(6) Schlussformel.}

\[
\boxed{
\begin{aligned}
&\text{Klassik: }
  (T^*Q,\omega),\ 
  (C^\infty(T^*Q),\{\ ,\ \}) \quad
  \longleftrightarrow\quad
  \text{Flüsse in }\mathrm{Ham}(T^*Q);\\[3pt]
&\text{Quantisierung: }
  f\mapsto\hat f,\ 
  \{f,g\}\mapsto\tfrac{1}{i\hbar}[\hat f,\hat g];\\[3pt]
&\text{Quantentheorie: }
  (\mathcal H,\mathcal U(\mathcal H),\mathfrak u(\mathcal H)),\ 
  U(\exp tX)=e^{t\,dU(X)};\\[3pt]
&\text{Beide Ebenen sind durch dieselbe Lie-Funktorialität verbunden.}
\end{aligned}
}
\]



\pagebreak


%------------------------------------------------------------
[Operatorwertige Distributionen und Quantenfelder]

[Operatorwertige Distribution]

Sei $\mathcal{H}$ ein Hilbertraum mit einem dichten, invarianten Unterraum 
$\mathcal{D}\subset\mathcal{H}$.
Eine \emph{operatorwertige Distribution} auf einer glatten Mannigfaltigkeit 
$M$ (typischerweise $M=\mathbb{R}^{1,3}$) ist eine lineare Abbildung
\[
  \Phi : \mathcal{S}(M) \longrightarrow \mathrm{Op}(\mathcal{D})
\]
von der Raumzeit-Testfunktionen (z.\,B.\ $\mathcal{S}(M)$, dem Raum der 
Schwarz-Funktionen) in die dicht definierten Operatoren auf $\mathcal{H}$ 
mit gemeinsamem Definitionsbereich $\mathcal{D}$, 
so dass für alle $\Psi_1,\Psi_2\in\mathcal{D}$ die Abbildung
\[
  f \ \longmapsto\ \langle \Psi_1,\Phi(f)\Psi_2\rangle
\]
eine (komplexwertige) Distribution auf $M$ ist.

Wir schreiben formal
\[
  \Phi(f) = \int_M \Phi(x)\,f(x)\,d^4x,
\]
wobei $\Phi(x)$ als \emph{formalwertiges Kernsymbol} der Distribution
verstanden wird.


[Quantisiertes (freies) Skalarfeld]

Eine operatorwertige Distribution 
\[
  \hat{\phi} : \mathcal{S}(M) \to \mathrm{Op}(\mathcal{D})
\]
heißt ein \emph{freies skalarsches Quantenfeld} der Masse $m>0$, 
wenn folgende Bedingungen erfüllt sind:
\begin{enumerate}
\item \textbf{Linearität und Selbstadjungiertheit:} 
      $\hat{\phi}$ ist reell-linear, 
      und es gilt $\hat{\phi}(f)^\dagger = \hat{\phi}(\overline{f})$ 
      auf $\mathcal{D}$.

\item \textbf{Klein--Gordon-Gleichung (distributionell):}
      \[
        \hat{\phi}\big((\Box + m^2)f\big) = 0
        \qquad \forall f\in\mathcal{S}(M).
      \]

\item \textbf{Kanonische Vertauschungsrelationen (CCR):}
      Für alle $f,g\in\mathcal{S}(M)$ gilt
      \[
        [\,\hat{\phi}(f),\hat{\phi}(g)\,]
        \;=\;
        i\,\sigma(f,g)\,\mathbf{1},
        \qquad
        \sigma(f,g)
        =\int_M f(x)\,(\Delta g)(x)\,d^4x,
      \]
      wobei $\Delta = E_{\mathrm R}-E_{\mathrm A}$ der kausale Propagator 
      der Klein--Gordon-Gleichung ist.

\item \textbf{Poincaré-Kovarianz:}
      Es existiert eine unitäre Darstellung 
      $U(a,\Lambda)$ der Poincaré-Gruppe auf $\mathcal{H}$, so dass
      \[
        U(a,\Lambda)\,\hat{\phi}(f)\,U(a,\Lambda)^{-1}
        = \hat{\phi}(f_{(a,\Lambda)}),
        \qquad
        f_{(a,\Lambda)}(x)
        = f(\Lambda^{-1}(x-a)).
      \]
\end{enumerate}


[Distributionelle Operationen]

Für eine operatorwertige Distribution $\Phi$ sind die folgenden
Operationen \emph{distributionell} definiert:
\begin{itemize}
\item \textbf{Ableitung:} 
      $\displaystyle \partial_\mu\Phi(f) := -\,\Phi(\partial_\mu f).$

\item \textbf{Kommutator-Kern:}
      $\displaystyle [\Phi(x),\Phi(y)] = i\,\Delta(x-y)\,\mathbf{1}$
      bedeutet, dass
      \[
        [\Phi(f),\Phi(g)] 
        = i\!\int\!\!\int f(x)\,\Delta(x-y)\,g(y)\,d^4x\,d^4y \;\mathbf{1}.
      \]

\item \textbf{Adjungierte:} 
      $\Phi(f)^\dagger=\Phi(\overline f)$,
      und für reelles $f$ ist $\Phi(f)$ wesentlich selbstadjungiert 
      auf $\mathcal{D}$ (Segal--Nelson).

\item \textbf{Zeitordnung:} 
      Das formal geschriebene 
      $T\{\Phi(x)\Phi(y)\}$ bezeichnet eine neu definierte 
      Distribution $T_2(x,y)$, deren Träger und Singularität 
      kausal fortgesetzt werden, so dass sie
      \[
        T_2(f,g)
        = \theta(x^0-y^0)\Phi(f)\Phi(g)
        + \theta(y^0-x^0)\Phi(g)\Phi(f)
      \]
      reproduziert, wo die rechte Seite sinnvoll fortgesetzt ist.
\end{itemize}


[Freies Skalarfeld im Impulsraum]

Die Fock-Darstellung der obigen Definition lautet formal
\[
  \hat{\phi}(x)
  = \int \frac{d^3p}{(2\pi)^3\,2E_{\mathbf p}}\;
    \Big(
      \hat{a}(\mathbf p)\,e^{-ip\cdot x}
      + \hat{a}^\dagger(\mathbf p)\,e^{ip\cdot x}
    \Big),
  \qquad E_{\mathbf p}=\sqrt{\mathbf p^2+m^2},
\]
wobei $\hat{a}(\mathbf p),\hat{a}^\dagger(\mathbf p)$
selbst operatorwertige Distributionen im Impulsraum sind und
\[
  [\hat{a}(\mathbf p),\hat{a}^\dagger(\mathbf q)]
  = (2\pi)^3\,2E_{\mathbf p}\,\delta^{(3)}(\mathbf p-\mathbf q).
\]
Der Fockraum ist $\mathcal{H}=\Gamma_s(L^2(\mathbb R^3,d^3p/(2\pi)^3 2E_{\mathbf p}))$,
und der dichte Definitionsbereich $\mathcal{D}$ besteht aus 
Endlich-Teilchen-Zuständen.


[Physikalische Bedeutung]

Die Feldoperatoren $\hat{\phi}(x)$ sind formal „lokalisiert“ an einem 
Raumzeitpunkt und erzeugen bzw.\ vernichten Feldquanten (Teilchen). 
Da die zugrunde liegenden Vertauschungsrelationen 
$\delta$-Kerne enthalten, können sie nur 
als \emph{operatorwertige Distributionen} verstanden werden.
Erst die gesmearte Größe 
\[
  \hat{\phi}(f)=\int\hat{\phi}(x)f(x)\,d^4x
\]
ist ein wohldefinierter Operator, der die mittlere Feldstärke 
in der Raumzeitregion $\mathrm{supp}(f)$ beschreibt. 
Diese Form ist die natürlichste mathematische Realisierung der 
Gleichzeitigkeit von
\begin{enumerate}
\item Quantenmechanik (Operatoren auf einem Zustandsraum),
\item Relativistischer Lokalität (raumzeitliche Stützmenge),
\item und Kausalität (Kommutatoren mit kausalem Träger).
\end{enumerate}

%------------------------------------------------------------

%========================================================
\textbf{Herleitung und Interpretation der Modenzerlegung des freien Skalarfeldes}


In diesem Abschnitt wird die Darstellung
\[
  \hat{\phi}(x)
  \;=\;
  \int \frac{d^3p}{(2\pi)^3\,2E_{\mathbf p}}\;
    \Big(
      \hat{a}(\mathbf p)\,e^{-ip\cdot x}
      + \hat{a}^\dagger(\mathbf p)\,e^{ip\cdot x}
    \Big),
  \qquad E_{\mathbf p}=\sqrt{\mathbf p^2+m^2},
\]
aus der strengen (smeared) Quantisierung hergeleitet und präzisiert, was es mathematisch und konzeptionell bedeutet, mit dieser Form zu rechnen. Zentral ist, dass \(\hat\phi(x)\) \emph{keine} punktweise definierte Operatorfunktion ist, sondern eine \emph{operatorwertige Distribution}; die obige Gleichung ist als Aussage über Distributionen zu lesen.

%--------------------------------------------------------
\textbf{Ausgangspunkt: Smeared Feldoperatoren und CCR}

Sei \(M=\mathbb{R}^{1,3}\) die Minkowski-Raumzeit, \(P:=\Box + m^2\) der Klein--Gordon-Operator und \(\Delta := E_{\mathrm R}-E_{\mathrm A}\) der kausale Propagator. Definiere den quotientierten Testfunktionsraum
\[
  \mathcal V \;:=\; \mathcal{S}(M)\big/ P\mathcal{S}(M),
\]
und die symplektische Form
\[
  \sigma([f],[g]) \;:=\; \int_{M} f(x)\,(\Delta g)(x)\,d^4x
  \;=\; \int_\Sigma \big(\phi_f\,\partial_n\phi_g - \phi_g\,\partial_n\phi_f\big)\,d\Sigma,
\]
wobei \(\phi_f := E f\) und \(\Sigma\) eine Cauchyfläche ist (Cauchy-Invarianz). Das \emph{freie skalare Quantenfeld} ist eine linear-distributionelle Abbildung
\[
  \hat\phi:\ \mathcal{S}(M)\to \mathrm{Op}(\mathcal D),\qquad
  f\mapsto \hat\phi(f),
\]
definiert auf einem dichten, invarianten Kern \(\mathcal D\subset\mathcal H\) (typisch: Endlich-Teilchen-Vektoren im Fockraum), die die \emph{CCR}
\[
  [\,\hat\phi(f),\hat\phi(g)\,] \;=\; i\,\sigma(f,g)\,\mathbf 1
\]
und die Feldgleichung \(\hat\phi(P f)=0\) (distributionell) erfüllt. Formal schreibt man
\[
  \hat\phi(f) \;=\; \int \hat\phi(x)\,f(x)\,d^4x,
\]
wobei \(\hat\phi(x)\) das \emph{formalwertige Kernsymbol} einer operatorwertigen Distribution ist.

%--------------------------------------------------------
\textbf{Einteilchenraum und Positiv-Frequenz-Projektion}

Der Einteilchen-Hilbertraum ist
\[
  \mathfrak h \;=\; L^2\!\Big(\mathbb{R}^3,\ \frac{d^3p}{(2\pi)^3\,2E_{\mathbf p}}\Big),
  \qquad E_{\mathbf p}=\sqrt{\mathbf p^2+m^2},
\]
und der bosonische Fockraum \(\mathcal H=\Gamma_s(\mathfrak h)=\bigoplus_{n=0}^\infty \mathrm{Sym}^n\mathfrak h\). Die Abbildung
\[
  K:\ \mathcal{S}(M)\to \mathfrak h,\qquad
  (Kf)(\mathbf p) \;:=\; \widetilde f(E_{\mathbf p},\mathbf p),
\]
liest die \emph{positive Frequenzkomponente} (Massenschale \(p^2=m^2,\ p^0>0\)) der Fouriertransformierten \(\widetilde f\) aus. Damit definiert man den smeared Feldoperator als \emph{Segal-Feld}
\[
  \hat\phi(f) \;=\; \frac{1}{\sqrt{2}}\Big( a^\dagger(Kf)+a(Kf) \Big)
  \quad\text{auf}\ \mathcal D,
\]
wobei \(a^\dagger(\varphi),a(\varphi)\) die Erzeuger/Vernichter auf \(\Gamma_s(\mathfrak h)\) sind. Auf \(\mathcal D\) gilt die CCR
\(
  [\hat\phi(f),\hat\phi(g)]=i\,\sigma(f,g)\,\mathbf 1.
\)

%--------------------------------------------------------
\textbf{Herleitung der Ortsdarstellung als Distribution}

Schreibe die Fourier-Inversion heuristisch als Paarung mit der \(\delta\)-Distribution:
\[
  \delta_{x_0}(x) = \delta(x-x_0),\qquad
  \widetilde{\delta_{x_0}}(p)=e^{ip\cdot x_0}.
\]
Setzt man \(f=\delta_{x_0}\) in die Segal-Formel, so ergibt sich formal
\[
  \hat\phi(x_0)
  \;=\; \hat\phi(\delta_{x_0})
  \;=\; \frac{1}{\sqrt{2}}\!\int\!\frac{d^3p}{(2\pi)^3\,2E_{\mathbf p}}
  \Big( a^\dagger(\mathbf p)\,e^{ip\cdot x_0} + a(\mathbf p)\,e^{-ip\cdot x_0}\Big).
\]
Damit folgt die \emph{Modenzerlegung als operatorwertige Distribution}
\begin{equation}\label{eq:mode-expansion}
  \hat\phi(x)
  \;=\; \int \frac{d^3p}{(2\pi)^3\,2E_{\mathbf p}}\;
  \big(\hat a(\mathbf p)e^{-ip\cdot x}+\hat a^\dagger(\mathbf p)e^{ip\cdot x}\big),
\end{equation}
wobei \(\hat a(\mathbf p),\hat a^\dagger(\mathbf p)\) selbst operatorwertige Distributionen in \(\mathbf p\) sind, die die (distributionellen) CCR
\[
  [\hat a(\mathbf p),\hat a^\dagger(\mathbf q)]
  \;=\; (2\pi)^3\,2E_{\mathbf p}\,\delta^{(3)}(\mathbf p-\mathbf q),
  \qquad
  [\hat a,\hat a]=[\hat a^\dagger,\hat a^\dagger]=0
\]
erfüllen. Die Gleichung \eqref{eq:mode-expansion} ist \emph{äquivalent} zur strengen smeared-Formel in dem Sinne, dass für alle \(f\in\mathcal S(M)\)
\[
  \hat\phi(f)
  \;=\; \int d^4x\, f(x)\,\hat\phi(x)
  \;=\; \int\!\frac{d^3p}{(2\pi)^3\,2E_{\mathbf p}}
  \Big( \hat a(\mathbf p)\,\widetilde f(E_{\mathbf p},\mathbf p)
       + \hat a^\dagger(\mathbf p)\,\overline{\widetilde f(E_{\mathbf p},\mathbf p)}\Big).
\]


%========================================================
\textbf{Einteilchenraum, Positiv-Frequenz-Projektion und Herleitung der Ortsdarstellung}
\label{sec:one-particle-K}

In diesem Abschnitt werden (i) der \emph{Einteilchenraum} $\mathfrak h$, (ii) die
\emph{Positiv-Frequenz-Projektion} und (iii) die \emph{Herleitung der
Ortsdarstellung} $\hat\phi(x)$ als operatorwertige Distribution präzise konstruiert.
Wir geben zwei äquivalente Blickwinkel: (A) die \emph{Fourier-/Massenschalen-Konstruktion}
und (B) die \emph{Cauchy-Daten-/Zeitentwicklungs-Konstruktion}. In beiden Fällen
entsteht ein kanonischer linearer Operator
\[
  K:\ \mathcal S(M)/P\mathcal S(M)\ \longrightarrow\ \mathfrak h
\]
(vgl. \emph{Time-slice}-Theorem), der die Positivfrequenz-Anteile der Testfunktionen
auf der Massenschale extrahiert. Aus $K$ folgt die Segal-Feldformel
$\hat\phi(f)=\tfrac1{\sqrt2}\big(a^\dagger(Kf)+a(Kf)\big)$ und daraus die
Modenzerlegung.

%--------------------------------------------------------
\textbf{(A) Fourier-/Massenschalen-Konstruktion des Einteilchenraums}
\label{subsec:fourier-construction}

Sei $M=\mathbb R^{1,3}$ (Signatur $(+,-,-,-)$), $P=\Box+m^2$.
Für $f\in\mathcal S(M)$ bezeichne $\widetilde f$ die Fouriertransformierte
\[
  \widetilde f(p)\ :=\ \int_{\mathbb R^4} e^{ip\cdot x} f(x)\,d^4x,
  \qquad p\cdot x := p_\mu x^\mu.
\]
Die \emph{positiv zeitartige Massenschale} ist
\[
  \mathcal C_m^+\ :=\ \{\,p\in\mathbb R^4\mid p^2=m^2,\ p^0>0\,\}
  \ \cong\ \mathbb R^3,\quad p=(E_{\mathbf p},\mathbf p),\ E_{\mathbf p}=\sqrt{\mathbf p^2+m^2}.
\]
Wir verwenden das Lorentz-invariante Maß
\[
  d\mu_m(p)\ :=\ \frac{d^3\mathbf p}{(2\pi)^3\,2E_{\mathbf p}}\quad\text{auf }\mathcal C_m^+.
\]
\emph{Einteilchenraum.}
Definiere den Einteilchen-Hilbertraum
\[
  \mathfrak h\ :=\ L^2\!\big(\mathcal C_m^+,d\mu_m\big)
  \ \cong\ L^2\!\Big(\mathbb R^3,\frac{d^3\mathbf p}{(2\pi)^3\,2E_{\mathbf p}}\Big).
\]
\emph{Positiv-Frequenz-Projektion und $K$.}
Definiere
\begin{equation}
  \label{eq:def-K-Fourier}
  K:\ \mathcal S(M)\ \longrightarrow\ \mathfrak h,\qquad
  (Kf)(\mathbf p)\ :=\ \widetilde f\!\big(E_{\mathbf p},\mathbf p\big).
\end{equation}
Liegt $f=Ph$ in $P\mathcal S(M)$, so verschwindet $\widetilde f$ auf der Massenschale
(da $\widetilde f(p)=(p^2-m^2)\,\widetilde h(p)$), also $K(P h)=0$. Damit faktorisiert $K$
kanonisch zum Quotienten
\[
  K:\ \mathcal S(M)/P\mathcal S(M)\ \longrightarrow\ \mathfrak h.
\]
\emph{Symplektische Form und Einteilchenskalarsprodukt.}
Mit $\sigma([f],[g])=\int f(x)\,\Delta g(x)\,d^4x$ gilt die Standardrelation
\[
  2\,\mathrm{Im}\,\langle Kf,Kg\rangle_{\mathfrak h}\ =\ \sigma([f],[g]),
\]
d.\,h.\ $K$ ist \emph{kompatibel} mit der symplektischen Struktur; die positiv-definite Form
$\mu([f],[g]):=\mathrm{Re}\,\langle Kf,Kg\rangle_{\mathfrak h}$ induziert die
quasifreie Vakuumstruktur.

%--------------------------------------------------------
\textbf{(B) Cauchy-Daten-/Zeitentwicklungs-Konstruktion}
\label{subsec:cauchy-construction}

Wähle eine Cauchy-Hyperfläche $\Sigma=\{t=0\}\cong\mathbb R^3$. Jede glatte Lösung der
KG-Gleichung wird eindeutig durch Cauchy-Daten $(\varphi,\pi)\in\mathcal S(\Sigma)\oplus\mathcal S(\Sigma)$
mit
\[
  \varphi(\mathbf x)=\phi(0,\mathbf x),\qquad \pi(\mathbf x)=\partial_t\phi(0,\mathbf x)
\]
bestimmt. Mit Fouriertransformationen $\widehat\varphi(\mathbf p)$, $\widehat\pi(\mathbf p)$ definiere
\[
  \mathcal F:\ \mathcal S(\Sigma)\oplus\mathcal S(\Sigma)\ \to\ \mathfrak h,\qquad
  \big(\varphi,\pi\big)\ \mapsto\ \frac{1}{\sqrt{2}}
  \Big(E_{\mathbf p}^{1/2}\,\widehat\varphi(\mathbf p)\ +\ i\,E_{\mathbf p}^{-1/2}\,\widehat\pi(\mathbf p)\Big).
\]
Dann ist $\mathcal F$ eine isometrische Einbettung bzgl. der Klein-Gordon-Energieform; die Zeitentwicklung
$U(t)=e^{-itH}$ wirkt in $\mathfrak h$ einfach als Multiplikator $e^{-itE_{\mathbf p}}$. Der Operator $K$
ergibt sich als Komposition: \emph{Erzeuge aus $f$ zunächst die Lösung $\phi_f=E f$, lese dann die Cauchy-Daten
bei $t=0$ ab und wende $\mathcal F$ an}. In Formeln:
\[
  K[f]\ =\ \mathcal F\big(\phi_f|_{t=0},\ \partial_t\phi_f|_{t=0}\big).
\]
Man prüft, dass dies mit \eqref{eq:def-K-Fourier} übereinstimmt.

%--------------------------------------------------------
\textbf{Segal-Feld und Modenoperatoren}
\label{subsec:segal-and-modes}

Mit $K$ ist das freie Feld als Segal-Feld definiert:
\begin{equation}
  \label{eq:Segal-field}
  \hat\phi(f)\ :=\ \frac{1}{\sqrt2}\Big(a^\dagger(Kf)+a(Kf)\Big),\qquad f\in\mathcal S(M),
\end{equation}
auf dem dichten Fock-Kern $\mathcal D$ (endlich-Teilchen-Vektoren), wobei
$a^\dagger(\psi),a(\psi)$ die bosonischen Erzeuger/Vernichter für $\psi\in\mathfrak h$ sind.
Für Testfunktionen in Impulsdarstellung, $f\mapsto Kf$, ist es natürlich, die
\emph{distributionswertigen Modenoperatoren} $\,\hat a(\mathbf p),\hat a^\dagger(\mathbf p)\,$ so zu verstehen, dass
\[
  a^\dagger(\psi)=\int \overline{\psi(\mathbf p)}\,\hat a^\dagger(\mathbf p)\,d\mu_m(\mathbf p),
  \qquad
  a(\psi)=\int \psi(\mathbf p)\,\hat a(\mathbf p)\,d\mu_m(\mathbf p),
\]
mit den (distributionellen) CCR
\[
  [\hat a(\mathbf p),\hat a^\dagger(\mathbf q)]\ =\ (2\pi)^3\,2E_{\mathbf p}\,\delta^{(3)}(\mathbf p-\mathbf q),
  \quad [\hat a,\hat a]=[\hat a^\dagger,\hat a^\dagger]=0.
\]

%--------------------------------------------------------
\textbf{Herleitung der Ortsdarstellung als operatorwertige Distribution}
\label{subsec:derivation-x}

Die formale Feldverteilung $\hat\phi(x)$ ist das \emph{Kernsymbol} von \eqref{eq:Segal-field}, d.\,h.
\[
  \hat\phi(f)\ =\ \int_{\mathbb R^4} \hat\phi(x)\,f(x)\,d^4x.
\]
Setzt man \eqref{eq:Segal-field} und die Definition \eqref{eq:def-K-Fourier} ein, erhält man
\[
  \hat\phi(f)
  = \frac{1}{\sqrt2}\!\int\!d\mu_m(\mathbf p)\,
    \Big(\hat a^\dagger(\mathbf p)\,\widetilde f(E_{\mathbf p},\mathbf p)
        + \hat a(\mathbf p)\,\overline{\widetilde f(E_{\mathbf p},\mathbf p)}\Big).
\]
Fourier-Inversion liefert (heuristisch mit $f=\delta_{x_0}$; streng über Approximationen der Eins)
\[
  \widetilde{\delta_{x_0}}(p)=e^{ip\cdot x_0}
  \quad\Rightarrow\quad
  \boxed{\;
  \hat\phi(x_0)\ =\ \int \frac{d^3\mathbf p}{(2\pi)^3\,2E_{\mathbf p}}
  \Big(\hat a(\mathbf p)\,e^{-ip\cdot x_0}
      + \hat a^\dagger(\mathbf p)\,e^{+ip\cdot x_0}\Big)\; }.
\]
Dies ist die übliche Modenzerlegung. Ihre \emph{Bedeutung} ist distributionell:
für jedes $f\in\mathcal S(M)$ gilt
\[
  \int \hat\phi(x)\,f(x)\,d^4x
  \;=\; \int d\mu_m(\mathbf p)\,
      \Big(\hat a(\mathbf p)\,\widetilde f(E_{\mathbf p},\mathbf p)
          + \hat a^\dagger(\mathbf p)\,\overline{\widetilde f(E_{\mathbf p},\mathbf p)}\Big),
\]
also ist die Gleichheit \emph{äquivalent} zur strengen Segal-Definition.

%--------------------------------------------------------
\textbf{Rolle der Variablen \texorpdfstring{$x$}{x} und \texorpdfstring{$p$}{p}}
\label{subsec:roles-x-p}

\emph{(i) $x$ als Lokalisierungs-Parameter.}
Das Symbol $x\in M$ ist \emph{kein} Punkt, an dem ein beschränkter Operator ausgewertet würde; vielmehr
kennzeichnet $x$ die \emph{Evaluation-Distribution} $\delta_x$, gegen die geschmiert wird. Das Feld
$\hat\phi(x)$ existiert nur als \emph{operatorwertige Distribution}; physikalisch entspricht
$\hat\phi(f)$ der Messung der „gemittelten Feldstärke“ in der Region $\mathrm{supp}(f)$.

\emph{(ii) $p$ als Spektralvariable der Translationen.}
$p=(E_{\mathbf p},\mathbf p)\in\mathcal C_m^+$ parametrisiert die irreduziblen Einteilchenmoden; in der Fockdarstellung
ist die unitäre Translationsdarstellung durch $P^\mu$ gegeben, und $p^\mu$ ist die Spektralvariable (Eigenwert) des
Einteilchen-Impulsoperators. Die Modenoperatoren $\hat a^\dagger(\mathbf p)$ erzeugen Einteilchenzustände
mit Vierimpuls $p^\mu$, und die Exponenten $e^{\pm ip\cdot x}$ implementieren die kovariante Raumzeit-Transformation.

\emph{(iii) Normierung und Maß.}
Die Wahl des Maßes $d\mu_m(\mathbf p)=(2\pi)^{-3}(2E_{\mathbf p})^{-1} d^3\mathbf p$ garantiert sowohl Lorentzinvarianz
der Zweipunktfunktion als auch die kanonischen (distributionellen) CCR:
$[\hat a(\mathbf p),\hat a^\dagger(\mathbf q)]=(2\pi)^3\,2E_{\mathbf p}\,\delta^{(3)}(\mathbf p-\mathbf q)$.

%--------------------------------------------------------
\textbf{Warum die Modenzerlegung legitim ist (Distributionensinn)}
\label{subsec:legitimacy}

\begin{itemize}
  \item \textbf{Distributioneller Sinn:} Die Gleichung
  $\hat\phi(x)=\int d\mu_m(\mathbf p)\, \big(\hat a(\mathbf p)e^{-ipx}+\hat a^\dagger(\mathbf p)e^{ipx}\big)$
  bedeutet, dass \emph{alle} physikalisch relevanten Größen (Erwartungswerte, $n$-Punkt-Funktionen) nach Smearing
  mit Testfunktionen identisch mit der Segal-Formel sind. Punktweise Produkte sind nicht definiert; normalgeordnete
  bzw. renormierte Produkte werden als Distributionen konstruiert.
  \item \textbf{Spektraldiagonalisierung:} In dieser Darstellung sind $H,P^i$ simultan diagonalisierbar (als
  $d\Gamma(E_{\mathbf p})$, $d\Gamma(p^i)$), was Streurechnungen und die Teilcheninterpretation transparent macht.
  \item \textbf{Implizites Smearing in der Praxis:} In Lagrange-Dichten, Feynman-Regeln und LSZ treten stets Integrale
  $\int d^4x\,(\cdots\hat\phi(x)\cdots)$ auf; das ist \emph{implizites} Smearing. Daher ist die Modenform die
  natürliche Rechenform, obwohl sie distributionell zu verstehen ist.
\end{itemize}

%--------------------------------------------------------
\textbf{Mini-Checkliste (Eigenschaften von $K$)}
\label{subsec:K-properties}

\begin{enumerate}
  \item $K$ ist linear und faktorisiert zum Quotienten $\mathcal S(M)/P\mathcal S(M)$.
  \item $2\,\mathrm{Im}\,\langle Kf,Kg\rangle_{\mathfrak h}=\sigma([f],[g])$ (Kompatibilität mit CCR).
  \item $K$ ist Poincaré-kovariant: $(Kf)^\Lambda(\mathbf p)=K(f\circ\Lambda^{-1})(\mathbf p)$
  unter geeigneter Wirkung (Details: Massenschalenabbildung).
  \item $K$ wählt die Positiv-Frequenz-Komponente (Massenschalenrestriktion $p^0=+E_{\mathbf p}$).
  \item In Cauchy-Daten-Form: $K[f]=\frac1{\sqrt2}(E_{\mathbf p}^{1/2}\widehat\varphi+iE_{\mathbf p}^{-1/2}\widehat\pi)$
  mit $(\varphi,\pi)$ den durch $f$ induzierten Cauchy-Daten.
\end{enumerate}
%========================================================


%--------------------------------------------------------
\textbf{Distributionelle Bedeutung und Legitimität der Rechnung}

\emph{(i) Operatorwertige Distribution.}
Die Gleichung \eqref{eq:mode-expansion} ist \emph{nicht} als Identität in \(\mathcal B(\mathcal H)\) zu verstehen, sondern als Identität \emph{von Distributionen mit Werten in unbeschränkten Operatoren} auf einem gemeinsamen dichten Bereich \(\mathcal D\). Aussagen wie
\[
  [\,\hat\phi(x),\hat\phi(y)\,] \;=\; i\,\Delta(x-y)\,\mathbf 1
\]
bedeuten streng:
\[
  [\,\hat\phi(f),\hat\phi(g)\,]
  \;=\; i\!\int\!\!\int f(x)\,\Delta(x-y)\,g(y)\,d^4x\,d^4y\;\mathbf 1
  \qquad\forall f,g\in\mathcal S(M).
\]

\emph{(ii) Ableitungen und Feldgleichung.}
Ableitungen sind distributionell definiert:
\[
  \partial_\mu \hat\phi(f) \;:=\; -\,\hat\phi(\partial_\mu f),\qquad
  \hat\phi(P f)=0\ \ \text{(Klein--Gordon-Gleichung distributionell)}.
\]

\emph{(iii) Erwartungswerte \& Wightman-Distributionen.}
Für einen Zustand \(\omega(\cdot)=\langle\Omega,\cdot\,\Omega\rangle\) sind
\[
  W_n(x_1,\dots,x_n) \;:=\; \langle \Omega,\hat\phi(x_1)\cdots \hat\phi(x_n)\Omega\rangle
\]
\emph{tempered Distributionen}, und alle physikalischen Größen (Propagatoren, Spektralmaße) werden als Distributionen verarbeitet (Fouriertransformation, Zeitordnung als kausale Fortsetzung, \emph{etc.}).

%--------------------------------------------------------
[Warum man praktisch mit \texorpdfstring{$\hat\phi(x)$}{phi(x)} „ohne Smearing“ rechnet]

\emph{(a) Implizites Smearing in physikalischen Rechnungen.}
In jeder realen Rechnung tritt \(\hat\phi(x)\) unter Integralen auf (z.\,B.\ Insertions \(\int d^4x\,\mathcal L_{\rm int}(\hat\phi(x))\), Fouriertransformationen, Feynman-Regeln). Das ist \emph{äquivalent} zum Smearing gegen Testfunktionen. Daher „spart“ man das Smearing in der Notation: es ist \emph{implizit} vorhanden.

\emph{(b) Energie-Impuls-Diagonalisierung.}
Die Darstellung \eqref{eq:mode-expansion} ist die Basis, in der Energie und Impuls diagonal sind:
\[
  \hat H \;=\; \int \frac{d^3p}{(2\pi)^3\,2E_{\mathbf p}}\,E_{\mathbf p}\,\hat a^\dagger(\mathbf p)\hat a(\mathbf p),\qquad
  \hat{\mathbf P} \;=\; \int \frac{d^3p}{(2\pi)^3\,2E_{\mathbf p}}\,\mathbf p\,\hat a^\dagger(\mathbf p)\hat a(\mathbf p),
\]
und \(\hat a^\dagger(\mathbf p)\) erzeugt ein Quant mit Vierimpuls \(p=(E_{\mathbf p},\mathbf p)\). Das macht die Streurechnung und die Verbindung zum Teilchenbild transparent.

\emph{(c) Distributionelle Korrektheit.}
\(\hat a(\mathbf p),\hat a^\dagger(\mathbf p)\) sind selbst operatorwertige Distributionen in \(\mathbf p\); Wellenpakete \(a^\dagger(\varphi)=\int \overline{\varphi(\mathbf p)}\,\hat a^\dagger(\mathbf p)\,d\mu(\mathbf p)\) (\(\varphi\in\mathfrak h\)) sind echte dicht definierte Operatoren. Erwartungswerte und Korrelatoren sind tempered; damit sind Ableitungen, Fouriertransformationen und kausale Fortsetzungen sauber als Distributionenoperationen definierbar.

\emph{(d) Produkte und Renormierung.}
Punktweise Produkte \(\hat\phi(x)^2\) sind \emph{nicht} definiert. In der Praxis verwendet man normalgeordnete (oder renormierte) Produkte, z.\,B.\ durch Point-Splitting:
\[
  :\!\hat\phi^2\!:(f)
  \;:=\; \lim_{y\to x}\int\!\!\int f(x)\,\Big(\hat\phi(x)\hat\phi(y)-\omega_2(x,y)\,\mathbf 1\Big)\,\delta(x-y)\,dx\,dy,
\]
wobei \(\omega_2\) die Zweipunktfunktion (Hadamard-Singularität) ist. Zeitordnung \(T\{\cdots\}\) wird als \emph{kausale Fortsetzung} einer zunächst auf \(x_i\neq x_j\) wohldefinierten Distribution konstruiert (Epstein--Glaser).

%--------------------------------------------------------
[Konzeptuelle Zusammenfassung]

\begin{itemize}
  \item \textbf{Mathematisch:} \eqref{eq:mode-expansion} ist die Fourier-\emph{Distribution} des freien Feldes auf der Massenschale \(p^2=m^2\). Sie ist äquivalent zur strengen smeared-Definition \(\hat\phi(f)=\frac1{\sqrt2}(a^\dagger(Kf)+a(Kf))\).
  \item \textbf{Physikalisch:} Das Feld ist eine Überlagerung quantisierter harmonischer Moden; \(\hat a^\dagger(\mathbf p)\) erzeugt, \(\hat a(\mathbf p)\) vernichtet ein Quant mit Impuls \(\mathbf p\).
  \item \textbf{Praktisch:} Rechnungen verwenden \(\hat\phi(x)\) stets unter Integralen (implizites Smearing). Energie/Impuls sind diagonal; Feynman-Regeln, Propagatoren und LSZ-Reduktion werden so direkt formuliert.
  \item \textbf{Vorsicht:} Punktprodukte erfordern Renormierungsrezepte (Normalordnung, Point-Splitting, kausale Fortsetzung); alle Aussagen sind distributionell und auf einem gemeinsamen dichten Operator-Domain \(\mathcal D\) zu verstehen.
\end{itemize}

%--------------------------------------------------------
[Anhang: Domänen- und Selbstadjungiertheits-Hinweise]

\begin{itemize}
  \item Der gemeinsame Kern \(\mathcal D\) kann als die Menge der Endlich-Teilchen-Zustände gewählt werden; er ist invariant unter \(a^\dagger(\varphi),a(\varphi)\) und unter \(\hat\phi(f)\).
  \item Für reelles \(f\) ist \(\hat\phi(f)\) auf \(\mathcal D\) symmetrisch und wesentlich selbstadjungiert (Segal--Nelson); daraus folgen die üblichen Spektraleigenschaften für \(H=\hat P^0\) und \(\mathbf P\).
  \item Mikrokausalität: \([\hat\phi(f),\hat\phi(g)]=0\), wenn \(\mathrm{supp}(f)\) und \(\mathrm{supp}(g)\) raumartig getrennt sind; entsprechend ist der Kommutator-Kern \(i\Delta(x-y)\) kausal getragen.
\end{itemize}
%========================================================



\pagebreak


\emph{Ziel.}
Wir erklären „von unten nach oben“, welche \emph{Lie-Gruppen}, \emph{Lie-Algebren} und \emph{Operatorräume}
im Quantisierungsdiagramm auftreten, \emph{warum} gerade diese Objekte gewählt werden
und wie sie durch Funktorialität miteinander verbunden sind.

%------------------------------------------------------------
\emph{(A) Klassische Ebene: Phasenraum, Flüsse, Poisson-Lie-Struktur}

\emph{(A1) Phasenraum als symplektische Mannigfaltigkeit.}
Sei $Q\simeq\mathbb R^n$ der Konfigurationsraum, der klassische Phasenraum
\[
\mathcal P:=T^*Q\simeq\mathbb R^{2n},\qquad
\omega=\sum_{i=1}^n dq^i\wedge dp_i
\]
ist symplektisch. Die \emph{Poisson-Klammer} auf $C^\infty(\mathcal P)$ lautet
\[
\{f,g\}
=\sum_{i=1}^n\Big(
\partial_{q^i}f\,\partial_{p_i}g-\partial_{p_i}f\,\partial_{q^i}g
\Big),
\qquad f,g\in C^\infty(\mathcal P),
\]
insbesondere $\{q^i,p_j\}=\delta^i_j$, $\{q^i,q^j\}=0$, $\{p_i,p_j\}=0$.

\emph{(A2) Lie-Algebra der Observablen \& Hamiltonsche Flüsse.}
Jede Observable $f\in C^\infty(\mathcal P)$ erzeugt ein Hamiltonsches Vektorfeld
$X_f:=\{f,\cdot\}$ und einen Fluss $\Phi^f_t=\exp(tX_f)$ in der Diffeomorphismengruppe von $\mathcal P$.
Die Abbildung
\[
C^\infty(\mathcal P)\longrightarrow \mathfrak X(\mathcal P),\quad
f\mapsto X_f
\]
ist ein \emph{Lie-Algebra-Homomorphismus}:
\[
[X_f,X_g]=X_{\{f,g\}}.
\]
Damit ist $(C^\infty(\mathcal P),\{\cdot,\cdot\})$ die \emph{klassische Lie-Algebra der Observablen}.

\emph{(A3) Warum diese Struktur?}
Die Poisson-Klammer ist die \emph{infinitesimale Kodierung} der Komposition klassischer Flüsse
(„erst $f$, dann $g$, minus umgekehrt“).
Sie ist der exakte klassische Gegenpart zum Quanten-Kommutator:
\[
\text{„Klammer = erste nichttriviale Abweichung der Komposition“}.
\]

%------------------------------------------------------------
\emph{(B) Brückenebene: Heisenberg-Struktur als minimaler Träger der Kanonik}

\emph{(B1) Heisenberg-Lie-Algebra (endlich viele Freiheitsgrade).}
Die lineare Hülle
\[
\mathfrak h_n:=\mathrm{span}\{q^i,p_i, \mathbf 1\;\mid\;i=1,\dots,n\}
\]
mit
\[
[q^i,p_j]_{\text{klass}}:=\{q^i,p_j\}=\delta^i_j,\qquad [q^i,q^j]_{\text{klass}}=0,\quad [p_i,p_j]_{\text{klass}}=0
\]
ist die \emph{kanonische} Lie-Algebra. Ihre \emph{Gruppenexponential} ist die \emph{Heisenberg-Gruppe} $H_n$.
Sie ist der kleinste (zentrale) Erweiterungsträger, auf dem die Kanonik konsistent exponentiiert (Weyl-Form, BCH).

\emph{(B2) Warum gerade $H_n$? (Minimalität \& Eindeutigkeit)}
\begin{itemize}
\item \emph{Minimalität:} $H_n$ ist die zentrale Erweiterung von $\mathbb R^{2n}$, die benötigt wird,
damit die kanonischen Relationen bei \emph{endlichen} Transformationen (Weyl-Operatoren) konsistent sind.
\item \emph{Eindeutigkeit (Stone–von Neumann):} Die irreduzible, stark stetige, unitäre Darstellung
der Weyl-Form der CCR ist (bis auf Unitäräquivalenz) eindeutig;
das erklärt, warum die Schrödinger-Darstellung kanonisch ist.
\end{itemize}

\emph{(B3) Zeittranslationsgruppe.}
Zusätzlich wirkt $G_t=\mathbb R$ (Zeittranslationen) mit Generator $H$ (Hamiltonfunktion klassisch,
Hamiltonoperator quantisiert). Sie liefert den dynamischen Fluss.

%------------------------------------------------------------
\emph{(C) Quantenebene: Operatorräume, unitäre Gruppe, Generatoren}

\emph{(C1) Operatorräume.}
\[
\mathcal U(\mathcal H)\subset \mathcal B(\mathcal H),\qquad
\mathfrak u(\mathcal H)=T_{\mathbf 1}\mathcal U(\mathcal H)=\{A^\dagger=-A\}.
\]
Selbstadjungierte Observablen sind $i\hbar\,\mathfrak u(\mathcal H)$:
\[
\mathfrak{herm}(\mathcal H)=\{\hat A=\hat A^\dagger\} \simeq i\hbar\,\mathfrak u(\mathcal H).
\]

\emph{(C2) Darstellung der Heisenberg-Gruppe.}
Eine unitäre Darstellung $U:H_n\to\mathcal U(\mathcal H)$ mit Differential
\[
dU:\mathfrak h_n\to \mathfrak u(\mathcal H),\qquad
dU(q^i)=\frac{1}{i\hbar}\hat q^i,\;\; dU(p_i)=\frac{1}{i\hbar}\hat p_i
\]
realisiert die \emph{kanonischen Vertauschungsrelationen (CCR)}:
\[
[\hat q^i,\hat p_j]=i\hbar\,\delta^i_j\,\mathbf 1,\qquad
[\hat q^i,\hat q^j]=0,\quad [\hat p_i,\hat p_j]=0.
\]
In der Schrödinger-Darstellung:
\[
(\hat q^i\psi)(q)=q^i\psi(q),\qquad
(\hat p_i\psi)(q)=-\,i\hbar\,\partial_{q^i}\psi(q),\qquad \mathcal H=L^2(\mathbb R^n).
\]

\emph{(C3) Zeitentwicklung als Darstellung von $G_t=\mathbb R$.}
\[
U_t(t)=\exp\!\Big(-\frac{i}{\hbar}t\,\hat H\Big),\qquad
\dot\psi(t)=-\frac{i}{\hbar}\hat H\psi(t).
\]
Damit fügt sich die Dynamik als weiterer Einparameter-Fluss in dasselbe Diagramm-Schema ein.

%------------------------------------------------------------
\emph{(D) Das \emph{Diagramm der Quantisierung}: Wer wirkt worauf und warum?}

\emph{(D1) Exponential-/Differential-Kompatibilität.}
\[
\begin{tikzcd}[column sep=huge,row sep=large]
\text{Heisenberg-Algebra } \mathfrak h_n \arrow{r}{\exp} \arrow{d}{dU}
& \text{Heisenberg-Gruppe } H_n \arrow{d}{U}
\\
\mathfrak u(\mathcal H) \arrow{r}{\exp}
& \mathcal U(\mathcal H)
\end{tikzcd}
\qquad
\begin{tikzcd}[column sep=huge,row sep=large]
\text{Zeit-Algebra } \mathbb R X_t \arrow{r}{\exp} \arrow{d}{dU}
& \text{Zeitgruppe } \mathbb R \arrow{d}{U_t}
\\
\mathfrak u(\mathcal H) \arrow{r}{\exp}
& \mathcal U(\mathcal H)
\end{tikzcd}
\]
Beide Quadrate \emph{kommutieren}:
\[
U(\exp Y)=\exp(dU(Y)),\qquad Y\in\mathfrak h_n\ \text{ oder }\ X_t.
\]

\emph{(D2) Poisson $\;\leadsto\;$ Kommutator (Lie-Funktorialität).}
\[
\boxed{\quad
\frac{1}{i\hbar}\,[\hat f,\hat g]\ \approx\ \widehat{\{f,g\}}
\quad}
\]
auf einem dichten Funktions/Kernraum. Für die Generatoren:
\[
\{q^i,p_j\}=\delta^i_j \;\leadsto\; [\hat q^i,\hat p_j]=i\hbar\,\delta^i_j.
\]

\emph{(D3) Gesamtfluss: klassisch vs. quantisiert.}
\[
\begin{tikzcd}[column sep=huge,row sep=large]
\text{Klassisch: } (C^\infty(\mathcal P),\{\ ,\ \}) \arrow{r}{f\mapsto X_f}
\arrow{d}{\text{Quantisierung}}
& \mathfrak X(\mathcal P) \arrow{r}{\exp} & \mathrm{Ham}(\mathcal P)
\\
\text{Quanten: } (\mathfrak{herm}(\mathcal H),\tfrac{1}{i\hbar}[\ ,\ ])
\arrow{r}{\hat f\mapsto dU(f)} & \mathfrak u(\mathcal H) \arrow{r}{\exp}
& \mathcal U(\mathcal H)
\end{tikzcd}
\]
\emph{Warum diese Räume?}
\begin{itemize}
\item $C^\infty(\mathcal P)$: trägt die \emph{Poisson}-Lie-Algebra der klassischen Observablen.
\item $\mathfrak h_n\subset C^\infty(\mathcal P)$: \emph{minimale} kanonische Generatoren (lineare Observablen) und zentrale Erweiterung.
\item $H_n$: exponentiierbare, globale Träger der Kanonik (Weyl-Form), kompatibel mit BCH.
\item $\mathfrak u(\mathcal H)$, $\mathcal U(\mathcal H)$: natürliche (Banach-)Lie-Strukturen für \emph{unitäre} (reversible) Dynamik.
\item $dU,\exp$: verbinden \emph{infinitesimale} und \emph{endliche} Transformationen funktoriell.
\end{itemize}

%------------------------------------------------------------
\emph{(E) Vergleich: QM vs. QFT (Rollen der Gruppen/Algebren)}

\emph{(E1) Endlich viele Freiheitsgrade (QM).}
\[
\mathfrak h_n \xrightarrow{\exp} H_n
\quad\text{und}\quad
\mathbb R X_t \xrightarrow{\exp} \mathbb R_t
\]
sind die \emph{zentralen} Gruppen:
$H_n$ (Kanondeformation) liefert CCR, $\mathbb R_t$ liefert Zeitflüsse.

\emph{(E2) Unendlich viele Freiheitsgrade (QFT).}
Die Rolle von $(q^i,p_i)$ wird durch \emph{Felder} $(\phi(\mathbf x),\Pi(\mathbf x))$ ersetzt:
\[
\{\phi(\mathbf x),\Pi(\mathbf y)\}=\delta^3(\mathbf x-\mathbf y)
\quad\leadsto\quad
[\hat\phi(\mathbf x),\hat\Pi(\mathbf y)]=i\hbar\,\delta^3(\mathbf x-\mathbf y),
\]
d.\,h.\ eine \emph{kontinuierliche Heisenberg-Algebra} (CCR mit kontinuierlichen Indizes),
deren exponentierte Weyl-Form auf dem Fockraum dargestellt wird.
Zeit- und Raumtranslationsgruppe (Poincaré-Untergruppen) liefern die dynamischen/unitären Flüsse.

%------------------------------------------------------------
\emph{(F) Fazit: Warum gerade diese Bausteine?}

\begin{itemize}
\item \textbf{Heisenberg-Algebra/-Gruppe} sind die \emph{minimalen} Träger der \emph{kanonischen Struktur}
(„$q$–$p$-Paare + zentrale Phase“), exponentierbar via BCH (Weyl-Operatoren).
\item \textbf{Unitäre Gruppen $\mathcal U(\mathcal H)$} und ihre \textbf{Lie-Algebren $\mathfrak u(\mathcal H)$}
tragen \emph{reversible Dynamik} und \emph{Spektraltheorie} (Selbstadjungiertheit, Stone).
\item \textbf{Zeit-/Raumtranslationsgruppen} integrieren die Generatoren (Hamilton-/Impulsoperatoren) zu \emph{globalen Flüssen}.
\item \textbf{Poisson $\to$ Kommutator} ist die \emph{Lie-funktoriale} Brücke:
sie bewahrt die „erste Abweichung von Komposition“ und erklärt die Struktur der CCR.
\item \textbf{Stone–von Neumann} fixiert (bis auf Unitäräquivalenz) die Schrödinger-Darstellung:
das erklärt die \emph{Eindeutigkeit} der kanonischen Quantisierung (QM).
\end{itemize}

\emph{Schlussformel (Sammelbox).}
\[
\boxed{
\begin{aligned}
&\text{Klassik: }(C^\infty(T^*Q),\{\ ,\ \})\ \leadsto\ \mathrm{Ham}(T^*Q)\\
&\text{Brücke: } \mathfrak h_n \xrightarrow{\exp} H_n,\quad
\text{CCR in Weyl-Form (BCH)}\\
&\text{Quant: } dU(\mathfrak h_n)\subset \mathfrak u(\mathcal H),\ 
[\hat q^i,\hat p_j]=i\hbar\,\delta^i_j,\ 
U(t)=e^{-\,\frac{i}{\hbar}t\hat H}.
\end{aligned}
}
\]



\pagebreak

[Kanonische Quantisierung als endlichdimensionale Lie-Algebra-Quantisierung]

\textbf{(1) Klassische Ebene: Phasenraum und Poisson-Algebra.}

In der klassischen Mechanik eines Teilchens mit Konfigurationsraum $Q\simeq\mathbb R^n$
ist der Phasenraum
\[
\mathcal P = T^*Q \;\simeq\; \mathbb R^{2n}
\]
mit Koordinaten $(q^i,p_i)$ ausgestattet mit der kanonischen symplektischen Form
\[
\omega = \sum_{i=1}^n dq^i \wedge dp_i.
\]
Daraus ergibt sich die \emph{Poisson-Klammer}
\[
\{f,g\}
=\sum_{i=1}^n
\Big(
\frac{\partial f}{\partial q^i}\frac{\partial g}{\partial p_i}
-\frac{\partial f}{\partial p_i}\frac{\partial g}{\partial q^i}
\Big),
\]
wobei $f,g\in C^\infty(\mathcal P)$ Observablen sind.
Insbesondere gilt für die Basisfunktionen
\[
\{q^i,p_j\}=\delta^i_j,
\qquad
\{q^i,q^j\}=0,
\qquad
\{p_i,p_j\}=0.
\]

Damit bildet der lineare Raum
\[
\mathfrak g_\text{class}=\operatorname{span}\{q^i,p_i\mid i=1,\ldots,n\}
\]
mit der Poisson-Klammer $\{\ ,\ \}$ eine Lie-Algebra, die \emph{Heisenberg-Lie-Algebra}.

\textbf{(2) Quantisierung als Funktor der Lie-Algebren.}

Die Quantisierung ist eine Zuordnung
\[
\mathfrak g_\text{class} \xrightarrow{\text{Quantisierung}}
\mathfrak u(\mathcal H),
\]
die die Lie-Klammer bis auf den Faktor $i\hbar$ erhält:
\[
\boxed{
\frac{1}{i\hbar}[\hat f,\hat g]\;\approx\;\widehat{\{f,g\}}.
}
\]

Auf den Generatoren gilt:
\[
q^i\longmapsto\hat q^i,
\qquad
p_i\longmapsto\hat p_i,
\]
mit den kanonischen Kommutatorrelationen
\[
[\hat q^i,\hat p_j]=i\hbar\,\delta^i_j,
\qquad
[\hat q^i,\hat q^j]=0,
\qquad
[\hat p_i,\hat p_j]=0.
\]

Diese bilden die \emph{Heisenberg-Lie-Algebra}
\[
\mathfrak h_n=\operatorname{span}\{\hat q^i,\hat p_i,\hat 1\}
\quad\text{mit}\quad
[\hat q^i,\hat p_j]=i\hbar\,\delta^i_j\,\hat 1.
\]

\textbf{(3) Diagrammatische Darstellung der Quantisierung.}

\[
\begin{tikzcd}[column sep=huge,row sep=large]
\text{klassische Variablen} 
\arrow{r}{\text{Poisson-Klammer }\{\ ,\ \}}
\arrow{d}{f\mapsto\hat f}
& \mathfrak g_\text{class}
\arrow{d}{\text{Quantisierung}}
\\
\text{Operatoren auf }\mathcal H
\arrow{r}{\frac{1}{i\hbar}[\ ,\ ]}
& \mathfrak u(\mathcal H)
\end{tikzcd}
\]

Damit wird die Poisson-Algebra $(C^\infty(T^*Q),\{\ ,\ \})$
durch die Operator-Algebra $(\mathcal B(\mathcal H),\frac{1}{i\hbar}[\ ,\ ])$ ersetzt.

\textbf{(4) Schrödinger-Darstellung.}

Die einfachste (und für Teilchenmechanik fundamentale) Darstellung ist die \emph{Schrödinger-Darstellung}:
\[
\mathcal H = L^2(\mathbb R^n),
\qquad
(\hat q^i\psi)(q)=q^i\psi(q),
\qquad
(\hat p_i\psi)(q)=-\,i\hbar\,\frac{\partial\psi}{\partial q^i}.
\]

Die kanonischen Kommutatorrelationen
\[
[\hat q^i,\hat p_j]\psi = i\hbar\,\delta^i_j\,\psi
\]
sind direkt erfüllt.

\textbf{(5) Hamiltonoperator und unitäre Flüsse.}

Für eine klassische Hamiltonfunktion $H(q,p)$
setzt man formal
\[
\hat H = H(\hat q,\hat p),
\]
z.\,B.
\[
\hat H = \frac{1}{2m}\hat p^2 + V(\hat q).
\]
Dann erzeugt $\hat H$ über die unitäre Einparametergruppe
\[
U(t)=\exp\!\Big(-\frac{i}{\hbar}t\hat H\Big)
\]
den Fluss der Zeitentwicklung im Hilbertraum:
\[
\psi(t)=U(t)\psi_0,
\qquad
\dot\psi(t)=-\frac{i}{\hbar}\hat H\psi(t).
\]
Die Schrödinger-Gleichung ist also die Darstellung der Lie-Gruppenwirkung
der Zeittranslationsgruppe $\mathbb R_t$ auf $\mathcal H$.

\textbf{(6) Physikalische Interpretation.}

\begin{itemize}
\item Die klassischen Observablen $(q,p)$ werden zu Operatoren $(\hat q,\hat p)$,
      die die algebraische Struktur der Poisson-Klammern auf der Operator-Ebene realisieren.
\item Die Quantisierung ist ein Funktor von der Poisson-Algebra auf die Operator-Algebra:
  \[
  \{f,g\}\;\longleftrightarrow\;\frac{1}{i\hbar}[\hat f,\hat g].
  \]
\item Die unitäre Gruppe $\mathcal U(\mathcal H)$ ersetzt den klassischen Phasenfluss auf $T^*Q$;
      ihre infinitesimale Version $\mathfrak u(\mathcal H)$ ersetzt die klassischen Hamiltonschen Vektorfelder.
\end{itemize}

\textbf{(7) Vergleich zur Feldquantisierung.}

\begin{center}
\begin{tabular}{lll}
 & \textbf{Quantenmechanik} & \textbf{Quantenfeldtheorie}\\
Koordinaten & $q^i$ (endlich viele) & $\phi(\mathbf x)$ (kontinuierlich viele)\\
Impulse & $p_i$ & $\Pi(\mathbf x)$\\
Klammer & $\{q^i,p_j\}=\delta^i_j$ & $\{\phi(\mathbf x),\Pi(\mathbf y)\}=\delta^3(\mathbf x-\mathbf y)$\\
Kommutator & $[\hat q^i,\hat p_j]=i\hbar\,\delta^i_j$ & $[\hat\phi(\mathbf x),\hat\Pi(\mathbf y)]=i\hbar\,\delta^3(\mathbf x-\mathbf y)$\\
Phasenraum & $T^*\mathbb R^n$ & unendlichdim.\ Phasenraum der Felder\\
Hilbertraum & $L^2(\mathbb R^n)$ & Fockraum $\mathcal H_\text{Fock}$\\
\end{tabular}
\end{center}

\textbf{(8) Diagrammatische Gesamtstruktur.}

\[
\begin{tikzcd}[column sep=huge,row sep=large]
\text{Klassische Ebene: }T^*Q
\arrow{r}{\text{Poisson-Fluss}}
\arrow{d}{\text{Quantisierung}}
& \mathfrak g_\text{class}\simeq\mathfrak h_n
\arrow{d}{\text{Heisenberg-Darstellung}}\\
\text{Quantenebene: }\mathcal H
\arrow{r}{\text{unitärer Fluss }U(t)=e^{-\,\frac{i}{\hbar}t\hat H}}
& \mathfrak u(\mathcal H)
\end{tikzcd}
\]

\textbf{(9) Fazit.}

\[
\boxed{
\begin{aligned}
&\text{Klassische Mechanik:}\;
(\mathcal P,\omega,\{\ ,\ \})\text{ mit Hamiltonschen Flüssen.}\\
&\text{Quantisierung:}\;
f\mapsto\hat f,\;\{f,g\}\mapsto\tfrac{1}{i\hbar}[\hat f,\hat g].\\
&\text{Quantenmechanik:}\;
(\mathcal H,\langle\cdot,\cdot\rangle,\hat H)\text{ mit unitären Flüssen.}
\end{aligned}
}
\]




\pagebreak

Diagrammatische Sicht auf Quantisierung


\textbf{(1) Ausgangspunkt: klassische Lie-Struktur.}

Sei $G$ eine Lie-Gruppe, $\mathfrak g$ ihre Lie-Algebra und $\mathfrak g^*$ der duale Raum.
Jedes $f\in C^\infty(\mathfrak g^*)$ kann als glatte Funktion auf den
Koordinaten der „klassischen Observablen“ verstanden werden.
Die \emph{Lie-Poisson-Struktur} auf $\mathfrak g^*$ ist definiert durch
\[
\{f,g\}(\xi)
=\langle \xi, [df_\xi,dg_\xi]\rangle,
\qquad
\xi\in\mathfrak g^*,\;
f,g\in C^\infty(\mathfrak g^*),
\]
wobei $df_\xi\in\mathfrak g$ der differentielle Kofaktor von $f$ bei $\xi$ ist.
Damit ist $(\mathfrak g^*,\{\cdot,\cdot\})$ eine \emph{Poisson-Mannigfaltigkeit};
ihre symplektischen Blätter sind die \emph{Koadjungierten Orbits} $G\cdot\xi\subset\mathfrak g^*$.

\textbf{(2) Diagrammatische Form der klassischen Ebene.}

\[
\begin{tikzcd}[column sep=huge,row sep=large]
\mathfrak g \arrow{r}{\exp}
\arrow{d}{\text{Koadjunktion}}
& G \arrow{d}{\text{Ad}^*} \\
\mathfrak g^* \arrow{r}{\text{Orbits in }\mathfrak g^*}
& G\text{-Orbits (klassische Phasenräume)}
\end{tikzcd}
\]

\begin{itemize}
\item Oben: Lie-Algebra und -Gruppe kodieren infinitesimale bzw.\ endliche Symmetrien.
\item Unten: deren duale Koadjungierten-Orbits tragen eine natürliche symplektische Form
  (Kostant–Kirillov–Souriau) und bilden die klassischen Phasenräume.
\end{itemize}

\textbf{(3) Quantisierung als Funktor zwischen Diagrammen.}

Quantisierung ist die Zuweisung
\[
\text{(klassische observablen Algebra) }
\big(C^\infty(\mathfrak g^*),\{\ ,\ \}\big)
\longrightarrow
\text{(Quanten-Operator-Algebra) }
\big(\mathfrak u(\mathcal H),\tfrac{1}{i\hbar}[\ ,\ ]\big),
\]
so dass die Lie-Klammer der Operatoren der Poisson-Klammer der Observablen entspricht:
\[
\boxed{
\frac{1}{i\hbar}[\hat f,\hat g]\;\approx\;\widehat{\{f,g\}}.
}
\]

Diagrammatisch:

\[
\begin{tikzcd}[column sep=huge,row sep=large]
\text{Klassische Seite:} & \mathfrak g^* \arrow{r}{\{\cdot,\cdot\}} \arrow{d}{\text{Quantisierung}}
& C^\infty(\mathfrak g^*) \arrow{d}{f\mapsto \hat f} \\
\text{Quantenseite:} & \mathcal H \arrow{r}{\tfrac{1}{i\hbar}[\cdot,\cdot]}
& \mathfrak u(\mathcal H)
\end{tikzcd}
\]

\textbf{(4) Coadjungierte Orbits und Darstellungen.}

Jeder Orbit $\mathcal O_\xi=G\cdot\xi\subset\mathfrak g^*$
trägt eine symplektische Form $\omega_\xi$.
Nach dem Satz von Kirillov–Kostant–Souriau gilt:

\[
\text{„Irreduzible unitäre Darstellungen von $G$
\quad$\leftrightarrow$\quad
Quantisierung der Koadjungierten-Orbits in $\mathfrak g^*$.“}
\]

Formal:
\[
\begin{tikzcd}[column sep=huge,row sep=large]
\text{Klassisch: } \mathcal O_\xi\subset \mathfrak g^* \arrow{r}{\text{Geometr. Quantisierung}}
& \text{Hilbertraum } \mathcal H_\xi \arrow{r}{U_\xi}
& \text{Irred. Darstellung von } G
\end{tikzcd}
\]

Damit interpretiert sich eine Darstellung $U_\xi$ als „quantisierte Version“ des klassischen Phasenraums $\mathcal O_\xi$.
Der Zusammenhang zwischen Lie-Klammer und Kommutator bleibt erhalten:
\[
[dU(X),dU(Y)]=dU([X,Y]),
\qquad
\text{klassisch: }\;\{f_X,f_Y\}=f_{[X,Y]}.
\]

\textbf{(5) Physikalische Interpretation.}

\begin{itemize}
\item \emph{Klassisch:} Zustände sind Punkte $\xi\in\mathfrak g^*$,
  Dynamik erfolgt entlang der Hamiltonschen Flüsse der Poisson-Klammer.
\item \emph{Quantisiert:} Zustände sind Strahlen $[\psi]\in\mathbb P(\mathcal H)$,
  Dynamik erfolgt durch unitäre Flüsse $U(t)=\exp(-\tfrac{i}{\hbar}t\hat H)$.
\item Die Übergangsregel
  \[
  f_X(\xi)=\langle \xi,X\rangle \quad\longleftrightarrow\quad \hat X=i\hbar\,dU(X)
  \]
  identifiziert die klassische Observablen mit den Generatoren der unitären Gruppe.
\end{itemize}

\textbf{(6) Diagrammatische Gesamtstruktur.}

\[
\begin{tikzcd}[column sep=huge,row sep=large]
\mathfrak g \arrow{r}{\exp} \arrow{d}{X\mapsto f_X(\xi)=\langle\xi,X\rangle}
& G \arrow{d}{\text{Koadjunktion}} \arrow[bend left=35]{dd}{U} \\
\mathfrak g^* \arrow{r}{\text{Orbit}} \arrow{d}{\text{Quantisierung}}
& \mathcal O_\xi \arrow{d}{\text{Geometr. Quantisierung}} \\
\mathcal H_\xi \arrow{r}{\text{Darstellung}}
& \mathcal U(\mathcal H_\xi)
\end{tikzcd}
\]

\textbf{(7) Zusammenfassung.}
\[
\boxed{
\begin{aligned}
&\text{Klassisch: Lie-Algebra }(\mathfrak g,\{\ ,\ \})\text{ wirkt auf }\mathfrak g^*,\\[2pt]
&\text{Quantisiert: }(\mathfrak u(\mathcal H),\tfrac{1}{i\hbar}[\ ,\ ])\text{ wirkt auf }\mathcal H,\\[2pt]
&\text{Quantisierung: } f_X\mapsto \hat X=i\hbar\,dU(X)
\text{ verbindet beide Ebenen,}\\[2pt]
&\text{so dass der Fluss der Poisson-Algebra zu unitären Flüssen in }\mathcal H\text{ wird.}
\end{aligned}
}
\]





[Feldquantisierung als kontinuierliche Lie-Algebra-Quantisierung]

\textbf{(1) Übergang von klassischen zu quantisierten Variablen.}

Im klassischen Feldformalismus sind die dynamischen Variablen Funktionen
\[
\phi:\mathbb R^{1,3}\to\mathbb R,\qquad
\Pi:\mathbb R^{1,3}\to\mathbb R,
\]
die auf jedem festen Zeitpunkt $t$ die Punkte $(\phi(\mathbf x,t),\Pi(\mathbf x,t))$ des
unendlichdimensionalen Phasenraums
\[
\mathcal P=\{(\phi(\cdot),\Pi(\cdot))\mid \phi,\Pi\in C^\infty(\mathbb R^3)\}
\]
bestimmen.

Darauf ist die \emph{Poisson-Struktur} gegeben durch
\[
\{F,G\}=\int d^3x\,
\Big(
\frac{\delta F}{\delta \phi(\mathbf x)}\frac{\delta G}{\delta \Pi(\mathbf x)}
-
\frac{\delta F}{\delta \Pi(\mathbf x)}\frac{\delta G}{\delta \phi(\mathbf x)}
\Big),
\]
so dass insbesondere
\[
\{\phi(\mathbf x,t),\Pi(\mathbf y,t)\}=\delta^3(\mathbf x-\mathbf y),
\qquad
\{\phi(\mathbf x,t),\phi(\mathbf y,t)\}=0=\{\Pi(\mathbf x,t),\Pi(\mathbf y,t)\}.
\]

Dies ist die \emph{kontinuierliche} Version der kanonischen Poisson-Algebra
\[
\mathfrak g_\text{class}
=\operatorname{span}\{\phi(\mathbf x),\Pi(\mathbf x)\}_{\mathbf x\in\mathbb R^3},
\quad
\{\cdot,\cdot\}:\mathfrak g_\text{class}\times\mathfrak g_\text{class}\to C^\infty(\mathbb R^3).
\]

\textbf{(2) Diagrammatische Quantisierung.}

Quantisierung bedeutet, dass man diese Poisson-Algebra in eine Operator-Lie-Algebra überführt:
\[
\begin{tikzcd}[column sep=huge,row sep=large]
\text{klassische Algebra }
(\mathfrak g_\text{class},\{\ ,\ \})
\arrow{r}{\text{Quantisierung}}
& \text{Operator-Algebra }
(\mathfrak g_\text{quant},\tfrac{1}{i\hbar}[\ ,\ ]).
\end{tikzcd}
\]

Explizit:
\[
\phi(\mathbf x)\longmapsto \hat\phi(\mathbf x),\qquad
\Pi(\mathbf x)\longmapsto \hat\Pi(\mathbf x),
\]
mit den Kommutatorregeln
\[
[\hat\phi(\mathbf x,t),\hat\Pi(\mathbf y,t)] = i\hbar\,\delta^3(\mathbf x-\mathbf y),
\quad
[\hat\phi(\mathbf x,t),\hat\phi(\mathbf y,t)]=0,
\quad
[\hat\Pi(\mathbf x,t),\hat\Pi(\mathbf y,t)]=0.
\]
Dies ist die \emph{kanonische Quantisierung}: dieselbe Funktorregel wie im endlichdimensionalen Fall,
nur dass die diskreten Indizes $(i,j)$ durch kontinuierliche Koordinaten $(\mathbf x,\mathbf y)$ ersetzt sind.

\textbf{(3) Lie-theoretischer Hintergrund.}

Die klassische Poisson-Klammer ist die Lie-Algebra-Struktur auf der
Raumzeit-abhängigen Observablenalgebra $C^\infty(\mathcal P)$.
Die Quantisierung ersetzt diese durch eine unitäre Darstellung
\[
dU:\mathfrak g_\text{class}\longrightarrow \mathfrak u(\mathcal H_\text{Fock}),
\]
deren Bild
\[
dU(\phi(\mathbf x))=\frac{1}{i\hbar}\hat\phi(\mathbf x),\qquad
dU(\Pi(\mathbf x))=\frac{1}{i\hbar}\hat\Pi(\mathbf x)
\]
den infinitesimalen Generatoren der unitären Gruppe der Fockraum-Darstellung entspricht.

\[
\begin{tikzcd}[column sep=huge,row sep=large]
\mathfrak g_\text{class} \arrow{r}{\text{Quantisierung}}
\arrow{d}{\text{Poisson-Klammer}}
& \mathfrak u(\mathcal H_\text{Fock}) \arrow{d}{\text{Kommutator }[\ ,\ ]}\\
C^\infty(\mathcal P) \arrow{r}{f\mapsto \hat f}
& \mathcal B(\mathcal H_\text{Fock})
\end{tikzcd}
\]

Das bedeutet:
\[
\frac{1}{i\hbar}[\hat f,\hat g]\;\approx\;\widehat{\{f,g\}}.
\]
Die Fockraumdarstellung $\mathcal H_\text{Fock}$ ist somit die „unitäre Repräsentation“
der unendlichdimensionalen Lie-Gruppe,
deren Lie-Algebra durch die klassischen Feld-Observablen gebildet wird.

\textbf{(4) Hamilton- und Impulsoperatoren.}

Die klassischen Größen
\[
H=\int d^3x\,\mathscr H(\phi,\Pi,\nabla\phi),\qquad
P^i=\int d^3x\,\Pi\,\partial^i\phi
\]
werden durch denselben Funktor zu Operatoren
\[
\hat H=\int d^3x\,\hat{\mathscr H}(\hat\phi,\hat\Pi,\nabla\hat\phi),\qquad
\hat P^i=\int d^3x\,\hat\Pi\,\partial^i\hat\phi.
\]
Sie sind die Generatoren der durch $U(a)$ beschriebenen Flüsse im Fockraum:
\[
U(a^0)=e^{-\,\frac{i}{\hbar}a^0\hat H},\qquad
U(\mathbf a)=e^{-\,\frac{i}{\hbar}\mathbf a\cdot\hat{\mathbf P}},
\]
d.h.\ die Zeit- und Raumtranslationen der Felder sind
\[
U(a^0)^\dagger\hat\phi(x)U(a^0)=\hat\phi(x+a^0e_0),\qquad
U(\mathbf a)^\dagger\hat\phi(x)U(\mathbf a)=\hat\phi(x+\mathbf a).
\]

\textbf{(5) Interpretative Bedeutung.}

\begin{itemize}
\item Die Poisson-Klammern der klassischen Felder definieren eine (unendlichdimensionale) Lie-Algebra
   der Observablen.
\item Die Quantisierung ist der Funktor, der diese algebraische Struktur in eine
   unitäre Darstellung auf dem Fockraum überführt.
\item Die Kommutatorrelationen sind die \emph{Lie-Klammern} in $\mathfrak u(\mathcal H_\text{Fock})$,
   die die Poisson-Klammern reproduzieren.
\item Damit ist die Feldquantisierung der Fall einer \emph{kontinuierlichen Lie-Gruppen-Darstellung},
   deren Generatoren $\hat\phi(x),\hat\Pi(x)$ durch die Felder an jedem Ort gegeben sind.
\end{itemize}

\textbf{(6) Diagrammatische Gesamtstruktur.}

\[
\begin{tikzcd}[column sep=huge,row sep=large]
\text{Klassische Variablen}
\arrow{r}{\text{Poisson-Klammer}}
\arrow{d}{\text{Quantisierung}}
& \mathfrak g_\text{class}\subset C^\infty(\mathcal P)
\arrow{d}{f\mapsto \hat f}
\\
\text{Operatoren auf }\mathcal H_\text{Fock}
\arrow{r}{\frac{1}{i\hbar}[\ ,\ ]}
& \mathfrak u(\mathcal H_\text{Fock})
\end{tikzcd}
\]

\textbf{(7) Fazit.}

\[
\boxed{
\begin{aligned}
&\text{Die Feldquantisierung ist die Anwendung der Lie-Funktorialität}\\
&(\mathfrak g_\text{class},\{\ ,\ \})\;\longmapsto\;
(\mathfrak u(\mathcal H_\text{Fock}),\tfrac{1}{i\hbar}[\ ,\ ])\\[3pt]
&\text{mit Generatoren }(\phi(\mathbf x),\Pi(\mathbf x))\mapsto(\hat\phi(\mathbf x),\hat\Pi(\mathbf x)).\\[3pt]
&\text{Die Hamilton- und Impulsoperatoren sind die zugehörigen Flussgeneratoren,}\\
&\text{die Translationen im Raumzeitparameterraum erzeugen.}
\end{aligned}
}
\]


\pagebreak



\textbf{Standardperspektiven auf physikalische Probleme in der Quantenfeldtheorie}

\emph{(1) Lie-Struktur als Symmetrie-Organisationsprinzip.}
Die fundamentale Idee:
\[
\text{Dynamik, Erhaltungsgrößen und Wechselwirkungen sind Ausdruck von Symmetrie.}
\]
Ausgehend vom Diagramm
\[
\mathfrak g \xrightarrow{dU} \mathfrak u(\mathcal H) \xrightarrow{\times i\hbar} \mathfrak{herm}(\mathcal H),
\]
nutzt man die Lie-Algebra, um:
\begin{itemize}
\item die Generatoren $\hat X=i\hbar dU(X)$ zu identifizieren (z.\,B. Energie, Impuls, Drehimpuls, Ladung);
\item Kommutatoren $[\hat X,\hat Y]=i\hbar\,\widehat{[X,Y]}$ als algebraische Relationen zwischen Observablen zu verstehen;
\item daraus Auswahlregeln, Erhaltungsgrößen und Dynamikbedingungen abzuleiten.
\end{itemize}
Diese Perspektive ist in der QFT universell:
\[
\text{(Symmetrie)}\quad \Longleftrightarrow\quad \text{(Erhaltungsgröße nach Noether)}.
\]

\emph{(2) Die Darstellungs-Perspektive (Wigner).}
In der relativistischen QFT:
\begin{itemize}
\item Die Poincaré-Gruppe $G=\mathcal P$ (Translations $\ltimes$ Lorentz) ist die fundamentale Symmetrie der Raumzeit.
\item Der physikalische Zustandsraum besteht aus einer \emph{unitären, irreduziblen Darstellung}
      $U:\mathcal P\to\mathcal U(\mathcal H)$.
\item Diese Darstellung ist durch die Casimir-Operatoren (Invarianten der Algebra)
      \[
      \hat P^2=m^2c^2,\qquad \hat W^\mu\hat W_\mu=-m^2s(s+1)\hbar^2
      \]
      charakterisiert — also Masse und Spin.
\end{itemize}
→ Diese Sichtweise führt direkt zur Klassifikation der Teilchenarten als irreduzible Darstellungen der Symmetriegruppe.

\emph{(3) Die Operator- und Kommutator-Perspektive.}
Hier arbeitet man \emph{im Hilbertraum selbst}:
\[
[\hat H,\hat A]=0 \;\Longrightarrow\; \text{$\hat A$ ist Erhaltungsgröße.}
\]
Man analysiert Strukturen wie:
\begin{itemize}
\item Kommutatoren als „infinitesimale“ Form der Lie-Symmetrie;
\item Spektralzerlegungen der Operatoren als physikalische Quantisierung der Observablen;
\item Algebren der Felder (z.\,B. kanonische Vertauschungsrelationen, CAR/CCR);
\item Repräsentationen dieser Algebren (Fock-Raum, Verteilungen).
\end{itemize}
→ In der QFT: „Quantisierung“ = Wahl einer *Darstellung der Vertauschungsalgebra*.

\emph{(4) Die Bündel- und Eich-Perspektive.}
Die Lie-Gruppen- und Darstellungsstruktur erweitert sich zu
\[
P(M,G) \;\to\; E=P\times_D W,
\]
mit Zusammenhang $A\in\Omega^1(M,\mathfrak g)$.
Die Lie-Algebra $\mathfrak g$ definiert:
\begin{itemize}
\item die internen Freiheitsgrade (Farben, Isospin, Chirality);
\item die Eichfeldstärke $F=dA+\tfrac{1}{2}[A,A]$ als geometrischen Ausdruck der Nichtkommutativität von $\mathfrak g$;
\item die kovariante Ableitung $D_\mu=\partial_\mu+A_\mu$ als lokale Umsetzung der Lie-Wirkung.
\end{itemize}
→ Die Feldgleichungen (Yang–Mills, Dirac, Maxwell) sind genau die \emph{kovarianten Operatorgleichungen}
unter dieser Lie-Struktur.

\emph{(5) Die Dynamik-Perspektive.}
Über den Hamiltonoperator $\hat H$ als Generator der Zeittranslation:
\[
U(t)=e^{-\,\frac{i}{\hbar}t\hat H},\qquad i\hbar\,\dot\psi=\hat H\psi.
\]
Man unterscheidet:
\begin{itemize}
\item Zeitentwicklung als Fluss in $\mathcal U(\mathcal H)$;
\item Heisenberg-Bild: Operatoren evolvieren, Zustände fix;
\item Schrödinger-Bild: Zustände evolvieren, Operatoren fix;
\item Wechselwirkungsbild: gemischte Beschreibung (zeitabhängige Lie-Algebra-Kurven $X(t)$).
\end{itemize}
→ Diese Sichtweise übersetzt alle physikalischen Prozesse in geometrische Flüsse in der Lie-Gruppe $\mathcal U(\mathcal H)$.

\emph{(6) Die Algebraische QFT-Perspektive.}
Hier abstrahiert man vollständig von konkreten Darstellungen:
\[
\mathcal A\subset \mathfrak B(\mathcal H)
\]
ist die $C^*$- oder von-Neumann-Algebra der Observablen;
$G$ wirkt durch Automorphismen
\[
\alpha_g:\mathcal A\to\mathcal A,\qquad \alpha_g(A)=U(g)AU(g)^{-1}.
\]
Dann:
\begin{itemize}
\item Zustände sind positive, normierte lineare Funktionale $\omega:\mathcal A\to\mathbb C$.
\item Das GNS-Konstrukt liefert die Darstellung $(\pi_\omega,\mathcal H_\omega,U_\omega)$.
\item Symmetrie, Invarianz, Dynamik werden auf der Ebene der Automorphismen beschrieben, unabhängig von einer festen Hilbertraum-Realisierung.
\end{itemize}
→ Dies ist die modernste Perspektive: Lie-Strukturen erscheinen als Automorphismengruppen der Observable-Algebra.

\emph{(7) Die Geometrisch-Analytische Perspektive.}
Auf den Jet- und Bündel-Strukturen aus Abschnitt~(12) aufbauend:
\begin{itemize}
\item Zustände sind Schnitte eines assoziierten Bündels $E=P\times_D W$.
\item Operatoren (Dirac, Laplace, Yang–Mills) sind kovariante Differentialoperatoren, 
  die aus der Lie-Algebra-Struktur des Zusammenhangs $A$ gebaut werden.
\item Invarianzbedingung $[dU(X),E]=0$ charakterisiert Forminvarianz und physikalische Symmetrie.
\end{itemize}
→ Diese Sichtweise dominiert in Geometrischer und Supersymmetrischer QFT.

\emph{(8) Synthese.}
Alle Perspektiven folgen aus demselben Diagramm:
\[
\boxed{
\begin{aligned}
\mathfrak g &\xrightarrow{dU}& \mathfrak u(\mathcal H)
 &\xrightarrow{\times i\hbar}& \mathfrak{herm}(\mathcal H)\\
 & &\text{Geometrie}&\text{→}&\text{Physik (Observablen)}
\end{aligned}
}
\]
\[
\text{„Lie-Struktur = Geometrie der Beschreibungsänderungen“}
\quad\Rightarrow\quad
\text{„Observablen = Generatoren dieser Änderungen“}.
\]

\emph{(9) Typische Anwendungspfade.}
Wenn ein QFT-Problem gegeben ist, prüfe:
\begin{enumerate}
\item \textbf{Welche Symmetriegruppe $G$?} (z.\,B.\ Poincaré, Lorentz, interne Gaugegruppe)
\item \textbf{Welche Lie-Algebra $\mathfrak g$?} (Generatoren, Relationen)
\item \textbf{Welche Darstellung $U$ oder Feldrepräsentation $D$?}
\item \textbf{Welche invarianten Operatoren / Casimirs?} (Masse, Spin, Ladung)
\item \textbf{Wie wirkt $dU(\mathfrak g)$ auf Observablen?} (Erhaltung, Kommutatoren)
\item \textbf{Wie entstehen Operatorgleichungen?} ($E[\psi]=0$, kovariant unter $G$)
\item \textbf{Welche Invarianzbedingung gilt?} ($[dU(X),E]=0$)
\item \textbf{Welche Quotienten oder Moduli?} ($\mathcal S/G$, physikalischer Raum der Zustände)
\end{enumerate}

\emph{(10) Zusammenfassung.}
\[
\boxed{
\begin{aligned}
&\text{Geometrie (Lie-Struktur):}\; G,\mathfrak g\\
&\text{Operatorisch (Darstellung):}\; U(G),dU(\mathfrak g)\\
&\text{Physikalisch (QFT):}\; \text{Felder, Observablen, Erhaltungsgrößen, Dynamik}\\[2pt]
&\Rightarrow \text{Jedes QFT-Problem kann als Frage über die Lie-Struktur von $G$ interpretiert werden.}
\end{aligned}
}
\]


\textbf{Standardperspektiven in Quantenmechanik und Quantenfeldtheorie I}

\emph{(1) Strukturrahmen.}
Beide Theorien sind spezifische Realisierungen des allgemeinen Diagramms
\[
\mathfrak g \xrightarrow{dU} \mathfrak u(\mathcal H)
\xrightarrow{\times i\hbar} \mathfrak{herm}(\mathcal H),
\]
mit jeweils unterschiedlichen Lie-Gruppen $G$ und Darstellungen $U$:
\begin{itemize}
\item Quantenmechanik (QM): endlich- oder separabel unendlichdimensionale $\mathcal H$, punktweise Zustände, wenige Freiheitsgrade.
\item Quantenfeldtheorie (QFT1): Fockräume, Felder als Operatoren, unendlich viele Freiheitsgrade.
\end{itemize}
Beide Theorien beruhen aber auf denselben fünf strukturellen Prinzipien:



\text{(i) Zustände in einem Hilbertraum,}

\text{(ii) Observablen als selbstadjungierte Operatoren,}

\text{(iii) Symmetrien als unitäre Darstellungen,}

\text{(iv) Dynamik als Fluss,}\quad

\text{(v) Invarianz als physikalisches Prinzip.}



\paragraph{(2) Quantenmechanik (QM).}
\emph{Standard-Symmetriegruppe:}
\[
G_{\text{QM}}=\mathrm{Galilei}\quad
(\text{bzw. bei Spin: ihre zentrale Erweiterung}).
\]
Ihre Lie-Algebra-Generatoren:
\[
\hat H\text{ (Zeittranslation)},\quad
\hat{\mathbf P}\text{ (Raumtranslation)},\quad
\hat{\mathbf J}\text{ (Rotation)},\quad
\hat{\mathbf K}\text{ (Galilei-Boost)}.
\]

\textbf{Perspektiven:}

\begin{itemize}
\item[(i)] \emph{Darstellungstheorie:}
  Zustände sind Vektoren in einer unitären, (im Idealfall irreduziblen) Darstellung
  $U:\mathrm{Galilei}\to\mathcal U(\mathcal H)$.
  Die Casimir-Operatoren (Invarianten der Algebra)
  \[
  \hat M,\quad \hat H-\frac{\hat{\mathbf P}^2}{2\hat M}
  \]
  klassifizieren irreduzible Darstellungen — d.\,h.\ Teilchentypen mit Masse $m$ und Spin $s$.

\item[(ii)] \emph{Dynamik:}
  Die Zeitentwicklung ist die Einparameter-Untergruppe
  \[
  U(t)=e^{-\,\frac{i}{\hbar}t\hat H},\qquad i\hbar\dot\psi=\hat H\psi.
  \]
  Das ist die Schrödinger-Gleichung, formal die Darstellung der Lie-Algebra $\mathfrak{u}(1)_t$ (Zeittranslationen).

\item[(iii)] \emph{Kommutatoralgebra:}
  \[
  [\hat x_i,\hat p_j]=i\hbar\,\delta_{ij},\quad
  [\hat H,\hat x_i]=\frac{i\hbar}{m}\hat p_i,\quad
  [\hat J_i,\hat J_j]=i\hbar\,\varepsilon_{ijk}\hat J_k.
  \]
  Diese Relationen sind genau die Lie-Klammern der entsprechenden Symmetriealgebra.

\item[(iv)] \emph{Messung und Spektralstruktur:}
  Selbstadjungierte Operatoren $\hat A$ (z.\,B. $\hat H,\hat L_z,\hat p_x$) bestimmen Spektralräume und Messwahrscheinlichkeiten.
  Stationäre Zustände sind Eigenvektoren von $\hat H$:
  \[
  \hat H\psi_n=E_n\psi_n,\qquad
  U(t)\psi_n=e^{-\,\frac{i}{\hbar}E_n t}\psi_n.
  \]
  → Dynamik als Fluss in $\mathcal U(\mathcal H)$.

\item[(v)] \emph{Invarianzprinzip:}
  Wenn $[\hat H,dU(X)]=0$, dann ist $dU(X)$ eine Erhaltungsgröße
  (Impuls bei Translationsinvarianz, Drehimpuls bei Rotationsinvarianz etc.).
\end{itemize}

\textbf{Schlagwort:}
\[
\boxed{\text{„Galilei-Symmetrie  $\Rightarrow$  Lie-Algebra der kanonischen Operatoren.“}}
\]

\paragraph{(3) Quantenfeldtheorie I (QFT1).}
\emph{Standard-Symmetriegruppe:}
\[
G_{\text{QFT1}}=\mathrm{Poincaré}=\mathbb R^{1,3}\rtimes SO^+(1,3).
\]
Ihre Lie-Algebra-Generatoren:
\[
\hat P_\mu \text{ (Translationen)},\qquad
\hat M_{\mu\nu}\text{ (Lorentz-Generatoren)}.
\]
Mit Vertauschungsrelationen:
\[
[\hat M_{\mu\nu},\hat M_{\rho\sigma}]
=i\hbar\left(\eta_{\nu\rho}\hat M_{\mu\sigma}-\eta_{\mu\rho}\hat M_{\nu\sigma}
+\eta_{\mu\sigma}\hat M_{\nu\rho}-\eta_{\nu\sigma}\hat M_{\mu\rho}\right),
\]
\[
[\hat P_\mu,\hat M_{\rho\sigma}]
=i\hbar(\eta_{\mu\rho}\hat P_\sigma-\eta_{\mu\sigma}\hat P_\rho),\qquad
[\hat P_\mu,\hat P_\nu]=0.
\]

\textbf{Perspektiven:}

\begin{itemize}
\item[(i)] \emph{Darstellungstheorie (Wigner-Klassifikation):}
  Irreduzible unitäre Darstellungen der Poincaré-Gruppe sind durch die Casimirs
  \[
  \hat P^2=m^2 c^2,\qquad
  \hat W^\mu \hat W_\mu=-m^2s(s+1)\hbar^2
  \]
  charakterisiert (Masse und Spin).
  Diese zwei Zahlen bestimmen die \emph{Teilchenarten}.
  Jeder Zustand im Fockraum transformiert gemäß $U(\Lambda,a)=e^{-\,\frac{i}{\hbar}(a^\mu \hat P_\mu+\tfrac{1}{2}\omega^{\mu\nu}\hat M_{\mu\nu})}$.

\item[(ii)] \emph{Felder als Operatorwertige Darstellungen:}
  Feldoperatoren $\Phi(x)$ sind Operatoren auf $\mathcal H$ mit
  \[
  U(\Lambda,a)\,\Phi(x)\,U(\Lambda,a)^{-1}
  =D(\Lambda)\,\Phi(\Lambda^{-1}(x-a)),
  \]
  wobei $D$ eine Darstellung der Lorentz-Gruppe auf internen Indizes ist.
  → Dies ist exakt dieselbe Lie-Gruppenwirkung wie in Abschnitt~(12) (Bündel-/Schnitt-Perspektive).

\item[(iii)] \emph{Dynamik:}
  Der Hamiltonoperator ist Generator der Zeittranslation:
  \[
  U(t)=e^{-\,\frac{i}{\hbar}t\hat H},\qquad \hat H=\int d^3x\,T^{00}(x).
  \]
  Das Feld erfüllt eine kovariante Operatorgleichung, z.\,B.
  \[
  (\Box+m^2)\Phi(x)=0\quad\text{(Klein–Gordon)},\qquad
  (i\gamma^\mu\partial_\mu-m)\psi(x)=0\quad\text{(Dirac)}.
  \]
  → Jede ist formal ein Ausdruck von $[dU(X),E]=0$ für $X\in\mathfrak{poincaré}$.

\item[(iv)] \emph{Kommutator- und Kausalstruktur:}
  Lokale Observablen erfüllen
  \[
  [\Phi(x),\Phi(y)]=0\quad\text{für spacelike getrennte Punkte }(x-y),
  \]
  was die Lie-Algebra-Struktur der Raumzeit-Symmetrie (Minkowski) widerspiegelt.

\item[(v)] \emph{Erhaltungsgrößen und Ströme:}
  Noether-Theorem:
  \[
  J^\mu(X)=\frac{\partial\mathcal L}{\partial(\partial_\mu\phi)}\,\delta_X\phi,\quad
  Q(X)=\int d^3x\,J^0(X),
  \]
  $Q(X)$ ist der Generator der infinitesimalen Symmetrie $X\in\mathfrak g$,
  und erfüllt
  \[
  [Q(X),Q(Y)]=i\hbar\,Q([X,Y]).
  \]
  → Dies ist der explizite Zusammenhang zwischen der Lie-Algebra und der Operatoralgebra der Felder.
\end{itemize}

\textbf{Schlagwort:}
\[
\boxed{\text{„Poincaré-Symmetrie  $\Rightarrow$  Klassifikation und Dynamik der Felder.“}}
\]

\emph{(4) Vergleichstabelle.}

\[
\begin{array}{lll}
\textbf{Strukturelement} & \textbf{QM (Galilei)} & \textbf{QFT1 (Poincaré)}\\
\text{Lie-Gruppe } G & \mathrm{Galilei} & \mathrm{Poincaré}\\
\text{Lie-Algebra } \mathfrak g & \{H,\mathbf P,\mathbf J,\mathbf K\} & \{P_\mu,M_{\mu\nu}\}\\
\text{Darstellung } U & \text{Einpartikel} & \text{Fock- oder Feldraum}\\
\text{Casimirs} & M,\,H-\tfrac{\mathbf P^2}{2M} & P^2,\,W^\mu W_\mu\\
\text{Generatoren} & \hat H,\hat{\mathbf P},\hat{\mathbf J} & \hat P_\mu,\hat M_{\mu\nu}\\
\text{Observablen} & \text{Selbstadjungierte Operatoren} & \text{Feldoperatoren / Integrale}\\
\text{Dynamik} & i\hbar\dot\psi=\hat H\psi & [\hat H,\Phi(x)]=i\hbar\,\dot\Phi(x)\\
\text{Invarianz} & [\hat H,dU(X)]=0 & [Q(X),E]=0\\
\text{Physikalische Zustände} & [\psi]\in\mathbb P(\mathcal H) & \text{Fock-Zustände / Verteilungen}\\
\text{Kovarianzgesetz} & U(g)\psi(x)=\psi(g^{-1}x) & U(\Lambda,a)\Phi(x)U^{-1}(\Lambda,a)=D(\Lambda)\Phi(\Lambda^{-1}(x-a))\\
\end{array}
\]

\emph{(5) Zusammenfassung.}
\[
\boxed{
\begin{aligned}
&\text{QM: Lie-Struktur der Galilei-Gruppe  }\Rightarrow\text{ kanonische Operatorrelationen.}\\[2pt]
&\text{QFT1: Lie-Struktur der Poincaré-Gruppe  }\Rightarrow\text{ Feldgleichungen und Teilchenklassifikation.}\\[2pt]
&\text{Beide: Darstellungen }U(G)\text{ liefern Dynamik, Invarianz und Erhaltungsgrößen.}
\end{aligned}
}
\]

\pagebreak



\section*{Chapter 5}


\section*{Lie Groups}
One approach to geometry is to view it as the study of invariance and symmetry. In our case, we are interested in studying symmetries of smooth manifolds, Riemannian manifolds, symplectic manifolds, etc. The usual way to deal with symmetry in mathematics is by the use of the notion of a transformation group. The wonderful thing for us is that the groups that arise in the study of geometric symmetries are often themselves smooth manifolds. Such "group manifolds" are called Lie groups.

In physics, Lie groups play a big role in connection with physical symmetries and conservation laws (Noether's theorem). Within physics, perhaps the most celebrated role played by Lie groups is in particle physics and gauge theory. In mathematics, Lie groups play a prominent role in harmonic analysis (generalized Fourier theory), group representations, differential equations, and in virtually every branch of geometry including Riemannian geometry, Cartan geometry, algebraic geometry, Kähler geometry, and symplectic geometry.

\subsection*{Definitions and Examples}
Definition 5.1. 


\begin{DEF}{LIE-L09-04-01}{Lie Gruppe}
A smooth manifold $G$ is called a Lie group if it is a group (abstract group) such that the multiplication map $\mu: G \times G \rightarrow G$ and the inverse map inv : $G \rightarrow G$, given respectively by $\mu(g, h)=g h$ and $\operatorname{inv}(g)=g^{-1}$, are $C^{\infty}$ maps. If the group is abelian, we sometimes opt to use the additive notation $g+h$ for the group operation.

We will usually denote the identity element of any Lie group by the same letter $e$. Exceptions include the case of matrix or linear groups where we use the letter $I$ or id. The map inv : $G \rightarrow G$ given by $g \mapsto g^{-1}$ is called inversion and is easily seen to be a diffeomorphism.
\end{DEF}

Example 5.2. 

\begin{EXA}{LIE-L09-04-02}{Vektorräume als Lie Gruppen}
$\mathbb{R}$ is a one-dimensional (abelian) Lie group, where the group multiplication is the usual addition + . Similarly, any real or complex vector space is a Lie group under vector addition.
\end{EXA}

Example 5.3. 

\begin{EXA}{LIE-L09-04-03}{$S^1$ als Lie Gruppen}
The circle $S^{1}=\left\{z \in \mathbb{C}:|z|^{2}=1\right\}$ is a 1 -dimensional (abelian) Lie group under complex multiplication. It is also traditional to denote this group by $U(1)$.
\end{EXA}

Example 5.4. 


\begin{EXA}{LIE-L09-04-04}{$\R^*$ als Lie Gruppen}
Let $\mathbb{R}^{*}=\mathbb{R} \backslash\{0\}, \mathbb{C}^{*}=\mathbb{C} \backslash\{0\}$ and $\mathbb{H}^{*}=\mathbb{H} \backslash\{0\}$ (here $\mathbb{H}$ is the quaternion division ring discussed in detail later). Then, using multiplication, $\mathbb{R}^{*}, \mathbb{C}^{*}$, and $\mathbb{H}^{*}$ are Lie groups. The Lie group $\mathbb{H}^{*}$ is not abelian.
\end{EXA}



\begin{EXA}{LIE-L09-04-05}{$\GL(n,\R \C)$ als Lie Gruppen}
The group of all invertible real $n \times n$ matrices is a Lie group denoted $\mathrm{GL}(n, \mathbb{R})$. A global chart on $\mathrm{GL}(n, \mathbb{R})$ is given by the $n^{2}$ functions $x_{j}^{i}$, where if $A \in \operatorname{GL}(n, \mathbb{R})$ then $x_{j}^{i}(A)$ is the $i j$-th entry of $A$. Showing that $\operatorname{GL}(n, \mathbb{R})$ is a Lie group is straightforward. Multiplication is clearly smooth. For the inversion map one appeals to the usual formula for the inverse of a matrix, $A^{-1}= \operatorname{adj}(A) / \operatorname{det}(A)$. Here $\operatorname{adj}(A)$ is the adjoint matrix (whose entries are the cofactors). This shows that $A^{-1}$ depends smoothly on the entries of $A$. Similarly, the group $\operatorname{GL}(n, \mathbb{C})$ of invertible $n \times n$ complex matrices is a Lie group.
\end{EXA}

Exercise 5.5. 

Let $H$ be a subgroup of $G$ and consider the cosets $g H, g \in G$. Recall that $G$ is the disjoint union of the cosets of $H$. Show that if $H$ is open, then so are all the cosets. Conclude that the complement $H^{c}$ is also open and hence $H$ is closed.

Theorem 5.6. 

\begin{THEO}{LIE-L09-04-06}{Lokale Umgebung um Einheitselement von ZSH Lie Gruppe generiert die Lie Gruppe}
If $G$ is a connected Lie group and $U$ is a neighborhood of the identity element $e$, then $U$ generates the group. In other words, every element of $g$ is a product of elements of $U$.
\end{THEO}

Proof. 


\begin{PROOF}{LIE-L09-04-07}{P: Lokale Umgebung um Einheitselement von ZSH Lie Gruppe generiert die Lie Gruppe}
First note that $V=\operatorname{inv}(U) \cap U$ is an open neighborhood of the identity with the property that $\operatorname{inv}(V)=V$. We say that $V$ is symmetric. We show that $V$ generates $G$. For any open $W_{1}$ and $W_{2}$ in $G$, the set $W_{1} W_{2}=\left\{w_{1} w_{2}: w_{1} \in W_{1}\right.$ and $\left.w_{2} \in W_{2}\right\}$ is an open set being a union of the open sets $\bigcup_{g \in W_{1}} g W_{2}$. Thus, in particular, the inductively defined sets
$$
V^{n}=V V^{n-1}, \quad n=1,2,3, \ldots
$$
are open. We have
$$
e \in V \subset V^{2} \subset \cdots V^{n} \subset \cdots
$$
It is easy to check that each $V^{n}$ is symmetric and so also is the union
$$
V^{\infty}:=\bigcup_{n=1}^{\infty} V^{n}
$$
Moreover, $V^{\infty}$ is not only closed under inversion, but also obviously closed under multiplication. Thus $V^{\infty}$ is an open subgroup. From Exercise 5.5, $V^{\infty}$ is also closed, and since $G$ is connected, we obtain $V^{\infty}=G$.
\end{PROOF}

In general, the connected component of a Lie group $G$ that contains the identity is a Lie group denoted $G_{0}$, and it is generated by any open neighborhood of the identity. We call $G_{0}$ the identity component of $G$.

Definition 5.7. 

\begin{DEF}{LIE-L09-04-08}{Links / Rechts Translation}
For a Lie group $G$ and a fixed element $g \in G$, the maps $L_{g}: G \rightarrow G$ and $R_{g}: G \rightarrow G$ are defined by
$$
\begin{aligned}
L_{g} x & =g x \text { for } x \in G, \\
R_{g} x & =x g \text { for } x \in G,
\end{aligned}
$$
and are called left translation and right translation (by $g$ ) respectively.

It is easy to see that $L_{g}$ and $R_{g}$ are diffeomorphisms with $L_{g}^{-1}=L_{g^{-1}}$ and $R_{g}^{-1}=R_{g^{-1}}$.
\end{DEF}



\begin{DEF}{LIE-L09-04-09}{Produkt Lie Gruppe}
If $G$ and $H$ are Lie groups, then so is the product manifold $G \times H$, where multiplication is $\left(g_{1}, h_{1}\right) \cdot\left(g_{2}, h_{2}\right)=\left(g_{1} g_{2}, h_{1} h_{2}\right)$. The Lie group $G \times H$ is called the product Lie group.
\end{DEF}


\begin{EXA}{LIE-L09-04-10}{$n$-Torus Gruppe}
For example, the product group $S^{1} \times S^{1}$ is called the 2 -torus group. More generally, the higher torus groups are defined by $T^{n}=S^{1} \times \cdots \times S^{1}$ ( $n$ factors).
\end{EXA}

Definition 5.8. 


\begin{DEF}{LIE-L09-04-11}{Lie Untergruppe}
Let $H$ be an abstract subgroup of a Lie group $G$. If $H$ is a Lie group such that the inclusion map $H \hookrightarrow G$ is an immersion, then we say that $H$ is a Lie subgroup of $G$.
\end{DEF}

Proposition 5.9. 

\begin{PROP}{LIE-L09-04-12}{Charakterisierung von Lie Untergruppen durch reguläre Untermannigfaltigkeiit}
If $H$ is an abstract subgroup of a Lie group $G$ that is also a regular submanifold, then $H$ is a closed Lie subgroup.
\end{PROP}

Proof. 

\begin{PROOF}{LIE-L09-04-13}{P: Charakterisierung von Lie Untergruppen durch reguläre Untermannigfaltigkeiit}
The multiplication and inversion maps, $H \times H \rightarrow H$ and $H \rightarrow H$, are the restrictions of the multiplication and inversion maps on $G$, and since $H$ is a regular submanifold, we obtain the needed smoothness of these maps. The harder part is to show that $H$ is closed. So let $x_{0} \in \bar{H}$ be arbitrary. Let ( $U, \mathrm{x}$ ) be a single-slice chart adapted to $H$ whose domain contains $e$. Let $\delta: G \times G \rightarrow G$ be the map $\delta\left(g_{1}, g_{2}\right)=g_{1}^{-1} g_{2}$, and choose an open set $V$ such that $e \in V \subset \bar{V} \subset U$. By continuity of the map $\delta$ we can find an open neighborhood $O$ of the identity element such that $O \times O \subset \delta^{-1}(V)$. Now if $\left\{h_{i}\right\}$ is a sequence in $H$ converging to $x_{0} \in \bar{H}$, then $x_{0}^{-1} h_{i} \rightarrow e$ and $x_{0}^{-1} h_{i} \in O$ for all sufficiently large $i$. Since $h_{j}^{-1} h_{i}=\left(x_{0}^{-1} h_{j}\right)^{-1} x_{0}^{-1} h_{i}$, we have that $h_{j}^{-1} h_{i} \in V$ for sufficiently large $i, j$. For any sufficiently large fixed $j$, we have
$$
\lim _{i \rightarrow \infty} h_{j}^{-1} h_{i}=h_{j}^{-1} x_{0} \in \bar{V} \subset U
$$
Since $U$ is the domain of a single-slice chart, $U \cap H$ is closed in $U$. Thus since each $h_{j}^{-1} h_{i}$ is in $U \cap H$, we see that $h_{j}^{-1} x_{0} \in U \cap H \subset H$ for all sufficiently large $j$. This shows that $x_{0} \in H$, and since $x_{0}$ was arbitrary, we are done.
\end{PROOF}


\begin{DEF}{LIE-L09-04-14}{Abgeschlossene Lie Gruppe}
By a closed Lie subgroup we shall always mean one that is a regular submanifold as in the previous theorem. It is a nontrivial fact that an abstract subgroup of a Lie group that is also a closed subset is automatically a closed Lie subgroup in this sense (see Theorem 5.81).
\end{DEF}

Example 5.10. 

\begin{EXA}{LIE-L09-04-15}{$S^1$ als Einbettung in $2$-Torus ist abgeschlossene Lie Gruppe}
$S^{1}$ embedded as $S^{1} \times\{1\}$ in the torus $S^{1} \times S^{1}$ is a closed subgroup.
\end{EXA}

Example 5.11. 

\begin{EXA}{LIE-L09-04-16}{$S^1$ als Einbettung oder dicht in $2$-Torus ist abgeschlossene Lie Gruppe}
Let $S^{1}$ be considered as the set of unit modulus complex numbers. The image in the torus $T^{2}=S^{1} \times S^{1}$ of the map $\mathbb{R}^{1} \rightarrow S^{1} \times S^{1}$ given by $t \mapsto\left(e^{i 2 \pi t}, e^{i 2 \pi a t}\right)$ is a Lie subgroup. This map is a homomorphism. If $a$ is a rational number, then the image is an embedded copy of $S^{1}$ wrapped around the torus several times depending on $a$. If $a$ is irrational, then the image is still a Lie subgroup but is now dense in $T^{2}$.
\end{EXA}

The last example is important since it shows that a Lie subgroup might actually be a dense subset of the containing Lie group.

\subsection*{Linear Lie Groups}


\begin{EXA}{LIE-L09-04-17}{Raum der Endomorphismen von Vektorräumen als glatte Mannigfaltigkeit}
Let V be an $n$-dimensional vector space over $\mathbb{F}$, where $\mathbb{F}=\mathbb{R}$ or $\mathbb{C}$. The space $L(\mathrm{V}, \mathrm{V})$ of linear maps from V to V is a vector space and therefore a smooth manifold. A global chart for $L(\mathrm{V}, \mathrm{V})$ may be obtained by first choosing a basis for V and then defining $n^{2}$ functions $\left\{x_{j}^{i}\right\}_{1 \leq i . j \leq n}$ by the rule that if $A \in L(\mathrm{V}, \mathrm{V})$, then $x_{j}^{i}(A)$ is the $i j$-th entry of the matrix that represents $A$ with respect to the chosen basis. If the field is $\mathbb{R}$, then these are the coordinate functions of a global chart. If the field is $\mathbb{C}$, then we simply take the real and imaginary parts of the $x_{j}^{i}$ and thereby obtain $2 n^{2}$ coordinate functions. The various choices of basis give compatible charts, and the reader may check that these are charts from the smooth structure that $L(\mathrm{V}, \mathrm{V})$ has by virtue of being a finite-dimensional vector space.
\end{EXA}

The determinant of an element $A \in L(\mathrm{V}, \mathrm{V})$ is given as the determinant of any matrix which represents $A$ with respect to some basis. 

\begin{EXA}{LIE-L09-04-18}{$\GL(V)$ als glatte offene Untermannigfaltigkeit von $E(V,V)$E}
The group $\mathrm{GL}(\mathrm{V})$ of all linear automorphisms of V is an open submanifold of $L(\mathrm{V}, \mathrm{V})$ given by the condition of nonvanishing determinant, and the restrictions of the coordinate functions just introduced provide a global chart for $\mathrm{GL}(\mathrm{V})$.
\end{EXA}


\begin{EXA}{LIE-L09-04-19}{$\GL(n,\mathbb{F})$ als Gruppe der invertierbaren Matrizen}
Let $\operatorname{GL}(n, \mathbb{F})$ denote the group of invertible matrices with entries from $\mathbb{F}$.
\end{EXA}

\begin{PROP}{LIE-L09-04-20}{$\GL(V)$ ist Lie Gruppe}
We obtain an isomorphism of $\operatorname{GL}(\mathrm{V})$ with the matrix group $\operatorname{GL}(n, \mathbb{F})$ by choosing a basis and then simply sending each element of $\mathrm{GL}(\mathrm{V})$ to its matrix representative with respect to that basis. If we write $\operatorname{mat}(A)$ for the matrix that represents $A \in \mathrm{GL}(\mathrm{V})$ with respect to a fixed basis, then $A \rightarrow \operatorname{mat}(A)$ is a group isomorphism and is clearly smooth. It follows that $\mathrm{GL}(\mathrm{V})$ is a Lie group, and for each choice of basis we have an isomorphism of Lie groups $\mathrm{GL}(\mathrm{V}) \cong \mathrm{GL}(n, \mathbb{F})$. In practice it is common to work with the matrix group $\mathrm{GL}(n, \mathbb{F})$ and its subgroups.
\end{PROP}



\begin{DEF}{LIE-L09-04-21}{Allgemeine lineare Gruppe}
The Lie group $\mathrm{GL}(\mathrm{V})$ is called the general linear group of V and is also denoted $\mathrm{GL}(\mathrm{V}, \mathbb{F})$ when we want to make the field apparent. The matrix group $\operatorname{GL}(n, \mathbb{F})$ is also referred to as a general linear (matrix) group and is often identified with $\operatorname{GL}\left(\mathbb{F}^{n}\right)$.


More specifically, $\operatorname{GL}(n, \mathbb{R})$ is called the real general linear group, and $\operatorname{GL}(n, \mathbb{C})$ is called the complex general linear group. Lie groups that are subgroups of GL(V) for some vector space V are referred to as linear Lie groups and are often realized as matrix subgroups of $\operatorname{GL}(n, \mathbb{F})$ for some $n$.
\end{DEF}

Definition 5.12. 

\begin{DEF}{LIE-L09-04-22}{Spezielle Lineare Gruppe $\SL(V)$}
Let V be an $n$-dimensional vector space over the field $\mathbb{F}$ which we take to be either $\mathbb{R}$ or $\mathbb{C}$. Then the group SL(V) defined by
$$
\mathrm{SL}(\mathrm{V})=\mathrm{SL}(\mathrm{V}, \mathbb{F}):=\{A \in \mathrm{GL}(\mathrm{V}): \operatorname{det}(A)=1\}
$$
is called the special linear group for V.
\end{DEF}


\begin{DEF}{LIE-L09-04-23}{Nicht degenerierte Bilinearformen}
A bilinear form $\beta: \mathrm{V} \times \mathrm{V} \rightarrow \mathbb{F}$ on an $\mathbb{F}$-vector space V is called nondegenerate if the maps $\mathrm{V} \rightarrow \mathrm{V}^{*}$ given by $\beta_{R}: v \mapsto \beta(v, \cdot)$ and $\beta_{L}: v \mapsto \beta(\cdot, v)$, are both linear isomorphisms. If V is finite-dimensional, then $\beta_{R}$ is an isomorphism if and only if $\beta_{L}$ is an isomorphism and then $\beta$ is nondegenerate provided it has the property that if $\beta(v, w)=0$ for all $w \in V$, then $v=0$.
\end{DEF}

Definition 5.13. 

\begin{DEF}{LIE-L09-04-24}{Reelles Skalarprodukt als symmetrische nicht-degenerierte BLF}
A (real) scalar product on a (real) finite-dimensional vector space V is a nondegenerate symmetric bilinear form $\beta: \mathrm{V} \times \mathrm{V} \rightarrow \mathbb{R}$. A (real) scalar product space is a pair ( $\mathrm{V}, \beta$ ) where V is a real vector space and $\beta$ is a scalar product. (As usual we refer to V itself as the scalar product space when $\beta$ is given)
\end{DEF}

Definition 5.14. 


\begin{DEF}{LIE-L09-04-25}{Sesquilineare und Hermitische Form}
Let V be a complex vector space. An $\mathbb{R}$-bilinear map $\beta$ : $\mathrm{V} \times \mathrm{V} \rightarrow \mathbb{C}$ that satisfies $\beta(a v, w)=\bar{a} \beta(v, w)$ and $\beta(v, a w)=a \beta(v, w)$ for all $a \in \mathbb{C}$ and $v, w \in \mathrm{V}$ is called a sesquilinear form. If also $\beta(v, w)=\overline{\beta(w, v)}$, we call $\beta$ a Hermitian form.

The sesquilinear conditions imposed are described by saying that $\beta$ is to be conjugate linear in the first slot and linear in the second slot.
\end{DEF}



\begin{DEF}{LIE-L09-04-26}{Hermitisches Skalarprodukt als nicht degenerierte hermitische Form}
If a Hermitian form is nondegenerate, we call it a Hermitian scalar product, and then ( $\mathrm{V}, \beta$ ) a Hermitian scalar product space.

Nondegeneracy for a sesquilinear form is defined as for bilinear forms. Many authors define sesquilinear and Hermitian forms to be conjugate linear in the second slot and linear in the first. Obviously, $\beta(v, v)$ is always real for a Hermitian form.
\end{DEF}

Definition 5.15. 

\begin{DEF}{LIE-L09-04-27}{(Semi) Positiv/Negativ definiete Skalarprodukt}
Let $\beta$ be a (real) scalar product or Hermitian scalar product on a vector space V . Then
\begin{enumerate}
\item $\beta$ is positive (resp. negative) definite if $\beta(v, v) \geq 0$ (resp. $\beta(v, v) \leq 0)$ for all $v \in \mathrm{V}$ and $\beta(v, v)=0 \Longrightarrow v=0$;

\item $\beta$ is positive (resp. negative) semidefinite if $\beta(v, v) \geq 0$ (resp. $\beta(v, v) \leq 0)$ for all $v \in \mathrm{V}$.
\end{enumerate}
\end{DEF}



\begin{DEF}{LIE-L09-04-28}{Innere Produkt}
What is called an inner product (a term we have already used) is a positive definite scalar product (or positive definite Hermitian scalar product in the complex case). In this book the term "inner product" always implies positive definiteness.
\end{DEF} 

\begin{DEF}{LIE-L09-04-29}{Orthonormale Basis für indefinite Skalarprodukte}
Let us generalize the notion of orthonormal basis for an inner product space to include indefinite scalar product spaces. A basis $\left(e_{1}, \ldots, e_{n}\right)$ for a scalar product space $(\mathrm{V}, \beta)$ is called an orthonormal basis if $\beta\left(e_{i}, e_{j}\right)=0$ when $i \neq j$, and $\beta\left(e_{i}, e_{i}\right)= \pm 1$ for all $i$. An orthonormal basis always exists for a finite-dimensional scalar product space.
\end{DEF}

Definition 5.16. 


\begin{DEF}{LIE-L09-04-30}{$\operatorname{Aut}(V,\beta)$ Automorphismenuntergruppe bzgl bilinear / sesquilinear Formen}
Let V be an $n$-dimensional vector space over the field $\mathbb{F}$ which we take to be either $\mathbb{R}$ or $\mathbb{C}$. If $\beta$ is a bilinear form or sesquilinear form on V, then $\operatorname{Aut}(\mathrm{V}, \beta)$ is the subgroup of $\operatorname{GL}(\mathrm{V})$ defined by
$$
\operatorname{Aut}(\mathrm{V}, \beta):=\{A \in \mathrm{GL}(\mathrm{V}): \beta(A v, A w)=\beta(v, w) \text { for all } v, w \in \mathrm{V}\}
$$
If $\beta$ is a scalar product, then the elements of $\operatorname{Aut}(\mathrm{V}, \beta)$ are called isometries of V.
\end{DEF}



Theorem 5.17. 

\begin{THEO}{LIE-L09-04-31}{$\operatorname{Aut}(V,\beta)$ und $\operatorname{SAut}(V,\beta)$ als abgeschlossene Lie Untergruppen von $\GL(V)$}
$\mathrm{SL}(\mathrm{V})$ is a closed Lie subgroup of $\mathrm{GL}(\mathrm{V})$. If $\beta$ is a bilinear or sesquilinear form as above, then $\operatorname{Aut}(\mathrm{V}, \beta)$ and $\operatorname{SAut}(\mathrm{V}, \beta):= \operatorname{Aut}(\mathrm{V}, \beta) \cap \mathrm{SL}(\mathrm{V})$ are closed Lie subgroups of $\mathrm{GL}(\mathrm{V})$. (The form $\beta$ is most often taken to be nondegenerate.)
\end{THEO}

Proof. 

\begin{PROOF}{LIE-L09-04-32}{P: $\operatorname{Aut}(V,\beta)$ und $\operatorname{SAut}(V,\beta)$ als abgeschlossene Lie Untergruppen von $\GL(V)$}
It is easy to check that the sets in question are subgroups. They are clearly closed. For example, $\operatorname{Aut}(\mathrm{V}, \beta)=\bigcap_{v, w} F_{v, w}$, where
$$
F_{v, w}:=\{A \in \mathrm{GL}(\mathrm{V}): \beta(A v, A w)=\beta(v, w)\}
$$
The fact that they are Lie subgroups follows from Theorem 5.81 below. However, as we shall see, most of the specific cases arising from various choices of $\beta$ can be proved to be Lie groups by other means. That they are Lie subgroups follows from Proposition 5.9 once we show that they are regular submanifolds of the appropriate group $\mathrm{GL}(\mathrm{V}, \mathbb{F})$.
\end{PROOF}

We will return to this later, once we have introduced another powerful theorem that will allow us to verify this without the use of Theorem 5.81. \(\square\)

Let $\operatorname{dim} \mathrm{V}=n$. After choosing a basis, $\mathrm{SL}(\mathrm{V})$ gives the matrix version $\operatorname{SL}(n, \mathbb{F}):=\left\{A \in M_{n \times n}(\mathbb{F}): \operatorname{det} A=1\right\}$. Notice that even when $\mathbb{F}=\mathbb{C}$, it may be that $\beta$ is only required to be $\mathbb{R}$-linear. Depending on whether $\mathbb{F}=\mathbb{C}$ or $\mathbb{R}$ and on the nature of $\beta$, the notation for the linear groups takes on special conventional forms introduced below. When choosing a basis in order to represent one of the groups associated to a form $\beta$ in a matrix version, it is usually the case that one uses a basis under which the matrix that represents $\beta$ takes on a canonical form.

Example 5.18 (The (semi) orthogonal groups). 

\begin{EXA}{LIE-L09-04-33}{$\Ogroup(V;\beta)$, $\Ogroup(p,q)$ als Semiorthogonale Gruppe assoziiert zu reellem Skalarprodukt $\beta$}
Let ($\mathrm{V},\beta$) be a real scalar product space. In this case we write $\operatorname{Aut}(\mathrm{V}, \beta)$ as $\mathrm{O}(\mathrm{V}, \beta)$ and refer to it as the semiorthogonal group associated to $\beta$. With respect to an appropriately ordered orthonormal basis, $\beta$ is represented by a diagonal matrix of the form
$$
\eta_{p, q}=\left[\begin{array}{cccccc}
1 & 0 & \cdots & & \cdots & 0 \\
0 & \ddots & 0 & & & \vdots \\
\vdots & 0 & 1 & \ddots & & \\
& & \ddots & -1 & 0 & \vdots \\
\vdots & & & 0 & \ddots & 0 \\
0 & \cdots & & \cdots & 0 & -1
\end{array}\right],
$$
where there are $p$ ones and $q$ minus ones down the diagonal. The group of matrices arising from $\mathrm{O}(\mathrm{V}, \beta)$ with such a choice of basis is denoted $\mathrm{O}(p, q)$ and consists exactly of the real matrices $Q$ satisfying $Q \eta_{p, q} Q^{t}=\eta_{p, q}$. These groups are called the semiorthogonal matrix groups. With such an orthonormal choice of basis as above, the bilinear form (scalar product) is given as a canonical form on $\mathbb{R}^{n}$ where $(p+q=n)$ :
$$
\langle x, y\rangle:=\sum_{i=1}^{p} x^{i} y^{i}-\sum_{i=p+1}^{n} x^{i} y^{i},
$$
and we have the alternative description
$$
\mathrm{O}(p, q)=\left\{Q \in \mathrm{GL}(n):\langle Q x, Q y\rangle=\langle x, y\rangle \text { for all } x, y \in \mathbb{R}^{n}\right\} .
$$
\end{EXA}


\begin{EXA}{LIE-L09-04-34}{$\Ogroup(V;\beta)$, $\Ogroup(n)$ als Orthogonale Gruppe assoziiert zu reellem positiv definit Skalarprodukt $\beta$}
If $\beta$ is positive definite, we then have $q=0$, and $\mathrm{O}(\mathrm{V}, \beta)$ is referred to as a real orthogonal group. We write $\mathrm{O}(n, 0)$ as $\mathrm{O}(n)$ and refer to it as the real orthogonal (matrix) group; $Q \in \mathrm{O}(n) \Longleftrightarrow Q^{t} Q=I$.
\end{EXA}



Example 5.19. 

\begin{EXA}{LIE-L09-04-35}{Komplexe Orthogonale Gruppe}
There are also complex orthogonal groups (not to be confused with unitary groups). In matrix representation, we have $\mathrm{O}(n, \mathbb{C}):= \left\{Q \in \mathrm{GL}(n, \mathbb{C}): Q^{t} Q=I\right\}$.
\end{EXA}

Example 5.20. 


\begin{EXA}{LIE-L09-04-36}{$\U(V;\beta)$, $\U(p,q)$ als Semiunitäre Gruppe assoziiert zu hermitischen Skalarprodukt $\beta$}
Let $(\mathrm{V}, \beta)$ be a Hermitian scalar product space. In this case, we write $\operatorname{Aut}(\mathrm{V},\beta,\mathbb{C})$ as $U(\mathrm{V},\beta)$ and refer to it as the semiunitary group associated to $\beta$. If $\beta$ is positive definite, then we call it a unitary group. Again we may choose a basis for V such that $\beta$ is represented by the Hermitian form on $\mathbb{C}^{n}$ given by
$$
\langle x, y\rangle:=\sum_{i=1}^{p} \bar{x}^{i} y^{i}-\sum_{i=p+1}^{p+q=n} \bar{x}^{i} y^{i}
$$
We then obtain the semiunitary matrix group
$$
\mathrm{U}(p, q)=\left\{A \in \mathrm{GL}(n, \mathbb{C}):\langle A x, A y\rangle=\langle x, y\rangle \text { for all } x, y \in \mathbb{R}^{n}\right\}
$$
We write $\mathrm{U}(n, 0)$ as $\mathrm{U}(n)$ and refer to it as the unitary (matrix) group. In particular, $U(1)=S^{1}=\{z \in \mathbb{C}:|z|=1\}$.
\end{EXA}



Definition 5.21. 

\begin{EXA}{LIE-L09-04-37}{Spezielle Orthogonale Gruppe $\SO(V,\beta)$ zu reellem Skalarprodukt}
For $(\mathrm{V},\beta)$ a real scalar product space, we have the special orthogonal group of ($\mathrm{V},\beta$) given by
$$
\mathrm{SO}(\mathrm{V}, \beta)=\mathrm{O}(\mathrm{V}, \beta) \cap \mathrm{SL}(\mathrm{V}, \mathbb{R})
$$
\end{EXA}


\begin{EXA}{LIE-L09-04-38}{Spezielle Unitäre Gruppe $\SU(V,\beta)$ zu hermitischen Skalarprodukt}
For ($\mathrm{V},\beta$) a Hermitian scalar product space, we have the special unitary group of ($\mathrm{V},\beta$) given by
$$
\mathrm{SU}(\mathrm{V}, \beta)=\mathrm{U}(\mathrm{V}, \beta) \cap \mathrm{SL}(\mathrm{V}, \mathbb{C})
$$
\end{EXA}


Definition 5.22. 

\begin{EXA}{LIE-L09-04-39}{Spezielle Orthogonale / Unitäre Matrixgruppe in $\R,\C$}
The group of $n \times n$ complex matrices of determinant one is the complex special linear matrix group $\operatorname{SL}(n, \mathbb{C})$. We also have the similarly defined real special linear group $\operatorname{SL}(n, \mathbb{R})$. The special orthogonal and special semiorthogonal matrix groups, $\mathrm{SO}(n)$ and $\mathrm{SO}(p, q)$, are the matrix groups defined by $\mathrm{SO}(n)=\mathrm{O}(n) \cap \mathrm{SL}(n, \mathbb{R})$ and $\mathrm{SO}(p, q)=\mathrm{O}(p, q) \cap \mathrm{SL}(n, \mathbb{R})$. The special unitary and special semiunitary matrix groups $\mathrm{SU}(n)$ and $\mathrm{SU}(p, q)$ are defined similarly.
\end{EXA}

The group $\mathrm{SO}(3)$ is the familiar matrix representation of the proper rotation group of Euclidean space and plays a prominent role in classical physics. Here "proper" refers to the fact that $\mathrm{SO}(3)$ does not contain any reflections. In the problems we ask the reader to show that $\mathrm{SO}(3)$ is the connected component of the identity in $\mathrm{O}(3)$.

Exercise 5.23. 

\begin{EXE}{LIE-L09-04-40}{$\SU(2)$ ist EZSH und $\SO(3)$ nicht}
Show that $\mathrm{SU}(2)$ is simply connected while $\mathrm{SO}(3)$ is not.
\end{EXE}


Example 5.24 (Symplectic groups). 


\begin{EXA}{LIE-L09-04-41}{Symplektische Gruppe $\operatorname{Sp}(\mathrm{V}, \mathbb{F})$ assoziert zu nichtdegenerierte, schrägsymmetrische $\mathbb{F}$-bilineare Form auf $2n$-dim Vektorraum}
We will describe both the real and the complex symplectic groups. Suppose that $\beta$ is a nondegenerate skewsymmetric $\mathbb{C}$-bilinear (resp. $\mathbb{R}$-bilinear) form on a $2 n$-dimensional complex (resp. real) vector space $V$. The group $\operatorname{Aut}(V, \beta)$ is called the complex (resp. real) symplectic group and is denoted by $\operatorname{Sp}(\mathrm{V}, \mathbb{C})$ (resp. $\operatorname{Sp}(\mathrm{V}, \mathbb{R})$ ). There exists a basis $\left\{f_{i}\right\}$ for V such that $\beta$ is represented in the canonical form by
$$
(v, w)=\sum_{i=1}^{n} v^{i} w^{n+i}-\sum_{j=1}^{n} v^{n+j} w^{j} .
$$
The symplectic matrix groups are given by
$$
\begin{aligned}
& \operatorname{Sp}(2 n, \mathbb{C}):=\left\{A \in M_{2 n \times 2 n}(\mathbb{C}):(A v, A w)=(v, w)\right\}, \\
& \operatorname{Sp}(2 n, \mathbb{R}):=\left\{A \in M_{2 n \times 2 n}(\mathbb{R}):(A v, A w)=(v, w)\right\},
\end{aligned}
$$
where ($v,w$) is given as above.
\end{EXA}


Exercise 5.25. 

For $\mathbb{F}=\mathbb{C}$ or $\mathbb{R}$, show that $A \in \operatorname{Sp}(2 n, \mathbb{F})$ if and only if $A^{t} J A=J$, where
$$
J=\left(\begin{array}{cc}
0 & I \\
-I & 0
\end{array}\right)
$$



Much of the above can be generalized somewhat more. Recall that the algebra of quaternions $\mathbb{H}$ is a copy of $\mathbb{R}^{4}$ endowed with a multiplication described as follows: First let a generic elements of $\mathbb{R}^{4}$ be denoted by $x= \left(x^{0}, x^{1}, x^{2}, x^{3}\right), y=\left(y^{0}, y^{1}, y^{2}, y^{3}\right)$, etc. Thus we are using $\{0,1,2,3\}$ as our index set. Let the standard basis be denoted by $\mathbf{1}, \mathbf{i}, \mathbf{j}, \mathbf{k}$. We define a multiplication by taking these basis elements as generators and insisting on the following relations:
$$
\begin{aligned}
\mathbf{i}^{2} & =\mathbf{j}^{2}=\mathbf{k}^{2}=-\mathbf{1}, \\
\mathbf{i} \mathbf{j} & =-\mathbf{j} \mathbf{i}=\mathbf{k}, \\
\mathbf{j} \mathbf{k} & =-\mathbf{k} \mathbf{j}=\mathbf{i}, \\
\mathbf{k i} & =-\mathbf{i} \mathbf{k}=\mathbf{j} .
\end{aligned}
$$
Of course, $\mathbb{H}$ is a vector space over $\mathbb{R}$ since it is just $\mathbb{R}^{4}$ with some extra structure. As a ring, $\mathbb{H}$ is a division algebra which is very much like a field, lacking only the property of commutativity. In particular, we shall see that every nonzero element of $\mathbb{H}$ has a multiplicative inverse. Elements of the form $a \mathbf{1}$ for $a \in \mathbb{R}$ are identified with the corresponding real numbers, and such quaternions are called real quaternions. By analogy with complex numbers, quaternions of the form $x^{1} \mathbf{i}+x^{2} \mathbf{j}+x^{3} \mathbf{k}$ are called imaginary quaternions. For a given quaternion $x=x^{0} \mathbf{1}+x^{1} \mathbf{i}+x^{2} \mathbf{j}+x^{3} \mathbf{k}$, the quaternion $x^{1} \mathbf{i}+x^{2} \mathbf{j}+x^{3} \mathbf{k}$ is called the imaginary part of $x$, and $x^{0} \mathbf{1}=x^{0}$ is called the real part of $x$. We also have a conjugation defined by
$$
x \mapsto \bar{x}:=x^{0} \mathbf{1}-x^{1} \mathbf{i}-x^{2} \mathbf{j}-x^{3} \mathbf{k} .
$$

Notice that $x \bar{x}=\bar{x} x$ is real and equal to $\left(x^{0}\right)^{2}+\left(x^{1}\right)^{2}+\left(x^{2}\right)^{2}+\left(x^{3}\right)^{2}$. We denote the positive square root of this by $|x|$ so that $\bar{x} x=|x|^{2}$.



Exercise 5.26. Verify the following for $x, y \in \mathbb{H}$ and $a, b \in \mathbb{R}$ :
$$
\begin{array}{cc}
\overline{a x+b y}=a \bar{x}+b \bar{y}, & \overline{(\bar{x})}=x \\
|x y|=|x||y|, & |\bar{x}|=|x| \\
& \overline{x y}=\bar{y} \bar{x}
\end{array}
$$
Now we can write down the inverse of a nonzero $x \in \mathbb{H}$ :
$$
x^{-1}=\frac{1}{|x|^{2}} \bar{x}
$$
Notice the strong analogy with complex number arithmetic.


Example 5.27. 

The set of unit quaternions is $U(1, \mathbb{H}):=\{|x|=1\}$. This set is closed under multiplication. As a manifold it is (diffeomorphic to) $S^{3}$. With quaternionic multiplication, $S^{3}=U(1, \mathbb{H})$ is a compact Lie group. Compare this to Example 5.3 where we saw that $U(1, \mathbb{C})=S^{1}$. For the future, we unify things by letting $U(1, \mathbb{R}):=\mathbb{Z}_{2}=S^{0} \subset \mathbb{R}$. In other words, we take the 0 -sphere to be the subset $\{-1,1\}$ with its natural structure as a multiplicative group.
$$
\begin{gathered}
U(1, \mathbb{H})=S^{3} \\
U(1, \mathbb{C})=S^{1} \\
U(1, \mathbb{R}):=\mathbb{Z}_{2}=S^{0}
\end{gathered}
$$


Exercise 5.28. 

Prove the assertions in the last example.


We now consider the $n$-fold product $\mathbb{H}^{n}$, which, as a real vector space (and a smooth manifold), is $\mathbb{R}^{4 n}$. However, let us think of elements of $\mathbb{H}^{n}$ as column vectors with quaternion entries. We want to treat $\mathbb{H}^{n}$ as a vector space over $\mathbb{H}$ with addition defined just as for $\mathbb{R}^{n}$ and $\mathbb{C}^{n}$, but since $\mathbb{H}$ is not commutative, we are not properly dealing with a vector space. In particular, we should decide whether scalars should multiply column vectors on the right or on the left. We choose to multiply on the right, and this could take some getting used to, but there is a good reason for our choice. This puts us into the category of right $\mathbb{H}$-modules were elements of $\mathbb{H}$ are the "scalars". The reader should have no trouble catching on, and so we do not make formal definitions at this time (but see Appendix D). For $v, w \in \mathbb{H}^{n}$ and $a, b \in \mathbb{H}$, we have
$$
\begin{aligned}
v(a+b) & =v a+v b \\
(v+w) a & =v a+w a \\
(v a) b & =v(a b)
\end{aligned}
$$
A map $A: \mathbb{H}^{n} \rightarrow \mathbb{H}^{n}$ is said to be $\mathbb{H}$-linear if $A(v a)=A(v) a$ for all $v \in \mathbb{H}^{n}$ and $a \in \mathbb{H}$. There is no problem with doing matrix algebra with matrices with quaternion entries, as long as one respects the noncommutativity of $\mathbb{H}$. For example, if $A=\left(a_{j}^{i}\right)$ and $B=\left(b_{j}^{i}\right)$ are matrices with quaternion entries, then writing $C=A B$ we have
$$
c_{j}^{i}=\sum a_{k}^{i} b_{j}^{k},
$$
but we cannot expect that $\sum a_{k}^{i} b_{j}^{k}=\sum b_{j}^{k} a_{k}^{i}$. For any $A=\left(a_{j}^{i}\right)$, the map $\mathbb{H}^{n} \rightarrow \mathbb{H}^{n}$ defined by $v \mapsto A v$ is $\mathbb{H}$-linear since $A(v a)=(A v) a$.



Definition 5.29. 

The set of all $m \times n$ matrices with quaternion entries is denoted $M_{m \times n}(\mathbb{H})$. The subset $\operatorname{GL}(n, \mathbb{H})$ is defined as the set of all $Q \in M_{m \times n}(\mathbb{H})$ such that the map $v \mapsto Q v$ is a bijection.

We will now see that $\operatorname{GL}(n, \mathbb{H})$ is a Lie group isomorphic to a subgroup of $\operatorname{GL}(2 n, \mathbb{C})$. First we define a map $\iota: \mathbb{C}^{2} \rightarrow \mathbb{H}$ as follows: For $\left(z_{1}, z_{2}\right) \in \mathbb{C}$ with $z_{1}=x^{0}+x^{1} \mathbf{i}$ and $z_{2}=x^{2}+x^{3} \mathbf{i}$, we let $\iota\left(z^{1}, z^{2}\right)=\left(x^{0}+x^{1} \mathbf{i}\right)+\left(x^{2}+x^{3} \mathbf{i}\right) \mathbf{j}$ where on the right hand side we interpret $\mathbf{i}$ as a quaternion. Note that $\left(x^{0}+x^{1} \mathbf{i}\right)+\left(x^{2}+x^{3} \mathbf{i}\right) \mathbf{j}=x^{0}+x^{1} \mathbf{i}+x^{2} \mathbf{j}+x^{3} \mathbf{k}$. It is easily shown that this map is an $\mathbb{R}$-linear bijection, and we use this map to identify $\mathbb{C}^{2}$ with $\mathbb{H}$. Another way of looking at this is that we identify $\mathbb{C}$ with the span of 1 and $\mathbf{i}$ in $\mathbb{H}$ and then every quaternion has a unique representation as $z^{1}+z^{2} \mathbf{j}$ for $z^{1}, z^{2} \in \mathbb{C} \subset \mathbb{H}$. We extend this idea to square quaternionic matrices; we can write every $Q \in M_{m \times n}(\mathbb{H})$ in the form $A+B \mathbf{j}$ for $A, B \in M_{m \times n}(\mathbb{C})$ in a unique way. This representation makes it clear that $M_{m \times n}(\mathbb{H})$ has a natural complex vector space structure, where the scalar multiplication is $z(A+B \mathbf{j})=z A+z B \mathbf{j}$. Direct computation shows that
$$
(A+B \mathbf{j})(C+D \mathbf{j})=(A C-B \bar{D})+(A D+B \bar{C}) \mathbf{j}
$$
for $A+B \mathbf{j} \in M_{m \times n}(\mathbb{H})$ and $C+D \mathbf{j} \in M_{n \times k}(\mathbb{H})$, where we have used the fact that for $Q \in M_{m \times n}(\mathbb{C})$ we have $Q \mathbf{j}=\mathbf{j} \bar{Q}$. From this it is not hard to show that the map $\vartheta_{m \times n}: M_{m \times n}(\mathbb{H}) \rightarrow M_{2 m \times 2 n}(\mathbb{C})$ given by
$$
\vartheta_{m \times n}: A+B \mathbf{j} \longmapsto\left(\begin{array}{cc}
A & B \\
-\bar{B} & \bar{A}
\end{array}\right)
$$
is an injective $\mathbb{R}$-linear map which respects matrix multiplication and thus is an $\mathbb{R}$-algebra isomorphism onto its image. We may identify $M_{m \times n}(\mathbb{H})$ with the subspace of $M_{2 m \times 2 n}(\mathbb{C})$ consisting of all matrices of the form $\left(\begin{array}{cc}A & B \\ -\bar{B} & \bar{A}\end{array}\right)$, where $A, B \in \mathbb{C}^{m \times n}$. In particular, if $m=n$, then we obtain an injective $\mathbb{R}$ linear algebra homomorphism $\vartheta_{n \times n}: M_{n \times n}(\mathbb{H}) \rightarrow M_{2 n \times 2 n}(\mathbb{C})$, and thus the image of this map in $M_{2 n \times 2 n}(\mathbb{C})$ is another realization of the matrix algebra $M_{n \times n}(\mathbb{H})$. If we specialize to the case of $n=1$, we get a realization of $\mathbb{H}$ as the set of all $2 \times 2$ complex matrices of the form $\left(-\bar{w} \frac{w}{\bar{z}}\right)$. This set of matrices is closed under multiplication and forms an algebra over the field $\mathbb{R}$. Let us denote this algebra of matrices by the symbol $\mathcal{R}^{4}$ since it is diffeomorphic to $\mathbb{H} \cong \mathbb{R}^{4}$. We now have an algebra isomorphism $\vartheta: \mathbb{H} \rightarrow \mathcal{R}^{4}$ under which the quaternions $\mathbf{1}, \mathbf{i}, \mathbf{j}$ and $\mathbf{k}$ correspond to the matrices
$$
\left(\begin{array}{ll}
1 & 0 \\
0 & 1
\end{array}\right),\left(\begin{array}{cc}
i & 0 \\
0 & -i
\end{array}\right),\left(\begin{array}{cc}
0 & 1 \\
-1 & 0
\end{array}\right) \text { and }\left(\begin{array}{cc}
0 & i \\
i & 0
\end{array}\right)
$$
respectively. Since $\mathbb{H}$ is a division algebra, each of its nonzero elements has a multiplicative inverse. Thus $\mathcal{R}^{4}$ must contain the matrix inverse of each of its nonzero elements. This can be seen directly:
$$
\left(\begin{array}{cc}
z & w \\
-\bar{w} & \bar{z}
\end{array}\right)^{-1}=\frac{1}{|z|^{2}+|w|^{2}}\left(\begin{array}{cc}
\bar{z} & -w \\
\bar{w} & z
\end{array}\right) .
$$
Consider again the group of unit quaternions $U(1, \mathbb{H})$. We have already seen that as a smooth manifold, $U(1, \mathbb{H})$ is $S^{3}$. However, under the isomorphism $\mathbb{H} \rightarrow \mathcal{R}^{4} \subset M_{2 \times 2}(\mathbb{C})$ just mentioned, $U(1, \mathbb{H})$ manifests itself as $\mathrm{SU}(2)$. Thus we obtain a smooth map $U(1, \mathbb{H}) \rightarrow \mathrm{SU}(2)$ that is a group isomorphism. We record this as a proposition:


Proposition 5.30. 

The map $U(1, \mathbb{H}) \rightarrow \mathrm{SU}(2)$ given by
$$
x=z+w \mathbf{j} \mapsto\left(\begin{array}{cc}
z & w \\
-\bar{w} & \bar{z}
\end{array}\right)
$$
where $x=x^{0}+x^{1} \mathbf{i}+x^{2} \mathbf{j}+x^{3} \mathbf{k}, z=x^{0}+x^{1} \mathbf{i}$ and $w=x^{2}+x^{3} \mathbf{i}$, is a group isomorphism. Thus $S^{3}=U(1, \mathbb{H}) \cong \mathrm{SU}(2)$.


Proof. 

The first equality has already been established. Notice that we then have
$$
|x|^{2}=|z|^{2}+|w|^{2}=\operatorname{det}\left(\begin{array}{cc}
z & w \\
-\bar{w} & \bar{z}
\end{array}\right)
$$
and so $x \in U(1, \mathbb{H})$ if and only if $(-\bar{w} \underset{\bar{z}}{w})$ has determinant one. But such matrices account for all elements of $\mathrm{SU}(2)$ (verify this). We leave it to the reader to check that the map $U(1, \mathbb{H}) \rightarrow \mathrm{SU}(2)$ is a group isomorphism.



Exercise 5.31. 

Show that $Q \in \mathrm{GL}(n, \mathbb{H})$ if and only if $\operatorname{det}\left(\vartheta_{n \times n}(Q)\right) \neq 0$.

The set of all elements of $\operatorname{GL}(2 n, \mathbb{C})$ which are of the form $\left(\begin{array}{cc}A & B \\ - & \bar{B}\end{array}\right)$ is a subgroup of $\mathrm{GL}(2 n, \mathbb{C})$ and in fact a Lie group. Using the last exercise, we see that we may identify $\operatorname{GL}(n, \mathbb{H})$ as a Lie group with this subgroup of $\mathrm{GL}(2 n, \mathbb{C})$. We want to find a quaternionic analogue of $\mathrm{U}(n, \mathbb{C})$, and so we define $b: \mathbb{H}^{n} \times \mathbb{H}^{n} \rightarrow \mathbb{H}$ by
$$
b(v, w)=\bar{v}^{t} w
$$
Explicitly, if
$$
v=\left[\begin{array}{c}
v^{1} \\
\vdots \\
v^{n}
\end{array}\right] \text { and } w=\left[\begin{array}{c}
w^{1} \\
\vdots \\
w^{n}
\end{array}\right],
$$
then
$$
b(v, w)=\left[\begin{array}{lll}
v^{1} & \cdots & v^{n}
\end{array}\right]\left[\begin{array}{c}
w^{1} \\
\vdots \\
w^{n}
\end{array}\right]=\sum \bar{v}^{i} w^{i} .
$$
Note that $b$ is obviously $\mathbb{R}$-bilinear. But if $a \in \mathbb{H}$, then we have $b(v a, w)= b(v, w) \bar{a}$ and $b(v, w a)=b(v, w) a$. Notice that we consistently use right multiplication by quaternionic scalars. Thus $b$ is the quaternionic analogue of an Hermitian scalar product.



Definition 5.32. 

We define $\mathrm{U}(n, \mathbb{H})$ :
$$
\mathrm{U}(n, \mathbb{H}):=\left\{Q \in \mathrm{GL}(n, \mathbb{H}): b(Q v, Q w)=b(v, w) \text { for all } v, w \in \mathbb{H}^{n}\right\}
$$
$\mathrm{U}(n, \mathbb{H})$ is called the quaternionic unitary group.


The group $\mathrm{U}(n, \mathbb{H})$ is sometimes called the symplectic group and is denoted $\operatorname{Sp}(n)$, but we will avoid this since we want no confusion with the symplectic groups we have already defined. The group $\mathrm{U}(n, \mathbb{H})$ is in fact a Lie group (Theorem 5.17 generalizes to the quaternionic setting). The image of $\mathrm{U}(n, \mathbb{H})$ in $M_{2 n \times 2 n}(\mathbb{C})$ under the map $\vartheta_{n \times n}$ is denoted $\operatorname{USp}(2 n, \mathbb{C})$. Since it is easily established that $\left.\vartheta_{n \times n}\right|_{\mathrm{U}(n, \mathbb{H})}$ is a group homomorphism, the image $\operatorname{USp}(2 n, \mathbb{C})$ is a subgroup of $\operatorname{GL}(2 n, \mathbb{C})$.


Exercise 5.33. 

Show that $\vartheta_{n \times n}\left(\bar{A}^{t}\right)=\left(\overline{\vartheta_{n \times n}(A)}\right)^{t}$. Show that $\operatorname{USp}(2 n, \mathbb{C})$ is a Lie subgroup of $\operatorname{GL}(2 n, \mathbb{C})$.

Exercise 5.34. 

Show that $\operatorname{USp}(2 n, \mathbb{C})=\mathrm{U}(2 n) \cap \operatorname{Sp}(2 n, \mathbb{C})$. Hint: Show that
$$
\vartheta_{n \times n}\left(M_{n \times n}(\mathbb{H})\right)=\left\{A \in \mathrm{GL}(2 n, \mathbb{C}): J A J^{-1}=\bar{A}\right\}
$$
where $J=\left(\begin{array}{cc}0 & \mathrm{id} \\ -\mathrm{id} & 0\end{array}\right)$. Next show that if $A \in \mathrm{U}(2 n)$, then $J A J^{-1}=\bar{A}$ if and only if $A^{t} J A=J$.





\pagebreak

\subsection*{Lie Group Homomorphisms}


Definition 5.35. 


\begin{DEF}{LIE-L09-04-42}{Lie Gruppen Homomorphismus}
Let $G$ and $H$ be Lie groups. A smooth map $f: G \rightarrow H$ that is a group homomorphism is called a Lie group homomorphism. A Lie group homomorphism is called a Lie group isomorphism in case it has an inverse that is also a Lie group homomorphism. A Lie group isomorphism $G \rightarrow G$ is called a Lie group automorphism of $G$.

If $f: G \rightarrow H$ is a Lie group homomorphism, then by definition $f\left(g_{1} g_{2}\right)= f\left(g_{1}\right) f\left(g_{2}\right)$ for all $g_{1}, g_{2} \in G$, and it follows that $f(e)=e$ and also that $f\left(g^{-1}\right)=f(g)^{-1}$ for all $g \in G$.
\end{DEF}



Example 5.36. 

\begin{EXA}{LIE-L09-04-43}{$\mathrm{SO}(n, \mathbb{R}) \hookrightarrow \operatorname{GL}(n, \mathbb{R})$ ist ein LG Homomorphismus}
The inclusion $\mathrm{SO}(n, \mathbb{R}) \hookrightarrow \operatorname{GL}(n, \mathbb{R})$ is a Lie group homomorphism.
\end{EXA}



Example 5.37. 

\begin{EXA}{LIE-L09-04-44}{$S^1 \mathrm{SO}(n)$ ist ein LG Homomorphismus}
The circle $S^{1} \subset \mathbb{C}$ is a Lie group under complex multiplication and the map
$$
z=e^{i \theta} \mapsto\left[\begin{array}{ccc}
\cos (\theta) & \sin (\theta) & 0 \\
-\sin (\theta) & \cos (\theta) & 0 \\
0 & 0 & I_{n-2}
\end{array}\right]
$$
is a Lie group homomorphism of $S^{1}$ into $\operatorname{SO}(n)$.
\end{EXA}


Example 5.38. 

\begin{EXA}{LIE-L09-04-45}{$U(1, \mathbb{H}) \rightarrow \mathrm{SU}(2)$ ist ein LG Homomorphismus}
The map $U(1, \mathbb{H}) \rightarrow \mathrm{SU}(2)$ of Proposition 5.30 is a Lie group isomorphism.
\end{EXA}

Example 5.39. 

\begin{EXA}{LIE-L09-04-46}{Konjugationsabbildung ist Lie Gruppen Automorphismus}
The conjugation map $C_{g}: G \rightarrow G$ given by $x \mapsto g x g^{-1}$ is a Lie group automorphism. Note that $C_{g}=L_{g} \circ R_{g^{-1}}$.
\end{EXA}

Proposition 5.40. 

\begin{PROP}{LIE-L09-04-47}{Identität von LG Homomorphismen für Übereinstimmung lokal um Einheitselement für Definitionsbereich ZSH}
Let $f_{1}: G \rightarrow H$ and $f_{2}: G \rightarrow H$ be Lie group homomorphisms that agree in a neighborhood of the identity. If $G$ is connected, then $f_{1}=f_{2}$.
\end{PROP}

Proof. 

\begin{PROOF}{LIE-L09-04-48}{P: Identität von LG Homomorphismen für Übereinstimmung lokal um Einheitselement für Definitionsbereich ZSH}
By Theorem 5.6, any $g \in G$ is a product of elements in the set on which $f_{1}$ and $f_{2}$ agree, so the homomorphism property forces $f_{1}=f_{2}$.
\end{PROOF}

Exercise 5.41. 


\begin{EXE}{LIE-L09-04-49}{Tagentenabbildung für Multiplikationsabbildung}
Show that the multiplication map $\mu: G \times G \rightarrow G$ has tangent map at $(e, e) \in G \times G$ given as $T_{(e, e)} \mu(v, w)=v+w$. Recall that we identify $T_{(e, e)}(G \times G)$ with $T_{e} G \times T_{e} G$.
\end{EXE}

Exercise 5.42. 

\begin{EXE}{LIE-L09-04-50}{Tagentenabbildung für Konjugationsabbildung}
$\operatorname{GL}(n, \mathbb{R})$ is an open subset of the vector space of all $n \times n$ matrices $M_{n \times n}(\mathbb{R})$. Using the natural identification of $T_{e} \mathrm{GL}(n, \mathbb{R})$ with $M_{n \times n}(\mathbb{R})$, show that
$$
T_{e} C_{g}(x)=g x g^{-1}
$$
where $g \in \mathrm{GL}(n, \mathbb{R})$ and $x \in M_{n \times n}(\mathbb{R})$.
\end{EXE}


Example 5.43. 

\begin{EXA}{LIE-L09-04-51}{LG Homomorphismus zwischen $\R$ und $S^1$}
The map $t \mapsto e^{i t}$ is a Lie group homomorphism from $\mathbb{R}$ to $S^{1} \subset \mathbb{C}$.
\end{EXA}

Definition 5.44. 


\begin{DEF}{LIE-L09-04-52}{1-parameter Untergruppe}
A Lie group homomorphism from the additive group $\mathbb{R}$ into a Lie group is called a one-parameter subgroup. (Note that despite the use of the word "subgroup", a one-parameter subgroup is actually a map.)
\end{DEF}

Example 5.45. 

\begin{EXA}{LIE-L09-04-53}{1-parameter Untergruppe von $2$-Torus}
We have seen that the torus $S^{1} \times S^{1}$ is a Lie group under multiplication given by 
$$\left(e^{i \tau_{1}}, e^{i \theta_{1}}\right)\left(e^{i \tau_{2}}, e^{i \theta_{2}}\right)=\left(e^{i\left(\tau_{1}+\tau_{2}\right)}, e^{i\left(\theta_{1}+\theta_{2}\right)}\right)$$ 
Every homomorphism of $\mathbb{R}$ into $S^{1} \times S^{1}$, that is, every one-parameter subgroup of $S^{1} \times S^{1}$, is of the form $t \mapsto\left(e^{t a i}, e^{t b i}\right)$ for some pair of real numbers $a, b \in \mathbb{R}$.
\end{EXA}

Example 5.46. 

\begin{EXA}{LIE-L09-04-54}{1-parameter Untergruppe von $\GL(3)$ und $\SO(3)$}
The map $R: \mathbb{R} \rightarrow \mathrm{SO}(3)$ given by
$$
t \mapsto\left(\begin{array}{ccc}
\cos t & -\sin t & 0 \\
\sin t & \cos t & 0 \\
0 & 0 & 1
\end{array}\right)
$$
is a one-parameter subgroup. Also, the map
$$
t \mapsto\left(\begin{array}{ccc}
\cos t & -\sin t & 0 \\
\sin t & \cos t & 0 \\
0 & 0 & e^{t}
\end{array}\right)
$$
is a one-parameter subgroup of GL(3).
\end{EXA}


Recall that an $n \times n$ complex matrix $A$ is called Hermitian (resp. skewHermitian) if $\bar{A}^{t}=A$ (resp. $\bar{A}^{t}=-A$ ). Let $\mathfrak{s u}(2)$ denote the vector space of skew-Hermitian matrices with zero trace. We will later identify $\mathfrak{s u}(2)$ as the "Lie algebra" of $\mathrm{SU}(2)$.

Example 5.47. 

\begin{EXA}{LIE-L09-04-55}{$Ad_g$ wirkt auf $\SO(3)$}
Recall that an $n \times n$ complex matrix $A$ is called Hermitian (resp. skewHermitian) if $\bar{A}^{t}=A$ (resp. $\bar{A}^{t}=-A$ ). Let $\mathfrak{s u}(2)$ denote the vector space of skew-Hermitian matrices with zero trace. Given $g \in \mathrm{SU}(2)$, we define the map $\operatorname{Ad}_{g}: \mathfrak{s u}(2) \rightarrow \mathfrak{s u}(2)$ by $\operatorname{Ad}_{g}: x \mapsto g x g^{-1}$. The skew-Hermitian matrices of zero trace can be identified with $\mathbb{R}^{3}$ by using the following matrices as a basis:
$$
\left(\begin{array}{cc}
0 & -i \\
-i & 0
\end{array}\right),\left(\begin{array}{cc}
0 & -1 \\
1 & 0
\end{array}\right),\left(\begin{array}{cc}
-i & 0 \\
0 & i
\end{array}\right) .
$$
These are just $-i$ times the Pauli matrices $\sigma_{1}, \sigma_{2}, \sigma_{3}$, and so the correspondence $\mathfrak{s u}(2) \rightarrow \mathbb{R}^{3}$ is given by $-x i \sigma_{1}-y i \sigma_{2}-i z \sigma_{3} \mapsto(x, y, z)$. Under this correspondence, the inner product on $\mathbb{R}^{3}$ becomes the inner product $(A, B)=\frac{1}{2} \operatorname{trace}\left(A \bar{B}^{t}\right)=-\frac{1}{2} \operatorname{trace}(A B)$. But then
$$
\begin{aligned}
\left(\operatorname{Ad}_{g} A, \operatorname{Ad}_{g} B\right) & =-\frac{1}{2} \operatorname{trace}\left(g A g^{-1} g B g^{-1}\right) \\
& =-\frac{1}{2} \operatorname{trace}(A B)=(A, B)
\end{aligned}
$$
So, $\operatorname{Ad}_{g}$ can be thought of as an element of $\mathrm{O}(3)$. More is true; $\operatorname{Ad}_{g}$ acts as an element of $\mathrm{SO}(3)$, and the map $g \mapsto \mathrm{Ad}_{g}$ is then a homomorphism from $\mathrm{SU}(2)$ to $\mathrm{SO}(\mathfrak{s u}(2)) \cong \mathrm{SO}(3)$. This is a special case of the adjoint map studied later. (This example is related to the notion of "spin". For more, see the online supplement.)
\end{EXA}



Definition 5.48. 

\begin{DEF}{LIE-L09-04-56}{(Universelle) Überdeckungsgruppe}
If a Lie group homomorphism $\wp: \widetilde{G} \rightarrow G$ is also a covering map then we say that $\widetilde{G}$ is a covering group and $\wp$ is a covering homomorphism. If $\widetilde{G}$ is simply connected, then $\widetilde{G}$ (resp. $\wp$ ) is called the universal covering group (resp. universal covering homomorphism) of $G$.
\end{DEF}

Exercise 5.49. 


\begin{EXE}{LIE-L09-04-57}{Überlagerungsraum mit induzierter Lie Gruppen Struktur }
Show that if $\wp: \widetilde{M} \rightarrow G$ is a smooth covering map and $G$ is a Lie group, then $\widetilde{M}$ can be given a unique Lie group structure such that $\wp$ becomes a covering homomorphism. (You may assume that $\widetilde{M}$ is paracompact.)
\end{EXE}

Example 5.50. 

\begin{EXA}{LIE-L09-04-58}{Möbiustransformationen}
The group Mob of Möbius transformations of the complex plane given by $T_{A}: z \mapsto \frac{a z+b}{c z+d}$ for $A=\left(\begin{array}{ll}a & b \\ c & d\end{array}\right) \in \mathrm{SL}(2, \mathbb{C})$ can be given the structure of a Lie group. The map $\wp: \mathrm{SL}(2, \mathbb{C}) \rightarrow$ Mob given by $\wp: A \mapsto T_{A}$ is onto but not injective. In fact, it is a (two fold) covering homomorphism. When do two elements of $\operatorname{SL}(2, \mathbb{C})$ map to the same element of Mob?
\end{EXA}




\pagebreak



\subsection*{Lie Algebras and Exponential Maps}


Definition 5.51. 


\begin{DEF}{LIE-L09-03-01}{Links / Rechts Invariante Vektorfelder}
A vector field $X \in \mathfrak{X}(G)$ is called left invariant if and only if 
$$\left(L_{g}\right)_{*} X=X$$ 
for all $g \in G$. A vector field $X \in \mathfrak{X}(G)$ is called right invariant if and only if 
$$\left(R_{g}\right)_{*} X=X$$ 
for all $g \in G$. The set of left invariant (resp. right invariant) vector fields is denoted $\mathfrak{X}^{L}(G)$ (resp. $\mathfrak{X}^{R}(G)$).

Recall that by definition 
$$\left(L_{g}\right)_{*} X=T L_{g} \circ X \circ L_{g}^{-1}$$ 
and so left invariance means that 
$$T L_{g} \circ X \circ L_{g}^{-1}=X$$ 
or that given any $x \in G$ we have 
$$T_{x} L_{g} \cdot X(x)= X(g x)$$ for all $g \in G$. 

Thus $X \in \mathfrak{X}(G)$ is left invariant if and only if the following diagram commutes for every $g \in G$ :\\
\includegraphics[scale=0.25, center]{2025_10_09_265d7e1a22c89f79c21eg-219}
There is a similar diagram for right invariance.
\end{DEF}


Lemma 5.52. 

\begin{LEM}{LIE-L09-03-02}{Raum der links Invarianten Vektorfelder ist abgeschlossen unter Lie-Klammer}
$\mathfrak{X}^{L}(G)$ is closed under the Lie bracket operation.
\end{LEM}


Proof. 

\begin{PROOF}{LIE-L09-03-03}{P: Raum der links Invarianten Vektorfelder ist abgeschlossen unter Lie-Klammer}
Suppose that $X, Y \in \mathfrak{X}^{L}(G)$. Then by Proposition 2.84 we have
$$
\left(L_{g}\right)_{*}[X, Y]=\left[L_{g *} X, L_{g *} Y\right]=[X, Y] .
$$
\end{PROOF}


\begin{DEF}{LIE-L09-03-04}{Assoziierte Links/Rechts Invariante Felder $L^v,R^v$ zu Tangentialvektoren in $T_e G$}
Given a vector $v \in T_{e} G$, we can define a smooth left (resp. right) invariant vector field $L^{v}$ (resp. $R^{v}$) such that $L^{v}(e)=v$ (resp. $\left.R^{v}(e)=v\right)$ by the simple prescription
$$
L^{v}(g)=T L_{g} \cdot v \quad\left(\text { resp. } R^{v}(g)=T R_{g} \cdot v\right) .
$$
A bit more precisely, $L^{v}(g)=T_{e}\left(L_{g}\right) \cdot v$. The proof that this prescription gives smooth invariant vector fields is left to the reader (see Problem 7). Given a vector in $T_{e} G$ there are various notations for denoting the corresponding left (or right) invariant vector field, and we shall have occasion to use some different notation later on. We will also write $L(v)$ for $L^{v}$ and $R(v)$ for $R^{v}$. 

The map $v \mapsto L(v)$ (resp. $v \mapsto R(v)$) is a linear isomorphism from $T_{e} G$ onto $\mathfrak{X}^{L}(G)$ (resp. $\left.\mathfrak{X}^{R}(G)\right)$.
\end{DEF}

Definition 2.78 (Lie algebra). 

\begin{DEF}{LIE-L09-03-05}{Lie Algebra}
A vector space $\mathfrak{a}$ (over a field $\mathbb{F}$ ) is called a Lie algebra if it is equipped with a bilinear map $\mathfrak{a} \times \mathfrak{a} \rightarrow \mathfrak{a}$ (a multiplication) denoted $(v, w) \mapsto[v, w]$ such that
$$
[v, w]=-[w, v]
$$
and such that we have the Jacobi identity
$$
[x,[y, z]]+[y,[z, x]]+[z,[x, y]]=0
$$
for all $x, y, z \in \mathfrak{a}$.
\end{DEF}


Definition 2.79. 

\begin{DEF}{LIE-L09-03-06}{Abelsche Lie Algebra}
A Lie algebra $\mathfrak{a}$ is called abelian (or commutative) if $[v, w]=0$ for all $v, w \in \mathfrak{a}$.
\end{DEF}

\begin{DEF}{LIE-L09-03-07}{Lie Unteralgebra}
A subspace $\mathfrak{h}$ of $\mathfrak{a}$ is called a Lie subalgebra if it is closed under the bracket operation, and it is called an ideal if $[v, w] \in \mathfrak{h}$ for any $v \in \mathfrak{a}$ and $w \in \mathfrak{h}$. (We indicate this by writing $[\mathfrak{a}, \mathfrak{h}] \subset \mathfrak{h}$.)
\end{DEF}


Exercise 5.53. Show that $v \mapsto L^{v}$ gives a linear isomorphism $T_{e} G \cong \mathfrak{X}^{L}(G)$. Similarly, $T_{e} G \cong \mathfrak{X}^{R}(G)$ by $v \mapsto R(v)$.



\begin{PROP}{LIE-L09-03-58}{Lineare Isomorphie $T_{e} G \cong \mathfrak{X}^{L}(G)$}
Sei $G$ eine Liegruppe mit neutralem Element $e$.
\begin{enumerate}\itemsep0.2em
\item Die Abbildung
\[
\Phi:\ T_eG\longrightarrow \mathfrak{X}^L(G),\qquad 
\Phi(v)=L^v,\ \ \text{wobei }(L^v)_g:=T_eL_g(v),
\]
ist ein linearer Isomorphismus.
\item Analog ist
\[
\Psi:\ T_eG\longrightarrow \mathfrak{X}^R(G),\qquad 
\Psi(v)=R^v,\ \ \text{wobei }(R^v)_g:=T_eR_g(v),
\]
ein linearer Isomorphismus.
\end{enumerate}
\end{PROP}



\begin{PROOF}{LIE-L09-03-59}{P: Lineare Isomorphie $T_{e} G \cong \mathfrak{X}^{L}(G)$}
Wir beweisen (1); (2) folgt wortgleich mit Rechtsübersetzungen.

\textbf{(i) Wohldefiniertheit und Linksinvarianz.}
Für $v\in T_eG$ sei $L^v$ durch $(L^v)_g:=T_eL_g(v)$ definiert.
Die Abbildung $g\mapsto (L^v)_g$ ist glatt, da $(g,w)\mapsto T_wL_g$ glatt ist
und $w=e$ festgehalten wird. Für $h,g\in G$ gilt
\[
T_gL_h\big((L^v)_g\big)=T_gL_h\,T_eL_g(v)=T_e(L_h\circ L_g)(v)=T_eL_{hg}(v)=(L^v)_{hg},
\]
also ist $L^v$ linksinvariant.

\textbf{(ii) Linearität.}
Für $a,b\in\mathbb{R}$ und $v_1,v_2\in T_eG$ gilt
\[
L^{a v_1+b v_2}_g
= T_eL_g(a v_1+b v_2)
= a\,T_eL_g(v_1)+b\,T_eL_g(v_2)
= a\,(L^{v_1})_g+b\,(L^{v_2})_g.
\]
Also ist $\Phi$ linear.

\textbf{(iii) Injektivität.}
Angenommen $L^v=0$. Dann $(L^v)_e=T_eL_e(v)=\mathrm{id}_{T_eG}(v)=v=0$.
Also ist $\ker\Phi=\{0\}$.

\textbf{(iv) Surjektivität.}
Sei $X\in\mathfrak{X}^L(G)$ linksinvariant. Setze $v:=X_e\in T_eG$.
Für jedes $g\in G$ gilt wegen Linksinvarianz
\[
X_g = T_eL_g(X_e) = T_eL_g(v) = (L^v)_g,
\]
also $X=L^v$. Damit ist $\Phi$ surjektiv.

Die inverse Abbildung ist gegeben durch die Auswertung am Neutralen,
\[
\mathfrak{X}^L(G)\longrightarrow T_eG,\qquad X\mapsto X_e,
\]
und ist offensichtlich linear. Somit ist $\Phi$ ein linearer Isomorphismus.
\end{PROOF}





We now restrict attention to the left invariant fields but keep in mind that essentially all of what we say for this case has analogies in the right invariant case. We will discover a conduit (the adjoint map) between the two cases. 


Proposition 5.54. 

\begin{PROP}{LIE-L09-03-08}{Raum der Linksinvarianten Vektorfelder über Lie Gruppe ist $n$-dim Lie Algebra}
If $G$ is a Lie group of dimension $n$, then $\mathfrak{X}^{L}(G)$ is an $n$-dimensional Lie algebra under the bracket of vector fields (see Definition 2.75).
\end{PROP}

\begin{PROOF}{LIE-L09-03-09}{P: Raum der Linksinvarianten Vektorfelder über Lie Gruppe ist $n$-dim Lie Algebra}
The linear isomorphism $T_{e} G \cong \mathfrak{X}^{L}(G)$ just discovered shows that $\mathfrak{X}^{L}(G)$ is, in fact, a vector space of finite dimension equal to the dimension of $G$. From this and Lemma 5.52 we immediately obtain the result.
\end{PROOF}



Using the isomorphism $T_{e} G \cong \mathfrak{X}^{L}(G)$, we can transfer the Lie algebra structure to $T_{e} G$. This is the content of the following:

Definition 5.55. 

\begin{DEF}{LIE-L09-03-10}{$T_eG$ von Lie gruppe als Lie Algebra}
For a Lie group $G$, define the bracket of any two elements $v, w \in T_{e} G$ by
$$
[v, w]:=\left[L^{v}, L^{w}\right](e) .
$$
With this bracket, the vector space $T_{e} G$ becomes a Lie algebra
\end{DEF}


\begin{DEF}{LIE-L09-03-11}{Lie Algebra einer Lie Gruppe}
We have two Lie algebras, $\mathfrak{X}^{L}(G)$ and $T_{e} G$, which are isomorphic by construction. The abstract Lie algebra isomorphic to either/both of them is often referred to as the Lie algebra of the Lie group $G$ and denoted variously by $\mathfrak{L}(G)$ or $\mathfrak{g}$. Of course, we are implying that $\mathfrak{L}(H)$ is denoted $\mathfrak{h}$ and $\mathfrak{L}(K)$ by $\mathfrak{k}$, etc. In some computations we will have to use a specific realization of $\mathfrak{g}$. Our default convention will be that $\mathfrak{g}=\mathfrak{L}(G):=T_{e} G$ with the bracket defined above.
\end{DEF}

Definition 5.56. 

\begin{DEF}{LIE-L09-03-12}{Strukturkonstanten der Lie Algebra}
Given a Lie algebra $\mathfrak{g}$, we can associate to every basis $v_{1}, \ldots, v_{n}$ for $\mathfrak{g}$, the structure constants $c_{i j}^{k}$ which are defined by
$$
\left[v_{i}, v_{j}\right]=\sum_{k} c_{i j}^{k} v_{k} \text { for } 1 \leq i, j, k \leq n .
$$
It follows from the skew symmetry of the Lie bracket and the Jacobi identity that the structure constants satisfy
\begin{enumerate}
\item $$c_{i j}^{k}=-c_{j i}^{k}$$
\item $$\sum_{k} c_{r s}^{k} c_{k t}^{i}+c_{s t}^{k} c_{k r}^{i}+c_{t r}^{k} c_{k s}^{i}=0$$
\end{enumerate}
The structure constants characterize the Lie algebra, and structure constants are sometimes used to actually define a Lie algebra once a basis is chosen.
\end{DEF} 


\begin{DEF}{LIE-L09-03-13}{Produkt Lie Algebra}
Let $\mathfrak{a}$ and $\mathfrak{b}$ be Lie algebras. For ($a_{1}, b_{1}$) and ($a_{2}, b_{2}$) elements of the vector space $\mathfrak{a} \times \mathfrak{b}$, define
$$
\left[\left(a_{1}, b_{1}\right),\left(a_{2}, b_{2}\right)\right]:=\left(\left[a_{1}, a_{2}\right],\left[b_{1}, b_{2}\right]\right) .
$$
With this bracket, $\mathfrak{a} \times \mathfrak{b}$ is a Lie algebra called the Lie algebra product of $\mathfrak{a}$ and $\mathfrak{b}$. Recall the definition of an ideal in a Lie algebra (Definition 2.79). The subspaces $\mathfrak{a} \times\{0\}$ and $\{0\} \times \mathfrak{b}$ are ideals in $\mathfrak{a} \times \mathfrak{b}$ that are clearly isomorphic to $\mathfrak{a}$ and $\mathfrak{b}$ respectively. We often identify $\mathfrak{a}$ with $\mathfrak{a} \times\{0\}$ and $\mathfrak{b}$ with $\{0\} \times \mathfrak{b}$.
\end{DEF}

Exercise 5.57. 

\begin{EXE}{LIE-L09-03-14}{Lie Algebra von Produkt Lie Gruppe}
Show that if $G$ and $H$ are Lie groups, then the Lie algebra $\mathfrak{g} \times \mathfrak{h}$ is (up to identifications) the Lie algebra of $G \times H$.
\end{EXE}

Definition 5.58. 

\begin{DEF}{LIE-L09-03-15}{Lie Algebra Homomorphismus}
Given two Lie algebras over a field $\mathbb{F}$, say ($\mathfrak{a},[,]_{\mathfrak{a}}$) and $\left(\mathfrak{b},[,]_{\mathfrak{b}}\right)$, an $\mathbb{F}$-linear map $\sigma$ is called a Lie algebra homomorphism if and only if
$$
\sigma\left([v, w]_{\mathfrak{a}}\right)=[\sigma v, \sigma w]_{\mathfrak{b}}
$$
for all $v, w \in \mathfrak{a}$. A Lie algebra isomorphism is defined in the obvious way. A Lie algebra isomorphism $\mathfrak{g} \rightarrow \mathfrak{g}$ is called an automorphism of $\mathfrak{g}$.
\end{DEF}

\begin{DEF}{LIE-L09-03-16}{Automorphismengruppe eine Lie Algebra ist Lie Untergruppe von $\GL(\mathfrak{g})$}
It is not hard to show that the set of all automorphisms of $\mathfrak{g}$, denoted $\operatorname{Aut}(\mathfrak{g})$, forms a Lie group (actually a Lie subgroup of $\operatorname{GL}(\mathfrak{g})$).
\end{DEF}



\begin{DEF}{LIE-L09-03-17}{Lie-Algebra von $\mathrm{GL}(V)$}
Sei $V$ ein endlichdimensionaler reeller Vektorraum.
\begin{enumerate}
\item Die allgemeine lineare Gruppe $\mathrm{GL}(V)$ ist die offene Teilmenge
\[
\mathrm{GL}(V)=\{A\in L(V,V)\mid \det A\neq 0\}
\]
des Vektorraums $L(V,V)$ aller linearen Endomorphismen von $V$.
\item Das Tangentialbündel von $\mathrm{GL}(V)$ kann mit
\[
T\mathrm{GL}(V)\cong \mathrm{GL}(V)\times L(V,V)
\]
identifiziert werden. Insbesondere gilt
\[
T_A\mathrm{GL}(V)=\{A\}\times L(V,V),
\qquad
T_I\mathrm{GL}(V)\cong L(V,V).
\]
\item Die Lie-Algebra von $\mathrm{GL}(V)$ ist definiert als
\[
\mathfrak{gl}(V):=T_I\mathrm{GL}(V)\cong L(V,V),
\]
versehen mit der Lie-Klammer, die aus der Lie-Algebra der linksinvarianten
Vektorfelder auf $\mathrm{GL}(V)$ induziert wird.
\end{enumerate}
\end{DEF}



\begin{PROP}{LIE-L09-03-18}{Kommutatorstruktur auf $\mathfrak{gl}(V)$}
Unter der natürlichen Identifikation $T_I\mathrm{GL}(V)\cong L(V,V)$
entspricht die Lie-Klammer der zugehörigen linksinvarianten
Vektorfelder dem \emph{Kommutator} von Endomorphismen:
\[
[A,B]:=A\circ B - B\circ A,
\qquad
A,B\in L(V,V).
\]
\end{PROP}



\begin{PROOF}{LIE-L09-03-19}{P: Kommutatorstruktur auf $\mathfrak{gl}(V)$}
Sei für $X\in L(V,V)$ das linksinvariante Vektorfeld
$\widetilde{X}$ auf $\mathrm{GL}(V)$ definiert durch
\[
\widetilde{X}(A)=A\circ X,\qquad A\in\mathrm{GL}(V).
\]
Dann gilt für $A,B\in L(V,V)$:
\[
[\widetilde{A},\widetilde{B}](C)
=
\widetilde{A}(\widetilde{B}(C))-\widetilde{B}(\widetilde{A}(C))
=
C\circ(A\circ B-B\circ A)
=
\widetilde{[A,B]}(C).
\]
Somit induziert die Lie-Klammer auf den linksinvarianten Vektorfeldern
den Kommutator $[A,B]=A\circ B-B\circ A$ auf $L(V,V)$.
\end{PROOF}






Let V be a finite-dimensional real vector space. Then $G L(\mathrm{V})$ is an open subset of the linear space $L(\mathrm{V}, \mathrm{V})$, and we identify the tangent bundle of $G L(\mathrm{V})$ with $G L(\mathrm{V}) \times L(\mathrm{V}, \mathrm{V})$ (recall Definition 2.58). The tangent space at $A \in G L(\mathrm{V})$ is then $\{A\} \times L(\mathrm{V}, \mathrm{V})$. Now the Lie algebra is $T_{I} G L(\mathrm{V})$, and it has a Lie algebra structure derived from the Lie algebra structure on $\mathfrak{X}^{L}(G L(\mathrm{V}))$. The natural isomorphism of $T_{I} G L(\mathrm{V})$ with $L(\mathrm{V}, \mathrm{V})$ puts a Lie algebra structure on $L(\mathrm{V}, \mathrm{V})$. 

We now show that the resulting bracket on $L(\mathrm{V}, \mathrm{V})$ is just the commutator bracket given by $[A, B]:=A \circ B-B \circ A$. In the following discussion we let $\widetilde{X}$ denote the left invariant vector field corresponding to 
$$X \in L(\mathrm{V}, \mathrm{V}) \cong T_{I}(G L(\mathrm{V}))$$ 
By definition we have $\widetilde{[A, B]}= [\widetilde{A}, \widetilde{B}]$. 


We will need some simple results from the following easy exercises.

Exercise 5.59. 


For a fixed $A \in G L(\mathrm{V})$, the map $L_{A}: G L(\mathrm{V}) \rightarrow G L(\mathrm{V})$ given by $A \mapsto A \circ B$ has tangent map given by $(A, X) \mapsto(A \circ B, A \circ X)$ where $(A, X) \in G L(\mathrm{V}) \times L(\mathrm{V}, \mathrm{V}) \cong T(G L(\mathrm{V}))$. Show that if $\widetilde{X}$ denotes the left invariant vector field corresponding to $X \in L(\mathrm{V}, \mathrm{V})$, then $\widetilde{X}(A)=(A, A X)$.




\begin{PROP}{LIE-L09-03-20}{Formel für linksinvariante Vektorfelder auf $\mathrm{GL}(V)$}
Sei $V$ endlichdimensional reell und identifiziere $T\mathrm{GL}(V)\cong \mathrm{GL}(V)\times L(V,V)$.
Für festes $A\in \mathrm{GL}(V)$ sei $L_A:\mathrm{GL}(V)\to \mathrm{GL}(V)$ die Linksmultiplikation $L_A(B)=A\circ B$.
Dann gilt für $X\in L(V,V)\cong T_I\mathrm{GL}(V)$ und das zu $X$ gehörige linksinvariante Vektorfeld $\widetilde X$:
\[
\widetilde X(A)\;=\; d(L_A)_I(X)\;=\;(A,\,A\circ X)\;\in T_A\mathrm{GL}(V)\cong \{A\}\times L(V,V).
\]
\end{PROP}



\begin{PROOF}{LIE-L09-03-21}{P: Formel für linksinvariante Vektorfelder auf $\mathrm{GL}(V)$}
Unter der Identifikation $T\mathrm{GL}(V)\cong \mathrm{GL}(V)\times L(V,V)$ ist die totale Ableitung von $L_A$ an $(B,Y)$ gegeben durch
\[
d(L_A)_{(B,Y)}\;:\;(B,Y)\longmapsto (A\circ B,\,A\circ Y).
\]
Insbesondere am Einselement $I\in \mathrm{GL}(V)$ liefert dies
\[
d(L_A)_I \;:\; T_I\mathrm{GL}(V)\cong L(V,V)\longrightarrow T_A\mathrm{GL}(V)\cong \{A\}\times L(V,V),\quad
X\longmapsto (A,\,A\circ X).
\]
Per Definition des linksinvarianten Vektorfeldes zu $X$ gilt
\[
\widetilde X(A)\;=\; d(L_A)_I(X).
\]
Kombiniert mit der obigen Formel ergibt sich
\[
\widetilde X(A)=(A,\,A\circ X),
\]
wie behauptet.
\end{PROOF}





We are going to consider functions on $G L(\mathrm{V})$ that are restrictions of linear functionals on the vector space $L(\mathrm{V}, \mathrm{V})$. We will not notationally distinguish the functional from its restriction. If $f$ is such a linear function and $X \in L(\mathrm{V}, \mathrm{V})$, let $f_{, X}$ be given by $f_{, X}(A):=f(A \circ X)$. It is easy to check that $f_{, X}$ is also a linear functional, and so for $Y \in L(\mathrm{V}, \mathrm{V})$ we also have $\left(f_{, X}\right)_{, Y}$. It is clear that $\left(f_{, X}\right)_{, Y}=f_{, Y \circ X}$ (notice the reversal of order).

Exercise 5.60. Show that the map $L(\mathrm{V}, \mathrm{V}) \rightarrow(L(\mathrm{V}, \mathrm{V}))^{*}$ given by $X \mapsto f_{, X}$ is linear over $\mathbb{R}$.

Exercise 5.61. Let $\widetilde{X}$ be the left invariant field corresponding to $X \in L(\mathrm{V}, \mathrm{V})$ as above. If $f$ is the restriction to $G L(\mathrm{V})$ of a linear functional on $L(\mathrm{V}, \mathrm{V})$ as above, then $\widetilde{X} f=f_{, X}$. Solution/Hint: $(\widetilde{X} f)(A)=\left.d f\right|_{A}\left(\widetilde{X}_{A}\right)= f(A \circ X)$.

Recall that if one picks a basis for V , then we obtain a global chart on $G L(\mathrm{V})$ which is given by coordinate functions $x_{j}^{i}$ defined by letting $x_{j}^{i}(A)$ be the $i j$-th entry of the matrix of $A$ with respect to the chosen basis. These coordinate functions are restrictions of linear functions on $L(\mathrm{V}, \mathrm{V})$. Using this we see that if $f_{, X}=f_{, Y}$ for all linear $f$, then $f(X)=f_{, X}(I)=f_{, Y}(I)= f(Y)$ for all linear $f$ and this in turn implies $X=Y$.


% Setup: lineare Funktionale und ihre Restriktionen

\begin{DEF}{LIE-L09-03-22}{Lineare Funktionale auf $GL(V)$ via Restriktion}
Sei $V$ endlichdimensional reell, $L(V,V)$ der Raum der Endomorphismen und $GL(V)\subset L(V,V)$ offen.
Eine \emph{lineare Funktion} $f$ auf $GL(V)$ sei stets die Restriktion eines linearen Funktionals
$f\in (L(V,V))^*$.
Für $X\in L(V,V)$ definieren wir
\[
f_{,X}:GL(V)\to\mathbb R,\qquad f_{,X}(A):=f(A\circ X).
\]
\end{DEF}



\begin{LEM}{LIE-L09-03-23}{Kompositionsregel für $f_{,X}$}
Für lineares $f$ und $X,Y\in L(V,V)$ gilt
\[
\bigl(f_{,X}\bigr)_{,Y}=f_{,\,Y\circ X}.
\]
\end{LEM}



\begin{PROOF}{LIE-L09-03-24}{P: Kompositionsregel für $f_{,X}$}
Für $A\in GL(V)$:
\[
\bigl(f_{,X}\bigr)_{,Y}(A)=f_{,X}(A\circ Y)=f\bigl((A\circ Y)\circ X\bigr)=f\bigl(A\circ (Y\circ X)\bigr)=f_{,\,Y\circ X}(A).
\]
\end{PROOF}



\begin{LEM}{LIE-L09-03-25}{Linearität von $X\mapsto f_{,X}$}
Die Abbildung
\[
T_f:L(V,V)\longrightarrow (L(V,V))^*,\qquad X\longmapsto f_{,X},
\]
ist $\mathbb R$-linear.
\end{LEM}



\begin{PROOF}{LIE-L09-03-26}{P: Linearität von $X\mapsto f_{,X}$}
Seien $X_1,X_2\in L(V,V)$ und $\alpha,\beta\in\mathbb R$. Für $A\in GL(V)$ gilt
\[
\bigl(\alpha f_{,X_1}+\beta f_{,X_2}\bigr)(A)=\alpha\,f(A\circ X_1)+\beta\,f(A\circ X_2)
=f\bigl(A\circ(\alpha X_1+\beta X_2)\bigr)=f_{,\,\alpha X_1+\beta X_2}(A).
\]
Also $T_f(\alpha X_1+\beta X_2)=\alpha T_f(X_1)+\beta T_f(X_2)$.
\end{PROOF}




\begin{PROP}{LIE-L09-03-27}{Linksinvariante Ableitung auf Restriktionen}
Sei $\widetilde X$ das linksinvariante Vektorfeld zu $X\in L(V,V)\cong T_I GL(V)$.
Für jedes lineare $f$ (als Restriktion) gilt
\[
\widetilde X f \;=\; f_{,X}.
\]
\end{PROP}



\begin{PROOF}{LIE-L09-03-28}{P: Linksinvariante Ableitung auf Restriktionen}
Mit der Identifikation $T GL(V)\cong GL(V)\times L(V,V)$ ist (vgl. Linksmultiplikation $L_A$)
\[
\widetilde X(A)=d(L_A)_I(X)=(A,\,A\circ X).
\]
Damit erhält man
\[
(\widetilde X f)(A)=\bigl.df\bigr|_A\bigl(\widetilde X(A)\bigr)=\bigl.df\bigr|_A(A\circ X)=f(A\circ X)=f_{,X}(A).
\]
\end{PROOF}



\begin{DEF}{LIE-L09-03-29}{Globale Koordinaten und Trennung von Punkten}
Fixiere eine Basis von $V$. Die Koordinatenfunktionen $x^i_j(A)$ (Matrixeinträge von $A$) sind
Restriktionen linearer Funktionale auf $L(V,V)$. Gilt $f_{,X}=f_{,Y}$ für alle linearen $f$, so folgt
aus der Auswertung bei $I$:
\[
f(X)=f_{,X}(I)=f_{,Y}(I)=f(Y)\quad\text{für alle linearen }f,
\]
also $X=Y$ (lineare Funktionale trennen Punkte in $L(V,V)$).
\end{DEF}







Proposition 5.62. 

\begin{PROP}{LIE-L09-03-30}{Induzierte Lie Algebra Klammer auf $L(\mathrm{V},\mathrm{V})$}
The Lie algebra bracket on $L(\mathrm{V},\mathrm{V})$ induced by the isomorphisms 
$$\mathfrak{X}_{L}(G L(\mathrm{V})) \cong T_{I} G L(\mathrm{V}) \cong L(\mathrm{V}, \mathrm{V})$$ is the commutator bracket
$$
[X, Y]:=X \circ Y-Y \circ X .
$$
\end{PROP}

Proof. 

\begin{PROOF}{LIE-L09-03-31}{P: Induzierte Lie Algebra Klammer auf $L(\mathrm{V},\mathrm{V})$}
Let $X, Y \in L(\mathrm{V}, \mathrm{V})$ and let $[X, Y]$ denote the bracket induced on $L(\mathrm{V}, \mathrm{V})$. Then for any linear $f$ we have
$$
\begin{aligned}
f_{,[X, Y]} & =\widetilde{[X, Y]} f=[\widetilde{X}, \widetilde{Y}] f=\widetilde{X} \widetilde{Y} f-\widetilde{Y} \widetilde{X} f=\widetilde{X}\left(f_{, Y}\right)-\widetilde{Y}\left(f_{, X}\right) \\
& =\left(f_{, Y}\right)_{, X}-\left(f_{, X}\right)_{, Y}=f_{, X \circ Y}-f_{, Y \circ X}=f_{,(X \circ Y-Y \circ X)}
\end{aligned}
$$
Since $f$ was an arbitrary linear functional, we conclude that $[X, Y]= X \circ Y-Y \circ X$.
\end{PROOF}


A choice of basis gives the algebra isomorphism $L(\mathrm{V}, \mathrm{V}) \cong M_{n \times n}(\mathbb{R})$ where $n=\operatorname{dim}(\mathrm{V})$. This isomorphism preserves brackets and restricts to a group isomorphism $G L(\mathrm{V}) \cong G L(n, \mathbb{R})$. 


\begin{PROP}{LIE-L09-03-32}{Koordinatendarstellung von $\mathrm{GL}(V)$ und $\mathfrak{gl}(V)$}
Sei $V$ ein endlichdimensionaler reeller Vektorraum der Dimension $n$ und
$(\mathbf e_1,\ldots,\mathbf e_n)$ eine feste Basis von $V$.
\begin{enumerate}
\item Die Wahl dieser Basis definiert einen kanonischen Isomorphismus von
Vektorräumen
\[
L(V,V)\;\cong\; M_{n\times n}(\mathbb R),
\qquad
A\longmapsto [A]_{\mathbf e},
\]
wobei $[A]_{\mathbf e}$ die darstellende Matrix von $A$ bezüglich der gewählten Basis ist.
\item Dieser Isomorphismus ist ein \emph{Algebraisomorphismus}, d.\,h.
\[
[A\circ B]_{\mathbf e}=[A]_{\mathbf e}\,[B]_{\mathbf e},
\qquad
[A,B]_{\mathbf e}=[A]_{\mathbf e}[B]_{\mathbf e}-[B]_{\mathbf e}[A]_{\mathbf e}.
\]
Somit überträgt sich die Lie-Algebra‐Struktur
\[
\mathfrak{gl}(V)\;\cong\;\mathfrak{gl}(n,\mathbb R).
\]
\item Einschränkung auf die invertierbaren Endomorphismen liefert einen
Gruppenisomorphismus
\[
\mathrm{GL}(V)\;\cong\;\mathrm{GL}(n,\mathbb R),
\qquad
A\longmapsto [A]_{\mathbf e},
\]
der mit der Matrizenmultiplikation verträglich ist.
\end{enumerate}
\end{PROP}



\begin{PROOF}{LIE-L09-03-33}{P: Koordinatendarstellung von $\mathrm{GL}(V)$ und $\mathfrak{gl}(V)$}
Die Abbildung $A\mapsto [A]_{\mathbf e}$ ordnet jedem Endomorphismus seine
darstellende Matrix zu. Diese Abbildung ist linear und erfüllt
$[A\circ B]_{\mathbf e}=[A]_{\mathbf e}[B]_{\mathbf e}$, also ist sie ein
Algebraisomorphismus zwischen $L(V,V)$ und $M_{n\times n}(\mathbb R)$.
Damit wird der Kommutator $[A,B]=A\circ B-B\circ A$ auf
$\mathfrak{gl}(V)$ zum Matrizenkommutator in $M_{n\times n}(\mathbb R)$.
Für die invertierbaren Endomorphismen gilt
$A\in \mathrm{GL}(V)\iff [A]_{\mathbf e}\in \mathrm{GL}(n,\mathbb R)$,
und die Multiplikationsstruktur bleibt erhalten.
\end{PROOF}






\begin{PROP}{LIE-L09-03-34}{Lie-Algebra von $\mathrm{GL}(n,\mathbb{R})$}
Die kanonische lineare Identifikation
\[
\mathfrak{gl}(n,\mathbb{R}) \;=\; T_I \mathrm{GL}(n,\mathbb{R}) 
\;\cong\; M_{n\times n}(\mathbb{R})
\]
induziert auf $M_{n\times n}(\mathbb{R})$ eine Lie-Algebra-Struktur,
deren Lie-Klammer durch den \emph{Matrizenkommutator}
\[
[A,B]\;:=\;AB-BA
\qquad
\text{für }A,B\in M_{n\times n}(\mathbb{R})
\]
gegeben ist.
\end{PROP}



\begin{PROOF}{LIE-L09-03-35}{P: Lie-Algebra von $\mathrm{GL}(n,\mathbb{R})$}
Das Tangentialbündel von $\mathrm{GL}(n,\mathbb{R})$ lässt sich als
$\mathrm{GL}(n,\mathbb{R})\times M_{n\times n}(\mathbb{R})$ identifizieren.
Für $X\in M_{n\times n}(\mathbb{R})$ ist das linksinvariante Vektorfeld
$\widetilde X(A)=A\,X$. Für $A,B\in M_{n\times n}(\mathbb{R})$ gilt daher
\[
[\widetilde A,\widetilde B](C)
=\widetilde A(\widetilde B(C))-\widetilde B(\widetilde A(C))
=C(AB-BA)
=\widetilde{[A,B]}(C).
\]
Damit entspricht die Lie-Klammer der linksinvarianten Vektorfelder auf
$\mathrm{GL}(n,\mathbb{R})$ genau dem Kommutator der Matrizen.
\end{PROOF}


It is now easy to arrive at the following corollary.

Corollary 5.63. 

The linear isomorphism of $\mathfrak{g l}(n, \mathbb{R})=T_{I} G L(n, \mathbb{R})$ with $M_{n \times n}(\mathbb{R})$ induces a Lie algebra structure on $M_{n \times n}(\mathbb{R})$ such that the bracket is given by $[A, B]:=A B-B A$.





From now on we follow the practice of identifying the Lie algebra $\mathfrak{g} \mathfrak{l}(V)$ of $G L(\mathrm{V})$ with the commutator Lie algebra $L(\mathrm{V}, \mathrm{V})$. Similarly for the matrix group, the Lie algebra of $G L(n, \mathbb{R})$ is taken to be $M_{n \times n}(\mathbb{R})$ with commutator bracket.

If $G \subset \mathrm{GL}(n)$ is some matrix group, then $T_{I} G$ may be identified with a linear subspace of $M_{n \times n}$. This linear subspace will be closed under the commutator bracket, and so we actually have an identification of Lie algebras: $\mathfrak{g}$ is identified with a subspace of $M_{n \times n}$. It is often the case that $G$ is defined by some matrix equation or equations. By differentiating these equations we find the defining equations for $\mathfrak{g}$ (as a subspace of $M_{n \times n}$ ). We first prove a general result for Lie algebras of closed subgroups, and then we apply this to some matrix groups. Recall that if $N$ is a submanifold of $M$ and $\iota: N \hookrightarrow M$ is the inclusion map, then we identify $T_{p} N$ with $T_{p} \iota\left(T_{p} N\right)$ and $T_{p} \iota$ is an inclusion. If $H$ is a closed Lie subgroup of a Lie group $G$, and $v \in \mathfrak{h}=T_{e} H$, then $v$ corresponds to a left invariant vector field on $G$ which is obtained by using the left translation in $G$. But $v$ also corresponds to a left invariant vector field on $H$ obtained from left translations in $H$. The notation we have been using so far is not sensitive to this distinction, so let us introduce an alternative notation.




\begin{DEF}{LIE-L09-03-38}{Identifikation von Lie-Algebren mit Kommutatorräumen}
Von nun an identifizieren wir die Lie-Algebra
\[
\mathfrak{gl}(V)\;=\;T_I \mathrm{GL}(V)
\]
mit dem Vektorraum $L(V,V)$ versehen mit der Kommutator-Klammer
\[
[A,B]=A\circ B - B\circ A.
\]
Analog identifizieren wir für den Matrixfall
\[
\mathfrak{gl}(n,\mathbb{R})\;=\;T_I \mathrm{GL}(n,\mathbb{R})
\]
mit $M_{n\times n}(\mathbb{R})$ und derselben Kommutatorstruktur.
\end{DEF}



\begin{PROP}{LIE-L09-03-36}{Lie-Algebra einer Matrixuntergruppe}
Sei $G\subset \mathrm{GL}(n,\mathbb{R})$ eine Matrixgruppe (d.\,h. ein Lie‐Untergruppe der allgemeinen linearen Gruppe).
Dann gilt:
\begin{enumerate}
\item Der Tangentialraum im Einselement,
\[
\mathfrak g\;:=\;T_I G,
\]
kann als linearer Unterraum von $M_{n\times n}(\mathbb{R})$ aufgefasst werden.
\item Dieser Unterraum ist unter der Kommutator‐Klammer abgeschlossen, d.\,h.
\[
[A,B]\in \mathfrak g \quad \text{für alle }A,B\in \mathfrak g,
\]
und somit ist $\mathfrak g$ eine Lie‐Unteralgebra von $\mathfrak{gl}(n,\mathbb{R})$.
\end{enumerate}
\end{PROP}



\begin{PROOF}{LIE-L09-03-37}{P: Lie-Algebra einer Matrixuntergruppe}
Da $G$ eine Lie‐Untergruppe von $\mathrm{GL}(n,\mathbb{R})$ ist, ist $G$ eine eingebettete Untermannigfaltigkeit
mit Inklusion $\iota:G\hookrightarrow \mathrm{GL}(n,\mathbb{R})$.
Der Tangentialraum im Einselement ist daher $T_I G=\mathrm{im}(T_I\iota)$, und wir identifizieren ihn
mit einem linearen Unterraum von $T_I\mathrm{GL}(n,\mathbb{R})\cong M_{n\times n}(\mathbb{R})$.

Seien $A,B\in \mathfrak g=T_I G$ und $\widetilde A,\widetilde B$ die zugehörigen linksinvarianten
Vektorfelder auf $G$. Ihr Lie‐Klammerfeld $[\widetilde A,\widetilde B]$ ist wieder linksinvariant und tangential an $G$,
da $G$ als Untergruppe unter der Gruppenmultiplikation abgeschlossen ist.
Daher gehört der Wert dieses Feldes im Einselement $I$ ebenfalls zu $T_I G$.
Somit ist $\mathfrak g$ unter der Kommutator‐Klammer abgeschlossen und daher eine Lie‐Unteralgebra.
\end{PROOF}



\begin{CONC}{LIE-L09-03-39}{Lie-Algebra als Ableitung von Gruppengleichungen}
Ist $G\subset \mathrm{GL}(n,\mathbb{R})$ durch eine oder mehrere glatte Gleichungen
\[
F_i(A)=0,\qquad F_i:\mathrm{GL}(n,\mathbb{R})\to \mathbb R,
\]
definiert, so erhält man die definierenden Gleichungen der Lie‐Algebra
$\mathfrak g=T_I G$ durch linearisieren:
\[
(dF_i)_I(X)=0,\qquad X\in \mathfrak g\subset M_{n\times n}(\mathbb{R}).
\]
Dies erlaubt die explizite Berechnung von $\mathfrak g$ für viele klassische Matrixgruppen
(z.\,B. $\mathrm{O}(n)$, $\mathrm{SO}(n)$, $\mathrm{SL}(n)$, $\mathrm{U}(n)$).
\end{CONC}






Notation 5.64. 

\begin{DEF}{LIE-L09-03-40}{Notation $v^G$}
For a Lie group $G$, we have the alternative notation $v^{G}$ for the left invariant vector field whose value at $e$ is $v$. If $H$ is a closed Lie subgroup of $G$ and $v \in T_{e} H$, then $v^{G} \in \mathfrak{X}^{L}(G)$ while $v^{H} \in \mathfrak{X}^{L}(H)$.
\end{DEF}

Proposition 5.65. 

\begin{PROP}{LIE-L09-03-41}{Linksinvariante Felder: Restriktion auf eine abgeschlossene Untergruppe}
Let $H$ be a closed Lie subgroup of a Lie group $G$. Let
$$
\widetilde{\mathfrak{X}^{L}}(H):=\left\{X \in \mathfrak{X}^{L}(G): X(e) \in T_{e} H\right\} .
$$
Then the restrictions of elements of $\widetilde{\mathfrak{X}^{L}}(H)$ to the submanifold $H$ are the elements of $\mathfrak{X}^{L}(H)$. This induces an isomorphism of Lie algebras of vector fields $\widetilde{\mathfrak{X}^{L}}(H) \cong \mathfrak{X}^{L}(H)$. For $v, w \in \mathfrak{h}$, we have
$$
[v, w]_{\mathfrak{h}}=\left[v^{H}, w^{H}\right]_{e}=\left[v^{G}, w^{G}\right]_{e}=[v, w]_{\mathfrak{g}}
$$
and so $\mathfrak{h}$ is a Lie subalgebra of $\mathfrak{g}$; the bracket on $\mathfrak{h}$ is the same as that inherited from $\mathfrak{g}$.
\end{PROP}

Proof. 

\begin{PROOF}{LIE-L09-03-42}{P: Linksinvariante Felder: Restriktion auf eine abgeschlossene Untergruppe}
The Lie bracket on $\mathfrak{h}=T_{e} H$ is given by $[v, w]:=\left[v^{H}, w^{H}\right](e)$. Notice that if $H$ is a closed Lie subgroup of a Lie group $G$, then for $h \in H$ we have left translation by $h$ as a map $G \rightarrow G$ and also as a map $H \rightarrow H$. The latter is the restriction of the former. To avoid notational clutter, let us denote\\
$\left.L_{h}\right|_{H}$ by $l_{h}$. If $\iota: H \hookrightarrow G$ is the inclusion, then we have $\iota \circ l_{h}=L_{h} \circ \iota$, and so $T \iota \circ T l_{h}=T L_{h} \circ T \iota$. If $v^{H} \in \mathfrak{X}^{L}(H)$, then we have
$$
\begin{aligned}
T_{h} \iota\left(v^{H}(h)\right) & =T_{h} \iota\left(T_{e} l_{h}(v)\right)=T \iota \circ T l_{h}(v) \\
& =T L_{h} \circ T \iota(v)=T_{e} L_{h}\left(T_{e} \iota(v)\right) \\
& =T_{e} L_{h}(v)=v^{G}(h)=\left(v^{G} \circ \iota\right)(h),
\end{aligned}
$$
so that $v^{H}$ and $v^{G}$ are $\iota$-related for any $v \in T_{e} H \subset T_{e} G$. Thus for $v, w \in T_{e} H$ we have
$$
[v, w]_{\mathfrak{h}}=\left[v^{H}, w^{H}\right]_{e}=\left[v^{G}, w^{G}\right]_{e}=[v, w]_{\mathfrak{g}},
$$
the formula we wanted. Next, notice that if we take $T_{h} \iota$ as an inclusion so that $T_{h} \iota\left(v^{H}(h)\right)=v^{H}(h)$ for all $h$, then we have really shown that if $v \in T_{e} H$, then $v^{H}$ is the restriction of $v^{G}$ to $H$. Also it is easy to see that
$$
\widetilde{\mathfrak{X}^{L}}(H)=\left\{v^{G}: v \in T_{e} H\right\},
$$
and so the restrictions of elements of $\widetilde{\mathfrak{X}^{L}}(H)$ are none other than the elements of $\mathfrak{X}^{L}(H)$. From what we have shown, the restriction map $\widetilde{\mathfrak{X}^{L}}(H) \rightarrow \mathfrak{X}^{L}(H)$ is given by $v^{G} \longmapsto v^{H}$ and is a surjective Lie algebra homomorphism. It also has kernel zero since if $v^{H}$ is the zero vector field, then $v=0$, which implies that $v^{G}$ is the zero vector field.
\end{PROOF} 



Because $[v, w]_{\mathfrak{h}}=[v, w]_{\mathfrak{g}}$, the inclusion $\mathfrak{h} \hookrightarrow \mathfrak{g}$ is a Lie algebra homomorphism. Examining the details of the previous proof we see that we have a commutative diagram\\
\includegraphics[scale=0.25, center]{2025_10_09_265d7e1a22c89f79c21eg-224(1)}\\
of Lie algebra homomorphisms where the top horizontal map is a restriction to $H$, the left diagonal map is $v \longmapsto v^{G}$ and the right diagonal map is $v \longmapsto v^{H}$. In practice, what this last proposition shows is that in order to find the Lie algebra of a closed subgroup $H \subset G$, we only need to find the subspace $\mathfrak{h}=T_{e} H$, since the bracket on $\mathfrak{h}$ is just the restriction of the bracket on $\mathfrak{g}$. The following is also easily seen to be a commutative diagram of Lie algebra homomorphisms:\\
\includegraphics[scale=0.25, center]{2025_10_09_265d7e1a22c89f79c21eg-224}\\
where both vertical maps are $v \longmapsto v^{G}$.



\begin{PROP}{BÜN-L09-02-29}{Inklusion und kommutative Diagramme von Lie-Algebrenhomomorphismen}
Da für alle $v,w\in \mathfrak h=T_eH$ gilt
\[
[v,w]_{\mathfrak h}=[v,w]_{\mathfrak g},
\]
ist die Inklusion $\mathfrak h\hookrightarrow \mathfrak g$ ein Lie-Algebrenhomomorphismus.
Aus den Details des vorhergehenden Beweises erhalten wir das folgende kommutative Diagramm:
\\
\includegraphics[scale=0.25, center]{2025_10_09_265d7e1a22c89f79c21eg-224(1)}\\
Hierbei ist die obere horizontale Abbildung die Restriktion auf $H$, die linke diagonale
$v\mapsto v^{G}$ und die rechte diagonale $v\mapsto v^{H}$.
In der Praxis zeigt diese Aussage, dass zur Bestimmung der Lie-Algebra einer abgeschlossenen
Untergruppe $H\subset G$ lediglich der Tangentialraum
\[
\mathfrak h=T_eH
\]
bestimmt werden muss, da die Klammer auf $\mathfrak h$ einfach die Einschränkung der Klammer auf
$\mathfrak g$ ist.

Ferner ergibt sich ein weiteres kommutatives Diagramm von Lie‐Algebrenhomomorphismen:
\\
\includegraphics[scale=0.25, center]{2025_10_09_265d7e1a22c89f79c21eg-224}\\
Hierbei sind beide vertikalen Abbildungen $v\mapsto v^{G}$.
\end{PROP}





The Lie algebra-Lie group correspondence works in the other direction too:

Theorem 5.66. 


\begin{THEO}{LIE-L09-03-43}{Assoziierte eindeutige ZSH Lie Untergruppe zu Lie Unteralgebra}
Let $G$ be a Lie group with Lie algebra $\mathfrak{g}$. If $\mathfrak{h}$ is a Lie subalgebra of $\mathfrak{g}$, then there is a unique connected Lie subgroup $H$ of $G$ whose Lie algebra is $\mathfrak{h}$.
\end{THEO}

The above theorem uses the Frobenius integrability theorem, so we defer the proof until we have that theorem in hand.

Since the Lie algebra of $\operatorname{GL}(n)$ is the set of all $n \times n$ matrices with the commutator bracket, and since we have just shown that the bracket for the Lie algebra of a subgroup is just the restriction of the bracket on the containing group, we see that the bracket on the Lie algebra of matrix subgroups can also be taken to be the commutator bracket if that Lie algebra is represented by the appropriate space of matrices. We record this as a proposition:


Proposition 5.67. 

\begin{PROP}{LIE-L09-03-44}{Assoziierte Lie Unteralgebra in $\operatorname{End}(V)$ zu linearer Lie Untergruppe von $\GL(V)$}
If $G \subset \mathrm{GL}(\mathrm{V})$ is a linear Lie group, then the Lie algebra of $G$ may be identified with a subalgebra of $\operatorname{End}(\mathrm{V})$ with the commutator bracket. A similar statement holds for matrix Lie algebras.
\end{PROP}


\begin{PROOF}{LIE-L09-03-45}{P: Assoziierte Lie Unteralgebra in $\operatorname{End}(V)$ zu linearer Lie Untergruppe von $\GL(V)$}
Da $G\subset \GL(V)$ eine lineare Lie‐Gruppe ist, ist die Inklusion
$\iota:G\hookrightarrow \GL(V)$ eine glatte Einbettung. Folglich ist die
Tangentialabbildung im Einselement $T_I\iota:\,T_I G\to T_I\GL(V)$ injektiv,
und wir identifizieren
\[
\mathfrak g:=T_I G \;\subset\; T_I\GL(V)\cong \operatorname{End}(V).
\]
Es bleibt zu zeigen, dass $\mathfrak g$ unter dem Kommutator abgeschlossen ist
und dass die induzierte Lie‐Klammer mit dem Kommutator $[A,B]=AB-BA$
auf $\operatorname{End}(V)$ übereinstimmt.

Für $X\in \operatorname{End}(V)\cong T_I\GL(V)$ bezeichne $\widetilde X$ das durch Linksmultiplikation
linksinvariante Vektorfeld auf $\GL(V)$, d.\,h.
\[
\widetilde X(A)=A\,X\in T_A\GL(V)\cong \{A\}\times \operatorname{End}(V).
\]
Ist $X\in \mathfrak g=T_I G$, so ist die Restriktion $\widetilde X|_{G}$ ein glattes
Vektorfeld auf $G$, das linksinvariant bzgl.\ der Gruppenstruktur von $G$ ist;
insbesondere ist $\widetilde X|_{G}$ tangential an $G$.

Für $A,B\in \mathfrak g$ betrachten wir die linksinvarianten Felder
$\widetilde A,\widetilde B$ auf $\GL(V)$. Dann gilt auf $\GL(V)$ (Standardrechnung)
\[
[\widetilde A,\widetilde B](C)=C\,(AB-BA)
= \widetilde{[A,B]}(C),
\]
also entspricht die Lie‐Klammer der linksinvarianten Felder dem Kommutator
in $\operatorname{End}(V)$. Da $G$ Untergruppe ist, ist $[\widetilde A|_G,\widetilde B|_G]$
wieder tangential an $G$; insbesondere liegt sein Wert im Einselement in $T_I G$.
Damit folgt $[A,B]\in \mathfrak g$ und die von $G$ induzierte Klammer auf
$\mathfrak g$ stimmt mit dem Kommutator in $\operatorname{End}(V)$ überein. Somit ist
$\mathfrak g$ eine Lie‐Unteralgebra von $(\operatorname{End}(V),[\cdot,\cdot])$.

Im Matrixfall liefert die kanonische Identifikation
$T_I\GL(n,\R)\cong M_{n\times n}(\R)$ die analoge Aussage: Für eine
Matrix‐Lie‐Untergruppe $G\subset \GL(n,\R)$ ist $\mathfrak g=T_I G$ ein
linearer Unterraum von $M_{n\times n}(\R)$, der unter dem Matrizenkommutator
abgeschlossen ist, also eine Lie‐Unteralgebra von $\mathfrak{gl}(n,\R)$.
\end{PROOF}



Example 5.68. 

\begin{EXA}{LIE-L09-03-46}{Lie Algebra von $\Ogroup(n)$}
Consider the orthogonal group $\mathrm{O}(n) \subset \operatorname{GL}(n)$. Given a curve of orthogonal matrices $Q(t)$ with $Q(0)=I$ and $\left.\frac{d}{d t}\right|_{t=0} Q(0)=A$, we compute by differentiating the defining equation $I=Q^{t} Q$ :
$$
\begin{aligned}
0 & =\left.\frac{d}{d t}\right|_{t=0} Q^{t} Q \\
& =\left(\left.\frac{d}{d t}\right|_{t=0} Q\right)^{t} Q(0)+Q^{t}(0)\left(\left.\frac{d}{d t}\right|_{t=0} Q\right) \\
& =A^{t}+A
\end{aligned}
$$
so that the space of skew-symmetric matrices is contained in the tangent space $T_{I} \mathrm{O}(n)$. But both $T_{I} \mathrm{O}(n)$ and the space of skew-symmetric matrices have dimension $n(n-1) / 2$, so they are equal. This means that we can identify the Lie algebra $\mathfrak{o}(n)=\mathfrak{L}(\mathrm{O}(n))$ with the space of skew-symmetric matrices with the commutator bracket. One can easily check that the commutator bracket of two such matrices is skew-symmetric, as expected.
\end{EXA}

\begin{CONC}{LIE-L09-03-47}{Berechnung von Lie Algebren in $\operatorname{GL}(n, \mathbb{C})$}
We have considered matrix groups as subgroups of $\mathrm{GL}(n)$, but it is often more convenient to consider subgroups of $\operatorname{GL}(n, \mathbb{C})$. Since $\operatorname{GL}(n, \mathbb{C})$ can be identified with a subgroup of $\operatorname{GL}(2 n)$, this is only a slight change in viewpoint. The essential parts of our discussion go through for $\operatorname{GL}(n, \mathbb{C})$ without any significant change.
\end{CONC}



Example 5.69. 

\begin{EXA}{LIE-L09-03-48}{Lie Algebra von $\U(n)$}
Consider the unitary group $\mathrm{U}(n) \subset \mathrm{GL}(n, \mathbb{C})$. Given a curve of unitary matrices $Q(t)$ with $Q(0)=I$ and $\left.\frac{d}{d t}\right|_{t=0} Q(0)=A$, we compute by differentiating the defining equation $I=\bar{Q}^{t} Q$. We have
$$
\begin{aligned}
0 & =\left.\frac{d}{d t}\right|_{t=0} \bar{Q}^{t} Q \\
& =\left(\left.\frac{d}{d t}\right|_{t=0} \bar{Q}\right)^{t} Q(0)+\bar{Q}^{t}(0)\left(\left.\frac{d}{d t}\right|_{t=0} Q\right) \\
& =\bar{A}^{t}+A
\end{aligned}
$$
Examining dimensions as before, we see that we can identify $\mathfrak{u}(n)$ with the space of skew-Hermitian matrices $\left(\bar{A}^{t}=-A\right)$ under the commutator bracket.
\end{EXA}

\begin{EXA}{LIE-L09-03-49}{Lie Algebren zu kanonischen Lie Untergruppen von $\GL(n,\C)$}
Along the lines of the above examples, each of the familiar matrix Lie groups has a Lie algebra presented as a matrix Lie algebra with commutator bracket. This subalgebra is defined in terms of simple conditions as in the examples above, and the chart below lists some of the more common examples.
\begin{center}
\begin{tabular}{lll}
Group & Lie algebra & \begin{tabular}{l}
Conditions defining \\
the Lie algebra \\
\end{tabular} \\
$\mathrm{SL}(n, \mathbb{R})$ & $\mathfrak{s l}(n, \mathbb{R})$ & Trace $(A)=0$ \\
$\mathrm{O}(n)$ & $\mathfrak{o}(n)$ & $A^{t}=-A$ \\
$S \mathrm{O}(n)$ & $\mathfrak{s o}(n)$ & $A^{t}=-A$ \\
$\mathrm{U}(n)$ & $\mathfrak{u}(n)$ & $\bar{A}^{t}=-A$ \\
$S \mathrm{U}(n)$ & $\mathfrak{s u}(n)$ & $\bar{A}^{t}=-A, \operatorname{Trace}(A)=0$ \\
$\mathrm{Sp}(2 n, \mathbb{F})$ & $\mathfrak{s p}(2 n, \mathbb{F})$ & $J A^{t} J=A$ \\
\end{tabular}
\end{center}
In the chart above the matrix $J$ is given by
$$
J=\left(\begin{array}{cc}
0 & I \\
-I & 0
\end{array}\right) .
$$
\end{EXA}


Exercise 5.70. Find the defining conditions for the Lie algebra of the semiorthogonal group $O(p, q)$.

We would like to relate Lie group homomorphisms to Lie algebra homomorphisms.

Proposition 5.71. 


\begin{PROP}{LIE-L09-03-50}{Lie Differential zu einem Lie Gruppen Homomorphismus}
Let $f: G_{1} \rightarrow G_{2}$ be a Lie group homomorphism. The map $T_{e} f: \mathfrak{g}_{1} \rightarrow \mathfrak{g}_{2}$ is a Lie algebra homomorphism called the Lie differential, which is denoted in this context by $d f: \mathfrak{g}_{1} \rightarrow \mathfrak{g}_{2}$.
\end{PROP}

Proof. 

\begin{PROOF}{LIE-L09-03-51}{P: Lie Differential zu einem Lie Gruppen Homomorphismus}
For $v \in \mathfrak{g}_{1}$ and $x \in G$, we have
$$
\begin{aligned}
T_{x} f \cdot L^{v}(x) & =T_{x} f \cdot\left(T_{e} L_{x} \cdot v\right) \\
& =T_{e}\left(f \circ L_{x}\right) \cdot v=T_{e}\left(L_{f(x)} \circ f\right) \cdot v \\
& =T_{e} L_{f(x)}\left(T_{e} f \cdot v\right)=T_{e} L_{f(x)}\left(T_{e} f \cdot v\right) \\
& =L^{d f(v)}(f(x)),
\end{aligned}
$$
so $L^{v} \sim_{f} L^{d f(v)}$. Thus by Proposition 2.84 we have that for any $v, w \in \mathfrak{g}_{1}$,
$$
L^{[v, w]} \sim_{f}\left[L^{d f(v)}, L^{d f(w)}\right] .
$$
In other words, $\left[L^{d f(v)}, L^{d f(w)}\right] \circ f=T f \circ L^{[v, w]}$, which at $e$ gives
$$
[d f(v), d f(w)]=[v, w]
$$
\end{PROOF}

Theorem 5.72. 

\begin{THEO}{LIE-L09-03-52}{Vollständigkeit der Invarianten Vektorfelder und Integralkurven als 1-Parameter Untergruppen}
Invariant vector fields are complete. The integral curves through the identity element are the one-parameter subgroups.
\end{THEO}

Proof. 

\begin{PROOF}{LIE-L09-03-53}{P: Vollständigkeit der Invarianten Vektorfelder und Integralkurven als 1-Parameter Untergruppen}
We prove the left invariant case since the right invariant case is similar. Let $X$ be a left invariant vector field and $c:(a, b) \rightarrow G$ be an integral curve of $X$ with $\dot{c}(0)=X(p)$. Let $a<t_{1}<t_{2}<b$ and choose an element $g \in G$ such that $g c\left(t_{1}\right)=c\left(t_{2}\right)$. Let $\Delta t=t_{2}-t_{1}$, and define $\bar{c}:(a+\Delta t, b+\Delta t) \rightarrow G$ by $\bar{c}(t)=g c(t-\Delta t)$. Then we have
$$
\begin{aligned}
\left.\frac{d}{d t}\right|_{t=0} \bar{c}(t) & =T L_{g} \cdot \dot{c}(t-\Delta t)=T L_{g} \cdot X(c(t-\Delta t)) \\
& =X(g c(t-\Delta t))=X(\bar{c}(t))
\end{aligned}
$$
and so $\bar{c}$ is also an integral curve of $X$. On the intersection ( $a+\Delta t, b$ ) of their domains, $c$ and $\bar{c}$ are equal since they are both integral curves of the same field and since $\bar{c}\left(t_{2}\right)=g c\left(t_{1}\right)=c\left(t_{2}\right)$. Thus we can concatenate the curves to get a new integral curve defined on the larger domain ( $a, b+\Delta t$ ). Since this extension can be done again for the same fixed $\Delta t$, we see that $c$ can be extended to $(a, \infty)$. A similar argument shows that we can extend in the negative direction to get the needed extension of $c$ to $(-\infty, \infty)$.

Next assume that $c$ is the integral curve with $c(0)=e$. The proof that $c(s+t)=c(s) c(t)$ proceeds by considering $\gamma(t)=c(s)^{-1} c(s+t)$. Then $\gamma(0)=e$ and also
$$
\begin{aligned}
\dot{\gamma}(t) & =T L_{c(s)^{-1}} \cdot \dot{c}(s+t)=T L_{c(s)^{-1}} \cdot X(c(s+t)) \\
& =X\left(c(s)^{-1} c(s+t)\right)=X(\gamma(t))
\end{aligned}
$$
By the uniqueness of integral curves we must have $c(s)^{-1} c(s+t)=c(t)$, which implies the result. Conversely, suppose $c: \mathbb{R} \rightarrow G$ is a one-parameter subgroup, and let $X_{e}=\dot{c}(0)$. There is a left invariant vector field $X$ such that $X(e)=X_{e}$, namely $X=L^{X_{e}}$. We must show that the integral curve through $e$ of the field $X$ is exactly $c$. But for this we only need that $\dot{c}(t)=X(c(t))$ for all $t$. We have $c(t+s)=c(t) c(s)$ or $c(t+s)=L_{c(t)} c(s)$. Thus
$$
\dot{c}(t)=\left.\frac{d}{d s}\right|_{0} c(t+s)=\left(T_{c(t)} L\right) \cdot \dot{c}(0)=X(c(t))
$$
\end{PROOF}


\begin{PROP}{LIE-L09-03-54}{Skalierungseigenschaft der Flüsse links-/rechtsinvarianter Felder}
Sei $G$ eine Lie-Gruppe und $v\in\mathfrak g=T_eG$. Bezeichne mit $L^{v}$ (resp.\ $R^{v}$)
das linksinvariante (resp.\ rechtsinvariante) Vektorfeld mit $L^{v}(e)=v$ (resp.\ $R^{v}(e)=v$),
und mit $\varphi_t^{L^{v}}$ (resp.\ $\varphi_t^{R^{v}}$) den zugehörigen Fluss.
Setze
\[
\varphi^{v}(t):=\varphi_t^{L^{v}}(e)\quad\text{(resp.\ }\psi^{v}(t):=\varphi_t^{R^{v}}(e)\text{)}.
\]
Dann gilt für alle $s,t\in\mathbb R$:
\[
\varphi^{v}(s t)=\varphi^{s v}(t)
\qquad\text{(resp.\ }\psi^{v}(s t)=\psi^{s v}(t)\text{)}.
\]
\end{PROP}



\begin{PROOF}{LIE-L09-03-55}{P: Skalierungseigenschaft der Flüsse links-/rechtsinvarianter Felder}
Wir zeigen die Aussage für das linksinvariante Feld; der rechtsinvariante Fall ist analog.
Fixiere $s\in\mathbb R$ und betrachte die Kurve
\[
\gamma(t):=\varphi^{v}(s t)=\varphi_{s t}^{L^{v}}(e).
\]
Dann gilt $\gamma(0)=e$ und nach der Kettenregel
\[
\dot\gamma(0)=\left.\frac{d}{dt}\right|_{t=0}\varphi^{v}(s t)
=\left.\frac{d}{du}\right|_{u=0}\varphi^{v}(u)\cdot \frac{du}{dt}(0)
=s\,\left.\frac{d}{du}\right|_{u=0}\varphi^{v}(u)=s\,v.
\]
Da $L^{v}$ linear in $v$ ist (d.\,h.\ $L^{s v}=s\,L^{v}$), erfüllt $\gamma$ die Anfangswertaufgabe
\[
\dot\gamma(t)=L^{s v}(\gamma(t)),\qquad \gamma(0)=e.
\]
Die eindeutige Lösung dieser ODE ist per Definition $t\mapsto \varphi_t^{L^{s v}}(e)=\varphi^{s v}(t)$.
Also $\varphi^{v}(s t)=\varphi^{s v}(t)$ für alle $t$. Der rechtsinvariante Fall folgt identisch
mit $R^{v}$ anstelle von $L^{v}$.
\end{PROOF}





Proposition 5.73. 


Let $v \in \mathfrak{g}=T_{e} G$. We have the corresponding left invariant field $L^{v}$ and flow $\varphi_{t}^{L^{v}}$. Then, with $\varphi^{v}(t):=\varphi_{t}^{L^{v}}(e)$, we have that
$$
\varphi^{v}(s t)=\varphi^{s v}(t)
$$
A similar statement holds with $R^{v}$ replacing $L^{v}$.


Proof. Let $u=s t$. We have that $\left.\frac{d}{d t}\right|_{t=0} \varphi^{v}(s t)=\left.\frac{d}{d u}\right|_{u=0} \varphi^{v}(u) \frac{d u}{d t}(0)=s v$ and so by uniqueness $\varphi^{v}(s t)=\varphi^{s v}(t)$.


Theorem 5.74. 

\begin{THEO}{LIE-L09-03-56}{Charakterisierung von 1-Parameter Untergruppen durch Flüsse}
Let $G$ be a Lie group. For a smooth curve $c: \mathbb{R} \rightarrow G$ with $c(0)=e$ and $\dot{c}(0)=v$, the following are all equivalent:
\begin{enumerate}
\item $c$ is a one-parameter subgroup with $\dot{c}(0)=v$.
\item $c(t)=\varphi_{t}^{L^{v}}(e)$ for all $t$.
\item $c(t)=\varphi_{t}^{R^{v}}(e)$ for all $t$.
\item $\varphi_{t}^{L^{v}}=R_{c(t)}$ for all $t$.
\item $\varphi_{t}^{R^{v}}=L_{c(t)}$ for all $t$.
\end{enumerate}
\end{THEO}


Proof. 

\begin{PROOF}{LIE-L09-03-57}{P: Charakterisierung von 1-Parameter Untergruppen durch Flüsse}
From Theorem 5.72 we already know that (ii) is equivalent to (i). The proof that (i) is equivalent to (iii) is analogous. Also, (iv) implies (i) since then $\varphi_{t}^{L^{v}}(e)=R_{c(t)}(e)=c(t)$. Now assuming (ii) we show (iv). We have $\dot{c}(t)=L^{v}(e)=v$ and
$$
\begin{aligned}
\left.\frac{d}{d t}\right|_{t=0} g c(t) & =\left.\frac{d}{d t}\right|_{t=0} L_{g}(c(t)) \\
& =T L_{g} v=L^{v}(g) \text { for any } g .
\end{aligned}
$$
In other words,
$$
\left.\frac{d}{d t}\right|_{t=0} R_{c(t)} g=L^{v}(g)
$$
for any $g$. But also, $R_{c(0)} g=e g=g$ and $\varphi_{0}^{L^{v}}(g)=g$ so by uniqueness we have $R_{c(t)} g=\varphi_{t}^{L^{v}}(g)$ for all $g$. We leave the remainder to the reader.

It follows from (iv) and (v) of the previous theorem that the bracket of any left invariant vector field with any right invariant vector field is zero since their flows commute (see Theorem 2.109).
\end{PROOF}







Definition 5.75 (Exponential map). 




\begin{DEF}{LIE-L09-02-01}{Exponential Abbildung von Lie Gruppen}
For any $v \in \mathfrak{g}=T_{e} G$, we have the corresponding left invariant field $L^{v}$ which has an integral curve $t \mapsto \varphi^{v}(t):= \varphi_{t}^{L^{v}}(e)$ through $e$. The map exp : $\mathfrak{g} \rightarrow G$ defined by exp : $v \mapsto \varphi^{v}(1)$ is referred to as the exponential map.

Thus for $s, t \in \mathbb{R}$ and $v \in \mathfrak{g}$ we have
$$
\begin{aligned}
\exp ((s+t) v) & =\exp (s v) \exp (t v) \\
\exp (-t v) & =(\exp (t v))^{-1}
\end{aligned}
$$
Note that we usually use the same symbol "exp" for the exponential map of any Lie group, but we may also write $\exp ^{G}$ to indicate that the group is $G$. By Proposition 5.73 we have
$$
\exp (t v)=\varphi^{t v}(1)=\varphi^{v}(t)=\varphi_{t}^{L^{v}}(e)
$$
\end{DEF}






Thus by Theorem 5.74 we obtain the following:

Proposition 5.76. 



\begin{PROP}{LIE-L09-02-02}{Exponentialkurve als Integralkurve linksinvarianter Felder}
Sei $G$ eine Lie-Gruppe und $v\in\mathfrak g=T_eG$.
Dann ist die Kurve
\[
\gamma_v:\mathbb R\to G,\qquad \gamma_v(t):=\exp(t\,v),
\]
eine \emph{eindimensionale Untergruppe} (ein Einparameteruntergruppe) von $G$
und zugleich die Integralkurve des linksinvarianten Vektorfeldes $L^{v}$
durch $e$:
\[
\gamma_v(0)=e,\qquad \dot{\gamma}_v(t)=L^{v}(\gamma_v(t))\ \ \text{für alle }t.
\]
\end{PROP}



\begin{PROOF}{LIE-L09-02-03}{P: Exponentialkurve als Integralkurve linksinvarianter Felder}
\emph{(1) Einparameteruntergruppe.)}
Die Exponentialabbildung $\exp:\mathfrak g\to G$ erfüllt $\exp(0)=e$ und
\[
\exp(t v)\,\exp(s v)=\exp((t+s)v)\quad\text{für alle }s,t\in\mathbb R,
\]
also ist $t\mapsto \exp(tv)$ ein Gruppenhomomorphismus $\mathbb R\to G$.
Damit ist $\gamma_v$ eine Einparameteruntergruppe.

\medskip
\emph{(2) Integralkurve von $L^v$.)}
Für $t\in\mathbb R$ gilt mit $L_g$ der Linksmultiplikation:
\[
\dot{\gamma}_v(t)
=\left.\frac{d}{ds}\right|_{s=0}\exp\bigl((t+s)v\bigr)
=\left.\frac{d}{ds}\right|_{s=0}\bigl(\exp(tv)\,\exp(sv)\bigr)
= T_e L_{\exp(tv)}\!\left(\left.\frac{d}{ds}\right|_{s=0}\exp(sv)\right).
\]
Da $\left.\frac{d}{ds}\right|_{s=0}\exp(sv)=v$ und $L^v$ durch
$L^v(g)=T_e L_g(v)$ definiert ist, folgt
\[
\dot{\gamma}_v(t)=T_e L_{\exp(tv)}(v)=L^v(\gamma_v(t)).
\]
Zudem ist $\gamma_v(0)=\exp(0)=e$. Damit ist $\gamma_v$ die eindeutige
Integralkurve von $L^v$ durch $e$.
\end{PROOF}





For $v \in \mathfrak{g}$, the map $\mathbb{R} \rightarrow G$ given by $t \mapsto \exp (t v)$ is the one-parameter subgroup that is the integral curve of $L^{v}$.



Lemma 5.77. 

\begin{LEM}{LIE-L09-02-04}{Exponentialabbildung ist glatt}
The map $\exp : \mathfrak{g} \rightarrow G$ is smooth.
\end{LEM}


Proof. 

\begin{PROOF}{LIE-L09-02-05}{P: Exponentialabbildung ist glatt}
Consider the map $\mathbb{R} \times G \times \mathfrak{g} \rightarrow G \times \mathfrak{g}$ given by
$$
(t, g, v) \mapsto(g \cdot \exp (t v), v)
$$
This map is easily seen to be the flow on $G \times \mathfrak{g}$ of the vector field $\widetilde{X}$ : $(g,v) \mapsto\left(L^{v}(g), 0\right)$ and so is smooth. The restriction of this smooth flow to the submanifold $\{1\} \times\{e\} \times \mathfrak{g}$ is $(1, e, v) \mapsto(\exp (v), v)$ and is also smooth. This clearly implies that exp is smooth also.
\end{PROOF}

Note that $\exp (0)=e$. In the following theorem, we use the canonical identification of the tangent space of $T_{e} G$ at the zero element (that is, $\left.T_{0}\left(T_{e} G\right)\right)$ with $T_{e} G$ itself.


Theorem 5.78. 


The tangent map of the exponential map $\exp : \mathfrak{g} \rightarrow G$ is the identity at $0 \in T_{e} G=\mathfrak{g}$, and $\exp$ is a diffeomorphism of some neighborhood of the origin onto its image in $G$,
$$
T_{e} \exp =\operatorname{id}: T_{e} G \rightarrow T_{e} G
$$


Proof. 

By Lemma 5.77, we know that $\exp : \mathfrak{g} \rightarrow G$ is a smooth map. Also, $\left.\frac{d}{d t}\right|_{0} \exp (t v)=v$, which means that the tangent map is $v \mapsto v$. If the reader thinks through the definitions carefully, he or she will discover that we have here used the identification of $\mathfrak{g}$ with $T_{0} \mathfrak{g}$.

From the definitions and Theorem 5.74 we have
$$
\begin{aligned}
\varphi_{t}^{L^{v}}(p) & =p \exp t v \\
\varphi_{t}^{L^{v}}(p) & =(\exp t v) p
\end{aligned}
$$
for all $v \in \mathfrak{g}$, all $t \in \mathbb{R}$ and all $p \in G$.


Proposition 5.79. 

For a (Lie group) homomorphism $f: G_{1} \rightarrow G_{2}$, the following diagram commutes:\\
\includegraphics[scale=0.25, center]{2025_10_09_265d7e1a22c89f79c21eg-230}

Proof. 

For $v$ in the Lie algebra of $G_{1}$, the curve $t \mapsto f\left(\exp ^{G_{1}}(t v)\right)$ is clearly a one-parameter subgroup. Also,
$$
\left.\frac{d}{d t}\right|_{0} f\left(\exp ^{G_{1}}(t v)\right)=d f(v)
$$
and so by uniqueness of integral curves $f\left(\exp ^{G_{1}}(t v)\right)=\exp ^{G_{2}}(t d f(v))$.



The Lie algebra of a Lie group and the group itself are closely related in many ways, and the exponential map is often the key to understanding the connection. One simple observation is that by Theorem 5.6, if $G$ is a connected Lie group, then for any open neighborhood $V \subset \mathfrak{g}$ of 0 the group generated by $\exp (V)$ is all of $G$.




\begin{THEO}{LIE-L09-02-06}{Differenzierbarkeit und Tangentialabbildung der Exponentialabbildung}
Sei $G$ eine Lie‐Gruppe mit Lie‐Algebra $\mathfrak g=T_eG$.
Dann gilt:
\begin{enumerate}
\item Die Exponentialabbildung
\[
\exp:\mathfrak g\longrightarrow G
\]
ist glatt.
\item Der Tangentialoperator im Ursprung ist die Identität:
\[
T_0\exp=\operatorname{id}_{\mathfrak g}:T_0\mathfrak g\longrightarrow T_eG.
\]
\item Es existiert eine Umgebung $U\subset\mathfrak g$ von $0$, sodass
$\exp|_U:U\to \exp(U)\subset G$ ein Diffeomorphismus auf sein Bild ist.
\end{enumerate}
\end{THEO}



\begin{PROOF}{LIE-L09-02-07}{P: Differenzierbarkeit und Tangentialabbildung der Exponentialabbildung}
Nach Lemma~5.77 ist $\exp$ eine glatte Abbildung.
Für $v\in\mathfrak g$ gilt
\[
\left.\frac{d}{dt}\right|_{t=0}\exp(tv)=v,
\]
wobei $\mathfrak g$ kanonisch mit $T_0\mathfrak g$ identifiziert wird.
Daraus folgt unmittelbar
\[
T_0\exp(v)=v,\qquad \text{also }T_0\exp=\operatorname{id}_{\mathfrak g}.
\]
Die Existenz einer Umgebung $U$ von $0$, auf der $\exp$ ein lokaler Diffeomorphismus
ist, folgt aus dem Satz über den Umkehrsatz für glatte Abbildungen, da
$T_0\exp=\operatorname{id}$ invertierbar ist.
\end{PROOF}



\begin{PROP}{LIE-L09-02-08}{Verhalten der Exponentialabbildung unter Lie-Gruppenhomomorphismen}
Seien $f:G_1\to G_2$ ein Lie‐Gruppenhomomorphismus und
$\mathfrak g_1=T_{e_1}G_1$, $\mathfrak g_2=T_{e_2}G_2$ die zugehörigen Lie‐Algebren.
Dann kommutiert das folgende Diagramm:
\[
\begin{array}{ccc}
\mathfrak g_1 & \xrightarrow{\quad df\quad} & \mathfrak g_2 \\[6pt]
\exp_{G_1}\!\!\downarrow &  & \downarrow\!\!\exp_{G_2}\\[6pt]
G_1 & \xrightarrow{\quad f\quad} & G_2
\end{array}
\]
d.\,h.
\[
f\bigl(\exp_{G_1}(v)\bigr)=\exp_{G_2}\bigl(df(v)\bigr)
\qquad\text{für alle }v\in\mathfrak g_1.
\]
\end{PROP}



\begin{PROOF}{LIE-L09-02-09}{P: Verhalten der Exponentialabbildung unter Lie-Gruppenhomomorphismen}
Für $v\in\mathfrak g_1$ ist die Kurve
$t\mapsto f\bigl(\exp_{G_1}(t v)\bigr)$ eine glatte Einparameteruntergruppe von $G_2$,
da $f$ ein Gruppenhomomorphismus ist.
Wir haben
\[
\left.\frac{d}{dt}\right|_{t=0}f\bigl(\exp_{G_1}(t v)\bigr)
=df(v)\in T_{e_2}G_2=\mathfrak g_2.
\]
Nach der Eindeutigkeit der Integralkurven linksinvarianter Felder folgt daher
\[
f\bigl(\exp_{G_1}(t v)\bigr)=\exp_{G_2}(t\,df(v))
\quad\text{für alle }t\in\mathbb R,
\]
und damit insbesondere für $t=1$ die behauptete Kommutativität.
\end{PROOF}



\begin{KORO}{LIE-L09-02-10}{Erzeugung eines zusammenhängenden Lie‐Gruppe durch die Exponentialabbildung}
Ist $G$ eine zusammenhängende Lie‐Gruppe und $V\subset\mathfrak g$ eine offene Umgebung von $0$,
so erzeugt die Teilmenge $\exp(V)\subset G$ die gesamte Gruppe $G$, d.\,h.
\[
\langle\exp(V)\rangle = G.
\]
\end{KORO}



\begin{PROOF}{LIE-L09-02-11}{P: Erzeugung eines zusammenhängenden Lie‐Gruppe durch die Exponentialabbildung}
Nach Theorem~5.6 ist eine zusammenhängende Lie‐Gruppe durch die Exponentialbilder
einer beliebigen Umgebung des Nullvektors in $\mathfrak g$ erzeugt.
\end{PROOF}




\begin{THEO}{LIE-L09-02-12}{Exponentialabbildung für lineare Lie‐Gruppen}
Sei $V$ ein endlichdimensionaler (reeller oder komplexer) Vektorraum mit Skalarprodukt
$\langle\cdot,\cdot\rangle$ und induzierter Norm $\|v\|=\sqrt{\langle v,v\rangle}$.
Wir identifizieren $\mathfrak{gl}(V)\cong L(V,V)$ und versehen diesen Raum mit der Operatornorm
\[
\|A\|:=\sup_{\|v\|\neq 0}\frac{\|Av\|}{\|v\|}.
\]
Dann gilt:
\begin{enumerate}
\item Die Reihe
\[
\exp(A):=\sum_{k=0}^{\infty}\frac{1}{k!}A^k
\]
konvergiert für jedes $A\in L(V,V)$ (absolut und gleichmäßig auf beschränkten Mengen)
und definiert die \emph{Exponentialabbildung}
\[
\exp:\mathfrak{gl}(V)\to \GL(V).
\]
\item Für festes $A$ ist die Kurve $\alpha(t)=\exp(tA)$ die eindeutige Lösung der Anfangswertaufgabe
\[
\alpha'(t)=A\,\alpha(t),\qquad \alpha(0)=I,
\]
d.\,h.\ $\alpha$ ist die Integralkurve des linksinvarianten Vektorfeldes, das durch $A$ bestimmt ist.
\item Für alle $s,t\in\mathbb R$ gilt
\[
\exp((s+t)A)=\exp(sA)\exp(tA),
\qquad
\exp(-tA)=(\exp(tA))^{-1}.
\]
\item Kommutieren $A,B\in L(V,V)$, so gilt
\[
\exp(A+B)=\exp(A)\exp(B).
\]
\end{enumerate}
\end{THEO}



\begin{PROOF}{LIE-L09-02-13}{P: Exponentialabbildung für lineare Lie‐Gruppen}
\emph{Zu (1).}
Zunächst gilt aus der Submultiplikativität der Norm $\|AB\|\le \|A\|\|B\|$
durch Induktion $\|A^k\|\le \|A\|^k$ für alle $k\ge 0$.  
Betrachten wir die Partialsummen
\[
s_N:=\sum_{k=0}^{N}\frac{1}{k!}A^k.
\]
Für $N>M$ folgt
\[
\|s_N-s_M\|
=\Bigl\|\sum_{k=M+1}^{N}\frac{1}{k!}A^k\Bigr\|
\le \sum_{k=M+1}^{N}\frac{1}{k!}\|A\|^k.
\]
Da die Reihe $\sum_{k=0}^{\infty}\frac{1}{k!}\|A\|^k$ konvergiert (Exponentialreihe in $\mathbb R$),
ist $\{s_N\}$ eine Cauchy-Folge in dem vollständigen normierten Raum $(L(V,V),\|\cdot\|)$.
Also existiert der Grenzwert $\exp(A):=\lim_{N\to\infty}s_N$.

\smallskip
\emph{Zu (2).}
Für $t\in\mathbb R$ definieren wir
\[
\exp(tA)=\sum_{k=0}^{\infty}\frac{t^k}{k!}A^k.
\]
Die Reihe konvergiert gleichmäßig für $t$ in kompakten Intervallen, also dürfen wir
termweise differenzieren:
\[
\frac{d}{dt}\exp(tA)
=\sum_{k=1}^{\infty}\frac{k\,t^{k-1}}{k!}A^k
=\sum_{k=1}^{\infty}\frac{t^{k-1}}{(k-1)!}A^k
=A\sum_{k=1}^{\infty}\frac{t^{k-1}}{(k-1)!}A^{k-1}
=A\,\exp(tA).
\]
Offensichtlich ist $\exp(0A)=I$. Somit löst $\alpha(t)=\exp(tA)$ die Anfangswertaufgabe
$\alpha'(t)=A\alpha(t)$, $\alpha(0)=I$.
Die Eindeutigkeit folgt aus der Standardtheorie linearer Differentialgleichungen in Banachräumen.

\smallskip
\emph{Zu (3).}
Seien $s,t\in\mathbb R$. Die Kurve $\alpha(t)=\exp(tA)$ erfüllt $\alpha'(t)=A\alpha(t)$ und $\alpha(0)=I$.
Da $A$ konstant ist, gilt
\[
\exp((s+t)A)=\exp(sA)\exp(tA),
\]
was sich direkt aus der Integralkurven-Eigenschaft ergibt:
$\exp((s+t)A)=\exp(sA)\alpha(t)$, und $\alpha(t)=\exp(tA)$ erfüllt dieselbe DGL.
Daraus folgt auch $\exp(-tA)=(\exp(tA))^{-1}$, da $\exp(-tA)\exp(tA)=\exp(0)=I$.

\smallskip
\emph{Zu (4).}
Falls $A$ und $B$ kommutieren, gilt $(A+B)^m=\sum_{j+k=m}\frac{m!}{j!\,k!}A^jB^k$.
Damit folgt
\[
\exp(A+B)
=\sum_{m=0}^{\infty}\frac{1}{m!}(A+B)^m
=\sum_{j,k=0}^{\infty}\frac{1}{j!\,k!}A^jB^k
=\Bigl(\sum_{j=0}^{\infty}\frac{1}{j!}A^j\Bigr)\Bigl(\sum_{k=0}^{\infty}\frac{1}{k!}B^k\Bigr)
=\exp(A)\exp(B).
\]
\end{PROOF}



\begin{KORO}{LIE-L09-02-14}{Logarithmusabbildung in der Nähe der Identität}
Mit der gleichen Norm sei
\[
U:=\{B\in\GL(V)\mid \|B-I\|<1\}.
\]
Dann ist die Potenzreihe
\[
\log(B):=\sum_{k=1}^{\infty}\frac{(-1)^{k+1}}{k}(B-I)^k
\]
konvergent und definiert eine glatte Abbildung
\[
\log:U\to \mathfrak{gl}(V).
\]
\end{KORO}


\begin{PROOF}{LIE-L09-02-17}{P: Logarithmusabbildung in der Nähe der Identität}
Für $\|B-I\|<1$ gilt wegen $\|(B-I)^k\|\le \|B-I\|^k$ und $\frac{1}{k}\le 1$:
\[
\sum_{k=1}^{\infty}\Bigl\|\frac{(-1)^{k+1}}{k}(B-I)^k\Bigr\|
\le \sum_{k=1}^{\infty}\|B-I\|^k
=\frac{\|B-I\|}{1-\|B-I\|}<\infty.
\]
Also konvergiert die Reihe absolut und gleichmäßig auf $\|B-I\|\le r<1$.
Termweise Differenzierbarkeit liefert die Glattheit von $\log$.
\end{PROOF}


\begin{PROP}{LIE-L09-02-15}{Wechselseitigkeit von $\exp$ und $\log$ in der Normumgebung}
Für $\|A\|<\log 2$ gilt $\|\exp(A)-I\|\le e^{\|A\|}-1<1$ und somit $\exp(A)\in U$.
In diesem Fall ist
\[
\log(\exp A)=A.
\]
Umgekehrt gilt für $\|B-I\|<1$:
\[
\exp(\log B)=B.
\]
\end{PROP}



\begin{PROOF}{LIE-L09-02-16}{P: Wechselseitigkeit von $\exp$ und $\log$ in der Normumgebung}
\emph{Zu (1).}
Aus der Exponentialreihe folgt
\[
\|\exp(A)-I\|
=\Bigl\|\sum_{k=1}^{\infty}\frac{1}{k!}A^k\Bigr\|
\le \sum_{k=1}^{\infty}\frac{1}{k!}\|A\|^k
=e^{\|A\|}-1.
\]
Falls $\|A\|<\log 2$, so ist $e^{\|A\|}-1<1$, d.h. $\exp(A)\in U$.
Da sowohl $\exp(A)$ als auch $\log(\exp(A))$ durch absolut konvergente Reihen definiert sind,
dürfen die Summen umgeordnet werden. Wie im komplexen Fall gilt für $|z|<\log 2$ die Identität
$\log(e^z)=z$; die identische formale Berechnung überträgt sich auf $A\in L(V,V)$
durch Funktorreihenentwicklung:
\[
\log(\exp A)=A+\Bigl(\frac{1}{2!}-\frac{1}{2}\Bigr)A^2
+\Bigl(\frac{1}{3!}-\frac{1}{2}+\frac{1}{3}\Bigr)A^3+\cdots=A.
\]

\smallskip
\emph{Zu (2).}
Für $\|B-I\|<1$ konvergieren sowohl $\log(B)$ als auch $\exp(\log(B))$ absolut.
Da $\log$ die Inverse der Exponentialreihe auf dieser Normkugel ist (vgl. Potenzreihen in Banachalgebren),
folgt durch Substitution
\[
\exp(\log B)=B.
\]
Formal ergibt sich dies durch dieselbe Reihe von Umordnungen wie im reellen bzw.\ komplexen Fall
$\exp(\log(1+z))=1+z$ für $|z|<1$.
\end{PROOF}


\begin{CONC}{LIE-L09-02-18}{ Beziehung zur abstrakten Exponentialabbildung}
Ist $H\subset G$ eine Lie‐Untergruppe, so ist die Inklusion $\iota:H\hookrightarrow G$
ein injektiver Lie‐Gruppenhomomorphismus. Nach Proposition~5.79 ist dann
\[
\exp_H(v)=\exp_G(v)\qquad\text{für alle }v\in\mathfrak h=T_eH.
\]
Somit ist die Exponentialabbildung einer linearen Lie‐Gruppe $G\subset\GL(V)$
die Einschränkung der in (1) definierten Abbildung
\[
\exp:\mathfrak{gl}(V)\to\GL(V),
\qquad
A\mapsto\sum_{k=0}^{\infty}\frac{1}{k!}A^k.
\]
\end{CONC}









If $H$ is a Lie subgroup of $G$, then the inclusion $\iota: H \hookrightarrow G$ is an injective homomorphism and Proposition 5.79 tells us that the exponential map on $\mathfrak{h} \subset \mathfrak{g}$ is the restriction of the exponential map on $\mathfrak{g}$. Thus, to understand the exponential map for linear Lie groups, we must understand the exponential map for the general linear group. Let V be a finite-dimensional vector space. It will be convenient to pick an inner product $\langle\cdot, \cdot\rangle$ on V and define the norm of $v \in \mathrm{V}$ by $\|v\|:=\sqrt{\langle v, v\rangle}$. In case V is a complex vector space, we use a Hermitian inner product. We put a norm on the set of linear transformations $L(\mathrm{V}, \mathrm{V})$ by

$$
\|A\|=\sup _{\|v\| \neq 0} \frac{\|A v\|}{\|v\|} .
$$

We have $\|A \circ B\| \leq\|A\|\|B\|$, which implies that $\left\|A^{k}\right\| \leq\|A\|^{k}$. If we use the identification of $\mathfrak{g l}(\mathrm{V})$ with $L(\mathrm{V}, \mathrm{V})$ (or equivalently the identification of $\mathfrak{g} \mathfrak{l}(n, \mathbb{R})$ with the vector space of $n \times n$ matrices $M_{n \times n}$ ), then the exponential map is given by a power series

$$
A \mapsto \exp (A)=\sum_{k=0}^{\infty} \frac{1}{k!} A^{k} .
$$

which can be seen from the following argument: The sequence of partial sums $s_{N}:=\sum_{k=0}^{N} \frac{1}{k!} A^{k}$ is a Cauchy sequence in the normed space $\mathfrak{g} \mathfrak{l}(\mathrm{V})$,

$$
\begin{aligned}
\left\|\sum_{k=0}^{N} \frac{1}{k!} A^{k}-\sum_{k=0}^{M} \frac{1}{k!} A^{k}\right\| & =\left\|\sum_{k=M+1}^{N} \frac{1}{k!} A^{k}\right\| \\
& \leq \sum_{k=M+1}^{N} \frac{1}{k!}\|A\|^{k}
\end{aligned}
$$

From this we see that

$$
\lim _{M, N \rightarrow \infty}\left\|\sum_{k=0}^{N} \frac{1}{k!} A^{k}-\sum_{k=0}^{M} \frac{1}{k!} A^{k}\right\|=0
$$

and so $\left\{s_{N}\right\}$ is a Cauchy sequence. Since $\mathfrak{g} \mathfrak{l}(\mathrm{V})$ together with the given norm is known to be complete, we see that $\sum_{k=0}^{\infty} \frac{1}{k!} A^{k}$ converges. For a fixed $A \in \mathfrak{g} \mathfrak{l}(\mathrm{V})$, the function $\alpha: t \mapsto \alpha(t)=\exp (t A)$ is the unique solution of the initial value problem

$$
\alpha^{\prime}(t)=A \alpha(t), \quad \alpha(0)=A
$$

This can be seen by differentiating term by term:

$$
\begin{aligned}
\frac{d}{d t} \exp (t A) & =\sum_{k=0}^{\infty} \frac{1}{k!} t^{k} A^{k}=\sum_{k=1}^{\infty} \frac{1}{(k-1)!} t^{k-1} A^{k} \\
& =A \sum_{k=1}^{\infty} \frac{1}{(k-1)!} t^{k-1} A^{k-1}=A \exp (t A)
\end{aligned}
$$

Under our identifications, this says that $\alpha$ is the integral curve corresponding to the left invariant vector field determined by $A$. Thus we have a concrete realization of the exponential map for $\mathfrak{g l}(\mathrm{V})$ and, by restriction, each Lie subgroup of $\mathfrak{g l}(\mathrm{V})$. Applying what we know about exponential maps in the abstract setting of a general Lie group we have in this concrete case $\exp ((s+t) A)=\exp (s A) \exp (t A)$ and $\exp (-t A)=(\exp (t A))^{-1}$. Let $A, B \in \mathfrak{g} \mathfrak{l}(\mathrm{V})$. Then

$$
\begin{aligned}
\exp (A) \exp (B) & =\left(\sum_{j=0}^{\infty} \frac{1}{j!} t^{j} A^{j}\right)\left(\sum_{k=0}^{\infty} \frac{1}{k!} t^{k} B^{k}\right) \\
& =\sum_{j=0}^{\infty} \sum_{k=0}^{\infty} \frac{1}{j!k!} t^{j+k} A^{j} B^{k}
\end{aligned}
$$

On the other hand, suppose that $A \circ B=B \circ A$. Then we have

$$
\begin{aligned}
\exp (A+B) & =\sum_{m=0}^{\infty} \frac{1}{m!} t^{m}(A+B)^{m}=\sum_{m=0}^{\infty} \frac{1}{m!}\left(\sum_{j+k=m}^{\infty} \frac{m!}{j!k!} A^{j} B^{k}\right) \\
& =\sum_{j=0}^{\infty} \sum_{k=0}^{\infty} \frac{1}{j!k!} t^{j+k} A^{j} B^{k}
\end{aligned}
$$

Thus in the case where $A$ commutes with $B$, we have

$$
\exp (A+B)=\exp (A) \exp (B)
$$

Since most Lie groups of interest in practice are linear Lie groups, it will pay to understand the exponential map a bit better in this case. Let V be a finite-dimensional vector space equipped with an inner product as before, and take the induced norm on $\mathfrak{g l}(\mathrm{V})$. By Problem 18 we can define a map $\log : U \rightarrow \mathfrak{g} \mathfrak{l}(\mathrm{V})$, where

$$
U=\{B \in \mathrm{GL}(\mathrm{V}):\|B\|<1\}
$$

by using the power series:

$$
\log B:=\sum_{k=1}^{\infty} \frac{(-1)^{k}}{k}(B-I)^{k}
$$

If we compute formally, then for $A \in \mathfrak{g} \mathfrak{l}(\mathrm{V})$,

$$
\begin{aligned}
\log (\exp A)= & \left(A+\frac{1}{2!} A^{2}\right)-\frac{1}{2}\left(A+\frac{1}{2!} A^{2}\right)^{2} \\
& +\frac{1}{3}\left(A+\frac{1}{2!} A^{2}\right)^{3}+\cdots \\
= & A+\left(\frac{1}{2!} A^{2}-\frac{1}{2} A^{2}\right)+\left(\frac{1}{3!} A^{3}-\frac{1}{2} A^{3}+\frac{1}{3} A^{3}\right)+\cdots
\end{aligned}
$$

We will argue that the above makes sense if $\|A\|<\log 2$ and that there must be cancellations in the last line so that $\log (\exp A)=A$. In fact, $\|\exp A-I\| \leq e^{\|A\|}-1$, and so the double series on the first line for $\log (\exp A)$ must converge absolutely if $e^{\|A\|}-1<1$ or if $\|A\|<\log 2$. This means that we may freely rearrange terms and expect the same cancellations as we find for the analogous calculation of $\log (\exp z)$ for complex $z$ with $|z|<\log 2$. But since $\log (\exp z)=z$ for such $z$, we have the desired conclusion. Similarly one may argue that

$$
\exp (\log B)=B \text { if }\|B-I\|<1
$$







Next, we prove a remarkable theorem that shows how an algebraic assumption can have implications in the differentiable category. First we need some notation.

Notation 5.80. If $S$ is any subset of a Lie group $G$, then we define

$$
S^{-1}=\left\{s^{-1}: s \in S\right\}
$$

and for any $x \in G$ we define

$$
x S=\{x s: s \in S\}
$$

Theorem 5.81. 


\begin{THEO}{LIE-L09-02-19}{Charakterisierung von abgeschlossenen Lie Untergruppen über Untergruppen}
An abstract subgroup $H$ of a Lie group $G$ is a (regular) submanifold and hence a closed Lie subgroup if and only if $H$ is a closed set in $G$.
\end{THEO}

Proof. 

\begin{PROOF}{LIE-L09-02-20}{P: Charakterisierung von abgeschlossenen Lie Untergruppen über Untergruppen}
First suppose that $H$ is a (regular) submanifold. Then $H$ is locally closed. That is, every point $x \in H$ has an open neighborhood $U$ such that $U \cap H$ is a relatively closed set in $H$. Let $U$ be such a neighborhood of the identity element $e$. We seek to show that $H$ is closed in $G$. Let $y \in \bar{H}$ and $x \in y U^{-1} \cap H$. Thus $x \in H$ and $y \in x U$. This means that $y \in \bar{H} \cap x U$, and hence $x^{-1} y \in \bar{H} \cap U=H \cap U$. So $y \in H$, and we have shown that $H$ is closed.

Conversely, suppose that $H$ is a closed abstract subgroup of $G$. Since we can always use left/right translation to translate any point to the identity, it suffices to find a single-slice chart at $e$. This will show that $H$ is a regular submanifold. The strategy is to first find $\mathfrak{L}(H)=\mathfrak{h}$ and then to exponentiate a neighborhood of $0 \in \mathfrak{h}$.

First choose any inner product on $T_{e} G$ so we may take norms of vectors in $T_{e} G$. Choose a small neighborhood $\widetilde{U}$ of $0 \in T_{e} G=\mathfrak{g}$ on which $\exp$ is a diffeomorphism, say $\exp : \widetilde{U} \rightarrow U$, and denote the inverse by $\log _{U}: U \rightarrow \widetilde{U}$. Define the set $\widetilde{H}$ in $\widetilde{U}$ by $\widetilde{H}=\log _{U}(H \cap U)$.

Claim. If $h_{n}$ is a sequence in $\widetilde{H}$ converging to zero and such that $u_{n}=h_{n} /\left\|h_{n}\right\|$ converges to $v \in \mathfrak{g}$, then $\exp (t v) \in H$ for all $t \in \mathbb{R}$.

Proof of the claim: Note that $t h_{n} /\left\|h_{n}\right\| \rightarrow t v$ while $\left\|h_{n}\right\|$ converges to zero. But since $\left\|h_{n}\right\| \rightarrow 0$, we must be able to find a sequence $k(n) \in \mathbb{Z}$ such that $k(n)\left\|h_{n}\right\| \rightarrow t$. From this we have $\exp \left(k(n) h_{n}\right)=\exp \left(k(n)\left\|h_{n}\right\| \frac{h_{n}}{\left\|h_{n}\right\|}\right) \rightarrow \exp (t v)$. But by the properties of $\exp$ proved previously, we have that $\exp \left(k(n) h_{n}\right)=\left(\exp \left(h_{n}\right)\right)^{k(n)}$. But also $\exp \left(h_{n}\right) \in H \cap U \subset H$ and so $\left(\exp \left(h_{n}\right)\right)^{k(n)} \in H$. Since $H$ is closed, we have
$$
\exp (t v)=\lim _{n \rightarrow \infty}\left(\exp \left(h_{n}\right)\right)^{k(n)} \in H
$$

Claim. If $W$ is the set of all $s v$ where $s \in \mathbb{R}$ and $v$ can be obtained as a limit $h_{n} /\left\|h_{n}\right\| \rightarrow v$ with $h_{n} \in \widetilde{H}$ and $h_{n} \rightarrow 0$, then $W$ is a vector space.

Proof of the claim: We just need to show that if $h_{n} /\left\|h_{n}\right\| \rightarrow v$ and $h_{n}^{\prime} /\left\|h_{n}^{\prime}\right\| \rightarrow w$ with $h_{n}^{\prime}, h_{n} \in \widetilde{H}$, then there is a sequence of elements $h_{n}^{\prime \prime}$ from $\widetilde{H}$ with $h_{n}^{\prime \prime} \rightarrow 0$ such that
$$
h_{n}^{\prime \prime} /\left\|h_{n}^{\prime \prime}\right\| \rightarrow \frac{v+w}{\|v+w\|}
$$
Using the previous claim, observe that
$$
h(t):=\log _{U}(\exp (t v) \exp (t w))=\left(\log _{U} \circ \mu\right)(\exp (t v), \exp (t w)) .
$$
Here $\mu$ is the group multiplication map. But, by the first claim, $\exp (t v)$ and $\exp (t w)$ are in $H$ for all $t$, and so $h(t)$ is in $\widetilde{H}$ for small $t$. By Exercise 5.41 and the fact that $T_{e} \log =\mathrm{id}$, we have that
$$
\lim _{t \downarrow 0} h(t) / t=h^{\prime}(0)=v+w .
$$
Thus,
$$
\frac{h(t)}{\|h(t)\|}=\frac{h(t) / t}{\|h(t) / t\|} \rightarrow \frac{v+w}{\|v+w\|} .
$$
Now just let $t_{n} \downarrow 0$ and let $h_{n}^{\prime \prime}:=h\left(t_{n}\right)$. Notice that by the first claim, $\exp (W) \subset H$.

Claim. Let $W$ be the set from the last claim. Then $\exp (W)$ contains an open neighborhood of $e$ in $H$.

Proof of the claim: Let $W^{\perp}$ be the orthogonal complement of $W$ with respect to the inner product chosen above. Then we have $T_{e} G=W^{\perp} \oplus W$. It is not difficult to show that the map $\Sigma: W \oplus W^{\perp} \rightarrow G$ defined by
$$
w+x \mapsto \exp (w) \exp (x)
$$
is a diffeomorphism in a neighborhood of the origin in $T_{e} G$. Denote this diffeomorphism by $\psi$. Now suppose that $\exp (W)$ does not contain an open neighborhood of $e$ in $H$. Then we may choose a sequence $h_{n} \rightarrow e$ such that $h_{n}$ is in $H$ but not in $\exp (W)$. But this means that we can choose a corresponding sequence $\left(w_{n}, x_{n}\right) \in W \oplus W^{\perp}$ with $\left(w_{n}, x_{n}\right) \rightarrow 0$ and $\exp \left(w_{n}\right) \exp \left(x_{n}\right) \in H$ and yet, $x_{n} \neq 0$. The space $W^{\perp}$ is closed and the unit sphere in $W^{\perp}$ is compact. After passing to a subsequence, we may assume that $x_{n} /\left\|x_{n}\right\| \rightarrow x \in W^{\perp}$, and of course $\|x\|=1$. Now $\exp \left(w_{n}\right) \in H$ since $\exp (W) \subset H$ and $H$ is at least an algebraic subgroup, so we see that $\exp \left(w_{n}\right) \exp \left(x_{n}\right) \in H$. Thus it must be that $\exp \left(x_{n}\right) \in H$ also and so $x_{n} \in \widetilde{H}$. But $x_{n} \rightarrow 0$ and $x_{n} /\left\|x_{n}\right\| \rightarrow 0$, and so, since we now know that $x_{n} \in \widetilde{H}$, we have that $x \in W$ by definition. This contradicts the fact that $\|x\|=1$ and $x \in W^{\perp}$. Thus $\exp (W)$ must contain a neighborhood of $e$ in $H$.

Finally, we let $O \subset \exp (W)$ be a neighborhood of $e$ in $H$. The set $O$ must be of the form $O=H \cap V$ for some open $V$ in $T_{e} G$ containing 0 . By shrinking $V$ further we obtain a diffeomorphism $\left.\psi\right|_{V}$. The inverse of this diffeomorphism, $\phi: U \rightarrow V$, has the required properties by construction: $\phi(H \cap U)=O \cap W$. We have actually constructed a chart with values in $T_{e} G$, but this is clearly good enough since we can choose a basis of $T_{e} G$ adapted to $W$.
\end{PROOF}

\pagebreak



\subsection*{The Adjoint Representation of a Lie Group}



Definition 5.82. 

\begin{DEF}{LIE-L09-05-01}{Konjugierte und Adjungierte Abbildung zu Elementen einer Lie Gruppe}
Fix an element $g \in G$. The map $C_{g}: G \rightarrow G$ defined by $C_{g}(x)=g x g^{-1}$ is a Lie group automorphism called the conjugation map, and the tangent map $T_{e} C_{g}: \mathfrak{g} \rightarrow \mathfrak{g}$, denoted $\mathrm{Ad}_{g}$, is called the adjoint map.
\end{DEF}

Proposition 5.83. 


\begin{PROP}{LIE-L09-05-02}{Konjugierte Abbildung (bzgl $g\in G$ ist Lie Gruppen Homomorphismus}
The map $C_{g}: G \rightarrow G$ is a Lie group homomorphism. The map $C: g \mapsto C_{g}$ is a Lie group homomorphism $G \rightarrow \operatorname{Aut}(G)$.
\end{PROP}

Proof. 

\begin{PROOF}{LIE-L09-05-03}{P: Konjugierte Abbildung (bzgl $g\in G$ ist Lie Gruppen Homomorphismus}
See Problem 4.
\end{PROOF}





Corollary 5.84. 

\begin{KORO}{LIE-L09-05-04}{Adjungierte Abbildung bzgl. $g\in G$ ist Lie Algebren Homomorphismus}
The map $\operatorname{Ad}_{g}: \mathfrak{g} \rightarrow \mathfrak{g}$ is a Lie algebra homomorphism.
\end{KORO}

\begin{PROOF}{LIE-L09-05-05}{P: Adjungierte Abbildung bzgl. $g\in G$ ist Lie Algebren Homomorphismus}
Using Proposition 5.71.
\end{PROOF}


Lemma 5.85. 


\begin{LEM}{LIE-L09-05-06}{Induzierte glatte Abbildung von $M$ nach $\GL(T_{y_0}N)$}
Let $f: M \times N \rightarrow N$ be a smooth map and define the partial map at $x \in M$ by $f_{x}(y)=f(x, y)$. Suppose that for every $x \in M$ the point $y_{0}$ is fixed by $f_{x}$ :
$$
f_{x}\left(y_{0}\right)=y_{0} \text { for all } x
$$
Then the map $A_{y_{0}}: x \mapsto T_{y_{0}} f_{x}$ is a smooth map from $M$ to $\operatorname{GL}\left(T_{y_{0}} N\right)$.
\end{LEM}


Proof. 

\begin{PROOF}{LIE-L09-05-07}{P: Induzierte glatte Abbildung von $M$ nach $\GL(T_{y_0}N)$}
It suffices to show that $A_{y_{0}}$ composed with an arbitrary coordinate function from an atlas of charts on $\operatorname{GL}\left(T_{y_{0}} N\right)$ is smooth. But GL $\left(T_{y_{0}} N\right)$ has an atlas consisting of a single chart. Namely, choose a basis $v_{1}, v_{2}, \ldots, v_{n}$ of $T_{y_{0}} N$ and let $v^{1}, v^{2}, \ldots, v^{n}$ be the dual basis of $T_{y_{0}}^{*} N$. Then $\chi_{j}^{i}: A \mapsto v^{i}\left(A v_{j}\right)$ is a typical coordinate function. Now we compose:
$$
\chi_{j}^{i} \circ A_{y_{0}}(x)=v^{i}\left(A_{y_{0}}(x) v_{j}\right)=v^{i}\left(T_{y_{0}} f_{x} \cdot v_{j}\right) .
$$
It is enough to show that $T_{y_{0}} f_{x} \cdot v_{j}$ is smooth in $x$. But this is just the composition of the smooth maps $M \rightarrow T M \times T N \cong T(M \times N) \rightarrow T N$ given by
$$
x \mapsto\left(0_{x}, v_{j}\right) \mapsto\left(\partial_{1} f\right)\left(x, y_{0}\right) \cdot 0_{x}+\left(\partial_{2} f\right)\left(x, y_{0}\right) \cdot v_{j}=T_{y_{0}} f_{x} \cdot v_{j}
$$
(Recall the discussion leading up to Lemma 2.28.)
\end{PROOF}


Proposition 5.86. 


\begin{DEF}{LIE-L09-05-10}{Adjungierte Darstellung von Lie Gruppe}
The map $\mathrm{Ad}: g \mapsto \mathrm{Ad}_{g}$ is a Lie group homomorphism $G \rightarrow \mathrm{GL}(\mathfrak{g})$, which is called the adjoint representation of $G$.
\end{DEF}

\begin{PROP}{LIE-L09-05-08}{Adjungierte Darstellung von Lie Gruppe}
The map $\mathrm{Ad}: g \mapsto \mathrm{Ad}_{g}$ is a Lie group homomorphism $G \rightarrow \mathrm{GL}(\mathfrak{g})$, which is called the adjoint representation of $G$.
\end{PROP}

Proof. 

\begin{PROOF}{LIE-L09-05-09}{P: Adjungierte Darstellung von Lie Gruppe}
We have
$$
\begin{aligned}
\operatorname{Ad}\left(g_{1} g_{2}\right) & =T_{e} C_{g_{1} g_{2}}=T_{e}\left(C_{g_{1}} \circ C_{g_{2}}\right) \\
& =T_{e} C_{g_{1}} \circ T_{e} C_{g_{2}}=\operatorname{Ad}_{g_{1}} \circ \operatorname{Ad}_{g_{2}},
\end{aligned}
$$
which shows that Ad is a group homomorphism. The smoothness follows from the previous lemma applied to the map $C:(g, x) \mapsto C_{g}(x)$.
\end{PROOF}

Recall that for $v \in \mathfrak{g}$ we have the associated left invariant vector field $L^{v}$ as well as the right invariant field $R^{v}$. Using this notation, we have

Lemma 5.87. 

\begin{LEM}{LIE-L09-05-11}{Beziehung zwischen Linksinvarainten und Rechtsinvarianten Vektorfeldern zu $g\in\mathfrak{g}$ via Adjungierte Abbildung}
Let $v \in \mathfrak{g}$. Then $L^{v}(x)=R^{\operatorname{Ad}_{x} v}$.
\end{LEM}


Proof. 

\begin{PROOF}{LIE-L09-05-12}{P: Beziehung zwischen Linksinvarainten und Rechtsinvarianten Vektorfeldern zu $g\in\mathfrak{g}$ via Adjungierte Abbildung}
We calculate as follows:
$$
\begin{aligned}
L^{v}(x) & =T_{e}\left(L_{x}\right) \cdot v=T\left(R_{x}\right) T\left(R_{x^{-1}}\right) T_{e}\left(L_{x}\right) \cdot v \\
& =T\left(R_{x}\right) T\left(R_{x^{-1}} \circ L_{x}\right) \cdot v=R^{\operatorname{Ad}(x) v} .
\end{aligned}
$$
\end{PROOF}

We now go one step further and take the differential of Ad.


Definition 5.88. 

\begin{DEF}{LIE-L09-05-13}{Adjungierte Abbildung von Lie Algebren zu Lie Gruppen}
For a Lie group $G$ with Lie algebra $\mathfrak{g}$, define the adjoint representation of $\mathfrak{g}$ as the map ad: $\mathfrak{g} \rightarrow \mathfrak{g} \mathfrak{l}(\mathfrak{g})$ given by
$$
\mathrm{ad}=T_{e} \mathrm{Ad}=d(\mathrm{Ad})
$$
\end{DEF}

Proposition 5.89. 

\begin{PROP}{LIE-L09-05-14}{Lie-Klammer Darstellung von Adjungierter Abbildung zu Lie Algebren}
$\operatorname{ad}(v) w=[v, w]$ for all $v, w \in \mathfrak{g}$.
\end{PROP}


Proof. 

\begin{PROOF}{LIE-L09-05-15}{P: Lie-Klammer Darstellung von Adjungierter Abbildung zu Lie Algebren}
Let $v^{1}, \ldots, v^{n}$ be a basis for $\mathfrak{g}$ so that $\operatorname{Ad}(x) w=\sum a_{i}(x) v^{i}$ for some functions $a_{i}$. Let $c$ be a curve with $\dot{c}(0)=v$. Then we have
$$
\begin{aligned}
\operatorname{ad}(v) w & =\left.\frac{d}{d t}\right|_{t=0} \operatorname{Ad}(c(t)) w=\left.\sum \frac{d}{d t}\right|_{t=0} a_{i}(c(t)) v^{i} \\
& =\sum\left(v a_{i}\right) v^{i}
\end{aligned}
$$
On the other hand, by Lemma 5.87 we have
$$
\begin{aligned}
L^{w}(x) & =R^{\operatorname{Ad}(x) w}=R\left(\sum a_{i}(x) v^{i}\right) \\
& =\sum a_{i}(x) R^{v^{i}}(x)
\end{aligned}
$$
From the fact that the bracket of any left invariant vector field with any right invariant vector field is zero, we have
$$
\left[L^{v}, L^{w}\right]=\left[L^{v}, \sum a_{i} R^{v^{i}}\right]=0+\sum L^{v}\left(a_{i}\right) R^{v^{i}}
$$
Finally, using the equation for $\operatorname{ad}(v) w$ derived above, we have
$$
\begin{aligned}
{[v, w] } & =\left[L^{v}, L^{w}\right](e) \\
& =\sum L^{v}\left(a_{i}\right)(e) R^{v^{i}}(e)=\sum L^{v}\left(a_{i}\right)(e) v^{i} \\
& =\sum\left(v a_{i}\right) v^{i}=\operatorname{ad}(v) w
\end{aligned}
$$
\end{PROOF}


The map ad : $\mathfrak{g} \rightarrow \mathfrak{g} \mathfrak{l}(\mathfrak{g})=\operatorname{End}\left(T_{e} G\right)$ is given as the tangent map at the identity of the map Ad which is a Lie group homomorphism. 

Thus by Proposition 5.71 we obtain the following:

Proposition 5.90. 

\begin{PROP}{LIE-L09-05-16}{Adjungierte Abbildung zur Lie Algebra ist Lie Algebra Homomorphismus}
ad: $\mathfrak{g} \rightarrow \mathfrak{g} \mathfrak{l}(\mathfrak{g})$ is a Lie algebra homomorphism.
\end{PROP}


Since ad is defined as the Lie differential of Ad, Proposition 5.79 tells us that the following diagram commutes for any Lie group $G$ :\\
\includegraphics[scale=0.25, center]{2025_10_09_265d7e1a22c89f79c21eg-237(1)}

On the other hand, for any $g \in G$, the map $C_{g}: x \mapsto g x g^{-1}$ is also a homomorphism, and so Proposition 5.79 applies again, giving the following commutative diagram:\\
\includegraphics[scale=0.25, center]{2025_10_09_265d7e1a22c89f79c21eg-237}

In other words,

$$
\exp \left(t \operatorname{Ad}_{g} v\right)=g \exp (t v) g^{-1}
$$

for any $g \in G, v \in \mathfrak{g}$ and $t \in \mathbb{R}$.

In the case of linear Lie groups $G \subset \mathrm{GL}(\mathrm{V})$, we have identified $\mathfrak{g}$ with a subspace of $\mathfrak{g} \mathfrak{l}(\mathrm{V})$, which is in turn identified with $L(\mathrm{V}, \mathrm{V})$. In this case, the exponential map is given by the power series as explained above. It is easy to show from the power series that $B \circ \exp (t A) \circ B^{-1}=\exp \left(t B \circ A \circ B^{-1}\right)$ for any $A \in \mathfrak{g l}(\mathrm{V})$ and $B \in \mathrm{GL}(\mathrm{V})$. In this special set of circumstances, we have

$$
\operatorname{Ad}_{B} A=B \circ A \circ B^{-1}
$$

This is seen as follows:

$$
\begin{aligned}
\operatorname{Ad}_{B} A & =\left.\frac{d}{d t}\right|_{t=0} B \circ \exp (t A) \circ B^{-1} \\
& =\left.\frac{d}{d t}\right|_{t=0} \exp \left(t B \circ A \circ B^{-1}\right)=B \circ A \circ B^{-1} .
\end{aligned}
$$

Earlier we noted that for a general Lie group we always have Ado exp $=$ exp $\circ \mathrm{ad}$. In the current context of linear Lie groups, this can be written as

$$
\exp (A) \circ B \circ \exp (-A)=\sum_{k=0}^{\infty} \frac{1}{k!}(\operatorname{ad}(A))^{k} B
$$

for any $A \in \mathfrak{g} \mathfrak{l}(\mathrm{V})$ and any $B \in \mathrm{GL}(\mathrm{V})$.



\begin{PROP}{LIE-L09-05-17}{Zusammenhang zwischen $\operatorname{Ad}$, $\operatorname{ad}$ und der Exponentialabbildung}
Da $\operatorname{ad}$ als das Lie-Differential von $\operatorname{Ad}$ definiert ist, kommutiert für jede Lie-Gruppe $G$
das folgende Diagramm:

\includegraphics[scale=0.25, center]{2025_10_09_265d7e1a22c89f79c21eg-237(1)}

Das heißt:
\[
\operatorname{Ad}_{\exp X} = \exp(\operatorname{ad}_X)
\qquad\text{für alle } X\in\mathfrak g.
\]
\end{PROP}



\begin{PROOF}{LIE-L09-05-18}{P: Zusammenhang zwischen $\operatorname{Ad}$, $\operatorname{ad}$ und der Exponentialabbildung}
Nach Proposition~5.79 ist $\operatorname{Ad}:G\to\GL(\mathfrak g)$ ein Lie-Gruppenhomomorphismus.
Für ein solches Homomorphismus $f:G_1\to G_2$ gilt allgemein:
\[
f(\exp^{G_1}X)=\exp^{G_2}\bigl((T_ef)(X)\bigr).
\]
Wenden wir dies auf $f=\operatorname{Ad}$ an, erhalten wir
\[
\operatorname{Ad}(\exp X)=\exp(\operatorname{ad}_X),
\]
wobei $\operatorname{ad}=T_e\operatorname{Ad}$. Dies zeigt die Kommutativität des Diagramms.

\smallskip
Ferner ist für jedes $g\in G$ auch
\[
C_g:G\to G,\qquad C_g(x)=g\,x\,g^{-1},
\]
ein Lie-Gruppenhomomorphismus. Da $\operatorname{Ad}_g=T_eC_g$, kommutiert ebenfalls
\[
\includegraphics[scale=0.25, center]{2025_10_09_265d7e1a22c89f79c21eg-237}
\]
und damit gilt für alle $g\in G$, $v\in\mathfrak g$ und $t\in\mathbb R$:
\[
\exp\!\bigl(t,\operatorname{Ad}_gv\bigr)
= g\,\exp(tv)\,g^{-1}.
\]

\smallskip
\emph{Spezialfall linearer Lie-Gruppen.}
Sei $G\subset\GL(V)$ eine lineare Lie-Gruppe mit Lie-Algebra
$\mathfrak g\subset\mathfrak{gl}(V)=L(V,V)$.
Dann gilt für $A\in\mathfrak{gl}(V)$ und $B\in\GL(V)$:
\[
B\exp(tA)B^{-1}=\exp\bigl(t(BAB^{-1})\bigr),
\quad\text{also}\quad
\operatorname{Ad}_BA=BAB^{-1}.
\]
Dies folgt durch Differentiation bei $t=0$:
\[
\operatorname{Ad}_BA=\left.\frac{d}{dt}\right|_{t=0}B\exp(tA)B^{-1}
=\left.\frac{d}{dt}\right|_{t=0}\exp(tBAB^{-1})=BAB^{-1}.
\]

\smallskip
Schließlich erhält man im Kontext linearer Lie-Gruppen die explizite Reihenentwicklung
\[
\exp(A)\,B\,\exp(-A)
=\sum_{k=0}^\infty\frac{1}{k!}\bigl(\operatorname{ad}(A)\bigr)^k(B),
\qquad A\in\mathfrak{gl}(V),\ B\in\GL(V),
\]
wobei $\operatorname{ad}(A)(B)=[A,B]=AB-BA$. Diese Reihe entspricht der Darstellung
\(\operatorname{Ad}_{\exp A}=e^{\operatorname{ad}_A}\).
\end{PROOF}








We end this section with a statement of the useful Campbell-BakerHausdorff (CBH) formula. A proof may be found in [Helg] or [Mich]. First notice that if $G$ is a Lie group with Lie algebra $\mathfrak{g}$, then for each $X \in \mathfrak{g}$ we have ad $X \in L(\mathfrak{g}, \mathfrak{g})$. We choose a norm on $\mathfrak{g}$, and then we have a natural operator norm and consequent notion of convergence in $L(\mathfrak{g}, \mathfrak{g})$. The analytic function $\frac{\ln z}{z-1}$ is defined by the power series
$$
\frac{\ln z}{z-1}:=\sum_{n=0}^{\infty} \frac{(-1)^{n}}{n+1}(z-1)^{n}
$$
which converges in a small ball about $z=1$. We may then use the corresponding power series to make sense of 
$$(\ln A)(A-1)^{-1}$$ 
for $A \in L(\mathfrak{g}, \mathfrak{g})$ sufficiently close to the identity map $I=\operatorname{id}_{\mathfrak{g}}$.

Theorem 5.91 (CBH Formula). Let $G$ be a Lie group and $\mathfrak{g}$ its Lie algebra. Let $f(z)=\frac{\ln z}{z-1}$ be defined by a power series as above. Then for sufficiently small $x, y \in \mathfrak{g}$,
$$
\exp (x) \exp (y)=\exp (C(x, y))
$$
where
$$
\begin{aligned}
C(x, y) & =y+\int_{0}^{1} f\left(e^{t \operatorname{ad} x} e^{t \operatorname{ad} y}\right) \cdot x d t \\
& =x+y+\sum_{n=1}^{\infty} \frac{(-1)^{n}}{n+1} \int_{0}^{1}\left(\sum_{k, l>0, k+l \geq 1} \frac{t^{k}}{k!l!}(\operatorname{ad} x)^{k}(\operatorname{ad} y)^{l}\right)^{n} x d t
\end{aligned}
$$
It follows from the above that
$$
C(x, y)=\sum_{n=1}^{\infty} C_{n}(x, y),
$$
where $C_{1}(x, y)=x+y, C_{2}(x, y)=\frac{1}{2}[x, y]$ and
$$
C_{3}(x, y)=\frac{1}{12}([[x, y], y]+[[y, x], x]) .
$$

In particular,
$$
\exp (t x) \exp (t y)=\exp \left(t(x+y)+O\left(t^{2}\right)\right)
$$
for any $x, y$ and sufficiently small $t$. One can show, using the above results, that $G$ is abelian if and only if $\mathfrak{g}$ is abelian (i.e. $[x, y]=0$ for all $x, y \in \mathfrak{g}$ ). The CBH formula shows that the Lie algebra structure on $\mathfrak{g}$ locally determines the multiplicative structure on $G$.







\begin{THEO}{LIE-L09-05-19}{Campbell--Baker--Hausdorff-Formel (CBH)}
Sei $G$ eine Lie-Gruppe mit Lie-Algebra $\mathfrak g$. Für jedes $X\in\mathfrak g$
sei $\operatorname{ad} X\in L(\mathfrak g,\mathfrak g)$ gegeben durch $\operatorname{ad}_X(Y)=[X,Y]$.
Wähle eine Norm auf $\mathfrak g$, sodass auf $L(\mathfrak g,\mathfrak g)$
die induzierte Operatornorm und eine Konvergenzstruktur definiert sind.

Die analytische Funktion
\[
f(z)=\frac{\ln z}{z-1}=\sum_{n=0}^\infty\frac{(-1)^n}{n+1}(z-1)^n
\]
konvergiert in einer Umgebung von $z=1$. Damit ist auch
\[
(\ln A)(A-I)^{-1}
\]
für $A\in L(\mathfrak g,\mathfrak g)$ in einer kleinen Umgebung der Identität wohldefiniert.

Dann gilt für hinreichend kleine $x,y\in\mathfrak g$ die \emph{Campbell--Baker--Hausdorff-Formel}
\[
\exp(x)\exp(y)=\exp(C(x,y)),
\]
wobei
\[
\begin{aligned}
C(x,y)
&=y+\int_0^1 f\!\left(e^{t\operatorname{ad} x}\,e^{t\operatorname{ad} y}\right)\cdot x\,dt\\
&=x+y+\sum_{n=1}^\infty\frac{(-1)^n}{n+1}\int_0^1
\left(\sum_{\substack{k,l>0\\k+l\ge1}}\frac{t^k}{k!\,l!}(\operatorname{ad} x)^k(\operatorname{ad} y)^l\right)^{\!n}\!x\,dt.
\end{aligned}
\]
Es folgt daraus die formale Reihenentwicklung
\[
C(x,y)=\sum_{n=1}^\infty C_n(x,y),
\]
mit den Anfangsgliedern
\[
C_1(x,y)=x+y,\qquad
C_2(x,y)=\tfrac{1}{2}[x,y],\qquad
C_3(x,y)=\tfrac{1}{12}\bigl([[x,y],y]+[[y,x],x]\bigr).
\]

Insbesondere gilt für kleines $t\in\mathbb R$:
\[
\exp(tx)\exp(ty)=\exp\!\left(t(x+y)+O(t^2)\right).
\]
\end{THEO}




\begin{KORO}{LIE-L09-05-20}{Abelsche Lie-Gruppen und Lie-Algebren}
Eine Lie-Gruppe $G$ ist genau dann abelsch, wenn ihre Lie-Algebra $\mathfrak g$
abelsch ist, d.\,h.
\[
[x,y]=0\quad\text{für alle }x,y\in\mathfrak g.
\]
\end{KORO}



\begin{CONC}{LIE-L09-05-21}{Interpretation der CBH-Formel}
Die Campbell--Baker--Hausdorff-Formel zeigt, dass die Lie-Algebra-Struktur
$\mathfrak g$ lokal die Multiplikationsstruktur auf $G$ bestimmt.
\end{CONC}






\pagebreak

\subsection*{The Maurer-Cartan Form}


Let $G$ be a Lie group, and for each $g \in G$, define $\omega_{G}(g): T_{g} G \rightarrow \mathfrak{g}$ and $\omega_{G}^{\text {right }}(g): T_{g} G \rightarrow \mathfrak{g}$ by
$$
\omega_{G}(g)\left(X_{g}\right)=T L_{g^{-1}} \cdot X_{g}
$$
and
$$
\omega_{G}^{\text {right }}(g)\left(X_{g}\right)=T R_{g^{-1}} \cdot X_{g} .
$$
The maps $\omega_{G}: g \mapsto \omega_{G}(g)$ and $\omega_{G}^{\text {right }}: g \mapsto \omega_{G}^{\text {right }}(g)$ are $\mathfrak{g}$-valued 1-forms called the left Maurer-Cartan form and right Maurer-Cartan form respectively. We can view $\omega_{G}$ and $\omega_{G}^{\text {right }}$ as maps $T G \rightarrow \mathfrak{g}$. For example, $\omega_{G}\left(X_{g}\right):=\omega_{G}(g)\left(X_{g}\right)$.


As we have seen, $\operatorname{GL}(n)$ is an open set in a vector space, and so its tangent bundle is trivial, $T \mathrm{GL}(n) \cong \mathrm{GL}(n) \times M_{n \times n}$ (recall Definition 2.58). A general abstract Lie group $G$ is not an open subset of a vector space, but we are still able to show that $T G$ is trivial. There are two such trivializations obtained from the Maurer-Cartan forms. These are $\operatorname{triv}_{L}: T G \rightarrow G \times \mathfrak{g}$ and $\operatorname{triv}_{R}: T G \rightarrow G \times \mathfrak{g}$ defined by

$$
\begin{aligned}
& \operatorname{triv}_{L}\left(v_{g}\right)=\left(g, \omega_{G}\left(v_{g}\right)\right) \\
& \operatorname{triv}_{R}\left(v_{g}\right)=\left(g, \omega_{G}^{\text {right }}\left(v_{g}\right)\right)
\end{aligned}
$$

for $v_{g} \in T_{g} G$. Observe that $\operatorname{triv}_{L}^{-1}(g, v)=L^{v}(g)$ and $\operatorname{triv}_{R}^{-1}(g, v)=R^{v}(g)$. It is easy to check that $\operatorname{triv}_{L}$ and $\operatorname{triv}_{R}$ are trivializations in the sense of Definition 2.58. 


Thus we have the following:

Proposition 5.92. The tangent bundle of a Lie group is trivial: $T G \cong G \times \mathfrak{g}$.

We will refer to $\operatorname{triv}_{L}$ and $\operatorname{triv}_{R}$ as the (left and right) Maurer-Cartan trivializations. 



\begin{DEF}{LIE-L09-06-01}{Linke und rechte Maurer--Cartan-Form}
Sei $G$ eine Lie-Gruppe mit Lie-Algebra $\mathfrak g=T_eG$.
Für jedes $g\in G$ definieren wir lineare Abbildungen
\[
\omega_G(g):T_gG\longrightarrow\mathfrak g,
\qquad
\omega_G^{\mathrm{right}}(g):T_gG\longrightarrow\mathfrak g
\]
durch
\[
\omega_G(g)(X_g):=T_gL_{g^{-1}}\cdot X_g,
\qquad
\omega_G^{\mathrm{right}}(g)(X_g):=T_gR_{g^{-1}}\cdot X_g.
\]
Die Abbildungen
\[
\omega_G:T G\longrightarrow\mathfrak g,\qquad
\omega_G^{\mathrm{right}}:T G\longrightarrow\mathfrak g
\]
sind $\mathfrak g$-wertige 1-Formen auf $G$ und heißen
\emph{linke} bzw. \emph{rechte Maurer--Cartan-Form}.
\end{DEF}



\begin{DEF}{LIE-L09-06-02}{Maurer--Cartan-Trivialisierungen des Tangentialbündels}
Mit Hilfe der Maurer--Cartan-Formen erhalten wir zwei kanonische
Trivialisierungen des Tangentialbündels $T G$:
\[
\begin{aligned}
\operatorname{triv}_L:T G&\longrightarrow G\times\mathfrak g,&
v_g&\longmapsto\bigl(g,\omega_G(v_g)\bigr),\\
\operatorname{triv}_R:T G&\longrightarrow G\times\mathfrak g,&
v_g&\longmapsto\bigl(g,\omega_G^{\mathrm{right}}(v_g)\bigr).
\end{aligned}
\]
Ihre Inversen sind gegeben durch
\[
\operatorname{triv}_L^{-1}(g,v)=L^v(g),\qquad
\operatorname{triv}_R^{-1}(g,v)=R^v(g),
\]
wobei $L^v$ (resp.\ $R^v$) das links- (resp.\ rechts-)invariante
Vektorfeld mit $L^v(e)=v$ (resp.\ $R^v(e)=v$) bezeichnet.
\end{DEF}



\begin{PROP}{LIE-L09-06-03}{Trivialität des Tangentialbündels einer Lie-Gruppe}
Das Tangentialbündel einer Lie-Gruppe $G$ ist trivial:
\[
T G \cong G\times\mathfrak g.
\]
Die Abbildungen $\operatorname{triv}_L$ und $\operatorname{triv}_R$
sind Bündelisomorphismen und heißen die \emph{linke} bzw. \emph{rechte
Maurer--Cartan-Trivialisierung}.
\end{PROP}






How do these two trivializations compare? There is no special reason to prefer left multiplication. We could have used right invariant vector fields as our means of producing the Lie algebra, and the whole theory would work "on the other side", so to speak. The bridge between left and right is the adjoint map:

Lemma 5.93 (Left-right lemma). 


\begin{LEM}{LIE-L09-06-04}{Links-Rechts Lemma für Trivialisierung}
For any $v \in \mathfrak{g}$, and $g \in G$ we have
$$
\operatorname{triv}_{R} \circ \operatorname{triv}_{L}^{-1}(g, v)=\left(g, \operatorname{Ad}_{g}(v)\right)
$$
\end{LEM}

Proof. 

\begin{PROOF}{LIE-L09-06-05}{P: Links-Rechts Lemma für Trivialisierung}
We compute:
$$
\begin{aligned}
& \operatorname{triv}_{R} \circ \operatorname{triv}_{L}^{-1}(g, v) \\
& =\left(g, T R_{g^{-1}} T L_{g} v\right)=\left(g, T\left(R_{g^{-1}} L_{g}\right) \cdot v\right) \\
& =\left(g, T C_{g} \cdot v\right)=\left(g, \operatorname{Ad}_{g}(v)\right)
\end{aligned}
$$
\end{PROOF}

\begin{CONC}{LIE-L09-06-06}{Prinzipielle Äquivalenz der Trivialisierungen}
Es besteht kein prinzipieller Grund, zwischen linker und rechter
Trivialisierung zu unterscheiden; beide liefern äquivalente Beschreibungen.
So könnte man die Lie-Algebra auch über rechtsinvariante Vektorfelder
definieren. Das obige Lemma zeigt, dass die Adjoint-Abbildung
die Brücke zwischen beiden Darstellungen bildet.
\end{CONC}


\begin{DEF}{LIE-L09-06-07}{Linke Identifikation des Tangentialbündels durch linke Maurer-Cartan Trivialisierung}
Im Folgenden identifizieren wir $T G$ mittels der \emph{linken
Maurer--Cartan-Trivialisierung}
\[
v_g\longmapsto(g,\omega_G(v_g))=(g,\,T L_g^{-1}v_g)
\]
mit $G\times\mathfrak g$, sofern nichts anderes angegeben ist.
\end{DEF}




\begin{CONC}{LIE-L09-06-08}{3 Trivialisierungen des Tangentialbündels von $\GL$}
Für die allgemeine lineare Gruppe existieren drei natürliche
Trivialisierungen des Tangentialbündels:
die Standardtrivialisierung von $T\GL(n)$
sowie die linken und rechten Maurer--Cartan-Trivialisierungen.
Die Standardtrivialisierung ist in der Tat die Beschränkung
der Maurer--Cartan-Trivialisierung der abelschen Lie-Gruppe
$(M_{n\times n},+)$ auf $T\GL(n)$.
\end{CONC}

It is often convenient to identify the tangent bundle $T G$ of a Lie group $G$ with $G \times \mathfrak{g}$. Of course we must specify which of the two trivializations described above is being invoked. Unless indicated otherwise we shall use the "left version" described above: $v_{g} \mapsto\left(g, \omega_{G}\left(v_{g}\right)\right)=\left(g, T L_{g}^{-1}\left(v_{g}\right)\right)$.

Warning: It must be realized that we now have three natural ways to trivialize the tangent bundle of the general linear group. In fact, the usual one which we introduced earlier is actually the restriction to $T \mathrm{GL}(n)$ of the Maurer-Cartan trivialization of the abelian Lie group ( $M_{n \times n},+$ ).



In order to use the (left) Maurer-Cartan trivialization as an identification effectively, we need to find out how a few basic operations look when this identification is imposed.

The picture obtained from using the trivialization produced by the left Maurer-Cartan form:



\begin{PROP}{LIE-L09-06-09}{Darstellung grundlegender Operationen unter der linken MC-Trivialisierung}
Sei $T G\cong G\times\mathfrak g$ durch die linke Maurer--Cartan-Form
identifiziert. Dann nehmen die folgenden Standardabbildungen
die nachstehenden Formen an:

(1) \textbf{Tangentialabbildung der Linkstranslation.} The tangent map of the left translation $T L_{g}: T G \rightarrow T G$ takes the form " $T L_{g}$ " : $(x, v) \mapsto(g x, v)$. Indeed, the following diagram commutes:
\[
\begin{tikzcd}[row sep=large, column sep=large]
T G \arrow[r, "T L_g"] \arrow[d] & T G \arrow[d] \\
G \times \mathfrak{g} \arrow[r, "{\text{``}T L_g\text{''}}"] & G \times \mathfrak{g}
\end{tikzcd}
\]
where elementwise we have
\[
\begin{tikzcd}
v_x \arrow[r, "TL_g"] \arrow[d] & TL_g \cdot v_x \arrow[d] \\
(x,\, TL_x^{-1} v_x) = (x,\, v) \arrow[r] & (gx,\, TL_{gx}^{-1} TL_g v_x) = (gx,\, v)
\end{tikzcd}
\]


(2) \textbf{Tangentialabbildung der Gruppenmultiplikation.} The tangent map of multiplication: This time we will invoke two identifications. First, group multiplication is a map $\mu: G \times G \rightarrow G$, and so on the tangent level we have a map $T(G \times G) \rightarrow T G$. Recall that we have a natural isomorphism $T(G \times G) \cong T G \times T G$ given by $T \pi_{1} \times T \pi_{2}:\left(v_{(x, y)}\right) \mapsto\left(T \pi_{1} \cdot v_{(x, y)}, T \pi_{2} \cdot v_{(x, y)}\right)$. If we also identify $T G$ with $G \times \mathfrak{g}$, then $T G \times T G \cong(G \times \mathfrak{g}) \times(G \times \mathfrak{g})$, and we end up with the following "version" of $T \mu$ :
$$
\begin{aligned}
& \text { " } T \mu \text { " }:(G \times \mathfrak{g}) \times(G \times \mathfrak{g}) \rightarrow G \times \mathfrak{g}, \\
& \text { " } T \mu \text { " }:((x, v),(y, w)) \mapsto\left(x y, T R_{y} v+T L_{x} w\right)
\end{aligned}
$$

(3) \textbf{Linke Maurer--Cartan-Form.} The (left) Maurer-Cartan form is a map $\omega_{G}: T G \rightarrow T_{e} G=\mathfrak{g}$, and so there must be a "version", " $\omega_{G}$ ", that uses the identification $T G \cong G \times \mathfrak{g}$. In fact, the map we seek is just projection:
$$
\text { " } \omega_{G} \text { " : }(x, v) \mapsto v .
$$

(4) \textbf{Rechte Maurer--Cartan-Form.} The right Maurer-Cartan form is a little more complicated since we are currently using the isomorphism $T G \cong G \times \mathfrak{g}$ obtained from the left Maurer-Cartan form. From Lemma 5.93 we obtain:
$$
\text { " } \omega_{G}^{\text {right }} \text { " }:(x, v) \mapsto \operatorname{Ad}_{g}(v) .
$$
The adjoint map is nearly the same thing as the right MaurerCartan form if we decide to use the left trivialization $T G \cong G \times \mathfrak{g}$ as an identification.


(5) \textbf{Vektorfelder.} A vector field $X \in \mathfrak{X}(G)$ should correspond to a section of the projection $G \times \mathfrak{g} \rightarrow G$ which must have the form $\overleftrightarrow{X}: x \mapsto\left(x, F^{X}(x)\right)$ for some smooth $\mathfrak{g}$-valued function $F^{X} \in C^{\infty}(G ; \mathfrak{g})$. It is an easy consequence of the definitions that $F^{X}(x)=\omega_{G}(X(x))= T L_{x}^{-1} \cdot X(x)$. Under this identification, a left invariant vector field becomes a constant section of $G \times \mathfrak{g}$. For example, if $X$ is left invariant, then the corresponding constant section is $x \mapsto(x, X(e))$.
\end{PROP}













Exercise 5.94. 

Refer to 2 above. Show that the map " $T \mu$ " defined so that the diagram below commutes is $((x, v),(y, w)) \mapsto\left(x y, T R_{y} v+T L_{x} w\right)$.
\[
\begin{tikzcd}
T(G \times G) \arrow[r, "T\mu"] \arrow[d] & TG \arrow[d] \\
(G \times \mathfrak{g}) \times (G \times \mathfrak{g}) \arrow[r, "{\text{``}T\mu\text{''}}"] & G \times \mathfrak{g}
\end{tikzcd}
\]



We have already seen that using available identifications the Lie algebra of a matrix groups and associated formulas take a concrete form. We now consider the Maurer-Cartan form in the case of matrix groups. Suppose that $G$ is a Lie subgroup of $\operatorname{GL}(n)$ and consider the coordinate functions $x_{j}^{i}$ on GL $(n)$ defined by $x_{j}^{i}(A)=a_{j}^{i}$ where $A=\left[a_{j}^{i}\right]$. We have the associated 1forms $d x_{j}^{i}$. Both the functions $x_{j}^{i}$ and the forms $d x_{j}^{i}$ restrict to $G$. Denoting these restrictions by the same symbols, the Maurer-Cartan form can be expressed as

$$
\omega_{G}=\left[x_{j}^{i}\right]^{-1}\left[d x_{j}^{i}\right] .
$$

One sometimes sees the shorthand, $g^{-1} d g$, which is a bit cryptic. Our goal is to understand this expression better. First, from a practical point of view we think of it as follows: An element $v_{g}$ of the tangent space $T_{g} G$ is also an element of $T_{g} \mathrm{GL}(n)$ and so can be expressed in terms of the vectors $\partial /\left.\partial x_{j}^{i}\right|_{g}$, say

$$
v_{g}=\left.\sum v_{j}^{i} \frac{\partial}{\partial x_{j}^{i}}\right|_{g} .
$$

Then,

$$
\omega_{G}\left(v_{g}\right)=\left[x_{j}^{i}\right]^{-1}\left[d x_{j}^{i}\right]\left(v_{g}\right)=\left[g_{j}^{i}\right]^{-1}\left[v_{j}^{i}\right],
$$

where $g=\left[g_{j}^{i}\right]$ and the matrix $\left[g_{j}^{i}\right]^{-1}\left[v_{j}^{i}\right]$ is interpreted as an element of the Lie algebra $\mathfrak{g}$. For instance, if $G=\mathrm{SO}(2)$, then the Maurer-Cartan form is given by

$$
\begin{aligned}
{\left[\begin{array}{cc}
x_{2}^{2} & -x_{2}^{1} \\
-x_{1}^{2} & x_{1}^{1}
\end{array}\right]\left[\begin{array}{cc}
d x_{1}^{1} & d x_{2}^{1} \\
d x_{1}^{2} & d x_{2}^{2}
\end{array}\right] } & =\left[\begin{array}{cc}
x_{2}^{2} d x_{1}^{1}-x_{2}^{1} d x_{1}^{2} & x_{2}^{2} d x_{2}^{1}-x_{2}^{1} d x_{2}^{2} \\
-x_{1}^{2} d x_{1}^{1}+x_{1}^{1} d x_{1}^{2} & -x_{1}^{2} d x_{2}^{1}+x_{1}^{1} d x_{2}^{2}
\end{array}\right] \\
& =\left[\begin{array}{cc}
x_{1}^{1} d x_{1}^{1}-x_{2}^{1} d x_{1}^{2} & x_{1}^{1} d x_{2}^{1}-x_{2}^{1} d x_{2}^{2} \\
x_{2}^{1} d x_{1}^{1}+x_{1}^{1} d x_{1}^{2} & x_{2}^{1} d x_{2}^{1}+x_{1}^{1} d x_{2}^{2}
\end{array}\right]
\end{aligned}
$$

since on $\operatorname{SO}(2)$ we have $x_{1}^{1}=x_{2}^{2}, x_{2}^{1}=-x_{1}^{2}$ and $x_{1}^{1} x_{2}^{2}-x_{2}^{1} x_{1}^{2}=1$. But this can be further simplified. If we let $v_{g} \in T_{g} \mathrm{SO}(2)$, then $\omega_{G}\left(v_{g}\right)$ is in $\mathfrak{s o}(2)$ and so must be antisymmetric. We conclude that

$$
\omega_{\mathrm{SO}(2)}=\left[\begin{array}{cc}
0 & x_{1}^{1} d x_{2}^{1}-x_{2}^{1} d x_{2}^{2} \\
-x_{1}^{1} d x_{2}^{1}+x_{2}^{1} d x_{2}^{2} & 0
\end{array}\right] .
$$

Let us try to understand things a bit more thoroughly. If $G$ is a subgroup of GL $(n)$, then consider the inclusion map $j: G \hookrightarrow M_{n \times n}$, where $M_{n \times n}$ is the set of $n \times n$ matrices. Then we have the differential $d j: T G \rightarrow M_{n \times n}$ and the two left multiplications $L_{g}: G \rightarrow G$ and $\mathcal{L}_{g}: M_{n \times n} \rightarrow M_{n \times n}$ for $g \in G$. We have the following commutative diagrams:
\[
\begin{tikzcd}[row sep=large, column sep=large]
M_{n\times n} \arrow[r, "\mathcal{L}_g"] & M_{n\times n} \\
G \arrow[u, "j"] \arrow[r, "L_g"] & G \arrow[u, "j"]
\end{tikzcd}
\quad
\begin{tikzcd}[row sep=large, column sep=large]
M_{n\times n} \arrow[r, "\mathcal{L}_g"] & M_{n\times n} \\
T_h G \arrow[u, "dj|_h"] \arrow[r, "T_h L_g"] & T_{gh} G \arrow[u, "dj|_{gh}"]
\end{tikzcd}
\]
for any $h \in G$. We have used the fact that since $\mathcal{L}_{g}$ is linear, $D \mathcal{L}_{g}(A)=\mathcal{L}_{g}$ for all $A \in M_{n \times n}$. Then for $v_{g} \in T_{g} G$ we have

$$
\begin{aligned}
\left.d j\right|_{e}\left(\omega_{G}\left(v_{g}\right)\right) & =\left.d j\right|_{e}\left(T L_{g^{-1}} v_{g}\right)=\left.\mathcal{L}_{g^{-1}} d j\right|_{g}\left(v_{g}\right) \\
& =\left.g^{-1} d j\right|_{g}\left(v_{g}\right)=\left(j \circ \pi\left(v_{g}\right)\right)^{-1} d j\left(v_{g}\right)
\end{aligned}
$$

where $\pi: T G \rightarrow G$ is the tangent bundle projection. Notice that the effect of $\left.d j\right|_{e}$ is simply to interpret elements of $T_{e} G$ as matrices, and so can be\\
suppressed from the notation. Taking this into account and applying a reasonable abbreviation for $\left(j \circ \pi\left(v_{g}\right)\right)^{-1} d j\left(v_{g}\right)$, we arrive at

$$
\omega_{G}=j^{-1} d j
$$

But notice that for a matrix group, $j$ is none other than the map $\left[x_{j}^{i}\right]: A \mapsto \left[x_{j}^{i}(A)\right]=\left[a_{j}^{i}\right]$, so that we are returned to the expression $\omega_{G}=\left[x_{j}^{i}\right]^{-1}\left[d x_{j}^{i}\right]=$ " $g^{-1} d g$ ". This tells us the meaning of $g$ in the expression $g^{-1} d g$.




\begin{PROP}{LIE-L09-06-10}{Maurer--Cartan-Form für Matrixgruppen}
Sei $G$ eine Lie-Untergruppe von $\GL(n)$.  
Bezeichne durch $x^i_j:\GL(n)\to\mathbb R$ die Koordinatenfunktionen
\[
x^i_j(A)=a^i_j,\qquad A=[a^i_j],
\]
und durch $dx^i_j$ die zugehörigen Differentialformen.  
Sowohl $x^i_j$ als auch $dx^i_j$ schränken sich auf $G$ ein;  
diese Einschränkungen bezeichnen wir mit denselben Symbolen.  

Dann kann die (linke) Maurer--Cartan-Form von $G$ geschrieben werden als
\[
\omega_G = [x^i_j]^{-1}[dx^i_j].
\]
In Kurzschreibweise wird dies häufig notiert als
\[
\omega_G = g^{-1}dg.
\]
\end{PROP}



\begin{PROOF}{LIE-L09-06-11}{P: Maurer--Cartan-Form für Matrixgruppen}
Ein Tangentialvektor $v_g\in T_gG$ kann in Koordinaten der Umgebung
in $\GL(n)$ dargestellt werden als
\[
v_g=\sum v^i_j\,
\left.\frac{\partial}{\partial x^i_j}\right|_g.
\]
Dann ergibt die Auswertung der Maurer--Cartan-Form
\[
\omega_G(v_g)
=[x^i_j]^{-1}[dx^i_j](v_g)
=[g^i_j]^{-1}[v^i_j],
\]
wobei die Matrix $[g^i_j]^{-1}[v^i_j]$ als Element der Lie-Algebra
$\mathfrak g$ zu interpretieren ist.
\end{PROOF}



\begin{EXA}{LIE-L09-06-12}{Beispiel: $\mathrm{SO}(2)$}
Für $G=\SO(2)$ mit
\[
g=
\begin{bmatrix}
x^1_1 & x^1_2\\
x^2_1 & x^2_2
\end{bmatrix},
\qquad
(x^1_1)^2+(x^2_1)^2=1,\quad x^1_1=x^2_2,\quad x^2_1=-x^1_2,
\]
ergibt sich
$$
\begin{aligned}
{\left[\begin{array}{cc}
x_{2}^{2} & -x_{2}^{1} \\
-x_{1}^{2} & x_{1}^{1}
\end{array}\right]\left[\begin{array}{cc}
d x_{1}^{1} & d x_{2}^{1} \\
d x_{1}^{2} & d x_{2}^{2}
\end{array}\right] } & =\left[\begin{array}{cc}
x_{2}^{2} d x_{1}^{1}-x_{2}^{1} d x_{1}^{2} & x_{2}^{2} d x_{2}^{1}-x_{2}^{1} d x_{2}^{2} \\
-x_{1}^{2} d x_{1}^{1}+x_{1}^{1} d x_{1}^{2} & -x_{1}^{2} d x_{2}^{1}+x_{1}^{1} d x_{2}^{2}
\end{array}\right] \\
& =\left[\begin{array}{cc}
x_{1}^{1} d x_{1}^{1}-x_{2}^{1} d x_{1}^{2} & x_{1}^{1} d x_{2}^{1}-x_{2}^{1} d x_{2}^{2} \\
x_{2}^{1} d x_{1}^{1}+x_{1}^{1} d x_{1}^{2} & x_{2}^{1} d x_{2}^{1}+x_{1}^{1} d x_{2}^{2}
\end{array}\right]
\end{aligned}
$$
Da $\mathfrak{so}(2)$ antisymmetrisch ist, folgt
\[
\omega_{\SO(2)}=
\begin{bmatrix}
0 & x^1_1dx^1_2-x^2_1dx^2_2\\[4pt]
-(x^1_1dx^1_2-x^2_1dx^2_2) & 0
\end{bmatrix}.
\]
\end{EXA}



\begin{PROP}{LIE-L09-06-13}{Herleitung über die Einbettung in den Matrizenraum}
Sei $G\subset\GL(n)$ mit Inklusion $j:G\hookrightarrow M_{n\times n}$.
Dann gilt für die zugehörigen Ableitungen
\[
dj:T G\longrightarrow T M_{n\times n}\cong M_{n\times n}.
\]
Für $g\in G$ betrachten wir die Linksabbildungen
\[
L_g:G\to G,\qquad \mathcal L_g:M_{n\times n}\to M_{n\times n},
\quad \mathcal L_g(A)=gA.
\]
Da $\mathcal L_g$ linear ist, gilt $D\mathcal L_g(A)=\mathcal L_g$ für alle $A$.
Somit kommutieren die Diagramme
\[
\begin{tikzcd}[row sep=large, column sep=large]
M_{n\times n}\arrow[r,"\mathcal L_g"] & M_{n\times n}\\
G\arrow[u,"j"]\arrow[r,"L_g"] & G\arrow[u,"j"]
\end{tikzcd}
\qquad
\begin{tikzcd}[row sep=large, column sep=large]
M_{n\times n}\arrow[r,"\mathcal L_g"] & M_{n\times n}\\
T_hG\arrow[u,"dj|_h"]\arrow[r,"T_hL_g"] & T_{gh}G\arrow[u,"dj|_{gh}"]
\end{tikzcd}
\]
für jedes $h\in G$.

Für $v_g\in T_gG$ erhalten wir daher
\[
\begin{aligned}
dj|_e\bigl(\omega_G(v_g)\bigr)
&=dj|_e(T L_{g^{-1}}v_g)
=\mathcal L_{g^{-1}}(dj|_g v_g)
=g^{-1}\,dj|_g(v_g)\\
&=(j\circ\pi(v_g))^{-1}dj(v_g),
\end{aligned}
\]
wobei $\pi:T G\to G$ die Bündelprojektion ist.  
Da $dj|_e$ lediglich die Identifikation von $T_eG$ mit $\mathfrak g$
bewirkt, können wir sie im Symbol unterdrücken und schreiben schließlich:
\[
\boxed{\;\omega_G = j^{-1}dj.\;}
\]
Da für Matrixgruppen $j(A)=[x^i_j(A)]=[a^i_j]$ gilt, folgt
\[
\omega_G = [x^i_j]^{-1}[dx^i_j]
= g^{-1}dg.
\]
Dies erklärt die Bedeutung des häufig verwendeten Ausdrucks
$g^{-1}dg$ für die Maurer--Cartan-Form in Matrixgruppen.
\end{PROP}



\pagebreak


\subsection*{Lie Group Actions}


The basic definitions for group actions were given earlier in Definitions 1.99 and 1.6.2. As before we give most of our definitions and results for left actions and ask the reader to notice that analogous statements can be made for right actions.

Definition 5.95. 

\begin{DEF}{LIE-L09-07-01}{Glatte Lie Gruppen Wirkung auf Mannigfaltigkeit}
Let $l: G \times M \rightarrow M$ be a left action, where $G$ is a Lie group and $M$ is a smooth manifold. If $l$ is a smooth map, then we say that $l$ is a (smooth) Lie group action.

As before, we also use any of the notations $g p, g \cdot p$ or $l_{g}(p)$ for $l(g, p)$.
\end{DEF}

We will need this notational flexibility. 

Recall that for $p \in M$, the orbit of $p$ is denoted $G p$ or $G \cdot p$, and the action is transitive if $G p=M$. Recall also that an action is effective if $l_{g}(p)=p$ for all $p$ only if $g=e$. For right actions $r: M \times G \rightarrow M$, similar definitions apply and we write $p g=r_{g}(p)=r(p, g)$. A right action corresponds to a left action by the rule $g p:=p g^{-1}$.



\begin{DEF}{LIE-L09-07-02}{Lie-Gruppenwirkung, Orbits und Effektivität}
Sei $G$ eine Lie-Gruppe und $M$ eine glatte Mannigfaltigkeit.
Eine \emph{(linke) glatte Aktion} von $G$ auf $M$ ist eine glatte Abbildung
\[
l : G\times M \longrightarrow M,\qquad (g,p)\longmapsto l_g(p)=g\cdot p,
\]
die folgenden Eigenschaften genügt:
\begin{enumerate}[label=(\roman*)]
\item $e\cdot p=p$ für alle $p\in M$, wobei $e$ das neutrale Element von $G$ ist;
\item $(g_1g_2)\cdot p = g_1\cdot(g_2\cdot p)$ für alle $g_1,g_2\in G$ und $p\in M$.
\end{enumerate}
\end{DEF}



\begin{DEF}{LIE-L09-07-03}{Orbits, Transitivität und Effektivität}
Für $p\in M$ heißt
\[
G\cdot p := \{g\cdot p \mid g\in G\}
\]
die \emph{Orbit} von $p$ unter der Aktion.
Eine Aktion heißt \emph{transitiv}, wenn
\[
G\cdot p = M
\quad\text{für jedes (oder ein) }p\in M.
\]
Sie heißt \emph{effektiv} (oder treu), wenn gilt:
\[
g\cdot p = p\text{ für alle }p\in M
\quad\Longrightarrow\quad g=e.
\]
\end{DEF}



\begin{DEF}{LIE-L09-07-04}{Rechte Wirkung und Zusammenhang zu linken Wirkung}
Eine \emph{(rechte) glatte Aktion} von $G$ auf $M$ ist eine glatte Abbildung
\[
r : M\times G \longrightarrow M,\qquad (p,g)\longmapsto r_g(p)=p\cdot g,
\]
mit den Eigenschaften:
\begin{enumerate}[label=(\roman*)]
\item $p\cdot e = p$ für alle $p\in M$;
\item $(p\cdot g_1)\cdot g_2 = p\cdot(g_1g_2)$ für alle $g_1,g_2\in G$ und $p\in M$.
\end{enumerate}
Zwischen linken und rechten Aktionen besteht die eindeutige Beziehung
\[
g\cdot p := p\cdot g^{-1},
\]
die eine rechte Aktion in eine linke überführt (und umgekehrt).
\end{DEF}




Definition 5.96. 

\begin{DEF}{LIE-L09-07-05}{Isotropiegruppe / Stabilisator von $p$ bzgl Lie gruppenwirkung}
Let $l$ be a Lie group action as above. For a fixed $p \in M$, the isotropy group of $p$ is defined to be
$$
G_{p}:=\{g \in G: g p=p\} .
$$
The isotropy group of $p$ is also called the stabilizer of $p$.
\end{DEF}


Exercise 5.97. 

\begin{EXE}{LIE-L09-07-06}{Stabilisator ist abgeschlossene Lie Untergruppe}
Show that $G_{p}$ is a closed subset and an abstract subgroup of $G$. This means that $G_{p}$ is a closed Lie subgroup.
\end{EXE}



Recalling the definition of a free action (Definition 1.100), it is easy to see that an action is free if and only if the isotropy subgroup of every point is the trivial subgroup consisting of the identity element alone.


\begin{DEF}{LIE-L09-07-07}{Freie Lie Gruppenwirkung und Isotropiegruppe}
Sei $l:G\times M\to M$, $(g,p)\mapsto g\cdot p$, eine glatte (linke) Aktion
einer Lie-Gruppe $G$ auf einer Mannigfaltigkeit $M$. Eine Aktion heißt \emph{frei}, wenn die Isotropiegruppe jedes Punktes trivial ist, d.\,h.
\[
G_p = \{e\}
\quad\text{für alle }p\in M.
\]
Äquivalent gilt:
\[
g\cdot p = p \text{ für ein }p\in M
\quad\Longrightarrow\quad g=e.
\]
\end{DEF}




Definition 5.98. 



\begin{DEF}{LIE-L09-07-08}{Invariante Teilmengen bzgl Lie Gruppenwirkung}
Suppose that we have a Lie group action of $G$ on $M$. If $N$ is a subset of $M$ and $G x \subset N$ for all $x \in N$, then we say that $N$ is an invariant subset. If $N$ is also a submanifold, then it is called an invariant submanifold.

In this definition, we include the possibility that $N$ is an open submanifold. If $N$ is an invariant subset of $M$, then it is easy to see that $g N=N$, where $g N=l_{g}(N)$ for any $g$. Furthermore, if $N$ is a submanifold, then the action restricts to a Lie group action $G \times N \rightarrow N$.

If $G$ is zero-dimensional, then by definition it is just a group with discrete topology, and we recover the definition of a discrete group action.
\end{DEF}





Example 5.99. 

\begin{EXA}{LIE-L09-07-09}{$L_g,R_g$ als Lie Gruppenwirkung}
The maps $G \times G \rightarrow G$ given by $(g, x) \mapsto L_{g} x$ and $(g, x) \mapsto R_{g} x$ are Lie group actions.
\end{EXA}

Example 5.100. 

\begin{EXA}{LIE-L09-07-10}{$\GL(V)$ als Lie Gruppenwirkung auf $V$}
In case $M=\mathbb{R}^{n}$, the Lie group $\operatorname{GL}(n, \mathbb{R})$ acts on $\mathbb{R}^{n}$ by matrix multiplication. Similarly, $\operatorname{GL}(n, \mathbb{C})$ acts on $\mathbb{C}^{n}$. More abstractly, $\mathrm{GL}(\mathrm{V})$ acts on the vector space V . This action is smooth since $A x$ depends smoothly (polynomially) on the components of $A$ and on the components of $x \in \mathbb{R}^{n}$.
\end{EXA}

Example 5.101. 

\begin{EXA}{LIE-L09-07-11}{Lie Untergruppen von $\GL(n,\R)$ als Lie Gruppenwirkung}
Any Lie subgroup of $\operatorname{GL}(n, \mathbb{R})$ acts on $\mathbb{R}^{n}$ also by matrix multiplication.
\end{EXA}



\begin{EXA}{LIE-L09-07-12}{$\Ogroup(n,\R)$ als Lie Gruppenwirkung}
For example, $\mathrm{O}(n, \mathbb{R})$ acts on $\mathbb{R}^{n}$. For every $x \in \mathbb{R}^{n}$, the orbit of $x$ is the sphere of radius $\|x\|$. This is trivially true if $x=0$. In general, if $\|x\| \neq 0$, then $\|g x\|=\|x\|$ for any $g \in \mathrm{O}(n, \mathbb{R})$. On the other hand, if $x, y \in \mathbb{R}^{n}$ and $\|x\|=\|y\|=r$, then let $\widehat{x}:=x / r$ and $\widehat{y}:=y / r$. Let $\widehat{x}:=e_{1}$ and $\widehat{y}=f_{1}$ and then extend to orthonormal bases $\left(e_{1}, \ldots, e_{n}\right)$ and $\left(f_{1}, \ldots, f_{n}\right)$. Then there exists an orthogonal matrix $S$ such that $S e_{i}=f_{i}$ for $i=1, \ldots, n$. In particular, $S \widehat{x}=\widehat{y}$, and so $S x=y$.
\end{EXA}

Exercise 5.102. 

\begin{EXE}{LIE-L09-07-13}{$\SO(n,\R)$ als transitive Lie Gruppenwirkung}
From the last example we can restrict the action of $\mathrm{O}(n, \mathbb{R})$ to a transitive action on $S^{n-1}$. Now $\operatorname{SO}(n, \mathbb{R})$ also acts on $S^{n-1}$ by restriction. Show that this action on $S^{n-1}$ is transitive as long as $n>1$.
\end{EXE}

Example 5.103. 

\begin{EXA}{LIE-L09-07-14}{$L_h$ als Lie Gruppenwirkung von Lie Untergruppe auf Lie Gruppe}
If $H$ is a Lie subgroup of a Lie group $G$, then we can consider $L_{h}$ for any $h \in H$ and thereby obtain a Lie group action of $H$ on $G$.
\end{EXA}


\begin{DEF}{LIE-L09-07-15}{Normale Untergruppe}
Recall that a subgroup $H$ of a group $G$ is called a normal subgroup if $g k g^{-1} \in K$ for any $k \in H$ and all $g \in G$. In other words, $H$ is normal if $g H g^{-1} \subset H$ for all $g \in G$, and it is easy to see that in this case we always have $g H g^{-1}=H$.
\end{DEF}

Example 5.104. 

\begin{EXA}{LIE-L09-07-16}{Lie Gruppe wirkt durch konjugierte Abbildung bzgl $g$ auf Lie Untergruppe}
If $H$ is a normal Lie subgroup of $G$, then $G$ acts on $H$ by conjugation:
$$
g \cdot h:=C_{g} h=g h g^{-1}
$$
Notice that the notation $g \cdot h$ cannot reasonably be abbreviated to $g h$ in this example.
\end{EXA}


\begin{EXA}{LIE-L09-07-17}{Äquivariante Abbildungen zwischen Mannigfaltigkeiten bzgl Gruppenwirkung und Lie Gruppen Wirkung}
Suppose that a Lie group $G$ acts on smooth manifolds $M$ and $N$. For simplicity, we take both actions to be left actions, which we denote by $l$ and $\lambda$, respectively. A map $f: M \rightarrow N$ such that 
$$f \circ l_{g}=\lambda_{g} \circ f$$ 
for all $g \in G$, is said to be an equivariant map (equivariant with respect to the given actions). This means that for all $g$ the following diagram commutes:
\[
\begin{tikzcd}[row sep=large, column sep=large]
M \arrow[r, "f"] \arrow[d, "l_g"'] & N \arrow[d, "\lambda_g"] \\
M \arrow[r, "f"] & N
\end{tikzcd}
\]
If $f$ is also a diffeomorphism, then $f$ is an equivalence of Lie group actions.
\end{EXA}


Example 5.105. 


\begin{EXA}{LIE-L09-07-18}{Induzierte Lie Gruppenwirkung durch Lie Gruppen Homomorphismus}
If $\phi: G \rightarrow H$ is a Lie group homomorphism, then we can define an action of $G$ on $H$ by $\lambda(g, h)=\lambda_{g}(h)=L_{\phi(g)} h$. In this situation, $\phi$ is equivariant with respect to the actions $\lambda$ and $L$ (left translation).
\end{EXA}



Example 5.106. 

\begin{EXA}{LIE-L09-07-19}{Äquivariante Wirkungen zwischen $\R^n$ und $T^n$}
Let $T^{n}=S^{1} \times \cdots \times S^{1}$ be the $n$-torus, where we identify $S^{1}$ with the complex numbers of unit modulus. Fix $k=\left(k_{1}, \ldots, k_{n}\right) \in \mathbb{R}^{n}$. Then $\mathbb{R}$ acts on $\mathbb{R}^{n}$ by 
$$\tau^{k}(t, x)=t \cdot x:=x+t k$$ 
On the other hand, $\mathbb{R}$ acts on $T^{n}$ by 
$$t \cdot\left(z^{1}, \ldots, z^{n}\right)=\left(e^{i t k_{1}} z^{1}, \ldots, e^{i t k_{n}} z^{n}\right)$$ 
The map $\mathbb{R}^{n} \rightarrow T^{n}$ given by $\left(x^{1}, \ldots, x^{n}\right) \mapsto\left(e^{i x^{1}}, \ldots, e^{i x^{n}}\right)$ is equivariant with respect to these actions.
\end{EXA}


Theorem 5.107 (Equivariant rank theorem). 

\begin{THEO}{LIE-L09-07-20}{Äquivariante Rang Theorem}
Suppose that $f: M \rightarrow N$ is smooth and that a Lie group $G$ acts on both $M$ and $N$ with the action on $M$ being transitive. If $f$ is equivariant, then it has constant rank. In particular, each level set of $f$ is a closed regular submanifold.
\end{THEO}

Proof. 

\begin{PROOF}{LIE-L09-07-21}{P: Äquivariante Rang Theorem}
Let the actions on $M$ and $N$ be denoted by $l$ and $\lambda$ respectively as before. Pick any two points $p_{1}, p_{2} \in M$. Since $G$ acts transitively on $M$, there is a $g$ with $l_{g} p_{1}=p_{2}$. By hypothesis, we have $f \circ l_{g}=\lambda_{g} \circ f$, which corresponds to the commutative diagram (5.3). Upon application of the tangent functor we have the commutative diagram
\[
\begin{tikzcd}[row sep=large, column sep=large]
T_{p_1} M \arrow[r, "T_{p_1} f"] \arrow[d, "T_{p_1} l_g"'] & T_{f(p_1)} N \arrow[d, "T_{f(p_1)} \lambda_g"] \\
T_{p_2} M \arrow[r, "T_{p_2} f"] & T_{f(p_2)} N
\end{tikzcd}
\]
Since the maps $T_{p_{1}} l_{g}$ and $T_{f\left(p_{1}\right)} \lambda_{g}$ are linear isomorphisms, we see that $T_{p_{1}} f$ must have the same rank as $T_{p_{2}} f$. Since $p_{1}$ and $p_{2}$ were arbitrary, we see that the rank of $f$ is constant on $M$. Apply Theorem C.5.
\end{PROOF}



There are several corollaries of this nice theorem. 


\begin{EXA}{LIE-L09-07-22}{$\Ogroup(n)$ als Levelset einer äquivarianten Abbildung von $\GL(n,\R)$ nach $\mathfrak{gl}(n,\R)$}
For example, we know that $\mathrm{O}(n, \mathbb{R})$ is the level set $f^{-1}(I)$, where $f: \mathrm{GL}(n, \mathbb{R}) \rightarrow \mathfrak{g} \mathfrak{l}(n, \mathbb{R})=M_{n \times n}$ is given by $f(A)=A^{T} A$. 

The group $\mathrm{O}(n, \mathbb{R})$ acts on itself via left translation, and we also let $\mathrm{O}(n, \mathbb{R})$ act on $\mathfrak{g} \mathfrak{l}(n, \mathbb{R})$ by $Q \cdot A:=Q^{T} A Q$ (adjoint\\
action). One checks easily that $f$ is equivariant with respect to these actions, and since the first action (left translation) is certainly transitive, we see that $\mathrm{O}(n, \mathbb{R})$ is a closed regular submanifold of $\mathrm{GL}(n, \mathbb{R})$.

It follows from Proposition 5.9 that $\mathrm{O}(n, \mathbb{R})$ is a closed Lie subgroup of $\operatorname{GL}(n, \mathbb{R})$.
\end{EXA}





Similar arguments apply for $\mathrm{U}(n, \mathbb{C}) \subset \mathrm{GL}(n, \mathbb{C})$ and other linear Lie groups. In fact, we have the following general corollary to Theorem 5.107 above.

Corollary 5.108. 

\begin{THEO}{LIE-L09-07-23}{Kern von Lie Gruppen Homomorphismus ist eine abgeschlossene Lie Untergruppe}
If $\phi: G \rightarrow H$ is a Lie group homomorphism, then the kernel $\operatorname{Ker}(h)$ is a closed Lie subgroup of $G$.
\end{THEO}

Proof. 

\begin{PROOF}{LIE-L09-07-24}{P: Kern von Lie Gruppen Homomorphismus ist eine abgeschlossene Lie Untergruppe}
Let $G$ act on itself and on $H$ as in Example 5.105. Then $\phi$ is equivariant, and $\phi^{-1}(e)=\operatorname{Ker}(h)$ is a closed Lie subgroup by Theorem 5.107 and Proposition 5.9.
\end{PROOF}


Corollary 5.109. 

\begin{KORO}{LIE-L09-07-25}{Isotriopiegruppe bzgl $g$ und Lie Gruppenwirkung ist abgeschlossene Lie Untergruppe}
Let $l: G \times M \rightarrow M$ be a Lie group action, and let $G_{p}$ be the isotropy subgroup of some $p \in M$. Then $G_{p}$ is a closed Lie subgroup of $G$.
\end{KORO}

Proof. 

\begin{PROOF}{LIE-L09-07-26}{P: Isotriopiegruppe bzgl $g$ und Lie Gruppenwirkung ist abgeschlossene Lie Untergruppe}
The orbit map $\theta_{p}: G \rightarrow M$ given by $\theta_{p}(g)=g p$ is equivariant with respect to left translation on $G$ and the given action on $M$. Thus by the equivariant rank theorem, $G_{p}$ is a regular submanifold of $G$, and then by Proposition 5.9 it is a closed Lie subgroup.
\end{PROOF}



Proper Actions and Quotients. At several points in this section, such as the proof of Proposition 5.111 below, we follow [Lee, John].


Definition 5.110. 

\begin{DEF}{LIE-L09-07-27}{Eigentliche Gruppenwirkung}
Let $l: G \times M \rightarrow M$ be a smooth (or merely continuous) group action. If the map $P: G \times M \rightarrow M \times M$ given by $(g, p) \mapsto\left(l_{g} p, p\right)$ is proper, we say that the action is a proper action.
\end{DEF}

It is important to notice that a proper action is not defined to be an action such that the defining map $l: G \times M \rightarrow M$ is proper. 

We now give a useful characterization of a proper action. For any subset $K \subset M$, let $g \cdot K:=\{g x: x \in K\}$.


Proposition 5.111. 

\begin{PROP}{LIE-L09-07-28}{Charakterisierung eigentlicher Gruppenwirkungen}
Let $l: G \times M \rightarrow M$ be a smooth (or merely continuous) group action. Then $l$ is a proper action if and only if the set
$$
G_{K}:=\{g \in G:(g \cdot K) \cap K \neq \emptyset\}
$$
is compact whenever $K$ is compact.
\end{PROP}


Proof. 

\begin{PROOF}{LIE-L09-07-29}{P: Charakterisierung eigentlicher Gruppenwirkungen}
Suppose that $l$ is proper so that the map $P$ is a proper map. Let $\pi_{G}$ be the first factor projection $G \times M \rightarrow G$. Then
$$
\begin{aligned}
G_{K} & =\{g: \text { there exists an } x \in K \text { such that } g x \in K\} \\
& =\{g: \text { there exists an } x \in M \text { such that } P(g, x) \in K \times K\} \\
& =\pi_{G}\left(P^{-1}(K \times K)\right)
\end{aligned}
$$
and so $G_{K}$ is compact.

Next we assume that $G_{K}$ is compact for all compact $K$. If $C \subset M \times M$ is compact, then letting $K=\pi_{1}(C) \cup \pi_{2}(C)$, where $\pi_{1}$ and $\pi_{2}$ are the first and second factor projections $M \times M \rightarrow M$ respectively, we have
$$
\begin{aligned}
P^{-1}(C) & \subset P^{-1}(K \times K) \subset\{(g, x): g x \in K \text { and } x \in K\} \\
& \subset G_{K} \times K
\end{aligned}
$$
Since $P^{-1}(C)$ is a closed subset of the compact set $G_{K} \times K$, it is compact. This means that $P$ is proper since $C$ was an arbitrary compact subset of $M \times M$.
\end{PROOF}


Using this proposition, one can show that Definition 1.106 for discrete actions is consistent with Definition 5.110 above.

Proposition 5.112. 


\begin{PROP}{LIE-L09-07-30}{Glatte Gruppenwirkungen mit kompaktem Definitionsbereich sind eigentlich}
If $G$ is compact, then any smooth action $l: G \times M \rightarrow M$ is proper.
\end{PROP}

Proof. 

\begin{PROOF}{LIE-L09-07-31}{P: Glatte Gruppenwirkungen mit kompaktem Definitionsbereich sind eigentlich}
Let $B \subset M \times M$ be compact. We find a compact subset $K \subset M$ such that $B \subset K \times K$ as in the proof of Proposition 5.111.

Claim: $P^{-1}(B)$ is compact. Indeed,
$$
\begin{aligned}
P^{-1}(B) & \subset P^{-1}(K \times K)=\bigcup_{k \in K} P^{-1}(K \times\{k\}) \\
& =\bigcup_{k \in K}\{(g, p):(g p, p) \in K \times\{k\}\} \\
& =\bigcup_{k \in K}\{(g, k): g p \in K\} \\
& \subset \bigcup_{k \in K}(G \times\{k\})=G \times K
\end{aligned}
$$
Thus $P^{-1}(B)$ is a closed subset of the compact set $G \times K$ and hence is compact.
\end{PROOF}



Exercise 5.113. 

\begin{EXE}{LIE-L09-07-34}{Eigenschaften von eigentlichen Gruppenwirkungen}
Prove the following:

(i) If $l: G \times M \rightarrow M$ is a proper action and $H \subset G$ is a closed subgroup, then the restricted action $H \times M \rightarrow M$ is proper.\\

(ii) If $N$ is an invariant submanifold for a proper action $l: G \times M \rightarrow M$, then the restricted action $G \times N \rightarrow N$ is also proper.
\end{EXE}



Let us consider a Lie group action $l: G \times M \rightarrow M$ that is both proper and free. The orbit map at $p$ is the map $\theta_{p}: G \rightarrow M$ given by $\theta_{p}(g)=g \cdot p$. It is easily seen to be smooth, and its image is obviously $G \cdot p$. In fact, since the action is free, each orbit map is injective. Also, $\theta_{p}$ is equivariant with respect to the left action of $G$ on itself and the action $l$:
$$
\begin{aligned}
\theta_{p}(g x) & =(g x) \cdot p=g \cdot(x \cdot p) \\
& =g \cdot \theta_{p}(x)
\end{aligned}
$$
for all $x, g \in G$. It follows from Theorem 5.107 (the equivariant rank theorem) that $\theta_{p}$ has constant rank, and since it is injective, it must be an immersion. Not only that, but it is a proper map. Indeed, for any compact $K \subset M$ the set $\theta_{p}^{-1}(K)$ is a closed subset of the set $G_{K \cup\{p\}}$, and since the latter set is compact by Theorem 5.111, $\theta_{p}^{-1}(K)$ is compact. By Exercise 3.9, $\theta_{p}$ is an embedding, so each orbit is a regular submanifold of $M$.

It will be very convenient to have charts on $M$ which fit the action of $G$ in a nice way. See Figure 5.1.

\begin{figure}[h]
\begin{center}
  \includegraphics[scale=0.25]{2025_10_09_265d7e1a22c89f79c21eg-248}

\caption{Figure 5.1. Action-adapted chart}
\end{center}
\end{figure}




\begin{DEF}{LIE-L09-07-35}{Orbitabbildung einer Lie-Gruppenwirkung}
Sei $l:G\times M\to M$ eine glatte Aktion einer Lie-Gruppe $G$ auf einer Mannigfaltigkeit $M$.
Für $p\in M$ heißt
\[
\theta_p:G\longrightarrow M,\qquad
\theta_p(g)=g\cdot p,
\]
die \emph{Orbitabbildung} in $p$.
Ihr Bild ist der Orbit $G\cdot p=\{\;g\cdot p\mid g\in G\;\}$.
\end{DEF}





\begin{PROP}{LIE-L09-07-36}{Eigenschaften der Orbitabbildung bei freien und properen Aktionen}
Sei $l:G\times M\to M$ eine \emph{freie} und \emph{propre} Aktion einer Lie-Gruppe $G$ auf einer Mannigfaltigkeit $M$.
Dann gilt für jedes $p\in M$:
\begin{enumerate}[label=(\roman*)]
\item Die Abbildung $\theta_p:G\to M$, $\theta_p(g)=g\cdot p$, ist glatt und $G$-äquivariant:
\[
\theta_p(gx)=(gx)\cdot p=g\cdot(x\cdot p)=g\cdot\theta_p(x)
\quad\text{für alle }g,x\in G.
\]
\item Da die Aktion frei ist, ist $\theta_p$ injektiv.
\item Nach dem äquivarianten Rangtheorem besitzt $\theta_p$ konstanten Rang; da sie injektiv ist, ist sie eine \emph{Immersion}.
\item Die Abbildung $\theta_p$ ist \emph{proper}:  
Ist $K\subset M$ kompakt, so ist
\[
\theta_p^{-1}(K)\subset G_{K\cup\{p\}}
\]
eine abgeschlossene Teilmenge der nach Theorem~5.111 kompakten Menge $G_{K\cup\{p\}}$; somit ist $\theta_p^{-1}(K)$ kompakt.
\item Nach Aufgabe~3.9 folgt aus (iii) und (iv), dass $\theta_p$ eine \emph{Einbettung} ist.
\end{enumerate}
Daraus folgt, dass jeder Orbit $G\cdot p$ eine eingebettete reguläre Untermannigfaltigkeit von $M$ ist.

$$\includegraphics[scale=0.25]{2025_10_09_265d7e1a22c89f79c21eg-248}$$
\end{PROP}







\begin{CONC}{LIE-L09-07-37}{Orbits als eingebettete Untermannigfaltigkeiten}
Die Eigenschaft, dass die Orbits eingebettete Untermannigfaltigkeiten sind,
ist eine der zentralen Konsequenzen freier und properer Aktionen.
Sie erlaubt es, lokale Karten auf $M$ zu wählen, die „an die Gruppenwirkung angepasst“ sind.
\end{CONC}






Definition 5.114. 

\begin{DEF}{LIE-L09-07-38}{Lie Gruppenwirkung angepasste lokale Karte von Mannigfaltigkeiten}
Let $M$ be an $n$-manifold and $G$ a Lie group of dimension $k$. If $l: G \times M \rightarrow M$ is a Lie group action, then an action-adapted chart on $M$ is a chart ( $U, \mathrm{x}$ ) such that\\
\begin{enumerate}
\item $\mathrm{x}(U)$ is a product open set $V_{1} \times V_{2} \subset \mathbb{R}^{k} \times \mathbb{R}^{n-k}=\mathbb{R}^{n}$
\item If an orbit has nonempty intersection with $U$, then that intersection has the form
$$
\left\{x^{k+1}=c^{1}, \ldots, x^{n}=c^{n-k}\right\}
$$
for some constants $c^{1}, \ldots, c^{n-k}$
\end{enumerate}
\end{DEF}


Theorem 5.115. 

\begin{THEO}{LIE-L09-07-39}{Für freie eigentliche Lie Gruppenwirkungen existieren angepasste Karten}
If $l: G \times M \rightarrow M$ is a free and proper Lie group action, then for every $p \in M$ there is an action-adapted chart centered at $p$.
\end{THEO}

Proof. 


\begin{PROOF}{LIE-L09-07-40}{P: Für freie eigentliche Lie Gruppenwirkungen existieren angepasste Karten}
Let $p \in M$ be given. Since $G \cdot p$ is a regular submanifold, we may choose a regular submanifold chart ( $W, \mathrm{y}$ ) centered at $p$ so that ( $G \cdot p$ ) $\cap W$ is exactly given by $y^{k+1}=\cdots=y^{n}=0$ in $W$. Let $S$ be the complementary slice in $W$ given by $y^{1}=\cdots=y^{k}=0$. Note that $S$ is a regular submanifold. The tangent space $T_{p} M$ decomposes as
$$
T_{p} M=T_{p}(G \cdot p) \oplus T_{p} S .
$$
Let $\varphi: G \times S \rightarrow M$ be the restriction of the action $l$ to the set $G \times S$. Also, let $i_{p}: G \rightarrow G \times S$ be the insertion map $g \mapsto(g, p)$ and let $j_{e}: S \rightarrow G \times S$ be the insertion map $s \mapsto(e, s)$. (See Figure 5.2.) These insertion maps are embeddings, and we have $\theta_{p}=\varphi \circ i_{p}$ and also $\varphi \circ j_{e}=\iota$, where $\iota$ is the inclusion $S \hookrightarrow M$. Now $T_{e} \theta_{p}\left(T_{e} G\right)=T_{p}(G \cdot p)$ since $\theta_{p}$ is an embedding. On the other hand, $T \theta_{p}=T \varphi \circ T i_{p}$, and so the image of $T_{(e, p)} \varphi$ must contain $T_{p}(G \cdot p)$. Similarly, from the composition $\varphi \circ j_{e}=\iota$ we see that the image of $T_{(e, p)} \varphi$ must contain $T_{p} S$. It follows that $T_{(e, p)} \varphi: T_{(e, p)}(G \times S) \rightarrow T_{p} M$ is surjective, and since $T_{(e, p)}(G \times S)$ and $T_{p} M$ have the same dimension, it is also injective.

By the inverse mapping theorem, there is a neighborhood $O$ of $(e, p)$ such that $\left.\varphi\right|_{O}$ is a diffeomorphism. By shrinking $O$ further if necessary we may assume that $\varphi(O) \subset W$. We may also arrange that $O$ has the form of a product $O=A \times B$ for $A$ open in $G$ and $B$ open in $S$. In fact, we can assume that there are diffeomorphisms $\alpha: I^{k} \rightarrow A$ and $\beta: I^{n-k} \rightarrow B$, where $I^{k}$ and $I^{n-k}$ are the open cubes in $\mathbb{R}^{k}$ and $\mathbb{R}^{n-k}$ given respectively by $I^{k}=(-1,1)^{k}$ and $I^{n-k}=(-1,1)^{n-k}$ and where $\alpha^{-1}(e)=0 \in \mathbb{R}^{k}$ and $\beta^{-1}(p)=0 \in \mathbb{R}^{n-k}$. Let $U:=\varphi(A \times B)$. The map $\varphi \circ(\alpha \times \beta)$ : $I^{k} \times I^{n-k} \rightarrow U$ is a diffeomorphism, and so its inverse is a chart. We now show that $B$ can be chosen small enough that the intersection of each orbit with $B$ is either empty or a single point. If this were not true, then there would be a sequence of open sets $\left\{B_{i}\right\}$ with compact closure and $\bar{B}_{i+1} \subset B_{i}$ (and with corresponding diffeomorphisms $\beta_{i}: I^{k} \rightarrow B_{i}$ as above), such that for every $i$ there is a pair of distinct points $p_{i}, p_{i}^{\prime} \in B_{i}$ with $g_{i} p_{i}=p_{i}^{\prime}$ for some sequence $\left\{g_{i}\right\} \subset G$. (We have used the fact that manifolds are first countable and normal.) This forces both $p_{i}$ and $p_{i}^{\prime}=g_{i} p_{i}$ to converge to $p$. From this we see that the set $K=\left\{\left(g_{i} p_{i}, p_{i}\right),(p, p)\right\} \subset M \times M$ is compact. Recall that by definition, the map $P:(g, x) \mapsto(g x, x)$ is proper. Since $\left(g_{i}, p_{i}\right)=P^{-1}\left(g_{i} p_{i}, p_{i}\right)$, we see that $\left\{\left(g_{i}, p_{i}\right)\right\}$ is a subset of the compact set $P^{-1}(K)$. Thus after passing to a subsequence, we have that $\left(g_{i}, p_{i}\right)$ converges to $(g, p)$ for some $g$ and hence $g_{i} \rightarrow g$ and $g_{i} p_{i} \rightarrow g p$. But this means that we have
$$
g p=\lim _{i \rightarrow \infty} g_{i} p_{i}=\lim _{i \rightarrow \infty} p_{i}^{\prime}=p
$$
and since the action is free, we conclude that $g=e$. However, this is impossible since it would mean that for large enough $i$ we would have $g_{i} \in A$, and this in turn would imply that
$$
\varphi\left(g_{i}, p_{i}\right)=l_{g_{i}}\left(p_{i}\right)=p_{i}^{\prime}=l_{e}\left(p_{i}^{\prime}\right)=\varphi\left(e, p_{i}^{\prime}\right)
$$
contradicting the injectivity of $\varphi$ on $A \times B$. Thus after shrinking $B$ we may assume that the intersection of each orbit with $B$ is either empty or a single point. One may now check that with $\mathrm{x}:=(\varphi \circ(\alpha \times \beta))^{-1}: U \rightarrow I^{k} \times I^{n-k} \subset \mathbb{R}^{n}$, we obtain a chart ($U,\mathrm{x}$) with the desired properties. Write $\mathrm{x}=(x, y)$, where $y$ takes values in $I^{n-k} \subset \mathbb{R}^{n-k}$ and $x$ takes values in $I^{k} \subset \mathbb{R}^{k}$. Each $y=c$ slice is of the form $\varphi(A \times\{q\}) \subset G q$ for $q \in B$ and so is contained in a single orbit. We see that the intersection of an orbit with $U$ must be
$$
\includegraphics[scale=0.25]{2025_10_09_265d7e1a22c89f79c21eg-250}
$$
a union of such slices. But each orbit can only intersect $B$ in at most one point, and so it is clear that each orbit intersects $U$ in one slice or does not intersect $U$ at all.
\end{PROOF}




Notice that we have actually constructed action-adapted charts that have image the cube $I^{n}$. For the next lemma, we continue with the convention that $I$ is the interval $(-1,1)$.

Lemma 5.116. 

\begin{LEM}{LIE-L09-07-41}{Konstruktion von angepassten Karten}
Let $\mathrm{x}:=(\varphi \circ(\alpha \times \beta))^{-1}: U \rightarrow I^{k} \times I^{n-k}=I^{n} \subset \mathbb{R}^{n}$ be an action-adapted chart map obtained as in the proof of Theorem 5.115 above. Then given any $p_{1} \in U$, there exists a diffeomorphism $\psi: I^{n} \rightarrow I^{n}$ such that $\psi \circ \mathrm{x}$ is an action-adapted chart centered at $p_{1}$ and with image $I^{n}$. Furthermore $\psi$ can be decomposed as $(a, b) \mapsto\left(\psi_{1}(a), \psi_{2}(b)\right)$, where $\psi_{1}: I^{k} \rightarrow I^{k}$ and $\psi_{2}: I^{n-k} \rightarrow I^{n-k}$ are diffeomorphisms.
\end{LEM}

Proof. 

\begin{PROOF}{LIE-L09-07-42}{P: Konstruktion von angepassten Karten}
We need to show that for any $a \in I^{n}$, there is a diffeomorphism $\psi: I^{n} \rightarrow I^{n}$ such that $\psi(a)=0$. Let $a^{i}$ be the $i$-th component of $a$. Let $\psi_{i}: I \rightarrow I$ be defined by
$$
\psi_{i}:=\phi^{-1} \circ t_{-\phi\left(a_{i}\right)} \circ \phi,
$$
where $t_{-c}(x):=x-c$ and $\phi:(-1,1) \rightarrow \mathbb{R}$ is the diffeomorphism $\phi: x \mapsto \tan \left(\frac{\pi}{2} x\right)$. The diffeomorphism we want is $\psi(x)=\left(\psi_{1}\left(x^{1}\right), \ldots, \psi_{1}\left(x^{n}\right)\right)$. The last part is obvious from the construction.
\end{PROOF}


If $l: G \times M \rightarrow M$ is a Lie group action, then there is a natural equivalence relation on $M$ whereby the equivalence classes are exactly the orbits of the action. The quotient space (or orbit space) is denoted $G \backslash M$, and we have the quotient map $\pi: M \rightarrow G \backslash M$. 

We put the quotient topology on $G \backslash M$ so that $A \subset G \backslash M$ is open if and only if $\pi^{-1}(A)$ is open in $M$. The quotient map is also open. Indeed, let $U \subset M$ be open. We want to show that $\pi(U)$ is open, and for this it suffices to show that $\pi^{-1}(\pi(U))$ is open. But $\pi^{-1}(\pi(U))$ is the union $\bigcup_{g} l_{g}(U)$, and this is open since each $l_{g}(U)$ is open.




\begin{DEF}{LIE-L09-07-43}{Orbitraum und Quotiententopologie einer Lie-Gruppenwirkung}
Sei $l:G\times M\to M$ eine glatte Aktion einer Lie-Gruppe $G$ auf einer Mannigfaltigkeit $M$.
Die durch die Aktion induzierte Äquivalenzrelation auf $M$ lautet:
\[
p\sim q
\quad:\Longleftrightarrow\quad
\exists g\in G\text{ mit } q=g\cdot p.
\]
Die Äquivalenzklassen dieser Relation sind genau die \emph{Orbits}
\[
G\cdot p=\{g\cdot p\mid g\in G\}.
\]
Der resultierende Quotientenraum heißt \emph{Orbitraum} oder \emph{Quotientraum der Aktion} und wird geschrieben als
\[
G\backslash M := M/\!\sim,
\]
mit der kanonischen Projektion
\[
\pi:M\longrightarrow G\backslash M,\qquad p\longmapsto G\cdot p.
\]
\end{DEF}



\begin{DEF}{LIE-L09-07-44}{Quotiententopologie auf $G\backslash M$}
Der Orbitraum $G\backslash M$ erhält die \emph{Quotiententopologie}, definiert durch:
\[
A\subset G\backslash M\text{ ist offen }
\iff
\pi^{-1}(A)\subset M\text{ ist offen.}
\]
\end{DEF}



\begin{LEM}{LIE-L09-07-45}{Offenheit der Quotientenabbildung}
Die Projektion $\pi:M\to G\backslash M$ ist eine offene Abbildung.
\end{LEM}



\begin{PROOF}{LIE-L09-07-46}{P: Offenheit der Quotientenabbildung}
Sei $U\subset M$ offen.  
Wir zeigen, dass $\pi(U)$ offen in $G\backslash M$ ist.  
Dazu genügt es zu zeigen, dass $\pi^{-1}(\pi(U))$ offen in $M$ ist.

Für $U\subset M$ gilt
\[
\pi^{-1}(\pi(U))
=\bigcup_{g\in G}l_g(U),
\]
wobei $l_g:M\to M$, $p\mapsto g\cdot p$, die Linksoperation bezeichnet.
Da jede Abbildung $l_g$ ein Diffeomorphismus ist, ist $l_g(U)$ offen,
und somit auch die Vereinigung $\bigcup_g l_g(U)$.  
Damit ist $\pi^{-1}(\pi(U))$ offen, und folglich ist $\pi(U)$ offen in $G\backslash M$.
\end{PROOF}





Proposition 5.117. 

\begin{PROP}{LIE-L09-07-47}{Induzierte Hausdorff / 2. Abzählbarkeit für Orbitraum}
Let $G$ act smoothly on $M$. Then $G \backslash M$ is a Hausdorff space if the set $\Gamma:=\{(g p, p): g \in G, p \in M\}$ is a closed subset of $M \times M$. If $M$ is second countable then $G \backslash M$ is also.
\end{PROP}

Proof. 

\begin{PROOF}{LIE-L09-07-48}{P: Induzierte Hausdorff / 2. Abzählbarkeit für Orbitraum}
Let $\underline{p}, \underline{q} \in G \backslash M$ with $\pi(p)=\underline{p}$ and $\pi(q)=\underline{q}$. If $\underline{p} \neq \underline{q}$, then $p$ and $q$ are not in the same orbit. This means that $(p, q) \notin \Gamma$, and so there must be a product open set $U \times V$ such that $(p, q) \in U \times V$ and $U \times V$ is disjoint from $\Gamma$. This means that $\pi(U)$ and $\pi(V)$ are disjoint neighborhoods of $\underline{p}$ and $\underline{q}$ respectively. Finally, if $\left\{U_{i}\right\}$ is a countable basis for the topology on $M$, then $\left\{\pi\left(U_{i}\right)\right\}$ is a countable basis for the topology on $G \backslash M$.
\end{PROOF}

Proposition 5.118. 

\begin{PROP}{LIE-L09-07-49}{Freie und eigentliche Lie Gruppenwirkung induziert Hausdorff Orbitraum}
If $l: G \times M \rightarrow M$ is a free and proper action, then $G \backslash M$ is Hausdorff.
\end{PROP}

Proof. 

\begin{PROOF}{LIE-L09-07-50}{P: Freie und eigentliche Lie Gruppenwirkung induziert Hausdorff Orbitraum}
To show that $G \backslash M$ is Hausdorff, we use the previous lemma. We must show that $\Gamma$ is closed. By Problem 2, proper continuous maps are closed. Thus $\Gamma=P(G \times M)$ is closed since $P$ is proper.
\end{PROOF}

We will shortly show that if the action is free and proper, then $G \backslash M$ has a smooth structure which makes the quotient map $\pi: M \rightarrow G \backslash M$ a submersion. Before coming to this let us note that if such a smooth structure exists, then it is unique and is determined by the smooth structure on $M$. Indeed, if $(G \backslash M)_{\mathcal{A}}$ is $G \backslash M$ with a smooth structure given by a maximal atlas $\mathcal{A}$ and similarly for $(G \backslash M)_{\mathcal{B}}$ for another atlas $\mathcal{B}$, then we have the following commutative diagram:
\[
\begin{tikzcd}[row sep=large, column sep=large]
& M \arrow[dl, "\pi"'] \arrow[dr, "\pi"] & \\
(G \backslash M)_A \arrow[rr, "\mathrm{id}"] & & (G \backslash M)_B
\end{tikzcd}
\]
Since $\pi$ is a surjective submersion, Proposition 3.26 applies to show that $(G \backslash M)_{\mathcal{A}} \xrightarrow{\text { id }}(G \backslash M)_{\mathcal{B}}$ is smooth as is its inverse. This means that $\mathcal{A}=\mathcal{B}$.




\begin{PROP}{LIE-L09-07-51}{Eindeutigkeit der glatten Struktur auf dem Orbitraum}
Sei $l:G\times M\to M$ eine glatte Aktion einer Lie-Gruppe $G$ auf einer glatten Mannigfaltigkeit $M$.
Angenommen, es existiert eine glatte Struktur auf dem Quotientenraum $G\backslash M$,
die die Projektion
\[
\pi:M\longrightarrow G\backslash M,\qquad p\longmapsto G\cdot p,
\]
zu einer glatten Submersion macht.
Dann ist diese glatte Struktur \emph{eindeutig} durch die glatte Struktur von $M$ bestimmt.
\end{PROP}



\begin{PROOF}{LIE-L09-07-52}{P: Eindeutigkeit der glatten Struktur auf dem Orbitraum}
Seien $(G\backslash M)_{\mathcal A}$ und $(G\backslash M)_{\mathcal B}$
zwei mögliche glatte Strukturen auf $G\backslash M$
mit zugehörigen maximalen Atlanten $\mathcal A$ bzw.\ $\mathcal B$.
In beiden Fällen ist die Projektion
\[
\pi: M \longrightarrow (G\backslash M)_{\mathcal A}
\quad\text{bzw.}\quad
\pi: M \longrightarrow (G\backslash M)_{\mathcal B}
\]
eine glatte surjektive Submersion.
Damit erhalten wir das kommutative Diagramm:
\[
\begin{tikzcd}[row sep=large, column sep=large]
& M \arrow[dl, "\pi"'] \arrow[dr, "\pi"] & \\
(G\backslash M)_{\mathcal A} \arrow[rr, "\mathrm{id}"] & & (G\backslash M)_{\mathcal B}
\end{tikzcd}
\]

Da $\pi$ in beiden Fällen eine surjektive Submersion ist,
kann Proposition~3.26 angewendet werden:
Sie besagt, dass eine Abbildung $f$ zwischen Mannigfaltigkeiten genau dann glatt ist,
wenn $f\circ\pi$ glatt ist.
Hieraus folgt, dass die Identität
\[
\mathrm{id}:(G\backslash M)_{\mathcal A}\longrightarrow (G\backslash M)_{\mathcal B}
\]
glatt ist, ebenso wie ihr Inverses.
Somit sind die beiden glatten Strukturen äquivalent, d.\,h.\ $\mathcal A=\mathcal B$.
\end{PROOF}





Theorem 5.119. 

\begin{THEO}{LIE-L09-07-53}{Eindeutig induzierte glatte Struktur auf Orbitraum}
If $l: G \times M \rightarrow M$ is a free and proper Lie group action, then there is a unique smooth structure on the quotient $G \backslash M$ such that
\begin{enumerate}
\item The induced topology is the quotient topology, and $G \backslash M$ is a smooth manifold
\item The projection $\pi: M \rightarrow G \backslash M$ is a submersion
\item $\operatorname{dim}(G \backslash M)=\operatorname{dim}(M)-\operatorname{dim}(G)$
\end{enumerate}
\end{THEO}

Proof. 

\begin{PROOF}{LIE-L09-07-54}{P: Eindeutig induzierte glatte Struktur auf Orbitraum}
Let $\operatorname{dim}(M)=n$ and $\operatorname{dim}(G)=k$. We have already shown that $G \backslash M$ is a Hausdorff space. All that is left is to exhibit an atlas such that the charts are homeomorphisms with respect to this quotient topology. Let $q \in G \backslash M$ and choose $p$ with $\pi(p)=q$. Let ( $U, \mathrm{x}$ ) be an action-adapted chart centered at $p$ and constructed exactly as in Theorem 5.115. Let $\pi(U)=V \subset G \backslash M$ and let $B$ be the slice $x^{1}=\cdots=x^{k}=0$. By construction $\left.\pi\right|_{B}: B \rightarrow V$ is a bijection. In fact, it is easy to check that $\left.\pi\right|_{B}$ is a homeomorphism, and $\sigma:=\left(\left.\pi\right|_{B}\right)^{-1}$ is the corresponding local section. Consider the map $\mathrm{y}=\pi_{2} \circ \mathrm{x} \circ \sigma$, where $\pi_{2}$ is the second factor projection $\pi_{2}: \mathbb{R}^{k} \times \mathbb{R}^{n-k} \rightarrow \mathbb{R}^{n-k}$. This is a homeomorphism since $\left.\left(\pi_{2} \circ \mathrm{x}\right)\right|_{B}$ is a homeomorphism and $\pi_{2} \circ \mathrm{x} \circ \sigma=\left.\left(\pi_{2} \circ \mathrm{x}\right)\right|_{B} \circ \sigma$. We now have a chart ($V, \mathrm{y}$).

Given two such charts ( $V, \mathrm{y}$ ) and ( $\bar{V}, \overline{\mathrm{y}}$ ), we must show that $\overline{\mathrm{y}} \circ \mathrm{y}^{-1}$ is smooth. The ( $V, \mathrm{y}$ ) and ( $\bar{V}, \overline{\mathrm{y}}$ ) are constructed from associated actionadapted charts ( $U, \mathrm{x}$ ) and ( $\bar{U}, \overline{\mathrm{x}}$ ) on $M$. Let $q \in V \cap \bar{V}$. As in the proof of Lemma 5.116, we may find diffeomorphisms $\psi$ and $\bar{\psi}$ such that ( $U, \psi \circ \mathrm{x}$ ) and $(\bar{U}, \bar{\psi} \circ \overline{\mathrm{x}})$ are action-adapted charts centered at points $p_{1} \in \pi^{-1}(q)$ and $p_{2} \in \pi^{-1}(q)$ respectively. Correspondingly, the charts ( $V, \mathrm{y}$ ) and ( $\bar{V}, \overline{\mathrm{y}}$ ) are modified to charts $\left(V, \mathrm{y}_{\psi}\right)$ and $\left(\bar{V}, \overline{\mathrm{y}}_{\bar{\psi}}\right)$ centered at $q$, where
$$
\begin{aligned}
& \mathrm{y}_{\psi}:=\pi_{2} \circ \mathrm{x}_{\psi} \circ \sigma^{\prime} \\
& \overline{\mathrm{y}}_{\bar{\psi}}:=\pi_{2} \circ \overline{\mathrm{x}}_{\psi} \circ \bar{\sigma}^{\prime}
\end{aligned}
$$
with $\mathrm{x}_{\psi}:=\psi \circ \mathrm{x}$ and similarly for $\overline{\mathrm{x}}_{\psi}$. Also, the sections $\sigma^{\prime}$ and $\bar{\sigma}^{\prime}$ are constructed to map into the zero slices of $\mathrm{x}_{\psi}$ and $\overline{\mathrm{x}}_{\psi}$. However, it is not hard to see that $\mathrm{y}_{\psi}$ and $\overline{\mathrm{y}}_{\bar{\psi}}$ are unchanged if we replace $\sigma^{\prime}$ and $\bar{\sigma}^{\prime}$ by the sections $\sigma$ and $\bar{\sigma}$ corresponding to the zero slices of x and y . Now, recall that $\psi$ was chosen to have the form $(a, b) \mapsto\left(\psi_{1}(a), \psi_{2}(b)\right)$. Using the above observations, one checks that $\mathrm{y}_{\psi} \circ \mathrm{y}^{-1}=\psi_{2}$ and similarly for $\overline{\mathrm{y}}_{\bar{\psi}} \circ \overline{\mathrm{y}}^{-1}$. From this it follows that the overlap map $\overline{\mathrm{y}}_{\bar{\psi}}^{-1} \circ \mathrm{y}_{\psi}$ will be smooth if and only if $\overline{\mathrm{y}}^{-1} \circ \mathrm{y}^{-1}$ is smooth. Thus we have reduced to the case where ( $U, \mathrm{x}$ ) and $(\bar{U}, \overline{\mathrm{x}})$ are centered at $p_{1} \in \pi^{-1}(q)$ and $p_{2} \in \pi^{-1}(q)$ respectively. This entails that both ( $V, \mathrm{y}$ ) and ( $\bar{V}, \overline{\mathrm{y}}$ ) are centered at $q \in V \cap \bar{V}$. If we choose a $g \in G$ such that $l_{g}\left(p_{1}\right)=p_{2}$, then by composing with the diffeomorphism $l_{g}$ we can reduce further to the case where $p_{1}=p_{2}$. Here we use the fact that $l_{g}$ takes the set of orbits to the set of orbits in a bijective manner and the special nature of our action-adapted charts with respect to these orbits. In this case, the overlap map $\overline{\mathrm{x}} \circ \mathrm{x}^{-1}$ must have the form $(a, b) \mapsto(f(a, b), g(b))$ for some smooth functions $f$ and $g$. It follows that $\overline{\mathrm{y}} \circ \mathrm{y}^{-1}$ has the form $b \mapsto g(b)$.

Finally, we give an argument that $G \backslash M$ is paracompact. Since $G \backslash M$ is Hausdorff and locally Euclidean, we will be done once we show that every connected component of $G \backslash M$ is second countable (see Proposition B.5). Let $Q_{0}$ be any such connected component of $G \backslash M$. Let $X_{0}$ be a connected component of $\pi^{-1}\left(Q_{0}\right)$. Then $X_{0}$ is a second countable manifold and open in $M$. We now argue that $\pi\left(X_{0}\right)=Q_{0}$. To this end we show that the connected set $\pi\left(X_{0}\right)$ is open and closed in $Q_{0}$. It is open since $\pi$ is an open map. Let $\bar{x}$ be in the closure of $\pi\left(X_{0}\right)$ and choose $x \in M$ with $\pi(x)=\bar{x}$. Let $U$ be the domain of an action-adapted chart centered at $x$ and let $S$ be the slice of $U$ that maps diffeomorphically onto an open neighborhood $O$ of $\bar{x}$. Now we may find $\bar{y} \in O \cap \pi\left(X_{0}\right)$ and a corresponding $y \in S$ such that $\pi(y)=\bar{y}$. Since $\bar{y}$ is in the image of $X_{0}$, there is a $y^{\prime} \in X_{0}$ such that $\pi\left(y^{\prime}\right)=\bar{y}$. But then there exists $g \in G$ such that $g y=y^{\prime}$. The open set $g U$ is diffeomorphic to $U$ under $l_{g}$ and hence is path connected. It contains $y^{\prime}$ and $g x$. Since $X_{0}$ is a path component, $g U \subset X_{0}$ and so $g x \in X_{0}$. But $\bar{x}=\pi(g x)$ so $\bar{x} \in \pi\left(X_{0}\right)$. We conclude that $\pi\left(X_{0}\right)$ is closed (and open) and connected. Hence $\pi\left(X_{0}\right)=Q_{0}$. Now since $X_{0}$ is a second countable manifold, we argue as in Proposition 5.117 that $Q_{0}$ is second countable. Conclusion: $G \backslash M$ is paracompact.
\end{PROOF}



Similar results hold for right actions. Some of the most important examples of proper actions are usually presented as right actions. In fact, principal bundle actions studied in Chapter 6 are usually presented as right actions. We shall also encounter situations where there are both a right and a left action in play.

Example 5.120. 

\begin{EXA}{LIE-L09-07-55}{Hopf Abbildung für $\mathbb{C}P^{n-1}$}
Consider $S^{2 n-1}$ as the subset of $\mathbb{C}^{n}$ given by $S^{2 n-1}= \left\{\xi \in \mathbb{C}^{n}:|\xi|=1\right\}$. Here $\xi=\left(z^{1}, \ldots, z^{n}\right)$ and $|\xi|=\sum \bar{z}^{i} z^{i}$. Let $S^{1}$ act on $S^{2 n-1}$ by $(a, \xi) \mapsto a \xi=\left(a z^{1}, \ldots, a z^{n}\right)$. This action is free and proper. The quotient is the complex projective space $\mathbb{C} P^{n-1}$,
\[
\begin{tikzcd}
S^{2n-1} \arrow[d] \\
\mathbb{C}P^{\,n-1}
\end{tikzcd}
\]
These maps (one for each $n$ ) are called the Hopf maps. In this context, $S^{1}$ is usually denoted by $U(1)$.
\end{EXA}

In what follows, we will consider the similar right action $S^{n} \times U(1) \rightarrow S^{n}$. In this case, we think of $\mathbb{C}^{n+1}$ as a set of column vectors, and the action is given by $(\xi, a) \mapsto \xi a$. Of course, since $U(1)$ is abelian, this makes essentially no difference, but in the next example we consider the quaternionic analogue where keeping track of order is important.

\begin{EXA}{LIE-L09-07-56}{Quaternionischer Projektiveer Raum $\mathbb{H}P^{n-1}$}
The quaternionic projective space $\mathbb{H} P^{n-1}$ is defined by analogy with $\mathbb{C} P^{n-1}$. The elements of $\mathbb{H} P^{n-1}$ are 1 -dimensional subspaces of the right $\mathbb{H}$-vector space $\mathbb{H}^{n}$. Let us refer to these as $\mathbb{H}$-lines. Each of these are of real dimension 4. Each element of $\mathbb{H}^{n} \backslash\{0\}$ determines an $\mathbb{H}$-line and the $\mathbb{H}$-line determined by $\left(\xi^{1}, \ldots, \xi^{n}\right)^{t}$ will be the same as that determined by $\left(\widetilde{\xi}^{1}, \ldots, \widetilde{\xi}^{n}\right)^{t}$ if and only if there is a nonzero element $a \in \mathbb{H}$ such that 
$$\left(\widetilde{\xi}^{1}, \ldots, \widetilde{\xi}^{n}\right)^{t}=\left(\xi^{1}, \ldots, \xi^{n}\right)^{t} a=\left(\xi^{1} a, \ldots, \xi^{n} a\right)^{t}$$ 
This defines an equivalence relation $\sim$ on $\mathbb{H}^{n} \backslash\{0\}$, and thus we may also think of $\mathbb{H} P^{n-1}$ as $\left(\mathbb{H}^{n} \backslash\{0\}\right) / \sim$. 

The element of $\mathbb{H} P^{n-1}$ determined by $\left(\xi^{1}, \ldots, \xi^{n}\right)^{t}$ is denoted by $\left[\xi^{1}, \ldots, \xi^{n}\right]$. 

Notice that the subset $\left\{\xi \in \mathbb{H}^{n}:|\xi|=1\right\}$ is $S^{4 n-1}$. Just as for the complex projective spaces, we observe that all such $\mathbb{H}$-lines contain points of $S^{4 n-1}$, and two points $\xi, \zeta \in S^{4 n-1}$ determine the same $\mathbb{H}$-line if and only if $\xi=\zeta a$ for some $a$ with $|a|=1$. Thus we can think of $\mathbb{H} P^{n-1}$ as a quotient of $S^{4 n-1}$. When viewed in this way, we also denote the equivalence class of $\xi=\left(\xi^{1}, \ldots, \xi^{n}\right)^{t} \in S^{4 n-1}$ by $[\xi]=\left[\xi^{1}, \ldots, \xi^{n}\right]$. The equivalence classes are clearly the orbits of an action as described in the following example.
\end{EXA}

Example 5.121. 

\begin{EXA}{LIE-L09-07-57}{Hopf Abbildung für $\mathbb{H}P^{n-1}$}
Consider $S^{4 n-1}$ regarded as the subset of $\mathbb{H}^{n}$ given by $S^{4 n-1}=\left\{\xi \in \mathbb{H}^{n}:|\xi|=1\right\}$. Here $\xi=\left(\xi^{1}, \ldots, \xi^{n}\right)^{t}$ and $|\xi|=\sum \bar{\xi}^{i} \xi^{i}$. Now we define a right action of $U(1, \mathbb{H})$ on $S^{4 n-1}$ by $(\xi, a) \mapsto \xi a=\left(\xi^{1} a, \ldots, \xi^{n} a\right)^{t}$. This action is free and proper. The quotient is the quaternionic projective space $\mathbb{H} P^{n-1}$, and we have the quotient map denoted by $\wp$,
\[
\begin{tikzcd}
S^{4n-1} \arrow[d, "\varphi"] \\
\mathbb{H}P^{n-1}
\end{tikzcd}
\]
This map is also referred to as a Hopf map. Recall that $\mathbb{Z}_{2}=\{1,-1\}$ acts on $S^{n-1}=\mathbb{R}^{n}$ on the right (or left) by multiplication, and the action is a (discrete) proper and free action with quotient $\mathbb{R} P^{n-1}$, and the examples above generalize this.
\end{EXA}

\begin{EXA}{LIE-L09-07-58}{Glatte Struktur auf $\mathbb{H}P^{n-1}$}
For completeness, we describe an atlas for $\mathbb{H} P^{n-1}$. View $\mathbb{H} P^{n-1}$ as the quotient $S^{4 n-1} / \sim$ described above. Let
$$
U_{k}:=\left\{[\xi] \in S^{4 n-1} \subset \mathbb{H}^{n}: \xi^{k} \neq 0\right\},
$$
and define $\varphi_{k}: U_{k} \rightarrow \mathbb{H}^{n-1} \cong \mathbb{R}^{4 n-1}$ by
$$
\varphi_{k}([\xi])=\left(\xi_{1} \xi_{k}^{-1}, \ldots, \widehat{1}, \ldots, \xi_{n} \xi_{k}^{-1}\right),
$$
where as before, the caret symbol ${ }^{\wedge}$ indicates that we have omitted the 1 in the $k$-th slot to obtain an element of $\mathbb{H}^{n-1}$. Notice that we insist that the $\xi_{k}^{-1}$ in this expression multiply from the right. The general pattern for the overlap maps become clear from the example $\varphi_{3} \circ \varphi_{2}^{-1}$. Here we have
$$
\begin{aligned}
\varphi_{3} \circ \varphi_{2}^{-1}\left(y_{1}, y_{3}, \ldots, y_{n}\right) & =\varphi_{3}\left(\left[y_{1}, 1, y_{3}, \ldots, y_{n}\right]\right) \\
& =\left(y_{1} y_{3}^{-1}, y_{3}^{-1}, y_{4} y_{3}^{-1}, \ldots, y_{n} y_{3}^{-1}\right)
\end{aligned}
$$
In the case $n=1$, we have an atlas of just two charts $\left\{\left(U_{1}, \varphi_{1}\right),\left(U_{2}, \varphi_{2}\right)\right\}$. In close analogy with the complex case we have $U_{1} \cap U_{2}=\mathbb{H} \backslash\{0\}$ and $\varphi_{1} \circ \varphi_{2}^{-1}(y)=\bar{y}^{-1}=\varphi_{2} \circ \varphi_{1}^{-1}(y)$ for $y \in \mathbb{H} \backslash\{0\}$.
\end{EXA}

Exercise 5.122. 


\begin{EXE}{LIE-L09-07-59}{$\mathbb{H}P^{n-1}\cong S^3, \mathbb{R} P^{1} \cong S^{1}, \mathbb{C} P^{1} \cong S^{2}, \mathbb{H} P^{1} \cong S^{3}$}
Show that by identifying $\mathbb{H}$ with $\mathbb{R}^{4}$ and modifying the stereographic charts on $S^{3} \subset \mathbb{R}^{4}$ we can obtain an atlas for $S^{3}$ with overlap maps of the same form as for $\mathbb{H} P^{1}$ given above. Use this to show that $\mathbb{H} P^{1} \cong S^{3}$.

Combining the last exercise with previous results we have
$$
\mathbb{R} P^{1} \cong S^{1}, \quad \mathbb{C} P^{1} \cong S^{2}, \quad \mathbb{H} P^{1} \cong S^{3}
$$
\end{EXE}


\pagebreak



\subsection*{Homogeneous Spaces}
Let $H$ be a closed Lie subgroup of a Lie group $G$. Then we have a right action of $H$ on $G$ given by right multiplication $r: G \times H \rightarrow G$. The orbits of this right action are exactly the left cosets of the quotient $G / H$. The action is clearly free, and we would like to show that it is also proper. Since we are now talking about a right action, and $G$ is the manifold on which we are acting, we need to show that the map $P_{\text {right }}: G \times H \rightarrow G \times G$ given by $(p, h) \mapsto(p, p h)$ is a proper map. The characterization of proper action becomes the condition that

$$
H_{K}:=\{h \in H:(K \cdot h) \cap K \neq \emptyset\}
$$

is a compact subset of $H$ whenever $K$ is compact in $G$. Let $K$ be any compact subset of $G$. It will suffice to show that $H_{K}$ is sequentially compact. To this end, let $\left(h_{i}\right)_{i \in \mathbb{N}}$ be a sequence in $H_{K}$. Then there must be sequences ( $a_{i}$ ) and ( $b_{i}$ ) in $K$ such that $a_{i} h_{i}=b_{i}$. Since $K$ is compact and hence sequentially compact, we can pass to subsequences $\left(a_{i(j)}\right)_{j \in \mathbb{N}}$ and $\left(b_{i(j)}\right)_{j \in \mathbb{N}}$ such that $\lim _{j \rightarrow \infty} a_{i(j)}=a$ and $\lim _{j \rightarrow \infty} b_{i(j)}=b$. Here $i \mapsto i(j)$ is a monotonic map on positive integers; $\mathbb{N} \rightarrow \mathbb{N}$. This means that $\lim _{j \rightarrow \infty} h_{i(j)}=\lim _{j \rightarrow \infty} a_{i(j)}^{-1} b_{i(j)}=a^{-1} b$. Thus the original sequence $\left\{h_{i}\right\}$ is shown to have a convergent subsequence. Since by Theorem 5.81, $H$ is an embedded submanifold, this sequence converges in the topology of $H$. We conclude that the right action is proper. Using Theorem 5.119 (or its analogue for right actions), we obtain the following Proposition.

Proposition 5.123. Let $H$ be a closed Lie subgroup of a Lie group $G$. Then,\\
(i) the right action $G \times H \rightarrow G$ is free and proper;\\
(ii) the orbit space is the left coset space $G / H$, and this has a unique smooth manifold structure such that the quotient map $\pi: G \rightarrow G / H$ is a surjection. Furthermore, $\operatorname{dim}(G / H)=\operatorname{dim}(G)-\operatorname{dim}(H)$.

If $K$ is a normal Lie subgroup of $G$, then the quotient is a group with multiplication defined by $\left[g_{1}\right]\left[g_{2}\right]=\left(g_{1} K\right)\left(g_{2} K\right)=g_{1} g_{2} K$. In this case, we may ask whether $G / K$ is a Lie group. If $K$ is closed, then we know from the considerations above that $G / K$ is a smooth manifold and that the quotient map is smooth. In fact, we have the following:

Proposition 5.124 (Quotient Lie groups). If $K$ is a closed normal subgroup of a Lie group $G$, then $G / K$ is a Lie group and the quotient map $G \rightarrow G / K$ is a Lie group homomorphism. Furthermore, if $f: G \rightarrow H$ is a surjective Lie group homomorphism, then $\operatorname{Ker}(f)$ is a closed normal subgroup, and the induced map $\widetilde{f}: G / \operatorname{Ker}(f) \rightarrow H$ is a Lie group isomorphism.

Proof. We have already observed that $G / K$ is a smooth manifold and that the quotient map is smooth. After taking into account what we know from standard group theory, the only thing we need to prove for the first part, is that the multiplication and inversion in the quotient are smooth. It is an easy exercise using Corollary 3.27 to show that both of these maps are smooth.

Consider a Lie group homomorphism $f$ as in the hypothesis of the proposition. It is standard that $\operatorname{Ker}(f)$ is a normal subgroup and it is clearly closed. It is also easy to verify fact that the induced $\widetilde{f}$ map is an isomorphism. One can then use Corollary 3.27 to show that the induced map $\widetilde{f}$ is smooth.

If a Lie group $G$ acts smoothly and transitively on $M$ (on the right or left), then $M$ is called a homogeneous space with respect to that action. Of course it is possible that a single group $G$ may act on $M$ in more than one way and so $M$ may be a homogeneous space in more than one way. We will give a few concrete examples shortly, but we already have an abstract example on hand.

Theorem 5.125. If $H$ is a closed Lie subgroup of a Lie group $G$, then the map $G \times G / H \rightarrow G / H$, given by $l:\left(g, g_{1} H\right) \rightarrow g g_{1} H$, is a transitive Lie group action. Thus $G / H$ is a homogeneous space with respect to this action.

Proof. The fact that $l$ is well-defined follows since if $g_{1} H=g_{2} H$, then $g_{2}^{-1} g_{1} \in H$, and so $g g_{2} H=g g_{2} g_{2}^{-1} g_{1} H=g g_{1} H$. We already know that $G / H$ is a smooth manifold and $\pi: G \rightarrow G / H$ is a surjective submersion.

We can form another submersion $\operatorname{id}_{G} \times \pi: G \times G \rightarrow G \times G / H$ making the following diagram commute:\\
\includegraphics[scale=0.25, center]{2025_10_09_265d7e1a22c89f79c21eg-257}

Here the upper horizontal map is group multiplication and the lower horizontal map is the action $l$. Since the diagonal map is smooth, it follows from Proposition 3.26 that $l$ is smooth. We see that $l$ is transitive by observing that if $g_{1} H, g_{2} H \in G / H$, then $l_{g_{2} g_{1}^{-1}}\left(g_{1} H\right)=g_{2} H$. \(\square\)

It turns out that up to appropriate equivalence, the examples of the above type account for all homogeneous spaces. Before proving this let us look at some concrete examples.

Example 5.126. Let $M=\mathbb{R}^{n}$ and let $G=\operatorname{Euc}(n, \mathbb{R})$ be the group of Euclidean motions. We realize $\operatorname{Euc}(n, \mathbb{R})$ as a matrix group

$$
\operatorname{Euc}(n, \mathbb{R})=\left\{\left[\begin{array}{cc}
1 & 0 \\
v & Q
\end{array}\right]: v \in \mathbb{R}^{n} \text { and } Q \in \mathrm{O}(n)\right\}
$$

The action of $\operatorname{Euc}(n, \mathbb{R})$ on $\mathbb{R}^{n}$ is given by the rule

$$
\left[\begin{array}{ll}
1 & 0 \\
v & Q
\end{array}\right] \cdot x=Q x+v,
$$

where $x$ is written as a column vector. Notice that this action is not given by a matrix multiplication, but one can use the trick of representing the points $x$ of $\mathbb{R}^{n}$ by the $(n+1) \times 1$ column vectors $\left[\begin{array}{c}1 \\ x\end{array}\right]$, and then we have $\left[\begin{array}{ll}1 & 0 \\ v & Q\end{array}\right]\left[\begin{array}{l}1 \\ x\end{array}\right]=\left[\begin{array}{c}1 \\ Q x+v\end{array}\right]$. The action is easily seen to be transitive.\\
Example 5.127. As in the previous example we take $M=\mathbb{R}^{n}$, but this time, the group acting is the affine group $\operatorname{Aff}(n, \mathbb{R})$ realized as a matrix group:

$$
\operatorname{Aff}(n, \mathbb{R})=\left\{\left[\begin{array}{cc}
1 & 0 \\
v & A
\end{array}\right]: v \in \mathbb{R}^{n} \text { and } A \in \mathrm{GL}(n, \mathbb{R})\right\} .
$$

The action is

$$
\left[\begin{array}{ll}
1 & 0 \\
v & A
\end{array}\right] \cdot x=A x+v,
$$

and this is again a transitive action.\\
Comparing these first two examples, we see that we have made $\mathbb{R}^{n}$ into a homogeneous space in two different ways. It is sometimes desirable to give different names and/or notations for $\mathbb{R}^{n}$ to distinguish how we are acting on the space. In the first example we might write $\mathbb{E}^{n}$ (Euclidean space), and\\
in the second case we write $\mathbb{A}^{n}$ and refer to it as affine space. Note that, roughly speaking, the action by $\operatorname{Euc}(n, \mathbb{R})$ preserves all metric properties of figures such as curves defined in $\mathbb{E}^{n}$. On the other hand, $\operatorname{Aff}(n, \mathbb{R})$ always sends lines to lines, planes to planes, etc.

Example 5.128. Let $M=\mathrm{H}:=\{z \in \mathbb{C}: \operatorname{Im} z>0\}$. This is the upper half-plane. The group acting on H will be $\mathrm{SL}(2, \mathbb{R})$, and the action is given by

$$
\left(\begin{array}{cc}
a & b \\
c & d
\end{array}\right) \cdot z=\frac{a z+b}{c z+d}
$$

This action is transitive.\\
Example 5.129. We have already seen in Example 5.102 that both $\mathrm{O}(n)$ and $\mathrm{SO}(n)$ act transitively on the sphere $S^{n-1} \subset \mathbb{R}^{n}$, so $S^{n-1}$ is a homogeneous space in at least two (slightly) different ways. Also, both $\operatorname{SU}(n)$ and $\mathrm{U}(n)$ act transitively on $S^{2 n-1} \subset \mathbb{C}^{n}$.

Example 5.130. Let $V_{n, k}^{\prime}$ denote the set of all $k$-frames for $\mathbb{R}^{n}$, where by a $k$-frame we mean an ordered set of $k$ linearly independent vectors. Thus an $n$-frame is just an ordered basis for $\mathbb{R}^{n}$. This set can easily be given a smooth manifold structure. This manifold is called the (real) Stiefel manifold of $k$-frames. The Lie group $\operatorname{GL}(n, \mathbb{R})$ acts (smoothly) on $V_{n, k}^{\prime}$ by $g \cdot\left(e_{1}, \ldots, e_{k}\right)=\left(g e_{1}, \ldots, g e_{k}\right)$. To see that this action is transitive, let $\left(e_{1}, \ldots, e_{k}\right)$ and $\left(f_{1}, \ldots, f_{k}\right)$ be two $k$-frames. Extend each to $n$-frames $\left(e_{1}, \ldots, e_{k}, \ldots, e_{n}\right)$ and $\left(f_{1}, \ldots, f_{k}, \ldots, f_{n}\right)$. Since we consider elements of $\mathbb{R}^{n}$ as column vectors, these two $n$-frames can be viewed as invertible $n \times n$ matrices $E$ and $F$. If we let $g:=E F^{-1}$, then $g E=F$, or $g \cdot\left(e_{1}, \ldots, e_{k}\right)= \left(g e_{1}, \ldots, g e_{k}\right)=\left(f_{1}, \ldots, f_{k}\right)$.

Example 5.131. Let $V_{n, k}$ denote the set of all orthonormal $k$-frames for $\mathbb{R}^{n}$, where by an orthonormal $k$-frame we mean an ordered set of $k$ orthonormal vectors. Thus an orthonormal $n$-frame is just an orthonormal basis for $\mathbb{R}^{n}$. This set can easily be given a smooth manifold structure and is called the Stiefel manifold of orthonormal $k$-frames. The group $\mathrm{O}(n, \mathbb{R})$ acts transitively on $V_{n, k}$ for reasons similar to those given in the last example.

Theorem 5.132. Let $M$ be a homogeneous space via the transitive action $l: G \times M \rightarrow M$, and let $G_{p}$ be the isotropy subgroup of a point $p \in M$. Recall that $G$ acts on $G / G_{p}$. If $G / G_{p}$ is second countable (in particular if $G$ is second countable), then there is an equivariant diffeomorphism $\phi$ : $G / G_{p} \rightarrow M$ such that $\phi\left(g G_{p}\right)=g \cdot p$.

Proof. We want to define $\phi$ by the rule $\phi\left(g G_{p}\right)=g \cdot p$, but we must show that this is well-defined. This is a standard group theory argument; if $g_{1} G_{p}=g_{2} G_{p}$, then $g_{1}^{-1} g_{2} \in G_{p}$, so that $\left(g_{1}^{-1} g_{2}\right) \cdot p=p$ or $g_{1} \cdot p=g_{2} \cdot p$.

This map is surjective by the transitivity of the action $l$. It is also injective since if $\phi\left(g_{1} G_{p}\right)=\phi\left(g_{2} G_{p}\right)$, then $g_{1} \cdot p=g_{2} \cdot p$ or $\left(g_{1}^{-1} g_{2}\right) \cdot p=p$, which by definition means that $g_{1}^{-1} g_{2} \in G_{p}$ and $g_{1} G_{p}=g_{2} G_{p}$. Notice that the following diagram commutes:\\
\includegraphics[scale=0.25, center]{2025_10_09_265d7e1a22c89f79c21eg-259}

From Corollary 3.27 we see that $\phi$ is smooth.\\
To show that $\phi$ is a diffeomorphism, it suffices to show that the rank of $\phi$ is equal to $\operatorname{dim} M$ or in other words that $\phi$ is a submersion. Since $\phi\left(g g_{1} G_{p}\right)=\left(g g_{1}\right) \cdot p=g \phi\left(g_{1} G_{p}\right)$, the map $\phi$ is equivariant and so has constant rank. By Lemma 3.28, $\phi$ is a submersion and hence in the present case a diffeomorphism. \(\square\)

Without the technical assumption on second countability, the proof shows that we still have that $\phi: G / G_{p} \rightarrow M$ is a smooth equivariant bijection.

Exercise 5.133. Show that if instead of the hypothesis of second countability in the last theorem we assume that $\theta_{p}$ has full rank at the identity, then $\phi: G / G_{p} \rightarrow M$ is a diffeomorphism.

Let $l: G \times M \rightarrow M$ be a left Lie group action and fix $p_{0} \in M$. Denote the projection onto cosets by $\pi$ and also write $\theta_{p_{0}}: g \mapsto g p_{0}$ as before. Then we have the following equivalence of maps:\\
\includegraphics[scale=0.25, center]{2025_10_09_265d7e1a22c89f79c21eg-259(1)}

Exercise 5.134. Let $G$ act on $M$ as above. Show that if $p_{2}=g p_{1}$ for some $g \in G$ and $p_{1}, p_{2} \in M$, then is a natural Lie group isomorphism $G_{p_{1}} \cong G_{p_{2}}$ and a natural equivariant diffeomorphism $G / G_{p_{1}} \cong G / G_{p_{2}}$.

We now look again at some of our examples of homogeneous spaces and apply the above theorem.

Example 5.135. Consider again Example 5.126. The isotropy group of the origin in $\mathbb{R}^{n}$ is the subgroup consisting of matrices of the form

$$
\left(\begin{array}{cc}
1 & 0 \\
0 & Q
\end{array}\right),
$$

where $Q \in \mathrm{O}(n)$. This group is clearly isomorphic to $\mathrm{O}(n, \mathbb{R})$, and so by the above theorem we have an equivariant diffeomorphism

$$
\mathbb{R}^{n} \cong \frac{\operatorname{Euc}(n, \mathbb{R})}{\mathrm{O}(n, \mathbb{R})}
$$

Example 5.136. Consider again Example 5.127. The isotropy group of the origin in $\mathbb{R}^{n}$ is the subgroup consisting of matrices of the form

$$
\left(\begin{array}{ll}
1 & 0 \\
0 & A
\end{array}\right),
$$

where $A \in \mathrm{GL}(n, \mathbb{R})$. This group is clearly isomorphic to $\mathrm{GL}(n, \mathbb{R})$, and so by the above theorem we have an equivariant diffeomorphism

$$
\mathbb{R}^{n} \cong \frac{\operatorname{Aff}(n, \mathbb{R})}{\operatorname{GL}(n, \mathbb{R})}
$$

It is important to realize that there is an implied action on $\mathbb{R}^{n}$, which is different from that in the previous example.

Example 5.137. Consider the action of $\operatorname{SL}(2, \mathbb{R})$ on the complex upper halfplane $H=\mathbb{C}_{+}$as in Example 5.128. We determine the isotropy subgroup for the point $i=\sqrt{-1}$. A matrix $A=\left(\begin{array}{ll}a & b \\ c & d\end{array}\right)$ is in this subgroup if and only if

$$
\frac{a i+b}{c i+d}=i
$$

This is true exactly if $b c-a d=1$ and $b d+a c=0$, and so the isotropy subgroup is $\mathrm{SO}(2, \mathbb{R})\left(\cong S^{1}=U(1, \mathbb{C})\right)$. Thus we have an equivariant diffeomorphism

$$
H=\mathbb{C}_{+} \cong \frac{\operatorname{SL}(2, \mathbb{R})}{\operatorname{SO}(2, \mathbb{R})}
$$

Example 5.138. From Example 5.129 we obtain

$$
\begin{aligned}
S^{n-1} & \cong \frac{\mathrm{O}(n)}{\mathrm{O}(n-1)}, & S^{n-1} & \cong \frac{\mathrm{SO}(n)}{\mathrm{SO}(n-1)} \\
S^{2 n-1} & \cong \frac{\mathrm{U}(n)}{\mathrm{U}(n-1)}, & S^{2 n-1} & \cong \frac{\mathrm{SU}(n)}{\mathrm{SU}(n-1)}
\end{aligned}
$$

Example 5.139. Let $\left(\mathbf{e}_{1}, \ldots, \mathbf{e}_{n}\right)$ be the standard basis for $\mathbb{R}^{n}$. Under the action of $\operatorname{GL}(n, \mathbb{R})$ on $V_{n, k}^{\prime}$ given in Example 5.130, the isotropy group of the $k$-frame $\left(\mathbf{e}_{k+1}, \ldots, \mathbf{e}_{n}\right)$ is the subgroup of $\operatorname{GL}(n, \mathbb{R})$ of the form

$$
\left(\begin{array}{cc}
A & 0 \\
0 & I
\end{array}\right) \text { for } A \in \mathrm{GL}(n-k, \mathbb{R})
$$

We identify this group with $\operatorname{GL}(n-k, \mathbb{R})$ and then we obtain

$$
V_{n, k}^{\prime} \cong \frac{\mathrm{GL}(n, \mathbb{R})}{\mathrm{GL}(n-k, \mathbb{R})}
$$

Example 5.140. A similar analysis leads to an equivariant diffeomorphism

$$
V_{n, k} \cong \frac{\mathrm{O}(n, \mathbb{R})}{\mathrm{O}(n-k, \mathbb{R})}
$$

where $V_{n, k}$ is the Stiefel manifold of orthonormal $k$-planes of Example 5.131. Notice that taking $k=1$ we recover Example 5.135.\\
Exercise 5.141. Show that if $k<n$, then we have $V_{n, k} \cong \frac{\mathrm{SO}(n, \mathbb{R})}{\mathrm{SO}(n-k, \mathbb{R})}$.\\
Next we introduce a couple of standard results concerning connectivity.\\
Proposition 5.142. Let $G$ be a Lie group acting smoothly on $M$. Let the action be a left (resp. right) action. If both $G$ and $M \backslash G$ (resp. $M / G$ ) are connected, then $M$ is connected.

Proof. Assume for concreteness that the action is a left action and that $G$ and $M \backslash G$ are connected. Suppose by way of contradiction that $M$ is not connected. Then there are disjoint open sets $U$ and $V$ whose union is $M$. Each orbit $G \cdot p$ is the image of the connected space $G$ under the orbit map $g \mapsto g \cdot p$ and so is connected. This means that each orbit must be contained in one and only one of $U$ and $V$. Since the quotient map $\pi$ is an open map, $\pi(U)$ and $\pi(V)$ are open, and from what we have just observed they must be disjoint and $\pi(U) \cup \pi(V)=M \backslash G$. This contradicts the assumption that $M \backslash G$ is connected.

Corollary 5.143. Let $H$ be a closed Lie subgroup of $G$. If both $H$ and $G / H$ are connected, then $G$ is connected.

Corollary 5.144. For each $n \geq 1$, the groups $\mathrm{SO}(n), \mathrm{SU}(n)$ and $\mathrm{U}(n)$ are connected while the group $\mathrm{O}(n)$ has exactly two components: $\mathrm{SO}(n)$ and the subset of $\mathrm{O}(n)$ consisting of elements with determinant -1 .

Proof. The groups $\mathrm{SO}(1)$ and $\mathrm{SU}(1)$ are both connected since they each contain only one element. The group $U(1)$ is the circle, and so it too is connected. We use induction. Suppose that $\mathrm{SO}(k), \mathrm{SU}(k)$ are connected for $1 \leq k \leq n-1$. We show that this implies that $\mathrm{SO}(n), \mathrm{SU}(n)$ and $\mathrm{U}(n)$ are connected. From Example 5.138 we know that $S^{n-1}=\mathrm{SO}(n) / \mathrm{SO}(n-1)$. Since $S^{n-1}$ and $\mathrm{SO}(n-1)$ are connected (the second by the induction hypothesis), we see that $\mathrm{SO}(n)$ is connected. The same argument works for $\mathrm{SU}(n)$ and $\mathrm{U}(n)$.

Every element of $\mathrm{O}(n)$ has determinant either 1 or -1 . The subset $\mathrm{SO}(n) \subset \mathrm{O}(n)$ is closed since it is exactly $\{g \in \mathrm{O}(n): \operatorname{det} g=1\}$. Fix an element $a_{0}$ with $\operatorname{det} a_{0}=-1$. It is easy to show that $a_{0} \mathrm{SO}(n)$ is exactly the set of elements of $\mathrm{O}(n)$ with determinant -1 so that $\mathrm{SO}(n) \cup a_{0} \mathrm{SO}(n)=\mathrm{O}(n)$ and $\mathrm{SO}(n) \cap a_{0} \mathrm{SO}(n)=\emptyset$. Indeed, by the multiplicative property of determinants, each element of $a_{0} \mathrm{SO}(n)$ has determinant -1 . But $a_{0} \mathrm{SO}(n)$\\
also contains every element of determinant -1 since for any such $g$ we have $g=a_{0}\left(a_{0}^{-1} g\right)$ and $a_{0}^{-1} g \in \mathrm{SO}(n)$. Since $\mathrm{SO}(n)$ and $a_{0} \mathrm{SO}(n)$ are complements of each other, they are also both open. Both sets are connected, since $g \mapsto a_{0} g$ is a diffeomorphism which maps the first to the second. Thus we see that $\mathrm{SO}(n)$ and $a_{0} \mathrm{SO}(n)$ are the connected components of $\mathrm{O}(n)$.

We close this chapter by relating the notion of a Lie group action with that of a Lie group representation. We give just a few basic definitions, some of which will be used in the next chapter.

Definition 5.145. A linear action of a Lie group $G$ on a finite-dimensional vector space V is a left Lie group action $\lambda: G \times \mathrm{V} \rightarrow \mathrm{V}$ such that for each $g \in G$ the map $\lambda_{g}: v \mapsto \lambda(g, v)$ is linear.

The map $G \rightarrow \mathrm{GL}(\mathrm{V})$ given by $g \mapsto \lambda(g):=\lambda_{g}$ is a Lie group homomorphism and will be denoted by the same letter $\lambda$ as the action so that $\lambda(g) v:=\lambda(g, v)$. A Lie group homomorphism $\lambda: G \rightarrow \mathrm{GL}(\mathrm{V})$ is called a representation of $G$. Given such a representation, we obtain a linear action by letting $\lambda(g, v):=\lambda(g) v$. Thus a linear action of a Lie group is basically the same as a Lie group representation. The kernel of the action is the kernel of the associated homomorphism (representation). An effective linear action is one such that the associated homomorphism has trivial kernel, which, in turn, is the same as saying that the representation is faithful. Two representations $\lambda: G \rightarrow \mathrm{GL}(\mathrm{V})$ and $\lambda^{\prime}: G \rightarrow \mathrm{GL}\left(\mathrm{V}^{\prime}\right)$ are equivalent if there exists a linear isomorphism $T: \mathrm{V} \rightarrow \mathrm{V}^{\prime}$ such that $T \circ \lambda_{g}=\lambda_{g}^{\prime} \circ T$ for all $g$.

Exercise 5.146. Show that if $\lambda: G \times \mathrm{V} \rightarrow \mathrm{V}$ is a map such that $\lambda_{g}: v \mapsto \lambda(g, v)$ is linear for all $g$, then $\lambda$ is smooth if and only if $\lambda_{g}: G \rightarrow \mathrm{GL}(\mathrm{V})$ is smooth for every $g \in G$. (Assume that V is finite-dimensional as usual.)

We have already seen one important example of a Lie group representation, namely, the adjoint representation. The adjoint representation came from first considering the action of $G$ on itself given by conjugation which leaves the identity element fixed. The idea can be generalized:

Theorem 5.147. Let $l: G \times M \rightarrow M$ be a (left) Lie group action. Suppose that $p_{0} \in M$ is a fixed point of the action ( $l_{g}\left(p_{0}\right)=p_{0}$ for all $g$ ). The map

$$
l^{\left(p_{0}\right)}: G \rightarrow \mathrm{GL}\left(T_{p_{0}} M\right)
$$

given by

$$
l^{\left(p_{0}\right)}(g):=T_{p_{0}} l_{g}
$$

is a Lie group representation.

Proof. Since

$$
\begin{aligned}
l^{\left(p_{0}\right)}\left(g_{1} g_{2}\right) & =T_{p_{0}}\left(l_{g_{1} g_{2}}\right)=T_{p_{0}}\left(l_{g_{1}} \circ l_{g_{2}}\right) \\
& =T_{p_{0}} l_{g_{1}} \circ T_{p_{0}} l_{g_{2}}=l^{\left(p_{0}\right)}\left(g_{1}\right) l^{\left(p_{0}\right)}\left(g_{2}\right),
\end{aligned}
$$

we see that $l^{\left(p_{0}\right)}$ is a homomorphism. We must show that $l^{\left(p_{0}\right)}$ is smooth. It will be enough to show that $g \mapsto \alpha\left(T_{p_{0}} l_{g} \cdot v\right)$ is smooth for any $v \in T_{p_{0}} M$ and any $\alpha \in T_{p_{0}}^{*} M$. This will follow if we can show that for fixed $v_{0} \in T_{p_{0}} M$, the map $G \rightarrow T M$ given by $g \mapsto T_{p_{0}} l_{g} \cdot v_{0}$ is smooth. This map is a composition

$$
G \rightarrow T G \times T M \cong T(G \times M) \xrightarrow{T l} T M,
$$

where the first map is $g \mapsto\left(0_{g}, v_{0}\right)$, which is clearly smooth. By Exercise 5.146 this implies that the map $G \times T_{p_{0}} M \rightarrow T_{p_{0}} M$ given by $(g, v) \mapsto T_{p_{0}} l_{g} \cdot v$ is smooth.

Definition 5.148. For a Lie group action $l: G \times M \rightarrow M$ with fixed point $p_{0}$, the representation $l^{\left(p_{0}\right)}$ from the last theorem is called the isotropy representation for the fixed point.

Now let us consider a transitive Lie group action $l: G \times M \rightarrow M$ and a point $p_{0}$. For notational convenience, denote the isotropy subgroup $G_{p_{0}}$ by $H$. Then $H$ acts on $M$ by restriction. We denote this action by $\lambda: H \times M \rightarrow M$,

$$
\lambda:(h, p) \mapsto h p \text { for } h \in H=G_{p_{0}} .
$$

Notice that $p_{0}$ is a fixed point of this action, and so we have an isotropy representation $\lambda^{\left(p_{0}\right)}: H \times T_{p_{0}} M \rightarrow T_{p_{0}} M$. On the other hand, we have another action $C: H \times G \rightarrow G$, where $C_{h}: G \rightarrow G$ is given by $g \mapsto h g h^{-1}$ for $h \in H$. The Lie differential of $C_{h}$ is the adjoint map $\operatorname{Ad}_{h}: \mathfrak{g} \rightarrow \mathfrak{g}$. The map $C_{h}$ fixes $H$, and $\operatorname{Ad}_{h}$ fixes $\mathfrak{h}$. Thus the map $\operatorname{Ad}_{h}: \mathfrak{g} \rightarrow \mathfrak{g}$ descends to a map $\widetilde{\operatorname{Ad}_{h}}: \mathfrak{g} / \mathfrak{h} \rightarrow \mathfrak{g} / \mathfrak{h}$. We are going to show that there is a natural isomorphism $T_{p_{0}} M \cong \mathfrak{g} / \mathfrak{h}$ such that for each $h \in H$ the following diagram commutes:\\
\includegraphics[scale=0.25, center]{2025_10_09_265d7e1a22c89f79c21eg-263}

One way to state the meaning of this result is to say that $h \mapsto \widetilde{\operatorname{Ad}_{h}}$ is a representation of $H$ on the vector space $\mathfrak{g} / \mathfrak{h}$, which is equivalent to the linear isotropy representation. The isomorphism $T_{p_{0}} M \cong \mathfrak{g} / \mathfrak{h}$ is given in the\\
following very natural way: Let $\xi \in \mathfrak{h}$ and consider $T_{e} \pi(\xi) \in T_{p_{0}} M$. We have

$$
T_{e} \pi(\xi)=\left.\frac{d}{d t}\right|_{t=0} \pi(\exp \xi t)=0
$$

since $\exp \xi t \in \mathfrak{h}$ for all $t$. Thus $\mathfrak{h} \subset \operatorname{Ker}\left(T_{e} \pi\right)$. On the other hand, $\operatorname{dim} \mathfrak{h}=\operatorname{dim} H=\operatorname{dim}\left(\operatorname{Ker}\left(T_{e} \pi\right)\right)$, so in fact $\mathfrak{h}=\operatorname{Ker}\left(T_{e} \pi\right)$ and we obtain an isomorphism $\mathfrak{g} / \mathfrak{h} \cong T_{p_{0}} M$ induced from $T_{e} \pi$. Let us see why the diagram (5.4) commutes. First, $L_{h}$ is well-defined as a map from $G / H$ to itself and the following diagram clearly commutes:\\
\includegraphics[scale=0.25, center]{2025_10_09_265d7e1a22c89f79c21eg-264(1)}

Using our equivariant diffeomorphism $\phi: G / H \rightarrow M$, we obtain an equivalent commutative diagram\\
\includegraphics[scale=0.25, center]{2025_10_09_265d7e1a22c89f79c21eg-264(3)}

Applying these maps to $\exp t \xi$ for $\xi \in \mathfrak{g}$, we have\\
\includegraphics[scale=0.25, center]{2025_10_09_265d7e1a22c89f79c21eg-264(2)}

Applying the tangent functor (looking at the differential), we get the commutative diagram\\
\includegraphics[scale=0.25, center]{2025_10_09_265d7e1a22c89f79c21eg-264}\\
and, taking quotients, this gives the desired commutative diagram (5.4).

\subsection*{Combining Representations}
We close this chapter with a bit about constructing new linear representations from old ones. Suppose that V is an $\mathbb{F}$-vector space and let $\mathcal{B}= \left(v_{1}, \ldots, v_{n}\right)$ be a basis for V . Then denoting the matrix representative of $\lambda_{g}$ with respect to $\mathcal{B}$ by $\left[\lambda_{g}\right]_{\mathcal{B}}$ we obtain a homomorphism $G \rightarrow \operatorname{GL}(n, \mathbb{F})$\\
given by $g \mapsto\left[\lambda_{g}\right]_{\mathcal{B}}$. In general, a Lie group homomorphism of a Lie group $G$ into $\operatorname{GL}(n, \mathbb{F})$ is called a matrix representation of $G$. We have already seen that any Lie subgroup of $\mathrm{GL}\left(\mathbb{R}^{n}\right)$ acts on $\mathbb{R}^{n}$ by matrix multiplication, and the corresponding homomorphism is the inclusion map $G \hookrightarrow \mathrm{GL}\left(\mathbb{R}^{n}\right)$. More generally, a Lie subgroup $G$ of GL(V) acts on V in the obvious way simply by employing the definition of GL(V) as a set of linear transformations of V . We call this the standard action of the linear Lie subgroup of GL(V) on V, and the corresponding homomorphism is just the inclusion map $G \hookrightarrow \mathrm{GL}(\mathrm{V})$. Choosing a basis, the subgroup corresponds to a matrix group, and the standard action becomes matrix multiplication on the left of $\mathbb{F}^{n}$, where the latter is viewed as a space of column vectors. This action of a matrix group on column vectors is also referred to as a standard action.

Given a representation $\lambda$ of $G$ in a vector space V , we have a dual representation $\lambda^{*}$ of $G$ in the dual space $\mathrm{V}^{*}$ by defining $\lambda^{*}(g):=\left(\lambda\left(g^{-1}\right)\right)^{*}$ : $\mathrm{V}^{*} \rightarrow \mathrm{V}^{*}$. Recall that if $L: \mathrm{V} \rightarrow \mathrm{V}$ is linear, then $L^{*}: \mathrm{V}^{*} \rightarrow \mathrm{V}^{*}$ is defined by $L^{*}(\alpha)(v)=\alpha(L v)$ for $\alpha \in \mathrm{V}^{*}$ and $v \in \mathrm{V}$. This dual representation is also sometimes called the contragredient representation (especially when $\mathbb{F}=\mathbb{R}$ ).

Now let $\lambda^{\mathrm{V}}$ and $\lambda^{\mathrm{W}}$ be representations of a lie group $G$ in $\mathbb{F}$-vector spaces V and W respectively. We can then form the direct sum representation $\lambda^{\mathrm{V}} \oplus \lambda^{\mathrm{W}}$ by $\left(\lambda^{\mathrm{V}} \oplus \lambda^{\mathrm{W}}\right)_{g}:=\lambda_{g}^{\mathrm{V}} \oplus \lambda_{g}^{\mathrm{W}}$ for $g \in G$, where we have $\left(\lambda_{g}^{\mathrm{V}} \oplus \lambda_{g}^{\mathrm{W}}\right)(v, w)=\left(\lambda_{g}^{\mathrm{V}} v, \lambda^{\mathrm{W}} w\right)$.

We will not pursue a serious study of Lie group representations but simply note that a major goal in the subject is the identification and classification of irreducible representations. A representation $\lambda: G \rightarrow \mathrm{GL}(\mathrm{V})$ is said to be irreducible if there is no nonzero proper subspace W of V such that $\lambda_{g}(\mathrm{~W}) \subset \mathrm{W}$ for all $g$. A large class of Lie groups known as semisimple Lie groups have the property that their representations break into direct sums of irreducible representations.

Example 5.149. A homogeneous polynomial of degree $d$ on $\mathbb{C}^{2}$ is a linear combination of monomials of total degree $d$. Let $H_{j}$ denote the vector space of homogeneous polynomials of degree $2 j$, where $j$ is a nonnegative "halfinteger" ( $j=k / 2$ for some nonnegative integer $k$ ). Define $\lambda_{j}: \mathrm{SU}(2) \rightarrow \mathrm{GL}\left(H_{j}\right)$ by $\left(\lambda_{j}(g) f\right)(z):=f\left(g^{-1} z\right)$ for $z=\left(z_{1}, z_{2}\right) \in \mathbb{C}^{2}$. Then $\lambda_{j}$ is an irreducible representation called the spin- $j$ representation of $\mathrm{SU}(2)$. The spin- $1 / 2$ representation turns out to be equivalent to the standard representation of $\operatorname{SU}(2)$ in $\mathbb{C}^{2}$. These spin representations play an important role in quantum physics.

One can also form the tensor product of representations. The definitions and basic facts about tensor products are given in the more general context of\\
module theory in Appendix D. Here we give a quick recounting of the notion of a tensor product of vector spaces, and then we define tensor products of representations. Given two vector spaces $V$ and $W$ over some field $\mathbb{F}$, consider the class $\mathcal{C}_{\mathrm{V} \times \mathrm{W}}$ consisting of all bilinear maps $\mathrm{V} \times \mathrm{W} \rightarrow \mathrm{X}$, where X varies over all $\mathbb{F}$-vector spaces, but V and W are fixed. We take members of $\mathcal{C}_{\mathrm{V} \times \mathrm{W}}$ as the objects of a category (see Appendix A). A morphism from, say, $\mu_{1}: \mathrm{V} \times \mathrm{W} \rightarrow \mathrm{X}$ to $\mu_{2}: \mathrm{V} \times \mathrm{W} \rightarrow \mathrm{Y}$ is defined to be a linear map $\ell: \mathrm{X} \rightarrow \mathrm{Y}$ such that the diagram\\
\includegraphics[scale=0.25, center]{2025_10_09_265d7e1a22c89f79c21eg-266(1)}\\
commutes.\\
There exists a vector space $\mathrm{T}_{\mathrm{V}, \mathrm{W}}$ together with a bilinear map $\otimes: \mathrm{V} \times \mathrm{W} \rightarrow \mathrm{T}_{\mathrm{V}, \mathrm{W}}$ that has the following universal property: For every bilinear map $\mu: \mathrm{V} \times \mathrm{W} \rightarrow \mathrm{X}$, there is a unique linear map $\widetilde{\mu}: \mathrm{T}_{\mathrm{V}, \mathrm{W}} \rightarrow \mathrm{X}$ such that the following diagram commutes:\\
\includegraphics[scale=0.25, center]{2025_10_09_265d7e1a22c89f79c21eg-266(2)}

If such a pair ( $T_{V, W}, \otimes$ ) exists with this property, then it is unique up to isomorphism in $\mathcal{C}_{\mathrm{V} \times \mathrm{W}}$. In other words, if $\widehat{\otimes}: \mathrm{V} \times \mathrm{W} \rightarrow \widehat{\mathrm{T}}_{\mathrm{V}, \mathrm{W}}$ also has this universal property, then there is a linear isomorphism $\mathrm{T}_{\mathrm{V}, \mathrm{W}} \cong \widehat{\mathrm{T}}_{\mathrm{V}, \mathrm{W}}$ such that the following diagram commutes:\\
\includegraphics[scale=0.25, center]{2025_10_09_265d7e1a22c89f79c21eg-266}

We refer to such universal object as a tensor product of V and W . We will indicate the construction of a specific tensor product that we denote by $\mathrm{V} \otimes \mathrm{W}$ with corresponding bilinear map $\otimes: \mathrm{V} \times \mathrm{W} \rightarrow \mathrm{V} \otimes \mathrm{W}$. The idea is\\
simple: We let $\mathrm{V} \otimes \mathrm{W}$ be the set of all linear combinations of symbols of the form $v \otimes w$ for $v \in \mathrm{V}$ and $w \in \mathrm{~W}$, subject to the relations

$$
\begin{aligned}
\left(v_{1}+v_{2}\right) \otimes w & =v_{1} \otimes w+v_{2} \otimes w, \\
v \otimes\left(w_{1}+w_{2}\right) & =v \otimes w_{1}+v \otimes w_{2}, \\
r(v \otimes w) & =r v \otimes w=v \otimes r w, \text { for } r \in \mathbb{F} .
\end{aligned}
$$

The map $\otimes$ is then simply $\otimes:(v, w) \rightarrow v \otimes w$. Somewhat more pedantically, let $F(\mathrm{V} \times \mathrm{W})$ denote the free vector space generated by the set $\mathrm{V} \times \mathrm{W}$ (the elements of $\mathrm{V} \times \mathrm{W}$ are treated as a basis for the space, and so the free space has dimension equal to the cardinality of the set $\mathrm{V} \times \mathrm{W}$ ). Next consider the subspace $R$ of $F(\mathrm{V} \times \mathrm{W})$ generated by the set of all elements of the form

$$
\begin{aligned}
& (a v, w)-a(v, w) \\
& (v, a w)-a(v, w) \\
& \left(v_{1}+v_{2}, w\right)-\left(v_{1}, w\right)+\left(v_{2}, w\right) \\
& \left(v, w_{1}+w_{2}\right)-\left(v, w_{1}\right)+\left(v, w_{2}\right)
\end{aligned}
$$

for $v_{1}, v_{2}, v \in \mathrm{V}, w_{1}, w_{2}, w \in \mathrm{~W}$, and $a \in \mathbb{F}$. Then we let $\mathrm{V} \otimes \mathrm{W}$ be defined as the quotient vector space $F(\mathrm{V} \times \mathrm{W}) / R$, and we have a corresponding quotient map $F(\mathrm{V} \times \mathrm{W}) \rightarrow \mathrm{V} \otimes \mathrm{W}$. The set $\mathrm{V} \times \mathrm{W}$ is contained in $F(\mathrm{V} \times \mathrm{W})$, and the map $\otimes: \mathrm{V} \times \mathrm{W} \rightarrow \mathrm{V} \otimes \mathrm{W}$ is then defined to be the restriction of the quotient map to $\mathrm{V} \times \mathrm{W}$. The image of $(v, w)$ under the quotient map is denoted by $v \otimes w$.

Tensor products of several vector spaces at a time are constructed similarly to be a universal space in a category of multilinear maps (Definition D.13). We may also form the tensor products two at a time and then use the easily proved fact that $(\mathrm{V} \otimes \mathrm{W}) \otimes \mathrm{U} \cong \mathrm{V} \otimes(\mathrm{W} \otimes \mathrm{U})$, which is then denoted by $\mathrm{V} \otimes \mathrm{W} \otimes \mathrm{U}$. Again the reader is referred to Appendix D for more about tensor products.

Elements of the form $v \otimes w$ generate $\mathrm{V} \otimes \mathrm{W}$, and in fact, if $\left(e_{1}, \ldots, e_{r}\right)$ is a basis for V and $\left(f_{1}, \ldots, f_{s}\right)$ is a basis for W , then

$$
\left\{e_{i} \otimes f_{j}: 1 \leq i \leq r, 1 \leq j \leq s\right\}
$$

is a basis for $\mathrm{V} \otimes \mathrm{W}$, which therefore has dimension $r s=\operatorname{dimV} \operatorname{dimW}$.\\
One more observation: If $A: \mathrm{V} \rightarrow \mathrm{X}$ and $B: \mathrm{W} \rightarrow \mathrm{Y}$ are linear maps, then we can define a linear map $A \otimes B: \mathrm{V} \otimes \mathrm{W} \rightarrow \mathrm{X} \otimes \mathrm{Y}$. First note that the map $(v, w) \mapsto A v \otimes B w$ is bilinear. Thus, by the universal property, there is a unique map $A \otimes B$ such that

$$
(A \otimes B)(v \otimes w)=A v \otimes B w
$$

for all $v \in \mathrm{V}, w \in \mathrm{~W}$.

Notice that if $A$ and $B$ are invertible, then $A \otimes B$ is invertible with $(A \otimes B)^{-1}(v \otimes w)=A^{-1} v \otimes B^{-1} w$. Pick bases for V and W as above and bases $\left\{e_{1}^{\prime}, \ldots, e_{r}^{\prime}\right\}$ and $\left\{f_{1}^{\prime}, \ldots, f_{s}^{\prime}\right\}$ for X and Y respectively. For $\tau \in \mathrm{V} \otimes \mathrm{W}$, we can write $\tau=\tau^{i j} e_{i} \otimes f_{j}$ using the Einstein summation convention. We have

$$
\begin{aligned}
A \otimes B(\tau) & =A \otimes B\left(\tau^{i j} e_{i} \otimes f_{j}\right) \\
& =\tau^{i j} A e_{i} \otimes B f_{j} \\
& =\tau^{i j} A_{i}^{k} e_{k}^{\prime} \otimes B_{j}^{l} f_{l}^{\prime} \\
& =\tau^{i j} A_{i}^{k} B_{j}^{l}\left(e_{k}^{\prime} \otimes f_{l}^{\prime}\right),
\end{aligned}
$$

so that the matrix of $A \otimes B$ is given by $(A \otimes B)_{i j}^{k l}=A_{i}^{k} B_{j}^{l}$.\\
Let $\lambda^{\mathrm{V}}$ and $\lambda^{\mathrm{W}}$ be representations of a Lie group $G$ in $\mathbb{F}$-vector spaces V and W , respectively. We can form a representation of $G$ in the tensor product space $\mathrm{V} \otimes \mathrm{W}$ by letting $\left(\lambda^{\mathrm{V}} \otimes \lambda^{\mathrm{W}}\right)_{g}:=\lambda_{g}^{\mathrm{V}} \otimes \lambda_{g}^{\mathrm{W}}$ for all $g \in G$. There is a variation on the tensor product that is useful when we have two groups involved. If $\lambda^{\mathrm{V}}$ is a representation of a Lie group $G_{1}$ in the $\mathbb{F}$-vector space V and $\lambda^{\mathrm{W}}$ is a representation of a Lie group $G_{2}$ in the $\mathbb{F}$-vector space W , then we can form a representation of the Lie group $G_{1} \times G_{2}$, also called the tensor product representation and denoted $\lambda^{\mathrm{V}} \otimes \lambda^{\mathrm{W}}$ as before. In this case, the definition is $\left(\lambda^{\mathrm{V}} \otimes \lambda^{\mathrm{W}}\right)_{\left(g_{1}, g_{2}\right)}:=\lambda_{g_{1}}^{\mathrm{V}} \otimes \lambda_{g_{2}}^{\mathrm{W}}$. Of course if it happens that $G_{1}=G_{2}$, then we have an ambiguity since $\lambda^{\mathrm{V}} \otimes \lambda^{\mathrm{W}}$ could be a representation of $G$ or of $G \times G$. One can usually determine which version is meant from the context. Alternatively, one can use pairs to denote actions so that an action $\lambda: G \times \mathrm{V} \rightarrow \mathrm{V}$ is denoted $(G, \lambda)$. Then the two tensor product representations would be ( $G \times G, \lambda^{\mathrm{V}} \otimes \lambda^{\mathrm{W}}$ ) and ( $G, \lambda^{\mathrm{V}} \otimes \lambda^{\mathrm{W}}$ ), respectively.

\section*{Problems}
(1) Verify that each of the groups described in Section 5.2 is (isomorphic to) a Lie subgroup of an appropriate Lie group of linear automorphisms.\\
(2) Show that proper continuous maps are closed maps.\\
(3) Show that $\mathrm{SL}(2, \mathbb{C})$ is simply connected and that $\wp: \mathrm{SL}(2, \mathbb{C}) \rightarrow$ Mob is a universal covering homomorphism. See Example 5.50.\\
(4) Prove Proposition 5.83.\\
(5) Let $\mathfrak{g}$ be a Lie algebra of a Lie group $G$. Show that the set of all automorphisms of $\mathfrak{g}$, denoted $\operatorname{Aut}(\mathfrak{g})$, forms a Lie group (actually a Lie subgroup of GL( $\mathfrak{g}$ )).\\
(6) Show that if we consider $\operatorname{SL}(2, \mathbb{R})$ as a subset of $\operatorname{SL}(2, \mathbb{C})$ in the obvious way, then $\mathrm{SL}(2, \mathbb{R})$ is a Lie subgroup of $\mathrm{SL}(2, \mathbb{C})$ and $\wp(\mathrm{SL}(2, \mathbb{R}))$ is a Lie subgroup of Mob. Show that if $T \in \wp(\mathrm{SL}(2, \mathbb{R}))$, then $T$ maps the upper half-plane of $\mathbb{C}$ onto itself (bijectively).\\
(7) Show that for $v \in T_{e} G$, the field defined by $g \mapsto L^{v}(g):=T L_{g} \cdot v$ is automatically smooth.\\
(8) Determine explicitly the map $T_{I}$ inv $: T_{I} \mathrm{GL}(n, \mathbb{R}) \rightarrow T_{I} \mathrm{GL}(n, \mathbb{R})$, where $\operatorname{inv}: \operatorname{GL}(n, \mathbb{R}) \rightarrow \operatorname{GL}(n, \mathbb{R})$ is defined by $\operatorname{inv}(A)=A^{-1}$.\\
(9) Let $H$ be the set of real $3 \times 3$ matrices of the form

$$
A=\left[\begin{array}{lll}
1 & a & b \\
0 & 1 & c \\
0 & 0 & 1
\end{array}\right]
$$

Find a global chart for $H$ and show that this and the usual matrix multiplication gives $H$ the structure of a Lie group.\\
(10) If $G$ is a connected Lie group and $h: G \rightarrow H$ is a Lie group homomorphism with discrete kernel $K$, then $K \subset Z(G)$, where $Z(G)=\{x \in G$ : $x g=g x$ for all $g \in G\}$ is the center of $G$.\\
(11) Show that for a Lie group $G$, the conjugation map $C_{g}: G \rightarrow G$ defined by $x \mapsto g x g^{-1}$ is a Lie group isomorphism. Show that the map $C: g \rightarrow \operatorname{Diff}(G)$ is a group homomorphism. Note that we have not defined any Lie group structure on $\operatorname{Diff}(G)$.\\
(12) Consider the map $T_{e} C_{g}: T_{e} G \rightarrow T_{e} G$. Show that $g \mapsto T_{e} C_{g}$ is a Lie group homomorphism from $G$ into $\mathrm{GL}\left(T_{e} G\right)$.\\
(13) Show that $\mathrm{SO}(3)$ is the connected component of the identity in $\mathrm{O}(3)$. Show that the special Lorentz group $\mathrm{SO}(1,3)$ is not connected. Show that the first entry of elements of $\operatorname{SO}(1,3)$ must have absolute value greater than or equal to 1 . Define $\operatorname{SO}(3,1)^{\uparrow}$ as the subset of $\operatorname{SO}(1,3)$ consisting of matrices with positive first entry (which must be greater than 1). Show that $\operatorname{SO}(3,1)^{\uparrow}$ is connected (and hence the connected component of the identity in $\mathrm{O}(1,3)$ ).\\
(14) Let $A \in \mathfrak{g} \mathfrak{l}(\mathrm{V})=L(\mathrm{V}, \mathrm{V})$ for some finite-dimensional vector space V . Show that if $A$ has eigenvalues $\left\{\lambda_{i}\right\}_{i=1, \ldots, n}$, then $\operatorname{ad}(A)$ has eigenvalues $\left\{\lambda_{j}-\lambda_{k}\right\}_{j, k=1, \ldots, n}$. Hint: Choose a basis for V such that $A$ is represented by an upper triangular matrix. Show that this induces a basis for $\mathfrak{g l}(\mathrm{V})$ such that with the appropriate ordering, $\operatorname{ad}(A)$ is upper triangular.\\
(15) Fix a nonzero vector $w \in \mathbb{R}^{3}$ with length $\theta=\|w\|$. Let $L_{w}: \mathbb{R}^{3} \rightarrow \mathbb{R}^{3}$ denote the linear transformation $v \mapsto w \times v$, where $\times$ is the vector cross product. Show that for any right handed orthonormal basis $\left\{e_{1}, e_{2}, e_{3}\right\}$\\
with $e_{3}$ parallel to $w$ we have

$$
\begin{aligned}
& \exp \left(L_{w}\right) e_{1}=\cos \theta e_{1}+\sin \theta e_{2} \\
& \exp \left(L_{w}\right) e_{2}=-\sin \theta e_{1}+\cos \theta e_{2} \\
& \exp \left(L_{w}\right) e_{3}=e_{3}
\end{aligned}
$$

(16) Let $L_{w}$ be as in the previous problem. Show that

$$
\exp L_{w}=I+\frac{\sin \theta}{\theta} L_{w}+\frac{1-\cos \theta}{\theta^{2}} L_{w}^{2},
$$

where $\frac{\sin \theta}{\theta}$ and $\frac{1-\cos \theta}{\theta^{2}}$ are defined in the obvious way using power series.\\
(17) Let $A, B \in \mathfrak{g l}(\mathrm{V})$, where V is a finite-dimensional vector space over the field $\mathbb{F}=\mathbb{R}$ or $\mathbb{C}$, and show that the following statements are equivalent:\\
(a) $[A, B]=0$.\\
(b) $\exp s A$ and $\exp t B$ commute for all $s, t \in \mathbb{F}$.\\
(c) $\exp (s A+t B)=\exp (s A) \exp (t B)$ for all $s, t \in \mathbb{F}$.\\
(18) Let V be a finite-dimensional normed space over the field $\mathbb{R}$ (resp. $\mathbb{C}$ ). Show that if $\sum_{n=0}^{\infty} a_{n} x^{n}$ is an absolutely convergent real (resp. complex) power series with radius of convergence $R$, then

$$
\sum_{n=0}^{\infty} a_{n} A^{n}
$$

converges (absolutely) in the normed space $\mathfrak{g l}(\mathrm{V})$ for $\|A\|<R$.\\
(19) Show that if $G$ is a connected Lie group, then $\pi_{1}(G)$ is abelian.\\
(20) Let $G$ be a Lie group and denote by $\mu: G \times G \rightarrow G$ the multiplication map.\\
(a) Identify $T(G \times G)$ with $T G \times T G$ in the usual way. Show that the tangent map $T \mu: T G \times T G \rightarrow T G$ defines a Lie group structure on $T G$ and show that if $\left(v_{g}, w_{h}\right) \in T_{g} G \times T_{h} G$, then

$$
T \mu\left(v_{g}, w_{h}\right)=T R_{h} v_{g}+T L_{g} w_{h},
$$

where $R_{h}$ and $L_{g}$ are right and left multiplications respectively.\\
(b) Show that under the isomorphism $G \times \mathfrak{g}$ with $T G$, the Lie group multiplication takes the form

$$
(g, A) \cdot(h, B)=\left(g h, \operatorname{Ad}_{h^{-1}} A+B\right)
$$

(21) Let $M$ be a smooth manifold, let $G$ be a Lie group, and let $r: M \times G \rightarrow M$ be a right Lie group action. Recall from the previous problem that $T G$ is naturally a Lie group.\\
(a) Show that $\operatorname{Tr}: T M \times T G \rightarrow T M$ defines a right Lie group action.\\
(b) For each $A \in \mathfrak{g}=T_{e} G$, define a vector field $\sigma(A)$ on $M$ by $\sigma(A)(p):=\left.\frac{d}{d t}\right|_{t=0} p \exp (A t)$. Show that for $A, B \in \mathfrak{g}$, we have

$$
\sigma([A, B])=[\sigma(A), \sigma(B)]
$$

(c) Show that if $A^{L}$ denote the left invariant vector field generated by $A$, then for $v_{p} \in T_{p} M$,

$$
\operatorname{Tr}\left(v_{p}, A^{L}(g)\right)=T R_{g}\left(v_{p}\right)+\sigma(A)(p g)
$$






\pagebreak


\section*{Chapter 6}



\section*{Fiber Bundles}


The notion of a bundle is basic in both topology and geometry. The reader need not master everything in this chapter before going on to later chapters and should skip forward rather than become too bogged down. The definition and basic examples of a vector bundle are most important. In this chapter we also introduce the more advanced notion of a structure group for a bundle. There is more than one approach to structure groups. We start out with an approach that takes the notion of a $G$-atlas as basic. This is essentially the approach of Steenrod $[\mathbf{S t}]$. One may also approach $G$ bundle structures by first introducing the notion of a principal $G$-bundle (see [Hus]). We discuss principal bundles near the end of this chapter.

\subsection*{General Fiber Bundles}
Definition 6.1. 

\begin{DEF}{BÜN-L09-02-01}{Faserbündel}
Let $F, M$, and $E$ be $C^{r}$ manifolds and let $\pi: E \rightarrow M$ be a $C^{r}$ map. The quadruple ( $E, \pi, M, F$ ) is called a (locally trivial) $C^{r}$ fiber bundle if for each point $p \in M$ there is an open set $U$ containing $p$ and a $C^{r}$ diffeomorphism $\phi: \pi^{-1}(U) \rightarrow U \times F$ such that the following diagram commutes:\\
\includegraphics[scale=0.25, center]{2025_10_09_265d7e1a22c89f79c21eg-272}
\end{DEF}

In differential geometry, attention is usually focused on $C^{\infty}$ fiber bundles (smooth fiber bundles), but the continuous case is also of interest. We will restrict ourselves to the smooth case, but the reader should keep in mind that most of the definitions and theorems have analogous $C^{0}$ versions where
\begin{figure}[h]
\begin{center}
  \includegraphics[scale=0.25]{2025_10_09_265d7e1a22c89f79c21eg-273}

\caption{Figure 6.1. Schematic for fiber bundle}
\end{center}
\end{figure}
the spaces are assumed merely to be sufficiently nice topological spaces and the maps are only assumed to be continuous. In this chapter, all maps and spaces will be smooth unless otherwise indicated.

Definition 6.2. 

\begin{DEF}{BÜN-L09-02-02}{Totaler Raum, Bündel-Projektion, Basisraum, Typische Faser, Faser über $p$}
If ( $E, \pi, M, F$ ) is a smooth fiber bundle, then $E$ is called the total space, $\pi$ is called the bundle projection, $M$ is called the base space and $F$ is called the typical fiber. For each $p \in M$, the set $E_{p}:=\pi^{-1}(p)$ is called the fiber over $p$.
\end{DEF}

Because the quadruple notation is cumbersome, it is common to denote a fiber bundle by a single symbol. For example, we could write $\xi=$ ( $E, \pi, M, F$ ). In the literature, it is common to see $E$ refer both to the total space and to the fiber bundle itself (an abuse of notation). The map $\pi$ is also a common way to reference the fiber bundle.

Example 6.3. 


\begin{EXA}{BÜN-L09-02-03}{Kreuzprodukt Mannigfaltigkeit als Faserbündel}
For smooth manifolds $M$ and $F$, we have the projection $\mathrm{pr}_{1}: M \times F \rightarrow M$. Then, ( $M \times F, \mathrm{pr}_{1}, M, F$ ) is a fiber bundle called a product bundle (or trivial bundle).
\end{EXA}

Exercise 6.4. 

\begin{EXE}{BÜN-L09-02-04}{Projektionen sind Submersionen und Fasern sind reguläre Untermannigfaltigkeiten und ZSH Charakterisierung}
Show that if $\xi=(E, \pi, M, F)$ is a (smooth) fiber bundle, then $\pi: E \rightarrow M$ is a submersion and each fiber $\pi^{-1}(p)$ is a regular submanifold which is diffeomorphic to $F$. Show that if both $F$ and $M$ are connected, then $E$ is connected.
\end{EXE}

There are various categories of bundles with corresponding notions of morphism. We give two very general definitions and modify them as needed.

Definition 6.5 (Bundle morphism (type I)). 

\begin{DEF}{BÜN-L09-02-05}{Kategorie der glatten Faserbündel über gleichem Basisraum}
Let $\xi_{1}=\left(E_{1}, \pi_{1}, M, F_{1}\right)$ and $\xi_{2}=\left(E_{2}, \pi_{2}, M, F_{2}\right)$ be smooth fiber bundles with the same base space $M$. A (type I) bundle morphism over $M$ from $\xi_{1}$ to $\xi_{2}$ is a smooth map $h: E_{1} \rightarrow E_{2}$ such that the following diagram commutes:\\
\includegraphics[scale=0.25, center]{2025_10_09_265d7e1a22c89f79c21eg-274(1)}
This type of morphism is also called an $M$-morphism or a morphism over $M$. If $h$ is also a diffeomorphism, then $h$ is called a bundle isomorphism over $M$ and in this case the bundles are said to be isomorphic (over $M$ ) or equivalent. A bundle isomorphism from a bundle to itself is called a bundle automorphism.
\end{DEF}


Definition 6.6 (Bundle morphism (type II)). 


\begin{DEF}{BÜN-L09-02-06}{Kategorie der glatten Faserbündel über verschiedene Basisräume}
Let $\xi_{1}=\left(E_{1}, \pi_{1}, M_{1}, F_{1}\right)$ and $\xi_{2}=\left(E_{2}, \pi_{2}, M_{2}, F_{2}\right)$ be smooth fiber bundles. A (type II) bundle morphism from $\xi_{1}$ to $\xi_{2}$ is a pair of smooth maps $\widehat{f}: E_{1} \rightarrow E_{2}$ and $f: M_{1} \rightarrow M_{2}$ such that the following diagram commutes:\\
\includegraphics[scale=0.25, center]{2025_10_09_265d7e1a22c89f79c21eg-274}
We write $(\widehat{f}, f): \xi_{1} \rightarrow \xi_{2}$ and say that $\widehat{f}$ is a bundle morphism along $f$. If both $\widehat{f}$ and $f$ are diffeomorphisms, then we call ( $\widehat{f}, f$ ) a bundle isomorphism. In this case, we say that the bundles are isomorphic over $f$.

Note that as a fiber preserving map, $\widehat{f}$ determines $f$ and so it is also proper to refer to $\widehat{f}$ as the bundle morphism and we sometimes say that $\widehat{f}$ is a bundle morphism along (or over) $f$.
\end{DEF}

Warning: The definitions of bundle morphism above are quite relaxed. There are a variety of definitions in the literature that require more than the definitions above, especially when structure groups (discussed below) are emphasized.

Definition 6.7. 

\begin{DEF}{BÜN-L09-02-07}{(Globale) glatte Schnitte von Faserbündeln}
A (global) smooth section of a fiber bundle $\xi=(E, \pi$, $M, F)$ is a smooth map $\sigma: M \rightarrow E$ such that $\pi \circ \sigma=\operatorname{id}_{M}$ (i.e., $\sigma(p) \in E_{p}$ ). A local smooth section over an open set $U$ is a smooth map $\sigma: U \rightarrow E$ such that $\pi \circ \sigma=\mathrm{id}_{U}$. The set of smooth sections of $\xi$ is denoted by $\Gamma(\xi)$ or sometimes by $\Gamma(E)$ or $\Gamma(\pi)$.
\end{DEF}

\begin{CONC}{BÜN-L09-02-08}{Natürliche Bijektion zwischen Raum der Schnitte als notwendiges Bedingung für Bündelisomorphie}
A very important point is that a fiber bundle may not have any global smooth sections. If two bundles are equivalent, via a bundle isomorphism $h$ (of type I), then there is a natural bijection between the spaces of sections given by $\sigma \mapsto h \circ \sigma$. This means that one quick way to conclude that two bundles are not equivalent is by showing that one bundle has global sections while the other does not.
\end{CONC}




\begin{DEF}{BÜN-L09-02-09}{Lokale Trivialisierung und Bündelkarten}
Sei $\pi:E\to M$ ein glattes Faserbündel mit typischer Faser $F$.
\begin{enumerate}
\item Eine \emph{lokale Trivialisierung} des Bündels über einer offenen Menge $U\subset M$
ist eine glatte Abbildung
\[
\phi:\pi^{-1}(U)\longrightarrow U\times F,
\]
die ein Diffeomorphismus ist und für die gilt
\[
\mathrm{pr}_1\circ\phi=\pi|_{\pi^{-1}(U)}.
\]
\item Jede solche Abbildung ist notwendigerweise von der Form
\[
\phi=(\pi,\Phi),\qquad
\Phi:\pi^{-1}(U)\to F,
\]
wobei $\Phi|_{E_p}:E_p\to F$ für jedes $p\in U$ ein Diffeomorphismus ist.
Die zweite Komponente $\Phi$ heißt die \emph{Hauptabbildung} oder der
\emph{principal part} der Trivialisierung.
\item Ein Paar $(U,\phi)$ mit einer lokalen Trivialisierung $\phi$ heißt eine
\emph{Bündelkarte}.
\end{enumerate}
\end{DEF}



\begin{DEF}{BÜN-L09-02-10}{Typ-I- und Typ-II-Bündelkarten}
\begin{enumerate}
\item Eine \emph{Typ~I-Bündelkarte} ist eine lokale Trivialisierung $(U,\phi)$ wie oben.
\item Eine \emph{Typ~II-Bündelkarte} ist ein Paar $(\phi,\mathrm{x})$, wobei
\[
\phi:\pi^{-1}(U)\longrightarrow V\times F
\quad\text{und}\quad
\mathrm{x}:U\longrightarrow V
\]
glatte Diffeomorphismen sind und das Diagramm
\[
\begin{array}{ccc}
\pi^{-1}(U) & \xrightarrow{\phi} & V\times F\\[3pt]
\pi\!\!\downarrow & & \downarrow\!\!\mathrm{pr}_1\\[3pt]
U & \xrightarrow{\mathrm{x}} & V
\end{array}
\]
kommutiert.
In der Regel ist $(U,\mathrm{x})$ eine Karte der Basis­mannigfaltigkeit.
\item Die beiden Typen sind äquivalent, da sich jede Typ-II-Karte durch
Komposition mit $(\mathrm{x}^{-1},\mathrm{id}_F)$ in eine Typ-I-Karte überführen lässt.
\end{enumerate}
\end{DEF}



\begin{DEF}{BÜN-L09-02-11}{Bündelatlas}
Eine Familie von Bündelkarten
\[
\{(U_\alpha,\phi_\alpha)\}_{\alpha\in A},
\qquad
\phi_\alpha:\pi^{-1}(U_\alpha)\to U_\alpha\times F,
\]
heißt ein \emph{Bündelatlas}, falls die offenen Mengen $\{U_\alpha\}_{\alpha\in A}$ den Basisraum $M$ überdecken.
\end{DEF}



\begin{DEF}{BÜN-L09-02-12}{Übergangsfunktionen}
Für zwei Bündelkarten $(U_\alpha,\phi_\alpha)$ und $(U_\beta,\phi_\beta)$
mit nichtleerem Schnitt $U_\alpha\cap U_\beta$
sei $\phi_\alpha=(\pi,\Phi_\alpha)$ und $\phi_\beta=(\pi,\Phi_\beta)$.
Dann ist
\[
\phi_\alpha\circ\phi_\beta^{-1}:
(U_\alpha\cap U_\beta)\times F
\longrightarrow
(U_\alpha\cap U_\beta)\times F,
\qquad
\phi_\alpha\circ\phi_\beta^{-1}(p,y)
=
(p,\Phi_{\alpha\beta}(p)(y)),
\]
wobei
\[
\Phi_{\alpha\beta}(p)
=
\Phi_\alpha|_{E_p}\circ
\Phi_\beta|_{E_p}^{-1}
\in\mathrm{Diff}(F).
\]
Die Abbildungen
$\Phi_{\alpha\beta}:U_\alpha\cap U_\beta\to\mathrm{Diff}(F)$
heißen \emph{Übergangsfunktionen} (oder \emph{transition maps}) des Bündels.
Sie erfüllen die folgenden Kokzykelbedingungen:
\[
\begin{array}{ll}
\Phi_{\alpha\alpha}(p)=e, & p\in U_\alpha,\\[2pt]
\Phi_{\alpha\beta}(p)=(\Phi_{\beta\alpha}(p))^{-1}, & p\in U_\alpha\cap U_\beta,\\[2pt]
\Phi_{\alpha\beta}(p)\circ\Phi_{\beta\gamma}(p)\circ\Phi_{\gamma\alpha}(p)=\mathrm{id}, & p\in U_\alpha\cap U_\beta\cap U_\gamma.
\end{array}
\]
\end{DEF}













\begin{CONC}{BÜN-L09-02-13}{Zusammenhang zwischen Bündelkarten und Übergangsabbildungen}
A bundle chart essentially gives a local type I bundle isomorphism, but it is sometimes more natural to consider bundle charts which are local type II isomorphisms. We will call these type II bundle charts. These are of the form $(\phi, \mathrm{x})$, where $\phi: \pi^{-1} U \rightarrow V \times F$ and $\mathrm{x}: U \rightarrow V$ are smooth diffeomorphisms such that the following diagram commutes:\\
\includegraphics[scale=0.25, center]{2025_10_09_265d7e1a22c89f79c21eg-275}

Usually, the pair ( $U, \mathrm{x}$ ) is a chart on the base manifold. The two types of bundle charts are equivalent since one may always compose a type II chart with ( $\mathrm{x}^{-1}, \mathrm{id}_{F}$ ) to obtain a type I bundle chart.

More restricted notions of bundle morphism can be obtained by making requirements such as that the induced maps on fibers $\left.\widehat{f}\right|_{\pi_{1}^{-1}(p)}: \pi_{1}^{-1}(p) \rightarrow \pi_{2}^{-1}(p)$ are $C^{\infty}$ diffeomorphisms.

The maps $\phi: \pi^{-1}(U) \rightarrow U \times F$ occurring in the definition of a fiber bundle are said to be local trivializations of the bundle. It is easy to see that such a local trivialization must be a map of the form $\phi=\left(\left.\pi\right|_{\pi^{-1}(U)}, \Phi\right)$ where $\Phi: \pi^{-1}(U) \rightarrow F$ is a smooth map with the property that $\left.\Phi\right|_{E_{p}}$ : $E_{p} \rightarrow F$ is a diffeomorphism. We will just write ( $\pi, \Phi$ ) instead of the more pedantic ( $\left.\pi\right|_{\pi^{-1}(U)}, \Phi$ ). Thus

$$
\phi(y)=(\pi, \Phi)(y):=(\pi(y), \Phi(y)) .
$$

The second component map $\Phi$ is called the principal part of the local trivialization. A pair ( $U, \phi$ ), where $\phi$ is a local trivialization over $U \subset M$, is called a bundle chart. (Clearly, a local trivialization and a bundle chart are essentially the same thing.) A family $\left\{\left(U_{\alpha}, \phi_{\alpha}\right)\right\}_{\alpha \in A}$ of bundle charts such that $\left\{U_{\alpha}\right\}_{\alpha \in A}$ is a cover of $M$ is said to be a bundle atlas. Given two such bundle charts ( $U_{\alpha}, \phi_{\alpha}$ ) and ( $U_{\beta}, \phi_{\beta}$ ), we have
$$
\phi_{\alpha}=\left(\pi, \Phi_{\alpha}\right): \pi^{-1}\left(U_{\alpha}\right) \rightarrow U_{\alpha} \times F
$$
and similarly for $\phi_{\beta}=\left(\pi, \Phi_{\beta}\right)$. If $U_{\alpha} \cap U_{\beta}$ is not empty, then $\pi^{-1}\left(U_{\alpha}\right) \cap \pi^{-1}\left(U_{\beta}\right)=\pi^{-1}\left(U_{\alpha} \cap U_{\beta}\right)$ is not empty and we have overlap maps
$$
\phi_{\alpha} \circ \phi_{\beta}^{-1}:\left(U_{\alpha} \cap U_{\beta}\right) \times F \rightarrow\left(U_{\alpha} \cap U_{\beta}\right) \times F .
$$
Since $\left.\Phi_{\alpha}\right|_{E_{p}}$ is a diffeomorphism for each $p \in U_{\alpha}$, the map $\left.\left.\Phi_{\alpha}\right|_{E_{p}} \circ \Phi_{\beta}\right|_{E_{p}} ^{-1}$ : $F \rightarrow F$ is a diffeomorphism for all $p \in U_{\alpha} \cap U_{\beta}$. We then obtain a map $\Phi_{\alpha \beta}: U_{\alpha} \cap U_{\beta} \rightarrow \operatorname{Diff}(F)$ defined by
$$
p \mapsto \Phi_{\alpha \beta}(p)=\left.\left.\Phi_{\alpha}\right|_{E_{p}} \circ \Phi_{\beta}\right|_{E_{p}} ^{-1} .
$$
It follows that
$$
\phi_{\alpha} \circ \phi_{\beta}^{-1}(p, y)=\left(p, \Phi_{\alpha \beta}(p)(y)\right) .
$$
The functions $\Phi_{\alpha \beta}: U_{\alpha} \cap U_{\beta} \rightarrow \operatorname{Diff}(F)$ are called transition maps or transition functions. Given a bundle atlas, the corresponding transition functions clearly satisfy the following "cocycle conditions":
$$
\begin{array}{ll}
\Phi_{\alpha \alpha}(p)=e & \text { for } p \in U_{\alpha} \\
\Phi_{\alpha \beta}(p)=\left(\Phi_{\beta \alpha}(p)\right)^{-1} & \text { for } p \in U_{\alpha} \cap U_{\beta} \\
\Phi_{\alpha \beta}(p) \circ \Phi_{\beta \gamma}(p) \circ \Phi_{\gamma \alpha}(p)=\text { id } & \text { for } p \in U_{\alpha} \cap U_{\beta} \cap U_{\gamma}
\end{array}
$$
for all $\alpha, \beta, \gamma$.


Notation 6.8. We will often denote $\Phi_{\alpha \beta}(p)(y)$ by $\left.\Phi_{\alpha \beta}\right|_{p}(y)$, which is, in many contexts, more transparent.


$\operatorname{Diff}(F)$ is a group, and we have a group action $\operatorname{Diff}(F) \times F \rightarrow F$ given by $(\psi, y) \mapsto \psi(y)$. However, $\operatorname{Diff}(F)$ is too big for many purposes, and we have certainly not attempted to give Diff( $F$ ) a Lie group structure. Even if we were to somehow extend the notion of Lie group sufficiently to include $\operatorname{Diff}(F)$, it would be infinite-dimensional and thereby take us out of the circle of ideas we have been developing. Because of this, the transition functions $\Phi_{\alpha \beta}$ above, which could be called "raw transition functions", might not be appropriate for our needs. We remedy this below by bringing Lie groups into the picture. First we give a simple example of a nontrivial bundle.
\end{CONC}



Example 6.9. 


\begin{EXA}{BÜN-L09-02-14}{Möbius Band Bündel}
The circle $S^{1}$ can be considered as a quotient $\mathbb{R} / \sim$ where $x$ is equivalent to $y$ if and only if $x-y$ is an integer multiple of $2 \pi$. For this example, we put an equivalence relation on $\mathbb{R} \times(-1,1)$ according to the prescription $(x, t) \sim\left(x+2 \pi n,(-1)^{n} t\right)$ for any integer $n$. The quotient $(\mathbb{R} \times(-1,1)) / \sim$ can easily be seen to be a smooth manifold and is none other than the familiar Möbius band which we denote by MB. Define a map $\pi: \mathrm{MB} \rightarrow \mathbb{R} / \sim=S^{1}$ by $\pi([x, t])=[x]$. We show that this is a fiber bundle by exhibiting an atlas consisting of three bundle charts. We call it the Möbius band bundle. We use three bundle charts instead of two in order that the overlaps be connected sets. Let $U_{1}=\{[x] \in \mathbb{R} / \sim:-2 \pi / 3< x<2 \pi / 3\}$ and $U_{2}=\{[x] \in \mathbb{R} / \sim: 0<x<4 \pi / 3\}$ and $U_{3}=\{[x] \in \mathbb{R} / \sim$ : $2 \pi / 3<x<2 \pi\}$. Then $U_{1} \cup U_{2} \cup U_{3}=\mathbb{R} / \sim=S^{1}$. For $i=2,3$ define $\phi_{i}: \pi^{-1}\left(U_{i}\right) \rightarrow U_{i} \times S^{1}$ by
$$
\phi_{i}([x, t])=([x], t),
$$
where ( $x, t$ ) is the unique representative of $[x, t]$ in the set $(0,2 \pi) \times(-1,1)$. For $\phi_{1}: \pi^{-1}\left(U_{1}\right) \rightarrow U_{1} \times S^{1}$, we define $\phi_{1}([x, t])=([x], t)$, where $(x, t)$ is the unique representative of $[x, t]$ in the set $(-2 \pi / 3,2 \pi / 3) \times(-1,1)$. One can check that $\phi_{2} \circ \phi_{3}^{-1}=\phi_{3} \circ \phi_{2}^{-1}=\mathrm{id}$ on the overlap $\pi^{-1}\left(U_{2}\right) \cap \pi^{-1}\left(U_{3}\right)$. Now consider the overlap $\pi^{-1}\left(U_{1}\right) \cap \pi^{-1}\left(U_{2}\right)$. If $[x, t] \in \pi^{-1}\left(U_{1}\right) \cap \pi^{-1}\left(U_{2}\right)$, then

\includegraphics[scale=0.25, center]{2025_10_09_265d7e1a22c89f79c21eg-277}

$[x, t]$ is uniquely represented by some $(x, t) \in(0,2 \pi / 3) \times(-1,1)$ and in view of the definitions we see that $\phi_{2} \circ \phi_{3}^{-1}=\mathrm{id}$ also. Finally we consider $\phi_{1} \circ \phi_{3}^{-1}$. If $[x, t] \in \pi^{-1}\left(U_{1}\right) \cap \pi^{-1}\left(U_{3}\right)$, then it has a unique representative ( $x, t$ ) in $(4 \pi / 3,2 \pi) \times(-1,1)$ and then $\phi_{3}^{-1}([x], t)=[x, t]$. For $\phi_{1}$, we need to represent $[x, t]$ properly. We use the fact that $[x, t]=[x-2 \pi,-t]$ and $(x-2 \pi,-t) \in (-2 \pi / 3,0) \times(-1,1)$ so that $\phi_{1}([x-2 \pi,-t])=([x-2 \pi],-t)=([x],-t)$. In short, we have $\phi_{1} \circ \phi_{3}^{-1}([x, t])=([x],-t)$. From these considerations and the fact that in general $\phi_{\alpha} \circ \phi_{\beta}^{-1}(p, y)=\left(p, \Phi_{\alpha \beta}(p)(y)\right)$ we see that
$$
\begin{aligned}
& \Phi_{12}(p)=\operatorname{id}_{(-1,1)} \in \operatorname{Diff}(-1,1) \text { for } p \in U_{1} \cap U_{2}, \\
& \Phi_{23}(p)=\operatorname{id}_{(-1,1)} \in \operatorname{Diff}(-1,1) \text { for } p \in U_{2} \cap U_{3}, \\
& \Phi_{13}(p)=-\operatorname{id}_{(-1,1)} \in \operatorname{Diff}(-1,1) \text { for } p \in U_{1} \cap U_{3} .
\end{aligned}
$$
A "twist" occurs on the overlap $\pi^{-1}\left(U_{1}\right) \cap \pi^{-1}\left(U_{3}\right)$. There is no way to construct an atlas for this bundle without having such a twist on at least one of the overlaps.

Notice that if we define an action $\lambda$ of $\mathbb{Z}_{2}=\{1,-1\}$ on the interval $(-1,1)$ by $\lambda(g, x) \mapsto g x$, then we can describe the transition functions in the last example by
$$
\Phi_{\alpha \beta}(p)(x)=\lambda\left(g_{\alpha \beta}(p), x\right),
$$
where $g_{\alpha \beta}: U_{\alpha} \cap U_{\beta} \rightarrow \mathbb{Z}_{2}$ is given by
$$
\begin{aligned}
& g_{12}=1 \text { on } U_{1} \cap U_{2}, \\
& g_{23}=1 \text { on } U_{2} \cap U_{3}, \\
& g_{13}=-1 \text { on } U_{1} \cap U_{3},
\end{aligned}
$$
and in this case the $g_{\alpha \beta}$ satisfy a cocycle condition like the $\Phi_{\alpha \beta}$. This is convenient since we understand $\mathbb{Z}_{2}$ very well. It is a zero-dimensional Lie group. Inspired by this, we seek to put Lie groups into the formalism. This will alleviate our concerns about the group $\operatorname{Diff}(F)$ mentioned above. We are also led to the theory of $G$-bundles that involves the group in subtle ways.
\end{EXA}





Definition 6.10. 

\begin{DEF}{BÜN-L09-02-15}{$G$-Kozyklus auf offener Überdeckung einer glatten Mannigfaltigkeit mit Lie Gruppe}
Let $\left\{U_{\alpha}\right\}$ be an indexed open cover of a smooth manifold $M$ and let $G$ be a Lie group. A $G$-cocycle on $\left\{U_{\alpha}\right\}$ is the assignment of a smooth map $g_{\alpha \beta}: U_{\alpha} \cap U_{\beta} \rightarrow G$ to every nonempty intersection $U_{\alpha} \cap U_{\beta}$ such that the cocycle conditions hold:
$$
\begin{aligned}
g_{\alpha \alpha}(p) & =e \text { for } p \in U_{\alpha} \\
g_{\alpha \beta}(p) & =\left(g_{\beta \alpha}(p)\right)^{-1} \text { for } p \in U_{\alpha} \cap U_{\beta} \\
g_{\alpha \beta}(p) g_{\beta \gamma}(p) g_{\gamma \alpha}(p) & =e \text { for } p \in U_{\alpha} \cap U_{\beta} \cap U_{\gamma}
\end{aligned}
$$
where $e$ is the identity in $G$. The family of maps $\left\{g_{\alpha \beta}\right\}$ forms a cocycle.
\end{DEF}



\begin{REM}{BÜN-L09-02-16}{Repräsentation von Wirkung von Übergangsabbildung durch Wirkung von Lie Gruppen}
The idea that we wish to pursue is that of representing the action of the raw transition maps by using Lie group actions. There is a subtle point here that the reader should not miss. Consider the following fact: If $\lambda$ : $G \times F \rightarrow F$ is a group action, then by letting $K=\left\{g: \lambda_{g}(p)=p\right.$ for all $p \in F\}$ (the kernel of the action) we obtain an effective action of $G / K$ on $F$. Things are not so simple on the global level of bundles as becomes clear when dealing with the notion of spin structure (See $[\mathbf{L}-\mathbf{M}]$ ). The best way to explain what is at stake is by the use of the notion of a principal bundle which we introduce later. Even before we get to that point, we will mention some things that will provide some idea as to why we need to be careful about ineffective actions.
\end{REM}

We start out assuming that the action is effective (see Definition 1.100):

Definition 6.11. 


\begin{DEF}{BÜN-L09-02-17}{($G,\lambda$)-Bündel Atlas}
Let $\xi=(E, \pi, M, F)$ be a fiber bundle and $G$ a Lie group. Suppose that we have an effective left action $\lambda: G \times F \rightarrow F$. Let $\left\{\left(\phi_{\alpha}, U_{\alpha}\right)\right\}$ be a bundle atlas for $\xi$. Suppose that for every nonempty intersection $U_{\alpha} \cap U_{\beta}$ there exists a smooth map $g_{\alpha \beta}: U_{\alpha} \cap U_{\beta} \rightarrow G$ such that $\lambda\left(g_{\alpha \beta}(p), y\right)= \left.\Phi_{\alpha \beta}\right|_{p}(y)$ for all $p \in U_{\alpha} \cap U_{\beta}$ and $y \in F$. Then the atlas $\left\{\left(\phi_{\alpha}, U_{\alpha}\right)\right\}$ is called a ( $G, \lambda$ )-bundle atlas. If the action $\lambda$ is understood or standard in some way, one also speaks of a $G$-bundle atlas.
\end{DEF}

\begin{CONC}{BÜN-L09-02-18}{Rolle von effektiver Wirkung für ($G,\lambda$)-Bündel Karte}
Because the action $\lambda$ in the above definition is assumed effective, it follows that the family $\left\{g_{\alpha \beta}\right\}$ satisfies the cocycle conditions of Definition 6.10. Notice that if we had not assume the action to be effective, then the maps $g_{\alpha \beta}$ would not be unique and may not satisfy a cocycle condition (although they would do so modulo the kernel of the action). Thus if we do not assume effectiveness, then we have to make the cocycle $\left\{g_{\alpha \beta}\right\}$ part of the definition. We return to this below.
\end{CONC}

The basic definition in the case of an effective action can be formulated as follows:

Definition 6.12. 


\begin{DEF}{BÜN-L09-02-19}{Effektive ($G,\lambda$)-Bündel und Bündel-Struktur}
Let $\xi=(E, \pi, M, F)$ be a fiber bundle and $G$ a Lie group. Suppose that we have an effective left action $\lambda: G \times F \rightarrow F$. Two $(G, \lambda)$ bundle atlases for $\xi$, say $\left\{\left(\phi_{\alpha}, U_{\alpha}\right)\right\}$ and $\left\{\left(\phi_{\alpha}^{\prime}, U_{\alpha}^{\prime}\right)\right\}$, are strictly equivalent if the union of the atlases is also a ( $G, \lambda$ )-bundle atlas. A strict equivalence class of atlases is referred to as an effective ( $G, \lambda$ )-bundle structure on $\xi$, and we say that $\xi$ together with this ( $G, \lambda$ )-bundle structure is an effective ( $G, \lambda$ )-bundle. Again, if the action is standard or understood, then it is common to speak of a $G$-bundle structure and refer to $\xi$ as a $G$-bundle.
\end{DEF}

\begin{CONC}{BÜN-L09-02-20}{Notation für Vereinigung von $(G,\lambda)$ Atlanten}
Actually, there is a tiny point to be made. To keep things neat we should always arrange that the indexing map $\alpha \mapsto\left(\phi_{\alpha}, U_{\alpha}\right)$ for any atlas is injective. Thus when taking the union of two atlases per the definition of strict equivalence, one may need to reindex so that the $g_{\alpha \beta}$ are notationally unambiguous. For example, if we take the union of an atlas $\left\{\left(\phi_{1}, U_{1}\right),\left(\phi_{2}, U_{2}\right)\right\}$ with an atlas $\left\{\left(\psi_{1}, V_{1}\right),\left(\psi_{2}, V_{2}\right)\right\}$, then how should the transition maps for the bigger atlas be denoted? What would $g_{11}$ mean? Sometimes the traditional indexing scheme is wisely dropped. Instead one uses the set of trivializing maps itself as the index set so that a chart is written as ( $\phi, U_{\phi}$ ) or ( $\psi, U_{\psi}$ ), and then one denotes transitions maps by $g_{\phi \psi}$ etc. The notion of maximal ( $G, \lambda$ )-bundle atlas is defined in the obvious way by direct analogy with the notion of maximal atlas for a smooth structure.
\end{CONC}

\begin{CONC}{BÜN-L09-02-21}{Effektive vs. Nicht effektive ($G,\lambda$)-Bündel}
Our main emphasis will be on the effective ( $G, \lambda$ )-bundles, but as mentioned above, if we wish to allow ineffective actions, then the notion of atlas should include the cocycle as part of the data. But even then we have to be careful. Indeed, for an ineffective action it is conceivable that there could be a different cocycle $\left\{g_{\alpha \beta}^{\prime}\right\}$ such that $\lambda\left(g_{\alpha \beta}^{\prime}(p), y\right)=\left.\Phi_{\alpha \beta}\right|_{p}(y)$ for all $y \in F$. Then $\left\{\left(\phi_{\alpha}, U_{\alpha}\right),\left(g_{\alpha \beta}\right), \lambda\right\}$ and $\left\{\left(\phi_{\alpha}, U_{\alpha}\right),\left(g_{\alpha \beta}^{\prime}\right), \lambda\right\}$ would be different (but possibly equivalent) $(G, \lambda)$-bundle atlases. For $\left\{\left(\phi_{\alpha}, U_{\alpha},\left(g_{\alpha \beta}\right), \lambda\right\}\right.$ and $\left\{\left(\phi_{j}^{\prime}, U_{j}^{\prime},\left(g_{i j}^{\prime}\right), \lambda\right\}\right.$ to define the same ( $G, \lambda$ )-bundle it must be the case that both of the cocycles are contained in a larger cocycle that gives the transitions for the atlas obtained as the union of the collection of charts from $\left\{\left(\phi_{\alpha}, U_{\alpha}\right),\left(g_{\alpha \beta}\right), \lambda\right\}$ and $\left\{\left(\phi_{j}^{\prime}, U_{j}^{\prime}\right),\left(g_{i j}^{\prime}\right), \lambda\right\}$. We handle ineffective actions this way so as to keep aligned with the notion of associated bundle introduced later.

If there is no chance of confusion, we will drop the adjective effective. The reader is warned that some standard expositions on fiber bundles allow ineffective actions right from the start, but in some cases assertions are made that would only be true in the effective case! It is interesting to note that in his famous book on the subject $[\mathbf{S t}]$, Norman Steenrod restricts himself to effective actions, although he announces this restriction in one easily overlooked sentence early in the book.
\end{CONC}

\begin{DEF}{BÜN-L09-02-22}{Übergangasabbildungen für ($G,\lambda$)-Bündel}
Notice that an alternative way to say that $\lambda\left(g_{\alpha \beta}(p), y\right)=\left.\Phi_{\alpha \beta}\right|_{p}(y)$ is $\phi_{\alpha} \circ \phi_{\beta}^{-1}(p, y)=\left(p, \lambda\left(g_{\alpha \beta}(p), y\right)\right)$ or
$$
\phi_{\alpha} \circ \phi_{\beta}^{-1}(p, y)=\left(p, g_{\alpha \beta}(p) \cdot y\right)
$$
(Actually, we shall at first avoid the notation $g \cdot y$ for the action of a group element $g$ on $y$ since we do not want the beginner to forget the role of the choice of action.) The maps $g_{\alpha \beta}$ are also called transition functions for the ( $G, \lambda$ )-bundle atlas. If the $G$ is literally a subgroup of $\operatorname{Diff}(F)$ and the action is simply $(\Phi, y) \mapsto \Phi(y)$, then the transition maps $g_{\alpha \beta}$ are simply the maps $\Phi_{\alpha \beta}$.
\end{DEF}

When dealing with $(G, \lambda)$-bundles there is a more specific notion of morphism and equivalence:

Definition 6.13. 


\begin{DEF}{BÜN-L09-02-23}{Kategorie der ($G,\lambda$)-Bündel}
Let $\xi_{1}=\left(E_{1}, \pi_{1}, M_{1}, F\right)$ be a ( $G, \lambda$ )-bundle with its $(G, \lambda)$-bundle structure determined by the strict equivalence class of the $(G, \lambda)$-atlas $\left\{\left(\varphi_{\alpha}, U_{\alpha}\right)\right\}_{\alpha \in A}$. Let $\xi_{2}=\left(E_{2}, \pi_{2}, M_{1}, F\right)$ be a $(G, \lambda)$-bundle with its $(G, \lambda)$-bundle structure determined by the strict equivalence class of the $(G, \lambda)$-atlas $\left\{\left(\psi_{\beta}, V_{\beta}\right)\right\}_{\beta \in B}$. Then a type II bundle morphism $(\widehat{h}, h): \xi_{1} \rightarrow \xi_{2}$ is called a ( $G, \lambda$ )-bundle morphism along $h$ if\\
\begin{enumerate}
\item $\widehat{h}$ carries each fiber of $E_{1}$ diffeomorphically onto the corresponding fiber of $E_{2}$;\\
\item whenever $U_{\alpha} \cap h^{-1}\left(V_{\beta}\right)$ is not empty, there is a smooth map $h_{\alpha \beta}$ : $U_{\alpha} \cap h^{-1}\left(V_{\beta}\right) \rightarrow G$ such that for each $p \in U_{\alpha} \cap h^{-1}\left(V_{\beta}\right)$ we have
$$
\left(\Psi_{\beta} \circ \widehat{h} \circ\left(\left.\Phi_{\alpha}\right|_{\pi_{1}^{-1}(p)}\right)^{-1}\right)(y)=\lambda\left(h_{\alpha \beta}(p), y\right) \text { for all } y \in F
$$
where as usual $\varphi_{\alpha}=\left(\pi_{1}, \Phi_{\alpha}\right)$ and $\psi_{\beta}=\left(\pi_{2}, \Psi_{\beta}\right)$.
\end{enumerate}
If $M_{1}=M_{2}$ and $h=\operatorname{id}_{M}$, then we call $\hat{h}$ a $(G, \lambda)$-bundle equivalence over $M$. (In this case, $\widehat{h}$ is a diffeomorphism.) Condition (ii) simply says that $\widehat{h}$ must be given by the action on each fiber when viewed in ( $G, \lambda$ )-charts. For this definition to be good it must be shown to be well-defined. That is, one must show that condition (ii) is independent of the choice of representatives $\left\{\left(\varphi_{\alpha}, U_{\alpha}\right)\right\}_{\alpha \in A}$ and $\left\{\left(\varphi_{i}, V_{i}\right)_{i \in J}\right\}$ of the strict equivalence classes of atlases that define the $(G, \lambda)$-bundle structures.
\end{DEF}


We leave this as an exercise. Later we will discover another, perhaps better, way to talk about equivalence of ( $G, \lambda$ )-bundles.

\begin{DEF}{BÜN-L09-02-24}{Triviale ($G,\lambda$)-Bündel}
The product bundle $\operatorname{pr}_{1}: M \times F \rightarrow M$ has a trivial $(G, \lambda)$-bundle structure for any $\lambda$ acting on $F$. Indeed, we just take the structure given by the single bundle chart ( $\operatorname{id}_{M \times F}, M$ ) where we have the resulting cocycle $\left\{g_{11}\right\}$, and where $g_{11}(x):=e$ for all $x$. In fact, if $\left\{U_{\alpha}\right\}_{\alpha \in A}$ is any open cover of $M$, then $\left\{\left(\operatorname{id}_{U_{\alpha} \times F}, U_{\alpha}\right)\right\}_{\alpha \in A}$ is a $(G, \lambda)$-atlas strictly equivalent to the atlas $\left\{\left(\operatorname{id}_{M \times F}, M\right)\right\}$ and so defining the same trivial $(G, \lambda)$-bundle. By trivial ( $G, \lambda$ )-bundle over $M$ we will mean either this product ( $G, \lambda$ )-bundle or one that is $(G, \lambda)$-equivalent to it.
\end{DEF}

Remark 6.14. 

\begin{CONC}{BÜN-L09-02-25}{($G,\lambda$)-Bündel-Äquivalenz im Fall von gleichem Basisraum und Projektionsabbildung}
The special case of a ( $G, \lambda$ )-bundle equivalence in the case that $E_{1}=E_{2}$ and $\pi_{1}=\pi_{2}$ is interesting but easily misunderstood. Suppose that $\left\{\left(\varphi_{\alpha}, U_{\alpha}\right)\right\}_{\alpha \in A}$ determines a $(G, \lambda)$-bundle structure on $(E, \pi, M, F)$ and suppose that $\left\{\left(\psi_{\alpha}, U_{\alpha}\right)\right\}_{\alpha \in A}$ determines another $(G, \lambda)$-bundle structure on ( $E, \pi, M, F$ ). Then these two ( $G, \lambda$ )-bundle structures might be ( $G, \lambda$ )bundle equivalent without being the very same. We would have ( $G, \lambda$ )bundle equivalence in case the two atlases give "cohomologous" cocycles. This is not the same as strict equivalence despite what one occasionally finds stated in the literature. Strict equivalence is the notion that defines the ( $G, \lambda$ )-bundle structure and is analogous to equivalence of manifold atlases which defines the notion of differentiable structure on a manifold. The other, weaker, notion of equivalence above is analogous to diffeomorphism and we know that two atlases on a manifold may define different smooth structures, which may or may not be diffeomorphic. For a more detailed discussion of these issues see the online supplement [Lee, Jeff].
\end{CONC}

We have defined fiber bundles, ( $G, \lambda$ )-bundles and various notions of morphism and equivalence. The question of classifying fiber bundles generally, or ( $G, \lambda$ )-bundles more specifically, is a huge part of topology which we do not have the space to discuss. The interested reader should consult [Hus], [Span] and [St].


$G$-bundle structures vs. $G$-structures. 

\begin{CONC}{BÜN-L09-02-26}{($G,\lambda$)- vs. G-Bündel-Struktur}
We have deliberately opted to use the term " $G$-bundle structure" rather than simply " $G$-structure", which could reasonably be taken to mean the same thing. Perhaps the reader is aware that there is a theory of $G$-structures on a smooth manifold (see [Stern]). One may rightly ask whether a $G$-structure on $M$ is nothing more than a $G$-bundle structure on the tangent bundle $T M$ where $G$ is a Lie subgroup of $G L(n)$ acting in the standard way. The answer is both yes and no. First, one could indeed say that a $G$-structure on $M$ is a kind of $G$-bundle structure on $T M$ even though the theory is usually fleshed out in terms of the frame bundle of $M$ (defined below). However, the notion of equivalence of two $G$-structures on $M$ is different from what we have given above. Roughly, $G$-structures on $M$ are equivalent in this sense if there is a diffeomorphism $\phi$ such that ( $T \phi, \phi$ ) is a type II bundle isomorphism that is also a $G$-bundle morphism along $\phi$.
\end{CONC}



\begin{DEF}{BÜN-L09-02-27}{Reduktion der Strukturgruppe}
Let $\xi=(E, \pi, M, F)$ have a ( $G, \lambda$ )-bundle structure given by a maximal $(G, \lambda)$-bundle atlas $\left\{\left(\phi_{\alpha}, U_{\alpha}\right)\right\}_{\alpha \in A}$. Suppose that there is a subatlas $\left\{\left(\phi_{j}, U_{j}\right)\right\}_{j \in J}, J \subset A$, such that the transition maps for the subatlas take values in a Lie subgroup $H$ of $G$. Then clearly $\xi$ has an ( $H,\left.\lambda\right|_{H}$ )-bundle structure, and this is called a reduction of the structure group. We say that the structure group is reducible.
\end{DEF}

The following theorem shows how we can build a fiber bundle using a cocycle and a choice of action.

Theorem 6.15 (Fiber bundle construction theorem). 

\begin{THEO}{BÜN-L09-02-28}{($G,\lambda$)-Bündel Konstruktion Theorem}
Let $M$ and $F$ be smooth manifolds and let $G$ be a Lie group. Let $\left\{U_{\alpha}\right\}_{\alpha \in A}$ be a cover of $M$ and $\left\{g_{\alpha \beta}\right\}$ a $G$-cocycle for the cover. For every action $\lambda: G \times F \rightarrow F$, there exists a fiber bundle with bundle-atlas $\left\{\left(U_{\alpha}, \phi_{\alpha}\right)\right\}$ satisfying $\phi_{\alpha} \circ \phi_{\beta}^{-1}(p, y)=$ ( $p, \lambda\left(g_{\alpha \beta}(p), y\right)$ ) on nonempty overlaps $U_{\alpha} \cap U_{\beta}$. Thus the resulting bundle has a $(G, \lambda)$-bundle structure.
\end{THEO}

Proof. 

\begin{PROOF}{BÜN-L09-02-29}{P: ($G,\lambda$)-Bündel Konstruktion Theorem}
On the union $\Sigma:=\bigcup_{\alpha}\{\alpha\} \times U_{\alpha} \times F$ define an equivalence relation such that $(\alpha, p, y) \in\{\alpha\} \times U_{\alpha} \times F$ is equivalent to $\left(\beta, p^{\prime}, y^{\prime}\right) \in\{\beta\} \times U_{\beta} \times F$ if and only if $p=p^{\prime}$ and $y=g_{\alpha \beta}(p) \cdot y^{\prime}$. Notice that $p=p^{\prime}$ is possible only in the case $U_{\alpha} \cap U_{\beta} \neq \emptyset$. The first member of the triple is only needed to make the union above disjoint. The cocycle conditions ensure that the equivalence relation is well-defined.

The total space of our bundle is then $E:=\Sigma / \sim$. The set $\Sigma$ is essentially the disjoint union of the product spaces $U_{\alpha} \times F$ and so has an obvious topology. We then give $E:=\Sigma / \sim$ the quotient topology. The bundle projection $\pi$ is induced by $(\alpha, p, v) \mapsto p$. To get our trivializations, we define
$$
\phi_{\alpha}(\epsilon):=(p, y) \text { for } \epsilon \in \pi^{-1}\left(U_{\alpha}\right)
$$
where ( $p, y$ ) is the unique member of $U_{\alpha} \times F$ such that $(\alpha, p, y) \in \epsilon$ (recall that $\epsilon$ is an equivalence class). The point here is that $\left(\alpha, p_{1}, y_{1}\right) \sim\left(\alpha, p_{2}, y_{2}\right)$ only if $\left(p_{1}, y_{1}\right)=\left(p_{2}, y_{2}\right)$. Now suppose $U_{\alpha} \cap U_{\beta} \neq \emptyset$. Then for $p \in U_{\alpha} \cap U_{\beta}$, the element $\phi_{\beta}^{-1}(p, y)$ is in $\pi^{-1}\left(U_{\alpha} \cap U_{\beta}\right)=\pi^{-1}\left(U_{\alpha}\right) \cap \pi^{-1}\left(U_{\beta}\right)$ and so $\phi_{\beta}^{-1}(p, y)= [(\beta, p, y)]$. But since $\phi_{\beta}^{-1}(p, y)$ is in $\pi^{-1}\left(U_{\alpha}\right)$, it must be equal to $\left[\left(\alpha, q, y_{2}\right)\right]$ for some $y_{2}$. This means that $p=q$ and $y_{2}=g_{\alpha \beta}(p) \cdot y=\lambda\left(g_{\alpha \beta}(p), y\right)$. Thus $\phi_{\beta}^{-1}(p, y)=\left[\alpha, p, g_{\alpha \beta}(p) \cdot y\right]$ and so
$$
\phi_{\alpha} \circ \phi_{\beta}^{-1}(p, y)=\left(p, g_{\alpha \beta}(p) \cdot y\right)
$$
We leave the existence of the smooth structure and the routine verification of the smoothness of these maps to the reader. That the topology induced is Hausdorff and paracompact is also easy to see.
\end{PROOF}



Example 6.16. 


\begin{EXA}{BÜN-L09-02-30}{($G,\lambda$)-Bündel Konstruktion auf $S^1$ (T1)}
Let $S^{1}$ be the circle realized as the unit complex numbers. We will construct a fiber bundle with typical fiber $F:=(-1,1) \subset \mathbb{R}$.

Let $U_{1}=\left\{e^{i \theta} \in S^{1}: 0<\theta<2 \pi\right\}$ and $U_{2}=\left\{e^{i \theta} \in S^{1}:-\pi<\theta<\pi\right\}$. Then $U_{1} \cap U_{2}$ is a disjoint union of open sets $V$ and $W$ where $1 \in V$ and $-1 \in W$. Now we define maps with values in the multiplicative group of two elements $\{1,-1\}=\mathbb{Z}_{2}$. Let $g_{11}(x)=1$ for all $x \in U_{1}, g_{22}(x)=1$ for all $x \in U_{1}$, and then let $g_{12}$ and $g_{21}$ be defined on $U_{1} \cap U_{2}$ by
$$
g_{12}(x):= \begin{cases}1 & \text { on } V, \\ -1 & \text { on } W,\end{cases}
$$
and $g_{21}:=g_{12}^{-1}$. Let $\mathbb{Z}_{2}$ act on the symmetric interval $(-1,1) \subset \mathbb{R}$ by multiplication (so that $-1 \cdot x:=-x$ ). Using the cocycle $\left\{g_{11}, g_{22}, g_{12}, g_{21}\right\}$ and this action on $\mathbb{R}$, Theorem 6.15 above gives a $\mathbb{Z}_{2}$-bundle which can be shown to be equivalent to the Möbius band bundle described in Example 6.9.
\end{EXA}

Example 6.17. 

\begin{EXA}{BÜN-L09-02-31}{($G,\lambda$)-Bündel Konstruktion auf $S^1$ (T2)}
Use the same cocycle on $S^{1}$ as in the last example but with typical fiber $S^{1}$ and action $\mathbb{Z}_{2} \times S^{1} \rightarrow S^{1}$ given by letting 1 act as the identity and -1 act by a rotation of $S^{1}$ by $\pi$, that is $-1 \cdot e^{i \theta}:=e^{i(\theta+\pi)}=-e^{i \theta}$. Then the bundle obtained is a $\mathbb{Z}_{2}$-bundle, sometimes called the twisted torus. It is of the upmost importance to realize that there is a fiber preserving diffeomorphism between the twisted torus and the trivial bundle $S^{1} \times S^{1} \xrightarrow{\mathrm{pr}_{1}} S^{1}$ and yet there is no $\mathbb{Z}_{2}$-bundle equivalence between the twisted torus and the trivial $\mathbb{Z}_{2}$-bundle $S^{1} \times S^{1} \xrightarrow{\mathrm{pr}_{1}} S^{1}$. Thus the twisted torus is trivial as a general bundle but not as a $\mathbb{Z}_{2}$-bundle (See Problem 5). This shows how much the involvement of the group matters.
\end{EXA}

\begin{CONC}{BÜN-L09-02-32}{($G,\lambda$)-Bündel Konstruktion für nicht effektive Wirkungen}
Notice that the previous theorem is true even without the assumption that the action is effective and we still end up with a genuine (ineffective) ( $G, \lambda$ )-bundle since there is no problem about the cocycle existing. A quotient by the kernel $K$ of the action would give an effective structure group action. Thinking of the action as a homomorphism $G \rightarrow \operatorname{Diff}(F)$, we have an induced homomorphism $G / K \rightarrow \operatorname{Diff}(F)$ such that the following diagram commutes:\\
\includegraphics[scale=0.25, center]{2025_10_09_265d7e1a22c89f79c21eg-283}
\end{CONC}

Definition 6.18. 


\begin{DEF}{BÜN-L09-02-33}{Pullback Bündel}
Let $\xi=(E, \pi, M, F)$ be a smooth fiber bundle and $f: N \rightarrow M$ a smooth map. The pull-back bundle $f^{*} \xi=\left(f^{*} E, \pi_{1}, M, F\right)$ (or induced bundle) is defined as follows: The total space $f^{*} E$ is the set
$$
f^{*} E:=\{(q, \epsilon) \in N \times E: f(q)=\pi(\epsilon)\} .
$$
Then we define $\pi_{1}$ as the restriction to $f^{*} E$ of the projection $\operatorname{pr}_{1}: N \times E \rightarrow$ N.
\end{DEF}

\begin{CONC}{BÜN-L09-02-34}{Induzierter Bündel Homomorphismus unter Pullback}
Notice that the second factor projection map $\operatorname{pr}_{2}: N \times E \rightarrow E$ restricts to a map $\widetilde{f}: f^{*} E \rightarrow E$ which is a bundle morphism over the map $f$ :\\
\includegraphics[scale=0.25, center]{2025_10_09_265d7e1a22c89f79c21eg-284}

The map $\widetilde{f}$ restricts to a diffeomorphism on each fiber. If $\phi=(\pi, \Phi)$ is a trivialization of the bundle $\xi$ over the open set $U$, then $\left(\pi_{1}, \bar{\Phi}\right):=\left(\pi_{1}, \Phi \circ \widetilde{f}\right)$ is a trivialization of $f^{*} \xi$ over the open set $f^{-1}(U)$. Thus a bundle atlas on $\xi$ induces a bundle atlas on $f^{*} \xi$. If $\left\{\Phi_{\alpha \beta}\right\}$ are (raw) transition maps for $\xi$ corresponding to a bundle atlas $\left\{\left(\pi, \Phi_{\alpha}\right)\right\}_{\alpha \in A}$, then $\left\{\Phi_{\alpha \beta} \circ f\right\}$ are transition maps for $f^{*} \xi$ corresponding to the atlas $\left\{\left(\pi_{1}, \Phi_{\alpha} \circ \widetilde{f}\right)\right\}_{\alpha \in A}$. In fact,
$$
\left.\bar{\Phi}_{\alpha}\right|_{q \times E_{f(q)}}(q, \epsilon)=\left.\Phi_{\alpha} \circ \widetilde{f}\right|_{q \times E_{f(q)}}(q, \epsilon)=\left.\Phi_{\alpha}\right|_{E_{f(q)}}(\epsilon)
$$
so the inverse is
$$
\left.\bar{\Phi}_{\alpha}\right|_{q \times E_{f(q)}} ^{-1}: y \mapsto\left(q,\left.\Phi_{\alpha}\right|_{E_{f(q)}} ^{-1}(y)\right) .
$$
Thus $\left.\left.\bar{\Phi}_{\alpha}\right|_{q \times E_{f(q)}} \circ \bar{\Phi}_{\beta}\right|_{q \times E_{f(q)}} ^{-1}$ is given by
$$
\left.\left.y \mapsto\left(q,\left.\Phi_{\beta}\right|_{E_{f(q)}} ^{-1}(y)\right) \mapsto \Phi_{\alpha}\right|_{E_{f(q)}} \circ \Phi_{\beta}\right|_{E_{f(q)}} ^{-1}(y),
$$
and so
$$
\begin{aligned}
\bar{\Phi}_{\alpha \beta}(q) & =\left.\left.\bar{\Phi}_{\alpha}\right|_{q \times E_{f(q)}} \circ \bar{\Phi}_{\beta}\right|_{q \times E_{f(q)}} ^{-1} \\
& =\left.\left.\Phi_{\alpha}\right|_{E_{f(q)}} \circ \Phi_{\beta}\right|_{E_{f(q)}} ^{-1}=\Phi_{\alpha \beta} \circ f(q)
\end{aligned}
$$

Furthermore, if $\left\{\left(\pi, \Phi_{\alpha}\right)\right\}_{\alpha \in A}$ is a $(G, \lambda)$-atlas for $\xi$, then since $\left.\Phi_{\alpha \beta}\right|_{p}(y)= \lambda\left(g_{\alpha \beta}(p), y\right)$, we have $\left.\bar{\Phi}_{\alpha \beta}\right|_{q}(y)=\Phi_{\alpha \beta}(f(q))(y)=\lambda\left(g_{\alpha \beta}(f(q)), y\right)$. We see that $f^{*} \xi$ has a $(G, \lambda)$-bundle structure with cocycles $g_{\alpha \beta} \circ f$. Note, however, that because of the composition with $f$, this structure may be reducible to a smaller group.
\end{CONC}


Definition 6.19. 


\begin{DEF}{BÜN-L09-02-35}{Schnitte entlang glatter Abbildungen zwischen Basisräumen}
Let $\xi=(E, \pi, M, F)$ be a smooth fiber bundle and let $f: N \rightarrow M$ be a smooth map. A section of $\xi$ along $f$ is a map $\sigma: N \rightarrow E$ such that $\pi \circ \sigma=f$. The set of sections along $f$ will be denoted $\Gamma_{f}(\xi)$ or $\Gamma_{f}(E)$.
\end{DEF}

\begin{DEF}{BÜN-L09-02-36}{Induzierter Schnitt unter Pullback}
If $\sigma: N \rightarrow E$ is a section of $\xi$ along $f$, then the map $\sigma^{\prime}: N \rightarrow f^{*} E$ given by $p \mapsto(p, \sigma(p))$ is a section of the pull-back bundle $f^{*} \xi$. It is not hard to show that all sections of $f^{*} \xi$ have this form.
\end{DEF}



Proposition 6.20. 


\begin{PROP}{BÜN-L09-02-37}{Natürliche Bijektion zwischen Mengen der Schnitten längs Abbildungen und Schnitte im Pullback}
Let $\xi=(E, \pi, M, F)$ be a smooth fiber bundle and $f: N \rightarrow M$ a smooth map. Then there is a natural bijection between $\Gamma_{f}(\xi)$ and $\Gamma\left(f^{*} \xi\right)$.
\end{PROP}

Proof. 


\begin{PROOF}{BÜN-L09-02-38}{P: Natürliche Bijektion zwischen Mengen der Schnitten längs Abbildungen und Schnitte im Pullback}
Let $s \in \Gamma\left(f^{*} \xi\right)$. For each $p \in N$, we have $s(p) \in\left(f^{*} E\right)_{p}$, which must have the form ( $p, y$ ) for some $y \in E_{f(p)}$. Therefore the smooth map $\sigma:=\mathrm{pr}_{2} \circ s: N \rightarrow E$ has the property that $\pi \circ \sigma=f$. Thus we obtain a map $\Gamma\left(f^{*} \xi\right) \rightarrow \Gamma_{f}(\xi)$ given by $s \mapsto \sigma$. But the inverse of this map is clearly $\sigma \mapsto\left(\mathrm{id}_{N}, \sigma\right)$.
\end{PROOF}


\pagebreak



\subsection*{Vector Bundles}


The tangent and cotangent bundles are examples of a general type of fiber bundle called a vector bundle. Roughly speaking, a vector bundle is a parametrized family of vector spaces. We shall need both complex vector bundles and real vector bundles, and so to facilitate definitions we let $\mathbb{F}$ denote either $\mathbb{R}$ or $\mathbb{C}$. Let V be a finite-dimensional $\mathbb{F}$-vector space. The simplest examples of vector bundles over a manifold $M$ are the product vector bundles which consist of a Cartesian product $M \times \mathrm{V}$ together with the projection onto the first factor $\operatorname{pr}_{1}: M \times \mathrm{V} \rightarrow M$. Each set of the form $\{x\} \times \mathrm{V} \subset M \times \mathrm{V}$ inherits an $\mathbb{F}$-vector space structure from that of V in the obvious way: $a(p, v)+b(p, w):=(p, a v+b w)$. We think of $M \times \mathrm{V}$ as copies of V parametrized by $M$.

Definition 6.21. 


\begin{DEF}{BÜN-L09-03-01}{Glatte Vektorbündel}
Let V be a finite-dimensional $\mathbb{F}$-vector space. A smooth $\mathbb{F}$-vector bundle with typical fiber V is a fiber bundle ( $E, \pi, M, \mathrm{V}$ ) such that:\\
\begin{enumerate}
\item for each $x \in M$ the set $E_{x}:=\pi^{-1}(x)$ has the structure of a vector space over the field $\mathbb{F}$, isomorphic to the fixed vector space $V$;\\
\item every $p \in M$ is in the domain of some bundle chart ( $U, \phi$ ), such that for each $x \in U$ the map $\left.\Phi\right|_{E_{x}}: E_{x} \rightarrow \mathrm{V}$ is a vector space isomorphism where $\phi=(\pi, \Phi)$.
\end{enumerate}
\end{DEF}


Definition 6.22 (Terminology). We refer to a vector bundle as a complex vector bundle (resp. real vector bundle) if $\mathbb{F}=\mathbb{C}$ (resp. $\mathbb{F}=\mathbb{R}$ ). "Vector bundle" shall mean either real or complex vector bundle as determined by the context.


\begin{DEF}{BÜN-L09-03-02}{Vektorbündel-Karten und Atlanten}
Sei $\pi:E\to M$ ein glattes \emph{Vektorbündel} mit typischer Faser $V$.
\begin{enumerate}
\item Eine \emph{Vektorbündelkarte} (oder \emph{lokale Vektorbündeltrivialisierung})
ist eine Bündelkarte $(U,\phi)$ der Form
\[
\phi:\pi^{-1}(U)\longrightarrow U\times V,
\]
bei der für jedes $p\in U$ die Einschränkung
\[
\phi|_{E_p}:E_p\longrightarrow\{p\}\times V
\]
eine lineare Abbildung der Vektorräume ist.
\item Eine Familie $\{(U_\alpha,\phi_\alpha)\}_{\alpha\in A}$ von Vektorbündelkarten,
deren offene Mengen $\{U_\alpha\}$ den Basisraum $M$ überdecken,
heißt ein \emph{Vektorbündelatlas} für $\pi:E\to M$.
\item Die Dimension der typischen Faser $V$ heißt der \emph{Rang} des Vektorbündels.
\end{enumerate}
\end{DEF}


A bundle chart ( $U, \phi$ ) of the sort described in the definition is called a vector bundle chart (VB-chart) or a local vector bundle trivialization (over $U$ ). In the setting of vector bundles, a local trivialization is assumed to be linear on fibers. A family $\left\{\left(U_{\alpha}, \phi_{\alpha}\right)\right\}$ of vector bundle charts such that $\left\{U_{\alpha}\right\}$ is an open cover of $M$ is called a vector bundle atlas for $\pi: E \rightarrow M$. The definition of a vector bundle guarantees that such an atlas exists. The dimension of the typical fiber $V$ is called the rank of the vector bundle.

Remark 6.23. 

\begin{CONC}{BÜN-L09-03-03}{Übersetzung von $V$-wertigen VB-Karten in $\R^n,\C^n$-wertige VB-Karten}
By choosing a basis for V , one gets an isomorphism with $\mathbb{C}^{k}$ or $\mathbb{R}^{k}$ as the case may be. Composing with this isomorphism we can convert the V-valued VB-charts into $\mathbb{C}^{k}$ - or $\mathbb{R}^{k}$-valued VB-charts. Thus we could have assumed from the start that we were dealing with one of these standard vector spaces, but it is not always natural to do so since our vector space may arise in a specific way (it could be a Lie algebra or perhaps a space of algebraic tensors) and may not have a preferred choice of basis.
\end{CONC}

Exercise 6.24. Show that the tangent and cotangent bundles of an $n$ manifold are vector bundles with typical fiber $\mathbb{R}^{n}$ (the cotangent bundle may be viewed as having typical fiber $\left(\mathbb{R}^{n}\right)^{*}$ ).

\begin{DEF}{BÜN-L09-03-04}{Raum der glatten Schnitte in VB als $C^\infty(M,\mathbb{F})$-Modul}
Let $\mathbb{F}$ be $\mathbb{R}$ or $\mathbb{C}$ as above. The space of smooth sections $\Gamma(\xi)$ of an $\mathbb{F}$-vector bundle $\xi=(E, \pi, M, \mathrm{V})$ has the structure of a module over the ring $C^{\infty}(M ; \mathbb{F})$; for $\sigma, \sigma_{1}, \sigma_{2} \in \Gamma(\xi)$ and $f \in C^{\infty}(M ; \mathbb{F})$, we define
$$
\begin{aligned}
\left(\sigma_{1}+\sigma_{2}\right)(p) & :=\sigma_{1}(p)+\sigma_{2}(p) \text { for all } p \in M, \\
f \sigma(p) & :=f(p) \sigma(p) \text { for all } p \in M .
\end{aligned}
$$
\end{DEF}

Definition 6.25. 


\begin{DEF}{BÜN-L09-03-05}{Kategorie der Vektorbündel}
Let $\xi_{1}$ and $\xi_{2}$ be $\mathbb{F}$-vector bundles with respective bundle projections $\pi_{1}$ and $\pi_{2}$. A bundle morphism $(\widehat{f}, f): \xi_{1} \rightarrow \xi_{2}$ is called a vector bundle morphism if the restrictions to fibers, $\left.\widehat{f}\right|_{\pi_{1}^{-1}(p)} \pi_{1}^{-1}(p) \rightarrow \pi_{2}^{-1}(f(p))$, are $\mathbb{F}$-linear. If $\xi_{1}$ and $\xi_{2}$ have the same base space $M$, then we obtain the definition of a vector bundle morphism over $M$ by specializing to the case $f=\operatorname{id}_{M}$. We then also have the corresponding notions of vector bundle isomorphism and automorphism (for both type I and II bundle morphisms).
\end{DEF}



\begin{DEF}{BÜN-L09-03-06}{Triviales Vektorbündel}
A vector bundle ( $E, \pi, M, \mathrm{V}$ ) is said to be trivial if it is vector bundle isomorphic to the product vector bundle $\operatorname{pr}_{1}: M \times \mathrm{V} \rightarrow M$. This happens exactly when there is a vector bundle trivialization over the entire manifold $M$, which we call a global vector bundle trivialization (a notion already introduced for tangent bundles).
\end{DEF}

Definition 6.26. 


\begin{DEF}{BÜN-L09-03-07}{Vektorunterbündel}
Let ( $E, \pi, M, \mathrm{V}$ ) be a rank $k$ vector bundle with typical fiber V and fix an $l$-dimensional subspace $\mathrm{V}^{\prime}$ of V . If $E^{\prime} \subset E$ is a submanifold with the property that for every $p \in M$ there is a VB-chart ( $U, \phi$ ) such that
$$
\phi\left(\pi^{-1}(U) \cap E^{\prime}\right)=U \times \mathrm{V}^{\prime} \subset U \times \mathrm{V}
$$
then ( $E^{\prime},\left.\pi\right|_{E^{\prime}}, M, \mathrm{V}^{\prime}$ ) is called a rank $l$ vector subbundle of ( $E, \pi, M, \mathrm{V}$ ). Charts with this property are said to be adapted to the subbundle.

The triple ( $E^{\prime},\left.\pi\right|_{E^{\prime}}, M, \mathrm{V}^{\prime}$ ) is a vector bundle and every adapted VBchart ( $U, \phi$ ) on $E$ gives rise to a chart on ( $E^{\prime},\left.\pi\right|_{E^{\prime}}, M, \mathrm{V}^{\prime}$ ); namely, ( $U, \phi^{\prime}$ ), where $\phi^{\prime}$ is the restriction of $\phi$ to $\pi^{-1}(U) \cap E^{\prime}=\left.\pi\right|_{E^{\prime}} ^{-1}(U)$. By picking a basis for $\mathrm{V}^{\prime}$ and extending to a basis for V one may take V to be $\mathbb{R}^{k}$ and $\mathrm{V}^{\prime}$ to be $\mathbb{R}^{l}$ embedded in $\mathbb{R}^{k}$ as $\mathbb{R}^{l} \times\{0\} \subset \mathbb{R}^{k}$.
\end{DEF}

Exercise 6.27. 


\begin{EXE}{BÜN-L09-03-08}{Konstruktion von Vektorunterbündel}
Let $E \rightarrow M$ be a vector bundle as above. Suppose that a subspace $E_{p}^{\prime}$ of $E_{p}$ is given for each $p \in M$ and consider the set $E^{\prime}= \bigcup_{p \in M} E_{p}^{\prime}$. Show that $E^{\prime}$ is the total space of a rank $l$ vector subbundle if and only if for each $p \in M$, there is an open neighborhood $U$ of $p$ on which smooth sections $\sigma_{1}, \ldots, \sigma_{l}$ are defined such that for each $q \in U$ the set $\left\{\sigma_{1}(q), \ldots, \sigma_{l}(q)\right\}$ is a basis of the subspace $E_{q}^{\prime}$.
\end{EXE}

\begin{DEF}{BÜN-L09-03-09}{Image und Kern von VB-Morphismen als Vekotunterbündel}
If $h: E_{1} \rightarrow E_{2}$ is a vector bundle morphism over $M$, then
$$
\operatorname{Ker} h:=\left.\bigcup_{p \in M} \operatorname{Ker} h\right|_{E_{1 p}}
$$
is a subset of $E_{1}$. This subset is not necessarily (the total space of) a subbundle, at least in the ordinary sense. However, if the rank $l$ of $\left.h\right|_{E_{1 p}}$ is independent of $p$, then we say that the bundle map has rank $l$ and, in this case, $\operatorname{Ker} h$ is a vector subbundle. Similarly, if $h$ has constant rank in this sense, then the image $\operatorname{Im} h$ is a vector subbundle of $E_{2}$.
\end{DEF} 


Both of these facts follow from

Proposition 6.28. 


\begin{PROP}{BÜN-L09-03-10}{Konstruktion von Vektorunterbündel aus Kern und Image}
Suppose that $h: E_{1} \rightarrow E_{2}$ is a vector bundle morphism over $M$ of constant rank $r$ and that $E_{1}$ and $E_{2}$ have typical fibers $\mathrm{V}_{1}$ and $\mathrm{V}_{2}$ respectively. Fix a rank $r$ linear map $A: \mathrm{V}_{1} \rightarrow \mathrm{V}_{2}$. Then for every $p \in M$ there is a VB-chart ( $U, \phi$ ) for $E_{1}$ with $p \in U$ and a $\operatorname{VB}$-chart ( $U, \psi$ ) for $E_{2}$ such that $\psi \circ h \circ \phi^{-1}: U \times \mathrm{V}_{1} \rightarrow U \times \mathrm{V}_{2}$ has the form
$$
(p, v) \mapsto(p, A v)
$$
It follows that $\operatorname{Ker} h$ is a vector subbundle with typical fiber $\operatorname{Ker} A$ and $\operatorname{Im} h$ is a vector subbundle of $E_{2}$ with typical fiber $\operatorname{Im} A$.
\end{PROP}


Proof. 

\begin{PROOF}{BÜN-L09-03-11}{P: Konstruktion von Vektorunterbündel aus Kern und Image}
Let us first make an observation. Notice that only the rank of the linear map $A: \mathrm{V}_{1} \rightarrow \mathrm{V}_{2}$ in this last proposition is important and we may replace $A$ by any linear map of the same rank. The reason for this is that if $B: \mathrm{V}_{1} \rightarrow \mathrm{V}_{2}$ is any other linear map with the same rank as $A$, then there exist linear isomorphisms $\alpha$ and $\beta$ such that $B=\beta A \alpha^{-1}$. In particular, if one has chosen bases and identified $\mathrm{V}_{1}$ with $\mathbb{R}^{k_{1}}$ and $\mathrm{V}_{2}$ with $\mathbb{R}^{k_{2}}$, then we may take $A$ to be a map of the form
$$
\left(x^{1}, \ldots, x^{k_{1}}\right) \mapsto\left(x^{1}, \ldots, x^{r}, 0, \ldots, 0\right),
$$
so that $\operatorname{Ker} A$ is a copy of $\mathbb{R}^{k_{1}-r}$ and $\operatorname{Im} h$ is a copy of $\mathbb{R}^{r}$.

What we need to prove is entirely local. Thus our task is to show that for any smooth map $h: U \times \mathbb{F}^{k_{1}} \rightarrow U \times \mathbb{F}^{k_{2}}$ of the form $(p, v) \rightarrow\left(p, h_{p} v\right)$, with $h_{p}$ a linear map of rank $r$ and where $p \mapsto h_{p}$ is smooth, we may find maps $\psi$ and $\phi$ such that $\psi \circ h \circ \phi^{-1}: U \times \mathbb{F}^{k_{1}} \rightarrow U \times \mathbb{F}^{k_{2}}$ is given by $\left(p, x^{1}, \ldots, x^{k_{1}}\right) \mapsto\left(p, x^{1}, \ldots, x^{r}, 0, \ldots, 0\right)$. Fix $p_{0} \in U$. There exist linear isomorphisms $\alpha: \mathbb{F}^{k_{1}} \rightarrow \mathbb{F}^{k_{1}}$ and $\beta: \mathbb{F}^{k_{2}} \rightarrow \mathbb{F}^{k_{2}}$ such that $\beta \circ h_{p} \circ \alpha^{-1}$ is given by a $k_{2} \times k_{1}$ matrix of the form
$$
\left[\begin{array}{cc}
A_{11}(p) & A_{12}(p) \\
A_{21}(p) & A_{22}(p)
\end{array}\right]
$$
and where $A_{11}\left(p_{0}\right)$ is an $r \times r$ matrix which is invertible. By shrinking $U$ if needed, we may assume that $A_{11}(p)$ is invertible for all $p \in U$. Thus we may as well assume from the start that $h_{p}$ is represented by a matrix of this form. Now consider the map $\phi_{p}: \mathbb{F}^{k_{1}} \rightarrow \mathbb{F}^{k_{1}}$ whose matrix is given by
$$
\left[\begin{array}{cc}
A_{11}(p) & A_{12}(p) \\
0 & I_{\left(k_{1}-r\right) \times\left(k_{1}-r\right)}
\end{array}\right]_{k_{1} \times k_{1}}
$$
Then $h_{p} \circ \phi_{p}^{-1}$ has a matrix of the form
$$
\left[\begin{array}{cc}
I_{r \times r} & 0 \\
A_{21}(p) A_{11}^{-1}(p) & C
\end{array}\right]_{k_{2} \times k_{1}}
$$
and since this matrix must have rank $r$, we see that $C=0$. Let $M_{p}:= A_{21}(p) A_{11}^{-1}(p)$ and let $\psi_{p}$ be the linear map $\mathbb{R}^{k_{2}} \rightarrow \mathbb{R}^{k_{2}}$ with matrix
$$
\left[\begin{array}{cc}
I_{r \times r} & 0 \\
-M_{p} & I_{\left(k_{2}-r\right) \times\left(k_{2}-r\right)}
\end{array}\right]
$$
Then $\psi_{p} \circ h_{p} \circ \phi_{p}^{-1}$ has the form $\left[\begin{array}{cc}I_{r \times r} & 0 \\ 0 & 0\end{array}\right]$. Now define $\phi(p, v)=\left(p, \phi_{p} v\right)$, $h(p, x):=\left(p, h_{p} x\right)$ and $\psi(p, v):=\left(p, \psi_{p} v\right)$ for $p \in U, x \in \mathbb{F}^{k_{1}}$, and $v \in \mathbb{F}^{k_{2}}$. Notice that $\psi_{p}, h_{p}$ and $\phi_{p}^{-1}$ each depend smoothly on $p$. The map $\psi \circ h \circ \phi^{-1}$ has the required form.
\end{PROOF}



Proposition 6.29. 


\begin{PROP}{BÜN-L09-03-12}{Charakterisierung von Faserbündel von Vektorbündel-Struktur auf Faserbündeln durch $\GL(V)$-Wirkung}
Let $\lambda_{0}: \mathrm{GL}(\mathrm{V}) \times \mathrm{V} \rightarrow \mathrm{V}$ be the standard action of $\mathrm{GL}(\mathrm{V})$ on the $\mathbb{F}$-vector space V . A fiber bundle with typical fiber V has an $\mathbb{F}$-vector bundle structure if and only if it admits a $\lambda_{0}$-bundle atlas ( $a \mathrm{GL}(\mathrm{V})$ bundle atlas). Furthermore, if $\lambda: \mathrm{GL}(\mathrm{V}) \times \mathrm{V} \rightarrow \mathrm{V}$ is any effective action which acts linearly, then any fiber bundle ( $E, \pi, M, \mathrm{V}$ ) that has a $\lambda$-atlas is a vector bundle in a natural way.
\end{PROP}

Proof. 

\begin{PROOF}{BÜN-L09-03-13}{P: Charakterisierung von Faserbündel von Vektorbündel-Struktur auf Faserbündeln durch $\GL(V)$-Wirkung}
That a vector bundle has a GL(V)-bundle structure follows directly from the definition. All that remains to show is the second part of the theorem, since this will imply the remainder of the first part. Let $\lambda: G \times \mathrm{V} \rightarrow$ V be any effective Lie group action which acts linearly and suppose that ( $E, \pi, M, \mathrm{V}$ ) has a $\lambda$-bundle structure. Let ( $U_{\alpha}, \phi_{\alpha}$ ) and ( $U_{\beta}, \phi_{\beta}$ ) be $\lambda$ compatible bundle charts and let $\phi_{\alpha}=\left(\pi, \Phi_{\alpha}\right)$ and $\phi_{\beta}=\left(\pi, \Phi_{\beta}\right)$. Fix $p \in U_{\alpha} \cap U_{\beta}$. For $v, w \in \mathrm{V}$, and $a, b \in \mathbb{F}$, we have
$$
\begin{aligned}
\Phi_{\alpha \beta}(p)(a v+b w) & =\lambda\left(g_{\alpha \beta}(p), a v+b w\right) \\
& =a \lambda\left(g_{\alpha \beta}(p), v\right)+b \lambda\left(g_{\alpha \beta}(p), w\right) \\
& =a \Phi_{\alpha \beta}(p)(v)+b \Phi_{\alpha \beta}(p)(w),
\end{aligned}
$$
which shows that $\Phi_{\alpha \beta}(p) \in \mathrm{GL}(\mathrm{V})$ for all $p \in U_{\alpha} \cap U_{\beta}$. We transfer the vector space structure from V to $E_{p}$ via $\left.\Phi_{\alpha}\right|_{E_{p}} ^{-1}$ and note that this is welldefined by Proposition 2.3. With this linear structure on the fibers it is easy to verify that ( $E, \pi, M, \mathrm{V}$ ) is a vector bundle.
\end{PROOF}


Theorem 6.30 (Vector bundle construction theorem). 


\begin{THEO}{BÜN-L09-03-14}{($G,\lambda$)-Vektorbündel Konstruktion Theorem}
Let $\left\{U_{\alpha}\right\}_{\alpha \in A}$ be a cover of $M$ and let $\left\{g_{\alpha \beta}\right\}$ be a $G$-cocycle for a Lie group $G$. If $G$ acts linearly on the vector space V (by say $\lambda$ ), then there exists a vector bundle over $M$ with a VB-atlas $\left\{\left(U_{\alpha}, \phi_{\alpha}\right)\right\}$ satisfying $\phi_{\alpha} \circ \phi_{\beta}^{-1}(p, v)=\left(p, g_{\alpha \beta}(p) \cdot v\right)$ ) on nonempty overlaps $U_{\alpha} \cap U_{\beta}$. In other words, there exists a vector bundle with $(G, \lambda)$-atlas.
\end{THEO}

Proof. 

\begin{PROOF}{BÜN-L09-03-15}{P: ($G,\lambda$)-Vektorbündel Konstruktion Theorem}
This is essentially a special case of Theorem 6.15. One only needs to check linearity of the $\phi_{\alpha}$ on fibers.
\end{PROOF}


\begin{DEF}{BÜN-L09-03-16}{Standard Übergangsabbildung und $\GL(V)$-Bündel-Struktur}
Since in the vector bundle case, the $\Phi_{\alpha \beta}$ arise directly from a VB-atlas and act by the standard action, we will call these standard transition maps, and the corresponding GL(V)bundle structure will be called the standard GL(V)-bundle structure. The standard GL(V)-bundle structure is the structure that a vector bundle has simply by virtue of being an $\mathbb{F}$-vector bundle with typical fiber V.
\end{DEF}

\begin{CONC}{BÜN-L09-03-17}{Relation zwischen $G$ in $G$-Bündel Struktur und $\GL(V)$ für Vektorbündel}
Perhaps some clarification is in order. In the case of a vector bundle, the raw transition maps $\Phi_{\alpha \beta}$ take values in the general linear group $\mathrm{GL}(\mathrm{V})$, which is a Lie group. They correspond to a $\lambda_{0}$-bundle structure where $\lambda_{0}$ is the standard linear action of $\mathrm{GL}(\mathrm{V})$ on V (the standard representation), and they automatically satisfy the cocycle condition. The more general transition maps that define a ( $G, \lambda$ )-bundle structure ( $G$-bundle structure) are $G$-valued. It is important to note that $G$ may be small compared to GL(V) and certainly need not be thought of as a subset of GL(V). For example, the tensor bundles have (possibly ineffective) GL(V)-bundle structures coming from tensor representations, but the tensor bundles themselves generally have rank greater than $k=\operatorname{dim}(\mathrm{V})$.
\end{CONC}



Remark 6.31. 

\begin{CONC}{BÜN-L09-03-18}{Anmerkung zu Darstellung und links lineare Wikrung}
We have previously mentioned that the notion of a representation is equivalent to that of a left linear action. When dealing with vector bundles it is perhaps more common to use the representation terminology and notation and this we shall do as convenient. So if $\lambda$ is a left linear action, then the map $G \rightarrow \mathrm{GL}(\mathrm{V})$ given by $g \mapsto \lambda(g):=\lambda_{g}$ is a representation of $G$. Conversely, if $\lambda$ is such a representation, we obtain a linear action by letting $\lambda(g, v):=\lambda(g) v$.
\end{CONC}

We already know what it means for two vector bundles over $M$ to be equivalent. Of course any two vector bundles that are equivalent in a natural way can be thought of as the same. Since we can and often do construct our bundles according to the above recipe, it will pay to know something about when two vector bundles over $M$ are isomorphic, based on their respective transition functions. Notice that the standard transition functions are easily recovered from every $(G, \lambda)$-atlas by the formula $\lambda\left(g_{\alpha \beta}(p)\right) y=\left.\Phi_{\alpha \beta}\right|_{p}(y)$.

Proposition 6.32. 


\begin{PROP}{BÜN-L09-03-19}{Charakterisierung von isomorphen Vektorbündeln durch standard Übergangsabbildungen}
Two vector bundles $\pi: E \rightarrow M$ and $\pi^{\prime}: E^{\prime} \rightarrow M$ with standard transition maps $\left\{\Phi_{\alpha \beta}: U_{\alpha} \cap U_{\beta} \rightarrow \mathrm{GL}(\mathrm{V})\right\}$ and $\left\{\Phi_{\alpha \beta}^{\prime}: U_{\alpha} \cap U_{\beta} \rightarrow\right. \mathrm{GL}(\mathrm{V})\}$ over the same cover $\left\{U_{\alpha}\right\}$ are isomorphic (over $M$ ) if and only if there are $\mathrm{GL}(\mathrm{V})$-valued functions $f_{\alpha}$ defined on each $U_{a}$ such that
$$
\Phi_{\alpha \beta}^{\prime}(x)=f_{\alpha}(x) \Phi_{\alpha \beta}(x) f_{\beta}^{-1}(x) \text { for } x \in U_{\alpha} \cap U_{\beta} .
$$
\end{PROP}

Proof. 

\begin{PROOF}{BÜN-L09-03-20}{P: Charakterisierung von isomorphen Vektorbündeln durch standard Übergangsabbildungen}
(Sketch) Given a vector bundle isomorphism $f: E \rightarrow E^{\prime}$ over $M$, let $f_{\alpha}(x):=\left.\Phi_{\alpha}^{\prime} \circ f \circ \Phi_{\alpha}\right|_{E_{x}}$. Check that this works. Conversely, given functions $f_{\alpha}$ satisfying equations (6.1), define $\widetilde{f}_{\alpha}: U_{\alpha} \times \mathrm{V} \rightarrow U_{\alpha} \times \mathrm{V}$ by $(x, v) \mapsto\left(x, f_{\alpha}(x) v\right)$. We define $f: E \rightarrow E^{\prime}$ by
$$
f(\epsilon):=\left(\left(\phi_{\alpha}^{\prime}\right)^{-1} \circ \widetilde{f}_{\alpha} \circ \phi_{\alpha}\right)(\epsilon) \text { for }\left.\epsilon \in E\right|_{U_{a}} .
$$
The conditions (6.1) insure that $f$ is well-defined on the overlaps $\left.E\right|_{U_{a}} \cap \left.E\right|_{U_{\beta}}=\left.E\right|_{U_{a} \cap U_{\beta}}$. One easily checks that this is a vector bundle isomorphism.
\end{PROOF}

We can use this construction to arrive at several common vector bundles.



Example 6.33. 


\begin{EXA}{BÜN-L09-03-21}{Konstruktion von (Ko)Tangentialbündel}
Given an atlas $\left\{\left(U_{\alpha}, \mathbf{x}_{\alpha}\right)\right\}$ for a smooth manifold $M$, we let $g_{\alpha \beta}(p)=T_{p} \mathrm{x}_{\alpha} \circ\left(T_{p} \mathrm{x}_{\beta}\right)^{-1}$ for all $p \in U_{\alpha} \cap U_{\beta}$. The bundle constructed according to the recipe of Theorem 6.30 is a vector bundle which is (naturally isomorphic to) the tangent bundle $T M$. If we let $g_{\alpha \beta}^{*}(p)=\left(T_{p} \mathrm{x}_{\beta} \circ\left(T_{p} \mathrm{x}_{\alpha}\right)^{-1}\right)^{*}$, then we obtain the cotangent bundle $T^{*} M$.
\end{EXA}



Proposition 6.34. 


\begin{PROP}{BÜN-L09-03-22}{Konstruktion eines globalen Schnitts im Vektorbündel}
Let $\pi: E \rightarrow M$ be an $\mathbb{F}$-vector bundle with typical fiber V and with VB-atlas $\left\{\left(U_{\alpha}, \phi_{\alpha}\right)\right\}$. Let $s_{\alpha}: U_{\alpha} \rightarrow \mathrm{V}$ be a collection of maps such that whenever $U_{\alpha} \cap U_{\beta} \neq \emptyset$, we have $s_{\alpha}(p)=\Phi_{\alpha \beta}(p) s_{\beta}(p)$ for all $p \in U_{\alpha} \cap U_{\beta}$. Then there is a global section $s$ such that $\left.s\right|_{U_{\alpha}}=s_{\alpha}$ for all $\alpha$.
\end{PROP}

Proof. 

\begin{PROOF}{BÜN-L09-03-23}{P: Konstruktion eines globalen Schnitts im Vektorbündel}
Let $\phi_{\alpha}:\left.E\right|_{U_{\alpha}} \rightarrow U_{\alpha} \times \mathrm{V}$ be the trivializations that give rise to the cocycle $\left\{\Phi_{\alpha \beta}\right\}$. Let $\gamma_{\alpha}(p):=\left(p, s_{\alpha}(p)\right)$ for $p \in U_{\alpha}$ and let $\left.s\right|_{U_{\alpha}}:=\phi_{\alpha}^{-1} \circ \gamma_{\alpha}$.

This gives a well-defined section $s$ because for $x \in U_{\alpha} \cap U_{\beta}$ we have
$$
\begin{aligned}
\phi_{\alpha}^{-1} \circ \gamma_{\alpha}(p) & =\phi_{\alpha}^{-1}\left(p, s_{\alpha}(p)\right) \\
& =\phi_{\alpha}^{-1}\left(p, \Phi_{\alpha \beta}(p) s_{\beta}(p)\right) \\
& =\phi_{\alpha}^{-1} \circ \phi_{\alpha} \circ \phi_{\beta}^{-1}\left(p, s_{\beta}(p)\right) \\
& =\phi_{\beta}^{-1}\left(p, s_{\beta}(p)\right)=\phi_{\beta}^{-1} \circ \gamma_{\beta}(p) .
\end{aligned}
$$
\end{PROOF}



\begin{EXA}{BÜN-L09-03-24}{Whitney Summen Bündel Konstruktion}
Suppose we have two vector bundles, $\pi_{1}: E_{1} \rightarrow M$ and $\pi_{2}: E_{2} \rightarrow M$. We give two constructions of the Whitney sum bundle $\pi_{1} \oplus \pi_{2}: E_{1} \oplus E_{2} \rightarrow M$. This is a globalization of the direct sum construction of vector spaces. In fact, the first construction simply takes $E_{1} \oplus E_{2}=\bigcup_{p \in M} E_{1 p} \oplus E_{2 p}$. Now, we have a vector bundle atlas $\left\{\left(\phi_{\alpha}, U_{\alpha}\right)\right\}$ for $\pi_{1}$ and a vector bundle atlas $\left\{\left(\psi_{\alpha}, U_{\alpha}\right)\right\}$ for $\pi_{2}$. Assume that both atlases have the same family of open sets (we can arrange this by taking a common refinement). Now let $\phi_{\alpha} \oplus \psi_{\alpha}$ : $\left(v_{p}, w_{p}\right) \mapsto\left(p, \operatorname{pr}_{2} \circ \phi_{\alpha}\left(v_{p}\right), \operatorname{pr}_{2} \circ \psi_{\alpha}\left(w_{p}\right)\right)$ for all $\left.\left(v_{p}, w_{p}\right) \in\left(E_{1} \oplus E_{2}\right)\right|_{U_{\alpha}}$. Then $\left\{\left(\phi_{\alpha} \oplus \psi_{\alpha}, U_{\alpha}\right)\right\}$ is a VB-atlas for $\pi_{1} \oplus \pi_{2}: E_{1} \oplus E_{2} \rightarrow M$.

Another method of constructing this bundle is to take the cocycle $\left\{g_{\alpha \beta}\right\}$ for $\pi_{1}$ and the cocycle $\left\{h_{\alpha \beta}\right\}$ for $\pi_{2}$ and then let $g_{\alpha \beta} \oplus h_{\alpha \beta}: U_{\alpha} \cap U_{\beta} \rightarrow \mathrm{GL}\left(\mathbb{F}^{k_{1}} \times \mathbb{F}^{k_{2}}\right)$ be defined by $\left(g_{\alpha \beta} \oplus h_{\alpha \beta}\right)(x)=g_{\alpha \beta}(x) \oplus h_{\alpha \beta}(x):(v, w) \mapsto \left(g_{\alpha \beta}(x) v, h_{\alpha \beta}(x) w\right)$. The maps $g_{\alpha \beta} \oplus h_{\alpha \beta}$ form a cocycle which determines a bundle by the construction of Proposition 6.30, which is (isomorphic to) $\pi_{1} \oplus \pi_{2}: E_{1} \oplus E_{2} \rightarrow M$.
\end{EXA}

\begin{DEF}{BÜN-L09-03-25}{Pullback Vektorbündel}
The pull-back of a vector bundle $\pi: E \rightarrow M$ by a smooth map $f: N \rightarrow M$ is naturally a vector bundle whose linear structure on each fiber $\left(f^{*} E\right)_{q}=\{q\} \times E_{p}$ is the obvious one induced from $E_{p}$. Put another way, we give the unique linear structure to each fiber that makes the bundle map $\tilde{f}: f^{*} E \rightarrow E$ linear on fibers. When given this vector bundle structure, we call $f^{*} E$ the pull-back vector bundle.
\end{DEF}

Example 6.35. 


\begin{EXA}{BÜN-L09-03-26}{Spezialfall Whitney Summen Bündel}
Let $\pi_{1}: E_{1} \rightarrow M$ and $\pi_{2}: E_{2} \rightarrow M$ be vector bundles and let $\triangle: M \rightarrow M \times M$ be the diagonal map $x \mapsto(x, x)$. From $\pi_{1}$ and $\pi_{2}$ one can construct a bundle $\pi_{E_{1} \times E_{2}}: E_{1} \times E_{2} \rightarrow M \times M$ by $\pi_{E_{1} \times E_{2}}\left(\epsilon_{1}, \epsilon_{2}\right):= \left(\pi_{1}\left(\epsilon_{1}\right), \pi_{2}\left(\epsilon_{2}\right)\right)$. The Whitney sum bundle $E_{1} \oplus E_{2}$ defined previously is naturally isomorphic to the pull-back $\triangle^{*} \pi_{E_{1} \times E_{2}}: \triangle^{*}\left(E_{1} \times E_{2}\right) \rightarrow M$ (Problem 10).
\end{EXA}

Exercise 6.36. 

\begin{EXE}{BÜN-L09-03-27}{$C^\infty(N;\mathbb{F})$-Modul durch Pullback}
Recall the space $\Gamma_{f}(\xi)$ from Definition 6.19. Show that if $\pi: E \rightarrow M$ is an $\mathbb{F}$-vector bundle and $f: N \rightarrow M$ is a smooth map, then both $\Gamma_{f}(\xi)$ and $\Gamma\left(f^{*} \xi\right)$ are modules over $C^{\infty}(N ; \mathbb{F})$, and that the natural correspondence between $\Gamma_{f}(\xi)$ and $\Gamma\left(f^{*} \xi\right)$ is a module isomorphism.
\end{EXE}

\begin{EXA}{BÜN-L09-03-28}{Nullschnitt in Vektorbündeln}
Every vector bundle has global sections. An obvious example is the zero section which maps each $x \in M$ to the zero element $0_{x}$ of the fiber $E_{x}$. The image of the zero section is also referred to as the zero section and is often identified with $M$. (Of course, the image of any global section is a submanifold diffeomorphic to the base manifold.)
\end{EXA}

We have the following simple analogue of Lemma 2.65:

Lemma 6.37. 

\begin{LEM}{BÜN-L09-03-29}{Existenz von globalen Schnitten zu lokalen Daten}
Let $\pi: E \rightarrow M$ be an $\mathbb{F}$-vector bundle with typical fiber V . If $v \in \pi^{-1}(p)$ then there exists a global section $\sigma \in \Gamma(\xi)$ such that $\sigma(p)=v$. Furthermore, if $s$ is a local section defined on $U$, and $V$ is an open set with compact closure with $V \subset \bar{V} \subset U$, then there is a section $\sigma \in \Gamma(\xi)$ such that $\sigma=s$ on $V$.
\end{LEM}

Proof. 

\begin{PROOF}{BÜN-L09-03-30}{P: Existenz von globalen Schnitten zu lokalen Daten}
Using a local trivialization one can easily get a local section $\sigma_{\text {loc }}$ defined near $p$ such that $\sigma(p)=v$. Now just use a cut-off function as in the proof of Lemma 2.65. For the second part we just choose a cut-off function $\beta$ with support in $U$ and such that $\beta=1$ on $\bar{V}$. Then $\beta s$ extends by zero to the desired global section.
\end{PROOF}

If a section of a vector bundle takes the zero value in some fiber we say that it vanishes at that point. Global smooth sections that never vanish do not always exist; such sections are called nowhere vanishing or nonvanishing. However, there is one case where it is easy to see that nonvanishing smooth sections exist:

Proposition 6.38. 


\begin{PROP}{BÜN-L09-03-31}{Globale nichtverschwindende Schnitte existieren auf zum trivialen VB äquivalenten Vektorbündeln}
Any bundle equivalent to a (trivial) product bundle must have a nowhere vanishing smooth global section.
\end{PROP}

It is a fact that the tangent bundle of $S^{2}$ does not have any such nowhere vanishing smooth sections. In other words, all smooth (or even continuous) vector fields on $S^{2}$ must vanish at some point. This is a result from algebraic topology called the "hairy sphere theorem" (see Theorem 10.15). If one fancifully imagines a vector field on a sphere to be hair, then the theorem suggests that one cannot comb the hair neatly "flat" without creating a cowlick somewhere. More generally, the analogous result holds for $S^{2 n}$ if $n \geq 1$.

Exercise 6.39. 

\begin{EXE}{BÜN-L09-02-53}{Nicht Existenz von nichtverschwindenden Schnitten im Böbiusband}
Modify either the construction of Example 6.9 or Example 6.16 to obtain a rank one vector bundle version of the Möbius band and give an argument proving that every global continuous section of this bundle must vanish somewhere.
\end{EXE}

Definition 6.40. 



\begin{DEF}{BÜN-L09-03-33}{Rahmen bei Basispunkt in Vektorbündel}
If $\xi=(E, \pi, M, \mathrm{V})$ is a vector bundle and $p \in M$, then a vector space basis for the fiber $E_{p}$ is called a frame at $p$.
\end{DEF}


Definition 6.41. 

\begin{DEF}{BÜN-L09-03-32}{(Lokales) Rahmenfeld von Vektorbündeln}
Let $\pi: E \rightarrow M$ be a rank $k$ vector bundle. A $k$-tuple $\sigma=\left(\sigma_{1}, \ldots, \sigma_{k}\right)$ of sections of $E$ over an open set $U$ is called a (local) frame field over $U$ if for all $p \in U,\left(\sigma_{1}(p), \ldots, \sigma_{k}(p)\right)$ is a frame at $p$.
\end{DEF}

If we choose a fixed basis $\left\{\mathbf{e}_{i}\right\}_{i=1, \ldots, k}$ for the typical fiber V , then a choice of a local frame field over an open set $U \subset M$ is equivalent to a local trivialization (a vector bundle chart). Namely, if $\phi$ is such a trivialization over $U$, then defining $\sigma_{i}(p)=\phi^{-1}\left(p, \mathbf{e}_{i}\right)$, we have that $\sigma_{\phi}=\left(\sigma_{1}, \ldots, \sigma_{k}\right)$ is a local frame over $U$. Conversely, if $\sigma_{\phi}=\left(\sigma_{1}, \ldots, \sigma_{k}\right)$ is a local frame over $U$, then every $v \in \pi^{-1}(U)$ has the form $v=\sum v^{i} \sigma_{i}(p)$ for a unique $p$ and unique numbers $v^{i}(p)$. Then the map $f: U \times \mathrm{V} \rightarrow \pi^{-1}(U)$ defined by $(p, v) \mapsto \sum v^{i} \sigma_{i}(p)$ is a diffeomorphism and its inverse $\phi=f^{-1}$ is a trivialization. Thus if there is a global frame field, then the vector bundle is trivial. 



\begin{PROP}{BÜN-L09-02-58}{Lokaler Rahmen $\Longleftrightarrow$ lokale Trivialisierung}
Sei $\pi:E\to M$ ein reelles (oder komplexes) Vektorbündel vom Rang $k$ mit typischer Faser $V$ und fixierter Basis $(\mathbf e_1,\dots,\mathbf e_k)$ von $V$. Für eine offene Menge $U\subset M$ sind folgende Daten äquivalent:
\begin{enumerate}
\item eine lokale Trivialisierung (Vektorbündelkarte) $\;\phi:\pi^{-1}(U)\xrightarrow{\;\cong\;} U\times V$,
\item ein lokaler Rahmen $\;\sigma=(\sigma_1,\dots,\sigma_k)$ von Schnitten $\sigma_i\in\Gamma(U,E)$, die punktweise eine Basis von $E_p$ bilden.
\end{enumerate}
Insbesondere existiert ein globaler Rahmen auf $M$ genau dann, wenn $E$ trivial ist, d.\,h.\ $E\cong M\times V$.
\end{PROP}



\begin{PROOF}{BÜN-L09-03-37}{P: Lokaler Rahmen $\Longleftrightarrow$ lokale Trivialisierung}
(1) $\Rightarrow$ (2): Sei $\phi$ wie in (1). Definiere für $p\in U$
\[
\sigma_i(p):=\phi^{-1}\bigl(p,\mathbf e_i\bigr)\in E_p.
\]
Dann ist $\sigma=(\sigma_1,\dots,\sigma_k)$ ein lokaler Rahmen über $U$, da $\phi_p:=\phi|_{E_p}:E_p\to\{p\}\times V\cong V$ linearer Isomorphismus ist und $(\mathbf e_i)$ Basis von $V$.

(2) $\Rightarrow$ (1): Sei umgekehrt $\sigma=(\sigma_1,\dots,\sigma_k)$ ein lokaler Rahmen über $U$. Für $p\in U$ und $v=\sum_{i=1}^k v^i\,\mathbf e_i\in V$ setze
\[
f:U\times V\longrightarrow \pi^{-1}(U),\qquad (p,v)\longmapsto \sum_{i=1}^k v^i\,\sigma_i(p).
\]
Die Abbildung $f$ ist glatt und für jedes $p$ ist $f_p:V\to E_p$ linearer Isomorphismus (weil $(\sigma_i(p))$ Basis von $E_p$ ist). Damit ist $f$ ein Vektorbündel-Isomorphismus über $U$. Der Inversenpfeil $\phi:=f^{-1}:\pi^{-1}(U)\to U\times V$ ist glatt; explizit gilt für $w\in E_p$ die eindeutige Darstellung $w=\sum_i a^i(p)\sigma_i(p)$ und
\[
\phi(w)=\bigl(p,\sum_i a^i(p)\,\mathbf e_i\bigr),
\]
wobei die Koeffizienten $a^i(p)$ glatt in $(p,w)$ variieren (z.\,B. mittels des dualen Rahmens in $E^*$). Somit ist $\phi$ eine lokale Trivialisierung.

Der letzte Satz folgt, indem man $U=M$ nimmt: Ein globaler Rahmen liefert obige Konstruktion global und damit einen Bündel-Isomorphismus $E\cong M\times V$; umgekehrt induziert jede globale Trivialisierung den kanonischen globalen Rahmen $\sigma_i(p)=\phi^{-1}(p,\mathbf e_i)$.
\end{PROOF}





\begin{DEF}{BÜN-L09-03-34}{Parallelisierbare Mannigfaltigkeiten}
A manifold $M$ is said to be parallelizable if $T M \rightarrow M$ has a global frame field; i.e. if the tangent bundle is trivial.
\end{DEF} 


\begin{EXA}{BÜN-L09-03-35}{Beispiele für Parallelisierbare Mannigfaltigkeiten}
For example, the hairy sphere theorem mentioned above implies that $S^{2}$ is not parallelizable. On the other hand, $S^{2} \times \mathbb{R}$ is parallelizable. It is easy to show that the torus $T^{2}=S^{1} \times S^{1}$ and its higher-dimensional analogues $T^{n}=S^{1} \times \cdots \times S^{1}$ are also parallelizable.
\end{EXA}





\begin{PROP}{BÜN-L09-03-36}{Lokaler Rahmen $\Longleftrightarrow$ lokale Trivialisierung}
Sei $\pi:E\to M$ ein reelles (oder komplexes) Vektorbündel vom Rang $k$ mit typischer Faser $V$ und fixierter Basis $(\mathbf e_1,\dots,\mathbf e_k)$ von $V$. Für eine offene Menge $U\subset M$ sind folgende Daten äquivalent:
\begin{enumerate}
\item eine lokale Trivialisierung (Vektorbündelkarte) $\;\phi:\pi^{-1}(U)\xrightarrow{\;\cong\;} U\times V$,
\item ein lokaler Rahmen $\;\sigma=(\sigma_1,\dots,\sigma_k)$ von Schnitten $\sigma_i\in\Gamma(U,E)$, die punktweise eine Basis von $E_p$ bilden.
\end{enumerate}
Insbesondere existiert ein globaler Rahmen auf $M$ genau dann, wenn $E$ trivial ist, d.\,h.\ $E\cong M\times V$.
\end{PROP}



\begin{PROOF}{BÜN-L09-03-37}{P: Lokaler Rahmen $\Longleftrightarrow$ lokale Trivialisierung}
(1) $\Rightarrow$ (2): Sei $\phi$ wie in (1). Definiere für $p\in U$
\[
\sigma_i(p):=\phi^{-1}\bigl(p,\mathbf e_i\bigr)\in E_p.
\]
Dann ist $\sigma=(\sigma_1,\dots,\sigma_k)$ ein lokaler Rahmen über $U$, da $\phi_p:=\phi|_{E_p}:E_p\to\{p\}\times V\cong V$ linearer Isomorphismus ist und $(\mathbf e_i)$ Basis von $V$.

(2) $\Rightarrow$ (1): Sei umgekehrt $\sigma=(\sigma_1,\dots,\sigma_k)$ ein lokaler Rahmen über $U$. Für $p\in U$ und $v=\sum_{i=1}^k v^i\,\mathbf e_i\in V$ setze
\[
f:U\times V\longrightarrow \pi^{-1}(U),\qquad (p,v)\longmapsto \sum_{i=1}^k v^i\,\sigma_i(p).
\]
Die Abbildung $f$ ist glatt und für jedes $p$ ist $f_p:V\to E_p$ linearer Isomorphismus (weil $(\sigma_i(p))$ Basis von $E_p$ ist). Damit ist $f$ ein Vektorbündel-Isomorphismus über $U$. Der Inversenpfeil $\phi:=f^{-1}:\pi^{-1}(U)\to U\times V$ ist glatt; explizit gilt für $w\in E_p$ die eindeutige Darstellung $w=\sum_i a^i(p)\sigma_i(p)$ und
\[
\phi(w)=\bigl(p,\sum_i a^i(p)\,\mathbf e_i\bigr),
\]
wobei die Koeffizienten $a^i(p)$ glatt in $(p,w)$ variieren (z.\,B. mittels des dualen Rahmens in $E^*$). Somit ist $\phi$ eine lokale Trivialisierung.

Der letzte Satz folgt, indem man $U=M$ nimmt: Ein globaler Rahmen liefert obige Konstruktion global und damit einen Bündel-Isomorphismus $E\cong M\times V$; umgekehrt induziert jede globale Trivialisierung den kanonischen globalen Rahmen $\sigma_i(p)=\phi^{-1}(p,\mathbf e_i)$.
\end{PROOF}





\begin{DEF}{BÜN-L09-03-38}{Reduktion der Strukturgruppe eines Vektorbündels}
Sei $\pi:E\to M$ ein glattes Vektorbündel vom Rang $k$ mit typischer Faser $V$
und Strukturgruppe $\mathrm{GL}(V)$, die in der Standarddarstellung
\[
\lambda_0:\mathrm{GL}(V)\longrightarrow \mathrm{GL}(V),\qquad
\lambda_0(A)(v)=Av
\]
auf $V$ wirkt.
\begin{enumerate}
\item Sei $G\subset\mathrm{GL}(V)$ eine Lie-Untergruppe und
$\lambda:=\lambda_0|_G$ die Einschränkung der Standarddarstellung.
Eine \emph{$(G,\lambda)$-Bündelstruktur} auf $E$
heißt eine \emph{Reduktion der Strukturgruppe} von $\mathrm{GL}(V)$ auf $G$,
wenn die Übergangsfunktionen eines Vektorbündelatlas
\[
g_{\alpha\beta}:U_\alpha\cap U_\beta\longrightarrow G
\]
Werte in $G$ annehmen und die Verträglichkeitsbedingungen
$g_{\alpha\alpha}=e$, $g_{\alpha\beta}=g_{\beta\alpha}^{-1}$,
$g_{\alpha\beta}g_{\beta\gamma}g_{\gamma\alpha}=e$ erfüllen.
\item Entsprechend gilt für die zugehörigen Bündeltrivialisierungen
$(U_\alpha,\phi_\alpha)$, $(U_\beta,\phi_\beta)$:
\[
\phi_\alpha\circ\phi_\beta^{-1}(p,v)
=
(p,\lambda(g_{\alpha\beta}(p))(v)),
\qquad p\in U_\alpha\cap U_\beta,\ v\in V.
\]
\item Seien $(\sigma_i)_i$ und $(\sigma_i')_i$ die lokalen Rahmenfelder über
$U_\alpha$ bzw. $U_\beta$, definiert durch
$\sigma_i(p)=\phi_\alpha^{-1}(p,\mathbf{e}_i)$ und
$\sigma_i'(p)=\phi_\beta^{-1}(p,\mathbf{e}_i)$.
Dann gilt auf $U_\alpha\cap U_\beta$:
\[
\sigma_i'(p)=\sum_j \lambda_i^{\,j}\!\big(g_{\alpha\beta}(p)\big)\,\sigma_j(p),
\]
wobei $(\lambda_i^{\,j}(g))$ die Matrixdarstellung von $\lambda(g)$
bezüglich der Basis $(\mathbf{e}_1,\ldots,\mathbf{e}_k)$ ist.
Die Matrixfunktion $(\lambda_i^{\,j}\circ g_{\alpha\beta})$ beschreibt
damit den Wechsel des lokalen Rahmens.
\end{enumerate}
\end{DEF}







\begin{CONC}{BÜN-L09-03-39}{Kam: Reduktion der Strukturgruppe eines Vektorbündels}
If $G$ is a Lie subgroup of $\mathrm{GL}(\mathrm{V})$ and $G$ acts in the standard way on V , that is, if the action is $\left.\lambda_{0}\right|_{G}$, the restriction of the standard action, then a $\left.\lambda_{0}\right|_{G}$-bundle structure is a reduction to the group $G$. Put another way, one has achieved such a reduction if one can find a cocycle of standard transition maps $\left\{\Phi_{\alpha \beta}\right\}$ arising from a vector bundle atlas which take values in $G$ (acting in the standard way on V ). By a slight extension, if $\lambda$ is a linear action on V , then an effective ( $G, \lambda$ )-bundle structure on $E$ can be considered as a reduction of the standard $\mathrm{GL}(\mathrm{V})$-bundle structure.

Let $\sigma_{i}(p)=\phi_{\alpha}^{-1}\left(p, \mathbf{e}_{i}\right)$ and $\sigma_{i}^{\prime}(p)=\phi_{\beta}^{-1}\left(p, \mathbf{e}_{i}\right)$ be frame fields corresponding to VB-charts ( $U_{\alpha}, \phi_{\alpha}$ ) and ( $U_{\beta}, \phi_{\beta}$ ) that lie in a ( $G, \lambda$ )-atlas and where $U_{\alpha} \cap U_{\beta}$ is nonempty. Then $\phi_{\alpha} \circ \phi_{\beta}^{-1}(p, v)=\left(p, \lambda\left(g_{\alpha \beta}(p)\right)(v)\right)$ for transition functions $g_{\alpha \beta}: U_{\alpha} \cap U_{\beta} \rightarrow G$. For each $g \in G$, there is a matrix $\left(\lambda_{i}^{j}(g)\right)$ which represents $\lambda_{g}$ with respect to ( $\mathbf{e}_{1}, \ldots, \mathbf{e}_{k}$ ). Unwinding definitions, we see that $\phi_{\alpha}^{-1}$ is linear on the vector space $\{p\} \times \mathrm{V}$. Using this we have
$$
\begin{aligned}
\sigma_{i}^{\prime}(p) & =\phi_{\beta}^{-1}\left(p, \mathbf{e}_{i}\right)=\phi_{\alpha}^{-1}\left(p, \lambda\left(g_{\alpha \beta}(p)\right)\left(\mathbf{e}_{i}\right)\right) \\
& =\phi_{\alpha}^{-1}\left(p, \sum_{j} \lambda_{i}^{j}\left(g_{\alpha \beta}(p)\right) \mathbf{e}_{j}\right)=\sum_{j} \lambda_{i}^{j}\left(g_{\alpha \beta}(p)\right) \phi_{\alpha}^{-1}\left(p, \mathbf{e}_{j}\right) \\
& =\sum_{j} \lambda_{i}^{j}\left(g_{\alpha \beta}(p)\right) \sigma_{j}(p)
\end{aligned}
$$
Thus the smooth matrix-valued function $\left(\lambda_{i}^{j} \circ g_{\alpha \beta}\right)$ defined on $U_{\alpha} \cap U_{\beta}$ gives the change of frame and embodies the transition map on $U_{\alpha} \cap U_{\beta}$. In practice, this is often a good way to look at things. Consider the common situation where $\mathrm{V}=\mathbb{R}^{k}$ and where $\lambda_{0}$ is the standard representation of $\mathrm{GL}\left(\mathbb{R}^{k}\right)$. If we identify $\operatorname{GL}\left(\mathbb{R}^{k}\right)$ with the matrix group $\operatorname{GL}(k)$, then $\left(\left(\lambda_{0}\right)_{i}^{j} \circ g_{\alpha \beta}(p)\right)$ is just the matrix $g_{\alpha \beta}(p)$ itself.
\end{CONC}

Metric differential geometry begins if we have a scalar product on the fibers of a vector bundle. We introduce the concept at this point so as to have an example of a reduction of the structure group.



Definition 6.42. 

\begin{CONC}{BÜN-L09-03-40}{Riemannisches Vektorbündel}
A Riemannian metric on a real vector bundle $\pi: E \rightarrow M$ is a map $p \mapsto g_{p}(\cdot, \cdot)$ which assigns to each $p \in M$ a positive definite scalar product $g_{p}(\cdot, \cdot)$ on the fiber $E_{p}$ that is smooth in the sense that $p \mapsto g_{p}\left(s_{1}(p), s_{2}(p)\right)$ is smooth for all smooth sections $s_{1}$ and $s_{2}$. A real vector bundle together with a Riemannian metric is referred to as a Riemannian vector bundle.
\end{CONC}

For example, a Riemannian metric on the tangent bundle of a smooth manifold is what one means by a Riemannian metric on the manifold. A smooth manifold with a Riemannian metric is called a Riemannian manifold and such will be studied later in this book. 



\begin{PROP}{BÜN-L09-03-41}{Wahl Riemannsche Metrik auf Vektorbündel $\Longleftrightarrow$ $\mathrm{O}(V)$-Reduktion}
Sei $\pi:E\to M$ ein reelles Vektorbündel vom Rang $k$ mit typischer Faser $V$.
Es sei auf $V$ ein festes Skalarprodukt $\langle\cdot,\cdot\rangle_V$ und eine ausgezeichnete
ON-Basis $(\mathbf e_1,\ldots,\mathbf e_k)$ gewählt. Dann gilt:

\begin{enumerate}
\item[(i)] Jede Riemannsche Bündelmetrik $g$ auf $E$ induziert eine Reduktion der
Strukturgruppe von $\mathrm{GL}(V)$ auf die orthogonale Gruppe
$\mathrm{O}(V)=\{A\in\mathrm{GL}(V)\mid \langle Av,Aw\rangle_V=\langle v,w\rangle_V\}$.
\item[(ii)] Umgekehrt induziert jede Reduktion der Strukturgruppe auf $\mathrm{O}(V)$
(eine $(\mathrm{O}(V),\lambda_0|_{\mathrm{O}(V)})$-Bündelstruktur in Standarddarstellung)
eindeutig eine Riemannsche Bündelmetrik $g$ auf $E$.
\end{enumerate}
\end{PROP}



\begin{PROOF}{BÜN-L09-03-42}{P: Wahl Riemannsche Metrik auf Vektorbündel $\Longleftrightarrow$ $\mathrm{O}(V)$-Reduktion}
\textbf{Zu (i).} Wähle einen beliebigen VB-Atlas $\{(U_\alpha,\phi_\alpha)\}$ mit
$\phi_\alpha:\pi^{-1}(U_\alpha)\to U_\alpha\times V$ linear auf den Fasern und
definiere die lokalen Rahmen
$\sigma_i^\alpha(p):=\phi_\alpha^{-1}(p,\mathbf e_i)$.
Für die gegebene Metrik $g$ auf $E$ wende auf die Basis
$(\sigma_1^\alpha(p),\ldots,\sigma_k^\alpha(p))$ punktweise den
Gram–Schmidt-Prozess an. Da $g$ glatt ist und die Ausgangsrahmen glatt von $p$
abhängen, erhält man einen \emph{glatten} ON-Rahmen
$e_j^\alpha(p)=\sum_i A_j^{\,i}(p)\,\sigma_i^\alpha(p)$ auf $U_\alpha$, wobei
$A(p)\in\mathrm{GL}(k)$ glatt in $p$ ist.

Ersetze nun $\phi_\alpha$ durch die Trivialisierung
$\phi'_\alpha$ mit
$\phi_\alpha^{\prime-1}(p,v)=\sum_i v^i\,e_i^\alpha(p)$
(\,$v=\sum_i v^i\mathbf e_i\in V$\,). Für zwei Karten $(U_\alpha,\phi'_\alpha)$
und $(U_\beta,\phi'_\beta)$ gilt auf $U_\alpha\cap U_\beta$ ein Rahmenwechsel
$e_j^\alpha(p)=\sum_i Q_j^{\,i}(p)\,e_i^\beta(p)$ mit einer glatten Matrixfunktion
$Q(p)\in \mathrm{O}(V)$, da beide Rahmen $g$-orthonormal sind.
Somit haben die Übergangsfunktionen der neuen Atlasfamilie Werte in $\mathrm{O}(V)$,
und wir erhalten eine Reduktion der Strukturgruppe auf $\mathrm{O}(V)$.

\medskip
\textbf{Zu (ii).} Sei umgekehrt ein $(\mathrm{O}(V),\lambda_0|_{\mathrm{O}(V)})$-Atlas
$\{(U_\alpha,\phi_\alpha)\}$ mit Übergangsfunktionen
$g_{\alpha\beta}:U_\alpha\cap U_\beta\to \mathrm{O}(V)$ gegeben.
Definiere auf jeder Faser $E_p$ eine Bilinearform $g_p$ durch:
Wähle $\alpha$ mit $p\in U_\alpha$ und setze für $u,v\in E_p$
\[
g_p(u,v):=\big\langle \Phi_\alpha(u),\,\Phi_\alpha(v)\big\rangle_V,
\quad\text{wobei }\ \phi_\alpha=(\pi,\Phi_\alpha).
\]
Dies ist wohldefiniert: Wählt man statt $\alpha$ eine Karte $\beta$, so gilt
$\Phi_\beta=\lambda_0(g_{\beta\alpha})\circ \Phi_\alpha$ mit
$g_{\beta\alpha}(p)\in \mathrm{O}(V)$, also
\[
\big\langle \Phi_\beta(u),\Phi_\beta(v)\big\rangle_V
=
\big\langle \lambda_0(g_{\beta\alpha})\Phi_\alpha(u),\lambda_0(g_{\beta\alpha})\Phi_\alpha(v)\big\rangle_V
=
\big\langle \Phi_\alpha(u),\Phi_\alpha(v)\big\rangle_V,
\]
da $g_{\beta\alpha}(p)$ das Skalarprodukt auf $V$ erhält.
Auf jedem $U_\alpha$ ist $g$ glatt, und die Glattheit verklebt über den
Überschneidungen, da die Übergänge orthogonal sind. Somit erhält man eine
glatte, positiv definite Bündelmetrik $g$ auf $E$.

Die Konstruktionen in (i) und (ii) sind zueinander invers (bis auf die
vorab fixierte Wahl von $\langle\cdot,\cdot\rangle_V$ und der ON-Basis von $V$).
\end{PROOF}







\begin{CONC}{BÜN-L09-03-43}{KOM: Wahl Riemannsche Metrik auf Vektorbündel $\Longleftrightarrow$ $\mathrm{O}(V)$-Reduktion}
If a rank $k$ real vector bundle $\pi: E \rightarrow M$ has a Riemannian metric, then it is convenient to assume that a fixed inner product is chosen on the typical fiber V and that a distinguished orthonormal basis $\left(\mathbf{e}_{1}, \ldots, \mathbf{e}_{k}\right)$ has been chosen. In most applications, V is $\mathbb{R}^{k}$, the inner product is the standard dot product, and the distinguished basis is the usual standard basis. Recall that the orthogonal group $\mathrm{O}(\mathrm{V})$ is the subgroup of $\mathrm{GL}(\mathrm{V})$ consisting of elements that preserve the inner product. Once the inner product and distinguished orthonormal basis are fixed, every choice of Riemannian metric on $E$ corresponds to a reduction of the standard structure group $\mathrm{GL}(\mathrm{V})$ to the subgroup $\mathrm{O}(\mathrm{V})$ as follows. 


Let us first show how a metric leads to a reduction. We start with an arbitrary VB-atlas $\left\{\left(U_{\alpha}, \phi_{\alpha}\right)\right\}$. Since we have fixed a basis for V, each chart ( $U_{\alpha}, \phi_{\alpha}$ ) defines a frame field ( $\sigma_{1}^{\alpha}, \ldots, \sigma_{k}^{\alpha}$ ) on $U_{\alpha}$ by $\sigma_{i}^{\alpha}(p):=\phi_{\alpha}^{-1}\left(p, \mathbf{e}_{i}\right)$ as explained above. One can perform a Gram-Schmidt process on the basis $\left(\sigma_{1}^{\alpha}(p), \ldots, \sigma_{k}^{\alpha}(p)\right)$ simultaneously for all $p \in U_{\alpha}$ so that we have a new orthonormal basis $\left(e_{1}^{\alpha}(p), \ldots, e_{k}^{\alpha}(p)\right)$ for each $p$, where $e_{j}^{\alpha}(p)=\sum A_{j}^{i}(p) \sigma_{i}^{\alpha}(p)$ for all $p \in U_{\alpha}$ and the matrix entries $A_{j}^{i}(p)$ depend smoothly on $p$. Thus $\left(e_{1}^{\alpha}, \ldots, e_{k}^{\alpha}\right)$ is a (smooth) local frame field called an orthonormal frame field. One then replaces the original chart ( $U_{\alpha}, \phi_{\alpha}$ ) by a new chart ( $U_{\alpha}, \phi_{\alpha}^{\prime}$ ) that is the inverse of the map $(p, v) \mapsto \sum v^{i} e_{i}(p)$, where $v=\sum v^{i} \mathbf{e}_{i}(p)$ in V . In other words, $\phi_{\alpha}^{-1}:(p, v) \mapsto \sum v^{i} \sigma_{i}(p)$ is replaced by
$$
\phi_{\alpha}^{\prime-1}:(p, v) \mapsto \sum v^{i} e_{i}(p)
$$
Make this replacement for each ( $U_{\alpha}, \phi_{\alpha}$ ) to obtain a new atlas $\left\{\left(U_{\alpha}, \phi_{\alpha}^{\prime}\right)\right\}$. Any two of these orthonormal frame fields, say $\left(e_{1}^{\alpha}, \ldots, e_{k}^{\alpha}\right)$ and $\left(e_{1}^{\beta}, \ldots e_{k}^{\beta}\right)$,\\
are related by
$$
e_{j}^{\alpha}(p)=\sum Q_{j}^{i}(p) e_{i}^{\beta}(p)
$$
for some smooth orthogonal matrix function $Q_{j}^{i}$. One now checks that the transition maps for this new atlas (which is still a subatlas for the maximal VB-atlas) take values in $\mathrm{O}(\mathrm{V})$. Indeed,
$$
\begin{aligned}
\left(p, \Phi_{\alpha \beta}^{\prime}(p)(v)\right) & =\phi_{\beta}^{\prime} \circ \phi_{\alpha}^{\prime-1}(p, v)=\phi_{\beta}^{\prime}\left(\sum v^{j} e_{j}^{\alpha}(p)\right) \\
& =\phi_{\beta}^{\prime}\left(\sum v^{j} Q_{j}^{i}(p) e_{i}^{\beta}(p)\right)=\left(p,\left.\Phi_{\beta}^{\prime}\right|_{p} \sum v^{j} Q_{j}^{i}(p) e_{i}^{\beta}(p)\right) \\
& =\left(p,\left.\sum v^{j} Q_{j}^{i}(p) \Phi_{\beta}^{\prime}\right|_{p} e_{i}^{\beta}(p)\right)=\left(p, \sum v^{j} Q_{j}^{i}(p) \mathbf{e}_{i}\right),
\end{aligned}
$$
from which we see that $\Phi_{\alpha \beta}^{\prime}(p)(v)=\Phi_{\alpha \beta}^{\prime}(p)\left(\sum v^{j} \mathbf{e}_{j}\right)=\sum v^{j} Q_{j}^{i}(p) \mathbf{e}_{i}$. Since $\left(Q_{j}^{i}(p)\right)$ is an orthogonal matrix for all $p$, we have $\Phi_{\alpha \beta}^{\prime}(p) \in \mathrm{O}(\mathrm{V})$ for all $p$. The converse is also true. Namely, a reduction to the structure group $\mathrm{O}(\mathrm{V})$ (acting in the standard way) is tantamount to the introduction of a Riemannian metric. The correspondence presumes the prior choice of inner product and distinguished orthonormal basis on V.
\end{CONC}

Exercise 6.43. Prove the converse statement referred to above.


Exercise 6.44. 

\begin{EXE}{BÜN-L09-03-44}{Hermitische Metrik auf Vektorbündel}
Let $E$ be a complex vector bundle of rank $k$. Define by analogy with Riemannian metric, the notion of a Hermitian metric on $E$ and show that every Hermitian metric on $E$ corresponds to a reduction of the standard $\mathrm{GL}(k, \mathbb{C})$-bundle structure to a $\mathrm{U}(n)$-bundle structure.
\end{EXE}

Proposition 6.45. 


\begin{PROP}{BÜN-L09-03-45}{Vektorbündel erlauben riemannsche / hermitische Metriken}
On every real vector bundle $E$ there can be defined a Riemannian metric. Similarly, on any complex vector bundle there exists a Hermitian metric.
\end{PROP}

Proof. 


\begin{PROOF}{BÜN-L09-03-46}{P: Vektorbündel erlauben riemannsche / hermitische Metriken}
We prove the Riemannian case; the Hermitian case is entirely analogous. The proof uses the fact that a strict convex combination of positive definite scalar products is a positive definite scalar product. This allows us to use a partition of unity argument. Endow V with an inner product. Let $\left\{\left(U_{\alpha}, \phi_{\alpha}\right)\right\}$ be a VB-atlas and let ( $U_{\alpha}, \phi_{\alpha}$ ) be a given VB-chart. On the trivial bundle $U_{\alpha} \times \mathrm{V} \rightarrow U_{\alpha}$ there certainly exists a Riemannian metric given on each fiber by $\langle(p, v),(p, w)\rangle_{\alpha}=\langle v, w\rangle$. We may transfer this to the bundle $\pi^{-1}\left(U_{\alpha}\right) \rightarrow U_{\alpha}$ by using the map $\phi_{\alpha}^{-1}$, thus obtaining a metric $g_{\alpha}$ on this restricted bundle over $U_{\alpha}$. We do this for every VB-chart in the atlas. The trick is to piece these together in a smooth way. For that, we take a smooth partition of unity ( $U_{\alpha}, \rho_{\alpha}$ ) subordinate to the cover $\left\{U_{\alpha}\right\}$. Let
$$
\mathrm{g}(p)=\sum \rho_{\alpha}(p) g_{\alpha}(p)
$$
The sum is finite at each $p \in M$ since the partition of unity is locally finite and the functions $\rho_{\alpha} \mathrm{g}_{\alpha}$ are extended to be zero outside of the corresponding $U_{\alpha}$. The fact that $\rho_{\alpha} \geq 0$ and $\rho_{\alpha}>0$ at $p$ for at least one $\alpha$ easily gives the result that g is positive definite at each $p$ and so it is a Riemannian metric on $E$. \(\square\)
\end{PROOF}

Example 6.46 (Tautological line bundle). 

\begin{EXA}{BÜN-L09-03-47}{Tautologisches Linien Bündel}
Recall that $\mathbb{R} P^{n}$ is the set of all lines through the origin in $\mathbb{R}^{n+1}$. Define the subset $\mathbb{L}\left(\mathbb{R} P^{n}\right)$ of $\mathbb{R} P^{n} \times \mathbb{R}^{n+1}$ consisting of all pairs ( $l, v$ ) such that $v \in l$ (think about this). This set together with the map $\pi_{\mathbb{R} P^{n}}: \mathbb{L}\left(\mathbb{R} P^{n}\right) \rightarrow \mathbb{R} P^{n}$ given by $(l, v) \mapsto l$, is a rank one vector bundle.
\end{EXA}

Example 6.47 (Tautological bundle). 


\begin{EXA}{BÜN-L09-03-48}{Tautologisches Bündel}
Let $G(n, k)$ denote the Grassmann manifold of $k$-planes in $\mathbb{R}^{n}$. Let $\gamma_{n, k}$ be the subset of $G(n, k) \times \mathbb{R}^{n}$ consisting of pairs ( $P, v$ ) where $P$ is a $k$-plane ( $k$-dimensional subspace) and $v$ is a vector in the plane $P$. The projection $\pi_{n, k}: \gamma_{n, k} \rightarrow G(n, k)$ is simply $(P, v) \mapsto P$. The result is a vector bundle $\left(\gamma_{n, k}, \pi_{n, k}, G(n, k), \mathbb{R}^{k}\right)$. We leave it to the reader to discover an appropriate VB-atlas (see Problem 12).
\end{EXA}

\begin{CONC}{BÜN-L09-03-49}{KOM: Tautologisches Bündel}
These tautological vector bundles are not just trivial bundles, and in fact their topology or twistedness (for large $n$ ) is of the utmost importance for classifying vector bundles (see $[\mathbf{B o - T u}]$ ). One may take the inclusions $\mathbb{R}^{n} \subset \mathbb{R}^{n+1} \subset \cdots \subset \mathbb{R}^{\infty}$ to construct inclusions $G(n, k) \subset G(n+1, k) \subset \cdots$ and $\gamma_{n, k} \subset \gamma_{n+1, k}$. Given a rank $k$ vector bundle $\pi: E \rightarrow M$, there is an $n$ such that $\pi: E \rightarrow M$ is (isomorphic to) the pull-back of $\gamma_{n, k}$ by some map $f: M \rightarrow G(n, k)$ :\\
\includegraphics[scale=0.25, center]{2025_10_09_265d7e1a22c89f79c21eg-296}
\end{CONC}

Exercise 6.48. To each point on a unit sphere in $\mathbb{R}^{n}$, attach the space of all vectors normal to the sphere at that point. Show that this normal bundle is in fact a (smooth) vector bundle. Generalize to define the normal bundle of a hypersurface in $\mathbb{R}^{n}$. When is such a normal bundle trivial?

Exercise 6.49. Fix a nonnegative integer $j$. Let $Y=\mathbb{R} \times(-1,1)$ and let $\left(x_{1}, y_{1}\right) \sim\left(x_{2}, y_{2}\right)$ if and only if $x_{1}=x_{2}+j k$ and $y_{1}=(-1)^{j k} y_{2}$ for some integer $k$. Show that $E:=Y / \sim$ is a vector bundle of rank 1 that is trivial if and only if $j$ is even. Prove or at least convince yourself that this is the Möbius band when $j$ is odd.


\pagebreak



\subsection*{Tensor Products of Vector Bundles}
Given two vector bundles $\pi_{1}: E_{1} \rightarrow M$ and $\pi_{2}: E_{2} \rightarrow M$ with respective typical fibers $V_{1}$ and $V_{2}$, we let

$$
E_{1} \otimes E_{2}:=\bigcup_{p \in M} E_{1 p} \otimes E_{2 p} \quad \text { (a disjoint union) }
$$

Then we have a projection map $\pi: E_{1} \otimes E_{2} \rightarrow M$ given by mapping any element in a fiber $E_{1 p} \otimes E_{2 p}$ to the base point $p$. We show how to construct a VB-atlas for $E_{1} \otimes E_{2}$ from an atlas on each of $E_{1}$ and $E_{2}$. The smooth structure and topology can be derived from the atlas as usual in such a way as to make all the relevant maps smooth. We leave the verification of this to the reader. The resulting bundle is the tensor product bundle. As usual we can assume that the atlases are based on the same open cover. Thus suppose that $\left\{\left(U_{\alpha}, \phi_{\alpha}\right)\right\}$ is a VB-atlas for $E_{1}$ while $\left\{\left(U_{\alpha}, \psi_{\alpha}\right)\right\}$ is a VB-atlas for $E_{2}$. Now let $\Phi_{\alpha} \otimes \Psi_{\alpha}:\left.\left(E_{1} \otimes E_{2}\right)\right|_{U_{\alpha}} \rightarrow \mathrm{V}_{1} \otimes \mathrm{V}_{2}$ be defined by $\left.\left(\Phi_{\alpha} \otimes \Psi_{\alpha}\right)\right|_{E_{1 p} \otimes E_{2 p}}:=\left.\left.\Phi_{\alpha}\right|_{E_{1 p}} \otimes \Psi_{\alpha}\right|_{E_{2 p}}$ for $p \in U_{\alpha}$. Then let

$$
\phi_{\alpha} \otimes \psi_{\alpha}:\left.\left(E_{1} \otimes E_{2}\right)\right|_{U_{\alpha}} \rightarrow U_{\alpha} \times\left(\mathrm{V}_{1} \otimes \mathrm{V}_{2}\right)
$$

be defined by $\phi_{\alpha} \otimes \psi_{\alpha}:=\left(\pi, \Phi_{\alpha} \otimes \Psi_{\alpha}\right)$. To clarify, the map $\left.\left.\Phi_{\alpha}\right|_{E_{1 p}} \otimes \Psi_{\alpha}\right|_{E_{2 p}}$ : $E_{1 p} \otimes E_{2 p} \rightarrow \mathrm{V}_{1} \otimes \mathrm{V}_{2}$ is the tensor product map of two linear maps as described at the end of Chapter 5. To see what the transition maps look like, we compute;

$$
\begin{aligned}
& \left.\left.\left(\Phi_{\alpha} \otimes \Psi_{\alpha}\right)\right|_{E_{1 p} \otimes E_{2 p}} \circ\left(\Phi_{\beta} \otimes \Psi_{\beta}\right)\right|_{E_{1 p} \otimes E_{2 p}} ^{-1} \\
& =\left.\left.\left.\left.\Phi_{\alpha}\right|_{E_{1 p}} \otimes \Psi_{\alpha}\right|_{E_{2 p}} \circ \Phi_{\beta}\right|_{E_{1 p}} ^{-1} \otimes \Psi_{\beta}\right|_{E_{2 p}} ^{-1} \\
& =\left(\left.\left.\Phi_{\alpha}\right|_{E_{1 p}} \circ \Phi_{\beta}\right|_{E_{1 p}} ^{-1}\right) \otimes\left(\left.\left.\Psi_{\alpha}\right|_{E_{1 p}} \circ \Psi_{\beta}\right|_{E_{1 p}} ^{-1}\right) \\
& =\Phi_{\alpha \beta}(p) \otimes \Psi_{\alpha \beta}(p)
\end{aligned}
$$

Thus the transition maps are given by $p \rightarrow \Phi_{\alpha \beta}(p) \otimes \Psi_{\alpha \beta}(p)$, which is a map from $U_{\alpha}$ to $\mathrm{GL}\left(\mathrm{V}_{1} \otimes \mathrm{V}_{2}\right)$. The group $\mathrm{GL}\left(\mathrm{V}_{1} \otimes \mathrm{V}_{2}\right)$ acts on $\mathrm{V}_{1} \otimes \mathrm{V}_{2}$ in a standard way, and this is the standard effective structure group of the bundle as we have just seen. However, it is also true that the bundle $E_{1} \otimes E_{2}$ has (ineffective) structure group $\mathrm{GL}\left(\mathrm{V}_{1}\right) \times \mathrm{GL}\left(\mathrm{V}_{2}\right)$ via a tensor product representation. Indeed, if $\iota_{1}$ denotes the standard representation of $\operatorname{GL}\left(\mathrm{V}_{1}\right)$ in $\mathrm{V}_{1}$ and $\iota_{2}$ denotes the standard representation of $\mathrm{GL}\left(\mathrm{V}_{2}\right)$ in $\mathrm{V}_{2}$, then we have a tensor product representation $\iota_{1} \otimes \iota_{2}$ of $\mathrm{GL}\left(\mathrm{V}_{1}\right) \times \mathrm{GL}\left(\mathrm{V}_{2}\right)$ in $\mathrm{V}_{1} \otimes \mathrm{V}_{2}$. This is usually not a faithful representation. Using the $\mathrm{GL}\left(\mathrm{V}_{1}\right) \times \mathrm{GL}\left(\mathrm{V}_{2}\right)-$ valued cocycle $p \mapsto h_{\alpha \beta}(p):=\left(\Phi_{\alpha \beta}(p), \Psi_{\alpha \beta}(p)\right)$, together with $\iota_{1} \otimes \iota_{2}$, we see that by definition

$$
\left(\iota_{1} \otimes \iota_{2}\right)_{h_{\alpha \beta}(p)}(\tau)=\left(\Phi_{\alpha \beta}(p) \otimes \Psi_{\alpha \beta}(p)\right)(\tau)
$$

Furthermore, if $\mathrm{V}_{1}=\mathrm{V}_{2}=\mathrm{V}$, then the tensor product representation is usually defined as a representation of $\mathrm{GL}(\mathrm{V})$ rather than $\mathrm{GL}(\mathrm{V}) \times \mathrm{GL}(\mathrm{V})$, and so $E_{1} \otimes E_{2}$ would have a ( $\mathrm{GL}(\mathrm{V}), \iota \otimes \iota$ )-bundle structure where $\iota$ is the standard representation. In this case $\iota \otimes \iota$ is still not a faithful representation since - idv is in the kernel. We can reconstruct the same vector bundle using any of these representation-cocycle pairs via Lemma 6.30. In fact, it is quite common that we have different representations by one group. Suppose that we have two faithful representations $\lambda_{1}$ and $\lambda_{2}$ of a Lie group $G$ acting on $\mathrm{V}_{1}$ and $\mathrm{V}_{2}$ respectively. If $\left\{g_{\alpha \beta}\right\}$ is a cocycle of transition maps, then we can use the pair $\left\{g_{\alpha \beta}, \lambda_{1}\right\}$ in Lemma 6.30 to form a vector bundle $E_{1}$ that has a ( $G, \lambda_{1}$ )-bundle structure by construction. Similarly, we can construct a vector bundle $E_{2}$ with ( $G, \lambda_{2}$ )-bundle structure. If we use $\lambda_{1} \otimes \lambda_{2}$ and the same cocycle $\left\{g_{\alpha \beta}\right\}$, then we obtain a bundle which, as a vector bundle, is $E_{1} \otimes E_{2}$. But by construction, it has a ( $G, \lambda_{1} \otimes \lambda_{2}$ )-structure (possibly ineffective). This is the case in the following exercise:

Exercise 6.50. Suppose that $E$ is a vector bundle with a $(G, \lambda)$-bundle structure given by a ( $G, \lambda$ )-atlas with a corresponding cocycle of transition functions. Show how one may use Theorem 6.30 to construct bundles isomorphic to $E^{*}, E \otimes E$ and $E \otimes E^{*}$ which will have a ( $G, \lambda^{*}$ )-bundle structure, a ( $G, \lambda \otimes \lambda$ )-bundle structure and a ( $G, \lambda \otimes \lambda^{*}$ )-bundle structure respectively.

\subsection*{Smooth Functors}
We have seen that various new vector bundles can be constructed starting with one or more vector bundles. Most of the operations of linear algebra extend to the vector bundle category. We can unify our thinking on these matters by introducing the notion of a $C^{\infty}$ functor (or smooth functor). With $\mathbb{F}=\mathbb{R}$ or $\mathbb{C}$, the set of all $\mathbb{F}$-vector spaces together with linear maps is a category that we denote by $\operatorname{Lin}(\mathbb{F})$. The set of morphisms from V to W is the space of $\mathbb{F}$-linear maps $L(\mathrm{V}, \mathrm{~W})$ (also denoted $\operatorname{Hom}(\mathrm{V}, \mathrm{W})$ ).

Definition 6.51. A covariant $C^{\infty}$ functor $\mathcal{F}$ of one variable on $\operatorname{Lin}(\mathbb{F})$ consists of a map, denoted again by $\mathcal{F}$, that assigns to every $\mathbb{F}$-vector space V an $\mathbb{F}$-vector space $\mathcal{F V}$, and a map, also denoted by $\mathcal{F}$, which assigns to every linear map $A \in L(\mathrm{V}, \mathrm{~W})$, a linear map $\mathcal{F} A \in L(\mathcal{F} \mathrm{V}, \mathcal{F} \mathrm{~W})$ such that\\
(i) $\mathcal{F}: L(\mathrm{V}, \mathrm{~W}) \rightarrow L(\mathcal{F} \mathrm{V}, \mathcal{F} \mathrm{~W})$ is smooth;\\
(ii) $\mathcal{F}\left(\mathrm{id}_{\mathrm{V}}\right)=\operatorname{id}_{\mathcal{F} \mathrm{V}}$ for all $\mathbb{F}$-vector spaces V ;\\
(iii) $\mathcal{F}(A \circ B)=\mathcal{F} A \circ \mathcal{F} B$ for all $A \in L(\mathrm{U}, \mathrm{V})$ and $B \in L(\mathrm{V}, \mathrm{~W})$ and vector spaces $\mathrm{U}, \mathrm{V}$ and W .

As an example we have the $C^{\infty}$ functor which assigns to each V the $k$-fold direct sum $\bigoplus_{k} \mathrm{V}=\mathrm{V} \oplus \cdots \oplus \mathrm{V}$ and to each linear map $A \in L(\mathrm{V}, \mathrm{~W})$\\
the map

$$
\bigoplus_{k} A: \bigoplus_{k} \mathrm{V} \rightarrow \bigoplus_{k} \mathrm{~W}
$$

given by $\bigoplus_{k} A\left(v_{1}, \ldots, v_{k}\right):=\left(A v_{1}, \ldots, A v_{k}\right)$. Similarly there is the functor which assigns to each V the $k$-fold tensor product $\otimes^{k} \mathrm{V}=\mathrm{V} \otimes \cdots \otimes \mathrm{V}$ and to each $A \in L(\mathrm{V}, \mathrm{~W})$ the map $\otimes^{k} A: \otimes^{k} \mathrm{V} \rightarrow \otimes^{k} \mathrm{~W}$ given on homogeneous elements by $\left(\otimes^{k} A\right)\left(v_{1} \otimes \cdots \otimes v_{k}\right):=A v_{1} \otimes \cdots \otimes A v_{k}$.

One can also consider $C^{\infty}$ covariant functors of several variables. For example, we may assign to each pair of vector spaces $(\mathrm{V}, \mathrm{W})$ the tensor product $\mathrm{V} \otimes \mathrm{W}$, and to each pair $(A, B) \in L\left(\mathrm{V}, \mathrm{V}^{\prime}\right) \times L\left(\mathrm{~W}, \mathrm{~W}^{\prime}\right)$, the map $A \otimes B: \mathrm{V} \otimes \mathrm{W} \rightarrow \mathrm{V}^{\prime} \otimes \mathrm{W}^{\prime}$.

There is also a similar notion of contravariant $C^{\infty}$ functor:\\
Definition 6.52. A contravariant $C^{\infty}$ functor $\mathcal{F}$ of one variable on $\operatorname{Lin}(\mathbb{F})$ consists of a map, denoted again by $\mathcal{F}$, which assigns to every $\mathbb{F}$ vector space V an $\mathbb{F}$-vector space $\mathcal{F} \mathrm{V}$, and a map, also denoted by $\mathcal{F}$, which assigns to every linear map $A \in L(\mathrm{V}, \mathrm{~W})$ a linear map $\mathcal{F} A \in L(\mathcal{F} \mathrm{~W}, \mathcal{F} \mathrm{V})$ (notice the reversal) such that\\
(i) $\mathcal{F}: L(\mathrm{V}, \mathrm{~W}) \rightarrow L(\mathcal{F} \mathrm{~W}, \mathcal{F} \mathrm{V})$ is smooth;\\
(ii) $\mathcal{F}\left(\mathrm{id}_{\mathrm{V}}\right)=\operatorname{id}_{\mathcal{F} \mathrm{V}}$ for all $\mathbb{F}$-vector spaces V ;\\
(iii) $\mathcal{F}(A \circ B)=\mathcal{F} B \circ \mathcal{F} A$ for all $A \in L(\mathrm{U}, \mathrm{V})$ and $B \in L(\mathrm{V}, \mathrm{~W})$ and vector spaces $\mathrm{U}, \mathrm{V}$ and W .

The map that assigns to each vector space its dual and to each map its dual map (transpose) is a contravariant $C^{\infty}$ functor $\mathcal{F}$. One may define the notion of a $C^{\infty}$ functor of several variables which may be covariant in some variables and contravariant in others. For example, consider the functor of two variables that assigns to each pair ( $\mathrm{V}, \mathrm{W}$ ) the space $\mathrm{V} \otimes \mathrm{W}^{*}$ and to each pair $(A, B) \in L\left(\mathrm{V}_{1}, \mathrm{V}_{2}\right) \times L\left(\mathrm{~W}_{1}, \mathrm{~W}_{2}\right)$ the map $A \otimes B^{*}: \mathrm{V}_{1} \otimes \mathrm{~W}_{2}^{*} \rightarrow \mathrm{V}_{2} \otimes \mathrm{~W}_{1}^{*}$.

Theorem 6.53. Let $\mathcal{F}$ be a $C^{\infty}$ functor of $m$ variables on $\operatorname{Lin}(\mathbb{F})$ and let $E_{1}, \ldots, E_{m}$ be $\mathbb{F}$-vector bundles with respective typical fibers $\mathrm{V}_{1}, \ldots, \mathrm{V}_{m}$. Then the set

$$
E:=\mathcal{F}\left(E_{1}, \ldots, E_{m}\right):=\bigcup_{p} \mathcal{F}\left(\left.E_{1}\right|_{p}, \ldots,\left.E_{m}\right|_{p}\right)
$$

together with the map $\pi: E \rightarrow M$ which takes elements of $\mathcal{F}\left(\left.E_{1}\right|_{p}, \ldots,\left.E_{m}\right|_{p}\right)$ to $p$ is naturally a vector bundle with typical fiber $\mathcal{F}\left(\mathrm{V}_{1}, \ldots, \mathrm{V}_{m}\right)$.

Proof. We will only prove the case of $m=2$ with covariant first variable and contravariant second variable. This should make it clear how the general case would go while keeping the notational complexity under control.

Given vector bundles $\pi_{1}: E_{1} \rightarrow M$ and $\pi_{2}: E_{2} \rightarrow M$, the total space of the constructed bundle is $\bigcup_{p} \mathcal{F}\left(\left.E_{1}\right|_{p},\left.E_{2}\right|_{p}\right)$ with the obvious projection which we call $\pi$. Let ( $\phi_{\alpha}, U_{\alpha}$ ) be a VB-atlas for $E_{1}$ and ( $\psi_{\alpha}, U_{\alpha}$ ) a VBatlas for $E_{2}$ (we have arranged that both atlases use the same cover by going to a common refinement as usual). For each $p$, let $E_{p}$ denote the fiber $\mathcal{F}\left(\left.E_{1}\right|_{p},\left.E_{2}\right|_{p}\right)$. Fix $\alpha$ and for each $p \in U_{\alpha}$ define $\left.\Theta_{\alpha}\right|_{p} \in L\left(E_{p}, \mathcal{F}\left(\mathrm{V}_{1}, \mathrm{V}_{2}\right)\right)$ by

$$
\left.\Theta_{\alpha}\right|_{p}:=\mathcal{F}\left(\left.\Phi_{\alpha}\right|_{p},\left.\Psi_{\alpha}\right|_{p} ^{-1}\right),
$$

where $\phi_{\alpha}=\left(\pi_{1}, \Phi_{\alpha}\right)$ and $\psi_{\alpha}=\left(\pi_{2}, \Psi_{\alpha}\right)$. Then define $\Theta_{\alpha}: \pi^{-1}\left(U_{\alpha}\right) \rightarrow \mathcal{F}\left(\mathrm{V}_{1}, \mathrm{V}_{2}\right)$ by $\Theta_{\alpha}(\epsilon)=\left.\Theta_{\alpha}\right|_{p}(\epsilon)$ whenever $\epsilon \in \mathcal{F}\left(\left.E_{1}\right|_{p},\left.E_{2}\right|_{p}\right)$. Next define

$$
\theta_{\alpha}=\left(\pi, \Theta_{\alpha}\right): \pi^{-1} U_{\alpha} \rightarrow U_{\alpha} \times \mathcal{F}\left(\mathrm{V}_{1}, \mathrm{V}_{2}\right) .
$$

The family $\left\{\left(\theta_{\alpha}, U_{\alpha}\right)\right\}$ is to be a VB-atlas for $E$. We check the transition maps:

$$
\begin{aligned}
\Theta_{\alpha \beta}(p) & =\left.\left.\Theta_{\alpha}\right|_{p} \circ \Theta_{\beta}^{-1}\right|_{p} \\
& =\mathcal{F}\left(\left.\Phi_{\alpha}\right|_{p},\left.\Psi_{\alpha}\right|_{p} ^{-1}\right) \circ \mathcal{F}\left(\left.\Phi_{\beta}\right|_{p},\left.\Psi_{\beta}\right|_{p} ^{-1}\right)^{-1} \\
& =\mathcal{F}\left(\left.\Phi_{\alpha}\right|_{p},\left.\Psi_{\alpha}\right|_{p} ^{-1}\right) \circ \mathcal{F}\left(\left.\Phi_{\beta}\right|_{p} ^{-1},\left.\Psi_{\beta}\right|_{p}\right) \\
& =\mathcal{F}\left(\left.\left.\Phi_{\alpha}\right|_{p} \circ \Phi_{\beta}\right|_{p} ^{-1},\left.\left.\Psi_{\beta}\right|_{p} \circ \Psi_{\alpha}\right|_{p} ^{-1}\right) \\
& =\mathcal{F}\left(\Phi_{\alpha \beta}(p), \Psi_{\beta \alpha}(p)\right) .
\end{aligned}
$$

(Remember that the functor is contravariant in the second variable.) Now we can see from the properties of $\Phi_{\alpha \beta}, \Psi_{\beta \alpha}$ and the definition of $C^{\infty}$ functor that $\mathcal{F}\left(\Phi_{\alpha \beta}(p), \Psi_{\beta \alpha}(p)\right) \in \mathrm{GL}\left(\mathcal{F}\left(\mathrm{V}_{1}, \mathrm{V}_{2}\right)\right)$ and the maps $\Theta_{\alpha \beta}: U_{\alpha} \cap U_{\beta} \rightarrow \operatorname{GL}\left(\mathcal{F}\left(\mathrm{V}_{1}, \mathrm{V}_{2}\right)\right)$ are smooth. \(\square\)

\subsection*{Hom}
Let $\xi_{1}:=\left(E_{1}, \pi_{1}, M, \mathrm{V}\right)$ and $\xi_{2}:=\left(E_{2}, \pi_{2}, M, \mathrm{V}\right)$ be smooth $\mathbb{F}$-vector bundles. The bundle whose fiber over $p \in M$ is $L\left(E_{1 p}, E_{2 p}\right)=\operatorname{Hom}\left(E_{1 p}, E_{2 p}\right)$ is denoted by $\operatorname{Hom}\left(\xi_{1}, \xi_{2}\right)$ or less precisely, by referring to the total space $\operatorname{Hom}\left(E_{1}, E_{2}\right)$. Here $\operatorname{Hom}\left(E_{1 p}, E_{2 p}\right)$ denotes $\mathbb{F}$-linear maps. If $f: E_{1} \rightarrow E_{2}$ is vector bundle homomorphism over $M$, then we may obtain a section $s$ of $\operatorname{Hom}\left(E_{1}, E_{2}\right)$ by defining $s: p \mapsto f \mid{ }_{E_{1 p}}$. Conversely, given $s \in \Gamma\left(\operatorname{Hom}\left(E_{1}, E_{2}\right)\right)$ we define $f: E_{1} \rightarrow E_{2}$ by requiring that

$$
\left.f\right|_{E_{1 p}}=s(p) .
$$

Thus every element of $\operatorname{Hom}\left(E_{1}, E_{2}\right)$ can be identified with a vector bundle homomorphism over $M$.

Exercise 6.54. Let $E_{1} \rightarrow M_{1}$ and $E_{2} \rightarrow M_{2}$ be smooth vector bundles. Show that the set of vector bundle homomorphisms along a smooth\\
map $g: M_{1} \rightarrow M_{2}$ is in natural bijection with the sections of the bundle $\operatorname{Hom}\left(E_{1}, g^{*} E_{2}\right)$.

Since $\Gamma\left(E_{1}\right)$ and $\Gamma\left(E_{2}\right)$ are $C^{\infty}(M, \mathbb{F})$ modules, we can look at the $C^{\infty}(M, \mathbb{F})$ module $\operatorname{Hom}\left(\Gamma\left(E_{1}\right), \Gamma\left(E_{2}\right)\right)$. Then we have

Proposition 6.55. Let $E_{1} \rightarrow M$ and $E_{2} \rightarrow M$ be smooth $\mathbb{F}$-vector bundles. Then $\Gamma\left(\operatorname{Hom}\left(E_{1}, E_{2}\right)\right)$ and $\operatorname{Hom}\left(\Gamma\left(E_{1}\right), \Gamma\left(E_{2}\right)\right)$ are naturally isomorphic as $C^{\infty}(M, \mathbb{F})$ modules.

Proof. To each section $s \in \Gamma\left(\operatorname{Hom}\left(E_{1}, E_{2}\right)\right)$ we assign a map $\phi_{s}: \Gamma\left(E_{1}\right) \rightarrow \Gamma\left(E_{2}\right)$ defined by the formula

$$
\phi_{s}(\sigma)(p)=s(p) \sigma(p) \text { for } \sigma \in \Gamma\left(E_{1}\right)
$$

Then we obtain a map $\Phi: \Gamma\left(\operatorname{Hom}\left(E_{1}, E_{2}\right)\right) \rightarrow \operatorname{Hom}\left(\Gamma\left(E_{1}\right), \Gamma\left(E_{2}\right)\right)$ which is defined by $\Phi: s \mapsto \phi_{s}$. The smoothness of $s$ and $\sigma$ implies the smoothness of $\phi_{s}(\sigma)$ and then the smoothness of $\phi_{s}$. The map $\phi_{s}$ is clearly in $\operatorname{Hom}\left(\Gamma\left(E_{1}\right), \Gamma\left(E_{2}\right)\right)$, and it is not difficult to check that $\Phi$ is also a module homomorphism.

Now suppose we are given a module homomorphism $\phi: \Gamma\left(E_{1}\right) \rightarrow \Gamma\left(E_{2}\right)$. Define $s \in \Gamma\left(\operatorname{Hom}\left(E_{1}, E_{2}\right)\right)$ by

$$
s(p)\left(v_{p}\right):=\phi(\sigma)(p) \text { for } v_{p} \in E_{1 p},
$$

where $\sigma$ is any section in $\Gamma\left(E_{1}\right)$ such that $\sigma(p)=v_{p}$. We need to show that this is well-defined. It suffices to show that if $\sigma(p)=0$ then $\phi(\sigma)(p)=0$. Let $\left(\chi_{1}, \ldots, \chi_{k}\right)$ be a local frame field for $E_{1}$ defined over an open set $U$. Choose $g \in C^{\infty}(M)$ with support in $U$ and $g(p)=1$. Define fields $X_{i}:=g \chi_{i}$ and extend by zero outside of $U$. Then there exist functions $f^{i}$ such that

$$
g \sigma=\sum_{i=1}^{k} f^{i} X_{i}
$$

and since $\sigma(p)=0$, we must have $f^{i}(p)=0$ for all $i$. Then we have

$$
\begin{aligned}
(\phi \sigma)(p) & =g(p)(\phi(\sigma))(p)=(g \phi(\sigma))(p) \\
& =\phi(g \sigma)(p)=\phi\left(\sum_{i=1}^{k} f^{i} X_{i}\right)(p) \\
& =\sum_{i=1}^{k}\left(f^{i} \phi\left(X_{i}\right)\right)(p)=\sum_{i=1}^{k} f^{i}(p) \phi\left(X_{i}\right)(p)=0 .
\end{aligned}
$$

The constructed map is easily checked to be the inverse of the map $\Phi: s \mapsto \phi_{s}$.

If ( $\sigma_{1}, \ldots, \sigma_{k_{1}}$ ) is a local frame field for $E_{1} \rightarrow M$ over an open set $U$, and if $\left(\phi_{1}, \ldots, \phi_{k_{2}}\right)$ a local frame field for $E_{2} \rightarrow M$ also over $U$, then we may construct a local frame field for $\operatorname{Hom}\left(E_{1}, E_{2}\right)$. The frame field is $\left\{e_{i}^{j}\right\}$ where $e_{i}^{j}\left(\sigma_{k}\right)=\delta_{k j} \phi_{i}$. If $\operatorname{Hom}\left(E_{1}, E_{2}\right)$ is identified with $E_{2} \otimes E_{1}^{*}$, then $e_{i}^{j}$ is $\phi_{i} \otimes \sigma^{j}$ where ( $\sigma^{1}, \ldots, \sigma^{k_{1}}$ ) is the dual frame field to ( $\sigma_{1}, \ldots, \sigma_{k_{1}}$ ). Suppose that over an open set $V$, we have frame fields $\left(\widetilde{\sigma}_{1}, \ldots, \widetilde{\sigma}_{k_{1}}\right)$ and $\left(\widetilde{\phi}_{1}, \ldots, \widetilde{\phi}_{k_{2}}\right)$. If $U \cap V$ is nonempty, then

$$
\begin{aligned}
& \left(\widetilde{\sigma}_{1}, \ldots, \widetilde{\sigma}_{k_{1}}\right)=\left(\sigma_{1}, \ldots, \sigma_{k_{1}}\right) C \\
& \left(\widetilde{\phi}_{1}, \ldots, \widetilde{\phi}_{k_{2}}\right)=\left(\phi_{1}, \ldots, \phi_{k_{2}}\right) D
\end{aligned}
$$

for smooth matrix-valued functions $C$ and $D$ defined on $U \cap V$. We obtain $\widetilde{e}_{i}^{j}=\widetilde{\phi}_{i} \otimes \widetilde{\sigma}^{j}=\sum_{r, s}\left(C^{-1}\right)_{s}^{j} e_{r}^{s} D_{i}^{r}$. If $A \in \Gamma\left(\operatorname{Hom}\left(E_{1}, E_{2}\right)\right)$ and on $U \cap V$ we have

$$
A=\sum \widetilde{A}_{j}^{i} \widetilde{e}_{i}^{j}=\sum A_{j}^{i} e_{i}^{j},
$$

then $\widetilde{A}_{j}^{i}=\sum_{r, s}\left(D^{-1}\right)_{r}^{i} A_{s}^{r} C_{j}^{s}$ on $U \cap V$.\\
Of special importance is $\operatorname{End}(E):=\operatorname{Hom}(E, E) \rightarrow M$ for a given vector bundle $E \rightarrow M$. Here if $\left(\widetilde{\sigma}_{1}, \ldots, \widetilde{\sigma}_{k}\right)=\left(\sigma_{1}, \ldots, \sigma_{k}\right) C$ is a change of frame field, then we take $\phi_{i}=\sigma_{i}$ and $\widetilde{\phi}_{i}=\widetilde{\sigma}_{i}$ so that the above rule specializes to

$$
\widetilde{A}_{j}^{i}=\sum_{r, s}\left(C^{-1}\right)_{r}^{i} A_{s}^{r} C_{j}^{s} \quad \text { on } \quad U \cap V
$$

which is the signature transformation law for sections of $\operatorname{End}(E)$. This bundle is important in the study of covariant derivatives and curvature. Notice how everything is formally similar to operations in linear algebra, but here we are dealing with fields and functions (often only locally defined).

\subsection*{Algebra Bundles}
Let $\mathbb{F}$ be $\mathbb{C}$ or $\mathbb{R}$. Recall that an $\mathbb{F}$-algebra is a vector space $V$ with a bilinear map $\mathrm{V} \times \mathrm{V} \rightarrow \mathrm{V}$ giving a product on V . Such a bilinear map is uniquely specified by the associated linear map $\mathrm{V} \otimes \mathrm{V} \rightarrow \mathrm{V}$ (see Definition D.17).

Definition 6.56. Let V be an $\mathbb{F}$-algebra. An $\mathbb{F}$-vector bundle $\xi=(E, \pi$, $M, \mathrm{V})$ is called an $\mathbb{F}$-algebra bundle if each fiber $E_{p}$ has an $\mathbb{F}$-algebra structure in such a way that the associated maps $E_{p} \otimes E_{p} \rightarrow E_{p}$ combine to give a vector bundle homomorphism $E \otimes E \rightarrow E$ and $\xi$ has a VB-atlas $\left\{\left(U_{\alpha}, \phi_{\alpha}\right)\right\}$ such that for each $\alpha,\left.\Phi_{\alpha}\right|_{E_{p}}: E_{p} \rightarrow \mathrm{V}$ is an algebra isomorphism for all $p \in U_{\alpha}$. Here $\phi_{\alpha}=\left(\pi, \Phi_{\alpha}\right)$ as usual.

Example 6.57. The endomorphism bundle $\operatorname{End}(E) \rightarrow M$ of a vector bundle ( $E, \pi, M, \mathrm{V}$ ) is the bundle whose fiber at $p$ is $\operatorname{End}\left(E_{p}\right)=\operatorname{Hom}\left(E_{p}, E_{p}\right)$. The space $\operatorname{End}\left(E_{p}\right)$ has an $\mathbb{F}$-algebra structure given by composition of linear maps.

Example 6.58. Let $\xi=(E, \pi, M, \mathrm{V})$ be an $\mathbb{F}$-vector bundle. The corresponding general linear Lie algebra bundle is the bundle whose fiber at $p$ is $L\left(E_{p}, E_{p}\right)$ but with the product being $(f, g) \mapsto[f, g]:=f \circ g-g \circ f$. When given this product, the fiber is denoted $\mathfrak{g} \mathfrak{l}\left(E_{p}\right)$ and is the Lie algebra of $\operatorname{GL}\left(E_{p}\right)$. The total space of the general linear Lie algebra bundle is then denoted by $\mathfrak{g} \mathfrak{l}(E)$. This provides an example of a Lie algebra bundle.

Recall that when a direct sum has an infinite number of summands, we require that each element have finite support. In other words, $\bigoplus_{i=1}^{\infty} V_{i}$ is the vector space of all formal sums $\sum_{i=1}^{\infty} v_{i}$, where $v_{i} \in \mathrm{V}_{i}$ and all but finitely many of the $v_{i}$ 's are zero.

Example 6.59. The definition of tensor algebra is given as Definition D. 46 of Appendix D. Let $\xi=(E, \pi, M, \mathrm{V})$ be an $\mathbb{F}$-vector bundle and consider the bundle $\bigotimes E \rightarrow M$ whose fiber at $p$ is the $\mathbb{F}$-tensor algebra $\bigotimes E_{p}= \mathbb{R} \oplus E_{p} \oplus\left(E_{p} \otimes E_{p}\right) \oplus \cdots$. If we consider the disjoint union

$$
\bigotimes E=\bigcup_{p \in M} \otimes E_{p},
$$

with the obvious projection $\bigotimes E \rightarrow M$, then it seems that we have a bundle with algebra fibers. However, each fiber $\bigotimes E_{p}$ is infinite-dimensional. On the other hand, we do at least have a nested sequence of vector bundles $\cdots \subset \bigotimes^{\leq k} E \subset \bigotimes^{\leq k+1} E \subset \cdots$, where

$$
\bigotimes^{\leq k} E=\mathbb{R} \oplus E \oplus(E \otimes E) \oplus \cdots \oplus\left(\bigotimes^{k} E\right)
$$

\subsection*{Sheaves}
Let $E \rightarrow M$ be a vector bundle. We have seen that $\Gamma(M, E)$ is a module over the smooth functions $C^{\infty}(M)$. It is important to realize that having a vector bundle at hand not only provides a module, but a family of modules $\{\Gamma(U, E)\}$ parametrized by the open subsets $U$ of $M$. How are these modules related to each other?

Consider a section $\sigma: M \rightarrow E$. Given any open set $U \subset M$, we may always produce the restricted section $\left.\sigma\right|_{U}: U \rightarrow E$. This gives us a family of sections; one for each open set $U$. Conversely, suppose that we have a family of sections $\sigma_{U}: U \rightarrow E$ where $U$ varies over the open sets (or just a cover of $M)$. When is it the case that such a family is just the family of restrictions of some section $\sigma: M \rightarrow E$ ? To help with these kinds of questions, and to provide a language that will occasionally be convenient, we will introduce another formalism. This is the formalism of sheaves and presheaves. The formalism of sheaf theory is used for studying the interplay between the local and the global, for example, using sheaf cohomology theory. It is especially\\
useful in complex geometry. Sheaf theory also provides a very good framework within which to develop the foundations of supergeometry, which is an extension of differential geometry that incorporates the important notion of "fermionic variables". A deep understanding of sheaf theory is not necessary for what we do here and it would be enough to acquire a basic familiarity with the definitions since we only want the convenience of the language.

Definition 6.60. A presheaf of abelian groups (resp. rings, etc.) on a manifold (or more generally a topological space) $M$ assigns to each open set an abelian group (resp. ring, etc.) $\mathcal{M}(U)$ and assigns to each nested pair $V \subset U$ of open sets a homomorphism $r_{V}^{U}: \mathcal{M}(U) \rightarrow \mathcal{M}(V)$ of abelian groups (resp. rings, etc.) such that\\
(i) $r_{W}^{V} \circ r_{V}^{U}=r_{W}^{U}$ whenever $W \subset V \subset U$;\\
(ii) $r_{V}^{V}=\operatorname{id}_{V}$ for all open $V \subset M$.

Definition 6.61. Let $\mathcal{M}$ be a presheaf and $\mathcal{R}$ a presheaf of rings over $M$. If for each open $U \subset M$ we have that $\mathcal{M}(U)$ is a module over the $\operatorname{ring} \mathcal{R}(U)$, and if multiplication commutes with restriction, that is, if the following diagram commutes for each nested pair $V \subset U$,\\
\includegraphics[scale=0.25, center]{2025_10_09_265d7e1a22c89f79c21eg-304}\\
then we say that $\mathcal{M}$ is a presheaf of modules over $\mathcal{R}$.\\
Definition 6.62. Let $\mathcal{M}_{1}$ and $\mathcal{M}_{2}$ be presheaves over $M$. A presheaf morphism $h: \mathcal{M}_{1} \rightarrow \mathcal{M}_{2}$ over $M$ is a collection of morphisms, $h_{U}: \mathcal{M}_{1}(U) \rightarrow \mathcal{M}_{2}(U)$, one for each open set and such that whenever $V \subset U$, the following diagram commutes:\\
\includegraphics[scale=0.25, center]{2025_10_09_265d7e1a22c89f79c21eg-304(1)}

Note that we have used the same notation for the restriction maps of both presheaves.

Definition 6.63. Let $\mathcal{M}$ be a presheaf. A family $\left\{s_{\alpha}\right\}$ with $s_{\alpha} \in \mathcal{M}\left(U_{\alpha}\right)$ is called consistent if $r_{U_{\alpha} \cap U_{\beta}}^{U} s_{\alpha}=r_{U_{\alpha} \cap U_{\beta}}^{U} s_{\beta}$ whenever $U_{\alpha} \cap U_{\beta} \neq \emptyset$.

Definition 6.64. We will call a presheaf $\mathcal{M}$ a sheaf if the following properties hold whenever $U=\bigcup_{U_{\alpha} \in \mathcal{U}} U_{\alpha}$ for some collection of open sets $\mathcal{U}$.\\
(i) If $s_{1}, s_{2} \in \mathcal{M}(U)$ and $r_{U_{\alpha}}^{U} s_{1}=r_{U_{\alpha}}^{U} s_{2}$ for all $U_{\alpha} \in \mathcal{U}$, then $s_{1}=s_{2}$.\\
(ii) Given a consistent family $\left\{s_{\alpha}: s_{\alpha} \in \mathcal{M}\left(U_{\alpha}\right)\right\}$, there exists $s \in \mathcal{M}(U)$ such that $r_{U_{\alpha}}^{U} s=s_{\alpha}$.\\
(iii) $\mathcal{M}(\emptyset)$ is the trivial group, (resp. ring, etc.).

If we need to indicate the space $M$ involved, we will write $\mathcal{M}_{M}$ instead of $\mathcal{M}$.

Definition 6.65. A morphism of sheaves is a morphism of the underlying presheaf.

The assignment $C^{\infty}(\cdot): U \mapsto C^{\infty}(U)$ is a sheaf of rings. This sheaf will also be denoted by $C_{M}^{\infty}$. The best and most important example of a sheaf of modules over $C^{\infty}(\cdot)$ is the assignment $\Gamma(\cdot, E): U \mapsto \Gamma(U, E)$ for some vector bundle $E \rightarrow M$, where by definition $r_{V}^{U}(s)=\left.s\right|_{V}$ for $s \in \Gamma(U, E)$. In other words, $r_{V}^{U}$ is just the restriction map. Let us denote this (pre)sheaf by $\Gamma_{E}: U \mapsto \Gamma_{E}(U):=\Gamma(U, E)$.

Exercise 6.66. For each open set $U$ in a manifold $M$, let $B(U)$ denote the ring of bounded smooth functions defined on $U$. Show that $U \mapsto B(U)$ defines a presheaf that is not a sheaf.

Many if not most of the constructions and operations we introduce for sections of vector bundles are really also operations appropriate to the (pre)sheaf category. Naturality with respect to restrictions is one of the features that is often not even mentioned (precisely because it seems obvious). This is the inspiration for a slight twist on our notation.

\begin{center}
\begin{tabular}{llll}
 & Global & Local & Sheaf \\
Functions on $M$ & $C^{\infty}(M)$ & $C^{\infty}(U)$ & $C_{M}^{\infty}$ \\
Vector fields on $M$ & $\mathfrak{X}(M)$ & $\mathfrak{X}(U)$ & $\mathfrak{X}_{M}$ \\
Sections of $E$ & $\Gamma(E)$ & $\Gamma(U, E)$ & $\Gamma_{E}$ \\
\end{tabular}
\end{center}

where $C_{M}^{\infty}: U \mapsto C_{M}^{\infty}(U):=C^{\infty}(U), \mathfrak{X}_{M}: U \mapsto \mathfrak{X}_{M}(U):=\mathfrak{X}(U)$ and so on.\\
Notation 6.67. For example, when we say that $D: C_{M}^{\infty} \rightarrow C_{M}^{\infty}$ is a derivation we mean that $D$ is actually a family of algebra derivations $D_{U}$ : $C_{M}^{\infty}(U) \rightarrow C_{M}^{\infty}(U)$ indexed by open sets $U$ such that we have naturality with respect to restrictions; i.e., diagrams of the form below for $V \subset U$ commute:\\
\includegraphics[scale=0.25, center]{2025_10_09_265d7e1a22c89f79c21eg-305}

It is easy to see that all of the following examples are sheaves. In each case, the maps $r_{V}^{U}$ are just the restriction maps.

Example 6.68 (Sheaf of holomorphic functions). Sheaf theory really shows its strength in complex analysis. This example is one of the most studied. However, we have not studied the notion of a complex manifold, and so this example is for those readers with some exposure to complex manifolds (see the online supplement). Let $M$ be a complex manifold and let $\mathcal{O}_{M}(U)$ be the algebra of holomorphic functions defined on $U$. Here too, $\mathcal{O}_{M}$ is a sheaf of modules over itself. Whereas the sheaf $\mathcal{C}_{M}^{\infty}$ always has global sections, the same is not true for $\mathcal{O}_{M}$. The sheaf-theoretic approach to the study of obstructions to the existence of global holomorphic functions has been very successful.

For a bit more on sheaves, see the online supplement [Lee, Jeff].

\subsection*{Principal and Associated Bundles}
Let $\pi: E \rightarrow M$ be a vector bundle with typical fiber V and for every $p \in M$ let $\mathrm{GL}\left(\mathrm{V}, E_{p}\right)$ denote the set of linear isomorphisms from V to $E_{p}$. If we choose a fixed basis ( $\mathbf{e}_{1}, \ldots, \mathbf{e}_{k}$ ) for V , then each frame ( $u_{1}, \ldots, u_{k}$ ) at $p$ gives an element $u \in \mathrm{GL}\left(\mathrm{V}, E_{p}\right)$ defined by

$$
u(v):=\sum v^{i} u_{i},
$$

where $v=\sum v^{i} \mathbf{e}_{i}$. We identify $u$ with ( $u_{1}, \ldots, u_{k}$ ) and refer to it as a frame. With this identification, notice that if $\sigma_{\alpha}:=\sigma_{\phi_{\alpha}}$ is the local frame field coming from a VB-chart ( $U_{\alpha}, \phi_{\alpha}$ ) as described above, then we have

$$
\sigma_{\alpha}(p)=\left.\Phi_{\alpha}\right|_{E_{p}} \text { for } p \in U_{\alpha}
$$

Now let

$$
F(E):=\bigcup_{p \in M} \mathrm{GL}\left(\mathrm{V}, E_{p}\right) \quad \text { (disjoint union) }
$$

It will shortly be clear that $F(E)$ is a smooth manifold and the total space of a fiber bundle. Let $\wp: F(E) \rightarrow M$ be the projection map defined by $\wp(u)=p$ for $u \in \mathrm{GL}\left(\mathrm{V}, E_{p}\right)$. Observe that GL(V) acts on the right of the set $F(E)$; the action $F(E) \times \mathrm{GL}(\mathrm{V}) \rightarrow F(E)$ is given by $r:(u, g) \mapsto u g=u \circ g$. If we pick a fixed basis for V as above, then we may view $g$ as a matrix and an element $u \in \mathrm{GL}\left(\mathrm{V}, E_{p}\right)$ as a basis $\left(u_{1}, \ldots, u_{k}\right)$. In this case, we have

$$
u g=\left(\sum u_{i} g_{1}^{i}, \ldots, \sum u_{i} g_{k}^{i}\right)
$$

It is easy to see that the orbit of a frame at $p$ is exactly the set $\wp^{-1}(p)= \mathrm{GL}\left(\mathrm{V}, E_{p}\right)$ and that the action is free. For each VB-chart ( $U, \phi$ ) for $E$, let $\sigma_{\phi}$ be the associated frame field. Define $f_{\phi}: U \times \mathrm{GL}(\mathrm{V}) \rightarrow \wp^{-1}(U)$\\
by $f_{\phi}(p, g)=\sigma_{\phi}(p) g$. It is easy to check that this is a bijection. Let $\widetilde{\phi}$ : $\wp^{-1}(U) \rightarrow U \times \mathrm{GL}(\mathrm{V})$ be the inverse of this map. We have $\widetilde{\phi}=(\wp, \widetilde{\Phi})$, where $\widetilde{\Phi}$ is uniquely determined by $\widetilde{\phi}$. Starting with a VB-atlas $\left\{\left(U_{\alpha}, \phi_{\alpha}\right)\right\}$ for $E$, we obtain a family $\left\{\widetilde{\phi}_{\alpha}: \wp^{-1}\left(U_{\alpha}\right) \rightarrow U_{\alpha} \times \mathrm{GL}(\mathrm{V})\right\}$ of trivializations which gives a fiber bundle atlas $\left\{\left(U_{\alpha}, \widetilde{\phi}_{\alpha}\right)\right\}$ for $F(E) \rightarrow M$ and simultaneously induces the smooth structure.

Definition 6.69. Let $\pi: E \rightarrow M$ be a vector bundle with typical fiber V . The fiber bundle ( $F(E), \wp, M, \mathrm{GL}(\mathrm{V})$ ) constructed above is called the linear frame bundle of $E$ and is usually denoted simply by $F(E)$. The frame bundle for the tangent bundle of a manifold $M$ is often denoted by $F(M)$ rather than by $F(T M)$.

Notation 6.70. It would be perhaps more appropriate to refer to the frame bundle of a vector bundle $\xi=(E, \pi, M, \mathrm{V})$ as $F(\xi)$ since the notation $F(E)$ is inconsistent with the notation $F(M)$ above. After all, $E$ is itself a manifold. Despite this we will continue with the dangerous notation.

Notice that any VB-atlas for $E$ induces an atlas on $F(E)$ according to our considerations above. We have

$$
\widetilde{\phi}_{\alpha} \circ \widetilde{\phi}_{\beta}^{-1}(p, g)=\widetilde{\phi}_{\alpha}\left(\sigma_{\beta}(p) g\right)=\widetilde{\phi}_{\alpha}\left(\left.\Phi_{\beta}\right|_{E_{p}} g\right)=\left(p, \Phi_{\alpha \beta}(p) g\right) .
$$

Thus the transition functions of $F(E)$ are given by the standard transition functions of $E$ acting by left multiplication on $\mathrm{GL}(\mathrm{V})$. The cocycle corresponding to the bundle atlas for $F(E)$ that we constructed from a VB-atlas $\left\{\left(U_{\alpha}, \phi_{\alpha}\right)\right\}$ for the original vector bundle $E$, is the very cocycle $\Phi_{\alpha \beta}$ deriving from this atlas on $E$.

We need to make one more observation concerning the right action of $\mathrm{GL}(\mathrm{V})$ on $F(E)$. Take a VB-chart for $E$, say ( $U, \phi$ ), and let us look again at the associated chart $(U, \widetilde{\phi})$ for $F(E)$. First, consider the trivial bundle $\mathrm{pr}_{1}: U \times \mathrm{GL}(\mathrm{V}) \rightarrow U$ and define the obvious right action on the total space $\widetilde{\phi}\left(\wp^{-1}(U)\right)=U \times \operatorname{GL}(\mathrm{V})$ by $\left(\left(p, g_{1}\right), g\right) \mapsto\left(p, g_{1} g\right):=\left(p, g_{1}\right) \cdot g$. Then this action is transitive on the fibers of this trivial bundle. Of course since GL(V) acts on $F(E)$ and preserves fibers, it also acts by restriction on $\wp^{-1}(U)$.

Proposition 6.71. The bundle map $\widetilde{\phi}: \wp^{-1}(U) \rightarrow U \times \mathrm{GL}(\mathrm{V})$ is equivariant with respect to the right actions described above.

Proof. We first look at the inverse:

$$
\widetilde{\phi}^{-1}\left(p, g_{1}\right) g=\sigma_{\phi}(p) g_{1} g=\widetilde{\phi}^{-1}\left(p, g_{1} g\right)
$$

To see this from the point of view of $\widetilde{\phi}$ rather than its inverse, take $u \in \wp^{-1}(U) \subset F(E)$ and let ( $p, g_{1}$ ) be the unique pair such that $u=\widetilde{\phi}^{-1}\left(p, g_{1}\right)$.

Then

$$
\widetilde{\phi}(u g)=\widetilde{\phi}\left(\widetilde{\phi}^{-1}\left(p, g_{1}\right) g\right)=\widetilde{\phi} \widetilde{\phi}^{-1}\left(p, g_{1} g\right)=\left(p, g_{1} g\right)=\left(p, g_{1}\right) \cdot g=\widetilde{\phi}(u) \cdot g .
$$

A section of $F(E)$ over an open set $U$ in $M$ is just a frame field over $U$. A global frame field is a global section of $F(E)$, and clearly a global section exists if and only if $E$ is trivial.

For a frame bundle $F(E)$, the following things stand out: The typical fiber is the structure group GL(V), and we constructed an atlas which showed that $F(E)$ has a GL(V)-bundle structure where the action is left multiplication. Furthermore there is a right action of $\mathrm{GL}(\mathrm{V})$ on the total space $F(E)$ which has the fibers as orbits. The charts constructed were of the form $(U, \widetilde{\phi})$, where $\widetilde{\phi}$ is equivariant in the sense that $\widetilde{\phi}(u g)=\widetilde{\phi}(u) g$, where if $\widetilde{\phi}(u)=\left(p, g_{1}\right)$, then $\left(p, g_{1}\right) g:=\left(p, g_{1} g\right)$ by definition. These facts motivate the concept of a principal bundle:

Definition 6.72. Let $\wp: P \rightarrow M$ be a smooth fiber bundle with typical fiber a Lie group $G$. The bundle ( $P, \wp, M, G$ ) is called a principal $G$-bundle if there is a smooth free right action of $G$ on $P$ such that\\
(i) The action preserves fibers; $\wp(u g)=\wp(u)$ for all $u \in P$ and $g \in G$.\\
(ii) For each $p \in M$, there exists a bundle chart ( $U, \phi$ ) with $p \in U$ and such that if $\phi=(\wp, \Phi)$, then

$$
\Phi(u g)=\Phi(u) g
$$

for all $u \in \wp^{-1}(U)$ and $g \in G$.\\
If the group $G$ is understood, then we may refer to ( $P, \wp, M, G$ ) simply as a principal bundle. Define a right action on $U \times G$ by $\left(p, g_{1}\right) g=$ ( $p, g_{1} g$ ). Then $U \times G \rightarrow U$ is a trivial principal bundle. If $\phi=(\wp, \Phi)$, then using this right action on $U \times G$ we have that $\phi(u g)=\phi(u) g$ if and only if $\Phi(u g)=\Phi(u) g$. Charts of the form described in (ii) of the definition are called principal bundle charts, and an atlas consisting of principal bundle charts is called a principal bundle atlas.

Proposition 6.73. If ( $P, \wp, M, G$ ) is a principal $G$-bundle, then the fibers are exactly the orbits of the right $G$-action.

Proof. The definition makes it clear that each orbit is contained in some fiber. Suppose that $u_{1}$ and $u_{2}$ are in the same fiber so that $\wp\left(u_{1}\right)=\wp\left(u_{2}\right)$. We wish to find a $g \in G$ such that $u_{1}=u_{2} g$. Let $g:=\Phi\left(u_{2}\right)^{-1} \Phi\left(u_{1}\right)$. Then $\Phi\left(u_{1}\right)=\Phi\left(u_{2}\right) g$ and so

$$
\begin{aligned}
\phi\left(u_{1}\right) & =\left(\wp\left(u_{1}\right), \Phi\left(u_{1}\right)\right)=\left(\wp\left(u_{2}\right), \Phi\left(u_{2}\right) g\right) \\
& =\left(\wp\left(u_{2} g\right), \Phi\left(u_{2} g\right)\right)=\phi\left(u_{2} g\right) .
\end{aligned}
$$

Since $\phi$ is bijective, we see that $u_{1}=u_{2} g$. The conclusion can be expressed by saying that the action is transitive on fibers.

The definition of a principal $G$-bundle above does not use the notion of a $G$-atlas or strict equivalence. The definition of principal bundle atlas is given after the notion of a principal bundle is already defined. However, once we have singled out this type of atlas, we can find out what the transition functions are and thereby make connection with earlier developments. Notice that if ( $\phi_{\alpha}, U_{\alpha}$ ) and ( $\phi_{\beta}, U_{\beta}$ ) are overlapping principal bundle charts with $\phi_{\alpha}=\left(\wp, \Phi_{\alpha}\right)$ and $\phi_{\beta}=\left(\wp, \Phi_{\beta}\right)$, then

$$
\Phi_{\alpha}(u g) \Phi_{\beta}(u g)^{-1}=\Phi_{\alpha}(u) g g^{-1} \Phi_{\beta}(u)^{-1}=\Phi_{\alpha}(u) \Phi_{\beta}(u)^{-1}
$$

so that the map $u \mapsto \Phi_{\alpha}(u) \Phi_{\beta}(u)^{-1}$ is constant on fibers. This means that there is a smooth function $g_{\alpha \beta}: U_{\alpha} \cap U_{\beta} \rightarrow G$ such that

$$
g_{\alpha \beta}(p)=\Phi_{\alpha}(u) \Phi_{\beta}(u)^{-1}
$$

where $u$ is any element in the fiber at $p$.\\
Lemma 6.74. Let ( $\phi_{\alpha}, U_{\alpha}$ ) and ( $\phi_{\beta}, U_{\beta}$ ) be overlapping principal bundle charts. For each $p \in U_{\alpha} \cap U_{\beta}$,

$$
\left.\Phi_{\alpha}\right|_{\wp^{-1}(p)} \circ\left(\left.\Phi_{\beta}\right|_{\wp^{-1}(p)}\right)^{-1}(g)=g_{\alpha \beta}(p) g
$$

where the $g_{\alpha \beta}$ are given as above.\\
Proof. Let $\left(\left.\Phi_{\beta}\right|_{\wp^{-1}(p)}\right)^{-1}(g)=u$. Then $g=\Phi_{\beta}(u)$ and so $\left.\Phi_{\alpha}\right|_{\wp^{-1}(p)} \circ \left.\Phi_{\beta}\right|_{\wp^{-1}(p)}(g)=\Phi_{\alpha}(u)$. On the other hand, $u \in \wp^{-1}(p)$ and so

$$
\begin{aligned}
g_{\alpha \beta}(p) g & =\Phi_{\alpha}(u) \Phi_{\beta}(u)^{-1} g=\Phi_{\alpha}(u) \Phi_{\beta}(u)^{-1} \Phi_{\beta}(u) \\
& =\Phi_{\alpha}(u)=\left.\left.\Phi_{\alpha}\right|_{\wp^{-1}(p)} \circ \Phi_{\beta}\right|_{\wp^{-1}(p)}(g)
\end{aligned}
$$

From this lemma we see that the structure group of a principal bundle is $G$ acting on itself by left translation. Conversely if $(P, \wp, M, G)$ is a fiber bundle with a $G$-atlas with $G$ acting by left translation, then ( $P, \wp, M, G$ ) is a principal bundle. To see this, we only need to exhibit the free right action. Let $u \in P$ and choose a chart ( $\phi_{\alpha}, U_{\alpha}$ ) from the $G$-atlas. Then let $u g:= \phi_{\alpha}^{-1}\left(p, \Phi_{\alpha}(u) g\right)$ where $p=\wp(u)$. We need to show that this is well-defined, so let $\left(\phi_{\beta}, U_{\beta}\right)$ be another such bundle chart with $p=\wp(u) \in U_{\beta} \cap U_{\alpha}$. Then if $u_{1}:=\phi_{\beta}^{-1}\left(p, \Phi_{\beta}(u) g\right)$, we have

$$
\phi_{\alpha}\left(u_{1}\right)=\phi_{\alpha} \phi_{\beta}^{-1}\left(p, \Phi_{\beta}(u) g\right)=\left(p, g_{\alpha \beta}(p) \Phi_{\beta}(u) g\right)=\left(p, \Phi_{\alpha}(u) g\right)
$$

so that $u_{1}=\phi_{\alpha}^{-1}\left(p, \Phi_{\alpha}(u) g\right)=u g$. It is easy to see that this action is free. Furthermore, since

$$
\phi_{\alpha}^{-1}\left(p, \Phi_{\alpha}(u g)\right)=\phi_{\alpha}^{-1} \circ \phi_{\alpha}(u g)=u g:=\phi_{\alpha}^{-1}\left(p, \Phi_{\alpha}(u) g\right),
$$

we see that $\Phi_{\alpha}(u g)=\Phi_{\alpha}(u) g$ as required by the definition of principal bundle. Obviously the frame bundles of vector bundles are examples of principal bundles.

Definition 6.75. Suppose that $\pi: E \rightarrow M$ is a rank $k$ vector bundle with typical fiber V. Let $G$ be a Lie subgroup of GL(V) and suppose that $\pi: E \rightarrow M$ has a $G$-bundle structure where $G$ acts on V by the standard action as a subgroup of $\mathrm{GL}(\mathrm{V})$. Thus we have a reduction of the structure group to $G$. Let $\mathcal{A}_{G}=\left\{\left(U_{\alpha}, \phi_{\alpha}\right)\right\}$ be the maximal $G$-atlas that defines this structure. For each $p \in M$, let

$$
F_{G}\left(E_{p}\right):=\left\{u \in \mathrm{GL}\left(\mathrm{V}, E_{p}\right): u=\phi^{-1}(p, \cdot) \text { for some }(U, \phi) \in \mathcal{A}_{G}\right\}
$$

Elements of $F_{G}\left(E_{p}\right)$ are called $G$-frames at $p$ (associated to $\mathcal{A}_{G}$ ).\\
The group $G$ acts on the right on $F_{G}\left(E_{p}\right)$ by $(u, g) \mapsto u \circ g$, and letting $F_{G}(E):=\bigcup_{p \in M} F_{G}\left(E_{p}\right)$ we see that $G$ acts on the right on $F_{G}(E)$. It is not hard to show that $F_{G}(E)$ has the structure of a principal $G$-bundle with this action. It is a subprincipal bundle of the frame bundle.

Definition 6.76. The bundle $F_{G}(E)$ is called the bundle of $G$-frames associated to the $G$-bundle structure on $\pi: E \rightarrow M$.

Actually, $G$ acts on the whole frame bundle of $E$ by the same formula and on the set of all frames $F\left(E_{p}\right)$ at a fixed $p$. Then, $F_{G}\left(E_{p}\right)$ is an orbit of this action on $F(E)$. In fact, one may simply define a $G$-bundle structure on a vector bundle to be a subbundle of the frame bundle such that $G$ acts transitively on each fiber. It is not hard to see that the notion of a bundle of $G$-frames gives us another way to describe the notion of reduction of the structure group. It is also important to notice that since there may be more than one $G$-bundle structure on $E$, the notation $F_{G}(E)$ is ambiguous. If one must deal with two different $G$-bundle structures on $E$, then one must resort to another notation such as $F_{G}^{1}(E)$ and $F_{G}^{2}(E)$.

Exercise 6.77. Show that a choice of metric on a vector bundle $E \rightarrow M$ is equivalent to a specification of a subbundle of the frame bundle such that $O(k)$ acts transitively on each fiber of the subbundle.

As another example of a principal bundle we also have the Hopf bundles described in the next example and the following exercise.

Example 6.78 (Hopf bundles). Recall the Hopf map $\wp: S^{2 n-1} \rightarrow \mathbb{C} P^{n-1}$ defined in Example 5.120. The quadruple ( $S^{2 n-1}, \wp, \mathbb{C} P^{n-1}, \mathrm{U}(1)$ ) is a principal fiber bundle. We have already defined the left action of $\mathrm{U}(1)$ on $S^{2 n-1}$ in Example 5.120. Since $\mathrm{U}(1)$ is abelian, we may take this action to also be a right action. Recall that in this context, we have $S^{2 n-1}=\left\{\xi \in \mathbb{C}^{n}:|\xi|=1\right\}$,\\
where for $\xi=\left(z^{1}, \ldots, z^{n}\right)$, we have $|\xi|^{2}=\sum \bar{z}^{i} z^{i}$. The right action of $\mathrm{U}(1)=S^{1}$ on $S^{2 n-1}$ is $(\xi, g) \mapsto \xi g=\left(z^{1} g, \ldots, z^{n} g\right)$. It is clear that $\wp(\xi g)=\wp(\xi)$. To finish the verification that ( $S^{2 n-1}, \wp, \mathbb{C} P^{n-1}, \mathrm{U}(1)$ ) is a principal bundle, we exhibit appropriate principal bundle charts. For each $k=1,2, \ldots, n$, we let $U_{k}:=\left\{\left[z^{1}, \ldots, z^{n}\right] \in \mathbb{C} P^{n-1}: z^{k} \neq 0\right\}$ and we let $\psi_{k}: \wp^{-1}\left(U_{k}\right) \rightarrow U_{k} \times \mathrm{U}(1)$ be defined by $\psi_{k}:=\left(\wp, \Psi_{k}\right)$, where

$$
\Psi_{k}(\xi)=\Psi_{k}\left(z^{1}, \ldots, z^{n}\right):=\left|z^{k}\right|^{-1} z^{k}
$$

We leave it to the reader to show that $\psi_{k}:=\left(\wp, \Psi_{k}\right)$ is a diffeomorphism. For $g \in \mathrm{U}(1)$, we have

$$
\Psi_{k}(\xi g)=\left|z^{k} g\right|^{-1}\left(z^{k} g\right)=\left|z^{k}\right|^{-1}\left(z^{k} g\right)=\left(\left|z^{k}\right|^{-1} z^{k}\right) g=\Psi_{k}(\xi) g
$$

as desired. Let us compute the transition cocycle $\left\{g_{i j}\right\}$. For $p=[\xi] \in U_{i} \cap U_{i}$, we have

$$
g_{i j}(p)=\Psi_{i}(\xi) \Psi_{j}(\xi)^{-1}=\left|z^{i}\right|^{-1} z^{i}\left(z^{j}\right)^{-1}\left|z^{j}\right| \in \mathrm{U}(1)
$$

Exercise 6.79. By analogy with the above example, show that we have principal bundles ( $S^{n-1}, \wp, \mathbb{R} P^{n-1}, \mathbb{Z}_{2}$ ) and ( $S^{4 n-1}, \wp, \mathbb{H} P^{n-1}, \mathrm{U}(1, \mathbb{H})$ ). Show that in the quaternionic case $g_{i j}(p)=\left|q^{i}\right|^{-1} q^{i}\left(q^{j}\right)^{-1}\left|q^{j}\right|$ for $p= \left[q^{1}, \ldots, q^{n}\right]$ and that the order matters in this case.

If ( $U, \phi$ ) is a principal bundle chart for a principal bundle ( $P, \wp, M, G$ ), then for each fixed $g \in G$, the map $\sigma_{\phi, g}: p \mapsto \phi^{-1}(p, g)$ is a smooth local section. The canonical choice for the fixed element is the identity $e$. Thus to each principal bundle chart ( $U, \phi$ ) we associate the local section $\sigma_{\phi}$ : $p \mapsto \phi^{-1}(p, e)$. Conversely if $\sigma: U \rightarrow P$ is a smooth local section, we let $f_{\sigma}: U \times G \rightarrow \wp^{-1}(U)$ be defined by $f_{\sigma}(p, g)=\sigma(p) g$. The proof of the next proposition shows that $f_{\sigma}$ is a diffeomorphism and its inverse $\phi: \wp^{-1}(U) \rightarrow U \times G$ defines a principal bundle chart. This is a principle worth emphasizing: Local sections of a principal bundle give rise to associated principal bundle charts, and conversely as described above.

Proposition 6.80. If $\wp: P \rightarrow M$ is a surjective submersion and a Lie group $G$ acts freely on $P$ so that for each $p \in M$ the orbit of $p$ is exactly $\wp^{-1}(p)$, then ( $P, \wp, M, G$ ) is a principal bundle.

Proof. Let us assume (without loss of generality) that the action is a right action since it can always be converted into such by group inversion if needed. We use Proposition 3.25: For each point $p \in M$, there is a local section $\sigma: U \rightarrow P$ on some neighborhood $U$ containing $p$. Consider the map $f_{\sigma}: U \times G \rightarrow \wp^{-1}(U)$ given by $f_{\sigma}(p, g)=\sigma(p) g$. One can check that this map is injective and has an invertible tangent map at each point of $U$. Now\\
let $\phi:=f_{\sigma}^{-1}$. Then we have $\phi=(\wp, \Phi)$ for a uniquely determined smooth map $\Phi: U \rightarrow G$. If $p=\wp(u)$, we have $\phi(u g)=(p, \Phi(u g))$ and so

$$
u g=\phi^{-1}(p, \Phi(u g))
$$

while

$$
\begin{aligned}
\phi^{-1}(p, \Phi(u) g) & =f_{\sigma}(p, \Phi(u) g) \\
& =\sigma(p)(\Phi(u) g)=(\sigma(p) \Phi(u)) g \\
& =f_{\sigma}(p, \Phi(u)) g=\phi^{-1}(p, \Phi(u)) g \\
& =u g=\phi^{-1}(p, \Phi(u g))
\end{aligned}
$$

Since $\phi^{-1}$ is a bijection, we have $\Phi(u g)=\Phi(u) g$. Thus the section $\sigma$ gives rise to a principal bundle chart ( $U, \phi$ ), where $\phi=(\pi, \Phi)$.

Combining this with our results on proper free actions from Chapter 5, we obtain the following corollary:

Corollary 6.81. If a Lie group $G$ acts properly and freely on $M$ (on the right), then ( $M, \pi, M / G, G$ ) is a principal bundle. In particular, if $H$ is a closed subgroup of a Lie group $G$, then ( $G, \pi, G / H, H$ ) is a principal bundle (with structure group $H$ ).\\

Definition 6.82. Let ( $P_{1}, \wp_{1}, M_{1}, G$ ) and ( $P_{2}, \wp_{2}, M_{2}, G$ ) be two principal $G$-bundles. A (type II) bundle morphism $\widetilde{f}: P_{1} \rightarrow P_{2}$ along a smooth map $f: M_{1} \rightarrow M_{2}$ is called a principal $G$-bundle morphism (along $f$ ) if

$$
\widetilde{f}(u \cdot g)=\widetilde{f}(u) \cdot g
$$

for all $g \in G$ and $u \in P$. If $M_{1}=M_{2}$ and $f=\operatorname{id}_{M}$, then we say that $\widetilde{f}$ is a principal $G$-bundle morphism over $M$.

Exercise 6.83. Show that if ( $P_{1}, \wp_{1}, M_{1}, G$ ) and ( $P_{2}, \wp_{2}, M_{2}, G$ ) are principal $G$-bundles and $\widetilde{f}: P_{1} \rightarrow P_{2}$ is a principal $G$-bundle morphism along a diffeomorphism $f$, then $\widetilde{f}$ is a diffeomorphism.

If $\widetilde{f}$ is a principal bundle morphism over $M$, then it is a diffeomorphism and hence a bundle equivalence (or bundle isomorphism over $M$ ) with the property $\widetilde{f}(u \cdot g)=\widetilde{f}(u) \cdot g$ for all $g \in G$ and $u \in P$. In this case, we call $\widetilde{f}$ a principal $G$-bundle equivalence and the two bundles are equivalent principal $G$-bundles over $M$. A principal $G$-bundle equivalence from a principal bundle to itself is called a principal bundle automorphism or also a (global) gauge transformation.

The classification problem for principal bundles (in the topological category) is regrettably beyond the scope of this volume and can be found in [Hus]. We can only offer the following comments: In the topological category, the notion of a Lie group is replaced by that of a topological group,\\
but the reader may continue to think of Lie groups. For every topological group $G$, there is a principal bundle $\xi(G)=\left(\left(E(G), \wp_{\infty}, B(G), G\right)\right.$ called a universal bundle with the property that all the homotopy groups of $E(G)$ are trivial. There is then a classification theorem that states that the equivalence classes of principal bundles with a fixed sufficiently nice ${ }^{1}$ base space $M$ are in one-to-one correspondence with the set of homotopy classes of maps from $M$ to $B(G)$. The correspondence is given by assigning to the homotopy class $[f]$, the pull-back principal bundle $f^{*} \xi(G)$.

The notion of a principal bundle morphism over $M$ can be generalized to the situation where we have two groups in play.\\
Definition 6.84. Let ( $P_{1}, \wp_{1}, M, G_{1}$ ) and ( $P_{2}, \wp_{2}, M, G_{2}$ ) be principal bundles and let $h: G_{1} \rightarrow G_{2}$ be a Lie group homomorphism. A bundle morphism over $M$ is called a principal bundle homomorphism with respect to $h$ if

$$
\widetilde{f}(u \cdot g)=\widetilde{f}(u) \cdot h(g)
$$

If $G_{1} \subset G_{2}$ and the homomorphism $h$ is the inclusion, then we call $\widetilde{f}$ a reduction of ( $P_{2}, \wp_{2}, M, G_{2}$ ) to ( $P_{1}, \wp_{1}, M, G_{1}$ ).

In the most common case of a reduction, $P_{1}$ is a submanifold of $P_{2}$ and $\widetilde{f}$ is the inclusion $P_{1} \hookrightarrow P_{2}$. For example, this is the case when one chooses a metric on a vector bundle and thereby obtains a bundle of orthonormal frames $F_{\mathrm{O}(n)}(E)$. The inclusion $F_{\mathrm{O}(n)}(E) \hookrightarrow F(E)$ is then a reduction, and we just say that $F_{\mathrm{O}(n)}(E)$ is a reduction of the frame bundle $F(E)$.

We have seen that a principal $G$-bundle atlas $\left\{\left(U_{\alpha}, \phi_{\alpha}\right)\right\}$ is associated to a cocycle $\left\{g_{\alpha \beta}\right\}$. From this cocycle and the left action of $G$ on itself we may construct a bundle which has $\left\{g_{\alpha \beta}\right\}$ as a transition cocycle. In fact, recall that in the construction we formed the total space by putting an equivalence relation on the set $\Sigma:=\bigcup_{\alpha}\{\alpha\} \times U_{\alpha} \times G$, where $(\alpha, p, g) \in\{\alpha\} \times U_{\alpha} \times G$ is equivalent to $\left(\beta, p^{\prime}, g^{\prime}\right) \in\{\beta\} \times U_{\beta} \times G$ if and only if $p=p^{\prime}$ and $g^{\prime}=g_{\beta \alpha}(p) \cdot g$. If we define a right action on the total space of the constructed bundle by $\left[\alpha, p, g_{1}\right] \cdot g=\left[\alpha, p, g_{1} g\right]$, then this is well-defined, smooth, and makes the constructed bundle a principal $G$-bundle equivalent to the original principal $G$-bundle.

Exercise 6.85. Prove the last assertion above.\\
Thus we see that $G$-cocycles on a smooth manifold $M$ give rise to principal $G$-bundles and conversely. If we start with two $G$-cocycles on $M$, then we may ask whether the principal $G$-bundles constructed from these cocycles are equivalent or not. First notice that the constructed bundles will have principal bundle atlases with the respective original transition cocycles. Thus we are led to the following related question: What conditions on

\footnotetext{${ }^{1} M$ should be a CW-complex.
}
the transition cocycles arising from principal bundle atlases on two principal $G$-bundles will ensure that the bundles are equivalent principal $G$-bundles? By restricting the trivializing maps to open sets of a common refinement, we obtain new atlases and so we may as well assume from the start that the respective principal bundle atlases are defined on the same cover of $M$.

Theorem 6.86. Let ( $P_{1}, \wp_{1}, M, G$ ) and ( $P_{2}, \wp_{2}, M, G$ ) be principal $G$ bundles with principal bundle atlases $\left\{\left(\phi_{\alpha}, U_{\alpha}\right)\right\}$ and $\left\{\left(\phi_{\alpha}^{\prime}, U_{\alpha}\right)\right\}$ respectively. Then ( $P_{1}, \wp_{1}, M, G$ ) is equivalent to ( $P_{2}, \wp_{2}, M, G$ ) if and only if there exists a family of (smooth) maps $\tau_{\alpha}: U_{\alpha} \rightarrow G$ such that $g_{\alpha \beta}^{\prime}(p)= \left(\tau_{\alpha}(p)\right)^{-1} g_{\alpha \beta}(p) \tau_{\beta}(p)$ for all $p \in U_{\alpha} \cap U_{\beta}$ and for all nonempty intersections $U_{\alpha} \cap U_{\beta}$. (Here $\left\{g_{\alpha \beta}\right\}$ is the cocycle associated to $\left\{\left(\phi_{\alpha}, U_{\alpha}\right)\right\}$ and $\left\{g_{\alpha \beta}^{\prime}\right\}$ is the cocycle associated to $\left\{\left(\phi_{\alpha}^{\prime}, U_{\alpha}\right)\right\}$.)

Sketch of proof. First suppose that $P_{1}$ and $P_{2}$ are equivalent principal $G$ bundles and let $\widetilde{f}: P_{1} \rightarrow P_{2}$ be an equivalence. Let $p \in U_{\alpha}$ and choose some $u \in \wp_{1}^{-1}(p)$, so $\widetilde{f}(u) \in \wp_{2}^{-1}(p)$. Write $\phi_{\alpha}=\left(\wp_{1}, \Phi_{\alpha}\right)$ and $\phi_{\alpha}^{\prime}=\left(\wp_{2}, \Phi_{\alpha}^{\prime}\right)$. One can easily show that $\Phi_{\alpha}(u)\left(\Phi_{\alpha}^{\prime}(\widetilde{f}(u))\right)^{-1}$ is an element of $G$ that is independent of the choice of $u \in \wp_{1}^{-1}(p)$. For each $\alpha$, define $\tau_{\alpha}: U_{\alpha} \rightarrow G$ by

$$
\tau_{\alpha}(p):=\Phi_{\alpha}(u)\left(\Phi_{\alpha}^{\prime}(\widetilde{f}(u))\right)^{-1}
$$

where $u \in \wp_{1}^{-1}(p)$. Suppose that $p \in U_{\alpha} \cap U_{\beta}$. Then we have $\left(\tau_{\alpha}(p)\right)^{-1}= \Phi_{\alpha}^{\prime}(\widetilde{f}(u))\left(\Phi_{\alpha}(u)\right)^{-1}$. Using the definitions of $g_{\alpha \beta}$ and $g_{\alpha \beta}^{\prime}$ (see equation (6.3)), we immediately have

$$
g_{\alpha \beta}^{\prime}(p)=\left(\tau_{\alpha}(p)\right)^{-1} g_{\alpha \beta}(p) \tau_{\beta}(p)
$$

Conversely, given the maps $\tau_{\alpha}: U_{\alpha} \rightarrow G$ satisfying

$$
g_{\alpha \beta}^{\prime}(p)=\left(\tau_{\alpha}(p)\right)^{-1} g_{\alpha \beta}(p) \tau_{\beta}(p)
$$

we define, for each $\alpha$, a map $f_{\alpha}: \wp_{1}^{-1}\left(U_{\alpha}\right) \rightarrow \wp_{2}^{-1}\left(U_{\alpha}\right)$ by

$$
f_{\alpha}(u):=\left(\phi_{\alpha}^{\prime}\right)^{-1}\left(p,\left(\tau_{\alpha}(p)\right)^{-1} \Phi_{\alpha}(u)\right) .
$$

Check that $f_{\alpha}(u)=f_{\beta}(u)$ when $\wp_{1}(u) \in U_{\alpha} \cap U_{\beta}$ so that there is a welldefined map $\widetilde{f}: P_{1} \rightarrow P_{2}$ such that $f_{\alpha}(u)=\widetilde{f}(u)$ whenever $\wp_{1}(u) \in U_{\alpha}$. Finally, check that $\widetilde{f}(u \cdot g)=\widetilde{f}(u) \cdot g$.

Let $\wp: P \rightarrow M$ be a principal $G$-bundle and suppose that we are given a smooth left action $\lambda: G \times F \rightarrow F$ on some smooth manifold $F$. Define a right action of $G$ on $P \times F$ according to

$$
(u, y) \cdot g:=\left(u g, g^{-1} y\right)=\left(u g, \lambda\left(g^{-1}, y\right)\right)
$$

Denote the orbit space of this action by $P \times_{\lambda} F$ (or $P \times_{G} F$ ) and let $\widetilde{\wp}$ denote the quotient map. Also denote the equivalence class of $(u, y)$ by $[u, y]$ so\\
$\widetilde{\wp}(u, y)=[u, y]$. One may check that there is a unique map $\pi: P \times_{\lambda} F \rightarrow M$ such that $\pi([u, y])=\wp(u)$, and so we have a commutative diagram:\\
\includegraphics[scale=0.25, center]{2025_10_09_265d7e1a22c89f79c21eg-315}

Next we show that ( $P \times_{\lambda} F, \pi, M, F$ ) is a fiber bundle (a ( $G, \lambda$ )-bundle). It is said to be associated to the principal bundle $P$. Bundles constructed in this way are called associated bundles. More precisely, if the action $\lambda$ is not effective, we should say that $P \times_{\lambda} F$ is weakly associated to $P$. This is what lies behind our previously introduced notion of "ineffective $(G, \lambda)$-bundle".

Theorem 6.87. Referring to the above diagram and notations, $P \times_{\lambda} F$ is a smooth manifold and the following hold:\\
(i) ( $P \times_{\lambda} F, \pi, M, F$ ) is a fiber bundle, and for every principal bundle atlas $\left\{\left(U_{\alpha}, \phi_{\alpha}\right)\right\}$, there is a corresponding bundle atlas $\left\{\left(U_{\alpha}, \widetilde{\phi}_{\alpha}\right)\right\}$ for $P \times_{\lambda} F$ such that\\
$\widetilde{\phi}_{\alpha} \circ \widetilde{\phi}_{\beta}^{-1}(p, y)=\left(p, \lambda\left(g_{\alpha \beta}(p), y\right)\right)$ if $p \in U_{\alpha} \cap U_{\beta}$ and $y \in F$,\\
where the $g_{\alpha \beta}$ are defined by equation (6.3).\\
(ii) $\left(P \times F, \widetilde{\wp}, P \times{ }_{\lambda} F, G\right)$ is a principal bundle with the right action given by

$$
(u, y) \cdot g:=\left(u g, g^{-1} y\right) .
$$

(iii) $P \times F \xrightarrow{\mathrm{pr}_{1}} P$ is a principal bundle morphism along $\pi$.

Proof. Let $\left\{\left(U_{\alpha}, \phi_{\alpha}\right)\right\}$ be a principal bundle atlas for $\wp: P \rightarrow M$. Note that $\widetilde{\wp}\left(\wp^{-1}\left(U_{\alpha}\right) \times F\right)=\pi^{-1}\left(U_{\alpha}\right)$. For each $\alpha$, define $\widetilde{\Phi}_{\alpha}: \pi^{-1}\left(U_{\alpha}\right) \rightarrow F$ by requiring that $\widetilde{\Phi}_{\alpha} \circ \widetilde{\wp}(u, y)=\Phi_{\alpha}(u) \cdot y$ for all $(u, y) \in \wp_{\sim}^{-1}\left(U_{\alpha}\right) \times F$ and then let $\widetilde{\phi}_{\alpha}:=\left(\pi, \widetilde{\Phi}_{\alpha}\right)$ on $\pi^{-1}\left(U_{\alpha}\right)$. We want to show that $\widetilde{\phi}_{\alpha}$ is bijective by defining an inverse for $\widetilde{\phi}_{\alpha}$. For every $p \in U_{\alpha}$, let $\sigma_{\alpha}(p):=\phi_{\alpha}^{-1}(p, e)$, where $e$ is the identity element in $G$. Then we have

$$
\sigma_{\alpha}(p) \cdot \Phi_{\alpha}(u)=\phi_{\alpha}^{-1}(p, e) \cdot \Phi_{\alpha}(u)=\phi_{\alpha}^{-1}\left(p, \Phi_{\alpha}(u)\right)=u
$$

Define $\eta_{\alpha}: U_{\alpha} \times F \rightarrow \pi^{-1}\left(U_{\alpha}\right)$ by $\eta_{\alpha}(p, y):=\widetilde{\wp}\left(\sigma_{\alpha}(p), y\right)$. We have

$$
\begin{aligned}
\eta_{\alpha} \circ \widetilde{\phi}_{\alpha}(\widetilde{\wp}(u, y)) & =\eta_{\alpha}\left(p, \Phi_{\alpha}(u) \cdot y\right)=\widetilde{\wp}\left(\sigma_{\alpha}(p), \Phi_{\alpha}(u) \cdot y\right) \\
& =\widetilde{\wp}\left(\sigma_{\alpha}(p) \cdot \Phi_{\alpha}(u), y\right)=\widetilde{\wp}(u, y)
\end{aligned}
$$

Thus $\eta_{\alpha}$ is a right inverse for $\widetilde{\phi}_{\alpha}$ and so $\widetilde{\phi}_{\alpha}$ is injective. It is easily checked that $\eta_{\alpha}$ is also a left inverse for $\phi_{\alpha}$. To see this first note that $\left(p, \Phi_{\alpha}\left(\sigma_{\alpha}(p)\right)\right)=$

$$
\begin{aligned}
& \phi_{\alpha}\left(\sigma_{a}(p)\right)=(p, e) \text { so } \Phi_{\alpha}\left(\sigma_{\alpha}(p)\right)=e . \text { Thus we have } \\
& \qquad \begin{aligned}
\widetilde{\phi}_{\alpha} \circ \eta_{\alpha}(p, y) & =\widetilde{\phi}_{\alpha}\left(\widetilde{\wp}\left(\sigma_{\alpha}(p), y\right)\right)=\left(p, \widetilde{\Phi}_{\alpha}\left(\widetilde{\wp}\left(\sigma_{\alpha}(p), y\right)\right)\right) \\
& =\left(p, \Phi_{\alpha}\left(\sigma_{\alpha}(p)\right) \cdot y\right)=(p, y)
\end{aligned}
\end{aligned}
$$

Thus $\widetilde{\phi}_{\alpha}$ is a bijection. Next we check the overlaps. We use Lemma 6.74;

$$
\begin{aligned}
\widetilde{\phi}_{\alpha} \circ \widetilde{\phi}_{\beta}^{-1}(p, y) & =\widetilde{\phi}_{\alpha} \circ \eta_{\beta}(p, y)=\widetilde{\phi}_{\alpha}\left(\widetilde{\wp}\left(\sigma_{\beta}(p), y\right)\right) \\
& =\left(p, \Phi_{\alpha}\left(\sigma_{\beta}(p)\right) \cdot y\right)=\left(p, \Phi_{\alpha}\left(\phi_{\beta}^{-1}(p, e)\right) \cdot y\right) \\
& \left.=\left(p,\left.\left.\Phi_{\alpha}\right|_{p} \circ \Phi_{\beta}\right|_{p} ^{-1}(e)\right) \cdot y\right) \\
& =\left(p, g_{\alpha \beta}(p) \cdot e \cdot y\right)=\left(p, g_{\alpha \beta}(p) y\right)
\end{aligned}
$$

This shows that the transitions mappings have the stated form and that the overlap maps $\widetilde{\phi}_{\alpha} \circ \widetilde{\phi}_{\beta}^{-1}$ are smooth. The family $\left\{\left(U_{\alpha}, \phi_{\alpha}\right)\right\}$ provides the induced smooth structure and is also a bundle atlas. Since $\phi_{\alpha} \circ \widetilde{\wp}(u, p)= \left(\pi, \widetilde{\Phi}_{\alpha}\right) \circ \widetilde{\wp}(u, p)=\left(\wp(u), \Phi_{\alpha}(u) y\right)$ in the domain of every bundle chart ( $U_{\alpha}, \phi_{\alpha}$ ), it follows that $\widetilde{\wp}$ is smooth.

We leave it to the reader to verify that ( $P \times F, \widetilde{\wp}, P \times{ }_{\lambda} F, G$ ) is a principal $G$-bundle. Notice that while the map $\operatorname{pr}_{1}: P \times F \rightarrow P$ is clearly a bundle map along $\pi$, we also have

$$
\operatorname{pr}_{1}((u, y) \cdot g)=\operatorname{pr}_{1}\left(\left(u \cdot g, g^{-1} y\right)\right)=u \cdot g=\operatorname{pr}_{1}(u, y) \cdot g
$$

and so $\mathrm{pr}_{1}$ is in fact a principal bundle morphism.\\
Clearly what we have is another way of looking at bundle construction. The principal bundle takes the place of the cocycle of transition maps.

Exercise 6.88. Construct a principal $\mathbb{Z}_{2}$-bundle $P$ and left actions $\lambda_{1}$ and $\lambda_{2}$ of $\mathbb{Z}_{2}$ on $S^{1}$ and $\mathbb{R}$ respectively, such that $P \times_{\lambda_{1}} S^{1}$ is the twisted torus and $P \times{ }_{\lambda_{2}} \mathbb{R}$ is the Mobius band line bundle.

We have seen that given a principal $G$-bundle, one may construct various fiber bundles with $G$-bundle structures. Let us look at the converse situation. Suppose that ( $E, \pi, M, F$ ) is a fiber bundle. Suppose that this bundle has a $(G, \lambda)$-atlas $\left\{\left(U_{\alpha}, \phi_{\alpha}\right)\right\}$ with associated $G$-valued cocycle of transition functions $\left\{g_{\alpha \beta}\right\}$. Using Theorem 6.15, one may construct a bundle with typical fiber $G$ by using left translation as the action. The resulting bundle is then a principal bundle ( $P, \wp, M, G$ ), and it turns out that $P \times_{\lambda} F$ is equivalent to the original bundle $E$.

If ( $E, \pi, M, \mathrm{V}$ ) is a vector bundle and we use the standard $\mathrm{GL}(\mathrm{V})$-cocycle $\left\{\Phi_{\alpha \beta}\right\}$ associated to a VB-atlas, then the principal bundle obtained by the above construction is (equivalent to) the linear frame bundle $F(E)$. Letting $\mathrm{GL}(\mathrm{V})$ act on V according to the standard action we have $F(E) \times{ }_{\mathrm{GL}(\mathrm{V})} \mathrm{V}$, which is equivalent to the original bundle ( $E, \pi, M, \mathrm{V}$ ). More generally, if\\
$\lambda: G \rightarrow \mathrm{GL}(\mathrm{V})$ is a Lie group representation, then by treating $\lambda$ as a linear action we can form $P \times_{\lambda} \mathrm{V}$.

Proposition 6.89. Let $P$ be a principal $G$-bundle and let $\lambda: G \rightarrow \mathrm{GL}(\mathrm{V})$ be a representation. Then $P \times_{\lambda} \mathrm{V}$ has a natural vector bundle structure with typical fiber V.

Proof. This follows from Theorem 6.87, but we can argue more directly. Let us denote the total space of $P \times_{\lambda} \mathrm{V}$ by $B$ and let $B_{p}$ be the fiber over some point $p \in M$. Then for each $u \in P_{p}$ there is a map $\psi_{u}:=[u, \cdot]: \mathrm{V} \rightarrow B_{p}$ given by $v \mapsto[u, v]$. We compare $\psi_{u}$ with $\psi_{u g}$ for $g \in G$ and $u \in P_{p}$. Since $[u g, v]=[u, \lambda(g) v]$ for all $v \in \mathrm{V}$, the following diagram commutes:\\
\includegraphics[scale=0.25, center]{2025_10_09_265d7e1a22c89f79c21eg-317(1)}

From this it follows that $\psi_{u}$ transfers the linear structure of V to $B_{p}$ independently of the choice of $u \in P_{p}$. We leave it to the reader to show that the local trivializations of $P \times_{\lambda} \mathrm{V}$ constructed as in the proof of Theorem 6.87 are linear on each fiber.

Example 6.90. Let $M$ be an $n$-manifold and let $F(M)$ be the frame bundle of $M$. Then, if $\lambda_{0}$ is the standard action of $\operatorname{GL}(n, \mathbb{R})$ on $\mathbb{R}^{n}$, we have the following vector bundle isomorphisms:

$$
\begin{aligned}
F(M) \times_{\lambda_{0}} \mathbb{R}^{n} & \cong T M \\
F(M) \times_{\lambda_{0}^{*}} \mathbb{R}^{n} & \cong T^{*} M \\
F(M) \times_{\lambda_{0} \otimes \lambda_{0}^{*}} \mathbb{R}^{n} & \cong T M \otimes T^{*} M
\end{aligned}
$$

If $E=P \times_{\lambda} \mathrm{V}$ is an associated vector bundle for $\lambda$ a representation, then we can map $P$ into the frame bundle of $E$. Indeed, the map is just $\psi: u \mapsto \psi(u)=\psi_{u}$, where $\psi_{u}:=[u, \cdot]$ as above. Furthermore, $\psi(u g)= \psi(u) \circ \lambda(g)$ and so we have a principal bundle morphism with respect to the homomorphism $\lambda$ :\\
\includegraphics[scale=0.25, center]{2025_10_09_265d7e1a22c89f79c21eg-317}

The map $\psi: P \rightarrow F(E)$ is only injective if the action $\lambda$ is effective.\\
Based on what we have seen above we can say that the theory of principal bundles and associated bundles is an alternative and "invariant" approach to $G$-bundles. By "invariant" we mean that the foundations can be laid out\\
without recourse to strict equivalence classes of $G$-atlases or the use of cocycles (of course, these notions can be brought in as convenient). According to this approach, the central notion is the principal bundle, and one recovers the other $G$-bundles of interest as associated bundles. Developing the theory in this way has the advantage that much can be accomplished without the direct need of bundle atlases. It is a more "intrinsic" approach. This approach seems to have originated with Ehresmann and is the approach followed by [Hus].

\section*{Problems}
(1) Show that $S^{n} \times \mathbb{R}$ and $S^{n} \times S^{1}$ are parallelizable.\\
(2) Let $X:=[0,1] \times \mathbb{R}^{n}$. Fix a linear isomorphism $L: \mathbb{R}^{n} \rightarrow \mathbb{R}^{n}$ and consider the quotient space $E=X / \sim$, where the equivalence relation is given by $(0, v) \sim(1, L v)$. Show that $E$ is the total space of a smooth vector bundle over the circle $S^{1}$.\\
(3) Exhibit the vector bundle charts for the pull-back bundle construction of Definition 6.18.\\
(4) Let $\xi=(E, \pi, M, F)$ be a $G$-bundle. Let $g_{\alpha \beta}$ be cocycles associated to a $G$-altas $\left\{\left(U_{\alpha}, \phi_{\alpha}\right)\right\}$ for $\xi$. Show that $\xi$ is $G$-equivalent to a product bundle if and only if there exist functions $\lambda_{\alpha}: U_{\alpha} \rightarrow G$ such that

$$
g_{\beta \alpha}(x)=\lambda_{\beta}(x) \lambda_{\alpha}^{-1}(x) \quad \text { for all } x \in U_{\alpha} \cap U_{\beta}
$$

and all $\alpha, \beta$.\\
(5) Show that the twisted torus of Example 6.17 is trivial as a fiber bundle but not trivial as a $\mathbb{Z}_{2}$-bundle. (Use Problem 4.)\\
(6) Show that the space of sections of a vector bundle over a compact base is a finitely generated module. Show that if the bundle is trivial, then the space of sections is a finitely generated free module.\\
(7) Let $E_{1} \rightarrow M_{1}$ and $E_{2} \rightarrow M_{2}$ be smooth vector bundles. Show that if $F: E_{1} \rightarrow E_{2}$ is a vector bundle homomorphism along a map $f: M_{1} \rightarrow M_{2}$ such that $F$ is an isomorphism of fibers, then $E_{1} \rightarrow M_{1}$ is isomorphic to the pull-back bundle $f^{*} E_{2} \rightarrow M_{1}$.\\
(8) Suppose that $\xi=(E, \pi, M)$ is a vector bundle with a positive definite metric. Show that the metric induces a vector bundle isomorphism $E \cong E^{*}$\\
(9) Show that the tangent bundle of the real projective plane is a vector bundle isomorphic to $\operatorname{Hom}\left(\mathbb{L}\left(\mathbb{R} P^{n}\right), \mathbb{L}\left(\mathbb{R} P^{n}\right)^{\perp}\right)$, where $\mathbb{L}\left(\mathbb{R} P^{n}\right) \rightarrow \mathbb{R} P^{n}$\\
is the tautological line bundle and $\mathbb{L}\left(\mathbb{R} P^{n}\right)^{\perp} \rightarrow \mathbb{R} P^{n}$ is the rank $n$ vector bundle whose fiber at $l \in \mathbb{R} P^{n}$ is $\left\{(l, v) \in \mathbb{R} P^{n} \times \mathbb{R}^{n+1}: v \perp l\right\}$.\\
(10) Recall Example 6.35. Show that the Whitney sum bundle $E_{1} \oplus E_{2}$ is naturally isomorphic to the pull-back $\Delta^{*} \pi_{E_{1} \times E_{2}}: \Delta^{*}\left(E_{1} \times E_{2}\right) \rightarrow M$.\\
(11) (a) Let $\pi: E \rightarrow M$ be an $\mathbb{F}$-vector bundle. We wish to show that $T \pi: T E \rightarrow T M$ is naturally a vector bundle. Consider the maps $\alpha: E \oplus E \rightarrow E$ and $\mu_{s}: E \rightarrow E$ for each $s \in \mathbb{F}$ given by

$$
\begin{aligned}
\alpha\left(v_{p}, w_{p}\right) & :=v_{p}+w_{p} \text { for } v_{p}, w_{p} \in E_{p} \\
\mu_{s}\left(e_{p}\right) & :=s e_{p} \text { for } e_{p} \in E_{p}
\end{aligned}
$$

Show that we may identify $T(E \oplus E)$ with the submanifold of $T E \times T E$ given by

$$
\{(v, w) \in T E \times T E: T \pi \cdot v=T \pi \cdot w\}
$$

Now suppose that for $v, w \in T E$ with $T \pi \cdot v=T \pi \cdot w$ we define $v \boxplus w:=T \alpha \cdot(v, w)$ and for $s \in \mathbb{F}$ and $v \in T E$ we define $s \odot v:= T \mu_{s} \cdot v$. Show that with these definitions of addition and scalar multiplication, $T \pi: T E \rightarrow T M$ is indeed an $\mathbb{F}$-vector bundle.\\
(b) Let $E$ be as above but assume for simplicity that $\mathbb{F}=\mathbb{R}$. Let $x^{1}, \ldots, x^{n}$ be coordinates on $U \subset M$. Suppose that $e_{1}, \ldots, e_{k}$ is a frame field over $U$. Let $\xi^{1}, \ldots, \xi^{n}$ be defined on $\left.E\right|_{U}$ by $y=\sum \xi^{i}(y) e^{i}(\pi(y))$ for any $y \in E$. Then, identifying $x^{i}$ with $x^{i} \circ \pi$, the functions $x^{1}, \ldots, x^{n}, \xi^{1}, \ldots, \xi^{n}$ are a coordinate system for $E$ defined on $\left.E\right|_{U}$ and such that the $\frac{\partial}{\partial \xi^{i}}$ are in the kernel of $T \pi$. Now if $v, w \in T E$ are such that $T \pi \cdot v=T \pi \cdot w$, then we may express $v$ and $w$ as

$$
v=\left.\sum_{i} a^{i} \frac{\partial}{\partial x^{i}}\right|_{y}+\left.\sum_{\alpha} b^{\alpha} \frac{\partial}{\partial \xi^{\alpha}}\right|_{y}
$$

and

$$
w=\left.\sum_{i} a^{i} \frac{\partial}{\partial x^{i}}\right|_{\widetilde{y}}+\left.\sum_{\alpha} \widetilde{b}^{\alpha} \frac{\partial}{\partial \xi^{\alpha}}\right|_{\widetilde{y}} .
$$

Here $y$ and $\widetilde{y}$ are the base points of $v$ and $w$, and the fact that the $a$ 's are the same for both $v$ and $w$ is a result of the condition $T \pi \cdot v=T \pi \cdot w$. Show that

$$
v \boxplus w=\left.\sum_{i} a^{i} \frac{\partial}{\partial x^{i}}\right|_{y+\widetilde{y}}+\left.\sum_{\alpha}\left(b^{\alpha}+\widetilde{b}^{\alpha}\right) \frac{\partial}{\partial \xi^{\alpha}}\right|_{y+\widetilde{y}},
$$

where in $v \boxplus w$, the $\boxplus$ refers to the addition described in part (a).\\
(12) Exhibit a VB-atlas for the tautological bundle of Example 6.47.\\
(13) Show that the tautological bundle over $\mathbb{R} P^{1}$ is a Möbius band.\\
(14) Let $P \rightarrow M$ be a principal bundle with group $G$. If $H$ is a Lie subgroup of $G$, then the quotient $P / H$ is an $H$-principal bundle. Show that $P / H \rightarrow M$ admits a global section if and only if the structure group of $P \rightarrow M$ is reducible to $H$.\\
(15) Show that the notions of smooth fiber bundle and vector bundle make sense when the base space is allowed to be a manifold with boundary. What issues arise if one considers allowing both the base space and typical fiber to have boundary?





\pagebreak


\section{Modules and Multilinearity}


A module is an algebraic object that shows up quite a bit in differential geometry and analysis (at least implicitly). A module is a generalization of a vector space where the field $\mathbb{F}$ is replaced by a ring or an algebra over a field. For the definition of ring and field consult any book on abstract algebra. The definition of algebra is given below. The modules that occur in differential geometry are almost always finitely generated projective modules over the algebra of $C^{r}$ functions, and these correspond to the spaces of $C^{r}$ sections of vector bundles. We give the abstract definitions but we ask the reader to keep two cases in mind. The first is just the vector spaces which are the fibers of vector bundles. In this case, the ring in the definition below is the field $\mathbb{F}$ (the real numbers $\mathbb{R}$ or the complex numbers $\mathbb{C}$ ), and the module is just a vector space. The second case, already mentioned, is where the ring is the algebra $C^{r}(M)$ for a $C^{r}$ manifold and the module is the set of $C^{r}$ sections of a vector bundle over $M$.

As we have indicated, a module is similar to a vector space with the differences stemming from the use of elements of a ring R as the scalars rather than the field of complex $\mathbb{C}$ or real numbers $\mathbb{R}$. For an element $v$ of a module V , one still has $0 v=0$, and if the ring is a ring with unity 1 , then we usually require $1 v=v$. Of course, every vector space is also a module since the latter is a generalization of the notion of vector space. We also have maps between modules, the module homomorphisms (see Definition D. 5 below), which make the class of modules and module homomorphisms into a category.

Definition D.1. 


\begin{DEF}{A-L09-04-01}{Links/Rechts Modul}
Let $R$ be a ring. A left $R$-module (or a left module over R ) is an abelian group ( $\mathrm{V},+$ ) together with an operation $\mathrm{R} \times \mathrm{V} \rightarrow \mathrm{V}$ written as $(a, v) \mapsto a v$ and such that
\begin{enumerate}
  \item $(a+b) v=a v+b v$ for all $a, b \in \mathrm{R}$ and all $v \in \mathrm{V}$;
  \item $a\left(v_{1}+v_{2}\right)=a v_{1}+a v_{2}$ for all $a \in \mathrm{R}$ and all $v_{2}, v_{1} \in \mathrm{V}$;\\
$3(a b) v=a(b v)$ for all $a, b \in \mathrm{R}$ and all $v \in \mathrm{V}$.
\end{enumerate}
A right R -module is defined similarly with the multiplication on the right so that
\begin{enumerate}
  \item $v(a+b)=v a+v b$ for all $a, b \in \mathrm{R}$ and all $v \in \mathrm{V}$;
  \item $\left(v_{1}+v_{2}\right) a=v_{1} a+v_{2} a$ for all $a \in \mathrm{R}$ and all $v_{2}, v_{1} \in \mathrm{V}$
  \item $v(a b)=(v a) b$ for all $a, b \in \mathrm{R}$ and all $v \in \mathrm{V}$.
\end{enumerate}
\end{DEF}


\begin{DEF}{A-L09-04-02}{Unitäres Modul}
If R has an identity and $1 v=v$ (or $v 1=v$ ) for all $v \in \mathrm{V}$, then we say that V is a unitary R-module.

If $R$ has an identity, then by " $R$-module" we shall always mean a unitary R-module unless otherwise indicated.
\end{DEF}

Also, if the ring is commutative (the usual case for us), then we may write $a v=v a$ and consider any right module as a left module and vice versa. Even if the ring is not commutative, we will usually stick to left modules in this appendix and so we drop the reference to "left" and refer to such as R-modules. We do also use right modules in the text. For example, we consider right modules over the quaternions.

Remark D.2. We shall often refer to the elements of $R$ as scalars.


Example D.3. 

\begin{EXA}{A-L09-04-03}{Abelsche Gruppe als $\Z$-Modul}
An abelian group ( $A,+$ ) is a $\mathbb{Z}$-module, and a $\mathbb{Z}$-module is none other than an abelian group. Here we take the product of $n \in \mathbb{Z}$ with $x \in A$ to be $n x:=x+\cdots+x$ if $n \geq 0$ and $n x:=-(x+\cdots+x)$ if $n<0$ (in either case we are adding $|n|$ terms).
\end{EXA}

Example D.4. 

\begin{EXA}{A-L09-04-04}{Matrix-Menge als R-Modul}
The set of all $m \times n$ matrices with entries that are elements of a commutative ring R is an R-module with scalar multiplication.
\end{EXA}

Definition D.5. 

\begin{DEF}{A-L09-04-05}{Modul Homomorphismus}
Let $\mathrm{V}_{1}$ and $\mathrm{V}_{2}$ be modules over a ring R . A map $L: \mathrm{V}_{1} \rightarrow \mathrm{V}_{2}$ is called a module homomorphism or a linear map if
$$
L\left(a v_{1}+b v_{2}\right)=a L\left(v_{1}\right)+b L\left(v_{2}\right)
$$
\end{DEF}


By analogy with the case of vector spaces we often characterize a module homomorphism $L$ by saying that $L$ is linear over R .

Example D.6. 

\begin{EXA}{A-L09-04-06}{$\operatorname{Hom}_R(V,M)$ / $L_R(V,M)$ als Modul bzgl kommutativen Ring}
The set of all module homomorphisms of a module V over a commutative ring R to another R -module M is also a (left) R -module in its own right and is denoted $\operatorname{Hom}_{\mathrm{R}}(\mathrm{V}, \mathrm{M})$ or $L_{\mathrm{R}}(\mathrm{V}, \mathrm{M})$ (we mainly use the latter). The scalar multiplication and addition in $L_{\mathrm{R}}(\mathrm{V}, \mathrm{M})$ are defined by
$$
\begin{aligned}
(f+g)(v) & :=f(v)+g(v) \text { for } f, g \in L_{\mathrm{R}}(\mathrm{V}, \mathrm{M}) \text { and all } v \in \mathrm{V} ; \\
(a f)(v) & :=a f(v) \text { for } a \in \mathrm{R} .
\end{aligned}
$$
Note that $((a b) f)(v):=(a b) f(v)=a(b f(v))=a((b f)(v))=(a(b f))(v)$. Also, in order to show that $a f$ is linear we argue as follows: $(a f)(c v)= a f(c v)=a c f(v)=c a f(v)=c(a f)(v)$. This argument fails if R is not commutative! Indeed, if R is not commutative, then $L_{\mathrm{R}}(\mathrm{V}, \mathrm{M})$ is not a module but rather only an abelian group.
\end{EXA}



Example D.7. 


\begin{EXA}{A-L09-04-07}{Vektorraum als Modul über Polynomring $\R[t]$}
Let V be a vector space and $\ell: \mathrm{V} \rightarrow \mathrm{V}$ a linear operator. Using $\ell$, we may consider V as a module over the ring of polynomials $\mathbb{R}[t]$ by defining the "scalar" multiplication by the rule
$$
p(t) v:=p(\ell) v
$$
for $p \in \mathbb{R}[t], v \in \mathrm{V}$. Here, if $p(t)=\sum a_{n} t^{n}$, then $p(\ell)$ is the linear map $\sum a_{n} \ell^{n}$.
\end{EXA}

Since the ring is usually fixed, we often omit mentioning the ring. In particular, we often abbreviate $L_{\mathrm{R}}(\mathrm{V}, \mathrm{W})$ to $L(\mathrm{V}, \mathrm{~W})$. Similar omissions will be made without further mention.

Remark D.8. 

\begin{DEF}{A-L09-04-08}{$L(V,W)$ im Fall von $\infty$-dim topologische Vektorräume}
If the modules are infinite-dimensional topological vector spaces such as Banach space, then we must distinguish between the bounded linear maps and simply linear maps. If $E$ and $F$ are infinite-dimensional Banach spaces, then $L(\mathrm{E} ; \mathrm{F})$ would normally denote bounded linear maps.
\end{DEF}



\begin{DEF}{A-L09-04-09}{Untermodul}
A submodule is defined in the obvious way as a subset $S \subset \mathrm{V}$ that is closed under the operations inherited from V so that $S$ itself is a module.
\end{DEF}

\begin{DEF}{A-L09-04-10}{Generiertes Untermodul}
The intersection of all submodules containing a subset $A \subset \mathrm{V}$ is called the submodule generated by $A$ and is denoted $\langle A\rangle$. In this case, $A$ is called a generating set. If $\langle A\rangle=\mathrm{V}$ for a finite set $A$, then we say that V is finitely generated.
\end{DEF}

\begin{DEF}{A-L09-04-11}{Quotienten Modul}
Let $S$ be a submodule of V and consider the quotient abelian group $\mathrm{V} / S$ consisting of cosets, that is, sets of the form $[v]:=v+S=\{v+x: x \in S\}$ with addition given by $[v]+[w]=[v+w]$. We define scalar multiplication by elements of the ring R by $a[v]:=[a v]$ for $a \in \mathrm{R}$. In this way, $\mathrm{V} / S$ is a module called a quotient module.
\end{DEF}

Many of the operations that exist for vector spaces have analogues in the module category. For example, 


\begin{DEF}{A-L09-04-12}{Produkt von Moduln}
If V and W are R-modules, then the set $\mathrm{V} \times \mathrm{W}$ can be made into an R -module by defining $\left(v_{1}, w_{1}\right)+\left(v_{2}, w_{2}\right):=\left(v_{1}+v_{2}, w_{1}+w_{2}\right)$ for $\left(v_{1}, w_{1}\right)$ and $\left(v_{2}, w_{2}\right)$ in $\mathrm{V} \times \mathrm{W}$\\
and
$$
a(v, w):=(a v, a w) \text { for } a \in \mathrm{R} \text { and }(v, w) \in \mathrm{V} \times \mathrm{W} .
$$
This module is sometimes written as $\mathrm{V} \oplus \mathrm{W}$, especially when taken together with the injections $\mathrm{V} \rightarrow \mathrm{V} \times \mathrm{W}$ and $\mathrm{W} \rightarrow \mathrm{V} \times \mathrm{W}$ given by $v \mapsto(v, 0)$ and $w \mapsto(0, w)$ respectively.
\end{DEF}

\begin{DEF}{A-L09-04-13}{Kern und Image von Modul Homomorphismen}
For any module homomorphism $L: \mathrm{V}_{1} \rightarrow \mathrm{V}_{2}$ we have the usual notions of kernel and image:
$$
\begin{aligned}
\operatorname{Ker} L & =\left\{v \in \mathrm{V}_{1}: L(v)=0\right\} \subset \mathrm{V}_{1} \\
\operatorname{Im}(L) & =L\left(\mathrm{V}_{1}\right)=\left\{w \in \mathrm{V}_{2}: w=L v \text { for some } v \in \mathrm{V}_{1}\right\} \subset \mathrm{V}_{2}
\end{aligned}
$$
These are submodules of $V_{1}$ and $V_{2}$ respectively.
\end{DEF}


On the other hand, modules are generally not as simple to study as vector spaces. For example, there are several notions of dimension. The following notions for a vector space all lead to the same notion of dimension. For a completely general module, these are all potentially different notions:

(1) The length $n$ of the longest chain of submodules
$$
0=\mathrm{V}_{n} \subsetneq \cdots \subsetneq \mathrm{V}_{1} \subsetneq \mathrm{V}
$$

(2) The cardinality of the largest linearly independent set (see below).\\

(3) The cardinality of a basis (see below).

For simplicity, in our study of dimension, let us now assume that R is commutative.


Definition D.9. 

\begin{DEF}{A-L09-04-14}{Lineare Unabhängigkeit für Moduln}
A set of elements $\left\{e_{1}, \ldots, e_{k}\right\}$ of a module are said to be linearly dependent if there exist ring elements $r_{1}, \ldots, r_{k} \in \mathrm{R}$ not all zero, such that $r_{1} e_{1}+\cdots+r_{k} e_{k}=0$. Otherwise, they are said to be linearly independent. We also speak of the set $\left\{e_{1}, \ldots, e_{k}\right\}$ as being a linearly independent set.
\end{DEF}

\begin{DEF}{A-L09-04-15}{Torsionselemente von Moduln}
So far so good, but it is important to realize that just because $e_{1}, \ldots, e_{k}$ are linearly dependent does not mean that we may write each of these $e_{i}$ as a linear combination of the others. It may even be that some single element $v$ forms a linearly dependent set since there may be a nonzero $r$ such that $r v=0$ (such a $v$ is said to be a torsion element).
\end{DEF}


\begin{DEF}{A-L09-04-16}{Rang eines Moduls}
If a linearly independent set $\left\{e_{1}, \ldots, e_{k}\right\}$ is maximal in size, then we say that the module has rank $k$.
\end{DEF}

Another strange possibility is that a maximal linearly independent set may not be a generating set for the module and hence may not be a basis in the sense to be defined below. The point is that although for an arbitrary $w \in \mathrm{V}$ we must have that $\left\{e_{1}, \ldots, e_{k}\right\} \cup\{w\}$ is linearly dependent and hence there must be a nontrivial expression $r w+ r_{1} e_{1}+\cdots+r_{k} e_{k}=0$, it does not follow that we may solve for $w$ since $r$ may not be an invertible element of the ring. In other words, it may not be a unit.

Definition D.10. 

\begin{DEF}{A-L09-04-17}{Basis eines Modul}
If $B$ is a generating set for a module V such that every element of V has a unique expression as a finite R -linear combination of elements of $B$, then we say that $B$ is a basis for V .
\end{DEF}

Definition D.11. 

\begin{DEF}{A-L09-04-18}{(Endlich erzeugte) freie Modul}
If an R-module has a basis, then it is referred to as a free module. If this basis is finite we indicate this by referring to the module as a finitely generated free module.
\end{DEF}

\begin{DEF}{A-L09-04-19}{Endlich dimensionales freies Modul}
It turns out that just as for vector spaces the cardinality of a basis for a finitely generated free module V is the same as that of every other basis for V . If a module over a (commutative) ring R has a basis, then the number of elements in the basis is called the dimension and must in this case be the same as the rank (the size of a maximal linearly independent set). Thus a finitely generated free module is also called a finite-dimensional free module.
\end{DEF}

Exercise D.12. Show that every finitely generated R -module is the homomorphic image of a finitely generated free module.

\begin{DEF}{A-L09-04-20}{Moduln sind Vektorräume für Körper}
If R is a field, then every module is free and is a vector space by definition. In this case, the current definitions of dimension and basis coincide with the usual ones.
\end{DEF}

\begin{DEF}{A-L09-04-21}{Ringe als endliche freie Moduln und Iso zu allg. endlichen freien Moduln}
The ring R is itself a free R-module with standard basis given by $\{1\}$. Also, $\mathrm{R}^{n}:=\mathrm{R} \oplus \cdots \oplus \mathrm{R}$ is a finitely generated free module with standard basis $\left\{\mathbf{e}_{1}, \ldots, \mathbf{e}_{n}\right\}$, where, as usual, $\mathbf{e}_{i}:=(0, \ldots, 1, \ldots, 0)$; the only nonzero entry is in the $i$-th position. Up to isomorphism, these account for all finitely generated free modules: If a module V is free with basis $e_{1}, \ldots, e_{n}$, then we have an isomorphism $\mathrm{R}^{n} \cong \mathrm{V}$ given by
$$
\left(r_{1}, \ldots, r_{n}\right) \mapsto r_{1} e_{1}+\cdots+r_{n} e_{n} .
$$
\end{DEF}

Definition D.13. 


\begin{DEF}{A-L09-04-22}{Multilineare Abbildungen für Moduln}
Let $\mathrm{V}_{i}, i=1, \ldots, k$, and W be modules over a ring R . A map $\mu: \mathrm{V}_{1} \times \cdots \times \mathrm{V}_{k} \rightarrow \mathrm{~W}$ is called multilinear ( $k$-multilinear) if for each $i, 1 \leq i \leq k$, and each fixed $\left(v_{1}, \ldots, \widehat{v}_{i}, \ldots, v_{k}\right) \in \mathrm{V}_{1} \times \cdots \times \widehat{\mathrm{V}_{i}} \times \cdots \times \mathrm{V}_{k}$ we have that the map
$$
v \mapsto \mu\left(v_{1}, \ldots, v_{i-1}, \underset{i-\text { th }}{v}, v_{i+1} \ldots, v_{k}\right),
$$
obtained by fixing all but the i-th variable, is a module homomorphism. In other words, we require that $\mu$ be R-linear in each slot separately. The set of all multilinear maps $\mathrm{V}_{1} \times \cdots \times \mathrm{V}_{k} \rightarrow \mathrm{~W}$ is denoted $L_{\mathrm{R}}\left(\mathrm{V}_{1}, \ldots, \mathrm{V}_{k} ; \mathrm{W}\right)$. If $\mathrm{V}_{1}=\cdots=\mathrm{V}_{k}=\mathrm{V}$, then we abbreviate this to $L_{\mathrm{R}}^{k}(\mathrm{V} ; \mathrm{W})$.
\end{DEF}

\begin{DEF}{A-L09-04-23}{$L_{\mathrm{R}}\left(\mathrm{V}_{1}, \ldots, \mathrm{V}_{k} ; \mathrm{W}\right)$ ist ein $R$-Modul für kommutative Ringe}
If R is commutative, then the space of multilinear maps $L_{\mathrm{R}}\left(\mathrm{V}_{1}, \ldots, \mathrm{V}_{k} ; \mathrm{W}\right)$ is itself an R -module in a fairly obvious way: If $a, b \in \mathrm{R}$ and $\mu_{1}, \mu_{2} \in L_{\mathrm{R}}\left(\mathrm{V}_{1}, \ldots, \mathrm{V}_{k} ; \mathrm{W}\right)$, then $a \mu_{1}+b \mu_{2}$ is defined pointwise in the usual way.
\end{DEF}

Note: For the remainder of this chapter, all modules will be taken to be over a fixed commutative ring R .

Let us agree to use the following abbreviation: $\mathrm{V}^{k}=\mathrm{V} \times \cdots \times \mathrm{V}$ ( $k$-fold Cartesian product).\\


Definition D.14. 

\begin{DEF}{A-L09-04-24}{Duales $R$-Modul}
The dual of an R-module V is the module $\mathrm{V}^{*}:=L_{\mathrm{R}}(\mathrm{V}, \mathrm{R})$ of all R-linear functionals on V.
\end{DEF}

\begin{DEF}{A-L09-04-25}{Reflexive $R$-Moduln}
Any element $w \in \mathrm{V}$ can be thought of as an element of $\mathrm{V}^{**}:=L_{\mathrm{R}}\left(\mathrm{V}^{*}, \mathrm{R}\right)$ according to $w(\alpha):=\alpha(w)$ where $\alpha \in \mathrm{V}^{*}$. 

This provides a map $\mathrm{V} \hookrightarrow \mathrm{V}^{* *}$, and if this map is an isomorphism, then we say that $V$ is reflexive. If $V$ is reflexive, then we are free to identify $V$ with $V^{* *}$.
\end{DEF}


Exercise D.15. 

\begin{PROP}{A-L09-04-26}{Endlich erzeugte Moduln sind reflexiv}
If V is a finitely generated free module, then V is reflexive.
\end{PROP}

For completeness, we include the definition of a projective module but what is important for us is that the finitely generated projective modules over $C^{\infty}(M)$ correspond to spaces of sections of smooth vector bundles. These modules are not necessarily free but are reflexive and have many other good properties such as being "locally free".\\


Definition D.16. 

\begin{DEF}{A-L09-04-27}{Projektive Moduln}
A module $V$ is projective if, whenever $V$ is a quotient of a module $W$, there exists a module $U$ such that the direct sum $V \oplus U$ is isomorphic to W.
\end{DEF}

Given two modules V and W over some commutative ring R , consider the class $\mathcal{C}_{\mathrm{V} \times \mathrm{W}}$ consisting of all bilinear maps $\mathrm{V} \times \mathrm{W} \rightarrow \mathrm{X}$ where X varies over all R -modules, but V and W are fixed. We take members of $\mathcal{C}_{\mathrm{V} \times \mathrm{W}}$ as the objects of a category (see Appendix A). A morphism from, say $\mu_{1}: \mathrm{V} \times \mathrm{W} \rightarrow \mathrm{X}_{1}$ to $\mu_{2}: \mathrm{V} \times \mathrm{W} \rightarrow \mathrm{X}_{2}$ is defined to be a homomorphism $\ell: \mathrm{X}_{1} \rightarrow \mathrm{X}_{2}$ such that $\mu_{2}=\ell \circ \mu_{1}$. There exists a vector space $\mathrm{T}_{\mathrm{V}, \mathrm{W}}$ together with a bilinear map $\otimes: \mathrm{V} \times \mathrm{W} \rightarrow \mathrm{T}_{\mathrm{V}, \mathrm{W}}$ that has the following universal property: For every bilinear map $\mu: \mathrm{V} \times \mathrm{W} \rightarrow \mathrm{X}$, there is a unique linear map $\widetilde{\mu}: \mathrm{T}_{\mathrm{V}, \mathrm{W}} \rightarrow \mathrm{X}$ such that $\mu=\widetilde{\mu} \circ \otimes$. If a pair ( $\mathrm{T}_{\mathrm{V}, \mathrm{W}}, \otimes$ ) with this property exists, then it is unique up to isomorphism in $\mathcal{C}_{\mathrm{V} \times \mathrm{W}}$. We refer to such a universal object as a tensor product of $V$ and $W$. We will indicate the construction of a specific tensor product that we denote by $\mathrm{V} \otimes \mathrm{W}$ with the corresponding map $\otimes: \mathrm{V} \times \mathrm{W} \rightarrow \mathrm{V} \otimes \mathrm{W}$. The idea is simple: We let $\mathrm{V} \otimes \mathrm{W}$ be the set of all linear combinations of symbols of the form $v \otimes w$ for $v \in \mathrm{V}$ and $w \in \mathrm{~W}$, subject to the relations

$$
\begin{aligned}
\left(v_{1}+v_{2}\right) \otimes w & =v_{1} \otimes w+v_{2} \otimes w, \\
v \otimes\left(w_{1}+w_{2}\right) & =v \otimes w_{1}+v \otimes w_{2}, \\
r(v \otimes w) & =r v \otimes w=v \otimes r w, \text { for } r \in \mathbb{F} .
\end{aligned}
$$

The map $\otimes$ is then simply $\otimes:(v, w) \rightarrow v \otimes w$. Let us generalize this idea to tensor products of several vector spaces at a time, and also, let us be a bit more pedantic about the construction. We seek a universal object for the category of $k$-multilinear maps of the form $\mu: \mathrm{V}_{1} \times \cdots \times \mathrm{V}_{k} \rightarrow \mathrm{~W}$ with $\mathrm{V}_{1}, \ldots, \mathrm{V}_{k}$ fixed.






\begin{DEF}{A-L09-04-28}{Universelle Eigenschaft des Tensorprodukts}
Seien $V$ und $W$ zwei Moduln über einem kommutativen Ring $R$.
Wir betrachten die Klasse
\[
\mathcal{C}_{V\times W}
:=\{\;\mu:V\times W\to X\;\mid\; X\text{ ist ein }R\text{-Modul, }\mu\text{ ist bilinear}\;\}.
\]
Die Objekte dieser Kategorie sind also bilineare Abbildungen
$\mu:V\times W\to X$ bei festem $V,W$,
wobei $X$ über alle $R$-Moduln variiert.

Ein \emph{Morphismus}
\[
\ell:\mu_1\longrightarrow \mu_2
\quad\text{für}\quad
\mu_1:V\times W\to X_1,\;\;
\mu_2:V\times W\to X_2
\]
ist definiert als ein $R$-Modulhomomorphismus $\ell:X_1\to X_2$ mit
\[
\mu_2=\ell\circ \mu_1.
\]

\smallskip
Ein Objekt $(T_{V,W},\otimes)$ in $\mathcal{C}_{V\times W}$ heißt
\emph{Tensorprodukt von $V$ und $W$},
wenn für jedes $\mu:V\times W\to X$ in $\mathcal{C}_{V\times W}$
ein eindeutiger $R$-linearer Homomorphismus
\[
\widetilde{\mu}:T_{V,W}\longrightarrow X
\]
existiert, so dass das Diagramm
\[
\begin{tikzcd}[row sep=large, column sep=large]
V\times W \arrow[r, "\otimes"] \arrow[dr, "\mu"'] & T_{V,W} \arrow[d, "\widetilde{\mu}"] \\
& X
\end{tikzcd}
\]
kommutiert, d.\,h. $\mu=\widetilde{\mu}\circ\otimes$ gilt.

\smallskip
Besitzt ein Paar $(T_{V,W},\otimes)$ diese Eigenschaft, so ist es bis auf eindeutigen Isomorphismus in der Kategorie $\mathcal{C}_{V\times W}$ bestimmt.


Die universelle Eigenschaft des Tensorprodukts charakterisiert $(T_{V,W},\otimes)$ eindeutig bis auf eindeutigen Isomorphismus.
Wir schreiben $T_{V,W}=V\otimes W$ und bezeichnen die zugehörige Abbildung
\[
\otimes:V\times W\longrightarrow V\otimes W,\qquad (v,w)\longmapsto v\otimes w,
\]
als \emph{kanonische bilineare Abbildung}.
\end{DEF}





\begin{DEF}{A-L09-04-29}{Explizite Konstruktion von $V\otimes W$}
Die Menge $V\otimes W$ wird als der Quotient des freien $R$-Moduls
über allen formalen Symbolen $v\otimes w$ ($v\in V,\,w\in W$)
nach den folgenden Relationen definiert:
\[
\begin{aligned}
(v_1+v_2)\otimes w &= v_1\otimes w + v_2\otimes w,\\
v\otimes (w_1+w_2) &= v\otimes w_1 + v\otimes w_2,\\
r(v\otimes w) &= (rv)\otimes w = v\otimes (rw),
\quad r\in R.
\end{aligned}
\]
Damit wird $V\otimes W$ ein $R$-Modul, und die Abbildung
\[
\otimes:(v,w)\longmapsto v\otimes w
\]
ist bilinear und erfüllt die universelle Eigenschaft.

Diese Konstruktion verallgemeinert sich unmittelbar auf das
\emph{$k$-fache Tensorprodukt}
\[
V_1\otimes V_2\otimes\cdots\otimes V_k,
\]
das durch dieselbe universelle Eigenschaft für multilineare Abbildungen
\[
\mu:V_1\times\cdots\times V_k\to W
\]
charakterisiert ist.
\end{DEF}







Definition D.17. 


\begin{DEF}{A-L09-04-30}{Universal für $k$-multilineare Abbildungen auf $V_1\times…\times V_k$}
A module $\mathrm{T}=\mathrm{T}_{\mathrm{V}_{1}, \ldots, \mathrm{V}_{k}}$ together with a multilinear map $\otimes: \mathrm{V}_{1} \times \cdots \times \mathrm{V}_{k} \rightarrow \mathrm{~T}$ is called universal for $k$-multilinear maps on $\mathrm{V}_{1} \times \cdots \times \mathrm{V}_{k}$ if for every multilinear map $\mu: \mathrm{V}_{1} \times \cdots \times \mathrm{V}_{k} \rightarrow \mathrm{~W}$ there is a unique linear map $\widetilde{\mu}: \mathrm{T} \rightarrow \mathrm{W}$ such that the following diagram commutes:
\[
\begin{tikzcd}[row sep=large, column sep=large]
V_1 \times \cdots \times V_k
  \arrow[r, "\mu"]
  \arrow[d, "\otimes"']
& W \\
T \arrow[ur, "\tilde{\mu}"']
\end{tikzcd}
\]
i.e. we must have $\mu=\widetilde{\mu} \circ \otimes$. If such a universal object exists, it will be called a tensor product of $\mathrm{V}_{1}, \ldots, \mathrm{V}_{k}$, and the module itself $\mathrm{T}=\mathrm{T}_{\mathrm{V}_{1}, \ldots, \mathrm{V}_{k}}$ is also referred to as a tensor product of the modules $\mathrm{V}_{1}, \ldots, \mathrm{V}_{k}$.
\end{DEF}

The tensor product is again unique up to isomorphism:

Proposition D.18. 


\begin{PROP}{A-L09-04-31}{Eindeutigkeit der Universalen von multilinearen Abbildungen auf Moduln bis auf Isomorphie}
If ($\mathrm{T}_{1},\otimes_{1}$) and ($\mathrm{T}_{2},\otimes_{2}$) are both universal for $k$-multilinear maps on $\mathrm{V}_{1} \times \cdots \times \mathrm{V}_{k}$, then there is a unique isomorphism $\Phi: \mathrm{T}_{1} \rightarrow \mathrm{~T}_{2}$ such that $\Phi \circ \otimes_{1}=\otimes_{2}$ :
\[
\begin{tikzcd}[row sep=large, column sep=large]
& V_1 \times \cdots \times V_k \arrow[dl, "\otimes_1"'] \arrow[dr, "\otimes_2"] & \\
T_1 \arrow[rr, "\Phi"] & & T_2
\end{tikzcd}
\]
\end{PROP}

Proof. 

\begin{PROOF}{A-L09-04-32}{P: Eindeutigkeit der Universalen von multilinearen Abbildungen auf Moduln bis auf Isomorphie}
By the assumption of universality, there are maps $\otimes_{1}$ and $\otimes_{2}$ such that $\Phi \circ \otimes_{1}=\otimes_{2}$ and $\bar{\Phi} \circ \otimes_{2}=\otimes_{1}$. We thus have $\bar{\Phi} \circ \Phi \circ \otimes_{1}=\otimes_{1}$, and by the uniqueness part of the definition of universality of $\otimes_{1}$ we must have $\bar{\Phi} \circ \Phi=\mathrm{id}$ or $\bar{\Phi}=\Phi^{-1}$.
\end{PROOF}

\begin{DEF}{A-L09-04-33}{Realisierung von Tensorprodukt von Moduln}
The usual specific realization of the tensor product of modules $\mathrm{V}_{1}, \ldots, \mathrm{V}_{k}$ is, roughly, the set of all linear combinations of symbols of the form $v_{1} \otimes \cdots \otimes v_{k}$ subject to the obvious multilinear relations:
$$
v_{1} \otimes \cdots \otimes a v_{i} \otimes \cdots \otimes v_{k}=a\left(v_{1} \otimes \cdots \otimes v_{i} \otimes \cdots \otimes v_{k}\right)
$$
and
$$
\begin{aligned}
& v_{1} \otimes \cdots \otimes\left(v_{i}+v_{i}^{\prime}\right) \otimes \cdots \otimes v_{k} \\
& \quad=v_{1} \otimes \cdots \otimes v_{i} \otimes \cdots \otimes v_{k}+v_{1} \otimes \cdots \otimes v_{i}^{\prime} \otimes \cdots \otimes v_{k} .
\end{aligned}
$$
This space is denoted by 
$$\bigotimes_{i=1}^{k} \mathrm{V}_{i}=\mathrm{V}_{1} \otimes \cdots \otimes \mathrm{V}_{k}$$
and called the tensor product of $\mathrm{V}_{1}, \ldots, \mathrm{V}_{k}$. Also, we will use $\mathrm{V}^{\otimes k}$ or $\otimes^{k} \mathrm{V}$ to denote $\mathrm{V} \otimes \cdots \otimes \mathrm{V}$ ( $k$-fold tensor product of V ).

The associated map $\otimes: \mathrm{V}_{1} \times \cdots \times \mathrm{V}_{k} \rightarrow \mathrm{V}_{1} \otimes \cdots \otimes \mathrm{V}_{k}$ is simply
$$
\otimes:\left(v_{1}, \ldots, v_{k}\right) \mapsto v_{1} \otimes \cdots \otimes v_{k}
$$
\end{DEF}


\begin{DEF}{A-L09-04-34}{Tensorproduktraum als Quotient von freiem Modul über direktes Produkt modulo generierte Untermodul von symbolischen Relationen}
We take 
$$\mathrm{V}_{1} \otimes \cdots \otimes \mathrm{V}_{k}:=F\left(\mathrm{V}_{1} \times\right. \left.\cdots \times \mathrm{V}_{k}\right) / U_{0}$$ 
where $F\left(\mathrm{V}_{1} \times \cdots \times \mathrm{V}_{k}\right)$ is the free module on the set $\mathrm{V}_{1} \times \cdots \times \mathrm{V}_{k}$ and $U_{0}$ is the submodule generated by the set of all elements of the form
$$
\left(v_{1}, \ldots, a v_{i}, \ldots, v_{k}\right)-a\left(v_{1}, \ldots, v_{i}, \ldots, v_{k}\right)
$$
and
$$
\begin{aligned}
& \left(v_{1}, \ldots,\left(v_{i}+v_{i}^{\prime}\right), \ldots, v_{k}\right) \\
& \quad-\left(v_{1}, \ldots, v_{i}, \ldots, v_{k}\right)-\left(v_{1}, \ldots, v_{i}^{\prime}, \ldots, v_{k}\right)
\end{aligned}
$$
where $v_{i}, v_{i}^{\prime} \in \mathrm{V}_{i}$, and $a \in \mathrm{R}$. Each element $\left(v_{1}, \ldots, v_{k}\right)$ of the set $\mathrm{V}_{1} \times \cdots \times \mathrm{V}_{k}$ is naturally identified with a generator of the free module $F\left(\mathrm{V}_{1} \times \cdots \times \mathrm{V}_{k}\right)$ and we have the obvious injection $$\mathrm{V}_{1} \times \cdots \times \mathrm{V}_{k} \hookrightarrow F\left(\mathrm{V}_{1} \times \cdots \times \mathrm{V}_{k}\right)$$ 
Its equivalence class is denoted $v_{1} \otimes \cdots \otimes v_{k}$ and the map $\otimes$ is then the composition 
$$\mathrm{V}_{1} \times \cdots \times \mathrm{V}_{k} \hookrightarrow F\left(\mathrm{V}_{1} \times \cdots \times \mathrm{V}_{k}\right) \rightarrow F\left(\mathrm{V}_{1} \times \cdots \times \mathrm{V}_{k}\right) / U_{0}$$
\end{DEF}

Proposition D.19. 


\begin{PROP}{A-L09-04-35}{Assoziierte Abbildung zu Tensorprodukt von Moduln ist universell}
$\otimes: \mathrm{V}_{1} \times \cdots \times \mathrm{V}_{k} \rightarrow \mathrm{V}_{1} \otimes \cdots \otimes \mathrm{V}_{k}$ is universal for multilinear maps on $\mathrm{V}_{1} \times \cdots \times \mathrm{V}_{k}$.
\end{PROP}

Exercise D.20. Prove the above proposition.


Proposition D.21. 

\begin{DEF}{A-L09-04-36}{Tensorprodukt von Modul Homomorphismen}
If $f: \mathrm{V}_{1} \rightarrow \mathrm{~W}_{1}$ and $g: \mathrm{V}_{2} \rightarrow \mathrm{~W}_{2}$ are module homomorphisms, then there is a unique homomorphism 
$$f \otimes g: \mathrm{V}_{1} \otimes \mathrm{V}_{2} \rightarrow \mathrm{W}_{1} \otimes \mathrm{~W}_{2}$$ 
the tensor product, which has the characterizing properties that $f \otimes g$ is linear and that 
$$(f \otimes g)\left(v_{1} \otimes v_{2}\right)=\left(f v_{1}\right) \otimes\left(g v_{2}\right)$$ 
for all $v_{1} \in \mathrm{V}_{1}, v_{2} \in \mathrm{V}_{2}$. Similarly, if $f_{i}: \mathrm{V}_{i} \rightarrow \mathrm{~W}_{i}$, we may obtain 
$$\otimes_{i} f_{i}: \otimes_{i=1}^{k} \mathrm{V}_{i} \rightarrow \bigotimes_{i=1}^{k} \mathrm{~W}_{i}$$
\end{DEF}

Proof. Exercise.


Definition D.22. 

\begin{DEF}{A-L09-04-37}{Einfache Tensoren}
Elements of $\bigotimes_{i=1}^{k} \mathrm{V}_{i}$ that may be written as $v_{1} \otimes \cdots \otimes v_{k}$ for some $v_{i}$ are called simple or decomposable.
\end{DEF}


Remark D.23. 

\begin{DEF}{A-L09-04-38}{Einfache Tensoren spannen das Tensorprodukt von Moduln auf}
It is clear from our specific realization of $\bigotimes_{i=1}^{k} V_{i}$ that elements in the image of $\otimes: \mathrm{V}_{1} \times \cdots \times \mathrm{V}_{k} \rightarrow \bigotimes_{i=1}^{k} \mathrm{V}_{i}$ span $\bigotimes_{i=1}^{k} \mathrm{V}_{i}$. I.e., decomposable elements span the space.
\end{DEF}

Exercise D.24. Not all elements are decomposable but the decomposable elements generate $\mathrm{V}_{1} \otimes \cdots \otimes \mathrm{V}_{k}$.

\begin{DEF}{A-L09-04-39}{Unterschied zwischen Tensorprodukt bzgl verschiedenen Ringen}
It may be that the $V_{i}$ are modules over more than one ring. For example, any complex vector space is module over both $\mathbb{R}$ and $\mathbb{C}$. Also, the module of\\
smooth vector fields $\mathfrak{X}_{M}(U)$ is a module over $C^{\infty}(U)$ and a module (actually a vector space) over $\mathbb{R}$. Thus it is sometimes important to indicate the ring involved, and so we write the tensor product of two R -modules V and W as $\mathrm{V} \otimes_{\mathrm{R}} \mathrm{W}$. For instance, there is a big difference between $\mathfrak{X}_{M}(U) \otimes_{C^{\infty}(U)} \mathfrak{X}_{M}(U)$ and $\mathfrak{X}_{M}(U) \otimes_{\mathbb{R}} \mathfrak{X}_{M}(U)$.
\end{DEF}

Lemma D.25. 

\begin{LEM}{A-L09-04-40}{Tensorprodukträume haben assoziative und kommutative Isomorphie}
There are the following natural isomorphisms:
\begin{enumerate}
\item $(\mathrm{V} \otimes \mathrm{W}) \otimes \mathrm{U} \cong \mathrm{V} \otimes(\mathrm{W} \otimes \mathrm{U}) \cong \mathrm{V} \otimes \mathrm{W} \otimes \mathrm{U}$, and under these isomorphisms, $(v \otimes w) \otimes u \longleftrightarrow v \otimes(w \otimes u) \longleftrightarrow v \otimes w \otimes u$.\\
\item $\mathrm{V} \otimes \mathrm{W} \cong \mathrm{W} \otimes \mathrm{V}$, and under this isomorphism $v \otimes w \longleftrightarrow w \otimes v$.
\end{enumerate}
\end{LEM}

Proof. 

\begin{PROOF}{A-L09-04-41}{P: Tensorprodukträume haben assoziative und kommutative Isomorphie}
We prove (1) and leave (2) as an exercise.

Elements of the form $(v \otimes w) \otimes u$ generate $(\mathrm{V} \otimes \mathrm{W}) \otimes \mathrm{U}$, so any map that sends $(v \otimes w) \otimes u$ to $v \otimes(w \otimes u)$ for all $v, w, u$ must be unique. Now we have compositions
$$
(\mathrm{V} \times \mathrm{W}) \times \mathrm{U}^{\otimes \times \mathrm{id}_{\mathrm{U}}}(\mathrm{V} \otimes \mathrm{~W}) \times \mathrm{U} \xrightarrow{\otimes}(\mathrm{V} \otimes \mathrm{~W}) \otimes \mathrm{U}
$$
and
$$
\mathrm{V} \times(\mathrm{W} \times \mathrm{U}) \xrightarrow{\mathrm{i}_{\mathrm{V}} \times \otimes} \mathrm{V} \times(\mathrm{W} \otimes \mathrm{U}) \xrightarrow{\otimes} \mathrm{V} \otimes(\mathrm{~W} \otimes \mathrm{U}) .
$$
It is a simple matter to check that these composite maps have the same universal property as the map $\mathrm{V} \times \mathrm{W} \times \mathrm{U} \xrightarrow{\otimes} \mathrm{V} \otimes \mathrm{W} \otimes \mathrm{U}$. The result now follows from the existence and essential uniqueness (Propositions D. 19 and D.18).
\end{PROOF}

\begin{DEF}{A-L09-04-42}{Assoziative Identifikation von Tensorprodukträumen }
We shall use the first isomorphism and the obvious generalizations to identify $\mathrm{V}_{1} \otimes \cdots \otimes \mathrm{V}_{k}$ with all legal parenthetical constructions such as $\left(\left(\left(\mathrm{V}_{1} \otimes \mathrm{V}_{2}\right) \otimes \cdots \otimes \mathrm{V}_{j}\right) \otimes \cdots\right) \otimes \mathrm{V}_{k}$ and so forth. In short, we may construct $\mathrm{V}_{1} \otimes \cdots \otimes \mathrm{V}_{k}$ by tensoring spaces two at a time. In particular, we assume the isomorphisms (as identifications)
$$
\left(\mathrm{V}_{1} \otimes \cdots \otimes \mathrm{V}_{k}\right) \otimes\left(\mathrm{W}_{1} \otimes \cdots \otimes \mathrm{~W}_{k}\right) \cong \mathrm{V}_{1} \otimes \cdots \otimes \mathrm{V}_{k} \otimes \mathrm{~W}_{1} \otimes \cdots \otimes \mathrm{~W}_{k},
$$
where $\left(v_{1} \otimes \cdots \otimes v_{k}\right) \otimes\left(w_{1} \otimes \cdots \otimes w_{k}\right)$ maps to $v_{1} \otimes \cdots \otimes v_{k} \otimes w_{1} \otimes \cdots \otimes w_{k}$.
\end{DEF}


Proposition D.26. 

\begin{PROP}{A-L09-04-43}{Natürlicher Isomorphismus $V\otimes R\cong V\cong R\otimes V$}
If V is an R-module, then we have natural isomorphisms
$$
\mathrm{V} \otimes \mathrm{R} \cong \mathrm{V} \cong \mathrm{R} \otimes \mathrm{V}
$$
given on decomposable elements as $v \otimes r \mapsto r v \mapsto r \otimes v$. (Recall that we are assuming that R is commutative.)
\end{PROP}


The proof is left to the reader. The following proposition gives a basic and often used isomorphism.

Proposition D.27. 

\begin{PROP}{A-L09-04-44}{$L_{\mathrm{R}}\left(\mathrm{~W}_{1} \otimes \cdots \otimes \mathrm{~W}_{k}, \mathrm{U}\right) \cong L\left(\mathrm{~W}_{1}, \ldots, \mathrm{~W}_{k} ; \mathrm{U}\right)$}
For R-modules $\mathrm{W}, \mathrm{V}, \mathrm{U}$, we have
$$
L_{\mathrm{R}}(\mathrm{~W} \otimes \mathrm{V}, \mathrm{U}) \cong L(\mathrm{~W}, \mathrm{V} ; \mathrm{U})
$$
More generally,
$$
L_{\mathrm{R}}\left(\mathrm{~W}_{1} \otimes \cdots \otimes \mathrm{~W}_{k}, \mathrm{U}\right) \cong L\left(\mathrm{~W}_{1}, \ldots, \mathrm{~W}_{k} ; \mathrm{U}\right)
$$
\end{PROP}

Proof. 

\begin{PROOF}{A-L09-04-45}{P: $L_{\mathrm{R}}\left(\mathrm{~W}_{1} \otimes \cdots \otimes \mathrm{~W}_{k}, \mathrm{U}\right) \cong L\left(\mathrm{~W}_{1}, \ldots, \mathrm{~W}_{k} ; \mathrm{U}\right)$}
This is more or less just a restatement of the universal property of $\mathrm{W} \otimes \mathrm{V}$. One should check that this association is indeed an isomorphism.
\end{PROOF}



Exercise D.28. 


\begin{PROP}{A-L09-04-46}{Dualmodul ist frei und Dualbasis für freie Moduln}
If W is free with basis $\left(f_{1}, \ldots, f_{n}\right)$, then $\mathrm{W}^{*}$ is also free and has a dual basis $\left(f^{1}, \ldots, f^{n}\right)$, that is, $f^{i}\left(f_{j}\right)=\delta_{j}^{i}$.
\end{PROP}



Theorem D.29. 

\begin{THEO}{A-L09-04-47}{Induzierte Basis von Tensorproduktraum aus Basis von freien Moduln}
If $\mathrm{V}_{1}, \ldots, \mathrm{V}_{k}$ are free R -modules and if $\left(e_{1}^{j}, \ldots, e_{n_{j}}^{j}\right)$ is a basis for $\mathrm{V}_{j}$, then the set of all decomposable elements of the form $e_{i_{1}}^{1} \otimes \cdots \otimes e_{i_{k}}^{k}$ is a basis for $\mathrm{V}_{1} \otimes \cdots \otimes \mathrm{V}_{k}$.
\end{THEO}

Proof. 

\begin{PROOF}{A-L09-04-48}{P: Induzierte Basis von Tensorproduktraum aus Basis von freien Moduln}
We prove this for the case of $k=2$. The general case is similar. We wish to show that if $\left(e_{1}, \ldots, e_{n_{1}}\right)$ is a basis for $\mathrm{V}_{1}$ and $\left(f_{1}, \ldots, f_{n_{2}}\right)$ is a basis for $\mathrm{V}_{2}$, then $\left\{e_{i} \otimes f_{j}\right\}$ is a basis for $\mathrm{V}_{1} \otimes \mathrm{V}_{2}$. Define $\phi_{l k}: \mathrm{V}_{1} \times \mathrm{V}_{2} \rightarrow \mathrm{R}$ by $\phi_{l k}\left(e_{i}, f_{j}\right)=\delta_{i}^{l} \delta_{j}^{k} 1$, where 1 is the identity in R and
$$
\delta_{i}^{l} \delta_{j}^{k} 1:= \begin{cases}1 & \text { if }(l, k)=(i, j) \\ 0 & \text { otherwise }\end{cases}
$$
Extend this definition bilinearly. These maps are linearly independent in $L\left(\mathrm{V}_{1}, \mathrm{V}_{2} ; \mathrm{R}\right)$ since if $\sum_{l k} a_{l k} \phi_{l k}=0$ in R , then for any $i, j$ we have
$$
\begin{aligned}
0 & =\sum_{l k} a_{l k} \phi_{l k}\left(e_{i}, f_{j}\right)=\sum_{l k} a_{l k} \delta_{i}^{l} \delta_{j}^{k} 1 \\
& =a_{i j}
\end{aligned}
$$
Thus $\operatorname{dim}\left(\mathrm{V}_{1} \otimes \mathrm{V}_{2}\right)=\operatorname{dim}\left(\left(\mathrm{V}_{1} \otimes \mathrm{V}_{2}\right)^{*}\right)=\operatorname{dim} L\left(\mathrm{V}_{1}, \mathrm{V}_{2} ; \mathrm{R}\right) \geq n_{1} n_{2}$. On the other hand, $\left\{e_{i} \otimes f_{j}\right\}$ spans the set of all decomposable elements and hence the whole space $\mathrm{V}_{1} \otimes \mathrm{V}_{2}$, so that $\operatorname{dim}\left(\mathrm{V}_{1} \otimes \mathrm{V}_{2}\right) \leq n_{1} n_{2}$ and it follows that $\left\{e_{i} \otimes f_{j}\right\}$ is a basis.
\end{PROOF}

Proposition D.30. 

\begin{PROP}{A-L09-04-49}{Identifikation von Tensorprodukträumen über Räume von Modul Homomorphismen}
There is a unique R-module map 
$$\iota: L\left(\mathrm{V}_{1}, \mathrm{~W}_{1}\right) \otimes \cdots \otimes L\left(\mathrm{V}_{k}, \mathrm{~W}_{k}\right) \rightarrow L\left(\mathrm{V}_{1} \otimes \cdots \otimes \mathrm{V}_{k}, \mathrm{~W}_{1} \otimes \cdots \otimes \mathrm{~W}_{k}\right)$$ 
such that if $f_{1} \otimes \cdots \otimes f_{k}$ is a (decomposable) element of $L\left(\mathrm{V}_{1}, \mathrm{~W}_{1}\right) \otimes \cdots \otimes L\left(\mathrm{V}_{k}, \mathrm{~W}_{k}\right)$ then
$$
\iota\left(f_{1} \otimes \cdots \otimes f_{k}\right)\left(v_{1} \otimes \cdots \otimes v_{k}\right)=f_{1}\left(v_{1}\right) \otimes \cdots \otimes f_{k}\left(v_{k}\right)
$$
If the modules are all finitely generated and free, then this is an isomorphism.
\end{PROP}

Proof. 

\begin{PROOF}{A-L09-04-50}{P: Identifikation von Tensorprodukträumen über Räume von Modul Homomorphismen}
If such a map exists, it must be unique since the decomposable elements span $L\left(\mathrm{V}_{1}, \mathrm{~W}_{1}\right) \otimes \cdots \otimes L\left(\mathrm{V}_{k}, \mathrm{~W}_{k}\right)$. To show the existence, we define a multilinear map
$$
\vartheta: L\left(\mathrm{V}_{1}, \mathrm{~W}_{1}\right) \times \cdots \times L\left(\mathrm{V}_{k}, \mathrm{~W}_{k}\right) \times \mathrm{V}_{1} \times \cdots \times \mathrm{V}_{k} \rightarrow \mathrm{~W}_{1} \otimes \cdots \otimes \mathrm{~W}_{k}
$$
by the recipe
$$
\left(f_{1}, \ldots, f_{k}, v_{1}, \ldots, v_{k}\right) \mapsto f_{1}\left(v_{1}\right) \otimes \cdots \otimes f_{k}\left(v_{k}\right)
$$
By the universal property there must be a linear map
$$
\widetilde{\vartheta}: \mathrm{V}_{1}^{*} \otimes \cdots \otimes \mathrm{V}_{k}^{*} \otimes \mathrm{V}_{1} \otimes \cdots \otimes \mathrm{V}_{k} \rightarrow \mathrm{~W}_{1} \otimes \cdots \otimes \mathrm{~W}_{k}
$$
such that $\widetilde{\vartheta} \circ \otimes=\vartheta$, where $\otimes$ is the universal map. Now define
$$
\begin{aligned}
\iota\left(f_{1}\right. & \left.\otimes \cdots \otimes f_{k}\right)\left(v_{1} \otimes \cdots \otimes v_{k}\right) \\
& :=\widetilde{\vartheta}\left(f_{1} \otimes \cdots \otimes f_{k} \otimes v_{1} \otimes \cdots \otimes v_{k}\right)
\end{aligned}
$$
The fact that $\iota$ is an isomorphism in case the $\mathrm{V}_{i}$ are all free follows easily from Exercise D. 28 and Theorem D.29.
\end{PROOF}

Since $\mathrm{R} \otimes \mathrm{R}=\mathrm{R}$, we obtain


Corollary D.31. 

\begin{KORO}{A-L09-04-51}{Identifikation von Tensorprodukträumen über Dualräume von Moduln}
There is a unique R-module map $\iota: \mathrm{V}_{1}^{*} \otimes \cdots \otimes \mathrm{V}_{k}^{*} \rightarrow \left(\mathrm{V}_{1} \otimes \cdots \otimes \mathrm{V}_{k}\right)^{*}$ such that if $\alpha_{1} \otimes \cdots \otimes \alpha_{k}$ is a (decomposable) element of $\mathrm{V}_{1}^{*} \otimes \cdots \otimes \mathrm{V}_{k}^{*}$, then
$$
\iota\left(\alpha_{1} \otimes \cdots \otimes \alpha_{k}\right)\left(v_{1} \otimes \cdots \otimes v_{k}\right)=\alpha_{1}\left(v_{1}\right) \cdots \alpha_{k}\left(v_{k}\right)
$$
If the modules are all finitely generated and free, then this is an isomorphism.
\end{KORO}


Corollary D.32. 

\begin{KORO}{A-L09-04-52}{Identifikation von Tensorprodukträumen $W\otimes V^*$}
There is a unique module map $\iota_{0}: \mathrm{W} \otimes \mathrm{V}^{*} \rightarrow L(\mathrm{V}, \mathrm{~W})$ such that if $v \otimes \beta$ is a (decomposable) element of $\mathrm{W} \otimes \mathrm{V}^{*}$, then
$$
\iota_{0}(w \otimes \beta)(v)=\beta(v) w .
$$
If V and W are finitely generated free modules, then this is an isomorphism.
\end{KORO}


Proof. 

\begin{PROOF}{A-L09-04-53}{P: Identifikation von Tensorprodukträumen $W\otimes V^*$}
If we associate to every $w \in \mathrm{~W}$ the map $w^{\text {map }} \in L(\mathrm{R}, \mathrm{W})$ given by $w^{\text {map }}(r):=r w$, then we obtain an isomorphism $\mathrm{W} \cong L(\mathrm{R}, \mathrm{W})$. Use this and then compose
$$
\begin{aligned}
\mathrm{W} \otimes \mathrm{V}^{*} & \rightarrow L(\mathrm{R}, \mathrm{~W}) \otimes L(\mathrm{V}, \mathrm{R}) \\
& \rightarrow L(\mathrm{R} \otimes \mathrm{V}, \mathrm{~W} \otimes \mathrm{R}) \cong L(\mathrm{V}, \mathrm{~W}) \\
\mathrm{V}^{*} \otimes \mathrm{~W} & \rightarrow L(\mathrm{V}, \mathrm{R}) \otimes L(\mathrm{R}, \mathrm{~W}) \\
& \rightarrow L(\mathrm{V} \otimes \mathrm{R}, \mathrm{R} \otimes \mathrm{~W}) \cong L(\mathrm{V}, \mathrm{~W}) .
\end{aligned}
$$
By combining Corollary D. 31 with Proposition D. 27 and taking $\mathrm{U}=\mathrm{R}$ we obtain the following assertion.
\end{PROOF}

Corollary D.33. 

\begin{KORO}{A-L09-04-54}{Identifikation von Tensorprodukträumen von Dualräumen mit Raum der multilinearen Abbildungen $L\left(\mathrm{V}_{1}, \ldots, \mathrm{V}_{k} ; \mathrm{R}\right)$}
There is a unique R-module map $\iota: \mathrm{V}_{1}^{*} \otimes \cdots \otimes \mathrm{V}_{k}^{*} \rightarrow L\left(\mathrm{V}_{1}, \ldots, \mathrm{V}_{k} ; \mathrm{R}\right)$ such that if $\alpha_{1} \otimes \cdots \otimes \alpha_{k}$ is a (decomposable) element of $\mathrm{V}_{1}^{*} \otimes \cdots \otimes \mathrm{V}_{k}^{*}$, then
$$
\iota\left(\alpha_{1} \otimes \cdots \otimes \alpha_{k}\right)\left(v_{1}, \ldots, v_{k}\right)=\alpha_{1}\left(v_{1}\right) \cdots \alpha_{k}\left(v_{k}\right) .
$$
If the modules are all finitely generated and free, then this is an isomorphism.
\end{KORO}


Theorem D.34. 

\begin{THEO}{A-L09-04-55}{$\bigotimes_{i=1}^{k} \mathrm{V}_{i} \times \bigotimes_{i=1}^{k} \mathrm{~W}_{i} \rightarrow \bigotimes_{i=1}^{k} \mathrm{U}_{i}$}
If $\varphi_{i}: \mathrm{V}_{i} \times \mathrm{W}_{i} \rightarrow \mathrm{U}_{i}$ are bilinear maps for $i=1, \ldots, k$, then there is a unique bilinear map
$$
\varphi: \bigotimes_{i=1}^{k} \mathrm{V}_{i} \times \bigotimes_{i=1}^{k} \mathrm{~W}_{i} \rightarrow \bigotimes_{i=1}^{k} \mathrm{U}_{i}
$$
such that for $v_{i} \in \mathrm{V}_{i}$ and $w_{i} \in \mathrm{~W}_{i}$,
$$
\varphi\left(v_{1} \otimes \cdots \otimes v_{k}, w_{1} \otimes \cdots \otimes w_{k}\right)=\varphi_{1}\left(v_{1}, w_{1}\right) \otimes \cdots \otimes \varphi_{k}\left(v_{k}, w_{k}\right)
$$
\end{THEO}

Proof. 

\begin{PROOF}{A-L09-04-56}{P: $\bigotimes_{i=1}^{k} \mathrm{V}_{i} \times \bigotimes_{i=1}^{k} \mathrm{~W}_{i} \rightarrow \bigotimes_{i=1}^{k} \mathrm{U}_{i}$}
We sketch the proof in the $k=2$ case. If $\varphi$ exists, it is unique since elements of the form $v_{1} \otimes v_{2}$ span $\mathrm{V}_{1} \otimes \mathrm{V}_{2}$ and similarly for $\mathrm{W}_{1} \otimes \mathrm{~W}_{2}$. Now by the universal property of tensor products, associated to $\varphi_{i}$ for $i=1,2$, we have unique linear maps $f_{i}: \mathrm{V}_{i} \otimes \mathrm{~W}_{i} \rightarrow \mathrm{U}_{i}$ with $f_{i} \circ \otimes=\varphi_{i}$. Then we obtain the linear map $f_{1} \otimes f_{2}:\left(\mathrm{V}_{1} \otimes \mathrm{~W}_{1}\right) \otimes\left(\mathrm{V}_{2} \otimes \mathrm{~W}_{2}\right) \rightarrow \mathrm{U}_{1} \otimes \mathrm{U}_{2}$. On the other hand we have the natural isomorphism $S:\left(\mathrm{V}_{1} \otimes \mathrm{V}_{2}\right) \otimes\left(\mathrm{W}_{1} \otimes \mathrm{~W}_{2}\right) \rightarrow \left(\mathrm{V}_{1} \otimes \mathrm{~W}_{1}\right) \otimes\left(\mathrm{V}_{2} \otimes \mathrm{~W}_{2}\right)$ induced by the obvious switching of factors of simple elements. Now define $\varphi:\left(\mathrm{V}_{1} \otimes \mathrm{V}_{2}\right) \times\left(\mathrm{W}_{1} \otimes \mathrm{~W}_{2}\right) \rightarrow \mathrm{U}_{1} \otimes \mathrm{U}_{2}$ by $\varphi(x, y)= \left(f_{1} \otimes f_{2}\right)(S(x \otimes y))$ for $x \in \mathrm{V}_{1} \otimes \mathrm{V}_{2}$ and $y \in \mathrm{~W}_{1} \otimes \mathrm{~W}_{2}$. Then we have
$$
\begin{aligned}
\varphi\left(v_{1} \otimes v_{2}, w_{1} \otimes w_{2}\right) & =\left(f_{1} \otimes f_{2}\right)\left(v_{1} \otimes w_{1} \otimes v_{2} \otimes w_{2}\right) \\
& =f_{1}\left(v_{1} \otimes w_{1}\right) \otimes f_{2}\left(v_{2} \otimes w_{2}\right) \\
& =\varphi_{1}\left(v_{1}, w_{1}\right) \otimes \varphi_{2}\left(v_{2}, w_{2}\right) .
\end{aligned}
$$
\end{PROOF}

Corollary D.35. 


\begin{KORO}{A-L09-04-57}{$\bigotimes_{i=1}^{k} \mathrm{V}_{i} \times \bigotimes_{i=1}^{k} \mathrm{~W}_{i} \rightarrow R$}
Let $\mathrm{V}_{i}$ and $\mathrm{W}_{i}$ be R -modules for $i=1, \ldots, k$. If $\varphi_{i}$ : $\mathrm{V}_{i} \times \mathrm{W}_{i} \rightarrow \mathrm{R}$ are bilinear maps, $i=1, \ldots, k$, then there is a unique bilinear map
$$
\varphi: \bigotimes_{i=1}^{k} \mathrm{V}_{i} \times \bigotimes_{i=1}^{k} \mathrm{~W}_{i} \rightarrow \mathrm{R}
$$
such that for $v_{i} \in \mathrm{V}_{i}$ and $w_{i} \in \mathrm{~W}_{i}$,
$$
\varphi\left(v_{1} \otimes \cdots \otimes v_{k}, w_{1} \otimes \cdots \otimes w_{k}\right)=\varphi_{1}\left(v_{1}, w_{1}\right) \cdots \varphi_{k}\left(v_{k}, w_{k}\right) .
$$
\end{KORO}

Proof. 

\begin{PROOF}{A-L09-04-58}{P: $\bigotimes_{i=1}^{k} \mathrm{V}_{i} \times \bigotimes_{i=1}^{k} \mathrm{~W}_{i} \rightarrow R$}
Imitate the proof of the previous theorem or just use the previous theorem together with the natural isomorphism $\mathrm{R} \otimes \cdots \otimes \mathrm{R} \cong \mathrm{R}$.
\end{PROOF}



\section*{D.1. R-Algebras}
Definition D.36. Let $R$ be a commutative ring. An (associative) R-algebra $\mathfrak{A}$ is a unitary R-module that is also a ring with identity $1_{\mathfrak{A}}$, where the ring addition and the module addition coincide and where $r\left(a_{1} a_{2}\right)=\left(r a_{1}\right) a_{2}= a_{1}\left(r a_{2}\right)$ for all $a_{1}, a_{2} \in \mathfrak{A}$ and all $r \in \mathrm{R}$.

As defined above, an algebra is associative. However, one can also define nonassociative algebras, and a Lie algebra is an example of such.

Definition D.37. Let $\mathfrak{A}$ and $\mathfrak{B}$ be R-algebras. A module homomorphism $h: \mathfrak{A} \rightarrow \mathfrak{B}$ that is also a ring homomorphism is called an R-algebra homomorphism. Epimorphism, monomorphism, and isomorphism are defined in the obvious way.

If a submodule $\mathfrak{I}$ of an algebra $\mathfrak{A}$ is also a two-sided ideal with respect to the ring structure on $\mathfrak{A}$, then $\mathfrak{A} / \mathfrak{I}$ is also an algebra.

Example D.38. Let $U$ be an open subset of a $C^{r}$ manifold $M$. The set of all smooth functions $C^{r}(U)$ is an $\mathbb{R}$-algebra ( $\mathbb{R}$ is the real numbers) with unity being the function constantly equal to 1 .

Example D.39. The set of all complex $n \times n$ matrices is an algebra over $\mathbb{C}$ with the product being matrix multiplication.

Example D.40. The set of all complex $n \times n$ matrices with real polynomial entries is an algebra over the ring of polynomials $\mathbb{R}[x]$.

Definition D.41. The set of all endomorphisms of an R -module W is an R-algebra denoted $\operatorname{End}_{R}(\mathrm{~W})$ and called the endomorphism algebra of W . Here, the sum and scalar multiplication are defined as usual and the product is composition. Note that for $r \in \mathrm{R}$

$$
r(f \circ g)=(r f) \circ g=f \circ(r g)
$$

where $f, g \in \operatorname{End}_{\mathrm{R}}(\mathrm{W})$.\\
Definition D.42. A set $A$ together with a binary operation $*: A \times A \rightarrow A$ is called a monoid if the operation is associative and there exists an element $e$ (the identity) such that $a * e=e * a=a$ for all $a \in A$.

With the operation of addition, $\mathbb{N}, \mathbb{Z}$ and $\mathbb{Z}_{2}$ are all commutative mondoids.

Definition D.43. Let $(A, *)$ be a monoid and R a ring. An $A$-graded R-algebra is an R-algebra with a direct sum decomposition $\mathfrak{A}=\sum_{i \in A} \mathfrak{A}_{i}$ such that $\mathfrak{A}_{i} \mathfrak{A}_{j} \subset \mathfrak{A}_{i * j}$. An $\mathbb{N}$-graded algebra is sometimes simply referred to as a graded algebra. A superalgebra is a $\mathbb{Z}_{2}$-graded algebra.

Definition D.44. Let $\mathfrak{A}=\sum_{i \in \mathbb{Z}} \mathfrak{A}_{i}$ and $\mathfrak{B}=\sum_{i \in \mathbb{Z}} \mathfrak{B}_{i}$ be $\mathbb{Z}$-graded algebras. An R-algebra homomorphism $h: \mathfrak{A} \rightarrow \mathfrak{B}$ is called a $\mathbb{Z}$-graded homomorphism if $h\left(\mathfrak{A}_{i}\right) \subset \mathfrak{B}_{i}$ for each $i \in \mathbb{Z}$.

We now construct the tensor algebra on a fixed R-module W. This algebra is important because using it we may construct by quotients many important algebras. Consider the following situation: $\mathfrak{A}$ is an R-algebra, W\\
an R -module, and $\phi: \mathrm{W} \rightarrow \mathfrak{A}$ is a module homomorphism. If $h: \mathfrak{A} \rightarrow \mathfrak{B}$ is an algebra homomorphism, then of course $h \circ \phi: \mathrm{W} \rightarrow \mathfrak{B}$ is an R -module homomorphism.

Definition D.45. Let W be an R -module. An R -algebra $\mathfrak{U}$ together with a map $\phi: \mathrm{W} \rightarrow \mathfrak{U}$ is called universal with respect to W if for any R -module homomorphism $\psi: \mathrm{W} \rightarrow \mathfrak{B}$ there is a unique algebra homomorphism $h$ : $\mathfrak{U} \rightarrow \mathfrak{B}$ such that $h \circ \phi=\psi$.

Again if such a universal object exists, it is unique up to isomorphism. We now exhibit the construction of this type of universal algebra. First we define $\bigotimes^{0} \mathrm{~W}:=\mathrm{R}$ and $\bigotimes^{1} \mathrm{~W}:=\mathrm{W}$. Then we define $\bigotimes^{k} \mathrm{~W}:=\mathrm{W}^{\otimes k}= \mathrm{W} \otimes \cdots \otimes \mathrm{W}$. The next step is to form the direct sum $\otimes \mathrm{W}:=\sum_{i=0}^{\infty} \otimes^{i} \mathrm{~W}$. In order to make this a $\mathbb{Z}$-graded algebra, we define $\bigotimes^{i} \mathrm{~W}:=0$ for $i<0$ and then define a product on $\bigotimes \mathrm{W}:=\sum_{i \in \mathbb{Z}} \bigotimes^{i} \mathrm{~W}$ as follows: We know that for $i, j>0$ there is an isomorphism $\mathrm{W}^{i \otimes} \otimes \mathrm{~W}^{\otimes j} \rightarrow \mathrm{~W}^{\otimes(i+j)}$ and so a bilinear map $\mathrm{W}^{i \otimes} \times \mathrm{W}^{\otimes j} \rightarrow \mathrm{~W}^{\otimes(i+j)}$ such that

$$
\left(w_{1} \otimes \cdots \otimes w_{i}\right) \times\left(w_{1}^{\prime} \otimes \cdots \otimes w_{j}^{\prime}\right) \mapsto w_{1} \otimes \cdots \otimes w_{i} \otimes w_{1}^{\prime} \otimes \cdots \otimes w_{j}^{\prime} .
$$

Similarly, we define $\bigotimes^{0} \mathrm{~W} \times \mathrm{W}^{\otimes i}=\mathrm{R} \times \mathrm{W}^{\otimes i} \rightarrow \mathrm{~W}^{\otimes i}$ by scalar multiplication. Also, $\mathrm{W}^{\otimes i} \times \mathrm{W}^{\otimes j} \rightarrow 0$ if either $i$ or $j$ is negative. Now we may use the symbol $\otimes$ to denote these multiplications without contradiction and put them together to form a product on $\bigotimes \mathrm{W}:=\sum_{i \in \mathbb{Z}} \bigotimes^{i} \mathrm{~W}$. It is now clear that

$$
\bigotimes^{i} \mathrm{~W} \times \bigotimes^{j} \mathrm{~W} \rightarrow \bigotimes^{i+j} \mathrm{~W}
$$

where we make the needed trivial definitions for the negative powers:

$$
\bigotimes^{i} \mathrm{~W}=0, \quad i<0
$$

Definition D.46. The algebra $\otimes \mathrm{W}$ is a graded algebra called the R-tensor algebra.






\pagebreak



\section*{Chapter 7}


\section*{Tensors}



In this chapter we shall employ the Einstein summation convention. 


\begin{DEF}{GA-L09-03-01}{Indizes Wahl}
For example, $\tau_{j}^{i k} \alpha^{j} v_{k}$ is taken to be shorthand for
$$
\sum_{k=1}^{m} \sum_{j=1}^{n} \tau_{j}^{i k} \alpha^{j} v_{k},
$$
where the range of summation is understood from the context. Normally, the repeated indices that are summed over occur once as a subscript and once as a superscript. For example, if $A=\left(a_{j}^{i}\right)$ is an $n \times m$ matrix and $B=\left(b_{j}^{i}\right)$ is an $m \times k$ matrix, where in this case we use upper indices to indicate rows and lower indices for columns, then $C=A B$ corresponds to
$$
c_{j}^{i}=\sum_{l=1}^{m} a_{l}^{i} b_{j}^{l} .
$$
This is reduced by the summation convention to $c_{j}^{i}=a_{l}^{i} b_{j}^{l}$. We will occasionally include the summation symbol $\sum$ for emphasis, or to meet the demands of clarity.
\end{DEF}

Tensor fields (often referred to simply as tensors) can be introduced in a rough and ready way by describing their local expressions in charts and then going on to explain how such expressions are transformed under a change of coordinates. With this approach one can gain proficiency with tensor calculations in short order, and this is usually the way physicists and engineers are introduced to tensors. However, since this approach hides much of the underlying algebraic and geometric structure, we will not pursue it here. Instead, we present tensors in terms of multilinear maps.

\subsection*{Some Multilinear Algebra}




It will be convenient to define the notion of an algebraic tensor on a vector space or module. The reader who has looked over the material in Appendix D will find this chapter easier to understand. In particular, we assume the definition of "multilinear" (Definition D.13). In this chapter, if we say that a module is finite-dimensional, ${ }^{1}$ we mean that it is free and finitely generated and thus has a basis. All modules in this chapter are assumed to be over a commutative ring with unity.

Definition 7.1. 

\begin{DEF}{GA-L09-03-02}{Algebraischer Moduln-wertiger Tensor}
Let $V$ and $W$ be modules over a commutative ring $R$ with unity. Then, an algebraic W -valued tensor on V is a multilinear mapping of the form
$$
\tau: \mathrm{V}_{1} \times \mathrm{V}_{2} \times \cdots \times \mathrm{V}_{m} \rightarrow \mathrm{~W},
$$
where each factor $V_{i}$ is either $V$ or $V^{*}$.
\end{DEF}


\begin{DEF}{GA-L09-03-03}{Ko/Kontravariante Klassifikation von Tensoren}
If the number of $V^{*}$ factors occurring is $r$ and the number of V factors is $s$, then we say that the tensor is $r$ contravariant and $s$-covariant. We also say that the tensor is of total type $\binom{r}{s}$.
\end{DEF}

The most common situation is where W is the ring R itself, in which case we often drop the adjective "R-valued". Notice that if $\tau: \mathrm{V}^{*} \times \mathrm{V} \times \mathrm{V}^{*} \times \mathrm{V} \rightarrow$ R is a tensor, then we can define a tensor $\widetilde{\tau}: \mathrm{V}^{*} \times \mathrm{V} \times \mathrm{V} \times \mathrm{V}^{*} \rightarrow \mathrm{R}$ by

$$
\widetilde{\tau}\left(\alpha_{1}, v_{1}, v_{2}, \alpha_{2}\right):=\tau\left(\alpha_{1}, v_{1}, \alpha_{2}, v_{2}\right) .
$$

Although these two tensors clearly contain the same information, they are nevertheless different. We indicate this with a more specific notation. We say that $\tau$ is a tensor of type ( $\begin{array}{cll}1 & & 1 \\ & 1 & \\ & 1\end{array}$ ), while $\widetilde{\tau}$ is of type ( $c^{1}{ }_{2}{ }^{1}$ ). More generally, a tensor might, for example, be specified to be of type

$$
\left(\begin{array}{ccccc}
r_{1} & r_{2} & & r_{a} \\
& s_{1} & s_{2} & & s_{b}
\end{array}\right) \text { or }\left(\begin{array}{cccc}
r_{1} & & & r_{a} \\
s_{1} & s_{2} & & s_{b}
\end{array}\right) .
$$

The general pattern should be clear. If $r=r_{1}+\cdots+r_{a}$ and $s=s_{1}+\cdots+s_{b}$, then the tensor would be of total type $\binom{r}{s}$, which we also write as $(r, s)$. The set of all tensors of fixed type (as above) is easily seen to be an R-module with the scalar multiplication and addition defined as is usual for spaces of functions. As another example, a multilinear map

$$
\Upsilon: \mathrm{V} \times \mathrm{V}^{*} \times \mathrm{V} \times \mathrm{V}^{*} \times \mathrm{V}^{*} \rightarrow \mathrm{~W}
$$

is a W -tensor which is of type $\left(\begin{array}{cc}1 & \\ 1 & \end{array}{ }^{2}\right.$ ) and total type ( $\begin{array}{l}3 \\ 2\end{array}$ ). The set of all W valued tensors on V of type $\left(\begin{array}{ccc}1 & & \end{array}{ }^{2}\right)$ is denoted $T_{1}{ }^{1}{ }_{1}{ }^{2}(\mathrm{V} ; \mathrm{W})$, and we have analogous notations for other types. In many, if not most, circumstances we agree to associate to each tensor of total type $\binom{r}{s}$, a unique element of $T^{r}{ }_{s}(\mathrm{V} ; \mathrm{W})$ by simply keeping the relative order among the V variables and

\footnotetext{${ }^{1}$ For modules, what we mean by dimension is what is usually called the rank.
}
among the $\mathrm{V}^{*}$ variables separately, but shifting all V variables to the right of the $\mathrm{V}^{*}$ variables. Following this procedure, we have, for example, the map
$$
T_{1}{ }^{1}{ }_{1}{ }^{2}(\mathrm{V} ; \mathrm{W}) \rightarrow T^{3}{ }_{2}(\mathrm{V} ; \mathrm{W}) .
$$

Maps like this will be called consolidation maps or consolidation isomorphisms.






\begin{DEF}{GA-L09-03-04}{Konsolidierungsabbildung}
Sei $V$ ein endlich erzeugter Modul über einem kommutativen Ring $R$,
und sei $W$ ein weiterer $R$-Modul.
Für eine gegebene Anordnung von Argumenten aus $V$ und $V^{*}$
sei
\[
\tau : (V^{*})^{\times s_{1}} \times V^{\times r_{1}}
        \times (V^{*})^{\times s_{2}} \times V^{\times r_{2}}
        \times \cdots \times (V^{*})^{\times s_{b}} \times V^{\times r_{a}}
        \;\longrightarrow\; W
\]
ein $W$-wertiger Tensor, der in dieser Reihenfolge multilinear ist.

\smallskip
Dann nennt man die Abbildung
\[
\mathrm{Cons}:\;
T_{r_{1}}{}^{s_{1}}{}_{r_{2}}{}^{s_{2}}\cdots{}_{r_{a}}{}^{s_{b}}(V;W)
\;\longrightarrow\;
T^{r}{}_{s}(V;W),
\qquad
r=\sum_{i=1}^{a} r_{i},\; s=\sum_{j=1}^{b} s_{j},
\]
die \emph{Konsolidierungsabbildung}, wobei
\[
(\mathrm{Cons}\,\tau)(\alpha_{1},\ldots,\alpha_{s},v_{1},\ldots,v_{r})
\;:=\;
\tau(\alpha_{1},\ldots,\alpha_{s_{1}},v_{1},\ldots,v_{r_{1}},
      \alpha_{s_{1}+1},\ldots,\alpha_{s_{2}},v_{r_{1}+1},\ldots,v_{r_{2}},
      \ldots).
\]
Sie entsteht also durch
\emph{Konsolidierung der Argumente}:
alle $\alpha_{i}\in V^{*}$ werden in ihrer relativen Reihenfolge
an den Anfang gestellt,
alle $v_{j}\in V$ in ihrer relativen Reihenfolge an das Ende.

\smallskip
Diese Abbildung ist $R$-linear und bijektiv; ihr Inverses
stellt die ursprüngliche Argumentanordnung wieder her.
Daher heißt $\mathrm{Cons}$ auch
\emph{Konsolidierungsisomorphismus}.
\end{DEF}




\begin{EXA}{GA-L09-03-05}{Beispiele Konsolidierungsabbildung}
Sei $\tau:V^{*}\times V\times V^{*}\times V\to R$ ein Tensor.
Dann kann man durch Vertauschung der Argumente einen weiteren Tensor definieren:
\[
\widetilde{\tau}(\alpha_1,v_1,v_2,\alpha_2)
:= \tau(\alpha_1,v_1,\alpha_2,v_2).
\]
Obwohl $\tau$ und $\widetilde{\tau}$ dieselbe Information enthalten,
sind sie formal verschiedene Tensoren,
da ihre Argumentanordnung unterschiedlich ist.
Wir unterscheiden sie notational durch ihre \emph{Typdarstellung}:
\[
\tau\text{ hat Typ }
\begin{array}{cll}
1 & & 1 \\
& 1 &
\end{array}
\quad\text{während}\quad
\widetilde{\tau}\text{ hat Typ }c^{1}{}_{2}{}^{1}.
\]
Allgemein kann ein Tensor z.\,B.\ vom Typ
\[
\left(\begin{array}{ccccc}
r_{1} & r_{2} & \cdots & r_{a}\\[2pt]
s_{1} & s_{2} & \cdots & s_{b}
\end{array}\right)
\quad\text{oder}\quad
\left(\begin{array}{cccc}
r_{1} & & & r_{a}\\[2pt]
s_{1} & s_{2} & & s_{b}
\end{array}\right)
\]
sein. Sind $r=r_1+\cdots+r_a$ und $s=s_1+\cdots+s_b$,
so ist der Tensor vom \emph{totalen Typ} $\binom{r}{s}$.

Eine multilineare Abbildung
\[
\Upsilon:V\times V^{*}\times V\times V^{*}\times V^{*}\longrightarrow W
\]
ist ein $W$-wertiger Tensor vom Typ
\[
\left(\begin{array}{cc}1 & \\[2pt] 1 & \end{array}\right)^{2},
\qquad\text{also vom totalen Typ }\binom{3}{2}.
\]
\end{EXA}











Definition 7.2. 

\begin{DEF}{GA-L09-03-06}{Konsolidierte $W$-Wertige Tensoren von totalem Typ $(r,s)$}
A tensor
$$
\tau: \underbrace{\mathrm{V}^{*} \times \mathrm{V}^{*} \times \cdots \times \mathrm{V}^{*}}_{r \text { times }} \times \underbrace{\mathrm{V} \times \mathrm{V} \times \cdots \times \mathrm{V}}_{s \text { times }} \rightarrow \mathrm{W},
$$
where all the V factors occur last, is said to be in consolidated form. The set of all such (consolidated) W-valued tensors on V will be denoted $T^{r}{ }_{s}(\mathrm{V} ; \mathrm{W})$. As a special case we have $T^{0}{ }_{1}(\mathrm{V} ; \mathrm{R})=\mathrm{V}^{*}$. We will often abbreviate $T^{r}{ }_{s}(\mathrm{V} ; \mathrm{R})$ to $T^{r}{ }_{s}(\mathrm{V})$.
\end{DEF}

For example, elements of $T_{2}^{3}(\mathrm{V} ; \mathrm{W})$ are in consolidated form, while tensors from $T_{1}{ }^{1}{ }_{1}{ }^{2}(\mathrm{V} ; \mathrm{W})$ are unconsolidated.\\
Remark 7.3. Some authors consolidate by putting all V arguments first. Also, sometimes it is appropriate to forgo the consolidation especially in connection with the "type changing" operations introduced later. Our policy will be to work with tensors in consolidated form whenever convenient.

Example 7.4. 

\begin{EXA}{GA-L09-03-07}{Kronecker Delta Tensor}
One always has the special tensor $\delta \in T_{1}^{1}(\mathrm{V} ; \mathrm{R})$ defined by
$$
\delta(a, v)=a(v)
$$
for $a \in \mathrm{V}^{*}$ and $v \in \mathrm{V}$. This tensor is sometimes referred to as the Kronecker delta tensor.
\end{EXA}

There is a natural map from V to $\mathrm{V}^{* *}$ given by $v \mapsto \widetilde{v}$, where $\widetilde{v}: \alpha \mapsto \alpha(v)$. If this map is an isomorphism, we say that V is a reflexive module and we identify V with $\mathrm{V}^{* *}$. Finite-dimensional vector spaces are reflexive.


Exercise 7.5. 

\begin{PROP}{GA-L09-03-08}{$C^\infty(M)$-Moduln der Schnitte in Vektorbündeln mit endl.-dim. Fasern sind reflexiv}
Show that the $C^{\infty}(M)$ module of sections of a vector bundle $E \rightarrow M$ is a reflexive module. (It is important here that we are only considering vector bundles with finite-dimensional fibers.)
\end{PROP}

We now consider the relationship between tensors as defined above and the abstract tensor product spaces described in Appendix D. We restrict our discussion to tensors in consolidated form since the implications for the general situation will be obvious. We specialize to the case of $R$-valued tensors where R is the ring. Recall that the $k$-th tensor power of an R module V is denoted by $\otimes^{k} \mathrm{V}:=\mathrm{V} \otimes \cdots \otimes \mathrm{V}$. We always have a module homomorphism

$$
\left(\bigotimes^{r} \mathrm{V}\right) \otimes\left(\bigotimes^{s} \mathrm{V}^{*}\right) \rightarrow T_{s}^{r}(\mathrm{V} ; \mathrm{R})
$$

whereby an element $u_{1} \otimes \cdots \otimes u_{r} \otimes \beta^{1} \otimes \cdots \otimes \beta^{s} \in\left(\otimes^{r} \mathrm{V}\right) \otimes\left(\otimes^{s} \mathrm{V}^{*}\right)$ corresponds to the multilinear map given by

$$
\left(\alpha^{1}, \ldots, \alpha^{r}, v_{1}, \ldots, v_{s}\right) \mapsto \alpha^{1}\left(u_{1}\right) \cdots \alpha^{r}\left(u_{r}\right) \beta^{1}\left(v_{1}\right) \cdots \beta^{s}\left(v_{s}\right) .
$$

We will identify $u_{1} \otimes \cdots \otimes u_{r} \otimes \beta^{1} \otimes \cdots \otimes \beta^{s}$ with this multilinear map. In particular, this entails identifying $v \in \mathrm{V}$ with the element $\widetilde{v} \in \mathrm{V}^{* *}$ where $\widetilde{v}: \alpha \mapsto \alpha(v)$. If V is a finite-dimensional vector space, then the map (7.1) is an isomorphism. In fact, it is also true that if V is the space of sections of some vector bundle over $M$ (with finite-dimensional fibers), then V is a $C^{\infty}(M)$-module and the map (7.1) is still an isomorphism. A tensor which can be written in the form $u_{1} \otimes \cdots \otimes \beta^{s}$ is called a simple or decomposable tensor. Note well that not all tensors are simple.

Remark 7.6. The reader should take careful notice of how we treat the orders of the factors: An element of $\mathrm{V} \otimes \mathrm{V}^{*} \otimes \mathrm{V}^{*}$ corresponds to a multilinear map $\mathrm{V}^{*} \times \mathrm{V} \times \mathrm{V} \rightarrow \mathrm{R}$ and not to a map $\mathrm{V} \times \mathrm{V}^{*} \times \mathrm{V}^{*} \rightarrow \mathrm{R}$.

Since the map (7.1) is not always an isomorphism for general modules, and since no analogous isomorphism exists in the case of tangent spaces to infinite-dimensional manifolds such as those discussed in $[\mathbf{L} \mathbf{1}]$, it becomes important to ask to what extent the map (7.1) is needed in differential geometry. Serge Lang has written a very fine differential geometry book for manifolds modeled on Banach spaces $[\mathbf{L} 1]$ without the help of such an isomorphism. In any case, we still can and will consider $u_{1} \otimes \cdots \otimes u_{r} \otimes \beta^{1} \otimes \cdots \otimes \beta^{s}$ to be an element of $T_{s}^{r}(\mathrm{V})$ as described above. Another thing to notice is that if (7.1) is an isomorphism for all $r$ and $s$, then in particular $\mathrm{V} \cong \mathrm{V}^{* *}$, that is, V must be reflexive. Corollary D. 33 of Appendix D states that for a finitely generated free module, being reflexive is enough to insure that (7.1) is an isomorphism for all $r$ and $s$. In the latter case, the consolidation maps introduced earlier can be described in terms of simple tensors. For example, the consolidation map

$$
T_{1}{ }_{2}{ }_{2}(\mathrm{V}) \rightarrow T^{2}{ }_{3}(\mathrm{V})
$$

is given on simple tensors by

$$
\alpha \otimes \mathrm{v} \otimes \mathrm{w} \otimes \beta \otimes \gamma \rightarrow \mathrm{v} \otimes \mathrm{w} \otimes \alpha \otimes \beta \otimes \gamma .
$$

Now let us consider the spaces $\mathrm{V} \otimes \mathrm{V}^{*}$ and $\mathrm{V} \otimes \mathrm{V}^{*} \otimes \mathrm{V}^{*}$. By a straightforward argument using the universal property of tensor product spaces, one can construct a bilinear map

$$
\left(V \otimes V^{*}\right) \times\left(V \otimes V^{*} \otimes V^{*}\right) \rightarrow V \otimes V^{*} \otimes V \otimes V^{*} \otimes V^{*}
$$

such that ( $\mathrm{v} \otimes \alpha, \mathrm{w} \otimes \beta_{1} \otimes \beta_{2}$ ) is mapped to $\mathrm{v} \otimes \alpha \otimes \mathrm{w} \otimes \beta_{1} \otimes \beta_{2}$. This corresponds to a product map

$$
\otimes: T^{1}{ }_{1}(\mathrm{V}) \times T^{1}{ }_{2}(\mathrm{V}) \rightarrow T^{1}{ }_{1}{ }^{1}{ }_{2}(\mathrm{V})
$$

such that $(S, T) \rightarrow S \otimes T$, where

$$
(S \otimes T)\left(\alpha_{1}, \mathrm{v}_{1}, \alpha_{2}, \mathrm{v}_{2}, \mathrm{v}_{3}\right)=S\left(\alpha_{1}, \mathrm{v}_{1}\right) T\left(\alpha_{2}, \mathrm{v}_{2}, \mathrm{v}_{3}\right) .
$$

The general pattern should be clear, but writing down the general case is notationally onerous. This product is the (unconsolidated) tensor product of tensors. Note carefully the order of the factors. 



\begin{DEF}{GA-L09-03-09}{Abbildung zwischen Tensorpotenzen und Tensorräumen}
Sei $V$ ein $R$-Modul über einem kommutativen Ring $R$.
Für ganze Zahlen $r,s\geq 0$ sei
\[
\otimes^{r}V := V\otimes\cdots\otimes V
\quad\text{und}\quad
\otimes^{s}V^{*} := V^{*}\otimes\cdots\otimes V^{*}
\]
das $r$- bzw.\ $s$-fache Tensorprodukt.

\smallskip
Dann existiert ein kanonischer $R$-Modulhomomorphismus
\[
\Phi_{r,s}:\;
(\otimes^{r}V)\otimes(\otimes^{s}V^{*})
\;\longrightarrow\;
T^{r}{}_{s}(V;R),
\]
definiert durch
\[
\Phi_{r,s}(u_{1}\otimes\cdots\otimes u_{r}
             \otimes\beta^{1}\otimes\cdots\otimes\beta^{s})
\;:\;
(\alpha^{1},\ldots,\alpha^{r},v_{1},\ldots,v_{s})
\longmapsto
\alpha^{1}(u_{1})\cdots\alpha^{r}(u_{r})
\beta^{1}(v_{1})\cdots\beta^{s}(v_{s}).
\]
Wir identifizieren jedes einfache Tensorprodukt
$u_{1}\otimes\cdots\otimes u_{r}\otimes\beta^{1}\otimes\cdots\otimes\beta^{s}$
mit dem entsprechenden multilinearen Abbild.

Diese Identifikation beinhaltet insbesondere,
dass man $v\in V$ mit dem Element
$\widetilde{v}\in V^{**}$ identifiziert, definiert durch
\[
\widetilde{v}(\alpha):=\alpha(v),\qquad \alpha\in V^{*}.
\]
\end{DEF}



\begin{DEF}{GA-L09-03-10}{Isomorphie für endlichdimensionale und freie Moduln}
Ist $V$ ein endlichdimensionaler $R$-Vektorraum
oder allgemeiner ein endlich erzeugter freier $R$-Modul,
so ist die Abbildung
\[
\Phi_{r,s}:(\otimes^{r}V)\otimes(\otimes^{s}V^{*})
\;\longrightarrow\; T^{r}{}_{s}(V;R)
\]
ein Isomorphismus von $R$-Moduln für alle $r,s$.

\smallskip
Insbesondere gilt dann $V\cong V^{**}$ (Reflexivität),
und die Konsolidierungsabbildungen können
in dieser Darstellung explizit auf einfachen Tensoren beschrieben werden.
\end{DEF}




\begin{EXA}{GA-L09-03-11}{Konsolidierung auf einfachen Tensoren}
Für den Fall
\[
T_{1}{}_{2}{}_{2}(V)\;\longrightarrow\;T^{2}{}_{3}(V)
\]
gilt auf einfachen Tensoren:
\[
\mathrm{Cons}(\alpha\otimes v\otimes w\otimes\beta\otimes\gamma)
\;=\;
v\otimes w\otimes\alpha\otimes\beta\otimes\gamma.
\]
\end{EXA}



\begin{DEF}{GA-L09-03-12}{Einfache Tensoren}
Ein Tensor der Form
\[
u_{1}\otimes\cdots\otimes u_{r}\otimes\beta^{1}\otimes\cdots\otimes\beta^{s}
\]
heißt \emph{einfacher} oder \emph{zerlegbarer Tensor}.
Allgemeine Tensoren sind Linearkombinationen solcher einfachen Tensoren,
doch nicht jeder Tensor ist einfach.
\end{DEF}





\begin{EXA}{GA-L09-03-13}{Zur Reihenfolge der Faktoren}
Ein Tensor
\[
v\otimes\alpha\otimes\beta \in V\otimes V^{*}\otimes V^{*}
\]
entspricht einer multilinearen Abbildung
\[
V^{*}\times V\times V \longrightarrow R,
\]
\emph{nicht} einer Abbildung
$V\times V^{*}\times V^{*}\to R$.
Die Reihenfolge der Tensorfaktoren legt also
die Argumentreihenfolge eindeutig fest.
\end{EXA}




\begin{CONC}{GA-L09-03-14}{Anwendungsbereich der Isomorphie}
Für allgemeine (nichtfreie oder unendlichdimensionale) Moduln
ist $\Phi_{r,s}$ im Allgemeinen kein Isomorphismus.
In der Differentialgeometrie spielt dies dennoch keine wesentliche Rolle:
man kann Tensorfelder auch dann definieren,
wenn $V$ kein reflexiver Modul ist.
Für endlich erzeugte freie Moduln genügt jedoch die Reflexivität von $V$,
um sicherzustellen, dass $\Phi_{r,s}$ ein Isomorphismus ist
(vgl.\ Korollar~D.33 im Anhang~D).
\end{CONC}




Seien $S\in T^{1}{}_{1}(V)$ und $T\in T^{1}{}_{2}(V)$.
Dann ist das Tensorprodukt
\[
S\otimes T\in T^{1}{}_{1}{}^{1}{}_{2}(V)
\]
definiert durch
\[
(S\otimes T)(\alpha_{1},v_{1},\alpha_{2},v_{2},v_{3})
= S(\alpha_{1},v_{1})\,T(\alpha_{2},v_{2},v_{3}).
\]
Dies entspricht der bilinearen Abbildung
\[
(V\otimes V^{*})\times(V\otimes V^{*}\otimes V^{*})
\longrightarrow V\otimes V^{*}\otimes V\otimes V^{*}\otimes V^{*},
\]
gegeben durch
\[
(v\otimes\alpha,\;w\otimes\beta_{1}\otimes\beta_{2})
\longmapsto
v\otimes\alpha\otimes w\otimes\beta_{1}\otimes\beta_{2}.
\]

Die Definition verallgemeinert sich auf beliebige Tensoren
$S\in T^{r_{1}}{}_{s_{1}}(V)$ und $T\in T^{r_{2}}{}_{s_{2}}(V)$,
wobei
\[
S\otimes T \in T^{r_{1}+r_{2}}{}_{s_{1}+s_{2}}(V).
\]
Die explizite Schreibweise wird allerdings
schnell unübersichtlich und ist meist entbehrlich.












To simplify the notation, the tensor product is often defined in a slightly different way when dealing with tensors which are in consolidated form:

Definition 7.7. 


\begin{DEF}{GA-L09-03-15}{Tensorprodukt von Tensoren}
For tensors $S \in T_{s_{1}}^{r_{1}}(\mathrm{V})$ and $T \in T_{s_{2}}^{r_{2}}(\mathrm{V})$, we define the (consolidated) tensor product $S \otimes T \in T^{r_{1}+r_{2}}{ }_{s_{1}+s_{2}}(\mathrm{V})$ by
$$
\begin{aligned}
& S \otimes T\left(\theta^{1}, \ldots, \theta^{r_{1}+r_{2}}, \mathrm{v}_{1}, \ldots, \mathrm{v}_{s_{1}+s_{2}}\right) \\
& :=S\left(\theta^{1}, \ldots, \theta^{r_{1}}, \mathrm{v}_{1}, \ldots, \mathrm{v}_{s_{1}}\right) T\left(\theta^{r_{1}+1}, \ldots, \theta^{r_{1}+r_{2}}, \mathrm{v}_{s_{1}+1}, \ldots, \mathrm{v}_{s_{1}+s_{2}}\right) .
\end{aligned}
$$
Whether or not a tensor product is the consolidated version will normally be clear from the context, and so we will drop the word "consolidated". We can also extend to products of several tensors at a time. While it is easy to see that the tensor product defined above is associative, it is not commutative since the order of the slots is an issue.
\end{DEF}

Let $T^{*}(\mathrm{V})$ denote the direct sum of all spaces of the form $T_{0}^{r}(\mathrm{V})$, where we take $T_{0}^{0}(\mathrm{V}):=\mathrm{R}$. The tensor product gives $T^{*}(\mathrm{V})$ the structure of an algebra over R as long as we make the definition that $r \otimes A:=r A$ for $r \in \mathrm{R}$.



\begin{DEF}{GA-L09-03-16}{Tensoralgebra eines $R$-Moduls}
Sei $V$ ein Modul über einem kommutativen Ring $R$.

\smallskip
Wir definieren die \emph{Tensoralgebra} von $V$ als die direkte Summe
\[
T^{*}(V)
:= \bigoplus_{r=0}^{\infty} T_{0}^{r}(V)
= R \;\oplus\; V \;\oplus\; (V\otimes V) \;\oplus\; (V\otimes V\otimes V) \;\oplus\; \cdots,
\]
wobei per Konvention $T_{0}^{0}(V):=R$ und $T_{0}^{r}(V):=\underbrace{V\otimes\cdots\otimes V}_{r\text{ Faktoren}}$ für $r\ge 1$ gilt.

\smallskip
Auf $T^{*}(V)$ definieren wir eine bilineare Multiplikation
\[
\otimes : T^{*}(V) \times T^{*}(V) \longrightarrow T^{*}(V),
\qquad
(A,B)\longmapsto A\otimes B,
\]
die komponentenweise durch das Tensorprodukt gegeben ist:
Für $A\in T_{0}^{r}(V)$ und $B\in T_{0}^{s}(V)$ setzen wir
\[
A\otimes B \in T_{0}^{r+s}(V).
\]

\smallskip
Damit $T^{*}(V)$ eine \emph{Algebra über $R$} wird,
wird zusätzlich festgelegt, dass für alle $r\in R$ und $A\in T^{*}(V)$ gilt:
\[
r\otimes A := rA \quad\text{und}\quad A\otimes r := rA.
\]

\smallskip
Mit dieser Multiplikation und Einselement $1_{R}\in T_{0}^{0}(V)$
wird $T^{*}(V)$ zu einer assoziativen, unitalen $R$-Algebra.
\end{DEF}





Proposition 7.8. 

\begin{PROP}{GA-L09-03-18}{Basis von $T_{s}^{r}(\mathrm{V})$ und Basisdarstellung von Tensoren}
Let V be a free R-module with basis ($e_{1},\ldots,e_{n}$) and corresponding dual basis $\left(e^{1},\ldots, e^{n}\right)$ for $\mathrm{V}^{*}$. Then the indexed set
$$
\left\{e_{i_{1}} \otimes \cdots \otimes e_{i_{r}} \otimes e^{j_{1}} \otimes \cdots \otimes e^{j_{s}}: i_{1}, \ldots, i_{r}, j_{1}, \ldots, j_{s}=1, \ldots, n\right\}
$$
is a basis for $T_{s}^{r}(\mathrm{V})$. If $\tau \in T_{s}^{r}(\mathrm{V})$, then
$$
\tau=\tau^{i_{1} \ldots i_{r}}{ }_{j_{1} \ldots j_{s}} e_{i_{1}} \otimes \cdots \otimes e_{i_{r}} \otimes e^{j_{1}} \otimes \cdots \otimes e^{j_{s}}
$$
where $\tau^{i_{1} \ldots i_{r}}{ }_{j_{1} \ldots j_{s}}=\tau\left(e^{i_{1}}, \ldots, e^{i_{r}}, e_{j_{1}}, \ldots, e_{j_{s}}\right)$.
\end{PROP}

Proof. 

\begin{PROOF}{GA-L09-03-19}{P: Basis von $T_{s}^{r}(\mathrm{V})$ und Basisdarstellung von Tensoren}
If $\tau \in T^{r}{ }_{s}(\mathrm{V} ; \mathrm{R})$ and we define $\tau^{i_{1} \ldots i_{r}}{ }_{j_{1} \ldots j_{s}}=\tau\left(e^{i_{1}}, \ldots, e^{i_{r}}, e_{j_{1}}, \ldots, e_{j_{s}}\right)$, then it is easy to check that
$$
\tau=\tau^{i_{1} \ldots i_{r}}{ }_{j_{1} \ldots j_{s}} e_{i_{1}} \otimes \cdots \otimes e_{i_{r}} \otimes e^{j_{1}} \otimes \cdots \otimes e^{j_{s}}
$$
and so, in particular, our indexed set spans $T_{s}^{r}(\mathrm{V} ; \mathrm{R})$. Indeed, if we denote the right hand side of the above equation by $\tau^{\prime}$, we obtain (using the\\
summation convention throughout)
$$
\begin{aligned}
& \tau^{\prime}\left(e^{k_{1}}, \ldots, e^{k_{r}}, e_{l_{1}}, \ldots, e_{l_{s}}\right) \\
& =\tau^{i_{1} \ldots i_{r}}{ }_{j_{1} \ldots j_{s}} e_{i_{1}}\left(e^{k_{1}}\right) \cdots e_{i_{r}}\left(e^{k_{r}}\right) e^{j_{1}}\left(e_{l_{1}}\right) \cdots e^{j_{s}}\left(e_{l_{s}}\right) \\
& =\tau^{i_{1} \ldots i_{r}}{ }_{j_{1} \ldots j_{s}} \delta_{i_{1}}^{k_{1}} \cdots \delta_{i_{r}}^{k_{r}} \delta_{l_{1}}^{j_{1}} \cdots \delta_{l_{s}}^{j_{s}}=\tau^{k_{1} \ldots k_{r}}{ }_{l_{1} \ldots l_{s}} \\
& =\tau\left(e^{k_{1}}, \ldots, e^{k_{r}}, e_{l_{1}}, \ldots, e_{l_{s}}\right)
\end{aligned}
$$
Thus $\tau^{\prime}$ and $\tau$ agree on basis elements, and by multilinearity $\tau^{\prime}=\tau$.\\
For independence, suppose $\tau^{i_{1} \ldots i_{r}}{ }_{j_{1} \ldots j_{s}} e_{i_{1}} \otimes \cdots \otimes e_{i_{r}} \otimes e^{j_{1}} \otimes \cdots \otimes e^{j_{s}}=0$ for some $n^{r+s}$ elements $\tau^{i_{1} \ldots i_{r}}{ }_{j_{1} \ldots j_{s}}$ of R . This is an equality of multilinear maps, and if we apply both sides to $\left(e^{k_{1}}, \ldots, e^{k_{r}}, e_{l_{1}}, \ldots, e_{l_{s}}\right)$, then we obtain $\tau^{k_{1} \ldots k_{r}}{ }_{l_{1} \ldots l_{s}}=0$. Since our choices were arbitrary, we see that all $n^{r+s}$ elements $\tau^{i_{1} \ldots i_{r}}{ }_{j_{1} \ldots j_{s}}$ are equal to 0.
\end{PROOF}


\begin{DEF}{GA-L09-03-20}{Darstellung von $(1,1)$-Tensoren}
As a special case we see that if $A \in T^{1}{ }_{1}(\mathrm{V} ; \mathrm{R})$, then $A=A^{i}{ }_{j} e_{i} \otimes e^{j}$, where $A^{i}{ }_{j}=A\left(e^{i}, e_{j}\right)$. This theorem is a special case of Theorem D. 29 of Appendix D , which we will also invoke below for spaces like $\mathrm{W} \otimes \mathrm{V}^{*}$.
\end{DEF}

If we are dealing with tensors that are not in consolidated form, it should still be clear how to obtain a basis. For example, $\left\{e_{i} \otimes e^{j} \otimes e_{k}\right\}$ is a basis for $T^{1}{ }_{1}{ }^{1}(\mathrm{V} ; \mathrm{R})$, and a typical element $A$ would have an expansion
$$
A=A_{j}^{i}{ }^{k} e_{i} \otimes e^{j} \otimes e_{k}
$$
Notice the purposeful staggered positioning of the indices in $A^{i}{ }_{j}{ }^{k}$.\\


Definition 7.9. 

\begin{DEF}{GA-L09-03-21}{Komponenten von Tensoren bzgl Basis}
The elements $A^{i_{1} \ldots i_{r}}{ }_{j_{1} \ldots j_{s}}$ from the previous proposition are called the components of $\tau$ with respect to the basis $e_{1}, \ldots, e_{n}$.
\end{DEF}

Example 7.10. 


\begin{EXA}{GA-L09-03-22}{Tensoren über Tangentialräumen}
If $\mathrm{V}=T_{p} M$ for some smooth manifold $M$, then we can use any basis of $T_{p} M$ we please. That said, we realize that if $p$ is in a coordinate chart ( $U, \mathrm{x}$ ), then the vectors $\left.\frac{\partial}{\partial x^{1}}\right|_{p}, \ldots,\left.\frac{\partial}{\partial x^{n}}\right|_{p}$ form a basis for $T_{p} M$, and we may form a basis for $T_{s}^{r}\left(T_{p} M\right)$ consisting of all tensors of the form
$$
\left.\left.\left.\left.\frac{\partial}{\partial x^{i_{1}}}\right|_{p} \otimes \cdots \otimes \frac{\partial}{\partial x^{i_{r}}}\right|_{p} \otimes d x^{j_{1}}\right|_{p} \otimes \cdots \otimes d x^{j_{s}}\right|_{p}
$$
For example, an element $A_{p}$ of $T^{1}{ }_{1}\left(T_{p} M\right)$ can be expressed in coordinate form as
$$
A_{p}=\left.\left.A^{i}{ }_{j} \frac{\partial}{\partial x^{i}}\right|_{p} \otimes d x^{j}\right|_{p}
$$
An element $A_{p}$ of $T_{s}^{r}\left(T_{p} M\right)$ can be written as
$$
A_{p}=\left.\left.\left.\left.A^{i_{1} \ldots i_{r}}{ }_{j_{1} \ldots j_{s}} \frac{\partial}{\partial x^{i_{1}}}\right|_{p} \otimes \cdots \otimes \frac{\partial}{\partial x^{i_{r}}}\right|_{p} \otimes d x^{j_{1}}\right|_{p} \otimes \cdots \otimes d x^{j_{s}}\right|_{p}
$$
and this is called the coordinate expression for $A_{p}$.
\end{EXA}



The components of a tensor depend on the basis chosen, and a different choice will give new components related to the first by a transformation law. This is the content of the following exercise:

Exercise 7.11. Let $e_{1}, \ldots, e_{n}$ be a basis for V and let $e^{1}, \ldots, e^{n}$ be the corresponding dual basis for $\mathrm{V}^{*}$. If $\bar{e}_{1}, \ldots, \bar{e}_{n}$ is another basis for V with

$$
\bar{e}_{i}=C_{i}^{k} e_{k},
$$

then the dual basis $\bar{e}^{1}, \ldots, \bar{e}^{n}$ is related to $e^{1}, \ldots, e^{n}$ by $\bar{e}^{i}=\left(C^{-1}\right)^{i}{ }_{k} e^{k}$, where $C=\left(C_{j}^{i}\right)$. Show that if $\tau^{i}{ }_{j k}$ are the components of $\tau$ with respect to the first basis (and its dual) and if $\bar{\tau}^{i}{ }_{j k}$ are the components with respect to the second basis, then

$$
\bar{\tau}^{i}{ }_{j k}=\tau^{a}{ }_{b c} C^{b}{ }_{j} C_{k}^{c}\left(C^{-1}\right)_{a}^{i} \quad(\text { sum over } a, b, c) .
$$

This is a transformation law. What is the analogous statement for $\tau \in T^{r}{ }_{s}(\mathrm{V} ; \mathrm{R})$ ?



\begin{DEF}{GA-L09-03-23}{Allgemeines Transformationsgesetz für Tensor-Komponenten}
Sei $V$ ein endlich erzeugter freier $R$-Modul (insbesondere ein endlichdimensionaler Vektorraum über einem Körper),
und sei $\{e_1,\dots,e_n\}$ eine Basis von $V$ mit Dualbasis $\{e^1,\dots,e^n\}$ von $V^*$.
Sei ferner $\{\bar e_1,\dots,\bar e_n\}$ eine zweite Basis von $V$, die durch eine
invertierbare Matrix $C=(C^i{}_j)\in \GL_n(R)$ gegeben ist via
\[
\bar e_i = C^k{}_i \, e_k
\qquad\text{(passive Basiswechsel).}
\]
Dann transformiert die Dualbasis gemäß
\[
\bar e^{\,i} = (C^{-1})^{i}{}_{k}\, e^{k}.
\]

Für einen konsolidierten $R$-wertigen Tensor $\tau\in T^{r}{}_{s}(V;R)$ mit Komponenten
\[
\tau=\tau^{i_1\cdots i_r}{}_{j_1\cdots j_s}\;
e_{i_1}\otimes\cdots\otimes e_{i_r}\otimes e^{j_1}\otimes\cdots\otimes e^{j_s},
\]
definieren die Komponenten $\bar\tau^{i_1\cdots i_r}{}_{j_1\cdots j_s}$ bezüglich
der Basis $\{\bar e_i\}$ und ihrer Dualbasis $\{\bar e^{\,i}\}$ das \emph{allgemeine Transformationsgesetz}:
\[
\boxed{\;
\bar\tau^{i_1\cdots i_r}{}_{j_1\cdots j_s}
=
(C^{-1})^{i_1}{}_{a_1}\cdots (C^{-1})^{i_r}{}_{a_r}\;
C^{b_1}{}_{j_1}\cdots C^{b_s}{}_{j_s}\;
\tau^{a_1\cdots a_r}{}_{b_1\cdots b_s}
\;}
\]
(Summation über alle wiederholten oberen/unteren Indizes).
Dabei gilt:
\begin{itemize}
\item Jeder \emph{kontravariant} (oben) erscheinende Index trägt einen Faktor
$(C^{-1})$.
\item Jeder \emph{kovariant} (unten) erscheinende Index trägt einen Faktor
$C$.
\end{itemize}

\paragraph{Spezialfälle.}
\begin{enumerate}
\item Vektor ($r=1,s=0$): $\;\bar v^{\,i}=(C^{-1})^{i}{}_{a}\,v^{a}$.
\item Kovektor ($r=0,s=1$): $\;\bar\alpha_{\,j}=C^{b}{}_{j}\,\alpha_{b}$.
\item $(1,1)$-Tensor: $\;\bar T^{\,i}{}_{j}=(C^{-1})^{i}{}_{a}\,C^{b}{}_{j}\,T^{a}{}_{b}$.
\end{enumerate}

\paragraph{Bemerkung (nicht-konsolidierte Typen).}
Für Tensoren in nicht-konsolidierter Argumentordnung (wechselnde $V$- und $V^*$-Eingaben)
gilt dieselbe Regel \emph{komponentenweise} an der jeweiligen Indexposition:
oben $\mapsto (C^{-1})$, unten $\mapsto C$, in der gegebenen Reihenfolge.

\paragraph{Kronecker-Schreibweise.}
In blockweiser Form kann man schreiben
\[
\mathrm{vec}(\bar\tau)
=
\bigl((C^{-1})^{\otimes r}\otimes C^{\otimes s}\bigr)\,\mathrm{vec}(\tau),
\]
wobei $\mathrm{vec}$ eine beliebige Vektorisierung der Komponenten bezeichnet.
\end{DEF}






Example 7.12. 


\begin{EXA}{GA-L09-03-24}{Beispiele für Tensorkomponenten}
It is easy to show that for any basis (with corresponding dual basis) as above, the Kronecker delta tensor $\delta$ has components $\delta^{i}{ }_{j}$, where $\delta_{j}^{i}=0$ if $i \neq j$ and $\delta_{i}^{i}=1$.

It is easy to show that if $S \in T_{2}^{1}(\mathrm{V})$ and $T \in T_{2}^{2}(\mathrm{V})$, then $S \otimes T$ has components given by
$$
(S \otimes T)^{a b c}{ }_{d e f g}=S^{a}{ }_{d e} T^{b c}{ }_{f g} .
$$
More generally, if $S \in T_{s_{1}}^{r_{1}}(\mathrm{V})$ and $T \in T_{s_{2}}^{r_{2}}(\mathrm{V})$, then
$$
(S \otimes T)^{a_{1} \ldots a_{r_{1}} \alpha_{1} \ldots \alpha_{r_{2}}}{ }_{b_{1} \ldots b_{s_{1}} \beta_{1} \ldots \beta_{s_{2}}}=S^{a_{1} \ldots a_{r_{1}}}{ }_{b_{1} \ldots b_{s_{1}}} T^{\alpha_{1} \ldots \alpha_{r_{2}}}{ }_{\beta_{1} \ldots \beta_{s_{2}}} .
$$
Notice the consolidation. If we choose not to employ consolidation, then the (unconsolidated) tensor product would be expressed differently in component form. For example,
$$
(S \otimes T)^{a}{ }_{d e}{ }_{f g}^{b c}=S^{a}{ }_{d e} T^{b c}{ }_{f g}
$$
This way of treating the position of the indices is a convention that is called (naturally enough) "positional notation". For more on positional notation, see $[\mathbf{P e}],[\mathbf{D o d}-\mathbf{P o s}]$ and $[\mathbf{S t e r n}]$.
\end{EXA}



If V is a finitely generated free module, then we have a natural isomorphism $\mathrm{V} \otimes \mathrm{V}^{*} \cong L(\mathrm{V}, \mathrm{V})$. We can be a bit more general. Let W be another finitely generated free module. Consider the map $\Psi: \mathrm{W} \times \mathrm{V}^{*} \rightarrow L(\mathrm{V}, \mathrm{~W})$ given by $(w, \alpha) \mapsto \Psi(w, \alpha)$, where
$$
\Psi(w, \alpha)(v):=\alpha(v) w
$$
This map is multilinear, and so using the universal property of tensor products we get a map
$$
\widehat{\Psi}: \mathrm{W} \otimes \mathrm{V}^{*} \rightarrow L(\mathrm{V}, \mathrm{~W})
$$
such that $\widehat{\Psi}: w \otimes \alpha \mapsto \Psi(w, \alpha) \in L(\mathrm{V}, \mathrm{~W})$. Let us show that this map is an isomorphism. A given element $A \in \mathrm{~W} \otimes \mathrm{V}^{*}$ may be written as $A=\sum w^{i} \otimes \alpha_{i}$, where the $w^{i}$ are linearly independent. Indeed, we just write $A$ in terms of a basis and collect terms. Then if $\widehat{\Psi}(A)=0$, we have $\sum \alpha_{i}(v) w^{i}=0$ for all $v$. Thus $\alpha_{i}(v)=0$ for all $v$. It follows that $A=0$. We see that $\widehat{\Psi}$ is injective. Both V and W are finite-dimensional, and $\mathrm{W} \otimes \mathrm{V}^{*}$ and $L(\mathrm{V}, \mathrm{~W})$ have the same dimension. Indeed, $L(\mathrm{V}, \mathrm{~W})$ is isomorphic to the space of $m \times n$ matrices with entries from R , where $m$ and $n$ are the dimensions (or ranks) of W and V respectively. If V and W were vector spaces, this would imply that $\widehat{\Psi}$ is onto. However, it remains true that $\widehat{\Psi}$ is onto in the case where W and V are finite-dimensional free modules, but we must argue differently. In Problem 2 we ask the reader to show that if $e_{1}, \ldots, e_{n}$ is a basis for V and $f_{1}, \ldots, f_{m}$ is a basis for W , then $\left\{\widehat{\Psi}\left(f_{i} \otimes e^{j}\right)\right\}$ is a basis for $L(\mathrm{V}, \mathrm{~W})$. Thus $\widehat{\Psi}: \mathrm{W} \otimes \mathrm{V}^{*} \rightarrow L(\mathrm{V}, \mathrm{~W})$ is an isomorphism. If $\tau=\tau_{j}^{i} f_{i} \otimes e^{j}$, then

$$
\widehat{\Psi}(\tau): v \mapsto \tau_{j}^{i} e^{j}(v) f_{i}=\tau_{j}^{i} v^{j} f_{i} .
$$

The components $\tau_{j}^{i}$ are exactly the entries of the matrix that represents $\widehat{\Psi}(\tau)$. When the above conditions hold, so that $\widehat{\Psi}$ is an isomorphism, we often identify $\mathrm{W} \otimes \mathrm{V}^{*}$ with $L(\mathrm{V}, \mathrm{~W})$ and write $\tau$ even when we mean $\widehat{\Psi}(\tau)$. We say that $\tau$ has two "interpretations". Under this identification

$$
(w \otimes \alpha)(v)=\alpha(v) w .
$$

In component form, the two interpretations of $\tau$ show up as

$$
\begin{aligned}
v & \mapsto w, \text { where } w^{i}=\tau_{j}^{i} v^{j}, \text { and } \\
(\alpha, v) & \mapsto \tau(\alpha, v), \text { where } \tau(\alpha, v)=\tau_{j}^{i} v^{j} \alpha_{i} .
\end{aligned}
$$

In particular, we identify $T_{1}^{1}(\mathrm{V})$ with $L(\mathrm{V}, \mathrm{V})$.


Remark 7.13 (Tensions of conventions). Both $\mathrm{W} \otimes \mathrm{V}^{*}$ and $\mathrm{V}^{*} \otimes \mathrm{~W}$ can be identified with $L(\mathrm{V}, \mathrm{V})$. For example, we may also interpret $\alpha \otimes w \in \mathrm{V}^{*} \otimes \mathrm{~W}$ as the map $(\alpha \otimes w)(v)=\alpha(v) w$. If the reader looks at how we have consolidated the spaces in the definition of $T_{s}^{r}(\mathrm{V})$, it will be apparent that we have preferred $\mathrm{W} \otimes \mathrm{V}^{*}$ over $\mathrm{V}^{*} \otimes \mathrm{~W}$. However, if one considers the case where the underlying ring is not commutative, it becomes clear that the isomorphism $\mathrm{V}^{*} \otimes \mathrm{~W} \cong L(\mathrm{V}, \mathrm{~W})$ is correct for left modules while the other is correct for right modules ( $\mathrm{V}^{*}$ is a right module if V is a left module, and vice versa). On the other hand, the identification $\mathrm{W} \otimes \mathrm{V}^{*} \cong L(\mathrm{V}, \mathrm{~W})$ is more natural for the conventions of matrix multiplication. For example, if we think of elements of $\mathbb{R}^{n}$ as column vectors and elements of $\left(\mathbb{R}^{n}\right)^{*}$ as row\\
vectors, then for $w \in \mathbb{R}^{n}$ and $\alpha \in\left(\mathbb{R}^{n}\right)^{*}$, the linear map $\alpha \otimes w$ is indeed given by the matrix $\alpha w$. The tension between standard matrix conventions and left modules is well known to algebraists. We think of modules over commutative rings as simultaneously both left and right modules.





\begin{DEF}{GA-L09-03-25}{Kanonsiche Abbildung $\,\widehat\Psi:W\otimes V^*\to L(V,W)$}
Seien $V,W$ endlich erzeugte freie $R$-Moduln über einem kommutativen Ring $R$.
Definiere
\[
\Psi: W\times V^*\longrightarrow L(V,W),\qquad
\Psi(w,\alpha)(v):=\alpha(v)\,w.
\]
Die Abbildung $\Psi$ ist bilinear. Nach der universellen Eigenschaft des Tensorprodukts
existiert genau eine $R$-lineare Abbildung
\[
\widehat\Psi:W\otimes V^*\longrightarrow L(V,W)
\]
mit $\widehat\Psi(w\otimes\alpha)=\Psi(w,\alpha)$.
\end{DEF}



\begin{PROP}{GA-L09-03-26}{$\widehat\Psi$ ist ein Isomorphismus (freier, endlich erzeugter Fall)}
Sind $V$ und $W$ endlich erzeugte freie $R$-Moduln, so ist
\[
\widehat\Psi:W\otimes V^*\xrightarrow{\;\cong\;} L(V,W)
\]
ein $R$-Modul-Isomorphismus.
\end{PROP}



\begin{PROOF}{GA-L09-03-27}{P: $\widehat\Psi$ ist ein Isomorphismus (freier, endlich erzeugter Fall)}
\emph{Injektivität.}
Schreibe ein $A\in W\otimes V^*$ als $A=\sum_{i=1}^k w^i\otimes \alpha_i$
mit $w^1,\dots,w^k$ $R$-linear unabhängig (durch Gruppieren bzgl.\ einer Basis).
Aus $\widehat\Psi(A)=0$ folgt
$\sum_i \alpha_i(v)\,w^i=0$ für alle $v\in V$, also $\alpha_i(v)=0$ für alle $v$ und alle $i$,
also $\alpha_i=0$ und somit $A=0$.

\smallskip
\emph{Surjektivität.}
Wähle Basen $\{e_1,\dots,e_n\}$ von $V$ und $\{f_1,\dots,f_m\}$ von $W$
sowie die Dualbasis $\{e^1,\dots,e^n\}$ von $V^*$.
Definiere $m\cdot n$ Elemente
\[
T_{i}{}^{j} := f_i\otimes e^{j}\in W\otimes V^*.
\]
Dann sind $\widehat\Psi(T_{i}{}^{j})\in L(V,W)$ genau die Abbildungen,
die $e_j\mapsto f_i$ und alle $e_{j'}$ mit $j'\neq j$ auf $0$ schicken.
Diese $\{\widehat\Psi(T_{i}{}^{j})\}$ bilden eine $R$-Basis von $L(V,W)$,
also ist $\widehat\Psi$ surjektiv. Zusammen mit der Injektivität ist $\widehat\Psi$ ein Isomorphismus.
\end{PROOF}




\begin{KORO}{GA-L09-03-28}{$\widehat\Psi$ ist ein Isomorphismus (freier, endlich erzeugter Fall, $V=W$)}
Für $V$ endlich erzeugt frei gilt der kanonische Isomorphismus
\[
V\otimes V^*\;\cong\; L(V,V)=\operatorname{End}_R(V).
\]
\end{KORO}




[Komponenten und zwei Interpretationen]

\begin{CONC}{GA-L09-03-29}{Komponenten und zwei Interpretationen}
Seien $\{e_j\}$ Basis von $V$ mit Dualbasis $\{e^j\}$ und $\{f_i\}$ Basis von $W$.
Für
\[
\tau=\tau^{\,i}{}_{j}\,f_i\otimes e^{j}\in W\otimes V^*
\]
ist
\[
\widehat\Psi(\tau)(v)=\tau^{\,i}{}_{j}\,e^{j}(v)\,f_i
\quad\text{(Matrixform: } w^i=\tau^{\,i}{}_{j} v^{j}\text{)}.
\]
Unter dem Isomorphismus identifizieren wir oft $\tau$ mit $\widehat\Psi(\tau)$
und schreiben kurz $(w\otimes\alpha)(v)=\alpha(v)\,w$.
Die Komponenten $\tau^{\,i}{}_{j}$ sind dabei genau die Einträge der Matrix,
die $\widehat\Psi(\tau)$ in den gewählten Basen repräsentiert.
\end{CONC}




\begin{CONC}{GA-L09-03-30}{Funktorialität/Naturalität von kanonischer Abbildung $\Psi$}
Die Konstruktion ist natürlich in $V$ und $W$:
Für $R$-lineare $a:W\to W'$ und $b:V\to V'$ kommutiert
\[
\begin{tikzcd}[column sep=large]
W\otimes V^* \arrow[r,"\widehat\Psi"] \arrow[d,"a\otimes b^*"'] & L(V,W) \arrow[d,"a\circ(\cdot)\circ b"] \\
W'\otimes (V')^* \arrow[r,"\widehat\Psi"] & L(V',W')
\end{tikzcd}
\]
wobei $b^*:(V')^*\to V^*$ das Dual von $b$ bezeichnet.
\end{CONC}





\begin{CONC}{GA-L09-03-31}{Kanonische Abbildung für Konventionen und nichtkommutativer Fall}
Auch $V^*\otimes W \cong L(V,W)$ via $(\alpha\otimes w)(v)=\alpha(v)\,w$.
Im kommutativen Fall sind beide Identifikationen üblich;
für nichtkommutative Ringe ist $V^*\otimes W\to L(V,W)$ die natürliche Identifikation
für \emph{linke} Moduln, während $W\otimes V^*\to L(V,W)$ zur rechten Modulkonvention passt.
Im üblichen Matrixkalkül (Spaltenvektoren $v$, Zeilenvektoren $\alpha$) entspricht
$\alpha\otimes w$ der Rang-1-Matrix $\alpha\,w$.
\end{CONC}




If $V$ is finite-dimensional, then one can make various other reinterpretations of tensors:

Example 7.14. 

Suppose that V is finite-dimensional. Then, elements of $T_{s}^{r}(\mathrm{V})$ can be interpreted as members of
$$
T_{s}^{0}\left(\mathrm{V} ; T_{0}^{r}(\mathrm{V})\right)
$$
according to the prescription that $\tau\left(v_{1}, \ldots, v_{s}\right)$ acts by
$$
\tau\left(v_{1}, \ldots, v_{s}\right)\left(\alpha^{1}, \ldots, \alpha^{r}\right):=\tau\left(\alpha^{1}, \ldots, \alpha^{r}, v_{1}, \ldots, v_{s}\right) .
$$
Similarly, elements of $T_{s}^{r}(\mathrm{V})$ can be interpreted as members of $T_{0}^{r}\left(\mathrm{V} ; T_{s}^{0}(\mathrm{V})\right)$.



\begin{DEF}{GA-L09-03-32}{Alternative Interpretationen von Tensoren als modulwertige Abbildungen}
Sei $V$ ein endlichdimensionaler freier $R$-Modul über einem kommutativen Ring $R$.

\smallskip
Für ganze Zahlen $r,s\ge 0$ bezeichne
\[
T^{r}{}_{s}(V)
:= \mathrm{Hom}_{R}\big((V^{*})^{\otimes s}\otimes V^{\otimes r},\,R\big)
\]
den Raum der $R$-wertigen Tensoren vom Typ $\binom{r}{s}$ auf $V$.
Wir zeigen nun, dass ein solcher Tensor auf zwei äquivalente Arten
als \emph{multilineare Abbildung mit modulwertigen Werten} aufgefasst werden kann.

\textbf{(1) Kovariant als $V$-wertige Abbildung.}
Ein Tensor $\tau\in T^{r}{}_{s}(V)$ kann als Element von
\[
T^{0}{}_{s}\big(V;\,T^{r}{}_{0}(V)\big)
:= \mathrm{Hom}_{R}\big(V^{\otimes s},\,T^{r}{}_{0}(V)\big)
\]
interpretiert werden, d.\,h. als multilineare Abbildung
\[
\tau: V^{\times s}\longrightarrow T^{r}{}_{0}(V)=V^{\otimes r}.
\]
Für feste Argumente $v_{1},\dots,v_{s}\in V$ erhält man ein Element
\[
\tau(v_{1},\dots,v_{s})\in T^{r}{}_{0}(V),
\]
welches selbst $r$-fach kontravariant ist.
Die Wirkung dieses Objekts auf $r$ kovariante Argumente
$\alpha^{1},\dots,\alpha^{r}\in V^{*}$ ist definiert durch
\[
\tau(v_{1},\dots,v_{s})
\big(\alpha^{1},\dots,\alpha^{r}\big)
:= \tau(\alpha^{1},\dots,\alpha^{r},v_{1},\dots,v_{s}).
\]

\textbf{(2) Kontravariant als $V^{*}$-wertige Abbildung.}
Ebenso kann $\tau$ als Element von
\[
T^{r}{}_{0}\big(V;\,T^{0}{}_{s}(V)\big)
:= \mathrm{Hom}_{R}\big((V^{*})^{\otimes r},\,T^{0}{}_{s}(V)\big)
\]
interpretiert werden, d.\,h. als multilineare Abbildung
\[
\tau:(V^{*})^{\times r}\longrightarrow T^{0}{}_{s}(V)=(V^{*})^{\otimes s}.
\]
Für feste kovariante Argumente $\alpha^{1},\dots,\alpha^{r}\in V^{*}$
erhält man ein Element
\[
\tau(\alpha^{1},\dots,\alpha^{r})\in T^{0}{}_{s}(V),
\]
welches $s$-fach kovariant ist.
Dieses wirkt auf $s$ kontravariante Argumente $v_{1},\dots,v_{s}\in V$ nach
\[
\tau(\alpha^{1},\dots,\alpha^{r})(v_{1},\dots,v_{s})
:= \tau(\alpha^{1},\dots,\alpha^{r},v_{1},\dots,v_{s}).
\]

\textbf{(3) Zusammenfassung.}
Beide Sichtweisen sind natürlich äquivalent und fassen dieselbe Struktur von
$\tau$ zusammen:
\[
T^{r}{}_{s}(V)
\;\cong\;
T^{0}{}_{s}\big(V;T^{r}{}_{0}(V)\big)
\;\cong\;
T^{r}{}_{0}\big(V;T^{0}{}_{s}(V)\big).
\]
Diese Isomorphismen sind durch die obigen Wirkungsdefinitionen eindeutig festgelegt.
\end{DEF}





Example 7.15. 


Let V be as in the previous example. Elements of $T_{s_{1}+s_{2}}^{0}(\mathrm{V})$ can be interpreted as members of
$$
T_{s_{1}}^{0}\left(\mathrm{V} ; T_{s_{2}}^{0}(\mathrm{V})\right)
$$
by the prescription
$$
\tau\left(v_{1}, \ldots, v_{s_{1}}\right)\left(u_{1}, \ldots, u_{s_{2}}\right):=\tau\left(v_{1}, \ldots, v_{s_{1}}, u_{1}, \ldots, u_{s_{2}}\right)
$$
One can easily see from the above examples that many reinterpretations are possible. One of the most common is where one interprets elements of $T_{2}^{0}(\mathrm{V}, \mathrm{R})$ as elements of $L\left(\mathrm{V}, \mathrm{V}^{*}\right)$ according to
$$
\tau(v)(u)=\tau(v, u) \text { for } u, v \in \mathrm{V}
$$




\begin{DEF}{GA-L09-03-33}{Iterative und funktionale Reinterpretation von Tensoren}
Sei $V$ ein endlichdimensionaler freier $R$-Modul. Für $s_1,s_2\ge 0$ bezeichne $T^{0}{}_{s}(V)=(V^{*})^{\otimes s}$ den Raum der $s$-fach kovarianten Tensoren auf $V$.

\textbf{(1) Iterierte Interpretation.}
Jedes $\tau\in T^{0}{}_{s_1+s_2}(V)$ kann kanonisch als Element von
\[
T^{0}{}_{s_1}\big(V;\,T^{0}{}_{s_2}(V)\big)
:= \mathrm{Hom}_{R}\!\big(V^{\otimes s_1},\,(V^{*})^{\otimes s_2}\big)
\]
interpretiert werden, d.\,h. als multilineare Abbildung
\[
\tau:V^{\times s_1}\longrightarrow T^{0}{}_{s_2}(V)
\]
mit der Vorschrift
\[
\tau(v_1,\dots,v_{s_1})(u_1,\dots,u_{s_2})
:=\tau(v_1,\dots,v_{s_1},u_1,\dots,u_{s_2}).
\]
Somit entstehen aus einem $(s_1+s_2)$-fach kovarianten Tensor
Abbildungen, deren Werte selbst wieder kovariante Tensoren sind.

\textbf{(2) Funktorielle Struktur.}
Die obige Identifikation definiert einen kanonischen Isomorphismus
\[
T^{0}{}_{s_1+s_2}(V)
\;\cong\;
T^{0}{}_{s_1}\big(V;\,T^{0}{}_{s_2}(V)\big),
\]
der natürlich in $V$ ist, d.\,h. $R$-lineare Abbildungen
$f:V\to W$ induzieren kommutative Diagramme zwischen beiden Seiten.

\textbf{(3) Beispiel: Bilineare Form als lineare Abbildung in das Dual.}
Ein besonders wichtiger Spezialfall ist $s_1=s_2=1$:
\[
T^{0}{}_{2}(V)\;\cong\; L(V,V^{*}).
\]
Jedes $\tau\in T^{0}{}_{2}(V)$ kann als lineare Abbildung
\[
\tau:V\longrightarrow V^{*},\qquad
\tau(v)(u):=\tau(v,u),
\]
aufgefasst werden. Ist $\tau$ zusätzlich symmetrisch bzw.\ antisymmetrisch,
so entspricht $\tau$ einer symmetrischen bzw.\ alternierenden Bilinearform auf $V$.

\textbf{(4) Bemerkung zu weiteren Iterationen.}
Durch wiederholte Anwendung dieser Identifikation können Tensoren beliebigen Typs
als mehrstufige Abbildungen interpretiert werden, deren Werte wiederum Tensoren sind,
z.\,B.
\[
T^{0}{}_{s_1+s_2+s_3}(V)
\;\cong\;
T^{0}{}_{s_1}\big(V;\,T^{0}{}_{s_2}\big(V;\,T^{0}{}_{s_3}(V)\big)\big),
\]
wobei jedes Argumentniveau die Kovarianzstruktur einer „inneren“ Tensorstufe beschreibt.
\end{DEF}




Exercise 7.16. Show that under the identification of $T_{1}^{1}(\mathrm{V}, \mathrm{R})$ with $L(\mathrm{V}, \mathrm{V})$ we can interpret the Kronecker delta tensor as the identity map.


Definition 7.17. 

\begin{DEF}{GA-L09-03-34}{Symmetrische ko/kontravariante Tensore}
A covariant tensor $\tau \in T_{s}^{0}(\mathrm{V}, \mathrm{~W})$ is said to be symmetric if
$$
\tau\left(v_{1}, \ldots, v_{s}\right)=\tau\left(v_{\sigma(1)}, \ldots, v_{\sigma(s)}\right)
$$
for all $v_{1}, \ldots, v_{s}$ and all permutations $\sigma$ of the letters $\{1,2, \ldots, s\}$. We define a symmetric contravariant tensor similarly.
\end{DEF}


Definition 7.18. 

\begin{DEF}{GA-L09-03-35}{Alternierende ko/kontravariante Tensore}
A covariant tensor $\tau \in T_{s}^{0}(\mathrm{V}, \mathrm{~W})$ is said to be alternating if
$$
\tau\left(v_{1}, \ldots, v_{s}\right)=\operatorname{sgn}(\sigma) \tau\left(v_{\sigma(1)}, \ldots, v_{\sigma(s)}\right)
$$
for all $v_{1}, \ldots, v_{s}$ and all permutations $\sigma$ of the letters $\{1,2, \ldots, s\}$, where $\operatorname{sgn}(\sigma)=1$ if $\sigma$ is an even permutation and -1 if it is an odd permutation. $\operatorname{sgn}(\sigma)$ is called the sign of the permutation $\sigma$; We define an alternating contravariant tensor similarly.
\end{DEF}

If $\ell: \mathrm{U} \rightarrow \mathrm{V}$ is a linear map, then the map $\ell^{*}: T^{0}{ }_{s}(\mathrm{V} ; \mathrm{W}) \rightarrow T^{0}{ }_{s}(\mathrm{U} ; \mathrm{W})$ is defined by
$$
\left(\ell^{*} \tau\right)\left(u_{1}, \ldots, u_{s}\right):=\tau\left(\ell\left(u_{1}\right), \ldots, \ell\left(u_{s}\right)\right) .
$$
$\ell^{*} \tau$ is called the pull-back of $\tau$ by $\ell$. It is easy to show that $\ell^{*}$ is linear. If we have linear maps $\ell: \mathrm{U}_{1} \rightarrow \mathrm{U}_{2}$ and $\lambda: \mathrm{U}_{2} \rightarrow \mathrm{V}$, then
$$
(\lambda \circ \ell)^{*}: T_{s}^{0}(\mathrm{V} ; \mathrm{W}) \rightarrow T_{s}^{0}\left(\mathrm{U}_{1} ; \mathrm{W}\right)
$$
and
$$
(\lambda \circ \ell)^{*}=\ell^{*} \circ \lambda^{*}
$$
Thus the pull-back $\ell \rightarrow \ell^{*}$ defines a contravariant functor in the category of W-valued covariant tensors. (Because of this, one might wish that covariant tensors were called contravariant and vice versa, and indeed some authors have reversed the traditional terminology.) 

Suppose that ( $e_{1}, \ldots, e_{n}$ ) is a basis for V and that $\left(f_{1}, \ldots, f_{m}\right)$ is a basis for W . If $\ell\left(e_{i}\right)=\sum \ell_{i}^{k} f_{k}$, then we have
$$
\begin{aligned}
\left(\ell^{*} \tau\right)_{i_{1} \ldots i_{s}} & =\left(\ell^{*} \tau\right)\left(e_{i_{1}}, \ldots, e_{i_{s}}\right)=\tau\left(\ell\left(e_{i_{1}}\right), \ldots, \ell\left(e_{i_{s}}\right)\right) \\
& =\tau\left(\ell_{i_{1}}^{k_{1}} f_{k_{1}}, \ldots, \ell_{i_{s}}^{k_{s}} f_{k_{s}}\right)=\ell_{i_{1}}^{k_{1}} \cdots \ell_{i_{s}}^{k_{s}} \tau\left(f_{k_{1}}, \ldots, f_{k_{s}}\right) \\
& =\tau_{k_{1} \ldots k_{s}} \ell_{i_{1}}^{k_{1}} \cdots \ell_{i_{s}}^{k_{s}}
\end{aligned}
$$
which gives the component form of the pull-back operation in terms of the matrix $\left(\ell_{i}^{k}\right)$.



\begin{DEF}{GA-L09-03-36}{Pull-back von kovarianten Tensoren}
Seien $U,V,W$ endlichdimensionale $R$-Moduln und sei
$\ell:U\to V$ eine $R$-lineare Abbildung.
Für jedes $s\ge0$ definiert $\ell$ eine Abbildung
\[
\ell^{*}:T^{0}{}_{s}(V;W)\longrightarrow T^{0}{}_{s}(U;W),
\]
genannt der \emph{Pull-back von $W$-wertigen kovarianten Tensoren} durch $\ell$,
durch die Vorschrift
\[
(\ell^{*}\tau)(u_{1},\dots,u_{s})
:=\tau\big(\ell(u_{1}),\dots,\ell(u_{s})\big),
\qquad
\tau\in T^{0}{}_{s}(V;W),\;u_{i}\in U.
\]
\end{DEF}



\begin{PROP}{GA-L09-03-37}{Eigenschaften des Pull-backs}
\begin{enumerate}
\item $\ell^{*}$ ist $R$-linear.
\item Für lineare Abbildungen $\ell:U_{1}\to U_{2}$ und
$\lambda:U_{2}\to V$ gilt die \emph{Funktorialität}
\[
(\lambda\circ\ell)^{*}=\ell^{*}\circ\lambda^{*}.
\]
Insbesondere gilt $\mathrm{id}_{V}^{*}=\mathrm{id}_{T^{0}{}_{s}(V;W)}$.
\item Somit definiert die Zuweisung
\[
V\longmapsto T^{0}{}_{s}(V;W),
\qquad
\ell\longmapsto\ell^{*},
\]
einen \emph{kontravarianten Funktor} von der Kategorie der $R$-Moduln
in die Kategorie der $R$-Moduln.
\end{enumerate}
\end{PROP}





\begin{PROOF}{GA-L09-03-38}{P: Eigenschaften des Pull-backs}
Für $u_{1},\dots,u_{s}\in U_{1}$ gilt
\[
\big((\lambda\circ\ell)^{*}\tau\big)(u_{1},\dots,u_{s})
=\tau\big(\lambda(\ell(u_{1})),\dots,\lambda(\ell(u_{s}))\big)
=(\ell^{*}(\lambda^{*}\tau))(u_{1},\dots,u_{s}).
\]
\end{PROOF}



[Komponentenform]

\begin{DEF}{GA-L09-03-39}{Komponentendarstellung des Pullbacks}
Seien $\{e_{1},\dots,e_{n}\}$ eine Basis von $U$ und
$\{f_{1},\dots,f_{m}\}$ eine Basis von $V$.
Schreibe die lineare Abbildung $\ell:U\to V$ in Komponentenform als
\[
\ell(e_{i})=\ell_{i}{}^{k}f_{k}.
\]
Für einen Tensor $\tau\in T^{0}{}_{s}(V;W)$ mit Komponenten
$\tau_{k_{1}\dots k_{s}}:=\tau(f_{k_{1}},\dots,f_{k_{s}})$ gilt dann
\[
\begin{aligned}
(\ell^{*}\tau)_{i_{1}\dots i_{s}}
&=(\ell^{*}\tau)(e_{i_{1}},\dots,e_{i_{s}})\\
&=\tau\big(\ell(e_{i_{1}}),\dots,\ell(e_{i_{s}})\big)\\
&=\tau\big(\ell_{i_{1}}{}^{k_{1}}f_{k_{1}},\dots,\ell_{i_{s}}{}^{k_{s}}f_{k_{s}}\big)\\
&=\ell_{i_{1}}{}^{k_{1}}\cdots \ell_{i_{s}}{}^{k_{s}}\,
\tau_{k_{1}\dots k_{s}}.
\end{aligned}
\]
Dies ist die übliche \emph{Transformationsformel} für die Komponenten
eines kovarianten Tensors unter der linearen Abbildung $\ell$.
\end{DEF}




[Intuitive Bedeutung]

\begin{CONC}{GA-L09-03-40}{Bedeutung des Pullbacks}
Der Pull-back \(\ell^{*}\tau\) „zieht“ den Tensor \(\tau\) von \(V\) auf \(U\)
zurück, indem er jedes seiner Argumente zuerst durch \(\ell\) nach \(V\)
abbildet, bevor \(\tau\) ausgewertet wird.
Dadurch werden kovariante Tensoren (die auf Vektoren wirken)
\emph{entgegengesetzt} zur Richtung von \(\ell\) transportiert,
weshalb der Pull-back ein \textbf{kontravarianter Funktor} ist.

[Sonderfälle]

\begin{itemize}
\item Für $s=1$ und $W=R$ ist $\ell^{*}$ einfach die Dualabbildung
$\ell^{*}:V^{*}\to U^{*}$, gegeben durch $(\ell^{*}\alpha)(u)=\alpha(\ell(u))$.
\item Für $s=2$ und $W=R$ zieht $\ell^{*}$ Bilinearformen auf $V$
auf Bilinearformen auf $U$ zurück:
\[
(\ell^{*}\tau)(u_{1},u_{2})
=\tau(\ell(u_{1}),\ell(u_{2})).
\]
\end{itemize}
\end{CONC}



Proposition 7.19. 

\begin{PROP}{GA-L09-03-41}{Verträglichkeit von Pullback mit Tensorprodukt}
If $\lambda \in L(\mathrm{V}, \mathrm{V}), \alpha \in T_{s_{1}}^{0}(\mathrm{V})$ and $\beta \in T_{s_{2}}^{0}(\mathrm{V})$, then
$$
\lambda^{*}(\alpha \otimes \beta)=\lambda^{*} \alpha \otimes \lambda^{*} \beta .
$$
\end{PROP}

Proof. 

\begin{PROOF}{GA-L09-03-42}{P: Verträglichkeit von Pullback mit Tensorprodukt}
We have
$$
\begin{aligned}
\lambda^{*}(\alpha \otimes \beta) & \left(u_{1}, \ldots, u_{s_{1}}, u_{s_{1}+1}, \ldots, u_{s_{1}+s_{2}}\right) \\
& =\alpha \otimes \beta\left(\lambda u_{1}, \ldots, \lambda u_{s_{1}}, \lambda u_{s_{1}+1}, \ldots, \lambda u_{s_{1}+s_{2}}\right) \\
& =\alpha\left(\lambda u_{1}, \ldots, \lambda u_{s_{1}}\right) \beta\left(\lambda u_{s_{1}+1}, \ldots, \lambda u_{s_{1}+s_{2}}\right) \\
& =\lambda^{*} \alpha\left(u_{1}, \ldots, u_{s_{1}}\right) \lambda^{*} \beta\left(u_{s_{1}+1}, \ldots, u_{s_{1}+s_{2}}\right) \\
& =\lambda^{*} \alpha \otimes \lambda^{*} \beta\left(u_{1}, \ldots, u_{s_{1}}, u_{s_{1}+1}, \ldots, u_{s_{1}+s_{2}}\right) .
\end{aligned}
$$
\end{PROOF}



Before going on to study tensor fields, we introduce one more notion from multilinear algebra referred to as contraction.

Definition 7.20. 


\begin{DEF}{GA-L09-03-43}{Kontraktion von Tensoren in $(k,l)$}
Let $\left(e_{1}, \ldots, e_{n}\right)$ be a basis for V and $\left(e^{1}, \ldots, e^{n}\right)$ the dual basis. If $\tau \in T^{r}{ }_{s}(\mathrm{V})$, then for $k \leq r$ and $l \leq s$, we define $C_{l}^{k} \tau \in T^{r-1}{ }_{s-1}(\mathrm{V})$ by
$$
\begin{aligned}
C_{l}^{k} \tau( & \left.\theta^{1}, \ldots, \theta^{r-1}, \mathrm{w}_{1}, \ldots, \mathrm{w}_{s-1}\right) \\
& :=\sum_{a=1}^{n} \tau\left(\theta^{1}, \ldots, e_{k \text {-th position }}^{a}, \ldots, \theta^{r-1}, \mathrm{w}_{1}, \ldots, e_{a\quad l \text {-th position }}, \ldots, \mathrm{w}_{s-1}\right) .
\end{aligned}
$$
This processes is called contraction.

Write the components of $\tau$ with respect to our basis as $\tau^{i_{1} \ldots i_{r}}{ }_{j_{1} \ldots j_{s}}$. If we pick out an upper index, say $i_{k}$, and also a lower index, say $j_{l}$, then we obtain the components of the contracted tensor $C_{l}^{k} \tau$ by:
$$
\left(C_{l}^{k} \tau\right)^{i_{1} \ldots{\hat{i_{k}}}^{\ldots} i_{r}}{ }_{j_{1} \ldots \hat{j}_{l} \ldots j_{s}}:=\tau^{i \ldots a \ldots i_{r}}{ }_{j_{1} \ldots a \ldots j_{s}}(\text { sum over } a) .
$$
Here the caret means omission. In practice, one often just writes
$$
\tau^{i_{1} \ldots \hat{i_{k}} \ldots i_{r}}{ }_{j_{1} \ldots \hat{j_{l}} \ldots j_{s}}
$$
instead of $\left(C_{l}^{k} \tau\right)^{i_{1} \ldots i_{k} \ldots i_{r}}{ }_{j_{1} \ldots \hat{j}_{l} \ldots j_{s}}$ as long as it has been made clear how the contraction was carried out.
\end{DEF}


For example, one often sees expressions like $R_{i j}:=R_{i r j}^{r}$. Notice that we always contract an upper index with a lower index. The map $C_{l}^{k}$ contracts the $k$-th upper index with the $l$-th lower index.

Consider a tensor of the form $\mathrm{v} \otimes \mathrm{w} \otimes \eta \otimes \theta \in T_{2}^{2}(\mathrm{V})$. One can show that
$$
C_{1}^{1}(\mathrm{v} \otimes \mathrm{w} \otimes \eta \otimes \theta)=\eta(\mathrm{v}) \mathrm{w} \otimes \theta .
$$
Similarly,
$$
C_{2}^{1}(\mathrm{v} \otimes \mathrm{w} \otimes \eta \otimes \theta)=\theta(\mathrm{v}) \mathrm{w} \otimes \eta
$$

In general, $C_{l}^{k}$ acts on simple tensors $\mathrm{v}_{1} \otimes \mathrm{v}_{2} \otimes \cdots \otimes \mathrm{v}_{r} \otimes \eta^{1} \otimes \eta^{2} \otimes \cdots \otimes \eta^{s}$ by an obvious extension of the above. Universal mapping properties can be invoked to give a basis free definition of contraction. Contraction generalizes the notion of the trace of a linear transformation.



\begin{DEF}{GA-L09-03-44}{Allgemeine Definition der Kontraktion auf einfachen Tensoren}
Sei $V$ ein endlichdimensionaler $R$-Modul und 
\[
v_{1},\dots,v_{r}\in V,\qquad
\eta^{1},\dots,\eta^{s}\in V^{*}.
\]
Dann bezeichnet man mit
\[
C_{l}^{k}:
T^{r}{}_{s}(V)\longrightarrow T^{r-1}{}_{s-1}(V),
\qquad 1\le k\le r,\;1\le l\le s,
\]
die \emph{Kontraktion der $k$-ten kontravarianten mit der $l$-ten kovarianten Komponente},
definiert durch die Vorschrift
\[
C_{l}^{k}
\big(v_{1}\otimes\cdots\otimes v_{r}
\otimes\eta^{1}\otimes\cdots\otimes\eta^{s}\big)
=\eta^{l}(v_{k})\;
v_{1}\otimes\cdots\widehat{v_{k}}\cdots\otimes v_{r}
\otimes
\eta^{1}\otimes\cdots\widehat{\eta^{l}}\cdots\otimes\eta^{s}.
\]
Dabei bedeutet das Symbol $\widehat{(\cdot)}$, dass der entsprechende Faktor 
im Tensorprodukt ausgelassen wird.  

Diese Vorschrift ist eine \textbf{lineare Erweiterung} der zuvor gegebenen Beispiele, 
etwa
\[
C_{1}^{1}(v\otimes w\otimes\eta\otimes\theta)
=\eta(v)\,w\otimes\theta,
\qquad
C_{2}^{1}(v\otimes w\otimes\eta\otimes\theta)
=\theta(v)\,w\otimes\eta,
\]
und wird durch Linearität eindeutig auf ganz $T^{r}{}_{s}(V)$ fortgesetzt.

\textbf{Interpretation.}
$C_{l}^{k}$ „zieht“ ein kontravariantes und ein kovariantes Bein des Tensors zusammen,
indem die entsprechende lineare Form auf den passenden Vektor angewendet wird.
Dies reduziert den Typ des Tensors von $(r,s)$ auf $(r-1,s-1)$.
\end{DEF}





\begin{EXA}{GA-L09-03-45}{Kontraktion über Tensorprodukte}
A common use of contraction involves first taking the tensor product of two tensors, and then performing a contraction of a contravariant slot of one with a covariant slot of the other. One often performs several contractions. For example, we may form a tensor that is given in components as
$$
\tau^{a c}{ }_{e f g}=S^{a}{ }_{k e} T^{k c}{ }_{f g} \quad(\text { sum over } k) .
$$
\end{EXA}

Evaluation of a tensor on its arguments is the result of a repeated contraction. For example, let V have a basis $e_{1}, \ldots, e_{n}$ and let $e^{1}, \ldots, e^{n}$ be the dual basis for $\mathrm{V}^{*}$ as above. If $v=v^{i} e_{i}, w=w^{i} e_{i}$, and $\alpha=\alpha_{i} e^{i}$, then for $\tau \in T_{2}^{1}(\mathrm{V})$ we easily deduce that
$$
\tau(\alpha, v, w)=\tau^{i}{ }_{j k} \alpha_{i} v^{j} w^{k},
$$
which is the result of a repeated contraction on the tensor $\tau \otimes \alpha \otimes v \otimes w$. More generally, if we express elements $v_{1}, \ldots, v_{s} \in \mathrm{V}$ and $\alpha_{1}, \ldots, \alpha_{r} \in \mathrm{V}^{*}$ in terms of our basis and its dual, then for $\tau \in T^{r}{ }_{s}(\mathrm{V})$, we have an analogous general expression for $\tau\left(\alpha_{1}, \ldots, \alpha_{r}, v_{1}, \ldots, v_{s}\right)$ in terms of the components of the tensor and its arguments.



\begin{DEF}{GA-L09-03-46}{Auswertung eines Tensors als wiederholte Kontraktion}
Sei $V$ ein endlichdimensionaler freier $R$-Modul mit Basis $(e_{1},\dots,e_{n})$
und dualer Basis $(e^{1},\dots,e^{n})$ von $V^{*}$.
Für $r,s\ge0$ sei $\tau\in T^{r}{}_{s}(V)=V^{\otimes r}\otimes(V^{*})^{\otimes s}$ ein Tensor
vom Typ $(r,s)$.

\textbf{(1) Definition.}
Die \emph{Auswertung eines Tensors} auf Argumenten
\[
(\alpha_{1},\dots,\alpha_{r},v_{1},\dots,v_{s})
\in (V^{*})^{r}\times V^{s}
\]
ist definiert als das Ergebnis einer \textbf{wiederholten Kontraktion}
des Tensors $\tau\otimes \alpha_{1}\otimes\cdots\otimes \alpha_{r}
\otimes v_{1}\otimes\cdots\otimes v_{s}$,
wobei jede kovariante Komponente von $\tau$ mit der entsprechenden
kontravarianten Komponente eines Arguments kontrahiert wird:
\[
\tau(\alpha_{1},\dots,\alpha_{r},v_{1},\dots,v_{s})
=
C^{r+s}_{s}\circ\cdots\circ C^{r+1}_{1}\big(
\tau\otimes \alpha_{1}\otimes\cdots\otimes \alpha_{r}
\otimes v_{1}\otimes\cdots\otimes v_{s}\big).
\]

\textbf{(2) Beispiel.}
Seien $v,w\in V$ und $\alpha\in V^{*}$, sowie
\[
\tau = \tau^{i}{}_{jk}\, e_{i}\otimes e^{j}\otimes e^{k} \in T^{1}{}_{2}(V).
\]
Dann gilt
\[
\tau(\alpha,v,w)
=
\tau^{i}{}_{jk}\,\alpha_{i}\,v^{j}\,w^{k},
\]
was genau dem Ergebnis einer wiederholten Kontraktion auf
$\tau\otimes \alpha\otimes v\otimes w$ entspricht.

\textbf{(3) Allgemeine Komponentendarstellung.}
Für $\tau\in T^{r}{}_{s}(V)$ mit Komponenten $\tau^{i_{1}\dots i_{r}}{}_{j_{1}\dots j_{s}}$
und Argumente
\[
\alpha_{p}=\alpha^{(p)}_{i} e^{i}\in V^{*},
\qquad
v_{q}=v_{(q)}^{i} e_{i}\in V,
\]
ergibt die Auswertung
\[
\tau(\alpha_{1},\dots,\alpha_{r},v_{1},\dots,v_{s})
=
\tau^{i_{1}\dots i_{r}}{}_{j_{1}\dots j_{s}}\,
\alpha^{(1)}_{i_{1}}\cdots \alpha^{(r)}_{i_{r}}\,
v_{(1)}^{j_{1}}\cdots v_{(s)}^{j_{s}}.
\]

\textbf{(4) Interpretation.}
Die Auswertung eines Tensors auf seinen Argumenten entspricht
dem sukzessiven Anwenden aller kovarianten Komponenten
auf kontravariante Argumente und umgekehrt.
Algebraisch ist dies eine \emph{vollständige Kontraktion} 
von $\tau$ mit den entsprechenden Argumenten, deren Ergebnis
ein Element von $R$ (ein Skalar) ist.
\end{DEF}



\pagebreak



\subsection*{Bottom-Up Approach to Tensor Fields}
There are two approaches to tensor fields on smooth manifolds that turn out to be equivalent (at least for finite-dimensional manifolds). We start with the "bottom-up" approach where we apply multilinear algebra first to individual tangent spaces. The second approach directly defines tensors on $M$ as tensors on the module $\mathfrak{X}(M)$.

Roughly speaking, a smooth ( $r, s$ )-tensor field on a manifold $M$ assigns to each $p \in M$ an element of $T_{s}^{r}\left(T_{p} M\right)$ in a smooth way. We are interested in making sense of smoothness for tensor fields, so we wish to view a tensor field as a section of an appropriate vector bundle (a tensor bundle). Let us start out being a bit more general by considering a real rank $k$ vector bundle $\xi=(E, \pi, M)$. For convenience, we take the typical fiber to be $\mathbb{R}^{k}$. Let $T^{r}{ }_{s}(E)=\bigcup_{p \in M} T^{r}{ }_{s}\left(E_{p}\right)$. We wish to construct a bundle $T^{r}{ }_{s}(\xi)$ which has $T^{r}{ }_{s}(E)$ as total space, $M$ as base space, and $T^{r}{ }_{s}\left(E_{p}\right)$ as fiber over $p$. If ( $U, \phi$ ) is a VB-chart for $\xi$, then we construct a VB-chart for $T_{s}^{r}(\xi)$ in the following way: Recall that $\phi$ has the form $\phi=(\pi, \Phi)$, where $\Phi: \pi^{-1} U \rightarrow \mathbb{R}^{k}$ and where $\Phi_{p}:=\left.\Phi\right|_{E_{p}}: E_{p} \rightarrow \mathbb{R}^{k}$ is a linear isomorphism for each $p$. We obtain a map $\Phi_{p}^{r, s}: T^{r}{ }_{s}\left(E_{p}\right) \rightarrow T^{r}{ }_{s}\left(\mathbb{R}^{k}\right)$ by

$$
\begin{aligned}
\left(\Phi_{p}^{r, s} \tau_{p}\right) & \left(\alpha_{1}, \ldots, \alpha_{r}, v_{1}, \ldots, v_{s}\right) \\
& :=\tau_{p}\left(\left(\Phi_{p}\right)^{*} \alpha_{1}, \ldots,\left(\Phi_{p}\right)^{*} \alpha_{r}, \Phi_{p}^{-1} v_{1}, \ldots, \Phi_{p}^{-1} v_{s}\right)
\end{aligned}
$$

These maps combine to give a map $\Phi^{r, s}: \pi^{-1} U \rightarrow T_{s}^{r}\left(\mathbb{R}^{k}\right)$ which is smooth (exercise). Our chart for $T_{s}^{r}(\xi)$ is

$$
\phi^{r, s}:=\left(\pi, \Phi^{r, s}\right): \pi^{-1} U \rightarrow U \times T_{s}^{r}\left(\mathbb{R}^{k}\right)
$$

If desired, one can choose, once and for all, an isomorphism $T_{s}^{r}\left(\mathbb{R}^{k}\right) \cong \mathbb{R}^{k^{r+s}}$. A VB-atlas $\left\{\left(U_{\alpha}, \phi_{\alpha}\right)\right\}$ for $\xi=(E, \pi, M)$ gives a VB-atlas $\left\{\left(U_{\alpha}, \phi_{\alpha}^{r, s}\right)\right\}$ for $T^{r}{ }_{s}(\xi)$.




\begin{DEF}{GA-L09-04-01}{Tensorbündel eines Vektorbündels}
Sei $\xi = (E,\pi,M)$ ein reelles Vektorbündel vom Rang $k$ über einer glatten Mannigfaltigkeit $M$.
Für jedes $p\in M$ bezeichne $E_p = \pi^{-1}(p)$ den Vektorraum in der Faser über $p$.

\textbf{(1) Faserweise Tensorräume.}
Für jedes $p\in M$ sei
\[
T^{r}{}_{s}(E_p)
:= E_p^{\otimes r} \otimes (E_p^{*})^{\otimes s}
\]
der Raum der $(r,s)$-Tensoren auf der Faser $E_p$.
Das disjunkte Vereinigungsbündel
\[
T^{r}{}_{s}(E)
:= \bigsqcup_{p\in M} T^{r}{}_{s}(E_p)
\]
bildet die \emph{Gesamtmenge} des Tensorbündels.

\textbf{(2) Lokale Trivialisierungen.}
Sei $(U,\phi)$ eine lokale Vektorbündelkarte für $\xi$, d.\,h.
\[
\phi = (\pi,\Phi): \pi^{-1}(U) \longrightarrow U \times \mathbb{R}^{k},
\]
wobei für jedes $p\in U$ die Abbildung
$\Phi_p := \Phi|_{E_p} : E_p \to \mathbb{R}^{k}$ ein linearer Isomorphismus ist.

Daraus wird eine kanonische Abbildung
\[
\Phi_p^{r,s} : T^{r}{}_{s}(E_p)
\longrightarrow
T^{r}{}_{s}(\mathbb{R}^{k})
\]
definiert durch
\[
\big(\Phi_p^{r,s}\tau_p\big)(\alpha_1,\dots,\alpha_r,v_1,\dots,v_s)
:=
\tau_p\!\big((\Phi_p)^{*}\alpha_1,\dots,(\Phi_p)^{*}\alpha_r,
\Phi_p^{-1}v_1,\dots,\Phi_p^{-1}v_s\big).
\]

Diese Abbildungen $\Phi_p^{r,s}$ hängen glatt von $p$ ab und liefern damit eine glatte Abbildung
\[
\Phi^{r,s} : \pi^{-1}(U) \longrightarrow T^{r}{}_{s}(\mathbb{R}^{k}).
\]

\textbf{(3) Lokale Koordinatenkarte des Tensorbündels.}
Die zugehörige lokale Trivialisierung des Tensorbündels ist
\[
\phi^{r,s} := (\pi,\Phi^{r,s}) :
\pi^{-1}(U) \longrightarrow U \times T^{r}{}_{s}(\mathbb{R}^{k}),
\]
und bildet damit eine Vektorbündelkarte für $T^{r}{}_{s}(\xi)$.

\textbf{(4) Atlas und typische Faser.}
Ein VB-Atlas $\{(U_\alpha,\phi_\alpha)\}$ für $\xi$
induziert den VB-Atlas $\{(U_\alpha,\phi_\alpha^{r,s})\}$ für $T^{r}{}_{s}(\xi)$.
Die typische Faser ist $T^{r}{}_{s}(\mathbb{R}^{k})$, die, falls gewünscht,
mit $\mathbb{R}^{k^{r+s}}$ identifiziert werden kann.

\textbf{(5) Tensorfeld.}
Ein \emph{glattes $(r,s)$-Tensorfeld} auf $M$
ist eine glatte Schnittabbildung
\[
\tau : M \longrightarrow T^{r}{}_{s}(\xi),
\qquad \pi\circ\tau = \mathrm{id}_M,
\]
die jedem Punkt $p\in M$ einen Tensor $\tau(p)\in T^{r}{}_{s}(E_p)$
auf der Faser zuordnet.
\end{DEF}





Exercise 7.21. 

\begin{PROP}{GA-L09-04-02}{Natürliche Isomorphie $T_{s}^{r}(\xi) \cong\left(\otimes^{r} E\right) \otimes\left(\otimes^{s} E^{*}\right)$}
There is a natural vector bundle isomorphism
$$
T_{s}^{r}(\xi) \cong\left(\otimes^{r} E\right) \otimes\left(\otimes^{s} E^{*}\right)
$$
\end{PROP}

We leave it to the interested reader to prove the following useful theorem.


Proposition 7.22. 


\begin{PROP}{GA-L09-04-03}{Charakterisierung von glatten Schnitten in Tensorbündeln}
Let $\xi=(E, \pi, M)$ be a vector bundle as above and let $\Upsilon: M \rightarrow T^{r}{ }_{s}(E)$ be a map which assigns to each $p \in M$ an element of $T_{s}^{r}\left(E_{p}\right)$. Then $\Upsilon$ is smooth if and only if
$$
p \mapsto \Upsilon(p)\left(\alpha_{1}(p), \ldots, \alpha_{r}(p), X_{1}(p), \ldots, X_{s}(p)\right)
$$
is smooth for all smooth sections $p \mapsto \alpha_{i}(p)$ and $p \mapsto X_{i}(p)$ of $E^{*} \rightarrow M$ and $E \rightarrow M$ respectively. The same statement is true if we use local sections.

The set of smooth sections of $\left.T^{r}{ }_{s}(\xi)\right)$ is denoted $\Gamma\left(T^{r}{ }_{s}(\xi)\right)$.
\end{PROP}


\begin{DEF}{GA-L09-04-04}{Auswertungen von glatten Tensorfeldern}
If $\Upsilon \in \Gamma\left(T^{r}{ }_{s}(\xi)\right)$, then for $X_{1}, \ldots, X_{s}$ any smooth sections of $E \rightarrow M$ and $\alpha_{1}, \ldots, \alpha_{r}$ smooth sections of $E^{*} \rightarrow M$, define $\Upsilon\left(\alpha_{1}, \ldots, \alpha_{r}, X_{1}, \ldots, X_{s}\right) \in C^{\infty}(M)$ by
$$
\Upsilon\left(\alpha_{1}, \ldots, \alpha_{r}, X_{1}, \ldots, X_{s}\right)(p):=\Upsilon_{p}\left(\alpha_{1}(p), \ldots, \alpha_{r}(p), X_{1}(p), \ldots, X_{s}(p)\right) .
$$
Now we have a map 
$$\Upsilon:\left(\Gamma E^{*}\right)^{k} \times(\Gamma E)^{l} \rightarrow C^{\infty}(M)$$ 
This map is clearly multilinear over $C^{\infty}(M)$, and we see that we can interpret elements of $\Gamma\left(T^{r}{ }_{s}(\xi)\right)$ as such maps when convenient. This extends the idea of thinking of a 1 -form $\alpha$ as a $C^{\infty}(M)$ linear map $\mathfrak{X}(M) \rightarrow C^{\infty}(M)$.
\end{DEF}

For $\tau \in \Gamma\left(T_{s_{1}}^{r_{1}}(\xi)\right)$ and $\eta \in \Gamma\left(T_{s_{2}}^{r_{2}}(\xi)\right)$, we define the (consolidated) tensor product $\tau \otimes \eta \in \Gamma\left(T^{r_{1}+r_{2}}{ }_{s_{1}+s_{2}}(\xi)\right)$ by $(\tau \otimes \eta)(p):=\tau_{p} \otimes \eta_{p}$. Thus
$$
\begin{aligned}
& (\tau \otimes \eta)(p)\left(\alpha^{1}, \ldots, \alpha^{r_{1}+r_{2}}, v_{1}, \ldots, v_{s_{1}+s_{2}}\right) \\
& =\tau\left(\alpha^{1}, \ldots, \alpha^{r_{1}}, v_{1}, \ldots, v_{s_{1}}\right) \eta\left(\alpha^{r_{1}+1}, \ldots, \alpha^{r_{1}+r_{2}}, v_{s_{1}+1}, \ldots, v_{s_{1}+s_{2}}\right)
\end{aligned}
$$
for all $\alpha^{i} \in E_{p}^{*}$ and $v_{i} \in E_{p}$.


Let $\left(s_{1}, \ldots, s_{k}\right)$ be a local frame field for $\xi$ over an open set $U$ and let $\sigma^{1}, \ldots, \sigma^{k}$ be the dual frame field of the dual bundle $E^{*} \rightarrow M$ so that $\sigma^{i}\left(s_{j}\right)=\delta_{j}^{i}$. Consider the set
$$
\left\{\sigma^{i_{1}} \otimes \cdots \otimes \sigma^{i_{r}} \otimes s_{j_{1}} \otimes \cdots \otimes s_{j_{s}}: i_{1}, \ldots, i_{r}, j_{1}, \ldots, j_{s}=1, \ldots, k\right\}
$$
If $\tau \in \Gamma\left(T^{r}{ }_{s}(\xi)\right)$, then we have functions 
$$\tau^{i_{1} \ldots i_{r}}{ }_{j_{1} \ldots j_{s}} \in C^{\infty}(U)$$ defined by 
$$\tau^{i_{1} \ldots i_{r}}{ }_{j_{1} \ldots j_{s}}=\tau\left(\sigma^{i_{1}}, \ldots, \sigma^{i_{r}}, s_{j_{1}}, \ldots, s_{j_{s}}\right)$$ 
It follows from Proposition 7.8 that $\tau$ (restricted to $U$ ) has the expansion
$$
\tau=\tau^{i_{1} \ldots i_{r}}{ }_{j_{1} \ldots j_{s}} \sigma^{i_{1}} \otimes \cdots \otimes \sigma^{i_{r}} \otimes s_{j_{1}} \otimes \cdots \otimes s_{j_{s}} .
$$
Also, applying equation (7.2) in each fiber $E_{p}$, we see that the component functions for $\tau \otimes \eta$ are given by
$$
\begin{aligned}
& (\tau \otimes \eta)^{i_{1} \ldots i_{r_{1}+r_{2}}}{ }_{j_{1} \ldots j_{s_{1}+s_{2}}} \\
& =\tau^{i_{1} \ldots i_{r_{1}}}{ }_{j_{1} \ldots j_{s_{1}}} \eta^{i_{r_{1}+1} \ldots i_{r_{2}}}{ }_{j_{s_{1}+1} \ldots j_{s_{2}}} .
\end{aligned}
$$
Here, and wherever convenient, we use the consolidated tensor product.




\begin{DEF}{GA-L09-04-05}{Tensorprodukt von Tensorfeldern}
Sei $\xi=(E,\pi,M)$ ein reelles Vektorbündel vom Rang $k$ über einer glatten Mannigfaltigkeit $M$.
Für $r_{1},s_{1},r_{2},s_{2}\in\mathbb{N}_{0}$ seien
\[
\tau \in \Gamma\!\left(T^{r_{1}}{}_{s_{1}}(\xi)\right),
\qquad
\eta \in \Gamma\!\left(T^{r_{2}}{}_{s_{2}}(\xi)\right)
\]
glatte Tensorfelder auf $\xi$.

\textbf{(1) Faserweise Definition.}
Das \emph{(konsolidierte) Tensorprodukt} von $\tau$ und $\eta$ ist das Tensorfeld
\[
\tau\otimes\eta \in \Gamma\!\left(T^{r_{1}+r_{2}}{}_{s_{1}+s_{2}}(\xi)\right),
\]
definiert durch die Vorschrift
\[
(\tau\otimes\eta)(p) := \tau_{p} \otimes \eta_{p}, \qquad p\in M,
\]
wobei das Tensorprodukt auf der Faser $E_{p}$ verstanden wird.

Für $\alpha^{i}\in E_{p}^{*}$ und $v_{i}\in E_{p}$ gilt explizit:
\[
\begin{aligned}
&(\tau\otimes\eta)(p)
  (\alpha^{1},\ldots,\alpha^{r_{1}+r_{2}},v_{1},\ldots,v_{s_{1}+s_{2}}) \\
&\quad=
\tau_{p}(\alpha^{1},\ldots,\alpha^{r_{1}},v_{1},\ldots,v_{s_{1}})
\;\cdot\;
\eta_{p}(\alpha^{r_{1}+1},\ldots,\alpha^{r_{1}+r_{2}},v_{s_{1}+1},\ldots,v_{s_{1}+s_{2}}).
\end{aligned}
\]

\textbf{(2) Lokale Darstellung.}
Sei $(s_{1},\ldots,s_{k})$ ein lokales Rahmenfeld des Bündels $\xi$ über einem offenen $U\subset M$,
und $(\sigma^{1},\ldots,\sigma^{k})$ das zugehörige duale Rahmenfeld des Dualbündels $E^{*}$, d.\,h.
\[
\sigma^{i}(s_{j})=\delta^{i}_{j}.
\]
Dann lässt sich ein Tensorfeld
\[
\tau \in \Gamma\!\left(T^{r}{}_{s}(\xi)\right)
\]
auf $U$ eindeutig in der Form entwickeln:
\[
\tau
=
\tau^{i_{1}\ldots i_{r}}{}_{j_{1}\ldots j_{s}}
\;
\sigma^{i_{1}}\otimes\cdots\otimes\sigma^{i_{r}}
\otimes
s_{j_{1}}\otimes\cdots\otimes s_{j_{s}},
\]
wobei die \emph{Komponentenfunktionen}
\[
\tau^{i_{1}\ldots i_{r}}{}_{j_{1}\ldots j_{s}}(p)
=
\tau_{p}(\sigma^{i_{1}},\ldots,\sigma^{i_{r}},s_{j_{1}},\ldots,s_{j_{s}})
\in C^{\infty}(U)
\]
glatte Funktionen sind.

\textbf{(3) Komponentenform des Tensorprodukts.}
Seien nun
\[
\tau \in \Gamma\!\left(T^{r_{1}}{}_{s_{1}}(\xi)\right),
\qquad
\eta \in \Gamma\!\left(T^{r_{2}}{}_{s_{2}}(\xi)\right),
\]
mit lokalen Komponenten
$\tau^{i_{1}\ldots i_{r_{1}}}{}_{j_{1}\ldots j_{s_{1}}}$ und
$\eta^{i_{1}\ldots i_{r_{2}}}{}_{j_{1}\ldots j_{s_{2}}}$.
Dann hat das Tensorprodukt $\tau\otimes\eta$
die Komponenten
\[
(\tau\otimes\eta)^{i_{1}\ldots i_{r_{1}+r_{2}}}{}_{j_{1}\ldots j_{s_{1}+s_{2}}}
=
\tau^{i_{1}\ldots i_{r_{1}}}{}_{j_{1}\ldots j_{s_{1}}}
\;
\eta^{i_{r_{1}+1}\ldots i_{r_{1}+r_{2}}}{}_{j_{s_{1}+1}\ldots j_{s_{1}+s_{2}}},
\]
wobei die Indizes gemäß der Konsolidierung konkateniert werden.

\textbf{(4) Eigenschaften.}
Das Tensorprodukt von Tensorfeldern erfüllt:
\begin{itemize}
\item[(i)] Bilinearität:
\[
(a\tau_{1}+b\tau_{2})\otimes\eta
= a(\tau_{1}\otimes\eta)+b(\tau_{2}\otimes\eta),
\qquad
\tau\otimes(a\eta_{1}+b\eta_{2})
= a(\tau\otimes\eta_{1})+b(\tau\otimes\eta_{2}).
\]
\item[(ii)] Assoziativität:
\[
(\tau\otimes\eta)\otimes\zeta
=\tau\otimes(\eta\otimes\zeta).
\]
\item[(iii)] Kompatibilität mit der Faserstruktur:
$(\tau\otimes\eta)(p)=\tau_{p}\otimes\eta_{p}$ für alle $p\in M$.
\end{itemize}
\end{DEF}




Notation 7.23. Whenever there is no chance of confusion, we will refer to $T^{r}{ }_{s}(\xi)$ by $T^{r}{ }_{s}(E) \rightarrow M$ or even just $T^{r}{ }_{s}(E)$ (the latter is the notation for the total space of the bundle).

In the case of the tangent bundle TM, we have special terminology and notation:

Definition 7.24. 


\begin{DEF}{GA-L09-04-06}{$(r,s)$-Tensorbündel auf Mannigfaltigkeiten}
The bundle $T^{r}{ }_{s}(T M) \rightarrow M$ is called the $(r, s)$-tensor bundle on $M$.

By Exercise 7.21, $T_{s}^{r}(T M) \cong\left(\bigotimes^{r} T M\right) \otimes\left(\bigotimes^{s} T^{*} M\right)$ and this natural isomorphism is taken as an identification so the latter bundle is also referred to as a tensor bundle.
\end{DEF}



We now restrict ourselves to the case of the tangent bundle of a manifold but note that much of what follows makes sense for general vector bundles.

Definition 7.25. 

\begin{DEF}{GA-L09-04-07}{Raum der $(r,s)$-Tensorfelder auf Mannigfaltigkeiten}
The space of sections $\Gamma\left(T^{r}{ }_{s}(T M)\right)$ is denoted by $\mathcal{T}_{s}^{r}(M)$ and its elements are referred to as $r$-contravariant $s$-covariant tensor fields or just type ( $r, s$ )-tensor fields. The space $\mathcal{T}_{0}^{r}(M)$ is denoted by $\mathcal{T}^{r}(M)$ and $\mathcal{T}_{s}^{0}(M)$ by $\mathcal{T}_{s}(M)$.

In summary, a smooth tensor field $A$ is a smooth assignment of a multilinear map on each tangent space of the manifold. Thus for each $p, A(p)$ is a multilinear map
$$
A(p):\left(T_{p}^{*} M\right)^{r} \times\left(T_{p} M\right)^{s} \rightarrow \mathbb{R}
$$
or in other words, an element of $T_{s}^{r}\left(T_{p} M\right)$. Elements of $T_{s}^{r}\left(T_{p} M\right)$ are called tensors at $p$. We also write $A_{p}$ for $A(p)$.
\end{DEF}

Example 7.26. 

\begin{EXA}{GA-L09-04-08}{Riemannische Metrik auf reellen Vektorbündeln}
In Definition 6.42, we introduced the notion of a Riemannian metric on a real vector bundle. We saw that such metrics always exist. The most important case is where the bundle is the tangent bundle $T M$ of a manifold $M$. In this case, we say that we have a Riemannian metric on $M$. Thus a Riemannian metric on $M$ is an element of $\mathcal{T}_{2}(M)$ which is symmetric and positive definite at each point.
\end{EXA}

\begin{EXA}{GA-L09-04-09}{Einschränkung von Tensorfeldern auf offene Umgebungen}
Of course the manifold in question could be an open submanifold $U$ of $M$ so we have $C^{\infty}(U)$-module $(r, s)$-tensor fields over that set denoted $\mathcal{T}_{s}^{r}(U)$. The open subsets are partially ordered by inclusion $V \subset U$ and the tensor fields on these are related by restriction. Let $r_{V}^{U}: \mathcal{T}_{s}^{r}(U) \rightarrow \mathcal{T}_{s}^{r}(V)$ denote the restriction map. The assignment $U \rightarrow \mathcal{T}_{s}^{r}(U)$ is an example of a presheaf and in fact a sheaf.
\end{EXA}

\begin{EXA}{GA-L09-04-10}{$TM$-wertige Tensoren}
We will also sometimes deal with tensors with values in $T M$ (or in $\left.T^{*} M\right)$. First note that the space $T^{r}{ }_{s}\left(T_{p} M ; T_{p} M\right)$ of all multilinear maps $\left(T_{p}^{*} M\right)^{r} \times\left(T_{p} M\right)^{s} \rightarrow T_{p} M$ is a vector space. The set $T_{s}^{r}(T M ; T M):= \bigcup_{p} T_{s}^{r}\left(T_{p} M ; T_{p} M\right)$ can be given a smooth vector bundle structure in a way that is closely analogous to $T_{s}^{r}(T M) \rightarrow M$.
\end{EXA}

Definition 7.27. 

\begin{DEF}{GA-L09-04-11}{Raum der $TM$-wertigen $(r,s)$ Tensorfelder}
The space of sections $\Gamma\left(T^{r}{ }_{s}(T M ; T M)\right)$ is denoted by $\mathcal{T}_{s}^{r}(M ; T M)$ and its elements are referred to as $r$-contravariant $s$-covariant $T M$-valued tensor fields. Similarly, we may define $T^{*} M$-valued tensor fields.
\end{DEF}

Note that $T M$-valued tensor fields can be associated in a natural manner with ordinary tensor fields. For example, if $A \in T_{2}^{0}(T M ; T M)$, then using the same letter $A$ by abuse of notation, we may define an element $A \in T_{2}^{1}(T M)$ by
$$
A_{p}\left(\theta_{p}, v_{p}, w_{p}\right)=\theta_{p}\left(A_{p}\left(v_{p}, w_{p}\right)\right) \text { for } \theta_{p} \in T_{p}^{*} M \text { and } v_{p}, w_{p} \in T_{p} M
$$
Many such reinterpretations are possible. For this reason, we shall stick to studying ordinary tensors and tensor fields in what follows.

We shall define several operations on spaces of tensor fields. We would like each of these to be natural with respect to restriction. We already have one such operation: the tensor product. If $A \in \mathcal{T}_{s_{1}}^{r_{1}}(U)$ and $B \in \mathcal{T}_{s_{2}}^{r_{2}}(U)$ and $V \subset U$, then $r_{V}^{U}(A \otimes B)=r_{V}^{U} A \otimes r_{V}^{U} B$. A ( $k, l$ )-tensor field $A$ may generally be expressed in a chart ( $U, \mathrm{x}$ ) as

$$
A=A^{i_{1} \ldots i_{r}}{ }_{j_{1} \ldots j_{s}} \frac{\partial}{\partial x^{i_{1}}} \otimes \cdots \otimes \frac{\partial}{\partial x^{i_{r}}} \otimes d x^{j_{1}} \otimes \cdots \otimes d x^{j_{s}}
$$

where $A^{i_{1} \ldots i_{r}}{ }_{j_{1} \ldots j_{s}}$ are functions, $\frac{\partial}{\partial x^{i}} \in \mathfrak{X}(U)$ and $d x^{j} \in \mathfrak{X}^{*}(U)$. Actually, it is the restriction of $\tau$ to $U$ that can be written in this way, but because of the naturality of all the operations we introduce, it is generally safe to use the same letter to denote a tensor field and its restriction to an open set such as a chart domain. It is easy to show that

$$
A^{i_{1} \ldots i_{r}}{ }_{j_{1} \ldots j_{s}}=A\left(d x^{i_{1}}, \ldots, d x^{i_{r}}, \frac{\partial}{\partial x^{j_{1}}}, \ldots, \frac{\partial}{\partial x^{j_{s}}}\right),
$$

and so the components of a smooth tensor field are $C^{\infty}$ for every choice of coordinates ( $U, \mathrm{x}$ ) by Proposition 7.22. Conversely, one can obviously define tensors that are not necessarily smooth sections of the appropriate tensor bundle, and then a tensor will be smooth exactly when its components with respect to every chart in an atlas are smooth. Evaluating the expression (7.4) above at a point $p \in U$ results in an expression such as that given at the end of Example 7.10.




\begin{DEF}{GA-L09-04-12}{Lokale Darstellung glatter Tensorfelder und Natürlichkeit von Operationen}
Sei $M$ eine glatte Mannigfaltigkeit und
\[
\mathcal{T}^{r}{}_{s}(U)
:= \Gamma\!\left(T^{r}{}_{s}(TM|_{U})\right)
\]
der Raum der glatten $(r,s)$-Tensorfelder auf einer offenen Menge $U\subset M$.

\textbf{(1) Natürlichkeit von Operationen unter Restriktion.}
Eine Operation $\mathcal{O}$ auf Tensorfeldern heißt \emph{natürlich bezüglich Restriktion}, wenn gilt:
\[
r^{U}_{V}\!\big(\mathcal{O}(A,B)\big)
= \mathcal{O}\!\big(r^{U}_{V}A,\,r^{U}_{V}B\big),
\qquad
\forall\, V\subset U.
\]
Beispielsweise ist das Tensorprodukt natürlich, da für alle
\[
A \in \mathcal{T}^{r_{1}}{}_{s_{1}}(U), \quad
B \in \mathcal{T}^{r_{2}}{}_{s_{2}}(U)
\]
und $V\subset U$ gilt:
\[
r^{U}_{V}(A\otimes B)
= (r^{U}_{V}A)\otimes(r^{U}_{V}B).
\]

\textbf{(2) Lokale Darstellung eines Tensorfeldes.}
Sei $(U,x)$ eine glatte Karte von $M$ mit lokalen Koordinaten
\[
x=(x^{1},\dots,x^{n}).
\]
Dann bilden
\[
\left\{\frac{\partial}{\partial x^{i}}\right\}_{i=1}^{n}
\subset \mathfrak{X}(U),
\qquad
\{dx^{j}\}_{j=1}^{n}
\subset \mathfrak{X}^{*}(U)
\]
lokale Rahmenfelder des Tangential- bzw. Kotangentialbündels über $U$.

Jedes glatte Tensorfeld
\[
A \in \mathcal{T}^{r}{}_{s}(U)
\]
kann eindeutig in der Form
\[
A
=
A^{i_{1}\ldots i_{r}}{}_{j_{1}\ldots j_{s}}\,
\frac{\partial}{\partial x^{i_{1}}}\otimes\cdots\otimes
\frac{\partial}{\partial x^{i_{r}}}
\otimes
dx^{j_{1}}\otimes\cdots\otimes dx^{j_{s}}
\tag{7.4}
\]
geschrieben werden, wobei die \emph{Komponentenfunktionen}
\[
A^{i_{1}\ldots i_{r}}{}_{j_{1}\ldots j_{s}}:U\to\mathbb{R}
\]
glatte Funktionen sind.

\textbf{(3) Berechnung der Komponenten.}
Die Komponenten ergeben sich durch Auswertung von $A$ auf den dualen Basisfeldern:
\[
A^{i_{1}\ldots i_{r}}{}_{j_{1}\ldots j_{s}}
=
A\!\left(dx^{i_{1}},\ldots,dx^{i_{r}},
\frac{\partial}{\partial x^{j_{1}}},\ldots,
\frac{\partial}{\partial x^{j_{s}}}\right).
\]
Somit sind die Komponenten eines glatten Tensorfeldes stets $C^{\infty}$-Funktionen
in jedem Koordinatensystem $(U,x)$.

\textbf{(4) Umkehrschluss.}
Eine Abbildung $A$ der Form \eqref{7.4}, deren Komponenten
$A^{i_{1}\ldots i_{r}}{}_{j_{1}\ldots j_{s}}$
in allen Karten eines Atlas glatt sind, definiert umgekehrt
ein glattes Tensorfeld,
d.\,h.
\[
A \in \Gamma\!\left(T^{r}{}_{s}(TM)\right).
\]

\textbf{(5) Bemerkung zur Notation.}
Da alle relevanten Operationen natürlich bezüglich Restriktion sind,
wird in der Praxis derselbe Buchstabe $A$ sowohl für das globale Tensorfeld
als auch für seine Einschränkung auf ein offenes Teilgebiet $U$
verwendet.
\end{DEF}


Exercise 7.28 (Transformation laws). Suppose that we have two charts ( $U, \mathrm{x}$ ) and ( $V, \overline{\mathrm{x}}$ ). If $A \in \mathcal{T}_{2}^{1}(M)$ has components $A_{j k}^{i}$ in the first chart and $\bar{A}_{j k}^{i}$ in the second chart, then on the overlap $U \cap V$ we have

$$
\bar{A}_{j k}^{i}=A_{b c}^{a} \frac{\partial \bar{x}^{i}}{\partial x^{a}} \frac{\partial x^{b}}{\partial \bar{x}^{j}} \frac{\partial x^{c}}{\partial \bar{x}^{k}},
$$

where

$$
d \bar{x}^{i}=\frac{\partial \bar{x}^{i}}{\partial x^{a}} d x^{a} \text { and } \frac{\partial}{\partial \bar{x}^{i}}=\frac{\partial x^{b}}{\partial \bar{x}^{i}} \frac{\partial}{\partial x^{b}} .
$$

This last exercise reveals the transformation law for tensor fields in $\mathcal{T}_{2}^{1}(M)$, and there is obviously an analogous law for tensor fields from $\mathcal{T}_{s}^{r}(M)$ for any values of $r$ and $s$. 



\begin{DEF}{GA-L09-04-13}{Allgemeines Transformationsgesetz für Tensorfeld-Komponenten}
Sei $M$ eine glatte $n$-dimensionale Mannigfaltigkeit und
\[
\mathcal{T}^{r}{}_{s}(M):=\Gamma\!\big(T^{r}{}_{s}(TM)\big)
\]
der Raum glatter $(r,s)$-Tensorfelder. 
Seien $(U,x)$ und $(V,\bar x)$ zwei Karten mit lokalen Koordinaten
\[
x=(x^{1},\dots,x^{n}),\qquad 
\bar x=(\bar x^{1},\dots,\bar x^{n}),
\]
und Jacobimatrizen
\[
J=(J^{i}{}_{a})=\Big(\frac{\partial \bar x^{i}}{\partial x^{a}}\Big), 
\qquad
J^{-1}=(K^{b}{}_{j})=\Big(\frac{\partial x^{b}}{\partial \bar x^{j}}\Big)
\]
auf $U\cap V$.
Die zugehörigen Basisfelder transformieren wie üblich:
\[
d\bar x^{i}=J^{i}{}_{a}\,dx^{a}, 
\qquad
\frac{\partial}{\partial \bar x^{j}}
=K^{b}{}_{j}\,\frac{\partial}{\partial x^{b}}.
\]

\textbf{(1) Lokale Darstellung.}
Ein Tensorfeld $A\in\mathcal{T}^{r}{}_{s}(M)$ besitzt in $(U,x)$ die Darstellung
\[
A
=
A^{i_{1}\ldots i_{r}}{}_{j_{1}\ldots j_{s}}\;
\frac{\partial}{\partial x^{i_{1}}}\otimes\cdots\otimes\frac{\partial}{\partial x^{i_{r}}}
\otimes
dx^{j_{1}}\otimes\cdots\otimes dx^{j_{s}},
\]
und in $(V,\bar x)$ entsprechend
\[
A
=
\bar A^{i_{1}\ldots i_{r}}{}_{j_{1}\ldots j_{s}}\;
\frac{\partial}{\partial \bar x^{i_{1}}}\otimes\cdots\otimes\frac{\partial}{\partial \bar x^{i_{r}}}
\otimes
d\bar x^{j_{1}}\otimes\cdots\otimes d\bar x^{j_{s}}.
\]

\paragraph{(2) Transformationsgesetz (allgemeiner Fall).}
Auf dem Überlapp $U\cap V$ hängen die Komponenten durch
\[
\boxed{\quad
\bar A^{i_{1}\ldots i_{r}}{}_{j_{1}\ldots j_{s}}
=
A^{a_{1}\ldots a_{r}}{}_{b_{1}\ldots b_{s}}\;
J^{i_{1}}{}_{a_{1}}\cdots J^{i_{r}}{}_{a_{r}}\;
K^{b_{1}}{}_{j_{1}}\cdots K^{b_{s}}{}_{j_{s}}
\quad}
\]
zusammen, d.\,h.\ jeder \emph{kontravariable} Index (oben) trägt einen Faktor $J=\partial\bar x/\partial x$,
und jeder \emph{kovariable} Index (unten) trägt einen Faktor $K=J^{-1}=\partial x/\partial\bar x$.

\textbf{(3) Beispiele.}
\begin{itemize}
\item Vektorfeld ($r=1,s=0$): $\ \bar X^{i}=J^{i}{}_{a}\,X^{a}$.
\item Kovektorfeld ($r=0,s=1$): $\ \bar\alpha_{\,j}=K^{b}{}_{j}\,\alpha_{b}$.
\item $(1,1)$-Tensor: $\ \bar T^{i}{}_{j}=J^{i}{}_{a}\,K^{b}{}_{j}\,T^{a}{}_{b}$.
\item $(1,2)$-Tensor: $\ \bar A^{i}{}_{jk}=J^{i}{}_{a}\,K^{b}{}_{j}\,K^{c}{}_{k}\,A^{a}{}_{bc}$.
\end{itemize}

\textbf{(4) Äquivalente Charakterisierungen.}
\begin{itemize}
\item \textbf{Basisbasierte Definition:}
Das Gesetz in (2) ist äquivalent zur Forderung, dass $A$ als multilineare Abbildung
\[
A:\underbrace{\mathfrak{X}^{*}(M)\times\cdots\times\mathfrak{X}^{*}(M)}_{r\text{ mal}}
\times
\underbrace{\mathfrak{X}(M)\times\cdots\times\mathfrak{X}(M)}_{s\text{ mal}}
\to C^{\infty}(M)
\]
unverändert bleibt, während die Basisfelder gemäß $d\bar x=J\,dx$ und 
$\partial/\partial\bar x=K\,\partial/\partial x$ transformieren.
\item \textbf{Bündelbeschreibung:}
Das Transformationsgesetz ist genau die durch die Übergangsfunktionen
$\mathrm{GL}(n,\mathbb{R})$-äquivarisierte Darstellung im Tensorbündel
$T^{r}{}_{s}(TM)$: die $r$-fache Standarddarstellung auf $TM$ und die $s$-fache
Dualdarstellung auf $T^{*}M$ wirken als
$J^{\otimes r}\otimes (J^{-1})^{\otimes s}$ auf den Komponenten.
\end{itemize}

\textbf{(5) Funktorielle/Naturalitäts-Eigenschaft.}
Bei einer Koordinatenverkettung $(x\to \bar x\to \tilde x)$ erfüllt das Gesetz die
Kokettenbedingung: die zugehörigen Jacobi-Matrizen multiplizieren sich wie
$J(x\to\tilde x)=J(\bar x\to\tilde x)\,J(x\to\bar x)$, und dementsprechend
transformieren die Komponenten via
\[
\tilde A
=
\big(J(\bar x\to\tilde x)^{\otimes r}\otimes (J(\bar x\to\tilde x)^{-1})^{\otimes s}\big)\,
\bar A
=
\big(J(x\to\tilde x)^{\otimes r}\otimes (J(x\to\tilde x)^{-1})^{\otimes s}\big)\,
A.
\]
Damit ist das Transformationsgesetz konsistent und unabhängig vom gewählten Zwischenkartensystem.
\end{DEF}




In some presentations, tensor fields are defined in terms of such transformations laws (see $[\mathbf{L}-\mathbf{R}]$ for this approach). It should be emphasized again that there are two slightly different ways of reading local expressions like the above. We may think of all of these functions as living on the manifold in the domain $U \cap V$. In this interpretation, we read the above as

$$
\bar{A}_{j k}^{i}(p)=A_{a b}^{l}(p) \frac{\partial x^{a}}{\partial \bar{x}^{j}}(p) \frac{\partial x^{b}}{\partial \bar{x}^{k}}(p) \frac{\partial \bar{x}^{i}}{\partial x^{l}}(p) \text { for each } p \in U \cap V .
$$

This is the default modern viewpoint. Alternatively, we could take $\frac{\partial x^{j}}{\partial \bar{x}^{m}}$ to be functions on $\overline{\mathrm{x}}(U \cap V)$ and write $\frac{\partial x^{j}}{\partial \bar{x}^{m}}\left(\bar{x}^{1}, \ldots, \bar{x}^{n}\right)$. Then, $\frac{\partial \bar{x}^{i}}{\partial x^{l}}$ would refer to $\frac{\partial \bar{x}^{i}}{\partial x^{l}} \circ \mathrm{x} \circ \overline{\mathrm{x}}^{-1}\left(\bar{x}^{1}, \ldots, \bar{x}^{n}\right)$ so that both sides of the equation are functions of variables which we abusively write as $\left(\bar{x}^{1}, \ldots, \bar{x}^{n}\right)$. The first version seems theoretically pleasing, but for specific calculations that use familiar coordinates such as polar coordinates, the second version is often convenient. For example, suppose that a tensor $\tau$ has components with respect to rectangular coordinates on $\mathbb{R}^{2}$ given by $A_{j k}$ and we wish to find the components in polar coordinates. For indexing purposes, we take $(x, y)=$ ( $u^{1}, u^{2}$ ) and ( $r, \theta$ ) = ( $v^{1}, v^{2}$ ). Then we have

$$
\bar{A}_{j k}=A_{a b} \frac{\partial u^{a}}{\partial v^{j}} \frac{\partial u^{b}}{\partial v^{k}},
$$

which can be read so that both sides are functions of $\left(v^{1}, v^{2}\right)$ by writing $u^{1}$ and $u^{2}$ as a function of $\left(v^{1}, v^{2}\right)$, etc. Of course, the charts are there to "identify" open sets in Euclidean space with open sets on the manifold so these viewpoints are really somehow the same after all. Using ( $x, y$ ) and $(r, \theta)$, the transformation (7.5) is given in matrix form as

$$
\left[\begin{array}{ll}
\bar{A}_{11} & \bar{A}_{12} \\
\bar{A}_{12} & \bar{A}_{22}
\end{array}\right]=\left[\begin{array}{cc}
\cos \theta & -r \sin \theta \\
\sin \theta & r \cos \theta
\end{array}\right]\left[\begin{array}{ll}
A_{11} & A_{12} \\
A_{12} & A_{22}
\end{array}\right]\left[\begin{array}{cc}
\cos \theta & \sin \theta \\
-r \sin \theta & r \cos \theta
\end{array}\right] .
$$






\begin{CONC}{GA-L09-04-14}{Zwei Sichtweisen auf das Transformationsgesetz}
In manchen Darstellungen werden Tensorfelder \emph{direkt} durch ihre Transformationsgesetze
definiert (siehe etwa $[\mathbf{L}\!-\!\mathbf{R}]$). 
Dabei ist zu beachten, dass lokale Gleichungen wie
\[
\bar A^{i}{}_{jk}
=
A^{l}{}_{ab}
\frac{\partial \bar x^{i}}{\partial x^{l}}
\frac{\partial x^{a}}{\partial \bar x^{j}}
\frac{\partial x^{b}}{\partial \bar x^{k}}
\]
zwei unterschiedliche, aber äquivalente Weisen der Interpretation zulassen.

\textbf{(1) Punktweise Interpretation auf der Mannigfaltigkeit.}
In der modernen Sichtweise liest man alle beteiligten Größen als Funktionen
auf der Mannigfaltigkeit selbst. Das Transformationsgesetz wird dann
für jedes $p\in U\cap V$ punktweise als
\[
\boxed{\;
\bar A^{i}{}_{jk}(p)
=
A^{l}{}_{ab}(p)
\frac{\partial x^{a}}{\partial \bar x^{j}}(p)
\frac{\partial x^{b}}{\partial \bar x^{k}}(p)
\frac{\partial \bar x^{i}}{\partial x^{l}}(p)
\;}
\]
verstanden.  
Dies ist die konzeptionell saubere Formulierung, bei der alle Terme
als glatte Funktionen $M\to\mathbb{R}$ betrachtet werden.

\textbf{(2) Koordinatenbasierte Interpretation auf $\bar x(U\cap V)\subset\mathbb{R}^{n}$.}
Alternativ kann man die partiellen Ableitungen
\[
\frac{\partial x^{j}}{\partial \bar x^{m}}
\]
als Funktionen auf dem Koordinatenraum $\bar x(U\cap V)$ auffassen, also als
\[
\frac{\partial x^{j}}{\partial \bar x^{m}}
\big(\bar x^{1},\ldots,\bar x^{n}\big).
\]
Dann gilt dieselbe Formel, wobei man nun
\[
\frac{\partial \bar x^{i}}{\partial x^{l}}
\text{ als }
\frac{\partial \bar x^{i}}{\partial x^{l}}
\circ x \circ \bar x^{-1}
\big(\bar x^{1},\ldots,\bar x^{n}\big)
\]
auffasst.  
Beide Seiten der Gleichung sind so Funktionen der Koordinaten
$\big(\bar x^{1},\ldots,\bar x^{n}\big)$.
Diese Schreibweise ist besonders praktisch für konkrete Rechnungen
in Standardkoordinaten.

\textbf{(3) Beispiel: Transformation in Polarkoordinaten.}
Betrachten wir $\mathbb{R}^{2}$ mit kartesischen Koordinaten $(x,y)$ 
und Polarkoordinaten $(r,\theta)$.
Seien $(x,y)=(u^{1},u^{2})$ und $(r,\theta)=(v^{1},v^{2})$.
Für ein $(0,2)$-Tensorfeld mit Komponenten $A_{jk}$ gilt dann
\[
\boxed{\;
\bar A_{jk}
=
A_{ab}\,
\frac{\partial u^{a}}{\partial v^{j}}\,
\frac{\partial u^{b}}{\partial v^{k}}.
\;}
\]
Dabei sind beide Seiten als Funktionen von $(v^{1},v^{2})=(r,\theta)$ zu lesen.

\textbf{(4) Matrixform der Transformation.}
In diesen Koordinaten lautet die Umrechnung explizit:
\[
\begin{bmatrix}
\bar A_{11} & \bar A_{12}\\[2pt]
\bar A_{12} & \bar A_{22}
\end{bmatrix}
=
\begin{bmatrix}
\cos\theta & -r\sin\theta\\[2pt]
\sin\theta & r\cos\theta
\end{bmatrix}
\begin{bmatrix}
A_{11} & A_{12}\\[2pt]
A_{12} & A_{22}
\end{bmatrix}
\begin{bmatrix}
\cos\theta & \sin\theta\\[2pt]
-r\sin\theta & r\cos\theta
\end{bmatrix}.
\]

\textbf{(5) Interpretation.}
Diese beiden Sichtweisen sind mathematisch äquivalent:
Koordinatenkarten identifizieren offene Teilmengen der Mannigfaltigkeit
mit offenen Teilmengen des $\mathbb{R}^{n}$, sodass es sich letztlich nur um
zwei äquivalente Darstellungsformen desselben Zusammenhangs handelt —
die erste ist konzeptionell, die zweite praktisch für Rechnungen.
\end{CONC}





We now introduce the pull-back of a covariant tensor field, which will play a big role in the next chapter.

Definition 7.29. 


If $f: M \rightarrow N$ is a smooth map and $\tau \in \mathcal{T}_{s}(N)$, then we define the pull-back $f^{*} \tau \in \mathcal{T}_{s}(M)$ by
$$
f^{*} \tau\left(v_{1}, \ldots, v_{s}\right)(p)=\tau\left(T f \cdot v_{1}, \ldots, T f \cdot v_{s}\right)
$$
for all $v_{1}, \ldots, v_{s} \in T_{p} M$ and any $p \in M$.


Notice the connection of this with the pull-back defined earlier in a purely algebraic context. It is not hard to see that $f^{*}: \mathcal{T}_{s}(N) \rightarrow \mathcal{T}_{s}(M)$ is linear over $\mathbb{R}$, and for any $h \in C^{\infty}(N)$ and $\tau \in \mathcal{T}_{s}(N)$ we have $f^{*}(h \tau)= (h \circ f) f^{*} \tau$. If $f: M \rightarrow N$ and $g: N \rightarrow P$ are smooth maps, then of course $(g \circ f)^{*}=f^{*} \circ g^{*}$.

Let us discover the local expression for pull-back. 

Choose a chart ( $U, \mathrm{x}$ ) on $M$ and a chart ( $V, \mathrm{y}$ ) on $N$ and assume that $f(U) \subset V$. Let us denote $\frac{\partial\left(y^{i} \circ f\right)}{\partial x^{j}}$ by $\frac{\partial y^{i}}{\partial x^{j}}$ for simplicity. We have $\left.T_{p} f \cdot \frac{\partial}{\partial x^{i}}\right|_{p}=\left.\sum \frac{\partial y^{k}}{\partial x^{i}}(p) \frac{\partial}{\partial y^{k}}\right|_{f(p)}$ and
$$
\begin{aligned}
\left(f^{*} \tau\right)_{i_{1} \ldots i_{s}}(p) & =\left(f^{*} \tau\right)\left(\left.\frac{\partial}{\partial x^{i_{1}}}\right|_{p}, \ldots,\left.\frac{\partial}{\partial x^{i_{s}}}\right|_{p}\right) \\
& =\tau\left(\left.T f \frac{\partial}{\partial x^{i_{1}}}\right|_{p}, \ldots,\left.T f \frac{\partial}{\partial x^{i_{s}}}\right|_{p}\right) \\
& =\tau\left(\left.\frac{\partial y^{k_{1}}}{\partial x^{i_{1}}}(p) \frac{\partial}{\partial y^{k_{1}}}\right|_{f(p)}, \ldots,\left.\frac{\partial y^{k_{s}}}{\partial x^{i_{s}}}(p) \frac{\partial}{\partial y^{k_{s}}}\right|_{f(p)}\right) \\
& =\tau\left(\left.\frac{\partial}{\partial y^{k_{1}}}\right|_{f(p)}, \ldots,\left.\frac{\partial}{\partial y^{k_{s}}}\right|_{f(p)}\right) \frac{\partial y^{k_{1}}}{\partial x^{i_{1}}}(p) \cdots \frac{\partial y^{k_{s}}}{\partial x^{i_{s}}}(p) \\
& =\tau_{k_{1} \ldots k_{s}}(f(p)) \frac{\partial y^{k_{1}}}{\partial x^{i_{1}}}(p) \cdots \frac{\partial y^{k_{s}}}{\partial x^{i_{s}}}(p)
\end{aligned}
$$
Thus we have
$$
\left(f^{*} \tau\right)_{i_{1} \ldots i_{s}}=\left(\tau_{k_{1} \ldots k_{s}} \circ f\right) \frac{\partial y^{k_{1}}}{\partial x^{i_{1}}} \cdots \frac{\partial y^{k_{s}}}{\partial x^{i_{s}}}
$$

This looks similar to a transformation law for a tensor, but here $f$ is not a change of coordinates and need not even be a diffeomorphism. 


In the case that $f: M \rightarrow N$ is a diffeomorphism, the notion of pull-back can be extended to contravariant tensors and tensors of mixed covariance. For such a diffeomorphism, let $\left(T f^{-1}\right)^{*}: T_{p}^{*} M \rightarrow T_{p}^{*} N$ denote the dual of the map $T f^{-1}: T_{p} N \rightarrow T_{p} M$.

Definition 7.31. If $f: M \rightarrow N$ is a diffeomorphism and $\tau$ is an ( $r, s$ )-tensor field on $N$, then define the pull-back $f^{*} \tau \in T^{r}{ }_{s}(M)$ by

$$
\begin{aligned}
& f^{*} \tau\left(a_{1}, \ldots, a_{r}, v_{1}, \ldots, v_{s}\right)(p) \\
& \quad:=\tau\left(\left(T f^{-1}\right)^{*} a_{1}, \ldots,\left(T f^{-1}\right)^{*} a_{r}, T f \cdot v_{1}, \ldots, T f \cdot v_{s}\right)
\end{aligned}
$$

for all $v_{1}, \ldots, v_{s} \in T_{p} M$ and $a_{1}, \ldots, a_{r} \in T_{p}^{*} M$ and any $p \in M$. The pushforward is then defined for $\tau \in T^{r}{ }_{s}(M)$ as $f_{*} \tau:=\left(f^{-1}\right)^{*} \tau$.





\begin{DEF}{GA-L09-04-15}{Pull-Back von Tensorfeldern}
Sei $f : M \to N$ eine glatte Abbildung zwischen Mannigfaltigkeiten.

\textbf{(1) Pull-Back kovarianter Tensorfelder.}  
Für $\tau \in \mathcal{T}_{s}(N) = \Gamma(T^{0}_{s}(TN))$ wird der \emph{Pull-Back} 
$f^{*}\tau \in \mathcal{T}_{s}(M)$ durch
\[
\boxed{\;
(f^{*}\tau)(v_{1},\ldots,v_{s})(p)
=
\tau\big(T_{p}f\cdot v_{1},\ldots, T_{p}f\cdot v_{s}\big),
\qquad
v_{i}\in T_{p}M,\;p\in M.
\;}
\]
definiert.  
Dies entspricht der natürlichen algebraischen Definition des Pull-Backs für multilineare Abbildungen.

\textbf{Eigenschaften.}
\begin{itemize}
\item $f^{*}$ ist $\mathbb{R}$-linear:
  \[
  f^{*}(a\tau + b\sigma) = a f^{*}\tau + b f^{*}\sigma.
  \]
\item Für $h\in C^{\infty}(N)$ und $\tau\in \mathcal{T}_{s}(N)$ gilt:
  \[
  f^{*}(h\tau) = (h\circ f)\,f^{*}\tau.
  \]
\item Für glatte Abbildungen $f:M\to N$ und $g:N\to P$ gilt die Kompositionsregel:
  \[
  (g\circ f)^{*} = f^{*}\circ g^{*}.
  \]
\end{itemize}

\textbf{(2) Lokale Darstellung.}  
Seien $(U,x)$ eine Karte auf $M$ und $(V,y)$ eine Karte auf $N$ mit $f(U)\subset V$.
Dann gilt in Komponentenform:
\[
(f^{*}\tau)_{i_{1}\ldots i_{s}}(p)
=
\tau_{k_{1}\ldots k_{s}}(f(p))
\frac{\partial y^{k_{1}}}{\partial x^{i_{1}}}(p)
\cdots
\frac{\partial y^{k_{s}}}{\partial x^{i_{s}}}(p),
\]
oder kompakter geschrieben:
\[
\boxed{\;
(f^{*}\tau)_{i_{1}\ldots i_{s}}
=
(\tau_{k_{1}\ldots k_{s}}\circ f)
\frac{\partial y^{k_{1}}}{\partial x^{i_{1}}}
\cdots
\frac{\partial y^{k_{s}}}{\partial x^{i_{s}}}.
\;}
\]
Choose a chart ( $U, \mathrm{x}$ ) on $M$ and a chart ( $V, \mathrm{y}$ ) on $N$ and assume that $f(U) \subset V$. Let us denote $\frac{\partial\left(y^{i} \circ f\right)}{\partial x^{j}}$ by $\frac{\partial y^{i}}{\partial x^{j}}$ for simplicity. We have $\left.T_{p} f \cdot \frac{\partial}{\partial x^{i}}\right|_{p}=\left.\sum \frac{\partial y^{k}}{\partial x^{i}}(p) \frac{\partial}{\partial y^{k}}\right|_{f(p)}$ and
$$
\begin{aligned}
\left(f^{*} \tau\right)_{i_{1} \ldots i_{s}}(p) & =\left(f^{*} \tau\right)\left(\left.\frac{\partial}{\partial x^{i_{1}}}\right|_{p}, \ldots,\left.\frac{\partial}{\partial x^{i_{s}}}\right|_{p}\right) \\
& =\tau\left(\left.T f \frac{\partial}{\partial x^{i_{1}}}\right|_{p}, \ldots,\left.T f \frac{\partial}{\partial x^{i_{s}}}\right|_{p}\right) \\
& =\tau\left(\left.\frac{\partial y^{k_{1}}}{\partial x^{i_{1}}}(p) \frac{\partial}{\partial y^{k_{1}}}\right|_{f(p)}, \ldots,\left.\frac{\partial y^{k_{s}}}{\partial x^{i_{s}}}(p) \frac{\partial}{\partial y^{k_{s}}}\right|_{f(p)}\right) \\
& =\tau\left(\left.\frac{\partial}{\partial y^{k_{1}}}\right|_{f(p)}, \ldots,\left.\frac{\partial}{\partial y^{k_{s}}}\right|_{f(p)}\right) \frac{\partial y^{k_{1}}}{\partial x^{i_{1}}}(p) \cdots \frac{\partial y^{k_{s}}}{\partial x^{i_{s}}}(p) \\
& =\tau_{k_{1} \ldots k_{s}}(f(p)) \frac{\partial y^{k_{1}}}{\partial x^{i_{1}}}(p) \cdots \frac{\partial y^{k_{s}}}{\partial x^{i_{s}}}(p)
\end{aligned}
$$
Thus we have
$$
\left(f^{*} \tau\right)_{i_{1} \ldots i_{s}}=\left(\tau_{k_{1} \ldots k_{s}} \circ f\right) \frac{\partial y^{k_{1}}}{\partial x^{i_{1}}} \cdots \frac{\partial y^{k_{s}}}{\partial x^{i_{s}}}
$$

Diese Formel ähnelt formal dem Transformationsgesetz für Tensoren, beschreibt aber eine andere Situation:  
$f$ muss keine Diffeomorphismus sein und kann Abbildungen zwischen verschieden-dimensionalen Mannigfaltigkeiten darstellen.

\textbf{(3) Erweiterung auf gemischte Tensoren (Diffeomorphismen).}  
Sei $f:M\to N$ ein Diffeomorphismus und $\tau\in \mathcal{T}^{r}_{s}(N)$ ein $(r,s)$-Tensorfeld.
Dann definieren wir den Pull-Back $f^{*}\tau\in \mathcal{T}^{r}_{s}(M)$ durch
\[
\boxed{
\begin{aligned}
(f^{*}\tau)
(a_{1},\ldots,a_{r},v_{1},\ldots,v_{s})(p)
&:=
\tau\big(
(Tf^{-1})^{*}a_{1},\ldots,(Tf^{-1})^{*}a_{r},\\[-1ex]
&\hspace{4.8em}
Tf\cdot v_{1},\ldots,Tf\cdot v_{s}
\big),
\end{aligned}
}
\]
für $a_{i}\in T_{p}^{*}M$, $v_{i}\in T_{p}M$ und $p\in M$.
Hier bezeichnet $(Tf^{-1})^{*}$ das Dual (Pull-Back) des Tangentenabbildung-Inversen:
\[
(Tf^{-1})^{*}: T_{p}^{*}M \to T_{f(p)}^{*}N.
\]

\textbf{(4) Pushforward.}  
Für ein $(r,s)$-Tensorfeld $\tau$ auf $M$ definieren wir den \emph{Pushforward} durch
\[
f_{*}\tau := (f^{-1})^{*}\tau.
\]

\textbf{(5) Bemerkung.}
Der Pull-Back eines kovarianten Tensorfeldes hängt nur von der Ableitung $Tf$ ab.
Für kontravariante Indizes ist die Invertierbarkeit von $Tf$ erforderlich — daher ist 
der allgemeine Pull-Back für gemischte Tensoren nur für Diffeomorphismen wohldefiniert.
\end{DEF}






Pull-back respects tensor products:

Exercise 7.30. Let $f: M \rightarrow N$ be as above. Show that for $\tau_{1} \in \mathcal{T}_{s_{1}}(N)$ and $\tau_{2} \in \mathcal{T}_{s_{2}}(N)$ we have $f^{*}\left(\tau_{1} \otimes \tau_{2}\right)=f^{*} \tau_{1} \otimes f^{*} \tau_{2}$.




\begin{PROP}{GA-L09-04-16}{Verträglichkeit des Pull-Backs mit dem Tensorprodukt}
Sei $f:M\to N$ glatt, $\tau_{1}\in\mathcal{T}_{s_{1}}(N)$ und $\tau_{2}\in\mathcal{T}_{s_{2}}(N)$.
Dann gilt
\[
f^{*}(\tau_{1}\otimes \tau_{2})
=
f^{*}\tau_{1}\ \otimes\ f^{*}\tau_{2}
\ \in\ \mathcal{T}_{s_{1}+s_{2}}(M).
\]
\end{PROP}



\begin{PROOF}{GA-L09-04-17}{P: Verträglichkeit des Pull-Backs mit dem Tensorprodukt}
Sei $p\in M$ und $v_{1},\dots,v_{s_{1}+s_{2}}\in T_{p}M$. Mit der Definition des Pull-Backs für kovariante Tensoren erhält man
\[
\begin{aligned}
\big(f^{*}(\tau_{1}\otimes \tau_{2})\big)(v_{1},\dots,v_{s_{1}+s_{2}})(p)
&= (\tau_{1}\otimes \tau_{2})\big(Tf\cdot v_{1},\dots,Tf\cdot v_{s_{1}+s_{2}}\big) \\[2pt]
&= \tau_{1}\big(Tf\cdot v_{1},\dots,Tf\cdot v_{s_{1}}\big)\;
   \tau_{2}\big(Tf\cdot v_{s_{1}+1},\dots,Tf\cdot v_{s_{1}+s_{2}}\big) \\[2pt]
&= \big(f^{*}\tau_{1}\big)(v_{1},\dots,v_{s_{1}})(p)\ \cdot\
   \big(f^{*}\tau_{2}\big)(v_{s_{1}+1},\dots,v_{s_{1}+s_{2}})(p) \\[2pt]
&= \big((f^{*}\tau_{1}\otimes f^{*}\tau_{2})\big)(v_{1},\dots,v_{s_{1}+s_{2}})(p).
\end{aligned}
\]
Da die Gleichheit punktweise und auf beliebigen Tangentialargumenten gilt, folgt
$f^{*}(\tau_{1}\otimes \tau_{2})=f^{*}\tau_{1}\otimes f^{*}\tau_{2}$ als Tensorfelder.
\end{PROOF}



\pagebreak

\subsection*{Top-Down Approach to Tensor Fields}
Specializing what we learned from the discussion following Proposition 7.22 to the case of the tangent bundle, we see that a tensor field gives us a $C^{\infty}(M)$-multilinear map based on the module $\mathfrak{X}(M)$. This observation leads\\
to an alternative definition of a tensor field over $M$. In this "top-down" view, we simply define an ( $r, s$ )-tensor field to be a $C^{\infty}(M)$-multilinear map

$$
\mathfrak{X}^{*}(M)^{r} \times \mathfrak{X}(M)^{s} \rightarrow C^{\infty}(M) .
$$

In this view, a tensor field is an element of $T^{r}{ }_{s}(\mathfrak{X}(M))$. For example, a global covariant 2-tensor field on a manifold $M$ is a map $\tau: \mathfrak{X}(M) \times \mathfrak{X}(M) \rightarrow C^{\infty}(M)$ such that

$$
\begin{aligned}
\tau\left(f_{1} X_{1}+f_{2} X_{2}, Y\right) & =f_{1} \tau\left(X_{1}, Y\right)+f_{2} \tau\left(X_{2}, Y\right), \\
\tau\left(Y, f_{1} X_{1}+f_{2} X_{2}\right) & =f_{1} \tau\left(Y, X_{1}\right)+f_{2} \tau\left(Y, X_{2}\right)
\end{aligned}
$$

for all $f_{1}, f_{2} \in C^{\infty}(M)$ and all $X_{1}, X_{2}, Y \in \mathfrak{X}(M)$. As we shall see, it turns out that such $C^{\infty}(M)$-multilinear maps determine tensor fields in the sense of the previous section.




\begin{DEF}{GA-L09-04-18}{Top-down-Definition von Tensorfeldern über $M$}
Sei $M$ eine glatte Mannigfaltigkeit.

\textbf{(1) Modulischer Zugang.}
Bezeichne mit $\mathfrak{X}(M)$ den $C^{\infty}(M)$-Modul der glatten Vektorfelder
und mit $\mathfrak{X}^{*}(M)$ dessen Dualmodul
$\mathfrak{X}^{*}(M):=\mathrm{Hom}_{C^{\infty}(M)}\!\big(\mathfrak{X}(M),C^{\infty}(M)\big)$.
Ein \emph{$(r,s)$-Tensorfeld (top-down)} ist eine $C^{\infty}(M)$-multilineare Abbildung
\[
A:\ \underbrace{\mathfrak{X}^{*}(M)\times\cdots\times\mathfrak{X}^{*}(M)}_{r\ \text{Faktoren}}
\times
\underbrace{\mathfrak{X}(M)\times\cdots\times\mathfrak{X}(M)}_{s\ \text{Faktoren}}
\longrightarrow C^{\infty}(M),
\]
die in jedem Argument $C^{\infty}(M)$-linear ist.
Den Raum aller solcher Abbildungen schreiben wir als
\[
T^{r}{}_{s}\big(\mathfrak{X}(M)\big)
:=\mathrm{Multilin}_{C^{\infty}(M)}
\!\Big(\big(\mathfrak{X}^{*}(M)\big)^{\times r}\times\big(\mathfrak{X}(M)\big)^{\times s},\,C^{\infty}(M)\Big).
\]

\textbf{(2) Beispiel $(0,2)$-Tensorfeld.}
Eine Abbildung $\tau:\mathfrak{X}(M)\times\mathfrak{X}(M)\to C^{\infty}(M)$ ist ein
kovariantes $2$-Tensorfeld genau dann, wenn
\[
\begin{aligned}
\tau\!\big(f_{1}X_{1}+f_{2}X_{2},\,Y\big)&=f_{1}\,\tau(X_{1},Y)+f_{2}\,\tau(X_{2},Y),\\
\tau\!\big(Y,\,f_{1}X_{1}+f_{2}X_{2}\big)&=f_{1}\,\tau(Y,X_{1})+f_{2}\,\tau(Y,X_{2})
\end{aligned}
\]
für alle $f_{1},f_{2}\in C^{\infty}(M)$ und $X_{1},X_{2},Y\in\mathfrak{X}(M)$ gilt.
\end{DEF}



\begin{PROP}{GA-L09-04-19}{Äquivalenz der beiden Definitionen von Tensorfeldern}
Die top-down-Definition ist äquivalent zur geometrischen
Definition als glatte Schnitte des Tensorbündels:
\[
T^{r}{}_{s}\big(\mathfrak{X}(M)\big)\ \cong\ \Gamma\!\big(T^{r}{}_{s}(TM)\big)=:\mathcal{T}^{r}{}_{s}(M).
\]
Der Isomorphismus ist $C^{\infty}(M)$-linear, natürlich in $M$, und wird lokal durch
die Identifikation mit der Standardbasis $\{\partial/\partial x^{i},\,dx^{j}\}$
gegeben (Komponentenfunktionen stimmen überein).
\end{PROP}




If we take a top-down approach to tensor fields, then we must work to recover the presheaf/sheaf aspects. Indeed, it is not obvious what is the relation between $T^{r}{ }_{s}(\mathfrak{X}(M))$ and $T^{r}{ }_{s}(\mathfrak{X}(U))$ for some proper open subset $U \subset M$. Indeed, thinking purely in terms of modules makes the issue clear. The module $\mathfrak{X}(M)$ is not the same module as $\mathfrak{X}(U)$ unless $U=M$. A priori, there is no immediate reason to think that a multilinear map with arguments from the module $\mathfrak{X}(M)$ should be able to take elements of $\mathfrak{X}(U)$ as arguments! For instance, from the top-down viewpoint, how can we insert coordinate fields $\frac{\partial}{\partial x^{i}}$ and $d x^{i}$ into an element of $T_{s}^{r}(\mathfrak{X}(M))$ to get coordinate expressions if the chart domain is not all of $M$ ? We address this in the next section indirectly by showing how the top-down approach gives back tensors as sections (the bottom-up approach). Another comment is that both $\mathfrak{X}(U)$ and $T^{r}{ }_{s}(\mathfrak{X}(U))$ are finite-dimensional free modules over the ring $C^{\infty}(U)$ whenever $U$ is a chart domain or, more generally, the domain of a frame field. The reason is that a local frame field and its dual frame field provide a module basis for $\mathfrak{X}(U)$ and $\mathfrak{X}^{*}(U)$ and the latter really is the dual of the first in the module sense. On the other hand, the $C^{\infty}(M)$-modules $\mathfrak{X}(M)$ and $T_{s}^{r}(\mathfrak{X}(M))$ are not generally free unless $M$ is parallelizable.



\begin{CONC}{GA-L09-04-20}{Lokale und globale Aspekte des Top-Down-Ansatzes}
Der sogenannte \emph{Top-Down-Ansatz} zur Definition von Tensorfeldern beschreibt Tensoren
ausschließlich in algebraischer Form als $C^{\infty}(M)$-multilineare Abbildungen
\[
T^{r}{}_{s}(\mathfrak{X}(M))
=\mathrm{Multilin}_{C^{\infty}(M)}
\!\Big((\mathfrak{X}^{*}(M))^{r}\times (\mathfrak{X}(M))^{s},\,C^{\infty}(M)\Big).
\]
Diese Sichtweise hat jedoch einen wichtigen Nachteil:
die Beziehung zwischen globalen und lokalen Tensorfeldern ist nicht unmittelbar ersichtlich.
Im Gegensatz zur geometrischen („Bottom-Up“-)Definition ist hier nicht offensichtlich,
wie sich lokale Einschränkungen oder der Zusammenhang zu Sheaf-Strukturen ergeben.

\textbf{(1) Problem der Einschränkung.}
Für eine echte offene Teilmenge $U\subsetneq M$ gilt
\[
\mathfrak{X}(U)\neq \mathfrak{X}(M)|_{U},
\]
d.\,h. der Modul der Vektorfelder auf $U$ ist ein anderer als der der globalen Vektorfelder.
Daher ist \emph{a priori} unklar, wie ein Element
$A\in T^{r}{}_{s}(\mathfrak{X}(M))$ mit Argumenten aus $\mathfrak{X}(U)$ ausgewertet werden soll.
Auch die Bildung lokaler Koordinatenausdrücke wie
$\frac{\partial}{\partial x^{i}}$ und $dx^{i}$ ist nicht unmittelbar möglich,
wenn die Kartendomäne $U$ nicht ganz $M$ umfasst.

\textbf{(2) Lokale Freiheits- und Modulstruktur.}
Trotzdem gilt: Für ein offenes $U\subset M$, das eine lokale Basis (Frame-Feld)
\[
\{s_{1},\dots,s_{k}\}\subset \mathfrak{X}(U)
\quad\text{und die duale Basis}\quad
\{\sigma^{1},\dots,\sigma^{k}\}\subset \mathfrak{X}^{*}(U)
\]
besitzt, ist $\mathfrak{X}(U)$ ein endlichdimensionaler freier
$C^{\infty}(U)$-Modul.  
Ebenso sind alle
\[
T^{r}{}_{s}(\mathfrak{X}(U))
\]
freie $C^{\infty}(U)$-Module, und $\mathfrak{X}^{*}(U)$ ist tatsächlich der Dualmodul
von $\mathfrak{X}(U)$ im algebraischen Sinn.

\textbf{(3) Globaler Gegensatz.}
Im Allgemeinen ist $\mathfrak{X}(M)$ kein freier $C^{\infty}(M)$-Modul,
es sei denn, $M$ ist \emph{parallelisierbar},
d.\,h. das Tangentialbündel $TM$ ist trivial:
\[
TM\cong M\times\mathbb{R}^{n}.
\]
Nur in diesem Fall bilden globale Frame-Felder eine Basis des Moduls $\mathfrak{X}(M)$,
und die $C^{\infty}(M)$-Multilinearität spiegelt exakt die Bündelstruktur wider.

\textbf{(4) Fazit.}
Der Top-Down-Ansatz ist algebraisch elegant und lokal vollständig,
verliert aber zunächst die sheaf-artige Lokalisierbarkeit der Tensorfelder.
In der folgenden Sektion wird gezeigt,
dass sich dieser Zugang durch geeignete Identifikationen
mit der geometrischen („Bottom-Up“) Definition als glatte Schnitte des Tensorbündels
vereinigen lässt, wodurch die Sheaf-Struktur wiedergewonnen wird.
\end{CONC}



\pagebreak

\subsection*{Matching the Two Approaches to Tensor Fields}


If we define a tensor field as we first did, that is, as a field of tensors in tangent spaces, then we immediately obtain a tensor as defined in the topdown approach. On the other hand, if $\tau$ is initially defined as a $C^{\infty}(M)$ multilinear map, then how should we recover a field of tensors on the tangent spaces? ${ }^{2}$ Answering this is our next goal.

\footnotetext{${ }^{2}$ This is exactly where things might not go so well if the manifold is not finite-dimensional. What we need is the existence of smooth cut-off functions. Some Banach manifolds support cut-off functions but not all do.
}

Proposition 7.32. 

\begin{PROP}{GA-L09-04-21}{Punktweise Abhängigkeit von Tensorfeldern}
Let $p \in M$ and $\tau \in T^{r}{ }_{s}(\mathfrak{X}(M))$. Let $\theta_{1}, \ldots, \theta_{r}$ and $\bar{\theta}_{1}, \ldots, \bar{\theta}_{r}$ be smooth 1-forms such that $\theta_{i}(p)=\bar{\theta}_{i}(p)$ for $1 \leq i \leq r$; also let $X_{1}, \ldots, X_{s}$ and $\bar{X}_{1}, \ldots, \bar{X}_{s}$ be smooth vector fields such that $X_{i}(p)=\bar{X}_{i}(p)$ for $1 \leq i \leq s$. Then we have that
$$
\tau\left(\theta_{1}, \ldots, \theta_{r}, X_{1}, \ldots, X_{s}\right)(p)=\tau\left(\bar{\theta}_{1}, \ldots, \bar{\theta}_{r}, \bar{X}_{1}, \ldots, \bar{X}_{s}\right)(p)
$$
\end{PROP}

Proof. 

\begin{PROOF}{GA-L09-04-22}{P: Punktweise Abhängigkeit von Tensorfeldern}
The result will follow easily if we can show that
$$
\tau\left(\theta_{1}, \ldots, \theta_{r}, X_{1}, \ldots, X_{s}\right)(p)=0
$$
whenever one of $\theta_{1}(p), \ldots, \theta_{r}(p), X_{1}(p), \ldots, X_{s}(p)$ is zero. We shall assume for simplicity of notation that $r=1$ and $s=2$. Now suppose that $X_{1}(p)=0$. If ( $U, \mathrm{x}$ ), with $\mathrm{x}=\left(x^{1}, \ldots, x^{n}\right)$, is a chart with $p \in U$, then $\left.X_{1}\right|_{U}=\sum \xi^{i} \frac{\partial}{\partial x^{i}}$ for some smooth functions $\xi^{i} \in C^{\infty}(U)$. Let $\beta$ be a cut-off function with support in $U$ and $\beta(p)=1$. Then for any smooth vector field $X$ defined on $U$ we can consider both $\beta X$ and $\beta^{2} X$ to be globally defined and zero outside of $U$. Similarly, if $f$ is a smooth function defined on $U$, then $\beta f$ can be taken to be globally defined on $M$ and zero outside of $U$. Now $\beta^{2} X_{1}=\sum\left(\beta \xi^{i}\right)\left(\beta \frac{\partial}{\partial x^{i}}\right)$. (Notice that in this last expression we have used $\beta$ to extend both the functions $\xi^{i}$ and the coordinate fields $\frac{\partial}{\partial x^{i}}$, which is why we used $\beta^{2}$ rather than just $\beta$.) Thus
$$
\begin{aligned}
\beta^{2} \tau\left(\theta_{1}, X_{1}, X_{2}\right) & =\tau\left(\theta_{1}, \beta^{2} X_{1}, X_{2}\right) \\
& =\tau\left(\theta_{1}, \beta^{2} \xi^{i} \frac{\partial}{\partial x^{i}}, X_{2}\right) \\
& =\tau\left(\theta_{1},\left(\beta \xi^{i}\right) \beta \frac{\partial}{\partial x^{i}}, X_{2}\right) \\
& =\beta \xi^{i} \tau\left(\theta_{1}, \beta \frac{\partial}{\partial x^{i}}, X_{2}\right)
\end{aligned}
$$
(Notice that the point of the above expression is that, at this moment, $\tau$ is defined only for global sections and is linear over global functions. For example, $\frac{\partial}{\partial x^{i}}$ is not a global section while $\beta \frac{\partial}{\partial x^{i}}$ is a global section.) Since $X_{1}(p)=0$, we must have $\xi^{i}(p)=0$ for all $i$. Also recall that $\beta(p)=1$. Plugging $p$ into the formula above we obtain
$$
\tau\left(\theta_{1}, X_{1}, X_{2}\right)(p)=0
$$
A similar argument holds when $X_{2}(p)=0$ or $\theta_{1}(p)=0$.

Assume that $\theta_{1}(p)=\bar{\theta}_{1}(p), X_{1}(p)=\bar{X}_{1}(p)$ and $X_{2}(p)=\bar{X}_{2}(p)$. Then we have
$$
\begin{aligned}
& \tau\left(\bar{\theta}_{1}, \bar{X}_{1}, \bar{X}_{2}\right)-\tau\left(\theta_{1}, X_{1}, X_{2}\right) \\
& \quad=\tau\left(\bar{\theta}_{1}-\theta_{1}, \bar{X}_{1}, \bar{X}_{2}\right)+\tau\left(\theta_{1}, \bar{X}_{1}-X_{1}, \bar{X}_{2}\right)+\tau\left(\theta_{1}, X_{1}, \bar{X}_{2}-X_{2}\right)
\end{aligned}
$$
Since $\bar{\theta}_{1}-\theta_{1}, \bar{X}_{1}-X_{1}$, and $\bar{X}_{2}-X_{2}$ are all zero at $p$, we obtain the result that $\tau\left(\bar{\theta}_{1}, \bar{X}_{1}, \bar{X}_{2}\right)(p)=\tau\left(\theta_{1}, X_{1}, X_{2}\right)(p)$.
\end{PROOF}

Thus we have a natural correspondence between $T^{r}{ }_{s}(\mathfrak{X}(M))$ and $\mathcal{T}_{s}^{r}(M)$ (the latter being smooth sections of the bundle $T^{r}{ }_{s}(T M) \rightarrow M$ ). For example, if $A \in \mathcal{T}_{3}^{1}(M)$, then we obtain an element of $T_{3}^{1}(\mathfrak{X}(M))$, also denoted by $A$, by defining a smooth function $A(\theta, X, Y, Z)$ for given fields ( $\theta, X, Y, Z$ ) by

$$
A(\theta, X, Y, Z)(p):=A(p)(\theta(p), X(p), Y(p), Z(p)) .
$$

Conversely, if $A \in T_{3}^{1}(\mathfrak{X}(M))$, then we can use the above proposition to define an element of $\mathcal{T}_{3}^{1}(M)$, which we denote by the same letter. Given $A \in T_{3}^{1}(\mathfrak{X}(M))$, define $A(p) \in T_{3}^{1}\left(T_{p} M\right)$ for each $p$ as follows: For $X_{p}, Y_{p}, Z_{p} \in T_{p} M$ and $\theta_{p} \in T_{p}^{*} M$ we let

$$
A(p)\left(\theta_{p}, X_{p}, Y_{p}, Z_{p}\right):=A(\theta, X, Y, Z)(p),
$$

where $\theta, X, Y, Z$ are any fields chosen so that $\theta(p)=\theta_{p}, X(p)=X_{p}, Y(p)= Y_{p}$, and $Z(p)=Z_{p}$. By Lemma 6.37 we can always find such extensions, and by Proposition 7.32 above, $A$ is well-defined. That $A$ so defined is smooth follows from Proposition 7.22. The general case should be clear and, all said, we end up with a natural isomorphism of $C^{\infty}(M)$-modules:

$$
T_{s}^{r}(\mathfrak{X}(M)) \cong \mathcal{T}_{s}^{r}(M)
$$





\begin{THEO}{GA-L09-04-23}{Kanonische Korrespondenz Top-Down $\longleftrightarrow$ Bottom-Up}
Sei $M$ eine glatte Mannigfaltigkeit. Bezeichne
\[
T^{r}{}_{s}\!\big(\mathfrak{X}(M)\big)
:= \mathrm{Multilin}_{C^{\infty}(M)}
\!\Big((\mathfrak{X}^{*}(M))^{r}\times(\mathfrak{X}(M))^{s},\,C^{\infty}(M)\Big)
\]
den $C^{\infty}(M)$-Modul der \emph{Top-Down}-Tensorfelder vom Typ $(r,s)$, und
\[
\mathcal{T}^{r}{}_{s}(M):=\Gamma\!\big(T^{r}{}_{s}(TM)\big)
\]
den $C^{\infty}(M)$-Modul der glatten Schnitte des \emph{Tensorbündels} (\emph{Bottom-Up}).
Dann gibt es einen kanonischen, natürlichen Isomorphismus von $C^{\infty}(M)$-Modulen
\[
\Phi:\ T^{r}{}_{s}\!\big(\mathfrak{X}(M)\big)\ \xrightarrow{\ \cong\ }\ \mathcal{T}^{r}{}_{s}(M),
\qquad
\Psi:=\Phi^{-1}.
\]
\textbf{Konstruktion und Eigenschaften}:
\begin{enumerate}
\item[\textbf{(A)}] \textbf{Abbildung \(\Psi:\mathcal{T}^{r}{}_{s}(M)\to T^{r}{}_{s}(\mathfrak{X}(M))\).}
Für $A\in\mathcal{T}^{r}{}_{s}(M)$ definiere
\[
(\Psi A)(\theta_{1},\dots,\theta_{r},X_{1},\dots,X_{s})
\ \in\ C^{\infty}(M)
\]
durch
\[
\big((\Psi A)(\theta_{1},\dots,\theta_{r},X_{1},\dots,X_{s})\big)(p)
:= A(p)\big(\theta_{1}(p),\dots,\theta_{r}(p),X_{1}(p),\dots,X_{s}(p)\big),
\]
wobei $\theta_{i}\in\mathfrak{X}^{*}(M)$ und $X_{j}\in\mathfrak{X}(M)$ beliebig sind.
Dann ist $\Psi A$ $C^{\infty}(M)$-multilinear und glatt (Komposition glatter Schnitte mit Evaluation).

\item[\textbf{(B)}] \textbf{Abbildung \(\Phi:T^{r}{}_{s}(\mathfrak{X}(M))\to \mathcal{T}^{r}{}_{s}(M)\).}
Sei $B\in T^{r}{}_{s}(\mathfrak{X}(M))$. Für jeden Punkt $p\in M$ definiere
\[
(\Phi B)(p)\ \in\ T^{r}{}_{s}(T_{p}M)
\]
durch folgende Punktweise-Formel:
Für $a_{1},\dots,a_{r}\in T_{p}^{*}M$ und $v_{1},\dots,v_{s}\in T_{p}M$ wähle beliebige glatte Fortsetzungen
\[
\theta_{i}\in\mathfrak{X}^{*}(M),\ \theta_{i}(p)=a_{i}, 
\qquad
X_{j}\in\mathfrak{X}(M),\ X_{j}(p)=v_{j},
\]
und setze
\[
(\Phi B)(p)(a_{1},\dots,a_{r},v_{1},\dots,v_{s})
:=\ B(\theta_{1},\dots,\theta_{r},X_{1},\dots,X_{s})(p).
\]
Diese Definition ist \emph{wohldefiniert}, d.\,h. unabhängig von den gewählten Fortsetzungen, 
und liefert glatt von $p$ abhängige Tensoren.

\item[\textbf{(C)}] \textbf{Wohldefiniertheit in (B) (Punktweise-Abhängigkeit).}
Geben $\theta_{i},\bar\theta_{i}$ bzw. $X_{j},\bar X_{j}$ an $p$ denselben Wert, so gilt
\[
B(\theta_{1},\dots,\theta_{r},X_{1},\dots,X_{s})(p)
=
B(\bar\theta_{1},\dots,\bar\theta_{r},\bar X_{1},\dots,\bar X_{s})(p),
\]
denn $B$ ist $C^{\infty}(M)$-multilinear und damit punktweise nur vom Keim der Argumente abhängig.
Somit hängt $(\Phi B)(p)$ nur von $(a_{i})$ und $(v_{j})$ ab.

\item[\textbf{(D)}] \textbf{Glattheit in (B).}
Wähle lokal ein Koordinatenfeld $(\partial/\partial x^{i},dx^{j})$ oder ein lokales Frame-Feld
$(s_{1},\dots,s_{n})$ mitsamt Dualfeld $(\sigma^{1},\dots,\sigma^{n})$ auf $U\subset M$.
Schreibe $B$ auf $U$ durch seine Komponenten
\[
B = B^{i_{1}\dots i_{r}}{}_{j_{1}\dots j_{s}}\ \sigma^{j_{1}}\otimes\cdots\otimes\sigma^{j_{s}}
\otimes s_{i_{1}}\otimes\cdots\otimes s_{i_{r}},
\]
wobei $B^{i_{1}\dots i_{r}}{}_{j_{1}\dots j_{s}}\in C^{\infty}(U)$.
Dann ist $(\Phi B)(p)$ durch diese glatten Komponenten gegeben; 
also ist $p\mapsto(\Phi B)(p)$ ein glatter Schnitt.

\item[\textbf{(E)}] \textbf{$C^{\infty}(M)$-Linearität und Multilinearität.}
Beide Abbildungen $\Phi$ und $\Psi$ sind $C^{\infty}(M)$-linear und respektieren 
Tensorprodukt, Kontraktion und (für Diffeomorphismen) Pull-Back/Push-Forward:
\[
\Phi(B_{1}\otimes B_{2})=\Phi(B_{1})\otimes\Phi(B_{2}),\qquad
\Psi(A_{1}\otimes A_{2})=\Psi(A_{1})\otimes\Psi(A_{2}),
\]
\[
\Phi\big(C^{k}_{\ell} B\big)=C^{k}_{\ell}\big(\Phi B\big),\qquad
\Psi\big(C^{k}_{\ell} A\big)=C^{k}_{\ell}\big(\Psi A\big),
\]
und für einen Diffeomorphismus $f:M\to N$ gilt die Natürlichkeit
\[
\Phi_{M}\circ f^{*}_{\mathrm{top}} = f^{*}_{\mathrm{bun}}\circ \Phi_{N},
\]
wobei $f^{*}_{\mathrm{top}}$ der Pull-Back im Top-Down-Sinn (via $Tf$, $(Tf^{-1})^{*}$) 
und $f^{*}_{\mathrm{bun}}$ der Bündel-Pull-Back ist.

\item[\textbf{(F)}] \textbf{Inversen-Eigenschaft.}
Für $A\in\mathcal{T}^{r}{}_{s}(M)$ und $B\in T^{r}{}_{s}(\mathfrak{X}(M))$ gilt
\[
\Phi(\Psi A)=A,
\qquad
\Psi(\Phi B)=B.
\]
Denn $\Psi A$ evaluiert $A$ auf Feldern, und $\Phi(\Psi A)$ stellt punktweise genau $A(p)$ wieder her;
umgekehrt liefert die punktweise Def. in (B) aus $B$ die identische $C^{\infty}(M)$-multilineare Abbildung nach (A).
\end{enumerate}
\end{THEO}



\begin{PROOF}{GA-L09-04-24}{P: Kanonische Korrespondenz Top-Down $\longleftrightarrow$ Bottom-Up}
\textbf{Zu (A):} Multilinearität folgt direkt aus der Multilinearität von $A(p)$ in jedem Argument und der $C^{\infty}(M)$-Linearität der Evaluation $X\mapsto X(p)$, $\theta\mapsto\theta(p)$. Glattheit folgt aus der glatten Abhängigkeit der Komponenten von $p$ (lokale Koordinatendarstellung).

\textbf{Zu (B)\,\&\,(C):} Existenz glatter Fortsetzungen folgt z.\,B. durch Lemma über lokale Erweiterungen (Partitions der Eins). Unabhängigkeit von der Wahl der Fortsetzungen ist eine Konsequenz der $C^{\infty}(M)$-Multilinearität: unterscheiden sich zwei Fortsetzungen in einer Umgebung von $p$ nur um ein Feld, das am Punkt $p$ verschwindet, so ist der Unterschied durch einen Faktor aus dem verschwindenden Ideal $I_{p}\subset C^{\infty}(M)$ multipliziert und verschwindet bei Auswertung an $p$.

\textbf{Zu (D):} In lokalen Koordinaten/Frames sind die Komponenten von $B$ glatt; daher ist die so definierte punktweise multilineare Abbildung $p\mapsto (\Phi B)(p)$ glatt.

\textbf{Zu (E):} Die Aussagen folgen aus den Definitionen und der Faserweise- bzw. Punktweise-Natur von Tensorprodukt, Kontraktion und (für Diffeomorphismen) Pull-Back/Push-Forward.

\textbf{Zu (F):} Direktes Nachrechnen: $\Psi A$ evaluiert $A$ auf Feldern und $\Phi$ stellt dies punktweise als Faser-Tensor wieder her (und umgekehrt). Damit sind $\Phi$ und $\Psi$ invers.

Alle Schritte sind kanonisch, somit ist der Isomorphismus natürlich in $M$.
\end{PROOF}






\begin{CONC}{GA-L09-04-25}{Modul-Isomorphismus und Lokalkompatibilität}
Die Abbildungen $\Phi$ und $\Psi$ sind Isomorphismen von $C^{\infty}(M)$-Modulen. 
Sie vertragen sich mit Einschränkungen auf offene Teilmengen $U\subset M$
(\emph{Sheaf-Naturalität}): Diagramme der Form
\[
\begin{tikzcd}
T^{r}{}_{s}(\mathfrak{X}(M)) \arrow[r,"\Phi_{M}"] \arrow[d,"r^{M}_{U}"'] 
& \mathcal{T}^{r}{}_{s}(M) \arrow[d,"r^{M}_{U}"] \\
T^{r}{}_{s}(\mathfrak{X}(U)) \arrow[r,"\Phi_{U}"]
& \mathcal{T}^{r}{}_{s}(U)
\end{tikzcd}
\]
pendeln. Insbesondere stimmen lokale Komponenten-Darstellungen unter beiden Sichtweisen überein.
\end{CONC}





Similar reasoning shows that there is a correspondence between fields of $T M$-valued tensors and $\mathfrak{X}(M)$-valued tensors on $\mathfrak{X}(M)$. For example, we have

$$
T_{2}^{0}(\mathfrak{X}(M) ; \mathfrak{X}(M)) \cong \mathcal{T}_{2}(M ; T M) .
$$

Elements of $T_{2}^{0}(\mathfrak{X}(M) ; \mathfrak{X}(M))$ are $C^{\infty}(M)$-bilinear maps $\mathfrak{X}(M) \times \mathfrak{X}(M) \rightarrow \mathfrak{X}(M)$, while elements of $\mathcal{T}_{2}(M ; T M)$ are sections of the bundle whose fiber at $p$ is $T_{2}^{0}\left(T_{p} M ; T_{p} M\right)$. So, if $A \in \mathcal{T}_{2}(M ; T M)$, then for each $p, A(p)$ is an $\mathbb{R}$-multilinear map $T_{p} M \times T_{p} M \rightarrow T_{p} M$. Similarly, we have

$$
T_{3}^{0}(\mathfrak{X}(M) ; \mathfrak{X}(M)) \cong \mathcal{T}_{3}(M ; T M) .
$$

In fact, later, when we define the curvature tensor on a semi-Riemannian manifold, it will initially be given as a multilinear map on modules of fields with values in $\mathfrak{X}(M)$. The correspondence is then invoked to get a tensor field (as defined in the bottom-up approach) with values in $T M$. It is just as easy to give a similar correspondence between $T^{r}{ }_{s}(\Gamma(\xi))$ and $\Gamma\left(T^{r}{ }_{s}(\xi)\right)$ for some vector bundle $\xi=(E, \pi, M)$, where we view $\Gamma(\xi)$ as a $C^{\infty}(M)$ module.




\begin{THEO}{GA-L09-04-26}{Kanonische Korrespondenz zwischen $\mathfrak{X}(M)$-wertigen und $TM$-wertigen Tensorfeldern}
Sei $M$ eine glatte Mannigfaltigkeit. Wir bezeichnen mit
\[
T^{r}{}_{s}\big(\mathfrak{X}(M);\,\mathfrak{X}(M)\big)
:= \mathrm{Multilin}_{C^{\infty}(M)}
\!\Big((\mathfrak{X}^{*}(M))^{r}\times(\mathfrak{X}(M))^{s},\,\mathfrak{X}(M)\Big)
\]
den $C^{\infty}(M)$-Modul der \textbf{Top-Down}-Tensorfelder mit Werten in $\mathfrak{X}(M)$
(somit also $\mathfrak{X}(M)$-wertige multilineare Abbildungen), und mit
\[
\mathcal{T}^{r}{}_{s}(M;TM):=\Gamma\!\big(T^{r}{}_{s}(TM;TM)\big)
\]
den $C^{\infty}(M)$-Modul der \textbf{Bottom-Up}-Tensorfelder mit Werten im Tangentialbündel $TM$,
also der glatten Schnitte des Bündels, dessen Faser über $p\in M$ der Raum
$T^{r}{}_{s}\big(T_{p}M;T_{p}M\big)$ ist.

Dann existiert ein kanonischer, natürlicher Isomorphismus von $C^{\infty}(M)$-Modulen
\[
\Phi:\ T^{r}{}_{s}\big(\mathfrak{X}(M);\mathfrak{X}(M)\big)
\ \xrightarrow{\ \cong\ }\ 
\mathcal{T}^{r}{}_{s}(M;TM),
\qquad
\Psi:=\Phi^{-1}.
\]

\textbf{Konstruktion.}
\begin{enumerate}
\item[\textbf{(A)}] 
\textbf{Abbildung $\Psi:\mathcal{T}^{r}{}_{s}(M;TM)\to T^{r}{}_{s}(\mathfrak{X}(M);\mathfrak{X}(M))$.}
Für ein glattes $TM$-wertiges Tensorfeld $A\in \mathcal{T}^{r}{}_{s}(M;TM)$
definiere
\[
(\Psi A)(\theta_{1},\dots,\theta_{r},X_{1},\dots,X_{s})
\in \mathfrak{X}(M)
\]
durch
\[
\big((\Psi A)(\theta_{1},\dots,\theta_{r},X_{1},\dots,X_{s})\big)(p)
:=A(p)\big(\theta_{1}(p),\dots,\theta_{r}(p),X_{1}(p),\dots,X_{s}(p)\big).
\]
Diese Definition ist $C^{\infty}(M)$-multilinear, da $A(p)$ in jedem Argument linear ist
und die Abbildung $p\mapsto A(p)(\dots)$ glatt variiert.

\item[\textbf{(B)}]
\textbf{Abbildung $\Phi:T^{r}{}_{s}(\mathfrak{X}(M);\mathfrak{X}(M))\to \mathcal{T}^{r}{}_{s}(M;TM)$.}
Sei $B\in T^{r}{}_{s}(\mathfrak{X}(M);\mathfrak{X}(M))$. Für $p\in M$ definiere
\[
(\Phi B)(p)\in T^{r}{}_{s}(T_{p}M;T_{p}M)
\]
durch
\[
(\Phi B)(p)(a_{1},\dots,a_{r},v_{1},\dots,v_{s})
:=B(\theta_{1},\dots,\theta_{r},X_{1},\dots,X_{s})(p),
\]
wobei $\theta_{i}\in\mathfrak{X}^{*}(M)$ und $X_{j}\in\mathfrak{X}(M)$ beliebige
glatte Fortsetzungen mit
\(
\theta_{i}(p)=a_{i},\ 
X_{j}(p)=v_{j}
\)
sind. 
Diese Definition ist wohldefiniert, da $B$ $C^{\infty}(M)$-multilinear ist
und der Wert an $p$ nur von den Werten der Argumente im Tangentialraum abhängt.

\item[\textbf{(C)}]
\textbf{Glattheit und Natürlichkeit.}
In lokalen Koordinaten $(x^{i})$ oder mit einem lokalen Frame $(s_{1},\dots,s_{n})$
für $TM$ lässt sich $B$ durch glatte Komponentenfunktionen darstellen.
Somit ist $p\mapsto (\Phi B)(p)$ ein glatter Schnitt des Bündels
$T^{r}{}_{s}(TM;TM)$. 
Beide Abbildungen sind $C^{\infty}(M)$-linear und respektieren
Einschränkungen auf offene Mengen $U\subset M$:
\[
r^{M}_{U}\circ \Phi_{M}=\Phi_{U}\circ r^{M}_{U},
\]
sowie Tensorprodukt und Kontraktion in jedem Faktor.

\item[\textbf{(D)}]
\textbf{Inversen-Eigenschaft.}
Für jedes $A\in\mathcal{T}^{r}{}_{s}(M;TM)$ und $B\in T^{r}{}_{s}(\mathfrak{X}(M);\mathfrak{X}(M))$ gilt
\[
\Phi(\Psi A)=A,\qquad
\Psi(\Phi B)=B.
\]
Somit sind $\Phi$ und $\Psi$ zueinander invers und liefern einen Isomorphismus
von $C^{\infty}(M)$-Modulen.
\end{enumerate}
\end{THEO}





\begin{KORO}{GA-L09-04-27}{Spezialfälle und Verallgemeinerung der $TM$-wertigen Korrespondenz}
\begin{enumerate}
\item Für $(r,s)=(0,2)$ erhalten wir den Isomorphismus
\[
T^{0}{}_{2}\big(\mathfrak{X}(M);\mathfrak{X}(M)\big)
\ \cong\ 
\mathcal{T}_{2}(M;TM),
\]
d.\,h. bilineare Abbildungen
$B:\mathfrak{X}(M)\times\mathfrak{X}(M)\to\mathfrak{X}(M)$
entsprechen genau den $TM$-wertigen $(0,2)$-Tensorfeldern.

\item Analog gilt
\[
T^{0}{}_{3}\big(\mathfrak{X}(M);\mathfrak{X}(M)\big)
\ \cong\ 
\mathcal{T}_{3}(M;TM),
\]
was insbesondere für die spätere Definition der Krümmungstensoren 
einer (semi-)Riemannschen Mannigfaltigkeit wesentlich ist:
Dort erscheint die Krümmung zunächst als 
$C^{\infty}(M)$-multilineare Abbildung
$\mathfrak{X}(M)\times\mathfrak{X}(M)\times\mathfrak{X}(M)\to\mathfrak{X}(M)$
und wird mittels dieser Korrespondenz als 
glatter Schnitt $R\in \Gamma(T^{1}{}_{3}(TM))$ interpretiert.

\item Die Konstruktion verallgemeinert sich unmittelbar auf beliebige Vektorbündel
$\xi=(E,\pi,M)$:
\[
T^{r}{}_{s}\big(\Gamma(\xi);\Gamma(\xi)\big)
\ \cong\
\Gamma\!\big(T^{r}{}_{s}(\xi;\xi)\big),
\]
wobei die rechte Seite den Modul der glatten Schnitte
des Bündels bezeichnet, dessen Faser über $p\in M$ der Raum
$T^{r}{}_{s}(E_{p};E_{p})$ ist.
\end{enumerate}
\end{KORO}




\begin{PROP}{GA-L09-04-28}{Charakterisierung von Tensoridentität auf lokal definierte, kommutierende Argumentfelder}
Seien $S,T\in\mathcal{T}^{r}{}_{s}(M)$ Tensorfelder gleichen Typs. 
Angenommen, für \emph{alle} lokal definierten Vektorfelder $X_{1},\dots,X_{s}$ mit 
paarweise verschwindenden Lie-Klammern $[X_{i},X_{j}]=0$ und \emph{alle} lokal definierten
$1$-Formen $\theta_{1},\dots,\theta_{r}$ gilt
\[
S(\theta_{1},\dots,\theta_{r},X_{1},\dots,X_{s})
\;=\;
T(\theta_{1},\dots,\theta_{r},X_{1},\dots,X_{s})
\quad\text{(als Funktion auf der gemeinsamen Definitionsmenge).}
\]
Dann folgt $S=T$ auf $M$.
\end{PROP}



\begin{PROOF}{GA-L09-04-29}{P: Charakterisierung von Tensoridentität auf lokal definierte, kommutierende Argumentfelder}
Es genügt, punktweise zu zeigen, dass $S_{p}=T_{p}$ als multilineare Abbildungen 
$T_{p}^{*}M{}^{\,r}\times T_{p}M{}^{\,s}\to\mathbb{R}$ übereinstimmen. 
Sei daher $p\in M$ fest sowie $a_{1},\dots,a_{r}\in T_{p}^{*}M$ und 
$v_{1},\dots,v_{s}\in T_{p}M$ beliebig. 

Wähle eine Karte $(U,x)$ um $p$ und schreibe $v_{j}=v_{j}^{k}\,\partial_{k}\rvert_{p}$. 
Definiere auf $U$ Vektorfelder mit \emph{konstanten} Koeffizienten
\[
Y_{j}\ :=\ v_{j}^{k}\,\frac{\partial}{\partial x^{k}}\quad(j=1,\dots,s).
\]
Dann gilt $Y_{j}(p)=v_{j}$ und wegen der konstanten Koeffizienten 
$[Y_{i},Y_{j}]=0$ auf $U$ (Linearkombination kommutierender Koordinatenfelder).
Ferner wähle lokale $1$-Formen $\theta_{i}$ mit $\theta_{i}(p)=a_{i}$.

Da $S$ und $T$ Tensorfelder sind, hängt die Auswertung am Punkt nur von den 
Werten der Argumente am Punkt ab (punktweise Abhängigkeit). Also
\[
S_{p}(a_{1},\dots,a_{r},v_{1},\dots,v_{s})
=
S(\theta_{1},\dots,\theta_{r},Y_{1},\dots,Y_{s})(p),
\]
und analog für $T$. Nach Voraussetzung ist 
$S(\theta,\!Y)(q)=T(\theta,\!Y)(q)$ für alle $q$ in einer Umgebung von $p$,
insbesondere am Punkt $p$. Somit
\[
S_{p}(a_{1},\dots,a_{r},v_{1},\dots,v_{s})
=
T_{p}(a_{1},\dots,a_{r},v_{1},\dots,v_{s}).
\]
Da $(a_{i})$ und $(v_{j})$ beliebig waren, folgt $S_{p}=T_{p}$ und damit $S=T$ auf $M$.

Die Konstruktion der kommutierenden Testfelder $Y_{j}$ benutzt nur, 
dass Koordinatenvektorfelder $\partial/\partial x^{k}$ vertauschen und 
dass Linearkombinationen mit \emph{konstanten} Koeffizienten ebenfalls vertauschen.
Dies reduziert Gleichheitsnachweise für Tensoren auf Auswertungen mit 
„koordinatenartigen“ Argumenten.
\end{PROOF}




Exercise 7.33. Exhibit the isomorphism (7.7) in detail.

Exercise 7.34. Suppose that $S$ and $T$ are tensors of the same type and we wish to show that they are equal. Then it is enough to check equality under the assumption that the vector fields inserted into the slots of $S$ and $T$ are locally defined and have vanishing Lie brackets. Hint: Think about coordinate vector fields.




We end this section with some warnings. It may seem that there is a simple way to obtain a pull-back by a smooth map $f: M \rightarrow N$ entirely from the top-down or module-theoretic view. In fact, one often sees expressions like
$$
f^{*} \tau\left(X_{1}, \ldots, X_{s}\right)=\tau\left(f_{*} X_{1}, \ldots, f_{*} X_{s}\right) \quad \text { (problematic expression!). }
$$

This looks cute, but invites misunderstanding. The left hand side takes fields $X_{1}, \ldots, X_{s}$ as arguments, while on the right hand side, if we consider $\tau$ as a multilinear map $\mathfrak{X}(N) \times \cdots \times \mathfrak{X}(N) \rightarrow C^{\infty}(N)$, then $f_{*} X_{i}$ must be fields. But the push-forward map $f_{*}$ is generally not defined on fields, and even if it were, the above expression would seem to be an equality of a function on $M$ with a function on $N$. Note that $\mathfrak{X}(M)$ is a $C^{\infty}(M)$-module, while $\mathfrak{X}(N)$ is a $C^{\infty}(N)$-module. The above expression may be taken to mean something like $f^{*} \tau\left(X_{1}, \ldots, X_{s}\right)(p)=\tau\left(T f \cdot X_{1}(p), \ldots, T f \cdot X_{s}(p)\right)$, but now the right hand side has tangent vectors as arguments, and we are back to the bottom-up approach! A correct statement is the following:

Proposition 7.35. 

Let $f: M \rightarrow N$ be a smooth map. Let $\tau$ be a ($0,s$)tensor field. If $\tau$ and $f^{*} \tau$ are interpreted as elements of $T_{s}^{0}(\mathfrak{X}(N))$ and $T_{s}^{0}(\mathfrak{X}(M))$ respectively, then
$$
f^{*} \tau\left(X_{1}, \ldots, X_{s}\right)=\tau\left(Y_{1}, \ldots, Y_{s}\right) \circ f
$$
whenever $Y_{i}$ is $f$-related to $X_{i}$ for $i=1, \ldots, s$.


Of course, we can use Definition 7.31 to make sense of both push-forward and pull-back in the case that $f$ is a diffeomorphism.



\begin{PROP}{GA-L09-04-30}{Allgemeines Funktorialitätsgesetz für Pull-Back von Tensorfeldern}
Sei $f:M\to N$ eine glatte Abbildung und $\tau\in\mathcal{T}^{r}{}_{s}(N)$ ein glattes $(r,s)$-Tensorfeld auf $N$.  
Dann gilt für beliebige $1$-Formfelder $\theta_{1},\dots,\theta_{r}\in\mathfrak{X}^{*}(M)$ und Vektorfelder $X_{1},\dots,X_{s}\in\mathfrak{X}(M)$ die folgende \textbf{Kompatibilitätsrelation}:
\[
\boxed{
\big(f^{*}\tau\big)(\theta_{1},\dots,\theta_{r},X_{1},\dots,X_{s})
\;=\;
\big(\tau\big((Tf^{-1})^{*}\theta_{1},\dots,(Tf^{-1})^{*}\theta_{r},\,Tf\!\cdot\!X_{1},\dots,Tf\!\cdot\!X_{s}\big)\big)\circ f
}
\]
wobei die Terme auf der rechten Seite wohldefiniert sind, falls $f$ ein Diffeomorphismus ist.  
Im Spezialfall eines rein kovarianten Tensors $\tau\in\mathcal{T}_{s}^{0}(N)$ genügt die schwächere Voraussetzung, dass Vektorfelder $Y_{i}\in\mathfrak{X}(N)$ existieren, die zu $X_{i}\in\mathfrak{X}(M)$ $f$-verwandt sind; dann gilt punktweise
\[
\boxed{
(f^{*}\tau)(X_{1},\ldots,X_{s})
=\big(\tau(Y_{1},\ldots,Y_{s})\big)\circ f.
}
\]
Diese Formeln drücken die natürliche Verträglichkeit des Pull-Back-Operators mit der Auswertung von Tensorfeldern auf $f$-verwandten Argumenten aus.
\end{PROP}


\pagebreak


\subsection*{Tensor Derivations}
We would like to be able to differentiate tensor fields. In particular, we would like to extend the Lie derivative to tensor fields. For this purpose we introduce the following definition, which will be useful not only for extending the Lie derivative, but also in several other contexts. Recall the presheaf of tensor fields $U \mapsto \mathcal{T}_{s}^{r}(U)$ on a manifold $M$.

Definition 7.36. A tensor derivation is a collection of maps $\left.\mathcal{D}_{s}^{r}\right|_{U}$ : $\mathcal{T}_{s}^{r}(U) \rightarrow \mathcal{T}_{s}^{r}(U)$, all denoted by $\mathcal{D}$ for convenience, such that\\
(1) $\mathcal{D}$ is a presheaf map for $\mathcal{T}_{s}^{r}$ considered as a presheaf of vector spaces over $\mathbb{R}$. In particular, for all open $U$ and $V$ with $V \subset U$ we have

$$
\left.\mathcal{D} A\right|_{V}=\mathcal{D}\left(\left.A\right|_{V}\right)
$$

for all $A \in \mathcal{T}_{s}^{r}(U)$, i.e., the restriction of $\mathcal{D} A$ to $V$ is just $\mathcal{D}\left(\left.A\right|_{V}\right)$.\\
(2) $\mathcal{D}$ commutes with contractions.\\
(3) $\mathcal{D}$ satisfies a derivation law. Specifically, for $A \in \mathcal{T}_{s}^{r}(U)$ and $B \in \mathcal{T}_{k}^{j}(U)$ we have

$$
\mathcal{D}(A \otimes B)=\mathcal{D} A \otimes B+A \otimes \mathcal{D} B
$$

For smooth $n$-manifolds, the conditions (2) and (3) imply that for $A \in \mathcal{T}_{s}^{r}(U), \alpha_{1}, \ldots, \alpha_{r} \in \mathfrak{X}^{*}(U)$ and $X_{1}, \ldots, X_{s} \in \mathfrak{X}(U)$, we have

$$
\begin{aligned}
& \mathcal{D}\left(A\left(\alpha_{1}, \ldots, \alpha_{r}, X_{1}, \ldots, X_{s}\right)\right) \\
&=(\mathcal{D} A)\left(\alpha_{1}, \ldots, \alpha_{r}, X_{1}, \ldots, X_{s}\right) \\
&+\sum_{i=1}^{r} A\left(\alpha_{1}, \ldots, \mathcal{D} \alpha_{i}, \ldots, \alpha_{r}, X_{1}, \ldots, X_{s}\right) \\
&+\sum_{i=1}^{s} A\left(\alpha_{1}, \ldots, \alpha_{r}, X_{1}, \ldots, \mathcal{D} X_{i}, \ldots, X_{s}\right)
\end{aligned}
$$

This follows by noticing that

$$
A\left(\alpha_{1}, \ldots, \alpha_{r}, X_{1}, \ldots, X_{s}\right)=C\left(A \otimes\left(\alpha_{1} \otimes \cdots \otimes \alpha_{r} \otimes X_{1} \otimes \cdots \otimes X_{s}\right)\right)
$$

(where $C$ is the repeated contraction) and then applying (2) and (3).\\
Note that $\mathcal{D}$ stands for a family of maps whose domains $\mathcal{T}_{s}^{r}(U)$ depend not only on $r$ and $s$, but also on $U$. The next proposition considers the situation where we only have derivations defined for $U=M$ (the global case).

Proposition 7.37. Let $M$ be a smooth manifold and suppose we have a map on globally defined tensor fields $\mathcal{D}: \mathcal{T}_{s}^{r}(M) \rightarrow \mathcal{T}_{s}^{r}(M)$ for all nonnegative integers $r, s$ such that (2) and (3) above hold for the case $U=M$. Then there is a unique induced tensor derivation that agrees with $\mathcal{D}$ on global sections, that is, on the various $\mathcal{T}_{s}^{r}(M)$.

Proof. We need to define $\mathcal{D}: \mathcal{T}_{s}^{r}(U) \rightarrow \mathcal{T}_{s}^{r}(U)$ for arbitrary open $U$ as a derivation. Let $\delta$ be a function in $C^{\infty}(U)$ that vanishes on a neighborhood $V$ of $p \in U$. We claim that $(\mathcal{D} \delta)(p)=0$. To see this, let $\beta$ be a cut-off function equal to 1 on a neighborhood of $p$ and zero outside of $V$. Then $\delta=(1-\beta) \delta$ and so

$$
\begin{aligned}
\mathcal{D} \delta(p) & =\mathcal{D}((1-\beta) \delta)(p) \\
& =\delta(p) \mathcal{D}(1-\beta)(p)+(1-\beta(p)) \mathcal{D} \delta(p)=0 .
\end{aligned}
$$

Now given $\tau \in \mathcal{T}_{s}^{r}(U)$, let $\beta$ be a cut-off function with support in $U$ and equal to 1 on a neighborhood of $p \in U$. Then $\beta \tau \in \mathcal{T}_{s}^{r}(M)$ after extending by zero. Define $(\mathcal{D} \tau)(p):=\mathcal{D}(\beta \tau)(p)$. To show that this is well-defined let $\beta_{2}$ be any other cut-off function with support in $U$ and equal to 1 on a neighborhood of $p \in U$. Then we have

$$
\begin{aligned}
& \mathcal{D}(\beta \tau)(p)-\mathcal{D}\left(\beta_{2} \tau\right)(p) \\
& \quad=\left(\mathcal{D}(\beta \tau)-\mathcal{D}\left(\beta_{2} \tau\right)\right)(p)=\mathcal{D}\left(\left(\beta-\beta_{2}\right) \tau\right)(p)=0
\end{aligned}
$$

where the last equality follows from our claim above with $\delta=\beta-\beta_{2}$. Thus $\mathcal{D}$ is well-defined on $\mathcal{T}_{s}^{r}(U)$. We now show that $\mathcal{D} \tau$ so defined is an element of $\mathcal{T}_{s}^{r}(U)$. Let ( $U^{\prime}, \mathrm{x}$ ) be a chart with $p \in U^{\prime} \subset U$. Then we can write $\left.\tau\right|_{U^{\prime}} \in \mathcal{T}_{s}^{r}\left(U^{\prime}\right)$ as

$$
\tau_{U^{\prime}}=\tau_{i_{1}, \ldots, i_{s}}^{j_{1}, \ldots, j_{r}} d x^{i_{1}} \otimes \cdots \otimes d x^{i_{s}} \otimes \frac{\partial}{\partial x^{j_{1}}} \otimes \cdots \otimes \frac{\partial}{\partial x^{j_{r}}}
$$

We can use this to show that $\mathcal{D} \tau$ as defined agrees with a global section in a neighborhood $O$ of $p$ and so must be a smooth section itself since the choice of $p \in U$ was arbitrary. To save on notation, let us take the case $r=1$, $s=1$. Then $\tau_{U_{\alpha}}=\tau_{j}^{i} d x^{j} \otimes \frac{\partial}{\partial x^{i}}$. Let $\beta$ be a cut-off function equal to 1 in the neighborhood $O$ of $p$ and zero outside of $U^{\prime}$. Extend each of the sections $\beta \tau_{j}^{i}$, $\beta d x^{j}$, and $\beta \frac{\partial}{\partial x^{i}}$ to global sections and apply $\mathcal{D}$ to $\beta^{3} \tau=\left(\beta \tau_{j}^{i}\right)\left(\beta d x^{j}\right) \otimes\left(\beta \frac{\partial}{\partial x^{i}}\right)$ to get

$$
\begin{aligned}
\mathcal{D}\left(\beta^{3} \tau\right)= & \mathcal{D}\left(\beta \tau_{j}^{i} \beta d x^{j} \otimes \beta \frac{\partial}{\partial x^{i}}\right) \\
= & \mathcal{D}\left(\beta \tau_{j}^{i}\right) \beta d x^{j} \otimes \beta \frac{\partial}{\partial x^{i}}+\beta \tau_{j}^{i} \mathcal{D}\left(\beta d x^{j}\right) \otimes \beta \frac{\partial}{\partial x^{i}} \\
& +\beta \tau_{j}^{i} \beta d x^{j} \otimes \mathcal{D}\left(\beta \frac{\partial}{\partial x^{i}}\right)
\end{aligned}
$$

By assumption, $\mathcal{D}$ takes smooth global sections to smooth global sections, so both sides of the above equation are smooth. On the other hand, we have $\mathcal{D}\left(\beta^{3} \tau\right)(q)=\mathcal{D}(\tau)(q)$ by definition and valid for all $q \in O$. Thus $\mathcal{D}(\tau)$ is smooth and is the restriction of a smooth global section. This gives a unique derivation $\mathcal{D}: \mathcal{T}_{s}^{r}(U) \rightarrow \mathcal{T}_{s}^{r}(U)$ for all $U$ satisfying the naturality conditions (1), (2) and (3). We leave it to the reader to check this last statement.

Exercise 7.38. Let $\mathcal{D}_{1}$ and $\mathcal{D}_{2}$ be two tensor derivations (so satisfying conditions (1), (2) and 3 of Definition 7.36) that agree on functions and vector fields. Then $\mathcal{D}_{1}=\mathcal{D}_{2}$. [Hint: If $\alpha \in \mathfrak{X}^{*}(U)=\mathcal{T}_{1}^{0}(U)$, we must have $\left(\mathcal{D}_{i} \alpha\right)(X)=\mathcal{D}_{i}(\alpha(X))-\alpha\left(\mathcal{D}_{i} X\right)$ for $i=1,2$. Then both $\mathcal{D}_{1}$ and $\mathcal{D}_{2}$ must obey formula (7.8) above.]

Theorem 7.39. If $\mathcal{D}_{U}$ can be defined on $C^{\infty}(U)$ and $\mathfrak{X}(U)$ for each open $U \subset M$ so that\\
(1) $\mathcal{D}_{U}(f g)=\left(\mathcal{D}_{U} f\right) g+f \mathcal{D}_{U} g$ for all $f, g \in C^{\infty}(U)$,\\
(2) $\left.\left(\mathcal{D}_{M} f\right)\right|_{U}=\left.\mathcal{D}_{U} f\right|_{U}$ for each $f \in C^{\infty}(M)$,\\
(3) $\mathcal{D}_{U}(f X)=\left(\mathcal{D}_{U} f\right) X+f \mathcal{D}_{U} X$ for all $f \in C^{\infty}(U)$ and $X \in \mathfrak{X}(U)$,\\
(4) $\left.\left(\mathcal{D}_{M} X\right)\right|_{U}=\left.\mathcal{D}_{U} X\right|_{U}$ for each $X \in \mathfrak{X}(M)$,\\
then there is a unique tensor derivation $\mathcal{D}$ that is equal to $\mathcal{D}_{U}$ on $C^{\infty}(U)$ and $\mathfrak{X}(U)$ for all $U$.

Sketch of proof. We wish to define $\mathcal{D}$ on $\mathfrak{X}^{*}(U)$ so that

$$
\mathcal{D}_{U}(\alpha \otimes X)=\mathcal{D}_{U} \alpha \otimes X+\alpha \otimes \mathcal{D}_{U} X
$$

By contraction we see that we must have $\left(\mathcal{D}_{U} \alpha\right)(X)=\mathcal{D}_{U}(\alpha(X))-\alpha\left(\mathcal{D}_{U} X\right)$, which we take as the definition. Then check that (7.9) holds. Now define $\mathcal{D}_{U}$ by formula (7.8) and verify that we really have a map $\mathcal{T}_{s}^{r}(U) \rightarrow \mathcal{T}_{s}^{r}(U)$. Check that $\mathcal{D}_{U}$ commutes with contraction $C: \mathcal{T}_{1}^{1}(U) \rightarrow C^{\infty}(U)$ for simple tensors $\alpha \otimes X \in \mathcal{T}_{1}^{1}(U)$. Use the fact that, locally, every element of $\mathcal{T}_{1}^{1}$ can be written as a sum of simple tensors. Next extend to $\mathcal{T}_{s}^{r}$ along the lines exemplified by the case of $\mathcal{T}_{2}^{1}(U)$ and the contraction $C_{2}^{1}$ as follows: For $\tau \in \mathcal{T}_{2}^{1}(U)$, we have

$$
\begin{aligned}
\left(\mathcal{D}_{U} C_{2}^{1} \tau\right)(X) & =\mathcal{D}_{U}\left(\left(C_{2}^{1} \tau\right)(X)\right)-\left(C_{2}^{1} \tau\right) \mathcal{D}_{U} X \\
& =\mathcal{D}_{U}(C(\tau(\cdot, X, \cdot)))-C\left(\tau\left(\cdot, \mathcal{D}_{U} X, \cdot\right)\right) \\
& =C \mathcal{D}_{U}\left(\tau(\cdot, X, \cdot)-\tau\left(\cdot, \mathcal{D}_{U} X, \cdot\right)\right) \\
& =C\left(\left(\mathcal{D}_{U} \tau\right)(\cdot, X, \cdot)\right)=\left(C_{2}^{1} \mathcal{D}_{U} \tau\right)(X)
\end{aligned}
$$

The general case would involve an inconvenient profusion of parentheses. Uniqueness follows from Exercise 7.38. Finally check by direct calculation that (3) of Definition 7.36 holds.

Corollary 7.40. The Lie derivative $\mathcal{L}_{X}$ can be extended to a tensor derivation for any $X \in \mathfrak{X}(M)$.

The last corollary extends the Lie derivative to tensor fields. It follows from formula (7.8) that we have

$$
\begin{aligned}
\left(\mathcal{L}_{X} S\right)\left(Y_{1}, \ldots, Y_{s}\right)= & X\left(S\left(Y_{1}, \ldots, Y_{s}\right)\right) \\
& -\sum_{i=1}^{s} S\left(Y_{1}, \ldots, Y_{i-1}, \mathcal{L}_{X} Y_{i}, Y_{i+1}, \ldots, Y_{s}\right)
\end{aligned}
$$

We now present a different way of extending the Lie derivative to tensor fields that is equivalent to what we have just done. First let $A \in \mathcal{T}_{s}^{r}(M)$\\
and recall that if $f: M \rightarrow M$ is a diffeomorphism, then we can define $f^{*} A \in \mathcal{T}_{s}^{r}(M)$ by

$$
\begin{aligned}
& \left(f^{*} A\right)(p)\left(\alpha^{1}, \ldots, \alpha^{r}, v_{1}, \ldots, v_{s}\right) \\
& \quad=A(f(p))\left(\left(T_{p} f^{-1}\right)^{*}\left(\alpha^{1}\right), \ldots,\left(T_{p} f^{-1}\right)^{*}\left(\alpha^{r}\right), T_{p} f\left(v_{1}\right), \ldots, T_{p} f\left(v_{s}\right)\right)
\end{aligned}
$$

for all $\alpha^{1}, \ldots, \alpha^{r} \in\left(T_{p} M\right)^{*}$ and $v_{1}, \ldots, v_{s} \in T_{p} M$. If $X$ is a vector field on $M$ (possibly locally defined), we can define

$$
\left(\mathcal{L}_{X} A\right)(p)=\left.\frac{d}{d t}\right|_{0}\left(\left(\varphi_{t}^{X}\right)^{*} A\right)(p)
$$

just as we did for vector fields. We leave it as a project for the reader to show that this definition agrees with our first definition of the Lie derivative of a tensor field.

The Lie derivative on tensor fields is natural with respect to diffeomorphisms in the sense that for any diffeomorphism $f: M \rightarrow N$ and any vector field $X$ we have

$$
\mathcal{L}_{f_{*} X} f_{*} A=f_{*} \mathcal{L}_{X} A
$$

This property is not shared by some other important derivations such as the covariant derivative, which we define later in this book.

Exercise 7.41. Show that the Lie derivative on tensor fields is natural with respect to diffeomorphisms in the above sense of equation (7.12) by using the fact that it is natural on functions and vector fields.

\subsection*{Metric Tensors}
We start out again considering some linear algebra that we wish to globalize. Thus the vector space $V$ that we discuss next should be thought of as a tangent space of a manifold or a fiber of some vector bundle.

We recall the following definitions: A symmetric bilinear form $g$ on a finite-dimensional vector space V is nondegenerate if and only if $g(v, w)=$ 0 for all $w \in \mathrm{V}$ implies that $v=0$. A (real) scalar product on a (real) finite-dimensional vector space V is a nondegenerate symmetric bilinear form $g: \mathrm{V} \times \mathrm{V} \rightarrow \mathbb{R}$. A scalar product space is a pair ( $\mathrm{V}, g$ ) where V is a vector space and $g$ is a scalar product. We say that $g$ is positive (resp. negative) definite if $g(v, v) \geq 0$ (resp. $g(v, v) \leq 0$ ) for all $v \in \mathrm{V}$ and $g(v, v)=0 \Longrightarrow v=0$. In case the scalar product is positive definite, we also refer to it as an inner product and the pair ( $\mathrm{V}, g$ ) as an inner product space. Otherwise we say that the scalar product is indefinite. A scalar product on V is sometimes called a metric tensor on V . We now need to introduce quite a few more definitions.

Definition 7.42. The index of a symmetric bilinear form $g$ on V is the dimension of the largest subspace $\mathrm{W} \subset \mathrm{V}$ such that the restriction $\left.g\right|_{W}$ is negative definite. The index is denoted $\operatorname{ind}(g)$.

Definition 7.43. Let ( $\mathrm{V}, g$ ) be a scalar product space. We say that $v$ and $w$ are mutually orthogonal if and only if $g(v, w)=0$. Furthermore, given two subspaces $W_{1}$ and $W_{2}$ of $V$ we say that $W_{1}$ is orthogonal to $W_{2}$ and write $\mathrm{W}_{1} \perp \mathrm{~W}_{2}$ if and only if every element of $\mathrm{W}_{1}$ is orthogonal to every element of $\mathrm{W}_{2}$.

Since, in general, $g$ is not necessarily positive definite or negative definite, there may be nonzero elements that are orthogonal to themselves.

Definition 7.44. Given a subspace $W$ of a scalar product space $V$, we define the orthogonal complement as $\mathrm{W}^{\perp}=\{v \in \mathrm{V}: g(v, w)=0$ for all $w \in \mathrm{~W}\}$.

Exercise 7.45. We always have $\operatorname{dim}(\mathrm{W})+\operatorname{dim}\left(\mathrm{W}^{\perp}\right)=\operatorname{dim}(\mathrm{V})$, but unless $g$ is definite, we may not have $\mathrm{W} \cap \mathrm{W}^{\perp}=\{0\}$.

Definition 7.46. A subspace W of a scalar product space ( $V, g$ ) is called nondegenerate if $\left.g\right|_{\mathrm{W}}$ is nondegenerate.

Lemma 7.47. A subspace $\mathrm{W} \subset(\mathrm{V}, g)$ is nondegenerate if and only if $\mathrm{V}= \mathrm{W} \oplus \mathrm{W}^{\perp}$ (inner direct sum).

Proof. This an easy exercise. One uses the standard fact that

$$
\operatorname{dim} W+\operatorname{dim} W^{\perp}=\operatorname{dim}\left(W+W^{\perp}\right)+\operatorname{dim}\left(W \cap W^{\perp}\right)
$$

It is a standard fact from linear algebra, already mentioned in Chapter 5, that if $g$ is a scalar product, then there exists a basis $e_{1}, \ldots, e_{n}$ for V such that the matrix representative of $g$ with respect to this basis is a diagonal matrix with ones or minus ones along the diagonal. Such a basis is called an orthonormal basis for $(\mathrm{V}, g)$. The number of minus ones appearing is the index $\operatorname{ind}(g)$ and so is independent of the orthonormal basis chosen. It is easy to see that the index $\operatorname{ind}(g)$ is zero if and only if $g$ is positive definite.

Definition 7.48. For each $v \in \mathrm{V}$ with $\langle v, v\rangle \neq 0$, let $\epsilon(v):=\operatorname{sgn}\langle v, v\rangle$. Then if $e_{1}, \ldots, e_{n}$ are orthonormal, we have $\epsilon_{i}=\epsilon(i):=\epsilon\left(e_{i}\right)$.

Thus if $e_{1}, \ldots, e_{n}$ is an orthonormal basis for ( $\mathrm{V}, g$ ), then $g\left(e_{i}, e_{j}\right)=\epsilon_{i} \delta_{i j}$, where $\epsilon_{i}=g\left(e_{i}, e_{i}\right)= \pm 1$ are the entries of the diagonal matrix $\operatorname{ind}(g)$ of which are equal to -1 and the remaining are equal to 1 . Let us refer to the list of $\pm 1$ 's given by $\left(\epsilon_{1}, \ldots, \epsilon_{n}\right)$ as the signature. We may arrange for the -1 's to come first by permuting the elements of the basis. For example, if $(-1,-1,1,1)$ is the signature, then the index is 2 .

Remark 7.49. From now on, whenever context allows, we shall always assume that by "orthonormal basis" we mean an orthonormal basis that is arranged so that the -1 's come first as described above. The convention of putting the minus signs first is not universal, and in fact we used the opposite convention in Chapter 5. The negative signs first convention is popular in relativity theory and semi-Riemannian geometry, but the reverse convention is perhaps more common in Lie group theory and quantum field theory. It makes no difference in the final analysis as long as one is consistent, but it can be confusing when comparing references in the literature.

Another difference between the theory of positive definite scalar products and indefinite scalar products is the appearance of the $\epsilon_{i}$ 's from the signature in formulas that would be familiar in the positive definite case. For example, we have the following:

Proposition 7.50. Let $e_{1}, \ldots, e_{n}$ be an orthonormal basis for $(\mathrm{V}, g)$. For any $v \in \mathrm{V}$, we have a unique expansion given by $v=\sum_{i} \epsilon_{i}\left\langle v, e_{i}\right\rangle e_{i}$.

Proof. The usual proof works. One just has to notice the appearance of the $\epsilon_{i}$ 's.

Definition 7.51. If $v \in \mathrm{V}$, then let $\|v\|$ denote the nonnegative number $|g(v, v)|^{1 / 2}$ and call this the (absolute or positive) length or norm of $v$.

Some authors call $g(v, v)$ or $g(v, v)^{1 / 2}$ the norm, which would make it possible for the norm to be negative or even complex-valued. We will avoid this.

Just as for positive definite inner product spaces, we call a linear isomorphism $\Phi:\left(\mathrm{V}_{1}, g_{1}\right) \rightarrow\left(\mathrm{V}_{2}, g_{2}\right)$ from one scalar product space to another an isometry if $g_{1}(v, w)=g_{2}(\Phi v, \Phi w)$. It is not hard to show that if such an isometry exists, then $g_{1}$ and $g_{2}$ have the same index and signature.

Let ( $\mathrm{V}_{i}, g_{i}$ ) be scalar product spaces for $i=1, \ldots, k$. By Corollary D. 35 of Appendix D, there is a unique bilinear form

$$
\varphi: \bigotimes_{i=1}^{k} \mathrm{V}_{i} \times \bigotimes_{i=1}^{k} \mathrm{V}_{i} \rightarrow \mathbb{R}
$$

such that for $v_{i} \in \mathrm{V}_{i}$ and $w_{i} \in \mathrm{~W}_{i}$,

$$
\varphi\left(v_{1} \otimes \cdots \otimes v_{k}, w_{1} \otimes \cdots \otimes w_{k}\right)=g_{1}\left(v_{1}, w_{1}\right) \cdots g_{k}\left(v_{k}, w_{k}\right)
$$

The form $\varphi$ is clearly symmetric and by Problem 3, it is a scalar product on $\bigotimes_{i=1}^{k} \mathrm{V}_{i}$ that is positive definite if each $g_{i}$ is positive definite. If $(\mathrm{V}, g)$ is a scalar product space, then we use the above to endow each of the various tensor spaces $T^{r}{ }_{s}(\mathrm{V})$ with a scalar product. First consider V*. Since $g$ is nondegenerate, there is a linear isomorphism $g_{b}: \mathrm{V} \rightarrow \mathrm{V}^{*}$ defined by

$$
g_{b}(v)(w)=g(v, w)
$$

Denote the inverse by $g^{\sharp}: \mathrm{V}^{*} \rightarrow \mathrm{V}$. We force this to be an isometry by defining the scalar product on $\mathrm{V}^{*}$ to be

$$
g^{*}(\alpha, \beta)=g\left(g^{\sharp}(\alpha), g^{\sharp}(\beta)\right) .
$$

Under this prescription, the dual basis $\left(e^{1}, \ldots, e^{n}\right)$ corresponding to an orthonormal basis $\left(e_{1}, \ldots, e_{n}\right)$ for V will also be orthonormal. The signatures (and hence the indexes) of $g^{*}$ and $g$ are the same.

Notation 7.52. When convenient, we shall also denote $g_{b}(v)$ by either b $v$ or $v^{b}$ and similarly for $g^{\sharp}(\alpha)$.

The above procedure now applies to give a scalar product on any tensor space $T_{s}^{r}(\mathrm{V})$. For example, consider $T_{1}^{1}(\mathrm{V})=\mathrm{V} \otimes \mathrm{V}^{*}$. Then there is a unique scalar product $g_{1}^{1}$ on $\mathrm{V} \otimes \mathrm{V}^{*}$ such that for $v_{1} \otimes \alpha_{1}$ and $v_{2} \otimes \alpha_{2} \in \mathrm{V} \otimes \mathrm{V}^{*}$ we have

$$
g_{1}^{1}\left(v_{1} \otimes \alpha_{1}, v_{2} \otimes \alpha_{2}\right)=g\left(v_{1}, v_{2}\right) g^{*}\left(\alpha_{1}, \alpha_{2}\right)
$$

One can then see that for orthonormal $e_{1}, \ldots, e_{n}$ we have that

$$
\left\{e_{i} \otimes e^{j}\right\}_{1 \leq i, j \leq n}
$$

is an orthonormal basis for $\left(T_{1}^{1}(\mathrm{V}), g_{1}^{1}\right)$. In general, if we endow $T_{s}^{r}(\mathrm{V})$ with a scalar product as above, then the natural basis for $T_{s}^{r}(\mathrm{V})$ formed from the orthonormal basis $\left(e_{1}, \ldots, e_{n}\right)$ (and its dual $\left(e^{1}, \ldots, e^{n}\right)$ ) will also be orthonormal.

Notation 7.53. In order to reduce notational clutter, let us reserve the option to denote all these scalar products coming from $g$ by the same letter $g$ or, even more conveniently, by $\langle\cdot, \cdot\rangle$. So, for example, $\left\langle v_{1} \otimes \alpha_{1}, v_{2} \otimes \alpha_{2}\right\rangle= \left\langle v_{1}, v_{2}\right\rangle\left\langle\alpha_{1}, \alpha_{2}\right\rangle$ by definition.

Exercise 7.54. Show that under the natural identification of $\mathrm{V} \otimes \mathrm{V}^{*}$ with $L(\mathrm{V}, \mathrm{V})$, the scalar product of linear transformations $A$ and $B$ is given by $\langle A, B\rangle=\operatorname{trace}\left(A^{t} B\right)$.

The maps $g_{b}$ and $g^{\sharp}$ are called musical isomorphisms. Let us see how things look in terms of components. Let $f_{1}, \ldots, f_{n}$ be an arbitrary basis of V and let $f^{1}, \ldots, f^{n}$ be the dual basis for $\mathrm{V}^{*}$. The components of $g$ are given by $g_{i j}:=g\left(f_{i}, f_{j}\right)$. So if $v=v^{i} f_{i}$ and $w=w^{i} f_{i}$, then

$$
g(v, w)=g\left(v^{i} f_{i}, w^{j} f_{j}\right)=v^{i} w^{j} g\left(f_{i}, f_{j}\right)=g_{i j} v^{i} w^{j}
$$

where we continue to use the Einstein summation convention. There must be a matrix $\left(A_{i j}\right)$ such that $b f_{i}=A_{k i} f^{k}$. On the other hand,

$$
\begin{aligned}
g_{i j} & =g\left(f_{i}, f_{j}\right)=\left(b f_{i}\right)\left(f_{j}\right)=A_{k i} f^{k}\left(f_{j}\right) \\
& =A_{k i} \delta_{j}^{k}=A_{j i}
\end{aligned}
$$

So we have

$$
b f_{i}=g_{i k} f^{k}
$$

Thus if $v=v^{i} f_{i}$, then $b v=v^{j} b f_{j}=v^{j} g_{j i} f^{i}$, so the components of $b v$ are $(b v)_{i}=v^{j} g_{j i}=g_{i j} v^{j}$. It is a common convention that if $v^{i}$ are the components of $v$ with respect to a basis, then the components of bv are denoted simply by lowering the index:

$$
v_{i}=g_{i j} v^{j}
$$

The map $g_{b}$ is called the flatting operator and the effect of this operator is sometimes described as "index lowering". If we write $\sharp f^{i}=g^{i j} f_{j}$ for some matrix $\left(g^{i j}\right)$, then $\left(g_{i j}\right)^{-1}=\left(g^{i j}\right)$ so that

$$
g^{i k} g_{k j}=\delta_{j}^{k}
$$

This follows from $f^{i}=b \sharp f^{i}=g^{i k} b f_{k}=g^{i k} g_{k j} f^{j}$. If $\omega \in \mathrm{V}^{*}$ is written as $\omega=\omega_{i} f^{i}$, then an easy calculation shows that the components of $\sharp \omega$ are given by

$$
\omega^{i}:=(\sharp \omega)^{i}=g^{i k} \omega_{k} .
$$

The isomorphism $g^{\sharp}$ is called the sharping operator and its effect is referred to as "index raising". The scalar product $g_{*}$ on $\mathrm{V}^{*}$ introduced above has components $g^{i j}$ with respect to the dual basis. Indeed,

$$
\begin{aligned}
g^{*}\left(f^{i}, f^{j}\right) & =g\left(\sharp f^{i}, \sharp f^{j}\right)=g\left(g^{i k} f_{k}, g^{j l} f_{l}\right) \\
& =g^{i k} g^{j l} g_{k l}=g^{j l} \delta_{l}^{i}=g^{j i}=g^{i j} .
\end{aligned}
$$

Next we see how to extend the notions of index raising and lowering to tensors. Suppose we have a tensor $A \in T_{2}^{2}(\mathrm{V})$ and we wish to obtain a new related tensor $A^{\prime}$ whose final slot takes elements of $\mathrm{V}^{*}$ rather than elements of V . Thus we want $A^{\prime}$ be of type $\left({ }^{2}{ }_{1}{ }^{1}\right)$. The trick is to define $A^{\prime}$ using $\sharp$ as follows:

$$
A^{\prime}(\omega, \eta, v, \alpha):=A(\omega, \eta, v, \sharp \alpha) .
$$

Let us compute the components of $A^{\prime}$ with respect to our basis. We have

$$
\begin{aligned}
\left(A^{\prime}\right)_{k}^{i j} & =A^{\prime}\left(f^{i}, f^{j}, f_{k}, f^{\ell}\right)=A\left(f^{i}, f^{j}, f_{k}, \sharp f^{\ell}\right) \\
& =A\left(f^{i}, f^{j}, f_{k}, g^{\ell r} f_{k r}\right)=g^{\ell r} A_{k r}^{i j}
\end{aligned}
$$

It is common to write $A^{i j}{ }_{k}{ }^{\ell}$ in place of $\left(A^{\prime}\right)^{i j}{ }_{k}{ }^{\ell}$ when the context makes the meaning clear. In other words, we use the same root symbol but reposition the indices. This is an instance of index raising. Similarly we might use the flatting operation to obtain a new tensor from $A$. For $\omega \in \mathrm{V}^{*}$ and $u, v, w \in \mathrm{V}$ we could define $A^{\prime}$ by

$$
A^{\prime}(u, \omega, v, w):=A(b u, \omega, v, w)
$$

and the components of $A^{\prime}$ would be given by

$$
\left(A^{\prime}\right)_{i}^{j}{ }_{k \ell}:=g_{i r} A_{k \ell}^{r j},
$$

and again it is common to see simply $A_{i}{ }^{j}{ }_{k \ell}$. This process of raising and lowering indices is called type changing. Notice that we often obtain tensors that are not in consolidated form. However, one may simply apply consolidation as desired. But, notice that the staggering of position makes the relation of $A_{i}{ }^{j}{ }_{k \ell}$ to the original tensor $A^{i j}{ }_{k \ell}$ clear. This is where positional index notation excels.

The above can be approached in a slightly different way. We can take tensor products of various combinations of $g_{b}, g^{\sharp}$ and the identity maps idv and $\operatorname{id}_{V^{*}}$. For example,\\
$g_{b} \otimes \mathrm{id}_{\mathrm{V}} \otimes g_{b} \otimes g^{\sharp} \otimes \mathrm{id}_{\mathrm{V}^{*}}: \mathrm{V} \otimes \mathrm{V} \otimes \mathrm{V} \otimes \mathrm{V}^{*} \otimes \mathrm{V}^{*} \rightarrow \mathrm{V}^{*} \otimes \mathrm{V} \otimes \mathrm{V}^{*} \otimes \mathrm{V} \otimes \mathrm{V}^{*}$.\\
Depending on convenience, this map might then be followed by the consolidation isomorphism

$$
\mathrm{V}^{*} \otimes \mathrm{V} \otimes \mathrm{V}^{*} \otimes \mathrm{V} \otimes \mathrm{V}^{*} \rightarrow \mathrm{V}^{*} \otimes \mathrm{V}^{*} \otimes \mathrm{V}^{*} \otimes \mathrm{V} \otimes \mathrm{V}
$$

Exercise 7.55. Show that the map $g_{\mathrm{b}} \otimes \operatorname{id}_{\mathrm{V}} \otimes g_{\mathrm{b}} \otimes g^{\sharp} \otimes \operatorname{id}_{\mathrm{V}^{*}}$ effects three iterated type changes and is given in component form by

$$
A^{i j k}{ }_{\ell m} \mapsto A_{i}{ }^{j}{ }_{k m}^{\ell}:=g_{i a} g_{k b} g^{\ell c} A^{a j b}{ }_{c m} .
$$

In the presence of a scalar product, a type-changed tensor is considered to be just a different manifestation of the original tensor. We say that it is metrically equivalent. The reader should expect to see some slight variability with regard to how index positioning and order of slots is handled when type changing is done. For example, there is nothing stopping us from raising the $a$-th lower index into the $b$-th upper position while keeping everything in consolidated form:

$$
\left(A^{\prime}\right)^{i_{1} \ldots i_{b-1} j i_{b+1} \ldots i_{r}}{ }_{j_{1} \ldots j_{a-1} j_{a+1} \ldots j_{s}}:=A^{i_{1} \ldots i_{r}}{ }_{j_{1} \ldots j_{a-1} m j_{a+1} \ldots j_{s}} g^{m j} .
$$

Invariantly, this is described as

$$
\begin{aligned}
& A^{\prime}\left(\alpha^{1}, \ldots, \alpha^{r+1}, v_{1}, \ldots, v_{s-1}\right) \\
& \quad:=A\left(\alpha^{1}, \ldots, \widehat{\alpha^{b}}, \ldots, \alpha^{r+1}, v_{1}, \ldots, v_{a-1}, g^{\sharp}\left(\alpha^{b}\right), v_{a+1}, \ldots, v_{s}\right) .
\end{aligned}
$$

Notice that if we raise all the lower indices and lower all the upper indices on a tensor, then we can "completely contract" against another tensor of the original type. We leave it to the reader to show that the result is the scalar product of tensors defined earlier. For example, let $\chi=\sum \chi_{i j} f^{i} \otimes f^{j}$ and $\tau=\sum \tau_{i j} f^{i} \otimes f^{j}$. We may apply two type changes to $\tau$ that are given\\
in component form as $\tau_{i j} \mapsto \tau^{i}{ }_{j} \mapsto \tau^{i j}$. In other words, $\tau^{i j}=g^{i k} g^{j l} \tau_{k l}$. Then we have

$$
\langle\chi, \tau\rangle=\chi_{i j} \tau^{i j}
$$

See Problem 18.\\
7.6.1. Metrics on manifolds. If $g \in \mathcal{T}_{2}(M)$ is nondegenerate, symmetric and positive definite at every tangent space, we call $g$ a Riemannian metric (tensor). If $g$ is a Riemannian metric, then we call the pair ( $M, g$ ) a Riemannian manifold. For example, in Chapter 4, we saw how a hypersurface in $\mathbb{R}^{n}$ inherits a Riemannian metric. This works just as well for a regular submanifold $M \subset \mathbb{R}^{n}$ of arbitrary codimension.

In Riemannian geometry, it is the metric that is the basis for generalizations of length, volume and so on. Motivated by a desire to generalize and to include the mathematics needed for general relativity, we also allow the metric to be indefinite. In this case, some nonzero tangent vectors $v$ might have zero or negative self-scalar product $\langle v, v\rangle$. If $g \in \mathcal{T}_{2}(M)$ is a symmetric tensor field, then we say that it is nondegenerate if $g_{p}$ is nondegenerate on $T_{p} M$ for every $p$. If furthermore $g_{p}$ has the same index for all $p$, then we say it has constant index.

Definition 7.56. If $g \in \mathcal{T}_{2}(M)$ is symmetric nondegenerate and has constant index on $M$, then we call $g$ a semi-Riemannian metric and ( $M, g$ ) a semi-Riemannian manifold or pseudo-Riemannian manifold. The index is called the index of ( $M, g$ ) and denoted $\operatorname{ind}(g)$ or ind $(M)$. The signature is also constant and so the manifold has a signature also. If the signature of a semi-Riemannian manifold (with $\operatorname{dim}(M) \geq 2$ ) is ( $-1,+1,+1,+1, \ldots$ ) (or according to some conventions $(1,-1,-1,-1, \ldots)$ ), then the manifold is called a Lorentz manifold.

The simplest semi-Riemannian manifolds are the spaces $\mathbb{R}_{\nu}^{n}$, which are the spaces $\mathbb{R}^{n}$ endowed with the scalar products given by

$$
\langle x, y\rangle_{\nu}=-\sum_{i=1}^{\nu} x^{i} y^{i}+\sum_{i=\nu+1}^{n} x^{i} y^{i} .
$$

Since ordinary Euclidean geometry does not use indefinite scalar products, we shall call the spaces $\mathbb{R}_{\nu}^{n}$ semi-Euclidean spaces when the index $\nu= \operatorname{ind}(g)$ is not zero. If we write just $\mathbb{R}^{n}$, then either we are not concerned with a scalar product at all, or the scalar product is assumed to be the usual inner product ( $\nu=0$ ). Thus a Riemannian metric is just the special case of index 0 . The space $\mathbb{R}_{1}^{4}$ is called the Minkowski space.

We will usually write $\left\langle X_{p}, Y_{p}\right\rangle$ or $g\left(X_{p}, Y_{p}\right)$ in place of $g(p)\left(X_{p}, X_{p}\right)$. Also, for a pair of vector fields $X$ and $Y$, we define the function $\langle X, Y\rangle$ which is given by $\langle X, Y\rangle(p)=\left\langle X_{p}, Y_{p}\right\rangle$. In local coordinates ( $x^{1}, \ldots, x^{n}$ ) on $U \subset M$,\\
we have that $\left.g\right|_{U}=g_{i j} d x^{i} \otimes d x^{j}$, where $g_{i j}=\left\langle\frac{\partial}{\partial x^{i}}, \frac{\partial}{\partial x^{j}}\right\rangle$. Thus if $X=X^{i} \frac{\partial}{\partial x^{i}}$ and $Y=Y^{i} \frac{\partial}{\partial x^{i}}$ on $U$, then

$$
\langle X, Y\rangle=g_{i j} X^{i} Y^{i}
$$

which is a smooth function defined on $U$. The expression $\langle X, Y\rangle=g_{i j} X^{i} Y^{i}$ means that for all $p \in U$ we have $\langle X(p), Y(p)\rangle=g_{i j}(p) X^{i}(p) Y^{i}(p)$. As we know, the functions $X^{i}$ and $Y^{i}$ are given by $X^{i}=d x^{i}(X)$ and $Y^{i}=d x^{i}(Y)$.

On a semi-Riemannian manifold, the musical isomorphisms are globalized in the obvious way to act on tensor fields. We simply apply the type change at each point in the domain of a given tensor field. For example, if $A$ is a tensor field of type ( 2,2 ), then we may obtain a new metrically equivalent tensor field $A^{\prime}$ of type $(1,3)$ by the rule that for any $p \in M$, we have $A^{\prime}(p)(\alpha, u, v, w):=A(p)(\alpha, b u, v, w)$ for $\alpha \in T_{p}^{*} M$ and $u, v, w \in T_{p} M$. If we choose a chart ( $U, \mathrm{x}$ ), then in terms of the coordinate frames and the corresponding $g_{i j}$, we have

$$
A^{\prime}=A^{i}{ }_{j k l} \frac{\partial}{\partial x^{i}} \otimes d x^{j} \otimes d x^{k} \otimes d x^{l}
$$

where

$$
A_{j k l}^{i}=g_{j a} A_{k l}^{i a}
$$

Of course, each of the possible conventions for consolidation and index position globalize accordingly. The reader is invited to compare our treatment with those found in $[\mathbf{P e}],[\mathbf{D o d}-\mathbf{P o s}],[\mathbf{O N 1}]$ and $[\mathbf{S t e r n}]$.

We have been using coordinate frame fields, but there is nothing preventing us from giving local components of tensor fields with respect to arbitrary smooth frame fields. For example, if we choose a frame field ( $E_{1}, \ldots, E_{n}$ ) and the corresponding dual frame field, then we may define $g_{i j}:=\left\langle E_{i}, E_{j}\right\rangle$ with the corresponding $g^{i j}$, and then local expressions analogous to those above hold with the $\frac{\partial}{\partial x^{i}}$ 's and $d x^{j}$ 's replaced by the $E_{i}$ 's and $E_{j}$ 's, where the components of the tensor are obtained by evaluating on these frame fields. In particular, if the frame field is an orthonormal frame field, then $g_{i j}= \pm 1$ for $i=j$ and $g_{i j}=0$ for $i \neq j$. This can result in a good deal of simplification.

We now say a few words about the appropriate notion of equivalence of semi-Riemannian manifolds.

Definition 7.57. Let ( $M, g$ ) and ( $N, h$ ) be two semi-Riemannian manifolds. A diffeomorphism $\Phi: M \rightarrow N$ is called an isometry if $\Phi^{*} h=g$. Thus for an isometry $\Phi: M \rightarrow N$ we have $g(v, w)=h(T \Phi \cdot v, T \Phi \cdot w)$ for all $v, w \in T M$. If $\Phi: M \rightarrow N$ is a local diffeomorphism such that $\Phi^{*} h=g$, then $\Phi$ is called a local isometry. If there is an isometry $\Phi: M \rightarrow N$, then we say that $(M, g)$ and $(N, h)$ are isometric.

Definition 7.58. The set of all isometries of a semi-Riemannian manifold $M$ to itself is a group called the isometry group. It is denoted by $\operatorname{Isom}(M)$.

The isometry group of a generic manifold is most likely trivial, but examples of manifolds with relatively large isometry groups are easy to find using Lie group theory. Also, Myers and Steenrod showed that the isometry group of a compact Riemannian manifold is a Lie group (see [My-St]). Recall from Chapter 5 that associated to $\mathbb{R}_{\nu}^{n}$ we have the matrix groups $\mathrm{O}(\nu, n-\nu)$ and $\mathrm{SO}(\nu, n-\nu)$. The isometry group of $\mathbb{R}_{\nu}^{n}$ is given by

$$
\begin{aligned}
& \operatorname{Iso}(\nu, n-\nu)=\left\{L: L(x)=Q x+x_{0}\right. \\
& \left.\qquad \text { for some } Q \in \mathrm{O}(\nu, n-\nu) \text { and } x_{0} \in \mathbb{R}_{\nu}^{n}\right\}
\end{aligned}
$$

This is the group of semi-Euclidean motions.\\
Example 7.59. We have seen that a regular submanifold of a Euclidean space $\mathbb{R}^{n}$ is a Riemannian manifold with the metric inherited from $\mathbb{R}^{n}$. In particular, the sphere $S^{n-1} \subset \mathbb{R}^{n}$ is a Riemannian manifold. Every isometry of $S^{n-1}$ is the restriction to $S^{n-1}$ of an isometry of $\mathbb{R}^{n}$ that fixes the origin.

Definition 7.60. Let $\widetilde{M}$ and $M$ be semi-Riemannian manifolds. If $\wp: \widetilde{M} \rightarrow M$ is a covering map such that $\wp$ is a local isometry, we call $\wp: \widetilde{M} \rightarrow M$ a semi-Riemannian covering.

If we have a local isometry $\phi: N \rightarrow M$, then any lift $\phi: N \rightarrow \widetilde{M}$ is also a local isometry (Problem 16). Deck transformations are lifts of the identity map $M \rightarrow M$, and so are diffeomorphisms which are local isometries. Thus deck transformations are, in fact, isometries. We conclude that the group of deck transformations of a semi-Riemannian cover is a subgroup of the group of isometries $\operatorname{Isom}(\widetilde{M})$.

Let us consider here the case of a discrete group $G$ and a discrete group action $\lambda: G \times M \rightarrow M$ that is smooth, proper, and free. We have already seen that the quotient space $M / G$ has a unique structure as a smooth manifold such that the projection $\kappa: M \rightarrow M / G$ is a covering. Let us now assume that $G$ acts by isometries so that $\lambda_{g}^{*}\langle\cdot, \cdot\rangle=\langle\cdot, \cdot\rangle$ for all $g \in G$. The tangent map $T \kappa: T_{p} M \rightarrow T_{\kappa(p)}(M / G)$ is onto. For $x \in M / G$, let $\bar{v}_{1}, \bar{v}_{2} \in T_{x}(M / G)$. Define $h_{x}\left(\bar{v}_{1}, \bar{v}_{2}\right)=\left\langle v_{1}, v_{2}\right\rangle$, where $v_{1}$ and $v_{2}$ are chosen at the same point and such that $T \kappa \cdot v_{i}=\bar{v}_{i}$. We wish to show that this is well-defined. Indeed, if $v_{i} \in T_{p} M$ and $w_{i} \in T_{q} M$ are such that $T_{p} \kappa \cdot v_{i}=T_{q} \kappa \cdot w_{i}=\bar{v}_{i}$ for $i=1,2$, then there is an isometry $\lambda_{g}$ with $\lambda_{g} p=q$. Furthermore, since $\lambda_{g}$ is a deck transformation and curves representing $v_{i}$ and $w_{i}$ must be related by this deck transformation, we also have $T_{p} \lambda_{g} \cdot v_{i}=w_{i}$. Thus

$$
\left\langle v_{1}, v_{2}\right\rangle=\left\langle T_{p} \lambda_{g} v_{1}, T_{p} \lambda_{g} v_{2}\right\rangle=\left\langle w_{1}, w_{2}\right\rangle
$$

which means $h_{x}$ is well-defined. It is easy to show that $x \mapsto h_{x}$ is smooth and defines a metric on $M / G$ with the same signature as that of $\langle\cdot, \cdot\rangle$ and that further, $\kappa^{*} h=\langle\cdot, \cdot\rangle$. In fact, we will use the same notation for either the metric on $M / G$ or on $M$.

Definition 7.61. A lattice of rank $k$ in $\mathbb{R}^{n}$ is a set of the form

$$
\Gamma:=\left\{x \in \mathbb{R}^{n}: x=\sum n_{i} f_{i} \text { where } n_{i} \in \mathbb{Z}\right\}
$$

where $f_{1}, \ldots, f_{k}$ are linearly independent elements of $\mathbb{R}^{n}$. The $f_{1}, \ldots, f_{k}$ are called the generators of the lattice.

The lattice $\mathbb{Z}^{n} \subset \mathbb{R}^{n}$ is the standard rank $n$ lattice, and it is generated by the standard basis. A lattice is a subgroup of $\mathbb{R}^{n}$ and so acts on $\mathbb{R}^{n}$ by $a \mapsto a+v$ for $v \in \Gamma$. This is a discrete, free and proper action, and so the quotient $\mathbb{R}^{n} / \Gamma$ provides a simple example of the above construction and so has a metric induced from $\mathbb{R}^{n}$. If the lattice has full rank $n$, then $\mathbb{R}^{n} / \Gamma$ is called a flat torus (or flat $n$-torus) and is diffeomorphic to the product of $n$ copies of the circle $S^{1}$. Each of these $n$-dimensional flat tori is locally isometric, but may not be globally isometric. To be more precise, suppose that $f_{1}, f_{2}, \ldots, f_{n}$ is a basis for $\mathbb{R}^{n}$ which is not necessarily orthonormal. Let $\Gamma_{f}$ be the lattice consisting of integer linear combinations of $f_{1}, f_{2}, \ldots, f_{n}$. Now suppose we have two such lattices $\Gamma_{f}$ and $\Gamma_{\bar{f}}$. When is $\mathbb{R}^{n} / \Gamma_{f}$ isometric to $\mathbb{R}^{n} / \Gamma_{\bar{f}}$ ? It may seem that, since these are clearly diffeomorphic and since they are locally isometric, they must be (globally) isometric. But this is not the case (see Problem 17). The study of the global geometry of flat tori is quite interesting and even has deep connections with fields outside of geometry such as arithmetic, which we shall not have the space to pursue.

We know from Chapter 4 that there are surfaces in $\mathbb{R}^{3}$ that are diffeomorphic to a torus $S^{1} \times S^{1}$. Such surfaces inherit a metric from the ambient space, but it turns out that the Riemannian surface obtained in this way cannot be isometric to one of the flat 2 -tori introduced here. In Chapter 13 we will see how each metric on a manifold gives rise to an associated curvature tensor. The reason the tori just introduced are referred to as flat is because (being locally isometric to some $\mathbb{R}^{n}$ ) they have vanishing curvature tensor.

If we have semi-Riemannian manifolds ( $M, g$ ) and ( $N, h$ ), then we can consider the product manifold $M \times N$ and the projections $\operatorname{pr}_{1}: M \times N \rightarrow M$ and $\mathrm{pr}_{2}: M \times N \rightarrow N$. The tensor $g \times h=\mathrm{pr}_{1}^{*} g+\mathrm{pr}_{2}^{*} h$ provides a semiRiemannian metric on the manifold $M \times N$, which is then called the semiRiemannian product of ( $M, g$ ) and ( $N, h$ ). Let ( $U_{1}, \mathrm{x}$ ) and ( $U_{2}, \mathrm{y}$ ) denote charts on $M$ and $N$ respectively. Then we may form a product chart for $M \times N$ defined on $U_{1} \times U_{2}$. The coordinate functions of this chart are given by $\widetilde{x}^{i}=x^{i} \circ \operatorname{pr}_{1}$ and $\widetilde{y}^{i}=y^{i} \circ \operatorname{pr}_{2}$. We have the associated frame fields $\frac{\partial}{\partial \widetilde{x}^{i}}$\\
and $\frac{\partial}{\partial \widetilde{y}^{2}}$. The components of $g \times h=\operatorname{pr}_{1}^{*} g+\operatorname{pr}_{2}^{*} h$ in these coordinates are discovered by choosing a point $\left(p_{1}, p_{2}\right) \in U_{1} \times U_{2}$ and then calculating. We have

$$
\begin{aligned}
g & \times h\left(\left.\frac{\partial}{\partial \widetilde{x}^{i}}\right|_{(p, q)},\left.\frac{\partial}{\partial \widetilde{y}^{j}}\right|_{(p, q)}\right) \\
& =\operatorname{pr}_{1}^{*} g\left(\left.\frac{\partial}{\partial \widetilde{x}^{i}}\right|_{(p, q)},\left.\frac{\partial}{\partial \widetilde{y}^{j}}\right|_{(p, q)}\right)+\operatorname{pr}_{2}^{*} h\left(\left.\frac{\partial}{\partial \widetilde{x}^{i}}\right|_{(p, q)},\left.\frac{\partial}{\partial \widetilde{y}^{j}}\right|_{(p, q)}\right) \\
& =g\left(\left.T \operatorname{pr}_{1} \frac{\partial}{\partial \widetilde{x}^{i}}\right|_{(p, q)},\left.T \operatorname{pr}_{1} \frac{\partial}{\partial \widetilde{y}^{j}}\right|_{(p, q)}\right)+h\left(\left.T \operatorname{pr}_{2} \frac{\partial}{\partial \widetilde{x}^{i}}\right|_{(p, q)},\left.T \operatorname{pr}_{2} \frac{\partial}{\partial \widetilde{y}^{j}}\right|_{(p, q)}\right) \\
& =g\left(\left.\frac{\partial}{\partial x^{i}}\right|_{p}, 0_{p}\right)+h\left(0_{q},\left.\frac{\partial}{\partial y^{j}}\right|_{q}\right) \\
& =0+0=0
\end{aligned}
$$

and (abbreviating a bit)

$$
g \times h\left(\left.\frac{\partial}{\partial \widetilde{x}^{i}}\right|_{(p, q)},\left.\frac{\partial}{\partial \widetilde{x}^{j}}\right|_{(p, q)}\right)=g\left(\left.\frac{\partial}{\partial x^{i}}\right|_{p},\left.\frac{\partial}{\partial x^{i}}\right|_{p}\right)+h\left(0_{q}, 0_{q}\right)=g_{i j}(p) .
$$

Similarly $g \times h\left(\frac{\partial}{\partial \widetilde{y}^{i}}, \frac{\partial}{\partial \widetilde{y}^{j}}\right)(p, q)=h_{i j}(q)$. In practice, the coordinate functions constructed above are often abusively denoted by $\left(x^{1}, \ldots, x^{n_{1}}, y^{1}, \ldots, y^{n_{2}}\right)$ and the frame field,s by $\frac{\partial}{\partial x^{1}}, \ldots, \frac{\partial}{\partial x^{n} 1}, \frac{\partial}{\partial y^{1}}, \ldots, \frac{\partial}{\partial y^{n} 2}$. So with respect to these coordinates, the matrix of $g \times h$ is of the form

$$
\left(\begin{array}{cc}
G & 0 \\
0 & H
\end{array}\right),
$$

where $G=\left(g_{i j} \circ \mathrm{pr}_{1}\right)$ and $H=\left(h_{i j} \circ \mathrm{pr}_{2}\right)$.\\
Notation 7.62. The product metric is often denoted by $g+h$ or in coordinates by $d s^{2}=g_{i j} d x^{i} d x^{j}+h_{k l} d y^{k} d y^{l}$.

Every smooth manifold that admits partitions of unity also admits at least one (in fact infinitely many) Riemannian metric. This includes all (finite-dimensional) paracompact manifolds. The reason for this is that the set of all Riemannian metric tensors is, in an appropriate sense, convex. We record this as a proposition.

Proposition 7.63. Every smooth (paracompact) manifold admits a Riemannian metric.

Proof. This is a special case of Proposition 6.45.\\
If $M$ is a regular submanifold of a Riemannian manifold ( $N, h$ ), then $M$ inherits a Riemannian metric $g:=\imath^{*} h$, where $\imath: M \hookrightarrow N$ is the inclusion map. We have already used this idea for submanifolds of the Euclidean\\
space $\mathbb{R}^{d}$. More generally, if $f: M \rightarrow N$ is an immersion, then ( $M, f^{*} h$ ) is a Riemannian manifold. In particular, if $f: M \rightarrow \mathbb{R}^{d}$ is an immersion, then we obtain a Riemannian metric on $M$. It turns out that every Riemannian metric on $M$ can be obtained in this way. Actually, more is true! For any Riemannian manifold ( $M, g$ ) there is an embedding $f: M \rightarrow \mathbb{R}^{d}$, for some $d$, such that $g:=\imath^{*} g_{0}$, where $g_{0}$ denotes the standard metric on $\mathbb{R}^{d}$. What this means is that $f(M)$ is a regular submanifold, and if we give $f(M)$ the metric induced from the ambient space $\mathbb{R}^{d}$, then $f$ becomes an isometry when viewed as a map into $f(M)$. We say that such an $f$ is an isometric embedding of $M$ into $\mathbb{R}^{d}$. In short, the result is that every Riemannian manifold can be isometrically embedded into some Euclidean space of sufficiently high dimension. This difficult theorem is called the Nash embedding theorem and is due to John Forbes Nash (see [Nash1] and [Nash2]). Note that $d$ must be quite large in general $\left(d=(\operatorname{dim} M)^{2}+\right. 5(\operatorname{dim} M)+3$ is sufficient). For an indefinite semi-Riemannian manifold ( $N, h$ ), the pull-back $f^{*} h$ by a smooth map $f: M \rightarrow N$ may not be a metric because there may be points $p \in M$ such that $T_{p} f\left(T_{p} M\right)$ is a degenerate subspace of $T_{f(p)} N$. In particular, not every embedding of a manifold $M$ into a semi-Euclidean space $\mathbb{R}_{\nu}^{d}$ (with $1<\nu<d$ ) induces a metric on $M$. Nevertheless, every metric on $M$ of any index can be obtained using an appropriate embedding into some $\mathbb{R}_{\nu}^{d}$ (see [Clark]).

\section*{Problems}
(1) Show that if $\tau \in V \otimes V^{*}$ has the same components $\tau_{j}^{i}$ with respect to every basis, then $\tau_{j}^{i}=a \delta_{j}^{i}$ for some $a \in \mathbb{R}$.\\
(2) If $e_{1}, \ldots, e_{n}$ is a basis for V and $f_{1}, \ldots, f_{m}$ is a basis for W , then $\left\{E_{j}^{i}\right\}_{i=1, \ldots, n} \quad j=1, \ldots, m$ is a basis for $L(\mathrm{V}, \mathrm{~W})$, where $E_{j}^{i}(v):=e^{i}(v) f_{j}$. Show this directly without assuming the isomorphism of $\mathrm{W} \otimes \mathrm{V}^{*}$ with $L(\mathrm{V}, \mathrm{~W})$.\\
(3) Let $\left(\mathrm{V}_{i}, g_{i}\right)$ be scalar product spaces for $i=1, \ldots, k$. By Corollary D. 35 of Appendix D, there is a unique bilinear form

$$
\varphi: \bigotimes_{i=1}^{k} \mathrm{V}_{i} \times \bigotimes_{i=1}^{k} \mathrm{V}_{i} \rightarrow \mathbb{R}
$$

such that for $v_{i} \in \mathrm{V}_{i}$ and $w_{i} \in \mathrm{~W}_{i}$,

$$
\varphi\left(v_{1} \otimes \cdots \otimes v_{k}, w_{1} \otimes \cdots \otimes w_{k}\right)=\varphi_{1}\left(v_{1}, w_{1}\right) \cdots \varphi_{k}\left(v_{k}, w_{k}\right)
$$

Show that $\varphi$ is nondegenerate and that it is positive definite if each $g_{i}$ is positive definite.\\
(4) Define $\tau: \mathfrak{X}(M) \times \mathfrak{X}(M) \rightarrow C^{\infty}(M)$ by $\tau(X, Y)=X Y f$. Show that $\tau$ does not define a tensor field.\\
(5) Let $b_{i j}$ and $b_{i j}^{\prime}$ be the components of a bilinear form $b$ with respect to bases $e_{1}, \ldots, e_{n}$ and $e_{1}^{\prime}, \ldots, e_{n}^{\prime}$ respectively. Show that in general $\operatorname{det}\left(b_{i j}\right)$ does not equal $\operatorname{det}\left(b_{i j}^{\prime}\right)$. Show that if $\operatorname{det}\left(b_{i j}\right)$ is nonzero, then the same is true of $\operatorname{det}\left(b_{i j}^{\prime}\right)$.\\
(6) Show that while a single algebraic tensor $\tau_{p}$ at a point on a manifold can always be extended to a smooth tensor field, it is not the case that one may always extend a (smooth) tensor field defined on an open subset to a smooth tensor field on the whole manifold.\\
(7) Let $\phi: \mathbb{R}^{2} \rightarrow \mathbb{R}^{2}$ be defined by $(x, y) \mapsto(x+2 y, y)$. Let $\tau:=x \frac{\partial}{\partial x} \otimes d y+ \frac{\partial}{\partial y} \otimes d y$. Compute $\phi_{*} \tau$ and $\phi^{*} \tau$.\\
(8) Prove Proposition 7.35.\\
(9) Let $\mathcal{D}$ be a tensor derivation on $M$ and suppose that in a local chart we have $\mathcal{D}\left(\frac{\partial}{\partial x^{i}}\right)=\sum D_{i}^{j} \frac{\partial}{\partial x^{j}}$ for smooth functions $D_{i}^{j}$. Show that $\mathcal{D}\left(d x^{j}\right)= -\sum D_{i}^{j} d x^{i}$. Let $X$ be a fixed vector field with components $X^{i}$ in our chart. Find the $D_{i}^{j}$ in the case that $\mathcal{D}=\mathcal{L}_{X}$.\\
(10) Let $A \in \mathcal{T}_{0}^{2}(M)$. Show that the component form of the Lie derivative with respect to a chart is given as

$$
\left(\mathcal{L}_{X} A\right)^{a b}=\frac{\partial A^{a b}}{\partial x^{h}} X^{h}-\frac{\partial X^{a}}{\partial x^{h}} A^{h b}-\frac{\partial X^{b}}{\partial x^{h}} A^{a h}
$$

(where we use the Einstein summation convention). Show that if $A \in \mathcal{T}_{2}^{0}(M)$, then the formula becomes

$$
\left(\mathcal{L}_{X} A\right)=\frac{\partial A_{a b}}{\partial x^{h}} X^{h}+\frac{\partial X^{h}}{\partial x^{a}} A_{h b}+\frac{\partial X^{h}}{\partial x^{b}} A_{a h} .
$$

Find a formula for $A \in \mathcal{T}_{2}^{2}(M)$.\\
(11) Show that our two definitions of the Lie derivative of a tensor field agree with each other.\\
(12) In some chart $(U,(x, y))$ on a 2 -manifold, let $A=x \frac{\partial}{\partial y} \otimes d x \otimes d y+\frac{\partial}{\partial x} \otimes d y \otimes d y$ and let $X=\frac{\partial}{\partial x}+x \frac{\partial}{\partial y}$. Compute the coordinate expression for $\mathcal{L}_{X} A$.\\
(13) Suppose that for every chart ( $U, \mathrm{x}$ ) in an atlas for a smooth $n$-manifold $M$ we have assigned $n^{3}$ smooth functions $\Gamma_{i j}^{k}$, which we call Christoffel symbols. Suppose that rather than obeying the transformation law expected for a tensor, we have the following horrible formula relating the Christoffel symbols $\Gamma_{i j}^{\prime k}$ on a chart ( $U^{\prime}, \mathrm{y}$ ) to the symbols $\Gamma_{i j}^{k}$ :

$$
\Gamma_{i j}^{\prime k}=\frac{\partial^{2} x^{l}}{\partial y^{i} \partial y^{j}} \frac{\partial y^{k}}{\partial x^{l}}+\Gamma_{r s}^{t} \frac{\partial x^{r}}{\partial y^{i}} \frac{\partial x^{s}}{\partial y^{j}} \frac{\partial y^{k}}{\partial x^{t}} \text { (sum). }
$$

Assume that such a transformation law holds between the Christoffel symbol functions for all pairs of intersecting charts. For any pair of vector fields $X, Y \in \mathfrak{X}(M)$, consider the functions $\left(D_{X} Y\right)^{k}$ given in every chart by the formula $\left(D_{X} Y\right)^{k}:=\frac{\partial Y^{k}}{\partial x^{h}} X^{h}+\Gamma_{i j}^{k} X^{i} Y^{j}$. Show that the local vector fields of the form $\left(D_{X} Y\right)^{k} \frac{\partial}{\partial x^{k}}$, defined on each chart, are the restrictions of a single global vector field $D_{X} Y$. Show that $D_{X}: Y \mapsto D_{X} Y$ is a derivation of $\mathfrak{X}(M)$ and that with $D_{X} f:=X f$ for smooth functions, we may extend to a tensor derivation; $D_{X}$ is called a covariant derivative with respect to $X$. There are many possible covariant derivatives.\\
(14) Continuing on the last problem, show that $D_{f X+g Y} \Upsilon=f D_{X} \Upsilon+g D_{Y} \Upsilon$ for all $f, g \in C^{\infty}(M)$ and $X, Y \in \mathfrak{X}(M)$ and $\Upsilon \in \mathcal{T}_{s}^{r}(M)$.\\
(15) Show that if $\frac{\partial}{\partial x^{1}}, \ldots, \frac{\partial}{\partial x^{n}}$ are coordinate vector fields from some chart, then $\left[\frac{\partial}{\partial x^{i}}, \frac{\partial}{\partial x^{j}}\right] \equiv 0$. Consider the vector fields $\frac{\partial}{\partial x}$ and $\frac{\partial}{\partial y}$ arising from standard coordinates on $\mathbb{R}^{2}$ and also the $\frac{\partial}{\partial r}$ and $\frac{\partial}{\partial \theta}$ from polar coordinates. Show that $\left[\frac{\partial}{\partial x}, \frac{\partial}{\partial r}\right]$ is not identically zero by explicit computation.\\
(16) Let $\widetilde{M} \rightarrow M$ be a semi-Riemannian cover. Prove that if we have a local isometry $\phi: N \rightarrow M$, then any lift $\phi: N \rightarrow \widetilde{M}$ is also a local isometry.\\
(17) Suppose we have two lattices $\Gamma_{f}$ and $\Gamma_{\bar{f}}$ in $\mathbb{R}^{n}$. Let $\mathbb{R}^{n}$ have the standard metric. Describe the induced metrics on $\mathbb{R}^{n} / \Gamma_{f}$ and $\mathbb{R}^{n} / \Gamma_{\bar{f}}$ and provide a necessary and sufficient condition for the existence of an isometry $\mathbb{R}^{n} / \Gamma_{f} \rightarrow \mathbb{R}^{n} / \Gamma_{\bar{f}}$.\\
(18) Show that if $\alpha, \beta \in T_{k}(\mathrm{V})$ where V is a scalar product space, then the scalar product on $T_{k}(\mathrm{V})$ is given in terms of index raising and contraction by

$$
\langle\alpha, \beta\rangle=\sum \alpha_{i_{1} \ldots i_{k}} \beta^{i_{1} \ldots i_{k}}
$$









\pagebreak




\section*{Chapter 8}
\section*{Differential Forms}
In one guise, a differential form is nothing but an alternating (antisymmetric) tensor field. What is new is the introduction of an antisymmetrized version of the tensor product and also a natural differential operator called the exterior derivative. We start off with some more multilinear algebra.

\subsection*{More Multilinear Algebra}



Definition 8.1. 


Let V and W be real finite-dimensional vector spaces. A $k$ multilinear map $\alpha: \mathrm{V} \times \cdots \times \mathrm{V} \rightarrow \mathrm{W}$ is called alternating if $\alpha\left(v_{1}, \ldots, v_{k}\right)=$ 0 whenever $v_{i}=v_{j}$ for some $i \neq j$. The space of all alternating $k$ multilinear maps into W will be denoted by $L_{\text {alt }}^{k}(\mathrm{V} ; \mathrm{W})$ or by $L_{\text {alt }}^{k}(\mathrm{V})$ if $\mathrm{W}=\mathbb{R}$. By convention, $L_{\text {alt }}^{0}(\mathrm{V} ; \mathrm{W})$ is taken to be W and in particular, $L_{\text {alt }}^{0}(\mathrm{V})=\mathbb{R}$.


Since we are dealing with the field $\mathbb{R}$ (which has characteristic zero), it is easy to see that alternating $k$-multilinear maps are the same as (completely) antisymmetric $k$-multilinear maps which are defined by the property that for any permutation $\sigma$ of the letters $1,2, \ldots, k$ we have

$$
\omega\left(v_{1}, v_{2}, \ldots, v_{k}\right)=\operatorname{sgn}(\sigma) \omega\left(v_{\sigma(1)}, v_{\sigma(2)}, \ldots, v_{\sigma(k)}\right) .
$$

Let us denote the group of permutations of the $k$ letters $1,2, \ldots, k$ by $S_{k}$. In what follows, we will occasionally write $\sigma_{i}$ in place of $\sigma(i)$.




\begin{DEF}{GA-L09-06-01}{Alternierende multilineare Abbildungen (Vektorräume)}
Seien $\mathrm{V}$ und $\mathrm{W}$ endlichdimensionale reelle Vektorräume.  
Eine $k$-lineare Abbildung
\[
\alpha:\;\mathrm{V}\times\cdots\times \mathrm{V}\longrightarrow \mathrm{W}
\]
heißt \textbf{alternierend}, falls
\[
\alpha(v_{1},\ldots,v_{k})=0
\quad\text{immer dann, wenn }v_{i}=v_{j}\text{ für ein Paar }i\neq j.
\]
Der reelle Vektorraum aller solcher alternierenden $k$-multilinearen Abbildungen wird bezeichnet durch
\[
L_{\mathrm{alt}}^{k}(\mathrm{V};\mathrm{W})
\quad\text{bzw.}\quad
L_{\mathrm{alt}}^{k}(\mathrm{V}) := L_{\mathrm{alt}}^{k}(\mathrm{V};\mathbb{R})
\text{ im Falle }\mathrm{W}=\mathbb{R}.
\]
Per Konvention gilt $L_{\mathrm{alt}}^{0}(\mathrm{V};\mathrm{W})=\mathrm{W}$, also insbesondere $L_{\mathrm{alt}}^{0}(\mathrm{V})=\mathbb{R}$.

\medskip
Da der zugrundeliegende Körper $\mathbb{R}$ charakteristisch $0$ besitzt, ist eine $k$-multilineare Abbildung genau dann alternierend, wenn sie \textbf{antisymmetrisch} ist, d.\,h.\ wenn für jede Permutation $\sigma\in S_{k}$ der Symmetriegruppe
\[
\boxed{
\omega(v_{1},v_{2},\ldots,v_{k})
=\operatorname{sgn}(\sigma)\,
\omega(v_{\sigma(1)},v_{\sigma(2)},\ldots,v_{\sigma(k)})
}
\]
gilt. Somit bilden alternierende und antisymmetrische Multilinearformen denselben Raum.
\end{DEF}





Definition 8.2. 

\begin{DEF}{GA-L09-06-02}{Antisymmetrisierungsabbildung für kovariante Tensoren}
The antisymmetrization map $\mathrm{Alt}^{k}: T_{k}^{0}(\mathrm{V}) \rightarrow L_{\text {alt }}^{k}(\mathrm{V})$ is defined by
$$
\operatorname{Alt}^{k}(\omega)\left(v_{1}, v_{2}, \ldots, v_{k}\right):=\frac{1}{k!} \sum_{\sigma \in S_{k}} \operatorname{sgn}(\sigma) \omega\left(v_{\sigma_{1}}, v_{\sigma_{2}}, \ldots, v_{\sigma_{k}}\right) .
$$
\end{DEF}

Lemma 8.3. 

\begin{LEM}{GA-L09-06-03}{Antisymmetrisierung und Tensorprodukt}
For $\alpha \in T_{k_{1}}^{0}(\mathrm{V})$ and $\beta \in T_{k_{2}}^{0}(\mathrm{V})$, we have
$$
\begin{aligned}
& \mathrm{Alt}^{k_{1}+k_{2}}\left(\mathrm{Alt}^{k_{1}} \alpha \otimes \beta\right)=\mathrm{Alt}^{k_{1}+k_{2}}(\alpha \otimes \beta), \\
& \mathrm{Alt}^{k_{1}+k_{2}}\left(\alpha \otimes \mathrm{Alt}^{k_{2}} \beta\right)=\mathrm{Alt}^{k_{1}+k_{2}}(\alpha \otimes \beta),
\end{aligned}
$$
and
$$
\operatorname{Alt}^{k_{1}+k_{2}}\left(\operatorname{Alt}^{k_{1}}(\alpha) \otimes \operatorname{Alt}^{k_{2}}(\beta)\right)=\operatorname{Alt}^{k_{1}+k_{2}}(\alpha \otimes \beta) .
$$
\end{LEM}

Proof. 

\begin{PROOF}{GA-L09-06-04}{P: Antisymmetrisierung und Tensorprodukt}
For a permutation $\sigma \in S_{k}$ and any $T \in T_{k}^{0}(\mathrm{V})$, let $\sigma T$ denote the element of $T_{k}^{0}(\mathrm{V})$ given by $(\sigma T)\left(v_{1}, \ldots, v_{k}\right):=T\left(v_{\sigma(1)}, \ldots, v_{\sigma(k)}\right)$. We then have $\operatorname{Alt}^{k}(\sigma T)=\operatorname{sgn}(\sigma) \operatorname{Alt}^{k}(T)$ as may easily be checked. Also, by definition Alt ${ }^{k}(T)=\sum \operatorname{sgn}(\sigma) \sigma T$. We have
$$
\begin{aligned}
\operatorname{Alt}^{k_{1}+k_{2}} & \left(\operatorname{Alt}^{k_{1}}(\alpha) \otimes \beta\right) \\
& =\operatorname{Alt}^{k_{1}+k_{2}}\left(\left(\frac{1}{k_{1}!} \sum_{\sigma \in S_{k_{1}}} \operatorname{sgn} \sigma(\sigma \alpha)\right) \otimes \beta\right) \\
& =\operatorname{Alt}^{k_{1}+k_{2}}\left(\frac{1}{k_{1}!} \sum_{\sigma \in S_{k_{1}}} \operatorname{sgn} \sigma(\sigma \alpha \otimes \beta)\right) \\
& =\frac{1}{k_{1}!} \sum_{\sigma \in S_{k_{1}}} \operatorname{sgn} \sigma \operatorname{Alt}^{k_{1}+k_{2}}(\sigma \alpha \otimes \beta) .
\end{aligned}
$$
Let us examine the expression $\operatorname{sgn} \sigma \operatorname{Alt}^{k_{1}+k_{2}}(\sigma \alpha \otimes \beta)$. If we extend each $\sigma \in S_{k_{1}}$ to a corresponding element $\sigma^{\prime} \in S_{k_{1}+k_{2}}$ by letting $\sigma^{\prime}(i)=\sigma(i)$ for $i \leq k_{1}$ and $\sigma^{\prime}(i)=i$ for $i>k_{1}$, then we have $\sigma \alpha \otimes \beta=\sigma^{\prime}(\alpha \otimes \beta)$ and also $\operatorname{sgn}(\sigma)=\operatorname{sgn}\left(\sigma^{\prime}\right)$. Thus $\operatorname{sgn} \sigma$ Alt $^{k_{1}+k_{2}}(\sigma \alpha \otimes \beta)=\operatorname{sgn} \sigma^{\prime}$ Alt $^{k_{1}+k_{2}} \sigma^{\prime}(\alpha \otimes \beta)$ and so
$$
\begin{aligned}
\operatorname{Alt}^{k_{1}+k_{2}} & \left(\operatorname{Alt}^{k_{1}} \alpha \otimes \beta\right) \\
& =\frac{1}{k_{1}!} \sum_{\sigma \in S_{k_{1}}} \operatorname{sgn} \sigma^{\prime} \operatorname{Alt}^{k_{1}+k_{2}}\left(\sigma^{\prime}(\alpha \otimes \beta)\right) \\
& =\frac{1}{k_{1}!} \sum_{\sigma \in S_{k_{1}}} \operatorname{sgn} \sigma^{\prime} \operatorname{sgn} \sigma^{\prime} \operatorname{Alt}^{k_{1}+k_{2}}(\alpha \otimes \beta) \\
& =\operatorname{Alt}^{k_{1}+k_{2}}(\alpha \otimes \beta) \frac{1}{k_{1}!} \sum_{\sigma \in S_{k_{1}}} 1=\operatorname{Alt}^{k_{1}+k_{2}}(\alpha \otimes \beta) .
\end{aligned}
$$
We arrive at $\mathrm{Alt}^{k_{1}+k_{2}}\left(\mathrm{Alt}^{k_{1}}(\alpha) \otimes \beta\right)=\mathrm{Alt}^{k_{1}+k_{2}}(\alpha \otimes \beta)$. In a similar way, $\operatorname{Alt}^{k_{1}+k_{2}}\left(\alpha \otimes \operatorname{Alt}^{k_{2}} \beta\right)=\operatorname{Alt}^{k_{1}+k_{2}}(\alpha \otimes \beta)$, and so the last part of the theorem follows.
\end{PROOF}




\begin{DEF}{GA-L09-06-05}{Äußeres / Keil Produkt }
Given $\omega \in L_{\text {alt }}^{k_{1}}(\mathrm{V})$ and $\eta \in L_{\text {alt }}^{k_{2}}(\mathrm{V})$, we define their exterior product or wedge product $\omega \wedge \eta \in L_{\text {alt }}^{k_{1}+k_{2}}(\mathrm{V})$ by the formula
$$
\omega \wedge \eta:=\frac{\left(k_{1}+k_{2}\right)!}{k_{1}!k_{2}!} \operatorname{Alt}^{k_{1}+k_{2}}(\omega \otimes \eta)
$$
Written out, this is
$$
\begin{aligned}
& \omega \wedge \eta\left(\mathrm{v}_{1}, \ldots, \mathrm{v}_{k_{1}}, \mathrm{v}_{k_{1}+1}, \ldots, \mathrm{v}_{k_{1}+k_{2}}\right) \\
& \quad=\frac{1}{k_{1}!k_{2}!} \sum_{\sigma \in S_{k_{1}+k_{2}}} \operatorname{sgn}(\sigma) \omega\left(\mathrm{v}_{\sigma_{1}}, \ldots, \mathrm{v}_{\sigma_{k_{1}}}\right) \eta\left(\mathrm{v}_{\sigma_{k_{1}+1}}, \ldots, \mathrm{v}_{\sigma_{k_{1}+k_{2}}}\right)
\end{aligned}
$$
$$
\begin{aligned}
& (\omega \wedge \eta)\left(\mathrm{v}_{1}, \ldots, \mathrm{v}_{k_{1}}, \mathrm{v}_{k_{1}+1}, \ldots, \mathrm{v}_{k_{1}+k_{2}}\right) \\
& \quad=\sum_{\left(k_{1}, k_{2}\right) \text {-shuffles } \sigma} \operatorname{sgn}(\sigma) \omega\left(\mathrm{v}_{\sigma_{1}}, \ldots, \mathrm{v}_{\sigma_{k_{1}}}\right) \eta\left(\mathrm{v}_{\sigma_{k_{1}+1}}, \ldots, \mathrm{v}_{\sigma_{k_{1}+k_{2}}}\right) .
\end{aligned}
$$
In the latter formula, we sum over all permutations such that $\sigma(1)<\sigma(2)< \cdots<\sigma\left(k_{1}\right)$ and $\sigma\left(k_{1}+1\right)<\sigma\left(k_{1}+2\right)<\cdots<\sigma\left(k_{1}+k_{2}\right)$. This kind of permutation is called a ( $k_{1}, k_{2}$ )-shuffle as indicated in the summation.

The most important case of (8.1) is for $\omega, \eta \in L_{\text {alt }}^{1}(\mathrm{V})$, in which case
$$
(\omega \wedge \eta)(v, w)=\omega(v) \eta(w)-\omega(w) \eta(v)
$$
\end{DEF}


Warning: The factor in front of Alt in the definition of the exterior product is a convention but not the only convention in use. This choice has an effect on many of the formulas to follow which differ by a factor from the corresponding formulas written by authors following other conventions.

It is an exercise in combinatorics that we also have
$$
\begin{aligned}
& (\omega \wedge \eta)\left(\mathrm{v}_{1}, \ldots, \mathrm{v}_{k_{1}}, \mathrm{v}_{k_{1}+1}, \ldots, \mathrm{v}_{k_{1}+k_{2}}\right) \\
& \quad=\sum_{\left(k_{1}, k_{2}\right) \text {-shuffles } \sigma} \operatorname{sgn}(\sigma) \omega\left(\mathrm{v}_{\sigma_{1}}, \ldots, \mathrm{v}_{\sigma_{k_{1}}}\right) \eta\left(\mathrm{v}_{\sigma_{k_{1}+1}}, \ldots, \mathrm{v}_{\sigma_{k_{1}+k_{2}}}\right) .
\end{aligned}
$$

In the latter formula, we sum over all permutations such that $\sigma(1)<\sigma(2)< \cdots<\sigma\left(k_{1}\right)$ and $\sigma\left(k_{1}+1\right)<\sigma\left(k_{1}+2\right)<\cdots<\sigma\left(k_{1}+k_{2}\right)$. This kind of permutation is called a ( $k_{1}, k_{2}$ )-shuffle as indicated in the summation.

The most important case of (8.1) is for $\omega, \eta \in L_{\text {alt }}^{1}(\mathrm{V})$, in which case
$$
(\omega \wedge \eta)(v, w)=\omega(v) \eta(w)-\omega(w) \eta(v)
$$

This clearly defines an antisymmetric bilinear map.


Proposition 8.4. 

\begin{PROP}{GA-L09-06-06}{Keil Produkt als antisymmetrische bilineare Abbildung}
For $\alpha \in L_{\text {alt }}^{k_{1}}(\mathrm{V}), \beta \in L_{\text {alt }}^{k_{2}}(\mathrm{V})$, and $\gamma \in L_{\text {alt }}^{k_{3}}(\mathrm{V})$, we have
\begin{enumerate}
\item $\wedge: L_{\text {alt }}^{k_{1}}(\mathrm{V}) \times L_{\text {alt }}^{k_{2}}(\mathrm{V}) \rightarrow L_{\text {alt }}^{k_{1}+k_{2}}(\mathrm{V})$ is $\mathbb{R}$-bilinear
\item $\alpha \wedge \beta=(-1)^{k_{1} k_{2}} \beta \wedge \alpha$
\item $\alpha \wedge(\beta \wedge \gamma)=(\alpha \wedge \beta) \wedge \gamma$
\end{enumerate}
\end{PROP}



Proof. 

\begin{PROOF}{GA-L09-06-07}{P: Keil Produkt als antisymmetrische bilineare Abbildung}
We leave the proof of (i) as an easy exercise. For (ii), we consider the special permutation $f$ given by $\left(f(1), f(2), \ldots, f\left(k_{1}+k_{2}\right)\right)=\left(k_{1}+1, \ldots\right.$, $\left.k_{1}+k_{2}, 1, \ldots, k_{1}\right)$. We have that $\alpha \otimes \beta=f(\beta \otimes \alpha)$. Also $\operatorname{sgn}(f)=(-1)^{k_{1} k_{2}}$. So we have
$$
\operatorname{Alt}^{k_{1}+k_{2}}(\alpha \otimes \beta)=\operatorname{Alt}^{k_{1}+k_{2}}(f(\beta \otimes \alpha))=(-1)^{k_{1} k_{2}} \operatorname{Alt}^{k_{1}+k_{2}}(\beta \otimes \alpha),
$$
which gives (ii).

For (iii), we compute
$$
\begin{aligned}
\alpha \wedge(\beta \wedge \gamma) & =\frac{\left(k_{1}+k_{2}+k_{3}\right)!}{k_{1}!\left(k_{2}+k_{3}\right)!} \operatorname{Alt}(\alpha \otimes(\beta \wedge \gamma)) \\
& =\frac{\left(k_{1}+k_{2}+k_{3}\right)!}{k_{1}!\left(k_{2}+k_{3}\right)!} \frac{\left(k_{2}+k_{3}\right)!}{k_{2}!k_{3}!} \operatorname{Alt}(\alpha \otimes \operatorname{Alt}(\beta \otimes \gamma)) \\
& =\frac{\left(k_{1}+k_{2}+k_{3}\right)!}{k_{1}!k_{2}!k_{3}!} \operatorname{Alt}(\alpha \otimes \operatorname{Alt}(\beta \otimes \gamma))
\end{aligned}
$$
By Lemma 8.3, we know that $\operatorname{Alt}(\alpha \otimes \operatorname{Alt}(\beta \otimes \gamma))=\operatorname{Alt}(\alpha \otimes(\beta \otimes \gamma))$, and so we arrive at
$$
\alpha \wedge(\beta \wedge \gamma)=\frac{\left(k_{1}+k_{2}+k_{3}\right)!}{k_{1}!k_{2}!k_{3}!} \operatorname{Alt}(\alpha \otimes(\beta \otimes \gamma))
$$
By a symmetric computation, we also have
$$
(\alpha \wedge \beta) \wedge \gamma=\frac{\left(k_{1}+k_{2}+k_{3}\right)!}{k_{1}!k_{2}!k_{3}!} \operatorname{Alt}((\alpha \otimes \beta) \otimes \gamma)
$$
and so by the associativity of the tensor product we obtain the result.
\end{PROOF}


Example 8.5. 


\begin{EXA}{GA-L09-06-08}{Beispiel für Keil-Produkt}
Let V have a basis $e_{1}, e_{2}, e_{3}$ with dual basis $e^{1}, e^{2}, e^{3}$. Let $\alpha=2 e^{1} \wedge e^{2}+e^{1} \wedge e^{2}$ and $\beta=e^{1}-e^{3}$. Then as a sample calculation we have
$$
\begin{aligned}
\alpha \wedge \beta & =\left(2 e^{1} \wedge e^{2}+e^{1} \wedge e^{3}\right) \wedge\left(e^{1}-e^{3}\right) \\
& =2 e^{1} \wedge e^{2} \wedge e^{1}+e^{1} \wedge e^{3} \wedge e^{1} \\
& =-2 e^{1} \wedge e^{2} \wedge e^{3}-e^{1} \wedge e^{3} \wedge e^{3} \\
& =-2 e^{1} \wedge e^{2} \wedge e^{3}
\end{aligned}
$$
where we have used that $e^{1} \wedge e^{2} \wedge e^{1}=-e^{1} \wedge e^{1} \wedge e^{2}=0$, etc.
\end{EXA}


Lemma 8.6. 


\begin{LEM}{GA-L09-06-09}{Determinante als Auswertung von Keilprodukten}
Let $\alpha^{1}, \ldots, \alpha^{k}$ be elements of $\mathrm{V}^{*}=L_{\text {alt }}^{1}(\mathrm{V})$ and let $v_{1}, \ldots, v_{k} \in \mathrm{V}$. Then we have
$$
\alpha^{1} \wedge \cdots \wedge \alpha^{k}\left(v_{1}, \ldots, v_{k}\right)=\operatorname{det} A
$$
where $A=\left(a_{j}^{i}\right)$ is the $k \times k$ matrix whose $i j$-th entry is $a_{j}^{i}=\alpha^{i}\left(v_{j}\right)$.
\end{LEM}


Proof. 

\begin{PROOF}{GA-L09-06-10}{P: Determinante als Auswertung von Keilprodukten}
From the proof of the last theorem we have
$$
\alpha \wedge(\beta \wedge \gamma)=\frac{\left(k_{1}+k_{2}+k_{3}\right)!}{k_{1}!k_{2}!k_{3}!} \operatorname{Alt}(\alpha \otimes(\beta \otimes \gamma))
$$
By inductive application of this we have
$$
\alpha^{1} \wedge \cdots \wedge \alpha^{k}=k!\operatorname{Alt}\left(\alpha^{1} \otimes \cdots \otimes \alpha^{k}\right)
$$
Thus
$$
\alpha^{1} \wedge \cdots \wedge \alpha^{k}\left(v_{1}, \ldots, v_{k}\right)=\sum_{\sigma} \operatorname{sgn}(\sigma) \alpha^{1}\left(v_{\sigma_{1}}\right) \cdots \alpha^{k}\left(v_{\sigma_{k}}\right)=\operatorname{det} A
$$
\end{PROOF}


\begin{DEF}{GA-L09-06-11}{Levi-Civita Symbole}
Let us define
$$
\epsilon_{j_{1} \ldots j_{k}}^{i_{1} \ldots i_{k}}=\left\{\begin{array}{cc}
1 & \text { if } j_{1}, \ldots, j_{k} \text { is an even permutation of } i_{1}, \ldots, i_{k} \\
-1 & \text { if } j_{1}, \ldots, j_{k} \text { is an odd permutation of } i_{1}, \ldots, i_{k} \\
0 & \text { otherwise. }
\end{array}\right.
$$
\end{DEF}


Then we have\\


Corollary 8.7. 


\begin{KORO}{GA-L09-06-12}{Levi-Civita Symbole als Auswertungen von Keilprodukten}
Let $e_{1}, \ldots, e_{n}$ be a basis for V and $e^{1}, e^{2}, \ldots, e^{n}$ the dual basis for $\mathrm{V}^{*}$. Then we have
$$
e^{i_{1}} \wedge \cdots \wedge e^{i_{k}}\left(e_{j_{1}}, \ldots, e_{j_{k}}\right)=\epsilon_{j_{1} \ldots j_{k}}^{i_{1} \ldots i_{k}}
$$
\end{KORO}

Theorem 8.8. 

\begin{THEO}{GA-L09-06-13}{Basis von $L_{\text {alt }}^{k}(\mathrm{V})$}
If $\left(e^{1}, e^{2}, \ldots, e^{n}\right)$ is a basis for $\mathrm{V}^{*}$, then the set of elements
$$
\left\{e^{i_{1}} \wedge e^{i_{2}} \wedge \cdots \wedge e^{i_{k}}: 1 \leq i_{1}<i_{2}<\cdots<i_{k} \leq n\right\}
$$
is a basis for $L_{\text {alt }}^{k}(\mathrm{V})$. Thus $\operatorname{dim}\left(L_{\text {alt }}^{k}(\mathrm{V})\right)=\binom{n}{k}=\frac{n!}{k!(n-k)!}$. In particular, $L_{\text {alt }}^{k}(\mathrm{V})=0$ for $k>n$.
\end{THEO}

\begin{PROOF}{GA-L09-06-14}{P: Basis von $L_{\text {alt }}^{k}(\mathrm{V})$}
Since any $\alpha \in L_{\text {alt }}^{k}(\mathrm{V})$ is also a member of $T_{k}^{0}(\mathrm{V})$, we may write
$$
\alpha=\sum \alpha_{i_{1} \ldots i_{k}} e^{i_{1}} \otimes \cdots \otimes e^{i_{k}},
$$
where $\alpha_{i_{1} \ldots i_{k}}=\alpha\left(e_{i_{1}}, \ldots, e_{i_{k}}\right)$. By Alt $(\alpha)=\alpha$ and the linearity of Alt we have
$$
\alpha=\sum \alpha_{i_{1} \ldots i_{k}} \operatorname{Alt}\left(e^{i_{1}} \otimes \cdots \otimes e^{i_{k}}\right)=\frac{1}{k!} \sum \alpha_{i_{1} \ldots i_{k}} e^{i_{1}} \wedge \cdots \wedge e^{i_{k}} .
$$
We conclude that the set of elements of the form $e^{i_{1}} \wedge \cdots \wedge e^{i_{k}}$ spans $L_{\text {alt }}^{k_{1}}(\mathrm{V})$. Furthermore, if we use the fact that both $\alpha_{i_{\sigma_{1}} \ldots i_{\sigma_{k}}}=\operatorname{sgn} \sigma \alpha_{i_{1} \ldots i_{k}}$ and $e^{i_{\sigma_{1}}} \wedge \cdots \wedge e^{i_{\sigma_{k}}}=\operatorname{sgn} \sigma e^{i_{1}} \wedge \cdots \wedge e^{i_{k}}$ for any permutation $\sigma \in S_{k}$, we see that we can permute the indices into increasing order and collect terms to get
$$
\alpha=\frac{1}{k!} \sum \alpha_{i_{1} \ldots i_{k}} e^{i_{1}} \wedge \cdots \wedge e^{i_{k}}=\sum_{i_{1}<i_{2}<\ldots<i_{k}} a_{i_{1} i_{2}, . ., i_{k}} e^{i_{1}} \wedge e^{i_{2}} \wedge \cdots \wedge e^{i_{k}},
$$
where in the last expression we sum only over strictly increasing indices. We can check that the set of $\binom{n}{k}$ elements of the form $e^{i_{1}} \wedge e^{i_{2}} \wedge \cdots \wedge e^{i_{k}}$ with $1 \leq i_{1}<i_{2}<\cdots<i_{k} \leq n$ is linearly independent as follows: Suppose $\alpha=\sum_{i_{1}<i_{2}<\cdots<i_{k}} a_{i_{1} i_{2}, . ., i_{k}} e^{i_{1}} \wedge e^{i_{2}} \wedge \cdots \wedge e^{i_{k}}=0$. Let $j_{1}<j_{2}<\cdots<j_{k}$. Then we have
$$
\begin{aligned}
0 & =\alpha\left(e_{j_{1}}, \ldots, e_{j_{k}}\right) \\
& =\sum_{i_{1}<i_{2}<\cdots<i_{k}} a_{i_{1} i_{2} \ldots i_{k}} e^{i_{1}} \wedge e^{i_{2}} \wedge \cdots \wedge e^{i_{k}}\left(e_{j_{1}}, \ldots, e_{j_{k}}\right) \\
& =\sum_{i_{1}<i_{2}<\cdots<i_{k}} a_{i_{1} i_{2} \ldots i_{k}} \epsilon_{j_{1} \ldots j_{k}}^{i_{1} \ldots i_{k}}=a_{j_{1} j_{2} \ldots j_{k}}
\end{aligned}
$$
We have used that $\epsilon_{j_{1} \ldots j_{k}}^{i_{1} \ldots i_{k}}$ is zero unless $j_{1} \ldots j_{k}$ is a permutation of $i_{1} \ldots i_{k}$, and then in this case, since both are increasing, we must have $i_{r}=j_{r}$ for $r=1, \ldots, k$. Thus we get $0=a_{j_{1} j_{2} \ldots j_{k}}$, and since the choice of $j$ 's was arbitrary, we have shown independence. Thus, we have proved the following theorem.
\end{PROOF}




Corollary 8.9. 


\begin{KORO}{GA-L09-06-15}{Charakterisierung von linear unabhängigen Kovektoren durch Keilprodukt}
If $\alpha^{1}, \ldots, \alpha^{k} \in \mathrm{V}^{*}$, then $\alpha^{1}, \ldots, \alpha^{k}$ are linearly independent if and only if
$$
\alpha^{1} \wedge \cdots \wedge \alpha^{k} \neq 0
$$
\end{KORO}

Proof. 

\begin{PROOF}{GA-L09-06-16}{P: Charakterisierung von linear unabhängigen Kovektoren durch Keilprodukt}
If $\alpha^{1}, \ldots, \alpha^{k}$ are independent, then there are elements $\alpha^{k+1}, \ldots, \alpha^{n}$ such that $\alpha^{1}, \ldots, \alpha^{n}$ is a basis for $\mathrm{V}^{*}$. Let $v_{1}, \ldots, v_{n}$ be the basis for V dual to the above basis. Then since $\alpha^{1} \wedge \cdots \wedge \alpha^{k}$ is a basis element for $L_{\text {alt }}^{k}(\mathrm{V})$, it cannot be zero.

For the other direction, suppose that, after rearranging if needed, we have
$$
\alpha^{1}=c_{2} \alpha^{2}+\cdots+c_{k} \alpha^{k}
$$
Then we would have
$$
\alpha^{1} \wedge \cdots \wedge \alpha^{k}=\left(c_{2} \alpha^{2}+\cdots+c_{k} \alpha^{k}\right) \wedge \alpha^{2} \wedge \cdots \wedge \alpha^{k}=0
$$
Thus we see that $\alpha^{1} \wedge \cdots \wedge \alpha^{k} \neq 0$ implies that $\alpha^{1}, \ldots, \alpha^{k}$ are linearly independent.
\end{PROOF}


Notation 8.10. In order to facilitate notation, we will sometimes abbreviate a sequence of $k$ integers, say $i_{1}, i_{2}, \ldots, i_{k}$, from the set $\{1,2, \ldots, \operatorname{dim}(\mathrm{V})\}$ as $I$, and $e^{i_{1}} \wedge e^{i_{2}} \wedge \cdots \wedge e^{i_{k}}$ will be written as $e^{I}$ and $\epsilon_{L}^{I}=\epsilon_{\ell_{1} \ldots \ell_{k}}^{i_{1} \ldots i_{k}}$. Also, if we require that $i_{1}<i_{2}<\cdots<i_{k}$, then we will write $\vec{I}$. We will freely use similar self-explanatory notation as we go along without further comment. For example, we may write

$$
\alpha=\sum a_{\vec{I}} e^{\vec{I}}
$$

to mean $\alpha=\sum_{i_{1}<i_{2}<\cdots<i_{k}} a_{i_{1} i_{2} \ldots i_{k}} e^{i_{1}} \wedge e^{i_{2}} \wedge \cdots \wedge e^{i_{k}}$.\\
Whenever we have $\alpha=\sum a_{\vec{I}} e^{\vec{I}}$, where $a_{\vec{I}}=a_{i_{1} \ldots i_{k}}$ and $i_{1}<\cdots<i_{k}$, we can define $a_{J}$ for any $k$-tuple of indices $J=\left(j_{1}, \ldots, j_{k}\right)$ by requiring that $a_{J}=\sum \epsilon_{J}^{\vec{I}} a_{\vec{I}}$. Then we have $a_{J}=0$ when the entries of $J$ are indices that are not distinct. Otherwise, $a_{J}=\epsilon_{J}^{\vec{J}} a_{\vec{J}}$ where $\vec{J}$ is $J$ rearranged in increasing order. But, it is also easy to see that $\alpha\left(e_{j_{1}}, \ldots, e_{j_{k}}\right)=\sum \epsilon_{J}^{\vec{I}} a_{\vec{I}}$. We then have

$$
\begin{aligned}
\alpha & =\sum a_{\vec{I}} e^{\vec{I}}=\frac{1}{k!} \sum a_{I} e^{I}=\frac{1}{k!} \sum a_{i_{1} \ldots i_{k}} e^{i_{1}} \wedge \cdots \wedge e^{i_{k}} \\
& =\sum a_{i_{1} \ldots i_{k}} e^{i_{1}} \otimes \cdots \otimes e^{i_{k}}
\end{aligned}
$$



\begin{DEF}{GA-L09-06-17}{Mehrindex- und Summationsnotation für alternierende Formen}
Zur Vereinfachung der Schreibweise verwenden wir folgende \textbf{Mehrindex-Notation}:

\begin{enumerate}
\item Eine geordnete $k$-Tupelmenge von Indizes $i_{1},i_{2},\dots,i_{k}\in\{1,\dots,\dim(\mathrm{V})\}$ wird durch einen Großbuchstaben
\[
I=(i_{1},i_{2},\ldots,i_{k})
\]
abgekürzt.

\item Für Basisformen $e^{1},\dots,e^{n}$ schreiben wir
\[
e^{i_{1}}\wedge e^{i_{2}}\wedge\cdots\wedge e^{i_{k}} = e^{I},
\quad\text{und analog}\quad
\epsilon_{L}^{I}=\epsilon_{\ell_{1}\ldots \ell_{k}}^{i_{1}\ldots i_{k}},
\]
wobei $\epsilon$ das Signum des Permutations-Tensors bezeichnet.

\item Falls die Indizes streng aufsteigend sind ($i_{1}<i_{2}<\cdots<i_{k}$), schreiben wir den Vektor $I$ als \emph{geordneten Mehrindex} $\vec{I}$.
\end{enumerate}

\medskip
Mit dieser Konvention kann eine alternierende $k$-Form $\alpha\in L_{\mathrm{alt}}^{k}(\mathrm{V})$ in kompakter Form geschrieben werden als
\[
\alpha = \sum a_{\vec{I}}\, e^{\vec{I}}
\quad\text{bzw. explizit}\quad
\alpha = \sum_{i_{1}<i_{2}<\cdots<i_{k}} a_{i_{1}i_{2}\ldots i_{k}}\, e^{i_{1}}\wedge e^{i_{2}}\wedge\cdots\wedge e^{i_{k}}.
\]

\textbf{Erweiterung der Koeffizienten auf alle Index-Tupel:}  
Für ein beliebiges $k$-Tupel $J=(j_{1},\dots,j_{k})$ definieren wir die Komponenten
\[
a_{J} := \sum \epsilon_{J}^{\vec{I}}\, a_{\vec{I}}.
\]
Damit gilt:
\begin{itemize}
\item $a_{J}=0$, falls die Einträge von $J$ nicht alle verschieden sind.
\item Ist $J$ eine Permutation eines geordneten $\vec{J}$, so gilt $a_{J}=\epsilon_{J}^{\vec{J}}\,a_{\vec{J}}$.
\end{itemize}
Weiterhin gilt für die Auswertung auf Basisvektoren
\[
\alpha(e_{j_{1}},\ldots,e_{j_{k}})=\sum \epsilon_{J}^{\vec{I}}\, a_{\vec{I}}.
\]

\textbf{Komponentendarstellung:}  
Aus dieser Notation ergibt sich
\[
\begin{aligned}
\alpha
&=\sum a_{\vec{I}}\,e^{\vec{I}}
=\frac{1}{k!}\sum a_{I}\,e^{I}
=\frac{1}{k!}\sum a_{i_{1}\ldots i_{k}}\, e^{i_{1}}\wedge\cdots\wedge e^{i_{k}} \\
&=\sum a_{i_{1}\ldots i_{k}}\, e^{i_{1}}\otimes\cdots\otimes e^{i_{k}},
\end{aligned}
\]
wobei die letzte Gleichung die kanonische Einbettung der alternierenden Form in den Tensorraum $T^{k}(\mathrm{V}^{*})$ ausdrückt.
\end{DEF}





Exercise 8.11. 


\begin{PROP}{GA-L09-06-18}{Darstellung von Keilprodukt}
For $\alpha=\frac{1}{k!} \sum a_{I} e^{I}$ and $\beta=\frac{1}{\ell!} \sum b_{J} e^{J}$ with $I$ and $J$ as above, we have
$$
\alpha \wedge \beta=\frac{1}{(k+\ell)!} \sum_{i_{1} \ldots i_{k+\ell}}(\alpha \wedge \beta)_{i_{1} \ldots i_{k+\ell}} e^{i_{1}} \wedge \cdots \wedge e^{i_{k+\ell}},
$$
where
$$
(\alpha \wedge \beta)_{i_{1} \ldots i_{k+\ell}}=\frac{1}{k!\ell!} \sum a_{j_{1} \ldots j_{k}} b_{j_{k+1} \ldots j_{k+\ell}} \epsilon_{i_{1} \ldots i_{k+\ell}}^{j_{1} \ldots j_{k+\ell}}
$$
\end{PROP}

Definition 8.12. 


\begin{DEF}{GA-L09-06-19}{Inneres Produkt von Vektor und antisymmetrischer multilinearer Abbildung}
Let $v \in \mathrm{V}$ and $\omega \in L_{\text {alt }}^{k}(\mathrm{V})$. Define $i_{v} \omega \in L_{\text {alt }}^{k-1}(\mathrm{V})$ by
$$
i_{v} \omega\left(v_{1}, \ldots, v_{k-1}\right):=\omega\left(v, v_{1}, \ldots, v_{k-1}\right)
$$
$i_{v} \omega$ is called the interior product of $v$ with $\omega$ or the contraction of $\omega$ by $v$. By convention $\imath_{v} a=0$ for $a \in L_{\text {alt }}^{0}(\mathrm{V}):=\mathbb{R}$. Thus we obtain a linear map $i_{v}: L_{\text {alt }}^{k}(\mathrm{V}) \rightarrow L_{\text {alt }}^{k-1}(\mathrm{V})$.
\end{DEF}


It is clear that $i_{v}$ depends linearly on $v$ and that $i_{v} i_{w} \omega=-i_{w} i_{v} \omega$ for all $v, w \in \mathrm{V}$ and hence $i_{v} \circ i_{v}=0$.

With $L_{\text {alt }}^{1}(\mathrm{V})=\mathrm{V}^{*}$ and $L_{\text {alt }}^{0}(\mathrm{V})=\mathbb{R}$, the sum

$$
L_{\mathrm{alt}}(\mathrm{V})=\bigoplus_{k=0}^{\operatorname{dim}(\mathrm{V})} L_{\mathrm{alt}}^{k}(\mathrm{V})
$$

is made into an $\mathbb{R}$-algebra via the exterior product just defined. Since $\alpha \wedge \beta \in L_{\text {alt }}^{k_{1}+k_{2}}(\mathrm{V})$ whenever $\alpha \in L_{\text {alt }}^{k_{1}}(\mathrm{V})$ and $\beta \in L_{\text {alt }}^{k_{2}}(\mathrm{V})$, the algebra $L_{\text {alt }}(\mathrm{V})$ is a graded algebra (see Definition D.43).

The contraction map $i_{v}$ for $v \in \mathrm{V}$ extends to a map $L_{\text {alt }}(\mathrm{V}) \rightarrow L_{\text {alt }}(\mathrm{V})$. Since $i_{v}\left(L_{\text {alt }}^{k}(\mathrm{V})\right) \subset L_{\text {alt }}^{k-1}(\mathrm{V})$, we say that $i_{v}$ is a map of degree -1 .



\begin{DEF}{GA-L09-06-20}{Kontraktion und graduierte Algebra Struktur}
Sei $\mathrm{V}$ ein endlichdimensionaler reeller Vektorraum.  
Für jedes $v\in\mathrm{V}$ bezeichne $i_{v}:L_{\mathrm{alt}}^{k}(\mathrm{V})\to L_{\mathrm{alt}}^{k-1}(\mathrm{V})$ die \textbf{Kontraktionsabbildung} (auch \textit{Innere Ableitung}) mit dem Vektor $v$.  
Sie ist durch die Eigenschaft bestimmt, dass für eine alternierende $k$-Form $\omega$ und Vektoren $v_{1},\dots,v_{k-1}\in \mathrm{V}$ gilt:
\[
(i_{v}\omega)(v_{1},\dots,v_{k-1})
:=\omega(v,v_{1},\dots,v_{k-1}).
\]
\textbf{Eigenschaften:}
\begin{enumerate}
\item \textbf{Linearität:} Die Abbildung $v\mapsto i_{v}$ ist linear in $v$, d.\,h.
\[
i_{a v + b w} = a\,i_{v} + b\,i_{w}
\quad\text{für alle }a,b\in\mathbb{R},\; v,w\in\mathrm{V}.
\]

\item \textbf{Antikommutativität:} Für alle $v,w\in\mathrm{V}$ gilt
\[
i_{v}\,i_{w} = -\, i_{w}\,i_{v},
\]
und insbesondere $i_{v}\circ i_{v}=0$.

\item \textbf{Erweiterung:}  
Mit
\[
L_{\mathrm{alt}}(\mathrm{V})
:= \bigoplus_{k=0}^{\dim(\mathrm{V})} L_{\mathrm{alt}}^{k}(\mathrm{V}),
\qquad
L_{\mathrm{alt}}^{0}(\mathrm{V})=\mathbb{R},\;
L_{\mathrm{alt}}^{1}(\mathrm{V})=\mathrm{V}^{*},
\]
wird die direkte Summe durch das äußere Produkt $\wedge$ zu einer \textbf{graduierten Algebra} über $\mathbb{R}$.

\item \textbf{Grad der Kontraktion:}  
Da $i_{v}$ die $k$-Formen in $(k-1)$-Formen überführt,
\[
i_{v}\big(L_{\mathrm{alt}}^{k}(\mathrm{V})\big)\subset L_{\mathrm{alt}}^{k-1}(\mathrm{V}),
\]
sagt man, $i_{v}$ ist eine \emph{Abbildung vom Grad $-1$} in der graduierten Algebra $L_{\mathrm{alt}}(\mathrm{V})$.
\end{enumerate}
\end{DEF}



Proposition 8.13. 

\begin{PROP}{GA-L09-06-21}{Eindeutige Charakterisierung der Kontraktionsabbildung}
For $v \in \mathrm{V}$, the contraction map $i_{v}$ satisfies the product rule
$$
i_{v}(\alpha \wedge \beta)=\left(\imath_{v} \alpha\right) \wedge \beta+(-1)^{k} \alpha \wedge\left(\imath_{v} \beta\right) \text { for } \alpha \in L_{\mathrm{alt}}^{k}(\mathrm{V})
$$
The map $i_{v}: L_{\text {alt }}(\mathrm{V}) \rightarrow L_{\text {alt }}(\mathrm{V})$ is the unique degree -1 map satisfying the above product rule and satisfying $i_{v} a=0$ for $a \in \mathbb{R}$ and
$$
i_{v} \theta=\theta(v) \text { for } \theta \in \mathrm{V}^{*}=L_{\mathrm{alt}}^{1}(\mathrm{V}) \text { and } v \in \mathrm{V} \text {. }
$$
\end{PROP}

Proof. 

\begin{PROOF}{GA-L09-06-22}{P: Eindeutige Charakterisierung der Kontraktionsabbildung}
Let $\alpha \in L_{\text {alt }}^{k}(\mathrm{V})(M)$ and $\beta \in L_{\text {alt }}^{\ell}(\mathrm{V})(M)$. In the following computation we use the permutation $\widetilde{\sigma}$ which is given by
$$
(2,3, \ldots, k+1,1, k+2, \ldots, k+\ell) \mapsto(1,2, \ldots, k+\ell) .
$$
The sign of $\widetilde{\sigma}$ is $(-1)^{k}$. We compute as follows:
$$
\begin{aligned}
\left(v_{v} \alpha \wedge\right. & \left.\beta+(-1)^{k} \alpha \wedge \imath_{v} \beta\right)\left(v_{2}, \ldots, v_{k+\ell}\right) \\
= & \frac{(k+\ell-1)!}{(k-1)!\ell!} \operatorname{Alt}\left(\imath_{v} \alpha \otimes \beta\right)\left(v_{2}, \ldots, v_{k+\ell}\right) \\
& +(-1)^{k} \frac{(k+\ell-1)!}{k!(\ell-1)!} \operatorname{Alt}\left(\alpha \otimes \imath_{v} \beta\right)\left(v_{2}, \ldots, v_{k+1}\right) \\
= & \sum_{\sigma} \frac{\operatorname{sgn}(\sigma)}{(k-1)!\ell!} \alpha\left(v, v_{\sigma_{2}}, \ldots, v_{\sigma_{k}}\right) \beta\left(v_{\sigma_{k+1}}, \ldots, v_{\sigma_{k+\ell}}\right) \\
& +(-1)^{k} \frac{1}{k!(\ell-1)!} \sum_{\sigma} \operatorname{sgn}(\sigma) \alpha\left(v_{\sigma_{2}}, \ldots, v_{\sigma_{(k+1)}}\right) \beta\left(v, v_{\sigma_{k+2}}, \ldots, v_{\sigma_{k+\ell}}\right) \\
= & \frac{1}{(k-1)!\ell!} \sum_{\sigma} \operatorname{sgn}(\sigma) \alpha\left(v, v_{\sigma_{2}}, \ldots, v_{\sigma_{k}}\right) \beta\left(v_{\sigma_{k+1}}, \ldots, v_{\sigma_{k+\ell}}\right) \\
& +(-1)^{k} \frac{1}{k!(\ell-1)!} \sum_{\sigma} \operatorname{sgn}(\sigma \widetilde{\sigma}) \alpha\left(v, v_{\sigma_{2}}, \ldots, v_{\sigma_{k+1}}\right) \beta\left(v_{\sigma_{k+1}}, \ldots, v_{\sigma_{k+\ell}}\right) \\
= & \frac{1}{(k-1)!\ell!} \sum_{\sigma} \operatorname{sgn}(\sigma) \alpha\left(v, v_{\sigma_{2}}, \ldots, v_{\sigma_{k}}\right) \beta\left(v_{\sigma_{k+1}}, \ldots, v_{\sigma_{k+\ell}}\right) \\
& +\frac{1}{k!(\ell-1)!} \sum_{\sigma} \operatorname{sgn}(\sigma) \alpha\left(v, v_{\sigma_{2}}, \ldots, v_{\sigma_{k+1}}\right) \beta\left(v_{\sigma_{k+1}}, \ldots, v_{\sigma_{k+\ell}}\right) \\
= & \frac{1}{(k-1)!\ell!} \sum_{\sigma} \alpha\left(v, v_{\sigma_{2}}, \ldots, v_{\sigma_{(k+1)}}\right) \beta\left(v_{\sigma_{k+1}}, \ldots, v_{\sigma_{k+\ell}}\right) \\
= & \alpha \wedge \beta\left(v, v_{2}, \ldots, v_{k+1}\right)=v_{v}(\alpha \wedge \beta)\left(v_{2}, \ldots, v_{k+1}\right)
\end{aligned}
$$
We leave the proof of the remaining statements of the theorem to the reader (see Problem 14).
\end{PROOF}


\begin{PROP}{GA-L09-06-23}{Darstellung der Kontraktion von Keilprodukten}
It follows from the above that if $\theta_{1}, \ldots, \theta_{k} \in \mathrm{V}^{*}$ and $v \in \mathrm{V}$, then
$$
i_{v}\left(\theta_{1} \wedge \cdots \wedge \theta_{k}\right)=\sum_{\ell=1}^{k}(-1)^{k+1} \theta_{\ell}(v) \theta_{1} \wedge \cdots \wedge \widehat{\theta_{\ell}} \wedge \cdots \wedge \theta_{k}
$$
where the caret over $\theta_{\ell}$ denotes omission.
\end{PROP}


Proposition 8.14. 


\begin{PROP}{GA-L09-06-24}{Verträglichkeit von Pullback und Keilprodukt}
If $\lambda \in L(\mathrm{V}, \mathrm{V}), \alpha \in L_{\text {alt }}^{k_{1}}(\mathrm{V})$ and $\beta \in L_{\text {alt }}^{k_{2}}(\mathrm{V})$, then
$$
\lambda^{*}(\alpha \wedge \beta)=\lambda^{*} \alpha \wedge \lambda^{*} \beta
$$
\end{PROP}

Proof. 

\begin{PROOF}{GA-L09-06-25}{P: Verträglichkeit von Pullback und Keilprodukt}
This follows from Proposition 1.83 and the definition of the exterior product.
\end{PROOF}

We look more closely at $L_{\text {alt }}^{n}(\mathrm{V})$ where $n=\operatorname{dimV}$. The dimension of $L_{\text {alt }}^{n}(\mathrm{V})$ is one, and any nonzero element of $L_{\text {alt }}^{n}(\mathrm{V})$ provides a basis. If $\lambda \in L(\mathrm{V}, \mathrm{V})$, then $\lambda^{*}: L_{\text {alt }}^{n}(\mathrm{V}) \rightarrow L_{\text {alt }}^{n}(\mathrm{V})$ is a linear transformation\\
between one-dimensional vector spaces, and so it must be multiplication by an element of $\mathbb{R}$. Thus there is a unique number $\operatorname{det}(\lambda) \in \mathbb{R}$ such that

$$
\lambda^{*} \omega=\operatorname{det}(\lambda) \omega
$$

for any $\omega \in L_{\text {alt }}^{n}(\mathrm{V})$. This number is called the determinant of $\lambda$. This provides a definition of determinant that does not involve a choice of basis. We will show that if $A=\left(a_{j}^{i}\right)$ represents $\lambda \in L(\mathrm{V}, \mathrm{V})$ with respect to a basis for V , then $\operatorname{det}(\lambda)=\operatorname{det}(A)$, where the determinant of a matrix is given by the standard definition. Let $\lambda \in L(\mathrm{V}, \mathrm{V})$ and suppose that $\lambda\left(e_{i}\right)=\sum a_{i}^{j} e_{j}$ for some basis $\left(e_{1}, \ldots, e_{n}\right)$ with dual $\left(e^{1}, \ldots, e^{n}\right)$. Then $A=\left(a_{j}^{i}\right)$ represents $\lambda$, and we have

$$
\begin{aligned}
\lambda^{*}\left(e^{1} \wedge \cdots \wedge e^{n}\right)\left(e_{1}, \ldots, e_{n}\right) & =\left(e^{1} \wedge \cdots \wedge e^{n}\right)\left(\lambda e_{1}, \ldots, \lambda e_{n}\right) \\
& =\operatorname{det}\left(e^{i}\left(\lambda e_{j}\right)\right)=\operatorname{det} A
\end{aligned}
$$

On the other hand, $\lambda^{*}\left(e^{1} \wedge \cdots \wedge e^{n}\right)=(\operatorname{det} \lambda)\left(e^{1} \wedge \cdots \wedge e^{n}\right)$, and since $e^{1} \wedge \cdots \wedge e^{n}\left(e_{1}, \ldots, e_{n}\right)=1$, it must be that $\operatorname{det}(\lambda)=\operatorname{det}(A)$.



\begin{PROP}{GA-L09-06-26}{Basisfreie Definition der Determinante über alternierende $n$-Formen}
Sei $\mathrm{V}$ ein endlichdimensionaler reeller Vektorraum mit $\dim(\mathrm{V}) = n$.  
Dann besitzt der Raum der alternierenden $n$-Formen
\[
L_{\mathrm{alt}}^{n}(\mathrm{V})
\]
die Dimension $1$, d.\,h. jede nichtverschwindende $n$-Form $\omega$ bildet eine Basis dieses Raumes.

\textbf{Definition.}  
Für eine lineare Abbildung $\lambda \in L(\mathrm{V}, \mathrm{V})$ induziert der Pullback
\[
\lambda^{*}: L_{\mathrm{alt}}^{n}(\mathrm{V}) \longrightarrow L_{\mathrm{alt}}^{n}(\mathrm{V})
\]
eine lineare Selbstabbildung eines eindimensionalen Vektorraums und ist daher durch Multiplikation mit einer reellen Zahl bestimmt.  
Es existiert also genau eine reelle Zahl $\det(\lambda)\in\mathbb{R}$ derart, dass
\[
\boxed{\;
\lambda^{*}\omega = \det(\lambda)\,\omega
\qquad\text{für alle }\omega\in L_{\mathrm{alt}}^{n}(\mathrm{V}).
\;}
\]
Diese Zahl wird \emph{Determinante} von $\lambda$ genannt.  
Dies ist eine basisfreie Definition der Determinante.

\textbf{Beweis der Übereinstimmung mit der klassischen Determinante.}  
Sei $(e_{1},\ldots,e_{n})$ eine Basis von $\mathrm{V}$ und $(e^{1},\ldots,e^{n})$ die zugehörige Dualbasis.  
Angenommen, $\lambda(e_{i})=\sum_{j} a^{j}_{i}\, e_{j}$, sodass $A=(a^{j}_{i})$ die darstellende Matrix von $\lambda$ ist.  
Dann gilt:
\[
\begin{aligned}
\lambda^{*}(e^{1}\wedge\cdots\wedge e^{n})(e_{1},\ldots,e_{n})
&= (e^{1}\wedge\cdots\wedge e^{n})(\lambda e_{1},\ldots,\lambda e_{n}) \\
&= \det\!\big(e^{i}(\lambda e_{j})\big)
= \det(A).
\end{aligned}
\]
Andererseits liefert die Definition der Determinante
\[
\lambda^{*}(e^{1}\wedge\cdots\wedge e^{n})
= \det(\lambda)\,(e^{1}\wedge\cdots\wedge e^{n}).
\]
Da $e^{1}\wedge\cdots\wedge e^{n}(e_{1},\ldots,e_{n})=1$, folgt unmittelbar
\[
\det(\lambda)=\det(A).
\]

\textbf{Bemerkung.}  
Diese Darstellung zeigt, dass $\det(\lambda)$ als das Skalierungsverhältnis verstanden werden kann, mit dem $\lambda$ orientierte Volumenelemente (d.\,h. $n$-Formen) transformiert.
\end{PROP}



Exercise 8.15. 

\begin{PROP}{GA-L09-06-27}{Eigenschaften der Basisunabhängigen Determinante}
If $\ell, \lambda \in L(\mathrm{V}, \mathrm{V})$, then\\
\begin{enumerate}
\item $\operatorname{det}(\ell \circ \lambda)=\operatorname{det} \ell \operatorname{det} \lambda$
\item $\operatorname{det}(\mathrm{id})=1$
\item if $\lambda \in \mathrm{GL}(\mathrm{V})$, then $\operatorname{det}\left(\lambda^{-1}\right)=(\operatorname{det} \lambda)^{-1}$
\end{enumerate}
\end{PROP}


We now briefly discuss orientations of real vector spaces. In what follows, let V be a real vector space with $n=\operatorname{dimV}$. The nonzero elements of $L_{\text {alt }}^{n}(\mathrm{V})$ are sometimes referred to as volume elements (although this term will also apply to a global object later on). Two nonzero elements $\omega_{1}$ and $\omega_{2}$ of $L_{\text {alt }}^{n}(\mathrm{V})$ are said to be equivalent if there is a scalar $c>0$ such that $\omega_{1}=c \omega_{2}$. The equivalence class of $\omega$ will be denoted $[\omega]$. There are clearly exactly two equivalence classes. An equivalence class is referred to as an orientation for $V$. An oriented vector space is a vector space with a choice of orientation and is sometimes written as a pair ( $\mathrm{V},[\omega]$ ).

An ordered basis $\left(e_{1}, \ldots, e_{n}\right)$ for an oriented real vector space $(\mathrm{V},[\omega])$ is said to be positive with respect to the orientation if $\omega\left(e_{1}, \ldots, e_{n}\right)>0$ for some and hence any $\omega \in[\omega]$. Equivalently, $\left(e_{1}, \ldots, e_{n}\right)$ is positive if $e^{1} \wedge \cdots \wedge e^{n} \in[\omega]$. Actually a choice of ordered basis for a real vector space determines an orientation with respect to which it is positive. Indeed, if $e^{1}, \ldots, e^{n}$ is dual to such a basis, then choose $[\omega]$ where $\omega=e^{1} \wedge \cdots \wedge e^{n}$.



\begin{DEF}{GA-L09-06-28}{Orientierung eines reellen Vektorraums}
Sei $\mathrm{V}$ ein endlichdimensionaler reeller Vektorraum mit $\dim(\mathrm{V})=n$.

\textbf{(1) Volumenelemente.}
Die nichtverschwindenden Elemente von $L_{\mathrm{alt}}^{n}(\mathrm{V})$ heißen
\emph{Volumenelemente} von $\mathrm{V}$.
Jedes Volumenelement $\omega\in L_{\mathrm{alt}}^{n}(\mathrm{V})\setminus\{0\}$ bestimmt eine Orientierungsklasse.

\textbf{(2) Äquivalenzrelation.}
Zwei nichtverschwindende $n$-Formen $\omega_{1},\omega_{2}\in L_{\mathrm{alt}}^{n}(\mathrm{V})$
heißen \emph{äquivalent}, wenn ein reeller Skalar $c>0$ existiert mit
\[
\omega_{1}=c\,\omega_{2}.
\]
Die Äquivalenzklasse einer Form $\omega$ wird mit $[\omega]$ bezeichnet.

\textbf{(3) Orientierung.}
Da die Relation genau zwei Äquivalenzklassen liefert,
heißen diese \emph{Orientierungen} des Vektorraums~$\mathrm{V}$.
Ein \emph{orientierter Vektorraum} ist ein Paar $(\mathrm{V},[\omega])$,
bestehend aus einem reellen Vektorraum $\mathrm{V}$ und einer Wahl einer Orientierungsklasse~$[\omega]$.

\textbf{(4) Positive Basen.}
Eine geordnete Basis $(e_{1},\dots,e_{n})$ von $\mathrm{V}$ heißt
\emph{positiv} bezüglich der Orientierung~$[\omega]$,
wenn für eine (und daher jede) Repräsentantin $\omega\in[\omega]$ gilt:
\[
\omega(e_{1},\dots,e_{n})>0.
\]
Äquivalent dazu ist $(e_{1},\dots,e_{n})$ positiv genau dann,
wenn $e^{1}\wedge\cdots\wedge e^{n}\in[\omega]$,
wobei $(e^{1},\dots,e^{n})$ die duale Basis ist.

\textbf{(5) Orientierung durch eine Basis.}
Jede geordnete Basis $(e_{1},\dots,e_{n})$ bestimmt kanonisch eine Orientierung von~$\mathrm{V}$,
nämlich die Äquivalenzklasse
\[
[\omega]\quad\text{mit}\quad \omega:=e^{1}\wedge\cdots\wedge e^{n},
\]
für die $(e_{1},\dots,e_{n})$ positiv ist.
\end{DEF}


Definition 8.16. 

\begin{DEF}{GA-L09-06-29}{Orientierungserhaltender Isomorphismus zwischen orientierten Vektorräumen}
Let $\left(V_{1},\left[\omega_{1}\right]\right)$ and $\left(V_{2},\left[\omega_{2}\right]\right)$ be oriented real vector spaces. A linear isomorphism $\lambda: \mathrm{V}_{1} \rightarrow \mathrm{V}_{2}$ is said to be orientation preserving if $\lambda^{*} \omega_{2}=c \omega_{1}$ for some $c>0$ and some, and hence any, choices $\omega_{1} \in\left[\omega_{1}\right]$ and $\omega_{2} \in\left[\omega_{2}\right]$.
\end{DEF}


\begin{DEF}{GA-L09-06-30}{Orientierungserhaltende Elemente in $\GL(V)$}
When one talks about an element $\lambda \in \mathrm{GL}(\mathrm{V})$ being orientation preserving, one means that $\operatorname{det} \lambda>0$ and this is tantamount to $\lambda:(\mathrm{V},[\omega]) \rightarrow (\mathrm{V},[\omega])$ being orientation preserving for any choice of orientation $[\omega]$.
\end{DEF}


8.1.1. 


The abstract Grassmann algebra. We now take a very abstract approach to constructing an algebra that will be seen to be isomorphic to $L_{\text {alt }}(\mathrm{V})$. We wish to construct a space that is universal with respect to alternating multilinear maps. We work in the category of real vector spaces, although much of what we do here makes sense for modules. Consider the tensor space $T^{k}(\mathrm{V}):=\bigotimes^{k} \mathrm{V}$ (take any realization of the abstract tensor product as in Definition D.17). Let A be the submodule of $T^{k}(\mathrm{V})$ generated by elements of the form

$$
v_{1} \otimes \cdots v_{i} \otimes \cdots \otimes v_{i} \cdots \otimes v_{k} .
$$

In other words, A is generated by simple tensors with two (or more) equal factors. Recall that associated to $T^{k}(\mathrm{V})$ we have the canonical map $\otimes$ : $\mathrm{V} \times \cdots \times \mathrm{V} \rightarrow T^{k}(\mathrm{V})$ defined so that $\otimes\left(v_{1}, \ldots, v_{k}\right)=v_{1} \otimes \cdots \otimes v_{k}$. We define the space of $k$-vectors to be

$$
\mathrm{V} \wedge \cdots \wedge \mathrm{V}:=\bigwedge^{k} \mathrm{V}:=T^{k}(\mathrm{V}) / \mathrm{A} .
$$

Let $\wedge_{k}: \mathrm{V} \times \cdots \times \mathrm{V} \rightarrow \bigwedge^{k} \mathrm{V}$ be the composition of the canonical map $\otimes$ with quotient map of $T^{k}(\mathrm{V})$ onto $\wedge^{k} \mathrm{V}$. This map turns out to be an alternating multilinear map. We will denote $\wedge_{k}\left(v_{1}, \ldots, v_{k}\right)$ by $v_{1} \wedge \cdots \wedge v_{k}$. Using the universal property of $T^{k}(\mathrm{V})$ as described in Appendix D, one can show that the pair ( $\bigwedge^{k} \mathrm{V}, \wedge_{k}$ ) is universal with respect to alternating $k$-multilinear maps: That is, given any alternating $k$-multilinear map $\alpha: \mathrm{V} \times \cdots \times \mathrm{V} \rightarrow \mathrm{W}$, there is a unique linear map $\alpha_{\wedge}: \bigwedge^{k} \mathrm{V} \rightarrow \mathrm{~W}$ such that $\alpha=\alpha_{\wedge} \circ \wedge_{k}$; that is, the following diagram commutes:\\
\includegraphics[scale=0.25, center]{2025_10_09_265d7e1a22c89f79c21eg-369}

Notice that we also have that $v_{1} \wedge \cdots \wedge v_{k}$ is the image of $\mathrm{v}_{1} \otimes \cdots \otimes \mathrm{v}_{k}$ under the quotient map $T^{k}(\mathrm{V}) \rightarrow \bigwedge^{k} \mathrm{V}$. Next we define $\Lambda \mathrm{V}:=\bigoplus_{k=0}^{\infty} \bigwedge^{k} \mathrm{V}$, which is a direct sum, and we take $\Lambda^{0} \mathrm{V}:=\mathbb{R}$. We impose on $\Lambda \mathrm{V}$ the multiplication generated by the rule

$$
\left(v_{1} \wedge \cdots \wedge v_{i}\right) \times\left(v_{1}^{\prime} \wedge \cdots \wedge v_{j}^{\prime}\right) \mapsto v_{1} \wedge \cdots \wedge v_{i} \wedge v_{1}^{\prime} \wedge \cdots \wedge v_{j}^{\prime} \in \bigwedge^{i+j} \mathrm{V} .
$$

The resulting graded algebra is called the Grassmann algebra or exterior algebra. (Of course, the definition of $\wedge$ here is different from what we defined previously.) If we need to have a $\mathbb{Z}$ grading rather than an $\mathbb{N}$ grading,\\
we may define $\Lambda^{k} \mathrm{V}:=0$ for $k<0$ and extend the multiplication in the obvious way. Elements of $\Lambda \mathrm{V}$ are called multivectors and specifically, elements of $\bigwedge^{k} \mathrm{V}$ are called $k$-multivectors.




\begin{DEF}{GA-L09-06-31}{Abstrakte Grassmann-Algebra (äußere Algebra)}
Sei $\mathrm{V}$ ein endlichdimensionaler reeller Vektorraum.  
Ziel ist der Aufbau einer \emph{universellen Algebra} für alternierende multilineare Abbildungen.

\textbf{(1) Tensorraum.}  
Für jedes $k\in\mathbb{N}$ sei
\[
T^{k}(\mathrm{V}) := \bigotimes^{k}\mathrm{V}
\]
der $k$-te Tensorraum über $\mathrm{V}$, und sei
\[
\otimes : \mathrm{V}^{k} \to T^{k}(\mathrm{V}), \quad 
(v_{1},\dots,v_{k}) \mapsto v_{1}\otimes\cdots\otimes v_{k}
\]
die kanonische multilineare Abbildung (vgl. Definition~D.17).

\textbf{(2) Alternierende Quotientenbildung.}  
Sei $A \subset T^{k}(\mathrm{V})$ der von allen Tensoren der Form
\[
v_{1}\otimes\cdots\otimes v_{i}\otimes\cdots\otimes v_{i}\otimes\cdots\otimes v_{k}
\]
erzeugte Unterraum, also der Unterraum, der alle einfachen Tensoren mit mindestens zwei identischen Faktoren enthält.  
Wir definieren den \emph{Raum der $k$-Vektoren} durch den Quotienten
\[
\boxed{\;
\bigwedge^{k}\mathrm{V} := T^{k}(\mathrm{V})/A.
\;}
\]
Der natürliche Quotientenhomomorphismus $q:T^{k}(\mathrm{V})\to \bigwedge^{k}\mathrm{V}$ erlaubt die Definition der kanonischen alternierenden Abbildung
\[
\wedge_{k} := q\circ \otimes :
\mathrm{V}^{k}\to \bigwedge^{k}\mathrm{V},\qquad
(v_{1},\dots,v_{k})\mapsto v_{1}\wedge\cdots\wedge v_{k}.
\]

\textbf{(3) Universelle Eigenschaft.}  
Das Paar $(\bigwedge^{k}\mathrm{V},\wedge_{k})$ ist \emph{universell}
unter allen alternierenden $k$-multilinearen Abbildungen:
Für jeden reellen Vektorraum $\mathrm{W}$ und jede alternierende multilineare Abbildung
\[
\alpha:\mathrm{V}^{k}\to \mathrm{W}
\]
existiert genau eine lineare Abbildung
\[
\alpha_{\wedge}:\bigwedge^{k}\mathrm{V}\to \mathrm{W}
\]
mit
\[
\alpha = \alpha_{\wedge}\circ \wedge_{k}.
\]
Dieses universelle Diagramm kommutiert:
\[
\begin{tikzcd}
\mathrm{V}^{k}\arrow[rr,"\wedge_{k}"]\arrow[dr,"\alpha"'] & &
\bigwedge^{k}\mathrm{V}\arrow[dl,"\exists!\,\alpha_{\wedge}"] \\
& \mathrm{W} &
\end{tikzcd}
\]

\textbf{(4) Direkte Summe und Algebra-Struktur.}  
Wir definieren die \emph{Grassmann- oder äußere Algebra} von~$\mathrm{V}$ als
\[
\boxed{\;
\Lambda \mathrm{V}
:= \bigoplus_{k=0}^{\infty} \bigwedge^{k}\mathrm{V},
\quad
\Lambda^{0}\mathrm{V}:=\mathbb{R}.
\;}
\]
Für $i,j\ge 0$ wird die Multiplikation durch
\[
\big(\,v_{1}\wedge\cdots\wedge v_{i}\,\big)
\cdot
\big(\,v'_{1}\wedge\cdots\wedge v'_{j}\,\big)
:= v_{1}\wedge\cdots\wedge v_{i}\wedge v'_{1}\wedge\cdots\wedge v'_{j}
\in \bigwedge^{i+j}\mathrm{V}
\]
definiert und linear auf ganz $\Lambda \mathrm{V}$ fortgesetzt.

\textbf{(5) Eigenschaften.}  
Damit wird $\Lambda \mathrm{V}$ eine \emph{gradierte Algebra} über~$\mathbb{R}$, d.\,h.
\[
\Lambda\mathrm{V} = \bigoplus_{k\ge 0}\Lambda^{k}\mathrm{V},
\qquad
\Lambda^{i}\mathrm{V}\cdot \Lambda^{j}\mathrm{V}\subseteq \Lambda^{i+j}\mathrm{V}.
\]
Für $k<0$ setzt man konventionsgemäß $\Lambda^{k}\mathrm{V}:=0$.
Elemente von $\Lambda^{k}\mathrm{V}$ heißen \emph{$k$-Vektoren},  
und Elemente von $\Lambda \mathrm{V}$ heißen allgemein \emph{Multivektoren}.

\textbf{(6) Bemerkung.}  
Diese abstrakte Konstruktion liefert eine basisfreie Definition der äußeren Algebra,
die kanonisch isomorph ist zu der algebraischen Struktur der alternierenden Formen
$L_{\mathrm{alt}}(\mathrm{V})$.
\end{DEF}







Notice that since $(v+w) \wedge(w+v)=0$, it follows that $v \wedge w=-w \wedge v$. In fact, any transposition of the factors in a simple element such as $v_{1} \wedge \cdots \wedge v_{k}$, introduces a change of sign:

$$
\begin{aligned}
& v_{1} \wedge \cdots \wedge v_{i} \wedge \cdots \wedge v_{j} \wedge \cdots \wedge v_{k} \\
& =-v_{1} \wedge \cdots \wedge v_{j} \wedge \cdots \wedge v_{i} \wedge \cdots \wedge v_{k}
\end{aligned}
$$

Lemma 8.17. 

\begin{LEM}{GA-L09-06-32}{Basis für Raum der $k$-Vektoren}
If V has dimension $n$, then $\bigwedge^{k} \mathrm{V}=0$ for $k>n$. If $e_{1}, \ldots, e_{n}$ is a basis for V , then the set
$$
\left\{e_{i_{1}} \wedge \cdots \wedge e_{i_{k}}: 1 \leq i_{1}<\cdots<i_{k} \leq n\right\}
$$
is a basis for $\bigwedge^{k} \mathrm{V}$ where we agree that $e_{i_{1}} \wedge \cdots \wedge e_{i_{k}}=1$ if $k=0$.
\end{LEM}


Proof. 

\begin{PROOF}{GA-L09-06-33}{P: Basis für Raum der $k$-Vektoren}
The first statement is easy and we leave it to the reader. We will show that the set above is indeed a basis. First note that $\Lambda^{n} \mathrm{V}$ is spanned by $e_{1} \wedge \cdots \wedge e_{n}$. To see that $e_{1} \wedge \cdots \wedge e_{n}$ is not zero we let
$$
\operatorname{det}: \mathrm{V} \times \cdots \times \mathrm{V} \rightarrow \mathbb{R}
$$
be the multilinear map given by representing the arguments as column vectors of components with respect to the given basis and then taking the determinant of the $n \times n$ matrix built from these column vectors. Then $\operatorname{det}\left(e_{1}, \ldots, e_{n}\right)=1$. But by the universal property above there is a linear map $\operatorname{det}_{\wedge}$ such that $\operatorname{det}=\operatorname{det}_{\wedge} \circ \wedge_{k}$ and so
$$
\begin{aligned}
\operatorname{det}_{\wedge}\left(e_{1} \wedge \cdots \wedge e_{n}\right) & =\operatorname{det}_{\wedge} \circ \wedge_{k}\left(e_{1}, \ldots, e_{n}\right) \\
& =\operatorname{det}\left(e_{1}, \ldots, e_{n}\right)=1 ;
\end{aligned}
$$
thus, we conclude that $e_{1} \wedge \cdots \wedge e_{n}$ is not zero (and is a basis for $\wedge^{n} \mathrm{V}$ ). Now it is easy to see that the elements of the form $e_{i_{1}} \wedge \cdots \wedge e_{i_{k}}$ span $\wedge^{k} \mathrm{V}$. To see that we have linear independence, suppose that
$$
\sum_{1 \leq i_{1}<\cdots<i_{k} \leq n} a_{i_{1} \ldots i_{k}} e_{i_{1}} \wedge \cdots \wedge e_{i_{k}}=0 .
$$
Now multiply both sides by $e_{j_{1}} \wedge \cdots \wedge e_{j_{n-k}}$, where $\left\{j_{1}, \ldots, j_{n-k}\right\}$ is the complement of $\left\{i_{1}, \ldots, i_{k}\right\}$. Then we obtain
$$
\pm a_{i_{1} \ldots i_{k}} e_{1} \wedge \cdots \wedge e_{n}=0
$$
from which we conclude that $a_{i_{1} \ldots i_{k}}=0$, and since $i_{1}, \ldots, i_{k}$ was arbitrary, we are done.
\end{PROOF}



Exercise 8.18. 


\begin{PROP}{GA-L09-06-34}{Charakterisierung von linearer Unabhängigkeit im Raum der $k$-Vektoren}
$v_{1}, \ldots, v_{k} \in \mathrm{V}$ are linearly independent if and only if $v_{1} \wedge \cdots \wedge v_{k} \neq 0$. Compare this with Corollary 8.9.
\end{PROP}



Definition 8.19. 


\begin{DEF}{GA-L09-06-35}{Einfache $k$-Vektoren}
An element $\xi \in \Lambda^{k} \mathrm{V}$ is called decomposable if $\xi= v_{1} \wedge \cdots \wedge v_{k}$ for some $v_{1}, \ldots, v_{k} \in \mathrm{V}$.
\end{DEF}

Let us gain a little practice dealing with multivectors by proving the following proposition:

Proposition 8.20. 

\begin{PROP}{GA-L09-06-36}{Darstellung von $2$-Vektoren}
Let $\xi \in \bigwedge^{2} \mathrm{V}$ with $\xi \neq 0$. Then there is a basis $v_{1}, \ldots, v_{n}$ of V such that
$$
\xi=v_{1} \wedge v_{2}+v_{3} \wedge v_{4}+\cdots+v_{2 r-1} \wedge v_{2 r}
$$
Furthermore, in this case the $r$-fold product $\xi \wedge \cdots \wedge \xi$ is nonzero and decomposable while the $r+1$-fold product is zero.
\end{PROP}

Proof. 

\begin{PROOF}{GA-L09-06-37}{P: Darstellung von $2$-Vektoren}
First we prove that there exists such a decomposition for $\xi$. Let $e_{1}, \ldots, e_{n}$ be a basis for V . We have
$$
\begin{aligned}
\xi=\sum_{i<j} a_{i j} e_{i} \wedge e_{j}= & a_{12} e_{1} \wedge e_{2}+a_{13} e_{1} \wedge e_{2}+\cdots+a_{1 n} e_{1} \wedge e_{n} \\
& +a_{23} e_{2} \wedge e_{3}+a_{24} e_{2} \wedge e_{4}+\cdots+a_{2 n} e_{2} \wedge e_{n}+\xi^{\prime}
\end{aligned}
$$
where $\xi^{\prime}$ does not involve $e_{1}$ or $e_{2}$. By renumbering if necessary we may assume that $a_{12}$ is not zero. But then if
$$
\begin{aligned}
& v_{1}:=e_{1}+\left(a_{23} / a_{12}\right) e_{3}+\cdots+\left(a_{2 n} / a_{12}\right) e_{n} \\
& v_{2}:=a_{12} e_{2}+a_{13} e_{3}+\cdots+a_{1 n} e_{n},
\end{aligned}
$$
we have that $v_{1}, v_{2}, e_{3}, \ldots, e_{n}$ are linearly independent and
$$
\xi=v_{1} \wedge v_{2}+v
$$
where $v$ does not involve $e_{1}$ or $e_{2}$. If $v=0$, then we are done. Otherwise we may repeat the process with
$$
v=\sum_{\substack{i<j \\ i, j \geq 3}} b_{i j} e_{i} \wedge e_{j}
$$
Clearly a simple induction gives
$$
\xi=v_{1} \wedge v_{2}+v_{3} \wedge v_{4}+\cdots+v_{2 r-1} \wedge v_{2 r}
$$
for some nonzero $v_{1}, \ldots, v_{2 r}$ such that $v_{1}, \ldots, v_{2 r}, e_{2 r+1}, \ldots, e_{n}$ is the desired basis.

Next we consider the $r$-fold product. If we set $\phi_{k}:=v_{2 k-1} \wedge v_{2 k}$ for $k=1, \ldots, r$, then
$$
\xi=\sum_{i=1}^{r} \phi_{i} \text { and } \phi_{i} \wedge \phi_{i}=0
$$
On the other hand, if $i \neq j$, then $\phi_{i} \wedge \phi_{j}=\phi_{j} \wedge \phi_{i}$ is plus or minus one times one of the elements in the linearly independent set
$$
\left\{2 v_{\ell_{1}} \wedge v_{\ell_{2}} \wedge v_{\ell_{3}} \wedge v_{\ell_{4}}: \ell_{1}<\ell_{2}<\ell_{3}<\ell_{4} \text { and } \ell_{i} \in\{1, \ldots, 2 r\}\right\} .
$$
Thus
$$
\xi^{2}=\xi \wedge \xi=2 \sum_{1 \leq i<j \leq r} \phi_{i} \wedge \phi_{i} .
$$
In general, it is easy to see that since all of the $\phi_{i}$ commute with each other, we have
$$
\xi^{k}=\sum_{1 \leq i_{1}, \ldots, i_{k} \leq r} \phi_{i_{1}} \wedge \phi_{i_{2}} \wedge \cdots \wedge \phi_{i_{k}}=k!\sum_{1 \leq i_{1}<\cdots<i_{k} \leq r} \phi_{i_{1}} \wedge \phi_{i_{2}} \wedge \cdots \wedge \phi_{i_{k}} .
$$
Then, it is easy to see that
$$
\xi^{r}=r!\phi_{1} \wedge \phi_{2} \wedge \cdots \wedge \phi_{r}=r!v_{1} \wedge v_{2} \wedge \cdots \wedge v_{2 r} \neq 0
$$
while $\xi^{r+1}=0$.
\end{PROOF}




The number $r$ from the previous proposition is called the rank of the element $\xi \in \bigwedge^{2} \mathrm{V}$ and the definition works just as well for elements of $\bigwedge^{2} \mathrm{V}^{*}$. It can be shown that if

$$
\xi=\sum_{i<j} a_{i j} e_{i} \wedge e_{j}=\frac{1}{2} \sum_{i, j} a_{i j} e_{i} \wedge e_{j},
$$

where ( $a_{i j}$ ) is an antisymmetric matrix, then the rank of $\xi$ is the rank of the matrix $\left(a_{i j}\right)$.



\begin{DEF}{GA-L09-06-38}{Rang eines bivectoriellen Elements}
Sei $\mathrm{V}$ ein endlichdimensionaler reeller Vektorraum und 
\[
\xi \in \bigwedge^{2}\mathrm{V}
\]
ein bivectorielles Element.  

\textbf{(1) Definition des Rangs.}  
Der \emph{Rang} von $\xi$ ist die kleinste ganze Zahl $r$, für die es Vektoren
$v_{1},\dots,v_{2r}\in\mathrm{V}$ gibt mit
\[
\xi = v_{1}\wedge v_{2} + v_{3}\wedge v_{4} + \cdots + v_{2r-1}\wedge v_{2r}.
\]
Der Rang misst somit die minimale Anzahl linear unabhängiger 2-Vektoren, deren Summe $\xi$ darstellt.
Die Definition gilt analog für Elemente von $\bigwedge^{2}\mathrm{V}^{*}$.

\textbf{(2) Darstellung in einer Basis.}  
Sei $(e_{1},\dots,e_{n})$ eine Basis von $\mathrm{V}$ und
\[
\xi = \sum_{i<j} a_{ij}\, e_{i}\wedge e_{j}
   = \frac{1}{2}\sum_{i,j} a_{ij}\, e_{i}\wedge e_{j},
\]
wobei die Koeffizienten $a_{ij}$ eine antisymmetrische Matrix $A=(a_{ij})$ bilden, d.\,h.
$a_{ij}=-a_{ji}$.

\textbf{(3) Zusammenhang mit der Matrizenrang.}  
Dann gilt:
\[
\boxed{\;
\operatorname{rank}(\xi) = \operatorname{rank}(A),
\;}
\]
d.\,h. der Rang des bivectoriellen Elements $\xi$ stimmt mit dem gewöhnlichen Rang der antisymmetrischen Matrix seiner Komponenten überein.

\textbf{(4) Bemerkung.}  
Diese Definition ist koordinatenunabhängig:
der Rang von $\xi$ hängt nur von der linearen Struktur von $\mathrm{V}$ ab,
nicht von der gewählten Basis.  
Er beschreibt die Dimension des kleinsten Unterraums $W\subseteq \mathrm{V}$,  
für den $\xi \in \bigwedge^{2}W$ gilt.
\end{DEF}



The following proposition follows easily from the universal properties of the exterior product.


Remark 8.21. There is a natural isomorphism

$$
L_{\mathrm{alt}}^{k}(\mathrm{V} ; \mathrm{W}) \cong L\left(\bigwedge^{k} \mathrm{V} ; \mathrm{W}\right)
$$

Lemma 8.22. In particular,

$$
L_{\mathrm{alt}}^{k}(\mathrm{V}) \cong\left(\bigwedge^{k} \mathrm{V}\right)^{*}
$$



\begin{THEO}{GA-L09-06-39}{Isomorphismus zwischen alternierenden Abbildungen und äußeren Potenzen}
Sei $\mathrm{V}$ ein endlichdimensionaler reeller Vektorraum und $\mathrm{W}$ ein weiterer reeller Vektorraum.  
Dann besteht ein kanonischer, natürlicher Isomorphismus
\[
\boxed{\;
L_{\mathrm{alt}}^{k}(\mathrm{V};\mathrm{W})
\;\cong\;
L\!\left(\bigwedge\nolimits^{k}\mathrm{V},\,\mathrm{W}\right),
\;}
\]
wobei $L_{\mathrm{alt}}^{k}(\mathrm{V};\mathrm{W})$ den Raum aller alternierenden $k$-multilinearen Abbildungen 
$\alpha:\mathrm{V}^{k}\to \mathrm{W}$ bezeichnet.

\textbf{Konstruktion:}  
Der Isomorphismus wird durch die universelle Eigenschaft der äußeren Potenz gegeben:  
Zu jeder alternierenden Abbildung $\alpha:\mathrm{V}^{k}\to \mathrm{W}$ existiert genau eine lineare Abbildung 
\[
\alpha_{\wedge}:\bigwedge\nolimits^{k}\mathrm{V}\to\mathrm{W},
\quad\text{sodass}\quad
\alpha = \alpha_{\wedge}\circ \wedge_{k},
\]
wobei $\wedge_{k}:\mathrm{V}^{k}\to \bigwedge^{k}\mathrm{V}$ der kanonische Alternator ist.

\textbf{Spezialfall:}  
Für $\mathrm{W}=\mathbb{R}$ ergibt sich der natürliche Isomorphismus
\[
\boxed{\;
L_{\mathrm{alt}}^{k}(\mathrm{V})
\;\cong\;
\left(\bigwedge\nolimits^{k}\mathrm{V}\right)^{*},
\;}
\]
das heißt: der Raum der alternierenden $k$-linearen Formen auf $\mathrm{V}$
ist der Dualraum der $k$-ten äußeren Potenz von~$\mathrm{V}$.
\end{THEO}




We would now like to embed $\Lambda^{k} \mathrm{V}^{*}$ into $\bigotimes^{k} \mathrm{V}^{*}$, and this involves a choice. For each $k$, let $A_{k}: \mathrm{V}^{*} \times \cdots \times \mathrm{V}^{*} \rightarrow \bigotimes^{k} \mathrm{V}^{*}$ be defined by

$$
A_{k}\left(\alpha_{1}, \ldots, \alpha_{k}\right):=\sum_{\sigma} \operatorname{sgn}(\sigma) \alpha_{\sigma_{1}} \otimes \cdots \otimes \alpha_{\sigma_{k}}
$$

By the universal property of $\bigwedge^{k} \mathrm{V}^{*}$ we obtain an induced map

$$
\widetilde{A_{k}}: \bigwedge^{k} \mathrm{V}^{*} \rightarrow \bigotimes^{k} \mathrm{V}^{*}
$$

If we identify $\bigotimes^{k} \mathrm{V}^{*}$ with $T_{k}(\mathrm{V})$, then we get a map $\widetilde{A_{k}}: \bigwedge^{k} \mathrm{V}^{*} \rightarrow T_{k}(\mathrm{V})$.\\
Proposition 8.23. The image of the map $\widetilde{A}_{k}: \bigwedge^{k} \mathrm{V}^{*} \rightarrow T_{k}(\mathrm{V})$ is a linear isomorphism with image equal to $L_{\text {alt }}^{k}(\mathrm{V})$ such that

$$
\widetilde{A_{k}}\left(\alpha_{1} \wedge \cdots \wedge \alpha_{k}\right)\left(v_{1}, \ldots, v_{k}\right)=\operatorname{det}\left(\alpha_{i}\left(v_{j}\right)\right)
$$

Proof. We leave the proof as Problem 5.\\
We combine these maps to obtain a linear isomorphism

$$
\widetilde{A}: \bigwedge \mathrm{V}^{*} \rightarrow L_{\mathrm{alt}}(\mathrm{V})
$$

Now both $L_{\text {alt }}(\mathrm{V})$ and $\Lambda \mathrm{V}^{*}$ have already independently been given the exterior algebra structures via their respective wedge products. One may check that $\widetilde{A}$ has been defined in such a way as to be an isomorphism of these algebras:

$$
\bigwedge \mathrm{V}^{*} \cong L_{\text {alt }}(\mathrm{V}) \text { (as exterior algebras). }
$$

Also notice that

$$
\bigwedge^{k} \mathrm{V}^{*} \cong L_{\mathrm{alt}}^{k}(\mathrm{V}) \cong\left(\bigwedge^{k} \mathrm{V}\right)^{*}
$$

which allows us to think of $\Lambda^{k} \mathrm{V}^{*}$ as dual to $\Lambda^{k} \mathrm{V}$ in such a way that

$$
\left(\alpha_{1} \wedge \cdots \wedge \alpha_{k}\right)\left(v_{1} \wedge \cdots \wedge v_{k}\right)=\operatorname{det}\left(\alpha_{i}\left(v_{j}\right)\right)
$$

Remark 8.24 (An identification). In what follows, we will identify $\wedge^{k} \mathrm{V}^{*}$ with $L_{\text {alt }}^{k}(\mathrm{V})$ and hence $\wedge \mathrm{V}^{*}$ with $L_{\text {alt }}(\mathrm{V})$ whenever convenient. In other words, we freely treat elements of $\Lambda^{k} \mathrm{V}^{*}$ as alternating multilinear forms.




\begin{THEO}{GA-L09-06-40}{Isomorphismus zwischen äußerer Potenz und Raum der alternierenden Formen}
Sei $\mathrm{V}$ ein endlichdimensionaler reeller Vektorraum. Für jedes $k\ge 1$ definiere die Abbildung
\[
A_{k}:\mathrm{V}^{*}\times\cdots\times\mathrm{V}^{*}\longrightarrow
\bigotimes\nolimits^{k}\mathrm{V}^{*}
\]
durch
\[
A_{k}(\alpha_{1},\ldots,\alpha_{k})
:=\sum_{\sigma\in S_{k}}\operatorname{sgn}(\sigma)\,
\alpha_{\sigma(1)}\otimes\cdots\otimes\alpha_{\sigma(k)}.
\]
Diese Abbildung ist $k$-multilinear und alternierend, und nach der universellen Eigenschaft der äußeren Potenz
existiert genau eine lineare Abbildung
\[
\widetilde{A_{k}}:\bigwedge\nolimits^{k}\mathrm{V}^{*}\longrightarrow
\bigotimes\nolimits^{k}\mathrm{V}^{*}\cong T_{k}(\mathrm{V})
\]
mit $\widetilde{A_{k}}\circ\wedge_{k}=A_{k}$.

\textbf{(1) Charakterisierung der Abbildung.}  
Für einfache Elemente gilt
\[
\boxed{\;
\widetilde{A_{k}}(\alpha_{1}\wedge\cdots\wedge\alpha_{k})(v_{1},\ldots,v_{k})
=\det\!\big(\alpha_{i}(v_{j})\big),
\;}
\]
wobei $v_{1},\ldots,v_{k}\in\mathrm{V}$.

\textbf{(2) Strukturresultat.}  
Die Abbildung $\widetilde{A_{k}}$ ist ein linearer Isomorphismus mit Bild
\[
\operatorname{im}(\widetilde{A_{k}})=L_{\mathrm{alt}}^{k}(\mathrm{V}),
\]
dem Raum aller alternierenden $k$-multilinearen Abbildungen $\mathrm{V}^{k}\to\mathbb{R}$.

\textbf{(3) Globale Version.}  
Die Abbildungen $\widetilde{A_{k}}$ kombinieren zu einem Algebraisomorphismus
\[
\widetilde{A}:\bigwedge\nolimits \mathrm{V}^{*}\;\longrightarrow\;
L_{\mathrm{alt}}(\mathrm{V}),
\]
der die jeweilige äußere Produktstruktur erhält, also
\[
\boxed{\;
\bigwedge \mathrm{V}^{*}\;\cong\;L_{\mathrm{alt}}(\mathrm{V})
\quad\text{(als Außenalgebren)}.
\;}
\]

\textbf{(4) Dualitätsidentifikation.}  
Insbesondere gilt der natürliche Isomorphismus
\[
\bigwedge\nolimits^{k}\mathrm{V}^{*}\cong
L_{\mathrm{alt}}^{k}(\mathrm{V})\cong
\left(\bigwedge\nolimits^{k}\mathrm{V}\right)^{*},
\]
wobei die Dualität durch die Beziehung
\[
\boxed{\;
(\alpha_{1}\wedge\cdots\wedge\alpha_{k})
(v_{1}\wedge\cdots\wedge v_{k})
=\det\!\big(\alpha_{i}(v_{j})\big)
\;}
\]
gegeben ist.

\textbf{(5) Bemerkung.}
Im Folgenden wird diese Identifikation stillschweigend verwendet:
Elemente von $\Lambda^{k}\mathrm{V}^{*}$ werden als alternierende $k$-Formen,
d.\,h. als Elemente von $L_{\mathrm{alt}}^{k}(\mathrm{V})$, aufgefasst.
\end{THEO}


\pagebreak


\subsection*{Differential Forms}

Let $M$ be an $n$-manifold. We now bundle together the various spaces $L_{\text {alt }}^{k}\left(T_{p} M\right)$. That is, we form the natural bundle $L_{\text {alt }}^{k}(T M)$ that has as its fiber at $p$, the space $L_{\text {alt }}^{k}\left(T_{p} M\right)$. Thus $L_{\text {alt }}^{k}(T M)=\bigcup_{p \in M} L_{\text {alt }}^{k}\left(T_{p} M\right)$.\\
Exercise 8.25. Exhibit the smooth structure and vector bundle structure on $L_{\text {alt }}^{k}(T M)=\bigcup_{p \in M} L_{\text {alt }}^{k}\left(T_{p} M\right)$. [Hint: Let ( $U, \mathrm{x}$ ) be a chart of $M$ and $x^{1}, \ldots, x^{n}$ the coordinate functions. Let $\widetilde{U}:=\bigcup_{p \in U} L_{\text {alt }}^{k}\left(T_{p} M\right)$ and let $d= \binom{n}{k}$. Then we have a map $\widetilde{U} \rightarrow U \times \mathbb{R}^{d}$ given by $\alpha_{p} \mapsto(p, a)$ where $a$ is the $d$-tuple of components (in some fixed order) of $\alpha_{p}$ given by its local coordinate representation.]




\begin{DEF}{GA-L09-06-41}{Vektorbündelstruktur auf dem Bündel alternierender $k$-Formen}
Sei $M$ eine $n$-dimensionale glatte Mannigfaltigkeit und $0\le k\le n$. Setze
\[
L_{\mathrm{alt}}^{k}(TM)\;:=\;\bigsqcup_{p\in M} L_{\mathrm{alt}}^{k}(T_{p}M),
\qquad
\pi: L_{\mathrm{alt}}^{k}(TM)\to M,\ \alpha_{p}\mapsto p .
\]
Dann trägt $L_{\mathrm{alt}}^{k}(TM)$ eine natürliche glatte Vektor\-bündelstruktur vom Rang
\(
d=\binom{n}{k}
\),
wobei die Faser über $p$ gleich $L_{\mathrm{alt}}^{k}(T_{p}M)$ ist. Diese Struktur ist durch folgende lokale Trivialisierungen gegeben.


\medskip
\textbf{Konstruktion der lokalen Trivialisierungen.}
Sei $(U,\mathbf{x})=(U,(x^{1},\ldots,x^{n}))$ eine Karte von $M$.
Schreibe $\widetilde U:=\bigsqcup_{p\in U}L_{\mathrm{alt}}^{k}(T_{p}M)$.
Für $p\in U$ sei $(\partial_{i}|_{p})_{i=1}^{n}$ die Koordinatenbasis von $T_{p}M$ und
\(
(dx^{i}|_{p})_{i=1}^{n}
\)
die duale Basis von $T_{p}^{*}M$. Für jeden streng geordneten Mehrindex $\vec I=(i_{1}<\cdots<i_{k})$ definiere die lokale Basis
\(
dx^{\vec I}|_{p}:=dx^{i_{1}}|_{p}\wedge\cdots\wedge dx^{i_{k}}|_{p}
\)
von $L_{\mathrm{alt}}^{k}(T_{p}M)\cong \Lambda^{k}T_{p}^{*}M$.
Jedes $\alpha_{p}\in L_{\mathrm{alt}}^{k}(T_{p}M)$ besitzt die eindeutige Entwicklung
\[
\alpha_{p}=\sum_{\vec I} a_{\vec I}(p)\, dx^{\vec I}|_{p},
\qquad \vec I\in\mathcal{I}_{k}:=\{(i_{1}<\cdots<i_{k})\subset\{1,\dots,n\}\}.
\]
Definiere
\[
\Phi_{U}:\ \widetilde U\longrightarrow U\times\mathbb{R}^{d},
\quad
\alpha_{p}\longmapsto\Big(p,\ (a_{\vec I}(p))_{\vec I\in\mathcal{I}_{k}}\Big).
\]
Dann ist $\Phi_{U}$ ein Faser\-weise linearer Bijektor (mit inverser Abbildung
$(p,(c_{\vec I}))\mapsto\sum_{\vec I} c_{\vec I}\,dx^{\vec I}|_{p}$)
und liefert eine Bündeltrivialisierung über $U$:
\(
\Phi_{U}:\pi^{-1}(U)\overset{\cong}{\longrightarrow} U\times\mathbb{R}^{d}.
\)

\medskip
\textbf{Überlappungsabbildungen.}
Seien $(U,\mathbf{x})$ und $(V,\overline{\mathbf{x}})$ zwei Karten mit $U\cap V\neq\emptyset$.
Bezeichne die zugehörigen Trivialisierungen mit $\Phi_{U}$ bzw. $\Phi_{V}$.
Die Übergangsfunktion
\[
\Phi_{V}\circ\Phi_{U}^{-1}:\ (U\cap V)\times\mathbb{R}^{d}\longrightarrow (U\cap V)\times\mathbb{R}^{d}
\]
hat die Form
\[
\big(p,(a_{\vec I})\big)\longmapsto \Big(p,\ \big(\bar a_{\vec J}\big)_{\vec J\in\mathcal{I}_{k}}\Big),
\]
wobei die Koeffizienten $\bar a_{\vec J}$ glatt in $p$ von den $a_{\vec I}$ durch
\emph{exteriore Potenz der Jacobi\-matrix} gegeben sind:
\[
\boxed{\;
\bar a_{\vec J}
\;=\;
\sum_{\vec I\in\mathcal{I}_{k}}
\big(\wedge^{k}J(\overline{\mathbf{x}}\circ\mathbf{x}^{-1})\big)_{\vec J}{}^{\ \vec I}(p)\; a_{\vec I}.
\;}
\]
Hier ist $J(\overline{\mathbf{x}}\circ\mathbf{x}^{-1})(p)=(\partial \bar x^{\mu}/\partial x^{\nu})(p)$
die $n\times n$-Jacobi\-matrix und
\(
\wedge^{k}J
\)
die $d\times d$-Matrix der $k\times k$-Unterminoren (d.\,h. die Darstellung der $k$-ten äußeren Potenz in den Basen
$(dx^{\vec I})$ bzw. $(d\bar x^{\vec J})$). Äquivalent: Die Komponenten transformieren nach
\[
d\bar x^{\vec J}
=\sum_{\vec I}
\big(\wedge^{k}J\big)_{\vec J}{}^{\ \vec I}\, dx^{\vec I},
\qquad
\bar a_{\vec J}
=\sum_{\vec I}
\big(\wedge^{k}J\big)_{\vec J}{}^{\ \vec I}\, a_{\vec I}.
\]

\medskip
\textbf{Schluss.}
Die Abbildungen $\{\Phi_{U}\}$ definieren einen glatten Vektor\-bündelatlas
mit glatten, linear-faserweisen Übergangsfunktionen, die die Kokzykelbedingungen erfüllen.
Damit ist $\pi:L_{\mathrm{alt}}^{k}(TM)\to M$ ein glattes Vektor\-bündel vom Rang $d=\binom{n}{k}$.
Über jedem $U$ identifiziert $\Phi_{U}$ das Bündel kanonisch mit $U\times\mathbb{R}^{d}$.

\medskip
\textbf{Notation.} Üblicherweise schreibt man $L_{\mathrm{alt}}^{k}(TM)\cong \Lambda^{k}T^{*}M$ und spricht vom
\emph{Bündel der $k$-Formen}. Seine glatten Schnitte sind die \emph{Differentialformen} vom Grad $k$ auf $M$.
\end{DEF}




Let the smooth sections of this bundle be denoted by

$$
\Omega^{k}(M)=\Gamma\left(M ; L_{\mathrm{alt}}^{k}(T M)\right),
$$

and sections over $U \subset M$ by $\Omega_{M}^{k}(U)$.

Definition 8.26. Elements of $\Omega^{k}(M)$ are called differential $k$-forms or just $k$-forms.

The space $\Omega^{k}(M)$ is a module over the algebra of smooth functions $C^{\infty}(M)=\mathcal{F}(M)$. If $n=\operatorname{dim} M$ then we have the direct sum

$$
\Omega(M)=\bigoplus_{k=0}^{n} \Omega^{k}(M),
$$

with a similar decomposition for any open $U \subset M$.


\begin{DEF}{GA-L09-06-42}{Differentialformen auf einer Mannigfaltigkeit}
Sei $M$ eine glatte $n$-dimensionale Mannigfaltigkeit und
\[
L_{\mathrm{alt}}^{k}(TM)=\bigsqcup_{p\in M} L_{\mathrm{alt}}^{k}(T_{p}M)
\]
das in Proposition~\ref{DFM-L08-25} konstruierte Bündel der alternierenden $k$-Formen (auch \emph{Bündel der $k$-Kovektoren} oder \emph{Bündel der $k$-Formen}).  

\textbf{Definition.}  
Die Menge der glatten Schnitte dieses Bündels wird bezeichnet mit
\[
\boxed{\;
\Omega^{k}(M):=\Gamma\!\left(M;L_{\mathrm{alt}}^{k}(TM)\right),
\;}
\]
und ihre Elemente heißen \emph{differentiale $k$-Formen} oder einfach \emph{$k$-Formen auf $M$}.
Für eine offene Teilmenge $U\subset M$ schreiben wir entsprechend
\[
\Omega^{k}_{M}(U):=\Gamma\!\left(U;L_{\mathrm{alt}}^{k}(TM)\right)
\]
für die $k$-Formen, die auf $U$ definiert sind.

\medskip
\textbf{Eigenschaften.}
\begin{enumerate}
\item $\Omega^{k}(M)$ ist ein Modul über der Algebra der glatten Funktionen
\[
C^{\infty}(M)=\mathcal{F}(M),
\qquad
(f\omega)(p):=f(p)\,\omega(p)
\quad\text{für }f\in C^{\infty}(M),\ \omega\in\Omega^{k}(M).
\]
\item Ist $n=\dim M$, so definiert die direkte Summe
\[
\boxed{\;
\Omega(M):=\bigoplus_{k=0}^{n}\Omega^{k}(M)
\;}
\]
den Raum aller glatten Differentialformen auf $M$ beliebigen Grades.
Analog gilt für jedes offene $U\subset M$
\[
\Omega(U)=\bigoplus_{k=0}^{n}\Omega^{k}_{M}(U).
\]
\end{enumerate}

\textbf{Bemerkung.}
$\Omega^{0}(M)$ ist identisch mit $C^{\infty}(M)$, während $\Omega^{n}(M)$ die glatten
Volumenformen auf $M$ beschreibt.
\end{DEF}




Exercise 8.27. Show that 

\begin{PROP}{GA-L09-06-43}{$ \bigoplus_{k=0}^{n} \Omega^{k}(M) \cong \Gamma\left(\bigoplus_{k=0}^{n} L_{\text {alt }}^{k}(T M)\right)$}
There is a module isomorphism
$$
\bigoplus_{k=0}^{n} \Omega^{k}(M) \cong \Gamma\left(\bigoplus_{k=0}^{n} L_{\text {alt }}^{k}(T M)\right) .
$$
\end{PROP}

Definition 8.28. Let $M$ be an $n$-manifold. The elements of $\Omega(M)$ are called differential forms on $M$. We identify $\Omega^{k}(M)$ with the obvious subspace of $\Omega(M)=\bigoplus_{k} \Omega^{k}(M)$. A differential form in $\Omega^{k}(M)$ is said to be homogeneous of degree $k$.

\begin{DEF}{GA-L09-06-44}{Darstellung von Differentialformen durch Summen von homogenen Komponenten}
If $\omega \in \Omega(M)$, then we can uniquely write $\omega=\sum_{k=1}^{n} \omega_{k}$ where $\omega_{k} \in \Omega^{k}(M)$ and the $\omega_{k}$ are called homogeneous components of $\omega$.
\end{DEF}

Definition 8.29. 

For $\omega \in \Omega^{k_{1}}(M)$, and $\eta \in \Omega^{k_{2}}(M)$, we define the exterior product $\omega \wedge \eta \in \Omega^{k_{1}+k_{2}}(M)$ by
$$
(\omega \wedge \eta)(p):=\omega(p) \wedge \eta(p) .
$$

It is easy to see that $\Omega(M)$ is a ring under the wedge product and, in fact, a $C^{\infty}(M)$-algebra. Whenever convenient, we may extend this to a sum over all $n \in \mathbb{Z}$ by defining (as before) $\Omega^{k}(M):=0$ for $k<0$ and $\Omega^{k}(M):=0$ if $k>\operatorname{dim}(M)$. We have made the trivial extension of $\wedge$ to a $\mathbb{Z}$-graded algebra by declaring that $\omega \wedge \eta=0$ if either $\eta$ or $\omega$ is homogeneous of negative degree. Thus $\Omega(M)$ is a graded algebra and is said to be graded commutative because $\alpha \wedge \beta=(-1)^{k \ell} \beta \wedge \alpha$ for $\alpha \in \Omega^{k}(M)$ and $\beta \in \Omega^{\ell}(M)$. Sometimes one see this notion referred to as skew-commutativity.



\begin{DEF}{GA-L09-06-45}{Äußeres Produkt und Gradstruktur auf $\Omega(M)$}
Sei $M$ eine glatte Mannigfaltigkeit der Dimension $n$.

\textbf{(1) Definition des äußeren Produkts.}  
Für $\omega\in\Omega^{k_{1}}(M)$ und $\eta\in\Omega^{k_{2}}(M)$ ist das
\emph{äußere Produkt} oder \emph{Wedge-Produkt}
\[
\boxed{\;
\omega\wedge\eta\in\Omega^{k_{1}+k_{2}}(M)
\;}
\]
definiert durch die punktweise Operation
\[
(\omega\wedge\eta)(p):=\omega(p)\wedge\eta(p),
\qquad p\in M,
\]
wobei rechts das äußere Produkt der alternierenden Formen
$\omega(p)\in L_{\mathrm{alt}}^{k_{1}}(T_{p}M)$ und
$\eta(p)\in L_{\mathrm{alt}}^{k_{2}}(T_{p}M)$ steht.

\medskip
\textbf{(2) Algebraische Struktur.}
\begin{itemize}
\item Das Produkt $\wedge$ ist bilinear über $C^{\infty}(M)$, d.\,h.
\[
(f\omega)\wedge\eta=f(\omega\wedge\eta)=\omega\wedge(f\eta),
\qquad f\in C^{\infty}(M).
\]
\item Für $\omega\in\Omega^{k}(M)$, $\eta\in\Omega^{\ell}(M)$ gilt die
\emph{graded (skew-)Kommutativität}
\[
\boxed{\;
\omega\wedge\eta=(-1)^{k\ell}\eta\wedge\omega.
\;}
\]
\item Das neutrale Element ist $1\in\Omega^{0}(M)=C^{\infty}(M)$, und für jedes $k>n$ ist $\Omega^{k}(M)=0$.
\end{itemize}

\medskip
\textbf{(3) Erweiterte Gradierung.}
Zur Vereinheitlichung der algebraischen Struktur kann man
$\Omega^{k}(M):=0$ für $k<0$ und $k>n$ setzen.
Damit wird
\[
\boxed{\;
\Omega(M)=\bigoplus_{k\in\mathbb{Z}}\Omega^{k}(M)
\;}
\]
eine $\mathbb{Z}$-gradierte, kommutative $C^{\infty}(M)$-Algebra, 
genannt \emph{externe Algebra der Differentialformen} auf $M$.

\medskip
\textbf{Bemerkung.}
Die Gradkommutativität ersetzt hier die gewöhnliche Kommutativität:
für ungerade Grade wechselt das Vorzeichen beim Vertauschen.
Diese Struktur ist grundlegend für die Definition des äußeren Differentialoperators und der de-Rham-Kohomologie.
\end{DEF}


Just as a tangent vector is the infinitesimal version of a (parametrized) curve through a point $p \in M$, so a covector at $p \in M$ is the infinitesimal version of a function defined near $p$. At this point one must be careful. It is true that for any single covector $\alpha_{p} \in T_{p} M$ there always exists a function $f$ such $d f_{p}=\alpha_{p}$. But as we saw in Chapter 2, if $\alpha \in \Omega^{1}(M)$, then it is not necessarily true that there is a function $f \in C^{\infty}(M)$ such that $d f=\alpha$. If $f_{1}, f_{2}, \ldots, f_{k}$ are smooth functions, then one way to picture $d f_{1} \wedge \cdots \wedge d f_{k}$ is by thinking of the intersecting family of level sets of the functions $f_{1}, f_{2}, \ldots, f_{k}$, which in some cases can be pictured as a sort of "egg crate", structure. For

\begin{figure}[h]
\begin{center}
  \includegraphics[scale=0.25]{2025_10_09_265d7e1a22c89f79c21eg-375}

\caption{Figure 8.1. 2-form as "flux tubes"}
\end{center}
\end{figure}

a 2 -form in a 3 -manifold, one obtains "flux tubes" as shown in Figure 8.1. The infinitesimal version of this is a sort of straightened out "linear egg crate structure," which may be thought of as existing in the tangent space at a point. This is the rough intuition for $\left.\left.d f_{1}\right|_{p} \wedge \cdots \wedge d f_{k}\right|_{p}$, and the $k$-form $d f_{1} \wedge \cdots \wedge d f_{k}$ is a field of such structures which somehow fit the level sets of the family $f_{1}, f_{2}, \ldots, f_{k}$. Of course, $d f_{1} \wedge \cdots \wedge d f_{k}$ is a very special kind of $k$-form. In general, a $k$-form over $U$ may not arise from a family of functions.

In local coordinates, calculation is often quite easy and formal. For example, in $\mathbb{R}^{3}$ with standard coordinates $x, y, z$, a simple wedge product calculation is as follows:

$$
\begin{aligned}
(x y d x+ & z d y+d z) \wedge(x d y+z d z) \\
= & x y d x \wedge x d y+x y d x \wedge z d z+z d y \wedge x d y \\
& +z d y \wedge z d z+d z \wedge x d y+d z \wedge z d z \\
= & x^{2} y d x \wedge d y+x y z d x \wedge d z+z^{2} d y \wedge d z+x d z \wedge d y \\
= & x^{2} y d x \wedge d y+x y z d x \wedge d z+\left(z^{2}-x\right) d y \wedge d z
\end{aligned}
$$

An equally trivial calculation shows that

$$
\left(x y z^{2} d x \wedge d y+d y \wedge d z\right) \wedge(d x+x d y+z d z)=\left(x y z^{3}+1\right) d x \wedge d y \wedge d z
$$




\begin{CONC}{GA-L09-06-46}{Intuitive Bedeutung von Differentialformen und lokale Berechnungen}
\textbf{(1) Geometrische Intuition.}  
Wie ein Tangentialvektor das \emph{infinitesimale Analogon einer Kurve} durch einen Punkt $p\in M$ ist, so kann man ein \emph{Kovektorfeld} (eine 1-Form) als das infinitesimale Analogon einer glatten Funktion in einer Umgebung von $p$ auffassen.  
Für jedes $\alpha_{p}\in T_{p}^{*}M$ existiert stets eine Funktion $f$ mit $df_{p}=\alpha_{p}$, aber im Allgemeinen gilt: Ist $\alpha\in\Omega^{1}(M)$, so existiert \emph{nicht notwendigerweise} ein $f\in C^{\infty}(M)$ mit $df=\alpha$ — die Form muss dazu \emph{exakt} sein.

Sind $f_{1},\ldots,f_{k}\in C^{\infty}(M)$ glatte Funktionen, so kann man die $k$-Form
\[
df_{1}\wedge\cdots\wedge df_{k}
\]
heuristisch als infinitesimale Beschreibung der Orientierung und „Dichte“ der gemeinsamen Niveaumengen der $f_{i}$ verstehen.  
Diese bilden lokal ein „Gitter“ aus Schnittebenen (eine „egg-crate“-Struktur).  
Für $k=2$ in einer 3-dimensionalen Mannigfaltigkeit lässt sich dies als ein System von \emph{Flusstuben} darstellen, siehe Abbildung~\ref{fig:flux2form}.  
Das infinitesimale Analogon ist die linearisierte Struktur
\[
(df_{1}|_{p})\wedge\cdots\wedge(df_{k}|_{p})\in\Lambda^{k}T_{p}^{*}M,
\]
und die $k$-Form $df_{1}\wedge\cdots\wedge df_{k}$ ist ein Feld solcher Strukturen, die glatt über $M$ variieren.
$$
\includegraphics[scale=0.25]{2025_10_09_265d7e1a22c89f79c21eg-375}
$$
\medskip
\textbf{(2) Lokale Berechnungen.}  
In lokalen Koordinaten sind Rechnungen mit Differentialformen formal und oft einfach.
Beispielsweise in $\mathbb{R}^{3}$ mit Standardkoordinaten $(x,y,z)$ gilt
\[
\begin{aligned}
(xy\,dx+z\,dy+dz)\wedge(x\,dy+z\,dz)
&=xy\,dx\wedge x\,dy+xy\,dx\wedge z\,dz+z\,dy\wedge x\,dy\\
&\quad +z\,dy\wedge z\,dz+dz\wedge x\,dy+dz\wedge z\,dz\\
&=x^{2}y\,dx\wedge dy+xyz\,dx\wedge dz+(z^{2}-x)\,dy\wedge dz.
\end{aligned}
\]
Ebenso erhält man
\[
\big(xyz^{2}\,dx\wedge dy+dy\wedge dz\big)\wedge(dx+x\,dy+z\,dz)
=(xyz^{3}+1)\,dx\wedge dy\wedge dz.
\]

\medskip
\textbf{(3) Interpretation.}  
Solche lokalen Rechnungen verdeutlichen:
\begin{itemize}
\item Wedge-Produkte kodieren die orientierte Volumendichte (z.\,B.\ $dx\wedge dy\wedge dz$ als orientiertes Volumenelement);
\item Koordinatenfunktionen und deren Differentiale spannen lokal das Bündel der Differentialformen auf;
\item Die Operationen mit $\wedge$ gehorchen den antisymmetrischen Rechenregeln:
$dx^{i}\wedge dx^{j}=-dx^{j}\wedge dx^{i}$, $dx^{i}\wedge dx^{i}=0$.
\end{itemize}
Diese Formalismen sind die Grundlage für die spätere Einführung des äußeren Differentials $d:\Omega^{k}(M)\to\Omega^{k+1}(M)$.
\end{CONC}






8.2.1. Pull-back of a differential form. 


Since we treat differential forms as alternating covariant tensor fields, we have a notion of pull-back already defined. It is easy to see that the pull-back of an alternating tensor field is also an alternating tensor field, and so given any smooth map $f: M \rightarrow N$,\\
we get a map $f^{*}: \Omega^{k}(N) \rightarrow \Omega^{k}(M)$. We recall here the definition:

$$
\left(f^{*} \eta\right)(p)\left(v_{1}, \ldots, v_{k}\right)=\eta_{f(p)}\left(T_{p} f \cdot v_{1}, \ldots, T_{p} f \cdot v_{k}\right),
$$

for tangent vectors $v_{1}, \ldots, v_{k} \in T_{p} M$. The pull-back extends in the obvious way to a map $f^{*}: \Omega(N) \rightarrow \Omega(M)$.

Proposition 8.30. Let $f: M \rightarrow N$ be a smooth map and let $\eta_{1}, \eta_{2} \in \Omega(N)$. Then we have

$$
f^{*}\left(\eta_{1} \wedge \eta_{2}\right)=f^{*} \eta_{1} \wedge f^{*} \eta_{2}
$$

Proof. This follows directly from Proposition 8.14.\\
Proposition 8.31. Let $f: M \rightarrow N$ and $g: N \rightarrow P$ be smooth maps. Then for any smooth differential form $\eta \in \Omega(P)$ we have $(f \circ g)^{*} \eta=g^{*}\left(f^{*} \eta\right)$. Thus $(f \circ g)^{*}=g^{*} \circ f^{*}$.

Proof. We prove only the case $\eta \in \Omega^{1}(N)$. The general case is entirely similar and is left to the reader. For $v \in T_{p} M$, we have

$$
\begin{aligned}
(f \circ g)^{*} \eta(v) & =\eta(T(f \circ g) v)=\eta(T f \circ T g(v)) \\
& =f^{*} \eta(T g \cdot v)=g^{*}\left(f^{*} \eta\right)(v),
\end{aligned}
$$

which completes the proof for the considered case.

From the above propositions we see that we have a contravariant functor from the category of smooth manifolds and smooth maps to the category of rings which assigns to each smooth manifold $M$ the space $\Omega(M)$ and to each smooth map $f$ the ring homomorphism $f^{*}$.

In case $S$ is a regular submanifold of $M$, we have the inclusion map $\iota: S \hookrightarrow M$, which maps $p \in S$ to the very same point $p \in M$. As mentioned before, it is natural to identify $T_{p} S$ with $T_{p} \iota\left(T_{p} S\right)$ for any $p \in S$. In other words, we often do not distinguish between a vector $v_{p}$ and $T_{p} \iota\left(v_{p}\right)$. Thus we view the tangent bundle of $S$ as a subset of $T M$. With this in mind we must realize that for any $\alpha \in \Omega^{k}(M)$ the form $\iota^{*} \alpha$ is just the restriction of $\alpha$ to vectors tangent to $S$. In particular, if $U \subset M$ is open and $\iota: U \hookrightarrow M$, then $\iota^{*} \alpha=\left.\alpha\right|_{U}$.




\begin{DEF}{GA-L09-06-47}{Pull-Back von Differentialformen und Funktorialität}
\textbf{(1) Definition des Pull-Backs.}  
Seien $M,N$ glatte Mannigfaltigkeiten und $f:M\to N$ eine glatte Abbildung.
Für eine $k$-Form $\eta\in\Omega^{k}(N)$ definieren wir den \emph{Pull-Back}
\[
\boxed{\;
(f^{*}\eta)(p)(v_{1},\ldots,v_{k})
:=\eta_{f(p)}\big(T_{p}f\cdot v_{1},\ldots,T_{p}f\cdot v_{k}\big),
\qquad v_{i}\in T_{p}M.
\;}
\]
Da $f^{*}\eta$ wieder alternierend in den $v_{i}$ ist, gilt
$f^{*}\eta\in\Omega^{k}(M)$.
Somit erhalten wir eine wohldefinierte Abbildung
\[
f^{*}:\Omega^{k}(N)\longrightarrow \Omega^{k}(M),
\qquad
\text{bzw. insgesamt } f^{*}:\Omega(N)\to\Omega(M).
\]

\medskip
\textbf{(2) Produktverträglichkeit.}
Für $\eta_{1}\in\Omega^{k_{1}}(N)$ und $\eta_{2}\in\Omega^{k_{2}}(N)$ gilt
\[
\boxed{\;
f^{*}(\eta_{1}\wedge\eta_{2})
=f^{*}\eta_{1}\wedge f^{*}\eta_{2}.
\;}
\]
\emph{Beweis:} folgt direkt aus der Punktweise-Definition und der analogen Aussage
für alternierende Tensorfelder (vgl.~Proposition~7.14).

\medskip
\textbf{(3) Funktorialität.}
Für glatte Abbildungen
$f:M\to N$ und $g:N\to P$
und jede $\eta\in\Omega(P)$ gilt
\[
\boxed{\;
(f\circ g)^{*}\eta
=g^{*}\big(f^{*}\eta\big),
\qquad\text{d.\,h.}\quad
(f\circ g)^{*}=g^{*}\circ f^{*}.
\;}
\]
\emph{Beweis:}
Für $\eta\in\Omega^{1}(P)$ und $v\in T_{p}M$ gilt
\[
(f\circ g)^{*}\eta(v)
=\eta(T(f\circ g)v)
=\eta(Tf(Tg(v)))
=f^{*}\eta(Tg(v))
=g^{*}(f^{*}\eta)(v).
\]
Der allgemeine Fall folgt durch Multilinearität.

\medskip
\textbf{(4) Funktorielle Interpretation.}
Die Zuweisung
\[
M\longmapsto\Omega(M),
\qquad
f\longmapsto f^{*},
\]
definiert einen \emph{kontravarianten Funktor}
von der Kategorie der glatten Mannigfaltigkeiten und glatten Abbildungen
in die Kategorie der kommutativen $C^{\infty}$-Algebren.

\medskip
\textbf{(5) Spezieller Fall: Inklusion einer Untermannigfaltigkeit.}
Ist $S\subset M$ eine reguläre Untermannigfaltigkeit
mit Inklusion $\iota:S\hookrightarrow M$, so gilt für $\alpha\in\Omega^{k}(M)$:
\[
\boxed{\;
(\iota^{*}\alpha)(p)
=\alpha_{p}\big|_{T_{p}S},
\qquad p\in S.
\;}
\]
Das heißt, $\iota^{*}\alpha$ ist einfach die \emph{Restriktion}
von $\alpha$ auf die Tangentialvektoren von $S$.
Im Spezialfall einer offenen Menge $U\subset M$ und
$\iota:U\hookrightarrow M$ gilt insbesondere:
\[
\iota^{*}\alpha=\alpha|_{U}.
\]
\end{DEF}



\begin{LEM}{GA-L09-06-48}{Lokale Darstellung des Pullbacks von Differentialformen}
Seien $(U,\mathrm{x})$ eine Karte von $M$ mit Koordinaten $(x^{1},\dots,x^{m})$ und
$(V,\mathrm{y})$ eine Karte von $N$ mit Koordinaten $(y^{1},\dots,y^{n})$, so dass
$f(U)\subset V$. Sei $\eta\in\Omega^{k}(N)$ mit lokaler Darstellung
\[
\eta=\sum_{\vec{J}} b_{\vec{J}}\; d y^{\vec{J}}
\qquad\text{auf }V,
\]
wobei die Summe über strikt aufsteigende Multiindizes
$\vec{J}=(j_{1}<\cdots<j_{k})$ läuft und $d y^{\vec{J}}:=d y^{j_{1}}\wedge\cdots\wedge d y^{j_{k}}$.
Dann gilt auf $U$ die lokale Formel für den Pullback
\[
f^{*}\eta
= \sum_{\vec{J}}\sum_{\vec{L}}
\bigl(b_{\vec{J}}\circ f\bigr)\;
\frac{\partial y^{\vec{J}}}{\partial x^{\vec{L}}}\;
d x^{\vec{L}},
\]
wobei $\vec{L}=(\ell_{1}<\cdots<\ell_{k})$ und $d x^{\vec{L}}:=d x^{\ell_{1}}\wedge\cdots\wedge d x^{\ell_{k}}$, sowie
\[
\frac{\partial y^{\vec{J}}}{\partial x^{\vec{L}}}
:=\frac{\partial\!\left(y^{j_{1}},\dots,y^{j_{k}}\right)}
{\partial\!\left(x^{\ell_{1}},\dots,x^{\ell_{k}}\right)}
=\det\!\begin{bmatrix}
\frac{\partial y^{j_{1}}}{\partial x^{\ell_{1}}} & \cdots & \frac{\partial y^{j_{1}}}{\partial x^{\ell_{k}}}\\
\vdots & & \vdots\\
\frac{\partial y^{j_{k}}}{\partial x^{\ell_{1}}} & \cdots & \frac{\partial y^{j_{k}}}{\partial x^{\ell_{k}}}
\end{bmatrix}\!\circ f.
\]

\textbf{Äquivalente Schreibweise mit Levi-Civita-Symbol.}
Schreibt man zunächst
\(
f^{*}\eta=\sum_{\vec{J}}(b_{\vec{J}}\circ f)\,
\bigl(\sum_{i_{1}}\frac{\partial y^{j_{1}}}{\partial x^{i_{1}}} d x^{i_{1}}\bigr)\wedge\cdots\wedge
\bigl(\sum_{i_{k}}\frac{\partial y^{j_{k}}}{\partial x^{i_{k}}} d x^{i_{k}}\bigr),
\)
so erhält man nach Ausmultiplizieren
\[
f^{*}\eta=\sum_{\vec{J}}(b_{\vec{J}}\circ f)\sum_{I}
\frac{\partial y^{j_{1}}}{\partial x^{i_{1}}}\cdots
\frac{\partial y^{j_{k}}}{\partial x^{i_{k}}}\;
d x^{i_{1}}\wedge\cdots\wedge d x^{i_{k}},
\]
und nach Umordnung der $dx$‘s mit Hilfe des Vorzeichentensors
$\epsilon_{\vec{L}}^{I}=\epsilon^{\,i_{1}\ldots i_{k}}_{\;\;\;\;\ell_{1}\ldots \ell_{k}}$
die obige determinantenförmige Darstellung:
\[
\sum_{I}\epsilon_{\vec{L}}^{I}
\frac{\partial y^{j_{1}}}{\partial x^{i_{1}}}\cdots
\frac{\partial y^{j_{k}}}{\partial x^{i_{k}}}
=\frac{\partial y^{\vec{J}}}{\partial x^{\vec{L}}}.
\]
\end{LEM}



The local expression for the pull-back is described as follows. Let us abbreviate: $\epsilon_{L}^{I}=\epsilon_{\ell_{1} \ldots \ell_{k}}^{i_{1} \ldots i_{k}}$ (recall the definition given by equation (8.2)). Let $(U, \mathrm{x})$ be a chart on $M$ and $(V, \mathrm{y})$ a chart on $N$ with $f(U) \subset V$. Then,\\
writing $\eta=\sum b_{\vec{J}} d y^{\vec{J}}$ and abbreviating $\frac{\partial\left(y^{j} \circ f\right)}{\partial x^{i}}$ to simply $\frac{\partial y^{j}}{\partial x^{i}}$, etc. we have

$$
\begin{aligned}
f^{*} \eta & =\sum_{b_{\vec{J}} \circ f} d\left(y^{j_{1}} \circ f\right) \wedge \cdots \wedge d\left(y^{j_{k}} \circ f\right) \\
& =\sum_{\vec{J}}\left(b_{\vec{J}} \circ f\right)\left(\sum_{i_{1}} \frac{\partial y^{j_{1}}}{\partial x^{i_{1}}} d x^{i_{1}}\right) \wedge \cdots \wedge\left(\sum_{i_{k}} \frac{\partial y^{j_{k}}}{\partial x^{i_{k}}} d x^{i_{k}}\right) \\
& =\sum_{\vec{J}}\left(b_{\vec{J}} \circ f\right) \sum_{I} \frac{\partial y^{j_{1}}}{\partial x^{i_{1}}} \cdots \frac{\partial y^{j_{k}}}{\partial x^{i_{k}}} d x^{i_{1}} \wedge \cdots \wedge d x^{i_{k}} \\
& =\sum_{\vec{J}}\left(b_{\vec{J}} \circ f\right) \sum_{\vec{L}}\left(\sum_{I} \epsilon_{\vec{L}}^{I} \frac{\partial y^{j_{1}}}{\partial x^{i_{1}}} \cdots \frac{\partial y^{j_{k}}}{\partial x^{i_{1}}}\right) d x^{\ell_{1}} \wedge \cdots \wedge d x^{\ell_{k}} \\
& =\sum_{\vec{J}} \sum_{\vec{L}} b_{\vec{J}} \circ f \frac{\partial\left(y^{j_{1}}, \ldots, y^{j_{k}}\right)}{\partial\left(x^{\ell_{1}}, . ., x^{\ell_{k}}\right)} d x^{\ell_{1}} \wedge \cdots \wedge d x^{\ell_{k}}
\end{aligned}
$$

where

$$
\frac{\partial y^{J}}{\partial x^{L}}:=\frac{\partial\left(y^{j_{1}}, \ldots, y^{j_{k}}\right)}{\partial\left(x^{\ell_{1}}, \ldots, x^{\ell_{k}}\right)}=\operatorname{det}\left[\begin{array}{ccc}
\frac{\partial y^{j_{1}}}{\partial x^{\ell_{1}}} & \cdots & \frac{\partial y^{j_{1}}}{\partial x^{\ell_{k}}} \\
\vdots & & \vdots \\
\frac{\partial y^{j_{k}}}{\partial x^{\ell_{1}}} & \cdots & \frac{\partial y^{j_{k}}}{\partial x^{\ell_{k}}}
\end{array}\right]
$$

\begin{EXA}{GA-L09-06-49}{Lokale Darstellung des Pullbacks von Differentialformen für $\operatorname{dim} M=2, \operatorname{dim} N=3, k=2$}
Since the above formula is a bit intimidating at first sight, we work out the case where $\operatorname{dim} M=2, \operatorname{dim} N=3$ and $k=2$. As a warm up, notice that since $d x^{i} \wedge d x^{i}=0$, we have

$$
\begin{aligned}
d\left(y^{2} \circ f\right) \wedge d\left(y^{3} \circ f\right) & =\sum \frac{\partial y^{2}}{\partial x^{i}} d x^{i} \wedge \frac{\partial y^{3}}{\partial x^{j}} d x^{j}=\sum \frac{\partial y^{2}}{\partial x^{i}} \frac{\partial y^{3}}{\partial x^{j}} d x^{i} \wedge d x^{j} \\
& =\frac{\partial y^{2}}{\partial x^{1}} \frac{\partial y^{3}}{\partial x^{2}} d x^{1} \wedge d x^{2}+\frac{\partial y^{2}}{\partial x^{2}} \frac{\partial y^{3}}{\partial x^{1}} d x^{2} \wedge d x^{1} \\
& =\left(\frac{\partial y^{2}}{\partial x^{1}} \frac{\partial y^{3}}{\partial x^{2}}-\frac{\partial y^{2}}{\partial x^{2}} \frac{\partial y^{3}}{\partial x^{1}}\right) d x^{1} \wedge d x^{2} \\
& =\frac{\partial\left(y^{2}, y^{3}\right)}{\partial\left(x^{1}, x^{2}\right)} d x^{1} \wedge d x^{2}
\end{aligned}
$$

Using similar expressions, we have

$$
\begin{aligned}
f^{*} \eta= & f^{*}\left(b_{23} d y^{2} \wedge d y^{3}+b_{13} d y^{1} \wedge d y^{3}+b_{12} d y^{1} \wedge d y^{2}\right) \\
= & b_{23} \circ f \sum \frac{\partial y^{2}}{\partial x^{i}} d x^{i} \wedge \frac{\partial y^{3}}{\partial x^{j}} d x^{j}+b_{13} \circ f \sum \frac{\partial y^{1}}{\partial x^{i}} d x^{i} \wedge \frac{\partial y^{3}}{\partial x^{j}} d x^{j} \\
& +b_{12} \circ f \sum \frac{\partial y^{1}}{\partial x^{i}} d x^{i} \wedge \frac{\partial y^{2}}{\partial x^{j}} d x^{j} \\
= & \left(b_{23} \circ f \frac{\partial\left(y^{2}, y^{3}\right)}{\partial\left(x^{1}, x^{2}\right)}+b_{13} \circ f \frac{\partial\left(y^{1}, y^{3}\right)}{\partial\left(x^{1}, x^{2}\right)}+b_{12} \circ f \frac{\partial\left(y^{1}, y^{2}\right)}{\partial\left(x^{1}, x^{2}\right)}\right) d x^{1} \wedge d x^{2}
\end{aligned}
$$
\end{EXA}

Remark 8.32. Notice that the space $\Omega^{0}(M)$ is just the space of smooth functions $C^{\infty}(M)$, and so unfortunately we now have several notations for the same space: $C^{\infty}(M)=\Omega^{0}(M)=\mathcal{T}_{0}^{0}(M)$.

\subsection*{Exterior Derivative}


Here we will define and study the exterior derivative $d$. For 0-forms, exterior differentiation is just the operation of taking the differential: $f \mapsto d f$. Let us start right out with giving an idea of what the exterior derivative looks like for $k$-forms defined on an open set in $U \subset \mathbb{R}^{n}$. Using standard rectangular coordinates $x^{i}$, all $k$-forms can be written as sums of terms of the form $f d x^{i_{1}} \wedge \cdots \wedge d x^{i_{k}}$ for some $f \in C^{\infty}(U)$. We know that the differential of a 0 -form is a 1-form: $d: f \mapsto \frac{\partial f}{\partial x^{i}} d x^{i}$. We inductively extend the definition of $d$ to an operator that takes $k$-forms to $k+1$ forms. We declare $d$ to be linear over real numbers and then define $d\left(f d x^{i_{1}} \wedge \cdots \wedge d x^{i_{k}}\right)=d f \wedge d x^{i_{1}} \wedge \cdots \wedge d x^{i_{k}}$. It is an easy exercise to show that if the latter formula holds for increasing indices $i_{1}<\cdots<i_{k}$, then it holds for all choices of indices. For example, if in $\mathbb{R}^{2}$ we have a 1 -form $\alpha=x^{2} d x+x y d y$, then

$$
\begin{aligned}
d \alpha & =d\left(x^{2} d x+x y d y\right) \\
& =d\left(x^{2}\right) \wedge d x+d(x y) \wedge d y \\
& =2 x d x \wedge d x+(y d x+x d y) \wedge d y \\
& =y d x \wedge d y
\end{aligned}
$$



\begin{DEF}{GA-L09-06-50}{Äußere Ableitung in $\mathbb{R}^{n}$}
Sei $U\subset\mathbb{R}^{n}$ offen mit Standardkoordinaten $(x^{1},\ldots,x^{n})$.  
Jede $k$-Form $\omega\in\Omega^{k}(U)$ kann eindeutig geschrieben werden als
\[
\omega=\sum_{i_{1}<\cdots<i_{k}} f_{i_{1}\ldots i_{k}}(x)\, d x^{i_{1}}\wedge\cdots\wedge d x^{i_{k}},
\qquad f_{i_{1}\ldots i_{k}}\in C^{\infty}(U).
\]
Die \emph{äußere Ableitung} oder \emph{exterior derivative}
\[
d:\Omega^{k}(U)\longrightarrow \Omega^{k+1}(U)
\]
wird durch folgende Regeln definiert:

\begin{enumerate}
\item[(1)] Für $f\in\Omega^{0}(U)=C^{\infty}(U)$ gilt
\[
d f = \sum_{i=1}^{n} \frac{\partial f}{\partial x^{i}}\, d x^{i}.
\]
\item[(2)] $d$ ist $\mathbb{R}$-linear: $d(\alpha\omega+\beta\eta)=\alpha\, d\omega+\beta\, d\eta$ für $\alpha,\beta\in\mathbb{R}$.
\item[(3)] Für einfache Terme gilt
\[
d\!\left(f\, d x^{i_{1}}\wedge\cdots\wedge d x^{i_{k}}\right)
:= d f \wedge d x^{i_{1}}\wedge\cdots\wedge d x^{i_{k}}.
\]
\item[(4)] In expliziter Indexschreibweise:
\[
d\omega=\sum_{i_{1}<\cdots<i_{k}} \sum_{j=1}^{n}
\frac{\partial f_{i_{1}\ldots i_{k}}}{\partial x^{j}}\,
d x^{j}\wedge d x^{i_{1}}\wedge\cdots\wedge d x^{i_{k}}.
\]
\end{enumerate}

\textbf{Eigenschaften.}
\begin{itemize}
\item $d$ ist eine Abbildung vom Grad $+1$ im graduierten Ring $\Omega(U)$.
\item Für $f\in C^{\infty}(U)$ und $\omega\in\Omega^{k}(U)$ gilt die \textbf{Leibniz-Regel}
\[
d(f\omega)=d f\wedge\omega+f\,d\omega.
\]
\item Es gilt stets $d\circ d=0$, also $d^{2}=0$.
\end{itemize}
\end{DEF}







We now develop the general theory on manifolds. For the next theorem we think of $\Omega_{M}$ as an assignment $\Omega_{M}: U \mapsto \Omega_{M}(U)=\Omega(U)$ for open $U \subset M$. Thus we are simply thinking in terms of (pre)sheaves. We will drop the subscript $M$ on $\Omega_{M}$ when there is no chance of confusion.

Definition 8.33. A (natural) graded derivation of degree $r$ on $\Omega_{M}$ is a family of maps, one for each open set $U \subset M$, denoted $\mathcal{D}_{U}: \Omega_{M}(U) \rightarrow \Omega_{M}(U)$, such that for each $U \subset M$,

$$
\mathcal{D}_{U}: \Omega_{M}^{k}(U) \rightarrow \Omega_{M}^{k+r}(U)
$$

and such that\\
(1) $\mathcal{D}_{U}$ is $\mathbb{R}$ linear;\\
(2) $\mathcal{D}_{U}(\alpha \wedge \beta)=\mathcal{D}_{U} \alpha \wedge \beta+(-1)^{k r} \alpha \wedge \mathcal{D}_{U} \beta$ for $\alpha \in \Omega^{k}(U)$ and $\beta \in \Omega_{M}(U) ;$\\
(3) $\mathcal{D}_{U}$ is natural with respect to restriction:\\
\includegraphics[scale=0.25, center]{2025_10_09_265d7e1a22c89f79c21eg-378}

As usual we will denote all of the maps by a single symbol $\mathcal{D}$. In summary, we have a map of (pre)sheaves $\mathcal{D}: \Omega_{M} \rightarrow \Omega_{M}$. Along the lines similar to our study of tensor derivations, one can show that a graded derivation of $\Omega_{M}$ is completely determined by, and can be defined by its action on $0-$ forms (functions) and 1 -forms. In fact, since every form can be locally built out of functions and exact 1 -forms, i.e. differentials, we only need to know the action on 0 -forms and exact 1 -form to determine the graded derivation. Recall that an element of $\Omega^{1}(U)$ is said to be exact if it is the differential of a smooth function.

Remark 8.34. If one has a map $\mathcal{D}: \Omega(M) \rightarrow \Omega(M)$ that satisfies (1) and (2) of the previous definition, then it is just called a graded derivation of degree $k$. But, when we meet derivations below, they will also be defined on open submanifolds and will all give natural derivations.



\begin{DEF}{GA-L09-06-51}{Natürliche graduierte Derivation auf dem Prägarben‐System der Differentialformen}
Sei $M$ eine glatte Mannigfaltigkeit und 
\[
\Omega_{M}:\ \mathrm{Open}(M)\longrightarrow \mathbf{GrAlg},
\qquad U\longmapsto \Omega(U)
\]
die Prägarbe der Differentialformen auf $M$,
wobei $\Omega^{k}(U)$ die glatten $k$-Formen auf $U$ bezeichnet und 
$\Omega(U)=\bigoplus_{k\ge0}\Omega^{k}(U)$ den graduierten kommutativen Ring der Formen auf $U$.

Eine \textbf{(natürliche) graduierte Derivation vom Grad $r$} auf $\Omega_{M}$
ist ein System von Abbildungen
\[
\mathcal{D}_{U}:\Omega(U)\longrightarrow \Omega(U),
\qquad U\subset M\ \text{offen},
\]
so dass für jedes offene $U\subset M$ gilt:
\begin{enumerate}
\item[(1)] \textbf{Gradbedingung:} $\mathcal{D}_{U}\bigl(\Omega^{k}(U)\bigr)\subseteq \Omega^{k+r}(U)$.
\item[(2)] \textbf{$\mathbb{R}$‐Linearität:} $\mathcal{D}_{U}(a\alpha+b\beta)=a\,\mathcal{D}_{U}\alpha+b\,\mathcal{D}_{U}\beta$
für alle $a,b\in\mathbb{R}$ und $\alpha,\beta\in\Omega(U)$.
\item[(3)] \textbf{Abgeleitete Produktregel:}
Für $\alpha\in\Omega^{k}(U)$ und $\beta\in\Omega(U)$ gilt
\[
\mathcal{D}_{U}(\alpha\wedge\beta)
=\mathcal{D}_{U}\alpha\wedge\beta
+(-1)^{kr}\,\alpha\wedge\mathcal{D}_{U}\beta.
\]
\item[(4)] \textbf{Natürlichkeit bezüglich Restriktion:}
Für offene Mengen $V\subset U$ kommutiert das folgende Diagramm:
\[
\begin{tikzcd}[column sep=large]
\Omega(U)\arrow{r}{\mathcal{D}_{U}}\arrow{d}[swap]{r^{U}_{V}} & \Omega(U)\arrow{d}{r^{U}_{V}}\\
\Omega(V)\arrow{r}{\mathcal{D}_{V}} & \Omega(V)
\end{tikzcd}
\]
Das heißt:
\[
\mathcal{D}_{V}\circ r^{U}_{V} = r^{U}_{V}\circ\mathcal{D}_{U}.
\]
\end{enumerate}

Ein solches System $\mathcal{D}=\{\mathcal{D}_{U}\}_{U\subset M}$
heißt dann eine \emph{natürliche graduierte Derivation vom Grad $r$}
auf der Prägarbe $\Omega_{M}$, oder kurz:
\[
\mathcal{D}:\Omega_{M}\longrightarrow \Omega_{M}
\quad\text{mit}\quad \deg(\mathcal{D})=r.
\]
\end{DEF}


\begin{CONC}{GA-L09-06-52}{Kommentare zu Natürliche graduierte Derivation auf dem Prägarben‐System}
\begin{enumerate}
\item Eine graduierte Derivation $\mathcal{D}$ ist vollständig bestimmt durch ihre Wirkung
auf $0$-Formen (glatte Funktionen) und auf $1$-Formen.
Da jede Form lokal als Kombination von glatten Funktionen und exakten $1$-Formen $d f$
aufgebaut werden kann, reicht die Kenntnis von $\mathcal{D}(f)$ und $\mathcal{D}(d f)$ aus,
um $\mathcal{D}$ eindeutig zu bestimmen.

\item Wenn eine Abbildung $\mathcal{D}:\Omega(M)\to\Omega(M)$ nur die Bedingungen (1) und (2) erfüllt,
so nennt man sie eine \emph{graduierte Derivation vom Grad $r$} auf dem algebraischen Objekt $\Omega(M)$.
Der Zusatz „natürlich“ bezieht sich auf die Erfüllung der Bedingung (4) für alle offenen Teilmengen $U\subset M$.

\item Wichtige Beispiele sind:
\begin{itemize}
\item Der \textbf{äußere Differentialoperator} $d:\Omega^{k}(M)\to\Omega^{k+1}(M)$ (Grad $r=+1$);
\item Die \textbf{Innere Ableitung} $i_{X}$ durch ein Vektorfeld $X$ (Grad $r=-1$);
\item Die \textbf{Lie-Ableitung} $\mathcal{L}_{X}$ (Grad $r=0$).
\end{itemize}
\end{enumerate}
\end{CONC}


Proposition 8.35. 

\begin{PROP}{GA-L09-06-53}{Induzierter natürliche graduierte Klammer-Derivation}
If $\mathcal{D}_{1}$ and $\mathcal{D}_{2}$ are (natural) graded derivations of degrees $r_{1}$ and $r_{2}$ respectively then the operator
$$
\left[\mathcal{D}_{1}, \mathcal{D}_{2}\right]:=\mathcal{D}_{1} \circ \mathcal{D}_{2}-(-1)^{r_{1} r_{2}} \mathcal{D}_{2} \circ \mathcal{D}_{1}
$$
is a (natural) graded derivation of degree $r_{1}+r_{2}$.
\end{PROP}


Proof. 

\begin{PROOF}{GA-L09-06-54}{P: Induzierter natürliche graduierte Klammer-Derivation}
See Problem 15.
\end{PROOF}


Lemma 8.36. 

\begin{LEM}{GA-L09-06-55}{Charakterisierung von identischen graduierten Derivationen}
Suppose $\mathcal{D}_{1}: \Omega_{M}^{k}(U) \rightarrow \Omega_{M}^{k+r}(U)$ and $\mathcal{D}_{2}: \Omega_{M}^{k}(U) \rightarrow \Omega_{M}^{k+r}(U)$ are defined for each open set $U \subset M$ and both satisfy (1), (2) and (3) of Definition 8.33 above. If $\mathcal{D}_{1}$ and $\mathcal{D}_{2}$ agree when applied to functions and exact forms, then $\mathcal{D}_{1}=\mathcal{D}_{2}$.
\end{LEM}

Proof. 

\begin{PROOF}{GA-L09-06-56}{P: Charakterisierung von identischen graduierten Derivationen}
By (3), if $\mathcal{D}_{1}$ and $\mathcal{D}_{2}$ agree on chart domains, then they agree globally. Let $x^{1}, \ldots, x^{n}$ be local coordinates on $U$. Then every element of $\Omega_{M}^{k}(U)$ is a sum of elements of the form $f d x^{i_{1}} \wedge \cdots \wedge d x^{i_{k}}$. But by (2) we have
$$
\begin{aligned}
\mathcal{D}_{1}\left(f d x^{i_{1}} \wedge \cdots \wedge d x^{i_{k}}\right) & =\mathcal{D}_{1} f \wedge d x^{i_{1}} \wedge \cdots \wedge d x^{i_{k}} \pm f \mathcal{D}_{1}\left(d x^{i_{1}} \wedge \cdots \wedge d x^{i_{k}}\right) \\
& =\mathcal{D}_{2} f \wedge d x^{i_{1}} \wedge \cdots \wedge d x^{i_{k}} \pm f \mathcal{D}_{1}\left(d x^{i_{1}} \wedge \cdots \wedge d x^{i_{k}}\right)
\end{aligned}
$$
The last term can be expanded using (2), and then the elements $\mathcal{D}_{1} d x^{i_{j}}$ can be replaced by $\mathcal{D}_{2} d x^{i_{j}}$. The result is equal to $\mathcal{D}_{2}\left(f d x^{i_{1}} \wedge \cdots \wedge d x^{i_{k}}\right)$.
\end{PROOF}

The differential $d$ defined by

$$
d f(X)=X f \text { for } X \in \mathfrak{X}_{M}(U) \text { and } f \in C^{\infty}(U)
$$

gives a map $\Omega_{M}^{0} \rightarrow \Omega_{M}^{1}$. Next we show that this map can be extended to a degree one graded derivation.

Theorem 8.37. Let $M$ be a smooth manifold. There is a unique degree one graded derivation $d: \Omega_{M} \rightarrow \Omega_{M}$ such that

$$
d \circ d=0
$$

and such that for each open $U \subset M$ and $f \in C^{\infty}(U)=\Omega_{M}^{0}(U)$, the 1-form df coincides with the usual differential. Furthermore, for any chart ( $U, \mathrm{x}$ ) for $M$, we have the following local formula:

$$
d \sum \alpha_{\vec{I}} d x^{\vec{I}}=\sum_{\vec{I}}\left(d \alpha_{\vec{I}}\right) \wedge d x^{\vec{I}}
$$

Proof. We define an operator $d_{\mathrm{x}}$ for each chart ( $U, \mathrm{x}$ ). For a 0 -form on $U$ (i.e. a smooth function), we just define $d_{\mathrm{x}} f$ to be the usual differential given by $d f=\sum \frac{\partial f}{\partial x^{i}} d x^{i}$. For $\alpha \in \Omega_{M}^{k}(U)$, we have $\alpha=\sum \alpha_{\vec{I}} d x^{\vec{I}}$ and we define $d_{\mathrm{x}} \alpha=\sum d \alpha_{\vec{I}} \wedge d x^{\vec{I}}$. To show the product rule ((2) of Definition 8.33), consider $\alpha=\sum \alpha_{\vec{I}} d x^{\vec{I}} \in \Omega_{M}^{k}(U)$ and $\beta=\sum \beta_{\vec{J}} d x^{\vec{J}} \in \Omega_{M}^{\ell}(U)$. Then

$$
\begin{aligned}
d_{\mathrm{x}}(\alpha \wedge \beta)= & d_{\mathrm{x}}\left(\sum \alpha_{\vec{I}} d x^{\vec{I}} \wedge \sum \beta_{\vec{J}} d x^{\vec{J}}\right) \\
= & d_{\mathrm{x}}\left(\sum \alpha_{\vec{I}} \beta_{\vec{J}} d x^{\vec{I}} \wedge d x^{\vec{J}}\right) \\
= & \sum\left(\left(d \alpha_{\vec{I}}\right) \beta_{\vec{J}}+\alpha_{\vec{I}}\left(d \beta_{\vec{J}}\right)\right) d x^{\vec{I}} \wedge d x^{\vec{J}} \\
= & \left(\sum_{\vec{I}} d \alpha_{\vec{I}} \wedge d x^{\vec{I}}\right) \wedge \sum_{\vec{J}} \beta_{\vec{J}} d x^{\vec{J}} \\
& +\sum_{\vec{I}} \alpha_{\vec{I}} d x^{\vec{I}} \wedge\left((-1)^{k} \sum_{\vec{J}} d \beta_{\vec{J}} \wedge d x^{\vec{J}}\right)
\end{aligned}
$$

since $d \beta_{\vec{J}} \wedge d x^{\vec{I}}=(-1)^{k} d x^{\vec{I}} \wedge d \beta_{\vec{J}}$ due to the $k$ interchanges of the basic differentials $d x^{i}$. This means that the product rule holds for each $d_{\mathrm{x}}$. For any function $f$, we have $d_{\mathrm{x}} d_{\mathrm{x}} f=d_{\mathrm{x}} d f=\sum_{i j}\left(\frac{\partial^{2} f}{\partial x^{i} \partial x^{j}}\right) d x^{i} \wedge d x^{j}=0$ since $\frac{\partial^{2} f}{\partial x^{i} \partial x^{j}}$ is symmetric in $i, j$ and $d x^{i} \wedge d x^{j}$ is antisymmetric in $i, j$. More generally, for any functions $f, g \in C^{\infty}(U)$ we have $d_{\mathrm{x}}(d f \wedge d g)=0$ because of the graded commutativity. Inductively we get $d_{\mathbf{x}}\left(d f_{1} \wedge d f_{2} \wedge \cdots \wedge d f_{k}\right)=0$ for any functions $f_{i} \in C^{\infty}(U)$. From this it follows that for any $\alpha=\sum \alpha_{\vec{I}} d x^{\vec{I}} \in \Omega_{M}^{k}(U)$ we have $d_{\mathrm{x}} d_{\mathrm{x}} \sum \alpha_{\vec{I}} d x^{\vec{I}}=d_{\mathrm{x}} \sum d \alpha_{\vec{I}} \wedge d x^{\vec{I}}=0$ since $d_{\mathrm{x}} d \alpha_{\vec{I}}=0$ and $d_{\mathrm{x}} d x^{\vec{I}}=0$. We have now defined, for each coordinate chart ( $U, \mathrm{x}$ ), an operator $d_{\mathrm{x}}$ that clearly has the desired properties on that chart. Consider two different charts ( $U, \mathrm{x}$ ) and ( $V, \mathrm{y}$ ) such that $U \cap V \neq \emptyset$. We need to show that $d_{\mathrm{x}}$ restricted to $U \cap V$ coincides with $d_{\mathrm{y}}$ restricted to $U \cap V$, but it is\\
clear that these restrictions of $d_{\mathrm{x}}$ and $d_{\mathrm{y}}$ satisfy the hypothesis of Lemma 8.36 and so they must agree on $U \cap V$.

It is now clear that the individual operators on coordinate charts fit together to give a well-defined operator with the desired properties.



\begin{DEF}{GA-L09-06-57}{Existenz und Eindeutigkeit des äußeren Differentialoperators auf Mannigfaltigkeiten}
Sei $M$ eine glatte Mannigfaltigkeit.  
Dann existiert genau ein Operator
\[
d:\Omega_{M}\longrightarrow\Omega_{M}
\]
mit folgenden Eigenschaften:
\begin{enumerate}
\item[(1)] $d$ ist eine natürliche graduierte Derivation vom Grad $+1$, d.\,h.
\[
d:\Omega^{k}(U)\longrightarrow\Omega^{k+1}(U)
\]
für alle offenen Mengen $U\subset M$, und es gilt die Produktregel
\[
d(\alpha\wedge\beta)=d\alpha\wedge\beta+(-1)^{k}\,\alpha\wedge d\beta
\quad\text{für }\alpha\in\Omega^{k}(U),\,\beta\in\Omega(U).
\]
\item[(2)] $d$ erfüllt die Nilpotenzbedingung
\[
d\circ d = 0.
\]
\item[(3)] Für glatte Funktionen $f\in C^{\infty}(U)=\Omega^{0}(U)$ gilt
\[
d f(X)=X(f)
\quad\text{für alle }X\in\mathfrak{X}(U).
\]
\end{enumerate}
Zudem gilt für jede Karte $(U,\mathrm{x})$ mit Koordinaten $(x^{1},\dots,x^{n})$
und für jede Form
\[
\alpha=\sum_{\vec{I}}\alpha_{\vec{I}}\, d x^{\vec{I}}
\quad (\alpha_{\vec{I}}\in C^{\infty}(U)),
\]
die lokale Formel
\[
d\alpha
=\sum_{\vec{I}}\bigl(d\alpha_{\vec{I}}\bigr)\wedge d x^{\vec{I}}
=\sum_{\vec{I}}\sum_{j=1}^{n}
\frac{\partial \alpha_{\vec{I}}}{\partial x^{j}}\,
d x^{j}\wedge d x^{\vec{I}}.
\]

\textbf{Beweis.}
Wir definieren für jede Karte $(U,\mathrm{x})$ einen Operator $d_{\mathrm{x}}$ wie folgt:
\begin{itemize}
\item Für $f\in C^{\infty}(U)$:
\[
d_{\mathrm{x}}f
=\sum_{i=1}^{n}\frac{\partial f}{\partial x^{i}}\,d x^{i}.
\]
\item Für $\alpha=\sum_{\vec{I}}\alpha_{\vec{I}}\,d x^{\vec{I}}\in\Omega^{k}(U)$:
\[
d_{\mathrm{x}}\alpha
:=\sum_{\vec{I}}(d\alpha_{\vec{I}})\wedge d x^{\vec{I}}.
\]
\end{itemize}
Man überprüft unmittelbar:
\begin{enumerate}
\item[(a)] $d_{\mathrm{x}}$ erfüllt die Produktregel
\(
d_{\mathrm{x}}(\alpha\wedge\beta)
=d_{\mathrm{x}}\alpha\wedge\beta+(-1)^{k}\alpha\wedge d_{\mathrm{x}}\beta.
\)
\item[(b)] $d_{\mathrm{x}}\circ d_{\mathrm{x}}=0$, da für $f\in C^{\infty}(U)$ gilt
\[
d_{\mathrm{x}}(d_{\mathrm{x}}f)
=\sum_{i,j}\frac{\partial^{2}f}{\partial x^{i}\partial x^{j}}\,d x^{i}\wedge d x^{j}
=0
\]
wegen der Symmetrie der zweiten partiellen Ableitungen und der Antisymmetrie des Keilprodukts.
\item[(c)] Für Funktionen $f,g$ gilt $d_{\mathrm{x}}(d f\wedge d g)=0$ und induktiv
$d_{\mathrm{x}}(d f_{1}\wedge\cdots\wedge d f_{k})=0$, also $d_{\mathrm{x}}^{2}=0$ auf allen $\Omega^{k}(U)$.
\end{enumerate}

Für zwei Karten $(U,\mathrm{x})$ und $(V,\mathrm{y})$ mit $U\cap V\neq\varnothing$
fallen die Operatoren $d_{\mathrm{x}}$ und $d_{\mathrm{y}}$ auf $U\cap V$ zusammen,
da sie beide die in Lemma 8.36 formulierten Naturbedingungen erfüllen.
Somit definieren die $d_{\mathrm{x}}$ global einen wohldefinierten Operator $d$ auf $\Omega_{M}$,
der alle obigen Eigenschaften besitzt.

Die Eindeutigkeit folgt daraus, dass jede graduierte Derivation vom Grad 1
durch ihre Wirkung auf $0$-Formen und exakte $1$-Formen bestimmt ist.
\end{DEF}



Definition 8.38. The degree one graded derivation just introduced is called the exterior derivative.

Another approach to the existence of the exterior derivative is to exhibit a global coordinate free formula. Let $\omega \in \Omega^{k}(M)$ and view $\omega$ as an alternating multilinear map on $\mathfrak{X}(M)$. Then for $X_{0}, X_{1}, \ldots, X_{k} \in \mathfrak{X}(M)$, define

$$
\begin{aligned}
& d \omega\left(X_{0}, X_{1}, \ldots, X_{k}\right) \\
& \quad=\sum_{0 \leq i \leq k}(-1)^{i} X_{i}\left(\omega\left(X_{0}, \ldots, \widehat{X}_{i}, \ldots, X_{k}\right)\right) \\
& \quad+\sum_{0 \leq i<j \leq k}(-1)^{i+j} \omega\left(\left[X_{i}, X_{j}\right], X_{0}, \ldots, \widehat{X}_{i}, \ldots, \widehat{X}_{j}, \ldots, X_{k}\right)
\end{aligned}
$$

One can check that $d \omega$ is an alternating $C^{\infty}(M)$-multilinear map on $\mathfrak{X}(M)$ and so defines a differential form of degree $k+1$. By applying this formula to coordinate fields one obtains the same local operator defined previously.



\begin{DEF}{GA-L09-06-58}{Koordinatenfreie Definition der äußeren Ableitung (Cartansche Formel)}
Sei $M$ eine glatte Mannigfaltigkeit und $\omega\in\Omega^{k}(M)$ eine glatte $k$-Form, die wir als alternierende $C^\infty(M)$-multilineare Abbildung
\[
\omega:\underbrace{\mathfrak{X}(M)\times\cdots\times\mathfrak{X}(M)}_{k\text{ Faktoren}}\to C^\infty(M)
\]
auffassen. Für $X_0,\dots,X_k\in\mathfrak{X}(M)$ definieren wir $d\omega\in\Omega^{k+1}(M)$ durch
\[
\boxed{%
\begin{aligned}
(d\omega)(X_0,\dots,X_k)
&=\sum_{i=0}^{k}(-1)^i\,X_i\!\bigl(\,\omega(X_0,\dots,\widehat{X_i},\dots,X_k)\,\bigr)\\
&\quad+\sum_{0\le i<j\le k}(-1)^{i+j}\,\omega\!\bigl([X_i,X_j],X_0,\dots,\widehat{X_i},\dots,\widehat{X_j},\dots,X_k\bigr).
\end{aligned}}
\]
Hier bezeichnet $\widehat{X_i}$ das Weglassen des Arguments $X_i$, und $[\,,\,]$ ist die Lie-Klammer auf $\mathfrak{X}(M)$.
\end{DEF}




\begin{PROP}{GA-L09-06-59}{Eigenschaften der koordinatenfreien Definition des äußeren Differentials}
Die in der Definition gegebene Vorschrift besitzt folgende Eigenschaften:
\begin{enumerate}
\item \textbf{Alternierung und $C^\infty$-Multilinearität.}
Für festes $\omega$ ist $d\omega:\mathfrak{X}(M)^{k+1}\to C^\infty(M)$ alternierend und $C^\infty(M)$-multilinear; somit ist $d\omega\in\Omega^{k+1}(M)$ wohldefiniert.

\item \textbf{Lokale Übereinstimmung.}
In jeder Karte $(U,x)$ gilt für $\omega=\sum_{\vec I}\omega_{\vec I}\,dx^{\vec I}$ die lokale Formel
\[
(d\omega)|_U=\sum_{\vec I} d\omega_{\vec I}\wedge dx^{\vec I},
\]
d.\,h. die koordinatenfreie Definition stimmt mit der in lokalen Koordinaten eingeführten überein.

\item \textbf{Graduierte Leibniz-Regel.}
Für $\alpha\in\Omega^k(M)$ und $\beta\in\Omega^\ell(M)$ gilt
\[
d(\alpha\wedge\beta)=d\alpha\wedge\beta+(-1)^k\,\alpha\wedge d\beta.
\]

\item \textbf{Nilpotenz.}
Es gilt $d\circ d=0$ auf $\Omega(M)$.

\item \textbf{Funktorialität (Natürlichkeit).}
Für eine glatte Abbildung $f:M\to N$ und $\eta\in\Omega(N)$ gilt
\[
f^{*}(d\eta)=d(f^{*}\eta).
\]

\item \textbf{Übereinstimmung auf Funktionen.}
Für $f\in C^\infty(M)$ und $X\in\mathfrak{X}(M)$ gilt
\[
(df)(X)=X(f).
\]
\end{enumerate}
\end{PROP}




\begin{PROOF}{GA-L09-06-60}{P: Eigenschaften der koordinatenfreien Definition des äußeren Differentials}
(1) $C^\infty$-Multilinearität folgt aus der Produktregel $X(h\cdot g)=X(h)\,g+h\,X(g)$ und der $C^\infty$-Linearität der Lie-Klammer in jedem Argument; Alternierung folgt aus der Alternierung von $\omega$ sowie der Antisymmetrie von $[\,,\,]$.  

(2) Einsetzen der Koordinatenvektorfelder $\partial_{x^i}$ und Benutzen $[\partial_{x^i},\partial_{x^j}]=0$ liefert die bekannte lokale Formel.  

(3) Folgt durch direktes Ausmultiplizieren in Definition und Sortieren der Terme nach Vorzeichen (Standardargument).  

(4) Aus Definition unter Verwendung der Jacobi-Identität für die Lie-Klammer sowie der Alternierung (klassische Rechnung).  

(5) Direkt aus der Definition und der Kettenregel für Tangentialabbildungen bzw.\ aus (2) in lokalen Koordinaten.  

(6) Setze $k=0$.
\end{PROOF}





Lemma 8.39. 

\begin{LEM}{GA-L09-06-61}{Äußeres Differential ist natürlich bzgl Pullback}
Given any smooth map $f: M \rightarrow N$, we have that $d$ is natural with respect to the pull-back:
$$
f^{*}(d \eta)=d\left(f^{*} \eta\right)
$$
\end{LEM}

Proof. 

\begin{PROOF}{GA-L09-06-62}{P: Äußeres Differential ist natürlich bzgl Pullback}
By Lemma 2.119 we know the result is true if $\eta$ is a 1 -form. Because $d$ is natural with respect to restriction, we need only prove the formula for a differential form defined in the domain of a chart ( $U, \mathrm{x}$ ). By linearity we may assume that $\eta=g d x^{i_{1}} \wedge \cdots \wedge d x^{i_{k}}$ since an arbitrary form on $U$ is a sum of forms of this type:
$$
\begin{aligned}
f^{*}(d \eta) & =f^{*}\left(d\left(g d x^{i_{1}} \wedge \cdots \wedge d x^{i_{k}}\right)\right)=f^{*}\left(d g \wedge d x^{i_{1}} \wedge \cdots \wedge d x^{i_{k}}\right) \\
& =d\left(f^{*} g\right) \wedge d\left(f^{*} x^{i_{1}}\right) \wedge \cdots \wedge d\left(f^{*} x^{i_{k}}\right) \\
& =d(g \circ f) \wedge d\left(x^{i_{1}} \circ f\right) \wedge \cdots \wedge d\left(x^{i_{k}} \circ f\right) \\
& =d\left((g \circ f) d\left(x^{i_{1}} \circ f\right) \wedge \cdots \wedge d\left(x^{i_{k}} \circ f\right)\right)=d\left(f^{*} \eta\right)
\end{aligned}
$$
\end{PROOF}

Definition 8.40. 

\begin{DEF}{GA-L09-06-63}{Geschlossene und Exakte Differentialform}
A smooth differential form $\alpha$ is called closed if $d \alpha=0$ and exact if $\alpha=d \beta$ for some differential form $\beta$.

Notice that since $d \circ d=0$, every exact form is closed.
\end{DEF}


In general, the converse is not true. The extent to which the converse fails is a topological property of the manifold. This is the point of the de Rham ${ }^{1}$ cohomology

\footnotetext{${ }^{1}$ Georges de Rham 1903-1990.
}
to be studied in detail in Chapter 10. Here we just give the following basic definitions. 


The set of closed forms of degree $k$ on a smooth manifold $M$ is the kernel of $d: \Omega^{k}(M) \rightarrow \Omega^{k+1}(M)$ and is denoted $Z^{k}(M)$. The set of exact $k$-forms is the image of the map $d: \Omega^{k-1}(M) \rightarrow \Omega^{k}(M)$ and is denoted $B^{k}(M)$. Since $d \circ d=0$, we have $B^{k}(M) \subset Z^{k}(M)$.

Definition 8.41. The $k$-th de Rham cohomology group (actually a vector space) is given by

$$
H^{k}(M)=\frac{Z^{k}(M)}{B^{k}(M)}
$$

In other words, we look at closed forms and identify any two whose difference is an exact form.

If $\alpha \in Z^{k}(M) \subset \Omega^{k}(M)$, then the equivalence class that contains $\alpha$ is denoted $[\alpha]$ and called the cohomology class of $\alpha$. If $f: M \rightarrow N$ is smooth and $[\beta] \in H^{k}(N)$, then it is easy to see that $f^{*} \beta$ is closed since $\beta$ is closed. Thus we obtain a cohomology class $\left[f^{*} \beta\right]$. Also, $\left[f^{*} \beta\right]$ depends only on the equivalence class of $\beta$. Indeed, if $\beta-\beta^{\prime}=d \eta$, then $f^{*} \beta-f^{*} \beta^{\prime}=d f^{*} \eta$. Thus we may define a linear map $f^{*}: H^{k}(N) \rightarrow H^{k}(M)$ by $f^{*}[\beta]:=\left[f^{*} \beta\right]$. We return to this topic later.



\begin{DEF}{GA-L09-07-01}{de\,Rham-Kohomologie}
Sei $M$ eine glatte Mannigfaltigkeit und $d:\Omega^{k}(M)\to\Omega^{k+1}(M)$ der äußere Differentialoperator.

\begin{enumerate}
\item[(1)] \textbf{Geschlossene Formen.}  
Die Menge der \emph{geschlossenen $k$-Formen} auf $M$ ist definiert als
\[
Z^{k}(M):=\ker\!\bigl(d:\Omega^{k}(M)\longrightarrow\Omega^{k+1}(M)\bigr)
=\bigl\{\alpha\in\Omega^{k}(M)\mid d\alpha=0\bigr\}.
\]

\item[(2)] \textbf{Exakte Formen.}  
Die Menge der \emph{exakten $k$-Formen} auf $M$ ist das Bild des äußeren Differentials auf $(k-1)$-Formen:
\[
B^{k}(M):=\operatorname{im}\!\bigl(d:\Omega^{k-1}(M)\longrightarrow\Omega^{k}(M)\bigr)
=\bigl\{\alpha\in\Omega^{k}(M)\mid \exists\,\beta\in\Omega^{k-1}(M):\alpha=d\beta\bigr\}.
\]
Da $d\circ d=0$, gilt stets $B^{k}(M)\subseteq Z^{k}(M)$.

\item[(3)] \textbf{de\,Rham-Kohomologiegruppe.}  
Die $k$-te \emph{de\,Rham-Kohomologiegruppe} (eigentlich ein reeller Vektorraum) wird definiert als das Quotientenobjekt
\[
H^{k}_{\mathrm{dR}}(M):=\frac{Z^{k}(M)}{B^{k}(M)}
=\frac{\ker(d:\Omega^{k}(M)\to\Omega^{k+1}(M))}
      {\operatorname{im}(d:\Omega^{k-1}(M)\to\Omega^{k}(M))}.
\]
Die Elemente von $H^{k}_{\mathrm{dR}}(M)$ heißen \emph{Kohomologieklassen}. Für $\alpha\in Z^{k}(M)$ bezeichnen wir die Klasse von $\alpha$ durch
\[
[\alpha]:=\alpha+B^{k}(M).
\]
\end{enumerate}
\end{DEF}




\begin{PROP}{GA-L09-07-02}{Functorialität des Pullbacks auf der de\,Rham-Kohomologie}
Sei $f:M\to N$ eine glatte Abbildung und $k\ge0$.
\begin{enumerate}
\item Für jede geschlossene $k$-Form $\beta\in Z^{k}(N)$ ist auch $f^{*}\beta$ geschlossen, da
\[
d(f^{*}\beta)=f^{*}(d\beta)=f^{*}(0)=0.
\]
\item Ist $\beta=d\eta$ exakt, so gilt $f^{*}\beta=d(f^{*}\eta)$, d.\,h.\ $f^{*}\beta$ ist exakt.
\end{enumerate}
Somit induziert $f^{*}$ einen wohldefinierten linearen Abbildungsoperator
\[
f^{*}:H^{k}_{\mathrm{dR}}(N)\longrightarrow H^{k}_{\mathrm{dR}}(M),
\qquad
f^{*}[\beta]:=[f^{*}\beta].
\]
\end{PROP}




\begin{DEF}{GA-L09-07-03}{de\,Rham-Kohomologie als kontravarianter Funktor}
Die Abbildung
\[
M\longmapsto H^{k}_{\mathrm{dR}}(M),\qquad 
(f:M\to N)\longmapsto f^{*}:H^{k}_{\mathrm{dR}}(N)\to H^{k}_{\mathrm{dR}}(M)
\]
definiert für jedes feste $k$ einen \textbf{kontravarianten Funktor}
\[
H^{k}_{\mathrm{dR}}:\mathbf{Man}\longrightarrow\mathbf{Vect}_{\mathbb{R}}
\]
von der Kategorie der glatten Mannigfaltigkeiten in die Kategorie der reellen Vektorräume.
\end{DEF}



\pagebreak

\subsection*{Vector-Valued and Algebra-Valued Forms}
Now we consider a straightforward generalization. Let V and W be real vector spaces so that we have $L_{\text {alt }}^{k}(\mathrm{V} ; \mathrm{W})$ (Definition 8.1). For $\omega \in L_{\text {alt }}^{k}(\mathrm{V} ; \mathrm{W})$ and $\eta \in L_{\text {alt }}^{\ell}(\mathrm{V} ; \mathrm{W})$, we define the exterior product using the same formula as before except that we use the tensor product so that $\omega \wedge \eta$ is an element of $L_{\text {alt }}^{k+\ell}(\mathrm{V} ; \mathrm{W} \otimes \mathrm{W})$ :\\
$(\omega \wedge \eta)\left(v_{1}, v_{2}, \ldots, v_{k}, v_{k+1}, v_{k+2}, \ldots, v_{k+\ell}\right)$

$$
=\frac{1}{k_{1}!k_{2}!} \sum_{\sigma \in S_{k_{1}+k_{2}}} \operatorname{sgn}(\sigma) \omega\left(v_{\sigma_{1}}, v_{\sigma_{2}}, \ldots, v_{\sigma_{k}}\right) \otimes \eta\left(v_{\sigma_{(k+1)}}, v_{\sigma_{(k+2)}}, \ldots, v_{\sigma_{(k+\ell)}}\right)
$$

We want to globalize this algebra. Let $M$ be a smooth $n$-manifold and consider the set

$$
L_{\mathrm{alt}}^{k}(T M ; \mathrm{W})=\bigcup_{p \in M} L_{\mathrm{alt}}^{k}\left(T_{p} M ; \mathrm{W}\right)
$$

This set can easily be given a rather obvious vector bundle structure. In this setting, it is convenient to identify $L_{\text {alt }}^{k}\left(T_{p} M ; \mathrm{W}\right)$ with $\mathrm{W} \otimes\left(\Lambda^{k} T_{p}^{*} M\right)$, so that our bundle be identified with the vector bundle $\mathrm{W} \otimes\left(\bigwedge^{k} T^{*} M\right)$ whose fiber at $p$ is $\mathrm{W} \otimes\left(\bigwedge^{k} T_{p}^{*} M\right)$. The $C^{\infty}(M)$-module of sections of this bundle is denoted $\Omega^{k}(M, \mathrm{~W})$. Elements of $\Omega^{k}(M, \mathrm{~W})$ are called (smooth) W -valued $k$-forms.

We obtain an exterior product $\Omega^{k}(M, \mathrm{~W}) \times \Omega^{\ell}(M, \mathrm{~W}) \rightarrow \Omega^{k+\ell}(M, \mathrm{~W} \otimes \mathrm{~W})$ as usual by $(\alpha \wedge \beta)(p):=\alpha(p) \wedge \beta(p)$.

We give an alternative definition of $\wedge$ and leave it to the assiduous reader to show that the result is the same. Let $\mathrm{w}_{1}, \ldots, \mathrm{w}_{m}$ be a basis for W and $\omega \in \Omega^{k}(M, \mathrm{~W})$ and $\eta \in \Omega^{\ell}(M, \mathrm{~W})$. We have

$$
\omega=\sum_{i=1}^{m} \mathrm{w}_{i} \omega^{i} \text { and } \eta=\sum_{j=1}^{m} \mathrm{w}_{j} \eta^{j}
$$

for some $\omega^{i} \in \Omega^{k}(M)$ and $\eta^{j} \in \Omega^{\ell}(M)$, where we write $\mathrm{w}_{i} \omega^{i}$ rather than $\mathrm{w}_{i} \otimes \omega^{i}$, etc. Then,

$$
\omega \wedge \eta=\sum_{i=1}^{m} \sum_{j=1}^{m} \mathrm{w}_{i} \otimes \mathrm{w}_{j} \omega^{i} \wedge \eta^{j}
$$

One can show that this definition is independent of the choices and is consistent with the basis free definition.

We still have a pull-back operation defined as before so that if $f: N \rightarrow M$ is smooth and $\omega=\sum_{i=1}^{m} \mathrm{w}_{i} \omega^{i} \in \Omega^{k}(M, \mathrm{~W})$ for smooth $k$-forms $\omega^{i}$, then

$$
f^{*} \omega=\sum_{i=1}^{m} f^{*}\left(\mathrm{w}_{i} \omega^{i}\right)=\sum_{i=1}^{m} \mathrm{w}_{i} f^{*} \omega^{i}
$$

is an element of $\Omega^{k}(N, \mathrm{~W})$. We also define a natural exterior derivative

$$
d: \Omega^{k}(M, \mathrm{~W}) \rightarrow \Omega^{k+1}(M, \mathrm{~W})
$$

as follows: If $\omega \in \Omega^{k}(M, \mathrm{~W})$, then choose a basis as above and write $\omega= \sum_{i=1}^{m} \mathrm{w}_{i} \omega^{i}$. Then

$$
d \omega:=\sum_{i=1}^{m} \mathrm{w}_{i} d \omega^{i}
$$

If $\omega=\sum_{i=1}^{m} \widetilde{\mathrm{w}}_{i} \widetilde{\omega}^{i}$ where $\widetilde{\mathrm{w}}_{i}=\sum \mathrm{w}_{j} A_{i}^{j}$, then we must have

$$
\sum_{i=1}^{m} \widetilde{\mathrm{w}}_{i} \widetilde{\omega}^{i}=\sum_{i=1}^{m}\left(\sum_{j=1}^{m} \mathrm{w}_{j} A_{i}^{j}\right) \widetilde{\omega}^{i}=\sum_{j=1}^{m} \mathrm{w}_{j}\left(\sum_{i=1}^{m} A_{i}^{j} \widetilde{\omega}^{i}\right)
$$

from which it follows that

$$
\omega^{j}=\sum_{i=1}^{m} A_{i}^{j} \widetilde{\omega}^{i}
$$

and so

$$
\sum_{j=1}^{m} \mathrm{w}_{j} d \omega^{j}=\sum_{j=1}^{m} \mathrm{w}_{j} \sum_{i=1}^{m} A_{i}^{j} d \widetilde{\omega}^{i}=\sum_{j=1}^{m} \widetilde{\mathrm{w}}_{i} d \widetilde{\omega}^{i}
$$

Thus the definition does not depend on the choices. It turns out that the invariant definition is

$$
\begin{aligned}
& d \omega\left(X_{0}, X_{1}, \ldots, X_{k}\right) \\
& \quad=\sum_{0 \leq i \leq k}(-1)^{i} X_{i}\left(\omega\left(X_{0}, \ldots, \widehat{X}_{i}, \ldots, X_{k}\right)\right) \\
& \quad+\sum_{0 \leq i<j \leq k}(-1)^{i+j} \omega\left(\left[X_{i}, X_{j}\right], X_{0}, \ldots, \widehat{X}_{i}, \ldots, \widehat{X}_{j}, \ldots, X_{k}\right)
\end{aligned}
$$

where now $\omega\left(X_{0}, \ldots, \widehat{X}_{i}, \ldots, X_{k}\right)$ is a W -valued function. (A vector field $X$ acts on a W -valued function as $X f:=d f(X)$.)

If W happens to be an algebra, then the algebra product $\mathrm{W} \times \mathrm{W} \rightarrow \mathrm{W}$ is bilinear, and so it gives rise to a linear map $m: \mathrm{W} \otimes \mathrm{W} \rightarrow \mathrm{W}$. We compose the exterior product with this map to get an exterior product $\stackrel{m}{\wedge}: \Omega^{k}(M, \mathrm{~W}) \times \Omega^{\ell}(M, \mathrm{~W}) \rightarrow \Omega^{k+\ell}(M, \mathrm{~W})$. If $\omega \in \Omega^{k}(M, \mathrm{~W})$ and $\eta \in \Omega^{\ell}(M, \mathrm{~W})$, then as above we may write

$$
\omega=\sum_{i=1}^{m} \mathrm{w}_{i} \omega^{i} \text { and } \eta=\sum_{j=1}^{m} \mathrm{w}_{j} \eta^{j}
$$

for some $\omega^{i} \in \Omega^{k}(M)$ and $\eta^{j} \in \Omega^{\ell}(M)$ and

$$
\omega \stackrel{m}{\wedge} \eta=\sum_{i=1}^{m} \sum_{j=1}^{m} m\left(\mathrm{w}_{i} \otimes \mathrm{w}_{j}\right) \omega^{i} \wedge \eta^{j}
$$

Using a dot for the multiplication, we also have

$$
\begin{aligned}
& (\omega \wedge \eta)\left(X_{1}, X_{2}, \ldots, X_{r}, X_{r+1}, X_{r+2}, \ldots, X_{r+s}\right) \\
& \quad=\frac{1}{k_{1}!k_{2}!} \sum_{\sigma \in S_{k_{1}+k_{2}}} \operatorname{sgn}(\sigma) \omega\left(X_{\sigma_{1}}, X_{\sigma_{2}}, \ldots, X_{\sigma_{r}}\right) \cdot \eta\left(X_{\sigma_{r+1}}, X_{\sigma_{r+2}}, \ldots, X_{\sigma_{r+s}}\right)
\end{aligned}
$$

In this case $d$ is also defined as before. A particularly important case is when W is a Lie algebra $\mathfrak{g}$ with bracket $[\cdot, \cdot]$. Then we write the resulting product $\stackrel{m}{\wedge}$ as $[., \cdot]^{\wedge}$ or just $[., \cdot]$ when there is no risk of confusion. Thus if $\omega, \eta \in \Omega^{1}(U, \mathfrak{g})$ are Lie algebra-valued 1-forms, then

$$
[\omega, \eta]^{\wedge}(X, Y)=[\omega(X), \eta(Y)]+[\eta(X), \omega(Y)]
$$

In particular, $\frac{1}{2}[\omega, \omega]^{\wedge}(X, Y)=[\omega(X), \omega(Y)]$, which might not be zero in general!

Example 8.42. The Maurer-Cartan forms are $\mathfrak{g}$-valued 1 -forms.

\subsection*{Bundle-Valued Forms}
It is convenient in several contexts to have on hand the notion of a differential form with values in a vector bundle. Let $\xi=(E, \pi, M)$ be a smooth real vector bundle of rank $r$. We can consider the vector bundle $L_{\text {alt }}^{k}(T M, E)$ over $M$ whose fiber at $p$ is $L_{\text {alt }}^{k}\left(T_{p} M, E_{p}\right)$. We identify $L_{\text {alt }}^{k}\left(T_{p} M, E_{p}\right)$ with $E_{p} \otimes \bigwedge^{k} T_{p}^{*} M$ and thus the bundle is identified with $E \otimes \bigwedge^{k} T^{*} M$.

Definition 8.43. Let $\xi=(E, \pi, M)$ be a smooth vector bundle. A smooth differential $k$-form with values in $\xi$ (or values in $E$ ) is a smooth section of the bundle $E \otimes \bigwedge^{k} T^{*} M$. These are denoted by $\Omega^{k}(M ; E)$.

Remark 8.44. The reader should avoid confusion between $\Omega^{k}(M ; E)$ and the space of sections $\Gamma\left(M, \bigwedge^{k} E\right)$.

Theorem 8.45. There is a natural $C^{\infty}(M)$ module isomorphism $\Omega^{k}(M ; E) \cong L_{\text {alt }}^{k}(\mathfrak{X}(M), \Gamma(E))$. If this isomorphism is taken as an identification, then

$$
\mu\left(X_{1}, \ldots, X_{k}\right)(p)=\mu\left(X_{1}(p), \ldots, X_{k}(p)\right)
$$

for $\mu \in \Omega^{k}(M ; E)$ and $X_{1}, \ldots, X_{k} \in L_{\text {alt }}^{k}(\mathfrak{X}(M), \Gamma(E))$.\\
Proof. We use Proposition 6.55. The reader should check each of the following:

$$
\begin{aligned}
\Omega^{k}(M ; E) & \cong \Gamma\left(E \otimes \bigwedge^{k} T^{*} M\right) \cong \Gamma\left(L_{\mathrm{alt}}^{k}(T M, E)\right) \\
& \cong \Gamma\left(\operatorname{Hom}\left(\bigwedge^{k} T M, E\right)\right) \cong \operatorname{Hom}\left(\Gamma\left(\bigwedge^{k} T M\right), \Gamma(E)\right) \\
& \cong \operatorname{Hom}\left(\bigwedge^{k} \mathfrak{X}(M), \Gamma(E)\right) \cong L_{\mathrm{alt}}^{k}(\mathfrak{X}(M), \Gamma(E))
\end{aligned}
$$

In order to get a grip on the meaning of $\Omega^{k}(M ; E)$, let us exhibit transition functions. For a vector bundle, knowing the transition functions is tantamount to knowing how local expressions with respect to a frame transform as we change frames. A local frame for $E \otimes \bigwedge^{k} T^{*} M$ can be given by combining a local frame for $E$ with a local frame for $\bigwedge^{k} T^{*} M$. Let ( $e_{1}, \ldots, e_{r}$ ) be a frame field for $E$ defined on an open set $U$. We may as well take $U$ to also be a chart domain for the manifold $M$. Then any local section of $\Omega^{k}(M ; E)$ defined on $U$ has the form

$$
s=\sum a_{\vec{I}}^{j} e_{j} \otimes d x^{\vec{I}}
$$

for some smooth functions $a_{\vec{I}}^{j}=a_{i_{1} \ldots i_{k}}^{j}$ defined in $U$. Then for a new local set up with frames $\left(f_{1}, \ldots, f_{r}\right)$ and $d y^{\vec{I}}=d y^{i_{1}} \wedge \cdots \wedge d y^{i_{k}}\left(i_{1}<\cdots<i_{k}\right)$ we have

$$
s=\sum \widetilde{a}_{\vec{I}}^{j} f_{j} \otimes d y^{\vec{I}}
$$

for some $\widetilde{a}_{\vec{I}}^{j}$ and the transformation law

$$
\widetilde{a}_{\vec{I}}^{j}=\sum a_{\vec{J}}^{i} C_{i}^{j} \frac{\partial x^{\vec{J}}}{\partial y^{\vec{I}}},
$$

where $C_{i}^{j}$ is defined by $f_{s} C_{j}^{s}=e_{j}$.\\
Exercise 8.46. Derive the above transformation law.\\
Note that if we write $\omega^{j}=\sum a_{\vec{I}}^{j} d x^{\vec{I}}$, then we may write $s$ in $U$ as

$$
s=\sum e_{j} \otimes \omega^{j} \text { for } \omega^{j} \in \Omega^{k}(U) .
$$

Example 8.47. If $E$ is a trivial product bundle $M \times \mathrm{V} \rightarrow M$, then $\Omega^{k}(M ; E)$ is canonically isomorphic to $\Omega^{k}(M ; \mathrm{V})$.

Example 8.48. For an $n$-manifold $M$, the bundle map $I: T M \rightarrow T M$ which is the identity on each fiber can be interpreted as an element of $\Omega^{1}(M ; T M)$.

Example 8.49. If $f: M \rightarrow N$ is a smooth map, then the tangent map $T f: T M \rightarrow T N$ can be interpreted as an element of $\Omega^{1}\left(M ; f^{*} T N\right)$, where $f^{*} T N$ is the pull-back bundle (Definition 6.18).

Now we want to define an important graded module structure on the direct sum

$$
\Omega(M ; E)=\bigoplus_{k} \Omega^{k}(M ; E) .
$$

This will be a module over the graded algebra $\Omega(M)$. The action of $\Omega(M)$ on $\Omega(M ; E)$ is given by maps $\wedge: \Omega^{k}(M) \times \Omega^{\ell}(M ; E) \rightarrow \Omega^{k+\ell}(M ; E)$, which in turn are defined by extending the following rule linearly:

$$
\omega^{1} \wedge\left(s \otimes \omega^{2}\right):=s \otimes \omega^{1} \wedge \omega^{2} \text { for } \omega^{1} \in \Omega^{k}(M), \omega^{2} \in \Omega^{\ell}(M) \text { and } s \in \Gamma(E)
$$

(with the same formula for local sections). Actually we can also define the analogous right multiplication $\wedge: \Omega^{\ell}(M ; E) \times \Omega^{k}(M) \rightarrow \Omega^{k+\ell}(M ; E)$,

$$
\left(s \otimes \omega^{1}\right) \wedge \omega^{2}:=s \otimes \omega^{1} \wedge \omega^{2}
$$

and then we have

$$
\eta \wedge \mu=(-1)^{k \ell} \mu \wedge \eta \text { for } \mu \in \Omega^{\ell}(M ; E), \eta \in \Omega^{k}(M) .
$$

If the vector bundle is actually an algebra bundle, say $\mathcal{A} \rightarrow M$, then we may turn $\mathcal{A} \otimes \wedge T^{*} M:=\sum_{p=0}^{n} \mathcal{A} \otimes \bigwedge^{p} T^{*} M$ into an algebra bundle whose sections can be multiplied: For $\omega^{1} \in \Omega^{k}(M), \omega^{2} \in \Omega^{\ell}(M)$, and $s_{1}, s_{2} \in \Gamma(\mathcal{A})$, define

$$
\left(s_{1} \otimes \omega^{1}\right) \circledast\left(s_{2} \otimes \omega^{2}\right):=s_{1} \cdot s_{2} \otimes \omega^{1} \wedge \omega^{2}
$$

where $\cdot$ is the product in $\mathcal{A}$. This extends linearly on (possibly locally defined) sections:

$$
\left(\sum a_{j}^{i} s_{i} \otimes \omega^{j}\right) \circledast\left(\sum b_{\ell}^{k} s_{k} \otimes \omega^{\ell}\right):=\sum a_{j}^{i} b_{\ell}^{k} s_{i} s_{k} \otimes \omega^{j} \wedge \omega^{\ell},
$$

where $a_{j}^{i}$ and $b_{\ell}^{k}$ are smooth functions. We obtain a product that is natural with respect to restriction to open sets. From this the sections $\Omega(M, \mathcal{A})= \Gamma\left(M, \mathcal{A} \otimes \wedge T^{*} M\right)$ become an algebra over the ring of smooth functions. We can think of elements of $\Omega^{k}(M, \mathcal{A})$ as $C^{\infty}(M)$ multilinear maps on $\mathfrak{X}(M)$. Then the invariant formula for $\phi \in \Omega^{k}(M, \mathcal{A})$ and $\psi \in \Omega^{\ell}(M, \mathcal{A})$ is

$$
\begin{aligned}
& \phi \circledast \psi\left(X_{1}, \ldots, X_{k+\ell}\right) \\
& \quad=\frac{1}{k!\ell!} \sum_{\sigma} \operatorname{sgn}(\sigma) \phi\left(X_{\sigma_{1}}, \ldots, X_{\sigma_{k}}\right) \cdot \psi\left(X_{\sigma_{(k+1)}}, \ldots, X_{\sigma_{(k+\ell)}}\right)
\end{aligned}
$$

for $X_{1}, \ldots, X_{k+\ell} \in \mathfrak{X}(M)$. Depending on the context, the symbol may be chosen to be the same as that for the product in $\mathcal{A}$, or it may just be $\wedge$ (especially if $\mathcal{A}$ is commutative), or it may be some hybrid symbol.

Another important example is where $\mathcal{A}=\operatorname{End}(E)$. Locally, say on $U$, sections $\mu_{1}$ and $\mu_{2}$ of $\Omega(M, \operatorname{End}(E))$ take the form $\mu_{1}=\sum A_{i} \otimes \alpha^{i}$ and $\mu_{2}=\sum B_{i} \otimes \beta^{i}$, where $A_{i}$ and $B_{i}$ are maps $U \rightarrow \operatorname{End}(E)$. Thus for each $x \in U$, the $A_{i}$ and $B_{i}$ evaluate to give $A_{i}(x), B_{i}(x) \in \operatorname{End}\left(E_{x}\right)$. The multiplication is then

$$
\left(\sum A_{i} \otimes \alpha^{i}\right) \wedge\left(\sum B_{j} \otimes \beta^{j}\right)=\sum_{i, j} A_{i} B_{j} \otimes \alpha^{i} \wedge \beta^{j}
$$

where the $A_{i} B_{j}: U \rightarrow \operatorname{End}(E)$ are local sections given by composition:

$$
A_{i} B_{j}: x \mapsto A_{i}(x) \circ B_{j}(x) .
$$

Exercise 8.50. Show that $\Omega(M, \operatorname{End}(E))$ acts on $\Omega(M, E)$ making $\Omega(M, E)$ a module over the algebra $\Omega(M, \operatorname{End}(E))$.

Perhaps it would help to think as follows: We have a cover of a manifold $M$ by open sets $\left\{U_{\alpha}\right\}$ that simultaneously locally trivialize both $E$ and $T^{*} M$. Then these also give local trivializations over these open sets of the bundles $\operatorname{End}(E)$ and $\Lambda T^{*} M$. Associated with each local trivialization is a frame field for $E \rightarrow M$, say $\left(e_{1}, \ldots, e_{r}\right)$, which allows us to associate with each section $\sigma \in \Omega^{k}(M, E)$ an $r$-tuple of $k$-forms $\sigma_{U}=\left(\sigma_{U}^{i}\right)$ for each $U$ such that $\sigma=\sum \sigma_{U}^{i} e_{i}$. Similarly, a section $A \in \Omega^{\ell}(M, \operatorname{End}(E))$ is equivalent to assigning to each open set $U \in\left\{U_{\alpha}\right\}$ a matrix of $\ell$-forms $A_{U}$. The algebra structure on $\Omega(M, \operatorname{End}(E))$ is then just matrix multiplication, where the entries are multiplied using the exterior product $A_{U} \wedge B_{U}$,

$$
\left(A_{U} \wedge B_{U}\right)_{j}^{i}=\sum A_{s}^{i} \wedge B_{j}^{s} .
$$

The module structure of the above exercise is given locally by $\sigma_{U} \mapsto A_{U} \wedge \sigma_{U}$. Where did the bundle go? The global topology is now encoded in the transformation laws, which tell us what the same section looks like when we change to a new frame field on an overlap $U_{\alpha} \cap U_{\beta}$. In this sense, the bundle is a combinatorial recipe for pasting together local objects.

Recall that with the product $[A, B]:=A \circ B-B \circ A$, the bundle $\operatorname{Hom}(E, E)$ is written as $\mathfrak{g l}(E)$ rather than $\operatorname{End}(E)$, and then our general construction gives $\Omega(M, \mathfrak{g l}(E))$ an algebra structure, whose product is denoted $[\phi, \psi]^{\wedge}$ or something similar.

\subsection*{Operator Interactions}
The Lie derivative acts on differential forms since the latter are, from one viewpoint, alternating tensor fields. When we apply the Lie derivative to a differential form, we get a differential form, so we should think about the Lie derivative in the context of differential forms.

Lemma 8.51. For any $X \in \mathfrak{X}(M)$ and any $f \in \Omega^{0}(M)$, we have $\mathcal{L}_{X} d f= d \mathcal{L}_{X} f$.

Proof. For a function $f$, we compute as

$$
\begin{aligned}
\left(\mathcal{L}_{X} d f\right)(Y) & =\left(\left.\frac{d}{d t}\right|_{0}\left(\varphi_{t}^{X}\right)^{*} d f\right)(Y)=\left.\frac{d}{d t}\right|_{0} d f\left(T \varphi_{t}^{X} \cdot Y\right)=\left.\frac{d}{d t}\right|_{0} Y\left(\left(\varphi_{t}^{X}\right)^{*} f\right) \\
& =Y\left(\left.\frac{d}{d t}\right|_{0}\left(\varphi_{t}^{X}\right)^{*} f\right)=Y\left(\mathcal{L}_{X} f\right)=d\left(\mathcal{L}_{X} f\right)(Y)
\end{aligned}
$$

where $Y \in \mathfrak{X}(M)$ is arbitrary.\\
We now have two ways to differentiate sections in $\Omega(M)$. First, there is the Lie derivative $\mathcal{L}_{X}: \Omega^{i}(M) \rightarrow \Omega^{i}(M)$, which turns out to be a graded derivation of degree zero,

$$
\mathcal{L}_{X}(\alpha \wedge \beta)=\mathcal{L}_{X} \alpha \wedge \beta+\alpha \wedge \mathcal{L}_{X} \beta
$$

We may apply $\mathcal{L}_{X}$ to elements of $\Omega(U)$ for $U \subset M$, and it is easy to see that we obtain a natural derivation in the sense of Definition 8.33.

Exercise 8.52. Prove the above product rule.\\
Second, there is the exterior derivative $d$ which is a graded derivation of degree 1. In order to relate the two operations, we need a third map, which, like the Lie derivative, is taken with respect to a given field $X \in \mathfrak{X}(M)$. This map is defined using the interior product given in Definition 8.12 by letting

$$
i_{X} \omega\left(X_{1}, \ldots, X_{i-1}\right)(p):=i_{X_{p}} \omega_{p}\left(X_{1}(p), \ldots, X_{i-1}(p)\right) .
$$

Alternatively, if $\omega \in \Omega^{i}(M)$ is viewed as a skew-symmetric multilinear map from $\mathfrak{X}(M) \times \cdots \times \mathfrak{X}(M)$ to $C^{\infty}(M)$, then we simply define

$$
i_{X} \omega\left(X_{1}, \ldots, X_{i-1}\right):=\omega\left(X, X_{1}, \ldots, X_{i-1}\right) .
$$

By convention $i_{X} f=0$ for $f \in C^{\infty}(M)$. We will call this operator the interior product or contraction operator. This operator is clearly linear over $\mathbb{R}$. Notice that for any $f \in C^{\infty}(M)$ we have

$$
i_{f X} \omega=f i_{X} \omega, \text { and } i_{X} d f=d f(X)=\mathcal{L}_{X} f
$$

Proposition 8.53. $i_{X}$ is a graded derivation of $\Omega(M)$ of degree -1 :

$$
i_{X}(\alpha \wedge \beta)=\left(i_{X} \alpha\right) \wedge \beta+(-1)^{k} \alpha \wedge\left(i_{X} \beta\right) \text { for } \alpha \in \Omega^{k}(M) \text {. }
$$

It is the unique degree -1 graded derivation of $\Omega(M)$ such that $i_{X} f=0$ for $f \in \Omega^{0}(M)$ and

$$
i_{X} \theta=\theta(X) \text { for } \theta \in \Omega^{1}(M) \text { and } X \in \mathfrak{X}(M) .
$$

Proof. This follows from Proposition 8.13.\\
Actually, $i_{X}$ is natural with respect to restriction, so it is a natural graded derivation in the sense of Definition 8.33. Formulas developed for the interior product in the vector space category also hold for vector fields and differential forms. For example, if $\theta_{1}, \ldots, \theta_{k} \in \Omega^{1}(M)$ and $X \in \mathfrak{X}(M)$, then

$$
i_{X}\left(\theta_{1} \wedge \cdots \wedge \theta_{k}\right)=\sum_{\ell=1}^{k}(-1)^{k+1} \theta_{\ell}(X) \theta_{1} \wedge \cdots \wedge \widehat{\theta}_{\ell} \wedge \cdots \wedge \theta_{k}
$$

Notation 8.54. Other notations for $i_{X} \omega$ include $\left.X\right\rfloor \omega$ and $\langle X, \omega\rangle$. These notations make the following theorem look more natural:

Theorem 8.55. The Lie derivative is a derivation with respect to the pairing $(X, \omega) \mapsto\langle X, \omega\rangle$. That is,

$$
\mathcal{L}_{X}\left(i_{Y} \omega\right)=i_{\mathcal{L}_{X}} Y \omega+i_{Y} \mathcal{L}_{X} \omega,
$$

or in alternative notations,

$$
\begin{aligned}
\left.\mathcal{L}_{X}(Y\rfloor \omega\right) & \left.\left.=\left(\mathcal{L}_{X} Y\right)\right\rfloor \omega+Y\right\rfloor\left(\mathcal{L}_{X} \omega\right), \\
\mathcal{L}_{X}\langle Y, \omega\rangle & =\left\langle\mathcal{L}_{X} Y, \omega\right\rangle+\left\langle Y, \mathcal{L}_{X} \omega\right\rangle .
\end{aligned}
$$

Proof. Exercise.\\
Now we can relate the Lie derivative, the exterior derivative and the contraction operator.

Theorem 8.56. Let $X \in \mathfrak{X}_{M}$. Then we have Cartan's formula,

$$
\mathcal{L}_{X}=d \circ i_{X}+i_{X} \circ d
$$

Proof. Both sides of the equation define derivations of degree zero (use Proposition 8.35). So by Lemma 8.36 we just have to check that they agree on functions and exact 1 -forms. On functions we have $i_{X} f=0$ and $i_{X} d f= X f=\mathcal{L}_{X} f$ so the formula holds. On differentials of functions we have

$$
\left(d \circ i_{X}+i_{X} \circ d\right) d f=\left(d \circ i_{X}\right) d f=d \mathcal{L}_{X} f=\mathcal{L}_{X} d f
$$

where we have used Lemma 8.51 in the last step. \(\square\)

As a corollary, we can extend Lemma 8.51:\\
Corollary 8.57. $d \circ \mathcal{L}_{X}=\mathcal{L}_{X} \circ d$.\\
Proof. We have

$$
\begin{aligned}
d \mathcal{L}_{X} \alpha & =d\left(d i_{X}+i_{X} d\right)(\alpha)=d i_{X} d \alpha \\
& =d i_{X} d \alpha+i_{X} d d \alpha=\left(\mathcal{L}_{X} \circ d\right) \alpha
\end{aligned}
$$ \(\square\)

Corollary 8.58. We have the following formulas:\\
(i) $i_{[X, Y]}=\mathcal{L}_{X} \circ i_{Y}-i_{Y} \circ \mathcal{L}_{X} ;$\\
(ii) $\mathcal{L}_{f X} \omega=f \mathcal{L}_{X} \omega+d f \wedge i_{X} \omega$ for all $\omega \in \Omega(M)$.

Proof. We leave (i) as Problem 9. For (ii), we compute:

$$
\begin{aligned}
\mathcal{L}_{f X} \omega & =i_{f X} d \omega+d\left(i_{f X} \omega\right)=i_{f X} d \omega+d\left(f i_{X} \omega\right) \\
& =f i_{X} d \omega+d f \wedge i_{X} \omega+f d\left(i_{X} \omega\right) \\
& =f\left(i_{X} d \omega+d\left(i_{X} \omega\right)\right)+d f \wedge i_{X} \omega=f \mathcal{L}_{X} \omega+d f \wedge i_{X} \omega
\end{aligned}
$$ \(\square\)

\subsection*{Orientation}
A vector bundle $E \rightarrow M$ is called oriented if every fiber $E_{p}$ is given a smooth choice of orientation. There are several equivalent ways to make a rigorous definition:

Proposition 8.59. Let $E \rightarrow M$ be a rank $k$ real vector bundle with typical fiber V . The following are equivalent:\\
(i) There is a smooth global section $\omega$ of the bundle $\Lambda^{k} E^{*} \cong L_{\text {alt }}^{k}(E) \rightarrow M$ such that $\omega$ is nowhere vanishing.\\
(ii) There is a smooth global section $s$ of the bundle $\bigwedge^{k} E \rightarrow M$ such that $s$ is nowhere vanishing.\\
(iii) The vector bundle has an atlas of VB-charts (local trivializations) such that the corresponding transition maps take values in $\mathrm{GL}^{+}(\mathrm{V})$, the group of positive determinant elements of $\mathrm{GL}(\mathrm{V})$. This means that the standard $\mathrm{GL}(\mathrm{V})$-structure on $E \rightarrow M$ can be reduced to $a \mathrm{GL}^{+}(\mathrm{V})$-structure (refer to Chapter 6).

Proof. We show that (i) is equivalent to (iii) and leave the rest as an easy exercise. Suppose that (i) holds and that $\omega$ is a nonvanishing section of $L_{\text {alt }}^{k}(E)$. Now fix a basis ( $\mathbf{e}_{1}, \ldots, \mathbf{e}_{k}$ ) on V and recall that with this basis fixed, each VB-chart ( $U, \phi$ ) for $E \rightarrow M$ corresponds to a local frame field. Indeed, we let $e_{i}(p):=\phi^{-1}\left(p, \mathbf{e}_{i}\right)$. The transition maps between two charts will have values in $\mathrm{GL}^{+}(\mathrm{V})$ exactly when the matrix function that relates the corresponding frame fields has positive determinant (check this). Given a VB-atlas we construct a new atlas. We retain those VB-charts ( $U, \phi$ ) whose corresponding frame field $e_{1}, \ldots, e_{k}$ satisfies $\omega\left(e_{1}, \ldots, e_{k}\right)>0$. For the charts for which $\omega\left(e_{1}, \ldots, e_{k}\right)<0$, we replace $e_{1}, \ldots, e_{k}$ by $-e_{1}, \ldots, e_{k}$, and the resulting chart will be included in our new atlas. Now if $e_{1}, \ldots, e_{k}$ and $f_{1}, \ldots, f_{k}$ are two frame fields coming from this atlas, then $f_{i}=\sum C_{i}^{j} e_{j}$ and

$$
\omega\left(f_{1}, \ldots, f_{k}\right)=(\operatorname{det} C) \omega\left(e_{1}, \ldots, e_{k}\right)
$$

We conclude that $\operatorname{det} C>0$.\\
Conversely, suppose the vector bundle has an atlas $\left\{\left(U_{\alpha}, \phi_{\alpha}\right)\right\}$ taking values in $\mathrm{GL}^{+}(\mathrm{V})$. We will use the frame fields coming from this atlas to construct a nowhere vanishing section of $L_{\text {alt }}^{k}(E) \rightarrow M$. If $e_{1}, \ldots, e_{k}$ and $f_{1}, \ldots, f_{k}$ are two frame fields coming from this atlas, then let $f^{1}, \ldots, f^{k}$ be dual to $f_{1}, \ldots, f_{k}$ and consider $f^{1} \wedge \cdots \wedge f^{k}$. We have $f_{i}=\sum C_{i}^{j} e_{j}$ and

$$
\left(f^{1} \wedge \cdots \wedge f^{k}\right)\left(e_{1}, \ldots, e_{k}\right)=\operatorname{det} C>0
$$

For each chart ( $U_{\alpha}, \phi_{\alpha}$ ) in our VB-atlas, let $f_{1}^{\alpha}, \ldots, f_{k}^{\alpha}$ be the corresponding frame field and let $\left(f_{\alpha}^{1}, \ldots, f_{\alpha}^{k}\right)$ be the dual frame field. Then let $\left\{\rho_{\alpha}\right\}$ be a partition of unity subordinate to the cover $\left\{U_{\alpha}\right\}$. Let

$$
\omega:=\sum_{\alpha} \rho_{\alpha} f_{\alpha}^{1} \wedge \cdots \wedge f_{\alpha}^{k}
$$

Then, $\omega$ is nowhere vanishing. To see this let $p \in M$ and suppose that $p \in U_{\beta}$ for some chart ( $U_{\beta}, \phi_{\beta}$ ) from the $\mathrm{GL}^{+}(\mathrm{V})$-valued atlas. Then

$$
\omega\left(f_{1}^{\beta}, \ldots, f_{k}^{\beta}\right)(p)=\sum_{\alpha} \rho_{\alpha}(p) \operatorname{det} C_{\beta \alpha}(p)>0
$$

where $C_{\beta \alpha}$ is the matrix that relates $f_{1}^{\beta}, \ldots, f_{k}^{\beta}$ and $f_{1}^{\alpha}, \ldots, f_{k}^{\alpha}$.\\
Definition 8.60. If any one (and hence all) of the conditions in Proposition 8.59 hold, then $E \rightarrow M$ is said to be orientable. A VB-atlas that satisfies (iii) will be called an oriented atlas.

If $E \rightarrow M$ is orientable as above, then two nowhere vanishing sections of $L_{\text {alt }}^{k}(E)$, say $\omega_{1}$ and $\omega_{2}$, are said to be equivalent if $\omega_{1}=f \omega_{2}$, where $f$ is a smooth positive function. We denote the equivalence class of such a nowhere vanishing $\omega$ by $[\omega]$.

Definition 8.61. An orientation for an orientable vector bundle of rank $k$ is an equivalence class $[\omega]$ of nowhere vanishing sections of $\Lambda^{k} E^{*}$. If such an orientation is chosen, then the vector bundle is said to be oriented by $[\omega]$.

Notice that if we have two oriented VB-atlases on a vector bundle, then we know what it means for them to determine the same $\mathrm{GL}^{+}(\mathrm{V})$-structure. This was the notion of strict equivalence from Chapter 6. The next exercise shows that the notion of a reduction to a $\mathrm{GL}^{+}(\mathrm{V})$-structure is equivalent to the notion of an orientation as we have defined it.

Exercise 8.62. Recall the construction of a nowhere vanishing section $\omega$ in the proof of Proposition 8.59. Show that the class $[\omega]$ does not depend on the partition of unity used in the construction. Show that if two oriented VBatlases determine the same $\mathrm{GL}^{+}(\mathrm{V})$-structure, then the constructed sections are equivalent and so determine the same orientation. Conversely, show that an orientation as we have defined it determines a unique reduction to a $\mathrm{GL}^{+}(\mathrm{V})$-structure on the vector bundle.

Let $E \rightarrow M$ be oriented by $[\omega]$. A frame $\left(v_{1}, \ldots, v_{k}\right)$ of fiber $E_{p}$ is positively oriented (or just positive) with respect to $[\omega]$ if and only if $\omega(p)\left(v_{1}, \ldots, v_{k}\right)>0$. This condition is independent of the choice of representative $\omega$ for the class $[\omega]$.

Definition 8.63. Let $\pi: E \rightarrow M$ be an oriented vector bundle. A frame field ( $f_{1}, \ldots, f_{k}$ ) over an open set $U$ is called a positively oriented frame field if $\left(f_{1}(p), \ldots, f_{k}(p)\right)$ is a positively oriented basis of $E_{p}$ for each $p \in U$.

Exercise 8.64. Let $\pi: E \rightarrow M$ be an oriented vector bundle. Show that if $M$ is connected, then there are exactly two possible orientations for the vector bundle.

Exercise 8.65. If $\pi_{1}: E_{1} \rightarrow M$ and $\pi_{2}: E_{2} \rightarrow M$ are orientable, then so is the Whitney sum $\pi_{1} \oplus \pi_{2}: E_{1} \oplus E_{2} \rightarrow M$.

Definition 8.66. A smooth manifold $M$ is said to be orientable if $T M$ is orientable. An orientation for the vector bundle $T M$ is also called an orientation for $M$. A manifold $M$ together with an orientation for $M$ is said to be an oriented manifold.

Definition 8.67. An atlas $\left\{\left(U_{\alpha}, \mathrm{x}_{\alpha}\right)\right\}$ for $M$ is said to be an oriented atlas if the associated frame fields are positively oriented. If this atlas is positively oriented with respect to an orientation on $M$ (an orientation $[\omega]$ of $T M$ ), then we call $\left\{\left(U_{\alpha}, \mathrm{x}_{\alpha}\right)\right\}$ a positively oriented atlas.

It follows from the definitions that an oriented atlas induces an orientation for which it is a positively oriented atlas. For this reason, a choice of\\
oriented atlas is equivalent to a choice of orientation, and so one often sees an orientation of an orientable manifold defined as simply being given by a choice of oriented atlas.

Now let $M$ be an $n$-manifold. Consider a top form, i.e. an $n$-form $\varpi \in \Omega^{n}(M)$, and assume that $\varpi$ is nowhere vanishing. Thus $M$ must be orientable. We call such a nonvanishing $\varpi$ a volume form for $M$, and every such volume form obviously determines an orientation for $M$. If $\varphi: M \rightarrow M$ is a diffeomorphism, then we must have that $\varphi^{*} \varpi=\delta \varpi$ for some $\delta \in C^{\infty}(M)$, which we will call the Jacobian determinant of $\varphi$ with respect to the volume element $\varpi$ :

$$
\varphi^{*} \varpi=J_{\varpi}(\varphi) \varpi .
$$

Clearly $J_{\varpi}(\varphi)$ is a nowhere vanishing smooth function.\\
Proposition 8.68. Let ( $M,[\varpi]$ ) be an oriented $n$-manifold. The sign of $J_{\varpi}(\varphi)$ is independent of the choice of volume form $\varpi$ in the orientation class $[\varpi]$.

Proof. Let $\varpi^{\prime} \in \Omega^{n}(M)$. We have $\varpi=a \varpi^{\prime}$ for some function $a$ that is never zero on $U$. Furthermore,

$$
J(\varphi) \varpi=\left(\varphi^{*} \varpi\right)=(a \circ \varphi)\left(\varphi^{*} \varpi^{\prime}\right)=(a \circ \varphi) J_{\varpi^{\prime}}(\varphi) \varpi^{\prime}=\frac{a \circ \varphi}{a} \varpi
$$

and since $\frac{a \circ \varphi}{a}>0$ and $\varpi$ is nonzero, the conclusion follows.\\
Let us consider a very important special case of this: Suppose that $\varphi$ : $U \rightarrow U$ is a diffeomorphism and $U \subset \mathbb{R}^{n}$. Then letting $\varpi_{0}=d u^{1} \wedge \cdots \wedge d u^{n}$ we have for any $x \in U$,

$$
\begin{aligned}
\varphi^{*} \varpi_{0}(x) & =\varphi^{*} d u^{1} \wedge \cdots \wedge \varphi^{*} d u^{n}(x) \\
& =\left(\left.\sum \frac{\partial\left(u^{1} \circ \varphi\right)}{\partial u^{i_{1}}}\right|_{x} d u^{i_{1}}\right) \wedge \cdots \wedge\left(\left.\sum \frac{\partial\left(u^{n} \circ \varphi\right)}{\partial u^{i_{n}}}\right|_{x} d u^{i_{n}}\right) \\
& =\operatorname{det}\left(\frac{\partial\left(u^{i} \circ \varphi\right)}{\partial u^{j}}(x)\right) \varpi_{0}(x)=J \varphi(x) \varpi_{0}(x)
\end{aligned}
$$

So in this case, $J_{\varpi_{0}}(\varphi)$ is just the usual Jacobian determinant of $\varphi$. More generally, let a nonvanishing top form $\varpi$ be defined on $M$ and let $\varpi^{\prime}$ be another such form defined on $N$. Then we say that a diffeomorphism $\varphi: M \rightarrow N$ is orientation preserving (or positive) with respect to the orientations determined by $\varpi$ and $\varpi^{\prime}$ if the unique function $J_{\varpi, \varpi^{\prime}}$ such that $\varphi^{*} \varpi^{\prime}=J_{\varpi, \varpi^{\prime}} \varpi$ is strictly positive on $M$.

Exercise 8.69. An open subset of an oriented manifold $M$ inherits the orientation from $M$ since we can just restrict a defining volume form. Show that a chart ( $U, \mathrm{x}$ ) on an oriented manifold is positive if and only if\\
$\mathrm{x}: U \rightarrow \mathrm{x}(U)$ is orientation preserving. Here, $\mathrm{x}(U)$ inherits its orientation from the ambient Euclidean space with its standard orientation.

We now construct a two-fold covering manifold $M^{\text {or }}$ for any manifold $M$. The orientation cover will itself always be orientable. Recall that the zero section of a vector bundle over $M$ is a submanifold of the total space diffeomorphic to $M$. Consider the vector bundle whose total space is $\bigwedge^{n} T^{*} M$ and remove the zero section to obtain

$$
\left(\bigwedge^{n} T^{*} M\right)^{\times}:=\left(\bigwedge^{n} T^{*} M\right) \backslash\{\text { zero section }\} .
$$

Define an equivalence relation on $\left(\bigwedge^{n} T^{*} M\right)^{\times}$by declaring $\nu_{1} \sim \nu_{2}$ if and only if $\nu_{1}$ and $\nu_{2}$ are in the same fiber and if $\nu_{1}=a \nu_{2}$ with $a>0$. The space of equivalence classes is denoted $M^{\text {or }}$ and we will show that it is a smooth manifold. Let $\mathrm{q}:\left(\bigwedge^{n} T^{*} M\right)^{\times} \rightarrow M^{\text {or }}$ be the quotient map and give $M^{\text {or }}$ the quotient topology. There is a unique smooth map $\pi_{o r}$ making the following diagram commute:\\
\includegraphics[scale=0.25, center]{2025_10_09_265d7e1a22c89f79c21eg-394}

It is easy to see that for each $p \in M$, the set $\pi_{\text {or }}^{-1}(p)$ contains exactly two elements, which are the two orientations of $T_{p} M$. We give the set $M^{\text {or }}$ a smooth structure. First let $\left[\mu_{0}\right]$ be the standard orientation of $\mathbb{R}^{n}$ and choose a fixed orientation reversing linear involution $r_{0}: \mathbb{R}^{n} \rightarrow \mathbb{R}^{n}$. Let $\left\{\left(U_{\alpha}, \mathrm{x}_{\alpha}\right)\right\}$ be an atlas for $M$. By composing some of the charts with $r_{0}$ and adding the resulting charts to the atlas, we may suppose that $\left\{\left(U_{\alpha}, \mathrm{x}_{\alpha}\right)\right\}$ has the property that for every chart ( $U, \mathrm{x}$ ) in the atlas there is a chart ( $U, \mathrm{y}$ ) in the atlas with the same domain such that $\mathrm{y} \circ \mathrm{x}^{-1}$ is orientation reversing. Let us say that such an atlas is "balanced". (The maximal atlas is obviously balanced.) Now for each chart ( $U, \mathrm{x}$ ), where $\mathrm{x}=\left(x^{1}, \ldots, x^{n}\right)$, define a map $\phi_{\mathrm{x}}: \mathrm{x}(U) \rightarrow M^{\text {or }}$ by

$$
\phi_{\mathbf{x}}(u):=\left[d x^{1}\left(\mathbf{x}^{-1}(u)\right) \wedge \cdots \wedge d x^{n}\left(\mathbf{x}^{-1}(u)\right)\right] \in M^{\text {or }} .
$$

Because we have assumed that the atlas is balanced, it is easy to see that each element of $M^{\text {or }}$ is in the image of some $\phi_{\mathrm{x}}$. The $\phi_{\mathrm{x}}$ are injective and so are bijections onto their images. These maps are local parameterizations of $M^{\text {or }}$ and the inverses of these maps are charts on $M^{\text {or }}$. Let us denote the chart arising from ( $U, \mathrm{x}$ ) by ( $\widetilde{U}, \widetilde{\mathrm{x}}$ ). Thus $\widetilde{U}=\phi_{\mathrm{x}}(U)$ and $\widetilde{\mathrm{x}}=\phi_{\mathrm{x}}^{-1}$. In this way the balanced atlas $\left\{\left(U_{\alpha}, \mathrm{x}_{\alpha}\right)\right\}$ gives an atlas on $M^{\text {or }}$. We must check that the overlaps are smooth. So let $\widetilde{U}_{\alpha} \cap \widetilde{U}_{\beta} \neq \emptyset$ and consider $\widetilde{\mathrm{x}}_{\beta} \circ \widetilde{\mathrm{x}}_{\alpha}^{-1}=\widetilde{\mathrm{x}}_{\beta} \circ \phi_{\mathrm{x}_{\alpha}}$. For $u \in \mathrm{x}_{\alpha}\left(U_{\alpha} \cap U_{\beta}\right)$ and $w=\mathrm{x}_{\beta} \circ \mathrm{x}_{\alpha}^{-1}(u)$ we have

$$
d x_{\alpha}^{1}\left(\mathrm{x}_{\alpha}^{-1}(u)\right) \wedge \cdots \wedge d x_{\alpha}^{n}\left(\mathrm{x}_{\alpha}^{-1}(u)\right)=\lambda d x_{\beta}^{1}\left(\mathrm{x}_{\beta}^{-1}(w)\right) \wedge \cdots \wedge d x_{\beta}^{n}\left(\mathrm{x}_{\beta}^{-1}(w)\right)
$$

for some $\lambda>0$. Indeed, we must have $\lambda=\operatorname{det}\left(D\left(\mathrm{x}_{\alpha} \circ \mathrm{x}_{\beta}^{-1}\right)(u)\right)$. Then,

$$
\begin{aligned}
\widetilde{\mathrm{x}}_{\beta} \circ \phi_{\mathrm{x}_{\alpha}}(u) & =\widetilde{\mathrm{x}}_{\beta}\left(\left[d x_{\alpha}^{1}\left(\mathrm{x}_{\alpha}(u)\right) \wedge \cdots \wedge d x_{\alpha}^{n}\left(\mathrm{x}_{\alpha}(u)\right)\right]\right) \\
& =\widetilde{\mathrm{x}}_{\beta}\left(\left[d x_{\beta}^{1}\left(\mathrm{x}_{\beta}(w)\right) \wedge \cdots \wedge d x_{\beta}^{n}\left(\mathrm{x}_{\beta}(w)\right)\right]\right) \\
& =w=\mathrm{x}_{\beta} \circ \mathrm{x}_{\alpha}^{-1}(u)
\end{aligned}
$$

Thus $\widetilde{\mathrm{x}}_{\beta} \circ \widetilde{\mathrm{x}}_{\alpha}^{-1}=\mathrm{x}_{\beta} \circ \mathrm{x}_{\alpha}^{-1}$ on $\widetilde{\mathrm{x}}_{\alpha}\left(\widetilde{U}_{\alpha} \cap \widetilde{U}_{\beta}\right)$. Note that for each $\alpha$ and $\beta$ the set $\widetilde{\mathrm{x}}_{\alpha}\left(\widetilde{U}_{\alpha} \cap \widetilde{U}_{\beta}\right)$ consists of exactly those connected components of $\mathrm{x}_{\alpha}\left(U_{\alpha} \cap U_{\beta}\right)$ on which $\operatorname{det}\left(D\left(\mathrm{x}_{\beta} \circ \mathrm{x}_{\alpha}^{-1}\right)\right)$ is positive. Thus $\widetilde{\mathrm{x}}_{\alpha}\left(\widetilde{U}_{\alpha} \cap \widetilde{U}_{\beta}\right)$ is open. One may check that the topology induced by this atlas coincides with the quotient topology and that the quotient map is smooth. Also, note that if $\widetilde{U}_{\alpha} \cap \widetilde{U}_{\beta} \neq \emptyset$, then $\operatorname{det}\left(D\left(\widetilde{\mathrm{x}}_{\beta} \circ \widetilde{\mathrm{x}}_{\alpha}^{-1}\right)\right)>0$ and so our atlas is oriented! We thus have a canonical orientation of $M^{\text {or }}$.

For each admissible chart ( $U, \mathrm{x}$ ), we obtain a section $U \rightarrow M^{\text {or }}$ given by

$$
p \mapsto\left[d x^{1}(p) \wedge \cdots \wedge d x^{n}(p)\right]
$$

It is easy to see that such sections are smooth. From the existence of these sections, one surmises that the connected components of the chart domains are evenly covered and so $\pi_{\text {or }}: M^{\text {or }} \rightarrow M$ is a two-fold covering map (but $M^{\text {or }}$ may not be connected). This two-fold covering map is called the orientation covering map. The space $M^{\text {or }}$ itself is called the orientation (double) cover.

Exercise 8.70. Let $R^{+}$be the multiplicative Lie group of positive real numbers. Show that $R^{+}$acts freely and properly on $\left(\bigwedge^{n} T^{*} M\right)^{\times}$and that the quotient manifold is $M^{\text {or }}$. Show that the canonical orientation on $M^{\text {or }}$ can be described as follows: Each $y \in M^{\mathrm{or}}$ is an orientation of $T_{\pi_{\mathrm{or}}(y)} M$. Since $T_{y} \pi_{\text {or }}: T_{y} M^{\text {or }} \rightarrow T_{\pi_{\text {or }}(y)} M$ is an isomorphism, we may transfer the orientation $y$ on $T_{\pi_{\mathrm{or}}(y)} M$ to an orientation for $T_{y} M^{\mathrm{or}}$. This gives a canonical orientation on each fiber of $T M^{\text {or }}$ which agrees with the orientation derived from the oriented atlas described above. How can we get a smooth global nonvanishing top form that induced this orientation?

Exercise 8.71. A manifold $M$ is orientable if and only if $M^{\text {or }}$ is disconnected. If $M$ is oriented, then $M^{\text {or }}$ has exactly two components and an orientation of $M$ corresponds to a choice of connected component of $M^{\text {or }}$.\\
8.7.1. Orientation induced on a boundary. Now let $M$ be a manifold with boundary. According to Problem 15, we can consider vector bundles over $M$. The definitions of orientable and orientation make sense for $M$. Here we wish to consider orientations of $\partial M$. If $[\omega]$ is an orientation for an orientable vector bundle $\pi: E \rightarrow M$, then $\left[\left.\omega\right|_{\partial M}\right]$ is an orientation on $\left.E\right|_{\partial M} \rightarrow \partial M$ where $\left.E\right|_{\partial M}=\pi^{-1}(\partial M)$. In particular, if $M$ is oriented by\\
$[\omega]$, then we may obtain an orientation on $\left.T M\right|_{\partial M} \rightarrow \partial M$. But what we would really like to obtain at this time is an orientation on $\partial M$. In other words, we need an orientation on the bundle $T(\partial M)$. First notice that by our conventions, an atlas for $M$ consists of charts that take values in half-spaces of the form $\mathbb{R}_{\lambda \geq c}^{n}$.

Exercise 8.72. Consider $v_{p} \in T_{p} M$ for some boundary point $p \in \partial M$. Let ( $U, \mathrm{x}$ ) be a half-space chart containing $p$ and taking values in $\mathbb{R}_{\lambda \geq c}^{n}$ and let ( $V, \mathrm{y}$ ) be another chart containing $p$, but taking values in $\mathbb{R}_{\mu \geq d}^{n}$. Then $d \mathrm{x} \circ d \mathrm{y}^{-1}$ maps $\mathbb{R}_{\mu<0}^{n}$ diffeomorphically onto $\mathbb{R}_{\lambda<0}^{n}$. Thus $\lambda \circ d \mathrm{x}\left(v_{p}\right)<0$ if and only if $\mu \circ d \mathrm{y}\left(v_{p}\right)<0$. Similarly, $\lambda \circ d \mathrm{x}\left(v_{p}\right)>0$ if and only if $\mu \circ d \mathrm{y}\left(v_{p}\right)>0$.

In light of the exercise above, the following definition makes sense:\\
Definition 8.73. Let $p \in \partial M$. A vector $\left.v_{p} \in T M_{p} \subset T M\right|_{\partial M}$ is called outward pointing if in some $\mathbb{R}_{\lambda \geq c}^{n}$-valued chart ( $U, \mathrm{x}$ ) we have $\lambda \circ d \mathrm{x}\left(v_{p}\right)<$ 0 . A smooth section $\chi$ of $\left.T M\right|_{\partial M}$ is called outward pointing if $\chi(p)$ is outward pointing for each $p$. Inward pointing is defined analogously.

To clarify the situation, let us consider the special case of an $\mathbb{R}_{u^{k} \leq 0^{-}}^{n}$ valued chart ( $U, \mathrm{x}$ ) for fixed $k \in\{1,2, \ldots, n\}$ (thus $\lambda=-u^{k}$ ). If $p \in U$, then $v_{p}$ is outward pointing exactly when $d x^{k}\left(v_{p}\right)>0$. In particular, $\left.\frac{\partial}{\partial x^{k}}\right|_{p}$ is outward pointing. For an $\mathbb{R}_{u^{k} \geq 0}^{n}$-valued chart ( $U, \mathrm{x}$ ), the reverse is true; $v_{p}$ is outward pointing exactly when $d x^{k}\left(v_{p}\right)<0$. We shall see that $\mathbb{R}_{u^{1} \leq 0}^{n}$ is special, so keep the corresponding criterion $d x^{1}\left(v_{p}\right)>0$ in mind.

\begin{figure}[h]
\begin{center}
  \includegraphics[scale=0.25]{2025_10_09_265d7e1a22c89f79c21eg-396}

\caption{Figure 8.2. Outward pointing}
\end{center}
\end{figure}

Notice that the notion of outward pointing on $\partial M$ does not depend on $M$ being orientable (see Figure 8.3). In fact, we have the following:

Lemma 8.74. Outward pointing sections of $\left.T M\right|_{\partial M}$ always exist.

\begin{figure}[h]
\begin{center}
  \includegraphics[scale=0.25]{2025_10_09_265d7e1a22c89f79c21eg-397}

\caption{Figure 8.3. Outward at boundary of Möbius band}
\end{center}
\end{figure}

Proof. We may assume that we are dealing with an $\mathbb{R}_{u^{1} \leq 0}^{n}$-valued atlas $\left\{\left(U_{\alpha}, \mathrm{x}_{\alpha}\right)\right\}$ for $M$. Let $\left\{\rho_{\alpha}\right\}$ be a partition of unity subordinate to $\left\{U_{\alpha}\right\}$. Define

$$
\chi:=\sum_{\alpha} \rho_{\alpha} \frac{\partial}{\partial x_{\alpha}^{1}} .
$$

Then for $p \in \partial M$ we have

$$
\chi(p)=\left.\sum_{\alpha} \rho_{\alpha}(p) \frac{\partial}{\partial x_{\alpha}^{1}}\right|_{p} .
$$

To see if $\chi(p)$ is outward pointing, let ( $U_{\beta}, \mathbf{x}_{\beta}$ ) be a given $\mathbb{R}_{u^{1} \leq 0}^{n}$-valued chart from the atlas. Then

$$
d y^{1}(\chi(p))=\sum_{\alpha} \rho_{\alpha}(p) d y^{1}\left(\left.\frac{\partial}{\partial x_{\alpha}^{1}}\right|_{p}\right)>0
$$

since $\rho_{\alpha}(p) \geq 0$ for all $\alpha$, we have $\rho_{\alpha}(p)>0$ for at least one $\alpha$, and $d y^{1}\left(\frac{\partial}{\partial x_{\alpha}^{1}}(p)\right)>0$ for all $\alpha$.

In the case that $M$ is oriented we can use the orientation plus the notion of outward to define an orientation on the boundary. Let ( $M,[\omega]$ ) be an oriented smooth manifold with boundary and suppose that $\chi$ is an outward pointing section of $\left.T M\right|_{\partial M}$. Let $\omega \in[\omega]$; then we define $i_{\chi} \omega \in \Omega^{n-1}(\partial M)$ by

$$
i_{\chi} \omega(p)\left(v_{2}, \ldots, v_{n}\right)=\omega(p)\left(\chi(p), v_{2}, \ldots, v_{n}\right)
$$

It is clear that $i_{\chi} \omega$ is nowhere vanishing. If also $\omega_{1} \in[\omega]$, then it is easy to see that $i_{\chi} \omega=f \iota_{\chi} \omega_{1}$ for a positive function $f$. Thus $\left[i_{\chi} \omega\right]$ depends only on $[\omega]$ and provides an orientation on $\partial M$.

Definition 8.75. The orientation on $\partial M$ defined above is called the induced orientation.

It follows that if $\left(f_{1}, f_{2}, \ldots, f_{n}\right)$ is a positively oriented frame field on an open set $U \subset M$ with nonempty intersection with $\partial M$ and if $f_{1}(p)$ is\\
outward pointing for all $p \in U \cap \partial M$, then $\left(f_{2}, \ldots, f_{n}\right)$ restricted to $U \cap \partial M$ is a positively oriented frame with respect to the induced orientation on $\partial M$.

The case $n=1$ needs some interpretation. Here $\partial M$ is a discrete set of points $\partial M=\left\{p_{1}, \ldots, p_{k}\right\}$, and an orientation is an assignment of +1 or -1 to each point. For $p_{i} \in \partial M$, we assign +1 if $\omega\left(p_{i}\right)(\chi(p))>0$ for some outward pointing vector $\chi(p)$.

If every chart in an atlas takes values in a fixed half-space $\mathbb{R}_{\lambda \geq c}^{n}$, then we say the atlas is $\mathbb{R}_{\lambda \geq c}^{n}$-valued. It is not true generally that an oriented $n$-manifold has an atlas of positively oriented charts with values in a fixed half-space. However, it is almost true:

Lemma 8.76. If $M$ is an oriented $n$-manifold with nonempty boundary and $n \geq 2$, then there is an atlas of positively oriented charts.

Proof. If ( $U, \mathrm{x}$ ) is not positively oriented, then replace $\mathrm{x}=\left(x^{1}, x^{2}, \ldots, x^{n}\right)$ by $\mathrm{y}=\left(y^{1}, y^{2}, \ldots, y^{n}\right):=\left(x^{1},-x^{2}, \ldots, x^{n}\right)$ to obtain a positively oriented chart. Notice that this does not work if $n=1$.

Exercise 8.77. Show that although, as a manifold with boundary, the interval $M=[0,1]$ is oriented in the standard way, there is no atlas consisting of positively oriented $\mathbb{R}_{u^{1} \leq 0}^{1}$-valued charts. Trace the problem to the fact that we have defined an atlas for a manifold with boundary as having charts with values in a fixed half-space (in this case $\mathbb{R}_{u^{1} \leq 0}^{1}$ ). Hint: There are no positively oriented charts containing 0 . Changing to $\mathbb{R}_{u^{1} \geq 0}^{1}$ pushes the problem to the other endpoint.

This last exercise exhibits a fact that is an annoyance if one wants to work with an atlas taking values in a fixed half-space such as the ever popular upper half-space $\mathbb{R}_{u^{n} \geq 0}^{n}$. It is an issue often overlooked in the literature.\\
Definition 8.78. A nice chart on a smooth manifold (possibly with boundary) is a chart ( $U, \mathrm{x}$ ) such that $\mathrm{x}(U)=\mathbb{R}_{u^{1} \leq 0}^{n}$ if $U \cap \partial M \neq \emptyset$ or where $\mathrm{x}(U)=\mathbb{R}^{n}$ if $U \cap \partial M=\emptyset$.

Lemma 8.79. The following assertions hold:\\
(i) Every smooth manifold has an atlas consisting of nice charts.\\
(ii) Every oriented smooth manifold without boundary has an atlas consisting of positively oriented nice charts.\\
(iii) Every oriented smooth manifold with boundary of dimension $n \geq 2$ has an atlas consisting of positively oriented nice charts.

Proof. (i) If ( $U, \mathrm{x}$ ) is a chart with range in the interior of the left half-space $\mathbb{R}_{u^{1} \leq 0}^{n}$, then we can find a ball $B$ inside $\mathrm{x}(U)$ in $\mathbb{R}_{u^{1}<0}^{n}$ and then we form a new chart on $\mathrm{x}^{-1}(B)$ with range $B$. But a ball is diffeomorphic to $\mathbb{R}^{n}$ (by\\
an orientation preserving diffeomorphism). If ( $U, \mathrm{x}$ ) is a chart with range meeting the boundary of the left half-space $\mathbb{R}_{u^{1} \leq 0}^{n}$, then we can find a halfball $B_{-}$in $\mathbb{R}_{u^{1} \leq 0}^{n}$ with center on $\mathbb{R}_{u^{1}=0}^{n}$. Reduce the chart domain as before to have range equal to this half-ball. But every half-ball is diffeomorphic to the half-space $\mathbb{R}_{u^{1} \leq 0}^{n}$ so we can proceed by composition as before.\\
(ii) For this, we repeat the procedure of (i) and notice that since $\mathbb{R}^{n}$ is diffeomorphic to $B$ by an orientation preserving diffeomorphism, the new nice chart will be positively oriented if ( $U, \mathrm{x}$ ) is positive.\\
(iii) If ( $U, \mathrm{x}$ ) is a chart with range in the interior of the left half-space $\mathbb{R}_{u^{1} \leq 0}^{n}$, we proceed as in (ii). If ( $U, \mathrm{x}$ ) is a chart with range meeting the boundary of the left half-space $\mathbb{R}_{u^{1} \leq 0}^{n}$ that is already positively oriented, then we proceed as in (i). If ( $U, \mathrm{x}$ ) is not positively oriented, then replace $\mathrm{x}=\left(x^{1}, x^{2}, \ldots, x^{n}\right)$ by $\mathrm{y}=\left(y^{1}, y^{2}, \ldots, y^{n}\right):=\left(x^{1},-x^{2}, \ldots, x^{n}\right)$ as in the proof of Lemma 8.76 and proceed as before.

Exercise 8.80. If $\mathrm{x}=\left(x^{1}, \ldots, x^{n}\right)$ is an oriented $\mathbb{R}_{u^{1} \leq 0^{-}}^{n}$-valued chart on a neighborhood of $p$ on the boundary of an oriented manifold with boundary, then the vectors $\frac{\partial}{\partial x^{2}}, \ldots, \frac{\partial}{\partial x^{n}}$ form a positive basis for $T_{p} \partial M$ with respect to the induced orientation on $\partial M$ and thus $\left(x^{2}, \ldots, x^{n}\right)$ restricts to $\partial M$ to give a positively oriented chart.

\subsection*{Invariant Forms}
Throughout this section $G$ is an $n$-dimensional Lie group.\\
Definition 8.81. An element $\omega \in \Omega^{k}(G)$ is called left invariant if $L_{g}^{*} \omega=\omega$ for all $g \in G$.

It is easy to see that a left invariant $k$-form is determined by its value $\omega(e)$ at the identity element. Now suppose that $\omega^{1}, \ldots, \omega^{n}$ are invariant 1 -forms such that $\omega^{1}(e), \ldots, \omega^{n}(e)$ is a basis of $T_{e}^{*} G$. The $\omega^{1}, \ldots, \omega^{n}$ are independent everywhere (why?). If $X_{1}, \ldots, X_{n}$ is the frame field dual to $\omega^{1}, \ldots, \omega^{n}$, then each $X_{i}$ is a left invariant vector field as may be easily checked. By definition, a left invariant form satisfies

$$
\omega(x)\left(v_{1}, \ldots, v_{k}\right)=\omega(g x)\left(T L_{g} v_{1}, \ldots, T L_{g} v_{k}\right)
$$

for all $x \in M$ and $g \in G$. This would make sense even if $\omega$ were not smooth, but $\omega\left(X_{1}, \ldots, X_{k}\right)$ is not only smooth but constant, so any left invariant form must be smooth after all. Every left invariant $k$-form $\omega$ can be written as

$$
\omega=\sum_{1 \leq i_{1}<\cdots<i_{k} \leq n} a_{i_{1} \ldots i_{k}} \omega^{i_{1}} \wedge \cdots \wedge \omega^{i_{k}}
$$

for unique constants $a_{i_{1} \ldots i_{k}}$.

The exterior derivative of a left invariant form is left invariant since for any $g \in G$, we have

$$
L_{g}^{*} d \omega=d L_{g}^{*} \omega=d \omega
$$

Furthermore for any left invariant vector fields $X, Y$ we have

$$
\begin{aligned}
d \omega(X, Y) & =X \omega(Y)-Y \omega(X)-\omega([X, Y]) \\
& =-\omega([X, Y])
\end{aligned}
$$

and so in particular

$$
\left.d \omega\right|_{e}(v, w)=-\omega_{e}([v, w])
$$

for any $v, w \in \mathfrak{g}$.\\
A form is called right invariant if $R_{g}^{*} \omega=\omega$ for all $g \in G$. Of course, if $\omega$ is right invariant, then so is $d \omega$. If inv : $x \mapsto x^{-1}$ denotes the inversion map on $G$, then

$$
\mathrm{inv} \circ R_{g}=L_{g^{-1}} \circ \mathrm{inv},
$$

and so

$$
R_{g}^{*} \circ \mathrm{inv}^{*}=\mathrm{inv}^{*} \circ L_{g^{-1}}^{*}
$$

As a consequence, if $\omega$ is left invariant, then

$$
R_{g}^{*} \circ \operatorname{inv}^{*} \omega=\operatorname{inv}^{*} \circ L_{g^{-1}}^{*} \omega=\operatorname{inv}^{*} \omega,
$$

so inv* $\omega$ is right invariant. Similarly, if $\omega$ is right invariant, then $\operatorname{inv}^{*} \omega$ is left invariant.

Lemma 8.82. $T_{e} \mathrm{inv}=-\mathrm{id}: \mathfrak{g} \rightarrow \mathfrak{g}$.\\
Proof. Let $v \in \mathfrak{g}$. Then the curve $t \mapsto \exp (t v)$ has tangent $v$. Thus $(\exp (t v))^{-1}$ has tangent $T_{e} \operatorname{inv} \cdot v$. But $(\exp (t v))^{-1}=\exp (-t v)$ so we must have $T_{e} \mathrm{inv} \cdot v=-v$.

Proposition 8.83. Let $\operatorname{inv}: x \mapsto x^{-1}$ be the inversion map on $G$.\\
(i) If $\omega \in \Omega^{k}(G)$, then $\left(\operatorname{inv}^{*} \omega\right)_{e}=(-1)^{k} \omega_{e}$.\\
(ii) If $\omega$ is left and right invariant, then $d \omega=0$.

Proof. (i) It suffices to assume that $\omega=f \omega^{i_{1}} \wedge \cdots \wedge \omega^{i_{k}}$ for certain 1 -forms $\omega^{i_{1}}, \ldots, \omega^{i_{k}}$ that may be taken to be left invariant. But the result will follow if we can show that $\left(\operatorname{inv}^{*} \omega\right)_{e}=-1 \omega_{e}$ for any 1 -form. So let $v \in \mathfrak{g}$. Then, by Lemma 8.82,

$$
\left(\operatorname{inv}^{*} \omega\right)_{e}(v)=\omega_{e}\left(T_{e} \operatorname{inv} \cdot v\right)=-\omega_{e}(v)
$$

(ii) From (i) we have

$$
\left(\operatorname{inv}^{*} \omega\right)_{e}=(-1)^{k} \omega_{e} .
$$

If $\omega$ is left and right invariant, then both $\operatorname{inv}^{*} \omega$ and $\omega$ are left invariant and this continues to hold globally:

$$
\operatorname{inv}^{*} \omega=(-1)^{k} \omega
$$

$d \omega$ is also left and right invariant, so

$$
\operatorname{inv}^{*} d \omega=(-1)^{k+1} d \omega
$$

On the other hand,

$$
\operatorname{inv}^{*} d \omega=d \operatorname{inv}^{*} \omega=(-1)^{k} d \omega
$$

so $(-1)^{k+1} d \omega=(-1)^{k} d \omega$ and then $d \omega=0$.\\
Corollary 8.84. If $G$ is abelian, then $\mathfrak{g}$ is abelian.\\
Proof. If $G$ is abelian, then every left invariant form is also right invariant, so $d \omega=0$ for all left invariant forms $\omega$. But then by equation (8.13), for any $v, w \in \mathfrak{g}$ we have $\omega_{e}([v, w])=0$ for all $\omega_{e} \in \mathfrak{g}^{*}$. Thus $[v, w]=0$ for any $v, w$.

If $X_{1}, \ldots, X_{n}$ are left invariant vector fields with $X_{1}(e), \ldots, X_{n}(e)$ dual to the basis $\omega^{1}(e), \ldots, \omega^{n}(e)$, then

$$
\left[X_{i}(e), X_{j}(e)\right]=\sum_{k} c_{i j}^{k} X_{k}(e) \text { for } 1 \leq i, j \leq n
$$

where $c_{i j}^{k}$ are the structure constants from Definition 5.56. We also have $\left[X_{i}, X_{j}\right]=\sum_{k} c_{i j}^{k} X_{k}$ for $1 \leq i, j \leq n$, so by the equation above and equation (8.13) we have for $i<j$,

$$
d \omega^{k}\left(X_{i}, X_{j}\right)=-\omega^{k}\left(\left[X_{i}, X_{j}\right]\right)=-c_{i j}^{k}
$$

Since $d \omega^{k}\left(X_{i}, X_{j}\right)$ gives the components of $d \omega^{k}$, we have

$$
d \omega^{k}=-\sum_{i<j} c_{i j}^{k} \omega^{i} \wedge \omega^{j}=-\frac{1}{2} \sum_{i, j} c_{i j}^{k} \omega^{i} \wedge \omega^{j}
$$

which are called the equations of structure, or the structural equations.

We now use the concept of Lie algebra-valued forms and the product $[,]^{\wedge}$ introduced earlier. Recall the left Maurer-Cartan form $\omega_{G}$. If $\omega^{1}, \ldots, \omega^{n}$ is a basis of left invariant 1-forms dual to $X_{1}, \ldots, X_{n}$, then (omitting tensor product signs),

$$
\omega_{G}=\sum_{i=1}^{n} X_{i}(e) \omega^{i}
$$

Indeed, if $X_{g}=\sum_{i=1}^{n} a^{i} X_{i}(g) \in T_{g} G$, then $X=\sum_{i=1}^{n} a^{i} X_{i}$ is the unique left invariant vector field such that $X(g)=X_{g}$ and we have

$$
\begin{aligned}
\left(\sum_{i=1} X_{i}(e) \omega^{i}\right)\left(X_{g}\right) & =\sum_{i=1}^{n} a^{i} X_{i}(e) \\
& =X(e)=T L_{g^{-1}} \cdot X(g)=\omega_{G}\left(X_{g}\right)
\end{aligned}
$$

Moreover,

$$
\begin{aligned}
d \omega_{G} & =\sum_{k=1}^{n} X_{k}(e) d \omega^{k}=-\sum_{k=1}^{n} X_{k}(e)\left(\sum_{i<j} c_{i j}^{k} \omega^{i} \wedge \omega^{j}\right) \\
& =-\sum_{k=1}^{n} \sum_{i<j} c_{i j}^{k} X_{k}(e) \omega^{i} \wedge \omega^{j}
\end{aligned}
$$

On the other hand we have

$$
\begin{aligned}
{\left[\omega_{G}, \omega_{G}\right]^{\wedge} } & =\left[\sum_{i=1}^{n} X_{i}(e) \omega^{i}, \sum_{i=1}^{n} X_{j}(e) \omega^{j}\right]^{\wedge} \\
& =\sum_{i, j=1}^{n}\left[X_{i}(e), X_{j}(e)\right] \omega^{i} \wedge \omega^{j} \\
& =\sum_{i, j=1}^{n} \sum_{k} c_{i j}^{k} X_{k} \omega^{i} \wedge \omega^{j}=\frac{1}{2} \sum_{k} \sum_{i<j}^{n} c_{i j}^{k} X_{k} \omega^{i} \wedge \omega^{j},
\end{aligned}
$$

so that an alternative and concise form of the equations of structure is the single equation

$$
d \omega_{G}=-\frac{1}{2}\left[\omega_{G}, \omega_{G}\right]^{\wedge}
$$

Using what we learned at the end of Section 8.4, this equation may be written as

$$
d \omega_{G}(X, Y)=-\left[\omega_{G}(X), \omega_{G}(Y)\right]
$$

for any $X, Y \in \mathfrak{X}(G)$. If $G$ is a matrix group, then the structural equation can be written as

$$
d \omega_{G}=-\omega_{G} \wedge \omega_{G}
$$

Indeed, if we abbreviate $\omega_{G}$ to just $\omega$, then we have $(d \omega)_{j}^{i}=d \omega_{j}^{i}$ and

$$
\begin{aligned}
{\left[\omega_{G}(X), \omega_{G}(Y)\right]_{j}^{i} } & =(\omega(X) \omega(Y)-\omega(Y) \omega(X))_{j}^{i} \\
& =\omega(X)_{k}^{i} \omega(Y)_{j}^{k}-\omega(Y)_{k}^{i} \omega(X)_{j}^{k} \\
& =\left(\omega_{k}^{i} \wedge \omega_{j}^{k}\right)(X, Y)
\end{aligned}
$$

Equation (8.15) is also called the Maurer-Cartan equation.

\section*{Problems}
(1) Let $S_{N}$ denote the group of permutations of the set $\{1,2, \ldots, N\}$. Now let $G_{k, \ell}$ be the subgroup of $S_{k+\ell}$ which consists of permutations that leave the sets $\{1, \ldots, k\}$ and $\{k+1, \ldots, k+\ell\}$ each invariant. A cross section of $G_{k, \ell}$ is a subset $K$ of $S_{k+\ell}$ that contains exactly one element from each coset in $S_{k+\ell} / G_{k, \ell}$. Show that for any such cross section we have

$$
\begin{aligned}
& \omega \wedge \eta\left(\mathrm{v}_{1}, \ldots, \mathrm{v}_{k}, \mathrm{v}_{k+1}, \ldots, \mathrm{v}_{k+\ell}\right) \\
& =\sum_{\sigma \in K} \operatorname{sgn}(\sigma) \omega\left(\mathrm{v}_{\sigma_{1}}, \ldots, \mathrm{v}_{\sigma_{k}}\right) \eta\left(\mathrm{v}_{\sigma_{k+1}}, \ldots, \mathrm{v}_{\sigma_{k+\ell}}\right)
\end{aligned}
$$

Also, show that the set of all ( $k, \ell$ )-shuffle permutations is a cross section of $G_{k, \ell}$.\\
(2) Show that if $v_{1}, \ldots, v_{k} \in \mathrm{V}$, then $v_{1} \wedge \cdots \wedge v_{k} \neq 0$ if and only if $v_{1}, \ldots, v_{k}$ are linearly independent.\\
(3) Let $f_{1}, \ldots, f_{n}$ be smooth functions on an open set in an $n$-manifold. Let $p$ be in their common domain. Then $d f_{1} \wedge \cdots \wedge d f_{n}$ is nonzero at $p$ if and only if $f_{1}, \ldots, f_{n}$ agree with the coordinate functions of a chart whose domain is a neighborhood of $p$.\\
(4) (a) Let $v \in \wedge^{k} \mathrm{V}$. Show that if $v \wedge w=0$ for all $w \in \wedge^{n-k} \mathrm{V}$, where $n=\operatorname{dimV}$, then $v=0$.\\
(b) Using part (a), show that more generally, if $v \wedge w=0$ for all $w \in \Lambda^{m} \mathrm{V}$, where $m \leq n-k$, then $v=0$. [Hints: If $v \wedge w=0$ for all $w \in \Lambda^{m} \mathrm{V}$, then $v \wedge(w \wedge x)=0$ for all $w \in \Lambda^{m} \mathrm{V}$ and all $x \in \bigwedge^{n-k-m} \mathrm{V}$. Elements of the form $w \wedge x$ as above span $\bigwedge^{n-k} \mathrm{V}$.]\\
(5) Prove Proposition 8.23.\\
(6) Show that the sphere is orientable.\\
(7) Prove (i) of Proposition 8.4.\\
(8) Prove Proposition 8.30.\\
(9) Prove equation (i) of Corollary 8.58.\\
(10) Prove Cartan's lemma: Let $k \leq n=\operatorname{dim} M$ and let $\omega_{1}, \ldots, \omega_{k}$ be 1forms on $M$ which are linearly independent at each point. Suppose that there are 1 -forms $\theta_{1}, \ldots, \theta_{k}$ such that

$$
\sum_{i=1}^{k} \theta_{i} \wedge \omega_{i}=0 \text { (identically). }
$$

Then there exists a symmetric $k \times k$ matrix of smooth functions ( $A_{i j}$ ) such that

$$
\theta_{i}=\sum_{j=1}^{k} A_{i j} \omega_{j} \text { for } i=1, \ldots, k
$$

(11) Let $M=\mathbb{R}^{3} \backslash\{0\}$ and let

$$
\omega=\frac{x d y \wedge d z+y d z \wedge d x+z d x \wedge d y}{\left(x^{2}+y^{2}+z^{2}\right)^{3 / 2}}
$$

Find $d \omega$ and determine whether $\omega$ is closed and if so, whether it is exact. Find the expression for $\omega$ in spherical coordinates.\\
(12) Show that every simply connected manifold is orientable.\\
(13) (a) Let V and W be vector spaces with $\operatorname{dim} \mathrm{V}=n$ and $\operatorname{dim} \mathrm{W}= m$. Show that if $A \in L(\mathrm{V}, \mathrm{~W})$, then there is a unique map $\wedge_{A}$ : $\wedge^{k} \mathrm{V} \rightarrow \wedge^{k} \mathrm{~W}$ such that $\wedge_{A}\left(v_{1} \wedge \cdots \wedge v_{k}\right)=A v_{1} \wedge \cdots \wedge A v_{k}$ whenever $v_{1}, \ldots, v_{k} \in \mathrm{V}$.\\
(b) Show that if $e_{1}, \ldots, e_{n}$ is a basis for V and $f_{1}, \ldots, f_{m}$ is a basis for W and if $A e_{i}=\sum a_{i}^{j} f_{j}$, then for $1 \leq i_{1}<\cdots<i_{k} \leq n$ we have

$$
\wedge_{A}\left(e_{i_{1}} \wedge \cdots \wedge e_{i_{k}}\right)=\sum_{1 \leq j_{1}<\cdots<j_{k} \leq m} a_{i_{1} \ldots i_{k}}^{j_{1} \ldots j_{k}} f_{j_{1}} \wedge \cdots \wedge f_{j_{k}}
$$

where

$$
a_{i_{1} \ldots i_{k}}^{j_{1} \ldots j_{k}}
$$

is the $k \times k$ minor determinant of the matrix ( $a_{j}^{i}$ ) given by

$$
a_{i_{1} \ldots i_{k}}^{j_{1} \ldots j_{k}}=\sum_{\sigma \in S_{k}} \operatorname{sgn}(\sigma) a_{i_{1}}^{\sigma\left(j_{1}\right)} \cdots a_{i_{k}}^{\sigma\left(j_{k}\right)}
$$

(14) Finish the proof of Proposition 8.13.\\
(15) Prove Proposition 8.35.

\section*{Chapter 9}
\section*{Integration and Stokes' Theorem}
Let $M$ be a smooth $n$-manifold possibly with boundary $\partial M$ and assume that $M$ is oriented and that $\partial M$ has the induced orientation (Definition 8.75).

Definition 9.1. The support of a differential form $\alpha \in \Omega(M)$ is the closure of the set $\{p \in M: \alpha(p) \neq 0\}$ and is denoted by $\operatorname{supp}(\alpha)$. The set of all $k$-forms that have compact support is denoted $\Omega_{c}^{k}(M)$, and the set of all $k$-forms with compact support contained in $U \subset M$ is denoted by $\Omega_{c}^{k}(U)$.

Let us consider the case of an $n$-form $\alpha$ on an open subset $U$ of $\mathbb{R}^{n}$. Let ( $u^{1}, \ldots, u^{n}$ ) be standard coordinates on $U$. We may write $\alpha=a d u^{1} \wedge \cdots \wedge d u^{n}$ for some function $a$. If $\alpha$ has compact support in $U$, we may define the integral $\int_{U} \alpha$ by

$$
\begin{aligned}
\int_{U} \alpha & =\int_{U} a d u^{1} \wedge \cdots \wedge d u^{n} \\
& :=\int_{U} a(u) d u^{1} \cdots d u^{n}
\end{aligned}
$$

where this latter integral is the Riemann (or Lebesgue) integral of $a(\cdot)$. For this we extend $a(\cdot)$ by zero to all of $\mathbb{R}^{n}$ and integrate over a sufficiently large closed $n$-cube containing the support of $a(\cdot)$ in its interior. We could have written $\left|d u^{1} \cdots d u^{n}\right|$ instead of $d u^{1} \cdots d u^{n}$ to emphasize that the order of the $d u^{i}$ does not matter as it does for $d u^{1} \wedge \cdots \wedge d u^{n}$. If $U$ is an open subset of the half-space $\mathbb{R}_{\lambda \geq c}^{n}$, then we define $\int_{U} \alpha$ by the same formula. If $\phi: V \rightarrow U$ is an orientation preserving diffeomorphism of open sets in $\mathbb{R}^{n}$, then $\operatorname{det} D \phi>0$. Let $u^{1}, \ldots, u^{n}$ denote standard coordinates on $U$, and let\\
$v^{1}, \ldots, v^{n}$ denote standard coordinates on $V$. Then by the classical change of variable formula,

$$
\begin{aligned}
\int_{U} \alpha & =\int_{U} a(u) d u^{1} \cdots d u^{n}=\int_{V} a \circ \phi(v)|\operatorname{det} D \phi| d v^{1} \cdots d v^{n} \\
& =\int_{V} a \circ \phi(v) \operatorname{det} D \phi d v^{1} \cdots d v^{n} \\
& =\int_{V} a \circ \phi(v) \operatorname{det} D \phi d v^{1} \wedge \cdots \wedge d v^{n}=\int_{V} \phi^{*} \alpha
\end{aligned}
$$

So

$$
\int_{U} \alpha=\int_{V} \phi^{*} \alpha
$$

Next consider an oriented $n$-manifold $M$ without boundary and let $\alpha \in \Omega^{n}(M)$. If $\alpha$ has compact support inside $U$ for some positively oriented chart $(U, \mathrm{x})$, then $\mathrm{x}^{-1}: \mathrm{x}(U) \rightarrow U$ and $\left(\mathrm{x}^{-1}\right)^{*} \alpha$ has compact support in $\mathrm{x}(U) \subset \mathbb{R}^{n}$. We define

$$
\int_{U} \alpha:=\int_{\mathrm{x}(U)}\left(\mathrm{x}^{-1}\right)^{*} \alpha
$$

The change of variables formula (9.1) shows that this definition is independent of the positively oriented chart chosen. In fact, if $(\mathrm{y}, V)$ is another chart and $\alpha$ has support inside $U \cap V$, then $\left(\mathrm{x}^{-1}\right)^{*} \alpha$ has support inside $\mathrm{x}(U \cap V)$ and $\left(\mathrm{y}^{-1}\right)^{*} \alpha$ has support in $\mathrm{y}(U \cap V)$. Thus since $\mathrm{x} \circ \mathrm{y}^{-1}$ is orientation preserving, we have

$$
\begin{aligned}
\int_{\mathrm{x}(U)}\left(\mathrm{x}^{-1}\right)^{*} \alpha & =\int_{\mathrm{x}(U \cap V)}\left(\mathrm{x}^{-1}\right)^{*} \alpha=\int_{\mathrm{y}(U \cap V)}\left(\mathrm{x} \circ \mathrm{y}^{-1}\right)^{*}\left(\mathrm{x}^{-1}\right)^{*} \alpha \\
& =\int_{\mathrm{y}(U \cap V)}\left(\mathrm{y}^{-1}\right)^{*} \alpha=\int_{\mathrm{y}(U)}\left(\mathrm{y}^{-1}\right)^{*} \alpha
\end{aligned}
$$

The same definition works fine in case $M$ has nonempty boundary, but there is a small technicality. Namely, suppose we wish to work only with positive charts taking values in a fixed half-space $\mathbb{R}_{\lambda \geq c}^{n}$. Then as long as $n \geq 2$ there is no problem, but if $n=1$, we are faced with the fact that there may be no positively oriented $\mathbb{R}_{\lambda \geq c}^{1}$-valued atlas at all even if $M$ is orientable. Some authors define manifold with boundary completely in terms of a fixed half-space and seem unaware of this little glitch. Since we allow multiple half-spaces, this is not a problem for us, but in any case, we could modify the definition slightly in a way that works in all dimensions and for any chart. Let $M$ be an oriented manifold with boundary. If $\alpha$ has compact support inside $U$ for some chart ( $U, \mathrm{x}$ ), then

$$
\int_{U} \alpha=\operatorname{sgn}(\mathrm{x}) \int_{\mathrm{x}(U)}\left(\mathrm{x}^{-1}\right)^{*} \alpha
$$

where

$$
\operatorname{sgn}(\mathrm{x})= \begin{cases}1 & \text { if }(U, \mathrm{x}) \text { is positively oriented } \\ -1 & \text { if }(U, \mathrm{x}) \text { is not positively oriented. }\end{cases}
$$

This could be taken as a definition. Once again, the standard change of variables formula shows that this definition is independent of the chart chosen.

Remark 9.2. Because we have used the $\operatorname{sgn}(\mathrm{x})$ factor in the definition, we can use arbitrary charts. But the manifold still must be oriented so that $\operatorname{sgn}(\mathrm{x})$ makes sense!

If $\alpha \in \Omega_{c}^{n}(M)$ has compact support but does not have support contained in some chart domain, then we choose a positively oriented atlas $\left\{\left(\mathrm{x}_{i}, U_{i}\right)\right\}$ for $M$ and a smooth partition of unity $\left\{\left(\rho_{i}, U_{i}\right)\right\}$ subordinate to the atlas, and consider the sum

$$
\sum_{i} \int_{U_{i}} \rho_{i} \alpha=\sum_{i} \int_{\mathbf{x}_{i}\left(U_{i}\right)}\left(\mathrm{x}_{i}^{-1}\right)^{*}\left(\rho_{i} \alpha\right) .
$$

Proposition 9.3. In the sum above, only a finite number of terms are nonzero. The sum is independent of the choice of atlas and smooth partition of unity $\left\{\left(\rho_{i}, U_{i}\right)\right\}$.

Proof. First, for any $p \in M$ there is an open set $O$ containing $p$ such that only a finite number of $\rho_{i}$ are nonzero on $O$. But $\alpha$ has compact support, and so a finite number of such open sets cover the support. This means that only a finite number of the $\rho_{i}$ are nonzero on this support. Now let $\left\{\left(\bar{x}_{i}, V_{i}\right)\right\}$ be another positive atlas and $\bar{\rho}_{i}$ a partition of unity subordinate to it. Then we have

$$
\begin{aligned}
\sum_{i} \int_{U_{i}} \rho_{i} \alpha & =\sum_{i} \int_{U_{i}}\left(\rho_{i} \sum_{j} \bar{\rho}_{j} \alpha\right) \\
& =\sum_{i} \sum_{j} \int_{U_{i} \cap V_{j}} \rho_{i} \bar{\rho}_{j} \alpha=\sum_{j} \int_{V_{j}}\left(\bar{\rho}_{j} \sum_{i} \rho_{i} \alpha\right) \\
& =\sum_{j} \int_{V_{j}} \bar{\rho}_{j} \alpha
\end{aligned}
$$

Since the sum (9.3) above is the same independently of the allowed choices, we make the following definition:

Definition 9.4. Let ( $M,[\varpi]$ ) be an oriented smooth manifold with or without boundary. Let $\alpha \in \Omega^{n}(M)$ have compact support. Choose a positively\\
oriented atlas $\left\{\left(\mathrm{x}_{i}, U_{i}\right)\right\}$ and a smooth partition of unity $\left\{\left(\rho_{i}, U_{i}\right)\right\}$ subordinate to $\left\{U_{i}\right\}$. Then we define

$$
\int_{(M,[\varpi])} \alpha:=\sum_{i} \int_{U_{i}} \rho_{i} \alpha .
$$

We usually omit the explicit reference to the orientation $[\varpi]$ and simply write $\int_{M} \alpha$.

Remark 9.5. Of course if we take (9.2) as a definition, then we may take

$$
\sum_{i} \int_{U_{i}} \rho_{i} \alpha=\sum_{i} \operatorname{sgn}\left(\mathrm{x}_{i}\right) \int_{\mathrm{x}_{i}\left(U_{i}\right)}\left(\mathrm{x}_{i}^{-1}\right)^{*}\left(\rho_{i} \alpha\right)
$$

and $\int_{(M,[\varpi])} \alpha=\sum_{i} \int_{U_{i}} \rho_{i} \alpha$ even for an arbitrary unoriented atlas. Once again, we still need $M$ to be oriented. In the online supplement we introduce twisted $n$-forms, and these may be integrated even on nonorientable manifolds!

In case $M$ is zero-dimensional, and therefore a discrete set of points $M=\left\{p_{1}, p_{2}, \ldots\right\}$, an orientation [ $\varpi$ ] is an assignment of +1 or -1 to each point. Then, if $\alpha=f \in \Omega^{0}(M)$, we have

$$
\int_{(M,[\varpi])} f=\sum \pm f\left(p_{i}\right),
$$

where we choose the $\pm$ according to the orientation at the point.

\subsection*{Stokes' Theorem}
In this section we take up the main theorem of the chapter. It is the fundamental theorem of exterior calculus known as Stokes' theorem. Our definition of integration works for any (positive) atlas, but we can use a specific atlas for theoretical purposes. We will employ $\mathbb{R}_{u^{1} \leq 0}^{n}$-valued atlases in this section. This is $\mathbb{R}_{\lambda \geq c}^{n}$ where $c=0$ and $\lambda=-u^{1}$. Let us consider two special cases of integration.

Case 1. This is the case of a compactly supported ( $n-1$ )-form on $\mathbb{R}^{n}$. Let $\omega_{j}=f d u^{1} \wedge \cdots \wedge \widehat{d u^{j}} \wedge \cdots \wedge d u^{n}$ be a smooth ( $n-1$ )-form with compact support in $\mathbb{R}^{n}$, where the caret symbol over the $d u^{j}$ means that this $j$-th\\
factor is omitted. Then we have

$$
\begin{aligned}
\int_{\mathbb{R}_{u^{1} \leq 0}^{n}} d \omega_{j} & =\int_{\mathbb{R}_{u^{1} \leq 0}^{n}} d\left(f d u^{1} \wedge \cdots \wedge \widehat{d u^{j}} \wedge \cdots \wedge d u^{n}\right) \\
& =\int_{\mathbb{R}_{u^{1} \leq 0}^{n}}\left(d f \wedge d u^{1} \wedge \cdots \wedge \widehat{d u^{j}} \wedge \cdots \wedge d u^{n}\right) \\
& =\int_{\mathbb{R}_{u^{1} \leq 0}^{n}}\left(\sum_{k} \frac{\partial f}{\partial u^{k}} d u^{k} \wedge d u^{1} \wedge \cdots \wedge \widehat{d u^{j}} \wedge \cdots \wedge d u^{n}\right) \\
& =\int_{\mathbb{R}_{u^{1} \leq 0}^{n}}(-1)^{j-1} \frac{\partial f}{\partial u^{j}} d u^{1} \wedge \cdots \wedge d u^{n} \\
& =\int_{\mathbb{R}^{n}}(-1)^{j-1} \frac{\partial f}{\partial u^{j}} d u^{1} \cdots d u^{n}=0
\end{aligned}
$$

by the fundamental theorem of calculus and the fact that $f$ has compact support. All compactly supported ( $n-1$ )-forms $\omega$ on $\mathbb{R}^{n}$ are sums of forms of this type and so we have

$$
\int_{\mathbb{R}^{n}} d \omega=\int_{\partial \mathbb{R}^{n}} \omega
$$

Case 2. This is the case of a compactly supported (n-1)-form on $\mathbb{R}_{u^{1} \leq 0}^{n}$. Let $\omega_{j}=f d u^{1} \wedge \cdots \wedge \widehat{d u^{j}} \wedge \cdots \wedge d u^{n}$ be a smooth ( $n-1$ )-form with compact support meeting $\partial \mathbb{R}_{u^{1} \leq 0}^{n}=\mathbb{R}_{u^{1}=0}^{n}=\{0\} \times \mathbb{R}^{n-1}$. Then, if $j \neq 1$,

$$
\begin{aligned}
\int_{\mathbb{R}_{u^{1} \leq 0}^{n}} d \omega_{j} & =\int_{\mathbb{R}_{u^{1} \leq 0}^{n}} d\left(f d u^{1} \wedge \cdots \wedge \widehat{d u^{j}} \wedge \cdots \wedge d u^{n}\right) \\
& =\int_{\mathbb{R}^{n-1}}(-1)^{j-1}\left(\int_{-\infty}^{\infty} \frac{\partial f}{\partial u^{j}} d u^{j}\right) d u^{1} \cdots \widehat{d u^{j}} \cdots d u^{n}=0
\end{aligned}
$$

If $j=1$ we have

$$
\begin{aligned}
\int_{\mathbb{R}_{u^{1} \leq 0}^{n}} d \omega_{1} & =\int_{\mathbb{R}^{n-1}}(-1)^{j-1}\left(\int_{-\infty}^{0} \frac{\partial f}{\partial u^{1}} d u^{1}\right) d u^{2} \wedge \cdots \wedge d u^{n} \\
& =\int_{\mathbb{R}^{n-1}} f\left(0, u^{2}, \ldots, u^{n}\right) d u^{2} \cdots d u^{n} \\
& =\int_{\partial \mathbb{R}_{u^{1} \leq 0}^{n}} f\left(0, u^{2}, \ldots, u^{n}\right) d u^{2} \wedge \cdots \wedge d u^{n}=\int_{\partial \mathbb{R}_{u^{1} \leq 0}^{n}} \omega_{1}
\end{aligned}
$$

For this last equality, it is important that we have set things up so that $d u^{2} \wedge \cdots \wedge d u^{n}$ be positive on $\partial \mathbb{R}_{u^{1} \leq 0}^{n}$ with the induced orientation. Since clearly $\int_{\partial \mathbb{R}_{u^{1} \leq 0}^{n}} \omega_{j}=0$ if $j \neq 1$ or if $\omega_{j}$ has support that does not meet $\partial \mathbb{R}_{u^{1} \leq 0}^{n}$, we see that in any case $\int_{\mathbb{R}_{u^{1} \leq 0}^{n}} d \omega_{j}=\int_{\partial \mathbb{R}_{u^{1} \leq 0}^{n}} \omega_{j}$. All compactly\\
supported ( $n-1$ )-forms $\omega$ on $\mathbb{R}_{u^{1} \leq 0}^{n}$ are sums of forms of this type, and so summing we have

$$
\int_{\mathbb{R}_{u^{1} \leq 0}^{n}} d \omega=\int_{\partial \mathbb{R}_{u^{1} \leq 0}^{n}} \omega
$$

Let us assume that $M$ is an oriented smooth manifold of dimension $n \geq 2$. Then there is a positively oriented atlas $\left\{\left(U_{\alpha}, \mathrm{x}_{\alpha}\right)\right\}_{\alpha \in A}$ consisting of nice charts so either $\mathrm{x}_{\alpha}: U_{\alpha} \cong \mathbb{R}^{n}$ or $\mathrm{x}_{\alpha}: U_{\alpha} \cong \mathbb{R}_{u^{1} \leq 0}^{n}$. Let $\left\{\rho_{\alpha}\right\}$ be a smooth partition of unity subordinate to $\left\{U_{\alpha}\right\}$. Notice that $\left\{\left.\rho_{\alpha}\right|_{U_{\alpha} \cap \partial M}\right\}$ is a partition of unity for the cover $\left\{U_{\alpha} \cap \partial M\right\}$ of $\partial M$. Next we apply the special cases above. For $\omega \in \Omega^{n-1}(M)$ with compact support, we have that

$$
\begin{aligned}
\int_{M} d \omega & =\int_{U_{\alpha}} \sum_{\alpha} d\left(\rho_{\alpha} \omega\right)=\sum_{\alpha} \int_{U_{\alpha}} d\left(\rho_{\alpha} \omega\right) \\
& =\sum_{\alpha} \int_{\mathrm{x}_{\alpha}\left(U_{\alpha}\right)}\left(\mathrm{x}_{\alpha}^{-1}\right)^{*} d\left(\rho_{\alpha} \omega\right)=\sum_{\alpha} \int_{\mathrm{x}_{\alpha}\left(U_{\alpha}\right)} d\left(\left(\mathrm{x}_{\alpha}^{-1}\right)^{*} \rho_{\alpha} \omega\right) \\
& =\sum_{\alpha} \int_{\partial\left\{\mathrm{x}_{\alpha}\left(U_{\alpha}\right)\right\}}\left(\left(\mathrm{x}_{\alpha}^{-1}\right)^{*} \rho_{\alpha} \omega\right)=\sum_{\alpha} \int_{\partial U_{\alpha}} \rho_{\alpha} \omega=\int_{\partial M} \omega,
\end{aligned}
$$

where we have used the fact that $\mathrm{x}_{\alpha}^{-1}: \mathrm{x}_{\alpha}\left(\partial U_{\alpha}\right) \rightarrow \partial U_{\alpha}$ is orientation preserving. We have now proved Stokes' theorem stated below for the $n>1$ case. The $n=1$ case is easily proved directly and amounts to the familiar fundamental theorem of calculus.

Theorem 9.6 (Stokes' theorem). Let $M$ be an oriented smooth manifold with boundary (possibly empty) and give $\partial M$ the induced orientation. Then for any $\omega \in \Omega_{c}^{n-1}(M)$ (i.e. with compact support) we have

$$
\int_{M} d \omega=\int_{\partial M} \omega
$$

Note that $\int_{\partial M} \omega:=0$ if $\partial M=\emptyset$. It can be shown that if ( $U, \mathrm{x}$ ) is a chart for a manifold such that $M \backslash U$ has measure zero, then

$$
\int_{M} \omega:=\int_{\mathrm{x}(U)}\left(\mathrm{x}^{-1}\right)^{*} \omega
$$

The definition of integration using partitions of unity is fine for the theoretical purposes we intend to pursue (such as cohomology), but for actually calculating integrals it is nearly useless. Consider the task of integrating the form $\omega:=z d y \wedge d z+x d x \wedge d y$ over the sphere $S^{2} \subset \mathbb{R}^{3}$. Technically speaking, we are integrating the restriction of $\omega$ to $S^{2}$. Let $\sigma(u, v)=(\cos u \cos v, \sin u \cos v, \sin v)$ for $(u, v) \in(0,2 \pi) \times(-\pi, \pi)$. This gives a parametrization of a portion of the sphere. Thus $\sigma$ plays the role of $\mathrm{x}^{-1}$. The image of our parametrization is the domain of the chart x , and the complement of the domain of the chart has measure zero in $S^{2}$.

Because of this it seems plausible that we would get the correct answer by calculating $\int_{(0,2 \pi) \times(-\pi, \pi)} \sigma^{*} \omega$. But $\sigma^{*} \omega$ does not have compact support in $(0,2 \pi) \times(-\pi, \pi)$, and so this is not justified in terms of the theory we have developed. Let us try it anyway. We pull the form back to the $u, v$ space and then integrate just as one normally would in calculus of several variables:

$$
\begin{aligned}
\int_{S^{2}} \omega & :=\int_{(0,2 \pi] \times(-\pi, \pi)} \sigma^{*} \omega \\
& =\int_{(0,2 \pi] \times(-\pi, \pi)}\left(z(u, v) \frac{\partial(y, z)}{\partial(u, v)}+x(u, v) \frac{\partial(x, y)}{\partial(u, v)}\right) d u d v \\
& =\int_{-\pi}^{\pi} \int_{0}^{2 \pi} \cos ^{2} v \cos u \sin v d u d v=0
\end{aligned}
$$

We are led to a nice integral in this case, but in general, proceeding in this way may lead to improper Riemann integrals, and anyway, it is not immediately clear how to connect this with our theory given in terms of partitions of unity. That the answer we obtained above is indeed the correct answer can be seen by applying Stokes' theorem:

$$
\int_{S^{2}} \omega \stackrel{\text { Stokes }}{=} \int_{B} d \omega=\int_{B} 0=0
$$

Here $B$ is the unit ball and we used the fact that $d \omega=0$.\\
A practical way of calculating integrals of forms is to break up the manifold in a nice way and add up the integrals over the pieces. In Problem 10 we ask the reader to prove the following theorem:

Theorem 9.7. Let $M$ be an oriented $n$-manifold with possibly nonempty boundary. Suppose that there are subsets $D_{i} \subset \mathbb{R}^{n}$ for $i=1, \ldots, K$ and smooth maps $\phi_{i}: D_{i} \rightarrow M$ such that the following assertions hold:\\
(i) Each $D_{i}$ is compact and has a boundary of measure zero.\\
(ii) Each $\phi_{i}$ restricts to an orientation preserving diffeomorphism of the interior of $D_{i}$ to the interior of $\phi_{i}\left(D_{i}\right)$.\\
(iii) $\phi_{i}\left(D_{i}\right) \cap \phi_{j}\left(D_{j}\right)$ is either empty or contains only boundary points of $\phi_{i}\left(D_{i}\right)$ and $\phi_{j}\left(D_{j}\right)$.\\
Then, for any smooth $n$-form $\omega$ with support in $\bigcup \phi_{i}\left(D_{i}\right)$, we have

$$
\int_{M} \omega=\sum_{i=1}^{K} \int_{D_{i}} \phi_{i}^{*} \omega
$$

\subsection*{Differentiating Integral Expressions; Divergence}
Suppose that $S \subset M$ is a $k$-dimensional regular submanifold with boundary $\partial S$ (possibly empty) and $\Phi_{t}$ is the flow of some complete vector field $X \in$\\
$\mathfrak{X}(M)$. In this case, $\Phi_{t}(S)$ is also a regular submanifold with boundary. We then consider $\frac{d}{d t} \int_{S} \Phi_{t}^{*} \eta$ for some $k$-form $\eta$ with compact support. We have

$$
\begin{aligned}
\frac{d}{d t} \int_{S} \Phi_{t}^{*} \eta & =\lim _{h \rightarrow 0} \frac{1}{h}\left[\int_{S} \Phi_{t+h}^{*} \eta-\int_{S} \Phi_{t}^{*} \eta\right] \\
& =\lim _{h \rightarrow 0}\left[\int_{S} \frac{1}{h} \Phi_{t}^{*}\left(\Phi_{h}^{*} \eta-\eta\right)\right] \\
& =\left[\int_{\Phi_{t} S} \lim _{h \rightarrow 0} \frac{1}{h}\left(\Phi_{h}^{*} \eta-\eta\right)\right]=\int_{\Phi_{t} S} \mathcal{L}_{X} \eta
\end{aligned}
$$

But also $\int_{S} \Phi_{t}^{*} \eta=\int_{\Phi_{t} S} \eta$. Thus we obtain the useful formula

$$
\frac{d}{d t} \int_{\Phi_{t} S} \eta=\int_{\Phi_{t} S} \mathcal{L}_{X} \eta
$$

As a special case ( $t=0$ ), we have $\left.\frac{d}{d t}\right|_{t=0} \int_{\Phi_{t} S} \eta=\int_{S} \mathcal{L}_{X} \eta$. We can go further using Cartan's formula $\mathcal{L}_{X}=i_{X} \circ d+d \circ i_{X}$. We get

$$
\begin{aligned}
\frac{d}{d t} \int_{\Phi_{t} S} \eta & =\frac{d}{d t} \int_{S} \Phi_{t}^{*} \eta=\int_{\Phi_{t} S} \mathcal{L}_{X} \eta=\int_{\Phi_{t} S} i_{X} d \eta+\int_{\Phi_{t} S} d i_{X} \eta \\
& =\int_{\Phi_{t} S} i_{X} d \eta+\int_{\partial\left(\Phi_{t} S\right)} i_{X} \eta
\end{aligned}
$$

This formula is particularly interesting in the case when $S=\Omega$ is an open submanifold of $M$ with compact closure and smooth boundary $\partial \Omega$, and where $\eta=\mu$ is a volume form on $M$. Let $\Omega_{t}:=\Phi_{t} \Omega$. We have $\frac{d}{d t} \int_{\Omega_{t}} \mu= \int_{\partial \Omega_{t}} i_{X} \mu$ and then

$$
\left.\frac{d}{d t}\right|_{t=0} \int_{\Omega_{t}} \mu=\int_{\partial \Omega} i_{X} \mu
$$

Definition 9.8. If $\mu$ is a volume form orienting a manifold $M$, then $\mathcal{L}_{X} \mu=$ ( $\operatorname{div} X$ ) $\mu$ for a unique function $\operatorname{div} X$ called the divergence of $X$ with respect to the volume form $\mu$.

We have

$$
\left.\frac{d}{d t}\right|_{t=0} \int_{\Omega_{t}} \mu=\int_{\partial \Omega} i_{X} \mu
$$

and

$$
\left.\frac{d}{d t}\right|_{t=0} \int_{\Omega_{t}} \mu=\int_{\Omega} \mathcal{L}_{X} \mu=\int_{\Omega}(\operatorname{div} X) \mu
$$

From this we conclude that

$$
\int_{\Omega}(\operatorname{div} X) \mu=\int_{\partial \Omega} i_{X} \mu
$$

This formula helps to give geometric meaning to $\operatorname{div} X$. The above formula is a version of Gauss' theorem and can be easily proven not just for domains\\
in a manifold with a volume form, but for a general manifold with boundary with a volume form:

Theorem 9.9. Let $M$ be a manifold with boundary and let $\mu$ be a volume form on $M$. In particular, $M$ is oriented by $[\mu]$. Then we have

$$
\int_{M}(\operatorname{div} X) \mu=\int_{\partial M} i_{X} \mu
$$

for any compactly supported smooth vector field $X$.\\
Proof. The proof uses the trick we used above:

$$
\begin{aligned}
\int_{M}(\operatorname{div} X) \mu & =\int_{M} \mathcal{L}_{X} \mu \\
& =\int_{M} d \iota_{X} \mu+\int_{M} i_{X} d \mu \\
& =\int_{M} d \iota_{X} \mu=\int_{\partial M} i_{X} \mu \quad \text { (Stokes' theorem) }
\end{aligned}
$$

Now let $E_{1}, \ldots, E_{n}$ be a local frame field on $U \subset M$ and $\theta^{1}, \ldots, \theta^{n}$ the dual frame field. Then for some smooth function $\rho$ we have

$$
\mu=\rho \theta^{1} \wedge \cdots \wedge \theta^{n}
$$

and a simple calculation gives

$$
\iota_{X}\left(\rho \theta^{1} \wedge \cdots \wedge \theta^{n}\right)=\sum_{k=1}^{n}(-1)^{j-1} \rho X^{k} \theta^{1} \wedge \cdots \wedge \widehat{\theta^{k}} \wedge \cdots \wedge \theta^{n}
$$

Then we have

$$
\begin{aligned}
\mathcal{L}_{X} \mu & =\mathcal{L}_{X}\left(\rho \theta^{1} \wedge \cdots \wedge \theta^{n}\right)=d \iota_{X}\left(\rho \theta^{1} \wedge \cdots \wedge \theta^{n}\right) \\
& =d \sum_{k=1}^{n}(-1)^{j-1} \rho X^{k} \theta^{1} \wedge \cdots \wedge \widehat{\theta^{k}} \wedge \cdots \wedge \theta^{n} \\
& =\sum_{k=1}^{n}(-1)^{j-1} d\left(\rho X^{k}\right) \wedge \theta^{1} \wedge \cdots \wedge \widehat{\theta^{k}} \wedge \cdots \wedge \theta^{n} \\
& =\sum_{k=1}^{n}(-1)^{j-1} \sum_{i=1}^{n}\left(\rho X^{k}\right)_{i} \theta^{i} \wedge \theta^{1} \wedge \cdots \wedge \widehat{\theta^{k}} \wedge \cdots \wedge \theta^{n} \\
& =\sum_{k=1}^{n}\left(\frac{1}{\rho}\left(\rho X^{k}\right)_{k}\right) \rho \theta^{1} \wedge \cdots \wedge \theta^{n}=\sum_{k=1}^{n}\left(\frac{1}{\rho}\left(\rho X^{k}\right)_{k}\right) \mu
\end{aligned}
$$

where $\left(\rho X^{k}\right)_{k}:=d\left(\rho X^{k}\right)\left(E_{k}\right)$. Thus we end up with a nice formula

$$
\operatorname{div} X=\sum_{k=1}^{n} \frac{1}{\rho}\left(\rho X^{k}\right)_{k}
$$

In particular, if $E_{k}=\frac{\partial}{\partial x^{k}}$ for some chart $(U, \mathrm{x})=\left(x^{1}, \ldots, x^{n}\right)$, then

$$
\operatorname{div} X=\sum_{k=1}^{n} \frac{1}{\rho} \frac{\partial}{\partial x^{k}}\left(\rho X^{k}\right) .
$$

If we were to replace the volume form $\mu$ by $-\mu$, then divergence with respect to that volume form would be given locally by $\sum_{k=1}^{n} \frac{1}{-\rho} \frac{\partial}{\partial x^{k}}\left(-\rho X^{k}\right)= \sum_{k=1}^{n} \frac{1}{\rho} \frac{\partial}{\partial x^{k}}\left(\rho X^{k}\right)$ and so the orientation seems superfluous! What if $M$ is not even orientable? In fact, since divergence is a local concept and orientation is a global concept, it seems that we should replace the volume form in the definition by something else that makes sense even on a nonorientable manifold. On nonorientable manifolds, pseudoforms can be used in place of forms for many purposes. Alternatively, many situations can be handled by going to the two-fold orientation cover. See the online supplement [Lee, Jeff] for a discussion.

\subsection*{Stokes' Theorem for Chains}
Here we give another version of Stokes' theorem that is very useful in connection with topology. We say that distinct points $p_{0}, \ldots, p_{k+1} \in \mathbb{R}^{n}$ are in general position if they are not contained in any $k$-dimensional affine subspace of $\mathbb{R}^{n}$. Equivalently, $p_{0}, \ldots, p_{k}$ are in general position if the vectors $p_{1}-p_{0}, \ldots, p_{k}-p_{0}$ are linearly independent. A set of the form

$$
\left\{\sum_{i=0}^{p} t_{i} p_{i}: 0 \leq t^{i} \leq 1 \text { and } \sum_{i=0}^{p} t_{i}=1\right\}
$$

for $p_{0}, \ldots, p_{k}$ in general position is called a geometric $k$-simplex. The set above is a closed convex set and is the smallest such set containing the points $p_{0}, \ldots, p_{k}$. The integer $k$ is the dimension of the simplex, and a geometric 0 -simplex is just a point. Let us denote the geometric $k$-simplex determined by the points $p_{0}, \ldots, p_{k}$ by $\left\langle p_{0}, \ldots, p_{k}\right\rangle$. If $\left\langle p_{0}, \ldots, p_{k}\right\rangle$ is such a geometric $k$-simplex and $\left\{p_{i_{1}}, \ldots, p_{i_{j}}\right\} \subset\left\{p_{0}, \ldots, p_{k+1}\right\}$, then $p_{i_{1}}, \ldots, p_{i_{j}}$ are in general position and the geometric $j$-simplex $\left\langle p_{i_{1}}, \ldots, p_{i_{j}}\right\rangle$ is called a face of $\left\langle p_{0}, \ldots, p_{k}\right\rangle$. In particular, the geometric ( $k-1$ )-simplices obtained by omitting one of $p_{0}, \ldots, p_{k}$ are called boundary faces. The $i$-th boundary face of $\left\langle p_{0}, \ldots, p_{k}\right\rangle$ is the simplex obtained by omitting the $p_{i}$ and is denoted $\partial_{i}\left\langle p_{0}, \ldots, p_{k}\right\rangle=\left\langle p_{0}, \ldots, \widehat{p_{i}}, \ldots, p_{k}\right\rangle$. Geometric simplices can be combined to create geometric simplicial spaces, but we take a slightly more flexible approach.

Definition 9.10. For $k \geq 0$, the set $\Delta^{k}:=\left\{a \in \mathbb{R}^{k}: a^{i} \geq 0\right.$ and $\left.\sum a^{i} \leq 1\right\}$ is called the standard $k$-simplex in $\mathbb{R}^{k}$. Note that $\Delta^{0}=\mathbb{R}^{0}$ consists of just a single point denoted 0 . In other words, $\Delta^{k}$ is $\left\langle e_{0}, \ldots, e_{k}\right\rangle$, where $e_{0}=(0, \ldots, 0), e_{1}=(1,0, \ldots, 0), e_{2}=(0,1, \ldots, 0), \ldots$, etc.

Definition 9.11. Let $M$ be an $n$-manifold. A $C^{r}$ map $\sigma: \Delta^{k} \rightarrow M$ is called a $C^{r}$ singular $k$-simplex. Let $G$ be an abelian group. A formal sum $c=\sum_{\sigma} c_{\sigma} \sigma$, where the sum is over all singular simplices $\sigma, c_{\sigma} \in G$, and $c_{\sigma}=0$ for all but finitely many simplices $\sigma$, is called a smooth singular $k$-chain with coefficients in $G$. A 0-simplex $\sigma: \Delta^{0} \rightarrow M$ is often identified with the image $\sigma(0)=p \in M$.

By definition, $c_{\sigma}=0$ for all but finitely many simplices; we say that a singular $k$-chain has finite support. The set of all $C^{r} k$-chains with coefficients in $G$ is an abelian group denoted $C_{k}(M, G)^{r}$, where we use additive notation for the group operation. Addition is given by

$$
\left(\sum_{\sigma} c_{\sigma} \sigma\right)+\left(\sum_{\sigma} c_{\sigma}^{\prime} \sigma\right)=\sum_{i=0}^{\infty}\left(c_{\sigma}+c_{\sigma}^{\prime}\right) \sigma .
$$

The cases $r=0$ and $r=\infty$ are the most important. We will restrict to the $r=\infty$ case and drop the superscript in $C_{k}(M, G)^{\infty}$. We now define the notion of the $i$-th face of a singular simplex and then the boundary of a simplex. First, if $q_{0}, \ldots, q_{k}$ are distinct points in $\mathbb{R}^{n}$, then there is a unique affine map $\mathbb{R}^{k} \rightarrow \mathbb{R}^{n}$ which maps $e_{i}$ to $q_{i}$ for $0 \leq i \leq k$. This map restricts to a map from $\Delta^{k}$ onto the convex hull of $q_{0}, \ldots, q_{k}$, which will be a geometric $k$-simplex if $q_{0}, \ldots, q_{k}$ are in general position. For given $q_{0}, \ldots, q_{k}$, denote this map by $\alpha\left(q_{0}, \ldots, q_{k}\right)$. For $0 \leq i \leq k$, let $f_{i}^{k}: \Delta^{k-1} \rightarrow \Delta^{k}$ be $\alpha\left(e_{0}, \ldots, \widehat{e}_{i}, \ldots, e_{k}\right)$. Explicitly, $f_{0}^{1}(0)=1, f_{1}^{1}(0)=0$ and for $k>0$,

$$
\begin{aligned}
& f_{0}^{k}\left(a^{1}, \ldots, a^{k-1}\right)=\left(1-\sum_{i=1}^{k-1} a^{i}, a^{1}, \ldots, a^{k-1}\right) \\
& f_{i}^{k}\left(a^{1}, \ldots, a^{k-1}\right)=\left(a^{1}, \ldots, a^{i-1}, 0, a^{i}, \ldots, a^{k-1}\right)
\end{aligned}
$$

Thus $f_{i}^{k}$ is a homeomorphism of $\Delta^{k-1}$ onto its image, which is the $i$-th boundary face of $\Delta^{k}$. Thus it parametrizes the face of $\Delta^{k}$ opposite $e_{i}$ while keeping track of the order of the vertices.

The $i$-th face of a singular $k$-simplex $\sigma$ is the singular ( $k-1$ )-simplex $\sigma^{i}$ given by

$$
\sigma^{i}:=\sigma \circ f_{i}^{k}
$$

Definition 9.12. The boundary operator $\partial_{k}: C_{k}(M, G) \rightarrow C_{k-1}(M, G)$ is defined on a simplex $\sigma$ by

$$
\partial_{k} \sigma:=\sum_{i=0}^{k}(-1)^{i} \sigma^{i}
$$

\begin{figure}[h]
\begin{center}
  \includegraphics[scale=0.25]{2025_10_09_265d7e1a22c89f79c21eg-417}

\caption{Figure 9.1. A face of $\sigma$}
\end{center}
\end{figure}

and then extended to be a group homomorphism $\partial_{k} \sum_{\sigma} c_{\sigma} \sigma=\sum_{\sigma} c_{\sigma} \partial_{k} \sigma$. If $\sigma$ is a 0 -simplex, then we define $\partial_{0} \sigma=0$ (the group identity) so that for a 0 -chain $\sum_{\sigma} c_{\sigma} \sigma$ we have $\partial_{k} \sum_{\sigma} c_{\sigma} \sigma=0$.

We will denote all the maps $\partial_{k}$ simply by $\partial$ for all $k$. It can be shown that

$$
\partial \circ \partial=0 .
$$

Exercise 9.13. Prove that $\partial \circ \partial=0$. Hint: First show that $f_{i}^{k} \circ f_{j}^{k-1}= f_{j}^{k} \circ f_{i-1}^{k-1}$ for $k \geq 1$ and $i>j$. The maps $f_{i}^{k}$ are defined above.

It is convenient to define $C_{k}(M, G)$ to be the trivial group $\{0\}$ for all $k<0$ and define $\partial_{k}=0$ for all $k<0$. Then $\partial \circ \partial=0$ remains true. The sequence of spaces and maps

$$
\cdots \xrightarrow{\partial} C_{k+1}(M, G) \xrightarrow{\partial} C_{k}(M, G) \xrightarrow{\partial} C_{k-1}(M, G) \xrightarrow{\partial} \cdots
$$

is called the singular chain complex with coefficients in $G$.\\
Let $Z_{k}(M, G)=\operatorname{Ker} \partial_{k}$ and $B_{k}(M, G):=\operatorname{Im} \partial_{k-1}$. Then because of equation (9.5) we have $B_{k}(M, G) \subset Z_{k}(M, G)$. The $k$-th singular homology group $H_{k}(M ; G)$ with coefficients in $G$ is defined by

$$
H_{k}(M ; G)=\frac{Z_{k}(M, G)}{B_{k}(M, G)}
$$

If $c \in Z_{k}(M, G)$, then the equivalence class of $c$ is denoted $[c]$. If $c_{1}$ and $c_{2}$ are $k$-chains in the same class, then $c_{1}=c_{2}+\partial c$ for some $c \in B_{k}(M, G)$, and in this case we say that $c_{1}$ and $c_{2}$ are homologous (or in the same homology class). The most important choices for $G$ are $\mathbb{R}$ and $\mathbb{Z}$.

Exercise 9.14. Check that if $G=\mathbb{R}$, then $C_{k}(M, \mathbb{R}), Z_{k}(M, \mathbb{R}), B_{k}(M, \mathbb{R})$ and $H_{k}(M ; \mathbb{R})$ are all vector spaces in a natural way and the boundary maps $\partial$ extend to linear maps.

We define the integral of a $k$-form $\alpha$ over a smooth singular $k$-simplex $\sigma$ as

$$
\int_{\sigma} \alpha:=\int_{\Delta^{k}} \phi^{*} \alpha
$$

and then for a chain $c=\sum_{\alpha} c_{\sigma} \sigma \in C_{k}(M, \mathbb{R})$ we can define

$$
\int_{c} \alpha=\sum_{\sigma} c_{\sigma} \int_{\sigma} \alpha
$$

If $\sigma$ is a 0 -simplex and if $f \in \Omega^{0}(M)=C^{\infty}(M)$, then $\int_{\sigma} f=f(0)$. We state without proof the following version of Stokes' theorem:

Theorem 9.15 (Stokes' theorem for chains). Let $M$ be a smooth manifold. For $c \in C_{k+1}(M, \mathbb{R})$, and $\alpha$ a $k$-form on $M$, we have

$$
\int_{c} d \alpha=\int_{\partial c} \alpha
$$

For a proof see $[\mathbf{W a r}]$. Notice that in this version, $M$ is not assumed to be orientable. Also, $\alpha$ need not have compact support since the chain $c$ has finite support. If $\sigma$ is a 1 -simplex and $f \in C^{\infty}(M)$, then the above reduces to

$$
\int_{\sigma} d f=f(\sigma(1))-f(\sigma(0))
$$

Now we define the de Rham map. First, if $\alpha \in Z^{k}(M) \subset \Omega^{k}(M)$, then we can define an element $I \alpha$ of $H^{k}(M ; \mathbb{R})$, the dual space of $H_{k}(M ; \mathbb{R})$ : For $[c] \in H_{k}(M ; \mathbb{R})$ represented by $c \in Z_{k}(M ; \mathbb{R})$ we define

$$
I \alpha([c])=\int_{c} \alpha
$$

This is well-defined since if $c+\partial c^{\prime} \in[c]$, then

$$
\begin{aligned}
I \alpha\left(c+\partial c^{\prime}\right) & =\int_{c+\partial c^{\prime}} \alpha \\
& =\int_{c} \alpha+\int_{\partial c^{\prime}} \alpha=\int_{c} \alpha+\int_{c^{\prime}} d \alpha \quad \text { (Stokes) } \\
& =\int_{c} \alpha \quad(\text { since } d \alpha=0)
\end{aligned}
$$

This gives a linear map $Z^{k}(M) \rightarrow H^{k}(M ; \mathbb{R})$. Now if $\alpha \in B^{k}(M)$ (image of $d)$, then $\alpha=d \beta$ and

$$
I \alpha([c])=\int_{c} \alpha=\int_{c} d \beta=\int_{\partial c} \beta=0(\text { since } \partial c=0)
$$

Thus we obtain a linear map called the de Rham map

$$
H_{\mathrm{deR}}^{k}(M) \rightarrow H^{k}(M ; \mathbb{R})
$$

where $H_{\mathrm{deR}}^{k}(M)=Z^{k}(M) / B^{k}(M)$ is the de Rham cohomology defined earlier. The content of the celebrated de Rham theorem is in part that this map is an isomorphism. The theorem is fairly difficult to prove.

Theorem 9.16 (de Rham). The de Rham map defined above is an isomorphism,

$$
H_{\mathrm{deR}}^{k}(M) \cong H^{k}(M ; \mathbb{R})
$$

Proof. See [Bo-Tu] or [War].\\
We have defined $H^{k}(M ; \mathbb{R})$ using smooth singular chains, but we could have used continuous chains. The result is isomorphic to $H^{k}(M ; \mathbb{R})$ as we have defined it (see [War] or [Bo-Tu]).

\subsection*{Differential Forms and Metrics}
Let ( $\mathrm{V}, g$ ) be a real scalar product space (not necessarily positive definite). We wish to induce a scalar product on $L_{\text {alt }}^{k}(\mathrm{V}) \cong \Lambda^{k}\left(\mathrm{V}^{*}\right)$. Even though elements of $L_{\text {alt }}^{k}(\mathrm{V}) \cong \Lambda^{k}\left(\mathrm{V}^{*}\right)$ can be thought of as tensors of type $(0, k)$ that just happen to be alternating, we will give a scalar product to this space in such a way that the basis

$$
\left\{e^{i_{1}} \wedge \cdots \wedge e^{i_{k}}\right\}_{i_{1}<\cdots<i_{k}}=\left\{e^{\vec{I}}\right\}
$$

is orthonormal if $e^{1}, \ldots, e^{n}$ is orthonormal. Recall that if $\alpha, \beta \in \mathrm{V}^{*}$, then by definition $\langle\alpha, \beta\rangle=g\left(\alpha^{\sharp}, \beta^{\sharp}\right)$. Now suppose that $\alpha=\alpha^{1} \wedge \alpha^{2} \wedge \cdots \wedge \alpha^{k}$ and $\beta=\beta^{1} \wedge \beta^{2} \wedge \cdots \wedge \beta^{k}$, where the $\alpha^{i}$ and $\beta^{i}$ are 1-forms. Then we want to have

$$
\begin{aligned}
\langle\alpha \mid \beta\rangle & =\left\langle\alpha^{1} \wedge \alpha^{2} \wedge \cdots \wedge \alpha^{k} \mid \beta^{1} \wedge \beta^{2} \wedge \cdots \wedge \beta^{k}\right\rangle \\
& =\operatorname{det}\left[\left\langle\alpha^{i}, \beta^{j}\right\rangle\right]
\end{aligned}
$$

where $\left[\left\langle\alpha^{i}, \beta^{j}\right\rangle\right]$ is an $n \times n$ matrix. Notice that we use $\langle\alpha \mid \beta\rangle$ rather than $\langle\alpha, \beta\rangle$ since the latter could be taken to be the natural inner product of $\alpha$ and $\beta$ as tensors - the latter differs from the first by a factor. We want to extend this bilinearly to all $k$-forms. We could just declare the basis (9.6) above to be orthonormal and thus define a scalar product. Of course, one must then show that this scalar product does not depend on the choice of basis. For completeness, we now show how to arrive at the appropriate scalar product using universal mapping properties and obtain some formulas. Fix $\beta^{1}, \beta^{2}, \ldots, \beta^{k} \in \mathrm{V}^{*}$ and consider the map

$$
\mu_{\beta^{1}, \beta^{2}, \ldots, \beta^{k}}: \mathrm{V}^{*} \times \cdots \times \mathrm{V}^{*} \rightarrow \mathbb{R}
$$

given by

$$
\mu_{\beta^{1}, \beta^{2}, \ldots, \beta^{k}}:\left(\alpha^{1}, \alpha^{2}, \ldots, \alpha^{k}\right) \mapsto \operatorname{det}\left[\left\langle\alpha^{i}, \beta^{j}\right\rangle\right] .
$$

Since this is an alternating multilinear map, we can use the universal property of $\wedge^{k}\left(\mathrm{V}^{*}\right)$ to see that this map defines a unique linear map

$$
\widetilde{\mu}_{\beta^{1}, \beta^{2}, \ldots, \beta^{k}}: \bigwedge^{k}\left(\mathrm{V}^{*}\right) \rightarrow \mathbb{R}
$$

such that

$$
\widetilde{\mu}_{\beta^{1}, \beta^{2}, \ldots, \beta^{k}}\left(\alpha^{1} \wedge \alpha^{2} \wedge \cdots \wedge \alpha^{k}\right)=\operatorname{det}\left[\left\langle\alpha^{i}, \beta^{j}\right\rangle\right] .
$$

Similarly, for fixed $\alpha \in \bigwedge^{k}\left(\mathrm{V}^{*}\right)$, the map $\widetilde{m}_{\alpha}:\left(\beta^{1}, \beta^{2}, \ldots, \beta^{k}\right) \mapsto \widetilde{\mu}_{\beta^{1}, \beta^{2}, \ldots, \beta^{k}}(\alpha)$ is alternating multilinear and so gives rise to a linear map

$$
\widetilde{m}_{\alpha}: \bigwedge^{k}\left(\mathrm{V}^{*}\right) \rightarrow \mathbb{R}
$$

such that

$$
\widetilde{m}_{\alpha}\left(\beta^{1} \wedge \beta^{2} \wedge \cdots \wedge \beta^{k}\right)=\widetilde{\mu}_{\beta^{1}, \beta^{2}, \ldots, \beta^{k}}(\alpha)
$$

Lemma 9.17. The map

$$
\langle\cdot \mid \cdot\rangle: \bigwedge^{k}\left(\mathrm{V}^{*}\right) \times \bigwedge^{k}\left(\mathrm{V}^{*}\right) \rightarrow \mathbb{R}
$$

defined by

$$
\langle\alpha \mid \beta\rangle:=\widetilde{m}_{\alpha}(\beta)
$$

is symmetric and bilinear. We have

$$
\begin{aligned}
& \left\langle\alpha^{1} \wedge \alpha^{2} \wedge \cdots \wedge \alpha^{k} \mid \beta^{1} \wedge \beta^{2} \wedge \cdots \wedge \beta^{k}\right\rangle \\
& =\operatorname{det}\left[\left\langle\alpha^{i}, \beta^{j}\right\rangle\right]
\end{aligned}
$$

for 1 -forms $\alpha^{1}, \alpha^{2}, \ldots, \alpha^{k}, \beta^{1}, \beta^{2}, \ldots, \beta^{k}$.\\
Proof. By construction, the map is linear in the second slot for each fixed $\alpha$. Fix $\beta$ and write $\beta$ as a sum $\beta=\sum b_{i_{1} \ldots i_{k}} \beta^{i_{1}} \wedge \cdots \wedge \beta^{i_{k}}$ in any way. Then

$$
\begin{aligned}
\widetilde{m}_{\alpha}(\beta) & =\sum b_{i_{1} \ldots i_{k}} \widetilde{m}_{\alpha}\left(\beta^{i_{1}} \wedge \cdots \wedge \beta^{i_{k}}\right) \\
& =\sum b_{i_{1} \ldots i_{k}} \widetilde{\mu}_{\beta^{1}, \beta^{2}, \ldots, \beta^{k}}(\alpha)
\end{aligned}
$$

But each map $\alpha \mapsto \widetilde{\mu}_{\beta^{1}, \beta^{2}, \ldots, \beta^{k}}(\alpha)$ is linear and so $\alpha \mapsto \widetilde{m}_{\alpha}(\beta)$ is also linear.\\
Now let $\varphi(\cdot, \cdot): \mathrm{W} \times \mathrm{W} \rightarrow \mathbb{R}$ be any bilinear map on a real vector space W . If $S \subset \mathrm{~W}$ spans W and if $\varphi\left(s_{1}, s_{2}\right)=\varphi\left(s_{2}, s_{1}\right)$ for all $s_{1}, s_{2} \in S$, then $\varphi$ is symmetric. The set of all elements $\beta \in \Lambda^{k}\left(\mathrm{V}^{*}\right)$ of the form $\beta=\beta^{1} \wedge \cdots \wedge \beta^{k}$ for 1 -forms $\beta^{1}, \ldots, \beta^{k}$, is a spanning set. Since

$$
\begin{aligned}
& \left\langle\alpha^{1} \wedge \alpha^{2} \wedge \cdots \wedge \alpha^{k} \mid \beta^{1} \wedge \beta^{2} \wedge \cdots \wedge \beta^{k}\right\rangle \\
& =\operatorname{det}\left[\left\langle\alpha^{i}, \beta^{j}\right\rangle\right]=\operatorname{det}\left[\left\langle\beta^{j}, \alpha^{i}\right\rangle\right] \\
& =\left\langle\beta^{1} \wedge \beta^{2} \wedge \cdots \wedge \beta^{k} \mid \alpha^{1} \wedge \alpha^{2} \wedge \cdots \wedge \alpha^{k}\right\rangle
\end{aligned}
$$

we conclude that $\langle\cdot \mid \cdot\rangle$ is symmetric.\\
The bilinear map defined in the previous lemma is a scalar product for the vector space $\Lambda^{k}\left(\mathrm{V}^{*}\right)$. Notice the vertical bar rather than a comma in\\
the notation. If $e^{1}, \ldots, e^{n}$ is an orthonormal basis for $\mathrm{V}^{*}$, then $\left\{e^{i_{1}} \wedge \cdots\right. \left.\wedge e^{i_{k}}\right\}_{i_{1}<\cdots<i_{k}}$ is orthonormal since

$$
\left\langle e^{i_{1}} \wedge \cdots \wedge e^{i_{k}} \mid e^{j_{1}} \wedge \cdots \wedge e^{j_{k}}\right\rangle=\left|\begin{array}{ccc}
\left\langle e^{i_{1}}, e^{j_{1}}\right\rangle & \left\langle e^{i_{1}}, e^{j_{2}}\right\rangle & \cdots \\
\left\langle e^{i_{2}}, e^{j_{1}}\right\rangle & \left\langle e^{i_{2}}, e^{j_{2}}\right\rangle & \cdots \\
\vdots & \vdots & \ddots
\end{array}\right|,
$$

which is zero if $i_{1}, \ldots, i_{k}$ is not a permutation of $j_{1}, \ldots, j_{k}$, while on the other hand, if $\left(j_{1}, \ldots, j_{k}\right)=\left(\sigma\left(i_{1}\right), \ldots, \sigma\left(i_{k}\right)\right)$, then every column of the above determinant has exactly one nonzero entry, which is either 1 or -1 . Abbreviating $e^{i_{1}} \wedge \cdots \wedge e^{i_{k}}$ to $e^{\vec{I}}$ and so on, we have in the orthonormal case

$$
\left\langle e^{\vec{I}} \mid e^{\vec{I}}\right\rangle= \pm 1
$$

So, if $\alpha=\sum \alpha_{\vec{I}} e^{\vec{I}}$, then $\left\langle e^{\vec{I}} \mid \alpha\right\rangle=\left\langle e^{\vec{I}} \mid e^{\vec{I}}\right\rangle \alpha_{\vec{I}}$ and

$$
\alpha_{\vec{I}}=\left\langle e^{\vec{I}} \mid e^{\vec{I}}\right\rangle\left\langle e^{\vec{I}} \mid \alpha\right\rangle= \pm\left\langle e^{\vec{I}} \mid \alpha\right\rangle .
$$

We see that the scalar product need not be positive definite. Unless we give a specific linear order to the basis $\left\{e^{\vec{I}}\right\}$, the signature is perhaps best thought of as the indexed set $\{\epsilon(\vec{I})\}_{\vec{I}}$ so that

$$
\alpha=\sum_{\vec{I}} \epsilon(\vec{I})\left\langle e^{\vec{I}} \mid \alpha\right\rangle e^{\vec{I}}
$$

We are very much interested in $\bigwedge^{k}\left(\mathrm{V}^{*}\right)$, but note that $\bigwedge^{k}(\mathrm{V})$ is also a scalar product space in the analogous way so that

$$
\left\langle v_{1} \wedge v_{2} \wedge \cdots \wedge v_{k} \mid w_{1} \wedge w_{2} \wedge \cdots \wedge w_{k}\right\rangle=\operatorname{det}\left[\left\langle v_{i}, w_{j}\right\rangle\right]
$$

Exercise 9.18. Let $\mathbb{R}_{1}^{4}$ denote Minkowski space. Determine the index of the scalar product on $\bigwedge^{2}\left(\mathbb{R}_{1}^{4}\right)$ described above.

We already have a scalar product denoted $\langle\cdot, \cdot\rangle$ for tensors. We wish to compare it with the scalar product $\langle\cdot \mid \cdot\rangle$ just defined on forms. For simplicity we consider the positive definite case. If $\alpha=\sum \alpha_{\vec{I}} e^{\vec{I}}$ and $\beta=\sum \beta_{\vec{I}} e^{\vec{I}}$ are $k$-forms considered as covariant tensor fields, then by Problem 18 of Chapter 7 we have

$$
\langle\alpha, \beta\rangle=\sum \alpha_{i_{1} \ldots i_{k}} \beta^{i_{1} \ldots i_{k}}=\sum \alpha_{I} \beta^{I} .
$$

Now let $e^{1}, \ldots, e^{n}$ be an orthonormal basis for $\mathrm{V}^{*}$. Then

$$
\langle\alpha, \beta\rangle=\sum \alpha_{I} \beta^{I}=\sum \alpha_{I} \beta_{I},
$$

and we have

$$
\begin{aligned}
\langle\alpha \mid \beta\rangle & =\sum_{\vec{I}, \vec{J}}\left\langle\alpha_{\vec{I}} e^{\vec{I}} \mid \beta_{\vec{J}} e^{\vec{J}}\right\rangle=\sum_{\vec{I}, \vec{J}} \alpha_{\vec{I}} \beta_{\vec{J}}\left\langle e^{\vec{I}} \mid e^{\vec{J}}\right\rangle \\
& =\sum_{\vec{I}} \alpha_{\vec{I}} \beta_{\vec{I}}=\frac{1}{k!} \sum_{I} \alpha_{I} \beta_{I}=\frac{1}{k!}\langle\alpha, \beta\rangle .
\end{aligned}
$$

We conclude that

$$
\langle\alpha \mid \beta\rangle=\frac{1}{k!}\langle\alpha, \beta\rangle,
$$

so the two scalar products differ by a factor of $k!$. Now let $\theta^{1}, \ldots, \theta^{n}$ be any basis for $\mathrm{V}^{*}$ (not necessarily orthonormal). For a given $\alpha \in \Lambda^{k}\left(\mathrm{V}^{*}\right)$, we can also write $\alpha=\frac{1}{k!} \sum \alpha_{I} e^{I}$ and then as a tensor

$$
\alpha=\frac{1}{k!} \sum \alpha_{i_{1} \ldots i_{k}} \theta^{i_{1}} \wedge \cdots \wedge \theta^{i_{k}}=\sum \alpha_{i_{1} \ldots i_{k}} \theta^{i_{1}} \otimes \cdots \otimes \theta^{i_{k}}
$$

But when $\alpha=\sum a_{\vec{I}} \theta^{\vec{I}}$ and $\beta=\sum b_{\vec{I}} \theta^{\vec{I}}$ are viewed as $k$-forms, we must have

$$
\langle\alpha \mid \beta\rangle=\frac{1}{k!}\langle\alpha, \beta\rangle=\frac{1}{k!} \sum a_{I} b^{I}=\sum a_{\vec{I}} b^{\vec{I}}
$$

Definition 9.19. We defined the scalar product on $\bigwedge^{k} \mathrm{V}^{*} \cong L_{\text {alt }}^{k}(\mathrm{V})$ by first using the above formula for exterior products of 1 -forms and then extending (bi) linearly to all of $\wedge^{k} \mathrm{V}^{*}$. We can also extend to the whole Grassmann algebra $\Lambda \mathrm{V}^{*}=\bigoplus \bigwedge^{k} \mathrm{V}^{*}$ by declaring forms of different degree to be orthogonal. We also have the obvious similar definition for $\bigwedge^{k} \mathrm{V}$ and $\Lambda \mathrm{V}$.

If we have an orthonormal basis $e^{1}, \ldots, e^{n}$ for $\mathrm{V}^{*}$, then $e^{1} \wedge \cdots \wedge e^{n} \in \bigwedge^{n} \mathrm{V}^{*}$. But $\bigwedge^{n} \mathrm{V}^{*}$ is one-dimensional, and if $\ell: \mathrm{V} \rightarrow \mathrm{V}$ is any isometry of $\left(\mathrm{V}^{*}, g^{*}\right)$, then the dual map $\ell^{*}: \mathrm{V}^{*} \rightarrow \mathrm{V}^{*}$ is an isometry of $(\mathrm{V}, g)$. Thus $\ell^{*} e^{1} \wedge \cdots \wedge \ell^{*} e^{n}= \pm e^{1} \wedge \cdots \wedge e^{n}$. In particular, for any permutation $\sigma$ of the letters $\{1,2, \ldots, n\}$ we have $e^{1} \wedge \cdots \wedge e^{n}=\operatorname{sgn}(\sigma) e^{\sigma(1)} \wedge \cdots \wedge e^{\sigma(n)}$.

For a given ordered basis $\left(e_{1}, \ldots, e_{n}\right)$ for V , with dual basis $\left(e^{1}, \ldots, e^{n}\right)$, we have

$$
\left\langle e^{1} \wedge \cdots \wedge e^{n} \mid e^{1} \wedge \cdots \wedge e^{n}\right\rangle=\epsilon_{1} \epsilon_{2} \cdots \epsilon_{n}= \pm 1
$$

Exercise 9.20. Let $\left(e_{1}, \ldots, e_{n}\right)$ be orthonormal. Show that the only elements $\omega$ of $\wedge^{n} \mathrm{V}^{*}$ with $|\langle\omega, \omega\rangle|=1$ are $e^{1} \wedge \cdots \wedge e^{n}$ and $-e^{1} \wedge \cdots \wedge e^{n}$.

Given a fixed orthonormal basis $\left(e_{1}, \ldots, e_{n}\right)$, all ordered orthonormal bases for V fall into two classes. Namely, those bases $\left(f_{1}, \ldots, f_{n}\right)$ for which $f^{1} \wedge \cdots \wedge f^{n}=e^{1} \wedge \cdots \wedge e^{n}$ and those for which $f^{1} \wedge \cdots \wedge f^{n}=-e^{1} \wedge \cdots \wedge e^{n}$. Thus for each orientation of V there is a corresponding element of $\Lambda^{n} \mathrm{V}^{*}$ called a metric volume element for ( $\mathrm{V}, g$ ). The metric volume element corresponding to a basis ( $e_{1}, \ldots, e_{n}$ ) is just $e^{1} \wedge \cdots \wedge e^{n}$. On the other hand,\\
we have seen that any nonzero top form $\omega$ determines an orientation. If we have an orientation given by a top form $\omega$, then obviously, $\left(e_{1}, \ldots, e_{n}\right)$ determines the same orientation if and only if $\omega\left(e_{1}, \ldots, e_{n}\right)>0$ since this means that $\omega=c e^{1} \wedge \cdots \wedge e^{n}$ for some $c>0$.

Proposition 9.21. Let an orientation be chosen on the scalar product space ( $\mathrm{V}, g$ ) and let $\mathcal{E}=\left(e_{1}, \ldots, e_{n}\right)$ be an oriented orthonormal frame so that vol $:=e^{1} \wedge \cdots \wedge e^{n}$ is the corresponding volume element. Then if $\mathcal{F}=\left(f_{1}, \ldots, f_{n}\right)$ is a positively oriented basis for V with dual basis $\mathcal{F}^{*}= \left(f^{1}, \ldots, f^{n}\right)$, then

$$
\operatorname{vol}=\sqrt{\left|\operatorname{det}\left(g_{i j}\right)\right|} f^{1} \wedge \cdots \wedge f^{n}
$$

where $g_{i j}=\left\langle f_{i}, f_{j}\right\rangle$. If $g$ is positive definite, then $\operatorname{det}\left(g_{i j}\right)>0$ and so

$$
\mathrm{vol}=\sqrt{\operatorname{det}\left(g_{i j}\right)} f^{1} \wedge \cdots \wedge f^{n}
$$

Proof. Let $e^{i}=\sum a_{j}^{i} f^{j}$. Then we have

$$
\begin{aligned}
\epsilon_{i} \delta^{i j} & = \pm \delta^{i j}=\left\langle e^{i}, e^{j}\right\rangle=\left\langle\sum a_{k}^{i} f^{k}, \sum a_{m}^{j} f^{m}\right\rangle \\
& =\sum_{k, m} a_{k}^{i} a_{m}^{j}\left\langle f^{k}, f^{m}\right\rangle=\sum_{k, m} a_{k}^{i} a_{m}^{j} g^{k m}
\end{aligned}
$$

so that $\pm 1=\operatorname{det}\left(\left[a_{k}^{i}\right]\right)^{2} \operatorname{det}\left(\left[g^{k m}\right]\right)=\left(\operatorname{det}\left(\left[a_{k}^{i}\right]\right)\right)^{2}\left(\operatorname{det}\left(g_{i j}\right)\right)^{-1}$. Thus,

$$
\pm \sqrt{\left|\operatorname{det}\left(g_{i j}\right)\right|}=\operatorname{det}\left(\left[a_{k}^{i}\right]\right)
$$

But since $\mathcal{E}$ and $\mathcal{F}$ are both positively oriented, we must have $\operatorname{det}\left(\left[a_{k}^{i}\right]\right)>0$ and so

$$
\sqrt{\left|\operatorname{det}\left(g_{i j}\right)\right|}=\operatorname{det}\left(\left[a_{k}^{i}\right]\right)
$$

On the other hand,

$$
\begin{aligned}
\operatorname{vol} & =e^{1} \wedge \cdots \wedge e^{n} \\
& =\left(\sum a_{k_{1}}^{1} f^{k_{1}}\right) \wedge \cdots \wedge\left(\sum a_{k_{1}}^{n} f^{k_{1}}\right)=\operatorname{det}\left(\left[a_{k}^{i}\right]\right) f^{1} \wedge \cdots \wedge f^{n}
\end{aligned}
$$

and the result follows. If $g$ is positive definite, then $\operatorname{det}\left(\left[a_{k}^{i}\right]\right)^{2} \operatorname{det}\left(\left[g^{k m}\right]\right)=1$ so $\operatorname{det}\left(g_{i j}\right)>0$.

Fix an orientation and let $\left(e_{1}, \ldots, e_{n}\right)$ be an orthonormal basis in that orientation class. Then we have the corresponding volume element vol $= e^{1} \wedge \cdots \wedge e^{n}$. Now we define the Hodge star operator $*: \bigwedge^{k} \mathrm{V}^{*} \rightarrow \bigwedge^{n-k} \mathrm{V}^{*}$ for $1 \leq k \leq n$, where $n=\operatorname{dim}(\mathrm{V})$.

Theorem 9.22. Let $(\mathrm{V}, g)$ be a scalar product space with $\operatorname{dim}(\mathrm{V})=n$ and the corresponding volume element vol. For each $k$, there is a unique linear\\
isomorphism $*: \bigwedge^{k} \mathrm{V}^{*} \rightarrow \bigwedge^{n-k} \mathrm{V}^{*}$ such that $\alpha \wedge * \beta=\langle\alpha \mid \beta\rangle$ vol for all $\alpha, \beta \in \bigwedge^{k} \mathrm{V}^{*}$.

Proof. Given $\gamma \in \bigwedge^{n-k} \mathrm{V}^{*}$, define a linear map $L_{\gamma}: \bigwedge^{k} \mathrm{V}^{*} \rightarrow \mathbb{R}$ by requiring that $L_{\gamma}(\alpha) \mathrm{vol}=\alpha \wedge \gamma$. By Problem 3, if $L_{\gamma}(\alpha)=0$ for all $\alpha$, then $\gamma=0$ and $\gamma \mapsto L_{\gamma}$ gives a one-to-one linear map $\Lambda^{k} \mathrm{V}^{*} \rightarrow\left(\bigwedge^{k} \mathrm{V}^{*}\right)^{*}$. But since $\operatorname{dim} \bigwedge^{k} \mathrm{V}^{*}=\operatorname{dim}\left(\bigwedge^{k} \mathrm{V}^{*}\right)^{*}$, this must be a linear isomorphism. Thus, for each $\beta \in \Lambda^{k} \mathrm{V}^{*}$, there is a unique element $* \beta$ such that

$$
L_{* \beta}(\alpha)=\langle\alpha \mid \beta\rangle
$$

for all $\alpha \in \Lambda^{k} \mathrm{V}^{*}$. This defines a map $*: \bigwedge^{k} \mathrm{V}^{*} \rightarrow \bigwedge^{n-k} \mathrm{V}^{*}$ such that

$$
\alpha \wedge * \beta=L_{* \beta}(\alpha) \mathrm{vol}=\langle\alpha \mid \beta\rangle \mathrm{vol}
$$

This map is easily seen to be linear. The equation above also shows that it is one-to-one and hence an isomorphism.

Proposition 9.23. Let ( $\mathrm{V}, g$ ) be a scalar product space with $\operatorname{dim}(\mathrm{V})= n$ and corresponding volume element vol. Let ( $e_{1}, \ldots, e_{n}$ ) be a positively oriented orthonormal basis with dual $\left(e^{1}, \ldots, e^{n}\right)$. Let $\sigma$ be a permutation of $(1,2, \ldots, n)$. On the basis elements $e^{\sigma(1)} \wedge \cdots \wedge e^{\sigma(k)}$ for $\wedge^{k} \mathrm{V}^{*}$ we have

$$
*\left(e^{\sigma(1)} \wedge \cdots \wedge e^{\sigma(k)}\right)=\epsilon_{\sigma_{1}} \epsilon_{\sigma_{2}} \cdots \epsilon_{\sigma_{k}} \operatorname{sgn}(\sigma) e^{\sigma(k+1)} \wedge \cdots \wedge e^{\sigma(n)}
$$

In other words, if we let $\left\{i_{k+1}, \ldots, i_{n}\right\}=\{1,2, \ldots, n\} \backslash\left\{i_{1}, \ldots, i_{k}\right\}$, then

$$
*\left(e^{i_{1}} \wedge \cdots \wedge e^{i_{k}}\right)= \pm\left(\epsilon_{i_{1}} \epsilon_{i_{2}} \cdots \epsilon_{i_{k}}\right) e^{i_{k+1}} \wedge \cdots \wedge e^{i_{n}}
$$

where we take the + sign if and only if $e^{i_{1}} \wedge \cdots \wedge e^{i_{k}} \wedge e^{i_{k+1}} \wedge \cdots \wedge e^{i_{n}}= e^{1} \wedge \cdots \wedge e^{n}$.

Proof. Formula (9.7) above actually defines a map $\wedge^{k} \mathrm{V}^{*} \rightarrow \bigwedge^{n-k} \mathrm{V}^{*}$. For this proof we denote this map by * and show it satisfies the same defining equations as the Hodge star. It is enough to check that the defining formula $\alpha \wedge * \beta=\langle\alpha \mid \beta\rangle$ vol holds for typical (orthonormal) basis elements $\alpha=e^{i_{1}} \wedge \cdots \wedge e^{i_{k}}$ and $\beta=e^{m_{1}} \wedge \cdots \wedge e^{m_{k}}$. We have

$$
\begin{aligned}
& \left(e^{m_{1}} \wedge \cdots \wedge e^{m_{k}}\right) \wedge *\left(e^{i_{1}} \wedge \cdots \wedge e^{i_{k}}\right) \\
& =e^{m_{1}} \wedge \cdots \wedge e^{m_{k}} \wedge\left( \pm e^{i_{k+1}} \wedge \cdots \wedge e^{i_{n 1}}\right)
\end{aligned}
$$

The last expression is zero unless $\left\{m_{1}, \ldots, m_{k}\right\} \cup\left\{i_{k+1}, \ldots, i_{n}\right\}=\{1,2$, $\ldots, n\}$, or in other words, unless $\left\{i_{1}, \ldots, i_{k}\right\}=\left\{m_{1}, \ldots, m_{k}\right\}$. But this is also true for

$$
\left\langle e^{m_{1}} \wedge \cdots \wedge e^{m_{k}} \mid e^{i_{1}} \wedge \cdots \wedge e^{i_{k}}\right\rangle \text { vol. }
$$

On the other hand if $\left\{i_{1}, \ldots, i_{k}\right\}=\left\{m_{1}, \ldots, m_{k}\right\}$, then both (9.8) and (9.9) give $\pm$ vol. So the lemma is proved up to a sign. We leave it to the reader to show that the definitions are such that the signs match.

Remark 9.24. In case the scalar product is positive definite, $\epsilon_{1}=\epsilon_{2}=\cdots =\epsilon_{n}=1$ and so the formulas are a bit less cluttered.

The Hodge star operators for each $k$ can be combined to give a Hodge star operator $\Lambda \mathrm{V}^{*} \rightarrow \Lambda \mathrm{V}^{*}$. We also note that the scalar product on $\Lambda^{0} \mathrm{V}^{*}=\mathbb{R}$ is just $\langle a \mid b\rangle=a b$. Then, we still have the formula $\alpha \wedge * \beta= \langle\alpha \mid \beta\rangle$ vol for any $\alpha, \beta$ in the algebra $\wedge \mathrm{V}^{*}$.

Proposition 9.25. Let V be as above with $\operatorname{dim}(\mathrm{V})=n$. The following identities hold for the star operator:\\
(1) $* 1=\mathrm{vol}$;\\
$(2) * \operatorname{vol}=(-1)^{\operatorname{ind}(g)}$;\\
(3) $* * \alpha=(-1)^{\operatorname{ind}(g)}(-1)^{k(n-k)} \alpha$ for $\alpha \in \bigwedge^{k} \mathrm{V}^{*}$;\\
(4) $\langle * \alpha \mid * \beta\rangle=(-1)^{\operatorname{ind}(g)}\langle\alpha \mid \beta\rangle$.

Proof. (1) and (2) follow directly from the definitions. For (3), it suffices to let $\alpha=e^{\sigma(1)} \wedge \cdots \wedge e^{\sigma(k)}$ for some permutation $\sigma \in S_{n}$. We first compute $*\left(e^{\sigma(k+1)} \wedge \cdots \wedge e^{\sigma(n)}\right)$. We must have $*\left(e^{\sigma(k+1)} \wedge \cdots \wedge e^{\sigma(n)}\right)=c e^{\sigma(1)} \wedge \cdots \wedge e^{\sigma(k)}$ for some constant $c$. On the other hand,

$$
\begin{aligned}
\epsilon_{\sigma(k+1)} \cdots \epsilon_{\sigma(n)} \mathrm{vol} & =\left\langle e^{\sigma(k+1)} \wedge \cdots \wedge e^{\sigma(n)} \mid e^{\sigma(k+1)} \wedge \cdots \wedge e^{\sigma(n)}\right\rangle \mathrm{vol} \\
& =\left(e^{\sigma(k+1)} \wedge \cdots \wedge e^{\sigma(n)}\right) \wedge *\left(e^{\sigma(k+1)} \wedge \cdots \wedge e^{\sigma(n)}\right) \\
& =\left(e^{\sigma(k+1)} \wedge \cdots \wedge e^{\sigma(n)}\right) \wedge c e^{\sigma(1)} \wedge \cdots \wedge e^{\sigma(k)} \\
& =(-1)^{k(n-k)} c \operatorname{sgn}(\sigma) \operatorname{vol}
\end{aligned}
$$

so that $c=\epsilon_{\sigma(k+1)} \cdots \epsilon_{\sigma(n)}(-1)^{k(n-k)} \operatorname{sgn}(\sigma)$. Using this and equation (9.7), we have

$$
\begin{aligned}
* *\left(e^{\sigma(1)} \wedge \cdots \wedge e^{\sigma(k)}\right)= & * \epsilon_{\sigma(1)} \epsilon_{\sigma(2)} \cdots \epsilon_{\sigma(k)} \operatorname{sgn}(\sigma) e^{\sigma(k+1)} \wedge \cdots \wedge e^{\sigma(n)} \\
= & \epsilon_{\sigma(1)} \epsilon_{\sigma(2)} \cdots \epsilon_{\sigma(k)} \epsilon_{\sigma(k+1)} \cdots \epsilon_{\sigma(n)}(\operatorname{sgn}(\sigma))^{2} \\
& \quad \times(-1)^{k(n-k)} e^{\sigma(1)} \wedge \cdots \wedge e^{\sigma(k)} \\
= & (-1)^{\operatorname{ind}(g)}(-1)^{k(n-k)} e^{\sigma(1)} \wedge \cdots \wedge e^{\sigma(k)}
\end{aligned}
$$

which implies the result. For (4) we compute as follows:

$$
\begin{aligned}
\langle * \alpha \mid * \beta\rangle \text { vol } & =* \alpha \wedge * * \beta=(-1)^{\operatorname{ind}(g)}(-1)^{k(n-k)} * \alpha \wedge \beta \\
& =(-1)^{\operatorname{ind}(g)} \beta \wedge * \alpha=(-1)^{\operatorname{ind}(g)}\langle\alpha \mid \beta\rangle \text { vol }
\end{aligned}
$$

We obtain a formula for the star operator that uses a basis that is not necessarily orthonormal. Recall $\epsilon_{\ell_{1} \ldots \ell_{k}}^{i_{1} \ldots i_{k}}$ defined by equation (8.2). If we let $\epsilon_{i_{1} \ldots i_{n}}:=\epsilon_{i_{1} \ldots i_{n}}^{12 \ldots n}$, then $\epsilon_{i_{1} \ldots i_{n}}$ is just the sign of the permutation $\binom{12 \ldots n}{i_{1} \ldots i_{n}}$. Using this notation, we have the following theorem.

Theorem 9.26. Let (V,g) be an oriented scalar product space with corresponding volume element vol. Let $e^{1}, \ldots, e^{n}$ be a positively oriented basis of $\mathrm{V}^{*}$. If $\omega=\frac{1}{k!} \sum \omega_{i_{1} \ldots i_{k}} e^{i_{i}} \wedge \cdots \wedge e^{i_{k}}$ (or $\left.\omega=\sum_{i_{1}<i_{2} \cdots<i_{k}} \omega_{i_{1} \ldots i_{k}} e^{i_{i}} \wedge \cdots \wedge e^{i_{k}}\right)$, then

$$
* \omega=\sqrt{\left|\operatorname{det}\left[g_{i j}\right]\right|} \sum_{j_{k+1}<\cdots<j_{n}} \omega^{j_{1} \ldots j_{k}} \epsilon_{j_{1} \ldots j_{k} j_{k+1} \ldots j_{n}} e^{j_{k+1}} \wedge \cdots \wedge e^{j_{n}} .
$$

Proof. Let us use the common abbreviation $g:=\operatorname{det}\left[g_{i j}\right]$. Let \& denote the operator defined so that $\boldsymbol{\operatorname { s i n }} \omega$ is given by the right hand side of equation (9.10) above. Our task is to show that $\boldsymbol{m}=*$. For this we need to show that $\alpha \wedge \boldsymbol{\beta}=\langle\alpha \mid \beta\rangle$ vol for all $\alpha, \beta \in \bigwedge^{k} \mathrm{V}^{*}$. First note that for fixed $j_{1}, \ldots, j_{n}$ we have $\epsilon_{j_{1} \ldots j_{n}} \epsilon_{i_{1} \ldots i_{n}}^{j_{1} \ldots j_{n}}=\epsilon_{i_{1} \ldots i_{n}}$ (no sum). We compute the coefficients of $\alpha \wedge \boldsymbol{\omega} \beta$ using formula (8.3):

$$
\begin{aligned}
(\alpha \wedge & \frac{1}{\beta)_{i_{1} \ldots i_{n}}} \\
= & \frac{1}{k!(n-k)!} \sum \alpha_{j_{1} \ldots j_{k}}\left(\boldsymbol{\omega}^{\beta} \beta\right)_{j_{k+1} \ldots j_{n}} \epsilon_{i_{1} \ldots i_{n}}^{j_{1} \ldots j_{n}} \\
& =\frac{1}{k!(n-k)!} \frac{1}{k!} \sum|g|^{1 / 2} \alpha_{j_{1} \ldots j_{k}} \beta^{m_{1} \ldots m_{k}} \epsilon_{m_{1} \ldots m_{k} j_{k+1} \ldots j_{n}} \epsilon_{i_{1} \ldots i_{n}}^{j_{1} \ldots j_{n}} \\
& =|g|^{1 / 2} \frac{1}{k!(n-k)!} \sum \alpha_{j_{1} \ldots j_{k}} \beta^{j_{1} \ldots j_{k}} \epsilon_{j_{1} \ldots j_{k} j_{k+1} \ldots j_{n}} \epsilon_{i_{1} \ldots i_{n}}^{j_{1} \ldots j_{n}} \\
& =|g|^{1 / 2} \frac{1}{k!} \sum \alpha_{j_{1} \ldots j_{k}} \beta^{j_{1} \ldots j_{k}} \epsilon_{i_{1} \ldots i_{n}}=(\langle\alpha \mid \beta\rangle \mathrm{vol})_{i_{1} \ldots i_{n}}
\end{aligned}
$$

and so $\alpha \wedge \boldsymbol{\mathfrak { m }} \beta=\langle\alpha \mid \beta\rangle$ vol; thus by uniqueness we have $\boldsymbol{\oplus}=*$.\\
The definitions we gave for volume element, star operator, and the musical isomorphisms all apply to each tangent space $T_{p} M$ of a semi-Riemannian manifold and these concepts globalize to the whole of $T M$ accordingly.

Let $M$ be oriented. The metric volume form induced by the metric tensor $g$ is defined to be the $n$-form vol such that $\operatorname{vol}_{p}$ is the metric volume form on $T_{p} M$ matching the orientation. If ( $U, \mathrm{x}$ ) is a positively oriented chart on $M$, then we have

$$
\left.\operatorname{vol}\right|_{U}=\sqrt{\left|\operatorname{det}\left[g_{i j}\right]\right|} d x^{1} \wedge \cdots \wedge d x^{n}
$$

If $f$ is a smooth function with compact support, then we may integrate $f$ over $M$ by using the volume form:

$$
\int_{M} f \mathrm{vol}
$$

The volume of a compact oriented Riemannian manifold is

$$
\operatorname{vol}(M):=\int_{M} \operatorname{vol}
$$

The volume of an open set $D \subset M$ is defined as\\
$\operatorname{vol}(D):=\sup \left\{\int_{M} f \operatorname{vol}: \operatorname{supp}(f) \subset D\right.$ with $f \in C_{c}^{\infty}(M)$ and $\left.0 \leq f \leq 1\right\}$,\\
where $C_{c}^{\infty}(M)$ denotes the space of smooth functions with compact support.\\
Example 9.27. The metric volume form $\operatorname{vol}_{\mathbb{R}^{n}}$ on $\mathbb{R}^{n}$ is $d x^{1} \wedge \cdots \wedge d x^{n}$, where $x^{1}, \ldots, x^{n}$ are standard coordinates. Also,

$$
\operatorname{vol}_{\mathbb{R}^{n}}\left(v_{1 p}, \ldots, v_{n p}\right)=\operatorname{det}\left(v_{1}, \ldots, v_{n}\right)
$$

where $v_{i p}=\left(p, v_{i}\right) \in T_{p} \mathbb{R}^{n}=\{p\} \times \mathbb{R}^{n}$.\\
Once again assume that $M$ is oriented by a metric volume form "vol". The star operator gives two types of maps, which are both denoted by $*$. Namely, the bundle maps * : $\bigwedge^{k} T^{*} M \rightarrow \bigwedge^{n-k} T^{*} M$, which have the obvious definition in terms of the star operators on each fiber, and the maps on forms $*: \Omega^{k}(M) \rightarrow \Omega^{n-k}(M)$, which are also induced in the obvious way. The definitions are set up so that $*$ is natural with respect to restriction and so that for any oriented orthonormal frame field $\left\{E_{1}, \ldots, E_{n}\right\}$ with dual $\left\{\theta^{1}, \ldots, \theta^{n}\right\}$ we have $*\left(\theta^{i_{1}} \wedge \cdots \wedge \theta^{i_{k}}\right)= \pm\left(\epsilon_{i_{1}} \epsilon_{i_{2}} \cdots \epsilon_{i_{k}}\right) \theta^{j_{1}} \wedge \cdots \wedge \theta^{j_{n-k}}$, where, as before, we use the + sign if and only if $\theta^{i_{1}} \wedge \cdots \wedge \theta^{i_{k}} \wedge \theta^{j_{1}} \wedge \cdots \wedge \theta^{j_{n-k}}= \theta^{1} \wedge \cdots \wedge \theta^{n}=$ vol. As expected we have

$$
\begin{aligned}
* 1 & =\mathrm{vol} \\
* \mathrm{vol} & =(-1)^{\operatorname{ind}(g)} \\
* * \alpha & =(-1)^{\operatorname{ind}(g)}(-1)^{k(n-k)} \alpha
\end{aligned}
$$

for $\alpha \in \Omega^{k}(M)$.\\
Example 9.28. On open sets in $\mathbb{R}^{3}$ the star operator associated with the standard metric and volume is given by

$$
\begin{array}{ccc}
f & \mapsto & f d x \wedge d y \wedge d z \\
f_{1} d x+f_{2} d y+f_{3} d z & \mapsto & f_{1} d y \wedge d z+f_{2} d z \wedge d x+f_{3} d x \wedge d y \\
g_{1} d y \wedge d z+g_{2} d z \wedge d x+g_{3} d x \wedge d y & \mapsto & g_{1} d x+g_{2} d y+g_{3} d z \\
f d x \wedge d y \wedge d z & \mapsto & f .
\end{array}
$$

For an oriented Riemannian manifold, div will always be the divergence with respect to the metric volume form (giving the orientation). From equation (9.4) we see that the local coordinate expression for the divergence of $X=\sum X^{k} \frac{\partial}{\partial x^{k}}$ is

$$
\operatorname{div} X=\sum_{k=1}^{n} \frac{1}{\sqrt{\operatorname{det} g}} \frac{\partial}{\partial x^{k}}\left(\sqrt{\operatorname{det} g} X^{k}\right)
$$

where $\operatorname{det} g=\operatorname{det}\left(g_{i j}\right)$.

Definition 9.29. The gradient of a smooth function $f$ on a Riemannian manifold ( $M, g$ ) is defined to be

$$
\operatorname{grad}(f):=\sharp d f
$$

and the Laplacian $\Delta$ is defined by

$$
\Delta f:=-\operatorname{div}(\operatorname{grad}(f))
$$

if $M$ is oriented. If $\Delta f=0$, then $f$ is called a harmonic function.\\
The minus sign in the above definition is a matter of convention, and this choice of sign is popular with differential geometers. Ironically, this choice of sign makes the Laplacian a positive operator.

Lemma 9.30. If $M$ is an oriented semi-Riemannian manifold with metric volume element vol, then $i_{X} \mathrm{vol}=* b X$ for $X \in \mathfrak{X}(M)$.

Proof. We show that for fixed $p \in M$ we have $i_{v_{1}}$ vol $=* b v_{1}$ for all $v_{1} \in T_{p} M$. First consider the case where $v$ is not a null vector and $\left\langle v_{1}, v_{1}\right\rangle=\epsilon_{1}= \pm 1$. Since the maps $\left(v_{2}, \ldots, v_{n}\right) \mapsto i_{v_{1}} \operatorname{vol}\left(v_{2}, \ldots, v_{n}\right)$ and $\left(v_{2}, \ldots, v_{n}\right) \mapsto \left(* b v_{1}\right)\left(v_{2}, \ldots, v_{n}\right)$ are both alternating multilinear maps, it suffices to check that they are equal on some basis. We extend to an orthonormal basis $\left(v_{1}, v_{2}, \ldots, v_{n}\right)$ with dual basis $e^{1}, \ldots, e^{n}$. Then we wish to show that $i_{v_{1}} \operatorname{vol}\left(v_{i_{1}}, \ldots, v_{i_{n-1}}\right)=\left(* b v_{1}\right)\left(v_{i_{1}}, \ldots, v_{i_{n-1}}\right)$ for $1 \leq i_{1}<\cdots<i_{n-1} \leq n$. But if $v_{i_{1}}=v_{1}$, then

$$
i_{v_{1}} \operatorname{vol}\left(v_{i_{1}}, \ldots, v_{i_{n-1}}\right)=\operatorname{vol}\left(v_{1}, v_{1}, \ldots\right)=0
$$

and

$$
\begin{aligned}
\left(* b v_{1}\right)\left(v_{i_{1}}, \ldots, v_{i_{n-1}}\right) & =\epsilon_{1}\left(* e^{1}\right)\left(v_{i_{1}}, \ldots, v_{i_{n-1}}\right) \\
& =\epsilon_{1}^{2} e^{2} \wedge \cdots \wedge e^{n}\left(v_{1}, \ldots\right)=0
\end{aligned}
$$

On the other hand, if $v_{i_{1}} \neq v_{1}$, then $\left(v_{i_{1}}, \ldots, v_{i_{n-1}}\right)=\left(v_{2}, \ldots, v_{n}\right)$ and

$$
i_{v_{1}} \operatorname{vol}\left(v_{2}, \ldots, v_{n}\right)=\operatorname{vol}\left(v_{1}, v_{2}, \ldots, v_{n}\right)=1
$$

while

$$
\begin{aligned}
\left(* b v_{1}\right)\left(v_{2}, \ldots, v_{n}\right) & =\epsilon_{1}\left(* e^{1}\right)\left(v_{2}, \ldots, v_{n}\right) \\
& =\epsilon_{1}^{2} e^{2} \wedge \cdots \wedge e^{n}\left(v_{2}, \ldots, v_{n}\right)=1
\end{aligned}
$$

Thus $i_{v_{1}}$ vol $=* b v_{1}$ for all nonnull vectors $v_{1}$. But since both sides are linear in $v_{1}$, this establishes the result.

Proposition 9.31. If $M$ is an oriented semi-Riemannian manifold with metric volume element vol, then $\operatorname{div} X=(-1)^{\operatorname{ind} g} * d * b X$ for $X \in \mathfrak{X}(M)$.

Proof. By Lemma 9.30 and Cartan's formula (8.56), we have $\mathcal{L}_{X}$ vol $= d i_{X} \operatorname{vol}=d * b X$ so that $* \operatorname{div} X \operatorname{vol}=* d * b X$ or

$$
\operatorname{div} X * \operatorname{vol}=* d * b X
$$

or

$$
(-1)^{\operatorname{ind} g} \operatorname{div} X=* d * b X .
$$

Proposition 9.32. If $M$ is an oriented semi-Riemannian manifold with metric volume element vol, then

$$
\operatorname{div} f X=\langle\operatorname{grad} f, X\rangle+f \operatorname{div} X .
$$

Proof. Write vol $=\mu$ and compute as follows:

$$
\begin{aligned}
(\operatorname{div} f X) \mu & =\mathcal{L}_{f X} \mu=d i_{f X} \mu+i_{f X} d \mu \\
& =d\left(f i_{X} \mu\right)=d f \wedge i_{X} \mu+f d i_{X} \mu \\
& =d f \wedge i_{X} \mu+f(\operatorname{div} X) \mu \\
& =-i_{X}(d f \wedge \mu)+i_{X} d f \wedge \mu+f(\operatorname{div} X) \mu \\
& =d f(X) \mu+f(\operatorname{div} X) \mu \\
& =(\langle X, \operatorname{grad} f\rangle+f(\operatorname{div} X)) \mu,
\end{aligned}
$$

where in the third line we use the fact that $i_{X}$ is a graded derivation and $d f \wedge \mu=0$.

Exercise 9.33. Show that the local expression for $\Delta f$ is

$$
\Delta f=-\frac{1}{\sqrt{\operatorname{det} g}} \sum_{j, k} \frac{\partial}{\partial x^{j}}\left(g^{j k} \sqrt{\operatorname{det} g} \frac{\partial f}{\partial x^{k}}\right) .
$$

\subsection*{Integral Formulas}
Definition 9.34. Let $S$ be a regular codimension one submanifold of an $n$-dimensional Riemannian manifold ( $M, g$ ). A smooth vector field $N$ along $S$ is called a unit normal field if $\langle N, N\rangle=1$ and if $N(p) \perp T_{p} S$ for every $p \in S$.

Let ( $M, g$ ) be oriented by a metric volume form vol and suppose that a regular ( $n-1$ )-dimensional submanifold $S$ of $M$ has a (global) unit normal field $N$ along $S$. If $N$ is extended to a field $\widetilde{N}$ on a neighborhood of $S$, then $i_{\widetilde{N}}$ vol is an ( $n-1$ )-form on this neighborhood. Clearly, the restriction of $i_{\widetilde{N}}$ vol to $S$ is independent of this extension, and so we can denote this restriction by $\left.i_{N} \mathrm{vol}\right|_{S}$. Then we have

$$
\left(\left.i_{N} \operatorname{vol}\right|_{S}\right)_{p}\left(v_{1}, \ldots, v_{n-1}\right)=\operatorname{vol}_{p}\left(N(p), v_{1}, \ldots, v_{n-1}\right)
$$

for $v_{1}, \ldots, v_{n-1} \in T_{p} S$. This is a volume form on $S$ which orients $S$ and is clearly the metric volume form for the induced metric $\iota^{*} g$ on $S$. Note that\\
$\left.i_{N} \operatorname{vol}\right|_{S}$ is sometimes denoted by $i_{N} \operatorname{vol}$, but this is an abuse of notation since $i_{N} \operatorname{vol}(p)\left(v_{1}, \ldots, v_{n-1}\right)$ makes sense for $v_{1}, \ldots, v_{n-1} \in T_{p} M$, but this is not what we want. On the other hand, $\left.\int_{S} i_{N} \operatorname{vol}\right|_{S}$ is $\int_{S} \imath^{*}\left(i_{\widetilde{N}} \operatorname{vol}\right)$, where $\imath: S \hookrightarrow M$ is the inclusion map. In turn, this means the same thing as $\int_{S} i_{N}$ vol by convention.

Proposition 9.35. Let $S$ be a regular codimension one submanifold of an $n$-dimensional oriented Riemannian manifold ( $M, g$ ) and let $\iota: S \hookrightarrow M$ be the inclusion. Let $N$ be a smooth unit normal vector field along $S$. Let $\operatorname{vol}_{M}$ be the volume form of $M$ and $\operatorname{vol}_{S}$ the volume form of $S$ with respect to the orientation induced by $N$ and the metric $\iota^{*} g$. If $X$ is a smooth vector field along $S$, then

$$
i_{X} \operatorname{vol}_{M}=\langle X, N\rangle \operatorname{vol}_{S} .
$$

Proof. Let $X^{\top}:=X-\langle X, N\rangle N$;

$$
\begin{aligned}
i_{X} \operatorname{vol}_{M} & =i_{X^{\top}} \operatorname{vol}_{M}+i_{\langle X, N\rangle N} \operatorname{vol}_{M} \\
& =i_{X^{\top}} \operatorname{vol}_{M}+\langle X, N\rangle i_{N} \operatorname{vol}_{M},
\end{aligned}
$$

and since $\left.i_{N} \operatorname{vol}_{M}\right|_{\partial M}=\operatorname{vol}_{S}$, it remains to show that $i_{X^{\top}} \operatorname{vol}_{M}=0$. But for $v_{1}, \ldots, v_{n-1} \in T_{p} S$,

$$
i_{X^{\top}} \operatorname{vol}_{M}\left(v_{1}, \ldots, v_{n-1}\right)=\operatorname{vol}_{M}\left(X^{\top}, v_{1}, \ldots, v_{n-1}\right)=0
$$

since $X^{\top}, v_{1}, \ldots, v_{n-1}$ cannot be linearly independent. \(\square\)

Now suppose that an oriented ( $M, g$ ) has a boundary $\partial M$. Since $\partial M$ is basically a codimension one regular submanifold of $M$, the above constructions can be applied to $\partial M$ so that $\partial M$ is Riemannian with induced metric. Since $\partial M$ has a global outward pointing vector field, we may normalize to obtain a global outward pointing unit normal field $N$ along $\partial M$. The orientation induced on $\partial M$ by the metric volume form $\left.i_{N} \operatorname{vol}_{M}\right|_{\partial M}$ is exactly the induced orientation on $\partial M$ described previously. We will denote $\left.i_{N} \operatorname{vol}_{M}\right|_{\partial M}$ by $\operatorname{vol}_{\partial M}$ when the orientation and normal field are understood.

Exercise 9.36. Let $\mathbf{x}: V \rightarrow U \subset M$ be a parametrization of an open subset of a surface $S$ in $\mathbb{R}^{3}$ and $N:=\left(\mathbf{x}_{u^{1}} \times \mathbf{x}_{u^{2}}\right) /\left\|\mathbf{x}_{u^{1}} \times \mathbf{x}_{u^{2}}\right\|$. Show that if $d S:=\left.i_{N} \operatorname{vol}\right|_{S}$, where vol is the usual volume form on $\mathbb{R}^{3}$, then $\mathbf{x}^{*} d S=\left\|\mathbf{x}_{u^{1}} \times \mathbf{x}_{u^{2}}\right\| d u^{1} \wedge d u^{2}$.

We now introduce a special ( $n-1$ )-form $\sigma$ on $\mathbb{R}^{n}$ (we use the same symbol for all $n$ if no confusion arises). This form is used in many constructions. If $x^{1}, \ldots, x^{n}$ denote standard coordinates, then

$$
\sigma:=\sum_{i=1}^{n}(-1)^{i-1} x^{i} d x^{1} \wedge \cdots \wedge \widehat{d x^{i}} \wedge \cdots \wedge d x^{n}
$$

where, as usual, the caret denotes omission. For example, on $\mathbb{R}^{3}$

$$
\sigma=x d y \wedge d z-y d x \wedge d z+z d x \wedge d y
$$

Proposition 9.37. If we use the outward pointing normal field $N$ on $S^{n-1}$ and the metric induced from $\mathbb{R}^{n}$, then the induced volume form on $S^{n-1}$ (the area form) is given by

$$
\operatorname{vol}_{S^{n-1}}=\iota^{*} \sigma
$$

where $\iota: S^{n-1} \hookrightarrow \mathbb{R}^{n}$ is inclusion.\\
Proof. Let $v_{i}=\sum v_{i}^{j} \partial /\left.\partial x^{j}\right|_{x} \in T_{x} S^{n-1}$. Identify elements of $T_{x} S^{n-1}$ with vectors in $\mathbb{R}^{n}$ so that $v_{i}$ is the column vector $\left(v_{i}^{1}, \ldots, v_{i}^{n}\right)^{T}$. Then we use expansion along the first column by cofactors to obtain

$$
\begin{aligned}
\operatorname{vol}_{S^{n-1}} & \left(v_{1}, \ldots, v_{n-1}\right)=\operatorname{det}\left(x, v_{1}, \ldots, v_{n-1}\right) \\
& =\sum_{i=1}^{n}(-1)^{i-1} x^{i} \operatorname{det}\left(M_{i}\right) \\
& =\sum_{i=1}^{n}(-1)^{i-1} x^{i}\left(d x^{1} \wedge \cdots \wedge \widehat{d x^{i}} \wedge \cdots \wedge d x^{n}\right)\left(v_{1}, \ldots, v_{n-1}\right)
\end{aligned}
$$

where $M_{i}$ is the submatrix obtained by deleting the first column and $i$-th row.

Exercise 9.38. Show that the volume form on $S^{n-1}(R)=\left\{\sum_{i=1}^{n}\left(x^{i}\right)^{2}=\right. \left.R^{2}\right\}$ is $\frac{1}{R} \iota^{*} \sigma$, where $\iota: S^{n-1}(R) \hookrightarrow \mathbb{R}^{n}$ is the inclusion map.

Theorem 9.39. Let $M$ be an oriented Riemannian $n$-manifold with boundary. For any $X \in \mathfrak{X}(M)$ with compact support,

$$
\int_{M} \operatorname{div} X \operatorname{vol}_{M}=\int_{\partial M}\langle X, N\rangle \operatorname{vol}_{\partial M}
$$

where $N$ is the outward pointing unit normal field on $\partial M$.\\
Proof. We have $\int_{M} \operatorname{div} X \operatorname{vol}_{M}=\int_{M} d\left(i_{X} \operatorname{vol}_{M}\right)=\int_{\partial M} i_{X} \operatorname{vol}_{M}$ which is equal to $\int_{\partial M}\langle X, N\rangle \operatorname{vol}_{\partial M}$ by Proposition 9.35.

Corollary 9.40. Let $M$ be a compact oriented Riemannian $n$-manifold with boundary and $X \in \mathfrak{X}(M)$. Then for $f \in C^{\infty}(M)$ we have

$$
\int_{M}\langle X, \operatorname{grad} f\rangle \operatorname{vol}_{M}+\int_{M} f \operatorname{div} X \operatorname{vol}_{M}=\int_{\partial M} f\langle X, N\rangle \operatorname{vol}_{\partial M}
$$

Proof. We have

$$
\int_{M}(\langle\operatorname{grad} f, X\rangle+f \operatorname{div} X) \operatorname{vol}_{M}=\int_{M} \operatorname{div} f X=\int_{\partial M} f\langle X, N\rangle \operatorname{vol}_{\partial M}
$$

Theorem 9.41. Let $M$ be an oriented Riemannian $n$-manifold with boundary and $f, g \in C^{\infty}(M)$. Then

$$
\int_{M}\langle\operatorname{grad} f, \operatorname{grad} g\rangle \operatorname{vol}_{M}-\int_{M} f \Delta g \operatorname{vol}_{M}=\int_{\partial M} f\langle\operatorname{grad} g, N\rangle \operatorname{vol}_{\partial M}
$$

and

$$
\int_{M}(-f \Delta g+g \Delta f) \operatorname{vol}_{M}=\int_{\partial M}[f\langle\operatorname{grad} g, N\rangle-g\langle\operatorname{grad} f, N\rangle] \operatorname{vol}_{\partial M}
$$

where $N$ is the outward unit normal field.

\section*{Proof.}
$$
\begin{aligned}
-\int_{M} f \Delta g \operatorname{vol}_{M} & =\int_{M} f \operatorname{div}(\operatorname{grad} g) \operatorname{vol}_{M} \\
& =\int_{\partial M} f\langle\operatorname{grad} g, N\rangle \operatorname{vol}_{\partial M}-\int_{M}\langle\operatorname{grad} g, \operatorname{grad} f\rangle \operatorname{vol}_{M}
\end{aligned}
$$

by Corollary 9.40. The second formula follows from the first by interchanging the roles of $f$ and $g$ and subtracting. \(\square\)

The integral equations of the previous theorem are called Green's first and second formulas respectively (when comparing with other versions in the literature, do not forget our convention: $\Delta f:=-\operatorname{divgrad} f$ ). If the boundary of $M$ is empty, then the boundary terms are zero, so that for instance $\int_{M} f \Delta g \operatorname{vol}_{M}=\int_{M} g \Delta f \operatorname{vol}_{M}=0$. We close this section with an application of Green's formulas.

Theorem 9.42 (Hopf). Let ( $M, g$ ) be a compact, connected and oriented Riemannian $n$-manifold without boundary. If $f \in C^{\infty}(M)$ and $\Delta f \geq 0$ (or $\Delta f \leq 0$ ), then $f$ is a constant function.

Proof. Suppose $\Delta f \geq 0$. Then integration gives

$$
0 \leq \int_{M} 1 \Delta f \operatorname{vol}_{M}=\int_{M} f \Delta 1 \operatorname{vol}_{M}=0
$$

which means that $\Delta f \equiv 0$ on $M$. Now use Green's second formula with $f=g$ to get

$$
\int_{M}\|\operatorname{grad} f\|^{2} \operatorname{vol}_{M}=\int_{M} f \Delta f \operatorname{vol}_{M}=0
$$

so that $\|\operatorname{grad} f\|^{2} \equiv 0$ on $M$. Thus $f$ is constant since $M$ is connected. \(\square\)

\subsection*{The Hodge Decomposition}
In this section we follow [War]. Let ( $M, g$ ) be an oriented semi-Riemannian $n$-manifold. (However, we soon restrict to the compact Riemannian case.) Let vol denote the metric volume form. For each $k$, the scalar products induced by $g_{p}$ on $\bigwedge^{k} T_{p}^{*} M$ for each $p$ combine to give a symmetric $C^{\infty}(M)$ bilinear map

$$
\langle\cdot \mid \cdot\rangle: \Omega^{k}(M) \times \Omega^{k}(M) \rightarrow C^{\infty}(M)
$$

with $\langle\alpha \mid \beta\rangle(p):=\langle\alpha(p) \mid \beta(p)\rangle_{p}=g_{p}(\alpha(p), \beta(p))$. We see that $\langle\alpha \mid \beta\rangle$ vol $= \alpha \wedge * \beta$. We can put a scalar product ( $\cdot \mid \cdot$ ) on the space $\Omega(M)=\sum_{k} \Omega^{k}(M)$ by letting $(\alpha \mid \beta)=0$ for $\alpha \in \Omega^{k_{1}}(M)$ and $\beta \in \Omega^{k_{2}}(M)$ with $k_{1} \neq k_{2}$ and letting

$$
(\alpha \mid \beta)=\int_{M} \alpha \wedge * \beta=\int_{M}\langle\alpha \mid \beta\rangle \text { vol if } \alpha, \beta \in \Omega^{k}(M) .
$$

Definition 9.43. Let ( $M, g$ ) be an oriented semi-Riemannian $n$-manifold. For each $k$ with $0 \leq k \leq n=\operatorname{dim} M$, the codifferential $\delta: \Omega^{k}(M) \rightarrow \Omega^{k-1}(M)$ is defined by

$$
\delta:=(-1)^{\operatorname{ind}(g)}(-1)^{n(k+1)+1} * d *
$$

and $\delta=0$ on $\Omega^{0}(M)$. These operators combine to give a linear map $\Omega(M) \rightarrow \Omega(M)$, which is also denoted $\delta$.

Remark 9.44. Notice that if $M$ is nonorientable, then $*$ is still defined locally by choosing an orientation valid on a chart domain. But * occurs twice in the formula for $\delta$, so any sign ambiguities cancel and thus $\delta$ is globally well-defined even if $M$ is nonorientable!

Proposition 9.45. The codifferential $\delta$ is the formal adjoint of the exterior derivative on $\Omega(M)$. That is,

$$
(d \alpha \mid \beta)=(\alpha \mid \delta \beta)
$$

for all $\alpha, \beta$ in $\Omega(M)$ with compact support in the interior of $M$.\\
Proof. It suffices to check that $(d \alpha \mid \beta)=(\alpha \mid \delta \beta)$ for $\alpha$ a $(k-1)$-form and $\beta$ a $k$-form. In this case, $d * \beta$ is an ( $n-k+1$ )-form and so

$$
\begin{aligned}
* \delta \beta & =*\left((-1)^{\operatorname{ind}(g)}(-1)^{n(k+1)+1} * d * \beta\right) \\
& =(-1)^{\operatorname{ind}(g)}(-1)^{n(k+1)+1} * * d * \beta \\
& =(-1)^{n(k+1)+1}(-1)^{(n-k+1)(k-1)} d * \beta=(-1)^{k^{2}} d * \beta
\end{aligned}
$$

where we have used that $n(k+1)+1+(n-k+1)(k-1)=2 k+2 k n-k^{2}= k^{2} \bmod 2$. Then we have

$$
\begin{aligned}
d(\alpha \wedge * \beta) & =d \alpha \wedge * \beta+(-1)^{k-1} \alpha \wedge d * \beta \\
& =d \alpha \wedge * \beta+(-1)^{k-1}(-1)^{k^{2}} \alpha \wedge * \delta \beta \\
& =d \alpha \wedge * \beta-\alpha \wedge * \delta \beta
\end{aligned}
$$

since $k-1+k^{2}=k(k+1)-1$, which is always odd. Now we integrate and use Stokes' theorem to get the desired result:

$$
\int_{M} d \alpha \wedge * \beta=\int_{M} \alpha \wedge * \delta \beta
$$

Obviously, $\delta$ is defined on $\Omega^{k}(U)$ for open sets $U$ in $M$ and $\delta$ is natural with respect to restriction.

Definition 9.46. For $0 \leq k \leq n$, the differential operator $\Delta: \Omega(M) \rightarrow \Omega(M)$ defined by $\Delta=\delta d+d \delta$ is called the Laplace-Beltrami operator. For each $k$, the restriction of $\Delta$ to $\Omega^{k}(M)$ is also called the Laplace-Beltrami operator (on $k$-forms). If $\Delta \omega=0$, we call $\omega$ a harmonic form.

Notice that by Remark 9.44 the Laplace-Beltrami operator is defined whether or not $M$ is orientable. On $\Omega^{0}(M)=C^{\infty}(M)$ the operator $\Delta$ reduces to the Laplacian defined earlier.

Proposition 9.47. The following assertions hold:\\
(i) For all $\alpha, \beta \in \Omega(M)$, we have

$$
(\Delta \alpha \mid \beta)=(d \alpha \mid d \beta)+(\delta \alpha \mid \delta \beta)
$$

(ii) $\Delta$ is formally self-adjoint. That is $(\Delta \alpha \mid \beta)=(\alpha \mid \Delta \beta)$ for all $\alpha, \beta \in \Omega(M)$ with compact support in the interior of $M$.

Proof. By Proposition 9.45 we have

$$
\begin{aligned}
(\Delta \alpha \mid \beta) & =(\delta d \alpha+d \delta \alpha \mid \beta)=(\delta d \alpha \mid \beta)+(d \delta \alpha \mid \beta)=(d \alpha \mid d \beta)+(\delta \alpha \mid \delta \beta) \\
& =(\alpha \mid \delta d \beta+d \delta \beta)=(\alpha \mid \Delta \beta)
\end{aligned}
$$

which proves (i). But by symmetry $(\Delta \alpha \mid \beta)=(d \alpha \mid d \beta)+(\delta \alpha \mid \delta \beta)=(\alpha \mid \Delta \beta)$, which proves (ii).

In the remainder of this section, we restrict attention to an oriented compact Riemannian $n$-manifold ( $M, g$ ) without boundary. In this case, $(\alpha \mid \beta)=\int_{M} \alpha \wedge * \beta$ defines a positive definite inner product on $\Omega(M)$. We write $\|\alpha\|:=(\alpha \mid \alpha)^{1 / 2}$.

Proposition 9.48. If ( $M, g$ ) is compact, oriented and Riemannian, then $\Delta \omega=0$ if and only if $d \omega=0$ and $\delta \omega=0$.

Proof. This follows from $(\Delta \omega \mid \omega)=(d \omega \mid d \omega)+(\delta \omega \mid \delta \omega)$ since $(\cdot \mid \cdot)$ is positive definite in the Riemannian case.

Now we wish to consider the equation $\Delta \omega=\alpha$ for $\alpha \in \Omega^{k}(M)$. First notice that if $\Delta \omega=\alpha$ holds, then for any $\beta \in \Omega^{k}(M)$ we have $(\Delta \omega \mid \beta)= (\omega \mid \Delta \beta)=(\alpha \mid \beta)$. Then since $|(\omega \mid \Delta \beta)|=|(\alpha \mid \beta)| \leq\|\alpha\|\|\beta\|$, the map $\ell_{\omega}$ : $\beta \mapsto(\omega \mid \Delta \beta)$ is a bounded linear functional. Employing a common idea from the theory of differential equations, we make the following definition:

Definition 9.49. A bounded linear functional $\ell: \Omega^{k}(M) \rightarrow \mathbb{R}$ is called a weak solution of the equation $\Delta \omega=\alpha$ if

$$
\ell(\Delta \beta)=(\alpha \mid \beta)
$$

for all $\beta \in \Omega^{k}(M)$.\\
Notice that if a weak solution $\ell$ is represented by $\omega$, then it must be an ordinary solution since in that case we have

$$
(\Delta \omega \mid \beta)=(\omega \mid \Delta \beta)=\ell(\Delta \beta)=(\alpha \mid \beta)
$$

for all $\beta$, which means that $\Delta \omega=\alpha$. We will use the following powerful regularity result (see [War] for a proof):

Theorem 9.50. Let $\alpha \in \Omega^{k}(M)$. If $\ell$ is a weak solution of $\Delta \omega=\alpha$, then there exists $\omega \in \Omega^{k}(M)$ such that

$$
\ell(\beta)=(\omega \mid \beta)
$$

for all $\beta \in \Omega^{k}(M)$ and hence

$$
\Delta \omega=\alpha
$$

We will also need a rather technical result whose proof can also be found in [War].

Proposition 9.51. Let $\left\{\alpha_{i}\right\}_{1}^{\infty}$ be a sequence in $\Omega^{k}(M)$. Suppose that for some $C>0$ we have

$$
\left\|\alpha_{i}\right\| \leq C \text { and }\left\|\Delta \alpha_{i}\right\| \leq C
$$

for all $i$. Then $\left\{\alpha_{i}\right\}_{1}^{\infty}$ has a Cauchy subsequence.\\
Let $\mathcal{H}^{k}=\left\{\omega \in \Omega^{k}(M): \Delta \omega=0\right\}$ and let $\left(\mathcal{H}^{k}\right)^{\perp}$ denote $\left\{\beta \in \Omega^{k}(M):\right. (\beta \mid \omega)=0$ for all $\left.\omega \in \mathcal{H}^{k}\right\}$.

Lemma 9.52. There exists a constant $C>0$ such that $\|\beta\| \leq C\|\Delta \beta\|$ for all $\beta \in\left(\mathcal{H}^{k}\right)^{\perp}$.

Proof. Suppose there is no such constant $C$. Then we may find a sequence $\left\{\beta_{i}\right\}_{1}^{\infty} \subset\left(\mathcal{H}^{k}\right)^{\perp}$ such that $\left\|\beta_{i}\right\|=1$ while $\lim _{i \rightarrow \infty}\left\|\Delta \beta_{i}\right\|=0$. By Proposition 9.51, $\left\{\beta_{i}\right\}_{1}^{\infty}$ has a Cauchy subsequence, which we may as well assume to be $\left\{\beta_{i}\right\}_{1}^{\infty}$. Hence we have that for any fixed $\theta \in \Omega^{k}(M)$, the sequence $\left\{\left(\beta_{i} \mid \theta\right)\right\}$ is Cauchy in $\mathbb{R}$ and so converges to some number. Now we define $\ell: \Omega^{k}(M) \rightarrow \mathbb{R}$ by

$$
\ell(\theta):=\lim _{i \rightarrow \infty}\left(\beta_{i} \mid \theta\right) .
$$

Then,

$$
\ell(\Delta \theta)=\lim _{i \rightarrow \infty}\left(\beta_{i} \mid \Delta \theta\right)=\lim _{i \rightarrow \infty}\left(\Delta \beta_{i} \mid \theta\right)=0
$$

and since $\ell$ is clearly bounded, it is a weak solution to $\Delta \omega=0$. By Theorem 9.50, there must exist $\beta \in \Omega^{k}(M)$ such that $\ell(\theta)=(\beta \mid \theta)$ for all $\theta \in \Omega^{k}(M)$. It follows that $(\beta \mid \theta)=\lim _{i \rightarrow \infty}\left(\beta_{i} \mid \theta\right)$ for all $\theta$ and so $\beta_{i} \rightarrow \beta$ since

$$
\begin{aligned}
\left\|\beta_{i}-\beta\right\|^{2} & =\left(\beta_{i}-\beta \mid \beta_{i}-\beta\right) \\
& =\left(\beta_{i} \mid \beta_{i}\right)-2\left(\beta_{i} \mid \beta\right)+(\beta \mid \beta) \rightarrow 0
\end{aligned}
$$

Since $\left\|\beta_{i}\right\|=1$, we must have $\|\beta\|=1$ and of course $\beta \in\left(\mathcal{H}^{k}\right)^{\perp}$. But by Theorem 9.50, $\Delta \beta=0$ so $\beta \in \mathcal{H}^{k} \cap\left(\mathcal{H}^{k}\right)^{\perp}=0$, which contradicts $\|\beta\|=1$. We conclude that $C$ exists after all.

Theorem 9.53 (Hodge decomposition). Let ( $M, g$ ) be compact and oriented (and without boundary). For each $k$ with $0 \leq k \leq n=\operatorname{dim} M$, the space of harmonic $k$-forms $\mathcal{H}^{k}$ is finite-dimensional. Furthermore, we have an orthogonal decomposition of $\Omega^{k}(M)$,

$$
\begin{aligned}
\Omega^{k}(M) & =\Delta\left(\Omega^{k}(M)\right) \oplus \mathcal{H}^{k} \\
& =d \delta\left(\Omega^{k}(M)\right) \oplus \delta d\left(\Omega^{k}(M)\right) \oplus \mathcal{H}^{k} \\
& =d\left(\Omega^{k-1}(M)\right) \oplus \delta\left(\Omega^{k+1}(M)\right) \oplus \mathcal{H}^{k}
\end{aligned}
$$

Proof. If $\mathcal{H}^{k}$ were infinite-dimensional, then it would contain an infinite orthonormal sequence $\left\{\omega_{i}\right\}_{1}^{\infty}$. In this case, we would have

$$
\left\|\omega_{i}-\omega_{j}\right\|^{2}=2
$$

for all $i, j$ with $i \neq j$. By Proposition 9.51, this sequence would contain a Cauchy subsequence. But this contradicts the above equation. Thus $\mathcal{H}^{k}$ must be finite-dimensional.

We now prove the orthogonal decomposition $\Omega^{k}(M)=\Delta\left(\Omega^{k}(M)\right) \oplus \mathcal{H}^{k}$. The other two decompositions can be derived from the first, and we leave\\
this as a problem for the reader (Problem 9). Choose an orthonormal basis $\omega_{1}, \ldots, \omega_{d}$ for $\mathcal{H}^{k}$. If $\alpha \in \Omega^{k}(M)$, then we may write

$$
\alpha=\beta+\sum_{i=1}^{d}\left(\omega_{i} \mid \alpha\right) \omega_{i}
$$

where $\beta \in\left(\mathcal{H}^{k}\right)^{\perp}$. It is easy to show that this decomposition is unique, so we have the orthogonal decomposition $\Omega^{k}(M)=\left(\mathcal{H}^{k}\right)^{\perp} \oplus \mathcal{H}^{k}$ and our task is to show that $\left(\mathcal{H}^{k}\right)^{\perp}=\Delta\left(\Omega^{k}(M)\right)$. Since $(\Delta \alpha \mid \omega)=(\alpha \mid \Delta \omega)=0$ whenever $\omega \in \mathcal{H}^{k}$, we see that $\Delta\left(\Omega^{k}(M)\right) \subset\left(\mathcal{H}^{k}\right)^{\perp}$. Now let $\alpha \in\left(\mathcal{H}^{k}\right)^{\perp}$ and define a linear functional $\ell$ on $\Delta\left(\Omega^{k}(M)\right)$ by

$$
\ell(\Delta \theta):=(\alpha \mid \theta) .
$$

This is well-defined since if $\Delta \theta_{1}=\Delta \theta_{2}$, then $\theta_{1}-\theta_{2} \in \mathcal{H}^{k}$ and so $\left(\alpha \mid \theta_{1}\right)- \left(\alpha \mid \theta_{2}\right)=\left(\alpha \mid \theta_{1}-\theta_{2}\right)=0$. We show that $\ell$ is bounded. Let $\phi:=\theta-H(\theta)$, where $H_{k}: \Omega^{k}(M) \rightarrow \mathcal{H}^{k}$ is the orthogonal projection. Then using Lemma 9.52 we have

$$
\begin{aligned}
|\ell(\Delta \theta)| & =|\ell(\Delta \phi)|=|(\alpha \mid \phi)| \leq\|\alpha\|\|\phi\| \\
& \leq C\|\alpha\|\|\Delta \phi\|=C\|\alpha\|\|\Delta \theta\| .
\end{aligned}
$$

By the Hahn-Banach theorem, the functional $\ell$ extends to a bounded functional $\widetilde{\ell}$ defined on all of $\Omega^{k}(M)$, which is then a weak solution of $\Delta \omega=\alpha$. By Theorem 9.50, there is an $\omega \in \Omega^{k}(M)$ with $\Delta \omega=\alpha$ so $\left(\mathcal{H}^{k}\right)^{\perp} \subset \Delta\left(\Omega^{k}(M)\right)$ and so $\left(\mathcal{H}^{k}\right)^{\perp}=\Delta\left(\Omega^{k}(M)\right)$.

In order to take full advantage of the Hodge decomposition, we now introduce a so-called "Green's operator" $G_{k}: \Omega^{k}(M) \rightarrow\left(\mathcal{H}^{k}\right)^{\perp}$. We simply define $G_{k}(\alpha)$ to be the unique solution of $\Delta \omega=\alpha-H_{k}(\alpha)$ where $H_{k}$ : $\Omega^{k}(M) \rightarrow \mathcal{H}^{k}$ is the orthogonal projection as above. The $G_{k}$ combine to give a map $G: \Omega(M) \rightarrow \bigoplus_{k}\left(\mathcal{H}^{k}\right)^{\perp}$ also called the Green's operator.

Lemma 9.54. Let $G_{k}: \Omega^{k}(M) \rightarrow\left(\mathcal{H}^{k}\right)^{\perp}$ be the Green's operator defined above for each $k$. Then\\
(i) $G_{k}$ is formally self-adjoint;\\
(ii) if $L: \Omega^{k}(M) \rightarrow \Omega^{r}(M)$ is linear and commutes with $\Delta$, then $G$ commutes with $L$ (that is, $L \circ G_{k}=G_{r} \circ L$ ). In particular, $G$ commutes with $d$ and $\delta$.

Proof. We have

$$
\begin{aligned}
\left(G_{k}(\alpha) \mid \beta\right) & =\left(G_{k}(\alpha) \mid \beta-H_{k}(\beta)\right)=\left(G_{k}(\alpha) \mid \Delta G_{k}(\beta)\right)=\left(\Delta G_{k}(\alpha) \mid G_{k}(\beta)\right) \\
& =\left(\alpha-H_{k}(\alpha) \mid G_{k}(\beta)\right)=\left(\alpha \mid G_{k}(\beta)\right),
\end{aligned}
$$

so $G$ is self-adjoint. For each $j$, let $\pi_{j}: \Omega^{j}(M) \rightarrow\left(\mathcal{H}^{j}\right)^{\perp}$ denote orthogonal projection (thus $\pi_{j}+H_{j}=\operatorname{id}_{\Omega^{j}}$ ). Now suppose that $L: \Omega^{k}(M) \rightarrow \Omega^{r}(M)$ is linear and commutes with $\Delta$. Notice that by definition we have $G_{k}= \left(\Delta \mid\left(\mathcal{H}^{k}\right)^{\perp}\right) \circ \pi_{k}$. The fact that $L \Delta=\Delta L$ implies that $L\left(\mathcal{H}^{k}\right) \subset \mathcal{H}^{r}$. Also, since $\Delta\left(\Omega^{k}(M)\right)=\left(\mathcal{H}^{k}\right)^{\perp}$, we have $L\left(\left(\mathcal{H}^{k}\right)^{\perp}\right) \subset\left(\mathcal{H}^{r}\right)^{\perp}$. Thus

$$
\begin{gathered}
L \circ \pi_{k}=\pi_{r} \circ L \\
L \circ\left(\Delta \mid\left(\mathcal{H}^{k}\right)^{\perp}\right)=\left.\left(\Delta \mid\left(\mathcal{H}^{r}\right)^{\perp}\right) \circ L\right|_{\left(\mathcal{H}^{r}\right)^{\perp}}
\end{gathered}
$$

and

$$
\left(\Delta \mid\left(\mathcal{H}^{r}\right)^{\perp}\right)^{-1} \circ L=\left.L\right|_{\left(\mathcal{H}^{r}\right)^{\perp}} \circ\left(\Delta \mid\left(\mathcal{H}^{k}\right)^{\perp}\right)^{-1} .
$$

It follows that $G$ commutes with $L$.\\
We note in passing that $G$ maps bounded sequences into sequences that have Cauchy subsequences. Indeed, suppose that $\left\{\alpha_{i}\right\}_{1}^{\infty} \subset \Omega^{k}(M)$ is a sequence with $\left\|\alpha_{i}\right\| \leq C$. If $\beta_{i}:=G\left(\alpha_{i}\right)$, then using Lemma 9.52 we have

$$
\left\|\beta_{i}\right\| \leq\left\|\Delta \beta_{i}\right\|=\left\|\alpha_{i}-H\left(\alpha_{i}\right)\right\| \leq\left\|\alpha_{i}\right\| \leq C
$$

and so by Proposition 9.51, $\left\{\beta_{i}\right\}_{1}^{\infty}$ has a Cauchy subsequence.\\
Theorem 9.55. Let ( $M, g$ ) be a compact Riemannian manifold (without boundary). Then each de Rham cohomology class contains a unique harmonic representative.

Proof. First assume that $M$ is orientable and fix an orientation. Let $\alpha \in \Omega^{k}(M)$ and use the Hodge decomposition and the definition of $G$ to obtain

$$
\alpha=\Delta G_{k}(\alpha)+H_{k}(\alpha)=d \delta G_{k}(\alpha)+\delta d G_{k}(\alpha)+H_{k}(\alpha) .
$$

Then since $G$ commutes with $d$, we have

$$
\alpha=d \delta G_{k}(\alpha)+\delta G_{k+1}(d \alpha)+H_{k}(\alpha) .
$$

So if $d \alpha=0$, then $\alpha-H_{k}(\alpha)=d \delta G_{k}(\alpha)$, and so $H_{k}(\alpha)$ represents the same cohomology class as $\alpha$. To show uniqueness, suppose that $\alpha_{1}$ and $\alpha_{2}$ are both harmonic and in the same class so that $\alpha_{2}-\alpha_{1}=d \beta$ for some $\beta$. Then we have $0=d \beta+\left(\alpha_{1}-\alpha_{2}\right)$. But $\alpha_{1}-\alpha_{2}$ is orthogonal to $d \beta$ since by Proposition 9.48

$$
\left(d \beta \mid \alpha_{1}-\alpha_{2}\right)=\left(\beta \mid \delta \alpha_{1}-\delta \alpha_{2}\right)=(\beta \mid 0)=0
$$

Next suppose that $M$ is nonorientable. If $\pi: M^{\text {or }} \rightarrow M$ is the 2fold orientation cover of $M$ (Section 8.7), then there is a unique metric on $M^{\text {or }}$ such that $\pi$ restricts to an isometry on sufficiently small open sets ( $\pi$ is a Riemannian covering). Now suppose that $c \in H^{k}(M)$ is a de Rham cohomology class and consider the class $\pi^{*} c \in H^{k}\left(M^{\text {or }}\right)$ (see the discussion after Definition 8.41). Since $M^{\text {or }}$ is orientable, $\pi^{*} c$ is uniquely represented\\
by a harmonic form $\widetilde{\eta}$ on $M^{\text {or }}$. If $\tau: M^{\text {or }} \rightarrow M^{\text {or }}$ is the involution which transposes the two points in each fiber, then it is easy to see that $\tau$ is an isometry so that $\tau^{*} \widetilde{\eta}$ is also harmonic. It follows from the identity $\pi \circ \tau=\pi$ that $\tau^{*} \widetilde{\eta}$ is also a representative of the cohomology class $\pi^{*} c$. Indeed, $\left[\tau^{*} \widetilde{\eta}\right]= \tau^{*}[\widetilde{\eta}]=\tau^{*} \pi^{*} c=\pi^{*} c$. So by uniqueness, we must have $\tau^{*} \widetilde{\eta}=\widetilde{\eta}$, which is exactly the condition that guarantees that $\widetilde{\eta}=\pi^{*} \eta$ for some $\eta \in \Omega^{k}(M)$. But $\pi$ is a local isometry and so by Problem 5, $\eta$ must be harmonic. We show that $\eta$ represents $c$. Suppose that $c=[\mu]$ for $\mu \in \Omega^{k}(M)$. Then both $\pi^{*} \mu$ and $\pi^{*} \eta$ represent the class $\pi^{*} c$, so $\pi^{*}(\eta-\mu)=d \beta$ for some $\beta$. But then, as for $\tau^{*} \widetilde{\eta}$ and $\widetilde{\eta}$ above, $\tau^{*} \beta$ is in the same class as $\beta$, which means that $d \beta=d \tau^{*} \beta$. Then

$$
\pi^{*}(\eta-\mu)=d \beta=\frac{1}{2}\left(d \beta+d \tau^{*} \beta\right)=d\left(\frac{1}{2}\left(\beta+\tau^{*} \beta\right)\right)
$$

Since $\tau$ is an involution, $\beta+\tau^{*} \beta$ is $\tau$-invariant, that is, $\tau^{*}\left(\beta+\tau^{*} \beta\right)=$ ( $\beta+\tau^{*} \beta$ ), so there exists $\alpha \in \Omega^{k}(M)$ with $\pi^{*} \alpha=\frac{1}{2}\left(\beta+\tau^{*} \beta\right)$. Thus

$$
\pi^{*}(\eta-\mu)=d \pi^{*} \alpha=\pi^{*} d \alpha
$$

and since $\pi^{*}$ is a local diffeomorphism, we must have $\eta-\mu=d \alpha$, which means that $[\eta]=[\mu]=c$. The uniqueness of the harmonic representative $\eta$ is clear.

Notice that if $\tau$ is the involution of $M^{\text {or }}$ introduced in the above proof, then $\pi^{*} \alpha=\tau^{*} \pi^{*} \alpha$ for any $\alpha \in \Omega^{k}(M)$.\\
Corollary 9.56. If $M$ is compact, then its cohomology spaces $H^{k}(M)$ are finite-dimensional.

Proof. Any manifold can be given a Riemannian metric, and so if $M$ is compact and oriented, then Theorems 9.53 and 9.55 combine to give the result. Now suppose that $M$ is not orientable. Let $\alpha \in \Omega^{k}(M)$ and suppose that $\pi^{*} \alpha=d \beta$ so that $\pi^{*}[\alpha]=0$. Then $\pi^{*} \alpha=\tau^{*} \pi^{*} \alpha=\tau^{*} d \beta=d \tau^{*} \beta$. Now $\pi^{*} \theta=\frac{1}{2}\left(\beta+\tau^{*} \beta\right)$ for some $\theta$ and so

$$
\pi^{*} \alpha=d\left(\frac{1}{2}\left(\beta+\tau^{*} \beta\right)\right)=d \pi^{*} \theta=\pi^{*} d \theta
$$

Since $\pi$ is a local diffeomorphism, we have $\alpha=d \theta$ or $[\alpha]=0$. Thus the map $\pi^{*}: H^{k}(M) \rightarrow H^{k}\left(M^{\text {or }}\right)$ has trivial kernel and is injective. The finitedimensionality of $H^{k}(M)$ now follows from that of $H^{k}\left(M^{\text {or }}\right)$.

The results above allow us to give a quick proof of Poincaré duality for de Rham cohomology. Choose an orientation for $M$. We define a bilinear pairing $H^{k}(M) \times H^{n-k}(M) \rightarrow \mathbb{R}$ as follows:

$$
(([\omega],[\eta]))=\int_{M} \omega \wedge \eta
$$

This is well-defined since if $\omega_{1}=\omega+d \alpha$ and $\eta_{1}=\eta+d \beta$ (with $\omega, \eta$ closed), then by Stokes' theorem

$$
\begin{aligned}
\int_{M} \omega_{1} \wedge \eta_{1} & =\int_{M} \omega \wedge \eta+\int_{M} d \alpha \wedge \eta+\int_{M} \omega \wedge d \beta+\int_{M} d \alpha \wedge d \beta \\
& =\int_{M} \omega \wedge \eta+\int_{M} d(\alpha \wedge \eta)-\int_{M} d(\omega \wedge \beta)+\int_{M} d(\alpha \wedge d \beta) \\
& =\int_{M} \omega \wedge \eta
\end{aligned}
$$

We wish to show that the pairing defined above is nondegenerate. Thus, given any nonzero $[\omega] \in H^{k}(M)$ we wish to produce a $[\eta] \in H^{n-k}(M)$ such that $(([\omega],[\eta])) \neq 0$. Choose a Riemannian metric and metric volume element for $M$. We may assume that $\omega$ is the harmonic representative for $[\omega]$. But it is easily checked that $\Delta$ commutes with $*$ and so $* \omega$ is also harmonic. In particular, $* \omega$ is closed and so represents a cohomology class $[* \omega]$. Then $[* \omega]$ is the desired class since

$$
(([\omega],[* \omega]))=\int_{M} \omega \wedge * \omega=\int_{M}\langle\omega \mid * \omega\rangle \text { vol }=\|\omega\|^{2}>0
$$

For each fixed $[\omega] \in H^{k}(M)$, we have the linear map $\ell_{[\omega]} \in\left(H^{n-k}(M)\right)^{*}$ given by $\ell_{[\omega]}([\eta]):=(([\omega],[\eta]))$. Since the pairing is nondegenerate, the map $[\omega] \mapsto \ell_{[\omega]}$ defines an isomorphism from $H^{k}(M)$ to $\left(H^{n-k}(M)\right)^{*}$. Thus we have proved the following:

Theorem 9.57 (Poincaré duality). If $M$ is an orientable compact $n$-manifold without boundary, then we have an isomorphism

$$
H^{k}(M) \cong\left(H^{n-k}(M)\right)^{*}
$$

coming from the pairing defined above.\\
We will take up Poincaré duality again in Chapter 10.

\subsection*{Vector Analysis on $\mathbb{R}^{3}$}
In $\mathbb{R}^{3}$, the 1 -forms may all be written (even globally) in the form $\theta=f_{1} d x+ f_{2} d y+f_{3} d z$ for some smooth functions $f_{1}, f_{2}$ and $f_{3}$ and all 2 -forms $\beta$ may be written $\beta=g_{1} d y \wedge d z+g_{2} d z \wedge d x+g_{3} d x \wedge d y$. The forms $d y \wedge d z, d z \wedge d x, d x \wedge d y$ constitute a basis (in the module sense) for the space of 2 -forms on $\mathbb{R}^{3}$ just as $d x, d y, d z$ form a basis for the 1 -forms. The single form $d x \wedge d y \wedge d z$ provides a module basis for the 3 -forms in $\mathbb{R}^{3}$. Suppose that $\mathbf{x}(u, v)$ parametrizes a surface $S \subset \mathbb{R}^{3}$ so that we have a map $\mathbf{x}: U \rightarrow \mathbb{R}^{3}$. Then the surface is\\
oriented by this parametrization, and the integral of $\beta$ over $S$ is

$$
\begin{aligned}
& \int_{S} \beta= \int_{S} g_{1} d y \wedge d z+g_{2} d z \wedge d x+g_{3} d x \wedge d y \\
&=\int_{U}\left(g_{1}(\mathbf{x}(u, v)) \frac{\partial(y, z)}{\partial(u, v)}+g_{2}(\mathbf{x}(u, v)) \frac{\partial(z, x)}{\partial(u, v)}\right. \\
&\left.+g_{3}(\mathbf{x}(u, v)) \frac{\partial(x, y)}{\partial(u, v)}\right) d u d v
\end{aligned}
$$

Here and in the following we disregard technical issues about integration (but recall Theorem 9.7).

Exercise 9.58. Find the integral of $\beta=x d y \wedge d z+d z \wedge d x+x z d x \wedge d y$ over the sphere oriented by the parametrization given by the usual spherical coordinates $\phi, \theta, \rho$.

If $\omega=h d x \wedge d y \wedge d z$ has support in a bounded open subset $U \subset \mathbb{R}^{3}$ which we may take to be given the usual orientation implied by the rectangular coordinates $x, y, z$, then

$$
\int_{U} \omega=\int_{U} h d x \wedge d y \wedge d z=\int_{U} h d x d y d z
$$

In order to relate differential forms on $\mathbb{R}^{3}$ to vector calculus on $\mathbb{R}^{3}$, we will need some ways to relate forms to vector fields. To a 1 -form $\theta=f_{1} d x+ f_{2} d y+f_{3} d z$, we can obviously associate the vector field $\sharp \theta=f_{1} \mathbf{i}+f_{2} \mathbf{j}+f_{3} \mathbf{k}$. But recall that this association depends on the notion of orthonormality provided by the dot product. If $\theta$ is expressed in say spherical coordinates $\theta=\widetilde{f}_{1} d \rho+\widetilde{f}_{2} d \theta+\widetilde{f}_{3} d \phi$, then it is not true that $\sharp \theta=\widetilde{f}_{1} \mathbf{i}+\widetilde{f}_{2} \mathbf{j}+\widetilde{f}_{3} \mathbf{k}$. Neither is it generally true that $\sharp \theta=\widetilde{f}_{1} \widehat{\rho}+\widetilde{f}_{2} \widehat{\theta}+\widetilde{f}_{3} \widehat{\phi}$, where $\widehat{\rho}, \widehat{\theta}, \widehat{\phi}$ are unit vector fields in the coordinate directions ${ }^{1}$ and where the $\widetilde{f}_{i}$ are just the $f_{i}$ expressed in polar coordinates. Rather, our general formulas give

$$
\sharp \theta=\widetilde{f}_{1} \widehat{\rho}+\widetilde{f}_{2} \frac{1}{\rho} \widehat{\theta}+\widetilde{f}_{3} \frac{1}{\rho \sin \theta} \widehat{\phi}
$$

In rectangular coordinates $x, y, z$, we have

$$
\begin{array}{llll}
\sharp: & d x & \rightarrow & \mathbf{i}, \\
\sharp: & d y & \rightarrow & \mathbf{j}, \\
\sharp: & d z & \rightarrow & \mathbf{k},
\end{array}
$$

while in spherical coordinates we have

$$
\begin{array}{cccc}
\sharp: & d \rho & \mapsto & \widehat{\rho}, \\
\sharp: & \rho d \theta & \mapsto & \widehat{\theta}, \\
\sharp: & \rho \sin \theta d \phi & \mapsto & \widehat{\phi} .
\end{array}
$$

\footnotetext{${ }^{1}$ Here $\theta$ is the polar angle ranging from 0 to $\pi$.
}As an example, we can derive the familiar formula for the gradient in spherical coordinates by first just writing $f$ in the new coordinates $f(\rho, \theta, \phi):= f(x(\rho, \theta, \phi), y(\rho, \theta, \phi), z(\rho, \theta, \phi))$ and then sharping the differential

$$
d f=\frac{\partial f}{\partial \rho} d \rho+\frac{\partial f}{\partial \theta} d \theta+\frac{\partial f}{\partial \phi} d \phi
$$

to get

$$
\operatorname{grad} f=\sharp d f=\frac{\partial f}{\partial \rho} \frac{\partial}{\partial \rho}+\frac{1}{\rho} \frac{\partial f}{\partial \theta} \frac{\partial}{\partial \theta}+\frac{1}{\rho \sin \theta} \frac{\partial f}{\partial \phi} \frac{\partial}{\partial \phi},
$$

where we have used

$$
\left(g^{i j}\right)=\left[\begin{array}{ccc}
1 & 0 & 0 \\
0 & \rho & 0 \\
0 & 0 & \rho \sin \theta
\end{array}\right]^{-1}=\left[\begin{array}{ccc}
1 & 0 & 0 \\
0 & \frac{1}{\rho} & 0 \\
0 & 0 & \frac{1}{\rho \sin \theta}
\end{array}\right] .
$$

In order to proceed to the point of including the curl and divergence of traditional vector calculus, we need a way to relate 2 -forms with vector fields. This part definitely depends on the fact that we are talking about forms in $\mathbb{R}^{3}$. We associate to a 2 -form $\eta$ the vector fields $\sharp(* \eta)$. Thus $g_{1} d y \wedge d z+g_{2} d z \wedge d x+g_{3} d x \wedge d y$ gives the vector field $X=g_{1} \mathbf{i}+g_{2} \mathbf{j}+g_{3} \mathbf{k}$.

Now we can see how the usual divergence of a vector field comes about. First flat the vector field, say $X=f_{1} \mathbf{i}+f_{2} \mathbf{j}+f_{3} \mathbf{k}$, to obtain b $X=f_{1} d x+ f_{2} d y+f_{3} d z$ and then apply the star operator to obtain $f_{1} d y \wedge d z+f_{2} d z \wedge d x+f_{3} d x \wedge d y$. Finally, we apply exterior differentiation to obtain

$$
\begin{aligned}
d\left(f_{1} d y\right. & \left.\wedge d z+f_{2} d z \wedge d x+f_{3} d x \wedge d y\right) \\
\quad=d f_{1} & \wedge d y \wedge d z+d f_{2} \wedge d z \wedge d x+d f_{3} \wedge d x \wedge d y \\
= & \left(\frac{\partial f_{1}}{\partial x} d x+\frac{\partial f_{1}}{\partial y} d y+\frac{\partial f_{1}}{\partial z} d z\right) \wedge d y \wedge d z+\text { the other two terms } \\
= & \frac{\partial f_{1}}{\partial x} d x \wedge d y \wedge d z+\frac{\partial f_{2}}{\partial x} d x \wedge d y \wedge d z+\frac{\partial f_{3}}{\partial x} d x \wedge d y \wedge d z \\
= & \left(\frac{\partial f_{1}}{\partial x}+\frac{\partial f_{2}}{\partial x}+\frac{\partial f_{3}}{\partial x}\right) d x \wedge d y \wedge d z
\end{aligned}
$$

Now we see the divergence appearing. In fact, if we apply the star operator one more time, we get the function $\operatorname{div} X=\frac{\partial f_{1}}{\partial x}+\frac{\partial f_{2}}{\partial x}+\frac{\partial f_{3}}{\partial x}$. We are thus led to $* d *(b X)=\operatorname{div} X$ which agrees with the definition of divergence given earlier for a general semi-Riemannian manifold.

What about the curl? For this, we just take $d(b X)$ to get

$$
\begin{aligned}
& d\left(f_{1} d x+f_{2} d y+f_{3} d z\right)=d f_{1} \wedge d x+d f_{2} \wedge d y+d f_{3} \wedge d z \\
& =\left(\frac{\partial f_{1}}{\partial x} d x+\frac{\partial f_{2}}{\partial y} d y+\frac{\partial f_{3}}{\partial z} d z\right) \wedge d x+\text { the obvious other two terms } \\
& =\left(\frac{\partial f_{2}}{\partial y}-\frac{\partial f_{3}}{\partial z}\right) d y \wedge d z+\left(\frac{\partial f_{3}}{\partial z}-\frac{\partial f_{1}}{\partial x}\right) d z \wedge d x+\left(\frac{\partial f_{1}}{\partial x}-\frac{\partial f_{2}}{\partial y}\right) d x \wedge d y
\end{aligned}
$$

and then apply the star operator and sharping to get back to vector fields obtaining

$$
\left(\frac{\partial f_{2}}{\partial y}-\frac{\partial f_{3}}{\partial z}\right) \mathbf{i}+\left(\frac{\partial f_{3}}{\partial z}-\frac{\partial f_{1}}{\partial x}\right) \mathbf{j}+\left(\frac{\partial f_{1}}{\partial x}-\frac{\partial f_{2}}{\partial y}\right) \mathbf{k}=\operatorname{curl} X .
$$

In short, we have

$$
\sharp * d(b X)=\operatorname{curl} X
$$

Exercise 9.59. Show that the fact that $d d=0$ leads to both of the following familiar facts:

$$
\begin{aligned}
& \operatorname{curl}(\operatorname{grad} f)=0 \\
& \operatorname{div}(\operatorname{curl} X)=0
\end{aligned}
$$

The 3-form $d x \wedge d y \wedge d z$ is the (oriented) volume element of $\mathbb{R}^{3}$. Every 3 -form is a function times this volume form, and integration of a 3 -form over a sufficiently nice subset (say a compact region) is given by $\int_{D} \omega= \int_{D} f d x \wedge d y \wedge d z=\int_{D} f d x d y d z$ (usual Riemann integral). Let us denote $d x \wedge d y \wedge d z$ by $\mathrm{d} V$. Of course, $\mathrm{d} V$ is not to be considered as the exterior derivative of some object $V$. Let $u^{1}, u^{2}, u^{3}$ be curvilinear coordinates on an open set $U \subset \mathbb{R}^{3}$ and let $g$ denote the determinant of the matrix $\left[g_{i j}\right]$ where $g_{i j}=\left\langle\frac{\partial}{\partial u^{i}}, \frac{\partial}{\partial u^{j}}\right\rangle$. Then $\mathrm{d} V=\sqrt{g} d u^{1} \wedge d u^{2} \wedge d u^{3}$. A familiar example is the case when ( $u^{1}, u^{2}, u^{3}$ ) are spherical coordinates $\rho, \theta, \phi$, in which case

$$
\mathrm{d} V=\rho^{2} \sin \theta d \rho \wedge d \theta \wedge d \phi
$$

and if $D \subset \mathbb{R}^{3}$ is parametrized by these coordinates, then

$$
\begin{aligned}
\int_{D} f \mathrm{~d} V & =\int_{D} f(\rho, \theta, \phi) \rho^{2} \sin \theta d \rho \wedge d \theta \wedge d \phi \\
& =\int_{D} f(\rho, \theta, \phi) \rho^{2} \sin \theta d \rho d \theta d \phi
\end{aligned}
$$

If we go to the trouble of writing Stokes' theorem for curves, surfaces and domains in $\mathbb{R}^{3}$ in terms of vector fields associated to the forms in an appropriate way, we obtain the following familiar theorems (using standard\\
notation):

$$
\begin{aligned}
\int_{c} \nabla f \cdot d \mathbf{r} & =f(\mathbf{r}(b))-f(\mathbf{r}(a)) \\
\iint_{S} \operatorname{curl}(\mathbf{X}) \times \mathbf{d S} & =\oint_{c} X \cdot d \mathbf{r}(\text { Stokes' theorem }) \\
\iiint_{D} \operatorname{div}(\mathbf{X}) d V & =\iint_{S} \mathbf{X} \cdot \mathbf{d S}(\text { Divergence theorem })
\end{aligned}
$$

Similar and simpler things can be done in $\mathbb{R}^{2}$ leading for example to the following version of Green's theorem for a planar domain $D$ with (oriented) boundary $c=\partial D$ :

$$
\int_{D}\left(\frac{\partial M}{\partial y}-\frac{\partial M}{\partial y}\right) d x \wedge d y=\int_{D} d(M d x+N d y)=\int_{c} M d x+N d y
$$

All of the standard integral theorems from vector calculus are special cases of the general Stokes' theorem.

\subsection*{Electromagnetism}
In this subsection we take a short trip into physics. Consider Maxwell's equations ${ }^{2}$ :

$$
\begin{aligned}
\nabla \cdot \mathbf{B} & =0 \\
\nabla \times \mathbf{E}+\frac{\partial \mathbf{B}}{\partial t} & =0 \\
\nabla \cdot \mathbf{E} & =\varrho \\
\nabla \times \mathbf{B}-\frac{\partial \mathbf{E}}{\partial t} & =\mathbf{j}
\end{aligned}
$$

Here $\mathbf{E}$ and $\mathbf{B}$, the electric and magnetic fields, are functions of space and time. We write $\mathbf{E}=\mathbf{E}(t, \mathbf{x}), \mathbf{B}=\mathbf{B}(t, \mathbf{x})$. The notation suggests that we have conceptually separated space and time as if we were stuck in the conceptual framework of the Galilean spacetime. Our purpose is to slowly discover how much better the theory becomes when we combine space and time in Minkowski spacetime $\mathbb{R}_{1}^{4}$. Recall that $\mathbb{R}_{1}^{4}$ is treated as a semi-Riemannian manifold, which is $\mathbb{R}^{4}$ endowed with the indefinite metric $\langle x, y\rangle_{\nu}=-x^{0} y^{0}+\sum_{i=1}^{3} x^{i} y^{i}$. Here the standard coordinates are conventionally denoted by $\left(x^{0}, x^{1}, x^{2}, x^{3}\right)$, and $x^{0}$ is to be thought of as a time coordinate and is also denoted by $t$ (we take units so that the speed of light $c$ is unity). In what follows, we let $\mathbf{r}=\left(x^{1}, x^{2}, x^{3}\right)$, so that $\left(x^{0}, x^{1}, x^{2}, x^{3}\right)=(t, \mathbf{r})$.

\footnotetext{${ }^{2}$ Actually, this is the form of Maxwell's equations after a certain convenient choice of units, and we are ignoring the somewhat subtle distinction between the two types of electric fields $E$ and $D$ and the two types of magnetic fields $B$ and $H$ and their relation in terms of dielectric constants.
}The electric field $\mathbf{E}$ is produced by the presence of charged particles. Under normal conditions a generic material is composed of a large number of atoms. To simplify matters, we will think of the atoms as being composed of just three types of particle; electrons, protons and neutrons. Protons carry a positive charge, electrons carry a negative charge and neutrons carry no charge. Normally, each atom will have a zero net charge since it will have an equal number of electrons and protons. If a relatively small percent of the electrons in a material body are stripped from their atoms and conducted away, then there will be a net positive charge on the body. In the vicinity of the body, there will be an electric field which exerts a force on charged bodies. Let us assume for simplicity that the charged body which has the larger, positive charge, is a point particle and stationary at $\mathbf{r}_{0}$ with respect to a rigid rectangular coordinate system that is stationary with respect to the laboratory. We must assume that our test particle carries a sufficiently small charge, so that the electric field that it creates contributes negligibly to the field we are trying to detect (think of a single electron). Let the test particle be located at $\mathbf{r}$. Careful experiments show that when both charges are positive, the force experienced by the test particle is directly away from the charged body located at $\mathbf{r}_{0}$ and has magnitude proportional to $q e / r^{2}$, where $r=\left|\mathbf{r}-\mathbf{r}_{0}\right|$ is the distance between the charged body and the test particle, and where $q$ and $e$ are positive numbers which represent the amount of charge carried by the stationary body and the test particle respectively. If the units are chosen in an appropriate way, we can say that the force $F$ is given by

$$
F=q e \frac{\mathbf{r}-\mathbf{r}_{0}}{\left|\mathbf{r}-\mathbf{r}_{0}\right|^{3}}
$$

By definition, the electric field at the location $\mathbf{r}$ of the test particle is

$$
\mathbf{E}=q \frac{\mathbf{r}-\mathbf{r}_{0}}{\left|\mathbf{r}-\mathbf{r}_{0}\right|^{3}}
$$

If the test particle has charge opposite to that of the source body, then one of $q$ or $e$ will be negative and the force is directed toward the source. The test particle could have been placed anywhere in space, and so the electric field is implicitly defined at each point in space and so gives a vector field on $\mathbb{R}^{3}$. If the charge is modeled as a smoothly distributed charge density $\rho$ which is nonzero in some region $U \subset \mathbb{R}^{3}$, then the total charge is given by integration $Q=\int_{U} \rho(t, \mathbf{r}) d V_{\mathbf{r}}$ and the field at $\mathbf{r}$ is now given by $\mathbf{E}(t, \mathbf{r})=\int_{U} \rho(t, \mathbf{y}) \frac{\mathbf{r}-\mathbf{y}}{|\mathbf{r}-\mathbf{y}|^{3}} d V_{\mathbf{y}}$. Since the source particle is stationary at $\mathbf{r}_{0}$, the electric field will be independent of time $t$. A magnetic field is produced by circulating charge (a current). If charge $e$ is located at $\mathbf{r}$ and moving with velocity $\mathbf{v}$ in a magnetic field $\mathbf{B}(t, \mathbf{r})$, then the force felt by the charge is $\mathbf{F}=e \mathbf{E}+\frac{e}{c} \mathbf{v} \times \mathbf{B}$, where $\mathbf{v}$ is the velocity of the test particle. The test\\
particle has to be moving to feel the magnetic part of the field! At this point it is worth pointing out that from the point of view of spacetime, we are not staying true to the spirit of differential geometry since a vector field should have a geometric reality that is independent of its expression in a coordinate system. But a change of inertial frame can make $\mathbf{B}$ zero. Only by treating $\mathbf{E}$ and $\mathbf{B}$ together as aspects of a single field can we obtain the proper view.

Our next task is to write Maxwell's equations in terms of differential forms. We already have a way to convert (time dependent) vector fields $\mathbf{E}$ and $\mathbf{B}$ on $\mathbb{R}^{3}$ into (time dependent) differential forms on $\mathbb{R}^{3}$. Namely, we use the flatting operation with respect to the standard metric on $\mathbb{R}^{3}$. For the electric field we have

$$
\mathbf{E}=E^{x} \partial_{x}+E^{y} \partial_{y}+E^{z} \partial_{z} \mapsto \mathcal{E}=E_{x} d x+E_{y} d y+E_{z} d z
$$

For the magnetic field we do something a bit different. Namely, we flat and then apply the star operator. In rectangular coordinates, we have

$$
\mathbf{B}=B^{x} \partial_{x}+B^{y} \partial_{y}+B^{z} \partial_{z} \mapsto \mathcal{B}=B_{x} d y \wedge d z+B_{y} d z \wedge d x+B_{z} d x \wedge d y
$$

If we stick to rectangular coordinates (as we have been), the matrix of the standard metric is just $I=\left(\delta_{i j}\right)$, and so we see that the above operations do not numerically change the components of the fields. Thus in any rectangular coordinate system we have

$$
E^{x}=E_{x}, \quad E^{y}=E_{y}, \quad E_{z}=E^{z}
$$

and similarly for the $B$ 's. It is not hard to check that in the static case where $\mathcal{E}$ and $\mathcal{B}$ are time independent, the first pair of (static) Maxwell's equations are equivalent to

$$
d \mathcal{E}=0 \text { and } d \mathcal{B}=0
$$

This is nice, but if we put time dependence back into the picture, we need to do a couple more things to get a nice viewpoint. So assume now that $\mathbf{E}$ and $\mathbf{B}$ and hence the forms $\mathcal{E}$ and $\mathcal{B}$ are time dependent, and let us view these as differential forms on spacetime $\mathbb{R}_{1}^{4}$. In fact, let us combine $\mathcal{E}$ and $\mathcal{B}$ into a single 2-form on $\mathbb{R}_{1}^{4}$ by setting

$$
\mathcal{F}=\mathcal{B}+\mathcal{E} \wedge d t
$$

Since $\mathcal{F}$ is a 2 -form, it can be written in the form $\mathcal{F}=\frac{1}{2} F_{\mu \nu} d x^{\mu} \wedge d x^{\nu}$, where $F_{\mu \nu}=-F_{\nu \mu}$ and where the Greek indices are summed over $\{0,1,2,3\}$. It is traditional in physics to let the Greek indices run over this set and to let Latin indices run over just the "space indices" $1,2,3$. We will follow this convention for a while. If we compare $\mathcal{F}=\mathcal{B}+\mathcal{E} \wedge d t$ with $\frac{1}{2} F_{\mu \nu} d x^{\mu} \wedge d x^{\nu}$,\\
we see that the $F_{\mu \nu}$ form an antisymmetric matrix which is none other than

$$
\left[\begin{array}{cccc}
0 & -E_{x} & -E_{y} & -E_{z} \\
E_{x} & 0 & B_{z} & -B_{y} \\
E_{y} & -B_{z} & 0 & B_{x} \\
E_{z} & B_{y} & -B_{x} & 0
\end{array}\right]
$$

Our goal now is to show that the first pair of Maxwell's equations are equivalent to the single differential form equation

$$
d \mathcal{F}=0
$$

Let $N$ be an $n$-manifold and let $M=(a, b) \times N$ for some interval $(a, b)$. Let the coordinate on $(a, b)$ be $t=x^{0}$ (time). Let $\left(x^{1}, \ldots, x^{n}\right)$ be a coordinate system on $N$. With the usual abuse of notation, $\left(x^{0}, x^{1}, \ldots, x^{n}\right)$ is a coordinate system on $(a, b) \times N$. One can easily show that the local expression $d \omega=\partial_{\mu} f_{\mu_{1} \ldots \mu_{k}} \wedge d x^{\mu} \wedge d x^{\mu_{1}} \wedge \cdots \wedge d x^{\mu_{k}}$ for the exterior derivative of a form $\omega=f_{\mu_{1} \ldots \mu_{k}} d x^{\mu_{1}} \wedge \cdots \wedge d x^{\mu_{k}}$ can be written as

$$
\begin{aligned}
d \omega= & \sum_{i=1}^{3} \partial_{i} \omega_{\mu_{1} \ldots \mu_{k}} \wedge d x^{i} \wedge d x^{\mu_{1}} \wedge \cdots \wedge d x^{\mu_{k}} \\
& +\partial_{0} \omega_{\mu_{1} \ldots \mu_{k}} \wedge d x^{0} \wedge d x^{\mu_{1}} \wedge \cdots \wedge d x^{\mu_{k}}
\end{aligned}
$$

where the $\mu_{i}$ sum over $\{0,1,2, \ldots, n\}$. Thus we may consider the spatial part $d_{S}$ of the exterior derivative operator $d$ on $(a, b) \times S=M$. That is, we think of a given form $\omega$ on ( $a, b) \times S$ as a time dependent form on $N$ so that $d_{S} \omega$ is exactly the first term in the expression (9.13) above. Then we may write $d \omega=d_{S} \omega+d t \wedge \partial_{t} \omega$ as a compact version of the expression (9.13). The part $d_{S} \omega$ contains no $d t$ 's. By definition $\mathcal{F}=\mathcal{B}+\mathcal{E} \wedge d t$ on $\mathbb{R} \times \mathbb{R}^{3}=\mathbb{R}_{1}^{4}$, and so

$$
\begin{aligned}
d \mathcal{F} & =d \mathcal{B}+d(\mathcal{E} \wedge d t) \\
& =d_{S} \mathcal{B}+d t \wedge \partial_{t} \mathcal{B}+\left(d_{S} \mathcal{E}+d t \wedge \partial_{t} \mathcal{E}\right) \wedge d t \\
& =d_{S} \mathcal{B}+\left(\partial_{t} \mathcal{B}+d_{S} \mathcal{E}\right) \wedge d t
\end{aligned}
$$

The part $d_{S} \mathcal{B}$ is the spatial part and contains no $d t$ 's. It follows that $d \mathcal{F}$ is zero if and only if both $d_{S} \mathcal{B}$ and $\partial_{t} \mathcal{B}+d_{S} \mathcal{E}$ are zero. Unraveling the definitions shows that the pair of equations $d_{S} \mathcal{B}=0$ and $\partial_{t} \mathcal{B}+d_{S} \mathcal{E}=0$ (which we just showed to be equivalent to $d \mathcal{F}=0$ ) are Maxwell's first two equations disguised in a new notation. In summary, we have

$$
d \mathcal{F}=0 \quad \Longleftrightarrow \quad \begin{gathered}
d_{S} \mathcal{B}=0 \\
\partial_{t} \mathcal{B}+d_{S} \mathcal{E}=0
\end{gathered} \quad \Longleftrightarrow \quad \begin{gathered}
\nabla \cdot \mathbf{B}=0 \\
\nabla \times \mathbf{E}+\frac{\partial \mathbf{B}}{\partial t}=0
\end{gathered}
$$

Below we rewrite the last pair of Maxwell's equations, where the advantage of combining time and space together manifests itself to an even greater degree. Let us first pause to notice an interesting aspect of the\\
first pair. Suppose that the electric and magnetic fields were really all along most properly thought of as differential forms. Then we see that the equation $d \mathcal{F}=0$ has nothing to do with the metric on Minkowski space at all. In fact, if $\phi: \mathbb{R}_{1}^{4} \rightarrow \mathbb{R}_{1}^{4}$ is any diffeomorphism at all, we have $d \mathcal{F}=0$ if and only if $d\left(\phi^{*} \mathcal{F}\right)=0$, and so the truth of the equation $d \mathcal{F}=0$ is really a differential topological fact; a certain form $\mathcal{F}$ is closed. The metric structure of Minkowski space is irrelevant. The same will not be true for the second pair. Even if we start out with the form $\mathcal{F}$ on spacetime it will turn out that the metric will necessarily be implicit in the differential forms version of the second pair of Maxwell's equations. In fact, what we will show is that if we use the star operator for the Minkowski metric, then the second pair can be rewritten as the single equation $* d * \mathcal{F}=\mathcal{J}$, where $\mathcal{J}$ is formed from $\mathbf{j}=\left(j^{1}, j^{2}, j^{3}\right)$ and $\rho$ as follows: First we form the 4 -vector field $J=\rho \partial_{t}+j^{1} \partial_{x}+j^{2} \partial_{y}+j^{3} \partial_{z}$ (called the 4 -current) and then using the flatting operation we produce $\mathcal{J}=-\rho d t+j^{1} d x+j^{2} d y+j^{3} d z=J_{0} d t+J_{1} d x+J_{2} d y+J_{3} d z$, which is the covariant form of the 4 -current.

We will only outline the passage from $* d * \mathcal{F}=\mathcal{J}$ to the pair $\nabla \cdot \mathbf{E}=\varrho$ and $\nabla \times \mathbf{B}-\frac{\partial \mathbf{E}}{\partial t}=\mathbf{j}$. Let $*_{S}$ be the operator one gets by viewing differential forms on $\mathbb{R}_{1}^{4}$ as time dependent forms on $\mathbb{R}^{3}$ and then acting by the star operator with respect to the standard metric on $\mathbb{R}^{3}$. The first step is to verify that the pair $\nabla \cdot \mathbf{E}=\varrho$ and $\nabla \times \mathbf{B}-\frac{\partial \mathbf{E}}{\partial t}=\mathbf{j}$ is equivalent to the pair ${ }_{{ }_{*} S} d_{S} *_{S} \mathcal{E}=\varrho$ and $-\partial_{t} \mathcal{E}+*_{S} d_{S} *_{S} \mathcal{B}=\jmath$, where $\jmath:=j^{1} d x+j^{2} d y+j^{3} d z$ and $\mathcal{B}$ and $\mathcal{E}$ are as before. Next we verify that

$$
* \mathcal{F}=*_{S} \mathcal{E}-*_{S} \mathcal{B} \wedge d t
$$

So the next goal is to get from $* d * \mathcal{F}=* \mathcal{J}$ to the pair $*_{S} d_{S} *_{S} \mathcal{E}=\varrho$ and $-\partial_{t} \mathcal{E}+*_{S} d_{S} *_{S} \mathcal{B}=\jmath$. The following exercise finishes things off.

Exercise 9.60. Show that $* d * \mathcal{F}=-\partial_{t} \mathcal{E}-*_{S} d_{S} *_{S} \mathcal{E} \wedge d t+*_{S} d_{S} *_{S} \mathcal{B}$ and then use this and what we did above to show that $* d * \mathcal{F}=\mathcal{J}$ is equivalent to the pair $*_{S} d_{S} *_{S} \mathcal{E}=\varrho$ and $-\partial_{t} \mathcal{E}+*_{S} d_{S} *_{S} \mathcal{B}=\jmath$.

We have arrived at the following formulation of Maxwell's equations:

$$
\begin{aligned}
d \mathcal{F} & =0 \\
* d * \mathcal{F} & =\mathcal{J}
\end{aligned}
$$

If we just think of this as a pair of equations to be satisfied by a 2 -form $\mathcal{F}$ where the 1 -form $\mathcal{J}$ is given, then this last version of Maxwell's equations makes sense on any semi-Riemannian manifold. In fact, on a Lorentz manifold that can be written as $(a, b) \times S=M$ with the product metric $-d t \otimes d t \times g$, for some Riemannian metric $g$ on $S$, we can write $\mathcal{F}=\mathcal{B}+\mathcal{E} \wedge d t$,\\
which allows us to identify the electric and magnetic fields in their covariant form.

\subsection*{Surface Theory Redux}
Consider a surface $M \subset \mathbb{R}^{3}$. In what follows, we take advantage of the natural identifications of the tangent spaces of $\mathbb{R}^{3}$ with $\mathbb{R}^{3}$ itself. Let $e_{1}, e_{2}$ be an oriented orthonormal frame field defined on some open subset $U$ of $M$ and let $e_{3}=e_{1} \times e_{2}$. We can think of each $e_{1}, e_{2}$ and $e_{3}$ as vector fields along $U$. Using the identifications mentioned above, we also think of them as $\mathbb{R}^{3}$-valued 0 -forms.

Note that the identity map $\operatorname{id}_{\mathbb{R}^{3}}$ may be considered as an $\mathbb{R}^{3}$-valued 0 form on $\mathbb{R}^{3}$. Let $I:=\iota^{*}\left(\operatorname{id}_{\mathbb{R}^{3}}\right)$ where $\iota: U \subset M \hookrightarrow \mathbb{R}^{3}$ is inclusion. Then $d I:=\iota^{*} d\left(\mathrm{id}_{\mathbb{R}^{3}}\right)=0$ since $\mathrm{id}_{\mathbb{R}^{3}}$ is constant. Thus, $I$ is an $\mathbb{R}^{3}$-valued 0 -form, and we may write

$$
I=e_{1} \theta^{1}+e_{2} \theta^{2},
$$

where ( $\theta^{1}, \theta^{2}$ ) is the frame field dual to ( $e_{1}, e_{2}$ ). Note that $\operatorname{vol}_{M}=\theta^{1} \wedge \theta^{2}$. We have the $\mathbb{R}^{3}$-valued 1 -forms $d e_{j}$ and

$$
d e_{j}=\sum_{k=1}^{3} e_{k} \omega_{j}^{k}
$$

for some matrix of 1 -forms $\left(\omega_{j}^{i}\right)$. Note that if $v \in T_{p} M$, then $d e_{j}(v)=\bar{\nabla}_{v} e_{j}$, where $\bar{\nabla}$ is the flat Levi-Civita derivative on $\mathbb{R}^{3}$. (In fact, we should point out that if $\mathbf{i}, \mathbf{j}, \mathbf{k}$ is the standard basis on $\mathbb{R}^{3}$, then any field $X$ along $U$ may be written as $X=f_{1} \mathbf{i}+f_{2} \mathbf{j}+f_{3} \mathbf{k}$ for some smooth functions $f_{1}, f_{2}, f_{3}$ defined on $U$ and may be considered as an $\mathbb{R}^{3}$-valued 1-form. Then $d X(v)=\bar{\nabla}_{v} X= d f_{1}(v) \mathbf{i}+d f_{2}(v) \mathbf{j}+d f_{3}(v) \mathbf{k}$.) For an arbitrary tangent vector $v$ we have

$$
\omega_{j}^{i}(v)=\left\langle d e_{j}(v), e_{i}\right\rangle=-\left\langle e_{j}, d e_{i}(v)\right\rangle=-\omega_{i}^{j}(v),
$$

and so it follows that the matrix ( $\omega_{j}^{i}$ ) is antisymmetric. We write $\mathbb{R}^{3}$-valued forms on $U$ as $\sum_{k=1}^{3} e_{k} \eta^{k}$ where $\eta^{k} \in \Omega(U)$. This is to conform with the order of matrix multiplication.

Theorem 9.61. Let $M \subset \mathbb{R}^{3}, e_{1}, e_{2}, \theta^{1}, \theta^{2}$ and $\left(\omega_{j}^{i}\right)$ be as above. Then the following structure equations hold:

$$
\begin{aligned}
& d \theta^{i}=-\sum_{k=1}^{2} \omega_{k}^{i} \wedge \theta^{k} \\
& \omega_{1}^{3} \wedge \theta^{1}+\omega_{2}^{3} \wedge \theta^{2}=0 \\
& d \omega_{j}^{i}=-\sum_{k=1}^{3} \omega_{k}^{i} \wedge \omega_{j}^{k}
\end{aligned}
$$

Proof. We calculate as follows:

$$
\begin{aligned}
0 & =d I=d \sum_{k=1}^{2} e_{k} \theta^{k}=\sum_{k=1}^{2} d e_{k} \wedge \theta^{k}+e_{k} \wedge d \theta^{k} \\
& =\sum_{k=1}^{2}\left(\sum_{j=1}^{3} \omega_{k}^{j} e_{j}\right) \wedge \theta^{k}+\sum_{k=1}^{2} e_{k} \wedge d \theta^{k} \\
& =\sum_{k=1}^{2}\left(\sum_{j=1}^{2} e_{j} \omega_{k}^{j}+e_{3} \omega_{k}^{3}\right) \wedge \theta^{k}+\sum_{j=1}^{2} e_{k} \wedge d \theta^{k} \\
& =\sum_{j=1}^{2} \sum_{k=1}^{2} e_{j} \omega_{k}^{j} \wedge \theta^{k}+\sum_{j=1}^{2} e_{j} \wedge d \theta^{j}+\sum_{k=1}^{2} e_{3} \omega_{k}^{3} \wedge \theta^{k} \\
& =\sum_{j=1}^{2} e_{j}\left(d \theta^{j}+\sum_{k=1}^{2} \omega_{k}^{j} \wedge \theta^{k}\right)+e_{3}\left(\sum_{k=1}^{2} \omega_{k}^{3} \wedge \theta^{k}\right)
\end{aligned}
$$

and it follows that $d \theta^{i}=-\sum_{k=1}^{2} \omega_{k}^{i} \wedge \theta^{k}$ and $\omega_{1}^{3} \wedge \theta^{1}+\omega_{2}^{3} \wedge \theta^{2}=0$.\\
Also,

$$
\begin{aligned}
0 & =d d e_{j}=d \sum_{k=1}^{3} e_{k} \omega_{j}^{k} \\
& =\sum_{k=1}^{3} d e_{k} \wedge \omega_{j}^{k}+\sum_{k=1}^{3} e_{k} \wedge d \omega_{j}^{k} \\
& =\sum_{k=1}^{3}\left(\sum_{i=1}^{3} e_{i} \omega_{k}^{i}\right) \wedge \omega_{j}^{k}+\sum_{i=1}^{3} e_{i} \wedge d \omega_{j}^{i} \\
& =\sum_{i=1}^{3} e_{i}\left(\sum_{k=1}^{3} \omega_{k}^{i} \wedge \omega_{j}^{k}+d \omega_{j}^{i}\right)
\end{aligned}
$$

and so $d \omega_{j}^{i}=-\sum_{k=1}^{3} \omega_{k}^{i} \wedge \omega_{j}^{k}$.\\
Notice that for $v \in T_{p} M$ we have $d e_{3}(v)=\bar{\nabla}_{v} e_{3}=-S(v)$ and so $I I(v, w)=\left\langle-d e_{3}(v), w\right\rangle$ (recall Definition 4.20). Therefore,

$$
I I\left(e_{i}, e_{j}\right)=\left\langle-d e_{3}\left(e_{i}\right), e_{j}\right\rangle=\left\langle-\sum_{k=1}^{3} e_{k} \omega_{3}^{k}\left(e_{i}\right), e_{j}\right\rangle=-\omega_{3}^{j}\left(e_{i}\right)
$$

and so

$$
\left[\begin{array}{cc}
I I\left(e_{1}, e_{1}\right) & I I\left(e_{1}, e_{2}\right) \\
I I\left(e_{2}, e_{1},\right) & I I\left(e_{2}, e_{2}\right)
\end{array}\right]=-\left[\begin{array}{cc}
\omega_{3}^{1}\left(e_{1}\right) & \omega_{3}^{2}\left(e_{1}\right) \\
\omega_{3}^{1}\left(e_{2}\right) & \omega_{3}^{2}\left(e_{2}\right)
\end{array}\right] .
$$

Since ( $e_{1}, e_{2}$ ) is an orthonormal frame field, the matrix of the first fundamental form is the identity matrix. It follows that the matrix which represents\\
the shape operator is the same as that which represents the second fundamental form. Thus

$$
K=\operatorname{det}\left[\begin{array}{cc}
I I\left(e_{1}, e_{1}\right) & I I\left(e_{1}, e_{2}\right) \\
I I\left(e_{2}, e_{1},\right) & I I\left(e_{2}, e_{2}\right)
\end{array}\right]=-\operatorname{det}\left[\begin{array}{cc}
\omega_{3}^{1}\left(e_{1}\right) & \omega_{3}^{2}\left(e_{1}\right) \\
\omega_{3}^{1}\left(e_{2}\right) & \omega_{3}^{2}\left(e_{2}\right)
\end{array}\right] .
$$

Notice that $d e_{3}=\sum_{k=1}^{3} e_{k} \omega_{3}^{k}$ reduces to $d e_{3}=e_{1} \omega_{3}^{1}+e_{2} \omega_{3}^{2}$ since $\omega_{3}^{3}=0$. We have

$$
\begin{aligned}
& \omega_{3}^{1}=\omega_{3}^{1}\left(e_{1}\right) \theta^{1}+\omega_{3}^{1}\left(e_{2}\right) \theta^{2}, \\
& \omega_{3}^{2}=\omega_{3}^{2}\left(e_{1}\right) \theta^{1}+\omega_{3}^{2}\left(e_{2}\right) \theta^{2},
\end{aligned}
$$

so that

$$
\begin{aligned}
\omega_{3}^{1} \wedge \omega_{3}^{2} & =-K \theta^{1} \wedge \theta^{2} \\
\omega_{3}^{1} \wedge \omega_{2}^{3} & =K \theta^{1} \wedge \theta^{2} \\
d \omega_{2}^{1} & =K \theta^{1} \wedge \theta^{2}
\end{aligned}
$$

Recall that for $v, w \in T_{x} S^{2}$ we have $\operatorname{vol}_{S^{2}}(v, w)=\langle v \times w, x\rangle$. If we consider $e_{3}$ as a map $e_{3}: U \subset M \rightarrow S^{2}$, then it is called the Gauss map. If $M$ is assumed oriented, then we may take $e_{3}$ to be equal to a global normal field. If $v, w \in T_{p} M$, then we have

$$
\begin{aligned}
e_{3}^{*} \operatorname{vol}_{S^{2}}(v, w) & =\left\langle d e_{3}(v) \times d e_{3}(w), e_{3}\right\rangle \\
& =\left\langle\left(\omega_{3}^{1}(v) e_{1}+\omega_{3}^{2}(v) e_{2}\right) \times\left(\omega_{3}^{1}(w) e_{1}+\omega_{3}^{2}(w) e_{2}\right), e_{3}\right\rangle \\
& =\omega_{3}^{1} \wedge \omega_{2}^{3}(v, w)=K \theta^{1} \wedge \theta^{2}(v, w)=K \operatorname{vol}_{M}(v, w)
\end{aligned}
$$

Thus we obtain

$$
e_{3}^{*} \operatorname{vol}_{S^{2}}=K \operatorname{vol}_{M}
$$

This shows that the Gauss curvature is a measure of distortion of the signed volume under the Gauss map. In particular, if $e_{3}: U \subset M \rightarrow S^{2}$ is orientation reversing at $p$, then the curvature at $p$ is negative. Let $p \in M$ with $p$ in the domain $U$ of $e_{3}$. If $A$ is a nice domain in $U$, then $e_{3}$ maps $A$ to a set in $S^{2}$. Without worrying about measure-theoretic technicalities, we have

$$
\begin{aligned}
\operatorname{vol}\left(e_{3}(A)\right) & =\int_{e_{3}(A)} \operatorname{vol}_{S^{2}} \\
& =\int_{A} e_{3}^{*} \operatorname{vol}_{S^{2}}=\int_{A} K \operatorname{vol}_{M}
\end{aligned}
$$

One can get an idea of the curvature near a point on a surface by visualizing the Gauss map (see Figure 9.2). For example, it is clear that the right circular cylinder has Gauss curvature zero since it maps every region of the cylinder onto a set with zero area. It is also fairly clear that the saddle surface in the diagram has negative curvature.

\begin{figure}[h]
\begin{center}
  \includegraphics[scale=0.25]{2025_10_09_265d7e1a22c89f79c21eg-452}

\caption{Figure 9.2. Gauss map}
\end{center}
\end{figure}

\section*{Problems}
(1) Let $\iota: S^{2} \hookrightarrow \mathbb{R}^{3} \backslash\{0\}$ be the inclusion map. Let $\tau=\iota^{*} \omega$ where

$$
\omega=\frac{x d y \wedge d z+y d z \wedge d x+z d x \wedge d y}{\left(x^{2}+y^{2}+z^{2}\right)^{3 / 2}} .
$$

Compute $\int_{S^{2}} \tau$ where $S^{2}$ is given the orientation induced by $\tau$ itself.\\
(2) Let $M$ be an oriented smooth compact manifold with boundary $\partial M$ and suppose that $\partial M$ has two connected components $N_{0}$ and $N_{1}$. Let $\imath_{i}: N_{i} \hookrightarrow M$ be the inclusion map for $i=0,1$. Suppose that $\alpha$ is a $p$-form with $\imath_{0}^{*} \alpha=0$ and $\beta$ an ( $n-p-1$ )-form with $\imath_{1}^{*} \beta=0$. Prove that in this case

$$
\int_{M} d \alpha \wedge \beta=(-1)^{p+1} \int_{M} \alpha \wedge d \beta
$$

(3) Consider the set up in the proof of Theorem 9.22 where $\gamma \in \Lambda^{n-k}$ and $L_{\gamma}: \bigwedge^{k} \mathrm{V}^{*} \rightarrow \mathbb{R}$ is defined so that $L_{\gamma}(\alpha)$ vol $=\alpha \wedge \gamma$. Show that $\gamma \mapsto L_{\gamma} \in\left(\bigwedge^{k} \mathrm{V}^{*}\right)^{*}$ is linear. Show that if $L_{\gamma}(\alpha)=0$ for all $\alpha \in \bigwedge^{k} \mathrm{V}^{*}$, then $\gamma=0$. Thus $\gamma \mapsto L_{\gamma}$ is injective (and hence an isomorphism). See the related Problem 4.\\
(4) Recall the notion of a manifold with corners as in Problem 21 in Chapter 1. Define orientation on manifolds with corners. Let $M$ be an $n$ manifold with corners. Show that if $C M$ is the set of corner points of $M$, then $M \backslash C M$ is a manifold with boundary. Develop integration theory and Stokes' theorem for manifolds with corners. [Hint: If $\omega$ is an ( $n-1$ )-form with compact support in the domain of a chart ( $U, \mathrm{x}$ ), then define

$$
\int_{\partial M} \omega:=\sum_{i=1}^{n} \int_{F_{i}}\left(\mathrm{x}^{-1}\right)^{*} \omega,
$$

where

$$
F_{i}:=\left\{x \in \overline{\mathbb{R}_{+}^{n}}: x^{i}=0\right\}
$$

is given the induced orientation as a subset of the boundary of $\left\{x \in \mathbb{R}^{n}\right.$ : $\left.x^{i} \geq 0\right\}$.]\\
(5) Show that if $f:(M, g) \rightarrow(N, h)$ is an isometry or a local isometry and $\omega$ is a $k$-form on $N$, then $f^{*} \omega$ is harmonic if and only if $\omega$ is harmonic.\\
(6) Let $M$ be a connected oriented compact Riemannian manifold with Laplace operator $\Delta$. A smooth nonzero function $f$ is called an eigenfunction for $\Delta$ with eigenvalue $\lambda$ if $\Delta f=\lambda f$.\\
(a) Show that zero is an eigenvalue and that all other eigenvalues are strictly positive.\\
(b) Show that if $\Delta f_{1}=\lambda_{1} f_{1}$ and $\Delta f_{2}=\lambda_{2} f_{2}$ for $\lambda_{1} \neq \lambda_{2}$, then $\left(f_{1} \mid f_{2}\right)=\int_{M} f_{1} f_{2} \operatorname{vol}_{M}=0$.\\
(a) Show that if $\wp: \widetilde{M} \rightarrow M$ is a smooth covering space of multiplicity $m$, then for any compactly supported $\omega \in \Omega^{n}(M)$ we have

$$
\int_{\widetilde{M}} \wp^{*} \omega=m \int_{M} \omega .
$$

(7) Let $M \subset \mathbb{R}^{n+1}$ be an oriented hypersurface. If $N$ is a positively oriented normal field along $M$ with $N=\sum N^{i} \partial / \partial x^{i}$, then the following formula gives the volume form corresponding to the induced metric on $M$ :

$$
\operatorname{vol}_{M}:=\sum_{i}(-1)^{i-1} N^{i} d x^{1} \wedge \cdots \wedge \widehat{d x^{i}} \wedge \cdots \wedge d x^{n+1}
$$

where $x^{1}, \ldots, x^{n+1}$ are the standard coordinates on $\mathbb{R}^{n+1}$ and the $d x^{i}$ are restricted to $M$.\\
(8) Let $U$ be a starshaped open set in $\mathbb{R}^{n}$ and let $x^{1}, \ldots, x^{n}$ be standard coordinates. Given

$$
\omega=\sum_{i_{1}<\cdots<i_{k}} \omega_{i_{1}, \ldots, i_{k}} d x^{i_{1}} \wedge \cdots \wedge d x^{i_{k}}
$$

define

$$
\begin{aligned}
I \omega:=\sum_{i_{1}<\cdots<i_{k}} \sum_{\alpha=1}^{k}( & -1)^{\alpha-1}\left(\int_{0}^{1} t^{k-1} \omega_{i_{1}, \ldots, i_{k}}(t x) d t\right) \\
& \times x^{i_{\alpha}} d x^{i_{1}} \wedge \cdots \wedge \widehat{d x^{i_{\alpha}}} \wedge \cdots \wedge d x^{i_{k}}
\end{aligned}
$$

where the caret means omission. Show that if $d \omega=0$, then $d I \omega=0$.\\
(9) Derive the decompositions

$$
\Omega^{k}(M)=d \delta\left(\Omega^{k}(M)\right) \oplus \delta d\left(\Omega^{k}(M)\right) \oplus \mathcal{H}^{k}
$$

and

$$
\Omega^{k}(M)=d\left(\Omega^{k-1}(M)\right) \oplus \delta\left(\Omega^{k+1}(M)\right) \oplus \mathcal{H}^{k}
$$

from

$$
\Omega^{k}(M)=\Delta\left(\Omega^{k}(M)\right) \oplus \mathcal{H}^{k}
$$

(10) Prove Theorem 9.7.

\section*{De Rham Cohomology}
We have already given the definition of de Rham cohomology (8.41) and the de Rham theorem which says that for a compact oriented smooth manifold $M$, the de Rham cohomology is isomorphic with the singular cohomology. In this chapter we introduce just a few of the most basic tools proper to the subject such as the Mayer-Vietoris sequence. We shall also introduce the compactly supported de Rham cohomology and revisit Poincaré duality. In this chapter all manifolds will be without boundary.

For a given $n$-manifold $M$, we have the sequence of maps called the de Rham complex

$$
0 \xrightarrow{d} \Omega^{0}(M) \xrightarrow{d} \Omega^{1}(M) \xrightarrow{d} \cdots \xrightarrow{d} \Omega^{n}(M) \xrightarrow{d} 0,
$$

and we have defined the de Rham cohomology groups (actually vector spaces) as the quotients

$$
H^{k}(M)=\frac{Z^{k}(M)}{B^{k}(M)}
$$

where $Z^{k}(M):=\operatorname{Ker}\left(d: \Omega^{k}(M) \rightarrow \Omega^{k+1}(M)\right)$ and $B^{k}(M):=\operatorname{Im}(d: \left.\Omega^{k-1}(M) \rightarrow \Omega^{k}(M)\right)$. The elements of $Z^{k}(M)$ are called closed $k$-forms or $k$-cocycles, and the elements of $B^{k}(M)$ are called exact $k$-forms or $k$ coboundaries. If $f: M \rightarrow N$ is smooth and $[\beta] \in H^{k}(N)$, then since $d$ commutes with pull-back, it is easy to see that $f^{*} \beta$ is closed because $\beta$ is closed. Thus we obtain a cohomology class $\left[f^{*} \beta\right]$. Also, $\left[f^{*} \beta\right]$ depends only on the equivalence class of $\beta$. Indeed, if $\beta-\beta^{\prime}=d \eta$, then $f^{*} \beta-f^{*} \beta^{\prime}=d f^{*} \eta$. Thus we may define a linear map $f^{*}: H^{k}(N) \rightarrow H^{k}(M)$ by $f^{*}[\beta]:=\left[f^{*} \beta\right]$.

Let us immediately review a simple situation from Section 2.10, which will help the reader better see what the de Rham cohomologies are all about.

Let $M=\mathbb{R}^{2} \backslash\{0\}$ and consider the 1-form

$$
\vartheta:=\frac{x d y-y d x}{x^{2}+y^{2}}
$$

We got this 1 -form by taking the exterior derivative of $\theta=\arctan (y / x)$. This function is not defined as a single-valued smooth function on all of $\mathbb{R}^{2} \backslash\{0\}$, but $\vartheta$ is well-defined on all of $\mathbb{R}^{2} \backslash\{0\}$. One can extend $\theta$ to be defined on the plane minus the ray $\{x \leq 0, y=0\}$ as follows:

$$
\theta= \begin{cases}\arctan (y / x) & \text { if } x>0 \\ \arccos \left(x / \sqrt{x^{2}+y^{2}}\right) & \text { if } x \leq 0 \text { and } y>0 \\ -\arccos \left(x / \sqrt{x^{2}+y^{2}}\right) & \text { if } x \leq 0 \text { and } y<0\end{cases}
$$

See Problem 1. However, this is still not defined on all of $\mathbb{R}^{2} \backslash\{0\}$. One may also check that $d \vartheta=0$ and so $\vartheta$ is closed. Using Proposition 2.125 and Example 2.126 of Section 2.10, it is easy to see that we have the following situation:\\
(1) $\vartheta:=\frac{x d y-y d x}{x^{2}+y^{2}}$ is a smooth 1 -form on $\mathbb{R}^{2} \backslash\{0\}$ with $d \vartheta=0$.\\
(2) There is no function $f$ defined on (all of) $\mathbb{R}^{2} \backslash\{0\}$ such that $\vartheta=d f$.\\
(3) From Section 2.10 we know that for any ball $B(p, \varepsilon)$ in $\mathbb{R}^{2} \backslash\{0\}$ there is a function $f \in C^{\infty}(B(p, \varepsilon))$ such that $\left.\vartheta\right|_{B(p, \varepsilon)}=d f$.\\
Assertion (1) says that $\vartheta$ is globally well-defined and closed, while (2) says that $\vartheta$ is not exact. Assertion (3) says that $\vartheta$ is what we might call locally exact or locally conservative. What prevents us from finding a (global) function with $\vartheta=d f$ ? Could the same kind of situation occur if the manifold is $\mathbb{R}^{2}$ ? The answer is no, and the difference between $\mathbb{R}^{2}$ and $\mathbb{R}^{2} \backslash\{0\}$ is that $H^{1}\left(\mathbb{R}^{2}\right)=0$ while $H^{1}\left(\mathbb{R}^{2} \backslash\{0\}\right) \neq 0$.

Exercise 10.1. Verify (1) and (3) above.\\
We recall from Chapter 2 that this example has something to do with path independence. In fact, if we could show that for a given 1 -form $\alpha$, the path integral $\int_{c} \alpha$ only depended on the beginning and ending points of the curve $c$, then we could define a function $f(x):=\int_{x_{0}}^{x} \alpha$, where $\int_{x_{0}}^{x} \alpha$ is just the path integral for any path beginning at a fixed $x_{0}$ and ending at $x$. With this definition one can show that $d f=\alpha$ and so $\alpha$ would be exact. In our example, the form $\vartheta$ is not exact, and so there must be a failure of path independence.

Exercise 10.2. A smooth fixed endpoint homotopy between a path $c_{0}$ : $[a, b] \rightarrow M$ and $c_{1}:[a, b] \rightarrow M$ is a one parameter family of paths $h_{s}$ such that $h_{0}=c_{0}$ and $h_{1}=c_{1}$ and such that the map $H(t, s):=h_{s}(t)$ is smooth on $[a, b] \times[0,1]$. Show that if $\alpha$ is an exact 1-form, then $\frac{d}{d s} \int_{h_{s}} \alpha=0$.

The de Rham complex and its cohomology can be viewed in terms of differential equations. For example, the task of finding a closed 1 -form $f d x+g d y$ on an open set $U \subset \mathbb{R}^{2}$ amounts to finding a solution of the differential equation

$$
\frac{\partial g}{\partial x}-\frac{\partial f}{\partial y}=0 .
$$

The "trivial" solutions are the exact forms since they are automatically closed. Thus the de Rham cohomology $H^{1}(U)$ is a space of solutions modulo the "uninteresting" solutions. Similar statements apply to the higher cohomology groups.

Since we have found a closed 1 -form on $\mathbb{R}^{2} \backslash\{0\}$ that is not exact, we know that $H^{1}\left(\mathbb{R}^{2} \backslash\{0\}\right) \neq 0$. We are not yet in a position to determine $H^{1}\left(\mathbb{R}^{2} \backslash\{0\}\right)$ completely. We will start out with even simpler spaces and eventually develop the machinery to bootstrap our way up to more complicated situations.

First, let $M=\{p\}$. That is, $M$ consists of a single point and is hence a 0-dimensional manifold. In this case,

$$
\Omega^{k}(\{p\})=Z^{k}(\{p\})= \begin{cases}\mathbb{R} & \text { for } k=0, \\ 0 & \text { for } k>0 .\end{cases}
$$

Furthermore, $B^{k}(\{p\})=0$ and so

$$
H^{k}(\{p\})= \begin{cases}\mathbb{R} & \text { for } k=0 \\ 0 & \text { for } k>0\end{cases}
$$

Next we consider the case $M=\mathbb{R}$. Here, $Z^{0}(\mathbb{R})$ is clearly just the constant functions and so is (isomorphic to) $\mathbb{R}$. On the other hand, $B^{0}(\mathbb{R})=$ 0 and so

$$
H^{0}(\mathbb{R})=\mathbb{R}
$$

Now since $d: \Omega^{1}(\mathbb{R}) \rightarrow \Omega^{2}(\mathbb{R})=0$, we see that $Z^{1}(\mathbb{R})=\Omega^{1}(\mathbb{R})$. If $g(x) d x \in \Omega^{1}(\mathbb{R})$, then letting

$$
f(x):=\int_{0}^{x} g(x) d x
$$

we get $d f=g(x) d x$. Thus, every $\Omega^{1}(\mathbb{R})$ is exact; $B^{1}(\mathbb{R})=\Omega^{1}(\mathbb{R})$. We are led to

$$
H^{1}(\mathbb{R})=0 .
$$

From this modest beginning we will be able to compute the de Rham cohomology for a large class of manifolds. Our first goal is to compute $H^{k}\left(\mathbb{R}^{n}\right)$ for all $k$. In order to accomplish this, we will need some preparation. The methods are largely algebraic, and so we will need to introduce a bit of "homological algebra".

Definition 10.3. Let $R$ be a commutative ring. A differential $R$-complex is a direct sum of modules $C=\bigoplus_{k \in \mathbb{Z}} C^{k}$ together with a linear map $d: C \rightarrow C$ such that $d \circ d=0$ and such that $d\left(C^{k}\right) \subset C^{k+1}$. Thus we have a sequence of linear maps

$$
\cdots \rightarrow C^{k-1} \xrightarrow{d} C^{k} \xrightarrow{d} C^{k+1} \rightarrow \cdots
$$

where we have denoted the restrictions $d_{k}=d \mid C^{k}$ all simply by the single letter $d$.

Let $A=\bigoplus_{k \in \mathbb{Z}} A^{k}$ and $B=\bigoplus_{k \in \mathbb{Z}} B^{k}$ be differential complexes. A map $f: A \rightarrow B$ is called a chain map if $f$ is a (degree 0 ) graded map such that $d \circ f=f \circ d$. In other words, if we let $f \mid A^{k}:=f_{k}$, then we require that $f_{k}\left(A^{k}\right) \subset B^{k}$ and that the following diagram commutes for all $k$ :\\
\includegraphics[scale=0.25, center]{2025_10_09_265d7e1a22c89f79c21eg-459}

Notice that if $f: A \rightarrow B$ is a chain map, then $\operatorname{Ker}(f)$ and $\operatorname{Im}(f)$ are complexes with $\operatorname{Ker}(f)=\bigoplus_{k \in \mathbb{Z}} \operatorname{Ker}\left(f_{k}\right)$ and $\operatorname{Im}(f)=\bigoplus_{k \in \mathbb{Z}} \operatorname{Im}\left(f_{k}\right)$. Thus the notion of exact sequence of chain maps may be defined in the obvious way.

Definition 10.4. The $k$-th cohomology of the complex $C=\bigoplus_{k \in \mathbb{Z}} C^{k}$ is

$$
H^{k}(C):=\frac{\operatorname{Ker}\left(d \mid C^{k}\right)}{\operatorname{Im}\left(d \mid C^{k-1}\right)}
$$

The elements of $\operatorname{Ker}\left(d \mid C^{k}\right)$ (also denoted $Z^{k}(C)$ ) are called $k$-cocycles, while the elements of $\operatorname{Im}\left(d \mid C^{k-1}\right)$ (also denoted $B^{k}(C)$ ) are called $k$-coboundaries.

If $f: A \rightarrow B$ is a chain map, then it is easy to see that there is a natural (degree 0 ) graded map $f_{*}: H(A) \rightarrow H(B)$ defined by $f_{*}([x]):=[f(x)]$ for $x \in A^{k}$.

Definition 10.5. An exact sequence of chain maps of the form

$$
0 \rightarrow A \xrightarrow{f} B \xrightarrow{g} C \rightarrow 0
$$

is called a short exact sequence.

Associated to every short exact sequence of chain maps is a long exact sequence of cohomology groups:\\
\includegraphics[scale=0.25, center]{2025_10_09_265d7e1a22c89f79c21eg-460(1)}

The maps $f_{*}$ and $g_{*}$ are the maps induced by $f$ and $g$, and the "coboundary map" or connecting homomorphism $d^{*}: H^{k}(C) \rightarrow H^{k+1}(A)$ is defined as follows: Let $c \in Z^{k}(C) \subset C^{k}$ represent the class $[c] \in H^{k}(C)$ so that $d c=0$. Starting with this we hope to end up with a well-defined element of $H^{k+1}(A)$, but we must refer to the following diagram with exact rows to explain how we arrive at a choice of representative of our class $d^{*}([c])$ :\\
\includegraphics[scale=0.25, center]{2025_10_09_265d7e1a22c89f79c21eg-460}

By the surjectivity of $g$ there is a $b \in B^{k}$ with $g(b)=c$. Also, since $g(d b)= d(g(b))=d c=0$, it must be that $d b=f(a)$ for some $a \in A^{k+1}$. The scheme of the process is

$$
c-\rightarrow b-\rightarrow a .
$$

Certainly $f(d a)=d(f(a))=d d b=0$, and so since $f$ is 1-1, we must have $d a=0$, which means that $a \in Z^{k+1}(A)$. We would like to define $d^{*}([c])$ to be $[a]$, but we must show that this is well-defined. Suppose that we repeat this process starting with $c^{\prime}=c+d \alpha$ for some $\alpha \in C^{k-1}$. In our first step, we find $b^{\prime} \in B^{k}$ with $g\left(b^{\prime}\right)=c^{\prime}$ and then $a^{\prime}$ with $f\left(a^{\prime}\right)=d b^{\prime}$. We wish to show that $[a]=\left[a^{\prime}\right]$. We have $g\left(b-b^{\prime}\right)=c-c^{\prime}=d \alpha$. But there must be an element $\beta \in B^{k-1}$ such that $g(\beta)=\alpha$. Now we have

$$
\begin{aligned}
g\left(b-b^{\prime}-d \beta\right) & =g(b)-g\left(b^{\prime}\right)-g(d \beta) \\
& =g(b)-g\left(b^{\prime}\right)-d g(\beta) \\
& =c-c^{\prime}-d \alpha=0 .
\end{aligned}
$$

By exactness at $B^{k}$, there must be a $\gamma \in A^{k}$ such that $f(\gamma)=b-b^{\prime}-d \beta$. Now we have

$$
\begin{aligned}
f\left(a-a^{\prime}-d \gamma\right) & =f(a)-f\left(a^{\prime}\right)-d f(\gamma) \\
& =d b-d b^{\prime}-d\left(b-b^{\prime}-d \beta\right)=0,
\end{aligned}
$$

and since $f$ is injective, we have $a-a^{\prime}-d \gamma=0$, which means that $[a]=\left[a^{\prime}\right]$. Thus our definition $d^{*}([c]):=[a]$ is independent of the choices. We leave it to the reader to check (if there is any doubt) that $d^{*}$, so defined, is linear.

Let us review the situation we have with de Rham cohomology noticing how it fits the abstract homological algebra above. We let $\Omega^{k}(M):=0$ for $k<0$ and we have a differential complex $d: \Omega(M) \rightarrow \Omega(M)$, where $d$ is the exterior derivative and $\Omega(M)$ is the direct sum of the $\Omega^{k}(M)$. In this case, $H^{k}(\Omega(M))=H^{k}(M)$ by definition. If $f: M \rightarrow N$ is a $C^{\infty}$ map, then we have $f^{*}: \Omega(N) \rightarrow \Omega(M)$. Since pull-back commutes with exterior differentiation $d$ and preserves the degree of differential forms, $f^{*}$ is a chain map. Thus we have the induced map on cohomology denoted by $f^{*}$ so that

$$
\begin{aligned}
& f^{*}: H^{*}(N) \rightarrow H^{*}(M), \\
& f^{*}:[\alpha] \mapsto\left[f^{*} \alpha\right],
\end{aligned}
$$

where we have used $H^{*}(M)$ to denote the direct sum $\bigoplus_{i} H^{i}(M)$. Notice that $f \mapsto f^{*}$ together with $M \mapsto H^{*}(M)$ is a contravariant functor, which means that for $f: M \rightarrow N$ and $g: N \rightarrow P$ we have

$$
(g \circ f)^{*}=f^{*} \circ g^{*}
$$

(This is why we have now put the stars in as superscripts.) In particular, if $\iota_{U}: U \rightarrow M$ is the inclusion of an open set $U$ in $M$, then $\iota_{U}^{*} \alpha$ is the same as the restriction of the form $\alpha$ to $U$. If $[\alpha] \in H^{*}(M)$, then $\iota_{U}{ }^{*}([\alpha]) \in H^{*}(U)$;

$$
\iota_{U}{ }^{*}: H^{*}(M) \rightarrow H^{*}(U) .
$$

Remark 10.6. If $\alpha \in Z^{k}(M)$ and $\beta \in Z^{l}(M)$, then for any $\alpha^{\prime} \in \Omega^{k-1}(M)$ and any $\beta^{\prime} \in \Omega^{l-1}(M)$ we have

$$
\begin{aligned}
\left(\alpha+d \alpha^{\prime}\right) \wedge \beta & =\alpha \wedge \beta+d \alpha^{\prime} \wedge \beta \\
& =\alpha \wedge \beta+d\left(\alpha^{\prime} \wedge \beta\right)-(-1)^{k-1} \alpha^{\prime} \wedge d \beta \\
& =\alpha \wedge \beta+d\left(\alpha^{\prime} \wedge \beta\right)
\end{aligned}
$$

and similarly $\alpha \wedge\left(\beta+d \beta^{\prime}\right)=\alpha \wedge \beta+d\left(\alpha \wedge \beta^{\prime}\right)$. Thus we may define a product $H^{k}(M) \times H^{l}(M) \rightarrow H^{k+l}(M)$ by $[\alpha] \wedge[\beta]:=[\alpha \wedge \beta]$. This gives $H^{*}(M)$ a graded ring structure.

\subsection*{The Mayer-Vietoris Sequence}
Suppose that $M=U \cup V$ for open sets $U$ and $V$. We have the following commutative diagram of inclusions:\\
\includegraphics[scale=0.25, center]{2025_10_09_265d7e1a22c89f79c21eg-462}

This gives rise to the Mayer-Vietoris short exact sequence

$$
0 \rightarrow \Omega(M) \xrightarrow{j^{*}} \Omega(U) \oplus \Omega(V) \xrightarrow{\partial i^{*}} \Omega(U \cap V) \rightarrow 0,
$$

where

$$
j^{*}(\omega):=\left(j_{1}^{*} \omega, j_{2}^{*} \omega\right)
$$

and

$$
\partial i^{*}(\alpha, \beta):=i_{2}^{*}(\beta)-i_{1}^{*}(\alpha)=\left.\beta\right|_{U \cap V}-\left.\alpha\right|_{U \cap V}
$$

Note that $j_{1}^{*} \omega \in \Omega(U)$ while $j_{2}^{*} \omega \in \Omega(V)$. Note also that $i_{1}^{*}(\alpha)=\left.\alpha\right|_{U \cap V}$ and $i_{1}^{*}(\beta)=\left.\beta\right|_{U \cap V}$ live in $\Omega(U \cap V)$. Let us write $j^{*}$ suggestively as $\left(j_{1}^{*}, j_{2}^{*}\right)$.

Let us show that this sequence is exact. First if $\left(j_{1}^{*}, j_{2}^{*}\right)(\omega):=\left(j_{1}^{*} \omega, j_{2}^{*} \omega\right) =(0,0)$, then $\left.\omega\right|_{U}=\left.\omega\right|_{V}=0$ and so $\omega=0$ on $M=U \cup V$. Thus $\left(j_{1}^{*}, j_{2}^{*}\right)$ is 1-1 and the exactness at $\Omega(M)$ is demonstrated. Next, if $\eta \in \Omega(U \cap V)$, then we take a smooth partition of unity $\left\{\rho_{U}, \rho_{V}\right\}$ subordinate to the cover $\{U, V\}$ and let $\xi:=\left(-\left(\rho_{V} \eta\right)^{U},\left(\rho_{U} \eta\right)^{V}\right)$, where $\left(\rho_{V} \eta\right)^{U}$ is the extension of the form $\left.\rho_{V}\right|_{U \cap V} \eta$ by zero to $U$, and $\left(\rho_{U} \eta\right)^{V}$ likewise extends $\left.\rho_{U}\right|_{U \cap V} \eta$ to $V$. This may look backwards at first, so note carefully that we use $\rho_{V}$, which has support in $V$, to get a function on $U$, and similarly $\rho_{U}$ is used to define $\left(\rho_{U} \eta\right)^{V}$ on $V$ ! Figure 10.1 shows the circle as the union of two open sets $U$ and $V$ and serves as a schematic of what we have in mind. In the figure $\eta \in \Omega^{0}$ is 1 on one connected component of $U \cap V$ and 0 on the other component. Now we have

$$
\begin{aligned}
& \partial i^{*}\left(-\left(\rho_{V} \eta\right)^{U},\left(\rho_{U} \eta\right)^{V}\right) \\
& \quad=\left.\left(\rho_{U} \eta\right)^{V}\right|_{U \cap V}+\left.\left(\rho_{V} \eta\right)^{U}\right|_{U \cap V} \\
& \quad=\left.\rho_{U} \eta\right|_{U \cap V}+\left.\rho_{V} \eta\right|_{U \cap V} \\
& \quad=\left(\rho_{U}+\rho_{V}\right) \eta=\eta
\end{aligned}
$$

Perhaps the notation is too pedantic. If we let the restrictions and extensions by zero take care of themselves, so to speak, then the idea is expressed by saying that $\partial i^{*}$ maps the element $\left(-\rho_{V} \eta, \rho_{U} \eta\right) \in \Omega(U) \oplus \Omega(V)$ to $\rho_{U} \eta- \left(-\rho_{V} \eta\right)=\eta \in \Omega(U \cap V)$. Thus we see that $\partial i^{*}$ is surjective.

It is easy to see that $\partial i^{*} \circ\left(j_{1}^{*}, j_{2}^{*}\right)=0$ so that $\operatorname{Im}\left(j_{1}^{*}, j_{2}^{*}\right) \subset \operatorname{Ker}\left(\partial i^{*}\right)$. Now let $(\alpha, \beta) \in \Omega(U) \oplus \Omega(V)$ and suppose that $\partial i^{*}(\alpha, \beta)=0$. This translates to $\left.\alpha\right|_{U \cap V}=\left.\beta\right|_{U \cap V}$, which means that there is a form $\omega \in \Omega(U \cup V)=\Omega(M)$ such that $\omega$ coincides with $\alpha$ on $U$ and with $\beta$ on $V$. Thus

$$
\left(j_{1}^{*}, j_{2}^{*}\right) \omega=\left(j_{1}^{*} \omega, j_{2}^{*} \omega\right)=(\alpha, \beta),
$$

so that $\operatorname{Ker}\left(\partial i^{*}\right) \subset \operatorname{Im}\left(j_{1}^{*}, j_{2}^{*}\right)$, which, together with the reverse inclusion, gives $\operatorname{Ker}\left(\partial i^{*}\right)=\operatorname{Im}\left(j_{1}^{*}, j_{2}^{*}\right)$.

Following the general algebraic pattern, the Mayer-Vietoris short exact sequence gives rise to the Mayer-Vietoris long exact sequence:\\
\includegraphics[scale=0.25, center]{2025_10_09_265d7e1a22c89f79c21eg-463}

Since our description of the coboundary map in the algebraic case was rather abstract, we will do well to take a closer look at the connecting homomorphism $d^{*}$ in the present context. Referring to the diagram below, $\eta \in \Omega^{k}(U \cap V)$ represents a cohomology class $[\eta] \in H^{k}(U \cap V)$ so that in particular $d \eta=0$ :\\
\includegraphics[scale=0.25, center]{2025_10_09_265d7e1a22c89f79c21eg-463(1)}

We will abbreviate the map $d \oplus d$ to just $d$ so that $d\left(\eta_{1}, \eta_{2}\right)=\left(d \eta_{1}, d \eta_{2}\right)$. By the exactness of the rows we can find a form $\xi \in \Omega^{k}(U) \oplus \Omega^{k}(V)$ which maps to $\eta$. In fact, we may take $\xi=\left(-\rho_{V} \eta, \rho_{U} \eta\right)$ as before. Since $d \eta=0$ and the diagram commutes, we see that $d \xi$ must map to 0 in $\Omega^{k+1}(U \cap V)$. This just tells us that $-d\left(\rho_{V} \eta\right)$ and $d\left(\rho_{U} \eta\right)$ agree on the intersection $U \cap V$. (Refer again to Figure 10.1.) Thus there is a well-defined form in $\Omega^{k+1}(M)$

\begin{figure}[h]
\begin{center}
  \includegraphics[scale=0.25]{2025_10_09_265d7e1a22c89f79c21eg-464}

\caption{Figure 10.1. Scheme for \$d\^{}\{*}\$ on the circle\}\end{center}
\end{figure}

which maps to $d \xi$. This global form $\gamma$ is given by

$$
\gamma=\left\{\begin{array}{cl}
-d\left(\rho_{V} \eta\right) & \text { on } U, \\
d\left(\rho_{U} \eta\right) & \text { on } V,
\end{array}\right.
$$

and then by definition $d^{*}[\eta]=[\gamma] \in H^{k+1}(M)$.\\
Exercise 10.7. Let the circle $S^{1}$ be parametrized by the angle $\theta$ in the usual way. Let $U$ be the part of a circle with $-\pi / 6<\theta<7 \pi / 6$ and let $V$ be given by $5 \pi / 6<\theta<13 \pi / 6$.\\
(a) Show that $H^{0}(U) \cong H^{0}(V) \cong \mathbb{R}$.\\
(b) Show that the "difference map" $H^{0}(U) \oplus H^{0}(V) \xrightarrow{d^{*}} H^{0}(U \cap V)$ has 1-dimensional image.\\
(c) What is the cohomology of $S^{1}$ ?

\subsection*{Homotopy Invariance}
Now we come to a result about the relation between $H^{*}(M)$ and $H^{*}(M \times \mathbb{R})$ which provides the leverage needed to compute the cohomology of some higher-dimensional manifolds based on that of lower-dimensional manifolds. One of our main goals in this section is to prove the homotopy invariance of de Rham cohomology and also the Poincaré lemma. We start with a theorem which is of interest in its own right. Let $M$ be an $n$-manifold and also consider the product manifold $M \times \mathbb{R}$. We then have the projection\\
$\pi: M \times \mathbb{R} \rightarrow M$, and for a fixed $a \in \mathbb{R}$ we have the section $s_{a}: M \rightarrow M \times \mathbb{R}$ given by $x \mapsto(x, a)$. We can apply the cohomology functor to this pair of maps to obtain the maps $\pi^{*}$ and $s_{a}^{*}$ :

$$
\begin{array}{cc}
M \times \mathbb{R} & H^{*}(M \times \mathbb{R}) \\
\pi\left(\prod_{M} s_{a}\right. & \pi^{*}()^{*} s_{a}^{*} \\
& H^{*}(M)
\end{array}
$$

Theorem 10.8. Given $M$ and the maps defined above, we have that $\pi^{*}$ : $H^{*}(M) \rightarrow H^{*}(M \times \mathbb{R})$ and $s_{a}^{*}: H^{*}(M \times \mathbb{R}) \rightarrow H^{*}(M)$ are mutual inverses for each a. In particular,

$$
H^{*}(M \times \mathbb{R}) \cong H^{*}(M)
$$

Proof. In what follows we let id denote the identity map on $\Omega(M \times \mathbb{R})$, $H^{*}(M)$ and $H^{*}(M \times \mathbb{R})$ as determined by the context. We already know that $\mathrm{id}=s_{a}^{*} \circ \pi^{*}$. The main idea of the proof is the use of a so-called homotopy operator, which in the present case is a degree -1 map $K: \Omega(M \times \mathbb{R}) \rightarrow \Omega(M \times \mathbb{R})$ with the property that

$$
\operatorname{id}-\pi^{*} \circ s_{a}^{*}= \pm(d \circ K-K \circ d)
$$

The point is that $d \circ K-K \circ d$ must send closed forms to exact forms. Thus on the level of cohomology $\pm(d \circ K-K \circ d)=0$, and hence id $-\pi^{*} \circ s_{a}^{*}$ must be the zero map. Thus we also have id $=\pi^{*} \circ s_{a}^{*}$ on $H^{*}(M)$, which gives the result.

So let us then prove equation (10.3) above. Let $\partial_{t}$ denote the obvious vector field that is related to the standard coordinate field on $\mathbb{R}$ under the projection $M \times \mathbb{R} \rightarrow \mathbb{R}$, and let $i_{\partial_{t}}$ denote interior product with respect to $\partial_{t}$. The map $K$ is defined on $\Omega^{k}(M \times \mathbb{R})$ by

$$
K(\omega)=(-1)^{k-1} \int_{a}^{t} \pi^{*} s_{\tau}^{*}\left(i_{\partial_{t}} \omega\right) d \tau=(-1)^{k-1} \int_{a}^{t}\left(s_{\tau} \circ \pi\right)^{*}\left(i_{\partial_{t}} \omega\right) d \tau
$$

More explicitly, for $v_{1}, \ldots, v_{k-1} \in T_{(q, t)} M \times \mathbb{R}$,

$$
\left.K(\omega)\right|_{(q, t)}\left(v_{1}, \ldots, v_{k-1}\right)=\int_{a}^{t} \omega\left(T s_{\tau} \circ T \pi\left(v_{1}\right), \ldots, T s_{\tau} \circ T \pi\left(v_{k-1}\right),\left.\partial_{t}\right|_{\tau}\right) d \tau
$$

This operator and $d$ are both local and linear over $\mathbb{R}$, and so if $\operatorname{id}-\pi^{*} \circ s_{a}^{*}= \pm(d \circ K-K \circ d)$ is true on charts, then it is true in general. Thus we may as well assume that $M$ is an open set $U$ in $\mathbb{R}^{n}$. For any $\omega \in \Omega^{k}(U \times \mathbb{R})$, we can find a pair of functions $f_{1}(x, t)$ and $f_{2}(x, t)$ such that

$$
\omega=f_{1}(x, t) \pi^{*} \alpha+f_{2}(x, t) \pi^{*} \beta \wedge d t
$$

for some forms $\alpha \in \Omega^{k}(U)$ and $\beta \in \Omega^{k-1}(U)$. This decomposition is unique in the sense that if $f_{1}(x, t) \pi^{*} \alpha+f_{2}(x, t) \pi^{*} \beta \wedge d t=0$, then $f_{1}(x, t) \pi^{*} \alpha=0$\\
and $f_{2}(x, t) \pi^{*} \beta \wedge d t=0$. In what follows we will abuse notation a bit. The standard coordinates on $U$ will be denoted by $x^{1}, \ldots, x^{n}$, and we write $\pi^{*} x^{i}$ simply as $x^{i}$ so that with $t$ the standard coordinate of $\mathbb{R}^{1}$, we have coordinates $\left(x^{1}, \ldots, x^{n}, t\right)$ on $U \times \mathbb{R}$. Using the decomposition above, one can see that $K\left(f_{1}(x, t) \pi^{*} \alpha\right)$ is zero and in general

$$
K: \omega \mapsto\left(\int_{a}^{t} f_{2}(x, \tau) d \tau\right) \times \pi^{*} \beta
$$

This map is our proposed homotopy operator.\\
Let us now check that $K$ has the required properties. It is clear from what we have said that we may check the action of $K$ separately on forms of the types $f_{1}(x, t) \pi^{*} \alpha$ and $f_{2}(x, t) \pi^{*} \beta \wedge d t$.

Case I (type $f(x, t) \pi^{*} \alpha$ ). If $\omega=f(x, t) \pi^{*} \alpha$, then $K \omega=0$ and so ( $d \circ K-K \circ d) \omega=-K(d \omega)$. Then

$$
\begin{aligned}
K(d \omega) & =K\left(d\left(f(x, t) \pi^{*} \alpha\right)\right) \\
& =K\left(d f(x, t) \wedge \pi^{*} \alpha+(-1)^{k} f(x, t) \pi^{*} d \alpha\right) \\
& =K\left(\sum \frac{\partial f_{1}}{\partial x^{i}} d x^{i} \wedge \pi^{*} \alpha+\frac{\partial f}{\partial t} d t \wedge \pi^{*} \alpha+(-1)^{k} f(x, t) \pi^{*} d \alpha\right) \\
& =(-1)^{k} K\left(\frac{\partial f}{\partial t}(x, t) \pi^{*} \alpha \wedge d t\right) \\
& =(-1)^{k} \int_{a}^{t} \frac{\partial f}{\partial t}(x, \tau) d \tau \times \pi^{*} \alpha=(-1)^{k}(f(x, t)-f(x, a)) \pi^{*} \alpha
\end{aligned}
$$

On the other hand, since

$$
\begin{aligned}
\pi^{*} s_{a}^{*} f(x, t) \pi^{*} \alpha & =\pi^{*}\left[f(x, a)\left(s_{a}^{*} \circ \pi^{*}\right) \alpha\right] \\
& =\pi^{*}\left[\left(f \circ s_{a}\right) \alpha\right]=\left(f \circ s_{a} \circ \pi^{*}\right) \pi^{*} \alpha=f(x, a) \pi^{*} \alpha
\end{aligned}
$$

we have

$$
\begin{aligned}
\left(\mathrm{id}-\pi^{*} \circ s_{a}^{*}\right) \omega & =\left(\mathrm{id}-\pi^{*} \circ s_{a}^{*}\right) f(x, t) \pi^{*} \alpha \\
& =f(x, t) \pi^{*} \alpha-f(x, a) \pi^{*} \alpha \\
& =(f(x, t)-f(x, a)) \pi^{*} \alpha
\end{aligned}
$$

as above. So in this case we get $(d \circ K-K \circ d) \omega= \pm\left(\mathrm{id}-\pi^{*} \circ s_{a}^{*}\right) \omega$.

Case II (type $\omega=f(x, t) \pi^{*} \beta \wedge d t$ ). We have

$$
\begin{aligned}
d \circ K(\omega)= & d \circ K\left(f(x, t) \pi^{*} \beta \wedge d t\right)=d\left(\pi^{*} \beta \int_{a}^{t} f(x, \tau) d \tau\right) \\
= & \pi^{*} d \beta\left(\int_{a}^{t} f(x, \tau) d \tau\right)+(-1)^{k-1} \pi^{*} \beta \wedge f(x, t) d t \\
& +(-1)^{k-1} \pi^{*} \beta \wedge \sum\left(\int_{a}^{t} \frac{\partial f}{\partial x^{i}}(x, \tau) d \tau\right) d x^{i}
\end{aligned}
$$

and

$$
\begin{aligned}
K \circ d \omega & =K d\left(f \pi^{*} \beta \wedge d t\right)=K d\left(\pi^{*} \beta \wedge f d t\right) \\
& =K\left(\pi^{*} d \beta \wedge f d t+(-1)^{k-1} \pi^{*} \beta \wedge d f \wedge d t\right) \\
& =K\left(\pi^{*} d \beta \wedge f d t+(-1)^{k-1} \pi^{*} \beta \wedge \sum \frac{\partial f}{\partial x^{i}} d x^{i} \wedge d t\right) \\
& =\left(\int_{a}^{t} f(x, \tau) d \tau\right) \pi^{*} d \beta+(-1)^{k-1} \pi^{*} \beta \wedge \sum\left(\int_{a}^{t} \frac{\partial f}{\partial x^{i}}(x, \tau) d \tau\right) d x^{i}
\end{aligned}
$$

Thus $(d \circ K-K \circ d) \omega=(-1)^{k-1} \pi^{*} \beta \wedge f(x, t) d t=(-1)^{k-1} \omega$. On the other hand, we also have ( $\mathrm{id}-\pi^{*} \circ s_{a}^{*}$ ) $\omega=\omega$ since $s_{a}^{*} d t=0$ and so ( $d \circ K-K \circ d)= \pm\left(\mathrm{id}-\pi^{*} \circ s_{a}^{*}\right)$, which is equation (10.3). As explained at the beginning of the proof, this implies that $\pi^{*}$ and $s_{a}^{*}$ are inverses and so $H^{*}(M \times \mathbb{R}) \cong H^{*}(M)$.

Corollary 10.9 (Poincaré lemma).

$$
H^{k}\left(\mathbb{R}^{n}\right)=H^{k}(\text { point })= \begin{cases}\mathbb{R} & \text { if } k=0 \\ 0 & \text { otherwise }\end{cases}
$$

Proof. One first verifies that the statement is true for $H^{k}($ point $)$. Then the remainder of the proof is a simple induction:

$$
\begin{aligned}
H^{k}(\text { point }) & \cong H^{k}(\text { point } \times \mathbb{R})=H^{k}(\mathbb{R}) \\
& \cong H^{k}(\mathbb{R} \times \mathbb{R})=H^{k}\left(\mathbb{R}^{2}\right) \\
& \cong \cdots \\
& \cong H^{k}\left(\mathbb{R}^{n-1} \times \mathbb{R}\right)=H^{k}\left(\mathbb{R}^{n}\right)
\end{aligned}
$$

Corollary 10.10 (Homotopy axiom). If $f: M \rightarrow N$ and $g: M \rightarrow N$ are homotopic, then the induced maps $f^{*}: H^{*}(N) \rightarrow H^{*}(M)$ and $g^{*}: H^{*}(N) \rightarrow H^{*}(M)$ are equal.

Proof. By extending the homotopy as in Exercise 1.77 we may assume that we have a map $F: M \times \mathbb{R} \rightarrow N$ such that

$$
\begin{aligned}
& F(x, t)=f(x) \text { for } t \geq 1 \\
& F(x, t)=g(x) \text { for } t \leq 0
\end{aligned}
$$

If $s_{1}(x):=(x, 1)$ and $s_{0}(x):=(x, 0)$, then $f=F \circ s_{1}$ and $g=F \circ s_{0}$, and so

$$
\begin{aligned}
f^{*} & =s_{1}^{*} \circ F^{*} \\
g^{*} & =s_{0}^{*} \circ F^{*}
\end{aligned}
$$

It is easy to check that $s_{1}^{*}$ and $s_{0}^{*}$ are one-sided inverses of $\pi^{*}$, where $M \times \mathbb{R} \rightarrow M$ is the projection as before. But we have shown that $\pi^{*}$ is an isomorphism. It follows that $s_{1}^{*}=s_{0}^{*}$ in cohomology, and then from the above we have $f^{*}=g^{*}$.

Homotopy plays a central role in algebraic topology, and so the last corollary is very important. Recall that if there exist maps $f: M \rightarrow N$ and $g: N \rightarrow M$ such that both $f \circ g$ and $g \circ f$ are defined and homotopic to $\mathrm{id}_{N}$ and $\mathrm{id}_{M}$ respectively, then $f$ (or $g$ ) is called a homotopy equivalence, and $M$ and $N$ are said to have the same homotopy type. In particular, if a topological space has the same homotopy type as a single point, then we say that the space is contractible. If we are dealing with smooth manifolds, we may take the maps to be smooth. In fact, any continuous map between two smooth manifolds is continuously homotopic to a smooth map. We shall use this fact often without comment. The following corollaries follow easily.

Corollary 10.11 (Homotopy invariance). If $M$ and $N$ are smooth manifolds which are of the same homotopy type, then $H^{*}(M) \cong H^{*}(N)$.

Corollary 10.12. If $M$ is a contractible $n$-manifold, then

$$
\begin{aligned}
& H^{0}(M) \cong \mathbb{R} \\
& H^{k}(M)=0 \text { for } 0<k \leq n
\end{aligned}
$$

Next consider the situation where $A$ is a subset of $M$ and $i: A \hookrightarrow M$ is the inclusion map. If there exists a map $r: M \rightarrow A$ such that $r \circ i=\operatorname{id}_{A}$, then we say that $r$ is a retraction of $M$ onto $A$. If $A$ is a regular submanifold of a smooth manifold $M$, then in case there is a retraction $r$ of $M$ onto $A$, we may assume that $r$ is smooth. If we can find a smooth retraction $r$ such that $i \circ r$ is smoothly homotopic to the identity $\operatorname{id}_{M}$, then we say that $r$ is a (smooth) deformation retraction, and this homotopy itself is also called a deformation retraction. In this case, $A$ is said to be a (smooth) deformation retract of $M$.

Corollary 10.13. If $A$ is a smooth deformation retract of $M$, then $A$ and $M$ have isomorphic cohomologies.

Exercise 10.14. Let $U_{+}$and $U_{-}$be open subsets of the sphere $S^{n} \subset \mathbb{R}^{n+1}$ given by

$$
\begin{aligned}
& U_{+}:=\left\{\left(x^{i}\right) \in S^{n}:-\varepsilon<x^{n+1} \leq 1\right\}, \\
& U_{-}:=\left\{\left(x^{i}\right) \in S^{n}:-1 \leq x^{n+1}<\varepsilon\right\},
\end{aligned}
$$

where $0<\varepsilon<1 / 2$. Show that there is a deformation retraction of $U_{+} \cap U_{-}$ onto the equator $x^{n+1}=0$ in $S^{n}$. Notice that the equator is a two point set in case $n=0$. Show that $S^{n}$ is not contractible. (See also Problem 6.)

Theorem 10.15 (Hairy sphere theorem). A nowhere vanishing smooth vector field exists on the sphere $S^{n}$ if and only if $n$ is odd.

Proof. If $X$ is a nowhere vanishing vector field on $S^{n}$, then define $\nu:= X /\|X\|$. We introduce a map $[0,1] \times S^{n} \rightarrow S^{n}$ by

$$
F(x, t):=(\cos \pi t) x+(\sin \pi t) \nu(x) .
$$

This is a smooth homotopy between the identity map and the antipodal map $a: x \mapsto-x$. Now if $\operatorname{vol}_{S^{n}}$ is the volume form on $S^{n}$, then $\left[\operatorname{vol}_{S^{n}}\right]$ is not zero in $H^{n}\left(S^{n}\right)$. Indeed, suppose that $\operatorname{vol}_{S^{n}}=d \omega$. Then $0 \neq \operatorname{vol}\left(S^{n}\right)= \int_{S^{n}}$ vol $=\int_{\emptyset} \omega=0$, which is a contradiction. Recall the special $n$-form on $\mathbb{R}^{n+1}$ given in Cartesian coordinates by

$$
\sigma:=\sum_{i=1}^{n+1}(-1)^{n-1} x^{i} d x^{1} \wedge \cdots \wedge \widehat{d x^{i}} \wedge \cdots \wedge d x^{n+1}
$$

It was shown previously that $\operatorname{vol}_{S^{n}}$ is the restriction of $\sigma$ to $S^{n}$. But if $n$ is even, then $n+1$ is odd, and it is clear from the above formula that $a^{*} \sigma=-\sigma$. Since $\operatorname{vol}_{S^{n}}$ is obtained from $\sigma$ taking restriction, we have $a^{*} \operatorname{vol}_{S^{n}}=-\operatorname{vol}_{S^{n}}$, and so also on the level of cohomology:

$$
a^{*}\left[\operatorname{vol}_{S^{n}}\right]=-\left[\operatorname{vol}_{S^{n}}\right] .
$$

But this is impossible since if the homotopy exists, we must have $a^{*}=\operatorname{id}^{*}$ on $H^{n}\left(S^{n}\right)$. We conclude that a nonvanishing vector field does not exist on $S^{n}$ if $n$ is even.

If $n$ is odd, then we can easily construct a nowhere vanishing vector field on $S^{n}$. For example, if $n=3$, then let

$$
X=x_{2} \frac{\partial}{\partial x_{1}}-x_{1} \frac{\partial}{\partial x_{2}}+x_{4} \frac{\partial}{\partial x_{3}}-x_{3} \frac{\partial}{\partial x_{4}},
$$

where $x_{1}, x_{2}, x_{3}, x_{4}$ are the standard coordinates on $\mathbb{R}^{4}$, and restrict $X$ to $S^{3}$. Notice that $X$ is indeed tangent to $S^{3}$. The generalization to higher odd dimensions should be clear.

A trivial consequence of the Poincaré lemma is that the cohomology spaces of a Euclidean space are finite-dimensional. Below we use an induction that shows that this is true for a large class of manifolds which include all compact manifolds. For this and later purposes, we introduce a technical condition. An open cover $\left\{U_{\alpha}\right\}_{\alpha \in A}$ of an $n$-manifold $M$ is called a good cover if for every choice $\alpha_{0}, \ldots, \alpha_{k}$ the set $U_{\alpha_{0}} \cap \cdots \cap U_{\alpha_{k}}$ is diffeomorphic to $\mathbb{R}^{n}$ (or empty).

Proposition 10.16. If $\left\{O_{\beta}\right\}_{\beta \in B}$ is any open cover of an $n$-manifold $M$, then there exists a good cover $\left\{U_{\alpha}\right\}_{\alpha \in A}$ such that each $U_{\alpha}$ is contained in some $O_{\beta}$.

Proof. The proof requires a result from Riemannian geometry (see Problem 2 of Chapter 13). We know that every manifold can be given a Riemannian metric. Each point of a Riemannian manifold has a neighborhood system made up of small geodesically convex sets. What matters to us now is that the intersection of any finite number of such sets is geodesically convex and hence diffeomorphic to $\mathbb{R}^{n}$. Actually, the assertion that an open geodesically convex set in a Riemannian manifold is diffeomorphic to $\mathbb{R}^{n}$ is common in the literature, but it is a more subtle issue than it may seem, and references to a complete proof are hard to find (but see [Grom]). Granted this claim, the result follows.

Many of the results we develop below will be proved for orientable manifolds that possess a finite good cover. This includes all orientable compact manifolds.

Theorem 10.17. If an $n$-manifold has a finite good cover, then its de Rham cohomology spaces are all finite-dimensional.

Proof. The proof is an induction argument that uses the Mayer-Vietoris sequence. Our $k$-th induction hypothesis is the following:\\
$\mathrm{P}(k)$ : Every smooth manifold that has a good cover consisting of $k$ open sets has finite-dimensional de Rham cohomologies.

Since we know that an open set diffeomorphic to $\mathbb{R}^{n}$ has finite-dimensional de Rham cohomology spaces, the statement $\mathrm{P}(1)$ is true by Corollary 10.9. Suppose that $\mathrm{P}(k)$ is true. We wish to show that this implies $\mathrm{P}(k+1)$. So suppose that $M$ has a good cover $\left\{U_{1}, \ldots, U_{k+1}\right\}$ and let $M_{k}:=U_{1} \cup \cdots \cup U_{k}$.

Since we are assuming $\mathrm{P}(k), M_{k}$ has finite-dimensional de Rham cohomology spaces. Note that $M_{k} \cap U_{k+1}$ has a good cover, which is just $\left\{U_{1} \cap U_{k+1}, \ldots, U_{k} \cap U_{k+1}\right\}$. From the long Mayer-Vietoris sequence we\\
have for a given $q$

$$
\longrightarrow H^{q}\left(M_{k} \cap U_{k+1}\right) \xrightarrow{d^{*}} H^{q}\left(M_{k+1}\right) \xrightarrow{\iota_{0}^{*}+\iota_{1}^{*}} H^{q+1}\left(M_{k}\right) \oplus H^{q+1}\left(U_{k+1}\right) \longrightarrow,
$$

which gives the exact sequence

$$
0 \longrightarrow \operatorname{Ker} d^{*} \xrightarrow{d^{*}} H^{q}\left(M_{k+1}\right) \xrightarrow{\iota_{0}^{*}+\iota_{1}^{*}} \operatorname{Im}\left(\iota_{0}^{*}+\iota_{1}^{*}\right) \longrightarrow 0 .
$$

Since $H^{q}\left(M_{k} \cap U_{k+1}\right)$ and $H^{q+1}\left(M_{k}\right) \oplus H^{q+1}\left(U_{k+1}\right)$ are finite-dimensional by hypothesis, the same is true of $\operatorname{Ker} d^{*}$ and $\operatorname{Im}\left(\iota_{0}^{*}+\iota_{1}^{*}\right)$. It follows that $H^{q}\left(M_{k+1}\right)$ is finite-dimensional and the induction is complete.

\subsection*{Compactly Supported Cohomology}
Let $\Omega_{c}(M)$ denote the algebra of compactly supported differential forms on a manifold $M$. Obviously, if $M$ is compact, then $\Omega_{c}(M)=\Omega(M)$, and so our main interest here is the case where $M$ is not compact. We now have a slightly different complex

$$
\cdots \xrightarrow{d} \Omega_{c}^{k}(M) \xrightarrow{d} \Omega_{c}^{k+1}(M) \xrightarrow{d} \cdots,
$$

which has a corresponding cohomology $H_{c}^{*}(M)$ called the de Rham cohomology with compact support. By definition

$$
H_{c}^{k}(M)=\frac{Z_{c}^{k}(M)}{B_{c}^{k}(M)}
$$

where $Z_{c}^{k}(M)$ is the vector space of closed $k$-forms with compact support and $B_{c}^{k}(M)$ is the space of all $k$-forms $d \omega$ where $\omega$ has compact support. Note carefully that $B_{c}^{k}(M)$ is not the set of exact $k$-forms with compact support. To drive the point home, consider $f \in C^{\infty}\left(\mathbb{R}^{n}\right)$ with compact support and with $f \geq 0$ and $f>0$ at some point. Then

$$
\omega=f d x^{1} \wedge \cdots \wedge d x^{n}
$$

is exact since every closed form on $\mathbb{R}^{n}$ is exact. However, $\omega$ cannot be $d \alpha$ for an $\alpha$ with compact support, since then we would have

$$
\int_{\mathbb{R}^{n}} \omega=\int_{\mathbb{R}^{n}} d \alpha=\int_{\partial \mathbb{R}^{n}} \alpha=0
$$

which contradicts the assumption that $f$ is nonnegative with $f>0$ at some point. This already shows that $H_{c}^{n}\left(\mathbb{R}^{n}\right) \neq 0$. Using a bump function with support inside a chart, one can similarly show that $H_{c}^{n}(M) \neq 0$ for any orientable manifold $M$.

Exercise 10.18. Let $M$ be an oriented $n$-manifold. If $\omega \in \Omega_{c}^{n}(M)$ and $\int_{M} \omega \neq 0$, then $[\omega] \neq 0$ in $H_{c}^{n}(M)$.\\
Exercise 10.19. Show that $H_{c}^{1}(\mathbb{R}) \cong \mathbb{R}$ and that $H_{c}^{0}(M)=0$ whenever $\operatorname{dim} M>0$ and $M$ is connected.

If we look at the behavior of differential forms under the operation of pull-back, we immediately realize that the pull-back of a differential form with compact support may not have compact support. In order to get desired functorial properties, we consider the class of smooth proper maps. Recall that a smooth map $f: P \rightarrow M$ is called a proper map if $f^{-1}(K)$ is compact whenever $K$ is compact. It is easy to verify that the set of all smooth manifolds together with proper smooth maps is a category and the assignments $M \mapsto \Omega_{c}(M)$ and $f \mapsto\left\{\alpha \mapsto f^{*} \alpha\right\}$ give a contravariant functor. In plain language, this means that if $f: P \rightarrow M$ is a proper map, then $f^{*}: \Omega_{c}(M) \rightarrow \Omega_{c}(P)$ and for two such maps we have $(f \circ g)^{*}=g^{*} \circ f^{*}$ as before, but now the assumption that $f$ and $g$ are proper maps is essential.

We will use a different functorial behavior associated with forms of compact support. The first thing we need is a new category (which is fortunately easy to describe). The category we have in mind has as objects the set of all open subsets of a fixed manifold $M$. The morphisms are the inclusion maps $j_{V, U}: V \hookrightarrow U$, which are only defined in case $V \subset U$. For any such inclusion $j_{V, U}$, we define a map $\left(j_{V, U}\right)_{*}: \Omega_{c}(V) \rightarrow \Omega_{c}(U)$ according to the following simple prescription: For $\alpha \in \Omega_{c}(V)$, let $\left(j_{V, U}\right)_{*} \alpha$ be the form in $\Omega_{c}(U)$ which is equal to $\alpha$ at all points in $V$ and equal to zero otherwise (extension by zero). Since the support of $\alpha$ is neatly inside the open set $V$, the extension $\left(j_{V, U}\right)_{*} \alpha$ is smooth. In what follows, we take this category whenever we employ the functor $\Omega_{c}$.

Let $U$ and $V$ be open subsets which together cover $M$. Recall the commutative diagram of inclusion maps (10.2). Now for each $k$, let us define a map $\left(-i_{1 *}, i_{2 *}\right): \Omega_{c}^{k}(U \cap V) \rightarrow \Omega_{c}^{k}(U) \oplus \Omega_{c}^{k}(V)$ by $\alpha \mapsto\left(-i_{1 *} \alpha, i_{2 *} \alpha\right)$. Now we also have the map $j_{1 *}+j_{2 *}: \Omega_{c}^{k}(V) \oplus \Omega_{c}(U) \rightarrow \Omega_{c}^{k}(M)$ given by $\left(\alpha_{1}, \alpha_{2}\right) \mapsto j_{1 *} \alpha_{1}+j_{2 *} \alpha_{2}$. Notice that if $\left\{\phi_{U}, \phi_{V}\right\}$ is a partition of unity subordinate to $\{U, V\}$, then for any $\omega \in \Omega_{c}^{k}(M)$ we can define $\omega_{U}:=\left.\phi_{U} \omega\right|_{U}$ and $\omega_{V}:=\left.\phi_{V} \omega\right|_{V}$ so that we have

$$
\left(j_{1 *}+j_{2 *}\right)\left(\omega_{U}, \omega_{V}\right)=j_{1 *} \omega_{U}+j_{2 *} \omega_{V}=\phi_{U} \omega+\phi_{V} \omega=\omega .
$$

Notice that $\omega_{U}$ and $\omega_{V}$ have compact support. For example,

$$
\operatorname{Supp}\left(\omega_{U}\right)=\operatorname{Supp}\left(\phi_{U} \omega\right) \subset \operatorname{Supp}\left(\phi_{U}\right) \cap \operatorname{Supp}(\omega)
$$

Since $\operatorname{Supp}\left(\phi_{U}\right) \cap \operatorname{Supp}(\omega)$ is compact, $\operatorname{Supp}\left(\omega_{U}\right)$ is also compact. Thus $j_{1 *}+j_{2 *}$ is surjective. We can associate to the diagram (10.2), the new sequence

$$
0 \rightarrow \Omega_{c}^{k}(V \cap U) \xrightarrow{\left(-i_{1 *}, i_{2 *}\right)} \Omega_{c}^{k}(V) \oplus \Omega_{c}^{k}(U) \xrightarrow{j_{1 *}+j_{2 *}} \Omega_{c}^{k}(M) \rightarrow 0 .
$$

This is the short Mayer-Vietoris sequence for differential forms with compact support.

Theorem 10.20. The sequence (10.4) is exact.

We have shown the surjectivity of $j_{1 *}+j_{2 *}$ above. The rest of the proof is also easy and left as Problem 4.

Corollary 10.21. There is a long exact sequence\\
\includegraphics[scale=0.25, center]{2025_10_09_265d7e1a22c89f79c21eg-473}

\section*{which is called the (long) Mayer-Vietoris sequence for cohomology with compact supports.}
Notice that we have denoted the connecting homomorphism by $d_{*}$. We will need to have a more explicit understanding of $d_{*}$ : If $[\omega] \in H_{c}^{k}(M)$, then using a partition of unity $\left\{\rho_{U}, \rho_{V}\right\}$ as above we write $\omega=j_{1 *} \omega_{U}+j_{2 *} \omega_{V}$ and then $d j_{1 *} \omega_{U}=-d j_{2 *} \omega_{V}$ on $U \cap V$. Then

$$
d_{*}[\omega]=-\left.d j_{1 *} \omega_{U}\right|_{U \cap V}=\left.d j_{2 *} \omega_{V}\right|_{U \cap V}
$$

Next we prove a version of the Poincaré lemma for compactly supported cohomology. For a given $n$-manifold $M$, we consider the projection $\pi$ : $M \times \mathbb{R} \rightarrow M$. We immediately notice that $\pi^{*}$ does not map $\Omega_{c}^{k}(M)$ into $\Omega_{c}^{k}(M \times \mathbb{R})$. What we need is a map $\pi_{*}: \Omega_{c}^{k}(M \times \mathbb{R}) \rightarrow \Omega_{c}^{k-1}(M)$ called integration along the fiber. Before giving the definition of $\pi_{*}$ we first note that every element of $\Omega_{c}^{k}(M \times \mathbb{R})$ is locally a sum of forms of the following types:

$$
\begin{array}{cc}
\text { Type I: } & f \pi^{*} \lambda \\
\text { Type II: } & f \pi^{*} \varphi \wedge d t
\end{array}
$$

where $\lambda \in \Omega_{c}^{k}(M), \varphi \in \Omega_{c}^{k-1}(M)$ and $f$ is a smooth function with compact support on $M \times \mathbb{R}$. By definition $\pi_{*}$ sends all forms of Type I to zero, and for Type II forms, we define

$$
\pi_{*}\left(\pi^{*} \varphi \cdot f \wedge d t\right)=\varphi \int_{-\infty}^{\infty} f(\cdot, t) d t
$$

By linearity this defines $\pi_{*}$ on all forms. In Problem 2 we ask the reader to show that $\pi_{*} \circ d=d \circ \pi_{*}$ so that $\pi_{*}$ is a chain map. Thus we get a map on\\
cohomology:

$$
\pi_{*}: H_{c}^{k}(M \times \mathbb{R}) \rightarrow H_{c}^{k-1}(M)
$$

Next choose $e \in \Omega_{c}^{1}(\mathbb{R})$ with $\int e=1$ and introduce the map $e_{*}: \Omega_{c}^{k}(M) \rightarrow \Omega_{c}^{k+1}(M \times \mathbb{R})$ by

$$
e_{*}: \omega \mapsto \pi^{*} \omega \wedge \pi_{2}^{*} e
$$

where $\pi_{2}: M \times \mathbb{R} \rightarrow \mathbb{R}$ is the projection on the second factor. It is easy to check that $e_{*}$ commutes with $d$, so once again we get a map on the level of cohomology:

$$
e_{*}: H_{c}^{k}(M) \rightarrow H_{c}^{k+1}(M \times \mathbb{R}) \text { for all } k
$$

Our immediate goal is to show that $e_{*} \circ \pi_{*}$ and $\pi_{*} \circ e_{*}$ are both identity operators on the level of cohomology. In fact, it is not hard to see that $\pi_{*} \circ e_{*}=\mathrm{id}$ already on $\Omega_{c}^{k}(M)$. We need to construct a homotopy operator $K$ between $e_{*} \circ \pi_{*}$ and id. The map $K: \Omega_{c}^{k}(M \times \mathbb{R}) \rightarrow \Omega_{c}^{k-1}(M \times \mathbb{R})$ is given by requiring that $K$ is linear, maps Type I forms to zero, and if $\omega=\pi^{*} \varphi \cdot f \wedge d t$ is Type II, then

$$
K(\omega)=K\left(\pi^{*} \varphi \cdot f \wedge d t\right):=\pi^{*} \varphi\left(\int_{-\infty}^{t} f(x, u) d u-T(t) \int_{-\infty}^{\infty} f(x, u) d u\right)
$$

where $T(t)=\int_{-\infty}^{t} e$. In Problem 3, we ask the reader to show that

$$
\mathrm{id}-e_{*} \circ \pi_{*}=(-1)^{k-1}(d K-K d)
$$

on $\Omega_{c}^{k}(M \times \mathbb{R})$. It follows that $\mathrm{id}=e_{*} \circ \pi_{*}$ on $H_{c}^{k}(M \times \mathbb{R})$, and so finally we have the following result:

Theorem 10.22. With notation as above the following maps are isomorphisms and mutual inverses:

$$
\begin{aligned}
& \pi_{*}: H_{c}^{k}(M \times \mathbb{R}) \rightarrow H_{c}^{k-1}(M), \\
& e_{*}: H_{c}^{k-1}(M) \rightarrow H_{c}^{k}(M \times \mathbb{R}) .
\end{aligned}
$$

Corollary 10.23 (Poincaré lemma for compactly supported cohomology). We have

$$
\begin{aligned}
H_{c}^{n}\left(\mathbb{R}^{n}\right) & =\mathbb{R} \\
H_{c}^{k}\left(\mathbb{R}^{n}\right) & =0 \text { if } k \neq n
\end{aligned}
$$

Proof. By Exercise 10.19 we have $H_{c}^{1}(\mathbb{R}) \cong \mathbb{R}$. But from the previous theorem $H_{c}^{n}\left(\mathbb{R}^{n}\right) \cong H_{c}^{n-1}\left(\mathbb{R}^{n-1}\right) \cong \ldots \cong H_{c}^{1}(\mathbb{R})$. Also, trivially, $H_{c}^{k}\left(\mathbb{R}^{n}\right)=$ 0 for $k>n$, and if $k<n$, then

$$
H_{c}^{k}\left(\mathbb{R}^{n}\right) \cong H_{c}^{k-1}\left(\mathbb{R}^{n-1}\right) \cong \ldots \cong H_{c}^{0}\left(\mathbb{R}^{n-k}\right)=0
$$

by Exercise 10.19.

The isomorphism $H_{c}^{n}\left(\mathbb{R}^{n}\right)=\mathbb{R}$ is given by repeated application of $\pi_{*}$, which is just integration over the last coordinate of $\mathbb{R}^{n}$.

Remark: Notice that $\mathbb{R}^{n}$ is homotopy equivalent to $\mathbb{R}^{n-1}$ and yet, $H_{c}^{n}\left(\mathbb{R}^{n}\right) \neq H_{c}^{n}\left(\mathbb{R}^{n-1}\right)$. This contrasts the compactly supported cohomology with the ordinary cohomology, since the latter is a homotopy invariant.

\subsection*{Poincaré Duality}
Looking back at our results for $\mathbb{R}^{n}$ we notice that $H^{k}\left(\mathbb{R}^{n}\right)=H_{c}^{n-k}\left(\mathbb{R}^{n}\right)^{*}$. This is a special case of Poincaré duality which is proved in this section. Recall that we have already met Poincaré duality for compact manifolds in Section 9.6. The version we develop here will be valid for noncompact manifolds also. So let $M$ be an oriented $n$-manifold which is not necessarily compact. For each $k$, we have a bilinear pairing

$$
\begin{aligned}
\Omega^{k}(M) \times \Omega_{c}^{n-k}(M) & \rightarrow \mathbb{R}, \\
\left(\omega_{1}, \omega_{2}\right) & \mapsto \int_{M} \omega_{1} \wedge \omega_{2},
\end{aligned}
$$

and as in Section 9.6 this defines a pairing on cohomology:

$$
\begin{aligned}
H^{k}(M) \times H_{c}^{n-k}(M) & \rightarrow \mathbb{R} \\
\left(\left[\omega_{1}\right],\left[\omega_{2}\right]\right) & \mapsto \int_{M} \omega_{1} \wedge \omega_{2}
\end{aligned}
$$

In turn, this provides a linear map $P D_{k}: H^{k}(M) \rightarrow\left(H_{c}^{n-k}(M)\right)^{*}$ defined as

$$
P D_{k}\left(\left[\omega_{1}\right]\right)\left(\left[\omega_{2}\right]\right):=\int_{M} \omega_{1} \wedge \omega_{2}
$$

Our goal is to show that this map is an isomorphism. We prove this for orientable manifolds with a finite cover, but the result remains valid more generally (see [Madsen]). Notice also that unless the cohomologies are finite-dimensional, we may have $H^{k}\left(\mathbb{R}^{n}\right)^{*} \neq H_{c}^{n-k}\left(\mathbb{R}^{n}\right)$. The reason is that for an infinite-dimensional vector spaces, $V^{*} \cong W$ does not imply $V \cong W^{*}$.

In what follows, we will denote all the duality maps $P D_{k}$ by the single designation $P D$. We use a theorem from algebra called the five lemma: Consider the following diagram of vector spaces (or abelian groups) and homomorphisms:\\
\includegraphics[scale=0.25, center]{2025_10_09_265d7e1a22c89f79c21eg-475}

Lemma 10.24 (Five lemma). Suppose that the above diagram commutes up to sign (so for example $h_{1}^{\prime} \circ \alpha= \pm \beta \circ h_{1}$ ). Suppose also that the rows are exact and that $\alpha, \beta, \delta$, and $\varepsilon$ are isomorphisms. Then the middle map $\gamma$ is also an isomorphism.

The five lemma is usually stated for the case when the diagram commutes, but a simple modification of the usual proof (see [L2]) gives the version above.

If $\{U, V\}$ is a cover of the oriented $n$-manifold $M$, then consider the following diagram:\\
\includegraphics[scale=0.25, center]{2025_10_09_265d7e1a22c89f79c21eg-476(1)}

Here the bottom row is the sequence of dual maps from the Mayer-Vietoris long exact sequence and is exact itself (easy check). The vertical maps are the duality maps $P D$ (or obvious direct sums of such).

Lemma 10.25. The above diagram commutes up to sign.

Proof. First let us consider the diagram\\
\includegraphics[scale=0.25, center]{2025_10_09_265d7e1a22c89f79c21eg-476}

In order to show that $(P D \oplus P D) \circ\left(j_{1}^{*}+j_{2}^{*}\right)=\left(\left(j_{1 *}\right)^{*}+\left(j_{2 *}\right)^{*}\right) \circ P D$, it is enough to show that $P D \circ j_{1}^{*}=\left(j_{1 *}\right)^{*} \circ P D$ and $P D \circ j_{2}^{*}=\left(j_{2 *}\right)^{*} \circ P D$. For a given $[\omega] \in H^{k}(M)$, the linear form $P D \circ j_{1}^{*}([\omega])$ takes an element $[\theta] \in H^{k}(U)$ to

$$
\int_{U} j_{1}^{*} \omega \wedge \theta .
$$

On the other hand, $\left(\left(j_{1 *}\right)^{*} \circ P D\right)([\omega])$ maps $[\theta]$ to

$$
\int_{M} \omega \wedge j_{1 *} \theta .
$$

But $\omega \wedge j_{1 *} \theta$ and $j_{1}^{*} \omega \wedge \theta$ both have support in $U$ and agree on this set. Therefore the above two integrals are equal. Similarly, $P D \circ j_{2}^{*}=\left(j_{2 *}\right)^{*} \circ P D$.

Next we show that the following square commutes up to sign:\\
\includegraphics[scale=0.25, center]{2025_10_09_265d7e1a22c89f79c21eg-477}

Let $\left\{\rho_{U}, \rho_{V}\right\}$ be a partition of unity subordinate to $\{U, V\}$ as before. If $[\omega] \in H^{k-1}(U \cap V)$, then $d^{*}[\omega]$ is represented by a form (which we denote by $\left.d^{*} \omega\right)$ that has the properties

$$
\begin{aligned}
\left.d^{*} \omega\right|_{U} & =-d\left(\rho_{V} \omega\right) \text { on } U \\
\left.d^{*} \omega\right|_{V} & =d\left(\rho_{U} \omega\right) \text { on } V
\end{aligned}
$$

On the other hand, if $[\tau] \in H_{c}^{n-k}(M)$, then $d_{*}[\tau]$ is represented by a form $d_{*} \tau$ which has the properties

$$
\begin{aligned}
-i_{1 *} d_{*} \tau & =d\left(\rho_{U} \omega\right) \text { on } U \\
i_{2 *} d_{*} \tau & =d\left(\rho_{V} \omega\right) \text { on } V
\end{aligned}
$$

Using the fact that $\omega$ is closed, we have

$$
\begin{aligned}
P D \circ d^{*}([\omega])([\tau]) & =\int_{M} d^{*} \omega \wedge \tau=\int_{U \cap V} d^{*} \omega \wedge \tau \\
& =\int_{U \cap V} d\left(\rho_{U} \omega\right) \wedge \tau=\int_{U \cap V} d \rho_{U} \wedge \omega \wedge \tau
\end{aligned}
$$

where have used the fact that $d^{*} \omega \wedge \tau$ has support in $U \cap V$. Meanwhile,

$$
\begin{aligned}
\left(d_{*}\right)^{*} \circ P D([\omega])([\tau]) & =\int_{U \cap V} \omega \wedge d_{*} \tau=-\int_{U \cap V} \omega \wedge d\left(\rho_{U} \tau\right) \\
& =-\int_{U \cap V} \omega \wedge d \rho_{U} \wedge \tau
\end{aligned}
$$

so the two integrals are equal up to sign. We leave the sign commutativity of the remaining square to the reader.

Theorem 10.26 (Poincaré duality). Let $M$ be an oriented $n$-manifold with a finite good cover. Then for each $k$ the map $P D_{k}: H^{k}(M) \rightarrow\left(H_{c}^{n-k}(M)\right)^{*}$ is an isomorphism.

Proof. We will prove that $P D_{k}$ is an isomorphism by induction. The inductive statement $P(N)$ is that $P D_{k}$ is an isomorphism for all orientable manifolds which have a good cover consisting of at most $N$ open sets. By the two Poincaré lemmas we already know that if $U$ is diffeomorphic to $\mathbb{R}^{n}$, then $H^{k}(U) \cong\left(H_{c}^{n-k}(U)\right)^{*}$, where both sides are either zero or isomorphic to $\mathbb{R}$. It is then an easy exercise to show that in fact $P D_{k}: H^{k}(U) \rightarrow\left(H_{c}^{n-k}(U)\right)^{*}$ is an isomorphism. This verifies the statement $P(1)$. Now assume that\\
$P(N)$ is true and suppose we are given an oriented manifold $M$ with a good cover $\left\{U_{1}, \ldots, U_{N+1}\right\}$. Let $M_{N}:=U_{1} \cup \cdots \cup U_{N}$ and use Lemma 10.25 with $U:=M_{N}, V:=U_{N+1}$ and $U \cup V=M$. Then by assumption $P D$ is an isomorphism for $U, V$ and $U \cap V$, so the hypotheses of Lemma 10.24 are satisfied. But then the middle homomorphism $H^{k}(M) \rightarrow H_{c}^{n-k}(M)^{*}$ is an isomorphism ( $k$ was arbitrary). \(\square\)

Corollary 10.27. If $M$ is a connected oriented $n$-manifold with finite good cover, then $H_{c}^{n}(M) \cong \mathbb{R}$. This isomorphism is given by integration over $M$.

This last corollary allows us to define the important concept of degree of a proper map. Let $M$ and $N$ be connected oriented $n$-manifolds and suppose that $f: M \rightarrow N$ is a proper map. Then a pull-back of a compactly supported form is compactly supported so that we obtain a map $f^{*}: H_{c}^{n}(N) \rightarrow H_{c}^{n}(M)$ and the following commutative diagram:\\
\includegraphics[scale=0.25, center]{2025_10_09_265d7e1a22c89f79c21eg-478}

The induced map $\operatorname{Deg}_{f}: \mathbb{R} \rightarrow \mathbb{R}$ must be multiplication by a real number. This number is called the degree of $f$ and is denoted $\operatorname{deg}(f)$. If $[\omega] \in H_{c}^{n}(N)$ is chosen so that $\int_{N} \omega=1$, then

$$
\int_{M} f^{*} \omega=\operatorname{deg}(f) .
$$

Such a form is called a generator and can be chosen to have support in an arbitrary open subset $V$ of $N$. To see this, just choose a chart ( $U, \mathrm{x}$ ) in $N$ with $U \subset \bar{U} \subset V$ and let $\phi$ be a cut-off function with $\operatorname{supp} \phi \subset U$. Then let $\omega:=\phi d x^{1} \wedge \cdots \wedge d x^{n}$ and then scale $\phi$ so that $\int_{N} \omega=1$. Of course, $\omega$ is closed, but also $[\omega] \neq 0$. Indeed, if we had $[\omega]=0$, then there would be an ( $n-1$ )-form $\eta$ with compact support and with $\omega=d \eta$. But then Stokes' theorem applies, and so $\int_{N} \omega=\int_{N} d \eta=\int_{\partial N} \eta=0$ since $\partial N=\emptyset$. This would contradict $\int_{N} \omega=1$.

Theorem 10.28. Let $M$ and $N$ be connected $n$-manifolds oriented by volume forms $\mu_{M}$ and $\mu_{N}$ respectively. If $f: M \rightarrow N$ is a proper map and $y \in N$ is a regular value, then

$$
\operatorname{deg}(f)=\sum_{x \in f^{-1}(y)} \operatorname{sign}_{x} f,
$$

where $\operatorname{sign}_{x} f=1$ if $\left(f^{*} \mu_{N}\right)(x)$ is a positive multiple of $\mu_{M}(x)$ and -1 otherwise. In particular, $\operatorname{deg}(f)$ is an integer.

Proof. Using Sard's Theorem 2.34 and the fact that $f$ is proper, we may choose a neighborhood $V$ of $y$ such that $f^{-1}(V)=\bigcup_{i=1}^{N} U_{i}$, where $U_{i} \cap U_{j}=\emptyset$ for $i \neq j$ and such that $\left.f\right|_{U_{i}}$ is a diffeomorphism onto $V$ for each $i$. Choose $[\omega] \in H_{c}^{n}(N)$ with $\int_{N} \omega=1$, and such that $\omega$ has support in $V$. We may arrange that $V$ and all $U_{i}$ are connected so that $\operatorname{sign}_{x} f$ is constant on each $U_{i}$. In this case, $\left.f\right|_{U_{i}}$ is orientation preserving or reversing according to whether $\operatorname{sign}_{x} f$ is 1 or -1 . Let $x_{i} \in U_{i}$ be the inverse image of $y$ in $U_{i}$. We have

$$
\begin{aligned}
\operatorname{deg}(f) & =\int_{M} f^{*} \omega=\sum_{i=1}^{N} \int_{U_{i}} f^{*} \omega=\sum_{i=1}^{N}\left(\operatorname{sign}_{x_{i}} f\right) \int_{V} \omega \\
& =\sum_{x \in f^{-1}(y)} \operatorname{sign}_{x} f
\end{aligned}
$$

We close this chapter with a quick explanation of the Poincaré dual of a submanifold. Let $Y$ be an oriented regular submanifold of an oriented $n$-manifold. Let the dimension of $Y$ be $n-k$. Then we obtain a linear map

$$
\begin{aligned}
\int_{Y} & : \Omega_{c}^{n-k}(M) \longrightarrow \mathbb{R} \\
\omega & \mapsto \int_{Y} \iota^{*} \omega
\end{aligned}
$$

where $\iota: Y \hookrightarrow M$ is the inclusion map. Since this map is zero on exact forms (by Stokes' theorem), it passes to the quotient giving a linear functional

$$
\int_{Y}: H_{c}^{n-k}(M) \longrightarrow \mathbb{R}
$$

By Poincaré duality $P D: H^{k}(M) \cong\left(H_{c}^{n-k}(M)\right)^{*}$, there must be a unique class $\left[\theta_{Y}\right] \in H^{k}(M)$ such that this linear form is given by

$$
\int_{Y}[\omega]=\int_{Y} \theta_{Y} \wedge \omega \text { for all }[\omega] \in H_{c}^{n-k}(M)
$$

This class (or a given representative $\theta_{Y}$ ) is called the Poincaré dual of the submanifold $Y$.

For more on this and other important topics such as the Thom isomorphism, the Leray-Hirsch theorem, and Čech cohomology see [Bo-Tu].

\section*{Problems}
(1) Show that formula (10.1) defines a smooth function $\theta$ on the open set $U:=\mathbb{R}^{2} \backslash\{x \leq 0, y=0\}$ and show that

$$
d \theta=\frac{x d y-y d x}{x^{2}+y^{2}} \text { on } U \text {. }
$$

(2) Show that $\pi_{*}: \Omega_{c}^{*}(M \times \mathbb{R}) \rightarrow \Omega_{c}^{*-1}(M)$ as defined in the discussion leading to Theorem 10.22 is a chain map.\\[0pt]
(3) Prove the homotopy formula (10.6) (or see [Bo-Tu], pages 38-39).\\
(4) Prove exactness of (10.4).\\
(5) Prove Corollary 10.11.\\
(6) Use Exercise 10.14 and the long Mayer-Vietoris sequence to show that

$$
H^{k}\left(S^{n}\right)= \begin{cases}\mathbb{R} & \text { if } k=0 \text { or } n, \\ 0 & \text { otherwise } .\end{cases}
$$

(7) Show that $\omega \in \Omega^{n}\left(S^{n}\right)$ is exact if and only if $\int_{S^{n}} \omega=0$.\\
(8) Determine the cohomology spaces for the punctured Euclidean spaces $\mathbb{R}^{n} \backslash\{0\}$. Show that if $B$ is a ball, then $\mathbb{R}^{n} \backslash \bar{B}$ has the same cohomology.\\
(9) Let $B_{x}$ be a ball centered at $x \in \mathbb{R}^{n}$. Show that a closed ( $n-1$ )-form on $\mathbb{R}^{n} \backslash \overline{B_{x}}$ is exact if and only $\int_{S} \iota^{*} \omega=0$ for some sphere centered at $x$. Here $\iota: S \hookrightarrow \mathbb{R}^{n}$ is the inclusion map. Show that this is true for all spheres centered at $x$.\\
(10) Let $n \geq 2$ and suppose that $\omega$ is a compactly supported $n$-form on $\mathbb{R}^{n}$ such that $\int_{\mathbb{R}^{n}} \omega=0$.\\
(a) Let $B_{1}$ and $B_{2}$ be open balls centered at the origin of $\mathbb{R}^{n}$ with $B_{1} \subset \bar{B}_{1} \subset B_{2}$ and $\operatorname{supp} \omega \subset B_{1}$. By the (first) Poincaré lemma there is a form $\alpha$ such that $d \alpha=\omega$. Show that $\int_{\partial B_{2}} \alpha=0$.\\
(b) Continuing from (a), show that $\alpha$ is exact on $\mathbb{R}^{n} \backslash \bar{B}$ (use Problem 9 above).\\
(c) Let $\rho \in \Omega^{n-2}\left(\mathbb{R}^{n} \backslash \bar{B}\right)$ be such that $d \rho=\alpha$ where $\alpha$ is as above. Show that there is a function $g$ such that if $\beta:=\alpha-d(g \rho)$, then $\beta$ is smooth, compactly supported and $d \beta=\omega$. Deduce that $H_{c}^{n}\left(\mathbb{R}^{n}\right)=0$.\\
(11) Find $H^{*}(M)$ where $M$ is $\mathbb{R}^{n} \backslash\{p, q\}$ for some distinct points $p, q \in \mathbb{R}^{n}$.

\section*{Distributions and Frobenius' Theorem}
The theory of distributions is a geometric formulation of the classical theory of certain systems of partial differential equations. The solutions are immersed submanifolds called integral manifolds. The Frobenius theorem gives necessary and sufficient conditions for the existence of such integral manifolds. One of the most important applications is to show that a subalgebra of the Lie algebra of a Lie group corresponds to a Lie subgroup. We give this application near the end of this chapter. We also develop a classic version of the Frobenius theorem and use it to prove the basic existence theorem for surfaces.

One can view the theory studied in this chapter as a higher-dimensional analogue of the study of vector fields and integral curves. If $X$ is a smooth vector field on an $n$-manifold $M$, then we know that integral curves exist through each point, and if $X$ never vanishes, then the integral curves are immersions. Nonvanishing vector fields do not even exist on many manifolds, and so we consider a somewhat different question. A one-dimensional distribution assigns to each $p \in M$ a one-dimensional subspace $E_{p}$ of the tangent space $T_{p} M$. We say that the assignment is smooth if for each $p \in M$ there is a smooth vector field $X$ defined in an open set $U$ containing $p$ such that $X(p)$ spans $E_{p}$ for each $p \in U$. It is easy to see that a one-dimensional distribution is essentially the same thing as a rank one subbundle of the tangent bundle.

A curve $c:(a, b) \rightarrow M$ with $\dot{c} \neq 0$ is called an integral curve of the one-dimensional distribution if $\dot{c}(t)=T_{t} c \cdot \frac{\partial}{\partial t}$ is contained in $E_{c(t)}$ for each

\begin{figure}[h]
\begin{center}
  \includegraphics[scale=0.25]{2025_10_09_265d7e1a22c89f79c21eg-483}

\caption{Figure 11.1. No integral manifold through the origin}
\end{center}
\end{figure}

$t \in(a, b)$. The curve is an immersion, and its restrictions to small enough subintervals are integral curves of vector fields that locally span the distribution. Since integral curves cannot cross themselves or each other, the curve is an injective immersion. The images of such curves are called integral manifolds (they are immersed submanifolds). Of course, every nonvanishing global vector field defines a one-dimensional distribution, but some onedimensional distributions do not arise in this way. For example, the tangent spaces to the fibers of the Möbius band bundle $\mathrm{MB} \rightarrow S^{1}$ form a distribution not spanned by any globally defined vector field (see Example 6.9). The integral manifolds are the fibers themselves.

\subsection*{Definitions}
We wish to generalize the above idea to higher-dimensional distributions.\\
Definition 11.1. A smooth rank $k$ distribution on an $n$-manifold $M$ is a (smooth) rank $k$ vector subbundle $E \rightarrow M$ of the tangent bundle.

Sometimes we refer to a distribution as a tangent distribution to distinguish it from the notion of a distributional function on a manifold. Recalling the definition of a vector subbundle (Definition 6.26) and the criteria given in Exercise 6.27, we see that a smooth rank $k$ distribution on an $n$-manifold $M$ gives a $k$-dimensional subspace $E_{p} \subset T_{p} M$ for each $p \in M$ such that for each fixed $p \in M$ there is a family of smooth vector fields $X_{1}, \ldots, X_{k}$ defined on some neighborhood $U_{p}$ of $p$ and such that $X_{1}(q), \ldots, X_{k}(q)$ are linearly independent and span $E_{q}$ for each $q \in U_{p}$. We say that such a set of vector fields locally spans the distribution. In other words, $X_{1}, \ldots, X_{k}$ can be viewed as a local frame field for the subbundle $E$.

Consider the punctured 3 -space $M=\mathbb{R}^{3} \backslash\{0\}$. The level sets of the function $\varepsilon:(x, y, x) \mapsto x^{2}+y^{2}+x^{2}$ are spheres whose union is all of $\mathbb{R}^{3} \backslash\{0\}$. Now define a distribution on $M$ by the rule that $E_{p}$ is the tangent space at\\
$p$ to the sphere containing $p$. Dually, we can define this distribution to be the one given by the rule

$$
E_{p}=\left\{v \in T_{p} M: d \varepsilon(v)=0\right\}
$$

Thus we see that a distribution of rank $n-1$ can arise as the family of tangent spaces to the level sets of nondegenerate smooth functions. More generally, rank $k$ distributions can arise from the level sets of sufficiently nondegenerate $\mathbb{R}^{n-k}$-valued functions. On the other hand, not all distributions arise in this way, even locally.

Definition 11.2. Let $E \rightarrow M$ be a distribution. An immersed submanifold $S$ of $M$ is called an integral manifold (or integral submanifold) for the distribution if for each $p \in S$ we have $T \iota\left(T_{p} S\right)=E_{p}$, where $\iota: S \hookrightarrow M$ is the inclusion map. The distribution is called integrable if each point of $M$ is contained in an integral manifold of the distribution. ${ }^{1}$

If $\iota: S \hookrightarrow M$ is the inclusion map, then we usually identify $T \iota\left(T_{p} S\right)$ with $T_{p} S$ so the condition which makes $S$ an integral manifold is stated simply as $E_{p}=T_{p} S$ for all $p \in S$.

Consider the distribution in $\mathbb{R}^{3}$ suggested by Figure 11.1. For each $p$, $E_{p}$ is the subspace of $T_{p} \mathbb{R}^{3}$ orthogonal to

$$
\left.y \frac{\partial}{\partial x}\right|_{p}-\left.x \frac{\partial}{\partial y}\right|_{p}+\left.\frac{\partial}{\partial z}\right|_{p}
$$

If there were an integral manifold though the origin, then it would clearly be tangent to the $x, y$ plane at the origin. But if one tries to imagine following a closed path on such an integral manifold around the $z$ axis, a problem appears. Namely, the curve would have to be tangent to the distribution, but this would clearly force the curve to spiral upward and never return to the same point.

A distribution as defined above is sometimes called a regular distribution to distinguish the concept from that of a singular distribution which is defined in a similar manner, but allows for the dimension of the subspaces to vary from point to point. Singular distributions are studied in the online supplement [Lee, Jeff].

Definition 11.3. Let $E \rightarrow M$ be a distribution on an $n$-manifold $M$ and let $X$ be a vector field defined on an open subset $U \subset M$. We say that $X$ lies in the distribution if $X(p) \in E_{p}$ for each $p$ in the domain of $X$; that is, if $X$ takes values in $E$. In this case, we sometimes write $X \in E$ (a slight abuse of notation).

\footnotetext{${ }^{1}$ In other contexts, integral manifold might be defined by the weaker condition $T \iota\left(T_{p} S\right) \subset E_{p}$. We will always require the stronger condition $T \iota\left(T_{p} S\right)=E_{p}$.
}Definition 11.4. If for every pair of locally defined vector fields $X$ and $Y$ with common domain that lie in a distribution $E \rightarrow M$, the bracket $[X, Y]$ also lies in the distribution, then we say that the distribution is involutive.

It is easy to see that a vector field $X$ lies in a distribution if and only if for any spanning local frame field $X_{1}, \ldots, X_{k}$ we have $X=\sum_{i=1}^{k} f^{i} X_{i}$ for smooth functions $f^{i}$ defined on the common domain of $X, X_{1}, \ldots, X_{k}$.

Locally, finding integral curves of vector fields is the same as solving a system of ordinary differential equations. The local model for finding integral manifolds of a distribution is that of solving a system of partial differential equations. Consider the following system of partial differential equations defined on some neighborhood $U$ of the origin in $\mathbb{R}^{2}$ :

$$
\begin{aligned}
& z_{x}=F(x, y, z) \\
& z_{y}=G(x, y, z)
\end{aligned}
$$

Here the functions $F$ and $G$ are assumed to be smooth and defined on $U \times \mathbb{R}$. We might attempt to find a solution defined near the origin by integrating twice. Suppose we want a solution with $z(0,0)=z_{0}$. First, solve $f^{\prime}(x)=F(x, 0, f(x))$ with $f(0)=z_{0}$. Then for each fixed $x$ in the domain of $f$ solve $\partial_{y} z(x, y)=G(x, y, z(x, y))$ with $z(x, 0)=f(x)$. This function $z(x, y)$ may not be a $C^{2}$ solution of our system of partial differential equations since, if it were, we would have $z_{x y}=z_{y x}$. But

$$
z_{x y}=\left(z_{x}\right)_{y}=\frac{\partial}{\partial y} F(x, y, z)=F_{y}+G F_{y}
$$

and similarly, $z_{y x}=G_{x}+F G_{x}$. Thus the equality of mixed partials gives a necessary integrability condition;

$$
F_{y}+G F_{y}=G_{x}+F G_{x} .
$$

We shall see that this condition is also sufficient for the local solvability of our system.

\section*{Exercise 11.5.}
(a) Show that solving the differential equations above is equivalent to finding the integral curves of the following pair of vector fields.

$$
X=\frac{\partial}{\partial x}-F(x, y, z) \frac{\partial}{\partial z} \text { and } Y=\frac{\partial}{\partial x}-G(x, y, z) \frac{\partial}{\partial z}
$$

(b) Show that the graph of a solution to the system (11.1) is an integral curve of the distribution spanned by $X$ and $Y$.\\
(c) Show that $[X, Y]=0$ is the same condition as (11.2).

There are no integrability conditions for finding integral curves of a vector field. But now we see that we should expect integrability conditions for the existence of integral manifolds of higher rank distributions.

\subsection*{The Local Frobenius Theorem}
We will relate the notion of involutivity to the existence of integral manifolds.\\
Lemma 11.6. Let $E \rightarrow M$ be a distribution (tangent subbundle) on an $n$ manifold $M$. Then $E \rightarrow M$ is involutive if and only if for every local frame field $X_{1}, \ldots, X_{k}$ for the subbundle $E \rightarrow M$ we have

$$
\left[X_{i}, X_{j}\right]=\sum_{s=1}^{k} c_{i j}^{s} X_{s}
$$

for some smooth functions $c_{i j}^{s}(1 \leq i, j, s \leq k)$.\\
Proof. It is clear that such functions exist if $E$ is involutive. Now, suppose that such functions exist. If $X$ and $Y$ lie in the distribution, then

$$
\begin{aligned}
& X=\sum_{i=1}^{k} f^{i} X_{i} \\
& Y=\sum_{j=1}^{k} g^{j} X_{j}
\end{aligned}
$$

for smooth functions $f^{i}$ and $g^{j}$. We wish to show that $[X, Y]$ lies in the distribution. This follows by direct calculation:

$$
\begin{aligned}
{[X, Y] } & =\left[\sum_{i} f^{i} X_{i}, \sum_{j} g^{j} X_{j}\right] \\
& =\sum_{i, j} f^{i} g^{j}\left[X_{i}, X_{j}\right]+\sum_{i, j} f^{i}\left(X_{i} g^{j}\right) X_{j}-\sum_{i, j} g^{j}\left(X_{j} f^{i}\right) X_{i}
\end{aligned}
$$

We will show that a distribution is involutive if and only if it is integrable. In fact, we shall see that the integral manifolds fit together nicely. In the following discussion we identify $\mathbb{R}^{k} \times \mathbb{R}^{n-k}$ with $\mathbb{R}^{n}$.

Definition 11.7. A rank $k$ distribution $E \rightarrow M$ on an $n$-manifold $M$ is called completely integrable if for each $p \in M$ there is a chart ( $U, \mathrm{x}$ ) centered at $p$ such that $\mathrm{x}(U)=V \times W \subset \mathbb{R}^{k} \times \mathbb{R}^{n-k}$ and such that each set of the form $\mathrm{x}^{-1}(V \times\{y\})$ for each $y \in W$ is an integral manifold of the distribution.

If $\mathrm{x}(p)=\left(a^{1}, \ldots, a^{n}\right)$, then

$$
\mathbf{x}^{-1}(V \times\{y\})=\left\{q \in U: x^{k+1}(q)=a^{k+1}, \ldots, x^{n}(q)=a^{n}\right\}
$$

where $y=\left(a^{k+1}, \ldots, a^{n}\right)$. For a chart as in the last definition we have

$$
T\left(\mathrm{x}^{-1}(V \times\{y\})\right)=\left.E\right|_{\mathrm{x}^{-1}(V \times\{y\})}
$$

for each $y \in W$. We can always arrange that $V$ and $W$ in the above definition are connected (say open cubes or open balls). In this case, we refer to such a chart as a distinguished chart for the distribution. The integral manifolds of the form $\mathrm{x}^{-1}(V \times\{y\})$ for a distinguished chart are called plaques. Notice that since $V$ and $W$ are assumed path connected for a distinguished chart, the plaques are path connected. The reader was asked to prove the following result in Problem 4 of Chapter 3, but we will state and prove the result here for the sake of completeness.

Lemma 11.8. Suppose that $f: N \rightarrow M$ is a one-one immersion and that $Y$ is a vector field on $M$ such that $Y(f(p)) \in T_{p} f\left(T_{p} N\right)$ for every $p \in N$. Then there is a unique vector field $X$ on $N$ that is $f$-related to $Y$.

Proof. It is clear that we can define a field $X$ such that $X(p)$ is the unique element of $T_{p} N$ with $Y(f(p))=T_{p} f \cdot X(p)$. The task is to show that $X$ so defined is smooth. For $p \in N$, let ( $U, \mathrm{x}$ ) be a chart centered at $p$ and ( $V, \mathrm{y}$ ) a chart centered around $f(p)$ such that $\mathrm{y} \circ f \circ \mathrm{x}^{-1}$ has the form

$$
\left(a^{1}, \ldots, a^{n}\right) \mapsto\left(a^{1}, \ldots, a^{n}, 0, \ldots, 0\right) .
$$

Then it is easy to check that

$$
\left.T_{p} f \cdot \frac{\partial}{\partial x^{i}}\right|_{p}=\left.\frac{\partial}{\partial y^{i}}\right|_{f(p)}
$$

So if

$$
Y=\sum Y^{i} \frac{\partial}{\partial y^{i}}
$$

for smooth functions $Y^{i}$, then

$$
X=\sum X^{i} \frac{\partial}{\partial x^{i}},
$$

where we must have $X^{i}=Y^{i} \circ f$. This shows that the $X^{i}$ are smooth and hence the field $X$ is smooth.

We apply the previous lemma to the situation where $S$ is an immersed submanifold of $M$ and $\iota: S \hookrightarrow M$ is the inclusion. In this case, if $Y$ is a vector field on $M$ which is tangent to $S$, then there is a unique smooth vector field on $S$, denoted $\left.Y\right|_{S}$, such that $\left.Y\right|_{S}(q)=Y(q)$ for all $q \in S$. The vector field $\left.Y\right|_{S}$ is called the restriction of $Y$ to $S$ and it is $\iota$-related to $Y$.

Theorem 11.9 (Local Frobenius theorem). A distribution is involutive if and only if it is completely integrable and if and only if it is integrable.

Proof. First, completely integrable clearly implies integrable. Now assume that $E \rightarrow M$ is a rank $k$ distribution which is integrable and let $\operatorname{dim}(M)=n$. Let $X$ and $Y$ be smooth vector fields defined on $U$ which both lie in the\\
distribution. Let $p \in U$ and let $S$ be an integral manifold that contains $p$. Then $X_{q}, Y_{q} \in T_{q} S$ for all $q \in U \cap S$. Thus the vector fields are tangent to $U \cap S$ and hence if $\iota: U \cap S \hookrightarrow U$ is the inclusion map, then $X$ and $Y$ are $\iota$-related to their restrictions to $U \cap S$ as in the previous lemma. This means that $\left[\left.X\right|_{U \cap S},\left.Y\right|_{U \cap S}\right]$ is $\iota$-related to $[X, Y]$ and hence $[X, Y]_{q} \in T_{q} S=E_{q}$ for all $q \in U \cap S$ and in particular $[X, Y]_{p} \in E_{p}$. Since $p$ was an arbitrary element of $U$, we conclude that the distribution is involutive.

Now suppose that $E \rightarrow M$ is involutive. Pick a point $p_{0} \in M$ and let $(U, \mathrm{x})$ be a chart centered at $p_{0}$. By shrinking $U$ if necessary we can find fields $X_{1}, \ldots, X_{k}$ defined on $U$ that span the distribution. Notice that if we have a chart of the type produced by Theorem 2.113, then we can easily arrange that the image of the chart is of the form $V \times W \subset \mathbb{R}^{k} \times \mathbb{R}^{n-k}$, and it is then easy to see that we have exactly the kind of chart needed (see the comments immediately following the proof of Theorem 2.113). Thus the strategy is to replace the frame field $X_{1}, \ldots, X_{k}$ by a commuting frame field.

Construct a chart $(U, \mathrm{x})$ centered at $p_{0}$ such that $X_{j}\left(p_{0}\right)=\left.\frac{\partial}{\partial x^{j}}\right|_{p_{0}}$ for $1 \leq j \leq k$. Let $\pi: \mathbb{R}^{n} \rightarrow \mathbb{R}^{k}$ be the map $\left(a^{1}, \ldots, a^{n}\right) \mapsto\left(a^{1}, \ldots, a^{k}\right)$ and let $f:=\pi \circ \mathrm{x}$. We have

$$
\left(T_{p_{0}} f\right)\left(X_{j}\left(p_{0}\right)\right)=\left.T_{p_{0}} f \frac{\partial}{\partial x^{j}}\right|_{p_{0}}=\left.\frac{\partial}{\partial t^{j}}\right|_{0},
$$

where $\left(t^{1}, \ldots, t^{k}\right)$ are standard coordinates on $\mathbb{R}^{k}$. Then $T_{p_{0}} f: T_{p_{0}} M \rightarrow T_{0} \mathbb{R}^{k}$ maps $E_{p_{0}}$ isomorphically onto $T_{0} \mathbb{R}^{k}$. It follows that $T_{p} f$ maps $E_{p}$ isomorphically onto $T_{f(p)} \mathbb{R}^{k}$ for all $p$ in some neighborhood of $p_{0}$. Hence, for each such $p$, there are unique vectors $Y_{1}(p), \ldots, Y_{k}(p)$ in $E_{p}$ with $T_{p} f\left(Y_{j}(p)\right) =\left.\frac{\partial}{\partial t^{t}}\right|_{f(p)}$. Thus we get vector fields $Y_{1}, \ldots, Y_{k}$ whose smoothness we leave to the reader to check. By shrinking $U$ if necessary, we may assume that these are defined on $U$. It is clear that $Y_{1}, \ldots, Y_{k}$ form a local frame for $E$ and $Y_{j}$ is $f$-related to $\frac{\partial}{\partial t^{j}}$. Thus $\left[Y_{i}, Y_{j}\right]$ is $f$-related to $\left[\frac{\partial}{\partial t^{i}}, \frac{\partial}{\partial t^{j}}\right]$ and in particular

$$
T_{p} f\left(\left[Y_{i}, Y_{j}\right]_{p}\right)=\left[\frac{\partial}{\partial t^{i}}, \frac{\partial}{\partial t^{j}}\right]_{0}=0
$$

Since the distribution $E$ is involutive, we have $\left[Y_{i}, Y_{j}\right]_{p} \in E_{p}$ for each $p$. Then since $T_{p} f$ is injective on $E_{p}$, we see that $\left[Y_{i}, Y_{j}\right]_{p}=0$ for all $p$. Thus we arrive at a local frame field $Y_{1}, \ldots, Y_{k}$ for $E$ with $\left[Y_{i}, Y_{j}\right]=0$.

By Theorem 2.113 and the comments above, we are done.

\subsection*{Differential Forms and Integrability}
One can use differential forms to effectively discuss distributions and involutivity. What is interesting is the way the exterior derivative comes into play.

The first thing to notice is that a smooth distribution is locally defined by 1 -forms. We need a bit of creative terminology:\\
Definition 11.10. If $\alpha^{1}, \ldots, \alpha^{j}$ are 1 -forms on an $n$-manifold $M$ and $S$ is a subset of $M$, then we say that $\alpha^{1}, \ldots, \alpha^{j}$ (defined at least on some open set containing $S$ ) are independent on $S$ if $\left\{\alpha^{1}(p), \ldots, \alpha^{j}(p)\right\}$ is an independent set in $T_{p}^{*} M$ for each $p \in S$.\\
Proposition 11.11. Let $E \rightarrow M$ be a rank $k$ distribution on an $n$-manifold $M$. Every $p \in M$ has an open neighborhood $U$ and 1 -forms $\alpha^{1}, \ldots, \alpha^{n-k}$ independent on $U$ such that

$$
E_{q}=\operatorname{Ker} \alpha^{1}(q) \cap \cdots \cap \operatorname{Ker} \alpha^{n-k}(q) \text { for all } q \in U
$$

Proof. Let $p \in M$ and let $X_{1}, \ldots, X_{k}$ span $E$ on a neighborhood $O$ of $p$. We may extend to a frame field $X_{1}, \ldots, X_{k}, \ldots, X_{n}$ defined on some possibly smaller neighborhood $U$ of $p$. Indeed, choose any other frame field $Y_{1}, \ldots, Y_{n}$ and then rearrange indices so that $X_{1}(p), \ldots, X_{k}(p), Y_{k+1}(p), \ldots, Y_{n}(p)$ are linearly independent. The vectors $X_{1}(q), \ldots, X_{k}(q), Y_{k+1}(q), \ldots, Y_{n}(q)$ must be independent for all $q$ in some neighborhood of $p$, which we may as well take to be $O$. Indeed, $q \mapsto X_{1}(q) \wedge \cdots \wedge X_{k}(q) \wedge Y_{k+1}(q) \wedge \cdots \wedge Y_{n}(q)$ is nonzero at $p$ and so remains nonzero in a neighborhood. In any case, let $X_{1}, \ldots, X_{k}, \ldots, X_{n}$ be our extended local frame and let $\theta^{1}, \ldots, \theta^{n}$ be the dual frame. Then it is easy to check that $v \in T_{q} M$ is in $E_{q}$ if and only if $\theta_{q}^{k+1}(v)=\cdots=\theta_{q}^{n}(v)=0$. Thus we just let ( $\alpha^{1}, \ldots, \alpha^{n-k}$ ) := $\left(\theta^{k+1}, \ldots, \theta^{n}\right)$.

The forms described in the previous proposition are sometimes called local defining forms for the distribution.

The set of all differential forms vanishing on a distribution has a convenient algebraic structure. Recall that $\Omega(M)$ is a ring under the wedge product. If $\mathcal{J}$ is the ideal in $\Omega(M)$ and $U \subset M$ is open, then $\left.\mathcal{J}\right|_{U}$ is the set of restrictions to $U$ of elements of $\mathcal{J}$ and is an ideal in $\Omega(U)$ since the wedge product is natural with respect to restrictions (that is, with respect to pull-back by inclusion maps).

Definition 11.12. Let $\mathcal{J}$ be an ideal in $\Omega(M)$; we say that $\mathcal{J}$ is locally generated by $n-k$ smooth 1 -forms if there is an open cover of $M$ such that for each open $U$ from the cover, there are 1 -forms $\alpha^{1}, \ldots, \alpha^{n-k}$ independent on $U$ such that $\left.\mathcal{J}\right|_{U}$ is an ideal of $\Omega(U)$ generated by $\alpha^{1}, \ldots, \alpha^{n-k}$.\\
Exercise 11.13. Show that if $\mathcal{J}$ is locally generated by $n-k$ smooth 1 forms as above, then it defines a unique smooth rank $k$ distribution $E$ on $M$ such that for each $U$ in the cover, and corresponding $\alpha^{1}, \ldots, \alpha^{n-k}$, we have

$$
E_{q}=\operatorname{Ker} \alpha^{1}(q) \cap \cdots \cap \operatorname{Ker} \alpha^{n-k}(q) \text { for all } q \in U .
$$

Definition 11.14. Let $E \rightarrow M$ be a rank $k$ distribution on an $n$-manifold $M$. Let $U$ be open in $M$. A $j$-form $\omega \in \Omega^{j}(U)$ is said to annihilate the distribution on $U$ if for each $p \in U$

$$
\omega_{p}\left(v_{1}, \ldots, v_{j}\right)=0 \text { whenever } v_{1}, \ldots, v_{j} \in E_{p} .
$$

An element $\omega \in \Omega(U)$ is said to annihilate the distribution if each homogeneous component of $\omega$ annihilates the distribution. The subset of all elements of $\Omega(U)$ that annihilate the distribution $E \rightarrow M$ is denoted $\mathcal{I}(U)$.

It is easy to check that $\mathcal{I}(U)$ is an ideal in $\Omega(U)$. Clearly, local defining forms for a distribution annihilate the distribution.

Lemma 11.15. The following conditions on a distribution $E \rightarrow M$ are equivalent:\\
(i) $E \rightarrow M$ is an involutive distribution.\\
(ii) For every open $U \subset M$, and 1 -form $\omega \in \Omega^{1}(U)$ that annihilates the distribution, the form $d \omega$ also annihilates the distribution on $U$.

Proof. The proof that these two conditions are equivalent follows easily from the formula

$$
d \omega(X, Y)=X(\omega(Y))-Y \omega(X)-\omega([X, Y])
$$

Suppose that (i) holds. If $\omega$ annihilates $E$ and $X, Y$ lie in $E$, then the above formula becomes

$$
d \omega(X, Y)=-\omega([X, Y])
$$

which shows that $d \omega$ vanishes on $E$ since $[X, Y] \in E$ by condition (i). Conversely, suppose that (ii) holds and that $X, Y \in E$. Let $\alpha^{1}, \ldots, \alpha^{n-k}$ be local defining 1 -forms. Then each form $d \alpha^{i}$ annihilates the distribution and also $\alpha^{j}(X)=\alpha^{j}(Y)=0$ for all $j$. Using equation (11.3) again, we have $\alpha^{i}([X, Y])=-d \alpha^{i}(X, Y)=0$ for all $i$ so that $[X, Y] \in E$.

Lemma 11.16. Let $E \rightarrow M$ be a rank $k$ distribution on the $n$-manifold $M$. $A j$-form $\omega$ annihilates $E$ if and only if, whenever $\alpha^{1}, \ldots, \alpha^{n-k}$ are local defining 1 -forms on some open $U$, we have

$$
\left.\omega\right|_{U}=\sum_{i=1}^{n-k} \alpha^{i} \wedge \beta^{i}
$$

for some $\beta^{1}, \ldots, \beta^{n-k} \in \Omega^{j-1}(U)$.\\
Proof. Recalling the definition of exterior product makes it clear that if $\omega$ satisfies (11.4), then it annihilates $E$ on $U$. If $\omega$ satisfies such an equation near each point for some generating 1 -forms, then $\omega$ annihilates $E$ on all of $M$.

Suppose that $\alpha^{1}, \ldots, \alpha^{n-k}$ are local defining forms for the distribution defined on an open set $U$ and let $p \in U$. Then on some neighborhood $V$ of $p$ we can extend the coframe field $\alpha^{1}, \ldots, \alpha^{n-k}$ to $\alpha^{1}, \ldots, \alpha^{n}$ in such a way that $\alpha^{1}, \ldots, \alpha^{n}$ is independent on $V$. Define $X_{1}, \ldots, X_{n}$ as dual to the frame $\left.\alpha^{1}\right|_{V}, \ldots,\left.\alpha^{n}\right|_{V}$ so that $X_{n-k+1}, \ldots, X_{n}$ span $E$ on $V$. Now for any $\omega \in \Omega^{j}(U)$, we have

$$
\left.\omega\right|_{V}=\left.\left.\sum_{i_{1}<\cdots<i_{j}} \omega_{i_{1} \cdots i_{j}} \alpha^{i_{1}}\right|_{V} \wedge \cdots \wedge \alpha^{i_{j}}\right|_{V}
$$

where $\omega_{i_{1} \cdots i_{j}}=\omega\left(X_{i_{1}}, \ldots, X_{i_{j}}\right)$ on $V$. Thus $\omega$ annihilates $E$ on $V$ if and only if $\omega\left(X_{i_{1}}, \ldots, X_{i_{j}}\right)=0$ on $V$ whenever $n-k+1 \leq i_{1}<\cdots<i_{j} \leq n$. So if $\omega$ annihilates $E$, then in the expansion of $\omega$ we need only include those $\omega_{i_{1} \cdots i_{j}}$ such that $i_{1}<\cdots<i_{j}$ and such that at least one index is less than or equal to $n-k$. But in the latter case we would have $i_{1} \leq n-k$. Thus we can write

$$
\begin{aligned}
\left.\omega\right|_{V} & =\left.\left.\sum_{\vec{I}: i_{1} \leq n-k} \omega_{i_{1} \cdots i_{j}} \alpha^{i_{1}}\right|_{V} \wedge \cdots \wedge \alpha^{i_{j}}\right|_{V} \\
& =\left.\sum_{i=1}^{n-k} \alpha^{i}\right|_{V} \wedge\left(\left.\left.\sum_{i_{2}<\cdots<i_{j}} \omega_{i i_{2} \cdots i_{j}} \alpha^{i_{2}}\right|_{V} \wedge \cdots \wedge \alpha^{i_{j}}\right|_{V}\right) \\
& :=\left.\sum_{i=1}^{n-k} \alpha^{i}\right|_{V} \wedge \beta_{V}^{i}
\end{aligned}
$$

This expression holds on $V \subset U$, but there is a cover of $U$ by such $V$ on which similar expressions hold. Denote this cover by $\mathcal{V}$. By using a partition of unity $\left\{\phi_{V}: V \in \mathcal{V}\right\}$ subordinate to this cover we can patch these expressions together. Notice that for any $V \in \mathcal{V}$ we have $\phi_{V}\left(\left.\alpha^{i}\right|_{V}\right)=\phi_{V} \alpha^{i}$ for $1 \leq i \leq n-k$, so

$$
\begin{aligned}
\left.\omega\right|_{U} & =\left.\sum_{V \in \mathcal{V}} \phi_{V} \omega\right|_{V}=\left.\sum_{V \in \mathcal{V}} \phi_{V} \sum_{i=1}^{n-k} \alpha^{i}\right|_{V} \wedge \beta_{V}^{i} \\
& =\sum_{V \in \mathcal{V}} \sum_{i=1}^{n-k}\left(\left.\phi_{V} \alpha^{i}\right|_{V}\right) \wedge \beta_{V}^{i} \\
& =\sum_{V \in \mathcal{V}} \sum_{i=1}^{n-k} \phi_{V} \alpha^{i} \wedge \beta_{V}^{i}=\sum_{i=1}^{n-k} \alpha^{i} \wedge \sum_{V \in \mathcal{V}} \phi_{V} \beta_{V}^{i} \\
& =\sum_{i=1}^{n-k} \alpha^{i} \wedge \beta^{i}
\end{aligned}
$$

for $\beta^{i}:=\sum_{V \in \mathcal{V}} \phi_{V} \beta_{V}^{i}$.

Corollary 11.17. A distribution $E \rightarrow M$ is involutive if and only if for every local defining set of 1 -forms $\alpha^{1}, \ldots, \alpha^{n-k}$ we have

$$
d \alpha^{i}=\sum_{j=1}^{n-k} \alpha^{j} \wedge \gamma_{j}^{i}
$$

for some 1 -forms $\gamma_{j}^{i}$ with $1 \leq i, j \leq n-k$.\\
Proof. Just use the above theorem together with Lemma 11.15.\\
Theorem 11.18. A distribution $E \rightarrow M$ on $M$ is involutive (and hence completely integrable) if and only if $d \mathcal{I}(E) \subset \mathcal{I}(E)$.

The condition $d \mathcal{I}(E) \subset \mathcal{I}(E)$ is expressed by saying that $\mathcal{I}(E)$ is a differential ideal.

Proof. Suppose that $d \mathcal{I}(E) \subset \mathcal{I}(E)$. If $\omega$ is any 1-form that annihilates $E$ on an open set $U$, then pick an arbitrary point $p \in U$ and a smooth cut-off function $\phi$ with support in $U$ and equal to one on a neighborhood of $p$. Then $d(\phi \omega) \in \mathcal{I}(E)$. But $d(\phi \omega)(p)=d \omega(p)$, and since $p$ was arbitrary, we see that $d \omega$ annihilates $E$ on $U$. The open set $U$ was also arbitrary so we can apply Lemma 11.15.

Now suppose that $E \rightarrow M$ is involutive and choose $\eta \in \mathcal{I}(E)$. Without loss of generality we may assume that $\eta$ is homogeneous of degree $j$. Notice that the case $j=0$ is trivial, while the case $j=1$ is Lemma 11.15. For $j \geq 2$, consider a local defining set of 1 -forms $\alpha^{1}, \ldots, \alpha^{n-k}$ on an open set $U$. Since by assumption $\eta$ annihilates $E$, we can use Lemma 11.16 and Corollary 11.17 to write

$$
\begin{aligned}
\left.d \eta\right|_{U} & =d\left(\sum_{i=1}^{n-k} \alpha^{i} \wedge \beta^{i}\right) \\
& =\sum_{i=1}^{n-k} d \alpha^{i} \wedge \beta^{i}-\sum_{i=1}^{n-k} \alpha^{i} \wedge d \beta^{i} \\
& =\sum_{i=1}^{n-k} \sum_{j} \alpha^{j} \wedge \gamma_{j}^{i} \wedge \beta^{i}-\sum_{i=1}^{n-k} \alpha^{i} \wedge d \beta^{i}
\end{aligned}
$$

Now use Lemma 11.16 again to conclude that $d \eta$ annihilates $E$ and so is in $\mathcal{I}(E)$.

We close this section with a convenient application of vector-valued forms.

Proposition 11.19. Let $\omega$ be a 1 -form on an $n$-manifold with values in a finite-dimensional vector space V . Then $E_{x}=\operatorname{Ker} \omega_{x}$ defines a distribution\\
$E$ if $\operatorname{dim} E_{x}=k$ is constant for $x \in M$. This distribution is involutive and hence integrable if and only if $d \omega(X, Y)=0$ whenever $X, Y$ lie in $E$.

Proof. The first statement is clear once we choose a basis $\mathbf{e}_{1}, \ldots, \mathbf{e}_{n}$ for V for then we can write $\omega=\sum \mathbf{e}_{i} \omega^{i}$ for some linearly independent smooth 1 -forms $\omega^{1}, \ldots, \omega^{k}$ which clearly define the distribution.

Choose local spanning fields $X_{1}, \ldots, X_{k}$ defined on an open set $U$. Then $E$ is integrable on $U$ if and only if

$$
\left[X_{i}, X_{j}\right] \in \operatorname{span}\left\{X_{1}, \ldots, X_{k}\right\} \text { for } 1 \leq i, j \leq k,
$$

which is true if and only if $\omega\left(\left[X_{i}, X_{j}\right]\right)=0$ for $1 \leq i, j \leq k$. But as before,

$$
\begin{aligned}
d \omega\left(X_{i}, X_{j}\right) & =X_{i}\left(\omega\left(X_{j}\right)\right)-X_{j}\left(\omega\left(X_{i}\right)\right)-\omega\left(\left[X_{i}, X_{j}\right]\right) \\
& =-\omega\left(\left[X_{i}, X_{j}\right]\right)
\end{aligned}
$$

Thus $E$ is integrable on $U$ if and only if $d \omega\left(X_{i}, X_{j}\right)=0$ for $1 \leq i, j \leq k$. This implies the result.

\subsection*{Global Frobenius Theorem}
First notice that if ( $U, \mathrm{x}$ ) is a distinguished chart for a distribution on an $n$-manifold with $\mathrm{x}(U)=V \times W \subset \mathbb{R}^{k} \times \mathbb{R}^{n-k}$, then the plaques $\mathrm{x}^{-1}(V \times\{y\})$ for each $y \in W$ are embedded integral manifolds. If $O \subset V$ and $y \in W$, then $\mathrm{x}^{-1}(O \times\{y\})$ is an open subset of the plaque $\mathrm{x}^{-1}(V \times\{y\})$ (in its submanifold topology). Clearly, all open subsets of a plaque are of this form.

Proposition 11.20. Let $E \rightarrow M$ be an integrable distribution of rank $k$ on the $n$-manifold $M$. Let ( $U, \mathrm{x}$ ) be a fixed distinguished chart for $E$ where $\mathrm{x}(U)=V \times W \subset \mathbb{R}^{k} \times \mathbb{R}^{n-k}$. If $S$ is a connected integral manifold for $E$, then $S \cap U$ is a countable disjoint union of connected open subsets of the plaques of ( $U, \mathrm{x}$ ). These open connected subsets are open in $S$ and embedded in $M$.

Proof. Since the inclusion $\iota: S \hookrightarrow M$ is an immersion, the sets $S \cap U= \iota^{-1}(U)$ are open in $S$. The set $S \cap U$ is a disjoint union of connected components each of which is open in $S$. Since $S$ is second countable, this union must be a countable union. Let $C$ be a connected component of $S \cap U$. We show first that $C \subset \mathrm{x}^{-1}(V \times\{y\})$ for some $y \in W$. Since ( $U, \mathrm{x}$ ) is a distinguished chart, the 1 -forms $d x^{k+1}, \ldots, d x^{n}$ are local defining forms for the distribution and so $d x^{k+1}=\cdots=d x^{n}=0$ on $S \cap U$ and hence on $C$. But $C$ is connected, so $x^{k+1}, \ldots, x^{n}$ are all constant on $C$, which puts it in a single plaque, say $\mathrm{x}^{-1}(V \times\{y\})$.

We know that $C \hookrightarrow M$ is smooth. Since $\mathrm{x}^{-1}(V \times\{y\})$ is embedded in $M$, the inclusion $C \hookrightarrow \mathrm{x}^{-1}(V \times\{y\})$ is also smooth by Corollary 3.15 and is thus an injective immersion. Since $C$ and $\mathrm{x}^{-1}(V \times\{y\})$ are manifolds of the\\
same dimension, the inclusion $C \hookrightarrow \mathrm{x}^{-1}(V \times\{y\})$ is a local diffeomorphism and hence an open map. It follows that $C \hookrightarrow \mathrm{x}^{-1}(V \times\{y\})$ is an embedding. Since the inclusion $C \hookrightarrow M$ is the composition of embeddings $C \hookrightarrow \mathrm{x}^{-1}(V \times \{y\}) \hookrightarrow M$, it is itself an embedding.

If $E \rightarrow M$ is an integrable distribution of rank $k$ on the $n$-manifold $M$ and $p \in M$, then we would like to consider the union

$$
N=\bigcup_{S \in \mathcal{F}(p)} S
$$

where $\mathcal{F}(p)$ is the family of all connected integral manifolds containing $p$. We would like to say that $N$ is the maximal integral manifold that contains $p$, but we must show that $N$ is indeed an integral manifold. Let $\left\{\left(U_{\alpha}, \mathrm{x}_{\alpha}\right)\right\}_{\alpha \in A}$ be an atlas of distinguished charts for the distribution where $\mathrm{x}_{\alpha}\left(U_{\alpha}\right)= V_{\alpha} \times W_{\alpha} \subset \mathbb{R}^{k} \times \mathbb{R}^{n-k}$. On each plaque $\mathrm{x}^{-1}\left(V_{\alpha} \times\{y\}\right)$ of ( $U_{\alpha}, \mathrm{x}_{\alpha}$ ) which is in $N$, we define

$$
\operatorname{pr}_{1} \circ \mathrm{x}_{\alpha}: \mathrm{x}^{-1}\left(V_{\alpha} \times\{y\}\right) \rightarrow V_{\alpha} \subset \mathbb{R}^{k}
$$

and this map is taken as a chart on $N$ whose domain is a plaque. Charts obtained in this way will be called plaque charts and the family of all such charts provides $N$ with an atlas. The question of the smoothness of overlaps is routine and is left to the reader. Thus we have a smooth atlas on $N$ which induces a topology such that each plaque, and indeed, each member of $\mathcal{F}(p)$, is open. The topology induced by this smooth structure is often called the leaf topology. It is also evident that $N$ is connected since it is the union of connected sets containing a fixed point. We can also see directly that $N$ is path connected. Indeed, if $q$ is in $N$, then it must be in some $S_{1} \in \mathcal{F}(p)$. But also $p \in S_{1}$ by definition, and so since $S_{1}$ is connected (and hence smoothly path connected), there is a smooth curve connecting $p$ and $q$. The questions that remain are whether the topology is Hausdorff and whether it is paracompact. Once again the proof that $N$ is Hausdorff is easy and is left to the reader. The harder question is paracompactness. But this follows from the fact that $N$ clearly satisfies property $\mathrm{W}(k)$ of Definition 3.18, and so $N$ is a weakly embedded submanifold by Proposition 3.20. On the other hand, the part of the proof of that proposition that established paracompactness appealed to induced Riemannian metrics and material about topologies induced by metrics that we still have not covered. For this reason we offer a direct and traditional proof. First, since $N$ is connected, the goal is to prove that it is second countable. The inclusion $N \hookrightarrow M$ is continuous, so $N$ is contained in a connected component of $M$. Thus we may assume without loss of generality that $M$ is second countable, and so also that the atlas $\left\{\left(U_{\alpha}, \mathrm{x}_{\alpha}\right)\right\}_{\alpha \in A}$ of distinguished charts is countable. This latter fact is one of the key points. Notice that each plaque is a regular

\begin{figure}[h]
\begin{center}
  \includegraphics[scale=0.25]{2025_10_09_265d7e1a22c89f79c21eg-495}

\caption{Figure 11.2. Slice accessible from point $p$}
\end{center}
\end{figure}

submanifold, so each such plaque is itself second countable. Observe that by construction, if a set $N \cap U_{\alpha}$ intersects a plaque in $U_{\alpha}$, then that plaque is actually contained in $N \cap U_{\alpha}$. Indeed, if a plaque of $U_{\alpha}$ intersects $N$, then there is an element $S \in \mathcal{F}(p)$ such that $S$ meets the plaque in an open set in the topology of the plaque (Proposition 11.20). But then the union of this plaque with $S$ is another element of $\mathcal{F}(p)$ and so occurs in the union that defines $N$. This reduces the problem to that of showing that $N$ contains only a countable number of such plaques.

We consider all sequences of the form $\left(U_{\alpha(1)}, P_{1}\right), \ldots,\left(U_{\alpha(m)}, P_{m}\right)$, where $P_{i}$ is a plaque for a chart ( $U_{\alpha(i)}, \mathrm{x}_{\alpha(i)}$ ) taken from our countable atlas of distinguished charts, such that $P_{m}=P$ is a plaque contained in $U_{\alpha(m)} \cap N$,

$$
p \in P_{1}, P_{m}=P, P_{i} \subset N \text { and } P_{i} \cap P_{i+1} \neq \emptyset \text { for } i=1, \ldots, m-1 .
$$

Notice that $m$ is not fixed, but must be finite. Let us call such a sequence a finite connection from $p$ to the plaque $P$. We first show that for each plaque $P$ contained in some $U_{\alpha} \cap N$, there is at least one such finite connection from $p$ to $P$. To see this, consider a point $q \in U_{\alpha} \cap N$ and suppose that $q \in P$. Let $c:[0,1] \rightarrow N$ be a smooth curve with $c(0)=p$ and $c(1)=q$. Now $c([0,1])$ is a compact set, so there is a sequence $t_{1}<t_{2}<\cdots<t_{m}$ such that for each $i, c\left(\left[t_{i}, t_{i+1}\right]\right)$ is contained in a distinguished chart, say $\left(U_{\alpha(i)}, \mathrm{x}_{\alpha(i)}\right)$. Of course, each $c\left(\left[t_{i}, t_{i+1}\right]\right)$ is also connected and so must be contained in a connected component of $U_{\alpha(i)} \cap N$ and thus in some plaque which we call $P_{i}$. Clearly we have created a finite connection from $p$ to $P$.

Now consider the family $\mathcal{C}$ of all finite connections of $p$ to the plaques contained in the various $U_{\alpha} \cap N$. Each finite connection can be mapped to the plaque on which it terminates, and we have just seen that this map is surjective. Thus the set of plaques of the various $U_{\beta} \cap N$ that are finitely connectable to $p$ must have cardinality less than or equal to that of $\mathcal{C}$. It may already seem that since the atlas is countable, $\mathcal{C}$ must also be countable. However, while the set of all finite sequences $U_{\alpha(1)}, \ldots, U_{\alpha(m)}$ is certainly\\
countable, there is still the question of how many ways there are to choose the $P_{i}$ inside each $U_{\alpha(i)}$ to build a connection of $p$ to a plaque. However, since each plaque $P_{i}$ is certainly an integral manifold, Proposition 11.20 tells us that it can only intersect $U_{\alpha(i+1)}$ in a countable number of plaques. Thus for a given sequence $U_{\alpha(1)}, \ldots, U_{\alpha(m)}$ there is only a countable number of ways to choose the corresponding $P_{i}$. It follows that $N$ contains only a countable number of plaque charts and so is a second countable connected integral manifold for the distribution. It is clearly maximal by construction.

One final thing to notice is that unless the rank of the distribution is equal to the dimension of the manifold, $M$ has an uncountable number of maximal integral manifolds, and they partition $M$. We have now proved the following:

Theorem 11.21. Let $E \rightarrow M$ be an integrable distribution of rank $k$ on the $n$-manifold $M$. There is a set of maximal integral manifolds for $E$ which partition $N$. These are weakly embedded submanifolds.

There is another slightly different way of looking at the above theorem. Namely, the set of all plaque charts puts a new smooth structure on the set $M$. We have seen that the overlaps are only nonempty if the plaques are on the same integral manifold. This smooth structure and topology gives us a new manifold which we might denote by $M(E)$. It has the same underlying set as $M$, but has a finer topology. The connected components of $M(E)$ are exactly the maximal integral manifolds.

Proposition 11.22. If $\mathcal{L}_{1}$ and $\mathcal{L}_{2}$ are connected integral manifolds of an integrable distribution, then either $\mathcal{L}_{1} \cap \mathcal{L}_{2}$ is empty or it is an open subset of the maximal integral manifold containing any point of $\mathcal{L}_{1} \cap \mathcal{L}_{2}$. The set $\mathcal{L}_{1} \cap \mathcal{L}_{2}$ is also open in the integral manifolds $\mathcal{L}_{1}$ and $\mathcal{L}_{2}$.

Proof. Assume that $\mathcal{L}_{1} \cap \mathcal{L}_{2}$ is not empty and let $p \in \mathcal{L}_{1} \cap \mathcal{L}_{2}$. Then there is a distinguished chart ( $U, \mathrm{x}$ ) for the distribution containing $p$ such that both $\mathcal{L}_{1} \cap U$ and $\mathcal{L}_{2} \cap U$ are plaques containing $p$. Clearly, the plaques must be the same. Since plaques are open by definition in the maximal integral manifold, and since $p$ was an arbitrary point in $\mathcal{L}_{1} \cap \mathcal{L}_{2}$, it follows that $\mathcal{L}_{1} \cap \mathcal{L}_{2}$ is open in this maximal integral manifold. The last statement is obvious. \(\square\)

Corollary 11.23. Suppose that $M$ and $N$ are smooth manifolds and that there is an integrable distribution $E$ on $M \times N$. Suppose that for some connected open subset $U$ of $M$ there are smooth functions $f_{1}: U \rightarrow N$ and $f_{2}: U \rightarrow N$ such that the graphs of $f_{1}$ and $f_{2}$ are integral manifolds. Then these graphs are either disjoint or equal. In the latter case, $f_{1}=f_{2}$.

Proof. Let $\Gamma_{f_{1}}$ and $\Gamma_{f_{2}}$ the graphs and suppose that $\Gamma_{f_{1}} \cap \Gamma_{f_{2}}$ is not empty. Let $\left(p_{0}, f_{1}\left(p_{0}\right)\right)=\left(p_{0}, f_{2}\left(p_{0}\right)\right) \in \Gamma_{f_{1}} \cap \Gamma_{f_{2}}$ and consider the set

$$
S=\left\{p \in U: f_{1}(p)=f_{2}(p)\right\}
$$

Then $S$ is clearly closed. We show that it is also open. Pick $p \in S$. Then $\Gamma_{f_{1}}$ and $\Gamma_{f_{2}}$ intersect at the point $\left(p, f_{1}(p)\right)=\left(p, f_{2}(p)\right)$. But since $\Gamma_{f_{1}}$ and $\Gamma_{f_{2}}$ are integral manifolds of an integrable distribution on $M \times N$, the previous proposition tells us that $\Gamma_{f_{1}} \cap \Gamma_{f_{2}}$ is open in $\Gamma_{f_{1}}$, and so, since $\mathrm{pr}_{1}: U \times N \rightarrow U$ restricts to a diffeomorphism on $\Gamma_{f_{1}}$ (as for all graphs of smooth functions), we see that $S=\operatorname{pr}_{1}\left(\Gamma_{f_{1}} \cap \Gamma_{f_{2}}\right)$ is an open set in $U$. Since $U$ is connected, we have $S=U$. Thus $f_{1}=f_{2}$ and the graphs are equal.

If we consider what we have discovered about integrable distributions, we arrive naturally at the following definition (and new terminology).

Definition 11.24. Let $M$ be an $n$-manifold. A $k$-dimensional foliation $\mathcal{F}_{M}$ of $M$ (or on $M$ ) is a partition of $M$ into a family of disjoint connected subsets $\left\{\mathcal{L}_{\alpha}\right\}_{\alpha \in A}$ such that for every $p \in M$, there is a chart ( $U, \mathrm{x}$ ) centered at $p$ of the form $\mathrm{x}: U \rightarrow V \times W \subset \mathbb{R}^{k} \times \mathbb{R}^{n-k}$ for connected $V$ and $W$ with the property that for each $\mathcal{L}_{\alpha}$ the connected components $\left(U \cap \mathcal{L}_{\alpha}\right)_{\beta}$ of $U \cap \mathcal{L}_{\alpha}$ are given by

$$
\varphi\left(\left(U \cap \mathcal{L}_{\alpha}\right)_{\beta}\right)=V \times\left\{c_{\alpha, \beta}\right\}
$$

where $c_{\alpha, \beta} \in W \subset F$ are constants. These charts are called distinguished charts for the foliation or foliation charts. The connected sets $\mathcal{L}_{\alpha}$ are called the leaves of the foliation while for a given chart as above, the connected components $\left(U \cap \mathcal{L}_{\alpha}\right)_{\beta}$ are called plaques.

It is easy to see that if $\mathcal{F}_{M}$ is a foliation as above, then there is a unique integrable distribution $E \rightarrow M$ on $M$ such that if $p$ is in the domain of a foliation chart ( $U, \mathrm{x}$ ), then

$$
E_{p}=\left.\left.\operatorname{Ker} d x\right|_{p} ^{k+1} \cap \cdots \cap \operatorname{Ker} d x\right|_{p} ^{n}
$$

Furthermore the leaves of the foliation are exactly the maximal integral manifolds of the distribution. We call this the distribution generated by the foliation. Combining this observation with Theorem 11.21 we obtain

Theorem 11.25. Every integrable rank $k$ distribution gives rise to a unique $k$-dimensional foliation whose leaves are the maximal integral manifolds. Conversely, a foliation generates an integrable distribution whose maximal integral manifolds are exactly the leaves of the original foliation.

Thus the theory of foliations is essentially the study of integrable distributions. Now we see that the leaves are weakly embedded submanifolds. The connected components $\left(U \cap \mathcal{L}_{\alpha}\right)_{\beta}$ of $U \cap \mathcal{L}_{\alpha}$ are of the form $C_{x}\left(U \cap \mathcal{L}_{\alpha}\right)$\\
for some $x \in \mathcal{L}_{\alpha}$ (recall Definition 3.17). An important point anticipated by our concern above is that a fixed leaf $\mathcal{L}_{\alpha}$ may intersect a given chart domain $U$ in many, even a countably infinite number of disjoint connected pieces no matter how small $U$ is taken to be. In fact, it may be that $U \cap \mathcal{L}_{\alpha}$ is dense in $U$. On the other hand, each $\mathcal{L}_{\alpha}$ is connected by definition. The usual first example of this behavior is given by the irrationally foliated torus. Here we take $M=T^{2}:=S^{1} \times S^{1}$ and let the leaves be given as the image of the immersions $\iota_{a}: t \mapsto\left(e^{i a \pi t}, e^{i \pi t}\right)$, where $a$ is a real number. If $a$ is irrational, then the image $\iota_{a}(\mathbb{R})$ is a (connected) dense subset of $S^{1} \times S^{1}$. On the other hand, even in this case there are infinitely many distinct leaves.

It may seem that a foliated manifold is just a special manifold, but from one point of view, a foliation is a generalization of a manifold. For instance, we can think of a manifold $M$ as a foliation where the points are the leaves. This is called the discrete foliation on $M$. At the other extreme a manifold may be thought of as a foliation with a single leaf $\mathcal{L}=M$ (the trivial foliation). We also have the following special cases:

Example 11.26. On a product manifold, say $M \times N$, we have two complementary foliations:

$$
\{\{p\} \times N\}_{p \in M}
$$

and

$$
\{M \times\{q\}\}_{q \in N}
$$

Example 11.27. Given any submersion $f: M \rightarrow N$, the level sets $\left\{f^{-1}(q)\right\}_{q \in N}$ form the leaves of a foliation.\\
Example 11.28. The fibers of any vector bundle foliate the total space.\\
Example 11.29 (Reeb foliation). Consider the strip in the plane given by $\{(x, y):|x| \leq 1\}$. For $a \in \mathbb{R} \cup\{ \pm \infty\}$, we form leaves $\mathcal{L}_{a}$ as follows:

$$
\begin{aligned}
\mathcal{L}_{a} & :=\{(x, a+f(x)):|x| \leq 1\} \text { for } a \in \mathbb{R}, \\
\mathcal{L}_{ \pm \infty} & :=\{( \pm 1, y):|y| \leq 1\}
\end{aligned}
$$

where $f(x):=\exp \left(\frac{x^{2}}{1-x^{2}}\right)-1$. By rotating this symmetric foliation about the $y$ axis we obtain a foliation of the solid cylinder. This foliation is such that translation of the solid cylinder $C$ in the $y$ direction maps leaves diffeomorphically onto leaves, and so we may let $\mathbb{Z}$ act on $C$ by $(x, y, z) \mapsto(x, y+n, z)$ and then $C / \mathbb{Z}$ is a solid torus with a foliation called the Reeb foliation.

Example 11.30. The one point compactification of $\mathbb{R}^{3}$ is homeomorphic to $S^{3} \subset \mathbb{R}^{4}$. Thus $S^{3}-\{p\} \cong \mathbb{R}^{3}$ and so there is a copy of the Reeb foliated solid torus inside $S^{3}$. The complement of a solid torus in $S^{3}$ is another solid torus. It is possible to put another Reeb foliation on this complement and thus foliate all of $S^{3}$. The only compact leaf is the torus that is the common boundary of the two complementary solid tori.

Theorem 11.31. Let $\mathcal{F}_{M}$ be a foliation on $M$ and $\mathcal{L}$ a leaf of the foliation. Then for any smooth map $f: N \rightarrow M$ with $f(N) \subset \mathcal{L}$ the map $f: N \rightarrow \mathcal{L}$ is smooth.

Proof. Just use Theorem 3.14.

\subsection*{Applications to Lie Groups}
First we fulfill the promise of proving Theorem 5.66, which we repeat here for convenience.

Theorem 11.32. Let $G$ be a Lie group with Lie algebra $\mathfrak{g}$. If $\mathfrak{h}$ is a Lie subalgebra of $\mathfrak{g}$, then there is a unique connected Lie subgroup $H$ of $G$ whose Lie algebra is $\mathfrak{h}$.

Proof. For $a \in G$, we let $\Delta_{a}$ denote the subspace of $T_{a} G$ that is the set of all $X(a)$ for left invariant vector fields $X$ with $X(e) \in \mathfrak{h}$. Thus $v_{a} \in \Delta_{a}$ if and only if $v_{a}=T_{e} L_{a} \cdot v$ for some $v \in \mathfrak{h}$. If $X_{1}, \ldots, X_{k}$ are left invariant fields such that $X_{1}(e), \ldots, X_{k}(e)$ is a basis for $\mathfrak{h}$, then $X_{1}, \ldots, X_{k}$ span the distribution which is therefore smooth. Since $\mathfrak{h}$ is a subalgebra, it follows that $\Delta: a \mapsto \Delta_{a}$ is an integrable distribution. Let $H$ be the maximal (connected) integral manifold containing $e$. Note that for any $b \in G$, we have $T_{e} L_{b}\left(\Delta_{a}\right)=\Delta_{b a}$ so that $T L_{b}: T G \rightarrow T G$ leaves the distribution invariant. Thus $L_{b}$ induces a permutation of the set of maximal integral manifolds and takes the maximal integral manifold through $a$ diffeomorphically to that through $b a$. Thus if $h \in H$, then $L_{h^{-1}}$ takes $H$ to the maximal integral manifold containing $e$, which is just $H$. In other words, $L_{h^{-1}}(H)=H$. This shows that $H$ is a subgroup. It remains to show that the multiplication map $\mu: H \times H \rightarrow H$ is smooth. But the multiplication map into $G, H \times H \rightarrow G$, is smooth and so by Theorem 11.31 the map $H \times H \rightarrow H$ is smooth.

Theorem 11.33. Let $G$ and $H$ be Lie groups with respective Lie algebras $\mathfrak{g}$ and $\mathfrak{h}$. If $h: \mathfrak{g} \rightarrow \mathfrak{h}$ is a Lie algebra homomorphism, then there is a neighborhood $U$ of the identity $e \in G$ and a smooth map $f: U \rightarrow H$ such that

$$
f(x y)=f(x) f(y)
$$

whenever $x, y \in U$ and $x y \in U$, and such that

$$
T_{e} f \cdot v=h(v)
$$

for every $v \in \mathfrak{g}$.\\
Proof. Let $\mathfrak{k} \subset \mathfrak{g} \times \mathfrak{h}$ be defined by

$$
\mathfrak{k}:=\{(v, h(v)): v \in \mathfrak{g}\} .
$$

The fact that $h$ is a homomorphism implies that $\mathfrak{k}$ is a Lie subalgebra of $\mathfrak{g} \times \mathfrak{h}$. Thus by Theorem 11.32, there is a connected Lie subgroup $K$ of $G \times H$ with Lie algebra $\mathfrak{k}$. Now let $\iota: K \hookrightarrow G \times H$ be inclusion and define a homomorphism $\rho: K \rightarrow G$ by

$$
\rho:=\operatorname{pr}_{1} \circ \iota
$$

If $v \in \mathfrak{g}$, then

$$
T \rho \cdot(v, h(v))=v
$$

and this means that $T \rho: T_{(e, e)} K \rightarrow T_{e} G$ is a linear isomorphism. Thus by the inverse mapping theorem, there is a neighborhood $V$ of $(e, e) \in K$ such that $\left.\rho\right|_{V}$ is a diffeomorphism onto an open neighborhood $U$ of $e \in G$. Define the homomorphism $\psi: K \rightarrow H$ by

$$
\psi:=\operatorname{pr}_{2} \circ \iota
$$

where $\mathrm{pr}_{2}: G \times H \rightarrow H$ is the second factor projection. Notice that $T_{e} \psi \cdot (v, h(v))=h(v)$. Now let

$$
f:=\left.\psi \circ \rho\right|_{V} ^{-1}
$$

A straightforward diagram chase argument shows that $f(x y)=f(x) f(y)$ if $x, y \in U$ and $x y \in U$.

If $v \in \mathfrak{g}$, then $T \rho(v, h(v))=v$ implies that $T\left(\left.\rho\right|_{V} ^{-1}\right) \cdot v=(v, h(v))$ so

$$
T f \cdot v=T \psi \circ T\left(\left.\rho\right|_{V} ^{-1}\right) \cdot v=T \psi \cdot(v, h(v))=h(v)
$$

Theorem 11.34. If $f_{1}: G \rightarrow H$ and $f_{2}: G \rightarrow H$ are Lie group homomorphisms with $d f_{1}=d f_{2}: \mathfrak{g} \rightarrow \mathfrak{h}$ and $G$ is connected, then $f_{1}=f_{2}$.

Proof. Let $h:=d f_{1}=d f_{2}$ and define $\mathfrak{k}:=\{(v, h(v)): v \in \mathfrak{g}\}$ and $K \subset G \times H$ as in the proof of the previous theorem. Now define $\theta: G \rightarrow G \times H$ by $\theta(x):=\left(x, f_{1}(x)\right)$. The image of $\theta$ is a subgroup $K_{1} \subset G \times H$. For $v \in \mathfrak{g}$, we have $T_{e} \theta \cdot v=(v, h(v))$ so the Lie algebra of $K_{1}$ must be $\mathfrak{k}$. Since $G$ is connected, we must have $K=K_{1}$ which implies that $f_{1}=f$ on $U$, where $f$ and $U$ are constructed from $h$ as in the last theorem. But equally, $f_{2}=h$ on $U$, and so by Proposition 5.40 we have $f_{1}=f_{2}$.

We say that two Lie groups $G$ and $H$ are locally isomorphic if there is a diffeomorphism $f$ from a neighborhood $U$ of the identity of $G$ onto a neighborhood $V$ of the identity of $H$ such that $f(x y)=f(x) f(y)$ whenever $x, y$ and $x y$ are contained in $U$.

Corollary 11.35. The following assertions hold:\\
(i) Two Lie groups with isomorphic Lie algebras are locally isomorphic.\\
(ii) A connected Lie group with abelian Lie algebra is abelian.

Proof. (i) Let $h: \mathfrak{g} \rightarrow \mathfrak{h}$ be a Lie algebra isomorphism. Then if $f$ is the map constructed in Theorem 11.33, then $f$ is a diffeomorphism on some possibly smaller neighborhood of the identity since $T_{e} f=h$ is an isomorphism.\\
(ii) By (i) a connected Lie group $G$ of dimension $n$ must be locally isomorphic to the (additive) abelian Lie group $\mathbb{R}^{n}$. But a neighborhood of the identity generates the whole group, and so $G$ is abelian.

\subsection*{Fundamental Theorem of Surface Theory}
In this section we state and outline the proof of a fundamental theorem concerning the existence of surfaces with prescribed first and second fundamental form. Our proof of the main theorem follows [Pa2]. To begin with, we need a few results about certain systems of partial differential equations. The first is equivalent to the local Frobenius theorem. For an open set $U \subset \mathbb{R}^{k} \times \mathbb{R}^{m}$, denote standard coordinates by $(x, z)=\left(x^{1}, \ldots, x^{k}, z^{1}, \ldots, z^{m}\right)$.

Proposition 11.36. Let $U$ be an open set in $\mathbb{R}^{k} \times \mathbb{R}^{m}$ and let ( $A_{j}^{i}$ ) be an $m \times k$ matrix of smooth functions on $U$. Then the following assertions are equivalent:\\
(i) For every $\left(x_{0}, z_{0}\right) \in U$, there is a neighborhood $V$ of $x_{0}$ in $\mathbb{R}^{k}$ and a unique smooth map $f: V \rightarrow \mathbb{R}^{m}$ with $f\left(x_{0}\right)=z_{0}$ such that

$$
\frac{\partial f^{i}}{\partial x^{j}}(x)=A_{j}^{i}(x, f(x)) \text { for all } i, j .
$$

(ii) The functions $A_{j}^{i}$ satisfy the following system of equations on $U$ :

$$
\frac{\partial A_{j}^{i}}{\partial x^{k}}+\sum_{l} A_{k}^{l} \frac{\partial A_{j}^{i}}{\partial z^{l}}=\frac{\partial A_{k}^{i}}{\partial x^{j}}+\sum_{l} A_{j}^{l} \frac{\partial A_{k}^{i}}{\partial z^{l}} \text { for all } i, j, k
$$

Proof. If (i) is true, then we obtain (ii) by equality of mixed partials of the $f^{i}$ and the chain rule (see the comments following the proof).

Conversely, consider the vector fields on $U$ defined by

$$
X_{j}:=\left.\frac{\partial}{\partial x^{j}}\right|_{p}+\left.\sum_{r} A_{j}^{r}(p) \frac{\partial}{\partial z^{r}}\right|_{p} .
$$

A bit of linear algebra tells us that these are everywhere independent. Let $E \rightarrow U$ be the distribution spanned by these fields. A straightforward check using (11.6) shows that

$$
\left[X_{i}, X_{j}\right]=0
$$

so there is an integral manifold through each $p$. Let $N$ be the integral manifold through $\left(x_{0}, z_{0}\right)$. Using the last $m$ coordinate functions of some distinguished chart, we obtain a map $\Phi: U^{\prime} \rightarrow \mathbb{R}^{m}$ for some connected open $U^{\prime} \subset U$ so that the level sets of $\Phi$ are integral manifolds. The tangent map\\
$T_{p} \Phi$ has kernel $E_{p}$ at each $p \in N$, and since $\left.\frac{\partial}{\partial z^{j}}\right|_{p}$ never lies in $E$, we see that

$$
\frac{\partial \Phi}{\partial z^{j}} \neq 0 \text { on } N \cap U^{\prime} \text { for all } j .
$$

In particular, this holds at ( $x_{0}, z_{0}$ ), and the implicit mapping theorem tells us that a neighborhood of $\left(x_{0}, z_{0}\right)$ in a plaque of $N$ is the graph of a function $f: V \rightarrow \mathbb{R}^{m}$ with $f\left(x_{0}\right)=z_{0}$. Define a function $F: V \rightarrow \mathbb{R}^{k} \times \mathbb{R}^{m}$ by $F(x):=(x, f(x))$. Writing $p=F(x)$, we see that for each $i$ the vector

$$
\left.T F \cdot \frac{\partial}{\partial x^{i}}\right|_{x}=\left.\frac{\partial}{\partial x^{i}}\right|_{p}+\left.\sum_{r} \frac{\partial f^{r}}{\partial x^{i}}(x) \frac{\partial}{\partial z^{r}}\right|_{p}
$$

is a linear combination of vectors $X_{j}$ defined above:

$$
\begin{aligned}
& \left.\frac{\partial}{\partial x^{i}}\right|_{f(x)}+\left.\sum_{r} \frac{\partial f^{r}}{\partial x^{i}}(x) \frac{\partial}{\partial z^{r}}\right|_{f(x)} \\
& \quad=\sum_{s=1}^{k} c_{i}^{s}\left(\left.\frac{\partial}{\partial x^{s}}\right|_{(x, f(x))}+\left.\sum_{r} A_{s}^{r}(p) \frac{\partial}{\partial z^{r}}\right|_{(x, f(x))}\right) .
\end{aligned}
$$

Collecting terms and comparing we see that $c_{i}^{s}=\delta_{i}^{s}$ and

$$
\frac{\partial f^{r}}{\partial x^{j}}(x)=A_{j}^{r}(x, f(x))
$$

It follows from Corollary 11.23 that $f$ is uniquely determined on a sufficiently small connected neighborhood of $(x, y)$.

It is often convenient to be able to come up with the integrability conditions for a given application without trying to match indexing and notation with the above theorem. The basis of the procedure is to set mixed partials equal to each other. We demonstrate this using the notation of the theorem. We start with

$$
\frac{\partial}{\partial x^{k}} A_{j}^{i}(x, f(x))=\frac{\partial}{\partial x^{j}} A_{k}^{i}(x, f(x)) .
$$

Apply the chain rule:

$$
\begin{aligned}
& \frac{\partial A_{j}^{i}}{\partial x^{k}}(x, f(x))+\sum_{r} \frac{\partial}{\partial z^{r}} A_{j}^{i}(x, f(x)) \frac{\partial f^{r}}{\partial x^{k}}(x) \\
& \quad=\frac{\partial A_{k}^{i}}{\partial x^{j}}(x, f(x))+\sum_{s} \frac{\partial}{\partial z^{s}} A_{k}^{i}(x, f(x)) \frac{\partial f^{s}}{\partial x^{j}}(x) .
\end{aligned}
$$

Finally, substitute back using the original equations (11.5) and replace all occurrences of $f(x)$ in the arguments with the independent variable $z$. We arrive at the integrability conditions:

$$
\frac{\partial A_{j}^{i}}{\partial x^{k}}(x, z)+\sum_{l} A_{k}^{l}(x, z) \frac{\partial A_{j}^{i}}{\partial z^{l}}(x, z)=\frac{\partial A_{k}^{i}}{\partial x^{j}}(x, z)+\sum_{l} A_{j}^{l}(x, z) \frac{\partial A_{k}^{i}}{\partial z^{l}}(x, z) .
$$

The convenience of this may not be clear yet, but we shall shortly demonstrate the usefulness of this method.

Proposition 11.37. Let $U$ be open in $\mathbb{R}^{k} \times \mathbb{R}^{m}$ and ( $A_{j}^{i}$ ) an $m \times k$ matrix of smooth functions on $U$. Let $\left(x_{0}, z_{0}\right) \in U$ and suppose that for some connected open set $V$, both $f_{1}: V \rightarrow U \subset \mathbb{R}^{m}$ and $f_{2}: V \rightarrow U \subset \mathbb{R}^{m}$ are solutions of

$$
\begin{aligned}
\frac{\partial f^{i}}{\partial x^{j}}(x) & =A_{j}^{i}(x, f(x)) \text { for all } i, j \\
f\left(x_{0}\right) & =z_{0}
\end{aligned}
$$

Then $f_{1}=f_{2}$.\\
Proof. This follows from Corollary 11.23 and the considerations in the proof of the previous proposition.

Lemma 11.38. Let $V$ be an open set in $\mathbb{R}^{k}$ and let ( $A_{j}^{i}$ ) be an $m \times k$ matrix of smooth functions on $V \times \mathbb{R}^{m}$ that are linear in the second argument and satisfy the integrability conditions (11.6) on $V \times \mathbb{R}^{m}$. Then for any $x_{0} \in V$ there is a ball $B_{x_{0}} \subset V$ such that for any $(a, b) \in B_{x_{0}} \times \mathbb{R}^{m}$ there is a solution defined on $B_{x_{0}}$ with $f(a)=b$.

Proof. Let $f_{i}$ be the solution with $f_{i}\left(x_{0}\right)=\mathrm{e}_{i}$, where $\mathrm{e}_{i}$ is the $i$-th standard basis vector of $\mathbb{R}^{m}$. Then $f_{1}, \ldots, f_{m}$ are defined and linearly independent on some ball $B_{x_{0}}$ containing $x_{0}$ and contained in the intersection of the domains of the $f_{i}$. Choose $(a, b) \in B_{x_{0}} \times \mathbb{R}^{m}$ and note that $b=\sum b^{r} f_{r}(a)$ for some uniquely defined numbers $b^{i}$. Now define $f=\sum b^{i} f_{i}$ on $B_{x_{0}}$. Then writing $f=\left(f^{1}, \ldots, f^{m}\right)$, we have for any $x \in V$,

$$
\begin{aligned}
A_{j}^{i}(x, f(x)) & =A_{j}^{i}\left(x, \sum b^{r} f_{r}(x)\right)=\sum b^{r} A_{j}^{i}\left(x, f_{r}(x)\right) \\
& =\sum b^{r} \frac{\partial f_{r}^{i}}{\partial x^{j}}(x)=\frac{\partial f^{i}}{\partial x^{j}}(x)
\end{aligned}
$$

and

$$
f(a)=\sum b^{r} f_{r}(a)=b
$$

Corollary 11.39. Let $V$ be a simply connected open set in $\mathbb{R}^{k}$ and let ( $A_{j}^{i}$ ) be an $m \times k$ matrix of smooth functions on $U=V \times \mathbb{R}^{m}$. Suppose that each $A_{j}^{i}$ is linear in its second argument. If

$$
\frac{\partial A_{j}^{i}}{\partial x^{k}}+\sum_{l} A_{k}^{l} \frac{\partial A_{j}^{i}}{\partial z^{l}}=\frac{\partial A_{k}^{i}}{\partial x^{j}}+\sum_{l} A_{j}^{l} \frac{\partial A_{k}^{i}}{\partial z^{l}} \quad \text { for all } i, j, k
$$

on $V \times \mathbb{R}^{m}$, then given any $\left(x_{0}, z_{0}\right) \in V \times \mathbb{R}^{m}$, there exists a unique smooth map $f: V \rightarrow \mathbb{R}^{m}$ such that

$$
\begin{aligned}
\frac{\partial f^{i}}{\partial x^{j}}(x) & =A_{j}^{i}(x, f(x)) \quad \text { for all } i, j, \\
f\left(x_{0}\right) & =z_{0} .
\end{aligned}
$$

Proof. Let $X_{j}:=\left.\frac{\partial}{\partial x^{j}}\right|_{p}+\left.\sum_{r} A_{j}^{r}(p) \frac{\partial}{\partial z^{r}}\right|_{p}$ be the fields that span an integrable distribution on $V \times \mathbb{R}^{m}$ as in Proposition 11.36. Let $L_{\left(x_{0}, z_{0}\right)}$ be the maximal integral manifold through the point $\left(x_{0}, z_{0}\right)$. Let $\wp$ denote the restriction of the projection $\operatorname{pr}_{1}: V \times \mathbb{R}^{m} \rightarrow V$ to $L_{\left(x_{0}, z_{0}\right)}$. Let $\left(a_{1}, b_{1}\right) \in L_{\left(x_{0}, z_{0}\right)}$ and consider the set

$$
F_{a_{1}}=\wp^{-1}\left(a_{1}\right) .
$$

By Lemma 11.38 above, there is a fixed open set $U$ containing $a_{1}$ such that for every $\left(a_{1}, b\right) \in F_{a_{1}}$ there is a solution $f_{b}: U \rightarrow \mathbb{R}^{m}$ with $f\left(a_{1}\right)=b$. By Corollary 11.23, the graphs of these solutions are all disjoint open sets in $L_{\left(x_{0}, z_{0}\right)}$ and $\wp$ restricts to a diffeomorphism on each such graph. Thus $\wp: L_{\left(x_{0}, z_{0}\right)} \rightarrow V$ is a smooth covering map. The local solutions guaranteed to exist by Proposition 11.36 are local sections of this covering. Thus since $V$ is simply connected, we know from Theorem 1.95 that there is a smooth lift $\widetilde{\wp}$ of $\mathrm{id}_{V}: V \rightarrow V$ such that $\widetilde{\wp}\left(x_{0}\right)=\left(x_{0}, z_{0}\right)$, which in this case means that we have a global section: $\wp \circ \widetilde{\wp}=\operatorname{id}_{V}$. Now let $f:=\operatorname{pr}_{2} \circ \widetilde{\wp}: V \rightarrow \mathbb{R}^{m}$. Then $\widetilde{\wp}(x)=(x, f(x))$ and $f(x)$ must be smooth and for every $a \in V$ the function $f$ must agree with the unique local solution which takes the value $f(a)$ at $a$.

We return to the situation studied in Chapter 4. Consider an immersion $\mathbf{x}: V \rightarrow \mathbb{R}^{3}$, where $V$ is an open set in $\mathbb{R}^{2}$ whose standard coordinates will be denoted by $u^{1}, u^{2}$. Let ( $f_{1}, f_{2}, f_{3}$ ) be the frame fields along $\mathbf{x}$ defined by

$$
\begin{aligned}
& f_{1}:=\mathbf{x}_{u^{1}}, \quad f_{2}:=\mathbf{x}_{u^{2}} \\
& f_{3}:=N=\mathbf{x}_{u^{1}} \times \mathbf{x}_{u^{2}} /\left\|\mathbf{x}_{u^{1}} \times \mathbf{x}_{u^{2}}\right\|,
\end{aligned}
$$

where $\mathbf{x}_{u^{1}}=\partial \mathbf{x} / \partial u^{1}$, etc. The first fundamental form is given by the matrix entries

$$
g_{i j}=\left\langle\mathbf{x}_{u^{i}}, \mathbf{x}_{u^{j}}\right\rangle \text { for } 1 \leq i, j \leq 2
$$

while the second fundamental form is given by the matrix entries

$$
\ell_{i j}=-\left\langle N_{u^{i}}, \mathbf{x}_{u^{j}}\right\rangle=\left\langle N, \mathbf{x}_{u^{i} u^{j}}\right\rangle=\left\langle f_{3}, \mathbf{x}_{u^{i} u^{j}}\right\rangle .
$$

Let us consider $\mathbf{f}=\left(f_{1}, f_{2}, f_{3}\right)$ as a matrix function of $\left(u^{1}, u^{2}\right)$ that takes values in GL(3). We have

$$
\left(f_{i}\right)_{u^{1}}=\sum_{r=1}^{3} P_{i}^{r} f_{r} \text { and }\left(f_{i}\right)_{u^{2}}=\sum_{r=1}^{3} Q_{i}^{r} f_{r}
$$

for some matrix functions $P$ and $Q$. In matrix notation, we have

$$
\begin{aligned}
\mathbf{f}_{u^{1}} & =\mathbf{f} P \\
\mathbf{f}_{u^{2}} & =\mathbf{f} Q
\end{aligned}
$$

and these are called the frame equations. For convenience, we define a $3 \times 3$ matrix function $G$ by $G_{i j}:=\left\langle f_{i}, f_{j}\right\rangle$ for $1 \leq i, j \leq 3$ so that

$$
G=\left(\begin{array}{ccc}
g_{11} & g_{12} & 0 \\
g_{21} & g_{22} & 0 \\
0 & 0 & 1
\end{array}\right)
$$

For any given $x=\sum_{i}^{3} x^{i} f_{i}$, we have

$$
\xi_{i}=\left\langle x, f_{i}\right\rangle=\left\langle\sum x^{k} f_{k}, f_{i}\right\rangle=x^{k} \sum\left\langle f_{k}, f_{i}\right\rangle=G_{k i} x^{k}
$$

and so

$$
\xi_{i}=\sum\left(G^{t}\right)_{i k} x^{k}=\sum G_{i k} x^{k}
$$

Now let $x=\left(f_{i}\right)_{u_{1}}=\sum f_{k} P_{i}^{k}$. Then $x^{k}=P_{i}^{k}$, so if we define $B_{i j}:= \left\langle f_{i},\left(f_{j}\right)_{u^{1}}\right\rangle$, then we have

$$
\begin{aligned}
\left\langle f_{i},\left(f_{1}\right)_{u^{1}}\right\rangle & =B_{i 1}=\sum G_{i k} P_{1}^{k} \\
\left\langle f_{i},\left(f_{2}\right)_{u^{1}}\right\rangle & =B_{i 2}=\sum G_{i k} P_{2}^{k} \\
\left\langle f_{i},\left(f_{3}\right)_{u^{1}}\right\rangle & =B_{i 3}=\sum G_{i k} P_{3}^{k}
\end{aligned}
$$

or $B=G P$. Similarly, if $C=\left(C_{i j}\right)=\left\langle f_{i},\left(f_{j}\right)_{u^{2}}\right\rangle$, then

$$
\begin{aligned}
\left\langle f_{i},\left(f_{1}\right)_{u^{2}}\right\rangle & =C_{i 1}=\sum G_{i k} Q_{1}^{k} \\
\left\langle f_{i},\left(f_{2}\right)_{u^{2}}\right\rangle & =C_{i 2}=\sum G_{i k} Q_{2}^{k} \\
\left\langle f_{i},\left(f_{3}\right)_{u^{2}}\right\rangle & =C_{i 3}=\sum G_{i k} Q_{3}^{k}
\end{aligned}
$$

or $C=G Q$. We arrive at

$$
\begin{aligned}
P & =G^{-1} B \\
Q & =G^{-1} C
\end{aligned}
$$

We denote the entries of $G^{-1}$ by $g^{i j}$ so that

$$
G^{-1}:=\left(\begin{array}{ccc}
g^{11} & g^{12} & 0 \\
g^{21} & g^{22} & 0 \\
0 & 0 & 1
\end{array}\right)
$$

Proposition 11.40. We have

$$
B=\left(\begin{array}{ccc}
\frac{1}{2}\left(g_{11}\right)_{u^{1}} & \frac{1}{2}\left(g_{11}\right)_{u^{2}} & -\ell_{11} \\
\left(g_{12}\right)_{u^{1}}-\frac{1}{2}\left(g_{11}\right)_{u^{2}} & \frac{1}{2}\left(g_{22}\right)_{u^{1}} & -\ell_{12} \\
\ell_{11} & \ell_{12} & 0
\end{array}\right)
$$

and

$$
C=\left(\begin{array}{ccc}
\frac{1}{2}\left(g_{11}\right)_{u^{2}} & \frac{1}{2}\left(g_{12}\right)_{u^{2}}-\frac{1}{2}\left(g_{22}\right)_{u^{1}} & -\ell_{12} \\
\left(g_{22}\right)_{u^{1}}-\frac{1}{2}\left(g_{11}\right)_{u^{2}} & \frac{1}{2}\left(g_{22}\right)_{u^{2}} & -\ell_{22} \\
\ell_{12} & \ell_{22} & 0
\end{array}\right) .
$$

In particular, $P$ and $Q$ can be written in terms of the matrix entries of the first and second fundamental forms.

Proof. The proof is just a calculation, and we only do part. For example, if $i$ is either 1 or 2 , then

$$
B_{i i}=\left\langle f_{i},\left(f_{i}\right)_{u^{1}}\right\rangle=\left\langle\mathbf{x}_{u^{i}}, \mathbf{x}_{u^{i} u^{1}}\right\rangle=\frac{1}{2}\left\langle\mathbf{x}_{u^{i}}, \mathbf{x}_{u^{i}}\right\rangle_{u^{1}}=\frac{1}{2}\left(g_{i i}\right)_{u^{1}}
$$

Similarly, for $i=1$ or 2 we have

$$
C_{i i}=\left\langle f_{i},\left(f_{i}\right)_{u^{2}}\right\rangle=\left\langle\mathbf{x}_{u^{i}}, \mathbf{x}_{u^{i} u^{2}}\right\rangle=\frac{1}{2}\left(g_{i i}\right)_{u^{2}}
$$

Now, $\frac{1}{2}\left(g_{12}\right)_{u^{1}}=\left\langle\mathbf{x}_{u^{1}}, \mathbf{x}_{u^{2}}\right\rangle_{u^{1}}=\left\langle\mathbf{x}_{u^{1} u^{1}}, \mathbf{x}_{u^{2}}\right\rangle+\frac{1}{2}\left(g_{11}\right)_{u^{2}}$ from above, and so

$$
B_{21}=\left\langle f_{2},\left(f_{1}\right)_{u^{1}}\right\rangle=\left\langle\mathbf{x}_{u^{1} u^{1}}, \mathbf{x}_{u^{2}}\right\rangle=\frac{1}{2}\left(g_{12}\right)_{u^{1}}-\frac{1}{2}\left(g_{11}\right)_{u^{2}} .
$$

The entries $B_{12}, C_{12}, C_{21}$ are calculated similarly. Next consider $B_{i 3}$ for $i=1$ or 2 . We have

$$
B_{i 3}=\left\langle f_{i},\left(f_{3}\right)_{u^{1}}\right\rangle=\left\langle\mathbf{x}_{u^{i}}, N_{u^{1}}\right\rangle=-\ell_{1 i}
$$

and

$$
0=\left\langle f_{3}, f_{i}\right\rangle_{u^{1}}=\left\langle\left(f_{3}\right)_{u^{1}}, f_{i}\right\rangle+\left\langle f_{3},\left(f_{i}\right)_{u^{1}}\right\rangle
$$

so

$$
B_{3 i}=-B_{i 3} .
$$

The entries $C_{3 i}$ and $C_{i 3}$ are obtained in the same way. Lastly, since $\left\langle f_{3}, f_{3}\right\rangle=1$,

$$
0=\frac{1}{2}\left\langle f_{3}, f_{3}\right\rangle_{u^{k}}=\left\langle\left(f_{3}\right)_{u^{k}}, f_{3}\right\rangle=\left\{\begin{array}{l}
B_{33} \text { if } k=1, \\
C_{33} \text { if } k=2 .
\end{array}\right.
$$

We record an observation to be used later:

$$
\begin{aligned}
& B+B^{t}=G_{u^{1}}, \\
& C+C^{t}=G_{u^{2}} .
\end{aligned}
$$

The frame equations (11.7) are a system to which Proposition 11.36 applies. Rather than trying to rewrite the equations in a form that matches that proposition we obtain the integrability conditions by setting

$$
\left(\mathbf{f}_{u^{1}}\right)_{u^{2}}=\left(\mathbf{f}_{u^{2}}\right)_{u^{1}}
$$

and then

$$
\mathbf{f}_{u^{2}} P+\mathbf{f} P_{u^{2}}=\mathbf{f}_{u^{1}} Q+\mathbf{f} Q_{u^{1}}
$$

Substituting from the frame equations we obtain

$$
\mathbf{f}\left(P_{u^{2}}-Q_{u^{1}}-(P Q-Q P)\right)=0
$$

Now $\mathbf{f}$ is a nonsingular matrix, so we have the equivalent integrability equation

$$
P_{u^{2}}-Q_{u^{1}}-(P Q-Q P)=0
$$

At this point we pause to appreciate an important fact. Namely, direct calculation reveals that these equations are equivalent to the combination of the Codazzi-Mainardi equation and the Gauss curvature equation, which we now see are integrability conditions (see Problem 7). We thus refer to the above integrability equations (11.15) as the Gauss-Codazzi equations with apologies to Gaspare Mainardi (1800-1879).

We now turn things around. Rather than assuming that we have a surface, we take the $\left(g_{i j}\right)$ and $\left(\ell_{i j}\right)$ as some given symmetric smooth matrixvalued functions defined on a connected open $V \subset \mathbb{R}^{2}$ with the assumption that ( $g_{i j}$ ) is positive definite. Furthermore, we now assume that $G, P$, and $Q$ are actually defined in terms of these by the formulas above, which we found to be true in the case where we started with a surface. We will show that we can obtain a surface with these as first and second fundamental form.

Theorem 11.41 (Fundamental existence theorem for surfaces). The following assertions hold:\\
(i) Let $V$ be an open set in $\mathbb{R}^{2}$ diffeomorphic to an open disk and let $\mathbf{x}: V \rightarrow \mathbb{R}^{3}$ be an immersion with the corresponding first and second fundamental forms given in matrix form as $\left(g_{i j}\right)$ and $\left(\ell_{i j}\right)$. Let $\mathbf{y}: V \rightarrow \mathbb{R}^{3}$ be another immersion with the corresponding forms $\left(\widetilde{g}_{i j}\right)$ and $\left(\widetilde{\ell}_{i j}\right)$. If

$$
\mathbf{y}=f \circ \mathbf{x}
$$

for some proper Euclidean motion $f: \mathbb{R}^{3} \rightarrow \mathbb{R}^{3}$, then

$$
\begin{aligned}
& \widetilde{g}_{i j}=g_{i j} \\
& \widetilde{\ell}_{i j}=\ell_{i j}
\end{aligned}
$$

Conversely, if the last equations hold, then $\mathbf{y}=f \circ \mathbf{x}$ for some Euclidean motion $f$.\\
(ii) Suppose that $\left(g_{i j}\right)$ and $\left(\ell_{i j}\right)$ are given symmetric matrix-valued functions defined on $V$ with $\left(g_{i j}\right)$ positive definite and suppose that $G$, $B$ and $C$ are defined in terms of the entries of $\left(g_{i j}\right)$ and $\left(\ell_{i j}\right)$ as in\\
formulas (11.8), (11.11), (11.12) and (11.13). Then if

$$
\begin{aligned}
& P=G^{-1} B \\
& Q=G^{-1} C
\end{aligned}
$$

and if $P_{u^{2}}-Q_{u^{1}}-(P Q-Q P)=0$, then there exists an embedding $\mathbf{x}: V \rightarrow \mathbb{R}^{3}$ such that $\left(g_{i j}\right)$ and $\left(\ell_{i j}\right)$ are the corresponding first and second fundamental forms.

Proof. We leave the proof of the first part of (i) to the reader, but note that it can be proved using direct calculation or it can be derived from Theorem 4.22.

For the rest of (i), note that by composing with a translation we may assume that both $\mathbf{x}$ and $\mathbf{y}$ map some fixed point $u \in V$ to the origin in $\mathbb{R}^{3}$. Let ( $f_{1}, f_{2}, f_{3}$ ) be the natural frame for $\mathbf{x}$ as above and let ( $\widetilde{f}_{1}, \widetilde{f}_{2}, \widetilde{f}_{3}$ ) be that of $\mathbf{y}$. By making a rotation we may assume that these two frames agree at $p$. But since we are assuming that $\widetilde{g}_{i j}=g_{i j}$ and $\widetilde{\ell}_{i j}=\ell_{i j}$, it follows that both frames satisfy the same frame equations and so by Proposition 11.37 they must agree on the connected set $V$. In particular, $\mathbf{x}_{u^{i}}=\mathbf{y}_{u^{i}}$ for $i=1,2$. Thus $\mathbf{x}$ and $\mathbf{y}$ only differ by a constant, which must be zero since $\mathbf{x}(u)=\mathbf{y}(u)$.

Next we consider (ii) where ( $g_{i j}$ ) and ( $\ell_{i j}$ ) are given. We want to construct a surface, but first we construct the frame for the desired surface. Since it is assumed that $g=\left(g_{i j}\right)$ is positive definite, $g$ and the extended matrix $G$ are both invertible and positive definite. Thus $P$ and $Q$ are well-defined. Since we assume that the integrability equations $P_{u^{2}}-Q_{u^{1}}- (P Q-Q P)=0$ hold, Theorem 11.36 tells us that we can solve the frame equations locally, near any point $u \in V$ and with any initial conditions $\mathbf{f}(u)=\mathbf{f}_{0}$ holding as desired. But the system is linear and our domain is diffeomorphic to a disk so we can use Corollary 11.39 to obtain a solution on all of $V$. Since $G$ is positive definite, we may choose these initial conditions so that

$$
\left\langle f_{i}(u), f_{j}(u)\right\rangle=G_{i j}(u) \quad(i j \text {-th entry of } G \text { at } u) .
$$

Having obtained the $f_{i}$ near $u$, we now wish to obtain a surface. This means solving the system

$$
\begin{aligned}
& \mathbf{x}_{u^{1}}=f_{1}, \\
& \mathbf{x}_{u^{2}}=f_{2},
\end{aligned}
$$

and this time the integrability conditions are derived from

$$
\left(f_{1}\right)_{u^{2}}=\left(f_{2}\right)_{u^{1}}
$$

Using the frame equations, we obtain integrability conditions

$$
\sum Q_{1}^{j} f_{j}=\sum P_{2}^{j} f_{j} .
$$

This just says that the second column of $P$ is equal to the first column of $Q$, which is true. Thus we can find $\mathbf{x}: V \rightarrow \mathbb{R}^{3}$ with $\mathbf{x}(u)=0$ so that (11.16) holds.

Next we show that $\left\langle f_{i}, f_{j}\right\rangle=G_{i j}$ on all of $V$. We compute as follows:

$$
\begin{aligned}
\left\langle f_{i}, f_{j}\right\rangle_{u^{1}} & =\left\langle\left(f_{i}\right)_{u^{1}}, f_{j}\right\rangle+\left\langle f_{i},\left(f_{j}\right)_{u^{1}}\right\rangle \\
& =\sum_{r} Q_{i}^{r}\left\langle f_{r}, f_{j}\right\rangle+\sum_{s} Q_{j}^{s}\left\langle f_{i}, f_{s}\right\rangle=\left(Q G+(Q G)^{t}\right)_{i j} \\
& =\left(G Q+(G Q)^{t}\right)_{i j}=\left(B+B^{t}\right)_{i j}=\left(G_{u^{1}}\right)_{i j}=\left(G_{i j}\right)_{u^{1}}
\end{aligned}
$$

by equations (11.14). Similarly for $u^{2}$. Thus $\left\langle f_{i}, f_{j}\right\rangle-G_{i j}$ is a constant, which must be zero since it is zero at $u$. From $\left\langle f_{i}, f_{j}\right\rangle=G_{i j}$ it follows that $\left\langle f_{3}, f_{3}\right\rangle=1$ and that $f_{1}, f_{2}$ are independent and orthogonal to $f_{3}$. It remains to show that $\left\langle\left(f_{3}\right)_{u^{i}}, f_{j}\right\rangle=-\ell_{i j}$. We compute as follows:

$$
\begin{aligned}
-\left\langle\left(f_{3}\right)_{u^{1}}, f_{j}\right\rangle & =\left(g^{11} \ell_{11}+g^{12} \ell_{12}\right)\left\langle f_{1}, f_{j}\right\rangle+\left(g^{12} \ell_{11}+g^{22} \ell_{12}\right)\left\langle f_{2}, f_{j}\right\rangle \\
& =\left(g^{11} \ell_{11}+g^{12} \ell_{12}\right) g_{1 j}+\left(g^{12} \ell_{11}+g^{22} \ell_{12}\right) g_{2 j} \\
& =\ell_{11}\left(g^{11} g_{1 j}+g^{12} g_{2 j}\right)+\ell_{12}\left(g^{21} g_{1 j}+g^{22} g_{2 j}\right) \\
& =\ell_{11} \delta_{j}^{1}+\ell_{12} \delta_{j}^{2}
\end{aligned}
$$

This shows that $\left\langle\left(f_{3}\right)_{u^{1}}, f_{j}\right\rangle=-\ell_{1 j}$ for $j=1,2$. The computation of $-\left\langle\left(f_{3}\right)_{u^{2}}, f_{j}\right\rangle$ is similar and left for the reader.

\subsection*{Local Fundamental Theorem of Calculus}
Recall the structure equations (8.15) satisfied by the Maurer-Cartan form $\omega_{G}$ for a Lie group $G$ :

$$
d \omega_{G}=-\frac{1}{2}\left[\omega_{G}, \omega_{G}\right]^{\wedge}
$$

If $v_{1}=X_{1}(e), \ldots, v_{n}=X_{n}(e)$ is a basis for the Lie algebra $\mathfrak{g}$ which extends to left invariant vector fields $X_{1}, \ldots, X_{n}$, then the above equation is equivalent to

$$
d \omega^{k}=-\sum_{i<j} c_{i j}^{k} \omega^{i} \wedge \omega^{j}=-\frac{1}{2} \sum_{i, j} c_{i j}^{k} \omega^{i} \wedge \omega^{j}
$$

where the $c_{i j}^{k}$ are the structure constants associated to $\omega^{1}, \ldots, \omega^{n}$, which is the left invariant frame field dual to $X_{1}, \ldots, X_{n}$. If $M$ is some $m$-manifold and $f: M \rightarrow G$ is a smooth map, then $\omega_{f}=f^{*} \omega_{G}$ is a $\mathfrak{g}$-valued 1 -form on $M$. By naturality we have

$$
d \omega_{f}=-\frac{1}{2}\left[\omega_{f}, \omega_{f}\right]^{\wedge}
$$

or equivalently

$$
d \omega_{f}^{k}=-\frac{1}{2} \sum_{i, j} c_{i j}^{k} \omega_{f}^{i} \wedge \omega_{f}^{j}
$$

where $\omega_{f}^{i}=f^{*} \omega^{i}$ for $i=1, \ldots, n$. The $\mathfrak{g}$-valued 1 -form $\omega_{f}$ is sometimes called the (left) Darboux derivative of $f$. The right Darboux derivative is defined similarly using the right Maurer-Cartan form.

If we think of a $\mathfrak{g}$-valued 1 -form on a manifold $M$ as a map $T M \rightarrow \mathfrak{g}$, then $\omega_{f}=f^{*} \omega_{G}=\omega_{G} \circ T f$. From this point of view we can understand why $\omega_{f}$ is a kind of derivative of $f$ by considering the special case where $G$ is a vector space V with its abelian (additive) Lie group structure. In this case, the Lie algebra is V itself and the Maurer-Cartan form is just the canonical map $\operatorname{pr}_{2}: T \mathrm{V}=\mathrm{V} \times \mathrm{V} \rightarrow \mathrm{V}$, and so for a smooth map $f: M \rightarrow \mathrm{V}$, the Darboux derivative is the differential $d f=\mathrm{pr}_{2} \circ T f$. Just as for the differential, the Darboux derivative embodies less information than the tangent map since the values that the map takes are "forgotten" and only tangential information is retained. Indeed, notice that if $L_{g}: G \rightarrow G$ is a left translation and $F:=L_{g} \circ f$, then

$$
\omega_{F}=F^{*} \omega_{G}=f^{*} L_{g}^{*} \omega_{G}=f^{*} \omega_{G}=\omega_{f}
$$

since $\omega_{G}$ is left invariant. Hence two smooth maps into $G$ that differ by left translation have the same (left) Darboux derivative. This generalizes the fact that two functions that differ by an additive constant have the same differential.

For a smooth 1-form $\vartheta=g d t$ on $\mathbb{R}$, we can always find a smooth function $f$ with $d f=\vartheta$ since by the Fundamental theorem of calculus one need only choose $f(t)=\int_{0}^{t} g(\tau) d \tau$. More generally, if $M$ is simply connected and $G$ is a vector space V , then the fact that $H^{1}(M)=0$ means that every V -valued 1-form is the differential of some smooth $f: M \rightarrow \mathrm{V}$. For a general $G$, if $\vartheta$ is a $\mathfrak{g}$-valued 1 -form on $M$, then we may ask for an $f$ such that $\vartheta=\omega_{f}$. But there is no reason to expect a general $\vartheta$ to satisfy the above structure equation that $\omega_{f}$ satisfies. Now if we choose a basis $\left\{v_{i}\right\}$ for $\mathfrak{g}$, then there must be 1 -forms $\vartheta^{1}, \ldots, \vartheta^{n} \in \Omega^{1}(M)$ such that

$$
\vartheta=\sum_{i=1}^{n} v_{i} \vartheta^{i}
$$

Then $d \vartheta=-\frac{1}{2}[\vartheta, \vartheta]^{\wedge}$ is equivalent to

$$
d \vartheta^{k}=-\frac{1}{2} \sum_{i, j} c_{i j}^{k} \vartheta^{i} \wedge \vartheta^{j}
$$

where $c_{i j}^{k}$ are the structure constants. As we said, this may or may not hold. These equations are the integrability conditions for the existence of an $f$ such that $\vartheta=\omega_{f}$. More precisely, we have the following theorem.

Theorem 11.42. Let $M$ be an $m$-manifold and $G$ an $n$-dimensional Lie group. If $\vartheta$ is a $\mathfrak{g}$-valued 1 -form on $M$ that satisfies the structural equation $d \vartheta=-\frac{1}{2}[\vartheta, \vartheta]^{\wedge}$, then for every $p_{0} \in M$ there is a neighborhood $U_{p_{o}}$ of $p_{0}$ such that given any $(a, b) \in U_{p_{o}} \times G$ there is a smooth function $f: U_{x_{o}} \rightarrow G$ with $f(a)=b$ and $\vartheta=\omega_{f}$.

Proof. Let $\operatorname{pr}_{1}: M \times G \rightarrow M$ and $\operatorname{pr}_{2}: M \times G \rightarrow G$ be the canonical projections and define a $\mathfrak{g}$-valued 1-form on $M \times G$ by

$$
\Omega:=\operatorname{pr}_{1}^{*} \vartheta-\operatorname{pr}_{2}^{*} \omega_{G}
$$

For each $(p, g) \in M \times G$, let $E_{(p, g)}=\operatorname{Ker} \Omega_{(p, g)}$. Now define a vector bundle homomorphism $T(M \times G) \rightarrow(M \times G) \times \mathfrak{g}$ by $v_{(p, g)} \mapsto\left((p, g), \Omega_{(p, g)}\left(v_{(p, g)}\right)\right)$. By Proposition 6.28, if this homomorphism has constant rank, then the kernel is a subbundle which clearly has fiber $E_{(p, g)}$ at ( $p, g$ ). By linear algebra, this is equivalent to showing that the dimension of $E_{(p, g)}$ is independent of $(p, g)$. This will follow if we show that $\left.T \operatorname{pr}_{1}\right|_{E_{(p, g)}}: E_{(p, g)} \rightarrow T_{p} M$ is an isomorphism for all ( $p, g$ ). If we identify $T_{(p, g)}(M \times G)$ with $T_{p} M \times T_{g} G$, then $T \operatorname{pr}_{1}$ is just the projection $(v, w) \mapsto v$ and similarly for $T \operatorname{pr}_{2}$. Now if $(v, w) \in E_{(p, g)}$ and $T \operatorname{pr}_{1} \cdot(v, w)=0$ then $v=0$. But, since $\vartheta(v)=\omega_{G}(w)$, we have $w=0$ also. Thus $\left.T \operatorname{pr}_{1}\right|_{E_{(p, g)}}$ is injective. It is also surjective since for any $v \in T_{p} M$, we clearly have $\left(v, T_{e} L_{g}(\vartheta(v)) \in E_{(p, g)}\right.$ and this has $v$ as its image.

Now we use Proposition 11.19 to show that $E$ is integrable:

$$
\begin{aligned}
d \Omega & =\operatorname{pr}_{1}^{*} d \vartheta-\operatorname{pr}_{2}^{*} d \omega_{G}=\operatorname{pr}_{1}^{*}\left(-\frac{1}{2}[\vartheta, \vartheta]^{\wedge}\right)-\operatorname{pr}_{2}^{*}\left(-\frac{1}{2}\left[\omega_{G}, \omega_{G}\right]^{\wedge}\right) \\
& =-\frac{1}{2}\left[\operatorname{pr}_{1}^{*} \vartheta, \operatorname{pr}_{1}^{*} \vartheta\right]^{\wedge}+\frac{1}{2}\left[\operatorname{pr}_{2}^{*} \omega_{G}, \operatorname{pr}_{2}^{*} \omega_{G}\right]^{\wedge}
\end{aligned}
$$

But $\operatorname{pr}_{1}^{*} \vartheta=\Omega+\operatorname{pr}_{2}^{*} \omega_{G}$, so

$$
\begin{aligned}
d \Omega & =-\frac{1}{2}\left[\left(\Omega+\operatorname{pr}_{2}^{*} \omega_{G}\right),\left(\Omega+\operatorname{pr}_{2}^{*} \omega_{G}\right)\right]^{\wedge}-\frac{1}{2}\left[\operatorname{pr}_{2}^{*} \omega_{G}, \operatorname{pr}_{2}^{*} \omega_{G}\right]^{\wedge} \\
& =-\frac{1}{2}[\Omega, \Omega]^{\wedge}-\frac{1}{2}\left[\Omega, \operatorname{pr}_{2}^{*} \omega_{G}\right]^{\wedge}-\frac{1}{2}\left[\operatorname{pr}_{2}^{*} \omega_{G}, \Omega\right]^{\wedge}
\end{aligned}
$$

which makes it clear that $d \Omega(X, Y)=0$ whenever $\Omega(X)=0$ and $\Omega(Y)=0$.\\
Now we use the leaves of the distribution to construct the solution. Let $x_{0} \in M$ and fix $g_{0} \in G$. Then let $L_{\left(x_{0}, g_{0}\right)}$ be the maximal integral manifold through $\left(x_{0}, g_{0}\right)$. The map $T \operatorname{pr}_{1} \mid E_{\left(p_{0}, g_{0}\right)}: E_{\left(p_{0}, g_{0}\right)} \rightarrow T_{p} M$ is an isomorphism so the inverse mapping theorem tells us that $\operatorname{pr}_{1} \mid L_{\left(p_{0}, g_{0}\right)}$ restricts to a diffeomorphism on some neighborhood $O$ of $\left(p_{0}, g_{0}\right)$ in $L_{\left(p_{0}, g_{0}\right)}$. Let $\Phi: U \rightarrow O \subset L_{\left(p_{0}, g_{0}\right)}$ denote the inverse of this diffeomorphism. Since $\operatorname{pr}_{1} \circ \Phi=\operatorname{id}_{U}$, there must be a smooth function $f$ such that $\Phi(p)=(p, f(p))$\\
for all $p \in U$. Observe that $\Phi^{*}(\Omega)=0$ since the image of $T \Phi$ lies in the distribution by construction. Thus we have

$$
\begin{aligned}
0 & =\Phi^{*}(\Omega)=\Phi^{*}\left(\operatorname{pr}_{1}^{*} \vartheta-\operatorname{pr}_{2}^{*} \omega_{G}\right) \\
& \left.=\Phi^{*} \operatorname{pr}_{1}^{*} \vartheta-\Phi^{*} \operatorname{pr}_{2}^{*} \omega_{G}\right)=\vartheta-f^{*} \omega_{G}
\end{aligned}
$$

or $\left.\vartheta\right|_{U}=f^{*} \omega_{G}=\omega_{f}$. Let $(a, b) \in U \times G$. Now we argue that we may modify $f$ so that we still have $\left.\vartheta\right|_{U}=\omega_{f}$ while now $f(a)=b$. In fact, if $f(a)=b_{1}$, then let $g=b_{1}^{-1} b$ and replace $f$ by $L_{g} \circ f$.

Proposition 11.43. Let $M$ be an $m$-manifold and $G$ an $n$-dimensional Lie group. Let $\vartheta$ be a $\mathfrak{g}$-valued 1 -form on $M$ and suppose that $f_{1}: U \rightarrow G$ and $f_{2}: U \rightarrow G$ are smooth maps such that $\left.\vartheta\right|_{U}=f_{i}^{*} \omega_{G}$ for $i=1,2$. Then if $U$ is connected and $f_{1}\left(p_{0}\right)=f_{2}\left(p_{0}\right)$ for some $p_{0} \in U$, then $f_{1}=f_{2}$.

Proof. Let $\Phi_{i}:=\left(\operatorname{id}_{U}, f_{i}\right)$. Then

$$
\begin{aligned}
\Phi_{i}^{*}(\Omega) & =\Phi_{i}^{*}\left(\operatorname{pr}_{1}^{*} \vartheta-\operatorname{pr}_{2}^{*} \omega_{G}\right) \\
& \left.=\Phi_{i}^{*} \operatorname{pr}_{1}^{*} \vartheta-\Phi_{i}^{*} \operatorname{pr}_{2}^{*} \omega_{G}\right)=\vartheta-f^{*} \omega_{G}=0,
\end{aligned}
$$

so $\Phi_{i}: U \rightarrow M \times G$ is an embedding for $i=1,2$ whose image is the graph of $f_{i}$ and an integral manifold of the distribution generated by $\Omega$ that contains ( $p_{0}, f\left(p_{0}\right)$ ). Corollary 11.23 applies to show that $\Phi_{1}(U)=\Phi_{2}(U)$ and $f_{1}=f_{2}$.

Corollary 11.44. Let $M, G, \vartheta$ and $\Omega$ be as above and suppose that $d \vartheta= -\frac{1}{2}[\vartheta, \vartheta]^{\wedge}$. Then the restriction of the map $\operatorname{pr}_{1}: M \times G \rightarrow M$ to any leaf of the distribution given by $\Omega$ is a covering map.

Proof. Let $L$ be a leaf and choose $p_{0} \in M$. Choose $\left(p_{0}, g_{0}\right) \in L$ and let $\Phi: U \rightarrow O \subset L_{\left(p_{0}, g_{0}\right)}=L$ be the diffeomorphism constructed as in the proof of the previous theorem where $U$ is connected. Now $\Phi=\left(\operatorname{id}_{U}, f\right)$ where $f\left(p_{0}\right)=g_{0}$. If ( $p_{0}, g_{1}$ ) is any other point in the leaf, then for $g=g_{1} g_{0}^{-1}$ the map $\Phi_{1}:=\left(\operatorname{id}_{U}, L_{g} \circ f\right)$ is a diffeomorphism $U \rightarrow O_{1}$ where $\left(p_{0}, g_{1}\right) \in O_{1}$ and $\operatorname{pr}_{1}\left(O_{1}\right)=U$. In fact, it is easy to see that $\Phi_{1} \circ \Phi^{-1}: O \rightarrow O_{1}$ is a diffeomorphism. If $\left(p_{0}, g_{1}\right)$ and $\left(p_{0}, g_{2}\right)$ are distinct points in the leaf, then we construct diffeomorphisms $\Phi_{1}: O \rightarrow O_{1}$ and $\Phi_{2}: O \rightarrow O_{2}$ as above, and Corollary 11.23 applies to show that $O_{1}$ and $O_{2}$ are disjoint and each map diffeomorphically onto $O$ under $\operatorname{pr}_{1}$. It is now clear that $M$ is evenly covered by $\operatorname{pr}_{1} \mid L$.

Corollary 11.45. Let $M$ be a simply connected $m$-manifold and $G$ an $n$ dimensional Lie group. If $\vartheta$ is a $\mathfrak{g}$-valued 1 -form on $M$ that satisfies the structural equation $d \vartheta=-\frac{1}{2}[\vartheta, \vartheta]^{\wedge}$, then for every $\left(p_{0}, g_{0}\right) \in M \times G$ there is a smooth function $f: M \rightarrow G$ with $f\left(p_{0}\right)=g_{0}$ and $\vartheta=\omega_{f}$.

Proof. Let $L$ be the leaf of the distribution determined by $\Omega$ that contains $\left(p_{0}, g_{0}\right)$. Since $M$ is simply connected, we can lift $\operatorname{id}_{M}: M \rightarrow M$ to a section $\Phi: M \rightarrow L$ of $\operatorname{pr}_{1} \mid L$ such that $\Phi\left(p_{0}\right)=g_{0}$. Once again $\Phi=\left(\operatorname{id}_{M}, f\right)$ for some smooth function $f$ with $f\left(p_{0}\right)=g_{0}$ and we argue as before to conclude that $\vartheta=\omega_{f}$ (this time globally).

\section*{Problems}
(1) Show that the following vector fields define a rank 2 distribution on $\mathbb{R}^{3}$ which is not involutive (and hence not integrable):

$$
\begin{aligned}
X & =\frac{\partial}{\partial x}+y \frac{\partial}{\partial z} \\
Y & =\frac{\partial}{\partial y}
\end{aligned}
$$

Draw a picture of the portion of this distribution that sits at points in the $x, y$-plane and try to see geometrically why the distribution is not integrable.\\
(2) Show that the distribution on $\mathbb{R}^{4}$ given by $X=\frac{\partial}{\partial y}+x \frac{\partial}{\partial z}$ and $Y= \frac{\partial}{\partial x}+y \frac{\partial}{\partial w}$, where $(x, y, z, w)$ are standard coordinates, has no integral manifolds.\\
(3) Let $\theta$ be a 1 -form. Show that a 2 -form $\eta$ is in the ideal generated by $\theta$ if and only if $\eta \wedge \theta=0$.\\
(4) Consider again the system of partial differential equations:

$$
\begin{aligned}
& z_{x}=F(x, y, z) \\
& z_{y}=G(x, y, z)
\end{aligned}
$$

Show that the graphs of solutions of these equations are integral manifolds of the distribution defined by the 1-form $\theta=d z-F d x-G d y$. Use Theorem 11.18 to deduce the integrability conditions for this system. [Hint: Use Problem 3.]\\
(5) Let $H$ be the Heisenberg group consisting of all matrices of the form

$$
A=\left[\begin{array}{ccc}
1 & x_{12} & x_{13} \\
0 & 1 & x_{23} \\
0 & 0 & 1
\end{array}\right]
$$

The $x_{i j}$ give global coordinates and a diffeomorphism with $\mathbb{R}^{3}$. Let $V_{12}, V_{13}, V_{23}$ be the left invariant vector fields on $H$ that have values at the identity with components with respect to the coordinate fields given by $(1,0,0),(0,1,0)$, and $(0,0,1)$, respectively. Let $\Delta_{\left\{V_{12}, V_{13}\right\}}$ and\\
$\Delta_{\left\{V_{12}, V_{23}\right\}}$ be the 2-dimensional distributions generated by the indicated pairs of vector fields. Show that $\Delta_{\left\{V_{12}, V_{13}\right\}}$ is integrable and $\Delta_{\left\{V_{12}, V_{23}\right\}}$ is not.\\
(6) Prove (i) of Theorem 11.41.\\
(7) Show that for a given surface, equations (11.15) are equivalent to a combination of the Codazzi-Mainardi equations and the Gauss curvature equations defined in Chapter 4.

\section*{Connections and Covariant Derivatives}
The terms "covariant derivative" and "connection" are sometimes treated as synonymous. In fact, a covariant derivative is sometimes called a Koszul connection. From one point of view, the central idea is that of measuring the rate of change of sections of bundles in the direction of a vector or vector field on the base manifold. Here the derivative viewpoint is prominent. From another related point of view, a connection provides an extra structure that gives a principled way of lifting curves from the base to the total space. The lifts are parallel sections along the curve. In this chapter, we will always take the typical fiber of an $\mathbb{F}$-vector bundle of rank $k$ to be $\mathbb{F}^{k}$. We let $\left(\mathbf{e}_{1}, \ldots, \mathbf{e}_{k}\right)$ be the standard basis of $\mathbb{F}^{k}$. Thus every vector bundle chart $(U, \phi)$ is associated with a frame field $\left(e_{1}, \ldots, e_{k}\right)$ where $e_{i}(x):=\phi^{-1}\left(x, \mathbf{e}_{i}\right)$.

\subsection*{Definitions}
Let $\pi: E \rightarrow M$ be a smooth $\mathbb{F}$-vector bundle of rank $k$ over a manifold $M$. A covariant derivative can either be defined as a map $\nabla: \mathfrak{X}(M) \times \Gamma(M, E) \rightarrow \Gamma(M, E)$ with certain properties from which one derives a well-defined map $\nabla: T M \times \Gamma(M, E) \rightarrow \Gamma(M, E)$ with nice properties or the other way around. We rather arbitrarily start with the first of these definitions.

Definition 12.1. A covariant derivative or Koszul connection on a smooth $\mathbb{F}$-vector bundle $E \rightarrow M$ is a map $\nabla: \mathfrak{X}(M) \times \Gamma(M, E) \rightarrow \Gamma(M, E)$ (where $\nabla(X, s)$ is written as $\nabla_{X} s$ ) satisfying the following four properties:\\
(i) $\nabla_{f X}(s)=f \nabla_{X} s$ for all $f \in C^{\infty}(M), X \in \mathfrak{X}(M)$ and $s \in \Gamma(M, E)$;\\
(ii) $\nabla_{X_{1}+X_{2}} s=\nabla_{X_{1}} s+\nabla_{X_{2}} s$ for all $X_{1}, X_{2} \in \mathfrak{X}(M)$ and $s \in \Gamma(M, E)$;\\
(iii) $\nabla_{X}\left(s_{1}+s_{2}\right)=\nabla_{X} s_{1}+\nabla_{X} s_{2}$ for all $X \in \mathfrak{X}(M)$ and $s_{1}, s_{2} \in \Gamma(M, E) ;$\\
(iv) $\nabla_{X}(f s)=(X f) s+f \nabla_{X} s$ for all $f \in C^{\infty}(M ; \mathbb{F}), X \in \mathfrak{X}(M)$ and $s \in \Gamma(M, E)$.

For a fixed $X \in \mathfrak{X}(M)$, the map $\nabla_{X}: \Gamma(M, E) \rightarrow \Gamma(M, E)$ is called the covariant derivative with respect to $X$.

As we will see below, this definition is enough to imply that $\nabla$ induces maps $\nabla^{U}: \mathfrak{X}(U) \times \Gamma(U, E) \rightarrow \Gamma(U, E)$, one for each open $U \subset M$, that are naturally related in a sense we make precise below (this is not necessarily true for infinite-dimensional manifolds). Furthermore, we also prove that for a fixed $p \in M$, the value $\left(\nabla_{X} s\right)(p)$ depends only on the value of $X$ at $p$ and only on the values of $s$ along any smooth curve $c$ representing $X_{p}$. Thus we get a well-defined map $\nabla: T M \times \Gamma(M, E) \rightarrow \Gamma(M, E)$ such that $\nabla_{v} s=\left(\nabla_{X} s\right)(p)$ for any extension of $v \in T_{p} M$ to a vector field $X$ with $X_{p}=v$. The resulting properties are\\
(i') $\nabla_{a v}(s)=a \nabla_{v} s$ for all $a \in \mathbb{R}, v \in T M$ and $s \in \Gamma(M, E) ;$\\
(ii') for all $p \in M$ we have $\nabla_{v_{1}+v_{2}} s=\nabla_{v_{1}} s+\nabla_{v_{2}} s$ for all $v_{1}, v_{2} \in T_{p} M$, and $s \in \Gamma(M, E)$;\\
(iii') $\nabla_{v}\left(s_{1}+s_{2}\right)=\nabla_{v} s_{1}+\nabla_{v} s_{2}$ for all $v \in T M$ and $s_{1}, s_{2} \in \Gamma(M, E)$;\\
(iv') for all $p \in M$ we have $\nabla_{v}(f s)=(v f) s(p)+f(p) \nabla_{v} s$ for all $v \in T_{p} M$, $s \in \Gamma(M, E)$ and $f \in C^{\infty}(M ; \mathbb{F}) ;$\\
(v') $p \mapsto \nabla_{X_{p}} s$ is smooth for all smooth vector fields $X$.\\
A map satisfying these properties is also called a covariant derivative (or Koszul connection). Note that it is easy to obtain a Koszul connection in the first sense since we just let $\left(\nabla_{X} s\right)(p):=\nabla_{X_{p}} s$.

Definition 12.2. A covariant derivative on the tangent bundle $T M$ of an $n$-manifold $M$ is usually referred to as a covariant derivative on $M$.

In Chapter 4 have already met the Levi-Civita connection $\bar{\nabla}$ on $\mathbb{R}^{n}$, which is, from the current point of view, a Koszul connection on the tangent bundle of $\mathbb{R}^{n}$. The definition of this connection takes advantage of the natural identification of tangent spaces which makes taking the difference quotient possible:

$$
\bar{\nabla}_{v} X=\lim _{t \rightarrow 0} \frac{X(p+t v)-X(p)}{t}
$$

In that same chapter we obtained, by a projection, a covariant derivative on (the tangent bundle of) any hypersurface in $\mathbb{R}^{n}$. A covariant derivative on\\
a submanifold of arbitrary codimension can be obtained in the same way. Let $M$ be a submanifold of $\mathbb{R}^{n}$ and let $X \in \mathfrak{X}(M)$ and $v \in T_{p} M$. We have

$$
\nabla_{v} X:=\left(\left.\frac{d}{d t}\right|_{0} X \circ c\right)^{\top} \in T_{p} M
$$

One may easily verify that $\nabla$, so defined, is a covariant derivative (in the second sense above).

Returning to the case of a general vector bundle, let us consider how covariant differentiation behaves with respect to restriction to open subsets of our manifold. Recall the restriction map $r_{V}^{U}: \Gamma(U, E) \rightarrow \Gamma(V, E)$ given by $r_{V}^{U}:\left.\sigma \mapsto \sigma\right|_{V}$ and where $V \subset U$.

Definition 12.3. A natural covariant derivative $\nabla$ on a smooth $\mathbb{F}$ vector bundle $E \rightarrow M$ is an assignment to each open set $U \subset M$ of a map $\nabla^{U}: \mathfrak{X}(U) \times \Gamma(U, E) \rightarrow \Gamma(U, E)$ written as $\nabla^{U}:(X, \sigma) \mapsto \nabla_{X}^{U} \sigma$ such that the following assertions hold:\\
(i) For every open $U \subset M$, the map $\nabla^{U}$ is a Koszul connection on the restricted bundle $\left.E\right|_{U} \rightarrow U$.\\
(ii) For nested open sets $V \subset U$, we have $r_{V}^{U}\left(\nabla_{X}^{U} \sigma\right)=\nabla_{r_{V}^{U} X}^{V} r_{V}^{U} \sigma$ (naturality with respect to restrictions).\\
(iii) For $X \in \mathfrak{X}(U)$ and $\sigma \in \Gamma(U, E)$ the value $\left(\nabla_{X}^{U} \sigma\right)(p)$ only depends on the value of $X$ at $p \in U$.

Here $\nabla_{X}^{U} \sigma$ is called the covariant derivative of $\sigma$ with respect to $X$. We will denote all of the maps $\nabla^{U}$ by the single symbol $\nabla$ when there is no chance of confusion. We have explicitly worked the naturality conditions (ii) and (iii) into the definition of a natural covariant derivative, so this definition is also appropriate for infinite-dimensional manifolds. The definition of Koszul connection did not specifically include these naturality features and was only defined for global sections. We shall now see that, in the case of finite-dimensional manifolds, a Koszul connection gives a natural covariant derivative for free.

Lemma 12.4. Suppose $\nabla: \mathfrak{X}(M) \times \Gamma(E) \rightarrow \Gamma(E)$ is a Koszul connection for the vector bundle $E \rightarrow M$. Then if for some open $U$ either $\left.X\right|_{U}=0$ or $\left.\sigma\right|_{U}=0$, then

$$
\left(\nabla_{X} \sigma\right)(p)=0 \text { for all } p \in U
$$

Proof. We prove the case of $\left.\sigma\right|_{U}=0$ and leave the case of $\left.X\right|_{U}=0$ to the reader.

Let $q \in U$. Then there is some relatively compact open set $V$ with $q \in \bar{V} \subset U$ and a smooth function $f$ that is identically one on $V$ and zero\\
outside of $U$. Thus $f \sigma \equiv 0$ on $M$ and so since $\nabla$ is linear, we have $\nabla(f \sigma) \equiv 0$ on $M$. Thus since (iv) of Definition 12.1 holds for global fields, we have

$$
\begin{aligned}
\nabla_{X}(f \sigma)(q) & =f(p)\left(\nabla_{X} \sigma\right)(q)+\left(X_{q} f\right) \sigma(q) \\
& =\left(\nabla_{X} \sigma\right)(q)=0
\end{aligned}
$$

Since $q \in U$ was arbitrary, we have the result.\\
We now define a natural covariant derivative derived from a given Koszul connection $\nabla$. Given any open set $U \subset M$, we define $\nabla^{U}: \mathfrak{X}(U) \times \Gamma\left(\left.E\right|_{U}\right) \rightarrow \Gamma\left(\left.E\right|_{U}\right)$ by

$$
\left(\nabla_{X}^{U} \sigma\right)(p):=\left(\nabla_{\widetilde{X}} \widetilde{\sigma}\right)(p), \quad p \in U,
$$

for $\widetilde{X} \in \mathfrak{X}(M)$ and $\widetilde{\sigma} \in \Gamma(E)$ chosen to be any sections which agree with $X$ and $\sigma$ on some open $V$ with $p \in V \subset \bar{V} \subset U$. By the above lemma this definition does not depend on the choices of $\widetilde{X}$ and $\widetilde{\sigma}$.

Proposition 12.5. Let $E \rightarrow M$ be a rank $k$ vector bundle and suppose that $\nabla: \mathfrak{X}(M) \times \Gamma(E) \rightarrow \Gamma(E)$ is a Koszul connection. If for each open $U$ we define $\nabla^{U}$ as in (12.1) above, then the assignment $U \mapsto \nabla^{U}$ is a natural covariant derivative as in Definition 12.3.

Proof. We must show that (i), (ii) and (iii) of Definition 12.3 hold. It is easily checked that (i) holds, that is, that $\nabla^{U}$ is a Koszul connection for each $U$. The demonstration that (ii) holds is also easy and we leave it for the reader to check. Now since $X \rightarrow \nabla_{X} \sigma$ is linear over $C^{\infty}(U)$, (iii) follows by familiar arguments ( $\nabla_{X} \sigma$ is linear over functions in the argument $X$ ).

Because of this last lemma, we may define $\nabla_{v_{p}} \sigma$ for $v_{p} \in T_{p} M$ by

$$
\nabla_{v_{p}} \sigma:=\left(\nabla_{X} \sigma\right)(p)
$$

where $X$ is any vector field with $X(p)=v_{p} \in T_{p} M$. We say that $\nabla_{X} \sigma$ is "tensorial" in the variable $X$. The result can be seen as a special case of Proposition 6.55. We now see that it is safe to write expressions not directly justified by the definition of Koszul connection. For example, if $X \in \mathfrak{X}(U)$ and $\sigma \in \Gamma(V, E)$, where $U \cap V \neq \emptyset$, then $\nabla_{X} \sigma$ is taken to be an element of $\Gamma(U \cap V, E)$ defined by

$$
\left(\nabla_{X} \sigma\right)(p)=\nabla_{X_{p}} \sigma:=\nabla_{X_{p}}^{U} \sigma \text { for all } p \in U \cap V
$$

This is a particularly useful convention when $U$ is the domain of a chart $(U, \mathrm{x})$ and $X=\frac{\partial}{\partial x^{i}}$ and also when $\sigma$ is a member of a frame field of the vector bundle defined on some open set.

In the same way that one extends a derivation on vector fields to a tensor derivation, one may show that a covariant derivative on a vector bundle induces naturally related covariant derivatives on all the multilinear bundles.

In particular, if $\pi^{*}: E^{*} \rightarrow M$ denotes the dual bundle to $\pi: E \rightarrow M$ we may define connections on $\pi^{*}: E^{*} \rightarrow M$ and on $\pi \otimes \pi^{*}: E \otimes E^{*} \rightarrow M$. We do this in such a way that for $s \in \Gamma(M, E)$ and $s^{*} \in \Gamma\left(M, E^{*}\right)$ we have

$$
\nabla_{X}^{E \otimes E^{*}}\left(s \otimes s^{*}\right)=\nabla_{X} s \otimes s^{*}+s \otimes \nabla_{X}^{E^{*}} s^{*}
$$

and

$$
\left(\nabla_{X}^{E^{*}} s^{*}\right)(s)=X\left(s^{*}(s)\right)-s^{*}\left(\nabla_{X} s\right)
$$

Of course, the second formula follows from the requirement that covariant differentiation commutes with contraction:

$$
\begin{aligned}
X\left(s^{*}(s)\right) & =\left(\nabla_{X} C\left(s \otimes s^{*}\right)\right)=C\left(\nabla_{X}^{E \otimes E^{*}}\left(s \otimes s^{*}\right)\right) \\
& =C\left(\nabla_{X} s \otimes s^{*}+s \otimes \nabla_{X}^{E^{*}} s^{*}\right)=s^{*}\left(\nabla_{X} s\right)+\left(\nabla_{X}^{E^{*}} s^{*}\right)(s),
\end{aligned}
$$

where $C$ denotes the contraction given by $s \otimes \alpha \mapsto \alpha(s)$. All this works like the tensor derivation extension procedure discussed previously, and we often write all of these covariant derivatives with the single symbol $\nabla$.

The bundle $E \otimes E^{*} \rightarrow M$ is naturally isomorphic to $\operatorname{End}(E)$, and by this isomorphism we get a connection on $\operatorname{End}(E)$. If we identify elements of $\Gamma(\operatorname{End}(E))$ with $\operatorname{End}(\Gamma E)$ (see Proposition 6.55), then we may use the following formula for the definition of the connection on $\operatorname{End}(E)$ :

$$
\left(\nabla_{X} A\right)(s):=\nabla_{X}(A(s))-A\left(\nabla_{X} s\right)
$$

Indeed, since $C: s \otimes A \mapsto A(s)$ is a contraction, we must have

$$
\begin{aligned}
\nabla_{X}(A(s)) & =C\left(\nabla_{X} s \otimes A+s \otimes \nabla_{X} A\right) \\
& =A\left(\nabla_{X} s\right)+\left(\nabla_{X} A\right)(s)
\end{aligned}
$$

Notice that if we fix $s \in \Gamma(E)$, then for each $p \in M$, we have an element $(\nabla s)(p)$ of $L\left(T_{p} M, E_{p}\right) \cong E \otimes T^{*} M$ given by

$$
(\nabla s)(p): v_{p} \mapsto \nabla_{v_{p}} s \text { for all } v_{p} \in T_{p} M
$$

Thus we obtain a section $\nabla s$ of $E \otimes T^{*} M$ given by $p \mapsto(\nabla s)(p)$, which is easily shown to be smooth. In this way, we can also think of a connection as giving a map

$$
\nabla: \Gamma(E) \rightarrow \Gamma\left(E \otimes T^{*} M\right)
$$

with the property that

$$
\nabla f s=f \nabla s+s \otimes d f
$$

for all $s \in \Gamma(E)$ and $f \in C^{\infty}(M)$. Since, by definition, $\Gamma(E)=\Omega^{0}(E)$ and $\Gamma\left(E \otimes T^{*} M\right)=\Omega^{1}(E)$, we really have a map $\Omega^{0}(E) \rightarrow \Omega^{1}(E)$. Later we will extend to a map $\Omega^{k}(E) \rightarrow \Omega^{k+1}(E)$ for all integral $k \geq 0$. Now if $X$ is a smooth vector field, then $X \mapsto \nabla_{X} s \in \Gamma(E)$, so we may also view $\nabla s$ as an element of $\operatorname{Hom}(\mathfrak{X}(M), \Gamma(E))$.

\subsection*{Connection Forms}
Let $\pi: E \rightarrow M$ be a rank $r$ vector bundle with a connection $\nabla$. Recall that a choice of a local frame field over an open set $U \subset M$ is equivalent to a trivialization of the restriction $\pi_{U}:\left.E\right|_{U} \rightarrow U$. Namely, if $\phi=(\pi, \Phi)$ is such a trivialization over $U$, then defining $e_{i}(x)=\phi^{-1}\left(x, \mathrm{e}_{i}\right)$, where ( $\mathrm{e}_{i}$ ) is the standard basis of $\mathbb{F}^{n}$, we have a frame field $\left(e_{1}, \ldots, e_{k}\right)$. We now examine the expression for a given covariant derivative from the viewpoint of such a local frame. It is not hard to see that for every such frame field there must be a matrix of 1 -forms $\omega=\left(\omega_{j}^{i}\right)_{1 \leq i, j \leq r}$ such that for $X \in \Gamma(U, E)$ we may write

$$
\nabla_{X} e_{j}=\sum_{i=1}^{k} \omega_{j}^{i}(X) e_{i} .
$$

The forms $\omega_{j}^{i}$ are called connection forms.\\
Proposition 12.6. If $s=\sum_{i} s^{i} e_{i}$ is the local expression of a section $s \in \Gamma(E)$ in terms of a local frame field ( $e_{1}, \ldots, e_{k}$ ), then the following local expression holds:

$$
\nabla_{X} s=\sum_{i}\left(X s^{i}+\sum_{r} \omega_{r}^{i}(X) s^{r}\right) e_{i}
$$

Proof. We simply compute:

$$
\begin{aligned}
\nabla_{X} s & =\nabla_{X}\left(\sum_{i} s^{i} e_{i}\right) \\
& =\sum_{i}\left(X s^{i}\right) e_{i}+\sum_{i} s^{i} \nabla_{X} e_{i} \\
& =\sum_{i}\left(X s^{i}\right) e_{i}+\sum_{i, j} s^{i} \omega_{i}^{j}(X) e_{j} \\
& =\sum_{i}\left(X s^{i}\right) e_{i}+\sum_{i, r} s^{r} \omega_{r}^{i}(X) e_{i} \\
& =\sum_{i}\left(X s^{i}+\sum_{r} \omega_{r}^{i}(X) s^{r}\right) e_{i} .
\end{aligned}
$$

So the $i$-th component of $\nabla_{X} s$ is

$$
\left(\nabla_{X} s\right)^{i}=X s^{i}+\sum_{r} \omega_{r}^{i}(X) s^{r}
$$

We may surely choose $U$ small enough that it is also the domain of a coordinate frame $\left\{\partial_{\mu}\right\}$ for $M$. Thus we have

$$
\nabla_{\partial_{\mu}} e_{j}=\sum_{k} \omega_{\mu j}^{k} e_{k}
$$

where $\omega_{\mu j}^{k}=\omega_{i}^{j}\left(\partial_{\mu}\right)$. We now have the local formula

$$
\nabla_{X} s=\sum_{i=1}^{k}\left(\sum_{\mu=1}^{n} X^{\mu} \partial_{\mu} s^{i}+\sum_{\mu=1}^{n} \sum_{r=1}^{k} X^{\mu} \omega_{\mu r}^{i} s^{r}\right) e_{i} .
$$

Or, using the summation convention,

$$
\nabla_{X} s=\left(X^{\mu} \partial_{\mu} s^{i}+X^{\mu} \omega_{\mu r}^{i} s^{r}\right) e_{i}
$$

Now suppose that we have two moving frames whose domains overlap, say $u=\left(e_{1}, \ldots, e_{k}\right)$ and $u^{\prime}=\left(e_{1}^{\prime}, \ldots, e_{k}^{\prime}\right)$. Let us examine how the matrix of 1-forms $\omega=\left(\omega_{i}^{j}\right)$ is related to the corresponding $\omega^{\prime}=\left(\omega_{i}^{\prime j}\right)$ defined in terms of the frame $u^{\prime}$. The change of frame is

$$
e_{i}^{\prime}=\sum_{j} g_{i}^{j} e_{j}
$$

which in matrix notation is

$$
u^{\prime}=u g
$$

for some smooth $g: U \cap U^{\prime} \rightarrow \mathrm{GL}(n)$. (We treat $u$ as a row vector of fields.) For a given moving frame, let $\nabla u:=\left[\nabla e_{1}, \ldots, \nabla e_{k}\right]$. Differentiating both sides of $u^{\prime}=u g$ and using matrix notation we have

$$
\begin{aligned}
& u^{\prime}=u g, \\
& \nabla u^{\prime}=\nabla(u g), \\
& u^{\prime} \omega^{\prime}=(\nabla u) g+u d g, \\
& u^{\prime} \omega^{\prime}=u g g^{-1} \omega g+u g g^{-1} d g, \\
& u^{\prime} \omega^{\prime}=u^{\prime} g^{-1} \omega g+u^{\prime} g^{-1} d g,
\end{aligned}
$$

and so we obtain the transformation law for connection forms:

$$
\omega^{\prime}=g^{-1} \omega g+g^{-1} d g
$$

\subsection*{Differentiation Along a Map}
Once again let $\pi: E \rightarrow M$ be a vector bundle with a Koszul connection $\nabla$. Consider a smooth map $f: N \rightarrow M$ and a section $\sigma: N \rightarrow E$ along $f$. Let $e_{1}, \ldots, e_{k}$ be a frame field defined over $U \subset M$. Since $f$ is continuous, $O=f^{-1}(U)$ is open and $e_{1} \circ f, \ldots, e_{k} \circ f$ are fields along $f$ defined on $O$. We may write $\sigma=\sum_{a=1}^{k} \sigma^{a} e_{a} \circ f$ for some functions $\sigma^{a}: O \subset N \rightarrow \mathbb{F}$.

For any $p \in O$ and $v \in T_{p} N$, we define

$$
\nabla_{v}^{f} \sigma:=\sum_{a=1}^{k}\left(\left(d \sigma^{a} \cdot v\right)+\sum_{r=1}^{k} \omega_{r}^{a}(T f(v)) \sigma^{r}(p)\right) e_{a}(f(p)) .
$$

A direct albeit tedious calculation shows that this result is independent of the choice of frame field. Thus we obtain a global map

$$
\nabla^{f}: T N \times \Gamma_{f}(E) \rightarrow \Gamma_{f}(E) .
$$

The map $\nabla^{f}: T N \times \Gamma_{f}(E) \rightarrow \Gamma_{f}(E)$ satisfies properties that qualify it as a covariant derivative along $f$. Namely, we have

Definition 12.7. Let $\pi: E \rightarrow M$ and $f: N \rightarrow M$ be as above. A covariant derivative along $f$ is a map $\nabla^{f}: T N \times \Gamma_{f}(E) \rightarrow \Gamma_{f}(E)$ that satisfies\\
(i) $\nabla^{f}: T N \times \Gamma_{f}(E) \rightarrow \Gamma_{f}(E)$ is fiberwise linear in the first argument:

$$
\nabla_{a u+b v}^{f} \sigma=\nabla_{a u}^{f} \sigma+\nabla_{b v}^{f} \sigma
$$

for all $\sigma \in \Gamma_{f}(E)$, scalars $a, b$, and $u, v \in T_{p} N$ ( $p$ is arbitrary).\\
(ii) $\nabla_{u}^{f} \sigma=\nabla_{u}^{f} \sigma_{1}+\nabla_{u}^{f} \sigma_{2}$ for any $u \in T N$ and any $\sigma_{1}, \sigma_{2} \in \Gamma_{f}(E)$.\\
(iii) For $v \in T_{p} N, h \in C^{\infty}(N ; \mathbb{F})$, and $\sigma \in \Gamma_{f}(E)$, we have

$$
\nabla_{v}^{f}(h \sigma)=h(p) \nabla_{v}^{f} \sigma+v(h) \sigma(p) .
$$

(iv) If $U \in \mathfrak{X}(N)$, then $p \mapsto \nabla_{U(p)}^{f} \sigma$ is smooth for all $\sigma \in \Gamma_{f}(E)$.\\
(v) If $g: P \rightarrow N$ and $f: N \rightarrow M$ are smooth, then $\nabla^{f \circ g}$ is related to $\nabla^{f}$ by the following chain rule:

$$
\nabla_{u}^{f \circ g}(\sigma \circ g)=\left(\nabla_{T g \cdot u}^{f} \sigma\right)
$$

for $u \in T P$; see the following diagram:\\
\includegraphics[scale=0.25, center]{2025_10_09_265d7e1a22c89f79c21eg-523}

In the next section we give a more geometric view of covariant differentiation. Among other things, we obtain a more geometric way of producing $\nabla^{f}$ that does not appeal to a frame field. Also, we usually omit the superscript on $\nabla^{f}$ since it will be clear from context.

As a special case, consider a section $\sigma$ of $E$ along a smooth curve $c$ : $(a, b) \rightarrow M$. Then we can form the operator $\nabla_{\frac{\partial}{\partial t}} \sigma$, which is also denoted by $\frac{D}{\partial t}$ or $\nabla_{\partial_{t}}$. We have the formula

$$
\nabla_{\frac{\partial}{\partial t}} \sigma=\sum_{i}\left(\frac{d \sigma^{i}}{d t}+\sum_{r}\left(\omega_{r}^{i} \circ \dot{c}\right) \sigma^{r}\right) e_{i}=\sum_{i}\left(\frac{d \sigma^{i}}{d t}+\sum_{r} c^{*} \omega_{r}^{i}\left(\frac{\partial}{\partial t}\right) \sigma^{r}\right) e_{i},
$$

where $e_{i} \circ c$ is abbreviated to $e_{i}$.

\subsection*{Ehresmann Connections}
One way to get a natural covariant derivative on a vector bundle is through the notion of connection in another sense. What we will define in this section is a special case of what is called an Ehresmann connection. An Ehresmann connection is a structure that can be defined in the category of general smooth fiber bundles, but the two most common instances are defined for vector bundles and for principal bundles, and in each case extra hypotheses are added to the definition. We work with vector bundles here and give an explanation of the principal bundle approach in the online supplement [Lee, Jeff].

We first describe the construction in words since the notation tends to obscure what is really a simple geometric idea. First, notice that an ordinary function of one variable is constant on an interval if its graph is horizontal over that interval. For sections of a vector bundle, there is no a priori notion of "horizontal", so we have no principled way to decide what sections should be considered constant. Consider a moving frame field $e_{1}, \ldots, e_{k}$ on some vector bundle. If we had a reason to declare these frame fields to be constant, then we could differentiate a general section $s=\sum s^{i} e_{i}$ in a direction $v$ by the rule $D_{v} s=\sum\left(v s^{i}\right) e_{i}$. So it seems clear that having a geometrically motivated way of picking out which sections should be constant will lead to a way of differentiating sections. The first step is to notice that we have a natural notion of vertical for a vector bundle $\pi: E \rightarrow M$. The tangent spaces of the fibers of $E$ are to be considered vertical. Thus the vertical subspace of $T_{y} E$ is $T_{y} E_{p} \subset T_{y} E$, where $\pi(y)=p$.

Now if we have a vector in $T_{y} E$ for some $y \in E$, we would like to project it onto the vertical direction. However, this entails having a complementary horizontal space along which we project. The idea is then to assume that there is a distribution on $E$, that is, a subbundle of $T E$, that is everywhere complementary to the vertical directions. Thus we obtain a subspace complementary to the vertical, which allows projection onto the vertical. Once we have this we can say that a section $s \in \Gamma(E)$ is horizontal (i.e. constant) along a curve $c: I \rightarrow M$ if $T s \cdot \dot{c}(t)$ has a vertical projection of zero for all $t$. Now we can define a covariant derivative as follows. If $v \in T_{p} M$ and $s$ is a smooth section, then we can project $T_{p} s \cdot v$ onto the vertical space $T_{s(p)} E_{p}$ obtaining an element $\nabla_{v} s$. But $E_{p}$ is a vector space, and so we may identify $T_{s(p)} E_{p}$ with $E_{p}$ and take $\nabla_{v} s$ to be an element of $E_{p}$. The resulting map $v \mapsto \nabla_{v} s$ turns out to define a covariant derivative. We will also consider sections along general smooth maps $f: N \rightarrow M$ and obtain covariant differentiation for sections along $f$.

There is a canonical "parallelism" on $\mathbb{R}^{n}$. Recall the canonical map $\jmath_{p}: \mathbb{R}^{n} \rightarrow T_{p} \mathbb{R}^{n}$. A vector $v_{1} \in T_{p} \mathbb{R}^{n}$ is parallel to a vector $v_{2} \in T_{q} \mathbb{R}^{n}$ exactly

\begin{figure}[h]
\begin{center}
  \includegraphics[scale=0.25]{2025_10_09_265d7e1a22c89f79c21eg-525}

\caption{Figure 12.1. Vertical space}
\end{center}
\end{figure}

when $\left(\jmath_{q} \circ{\jmath_{p}^{-1}}^{-1}\right)\left(v_{1}\right)=v_{2}$. The map $\jmath_{q} \circ \jmath_{p}^{-1}$ is a special example of what is called "parallel translation" and in this case establishes what is sometimes called "distant parallelism". In the presence of a connection, we obtain a map between fibers of a vector bundle which is called parallel translation. But this map may depend on a choice of smooth curve connecting the base points. Locally, this path dependence of parallel translation is due to the curvature of the connection.

We now proceed more formally. We give the definition of vertical bundle not just for a vector bundle, but also for a general fiber bundle. First a lemma:

Lemma 12.8. Let $E \xrightarrow{\pi} M$ be a fiber bundle with typical fiber a $k$-manifold $F$. Fix $p \in M$ and let $\imath: E_{p} \hookrightarrow E$ be the inclusion. For all $y \in E_{p}$, we have

$$
T_{y} \iota\left(T_{y} E_{p}\right)=\operatorname{Ker}\left[T_{y} \pi: T_{y} E \rightarrow T_{p} M\right]=\left(T_{y} \pi\right)^{-1}\left(0_{p}\right) \subset T_{y} E,
$$

where $0_{p} \in T_{p} M$ is the zero vector. If $\varphi: E_{p} \rightarrow F$ is a diffeomorphism and $(V, \mathrm{x})$ a chart on $F$, then for all $y \in \varphi^{-1}(V), d \mathrm{x} \circ T_{y} \varphi$ maps $T_{y} E_{p}$ isomorphically onto $\mathbb{R}^{k}$.

Proof. $\pi \circ \imath \circ \gamma$ is constant for each smooth curve $\gamma$ in $N$, so $T \pi \cdot(T \imath \cdot \dot{\gamma}(0))= 0_{p}$. Thus $T_{y} \imath\left(T_{y} E_{p}\right) \subset\left(T_{y} \pi\right)^{-1}\left(0_{p}\right)$. On the other hand,

$$
\operatorname{dim}\left(\left(T_{y} \pi\right)^{-1}\left(0_{p}\right)\right)=\operatorname{dim} E-\operatorname{dim} M=\operatorname{dim} F=\operatorname{dim} T_{y} E_{p},
$$

so $\left(T_{y} \imath\right)\left(E_{p}\right)=\left(T_{y} \pi\right)^{-1}\left(0_{p}\right)$. The rest is clear since $d \mathrm{x}$ is just $T \mathrm{x}$ followed by the projection $T \mathrm{x}(U)=\mathrm{x}(U) \times \mathbb{R}^{k} \rightarrow \mathbb{R}^{k}$.

Note: As usual we identify $T_{y} \imath\left(T_{y} E_{p}\right)$ with $T_{y} E_{p}$.

Definition 12.9. Let $\pi: E \rightarrow M$ be a fiber bundle with typical fiber $F$ and $\operatorname{dim} F=k$. Let $\mathcal{V}_{y} E:=\left(T_{y} \pi\right)^{-1}\left(0_{p}\right)$ where $\pi(y)=p$. The vertical bundle on $\pi: E \rightarrow M$ is the real vector bundle $\pi_{\mathcal{V}}: \mathcal{V} E \rightarrow E$ with total space defined by the disjoint union

$$
\mathcal{V} E:=\bigcup_{y \in E} \mathcal{V}_{y} E \subset T E .
$$

The projection map is defined by the restriction $\pi_{\mathcal{V}}:=\left.\pi_{T E}\right|_{\mathcal{V} E}$. A vector bundle atlas on $\mathcal{V} E$ is given by vector bundle charts of the form

$$
\left(\pi_{\mathcal{V}}, d \mathrm{x} \circ T \Phi\right): \pi_{\mathcal{V}}^{-1}\left(\pi^{-1}(U) \cap \Phi^{-1}(V)\right) \rightarrow\left(\pi^{-1}(U) \cap \Phi^{-1}(V)\right) \times \mathbb{R}^{k},
$$

where $\phi=(\pi, \Phi)$ is a bundle chart on $E$ over $U$ and $(V, \mathrm{x})$ a chart in $F$.\\
For the following, refer to the diagram:\\
\includegraphics[scale=0.25, center]{2025_10_09_265d7e1a22c89f79c21eg-526(1)}\\
\includegraphics[scale=0.25, center]{2025_10_09_265d7e1a22c89f79c21eg-526}

Exercise 12.10. Prove: Let $f: N \rightarrow M$ be a smooth map and $\pi: E \rightarrow M$ a fiber bundle with typical fiber $F$. Then $\mathcal{V} f^{*} E \rightarrow f^{*} E$ is bundle isomorphic to $\widetilde{f}^{*} \mathcal{V} E \rightarrow f^{*} E$ where $\widetilde{f}:=\left.\operatorname{pr}_{2}\right|_{f^{*} E}: f^{*} E \rightarrow E$ and $\operatorname{pr}_{2}: M \times E \rightarrow E$.

Now consider the pull-back bundle $f^{*} E$ where $f$ is as above. Since $f^{*} E$ is the submanifold of $N \times E$ defined by the condition that ( $q, y$ ) is in $f^{*} E$ if and only if $f(q)=\pi(y)$, a curve in $f^{*} E$ must be of the form $\left(c_{1}, c_{2}\right)$, where $c_{1}$ is a curve in $N$ and $c_{2}$ is a curve in $E$, and it must also be the case that $f \circ c_{1}=\pi \circ c_{2}$. Now if $\mathrm{pr}_{1}$ and $\mathrm{pr}_{2}$ are the first and second factor projections from $N \times E$, then ( $T \mathrm{pr}_{1}, T \mathrm{pr}_{2}$ ) gives a vector bundle isomorphism of the bundle $T(N \times E) \rightarrow N \times E$ with the bundle $T N \times T E \rightarrow N \times E$, and so we expect that under this isomorphism $T\left(f^{*} E\right)$ corresponds to a subbundle of $T N \times T E \rightarrow N \times E$.

Exercise 12.11. Show that under the bundle isomorphism $\left(T \operatorname{pr}_{1}, T \operatorname{pr}_{2}\right)$ : $T(N \times E) \cong T N \times T E$, the tangent bundle $T\left(f^{*} E\right)$ corresponds to $\{(v, w) \in T N \times T E: T f \cdot v=T \pi \cdot w\}$. Under this isomorphism $\left(\mathcal{V} f^{*} E\right)_{(q, y)}$ corresponds to $\left\{0_{q}\right\} \times \mathcal{V}_{y} E$.

We need to make an observation. If V is a complex vector bundle and $x \in \mathrm{V}$, then the tangent space $T_{x} \mathrm{V}$ has a natural complex structure. Indeed, the map $\lambda_{x}: \mathrm{V} \rightarrow T_{x} \mathrm{V}$ is used to transfer the complex structure of V to that of $T_{x} \mathrm{V}$. In particular, if $\pi: E \rightarrow M$ is a complex vector bundle, then we may view $T_{y} E_{p}=\left(T_{y} \pi\right)^{-1}\left(0_{p}\right)$ as a complex vector space. Thus $\mathcal{V} E$ is a complex vector bundle.

The vertical vector bundle $\mathcal{V} E$ is isomorphic to the vector bundle $\pi^{*} E$ over $E$ (we say that $\mathcal{V} E$ is isomorphic to $E$ along $\pi$ ). To see this, note that if $(v, w) \in \pi^{*} E$, then $\pi(v+t w)$ is constant in $t$. From this we see that the map from $\pi^{*} E$ to $T E$ given by $(v, w) \mapsto d /\left.d t\right|_{0}(v+t w)$, maps into $\mathcal{V} E$. It is easy to see that this map is a vector bundle isomorphism:

$$
\begin{aligned}
& \jmath: \pi^{*} E \cong \mathcal{V} E \\
& \jmath:(y, w) \mapsto \jmath_{y} w:=\left.\frac{d}{d t}\right|_{0}(y+t w)=w_{y} .
\end{aligned}
$$

The meaning of the symbol $y$ depends on the bundle we have in mind, but it is consistent with our previous use in the sense that if $E$ is the tangent bundle of an open set in a vector space, then $j_{y}$ is the canonical isomorphism as before.

Definition 12.12. A (linear Ehresmann) connection on a vector bundle $\pi: E \rightarrow M$ is a smooth distribution $\mathcal{H}$ on the total space $E$ such that\\
(i) $\mathcal{H}$ is complementary to the vertical bundle:

$$
T E=\mathcal{H} \oplus \mathcal{V} E
$$

(ii) $\mathcal{H}$ is homogeneous: $T_{y} \mu_{r}\left(\mathcal{H}_{y}\right)=\mathcal{H}_{r y}$ for all $y \in E, r \in \mathbb{R}$, where $\mu_{r}: E \rightarrow E$ is the multiplication map given by $\mu_{r}: y \mapsto r y$.\\
The subbundle $\mathcal{H}$ is called the horizontal distribution (or horizontal subbundle).

The statement $T E=\mathcal{H} \oplus \mathcal{V} E$ means that for every $y \in E$ we have the internal direct sum decomposition $T_{y} E=\mathcal{H}_{y} \oplus \mathcal{V}_{y} E$. Any $v \in T E$ has a corresponding decomposition $v=\mathrm{p}_{\mathrm{v}} v+\mathrm{p}_{\mathrm{h}} v$. Here, $\mathrm{p}_{\mathrm{v}}: v \mapsto \mathrm{p}_{\mathrm{v}} v$ and $\mathrm{p}_{\mathrm{h}}: v \mapsto \mathrm{p}_{\mathrm{h}} v$ are the obvious projections onto $\mathcal{V}_{y} E$ and $\mathcal{H}_{y}$ referred to respectively as the horizontal and vertical projections. Note well that without a choice of horizontal distribution $\mathcal{H}$ there is no vertical projection $\mathrm{p}_{\mathrm{v}}=1-\mathrm{p}_{\mathrm{h}}$. For $y \in E$, an individual element $w \in T_{y} E$ is horizontal if $w \in \mathcal{H}_{y}$ and vertical if $w \in \mathcal{V}_{y} E$. A vector field $\widetilde{X} \in \mathfrak{X}(E)$ is said to be a horizontal vector field (resp. vertical vector field) if $\widetilde{X}(y) \in \mathcal{H}_{y}$ (resp. $\widetilde{X}(y) \in \mathcal{V}_{y} E$ ) for all $y \in E$. On the right hand side of Figure 12.2 we see a schematic representation of the field of horizontal spaces. Notice that these spaces are tangent to the zero section (which we identify with $M$ ). This must be the case and is a consequence of the homogeneity condition (ii) from the definition. On the left hand side we see a particular tangent space to $E$ together with a vector and its vertical projection.

Exercise 12.13. Show that $\mathcal{H} \cong \pi^{*} T M$.\\
Theorem 12.14. Every vector bundle admits a connection.

\begin{figure}[h]
\begin{center}
  \includegraphics[scale=0.25]{2025_10_09_265d7e1a22c89f79c21eg-528(1)}

\caption{Figure 12.2. Horizontal distribution}
\end{center}
\end{figure}

Proof. First notice that we may easily define a connection on a trivial bundle $\operatorname{pr}_{1}: M \times \mathrm{V} \rightarrow M$. Given a fixed $v \in \mathrm{V}$, let $i_{v}: M \rightarrow M \times \mathrm{V}$ be defined by $i_{v}(p):=(p, v)$. Next define $\mathcal{H}_{(p, v)}:=T i_{v}\left(T_{p} M\right)$ for each $p$. We then have $T \operatorname{pr}_{1}\left(\mathcal{H}_{(p, v)}\right)=T_{p} M$. Since for any scalar $a$ we have $\mu_{a} \circ i_{v}=i_{a v}$, we also have $T \mu_{a} \circ T i_{v}=T i_{a v}$ so that

$$
T \mu_{a}\left(\mathcal{H}_{(p, v)}\right)=T \mu_{a}\left(T i_{v}\left(T_{p} M\right)\right)=T i_{a v}\left(T_{p} M\right)=\mathcal{H}_{(p, a v)}=\mathcal{H}_{a(p, v)} .
$$

For a general vector bundle $\pi: E \rightarrow M$, let $\left\{U_{\alpha}\right\}$ be a locally finite cover of $M$ such that the bundle is trivial over each $U_{\alpha}$. Then we may choose a connection $\mathcal{H}^{\alpha}$ on each $\pi^{-1}\left(U_{\alpha}\right)$. Let $\left\{\rho_{\alpha}\right\}$ be a partition of unity subordinate to $\left\{U_{\alpha}\right\}$. For each $y \in E$, define $L_{y}: T_{\pi(y)} M \rightarrow T_{y} E$ by

$$
L_{y}(v):=\sum_{\left\{\alpha: \pi(y) \in U_{a}\right\}} \rho_{\alpha}(\pi(y)) w_{\alpha},
$$

where for each $\alpha$, the vector $w_{\alpha}$ is the unique vector in $\mathcal{H}^{\alpha}$ such that $T \pi \cdot w_{\alpha}= v$. It is easy to check that $L_{y}$ is linear and satisfies $T_{y} \pi \circ L_{y}=\operatorname{id}_{T_{p} M}$. The distribution we seek is then defined by $L_{y}\left(T_{\pi(y)} M\right)$ for each $y$. We leave it to the reader to check that this distribution is smooth and satisfies the required conditions. \(\square\)

Theorem 12.15. If $\mathcal{H}$ is a connection on $\pi: E \rightarrow M$ and $f: N \rightarrow M$ a smooth map, then the distribution $f^{*} \mathcal{H}=(T \widetilde{f})^{-1}(\mathcal{H})$, where $\widetilde{f}:=\left.\operatorname{pr}_{2}\right|_{f^{*} E}$, defines a connection on the pull-back bundle $f^{*} E \rightarrow N$. This is referred to as the pull-back connection:\\
\includegraphics[scale=0.25, center]{2025_10_09_265d7e1a22c89f79c21eg-528}

Proof. First note that $\widetilde{f}$ is the restriction of the projection $N \times E \rightarrow E$ and so $f^{*} \mathcal{H}:=(T \widetilde{f})^{-1}(\mathcal{H})$ can be defined: For $(q, y) \in f^{*} E$, we let $\left(f^{*} \mathcal{H}\right)_{(q, y)}= \left(T_{(q, y)} \widetilde{f}\right)^{-1} \mathcal{H}_{y}$. By Exercise 12.11 we can identify $T\left(f^{*} E\right)$ with $\{(v, w) \in T N \times T E: T f \cdot v=T \pi \cdot w\}$. Under this isomorphism, $\left(\mathcal{V} f^{*} E\right)_{(q, y)}$ corresponds to $\left\{0_{q}\right\} \times \mathcal{V}_{y} E$ while $\left(f^{*} \mathcal{H}\right)_{(q, y)}$ corresponds to $\{(v, w) \in T N \times \mathcal{H}: T f \cdot u= T \pi \cdot v\}$. Thus we have the decomposition

$$
\left(v, \mathrm{p}_{\mathrm{h}} w\right)+\left(0, \mathrm{p}_{\mathrm{v}} w\right)
$$

which is easily seen to be unique. Hence the distribution $f^{*} \mathcal{H}$ is complementary to $\mathcal{V} f^{*} E$. Under the same identification, multiplication $m_{a}$ on $f^{*} E$ is $m_{a}(q, y)=\left(q, \mu_{a} y\right)$ and $T m_{a}(v, w)=\left(v, T \mu_{a} w\right)$, which makes it clear that the second defining condition for a connection holds for $f^{*} \mathcal{H}$ since it holds for $\mathcal{H}$.

Given a vector $v \in T_{p} M$ and a choice of $y \in E_{p}$, there is a unique vector $v_{y} \in \mathcal{H}_{y} \subset T_{y} E$ such that $T_{y} \pi \cdot v_{y}=v$. This vector is called the horizontal lift of $v$ to $T_{y} E$. The idea works for fields too. Given a vector field $X \in \mathfrak{X}(M)$, there is a unique vector field $\widetilde{X} \in \mathfrak{X}(E)$ such that $\widetilde{X}(y)$ is horizontal for all $y \in M$ and $T_{y} \pi \cdot \widetilde{X}(y)=X(\pi(y))$. Thus $\widetilde{X} \in \Gamma(\mathcal{H}) \subset \mathfrak{X}(E)$. This horizontal vector field on the total space $E$ is called the horizontal lift of $X$. Clearly $\widetilde{X}$ is $\pi$-related to $X$.

Proposition 12.16. Let $\pi: E \rightarrow M$ be a vector bundle with connection $\mathcal{H}$. If $X, Y \in \mathfrak{X}(M)$ and $f \in C^{\infty}(M)$, then\\
(i) $a \widetilde{X+b} Y=a \widetilde{X}+b \widetilde{Y}$ for all $a, b \in \mathbb{R}$;\\
(ii) $\widetilde{f X}=(f \circ \pi) \widetilde{X}$;\\
(iii) $\widetilde{[X, Y]}=\mathrm{p}_{\mathrm{h}}[\widetilde{X}, \widetilde{Y}]$.

Proof. (i) and (ii) are obvious. (iii) follows from the easy to check equalities $\pi_{*}[\widetilde{X}, \widetilde{Y}]=[X, Y]=\pi_{*} \widehat{[X, Y]}$ and the uniqueness of horizontal lifts.

Definition 12.17. Let $\sigma: N \rightarrow E$ be a section of $E$ along a map $f: N \rightarrow M$. We say that $\sigma$ is a parallel section if $T \sigma \cdot v$ is horizontal for all $v \in T N$. If $s$ is a section of $E$ and $c: I \rightarrow E$ is a curve, then we say that $s$ is parallel along $c$ provided $s \circ c$ is parallel.

Exercise 12.18. Recall that a section of $f^{*} E$ must have the form $s: q \mapsto \left(q, \sigma_{s}(q)\right)$, where $\sigma$ is a section along $f$. Show that if $s$ is parallel with respect to the pull-back connection on $f^{*} E$, then $\sigma_{s}$ is parallel.

Exercise 12.19. Let $[0, b]$ be an interval and let $t$ denote the standard coordinate on $[0, b]$. Suppose that $\pi: E \rightarrow[0, b]$ is a vector bundle over $[0, b]$ with connection. Let $\widetilde{\partial}$ denote the horizontal lift of $\partial / \partial t$.\\
(a) Show that if $c:[0, a] \rightarrow E$ is an integral curve of $\widetilde{\partial}$, then $c(a) \in E_{a}$. [Hint: Show that $\pi \circ c$ is an integral curve of $\partial / \partial t$.]\\
(b) Let $0 \leq t_{0}<b$. Show that there is a fixed $\epsilon>0$ depending only on $t_{0}$ such that all maximal integral curves of $\widetilde{\partial}$ originating in the fixed fiber $E_{t_{0}}$ are defined at least on $\left[t_{0}, \epsilon\right)$. [Hint: Endow $E$ with a bundle metric and consider all integral curves originating in the unit sphere in $E_{t_{0}}$ and then use property (ii) of Definition 12.12.]\\
(c) Using (a) and (b), show that integral curves of $\widetilde{\partial}$ all have domain equal to $[0, b]$.

Finding a horizontal lift of a vector field in $\mathfrak{X}(M)$ is trivial and automatic. On the other hand, finding a parallel section along a given map with prescribed value $\sigma(q)$ for some $q$ is generally nontrivial and may not exist. However, we have the following

Theorem 12.20. Let $\pi: E \rightarrow M$ be a (smooth) vector bundle with a connection $\mathcal{H}$. Suppose that $c:[a, b] \rightarrow M$ is a smooth curve. For each $u \in E_{c(a)}$, there is a unique parallel section $\sigma_{c, u}$ along $c$ such that $\sigma_{c, u}(a)=u$. If $P_{c}: E_{c(a)} \rightarrow E_{c(b)}$ denotes the map which takes $u \in E_{c(a)}$ to $\sigma_{c, u}(b)$, then $P_{c}$ is linear.

Proof. Without loss of generality we may take $a=0$. Let $\partial=\partial / \partial t$ denote the standard coordinate vector field on the interval $[0, b]$ and let $\widetilde{\partial}$ denote its horizontal lift in the pull-back bundle $c^{*} E$ with respect to the pull-back connection $c^{*} \mathcal{H}$. Let $c_{u}$ denote the maximal integral curve of $\widetilde{\partial}$ in $c^{*} E$ with $c_{u}(0)=(0, u) \in c^{*} E$. We have

$$
\frac{d}{d t}\left(\operatorname{pr}_{1} \circ c_{u}\right)=T \operatorname{pr}_{1} \circ \dot{c}_{u}=T \operatorname{pr}_{1} \circ \tilde{\partial} \circ c_{u}=\partial \circ \operatorname{pr}_{1} \circ c_{u}
$$

Thus $\operatorname{pr}_{1} \circ c_{u}$ is an integral curve of $\partial=\partial / \partial t$ and so $\operatorname{pr}_{1} \circ c_{u}(t)=t$. From this we see that $c_{u}(t)=\left(t, \mathrm{pr}_{2} \circ c_{u}\right)$. By Exercise 12.19 we know that $c_{u}$ is defined on $[0, b]$ and that $c_{u}(b) \in\left(c^{*} E\right)_{b}$. Let $\sigma_{c, u}:=\operatorname{pr}_{2} \circ c_{u}$ on $[0, b]$. Then $\sigma_{c, u}$ is a section of $E \rightarrow M$ along $c$ which is parallel since $\dot{c}_{u}$ is horizontal (see Exercise 12.18). We define $P_{c} u:=\sigma_{c, u}(b)$ for $u \in E_{c(0)}$. The uniqueness follows from the uniqueness of integral curves and we leave this to the reader. For any $r \in \mathbb{R}$, the field $r \sigma_{c, u}$ is parallel. Indeed, $\left(r \sigma_{c, u}\right)^{\cdot}=T \mu_{r} \circ \dot{\sigma}_{c, u}$ is horizontal since $T \mu_{r}$ preserves $\mathcal{H}$. But then $P_{c}(r u)=r P_{c}(u)$ so $P_{c}$ is homogeneous. We aim to show that $P_{c}=\jmath_{0}^{-1} \circ T P_{c} \circ \jmath_{0}$, which is a composition of linear maps. Let $v_{0} \in T_{0} E_{c(0)}$ so that $v_{0}=\dot{\gamma}(0)$, where $\gamma$ is defined by $\gamma(t)=t v$ for an appropriate $v \in E_{c(0)}$. Then $j_{0}^{-1} v_{0}=v$. We have

$$
T_{0} P_{c} v_{0}=\frac{d}{d t}\left(P_{c} \circ \gamma\right)(0)
$$

But also $P_{c} \circ \gamma(t)=P_{c}(t v)=t P_{c}(v)$ so that

$$
T_{0} P_{c} v_{0}=\jmath_{0}\left(P_{c}(v)\right)=\jmath_{0} \circ P_{c} \circ \jmath_{0}^{-1} v_{0}
$$

Thus $\jmath_{0} \circ P_{c} \circ \jmath_{0}^{-1}=T_{0} P_{c}$ or $P_{c}=\jmath_{0}^{-1} \circ T_{0} P_{c} \circ \jmath_{0}$, and we conclude that $P_{c}$ is linear.

Since $P_{c}$ has inverse $P_{c^{\leftarrow}}$ where $c^{\leftarrow}(t):=c(b-t)$, we see that it is a linear isomorphism.

Definition 12.21. Let $c:[a, b] \rightarrow M$ be as in the theorem above. The map $P_{c}$ is called parallel translation along $c$ from $c(a)$ to $c(b)$. Let $c$ be any smooth curve in $M$. For $t_{1}$ and $t_{2}$ in the domain of $c$, let $P(c)_{t_{1}}^{t_{2}}: E_{c\left(t_{1}\right)} \rightarrow E_{c\left(t_{2}\right)}$ be defined as $P(c)_{t_{1}}^{t_{2}}:=P_{c \mid\left[t_{1}, t_{2}\right]}$ if $t_{2} \geq t_{1}$ and $P(c)_{t_{1}}^{t_{2}}:=P_{c \mid\left[t_{2}, t_{1}\right]}^{-1}$ if $t_{1} \geq t_{2}$.

We also say that the curve $\sigma_{c, u}$ of Theorem 12.20 is a parallel lift or horizontal lift of the curve $c$. The map $P_{c}$ is also sometimes called parallel transport. Suppose that $c:[a, b] \rightarrow M$ is a (continuous) piecewise smooth curve (Definition 2.121). Thus we may find a monotonic sequence $t_{0}, t_{1}, \ldots, t_{j}=t$ such that $c_{i}:=\left.c\right|_{\left[t_{i-1}, t_{i}\right]}$ (or $\left.c\right|_{\left[t_{i}, t_{i-1}\right]}$ ) is smooth. ${ }^{1}$ In this case we define

$$
P(c)_{t_{0}}^{t}:=P(c)_{t_{j-1}}^{t} \circ \cdots \circ P(c)_{t_{0}}^{t_{1}}
$$

Exercise 12.22. The map $P(c)_{t_{0}}^{t}: E_{c\left(t_{0}\right)} \rightarrow E_{c(t)}$ is a linear isomorphism for all $t$ with inverse $P(c)_{t}^{t_{0}}$.

Parallel transport behaves nicely with respect to reparametrization. This is the content of the next theorem, which we ask the reader to prove in Problem 1.\\
\includegraphics[scale=0.25, center]{2025_10_09_265d7e1a22c89f79c21eg-531}

Theorem 12.23. Let $c:[a, b] \rightarrow M$ be a smooth curve. If $r:\left[a^{\prime}, b^{\prime}\right] \rightarrow[a, b]$ is a smooth map with $d r / d t>0$, then for $\gamma:=c \circ r$ we have $P_{c}=P_{\gamma}$.

\footnotetext{${ }^{1}$ It may be that $t<t_{0}$.
}It is often convenient to use special vector bundle charts constructed using parallel translation. Let ( $U, \mathrm{x}$ ) be a chart on $M$ centered at $p \in M$ and such that $\mathrm{x}(U)$ is a ball $B_{R}(0)$ of radius $R$. Consider the family of curves $c_{\mathbf{u}}:[0, R] \rightarrow M$ given for each $\mathbf{u} \in S^{n-1}(R)=\partial B_{R}(0)$ by

$$
c_{\mathbf{u}}(t):=\mathrm{x}^{-1}(t \mathbf{u})
$$

Define a frame field $\left(e_{1}, \ldots, e_{k}\right)$ for $E$ over $U$ as follows: Let $\left(e_{1}(0), \ldots, e_{k}(0)\right)$ be an ordered basis of $E_{p}$. If $q \in U$, then let $e_{i}(q)$ be defined as

$$
e_{i}(q):=P\left(c_{\mathbf{u}}\right)_{0}^{t}\left(e_{i}(0)\right),
$$

where ( $\mathbf{u}, t$ ) is the unique element of $S^{n-1}(R) \times[0,1]$ such that $c_{\mathbf{u}}(t)=q$. We say that $\left(e_{1}, \ldots, e_{k}\right)$ is radially parallel with respect to the spherical chart ( $U, \mathrm{x}$ ). Notice that by composing with a dilation of $\mathbb{R}^{n}$ if necessary, we may choose $R$ to be any positive number. Also, if we use a radially parallel frame field to define a vector bundle chart $\phi=(\pi, \Phi)$ in the usual way, then $\Phi$ will be constant along each curve $c_{\mathbf{u}}$.

We still have more to learn about the geometric meaning of a connection. We will eventually be led to the notion of curvature. At this point let us just consider what it means when a connection is integrable as a distribution.

Definition 12.24. A connection on a vector bundle is called flat if it is integrable as a distribution.

Let us agree to call a connection on a vector bundle $E$ a trivial connection if given any $u \in E$ there is a parallel section $s$ such that $s(\pi(u))=u$. Finding such sections is not always possible (even locally). In fact, we have the following characterization.

Theorem 12.25. Let $\pi: E \rightarrow M$ be a vector bundle with connection $\mathcal{H}$. The following assertions are equivalent:\\
(i) For any simply connected open set $U \subset M$, the restriction of $\mathcal{H}$ to $\pi^{-1}(U)$ is a trivial connection on the restricted vector bundle $\pi^{-1}(U) \rightarrow U$.\\
(ii) $\mathcal{H}$ is flat.

Proof. If (i) holds, then given any $u \in E$ there is a parallel section $s$ defined on $U$. Then it is easy to see that $s(U)$ is an integral manifold of the distribution $\mathcal{H}$. Since we can find such an integral manifold through any $u$, we see that $\mathcal{H}$ is integrable (i.e. flat).

If $\mathcal{H}$ is flat, then certainly $\left.\mathcal{H}\right|_{U}$ is flat for any open $U$. Suppose that $U$ is simply connected. By the Frobenius theorem there is a maximal integral manifold $L_{u}$ of $\left.\mathcal{H}\right|_{U}$ through any $u \in \pi^{-1}(U)$. By Theorem 12.20, we see that any smooth path in $U$ can be lifted uniquely to $L_{u}$. This is enough to\\
imply that $\left.\pi\right|_{L_{u}}: L_{u} \rightarrow U$ is a covering space (see [Span], 2.4.10). Since $U$ is simply connected, the lifting theorem for covering spaces (Theorem 1.95) implies that $\left.\pi\right|_{L_{u}}$ has an inverse. The desired parallel section is then $s:=\left(\left.\pi\right|_{L_{u}}\right)^{-1}$.

We now come to the task of relating covariant derivatives with Ehresmann connections on vector bundles. Denote the vector bundle isomorphism from $\mathcal{V} E$ to $E$ along $\pi$ by p:

$$
\begin{aligned}
& \mathrm{p}: \mathcal{V} E \rightarrow E \\
& \mathrm{p}: w_{y} \mapsto w
\end{aligned}
$$

For each $y$, the map p just gives the canonical identification of $T_{y} E_{p}$ with $E_{p}$, and on each fiber, it is the inverse of $\jmath$. If we have a connection on $\pi: E \rightarrow M$, then we have an associated connector (or connection map), which is the map $\kappa: T E \rightarrow E$ defined by

$$
\kappa(v):=\mathrm{p}\left(\mathrm{p}_{\mathrm{v}} v\right)=\jmath_{y}^{-1}\left(\mathrm{p}_{\mathrm{v}} v\right)
$$

for $v \in T_{y} E$. The connector is a vector bundle homomorphism along the map $\pi: E \rightarrow M$ :\\
\includegraphics[scale=0.25, center]{2025_10_09_265d7e1a22c89f79c21eg-533}

An interesting fact is that given the appropriate definition of vector space structure on the fibers, $T E$ is also a vector bundle over $T M$ via the map $T \pi$ (Problem 11 from Chapter 6). Recall that the addition and scalar multiplication on a fiber $T \pi^{-1}(x)$ of this bundle are defined by

$$
\begin{aligned}
& u \boxplus v:=T \alpha \cdot(u, v) \text { for } u, v \in T E \text { with } T \pi \cdot u=T \pi \cdot v=x, \\
& c \odot v:=T \mu_{c} \cdot v \text { for } v \in T E \text { and } c \in \mathbb{F},
\end{aligned}
$$

where $\alpha\left(y_{1}, y_{2}\right):=y_{1}+y_{2}$ for $\left(y_{1}, y_{2}\right) \in E \oplus E$ and $\mu_{c} y:=c y$ for $y \in E$ and $c \in \mathbb{F}$.

Lemma 12.26. Suppose that $f: \mathbb{R}^{K} \rightarrow \mathbb{R}^{k}$ is a smooth map such that $f(a v)=a f(v)$ for all $v \in \mathbb{R}^{K}$ and $a \in \mathbb{R}$. Then $f$ is linear. Similarly, if $f: \mathbb{C}^{K} \rightarrow \mathbb{C}^{k}$ is a smooth map such that $f(a v)=a f(v)$ for all $v \in \mathbb{C}^{K}$ and $a \in \mathbb{C}$, then $f$ is linear.

Proof. Let $f: \mathbb{R}^{K} \rightarrow \mathbb{R}^{k}$ as in the statement. Then $D f(0) v=\left.\frac{d}{d t}\right|_{t=0} f(t v) =\left.\frac{d}{d t}\right|_{t=0} t f(v)=f(v)$. Thus $f=D f(0)$ and so $f$ is linear. If $f: \mathbb{C}^{K} \rightarrow \mathbb{C}^{k}$ is a smooth map such that $f(a v)=a f(v)$ for all $v \in \mathbb{C}^{K}$ and $a \in \mathbb{C}$ then the first part shows that $f: \mathbb{C}^{K} \rightarrow \mathbb{C}^{k}$ is linear over $\mathbb{R}$. But by hypothesis $f(i v)=i f(v)$, so $f$ is actually complex linear.

Corollary 12.27. Let $\pi_{1}: E_{1} \rightarrow M_{1}$ and $\pi_{2}: E_{2} \rightarrow M_{2}$ be $\mathbb{F}$-vector bundles. Let $\widehat{f}: E_{1} \rightarrow E_{2}$ be a fiber bundle morphism over $f: M_{1} \rightarrow M_{2}$. If $\widehat{f}$ is homogeneous on each fiber, so that $\widehat{f}(a v)=a \widehat{f}(v)$ for all $v \in E_{1}$ and $a \in \mathbb{F}$, then $\widehat{f}$ is linear on fibers, and so it is a vector bundle morphism.

Proof. Use vector bundle charts and Lemma 12.26. \(\square\)

Lemma 12.28. Let $\mu_{r}: E \rightarrow E$ be multiplication by $r$. Then for any $p \in M$ and $y, w \in E_{p}$ we have

$$
T \mu_{r}\left(\jmath_{y} w\right)=\jmath_{r y}(r w)=r \jmath_{r y} w .
$$

Proof. We have

$$
\begin{aligned}
T \mu_{r}\left(\jmath_{y} w\right) & =\left.\frac{d}{d t}\right|_{t=0} \mu_{r}(y+t w)=\left.\frac{d}{d t}\right|_{t=0}(r y+t r w) \\
& =\jmath_{r y}(r w)=r \jmath_{r y} w
\end{aligned}
$$ \(\square\)

The connector $\kappa$ gives a vector bundle homomorphism along $\pi_{T M}$ : $T M \rightarrow M$. More precisely, we have

Theorem 12.29. If $\kappa$ is the connector for a connection on a vector bundle $\pi: E \rightarrow M$, then $\kappa$ gives a vector bundle homomorphism from the bundle $T \pi: T E \rightarrow T M$ to the bundle $\pi: E \rightarrow M$ along the map $\pi_{T M}: T M \rightarrow M$,\\
\includegraphics[scale=0.25, center]{2025_10_09_265d7e1a22c89f79c21eg-534}

Proof. It is easy to check that the above diagram commutes. Thus we must show that $\kappa$ is linear on fibers. Let $Y_{y} \in(T \pi)^{-1}\left(X_{p}\right)$ where $\pi(y)= p$. We may write $Y_{y}=H_{y}+V_{y}$, where $H_{y}$ and $V_{y}$ are horizontal and vertical respectively. Since $X_{p}=T \pi\left(Y_{y}\right)=T \pi\left(H_{y}\right)$, we see that $H_{y}$ is the horizontal lift $\widetilde{X}_{y}$ of $X_{p}$ to $T_{y} E$. Also, $V_{y}=\jmath_{y} w$ for a unique $w \in E_{p}$. So the decomposition may be written as $Y_{y}=\widetilde{X}_{y}+\jmath_{y} w$. By definition we have

$$
\kappa\left(Y_{y}\right)=w .
$$

Using the homogeneity property of the horizontal distribution together with Lemma 12.28 above we have

$$
T \mu_{r}\left(Y_{y}\right)=T \mu_{r}\left(\widetilde{X}_{y}\right)+T \mu_{r}\left(\jmath_{y} w\right)=\widetilde{X}_{r y}+\jmath_{r y} r w .
$$

Thus $\kappa\left(T \mu_{r}\left(Y_{y}\right)\right)=r w$. Therefore, we have $\kappa\left(T \mu_{r}\left(Y_{y}\right)\right)=r \kappa\left(Y_{y}\right)$. If we denote the scaling operation for the vector bundle structure on $T E \rightarrow T M$ using $\odot$ as before, then this last statement is just $\kappa\left(r \odot Y_{y}\right)=r \kappa\left(Y_{y}\right)$. The result now follows from the Corollary 12.27. \(\square\)

Once we have a connection on our bundle then the addition in the vector bundle $T E \rightarrow T M$ can be described in a convenient form. Any element of $T E$ that lies over the same element $X_{p} \in T M$ can be written in the form $\widetilde{X}_{y}+\jmath_{y} w$ for some $w \in E_{p}$, where $\widetilde{X}_{y}$ is the horizontal lift of $X_{p}$ to the point $y$. Let $\widetilde{X}_{y_{1}}+\jmath_{y_{1}} w_{1}$ and $\widetilde{X}_{y_{2}}+\jmath_{y_{2}} w_{2}$ be two such expressions. Then the sum of these two elements using the addition in the vector bundle $T E \rightarrow T M$ is given by $\widetilde{X}_{y_{1}+y_{2}}+\jmath_{y_{1}+y_{2}}\left(w_{1}+w_{2}\right)$, where $\widetilde{X}_{y_{1}+y_{2}}$ is the horizontal lift of $X_{p}$ to the point $y_{1}+y_{2}$.\\
Exercise 12.30. Prove the last statement.\\
Exercise 12.31. Deduce from the previous theorem that ( $\pi_{T E}, \kappa$ ) : $T E \rightarrow E \oplus E$ is a vector bundle isomorphism along the tangent bundle projection $\pi_{T M}: T M \rightarrow M$.

Using this notion of connection with associated connector $\kappa$ we can get a covariant derivative. If $\mathcal{H}$ is a connection on a vector bundle $E \rightarrow M$ with connector $\kappa$, then for a section $\sigma \in \Gamma_{f}(E)$ along a smooth map $f: N \rightarrow M$ we make the definition

$$
\nabla_{v}^{f} \sigma:=\kappa\left(T_{p} \sigma \cdot v\right) \quad \text { for } v \in T_{p} N
$$

If $V$ is a vector field on $N$, then $\left(\nabla_{V}^{f} \sigma\right)(p):=\nabla_{V(p)}^{f} \sigma$.\\
Theorem 12.32. Let $\pi: E \rightarrow M$ be a vector bundle. Suppose that the vector bundle is endowed with a connection $\mathcal{H}$ and associated connector $\kappa$. Then for each smooth map $f: N \rightarrow M$, the map $\nabla^{f}$ defined above is a covariant derivative along $f$ (Definition 12.7). If $f=\operatorname{id}_{M}$, then we obtain a Koszul connection.

Conversely, if $\nabla$ is a Koszul connection on $\pi: E \rightarrow M$, then we may define a connection by

$$
\mathcal{H}_{y}:=\left\{T s \cdot u-\jmath_{y} \nabla_{u} s: s \in \Gamma(E), s(\pi(y))=y, u \in T_{\pi(y)} M\right\} .
$$

The resulting connection, in turn, gives back the Koszul connection according to $\nabla_{v} s:=\kappa\left(T_{p} s \cdot v\right)$ for $v \in T_{p} M$.

Proof. We write $\nabla$ for $\nabla^{f}$ when no confusion is likely. We start by simply noting that (i) and (iv) of Definition 12.7 are easily seen to be true and leave these as an exercise. Notice that if $g: P \rightarrow N$ and $f: N \rightarrow M$ are smooth and $u \in T P$, then for each $\sigma \in \Gamma_{f}(E)$, we have $\nabla_{u}(s \circ g)=\nabla_{T g \cdot u} s$. Indeed,

$$
\nabla_{u}(\sigma \circ g)=(\kappa \circ T(\sigma \circ g)) u=\kappa(T \sigma(T g \cdot u))=\nabla_{T g \cdot u} \sigma
$$

This gives (v) of Definition 12.7. Next we claim that if $\sigma_{1}, \sigma_{2} \in \Gamma_{f}(E)$ and $u \in T_{p} N$, then

$$
\nabla_{u}\left(\sigma_{1}+\sigma_{2}\right):=\nabla_{u} \sigma_{1}+\nabla_{u} \sigma_{2},
$$

which implies (ii) of Definition 12.7.

Proof of the claim: Recall the definition of addition in the vector bundle $T \pi: T E \rightarrow T M$ in terms of the tangent lift of $\alpha:(u, v) \mapsto u+v$. If $u=\dot{\gamma}(0)$ for a smooth curve $\gamma$ in $N$ with $\gamma(0)=p$, then

$$
\begin{aligned}
T \sigma_{1} \cdot u \boxplus T \sigma_{2} \cdot u & :=T \alpha\left(T \sigma_{1} \cdot u, T \sigma_{2} \cdot u\right)=\left.\frac{d}{d t}\right|_{0}\left(\sigma_{1} \circ \gamma+\sigma_{2} \circ \gamma\right) \\
& =\left.\frac{d}{d t}\right|_{0}\left(\sigma_{1}+\sigma_{2}\right) \circ \gamma=T\left(\sigma_{1}+\sigma_{2}\right) \cdot u
\end{aligned}
$$

Now we use the above together with the fact that $\kappa$ is a bundle homomorphism along $\pi_{T M}$. We have

$$
\nabla_{u}\left(\sigma_{1}+\sigma_{2}\right)=\kappa\left(T\left(\sigma_{1}+\sigma_{2}\right) u\right)=\kappa\left(T \sigma_{1} u \boxplus T \sigma_{2} u\right)=\nabla_{u} \sigma_{1}+\nabla_{u} \sigma_{2}
$$

Claim: If $h \in C^{\infty}(N, \mathbb{F})$ and $\sigma \in \Gamma_{f}(E)$, then $\nabla_{u} h \sigma=u(h) \sigma(p)+ h(p) \nabla_{u} \sigma$, which is (iii) of Definition 12.7.

Proof: Let $u \in T_{p} N$ and let $\sigma: N \rightarrow E$ be a section along a smooth map $f: N \rightarrow M$. We begin by calculating $T \mu: T \mathbb{R} \times T E \rightarrow T E$, where $\mu: \mathbb{R} \times E \rightarrow E$ is scalar multiplication in the vector bundle $E \rightarrow M$. So let $(a, y) \in \mathbb{R} \times E$ and $\left(\left.b \frac{d}{d t}\right|_{a}, v_{y}\right) \in T_{a} \mathbb{R} \times T_{y} E$. To find $T \mu \cdot\left(\left.b \frac{d}{d t}\right|_{a}, v_{y}\right)$ we calculate $T \mu \cdot\left(0_{a}, v_{y}\right)$ and $T \mu \cdot\left(\left.b \frac{d}{d t}\right|_{a}, 0_{y}\right)$ separately. Let $c$ be a smooth curve with $c(0)=y$ and $\dot{c}(0)=v_{y}$. Then

$$
\begin{aligned}
T \mu \cdot\left(0_{a}, v_{y}\right) & =\left.\frac{d}{d t}\right|_{0} \mu(a, c(t))=\left.\frac{d}{d t}\right|_{0} \mu_{a}(c(t)) \\
& =T \mu_{a} \cdot v_{y}=: a \odot v_{y}
\end{aligned}
$$

where $\odot$ refers to multiplication in the vector bundle structure of $T E \rightarrow T M$. Next, let $c$ be the curve in the manifold $\mathbb{R}$ given by $c(t):=a+t b$ so that $\dot{c}(0)=\left.b \frac{d}{d t}\right|_{a}$. Then

$$
\begin{aligned}
T \mu \cdot\left(\left.b \frac{d}{d t}\right|_{a}, 0_{y}\right) & =\left.\frac{d}{d t}\right|_{0} \mu(c(t), y)=\left.\frac{d}{d t}\right|_{0}((a+b t) y) \\
& =\left.\frac{d}{d t}\right|_{0}(a y+t b y)=\jmath_{a y}(b y)
\end{aligned}
$$

We now have

$$
T \mu \cdot\left(\left.b \frac{d}{d t}\right|_{a}, v_{y}\right)=a \odot v_{y}+\jmath_{a y}(b y)
$$

Next we let $c$ be a curve in $N$ with $c(0)=p$ and $\dot{c}(0)=u \in T_{p} N$. Then

$$
\begin{aligned}
T_{p} h \cdot u & =\left.\frac{d}{d t}\right|_{0} h(c(t))=\left.(h \circ c)^{\prime}(t) \frac{d}{d t}\right|_{h(c(0))} \\
& =\left.u[h] \frac{d}{d t}\right|_{h(p)} .
\end{aligned}
$$

Now let $h \times \sigma: N \rightarrow \mathbb{R} \times E$ denote the map defined by $(h \times \sigma)(x)= (h(x), \sigma(x))$ and let $\widetilde{\sigma}:=h \sigma$. Then using the linearity of $\kappa$ and what we developed above we have

$$
\begin{aligned}
\nabla_{u} h \sigma & =\kappa(T \widetilde{\sigma} \cdot u)=\kappa[T \mu \circ T(h \times \sigma)(u)] \\
& =\kappa T \mu\left(\left.u[h] \frac{d}{d t}\right|_{h(p)}, T_{p} \sigma \cdot u\right) \\
& =\kappa\left\{h(p) \odot(T \sigma \cdot u)+\jmath_{h}(p) \sigma(p)(u[h] \sigma(p))\right\} \\
& =h(p) \kappa(T \sigma \cdot u)+u[h] \sigma(p)=h(p) \nabla_{u} \sigma+u[h] \sigma(p) .
\end{aligned}
$$

The proof that

$$
\mathcal{H}_{y}:=\left\{\jmath_{y} \nabla_{u} s-T s \cdot u \mid s \in \Gamma(E), u \in T_{\pi(y)} M\right\}
$$

defines a connection that gives back the original covariant derivative is left to the reader as Problem 9.

Example 12.33. In Chapter 4 we saw that each hypersurface in $\mathbb{R}^{n+1}$ has a natural covariant derivative (the Levi-Civita covariant derivative). This means that there is a corresponding horizontal distribution. Let us consider the canonical connection on $S^{n} \subset \mathbb{R}^{n+1}$. We may identify $T S^{n}$ with the subset of $S^{n} \times \mathbb{R}^{n+1}$ given by $\left\{(p, u) \in S^{n} \times \mathbb{R}^{n+1}:\langle p, u\rangle=0\right\}$. Under this identification, the velocity of a curve $c$ has the form $\dot{c}=\left(c, c^{\prime}\right)$. Now let us find a representation of $T\left(T S^{n}\right)$ as a subset of $T S^{n} \times \mathbb{R}^{2 n+2}$. A section of $T S^{n}$ along a curve $c$ has the form $\sigma: t \mapsto(c(t), x(t))$ and so we take $\dot{\sigma}=\left(c, x, c^{\prime}, x^{\prime}\right) \in T S^{n} \times \mathbb{R}^{2 n+2}$. Note that since $\langle c, x\rangle=0$, we have $\left\langle c, x^{\prime}\right\rangle=-\left\langle c^{\prime}, x\right\rangle$. Also, since $\langle c, c\rangle=1$, we have $\left\langle c, c^{\prime}\right\rangle=0$. It follows that we may make the identification\\
$T\left(T S^{n}\right)=\left\{((p, u),(v, w)) \in T S^{n} \times \mathbb{R}^{2 n+2}:\langle p, v\rangle=0,\langle p, w\rangle+\langle u, v\rangle=0\right\}$.\\
The vertical space $\mathcal{V}_{(p, u)}$ in $T_{(p, u)}\left(T S^{n}\right)$ is

$$
\mathcal{V}_{(p, u)}=\left\{((p, u),(0, w)) \in T_{(p, u)}\left(T S^{n}\right)\right\}
$$

Recall that in Chapter 4 we obtained the Levi-Civita covariant derivative by projection of the ambient derivative back into the tangent space. In the present context, this means that $((p, u),(v, w))$ is horizontal if $w$ is a multiple of $p$. But since we also have $\langle p, w\rangle+\langle u, v\rangle=0$, the criterion for being horizontal becomes $w=-\langle u, v\rangle p$. A little thought reveals that we should take $\mathcal{H}$ to be defined by

$$
\mathcal{H}_{(p, u)}:=\left\{((p, u),(v,-\langle u, v\rangle p)) \in T_{(p, u)} T S^{n}\right\} .
$$

Exercise 12.34. Let $p, q \in S^{2} \subset \mathbb{R}^{3}$ be such that $\langle p, q\rangle=0$. The great circle containing $p$ and $q$ can be parametrized as

$$
c(t):=(\cos t) p+(\sin t) q
$$

\begin{figure}[h]
\begin{center}
  \includegraphics[scale=0.25]{2025_10_09_265d7e1a22c89f79c21eg-538}

\caption{Figure 12.3. Parallel transport around path shows holonomy}
\end{center}
\end{figure}

for $0 \leq t \leq 2 \pi$. Using the notation of the previous example, show that if $\sigma=(c, x)$ is a parallel section along $c$, then both $\langle x, x\rangle$ and $\left\langle c^{\prime}, x\right\rangle$ are constant along $c$.

Figure 12.3 shows parallel translation of a vector at the north pole around a loop comprised of segments of great circles. Notice that parallel translation around this loop results in a map which is not the identity map. In general, if $\pi: E \rightarrow M$ is a vector bundle with connection and $c:[a, b] \rightarrow M$ is a piecewise smooth closed curve with $p=c(a)=c(b)$, then we obtain an isomorphism $P_{c} \in \mathrm{GL}\left(E_{p}\right)$. This map is called the holonomy of the curve.

Definition 12.35. The subset $G_{p} \subset \mathrm{GL}\left(E_{p}\right)$ consisting of the maps $P_{c}$ as $c$ ranges over all piecewise smooth closed curves $c$ with $p=c(a)=c(b)$ is called the holonomy group at $p$ for the connection.

It is easy to see that any piecewise smooth curve $c:[a, b] \rightarrow M$ induces a group isomorphism between $G_{c(a)}$ and $G_{c(b)}$ given by

$$
g \mapsto P_{c} \circ g \circ P_{c}^{-1} .
$$

Thus for a connected manifold we may speak of the holonomy group of the manifold, and this is well-defined up to group isomorphism.

One may use parallel translation to recover the covariant derivative and this is a very natural viewpoint:

Theorem 12.36. Let $\pi: E \rightarrow M$ be a vector bundle with connection $\mathcal{H}$. Let $f: N \rightarrow M$ be a smooth map. If $u \in T_{p} N$ and $\sigma$ is a section of $E$ along\\
$f$, then for any smooth curve $c:(-\varepsilon, \varepsilon) \rightarrow N$ with $\dot{c}(0)=u$ we have

$$
\nabla_{u}^{f} \sigma=\lim _{t \rightarrow 0} \frac{\left(P(f \circ c)_{0}^{t}\right)^{-1} \sigma(c(t))-\sigma(c(0))}{t}
$$

where $P(f \circ c)_{0}^{t}$ is the parallel transport along $f \circ c$ from $(f \circ c)(0)$ to $(f \circ c)(t)$.

Proof. Use parallel transport to obtain parallel frame fields $e_{1}, \ldots, e_{k}$ along $f \circ c$ with $e_{1}(0), \ldots, e_{k}(0)$ a basis of $E_{(f \circ c)(0)}$. Then we may write

$$
\sigma \circ c=\sum \sigma^{i} e_{i}
$$

for unique smooth functions $\sigma^{i}$ defined on $(-\varepsilon, \varepsilon)$. Then

$$
\begin{aligned}
\left(P(f \circ c)_{0}^{t}\right)^{-1} \sigma \circ c(t) & =\left(P(f \circ c)_{0}^{t}\right)^{-1} \sum \sigma^{i}(t) e_{i}(t) \\
& =\sum \sigma^{i}(t) e_{i}(0)
\end{aligned}
$$

Let $D_{0}$ denote $\left.\frac{\partial}{\partial t}\right|_{0}$. Then using the chain rule for connections and the fact that $\nabla_{\partial_{t}} e_{i}(t)=0$ we have

$$
\begin{aligned}
\nabla_{u} \sigma & =\nabla_{\dot{c}(0)} \sigma=\nabla_{T c \cdot D_{0}} \sigma \\
& =\nabla_{D_{0}}(\sigma \circ c)=\nabla_{D_{0}} \sum \sigma^{i}(t) e_{i}(t) \\
& =\left.\sum \frac{d \sigma^{i}}{d t}\right|_{0} e_{i}(0)=\left.\frac{d}{d t}\right|_{0}\left[\left(P(f \circ c)_{0}^{t}\right)^{-1} \sigma \circ c(t)\right] \\
& =\lim _{t \rightarrow 0} \frac{\left(P(f \circ c)_{0}^{t}\right)^{-1} \sigma(c(t))-\sigma(c(0))}{t}
\end{aligned}
$$

The most important special cases are the case of $f=\operatorname{id}_{M}$, where we recover a basic Koszul connection, and also the case where $f$ itself is a curve $c: I \rightarrow M$. However, recall that we often omit the superscripts indicating maps on the various covariant derivative operators.

Corollary 12.37. Let $\pi: E \rightarrow M$ be a vector bundle with connection $\mathcal{H}$. If $v \in T_{p} M$ and $s$ is a section of $E$, then for any smooth curve $c:(-\varepsilon, \varepsilon) \rightarrow M$ with $\dot{c}(0)=v$ we have

$$
\nabla_{v} s=\lim _{t \rightarrow 0} \frac{\left(P(c)_{0}^{t}\right)^{-1} s(c(t))-s(c(0))}{t}
$$

where $P(c)_{0}^{t}$ is the parallel transport along $c$ from $c(0)$ to $c(t)$.\\
Now recall that if $c: I \rightarrow M$ is a smooth curve and $\sigma: I \rightarrow E$ is a section along $c$, we define $\nabla_{\frac{\partial}{\partial t}} \sigma$ to be a section along $c$ given by

$$
\left(\nabla_{\frac{\partial}{\partial t}} \sigma\right)(t):=\nabla_{\left.\frac{\partial}{\partial t}\right|_{t}} \sigma .
$$

We have\\
Corollary 12.38. Let $\pi: E \rightarrow M$ be a vector bundle with connection $\mathcal{H}$. Let $c: I \rightarrow M$ be a smooth curve and suppose that $\sigma: I \rightarrow E$ is a section along c. Then we have

$$
\left(\nabla_{\frac{\partial}{\partial t}} \sigma\right)(t)=\lim _{\epsilon \rightarrow 0} \frac{\left(P(c)_{t}^{t+\epsilon}\right)^{-1} \sigma(t+\epsilon)-\sigma(t)}{\epsilon} .
$$

Exercise 12.39. Derive the above two corollaries from Theorem 12.36.\\
Corollary 12.40. Let $\pi: E \rightarrow M$ be a vector bundle with connection $\mathcal{H}$. Let $f: N \rightarrow M$ be a smooth map. If $\sigma$ is a section of $E$ along $f$, then $\sigma$ is parallel if and only if $\nabla_{u}^{f} \sigma=0$ for all $u \in T N$. In particular, a section $s \in \Gamma(E)$ is parallel along a curve $c$ if and only if $\nabla_{\partial / \partial t} s \circ c \equiv 0$.

If one approaches parallelism solely via covariant derivatives, then one could take $\nabla_{u}^{f} \sigma \equiv 0$ as the definition of parallel. This is a common equivalent approach.

Warning: There is a subtle point to be made here. If $s$ is a section of $E$ and $c:(-\varepsilon, \varepsilon) \rightarrow M$ is a curve, then $\nabla_{\dot{c}(t)} s$ is zero whenever $\dot{c}(t)=0$, but for a section $\sigma: I \rightarrow E$ along $c$ it is possible that $\left(\nabla_{\partial / \partial t} \sigma\right)(t)$ is nonzero even when $\dot{c}(t)=0$. For example, if $c_{p}: I \rightarrow M$ is a constant curve taking the fixed value $p \in M$, and $\sigma: I \rightarrow T_{p} M$ is a smooth map, then $\sigma$ can be considered a section along $c_{p}$ and then $\nabla_{\partial / \partial t} \sigma=\sigma^{\prime}$ (the ordinary derivative of a curve in the vector space $T_{p} M$ ).

Exercise 12.41. Show that $\sigma \in \Gamma(M, E)$ is a parallel section if and only if $\sigma \circ c$ is parallel along $c$ for every curve $c: I \rightarrow M$.

Exercise 12.42. Show that if $t \mapsto \sigma(t)$ is a curve in $E_{p}$, then we can consider $\sigma$ as a section along the constant map $c_{p}: t \mapsto p$ and then $\nabla_{\partial_{t}} \sigma(t)=\sigma^{\prime}(t) \in E_{p}$.

Exercise 12.43. Let $\nabla$ be a connection on $E \rightarrow M$ and let $\alpha:[a, b] \rightarrow M$ and $\beta:[a, b] \rightarrow E_{\alpha\left(t_{0}\right)}$. If $X(t):=P(\alpha)_{t_{0}}^{t}(\beta(t))$, then

$$
\left(\nabla_{\partial_{t}} X\right)(t)=P(\alpha)_{t_{0}}^{t}\left(\beta^{\prime}(t)\right) .
$$

Note: $\beta^{\prime}(t) \in E_{\alpha\left(t_{0}\right)}$ for all $t$. [Hint: Use a parallel frame field along the curve.]

\subsection*{Curvature}
An important fact about covariant derivatives is that they do not need to commute. If $\sigma: M \rightarrow E$ is a section and $X \in \mathfrak{X}(M)$, then $\nabla_{X} \sigma$ is a section also, and so we may take its covariant derivative $\nabla_{Y} \nabla_{X} \sigma$ with respect to some $Y \in \mathfrak{X}(M)$. In general, $\nabla_{Y} \nabla_{X} \sigma \neq \nabla_{X} \nabla_{Y} \sigma$. A measure of this\\
lack of commutativity is the curvature operator that is defined for a pair $X, Y \in \mathfrak{X}(M)$ to be the map $F(X, Y): \Gamma(E) \rightarrow \Gamma(E)$ given by

$$
F(X, Y) \sigma:=\nabla_{X} \nabla_{Y} \sigma-\nabla_{Y} \nabla_{X} \sigma-\nabla_{[X, Y]} \sigma
$$

i.e. $F(X, Y):=\left[\nabla_{X}, \nabla_{Y}\right]-\nabla_{[X, Y]}$.\\
Theorem 12.44. For fixed $\sigma$, the map $(X, Y) \mapsto F(X, Y) \sigma$ is $C^{\infty}(M)$ bilinear and antisymmetric. Also, $F(X, Y): \Gamma(E) \rightarrow \Gamma(E)$ is a $C^{\infty}(M)$ module homomorphism; that is, it is linear over the smooth functions,

$$
F(X, Y)(f \sigma)=f F(X, Y)(\sigma)
$$

Proof. We leave the proof of the first part as an easy exercise. For the second part, we calculate:

$$
\begin{aligned}
F(X, Y)(f \sigma)= & \nabla_{X} \nabla_{Y} f \sigma-\nabla_{Y} \nabla_{X} f \sigma-\nabla_{[X, Y]} f \sigma \\
= & \nabla_{X}\left(f \nabla_{Y} \sigma+(Y f) \sigma\right)-\nabla_{Y}\left(f \nabla_{X} \sigma+(X f) \sigma\right) \\
& -f \nabla_{[X, Y]} \sigma-([X, Y] f) \sigma \\
= & f \nabla_{X} \nabla_{Y} \sigma+(X f) \nabla_{Y} \sigma+(Y f) \nabla_{X} \sigma+X(Y f) \\
& -f \nabla_{Y} \nabla_{X} \sigma-(Y f) \nabla_{X} \sigma-(X f) \nabla_{Y} \sigma-Y(X f) \\
& -f \nabla_{[X, Y]} \sigma-([X, Y] f) \sigma \\
= & f\left[\nabla_{X}, \nabla_{Y}\right]-f \nabla_{[X, Y]} \sigma=f F(X, Y) \sigma .
\end{aligned}
$$

Thus we also have $F$ as a map $F: \mathfrak{X}(M) \times \mathfrak{X}(M) \rightarrow \Gamma(M, \operatorname{End}(E))$. But $F$ is $C^{\infty}(M)$ bilinear in the first two slots also.

Exercise 12.45. Show that $F$ is $C^{\infty}(M)$ bilinear in the first two slots. Let $v_{p}, w_{p} \in T_{p} M$ and $\sigma_{p} \in E_{p}$ be given. Let $X, Y \in \mathfrak{X}(M)$ and $\sigma \in \Gamma(E)$ such that $X(p)=v_{p}, Y(p)=w_{p}$ and $\sigma(p)=\sigma_{p}$. Argue that $F\left(X_{p}, Y_{p}\right) \sigma_{p}$ is well-defined and that the map $\left(v_{p}, w_{p}, \sigma_{p}\right) \mapsto F\left(v_{p}, w_{p}\right) \sigma_{p}$ is multilinear.

In light of this exercise we see that for each $v_{p}, w_{p} \in T_{p} M, F\left(v_{p}, w_{p}\right)$ is a linear map $E_{p} \rightarrow E_{p}$. In other words, $F\left(X_{p}, Y_{p}\right) \in \operatorname{End}\left(E_{p}\right)$ (a.k.a. $\left.L\left(E_{p}, E_{p}\right)\right)$. But $F\left(X_{p}, Y_{p}\right)$ is antisymmetric, and so we obtain, for each $p$, an element of $\operatorname{End}\left(E_{p}\right) \otimes \wedge^{2} T_{p}^{*} M$. Finally, we see that $F$ may be considered as a section of $\operatorname{End}(E) \otimes \wedge^{2} T^{*} M$. The space of such sections is denoted $\Omega^{2}(M, \operatorname{End}(E))$.

$$
\begin{array}{|l|}
\hline \text { Curvature is an } \operatorname{End}(E) \text {-valued 2-form. } \\
\hline
\end{array}
$$

We will have many occasions to consider differentiation along maps, and so we should also look at the curvature operator for sections along maps.

\begin{figure}[h]
\begin{center}
  \includegraphics[scale=0.25]{2025_10_09_265d7e1a22c89f79c21eg-542}

\caption{Figure 12.4. Parallel translation around loop}
\end{center}
\end{figure}

Let $f: N \rightarrow M$ be a smooth map and $\sigma$ a section of $E \rightarrow M$ along $f$. For $U, V \in \mathfrak{X}(N)$, we define a map $F^{f}(U, V): \Gamma_{f}(E) \rightarrow \Gamma_{f}(E)$ by

$$
F^{f}(U, V) \sigma:=\nabla_{U} \nabla_{V} \sigma-\nabla_{V} \nabla_{U} \sigma-\nabla_{[U, V]} \sigma
$$

for all $\sigma \in \Gamma_{f}(E)$.\\
Note that if $s \in \Gamma(E)$, then $s \circ f \in \Gamma_{f}(E)$, and if $U \in \mathfrak{X}(N)$, then $T f \circ U \in \Gamma_{f}(T M)=: \mathfrak{X}_{f}(M)$ is a vector field along $f$. As a matter of notation, we let $F(T f \circ U, T f \circ V) s$ denote the map $p \mapsto F\left(T f \cdot U_{p}, T f \cdot V_{p}\right) s(p)$, which makes sense because of Theorem 12.44. Thus $F(T f \circ U, T f \circ V) s \in \Gamma_{f}(E)$. Then we have the following useful fact:

Proposition 12.46. Let $X \in \Gamma_{f}(E)$ and $U, V \in \mathfrak{X}(N)$. Then

$$
F^{f}(U, V) X=F(T f \circ U, T f \circ V) X .
$$

Proof. Exercise! \(\square\)

The next theorem is an important step in understanding the geometric meaning of curvature in terms of parallel transport. Refer to Figure 12.4.

Theorem 12.47. Let $\pi: E \rightarrow M$ be a vector bundle with connection. Let $u, v \in T_{p} M$ and $y \in E_{p}$. Let $U$ be a neighborhood of the origin in $\mathbb{R}^{2}$ and denote standard coordinate frame fields on $\mathbb{R}^{2}$ by $\partial_{1}, \partial_{2}$ and their values at the origin by $\partial_{1}(0)$ and $\partial_{2}(0)$. Suppose that we have a smooth map $f: U \rightarrow M$ with $f(0)=p, T_{p} f \cdot \partial_{1}(0)=u$ and $T_{p} f \cdot \partial_{2}(0)=v$. For small $s, t$, let $y_{s, t}$ denote the element of $E_{p}$ obtained by parallel translation of $y$ around the piecewise smooth loop obtained by successively tracing the following four curves:\\
(1) $c_{1}^{s, t}: \sigma \mapsto f(\sigma, 0), 0 \leq \sigma \leq s ;$\\
(2) $c_{2}^{s, t}: \tau \mapsto f(s, \tau), 0 \leq \tau \leq t ;$\\
(3) $c_{3}^{s, t}: \sigma \mapsto f(s-\sigma, t), 0 \leq \sigma \leq s$;\\
(4) $c_{4}^{s, t}: \tau \mapsto f(0, t-\tau), 0 \leq \tau \leq t$.

Then,

$$
F(u, v) y=-\lim _{s, t \rightarrow 0} \frac{y_{s, t}-y}{s t}
$$

Proof. Let $Y$ be a section along $f$ defined as follows: If we choose $\epsilon>0$ sufficiently small, then for each $(s, t) \in[0, \epsilon] \times[0, \epsilon]$ the four curves described above will be defined. For each such $(s, t)$, let $Y(s, t)$ be the result of parallel translation of $y$ along the first two curves $c_{1}^{s, t}$ and $c_{2}^{s, t}$. We leave it to the reader to argue that $Y$ is smooth. Notice that $y=Y(0,0)$. By Proposition 12.46 we have

$$
F(u, v) y=F^{f}\left(\partial_{1}(0), \partial_{2}(0)\right) y=\nabla_{\partial_{1}(0)} \nabla_{\partial_{2}} Y-\nabla_{\partial_{2}(0)} \nabla_{\partial_{1}} Y
$$

since $\left[\partial_{1}, \partial_{2}\right]=0$. Since $Y$ is parallel along the curves $c_{2}^{s, t}$, we have $\nabla_{\partial_{2}} Y=0$ and so

$$
F(u, v) y=-\nabla_{\partial_{2}(0)} \nabla_{\partial_{1}} Y
$$

If we let $P_{t}$ denote parallel translation along $\tau \mapsto f(0, \tau)$ from $p$ to $f(0, t)$, then\\
$\nabla_{\partial_{2}(0)} \nabla_{\partial_{1}} Y=\lim _{t \rightarrow 0} \frac{P_{t}^{-1}\left(\nabla_{\partial_{1}} Y\right)(0, t)-\left(\nabla_{\partial_{1}} Y\right)(0,0)}{t}=\lim _{t \rightarrow 0} \frac{P_{t}^{-1}\left(\nabla_{\partial_{1}(0, t)} Y\right)}{t}$.\\
On the other hand, if $P_{s, t}$ denotes parallel translation along $\sigma \mapsto f(\sigma, t)$ from $f(0, t)$ to $f(s, t)$, then

$$
\nabla_{\partial_{1}(0, t)} Y=\lim _{s \rightarrow 0} \frac{P_{s, t}^{-1} Y(s, t)-Y(0, t)}{s}
$$

Putting things together we have

$$
F(u, v) y=-\lim _{s, t \rightarrow 0} \frac{P_{t}^{-1} P_{s, t}^{-1} Y(s, t)-P_{t}^{-1} Y(0, t)}{s t}=-\lim _{s, t \rightarrow 0} \frac{y_{s, t}-y}{s t} .
$$

Next we will try to understand curvature in terms of the horizontal distribution. Let us define an object that in some sense measures how far the horizontal distribution is from being integrable. First note that the horizontal lifts of local frame fields on the base manifold $M$ give vector fields on the total space that span the distribution. If the bracket of these lifted fields were always horizontal, then the distribution would be integrablethe connection would be flat. This suggests the following: For every pair of horizontal fields $Z_{1}, Z_{2}$ on $E$, let

$$
C\left(Z_{1}, Z_{2}\right):=\mathrm{p}_{\mathrm{v}}\left(\left[Z_{1}, Z_{2}\right]\right)
$$

Lemma 12.48. Let $\Gamma(\mathcal{H})$ denote the $C^{\infty}(E)$-module of horizontal fields and let $\Gamma(\mathcal{V})$ denote the $C^{\infty}(E)$-module of vertical fields. The map $C: \Gamma(\mathcal{H}) \times \Gamma(\mathcal{H}) \rightarrow \Gamma(\mathcal{V})$ given by $\left(Z_{1}, Z_{2}\right) \mapsto C\left(Z_{1}, Z_{2}\right)$ is a module homomorphism.

Proof. Let $f \in C^{\infty}(E)$. Then we have

$$
\begin{aligned}
C\left(f Z_{1}, Z_{2}\right) & =\mathrm{p}_{\mathrm{v}}\left(\left[f Z_{1}, Z_{2}\right]\right)=\mathrm{p}_{\mathrm{v}}\left(f\left[Z_{1}, Z_{2}\right]-\left(Z_{2} f\right) Z_{2}\right) \\
& =f \mathrm{p}_{\mathrm{v}}\left(\left[Z_{1}, Z_{2}\right]\right)-\mathrm{p}_{\mathrm{v}}\left(\left(Z_{2} f\right) Z_{2}\right)=f \mathrm{p}_{\mathrm{v}}\left(\left[Z_{1}, Z_{2}\right]\right) \\
& =f C\left(Z_{1}, Z_{2}\right)
\end{aligned}
$$

The analogous result for the second factor follows because $C$ is clearly skewsymmetric. Additivity is obvious.

Corollary 12.49. For $y \in E$, the value $C\left(Z_{1}, Z_{2}\right)(y)$ depends only on the values $Z_{1}(y)$ and $Z_{2}(y)$. Also, $C$ is well-defined on locally defined horizontal fields.

Proof. The corollary says that $C$ is "tensorial". The proof is just like the proof of Theorem 7.32 and can also be derived from Proposition 6.55.

Because of this corollary, we may think of $C$ as a map $\mathcal{H} \times \mathcal{H} \rightarrow \mathcal{V}$.\\
Theorem 12.50. Let $\pi: E \rightarrow M$ be a vector bundle with connection. Then for $u, v \in T_{p} M$ and $y \in E_{p}$ we have

$$
F(u, v) y=-\jmath_{y}^{-1}\left(C\left(u^{y}, v^{y}\right)\right)
$$

where $u^{y}, v^{y}$ denote the horizontal lifts of $u$ and $v$ to the point $y$. Here $\jmath_{y}: T_{y} E_{p} \rightarrow E_{p}$ is just the canonical map as usual.

Proof. Choose $U, V \in \mathfrak{X}(M)$ with $U(p)=u$ and $V(p)=v$. We may assume that $[U, V]=0$ near $p$. Let $\varphi_{s}$ and $\psi_{t}$ be local flows of $U$ and $V$ respectively. Also, let $\widetilde{\varphi}_{s}$ and $\widetilde{\psi}_{t}$ be the flows of the horizontal lifts $\widetilde{U}$ and $\widetilde{V}$. Observe that since $\varphi_{s} \circ \pi=\pi \circ \widetilde{\varphi}_{s}$, we have that $\widetilde{\varphi}_{s}(y)$ is the parallel translation of $y$ along the curve $\sigma \mapsto \varphi_{\sigma} \circ \pi(y), 0 \leq \sigma \leq s$. Similarly, $\widetilde{\psi}_{t}(y)$ is the parallel translation of $y$ along the curve $\tau \mapsto \psi_{\tau} \circ \pi(y), 0 \leq \tau \leq t$.

Now consider the curve $c$ defined for small $t$ by

$$
c(t):=\widetilde{\psi}_{-\sqrt{t}} \circ \widetilde{\varphi}_{-\sqrt{t}} \circ \widetilde{\psi}_{\sqrt{t}} \circ \widetilde{\varphi}_{\sqrt{t}} .
$$

Then, by Theorem 12.47, we have

$$
F(u, v) y=-\lim _{s, t \rightarrow 0} \frac{y_{s, t}-y}{s t}=-\lim _{t \rightarrow 0^{+}} \frac{y_{\sqrt{t}, \sqrt{t}}-y}{t}=-\lim _{t \rightarrow 0^{+}} \frac{c(t)-c(0)}{t} .
$$

But by Theorem 2.112,

$$
-\lim _{t \rightarrow 0^{+}} \frac{c(t)-c(0)}{t}=-\jmath_{y}^{-1} \dot{c}(0)=-\jmath_{y}^{-1}[\widetilde{U}, \widetilde{V}](y) .
$$

Finally, notice that $[U, V](p)=0$ so $[\widetilde{U}, \widetilde{V}](y)$ is vertical and is equal to $\mathrm{p}_{\mathrm{v}}([\widetilde{U}, \widetilde{V}])(y)=C(\widetilde{U}, \widetilde{V})(y)=C\left(u^{y}, v^{y}\right)$.

Unwinding the definitions, the result of the previous theorem can be written as

$$
(F(U, V) s)(p)=-\kappa\left([\widetilde{U}, \widetilde{V}]_{s(p)}\right)
$$

where $s \in \Gamma(E)$ and $\kappa$ is the connector associated to the connection.\\
Theorem 12.51. Let $\pi: E \rightarrow M$ be a vector bundle with connection. If the curvature $F$ vanishes, then the connection is flat. In particular, given any $y \in E$, there is a locally defined parallel section $s$ with $s(p)=y$ where $\pi(y)=p$.

Proof. Let $Z_{1}, Z_{2}$ be locally defined horizontal vector fields. Then for any $y$ in the common domain of $Z_{1}, Z_{2}$, let $u, v$ be the images of $Z_{1}(y), Z_{2}(y)$ under $T \pi$. Thus $Z_{1}(y)$ and $Z_{2}(y)$ are the horizontal lifts of $u$ and $v$. Then we have

$$
-\jmath_{y}^{-1}\left(\mathrm{p}_{\mathrm{v}}\left(\left[Z_{1}, Z_{2}\right]_{y}\right)\right)=F(u, v) y=0
$$

But $J_{y}^{-1}$ is an isomorphism, and so $\mathrm{p}_{\mathrm{v}}\left(\left[Z_{1}, Z_{2}\right]_{y}\right)=0$, which means that $\left[Z_{1}, Z_{2}\right]_{y}$ is horizontal. Since $y$ was arbitrary, $\left[Z_{1}, Z_{2}\right]=0$. We conclude that $\mathcal{H}$ is integrable by the Frobenius theorem. Now apply Theorem 12.25.

\subsection*{Connections on Tangent Bundles}
A connection on the tangent bundle $T M$ of a smooth manifold $M$ is called a linear connection on $M$. It is traditional to use special notation for the components of the connection form when using coordinate frames $\left\{\frac{\partial}{\partial x^{i}}\right\}$. In this case, formula (12.3) becomes

$$
\nabla_{\frac{\partial}{\partial x^{i}}} \frac{\partial}{\partial x^{j}}=\sum_{k} \Gamma_{i j}^{k} \frac{\partial}{\partial x^{k}}
$$

The $\Gamma_{i j}^{k}$ are a special case of the components of the connection forms. If $X=X^{j} \frac{\partial}{\partial x^{j}}$ and $Y=Y^{j} \frac{\partial}{\partial x^{j}}$, then

$$
\nabla_{X} Y=\left(\frac{\partial Y^{k}}{\partial x^{j}} X^{j}+\Gamma_{i j}^{k} X^{i} Y^{j}\right) \frac{\partial}{\partial x^{k}} \quad \text { (summation convention), }
$$

which is formula (12.4) in the current context. The functions $\Gamma_{i j}^{k}$ are called

\section*{Christoffel symbols.}
It is a consequence of Proposition 7.37 that for each $X \in \mathfrak{X}(M)$ there is a unique tensor derivation $\nabla_{X}$ on $\mathcal{T}_{s}^{r}(M)$ such that $\nabla_{X}$ commutes with contraction and coincides with the given covariant derivative on $\mathfrak{X}(M)$ (also denoted $\nabla_{X}$ ) and with $\mathcal{L}_{X}$ on $C^{\infty}(M)$.

To describe the covariant derivative on tensors more explicitly, consider $\omega \in \mathcal{T}_{1}^{0}(M)$. Since we have the contraction $Y \otimes \omega \mapsto C(Y \otimes \omega)=\omega(Y)$, we should have

$$
\begin{aligned}
\nabla_{X} \omega(Y) & =\nabla_{X} C(Y \otimes \omega) \\
& =C\left(\nabla_{X}(Y \otimes \omega)\right) \\
& =C\left(\nabla_{X} Y \otimes \omega+Y \otimes \nabla_{X} \omega\right) \\
& =\omega\left(\nabla_{X} Y\right)+\left(\nabla_{X} \omega\right)(Y)
\end{aligned}
$$

So we define $\left(\nabla_{X} \omega\right)(Y):=\nabla_{X}(\omega(Y))-\omega\left(\nabla_{X} Y\right)$. This implies that for a local frame field $E_{1}, \ldots, E_{n}$ with dual frame field $\theta^{1}, \ldots, \theta^{n}$ we have

$$
\nabla_{X} \theta^{i}=-\sum_{k} \omega_{j}^{i}(X) \theta^{j}
$$

where $\omega_{j}^{i}$ are the connection forms satisfying $\nabla_{X} E_{j}=\sum_{k} \omega_{j}^{i}(X) E_{i}$. More generally, if $\Upsilon \in \mathcal{T}_{s}^{r}$, then

$$
\begin{aligned}
\left(\nabla_{X} \Upsilon\right) & \left(\omega_{1}, \ldots, \omega_{r}, Y_{1}, \ldots, Y_{s}\right) \\
= & X\left(\Upsilon\left(\omega_{1}, \ldots, \omega_{r}, Y_{1}, \ldots, Y_{s}\right)\right)-\sum_{j=1}^{r} \Upsilon\left(\omega_{1}, . \nabla_{X} \omega_{j}, \ldots, \omega_{r}, Y_{1}, \ldots\right) \\
& -\sum_{i=1}^{s} \Upsilon\left(\omega_{1}, \ldots, \omega_{r}, Y_{1}, \ldots, \nabla_{X} Y_{i}, \ldots\right)
\end{aligned}
$$

Definition 12.52. The covariant differential of a tensor field $\Upsilon \in \mathcal{T}_{l}^{k}$ is denoted by $\nabla \Upsilon$ and is defined to be the element of $\mathcal{T}_{l+1}^{k}$ given by

$$
\nabla \Upsilon\left(\omega^{1}, \ldots, \omega^{1}, X, Y_{1}, \ldots, Y_{s}\right):=\nabla_{X} \Upsilon\left(\omega^{1}, \ldots, \omega^{1}, Y_{1}, \ldots, Y_{s}\right)
$$

For any fixed frame field $E_{1}, \ldots, E_{n}$, we denote the components of $\nabla \Upsilon$ by $\nabla_{i} \Upsilon_{j_{1} \ldots j_{s}}^{i_{1} \ldots i_{r}}$.\\
Remark 12.53. We have placed the new variable at the beginning as suggested by our notation $\nabla_{i} \Upsilon_{j_{1} \ldots j_{s}}^{i_{1}}$ for the components of $\nabla \Upsilon$ but in opposition to the equally common notation $\Upsilon_{j_{1} \ldots j_{s} ; i}^{i_{1}}$. This has the advantage of meshing well with exterior differentiation, making the statement of Theorem 12.56 as simple as possible.

The reader is asked in Problem 2 to show that $T(X, Y):=\nabla_{X} Y-\nabla_{Y} X- [X, Y]$ defines a tensor and so $T\left(X_{p}, Y_{p}\right)$ is well-defined for $X_{p}, Y_{p} \in T_{p} M$. This tensor is called the torsion tensor. If $T$ vanishes identically, then we say that the connection is torsion free. We have already noted that the

Levi-Civita connection defined in Chapter 4 for a hypersurface is torsion free.

We shall have more to say about torsion further on, but for now, notice that if the connection is torsion free, then

$$
0=\nabla_{\partial_{i}} \partial_{j}-\nabla_{\partial_{j}} \partial_{i}=\sum_{k}\left(\Gamma_{i j}^{k}-\Gamma_{j i}^{k}\right) \partial_{k}
$$

and so the Christoffel symbols are symmetric with respect to the lower two indices;

$$
\Gamma_{i j}^{k}=\Gamma_{j i}^{k}
$$

However, even for torsion free connections, this is generally not true for $\omega_{\mu j}^{k}$ defined by $\nabla_{\partial_{\mu}} e_{j}=\sum \omega_{\mu j}^{k} e_{k}$ unless the elements $e_{1}, \ldots, e_{n}$ of the frame field have pairwise vanishing Lie brackets. But this amounts to saying that they are locally coordinate frame fields for some chart.

\subsection*{Comparing the Differential Operators}
On a smooth manifold we have the Lie derivative $\mathcal{L}_{X}: \mathcal{T}_{s}^{r}(M) \rightarrow \mathcal{T}_{s}^{r}(M)$ and the exterior derivative $d: \Omega^{k}(M) \rightarrow \Omega^{k+1}(M)$, and in case we have a torsion free covariant derivative $\nabla$, that makes three differential operators which we would like to compare. To this end, we restrict attention to purely covariant tensor fields $\mathcal{T}_{s}^{0}(M)$.

The extended covariant derivative on tensor fields $\nabla_{X}: \mathcal{T}_{s}^{0}(M) \rightarrow \mathcal{T}_{s}^{0}(M)$ respects the subspace consisting of alternating tensors, and so for each $k$ we have a map

$$
\nabla_{X}: L_{\mathrm{alt}}^{k}(M) \rightarrow L_{\mathrm{alt}}^{k}(M)
$$

and these combine to give a degree preserving map

$$
\nabla_{X}: L_{\text {alt }}(M) \rightarrow L_{\text {alt }}(M)
$$

In other notation,

$$
\nabla_{X}: \Omega(M) \rightarrow \Omega(M)
$$

It is also easily seen that not only do we have $\nabla_{X}(\alpha \otimes \beta)=\nabla_{X} \alpha \otimes \beta+\alpha \otimes \nabla_{X} \beta$, but also

$$
\nabla_{X}(\alpha \wedge \beta)=\nabla_{X} \alpha \wedge \beta+\alpha \wedge \nabla_{X} \beta
$$

Recall (12.8) and the similar formula (7.10) for the Lie derivative. If $\nabla$ is torsion free so that $\mathcal{L}_{X} Y_{i}=\left[X, Y_{i}\right]=\nabla_{X} Y_{i}-\nabla_{Y_{i}} X$, then we obtain the following modification of formula (7.10) which incorporates $\nabla$.

Proposition 12.54. For a torsion free connection, we have the following equality for any $S \in \mathcal{T}_{s}^{0}(M)$ :

$$
\begin{aligned}
\left(\mathcal{L}_{X} S\right)\left(Y_{1}, \ldots, Y_{s}\right)= & \left(\nabla_{X} S\right)\left(Y_{1}, \ldots, Y_{s}\right) \\
& +\sum_{i=1}^{s} S\left(Y_{1}, \ldots, Y_{i-1}, \nabla_{Y_{i}} X, Y_{i+1}, \ldots, Y_{s}\right)
\end{aligned}
$$

Proof. See Problem 10.\\
Corollary 12.55. $\mathcal{L}_{X} S-\nabla_{X} S$ is a tensor.\\
For $\omega \in \Omega^{k}(M)$, we have that $\nabla \omega$ is a covariant tensor field but now not necessarily alternating. ${ }^{2}$ We search for a way to fix this. By antisymmetrizing, we get a map $\Omega^{k}(M) \rightarrow \Omega^{k+1}(M)$ which turns out to be none other than our old friend the exterior derivative as will be shown below.

Theorem 12.56. If $\nabla$ is a torsion free covariant derivative on $M$, then

$$
d=(k+1) \text { Alt } \circ \nabla
$$

or in other words, if $\omega \in \Omega^{k}(M)$, then

$$
d \omega\left(X_{0}, X_{1}, \ldots, X_{k}\right)=\sum_{i=0}^{k}(-1)^{i}\left(\nabla_{X_{i}} \omega\right)\left(X_{0}, \ldots, \widehat{X}_{i}, \ldots, X_{k}\right)
$$

Proof. First we show that

$$
(\text { Alt } \circ \nabla \omega)\left(X_{0}, X_{1}, \ldots, X_{k}\right)=\frac{1}{k+1} \sum_{i=0}^{k}(-1)^{i}\left(\nabla_{X_{i}} \omega\right)\left(X_{0}, \ldots, \widehat{X}_{i}, \ldots, X_{k}\right)
$$

Consider the subgroup $H$ consisting of the permutations of $\{0,1, \ldots, k\}$ that fix 0 . The cosets of $H$ are $H=H_{0}, H_{1}, \ldots, H_{k}$ where $H_{j}$ is the set of permutations that send 0 to $j$. If we decompose the group of permutations into these cosets, then

$$
\begin{aligned}
(\text { Alt } \circ \nabla \omega)\left(X_{0}, X_{1}, \ldots, X_{k}\right) & =\frac{1}{(k+1)!} \sum_{\sigma} \operatorname{sgn}(\sigma)\left(\nabla_{X_{\sigma_{0}}} \omega\right)\left(X_{\sigma_{1}}, \ldots, X_{\sigma_{k}}\right) \\
& =\frac{1}{(k+1)!} \sum_{i=0}^{k} \sum_{\sigma \in H_{i}} \operatorname{sgn}(\sigma)\left(\nabla_{X_{\sigma_{0}}} \omega\right)\left(X_{\sigma_{1}}, \ldots, X_{\sigma_{k}}\right) \\
& =\frac{k!}{(k+1)!} \sum_{i=0}^{k}(-1)^{i}\left(\nabla_{X_{i}} \omega\right)\left(X_{0}, \ldots, \widehat{X}_{i}, \ldots, X_{k}\right)
\end{aligned}
$$

\footnotetext{${ }^{2}$ However, $\nabla \omega$ can be viewed as being in $\Omega^{k}(M) \otimes_{C^{\infty}} \Omega^{1}(M)$.
}To finish we calculate as follows:

$$
\begin{aligned}
& d \omega\left(X_{0}, X_{1}, \ldots, X_{k}\right)=\sum_{i=0}^{k}(-1)^{i} X_{i}\left(\omega\left(X_{0}, \ldots, \widehat{X}_{i}, \ldots, X_{k}\right)\right) \\
& \quad+\sum_{1 \leq r<s \leq k}(-1)^{r+s} \omega\left(\left[X_{r}, X_{s}\right], X_{0}, \ldots, \widehat{X}_{r}, \ldots, \widehat{X}_{s}, \ldots, X_{k}\right) \\
& \quad=\sum_{i=0}^{k}(-1)^{i} X_{i}\left(\omega\left(X_{0}, \ldots, \widehat{X}_{i}, \ldots, X_{k}\right)\right) \\
& \quad+\sum_{1 \leq r<s \leq k}(-1)^{r+s} \omega\left(\nabla_{X_{r}} X_{s}-\nabla_{X_{s}} X_{r}, X_{0}, \ldots, \widehat{X}_{r}, \ldots, \widehat{X}_{s}, \ldots, X_{k}\right)
\end{aligned}
$$

which is equal to

$$
\begin{aligned}
\sum_{i=0}^{k}(-1)^{i} X_{i}\left(\omega\left(X_{0}, \ldots, \widehat{X_{i}}, \ldots, X_{k}\right)\right) \\
\quad+\sum_{1 \leq r<s \leq k}(-1)^{r+1} \omega\left(X_{0}, \ldots, \widehat{X_{r}}, \ldots, \nabla_{X_{r}} X_{s}, \ldots, X_{k}\right) \\
\quad-\sum_{1 \leq r<s \leq k}(-1)^{s} \omega\left(X_{0}, \ldots, \nabla_{X_{s}} X_{r}, \ldots, \widehat{X}_{s}, \ldots, X_{k}\right) \\
\left.\quad=\sum_{i=0}^{k}(-1)^{i} \nabla_{X_{i}} \omega\left(X_{0}, \ldots, \widehat{X}_{i}, \ldots, X_{k}\right) \quad \text { (by using } 12.8\right)
\end{aligned}
$$

\subsection*{Higher Covariant Derivatives}
Now let us suppose that we have a connection $\nabla^{E_{i}}$ on every vector bundle $E_{i} \rightarrow M$ in some family $\left\{E_{i}\right\}_{i \in I}$. We then also have the connections $\nabla^{E_{i}^{*}}$ induced on the duals $E_{i}^{*} \rightarrow M$. By demanding a product formula be satisfied as usual we can form a related family of connections on all bundles formed from tensor products of the bundles in the family $\left\{E_{i}, E_{i}^{*}\right\}_{i \in I}$. In this situation, it might be convenient to denote any and all of these connections by a single symbol as long as the context makes confusion unlikely. In particular, we have the following common situation: By the definition of a connection we have that $X \mapsto \nabla_{X}^{E} \sigma$ is $C^{\infty}(M)$ linear and so $\nabla^{E} \sigma$ is a section of the bundle $T^{*} M \otimes E$. We can use the Levi-Civita connection or any torsion free connection $\nabla$ on $M$ together with $\nabla^{E}$ to define a connection on $E \otimes T^{*} M$. To get a clear picture of this connection, we first notice that a section $\xi$ of the bundle $E \otimes T^{*} M$ can be written locally in terms of a local frame field $\left\{\theta^{i}\right\}$ on $T^{*} M$ and a local frame field $\left\{e_{i}\right\}$ on $E$. Namely, we may write $\xi=\sum \xi_{j}^{i} e_{i} \otimes \theta^{j}$. Then the connection $\nabla^{E \otimes T^{*} M}$ on $E \otimes T^{*} M$ is defined so that a product rule holds. Let $\gamma_{\nu}^{\mu}$ and $\omega_{j}^{i}$ be the connection forms for $\nabla$\\
and $\nabla^{E}$ respectively, so that locally, using the summation convention, we have

$$
\begin{aligned}
\nabla_{X}^{E \otimes T^{*} M} \xi & =\nabla_{X}^{E}\left(\xi_{\mu}^{i} e_{i}\right) \otimes \theta^{\mu}+\xi_{\mu}^{i} e_{i} \otimes \nabla_{X} \theta^{\mu} \\
& =\left(X \xi_{\mu}^{i} e_{i}+\xi_{\mu}^{i} \nabla_{X}^{E} e_{i}\right) \otimes \theta^{\mu}+\xi_{\mu}^{i} e_{i} \otimes \nabla_{X} \theta^{\mu} \\
& =\left(X \xi_{\nu}^{r} e_{r}+e_{r} \omega_{i}^{r}(X) \xi_{\nu}^{i}\right) \otimes \theta^{\nu}-\gamma_{\nu}^{\mu}(X) \xi_{\mu}^{r} e_{r} \otimes \theta^{\nu} \\
& =\left(X \xi_{\mu}^{r}+\omega_{i}^{r}(X) \xi_{\nu}^{i}-\gamma_{\nu}^{\mu}(X) \xi_{\mu}^{r}\right) e_{r} \otimes \theta^{\nu}
\end{aligned}
$$

Now let $\xi=\nabla^{E} \sigma$ for a given $\sigma \in \Gamma(E)$. The map $X \mapsto \nabla_{X}^{E \otimes T^{*} M}\left(\nabla^{E} \sigma\right)$ is $C^{\infty}(M)$ linear, and $\nabla^{E \otimes T^{*} M}\left(\nabla^{E} \sigma\right)$ is an element of $\Gamma\left(E \otimes T^{*} M \otimes T^{*} M\right)$, which can again be given the obvious connection. The process continues, and denoting all the connections by the same symbol we may consider the $k$-th covariant derivative $\nabla^{k} \sigma \in \Gamma\left(E \otimes T^{*} M^{\otimes k}\right)$ for each $\sigma \in \Gamma(E)$.

It is sometimes convenient to change viewpoints just slightly and define the covariant derivative operators $\nabla_{X_{1}, X_{2}, \ldots, X_{k}}: \Gamma(E) \rightarrow \Gamma(E)$. The definition is given inductively as

$$
\begin{aligned}
\nabla_{X_{1}, X_{2}} \sigma:= & \nabla_{X_{1}} \nabla_{X_{2}} \sigma-\nabla_{\nabla_{X_{1}} X_{2}} \sigma, \\
\nabla_{X_{1}, \ldots, X_{k}} \sigma:= & \nabla_{X_{1}}\left(\nabla_{X_{2}, X_{3}, \ldots, X_{k}} \sigma\right) \\
& -\nabla_{\nabla_{X_{1}} X_{2}, X_{3}, \ldots, X_{k}} \sigma-\cdots-\nabla_{X_{2}, X_{3}, \ldots, \nabla_{X_{1}} X_{k}} \sigma .
\end{aligned}
$$

Then we have the following convenient formula, which is true by definition:

$$
\nabla^{(k)} \sigma\left(X_{1}, \ldots, X_{k}\right)=\nabla_{X_{1}, \ldots, X_{k}} \sigma
$$

Warning: $\nabla_{\partial_{i}} \nabla_{\partial_{j}} \tau$ is a section of $E$, but it is not the same section as $\nabla_{\partial_{i} \partial_{j}} \tau$ since in general

$$
\nabla_{\partial_{i}} \nabla_{\partial_{j}} \tau \neq \nabla_{\partial_{i}} \nabla_{\partial_{j}} \sigma-\nabla_{\nabla_{\partial_{i}} \partial_{j}} \sigma
$$

Now we have that

$$
\begin{aligned}
\nabla_{X_{1}, X_{2}} \sigma-\nabla_{X_{2}, X_{1}} \sigma & =\nabla_{X_{1}} \nabla_{X_{2}} \sigma-\nabla_{\nabla_{X_{1}} X_{2}} \sigma-\left(\nabla_{X_{2}} \nabla_{X_{1}} \sigma-\nabla_{\nabla_{X_{2}} X_{1}} \sigma\right) \\
& =\nabla_{X_{1}} \nabla_{X_{2}} \sigma-\nabla_{X_{2}} \nabla_{X_{1}} \sigma-\nabla_{\left(\nabla_{X_{1}} X_{2}-\nabla_{X_{2}} X_{1}\right)} \sigma \\
& =F\left(X_{1}, X_{2}\right) \sigma-\nabla_{T\left(X_{1}, X_{2}\right)} \sigma .
\end{aligned}
$$

So, if $\nabla$ (the connection on the base $M$ ) is torsion free, then we recover the curvature

$$
\nabla_{X_{1}, X_{2}} \sigma-\nabla_{X_{2}, X_{1}} \sigma=F\left(X_{1}, X_{2}\right) \sigma
$$

One thing that is quite important to realize, is that $F$ depends only on the connection on $E$, while the operators $\nabla_{X_{1}, X_{2}}$ involve a torsion free connection on the tangent bundle TM.

\subsection{Exterior Covariant Derivative}
The exterior covariant derivative essentially antisymmetrizes the higher covariant derivatives just defined in such a way that the dependence on the auxiliary torsion free linear connection on the base cancels out. Of course this means that there must be a definition that does not involve this connection on the base at all. We give the definitions below, but first we point out a little more algebraic structure. We can give the space of $E$-valued forms $\Omega(M, E)$ the structure of a $(\Omega(M), \Omega(M))$-"bi-module". This means that we define a product $\wedge: \Omega(M) \times \Omega(M, E) \rightarrow \Omega(M, E)$ and another product $\wedge: \Omega(M, E) \times \Omega(M) \rightarrow \Omega(M, E)$ which are compatible in the sense that

$$
(\alpha \wedge \omega) \wedge \beta=\alpha \wedge(\omega \wedge \beta)
$$

for $\omega \in \Omega(M, E)$ and $\alpha \in \Omega(M)$. These products are defined by extending linearly the rules

$$
\begin{aligned}
\alpha \wedge(\sigma \otimes \omega) & :=\sigma \otimes \alpha \wedge \omega \text { for } \sigma \in \Gamma(E) \text { and } \alpha, \omega \in \Omega(M), \\
(\sigma \otimes \omega) & \wedge \alpha:=\sigma \otimes \omega \wedge \alpha \text { for } \sigma \in \Gamma(E) \text { and } \alpha, \omega \in \Omega(M), \\
\alpha \wedge \sigma & =\sigma \wedge \alpha=\sigma \otimes \alpha \text { for } \alpha \in \Omega(M) \text { and } \sigma \in \Gamma(E) .
\end{aligned}
$$

More precisely, we extend by using the universal properties of multilinear products. It follows that

$$
\alpha \wedge \omega=(-1)^{k l} \omega \wedge \alpha \text { for } \alpha \in \Omega^{k}(M) \text { and } \omega \in \Omega^{l}(M, E) .
$$

In the situation where $\sigma \in \Gamma(E):=\Omega^{0}(M, E)$ and $\omega \in \Omega(M)$, we have all three of the following conventional equalities:

$$
\sigma \omega=\sigma \wedge \omega=\sigma \otimes \omega \text { (special case). }
$$

The proof of the following theorem is analogous to the proof of the existence of the exterior derivative.

Theorem 12.57. Given a connection $\nabla$ on a vector bundle $\pi: E \rightarrow M$, there exists a unique operator $d^{\nabla}: \Omega(M, E) \rightarrow \Omega(M, E)$ such that\\
(i) $d^{\nabla}\left(\Omega^{k}(M, E)\right) \subset \Omega^{k+1}(M, E)$;\\
(ii) For $\alpha \in \Omega^{k}(M)$ and $\omega \in \Omega^{\ell}(M, E)$ we have

$$
\begin{aligned}
& d^{\nabla}(\alpha \wedge \omega)=d \alpha \wedge \omega+(-1)^{k} \alpha \wedge d^{\nabla} \omega \\
& d^{\nabla}(\omega \wedge \alpha)=d^{\nabla} \omega \wedge \alpha+(-1)^{\ell} \omega \wedge d \alpha
\end{aligned}
$$

(iii) $d^{\nabla} \sigma=\nabla \sigma$ for $\sigma \in \Gamma(E)$.

In particular, if $\alpha \in \Omega^{k}(M)$ and $\sigma \in \Omega^{0}(M, E)=\Gamma(E)$, we have

$$
d^{\nabla}(\sigma \otimes \alpha)=d^{\nabla} \sigma \wedge \alpha+\sigma \otimes d \alpha
$$

It can be shown that if we use $\Omega^{k}(M ; E) \cong L_{\text {alt }}^{k}(\mathfrak{X}(M), \Gamma(E))$, then we have the following formula:

$$
\begin{aligned}
d^{\nabla} \omega\left(X_{0}, \ldots, X_{k}\right) & =\sum_{0 \leq i \leq k}(-1)^{i} \nabla_{X_{i}}\left(\omega\left(X_{0}, \ldots, \widehat{X_{i}}, \ldots, X_{k}\right)\right) \\
& +\sum_{0 \leq i<j \leq k}(-1)^{i+j} \omega\left(\left[X_{i}, X_{j}\right], X_{0}, \ldots, \widehat{X_{i}}, \ldots, \widehat{X_{j}}, \ldots, X_{k}\right)
\end{aligned}
$$

Definition 12.58. The operator $d^{\nabla}$ whose existence is given by the previous theorem is called the exterior covariant derivative.

Exercise 12.59. We shall sometimes use the notation $d^{E}$ rather than $d^{\nabla}$. This is especially useful when two possibly unrelated connections on possibly different bundles are involved in the discussion. For example, we may deal with a connection $\nabla^{T M}$ on $T M$ and at the same time deal with a connection $\nabla^{E}$ on $E \rightarrow M$. Then it is conceivable that we may have to work with both $d^{E}$ and $d^{T M}$. Failure to notice these notational possibilities can result in serious confusion.

Note that we have the following special case for $\mu \in \Omega^{1}(M, E)$ :

$$
\begin{aligned}
d^{\nabla} \mu(X, Y) & =\nabla_{X}(\mu(Y))-\nabla_{Y}(\mu(X))-\mu([X, Y]) \\
& =d^{\nabla}(\mu(Y))(X)-d^{\nabla}(\mu(X))(Y)-\mu([X, Y]) .
\end{aligned}
$$

Example 12.60. Let us consider the case where $E=T M$ with connection $\nabla^{T M}$. Bundle maps $T M \rightarrow T M$ may be regarded as elements of $\Omega^{1}(M ; T M)$ (think about this). In particular, the identity map $\operatorname{id}_{T M}(v)=v$ can be viewed as a $T M$-valued 1-form. If we take the exterior covariant differential of $\theta:=\operatorname{id}_{T M}$, we obtain an element $d^{T M} \theta$ of $\Omega^{2}(M ; T M)$. For $X, Y \in \mathfrak{X}(M)$, we compute

$$
\begin{aligned}
d^{T M} \theta(X, Y) & =d^{T M}(\theta(Y))(X)-d^{T M}(\theta(X))(Y)-\theta([X, Y]) \\
& =d^{T M}(Y)(X)-d^{T M}(X)(Y)-[X, Y] \\
& =\nabla_{X} Y-\nabla_{Y} X-[X, Y]=T(X, Y)
\end{aligned}
$$

So we see that $d^{T M} \theta$ gives the torsion of the connection $\nabla^{T M}$.\\
If $\nabla^{T M}$ is a torsion free covariant derivative on $M$ and $\nabla^{E}$ is a connection on the vector bundle $E \rightarrow M$, then as before we get a covariant derivative $\nabla$ on all the bundles $E \otimes \bigwedge^{k} T^{*} M$ and as for the ordinary exterior derivative we have the formula

$$
d^{\nabla}=d^{E}=(k+1) \text { Alt } \circ \nabla
$$

or in other words, if $\omega \in \Omega^{k}(M, E)$, then

$$
d^{\nabla} \omega\left(X_{0}, X_{1}, \ldots, X_{k}\right)=\sum_{i=0}^{k}(-1)^{i}\left(\nabla_{X_{i}} \omega\right)\left(X_{0}, \ldots, \widehat{X}_{i}, \ldots, X_{k}\right)
$$

Let us look at the local expressions with respect to a moving frame. Let $\left\{e_{i}\right\}$ be a frame field for $E$ on $U \subset M$. Then locally, for $\omega \in \Omega^{k}(M, E)$, we may write

$$
\omega=\sum e_{i} \otimes \omega^{i},
$$

where $\omega^{i} \in \Omega^{k}(M)$. Using the summation convention, we have

$$
\begin{aligned}
d^{\nabla} \omega & =d^{\nabla}\left(e_{i} \otimes \omega^{i}\right)=e_{i} \otimes d \omega^{i}+d^{\nabla} e_{i} \wedge \omega^{i} \\
& =e_{j} \otimes d \omega^{j}+\omega_{i}^{j} e_{j} \wedge \omega^{i} \\
& =e_{j} \otimes d \omega^{j}+e_{j} \otimes \omega_{i}^{j} \wedge \omega^{i} \\
& =e_{j} \otimes\left(d \omega^{j}+\omega_{i}^{j} \wedge \omega^{i}\right)
\end{aligned}
$$

Thus the "( $k+1$ )-form coefficients" of $d^{\nabla} \omega$ with respect to the frame $\left\{e_{j}\right\}$ are given by $d \omega^{j}+\omega_{i}^{j} \wedge \omega^{i}$.

We can extend our wedge operation to a bilinear operation between $\Omega(M, \operatorname{End}(E))$ and $\Omega(M, E)$ in such a way that

$$
(A \otimes \alpha) \wedge(\sigma \otimes \beta)=A(\sigma) \otimes \alpha \wedge \beta .
$$

To understand what is going on a little better, let us consider how the action of $\Omega(M, \operatorname{End}(E))$ on $\Omega(M, E)$ comes about from a slightly different point of view. We can identify $\Omega(M, \operatorname{End}(E))$ with $\Omega\left(M, E \otimes E^{*}\right)$, and then the action of $\Omega\left(M, E \otimes E^{*}\right)$ on $\Omega(M, E)$ is given by tensoring and then contracting: For $\alpha, \beta \in \Omega(M), s, \sigma \in \Gamma(E)$, and $s^{*} \in \Gamma(E)$, we have

$$
\begin{aligned}
\left(s \otimes s^{*} \otimes \alpha\right) \wedge(\sigma \otimes \beta) & :=C\left(s \otimes s^{*} \otimes \sigma \otimes(\alpha \wedge \beta)\right) \\
& =s^{*}(\sigma) s \otimes(\alpha \wedge \beta)
\end{aligned}
$$

From $\nabla^{E}$ we get related connections on $E^{*}, E \otimes E^{*}$ and $E \otimes E^{*} \otimes E$. The connection on $E \otimes E^{*}$ is also a connection on $\operatorname{End}(E)$. Of course, we have

$$
\nabla_{X}^{\operatorname{End}(E) \otimes E}(L \otimes \sigma)=\left(\nabla_{X}^{\operatorname{End}(E)} L\right) \otimes \sigma+L \otimes \nabla_{X}^{E} \sigma
$$

and after contraction,

$$
\begin{aligned}
\nabla_{X}^{E}(L(\sigma)) & =\nabla_{X}^{E} C(L \otimes \sigma)=C\left(\nabla_{X}^{\operatorname{End}(E) \otimes E} L \otimes \sigma+L \otimes \nabla_{X}^{E} \sigma\right) \\
& =\left(\nabla_{X}^{\operatorname{End}(E)} L\right)(\sigma)+L\left(\nabla_{X}^{E} \sigma\right)=\left(\nabla_{X}^{\operatorname{End}(E)} L\right)(\sigma)+L\left(\nabla_{X}^{E} \sigma\right) .
\end{aligned}
$$

But $X \mapsto L\left(\nabla_{X}^{E} \sigma\right)$ is just $L \wedge \nabla^{E} \sigma$. So we have

$$
\nabla^{E}(L(\sigma))=\left(\nabla^{\operatorname{End}(E)} L\right)(\sigma)+L \wedge \nabla^{E} \sigma
$$

The connections on $\operatorname{End}(E)$ and $\operatorname{End}(E) \otimes E$ give the corresponding exterior covariant derivative operators

$$
d^{\operatorname{End}(E)}: \Omega^{k}(M, \operatorname{End}(E)) \rightarrow \Omega^{k+1}(M, \operatorname{End}(E))
$$

and

$$
d^{\operatorname{End}(E) \otimes E}: \Omega^{k}(M, \operatorname{End}(E) \otimes E) \rightarrow \Omega^{k+1}(M, \operatorname{End}(E) \otimes E) .
$$

It must be expected that many readers will find this formalism a bit daunting but it is not really as bad as it seems. Local calculations turn out to be quite natural as we shall see below. We encourage the reader to seek as much exposure as possible. To this end, we recommend $[\mathbf{B a M u}]$, [Poor], and [Dar]. Also, part of what might be intimidating is really just the bulkiness of the notations. The following notational convention makes things more palatable:

Notation 12.61. Let us now agree that, whenever convenient, we may write simply $\nabla$ for $\nabla^{E^{*}}, \nabla^{\operatorname{End}(E)}, \nabla^{\operatorname{End}(E) \otimes E}$, etc. and $d^{\nabla}$ for what we have called $d^{E^{*}}, d^{\operatorname{End}(E)}$ etc.

Proposition 12.62. For $\Phi \in \Omega^{k}(M, \operatorname{End}(E))$ and $\omega \in \Omega^{k}(M, E)$, we have

$$
d^{\nabla}(\Phi \wedge \omega)=d^{\nabla} \Phi \wedge \omega+(-1)^{k} \Phi \wedge d^{\nabla} \omega
$$

Proof. We have

$$
\begin{aligned}
d^{E}((A \otimes \alpha) \wedge & (\sigma \otimes \beta)):=d^{E}(A(\sigma) \otimes(\alpha \wedge \beta)) \\
= & \nabla^{E}(A(\sigma)) \wedge(\alpha \wedge \beta)+A(\sigma) \otimes d(\alpha \wedge \beta) \\
= & (-1)^{k}\left(\alpha \wedge \nabla^{E}(A(\sigma)) \wedge \beta\right)+A(\sigma) \otimes d(\alpha \wedge \beta) \\
= & (-1)^{k}\left(\alpha \wedge\left\{\nabla^{\operatorname{End}(E)} A \wedge \sigma+A \wedge \nabla^{E} \sigma\right\} \wedge \beta\right) \\
& +A(\sigma) \otimes d \alpha \wedge \beta+(-1)^{k} A(\sigma) \otimes \alpha \wedge d \beta \\
= & \left(\nabla^{\operatorname{End}(E)} A\right) \wedge \alpha \wedge \sigma \wedge \beta+(-1)^{k} A \wedge \alpha \wedge\left(\nabla^{E} \sigma\right) \wedge \beta \\
& +A(\sigma) \otimes d \alpha \wedge \beta+(-1)^{k} A(\sigma) \otimes \alpha \wedge d \beta \\
= & \left(\nabla^{\operatorname{End}(E)} A \wedge \alpha+A \otimes d \alpha\right) \wedge(\sigma \otimes \beta) \\
& +(-1)^{k}(A \otimes \alpha) \wedge\left(\nabla^{E} \sigma \wedge \beta+\sigma \otimes d \beta\right) \\
= & d^{\operatorname{End}(E)}(A \otimes \alpha) \wedge(\sigma \otimes \beta)+(-1)^{k}(A \otimes \alpha) \wedge d^{E}(\sigma \otimes \beta)
\end{aligned}
$$

By linearity we conclude that for $\Phi \in \Omega^{k}(M, \operatorname{End}(E))$ and $\omega \in \Omega(M, E)$ we have

$$
d^{E}(\Phi \wedge \omega)=d^{\operatorname{End}(E)} \Phi \wedge \omega+(-1)^{k} \Phi \wedge d^{E} \omega
$$

So in light of our notational conventions we are done.

Remark 12.63. In the literature, it seems that the different natures of $\left(d^{\nabla}\right)^{k}$ and $(\nabla)^{k}$ are not always appreciated. For example, the higher derivatives given by $\left(d^{\nabla}\right)^{k}$ are not appropriate for defining $k$-th order Sobolev spaces since $\left(d^{\nabla}\right)^{2}$ is zero for any flat connection.

\subsection*{Curvature Again}
The space $\Omega(M, \operatorname{End}(E))$ is an algebra over $C^{\infty}(M)$ where the multiplication is according to $\left(L_{1} \otimes \omega_{1}\right) \wedge\left(L_{2} \otimes \omega_{2}\right)=\left(L_{1} \circ L_{2}\right) \otimes \omega_{1} \wedge \omega_{2}$. Now $\Omega(M, E)$ is a module over this algebra because we can multiply (using the symbol $\wedge$ ) as follows:

$$
(L \otimes \alpha) \wedge(\sigma \otimes \beta)=L \sigma \otimes(\alpha \wedge \beta)
$$

As usual, this definition on simple elements is sufficient. If $X, Y \in \mathfrak{X}(M)$ and $\Psi \in \Omega^{2}(M, \operatorname{End}(E))$, then using $\Omega^{2}(M ; E) \cong L_{\text {alt }}^{2}(\mathfrak{X}(M), \Gamma(E))$ we have

$$
(\Psi \wedge \sigma)(X, Y)=\Psi(X, Y) \sigma
$$

for any $\sigma \in \Omega(M, E)$.\\
Proposition 12.64. The map $d^{\nabla} \circ d^{\nabla}: \Omega^{k}(M, E) \rightarrow \Omega^{k+2}(M, E)$ is given by the action of $F$, the curvature 2 -form of $\nabla^{E}$,

$$
d^{\nabla} \circ d^{\nabla} \mu=F \wedge \mu \text { for } \mu \in \Omega^{k}(M, E)
$$

Proof. Let us check the case where $k=0$ first. From formula (12.10) above we have for $\sigma \in \Omega^{0}(M, E)$,

$$
\begin{aligned}
\left(d^{\nabla} \circ d^{\nabla} \sigma\right)(X, Y) & =\nabla_{X}\left(d^{\nabla} \sigma(Y)\right)-\nabla_{Y}\left(d^{\nabla} \sigma(X)\right)-\sigma([X, Y]) \\
& =\nabla_{X} \nabla_{Y} \sigma-\nabla_{X} \nabla_{Y} \sigma-\sigma([X, Y]) \\
& =F(X, Y) \sigma=(F \wedge \sigma)(X, Y)
\end{aligned}
$$

More generally, we just check $d^{\nabla} \circ d^{\nabla}$ on elements of the form $\mu=\sigma \otimes \theta$ :

$$
\begin{aligned}
d^{\nabla} \circ d^{\nabla} \mu & =d^{\nabla} \circ d^{\nabla}(\sigma \otimes \theta) \\
& =d^{\nabla}\left(d^{\nabla} \sigma \wedge \theta+\sigma \otimes d \theta\right) \\
& =\left(d^{\nabla} d^{\nabla} \sigma\right) \wedge \theta-d^{\nabla} \sigma \wedge d \theta+d^{\nabla} \sigma \wedge d \theta+0 \\
& =(F \wedge \sigma) \wedge \theta=F \wedge(\sigma \wedge \theta)=F \wedge(\sigma \otimes \theta) \\
& =F \wedge \mu
\end{aligned}
$$

Let us take a look at how curvature appears in a local frame field. As before restrict to an open set $U$ on which $e_{U}=\left(e_{1}, \ldots, e_{r}\right)$ is a given local frame field and then write a typical element $s \in \Omega^{k}(M, E)$ as $s=e \eta$, where\\
$\eta=\left(\eta^{1}, \ldots, \eta^{r}\right)^{t}$ is a column vector of smooth $k$-forms. With $\omega_{U}=\left(\omega_{j}^{i}\right)$ the connection forms we have

$$
\begin{aligned}
d^{\nabla} d^{\nabla} s & =d^{\nabla} d^{\nabla}\left(e_{U} \eta_{U}\right)=d^{\nabla}\left(e_{U} d \eta_{U}+d^{\nabla} e_{U} \wedge \eta_{U}\right) \\
& =d^{\nabla}\left(e_{U}\left(d \eta_{U}+\omega_{U} \wedge \eta_{U}\right)\right) \\
& =d^{\nabla} e_{U} \wedge\left(d \eta_{U}+\omega_{U} \wedge \eta_{U}\right)+e_{U} \wedge d^{\nabla}\left(d \eta_{U}+\omega_{U} \wedge \eta_{U}\right) \\
& =e_{U} \omega_{U} \wedge\left(d \eta_{U}+\omega_{U} \wedge \eta_{U}\right)+e_{U} \wedge d^{\nabla}\left(d \eta_{U}+\omega_{U} \wedge \eta_{U}\right) \\
& =e_{U} \omega_{U} \wedge\left(d \eta_{U}+\omega_{U} \wedge \eta_{U}\right)+e_{U} \wedge d \omega_{U} \wedge \eta_{U}-e_{U} \wedge \omega_{U} \wedge d \eta_{U} \\
& =e_{U} d \omega_{U} \wedge \eta_{U}+e_{U} \omega_{U} \wedge \omega_{U} \wedge \eta_{U} \\
& =e_{U}\left(d \omega_{U}+\omega_{U} \wedge \omega_{U}\right) \wedge \eta_{U}
\end{aligned}
$$

The matrix $d \omega_{U}+\omega_{U} \wedge \omega_{U}$ represents a section of $\operatorname{End}(E)_{U} \otimes \Lambda^{2} T^{*} U$. In fact, we will now check that these local sections paste together to give a global section of $\operatorname{End}(E) \otimes \bigwedge^{2} T^{*} M$, or in other words, an element of $\Omega^{2}(M, \operatorname{End}(E))$, which is clearly the curvature form: $F:(X, Y) \mapsto F(X, Y) \in \Gamma(\operatorname{End}(E))$. Let $F_{U}=d \omega_{U}+\omega_{U} \wedge \omega_{U}$ and let $F_{V}=d \omega_{V}+\omega_{V} \wedge \omega_{V}$ be the corresponding form for a different moving frame $e_{U}:=e_{V} g$, where $g: U \cap V \rightarrow \mathrm{GL}\left(\mathbb{F}^{r}\right)$, and $r$ is the rank of $E$. What we need to verify is the transformation law

$$
F_{V}=g^{-1} F_{U} g
$$

which we met earlier in equation (6.2). Recall that $\omega_{V}=g^{-1} \omega_{U} g+g^{-1} d g$. Using $d\left(g^{-1}\right)=-g^{-1} d g g^{-1}$, we have

$$
\begin{aligned}
F_{V}= & d \omega_{V}+\omega_{V} \wedge \omega_{V} \\
= & d\left(g^{-1} \omega_{U} g+g^{-1} d g\right) \\
& +\left(g^{-1} \omega_{U} g+g^{-1} d g\right) \wedge\left(g^{-1} \omega_{U} g+g^{-1} d g\right) \\
= & d\left(g^{-1}\right) \wedge \omega_{U} g+g^{-1} d \omega_{U} g-g^{-1} \omega_{U} \wedge d g+d\left(g^{-1}\right) \wedge d g \\
& +g^{-1} \omega_{U} \wedge \omega_{U} g+g^{-1} d g g^{-1} \wedge \omega_{U} g \\
& +g^{-1} \omega_{U} \wedge d g+g^{-1} d g \wedge g^{-1} d g \\
= & g^{-1} d \omega_{U} g+g^{-1} \omega_{U} \wedge \omega_{U} g=g^{-1} F_{U} g
\end{aligned}
$$

where we have used that $g^{-1} d g \wedge g^{-1} d g=d\left(g^{-1} g\right)=0$.

\subsection*{The Bianchi Identity}
In this section we give several versions of the so-called Bianchi identity for a connection $\nabla$ on a vector bundle $E \rightarrow M$. Perhaps the simplest version to understand is the following: If $U, V, W \in \mathfrak{X}(M)$, then\\
$\left[\nabla_{U},\left[\nabla_{V}, \nabla_{W}\right]\right]+\left[\nabla_{V},\left[\nabla_{W}, \nabla_{U}\right]\right]+\left[\nabla_{W},\left[\nabla_{U}, \nabla_{V}\right]\right]=0$ (Bianchi identity),\\
where $\left[\nabla_{U}, \nabla_{V}\right]:=\nabla_{U} \circ \nabla_{V}-\nabla_{V} \circ \nabla_{U}$. This identity follows trivially once we observe that the set of linear operators on any vector space is a Lie algebra under the commutator bracket operation $[A, B]:=A \circ B-B \circ A$. So, in this form, the Bianchi identity is just an instance of the Jacobi identity. Let ( $U, \mathrm{x}$ ) be a chart on $M$, and let $F_{\mu \nu}:\left.\left.\Gamma(E)\right|_{U} \rightarrow \Gamma(E)\right|_{U}$ be the local curvature operator defined by

$$
F_{\mu \nu}=\left[\nabla_{\partial_{\mu}}, \nabla_{\partial_{v}}\right]=F\left(\partial_{\mu}, \partial_{\nu}\right)
$$

where $\partial_{\mu}=\partial / \partial x^{\mu}$, etc. Then we have the following version of the Bianchi identity:

$$
\left[\nabla_{\partial_{\mu}}, F_{\nu \lambda}\right]+\left[\nabla_{\partial_{\nu}}, F_{\lambda \mu}\right]+\left[\nabla_{\partial_{\lambda \nu}}, F_{\mu \nu}\right]=0 \quad \text { (Bianchi identity). }
$$

Another revealing form of the Bianchi identity depends on our discussions in the last section and in particular on Proposition 12.62. We have, for any $E$-valued form $\eta$,

$$
\left(d^{\nabla}\right)^{3} \eta=d^{\nabla}\left(\left(d^{\nabla}\right)^{2} \eta\right)=d^{\nabla}(F \wedge \eta)=d^{\nabla} F \wedge \eta+F \wedge d^{\nabla} \eta
$$

But equally,

$$
\left(d^{\nabla}\right)^{3} \eta=\left(d^{\nabla}\right)^{2}\left(d^{\nabla} \eta\right)=F \wedge d^{\nabla} \eta
$$

so it must be that $d^{\nabla} F \wedge \eta=0$ for any $\eta$ and so we obtain

$$
d^{\nabla} F=0 \quad \text { (Bianchi identity). }
$$

Exercise 12.65. Use a calculation in local coordinates and with respect to a local frame $e_{1}, \ldots, e_{k}$ to show that the above versions of the Bianchi identity are equivalent.

\subsection*{2. $G$-Connections}
If the vector bundle $\pi: E \rightarrow M$ has a $G$-bundle structure, then there should be certain connections that respect this structure. These are the $G$-connections. It is time to confess that we have come face to face with a weakness of our approach to connections. It is much easier and more natural to define the notion of $G$-connection if one first defines connections in terms of horizontal distributions on principal bundles. We treat this in the online supplement [Lee, Jeff]. However, we can still say what a $G$-connection should be from the current point of view without too much trouble.

Recall that if $\pi: E \rightarrow M$ is a rank $k$ vector bundle with typical fiber V that has a $G$-bundle structure where $G$ acts on V by the standard action as a subgroup of $\mathrm{GL}(\mathrm{V})$, then we have the bundle of $G$-frames

$$
F_{G}(E):=\bigcup_{p \in M} F_{G}\left(E_{p}\right)
$$

where $F_{G}\left(E_{p}\right):=\left\{u \in \mathrm{GL}\left(\mathrm{V}, E_{p}\right): u=\phi^{-1}(p, \cdot)\right.$ for some $\left.(U, \phi) \in \mathcal{A}_{G}\right\}$ and $\mathcal{A}_{G}$ is the maximal $G$-atlas defining the structure. The elements of\\
$F_{G}\left(E_{p}\right)$ are called $G$-frames. Having fixed a basis ( $\mathbf{e}_{1}, \ldots, \mathbf{e}_{k}$ ) for V , each element of $F_{G}\left(E_{p}\right)$ can be identified as a basis ( $u_{1}, \ldots, u_{k}$ ) for $E_{p}$ according to $u: x \mapsto \sum x^{i} u_{i}$. If $F\left(E_{p}\right)$ is the set of all frames at $p$, then $F_{G}\left(E_{p}\right) \subset F\left(E_{p}\right)$ and $F_{G}(E)$ is a subbundle of $F(E)$. Now let $c:[a, b] \rightarrow M$ be a smooth curve. If $\left(u_{1}, \ldots, u_{k}\right)$ is a frame, then $\left(P_{c} u_{1}, \ldots, P_{c} u_{k}\right)$ is also a frame. Thus we get a map $P_{c}: F\left(E_{c(a)}\right) \rightarrow F\left(E_{c(a)}\right)$. The following is then a workable definition of $G$-connection:

Definition 12.66. Let $\pi: E \rightarrow M$ be a vector bundle with a $G$-bundle structure as above. A connection on the bundle is called a $G$-connection if parallel transport $P_{c}$ takes $G$-frames to $G$-frames for all piecewise smooth curves $c$.

A $G$-frame field is a frame field which is a $G$-frame at each point of its domain.

Exercise 12.67. Show that if $\nabla$ is the covariant derivative associated to a $G$-connection on $E$, then for each $G$-frame field the associated connection forms take values in the Lie algebra of $G$ thought of as a matrix subgroup of $\operatorname{GL}(n, \mathbb{F})$.

Let $\pi: E \rightarrow M$ be a vector bundle with metric $h=\langle\cdot, \cdot\rangle$ and standard fiber $\mathbb{R}^{k}$. The metric $h$ gives a reduction of the structure group to $O(n)$. In this case, an $O(n)$-frame is an orthonormal frame. A connection on $E$ is an $O(n)$-connection if and only if the associated covariant derivative $\nabla$ satisfies

$$
v\left\langle s_{1}, s_{2}\right\rangle=\left\langle\nabla_{v} s_{1}, s_{2}\right\rangle+\left\langle s_{1}, \nabla_{v} s_{2}\right\rangle
$$

for all $s_{1}, s_{2} \in \Gamma(E)$ and all $v \in T M$.\\
Exercise 12.68. Prove this last assertion.\\
One may also consider metrics which are nondegenerate but not necessarily positive definite. In this case, the structure group reduces to one of the semiorthogonal groups $O(k, n-k)$.

Definition 12.69. A covariant derivative satisfying (12.11) above for some metric $h$ on a vector bundle $E \rightarrow M$ is called a metric covariant derivative (or metric connection).

A simple partition of unity argument shows that if $h$ is a given metric, then there exists a (nonunique) metric connection for $h$ (Problem 11). For a metric connection, parallel transport is an isometry (Problem 12). Furthermore if $h$ is the metric and $\nabla$ is metric with respect to $h$, then it is easy to check that $\nabla h=0$.

Proposition 12.70. Suppose that $h=\langle\cdot, \cdot\rangle$ is a metric on a vector bundle $E \rightarrow M$. If $\nabla$ is a metric covariant derivative, then the corresponding curvature satisfies

$$
\left\langle F(X, Y) \sigma_{1}, \sigma_{2}\right\rangle+\left\langle\sigma_{1}, F(X, Y) \sigma_{2}\right\rangle=0
$$

for all $X, Y \in \mathfrak{X}(M)$ and all $\sigma_{1}, \sigma_{2} \in \Gamma(E)$.

Proof. Let $X$ and $Y$ be fixed vector fields. Without loss of generality we may assume that $[X, Y]=0$ (recall Exercise 7.34). It is also enough to show that $\langle F(X, Y) \sigma, \sigma\rangle=0$ for all $\sigma$. We have

$$
\begin{aligned}
\langle F(X, Y) \sigma, \sigma\rangle= & \left\langle\nabla_{X} \nabla_{Y} \sigma, \sigma\right\rangle-\left\langle\sigma, \nabla_{X} \nabla_{Y} \sigma\right\rangle \\
= & X\left\langle\nabla_{Y} \sigma, \sigma\right\rangle-\left\langle\nabla_{Y} \sigma, \nabla_{X} \sigma\right\rangle \\
& -Y\left\langle\nabla_{X} \sigma, \sigma\right\rangle-\left\langle\nabla_{X} \sigma, \nabla_{Y} \sigma\right\rangle \\
= & \frac{1}{2}(X Y\langle\sigma, \sigma\rangle-Y X\langle\sigma, \sigma\rangle)=0
\end{aligned}
$$

since $[X, Y]=0$.

We shall study metric connections on tangent bundles in the next chapter.

\section*{Problems}
(1) Prove Theorem 12.23.\\
(2) Let $M$ have a linear connection $\nabla$ and let $T(X, Y):=\nabla_{X} Y-\nabla_{Y} X$. Show that $T$ is $C^{\infty}(M)$-bilinear (tensorial). The resulting tensor is called the torsion tensor for the connection.\\
(3) Show that a holonomy group is indeed a group. Show that the holonomy at any point of a sphere is isomorphic to the special orthogonal group $\mathrm{SO}(2)$.

\section*{Second order differential equations and sprays}
(4) A second order differential equation on a smooth manifold $M$ is a vector field on $T M$, that is, a section $X$ of the bundle $T T M$ (second tangent bundle) such that every integral curve $\alpha$ of $X$ is the velocity curve of its projection on $M$. In other words, $\alpha=\dot{\gamma}$ where $\gamma:=\pi_{T M} \circ \alpha$. A solution curve $\gamma: I \rightarrow M$ for a second order differential equation $X$ is, by definition, a curve with $\ddot{\alpha}(t)=X(\dot{\alpha}(t))$ for all $\tau \in I$.

In the case $M=\mathbb{R}^{n}$, show that this concept corresponds to the usual system of equations of the form

$$
\begin{aligned}
y^{\prime} & =v \\
v^{\prime} & =f(y, v),
\end{aligned}
$$

which is the reduction to a first order system of the second order equation $y^{\prime \prime}=f\left(y, y^{\prime}\right)$. What is the vector field on $T \mathbb{R}^{n}=\mathbb{R}^{n} \times \mathbb{R}^{n}$ which corresponds to this system?

Notation: For a second order differential equation $X$, the maximal integral curve through $v \in T M$ will be denoted by $\alpha_{v}$ and its projection will be denoted by $\gamma_{v}:=\pi_{T M} \circ \alpha$.\\
(5) A spray on $M$ is a second order differential equation, that is, a section $X$ of $T T M$ as in the previous problem, such that for $v \in T M$ and $s \in \mathbb{R}$, a number $t \in \mathbb{R}$ belongs to the domain of $\gamma_{s v}$ if and only if $s t$ belongs to the domain of $\gamma_{v}$, and in this case

$$
\gamma_{s v}(t)=\gamma_{v}(s t) .
$$

Show that there are infinitely many sprays on any smooth manifold $M$.\\[0pt]
[Hint: (i) Show that a vector field $X \in \mathfrak{X}(T M)$ is a spray if and only if $T \pi \circ X=\mathrm{id}_{T M}$, where $\pi=\pi_{T M}$ :\\
\includegraphics[scale=0.25, center]{2025_10_09_265d7e1a22c89f79c21eg-560}\\
(ii) Show that $X \in \mathfrak{X}(T M)$ is a spray if and only if for any $s \in \mathbb{R}$ and any $v \in T M$,

$$
X_{s v}=T \mu_{s}\left(s X_{v}\right),
$$

where $\mu_{s}: v \mapsto s v$ is the multiplication map.\\
(iii) Show the existence of a spray on an open ball in a Euclidean space.\\
(iv) Show that if $X_{1}$ and $X_{2}$ both satisfy one of the two characterizations of a spray above, then so does any convex combination of $X_{1}$ and $X_{2}$.]\\
(6) Show that if one has a linear connection on $M$, then there is a spray whose solutions are the geodesics of the connection.\\
(7) Show that given a spray on a manifold there is a (not unique) linear connection $\nabla$ on the manifold such that $\gamma: I \rightarrow M$ is a solution curve of the spray if and only if $\nabla_{\partial_{t}} \dot{\gamma}=0$ for all $t \in I$. Note: $\gamma$ is called a geodesic for the linear connection or the spray. Does the stipulation that the connection be torsion free force uniqueness?\\
(8) Let $X$ be a spray on $M$. Equivalently, we may start with a connection on $M$ (i.e. a connection on $T M$ ) which induces a spray. Show that the set $O_{X}:=\left\{v \in T M: \gamma_{v}(1)\right.$ is defined $\}$ is an open neighborhood of the zero section of $T M$.\\
(9) Finish the proof of Theorem 12.32.\\
(10) Prove formula (12.9).\\
(11) Let $h$ be a metric on a vector bundle $E \rightarrow M$. Show that there exists a metric connection for $h$. [Hint: Use orthonormal frames defined on each open set of a locally finite cover.]\\
(12) Show that parallel transport along curves with respect to a metric connection in a vector bundle with metric is an isometry of scalar product spaces.




\pagebreak
%\printbibliography

\clearpage
\end{document}